% !TEX program = xelatex
% 中文自动译本：Joel H. Ferziger, Milovan Perić, Robert L. Street,
% Computational Methods for Fluid Dynamics, Fourth Edition (Springer, 2020).
\documentclass[11pt,a4paper,openany,fontset=none]{ctexbook}
\usepackage[margin=2.55cm,headheight=15pt]{geometry}
\usepackage{microtype}
\usepackage{fancyhdr}
\usepackage{hyperref}
\usepackage{xurl}
\usepackage{titlesec}
\usepackage{amsmath,amssymb,amsfonts,bm,mathtools}
\usepackage[export]{adjustbox}
\usepackage{enumitem}
\usepackage{graphicx}
\usepackage{caption}
\usepackage{tikz}
\setmainfont{Arial Unicode MS}
\setCJKmainfont{Songti SC}
\setCJKsansfont{Songti SC}
\setCJKmonofont{Songti SC}
\setlength{\parindent}{2em}
\setlength{\parskip}{0.65em}
\setlength{\emergencystretch}{3em}
\linespread{1.30}\selectfont
\xeCJKsetup{CJKglue = \hskip 0.075em plus 0.04em minus 0.03em}
\setlist[itemize]{leftmargin=2.8em,itemsep=0.3em,topsep=0.45em}
\titleformat{\chapter}[display]{\bfseries\Huge\centering}{第\thechapter 章}{0.75em}{}
\titleformat{\section}{\large\bfseries}{}{0pt}{}
\titleformat{\subsection}{\normalsize\bfseries}{}{0pt}{}
\titleformat{\subsubsection}{\normalsize\bfseries}{}{0pt}{}
\titlespacing*{\section}{0pt}{1.45em}{0.7em}
\titlespacing*{\subsection}{0pt}{1.2em}{0.55em}
\pagestyle{fancy}
\fancyhf{}
\fancyhead[C]{\small\nouppercase{\leftmark}}
\fancyfoot[C]{\thepage}
\renewcommand{\headrulewidth}{0.3pt}
\hypersetup{
  colorlinks=true,
  linkcolor=black,
  urlcolor=blue,
  pdftitle={计算流体动力学的计算方法（第四版）},
  pdfauthor={Joel H. Ferziger, Milovan Perić, Robert L. Street}
}
\graphicspath{{./}}
\title{计算流体动力学的计算方法\\[0.5em]\Large 第四版·中文自动译文}
\author{Joel H. Ferziger、Milovan Perić、Robert L. Street 著\\
\small 译文由自动化工具生成，供学习与检索使用}
\date{原著第四版：2020\\生成日期：\today}
% 统一译本版式：避免公式、图片和题注跨页或侵入页眉、页脚。
\usepackage{needspace}
\clubpenalty=10000
\widowpenalty=10000
\displaywidowpenalty=10000
\raggedbottom
\setlength{\abovedisplayskip}{0.8\baselineskip plus 0.2\baselineskip minus 0.1\baselineskip}
\setlength{\belowdisplayskip}{0.8\baselineskip plus 0.2\baselineskip minus 0.1\baselineskip}
\setlength{\abovedisplayshortskip}{0.55\baselineskip}
\setlength{\belowdisplayshortskip}{0.55\baselineskip}
\newenvironment{sourceblock}
  {\par\addvspace{0.65\baselineskip}\noindent\begin{minipage}{\linewidth}\centering}
  {\end{minipage}\par\addvspace{0.45\baselineskip}}
\captionsetup{font=small,labelfont=bf,justification=centering,singlelinecheck=false,skip=0.45\baselineskip}
\begin{document}
\frontmatter
\maketitle
\chapter*{说明}
\markboth{说明}{说明}
本文件是 Joel H. Ferziger、Milovan Perić 与 Robert L. Street 所著
\emph{Computational Methods for Fluid Dynamics}（第四版）的中文自动译本。
正文按章节与小节的正常阅读顺序重新排版，不与原书逐页对齐。

原 PDF 使用专用数学字体，直接抽取时部分积分号、矩阵括号和上下标会错序。
为保证公式可靠，本译本把显示公式以及图形以原 PDF 的矢量裁片嵌入；
正文中的变量和较短表达式则保留在译文中。涉及推导、符号定义、图表数据与数值实现时，
请同时核对原著。原著版权归作者与 Springer Nature 所有。

\tableofcontents
\mainmatter
\chapter*{前言}\addcontentsline{toc}{chapter}{前言}
\markboth{前言}{前言}
计算流体动力学（通常简称为“CFD”）继续显著扩展。有许多软件包可以解决流体流动问题；成千上万的工程师正在广泛的行业和研究领域使用它们。市场每年以 15\% 左右的速度明显增长。如今，CFD 代码已被许多行业接受为设计工具，不仅用于解决问题，还可以帮助设计和优化各种产品，并作为研究工具。尽管自 1996 年本书第一版出版以来，商业 CFD 工具的用户友好性已大大提高，但为了使其有效且可靠的应用，用户仍然需要在流体力学和 CFD 方法方面拥有扎实的背景。我们假设读者熟悉理论流体力学，因此我们尝试提供有关另一个部分的有用信息 - 流体动力学的计算方法。

本书基于作者过去在斯坦福大学、埃尔兰根-纽伦堡大学和汉堡-哈尔堡技术大学的课程以及一些短期课程中提供的材料。它反映了作者在开发数值方法、编写 CFD 代码以及使用它们解决工程和地球物理问题方面的经验。示例中使用的许多代码，从涉及矩形网格的简单代码到使用非正交网格和多重网格方法的代码，都可供感兴趣的读者使用；附录中给出了如何通过互联网访问它们的信息。这些代码说明了书中描述的一些方法；它们可以扩展并适应许多流体机械问题的解决。学生应该尝试修改它们（例如，实现不同的边界条件、插值方案、微分和积分近似等）。这很重要，因为一个人在编写和/或运行一种方法之前并不真正了解该方法。我们了解到，许多研究人员过去曾使用这些代码作为其研究项目的基础。

本书偏爱有限体积法，尽管我们希望对有限差分法进行足够详细的描述。有限元方法没有详细介绍，因为已经有很多关于该主题的书籍。

每个主题的基本思想都以读者可以理解的方式描述；在可能的情况下，我们避免了冗长的数学分析。通常，对一个想法或方法的一般描述之后是代表该类型更好方法的一个或两个数值方案的更详细描述（包括必要的方程）；简要描述了其他可能的方法和扩展。我们试图强调方法的共同要素而不是它们的差异，并为构建变体提供基础。

我们非常重视估计数值误差的必要性；本书中几乎所有的例子都附有错误分析。尽管一个问题的定性错误的解可能看起来合理（甚至可能是另一个问题的良好解），但接受它的后果可能很严重。另一方面，有时如果小心处理，精度相对较低的解可能是有价值的。商业代码的工业用户在相信结果之前需要学会判断结果的质量。同样，研究人员也面临着同样的挑战。我们希望本书能够帮助人们认识到数值解始终是近似的并且需要进行适当的评估。

我们试图涵盖现代方法的横截面，包括任意多面体和重叠网格、多重网格方法和并行计算、移动网格和自由表面流的方法、湍流的直接和大涡模拟等。显然，我们无法详细涵盖所有这些主题，但我们希望本文包含的信息将为读者提供该主题的有用的一般知识；那些对特定主题进行更详细研究感兴趣的人会找到进一步阅读的建议。

本书上一版与当前版本之间相隔时间较长的原因是乔尔·H·费齐格 (Joel H. Ferziger) 在 2004 年突然去世。尽管上一版的剩余合著者在鲍勃街找到了继续该项目的优秀合作伙伴，但由于各种原因（但主要是时间不够），新版花了一段时间才完成。新的合著者带来了新的专业知识，时间的流逝也要求对大部分章节进行重大修改。值得注意的是，前任章。第 7 章涉及求解纳维-斯托克斯方程的方法已完全重写并分为两章。广泛用于大涡模拟的分数阶方法已经被更详细地描述，并且已经导出了新的隐式版本。基于分数步法的新代码已添加到该集中，读者可以从专门为此目的创建的网站 (\url{www.cfd-peric.de}) 下载这些代码。后面章节中的大多数示例都已使用商业软件重新计算；这些示例的模拟文件以及一些说明也可以从上述网站下载。

虽然我们已尽一切努力避免打字、拼写和其他错误，但毫无疑问，读者仍然会发现一些错误。如果您将可能发现的任何错误告知我们，以及对本书未来版本的改进提出意见和建议，我们将不胜感激。为此，作者的电子邮件地址如下。对某些示例的更正以及其他扩展报告将可以在上述网站上下载。

也希望作品没有被引用的同事见谅，因为任何遗漏都是无意的。

我们要感谢所有现在和以前的学生、同事和朋友，他们以各种方式帮助我们完成这项工作；完整的名单太长，无法在此列出。我们不能避免提及的名字包括（按字母顺序排列）博士。 Steven Arm field、David Briggs、Fotini (Tina) Katapodes Chow、Ismet Demirdžić、Gene Golub、Sylvain Lardeau、Ž eljko Lilek、Samir Muzaferija、Joseph Oliger、Eberhard Schreck、Volker Seidl 和 Kishan Shah。我们还非常感谢那些创建并提供 TEX、L A TEX、Linux、X fi g、Gnuplot 和其他使我们的工作变得更加轻松的工具的人们所提供的帮助。特别感谢 Rafael Ritterbusch 在第 1 章中提供了流固耦合示例。 13.

在我们的努力过程中，我们的家人给予了我们巨大的支持；我们特别感谢 Eva Ferziger、Anna James、Robinson 和 Kerstin Perić 以及 Norma Street。

地理遥远的同事之间的最初合作是通过亚历山大·冯·洪堡基金会（JHF）和 Deutsche Forschungsgemeinschaft（德国国家研究基金会，MP）的资助和奖学金实现的。没有他们的支持，这部作品就不可能问世，我们对他们表示深深的谢意。其中一位作者 (MP) 特别感谢已故 CD-adapco 前总裁 Peter S. MacDonald 的支持，并感谢西门子的经理（Jean-Claude Ercollanely、Deryl Sneider 和 Sven Enger），他们提供了支持和软件 Simcenter STAR-CCM+ 以在 Chaps 中创建示例。 9 – 13 1 . RLS 非常感谢有机会为他的好朋友兼研究同事 Joel Ferziger 的工作继续进行做出贡献。

德国杜伊斯堡 Milovan Peri ć milovan.peric@t-online.de 美国斯坦福 Joel H. Ferziger 美国斯坦福 Robert L. Street street@stanford.edu
1 节中的例子。 9.12.2、10.3.5、10.3.8、12.2.2、12.5.2、12.6.4 以及所有章节。 13（除非明确指定另一个来源）进行了模拟，并且图像是使用 Simcenter STAR-CCM+（Siemens Industry Software NV 及其任何附属公司的商标或注册商标）创建的。

\chapter*{缩略语}\addcontentsline{toc}{chapter}{缩略语}
\markboth{缩略语}{缩略语}
1D 一维 2D 二维 3D 三维 ADI 交替方向隐式 ALM 执行器线模型 BDS 后向差分格式 CDS 中心差分格式 CFD 计算流体动力学 CG 共轭梯度法 CGSTAB CG 稳定 CM 控制质量 CV 控制体积 CVFEM 基于控制体积的有限元方法 DDES 延迟分离涡模拟 DES 分离涡模拟 DNS 直接数值模拟EARSM 显式代数雷诺应力模型 EB 椭圆混合 ENO 本质无振荡 FAS 完全逼近格式 FD 有限差分 FDS 前向差分格式 FE 有限元 FFT 快速傅里叶变换 FMG 完全多重网格方法 FV 有限体积 GC 全局通信 GS 高斯-赛德尔方法 ICCG 不完全 Cholesky 方法预处理的 CG IDDES 改进的延迟分离涡模拟
IFSM 隐式分步法 ILES 隐式大涡模拟 ILU 不完全下上分解 LC 本地通信 LES 大涡模拟 LU 下上分解 MAC 标记单元 MG 多重网格 MPI 消息传递接口 ODE 常微分方程 PDE 偏微分方程 PVM 并行虚拟机 RANS 雷诺平均 纳维-斯托克斯 rms 均方根 rpm 每分钟转数 RSFS解析子滤波器尺度 RSM 雷诺应力模型 SBL 稳定边界层 SCL 空间守恒定律 SFS 子滤波器尺度 SGS 子网格尺度 SIP 强隐式过程 SOR 逐次过松弛 SST 剪应力传递 TDMA 三对角矩阵算法 TRANS 瞬态 RANS TVD 总变差递减 UDS 迎风差分方案 URANS 非定常 RANS VLES 超大涡模拟 VOF 流体体积

\chapter{流体流动的基本概念}
\section*{1.1 简介}\phantomsection\addcontentsline{toc}{section}{1.1 简介}
流体是其分子结构无法抵抗外部剪切力的物质：即使是最小的力也会导致流体颗粒变形。尽管液体和气体之间存在显著区别，但这两种类型的流体都遵循相同的运动定律。在大多数感兴趣的情况下，流体可以被视为连续体，即连续物质。

流体流动是由外力的作用引起的。常见的驱动力包括压差、重力、剪切力、旋转和表面张力。它们可以分为表面力（例如，由于海洋上方的风所产生的剪切力或刚性壁相对于流体的运动产生的压力和剪切力）和体积力（例如，重力和旋转引起的力）。

虽然所有流体在力的作用下表现相似，但它们的宏观特性却有很大差异。如果要研究流体运动，就必须了解这些属性；简单流体最重要的属性是密度和粘度。其他因素，例如普朗特数、比热和表面张力，仅在某些条件下（例如，当存在较大温差时）影响流体流动。流体特性是其他物理变量（例如温度和压力）的函数；尽管可以根据统计力学或动力学理论来估计其中一些，但它们通常是通过实验室测量获得的。

流体力学是一个非常广阔的领域。需要一个小型图书馆来涵盖其中可能包含的所有主题。在本书中，我们将主要关注工程师感兴趣的流动，但即使这是一个非常广泛的领域（例如，从风力涡轮机到燃气涡轮机，从纳米级到空中客车级，从暖通空调到人体血流）。不过，我们可以尝试对可能遇到的问题类型进行分类。这个方案的一个更数学化但不太完整的版本可以见第 1.8 节.

流动的速度通过多种方式影响其特性。在足够低的速度下，流体的惯性可以被忽略，并且我们有爬行流。这种状态对于含有小颗粒（悬浮液）的流动、通过多孔介质或狭窄通道（涂层技术、微型设备）的流动非常重要。随着速度的增加，惯性变得重要，但每个流体粒子都遵循平滑的轨迹；则称该流动为层流。速度的进一步增加可能会导致不稳定，最终产生一种更加随机的流动，称为湍流；层流-湍流转变过程本身就是一个重要领域。最后，流体中的流速与声速（马赫数）的比值决定是否需要考虑运动动能与内部自由度之间的交换。对于小马赫数，Ma < 0.3、流动可认为是不可压缩的；否则，它是可压缩的。如果 Ma < 1，则流动称为亚音速流动；当Ma > 1时，流动是超音速的，并且可能产生激波。最后，对于 Ma > 5，压缩可能会产生足够高的温度来改变流体的化学性质；这种流动被称为高超音速。这些区别影响问题的数学性质，从而影响解决方法。请注意，我们根据马赫数将流动称为可压缩或不可压缩，即使可压缩性是流体的属性。这是常见的术语，因为低马赫数下的可压缩流体的流动基本上是不可压缩的。

现在，工程师通常处理海洋和大气中的地球物理流。在那里，流体密度对压力作出响应，因此即使在没有运动的情况下，流体在许多情况下也可以有效压缩。然而，除了涉及深海的问题外，海水中的声速非常大，并且海水可以被视为不可压缩，尽管其密度取决于海洋温度和盐浓度。气氛截然不同。在那里，压力和空气密度随着海拔高度呈指数下降，因此流体可能需要被视为可压缩的，除了靠近地球表面的大气边界层。

在许多流动中，粘度的影响仅在壁附近很重要，因此域的最大部分中的流动可以被认为是无粘性的。在我们在本书中讨论的流体中，牛顿粘度定律是一个很好的近似，并且将专门使用它。遵守牛顿定律的流体称为牛顿流体；非牛顿流体对于某些工程应用很重要，但这里不予讨论。

许多其他现象影响流体流动。其中包括导致传热的温度差异和产生浮力的密度差异。它们以及溶质浓度的差异可能会显著影响流动，甚至是流动的唯一原因。相变（沸腾、冷凝、熔化和凝固）一旦发生，总是会导致流动的重要变化并产生多相流。其他属性（例如粘度、表面张力等）的变化也可能在确定流动性质方面发挥重要作用。除了少数例外，本书不会考虑这些影响。

在本章中，控制流体流动和相关现象的基本方程将以几种形式呈现：（i）坐标无关形式，它可以专门用于各种坐标系，（ii）有限控制体积的积分形式，它作为一类重要数值方法的起点，以及（iii）笛卡尔参考系中的微分（张量）形式，它是另一种重要方法的基础。这里仅简要概述用于推导这些方程的基本守恒定律和定律；更详细的推导可以在许多流体力学标准文本中找到（例如，Bird et al. 2006；White 2010）。假设读者对流体流动和相关现象的物理学有所了解，因此我们将重点关注控制方程的数值求解技术。

\section*{1.2 守恒原理}\phantomsection\addcontentsline{toc}{section}{1.2 守恒原理}
守恒定律可以通过考虑给定的物质或控制质量 (CM) 及其广泛的属性（例如质量、动量和能量）来推导。这种方法用于研究固体动力学，其中 CM（有时称为系统）很容易识别。然而，在流体流动中，很难跟踪一块物质。处理我们称为控制体积（CV）的特定空间区域内的流动比处理快速穿过感兴趣区域的物质块内的流动更方便。这种分析方法称为控制量法。

我们将主要关注两个广泛的性质：质量和动量。这些和其他性质的守恒方程具有共同的项，将首先考虑这些共同的项。

广泛属性的守恒定律将该属性在给定控制质量中的数量变化率与外部确定的影响联系起来。对于在工程利益的流动中既不产生也不破坏的质量，守恒方程可以写成：

\begin{sourceblock}
\includegraphics[page=19,trim={53.00bp 272.32bp 53.87bp 365.70bp},clip,width=0.92\textwidth,max height=0.38\textheight]{Ferziger4th.pdf}
\end{sourceblock}
另一方面，动量可以通过力的作用而改变，其守恒定律就是牛顿第二运动定律：

\begin{sourceblock}
\begin{tikzpicture}
\node[inner sep=0,anchor=south west] (image) at (0,0) {\includegraphics[page=19,trim={53.00bp 204.66bp 53.87bp 433.09bp},clip,width=0.92\textwidth,max height=0.38\textheight]{Ferziger4th.pdf}};
\begin{scope}[x={(image.south east)},y={(image.north west)}]
\fill[white] (0.40054,0.53751) rectangle (0.48217,0.95773);
\end{scope}
\end{tikzpicture}
\end{sourceblock}
其中 t 代表时间，m 代表质量，v 代表速度，f 代表作用在控制质量上的力。 1
我们将把这些定律转化为控制体积的形式，以便在本书中使用。基本变量将是密集的而不是广泛的属性；前者是与所考虑的物质数量无关的属性。例如密度 ρ（每单位体积的质量）和速度 v（每单位质量的动量）。

如果 φ 是任何守恒的密集性质（对于质量守恒， φ = 1；对于动量守恒， φ = v；对于标量守恒， φ 表示每单位质量的守恒性质），则相应的广延性质可以表示为：

\begin{sourceblock}
\begin{tikzpicture}
\node[inner sep=0,anchor=south west] (image) at (0,0) {\includegraphics[page=20,trim={53.00bp 551.46bp 53.87bp 85.47bp},clip,width=0.92\textwidth,max height=0.38\textheight]{Ferziger4th.pdf}};
\begin{scope}[x={(image.south east)},y={(image.north west)}]
\fill[white] (0.42650,0.56316) rectangle (0.45468,0.95892);
\end{scope}
\end{tikzpicture}
\end{sourceblock}
其中V CM 代表CM 所占据的体积。使用这个定义，控制体积的每个守恒方程的左侧可以写成 2：

\begin{sourceblock}
\begin{tikzpicture}
\node[inner sep=0,anchor=south west] (image) at (0,0) {\includegraphics[page=20,trim={53.00bp 474.01bp 53.87bp 156.30bp},clip,width=0.92\textwidth,max height=0.38\textheight]{Ferziger4th.pdf}};
\begin{scope}[x={(image.south east)},y={(image.north west)}]
\fill[white] (0.16746,0.64387) rectangle (0.18725,0.96651);
\fill[white] (0.16195,0.25147) rectangle (0.19068,0.57717);
\fill[white] (0.19092,0.03349) rectangle (0.23125,0.29082);
\end{scope}
\end{tikzpicture}
\end{sourceblock}
其中 V CV 是 CV 体积，S CV 是包围 CV 的表面，n 是与 S CV 正交且指向外的单位矢量，v 是流体速度，v s 是 CV 表面移动的速度。对于我们大多数时候要考虑的固定 CV，v s = 0 并且右侧的一阶导数成为局部（偏）导数。该方程指出，控制质量中的属性量的变化率 是控制体积内属性的变化率加上由于相对于 CV 边界的流体运动而通过 CV 边界的净通量。最后一项通常称为 穿过 CV 边界的对流（有时称为平流）通量。如果 CV 移动，使其边界与控制质量的边界重合，则 v = v s 并且根据要求此项将为零。

该方程的详细推导在许多流体动力学教科书中都有给出（例如，Bird et al. 2006；Street et al. 1996；Pritchard 2010），这里不再重复。质量、动量和标量守恒方程将在接下来的三节中介绍。为了方便起见，将考虑固定简历； V 代表 CV 体积，S 代表表面积。

\section*{1.3 质量守恒}\phantomsection\addcontentsline{toc}{section}{1.3 质量守恒}
通过设置 φ = 1，质量守恒（连续性）方程的积分形式直接从控制体积方程得出：

\begin{sourceblock}
\begin{tikzpicture}
\node[inner sep=0,anchor=south west] (image) at (0,0) {\includegraphics[page=20,trim={53.00bp 151.92bp 53.87bp 484.76bp},clip,width=0.92\textwidth,max height=0.38\textheight]{Ferziger4th.pdf}};
\begin{scope}[x={(image.south east)},y={(image.north west)}]
\fill[white] (0.30677,0.56687) rectangle (0.32749,0.95927);
\fill[white] (0.30123,0.08045) rectangle (0.33194,0.48201);
\end{scope}
\end{tikzpicture}
\end{sourceblock}
通过将高斯散度定理应用于对流项，我们可以将表面积分转化为体积积分。允许控制体积变得无穷小导致连续性方程的微分无坐标形式：

\footnotetext[2]{该方程通常称为控制体积方程或雷诺输运定理。}
\begin{sourceblock}
\includegraphics[page=21,trim={53.00bp 584.10bp 53.87bp 53.75bp},clip,width=0.92\textwidth,max height=0.38\textheight]{Ferziger4th.pdf}
\end{sourceblock}
通过提供该系统中的散度算子的表达式，可以将该形式转换为特定于给定坐标系的形式。笛卡尔坐标系、圆柱坐标系和球坐标系等常见坐标系的表达式可以在许多教科书中找到（例如，Bird et al. 2006）；给出了适用于一般非正交坐标系的表达式，例如，Aris (1990) 或 Chen 等人。 （2004a）。下面我们以张量和扩展表示法呈现笛卡尔形式。在这里以及在本书中，我们将采用爱因斯坦约定，即每当相同的索引在任何项中出现两次时，就暗示了对该索引范围的求和：

\begin{sourceblock}
\includegraphics[page=21,trim={53.00bp 431.84bp 53.87bp 205.18bp},clip,width=0.92\textwidth,max height=0.38\textheight]{Ferziger4th.pdf}
\end{sourceblock}
其中 x i (i = 1 , 2 , 3) 或 ( x , y , z ) 是笛卡尔坐标， u i 或 ( u x , u y , u z ) 是速度矢量 v 的笛卡尔分量。笛卡尔形式的守恒方程经常被使用，本工作就是这种情况。非正交坐标系中的微分守恒方程将在第 9 章中介绍.

\section*{1.4 动量守恒}\phantomsection\addcontentsline{toc}{section}{1.4 动量守恒}
有多种推导动量守恒方程的方法。一种方法是使用第 1 节中描述的控制体积方法。 1.2；在此方法中，使用方程（1.2）和（1.4）并用 v 代替 φ，例如，对于包含固定流体的空间体积：

\begin{sourceblock}
\begin{tikzpicture}
\node[inner sep=0,anchor=south west] (image) at (0,0) {\includegraphics[page=21,trim={53.00bp 242.09bp 53.87bp 394.60bp},clip,width=0.92\textwidth,max height=0.38\textheight]{Ferziger4th.pdf}};
\begin{scope}[x={(image.south east)},y={(image.north west)}]
\fill[white] (0.00000,0.89643) rectangle (0.09338,1.00000);
\fill[white] (0.09609,0.89643) rectangle (0.12583,1.00000);
\fill[white] (0.12854,0.89643) rectangle (0.20818,1.00000);
\fill[white] (0.26896,0.56672) rectangle (0.28968,0.95925);
\fill[white] (0.26343,0.08048) rectangle (0.29414,0.48217);
\fill[white] (0.31702,0.04075) rectangle (0.33465,0.33175);
\end{scope}
\end{tikzpicture}
\end{sourceblock}
为了用强度属性来表达右侧，必须考虑可能作用在 CV 中的流体上的力：

\begin{itemize}
\item 表面力（压力、法向应力和剪切应力、表面张力等）；
\item 体积力（重力、离心力、科里奥利力、电磁力等）。
\end{itemize}
从分子的角度来看，由于压力和应力而产生的表面力是表面上的微观动量通量。如果这些通量不能用方程所控制的守恒性质（密度和速度）来表示，则方程组不是封闭的；也就是说，方程少于因变量，并且无法求解。通过做出某些假设可以避免这种可能性。最简单的假设是流体是牛顿流体；幸运的是，牛顿模型适用于许多实际流体。

对于牛顿流体，应力张量 T 是动量传输的分子速率，可以写成：

\begin{sourceblock}
\begin{tikzpicture}
\node[inner sep=0,anchor=south west] (image) at (0,0) {\includegraphics[page=22,trim={53.00bp 583.46bp 53.87bp 50.00bp},clip,width=0.92\textwidth,max height=0.38\textheight]{Ferziger4th.pdf}};
\begin{scope}[x={(image.south east)},y={(image.north west)}]
\fill[white] (0.29143,0.25734) rectangle (0.31453,0.61108);
\fill[white] (0.31853,0.26438) rectangle (0.34671,0.61842);
\fill[white] (0.35020,0.26438) rectangle (0.37841,0.61842);
\end{scope}
\end{tikzpicture}
\end{sourceblock}
其中 μ 是动态粘度，I 是单位张量，p 是静压力，D 是应变（变形）张量率：

\begin{sourceblock}
\begin{tikzpicture}
\node[inner sep=0,anchor=south west] (image) at (0,0) {\includegraphics[page=22,trim={53.00bp 520.70bp 53.87bp 123.10bp},clip,width=0.92\textwidth,max height=0.38\textheight]{Ferziger4th.pdf}};
\begin{scope}[x={(image.south east)},y={(image.north west)}]
\fill[white] (0.41886,0.36661) rectangle (0.43865,0.88406);
\fill[white] (0.35353,0.05998) rectangle (0.37997,0.57744);
\fill[white] (0.38349,0.07028) rectangle (0.41170,0.58774);
\fill[white] (0.63432,0.07028) rectangle (0.64737,0.58774);
\fill[white] (0.92421,0.05372) rectangle (1.00000,0.57162);
\end{scope}
\end{tikzpicture}
\end{sourceblock}
这两个方程可以用笛卡尔坐标中的指数表示法写成如下：

\begin{sourceblock}
\begin{tikzpicture}
\node[inner sep=0,anchor=south west] (image) at (0,0) {\includegraphics[page=22,trim={53.00bp 421.53bp 53.87bp 178.12bp},clip,width=0.92\textwidth,max height=0.38\textheight]{Ferziger4th.pdf}};
\begin{scope}[x={(image.south east)},y={(image.north west)}]
\fill[white] (0.00000,0.92946) rectangle (0.10159,1.00000);
\fill[white] (0.26556,0.68642) rectangle (0.29158,0.87261);
\fill[white] (0.29014,0.68642) rectangle (0.30078,0.81531);
\fill[white] (0.30785,0.70266) rectangle (0.33603,0.87667);
\fill[white] (0.33955,0.70266) rectangle (0.36773,0.87667);
\fill[white] (0.42571,0.24004) rectangle (0.44553,0.41405);
\fill[white] (0.34189,0.12423) rectangle (0.37411,0.31042);
\fill[white] (0.37266,0.12423) rectangle (0.38331,0.25327);
\fill[white] (0.39038,0.14062) rectangle (0.41856,0.31448);
\fill[white] (0.56505,0.01805) rectangle (0.60809,0.20815);
\end{scope}
\end{tikzpicture}
\end{sourceblock}
其中 δ i j 是克罗内克符号（如果 i = j，则 δ i j = 1，否则 δ i j = 0）。对于不可压缩流，根据连续性方程，式（1.11）括号中的第二项为零。文献中经常使用以下符号来描述应力张量的粘性部分：

\begin{sourceblock}
\includegraphics[page=22,trim={53.00bp 327.83bp 53.87bp 310.29bp},clip,width=0.92\textwidth,max height=0.38\textheight]{Ferziger4th.pdf}
\end{sourceblock}
对于非牛顿流体，应力张量和速度之间的关系通常由一组偏微分方程定义，整个问题要复杂得多；例如，参见 Bird 和 Wiest (1995)。对于使用与上述相同类型的本构关系描述的非牛顿流体类别，但仅需要可变粘度（通常是速度梯度和温度的非线性函数）或与第 1 章中描述的雷诺应力模型相当的应力模型。如图10所示，可以使用用于牛顿流体的相同方法。 3 然而，不同类型的非牛顿流体需要不同的本构方程，如 Bird 和 Wiest (1995) 所示；反过来，这些可能需要专门的解决方法。这个主题很复杂，仅在第 1 章中简要提及。 13.

当体积力（每单位质量）由 b 表示时，动量守恒方程的积分形式变为：

\begin{sourceblock}
\begin{tikzpicture}
\node[inner sep=0,anchor=south west] (image) at (0,0) {\includegraphics[page=22,trim={53.00bp 139.62bp 53.87bp 497.06bp},clip,width=0.92\textwidth,max height=0.38\textheight]{Ferziger4th.pdf}};
\begin{scope}[x={(image.south east)},y={(image.north west)}]
\fill[white] (0.11630,0.56653) rectangle (0.13702,0.95927);
\fill[white] (0.18689,0.32043) rectangle (0.22217,0.72200);
\fill[white] (0.22779,0.31704) rectangle (0.26734,0.71317);
\fill[white] (0.27374,0.32960) rectangle (0.30192,0.72200);
\fill[white] (0.11080,0.08079) rectangle (0.14150,0.48235);
\fill[white] (0.16439,0.04073) rectangle (0.18202,0.33198);
\end{scope}
\end{tikzpicture}
\end{sourceblock}
通过将高斯散度定理应用于对流和扩散，很容易获得动量守恒方程（1.14）的无坐标矢量形式

\footnotetext[3]{例如，血液在高剪切速率下可以被视为牛顿流体（Tokuda 等人，2008 年），但在其他情况下具有可变粘度（Perktold 和 Rappitsch，1995 年）。}
\begin{sourceblock}
\begin{tikzpicture}
\node[inner sep=0,anchor=south west] (image) at (0,0) {\includegraphics[page=23,trim={53.00bp 573.14bp 53.87bp 57.16bp},clip,width=0.92\textwidth,max height=0.38\textheight]{Ferziger4th.pdf}};
\begin{scope}[x={(image.south east)},y={(image.north west)}]
\fill[white] (0.00000,0.64397) rectangle (0.05077,0.96652);
\fill[white] (0.05344,0.64397) rectangle (0.13317,0.96652);
\end{scope}
\end{tikzpicture}
\end{sourceblock}
第 i 个笛卡尔分量的相应方程为：

\begin{sourceblock}
\includegraphics[page=23,trim={53.00bp 517.28bp 53.87bp 120.58bp},clip,width=0.92\textwidth,max height=0.38\textheight]{Ferziger4th.pdf}
\end{sourceblock}
由于动量是矢量，因此它通过 CV 边界的对流和扩散通量是二阶张量（ ρ vv 和 T ）与表面矢量 n d S 的标量积。上述方程的积分形式为：

\begin{sourceblock}
\begin{tikzpicture}
\node[inner sep=0,anchor=south west] (image) at (0,0) {\includegraphics[page=23,trim={53.00bp 436.41bp 53.87bp 200.27bp},clip,width=0.92\textwidth,max height=0.38\textheight]{Ferziger4th.pdf}};
\begin{scope}[x={(image.south east)},y={(image.north west)}]
\fill[white] (0.10072,0.56687) rectangle (0.12144,0.95927);
\fill[white] (0.17128,0.29294) rectangle (0.21383,0.72234);
\fill[white] (0.22256,0.31704) rectangle (0.26211,0.71317);
\fill[white] (0.26854,0.32960) rectangle (0.29672,0.72234);
\fill[white] (0.09522,0.08045) rectangle (0.12592,0.48201);
\fill[white] (0.14881,0.04073) rectangle (0.16644,0.33164);
\end{scope}
\end{tikzpicture}
\end{sourceblock}
其中（参见方程（1.9）和（1.10））：

\begin{sourceblock}
\begin{tikzpicture}
\node[inner sep=0,anchor=south west] (image) at (0,0) {\includegraphics[page=23,trim={53.00bp 379.35bp 53.87bp 251.40bp},clip,width=0.92\textwidth,max height=0.38\textheight]{Ferziger4th.pdf}};
\begin{scope}[x={(image.south east)},y={(image.north west)}]
\fill[white] (0.08358,0.21334) rectangle (0.10493,0.56315);
\fill[white] (0.11197,0.24385) rectangle (0.14015,0.57078);
\fill[white] (0.14367,0.24385) rectangle (0.16653,0.57078);
\fill[white] (0.16668,0.21334) rectangle (0.21786,0.57078);
\fill[white] (0.22322,0.24385) rectangle (0.25143,0.57078);
\fill[white] (0.25326,0.24385) rectangle (0.27612,0.57078);
\fill[white] (0.27630,0.23651) rectangle (0.35516,0.61430);
\fill[white] (0.36165,0.24385) rectangle (0.37480,0.57078);
\fill[white] (0.37663,0.21334) rectangle (0.39570,0.56315);
\fill[white] (0.40105,0.24385) rectangle (0.42923,0.57078);
\end{scope}
\end{tikzpicture}
\end{sourceblock}
这里 b i 代表物体力的第 i 个分量，上标 T 表示转置， i i 是坐标 x i 方向上的笛卡尔单位向量。在笛卡尔坐标系中，上述表达式可以写为：

\begin{sourceblock}
\begin{tikzpicture}
\node[inner sep=0,anchor=south west] (image) at (0,0) {\includegraphics[page=23,trim={53.00bp 294.46bp 53.87bp 335.58bp},clip,width=0.92\textwidth,max height=0.38\textheight]{Ferziger4th.pdf}};
\begin{scope}[x={(image.south east)},y={(image.north west)}]
\fill[white] (0.32084,0.42271) rectangle (0.36547,0.77285);
\fill[white] (0.21287,0.22909) rectangle (0.23423,0.57202);
\fill[white] (0.24126,0.25900) rectangle (0.26944,0.57950);
\fill[white] (0.27296,0.25900) rectangle (0.29585,0.57950);
\fill[white] (0.31982,0.05568) rectangle (0.35630,0.38338);
\fill[white] (0.35579,0.03324) rectangle (0.36644,0.27064);
\end{scope}
\end{tikzpicture}
\end{sourceblock}
向量场可以用多种不同的方式表示。定义向量的基本向量可以是局部的或全局的。在曲线坐标系中，当边界复杂时通常需要曲线坐标系（参见第 9 章），人们可以选择协变基或逆变基，参见图 1.1。前者用其沿局部坐标的分量来表示向量；后者使用垂直于坐标表面的投影。在笛卡尔系统中，两者变得相同。而且，基向量可以是无量纲的或有量纲的。包括所有这些选项，可能有 70 多种不同形式的动量方程。从数学上来说，一切都是等价的；从数字的角度来看，有些比其他更难处理。

如果所有项都具有向量或张量散度的形式，则动量方程被称为“强守恒形式”。仅当使用固定方向的分量时，这对于方程的分量形式才是可能的。坐标定向矢量分量随坐标方向转动，需要“表观力”来产生转动；这些力量在上面定义的意义上是非保守的。例如，在柱坐标系中，径向和圆周方向发生变化，因此空间恒定矢量（例如均匀速度场）的分量随 r 和 θ 变化，并且在以下位置处是奇异的：

\begin{sourceblock}
\includegraphics[page=24,trim={53.00bp 478.75bp 53.87bp 52.00bp},clip,width=\textwidth,max height=0.38\textheight]{Ferziger4th.pdf}
\par\vspace{0.45\baselineskip}
\small 图 1.1 通过不同分量表示向量： u i ——笛卡尔分量； vi——逆变分量； vi — 协变分量 [ v A = v B , ( u i ) A = ( u i ) B , ( vi i ) A̸ = ( vi ) B , ( vi i ) A̸ = ( vi i ) B ]
\end{sourceblock}
坐标原点。为了解释这一点，这些分量的方程包含离心力和科里奥利力项。

\begin{center}\small 图 1.1 显示了向量 v 及其逆变、协变和笛卡尔分量。显然，即使向量 v 保持不变，逆变分量和协变分量也会随着基向量的变化而变化。我们将在第 1 章中讨论速度分量的选择对数值求解方法的影响。 9.\end{center}
当与有限体积方法一起使用时，方程的强守恒形式会自动确保计算中的全局动量守恒。这是守恒方程的一个重要性质，其在数值解中的保存同样重要。保留这一性质有助于保证数值方法在求解过程中不会发散，可以被视为一种“可实现性”。

对于某些流动，解析空间可变方向上的动量是有利的。例如，线涡旋中的速度在柱坐​​标系中只有一个分量 u θ 而在笛卡尔坐标系中只有两个分量。在极柱坐标系中分析时，无旋流的轴对称流是二维 (2D)，但在使用笛卡尔坐标系时是三维 (3D)。一些使用非正交坐标的数值技术需要使用逆变速度分量。然后，方程包含所谓的“曲率项”，这些项很难精确计算，因为它们包含难以近似的坐标变换的二阶导数。

在本书中，我们将根据笛卡尔分量来处理速度矢量和应力张量，并且我们将使用笛卡尔动量方程的保守形式。

式(1.16)是强守恒形式。利用连续性方程可以得到该方程的非保守形式；因为

\begin{sourceblock}
\includegraphics[page=24,trim={53.00bp 84.91bp 53.87bp 565.77bp},clip,width=0.92\textwidth,max height=0.38\textheight]{Ferziger4th.pdf}
\end{sourceblock}
由此可见：

\begin{sourceblock}
\begin{tikzpicture}
\node[inner sep=0,anchor=south west] (image) at (0,0) {\includegraphics[page=25,trim={53.00bp 573.14bp 53.87bp 64.71bp},clip,width=0.92\textwidth,max height=0.38\textheight]{Ferziger4th.pdf}};
\begin{scope}[x={(image.south east)},y={(image.north west)}]
\fill[white] (0.00000,0.81584) rectangle (0.02081,1.00000);
\fill[white] (0.02346,0.81584) rectangle (0.11738,1.00000);
\fill[white] (0.12009,0.81584) rectangle (0.17817,1.00000);
\end{scope}
\end{tikzpicture}
\end{sourceblock}
t i 中包含的压力项也可以写为

\begin{sourceblock}
\includegraphics[page=25,trim={53.00bp 528.31bp 53.87bp 122.38bp},clip,width=0.92\textwidth,max height=0.38\textheight]{Ferziger4th.pdf}
\end{sourceblock}
压力梯度则被视为体积力；这相当于对压力项的非保守处理。方程的非保守形式通常用于有限差分方法，因为它更简单。在极细网格的限制下，所有方程形式和数值求解方法都给出相同的解；然而，在粗网格上，非保守形式会引入额外的错误，这可能会变得很重要。

如果将应力张量粘性部分的表达式 Eq.( 1.13 ) 代入以指数表示法和笛卡尔坐标编写的 Eq.( 1.16 )，并且如果重力是唯一的体积力，则有：

\begin{sourceblock}
\includegraphics[page=25,trim={53.00bp 373.06bp 53.87bp 263.95bp},clip,width=0.92\textwidth,max height=0.38\textheight]{Ferziger4th.pdf}
\end{sourceblock}
其中 g i 是重力加速度 g 在笛卡尔坐标 x i 方向上的分量。对于密度和重力恒定的情况，ρ g 项可以写为 ∇ ( ρ g · r )，其中 r 是位置向量，r = x i i i （通常假设重力作用于负 z 方向，即 g = g z k，g z 为负；在这种情况下 g · r = g z z ）。那么 − ρ g z z 是静水压力，并且为了更有效地进行数值求解，将 ̃ p = p − ρ g z z 定义为水头并使用它代替压力是很方便的。然后项 ρ g i 从上面的方程中消失。如果需要实际压力，只需在 ̃ p 上加上 ρ g z z 即可。

由于方程中仅出现压力梯度，因此除了可压缩流（包括大气和一些海洋流）和自由表面暴露于大气的流外，压力的绝对值并不重要。

在变密度流中（本书中考虑的所有流中重力的变化都可以忽略不计），可以将 ρ g i 项分为两部分： ρ 0 g i + ( ρ − ρ 0 ) g i，其中 ρ 0 是参考密度。然后，第一部分可以包含压力，并且如果密度变化仅保留在引力项中，则我们有布辛涅斯克近似，请参见第 1 节。 1.7.

\section*{1.5 标量守恒}\phantomsection\addcontentsline{toc}{section}{1.5 标量守恒}
\begin{sourceblock}
\begin{tikzpicture}
\node[inner sep=0,anchor=south west] (image) at (0,0) {\includegraphics[page=26,trim={53.00bp 582.93bp 53.87bp 53.75bp},clip,width=0.92\textwidth,max height=0.38\textheight]{Ferziger4th.pdf}};
\begin{scope}[x={(image.south east)},y={(image.north west)}]
\fill[white] (0.25693,0.56687) rectangle (0.27765,0.95927);
\fill[white] (0.25140,0.08045) rectangle (0.28211,0.48201);
\fill[white] (0.30502,0.04073) rectangle (0.32265,0.33164);
\end{scope}
\end{tikzpicture}
\end{sourceblock}
其中 f φ 表示 φ 通过对流以外的机制以及标量的任何源或汇进行的传输。扩散传输始终存在（即使在静止的流体中），并且通常通过梯度近似来描述，例如热扩散的傅立叶定律和质量扩散的菲克定律：

\begin{sourceblock}
\begin{tikzpicture}
\node[inner sep=0,anchor=south west] (image) at (0,0) {\includegraphics[page=26,trim={53.00bp 491.65bp 53.87bp 150.55bp},clip,width=0.92\textwidth,max height=0.38\textheight]{Ferziger4th.pdf}};
\begin{scope}[x={(image.south east)},y={(image.north west)}]
\fill[white] (0.36653,0.39474) rectangle (0.37967,0.87761);
\fill[white] (0.38111,0.59190) rectangle (0.39642,0.94987);
\end{scope}
\end{tikzpicture}
\end{sourceblock}
其中 是量 φ 的扩散率。一个例子是能量方程，对于大多数工程流程，可以写成：

\begin{sourceblock}
\begin{tikzpicture}
\node[inner sep=0,anchor=south west] (image) at (0,0) {\includegraphics[page=26,trim={53.00bp 394.86bp 53.87bp 214.90bp},clip,width=0.92\textwidth,max height=0.38\textheight]{Ferziger4th.pdf}};
\begin{scope}[x={(image.south east)},y={(image.north west)}]
\fill[white] (0.15071,0.77368) rectangle (0.17143,0.97872);
\fill[white] (0.14517,0.51969) rectangle (0.17588,0.72951);
\fill[white] (0.19877,0.49876) rectangle (0.21639,0.65076);
\fill[white] (0.36051,0.02128) rectangle (0.37814,0.17347);
\end{scope}
\end{tikzpicture}
\end{sourceblock}
其中 h = p / ρ + e 是单位质量的焓或比焓，它是系统总能量的度量，e 是内能。另外，T为温度，k为热导率，k=μcp/Pr，S为应力张量的粘性部分，S=T+pI。 Pr 是普朗特数，定义为动量扩散率与热扩散率之比，cp 是恒压比热。源项代表压力和粘性力所做的功；在不可压缩流动中它可以被忽略。通过考虑具有恒定比热的流体来实现进一步的简化，在这种情况下，温度结果的对流/扩散方程：

\begin{sourceblock}
\begin{tikzpicture}
\node[inner sep=0,anchor=south west] (image) at (0,0) {\includegraphics[page=26,trim={53.00bp 241.55bp 53.87bp 395.13bp},clip,width=0.92\textwidth,max height=0.38\textheight]{Ferziger4th.pdf}};
\begin{scope}[x={(image.south east)},y={(image.north west)}]
\fill[white] (0.19113,0.56687) rectangle (0.21185,0.95927);
\fill[white] (0.18559,0.08045) rectangle (0.21633,0.48201);
\fill[white] (0.23922,0.04073) rectangle (0.25684,0.33164);
\end{scope}
\end{tikzpicture}
\end{sourceblock}
物质浓度方程具有相同的形式，其中 T 替换为浓度 c，Pr 替换为 Sc（施密特数），即动量扩散率与物种扩散率之比。

以一般形式编写守恒方程很有用，因为上述所有方程都有共同的项。然后就可以用一般的方式进行离散化和分析；必要时，可以单独处理方程特有的项。

通用守恒方程的积分形式直接由方程（1.22）和（1.23）得出：

\begin{sourceblock}
\begin{tikzpicture}
\node[inner sep=0,anchor=south west] (image) at (0,0) {\includegraphics[page=26,trim={53.00bp 88.23bp 53.87bp 548.45bp},clip,width=0.92\textwidth,max height=0.38\textheight]{Ferziger4th.pdf}};
\begin{scope}[x={(image.south east)},y={(image.north west)}]
\fill[white] (0.09038,0.56687) rectangle (0.11110,0.95927);
\fill[white] (0.16093,0.32994) rectangle (0.19907,0.72234);
\fill[white] (0.20427,0.31738) rectangle (0.24382,0.71317);
\fill[white] (0.25026,0.32994) rectangle (0.27844,0.72234);
\fill[white] (0.08484,0.08079) rectangle (0.11555,0.48235);
\fill[white] (0.13847,0.04073) rectangle (0.15609,0.33198);
\end{scope}
\end{tikzpicture}
\end{sourceblock}
其中 q φ 是 φ 的源或汇。该方程的无坐标向量形式为： ∂ ( ρφ )

\begin{sourceblock}
\begin{tikzpicture}
\node[inner sep=0,anchor=south west] (image) at (0,0) {\includegraphics[page=27,trim={53.00bp 561.19bp 53.87bp 76.67bp},clip,width=0.92\textwidth,max height=0.38\textheight]{Ferziger4th.pdf}};
\begin{scope}[x={(image.south east)},y={(image.north west)}]
\fill[white] (0.00000,0.81577) rectangle (0.03248,1.00000);
\end{scope}
\end{tikzpicture}
\end{sourceblock}
在笛卡尔坐标和张量表示法中，通用守恒方程的微分形式为：

\begin{sourceblock}
\begin{tikzpicture}
\node[inner sep=0,anchor=south west] (image) at (0,0) {\includegraphics[page=27,trim={53.00bp 492.09bp 53.87bp 144.93bp},clip,width=0.92\textwidth,max height=0.38\textheight]{Ferziger4th.pdf}};
\begin{scope}[x={(image.south east)},y={(image.north west)}]
\fill[white] (0.24150,0.56113) rectangle (0.31958,0.95810);
\fill[white] (0.35865,0.55254) rectangle (0.42268,0.95879);
\fill[white] (0.42289,0.52438) rectangle (0.46617,0.95879);
\fill[white] (0.32502,0.32143) rectangle (0.35320,0.71841);
\fill[white] (0.26385,0.06902) rectangle (0.29456,0.47527);
\end{scope}
\end{tikzpicture}
\end{sourceblock}
首先将描述该通用守恒方程的数值方法。连续性和动量方程（通常称为纳维-斯托克斯方程）的特殊特征将在后面描述，作为通用方程方法的扩展。

\section*{1.6 方程的无量纲形式}\phantomsection\addcontentsline{toc}{section}{1.6 方程的无量纲形式}
流动的实验研究通常在模型上进行，结果以无量纲形式显示，从而可以缩放到真实的流动条件。在数值研究中也可以采用相同的方法。通过使用适当的归一化，可以将控制方程转换为无量纲形式。例如，速度可以通过参考速度 v 0 标准化，空间坐标通过参考长度 L 0 标准化，时间通过某个参考时间 t 0 标准化，压力通过 ρ v 2 标准化
0，温度通过某个参考温差T 1 - T 0。无量纲变量则为：

\begin{sourceblock}
\includegraphics[page=27,trim={53.00bp 242.99bp 53.87bp 392.37bp},clip,width=0.92\textwidth,max height=0.38\textheight]{Ferziger4th.pdf}
\end{sourceblock}
如果流体属性恒定，则连续性、动量和温度方程的无量纲形式为：

\begin{sourceblock}
\includegraphics[page=27,trim={53.00bp 94.90bp 53.87bp 461.03bp},clip,width=0.92\textwidth,max height=0.38\textheight]{Ferziger4th.pdf}
\end{sourceblock}
方程中出现以下无量纲数：

\begin{sourceblock}
\includegraphics[page=28,trim={53.00bp 583.12bp 53.87bp 53.51bp},clip,width=0.92\textwidth,max height=0.38\textheight]{Ferziger4th.pdf}
\end{sourceblock}
分别称为斯特鲁哈尔数、雷诺数和弗劳德数。 γ i 是归一化重力加速度矢量在 x i 方向上的分量。

对于自然对流，通常使用 Boussinesq 近似，在这种情况下动量方程的最后一项变为：

Ra Re 2 Pr T ＊ γ i ,
其中 Ra 是瑞利数，定义为：

\begin{sourceblock}
\includegraphics[page=28,trim={53.00bp 433.96bp 53.87bp 203.14bp},clip,width=0.92\textwidth,max height=0.38\textheight]{Ferziger4th.pdf}
\end{sourceblock}
这里 Gr 是另一个无量纲数，称为格拉霍夫数，β 是热膨胀系数。

归一化量的选择在简单流程中是显而易见的； v 0 是平均速度，L 0 是几何长度尺度； T 0 和T 1 是冷壁温度和热壁温度。如果几何形状复杂、流体属性不稳定或边界条件不稳定，则描述流动所需的无量纲参数的数量可能会变得非常大，并且无量纲方程可能不再有用。

无量纲方程可用于分析研究和确定方程中各项的相对重要性。例如，它们表明，通道或管道中的稳态流动仅取决于雷诺数；然而，如果几何形状发生变化，流动也会受到边界形状的影响。因为我们对计算复杂几何形状的流动感兴趣，所以我们将在本书中使用输运方程的维度形式。

\section*{1.7 简化数学模型}\phantomsection\addcontentsline{toc}{section}{1.7 简化数学模型}
质量和动量的守恒方程比看起来更复杂。它们是非线性的、耦合的并且难以求解。现有的数学工具很难证明特定边界条件下存在唯一解。经验表明，纳维-斯托克斯方程可以准确地描述牛顿流体的流动。只有在少数情况下——大多数是简单几何形状中完全发展的流动，例如在管道中、平行板之间等——才有可能获得纳维-斯托克斯方程的解析解。这些流动对于研究流体动力学基础很重要，但其实际意义有限。

在所有可能有这种解的情况下，方程中的许多项都为零。对于其他流，有些术语并不重要，我们可能会忽略它们；这种简化引入了一个错误。在大多数情况下，即使是简化的方程也无法解析求解；必须使用数值方法。计算工作可能比完整方程小得多，这是简化的理由。我们在下面列出了一些可以简化运动方程的流动类型。

\subsection*{1.7.1 不可压缩流}\phantomsection\addcontentsline{toc}{subsection}{1.7.1 不可压缩流}
下列各节中提出的质量和动量守恒方程。 1.3和1.4是最通用的；他们假设所有流体和流动特性都随空间和时间而变化。在许多应用中，流体密度可以假设为恒定。这不仅适用于液体流（其压缩性在大多数情况下确实可以忽略不计），而且适用于马赫数低于 0.3 的气体流。这种流动被认为是不可压缩的。 4 如果流动也是等温的，则粘度也是恒定的。在这种情况下，质量和动量守恒方程 (1.6) 和 (1.16) 简化为：

\begin{sourceblock}
\begin{tikzpicture}
\node[inner sep=0,anchor=south west] (image) at (0,0) {\includegraphics[page=29,trim={53.00bp 320.31bp 53.87bp 295.90bp},clip,width=0.92\textwidth,max height=0.38\textheight]{Ferziger4th.pdf}};
\begin{scope}[x={(image.south east)},y={(image.north west)}]
\fill[white] (0.00000,0.97657) rectangle (0.03579,1.00000);
\end{scope}
\end{tikzpicture}
\end{sourceblock}
其中 ν = μ / ρ 是运动粘度。这种简化通常没有多大价值，因为方程很难求解。然而，它确实有助于数值求解。

\subsection*{1.7.2 无粘（欧拉）流}\phantomsection\addcontentsline{toc}{subsection}{1.7.2 无粘（欧拉）流}
在远离固体表面的流动中，粘度的影响通常很小。如果完全忽略粘性效应，即，如果我们假设应力张量简化为 T = − p I，则纳维-斯托克斯方程就简化为欧拉方程。连续性方程与(1.6)相同，动量方程为：

\begin{sourceblock}
\includegraphics[page=29,trim={53.00bp 131.23bp 53.87bp 506.62bp},clip,width=0.92\textwidth,max height=0.38\textheight]{Ferziger4th.pdf}
\end{sourceblock}
4 在某些情况下，例如非常高的压力或在深海中，需要考虑液体的可压缩性。同样，正如第 1 节中所指出的。 1.1，在模拟大气时，即使马赫数非常小，也可能需要使用流动方程的可压缩版本。

由于流体被假定为无粘性，因此它不能粘附在壁上，并且在固体边界处可能发生滑移。欧拉方程通常用于研究高马赫数下的可压缩流。在高速下，雷诺数非常高，粘性和湍流效应仅在靠近壁的小区域中很重要。这些流动通常可以使用欧拉方程很好地预测。

尽管欧拉方程不容易求解，但由于无需求解壁附近的边界层，因此可以使用较粗的网格。然而，正如 Hirsch (2007) 在其引言中所报告的那样，自 20 世纪 90 年代中期以来，求解方法和计算机能力的发展使得对整个飞机、船舶、车辆等以及通过多级压缩机和泵等的流动进行完整的三维纳维-斯托克斯模拟成为可能。自 20 世纪 80 年代初以来，完整的欧拉模拟已经成为可能。如今，工程师使用最有效的工具来完成任务，并且欧拉方程仍然是基本工具集的一部分（例如，参见 Wie 等人，2010）。

有许多方法被设计用来求解可压缩欧拉方程。其中一些将在第 1 章中简要描述。 11.有关这些方法的更多详细信息，请参阅 Hirsch (2007)、Fletcher (1991)、Knight (2006) 和 Tannehill 等人的书籍。 （1977）等。本书中描述的求解方法也可用于求解可压缩欧拉方程，正如我们将在第 1 章中看到的那样。如图 11 所示，它们的性能与为可压缩流设计的特殊方法一样好。

\subsection*{1.7.3 势流}\phantomsection\addcontentsline{toc}{subsection}{1.7.3 势流}
最简单的流动模型之一是势流。假设流体是无粘性的（如欧拉方程）；然而，对流动施加了一个附加条件——速度场必须是无旋的，即：

\begin{sourceblock}
\includegraphics[page=30,trim={53.00bp 260.27bp 53.87bp 391.54bp},clip,width=0.92\textwidth,max height=0.38\textheight]{Ferziger4th.pdf}
\end{sourceblock}
根据这个条件，存在速度势，因此速度矢量可以定义为 v = − ∇。不可压缩流的连续性方程 ∇ · v = 0，则变为势能的拉普拉斯方程：

\begin{sourceblock}
\begin{tikzpicture}
\node[inner sep=0,anchor=south west] (image) at (0,0) {\includegraphics[page=30,trim={53.00bp 188.54bp 53.87bp 463.27bp},clip,width=0.92\textwidth,max height=0.38\textheight]{Ferziger4th.pdf}};
\begin{scope}[x={(image.south east)},y={(image.north west)}]
\fill[white] (0.02313,0.91766) rectangle (0.03627,1.00000);
\end{scope}
\end{tikzpicture}
\end{sourceblock}
然后可以对动量方程进行积分，得到伯努利方程，这是一个一旦势能已知就可以求解的代数方程。因此，势流由标量拉普拉斯方程描述。后者无法针对任意几何形状进行解析求解，尽管有简单的解析解（均匀流、源流、汇流、涡流），但也可以将它们组合起来创建更复杂的流，例如围绕圆柱体的流。

对于每个速度势，还可以定义相应的流函数。速度矢量与流线（恒定流函数的线）相切；流线与恒势线正交，因此这些线族形成正交流网。

有时，潜在流量并不十分现实。例如，应用于物体周围流动的势理论会导致达朗贝尔悖论，即物体在势流中既不会受到阻力也不会受到升力（参见，例如，Street et al. 1996，或 Kundu 和 Cohen 2008）。然而，势流理论在多孔介质流动中具有许多应用，并且基于势流理论的计算方法被用于许多领域，例如造船（用于波浪阻力、螺旋桨性能、浮体运动等）。用于预测势流的数值方法通常基于边界元法或面板法（Hess 1990，Kim et al. 2018）；还有针对特定应用开发的特殊方法。这些内容不会在本书中涉及，但有兴趣的读者可以在 Wrobel (2002) 或《Engineering Analysis with Boundary Elements》杂志上找到相关信息。

\subsection*{1.7.4 蠕动（斯托克斯）流}\phantomsection\addcontentsline{toc}{subsection}{1.7.4 蠕动（斯托克斯）流}
当流速非常小时，流体非常粘稠，或者几何尺寸非常小（即雷诺数较小时），纳维-斯托克斯方程中的对流（惯性）项作用较小，可以忽略不计（参见动量方程的无量纲形式，方程（1.30））。然后，流动由粘性力、压力和体积力主导，称为爬行流。如果流体性质可以被认为是恒定的，则动量方程变为线性；它们通常被称为斯托克斯方程。由于速度较低，不稳定项也可以忽略不计，这是一个很大的简化。连续性方程与式（1.34）相同，而动量方程则变为：

\begin{sourceblock}
\includegraphics[page=31,trim={53.00bp 248.79bp 53.87bp 389.70bp},clip,width=0.92\textwidth,max height=0.38\textheight]{Ferziger4th.pdf}
\end{sourceblock}
爬行流存在于多孔介质、涂层技术、微型设备等中。

\subsection*{1.7.5 布辛涅斯克近似}\phantomsection\addcontentsline{toc}{subsection}{1.7.5 布辛涅斯克近似}
在伴随传热的流动中，流体特性通常是温度的函数。这些变化可能很小，但却是流体运动的原因。如果密度变化不大，则可以将非定常项和对流项中的密度视为常数，而仅在引力项中将其视为变量。这称为 Boussinesq 近似。人们通常假设密度随温度线性变化。如果在压力项中包括体积力对平均密度的影响，如第 1 节所述。 1.4，余项可表示为：

\begin{sourceblock}
\includegraphics[page=32,trim={53.00bp 593.89bp 53.87bp 56.79bp},clip,width=0.92\textwidth,max height=0.38\textheight]{Ferziger4th.pdf}
\end{sourceblock}
其中 β 是体积膨胀系数。如果温差低于 2°（水为 2°，空气为 15°），则该近似值会引入 1\% 左右的误差。当温差较大时，误差可能会更大；该解甚至可能在质量上是错误的（例如，参见 Bückle 和 Peri ́c 1992）。

\subsection*{1.7.6 边界层近似}\phantomsection\addcontentsline{toc}{subsection}{1.7.6 边界层近似}
当流动具有主要方向（即没有反向流动或再循环）并且几何形状的变化是渐进的时，流动主要受到上游发生的情况的影响。例如通道和管道中的流动以及平面或轻微弯曲的实体壁上的流动。这种流称为薄剪切层流或边界层流。对于此类流动，纳维-斯托克斯方程可以简化如下：

\begin{itemize}
\item 主要流动方向上的动量扩散远小于对流传输，可以忽略不计。
\item 主流方向上的速度分量远大于其他方向上的分量；
\item 流动方向上的压力梯度远小于主流动方向上的压力梯度。
\end{itemize}
二维边界层方程简化为：

\begin{sourceblock}
\includegraphics[page=32,trim={53.00bp 262.59bp 53.87bp 371.80bp},clip,width=0.92\textwidth,max height=0.38\textheight]{Ferziger4th.pdf}
\end{sourceblock}
必须与连续性方程一起求解；垂直于主流动方向的动量方程简化为 ∂ p / ∂ x 2 = 0。作为 x 1 函数的压力必须通过计算边界层外部的流动（通常假设为势流）来提供，因此边界层方程本身并不是流动的完整描述。可以使用类似于求解具有初始条件的常微分方程的步进技术来求解简化方程。这些技术在空气动力学中得到了广泛的应用。这些方法非常有效，但只能应用于没有分离的流。

\subsection*{1.7.7 复杂流动现象的建模}\phantomsection\addcontentsline{toc}{subsection}{1.7.7 复杂流动现象的建模}
流量，非常重要。由于精确描述通常是不切实际的，因此通常使用半经验模型来表示这些现象。例如湍流模型（将在第 10 章中详细讨论）、燃烧模型、多相模型等。这些模型以及上述简化都会影响解的准确性。各种近似值引入的误差可能会相互增加或抵消；因此，在从使用模型的计算中得出结论时需要小心。由于数值解中各种误差的重要性，我们将重点关注这个主题。将在遇到错误时定义和描述错误类型。

\section*{1.8 流的数学分类}\phantomsection\addcontentsline{toc}{section}{1.8 流的数学分类}
双自变量的拟线性二阶偏微分方程可分为双曲型、抛物型和椭圆型三种类型。这种区别是基于特性和曲线的性质，沿着这些曲线承载有关解的信息。这种类型的每个方程都有两组特征（例如，参见 Street 1973）。

在双曲情况下，特征是真实且明显的。这意味着信息以有限的速度在两组方向上传播。一般来说，信息传播是沿着特定方向进行的，因此每个特征的初始点都需要给出一个数据；因此，这两组特征需要两个初始条件。如果存在横向边界，通常每个点只需要一个条件，因为一个特征是将信息带出域，一个特征将信息带入。但是，这一规则也有例外。

在抛物线方程中，特征退化为单个实数集。因此，通常只需要一个初始条件。在横向边界处，每一点都需要一个条件。

最后，在椭圆的情况下，没有真实的特征；这两组特征是复杂的（想象的）并且截然不同。因此，信息传播没有特殊的方向。事实上，信息在各个方向上的传播基本上同样顺利。一般来说，边界上的每一点都需要一个边界条件，并且解的域通常是封闭的，尽管域的一部分可能延伸到无穷大。不稳态问题从来都不是椭圆形的。

方程性质上的这些差异反映在求解方程的方法上。数值方法应尊重所求解方程的性质，这是一条重要的一般规则。

纳维-斯托克斯方程是四个自变量的非线性二阶方程组。因此，分类方案并不直接适用于它们。尽管如此，纳维-斯托克斯方程确实具有上述许多性质，并且用于求解两个自变量二阶方程的许多思想也适用于它们，但必须小心谨慎。

\subsection*{1.8.1 双曲流}\phantomsection\addcontentsline{toc}{subsection}{1.8.1 双曲流}
首先，考虑不稳定非粘性可压缩流的情况。可压缩流体可以支持声波和激波，这些方程本质上具有双曲线特征也就不足为奇了。大多数用于求解这些方程的方法都是基于方程是双曲的这一思想，并且只要足够小心，它们就可以很好地发挥作用。这些是上面简要提到的方法。

对于稳定的可压缩流，特征取决于流的速度。如果流动是超音速的，则方程是双曲线的，而亚音速流动的方程基本上是椭圆形的。这导致了我们将在下面进一步讨论的困难。

然而，应该指出的是，粘性可压缩流的方程仍然更加复杂。他们的性格是上述所有类型元素的混合体。它们不太适合分类方案，并且它们的数值方法更难以构建。

\subsection*{1.8.2 抛物线流}\phantomsection\addcontentsline{toc}{subsection}{1.8.2 抛物线流}
上面简要描述的边界层近似产生一组基本上具有抛物线特征的方程。信息仅在这些方程中向下游传播，并且可以使用适合抛物线方程的方法来求解。

但请注意，边界层方程需要指定通常通过求解势流问题获得的压力。亚音速势流由椭圆方程控制（在不可压缩极限下，拉普拉斯方程就足够了），因此整个问题实际上具有混合抛物线-椭圆特征。

\subsection*{1.8.3 椭圆流}\phantomsection\addcontentsline{toc}{subsection}{1.8.3 椭圆流}
当流具有再循环区域（即，与流的主方向相反的意义上的流）时，信息可以向上游和下游传播。因此，不能仅在流的上游端应用条件。然后问题就具有了椭圆特征。这种情况发生在亚音速（包括不可压缩）流动中，使得方程的求解成为一项非常困难的任务。

应该指出的是，非稳态不可压缩流实际上具有椭圆形和抛物线形特征的组合。前者源于信息在空间中双向传播的事实，而后者源于信息只能在时间上向前流动的事实。此类问题称为不完全抛物线问题。

\subsection*{1.8.4 混合流类型}\phantomsection\addcontentsline{toc}{subsection}{1.8.4 混合流类型}
正如我们刚才所看到的，单个流可以用不纯粹是一种类型的方程来描述。另一个重要的例子发生在稳态跨音速流中，即同时包含超音速和亚音速区域的稳态可压缩流。超音速区域是双曲形的，而亚音速区域是椭圆形的。因此，可能有必要根据局部流的性质来改变逼近方程的方法。更糟糕的是，在求解方程之前无法确定区域。

\section*{1.9 本书计划}\phantomsection\addcontentsline{toc}{section}{1.9 本书计划}
第 2 章给出了数值求解方法的介绍。讨论了数值方法的优点和缺点，并概述了计算方法的可能性和局限性。接下来是数值求解方法的组成部分及其属性的描述。最后，对基本计算方法（有限差分、有限体积和有限元）进行了简要说明。

第 3 章介绍有限差分 (FD) 方法。在这里，我们提出使用泰勒级数展开和多项式拟合来逼近一阶、二阶和混合导数的方法。讨论了高阶方法的推导以及非线性项和边界的处理。还关注网格不均匀性对截断误差的影响以及离散化误差的估计。这里还简要描述了光谱方法。

第 4 章介绍了有限体积（FV）方法，包括表面和体积积分的近似以及使用插值来获得除单元中心之外的位置处的变量值和导数。还描述了使用延迟校正方法开发高阶方案和简化所得代数方程。特别关注插值和积分近似引起的离散化误差的分析。最后，讨论了各种边界条件的实现。

第 3、4 章描述基本 FD 和 FV 方法，并将其用于结构化笛卡尔网格。这种限制使我们能够将与几何复杂性相关的问题与离散化技术背后的概念分开。复杂几何形状的处理稍后在第 9 章中介绍.

在第 5 章中，我们描述了求解由离散化产生的代数方程组的方法。简要描述了直接方法，但本章的主要部分致力于迭代求解技术。不完全下上分解、共轭梯度和多重网格方法受到特别关注。还描述了解决耦合和非线性系统的方法，包括欠松弛和收敛标准的问题。

第 6 章专门讨论时间积分方法。首先，介绍了求解常微分方程的方法，包括基本方法、预测校正方法和多点方法以及龙格-库塔方法。接下来描述这些方法在非定常输运方程中的应用，包括稳定性和准确性的分析。

第 7、8 章讨论纳维-斯托克斯方程的复杂性和不可压缩流的特殊特征。详细描述了使用分数步和 SIMPLE 算法的交错和同位变量布置、压力方程以及不可压缩流的压力-速度耦合。还描述了其他方法（PISO 算法、流函数涡度、人工压缩性）。足够详细地描述了交错和共置笛卡尔网格的求解方法，以便能够编写计算机代码；这些代码可以在互联网上找到。最后，提出并讨论了使用基于分数步和 SIMPLE 算法的所提供代码计算的稳态和非稳态层流的一些说明性示例，包括迭代和离散化误差的评估。

第9章专门讨论复杂几何形状的处理。讨论了网格类型的选择、复杂几何形状的网格生成方法、网格属性、速度分量和变量排列。重新审视了 FD 和 FV 方法，并讨论了复杂几何形状的特殊特征（如非正交、块结构和非结构网格、非共形网格接口、任意形状的控制体、重叠网格等）。特别关注压力校正方程和边界条件。再次，提出并讨论了使用基于分数步和 SIMPLE 算法提供的代码计算的稳态和非稳态、二维和三维层流的一些说明性示例；还包括离散化误差的评估以及使用不同网格类型（修剪的笛卡尔和任意多面体）获得的结果的比较。

第 10 章讨论湍流的计算。我们讨论湍流的本质及其模拟的三种方法：直接模拟和大涡模拟以及基于雷诺平均纳维-斯托克斯方程的方法。描述了后两种方法中一些广泛使用的模型，包括与边界条件相关的细节。介绍了这些方法的应用示例，包括它们的性能比较。

第 11 章讨论可压缩流。简要讨论了为可压缩流设计的方法。描述了基于分数阶方法和不可压缩流到可压缩流的SIMPLE算法的压力校正方法的扩展。还讨论了处理激波的方法（例如，网格适应、总变差递减和本质无振荡方案）。描述了各种类型可压缩流（亚音速、跨音速和超音速）的边界条件。最后，给出并讨论了应用实例。

第12章专门讨论准确性和效率的提高。首先描述多重网格算法所提供的提高的效率，然后是示例。自适应网格方法和局部网格细化是另一节的主题。最后讨论了并行化。特别关注基于空间和时间域分解的隐式方法的并行处理，并分析并行处理的效率。使用示例计算来证明这些点。

最后，第 13 章讨论若干特殊问题。这些包括共轭传热、自由表面流动、需要移动网格的移动边界的处理、空化和流体-结构相互作用的模拟。简要讨论了传热传质、两相和化学反应流动中的特殊效应。

我们以简短的说明结束这一介绍性章节。计算流体动力学（CFD）可以被视为流体动力学或数值分析的一个子领域。 CFD 能力要求从业者在这两个领域都具有相当扎实的背景。那些在某一领域是专家但认为另一领域没有必要的人会产生糟糕的结果。我们希望读者注意到这一点并获得必要的背景知识。

\chapter{数值方法导论}
\section*{2.1 解决流体动力学问题的方法}\phantomsection\addcontentsline{toc}{section}{2.1 解决流体动力学问题的方法}
正如第一章所述，流体力学方程（已知一个多世纪了）只能求解有限数量的流动。已知的解在帮助理解流体流动方面非常有用，但很少可以直接用于工程分析或设计。工程师传统上被迫使用其他方法。

在最常见的方法中，使用方程的简化。这些通常基于近似值和尺寸分析的组合；几乎总是需要经验输入。例如，量纲分析表明，物体上的阻力可以表示为：

\begin{sourceblock}
\begin{tikzpicture}
\node[inner sep=0,anchor=south west] (image) at (0,0) {\includegraphics[page=38,trim={53.00bp 277.42bp 53.87bp 369.15bp},clip,width=0.92\textwidth,max height=0.38\textheight]{Ferziger4th.pdf}};
\begin{scope}[x={(image.south east)},y={(image.north west)}]
\fill[white] (0.40460,0.07869) rectangle (0.44144,0.71487);
\fill[white] (0.44644,0.13797) rectangle (0.47462,0.72867);
\fill[white] (0.47802,0.07869) rectangle (0.57850,0.80327);
\fill[white] (0.58364,0.13797) rectangle (0.59838,0.72867);
\fill[white] (0.93937,0.11906) rectangle (1.00000,0.70976);
\end{scope}
\end{tikzpicture}
\end{sourceblock}
其中 S 是身体呈现的流动面积，v 是流速，ρ 是流体密度；参数C D 称为阻力系数。它是问题的其他无量纲参数的函数，并且几乎总是通过关联实验数据获得。当系统可以用一两个参数描述时，这种方法非常成功，因此排除了复杂几何形状（只能用许多参数描述）的应用。

通过注意到对于许多流，纳维-斯托克斯方程的无量纲化将雷诺数作为唯一的独立参数，得出了一种相关的方法。如果身体形状保持固定，则可以从具有该形状的比例模型上进行的实验中获得所需的结果。通过仔细选择流体和流动参数或通过雷诺数外推来获得所需的雷诺数；后者可能很危险。这些方法非常有价值，即使在今天也是实际工程设计的主要方法。

问题在于，许多流量需要多个无量纲参数来满足其规格，并且可能无法建立正确缩放实际流量的实验。例如飞机或船舶周围的流动。为了用较小的模型实现相同的雷诺数，必须增加流体速度。对于飞机来说，如果使用相同的流体（空气），这可能会给出过高的马赫数；人们试图找到一种允许两个参数匹配的流体。对于船舶来说，问题是同时匹配雷诺数和弗劳德数，这几乎是不可能的。

在其他情况下，实验即使不是不可能，也是非常困难的。例如，测量设备可能会干扰流动或者流动可能无法接近（例如，晶体生长装置中液体硅的流动）。有些量根本无法用现有技术测量，或者测量精度不够。

实验是测量全局参数（例如阻力、升力、压降或传热系数）的有效方法。在很多情况下，细节很重要；了解是否发生流动分离或壁温是否超过某个极限可能是至关重要的。由于技术改进和竞争需要更仔细地优化设计，或者当新的高科技应用需要预测数据库不足的流量时，实验开发可能过于昂贵和/或耗时。找到合理的替代方案至关重要。

随着电子计算机的诞生，出现了一种替代方法（或者至少是一种补充方法）。尽管偏微分方程数值解法的许多关键思想都是一个多世纪前建立的，但在计算机出现之前它们几乎没有用处。自 20 世纪 50 年代以来，计算机的性能成本比以惊人的速度增长，并且没有任何放缓的迹象。虽然 20 世纪 50 年代制造的第一台计算机每秒只能执行几百次运算，但截至 2017 年 6 月，TOP500 列表 (\url{https://www.top500.org}) 上排名第一的计算机的测量性能峰值为 93 Pflops（petaflops = 每秒 10 15 次浮点运算）；它拥有超过 10 6 个处理核心，内存大小为 1.3 PB（PB = 10 15 字节）。即使笔记本电脑也具有多核芯片和多个处理器，GPU 也可用于完成大规模并行计算（Thibault 和 Senocak 2009；Senocak 和 Jacobsen 2010）。存储数据的能力也显著增强：20 年前只有超级计算机才能找到 10 GB（10GB = 10 10 字节或字符）容量的硬盘，现在笔记本电脑拥有 1 TB 硬盘。智能手机大小的备份磁盘可容纳 500GB 或更多。 1980 年，一台价值数百万美元、占据一个大房间、需要长期维护和操作人员的机器现在可以在笔记本电脑上使用！很难预测未来会发生什么，但可以肯定的是，价格实惠的计算机的计算速度和内存将进一步显著提高。

几乎不需要想象力就能理解计算机使流体流动的研究变得更容易、更有效。一旦人们认识到计算机的这种力量，人们对数值技术的兴趣就会急剧增加。该领域称为计算流体动力学 (CFD)。其中包含许多子专业。几十年来，CFD 已从一个专门的研究领域发展成为一种强大的工具，包含在大学课程中，几乎应用于每个行业，并被研究人员用来研究流体流动的本质。

我们将仅详细讨论解决描述流体流动和相关现象的方程的一小部分方法；将简要提及其他内容，并在适当时给出其他文献的参考文献。

\section*{2.2 什么是CFD？}\phantomsection\addcontentsline{toc}{section}{2.2 什么是CFD？}
正如我们在第一章中看到的。如图1所示，流动和相关现象可以通过偏微分（或积分微分）方程来描述，除非特殊情况，否则不能通过解析求解。为了获得数值上的近似解，我们必须使用离散化方法，通过代数方程组来近似微分方程，然后可以在计算机上求解。近似值应用于空间和/或时间中的小域，因此数值解提供空间和时间中离散位置的结果。正如实验数据的准确性取决于所用工具的质量一样，数值解的准确性也取决于所用离散化的质量。

计算流体动力学的广泛领域中包含的活动涵盖了从完善的工程设计方法的自动化到使用纳维-斯托克斯方程的详细解作为复杂流动性质的实验研究的替代品等一系列活动。一方面，人们可以购买管道系统的设计包，在个人计算机或工作站上在几秒钟或几分钟内解决问题。另一方面，有些代码在最大的超级计算机上可能需要数百小时。该范围与流体力学领域本身一样大，因此不可能在一部著作中涵盖所有 CFD。此外，该领域发展如此之快，以至于我们面临着在短时间内过时的风险。

本书中我们不会讨论自动化的简单方法。基础教科书和本科课程涵盖了它们的基础，并且可用的程序包相对容易理解和使用。

我们将关注旨在求解二维或三维流体运动方程的方法。这些是非标准应用程序中使用的方法，我们指的是在教科书或手册中找不到解（或至少是良好的近似值）的应用程序。虽然这些方法从一开始就被用于高科技工程（例如航空航天），但它们在几何形状复杂或某些重要特征（例如污染物浓度的预测）无法通过标准方法处理的工程领域中得到更频繁的使用。

CFD 已应用于机械、过程、化学、土木和环境工程，并且是大气科学各个方面（从天气预报到气候变化）的主要组成部分。这些领域的优化可以大大节省设备和能源成本，改善洪水和风暴的预测，并减少环境污染。

\section*{2.3 数值方法的可能性和局限性}\phantomsection\addcontentsline{toc}{section}{2.3 数值方法的可能性和局限性}
我们已经注意到与实验工作相关的一些问题。其中一些问题可以通过 CFD 轻松解决。例如，如果我们想在风洞中模拟移动汽车周围的流动，我们需要固定汽车模型并向其吹气，但地板必须以空气速度移动，这是很难做到的。在数值模拟中这并不难做到。其他类型的边界条件很容易在计算中规定；例如，流体的温度或不透明度不会造成问题。如果我们准确地求解非定常三维纳维-斯托克斯方程（如直接模拟湍流），我们就可以获得完整的数据集，从中可以导出任何具有物理意义的量。

这听起来好得令人难以置信。事实上，CFD 的这些优势是以能够准确求解纳维-斯托克斯方程为条件的，而这对于大多数工程兴趣流程来说是极其困难的。我们将在第 1 章中看到。 10 为什么获得高雷诺数流的纳维-斯托克斯方程的精确数值解如此困难。

如果我们无法获得所有流程的准确解，我们必须确定我们可以产生什么，并学会分析和判断结果。首先，我们必须记住，数值结果始终是近似值。计算结果与“现实”之间存在差异是有原因的，即用于产生数值解的过程的每个部分都会产生误差：

\begin{itemize}
\item 微分方程可能包含近似值或理想化，如第1 节中所讨论的。 1.7；
\item 在离散化过程中进行近似；
\item 在求解离散方程时，使用迭代方法。除非运行很长时间，否则不会产生离散方程的精确解。
\end{itemize}
当控制方程准确已知时（例如，不可压缩牛顿流体的纳维-斯托克斯方程），原则上可以实现任何所需精度的解。然而，对于许多现象（例如湍流、燃烧和多相流），精确的方程要么不可用，要么数值求解不可行。这使得引入模型成为必要。即使我们精确地求解方程，该解也不能正确表示现实。为了验证模型，我们必须依赖实验数据。即使可以进行精确的治疗，通常也需要模型来降低成本。

可以通过使用更精确的插值或近似或将近似应用于更小的区域来减少离散化误差，但这通常会增加获得解的时间和成本。通常需要妥协。我们将详细介绍一些方案，但也会指出创建更准确近似值的方法。

求解离散方程时还需要进行折衷。获得准确解的直接求解器很少使用，因为它们成本太高。迭代方法更常见，但需要考虑由于过早停止迭代过程而导致的错误。

本书将强调错误及其估计。我们将给出许多例子的误差估计；分析和估计数值误差的必要性怎么强调都不为过。

使用非定常流的矢量、轮廓或其他类型的绘图或电影（视频）对数值解进行可视化对于结果的解释非常重要。它无疑是解释计算产生的大量数据的最有效方法。然而，存在这样的危险：错误的解可能看起来不错，但可能与实际的边界条件、流体特性等不相符！作者遇到了不正确的数值产生的流动特征，这些特征可能并且已经被解释为物理现象。商业 CFD 代码的工业用户应特别小心，因为推销员的乐观态度是有口皆碑的。精彩的彩色图片给人留下了深刻的印象，但如果在数量上不正确，则毫无价值。结果必须经过非常严格的检验才能被相信。

\section*{2.4 数值求解方法的组成部分}\phantomsection\addcontentsline{toc}{section}{2.4 数值求解方法的组成部分}
考虑到本书不仅适合商业代码的用户，也适合开发新代码的年轻研究人员，因此我们将在这里介绍数值求解方法的重要组成部分。更多细节将在后续章节中介绍。

\subsection*{2.4.1 数学模型}\phantomsection\addcontentsline{toc}{subsection}{2.4.1 数学模型}
任何数值方法的起点都是数学模型，即偏微分或积分微分方程组和边界条件。第 1 章介绍了一些用于流动预测的方程组。 1.人们为目标应用选择合适的模型（不可压缩、无粘性、湍流、二维或三维等）。正如已经提到的，该模型可能包括精确守恒定律的简化。求解方法通常是为一组特定的方程设计的。试图产生一种通用求解方法，即适用于所有流程的方法，即使不是不可能，也是不切实际的，并且与大多数通用工具一样，它们通常对于任何一种应用都不是最佳的。

\subsection*{2.4.2 离散化方法}\phantomsection\addcontentsline{toc}{subsection}{2.4.2 离散化方法}
选择数学模型后，必须选择一种合适的离散化方法，即通过代数方程组对空间和时间上的一组离散位置处的变量来逼近微分方程的方法。

方法有很多，但其中最重要的是：有限差分 (FD)、有限体积 (FV) 和有限元 (FE) 方法。本章最后介绍了这三种离散化方法的重要特征。 CFD 中使用了谱方案、边界元方法和格子玻尔兹曼方法等其他方法，但它们的使用仅限于特殊类别的问题。

如果网格非常精细，每种类型的方法都会产生相同的解。但是，某些方法比其他方法更适合某些类别的问题。偏好往往取决于开发者的态度。稍后我们将讨论各种方法的优缺点。

\subsection*{2.4.3 坐标和基矢量系统}\phantomsection\addcontentsline{toc}{subsection}{2.4.3 坐标和基矢量系统}
章里有提到过。从图 1 可以看出，守恒方程可以写成多种不同的形式，具体取决于所使用的坐标系和基向量。例如，可以选择笛卡尔坐标系、柱坐标系、球坐标系、曲线正交坐标系或非正交坐标系，它们可以是固定的或移动的。选择取决于目标流，并且可能影响要使用的离散化方法和网格类型。

还必须选择定义向量和张量的基础（固定或可变、协变或逆变等）。根据这种选择，速度矢量和应力张量可以用例如笛卡尔、协变或逆变、物理或非物理坐标导向分量来表示。在本书中，我们将专门使用笛卡尔分量，其原因在第 1 章中进行了解释。 9.

\subsection*{2.4.4 数值网格}\phantomsection\addcontentsline{toc}{subsection}{2.4.4 数值网格}
要计算变量的离散位置由数值网格定义，数值网格本质上是要解决问题的几何域的离散表示。它将解域划分为有限数量的子域（元素、控制体积等）。一些可用的选项如下：

\begin{itemize}
\item 结构化（规则）网格——规则或结构化网格由网格线族组成，其属性是单个族的成员不会相互交叉，并且其他族的每个成员仅交叉一次。这允许给定集合的行连续编号。域内任何网格点（或控制体积）的位置由一组两个（2D）或三个（3D）索引唯一标识，例如（ i， j， k ）。这是最简单的网格结构，因为它在逻辑上相当于笛卡尔网格。每个点在二维上有四个最近邻居，在三维上有六个最近邻居；点 P 的每个邻居的索引之一（索引 i , j , k ）不同
\end{itemize}
\begin{sourceblock}
\includegraphics[page=44,trim={53.00bp 489.25bp 53.87bp 52.00bp},clip,width=\textwidth,max height=0.38\textheight]{Ferziger4th.pdf}
\par\vspace{0.45\baselineskip}
\small 图 2.1 二维结构化非正交网格示例，设计用于计算交错管束对称段中的流量
\end{sourceblock}
与 P 的相应索引相差±1。结构化二维网格的示例如图 2.1 所示。这种邻接连通性简化了编程，并且代数方程组的矩阵具有规则结构，可用于开发求解技术。事实上，有大量仅适用于结构化网格的高效求解器（参见第 5 章）。结构化网格的缺点是它们只能用于几何简单的解域。另一个缺点是可能难以控制网格点的分布：出于精度原因将点集中在一个区域中会在解域的其他部分中产生不必要的小间距并浪费资源。这个问题在 3D 几何中被夸大了。又长又细的细胞也可能对会聚产生不利影响。结构化网格可以是H型、O型或C型；这些名称源自求解域边界的拓扑。图 2.1 显示了一个 H 型网格，当映射到矩形时，它具有明显的东、西、北、南边界。在 O 网格中，两个相对的边界相互连接（例如，从东到西，或从南到北）。一个例子是围绕圆柱体的网格（见图 9.22）：如果圆柱体壁是南边界，外边缘是北边界，则西边界和东边界合并以在圆柱体周围创建无尽的网格线（它们可能是圆形的，也可能不是圆形的）。索引 i 的计数从代表东西边界之间的界面的任意径向线开始。网格在该界面上可能是共形的，也可能不是共形的；如何处理这些接口将在第 1 章中描述。 9.在 C 网格中，一个边界部分落在自身上。一个例子是围绕机翼的网格：一组网格线围绕箔片，而另一组则（几乎）与其正交。例如，如果翼型壁是南边界，那么它从后缘延伸到翼片后面的一定距离；出口是从后缘延伸的双南线下方的西边界和上方的东边界，而北边界覆盖解域的底部、左侧和顶部。对于 O 网格，南边界两部分接触的界面可以是共形的，也可以是非共形的。

\begin{sourceblock}
\includegraphics[page=45,trim={53.00bp 487.25bp 53.87bp 52.00bp},clip,width=\textwidth,max height=0.38\textheight]{Ferziger4th.pdf}
\par\vspace{0.45\baselineskip}
\small 图 2.2 具有共形界面的 2D 块结构网格示例
\end{sourceblock}
\begin{sourceblock}
\includegraphics[page=45,trim={53.00bp 342.83bp 53.87bp 192.44bp},clip,width=\textwidth,max height=0.38\textheight]{Ferziger4th.pdf}
\par\vspace{0.45\baselineskip}
\small 图 2.3 具有非共形界面的 2D 块结构网格示例
\end{sourceblock}
\begin{itemize}
\item 块结构网格——在块结构网格中，求解域有两级（或更多级）细分。在粗略层面上，有些块是域中相对较大的部分；它们的结构可能不规则，并且可能重叠也可能不重叠。在精细级别（每个块内）定义了结构化网格。块接口需要特殊处理。第 1 章描述了一些此类方法。 9.图 2.2 显示了具有共形界面的块结构网格；它是针对图 2.1 中使用的相同流动和几何形状而设计的。块界面最初是规则表面（圆柱体或平面段），但通过一些旨在提高网格质量的平滑操作，它们变成了曲面。通过良好的块结构设计，我们可以创建一个高质量的网格，但这需要时间。事实上，为中等复杂的几何形状生成这种类型的高质量网格需要工程师一两周的时间，这并不罕见。图 2.3 显示了具有非共形界面的块结构网格；它与图 2.2 中的网格类似，只是管周围的块中的网格更细。这种网格比以前的网格更灵活，因为它允许在需要更高分辨率的区域（例如，想要准确捕获传热的管道周围）使用更精细的网格。非共形界面可以以完全保守的方式或使用“悬挂节点”来处理，如下所示
\end{itemize}
\begin{sourceblock}
\includegraphics[page=46,trim={53.00bp 487.25bp 53.87bp 52.00bp},clip,width=\textwidth,max height=0.38\textheight]{Ferziger4th.pdf}
\par\vspace{0.45\baselineskip}
\small 图 2.4 复合二维网格，设计用于图 2.1 中使用的相同流动和几何形状
\end{sourceblock}
在第 1 章中讨论过。 9.这里网格平滑只能应用于每个块内；界面通常保持其原始形状，通过比较图 2 和 3 中的网格可以看出。 2.2 和 2.3。编程比上述网格类型更困难。结构化网格的求解器可以按块应用，并且可以使用这些网格处理更复杂的流域。局部细化可以按块进行（即，网格可以在某些块中进行细化）。具有重叠块的块结构网格有时称为复合网格或嵌合网格。图 2.4 显示了这样一个网格。在重叠区域中，一个块的边界条件是通过对另一（重叠）块的解进行插值来获得的。这些网格的缺点是不容易在块边界处实施保护。这种方法的优点是更容易处理复杂的域，并且可以用来跟踪移动的物体：一个块附着在物体上并随之移动，而静止的网格覆盖周围环境。许多作者在 20 世纪 80 年代末和 90 年代初开发了使用重叠结构化网格的方法（Tu 和 Fuchs 1992；Perng 和 Street 1991；Hinatsu 和 Ferziger 1991；Zang 和 Street 1995；Hubbard 和 Chen 1994，1995）。一些半商业代码使用这种方法，例如 SHIPFLOW、CFDShip-Iowa 和 OVERFLOW (NASA)。在世纪之交，人们再次对这种方法产生了浓厚的兴趣，特别是对于处理移动物体周围的流动。 Hadži ́c (2005) 开发了一种这样的方法，并且商业软件 STAR-CCM+ 中提供了适用于任意多面体网格的版本。它将在第 1 章中进行更详细的描述。 9. • 非结构化网格——对于非常复杂的几何形状，最灵活的网格类型是可以适应任意解域边界的网格。原则上，此类网格可以与任何离散化方案一起使用，但它们最适合有限体积和有限元方法。元件或控制体积可以具有任何形状；邻居元素或节点的数量也没有限制。在实践中，最常使用2D中由三角形、四边形或任意多边形组成的网格，以及3D中由四面体、六面体或任意多面体组成的网格。图 2.5 显示了沿壁棱柱层的典型非结构化网格的三个示例。

\begin{sourceblock}
\includegraphics[page=47,trim={53.00bp 257.25bp 53.87bp 52.00bp},clip,width=\textwidth,max height=0.38\textheight]{Ferziger4th.pdf}
\par\vspace{0.45\baselineskip}
\small 图 2.5 非结构化网格的三个示例：四面体（上）、多面体（中）和修剪六面体（下），沿壁有棱柱层并进行局部网格细化
\end{sourceblock}
近年来，由任意多面体控制体组成的网格开始流行，因为它们比四面体网格具有更好的性能，并且比由六面体组成的非结构化网格更容易自动生成。这样的网格可以通过商业网格生成工具自动生成。如果需要，可以优化网格的正交性，轻松控制纵横比，并且可以轻松地局部细化网格。灵活性的优点被数据结构不规则性的缺点所抵消。需要明确指定节点位置和邻居连接。代数方程组的矩阵不再具有规则的对角结构；需要通过对点重新排序来减少带宽。由于使用间接寻址，代数方程系统的求解器通常比结构化网格的求解器慢。例如，这使得使用图形处理单元（GPU）进行计算的效率低于结构化网格的情况；对于矢量处理器的情况也是如此。非结构化网格通常与有限元方法一起使用，并且越来越多地与有限体积方法一起使用。非结构化网格的计算机代码更加灵活。当局部细化网格或使用不同形状的元素或控制体积时，不需要更改它们。随着复杂几何形状的自动网格生成方法的出现，非结构化网格已成为行业中的规则而不是例外。为了保留近壁区域结构化网格的优点，大多数现代网格生成工具还可以沿边界创建棱柱层。然后，网格沿着壁是非结构化的，但是是分层的（结构化的）并且在壁法线方向上几乎正交，从而允许更准确地近似边界层。本书提出的有限体积方法适用于非结构化网格；更多细节将在第 1 章中给出。 9.

本书不会详细介绍网格生成方法。第 1 章简要讨论了网格属性和一些基本的网格生成方法。 9.有大量文献致力于生成块结构和非结构网格，感兴趣的读者可以参考 Thompson 等人的书籍。 ( 1985 ) 和 Arcilla 等人。 （1991）。然而，商业电网生成和优化工具中使用的许多方法并未在公共文献中描述；网格生成在某种程度上是一门艺术，用于处理特殊情况的许多步骤无法用数学方法描述。网格质量问题——这在非结构化网格的情况下尤其重要——将在第 1 章中处理。 12.

\subsection*{2.4.5 有限近似}\phantomsection\addcontentsline{toc}{subsection}{2.4.5 有限近似}
选择网格类型后，必须选择离散化过程中使用的近似值。在有限差分方法中，必须选择网格点处导数的近似值。在有限体积方法中，必须选择近似表面积分和体积积分的方法。在有限元方法中，必须选择形状函数（元素）和加权函数。

有多种可能性可供选择；本书中介绍了一些最常用的方法，有些只是简单提及，还有更多可以创建。该选择会影响近似的准确性。它还影响求解方法的开发、编码、调试的难度以及编码的速度。更准确的近似涉及更多节点并给出更完整的系数矩阵。增加的内存需求可能需要使用更粗糙的网格，部分抵消了更高精度的优势。必须在简单性、易于实施、准确性和计算效率之间做出折衷。本书中介绍的二阶方法是在考虑到这种折衷的情况下选择的。

\subsection*{2.4.6 求解方法}\phantomsection\addcontentsline{toc}{subsection}{2.4.6 求解方法}
离散化产生了一个大型的非线性代数方程组。解决方法取决于问题。对于非定常流动，使用基于常微分方程（时间行进）初值问题的方法。在每个时间步都必须解决一个椭圆问题。稳态流问题通常通过伪时间推进或等效迭代方案来解决。由于方程是非线性的，因此使用迭代方案来求解它们。这些方法使用方程的连续线性化，并且所得线性系统几乎总是通过迭代技术求解。求解器的选择取决于网格类型和每个代数方程中涉及的节点数量。一些求解器将在第 1 章中介绍。 5.

\subsection*{2.4.7 收敛准则}\phantomsection\addcontentsline{toc}{subsection}{2.4.7 收敛准则}
最后，需要设定迭代方法的收敛标准。通常，迭代有两个级别：内部迭代，在其中求解线性方程；以及外部迭代，处理方程的非线性和耦合。从准确性和效率的角度来看，决定何时停止每个级别的迭代过程都很重要。这些问题将在第 4 章中讨论。 5和12。

\section*{2.5 数值解法的性质}\phantomsection\addcontentsline{toc}{section}{2.5 数值解法的性质}
求解方法应具有一定的性质。在大多数情况下，不可能分析完整的解决方法。分析该方法的组成部分；如果组件不具备所需的属性，则完整的方法也不具备所需的属性，但反之则不一定成立。最重要的属性总结如下。

\subsection*{2.5.1 一致性}\phantomsection\addcontentsline{toc}{subsection}{2.5.1 一致性}
当网格间距趋于零时，离散化应该变得精确。离散方程与精确方程之间的差异称为截断误差。通常通过将离散近似中的所有节点值替换为关于单点的泰勒级数展开来估计。结果，我们恢复了原始微分方程以及代表截断误差的余数。为了使方法保持一致，当网格间距 t → 0 和/或 x i → 0 时，截断误差必须变为零。截断误差通常与网格间距 x i 和/或时间步长 t 的幂成正比。如果最重要的项与 ( x ) n 或 ( t ) n 成比例，我们将该方法称为 n 阶近似；为了一致性，需要 n > 0。理想情况下，所有项都应以相同精度的近似值进行离散化；然而，某些项（例如，高雷诺数流动中的对流项或低雷诺数流动中的扩散项）可能在特定流动中占主导地位，并且比其他项更准确地处理它们可能是合理的。

一些离散化方法会导致截断误差，该误差是 x i 与 t 之比的函数，反之亦然。在这种情况下，一致性要求只是有条件地满足：x i 和 t 必须以允许适当比率变为零的方式减少。在接下来的两章中，我们将演示几种离散化方案的一致性。

即使近似一致，也并不一定意味着离散方程组的解在小步长的限制下会成为微分方程的精确解。为此，求解方法必须稳定；这在下面定义。

\subsection*{2.5.2 稳定性}\phantomsection\addcontentsline{toc}{subsection}{2.5.2 稳定性}
如果数值求解过程中出现的误差不被放大，则称数值求解方法是稳定的。对于时间问题，只要精确方程的解有界，稳定性就能保证该方法产生有界解。对于迭代方法，稳定的方法是不发散的方法。稳定性可能很难研究，特别是当存在边界条件和非线性时。因此，通常研究无边​​界条件的常系数线性问题方法的稳定性。经验表明，以这种方式获得的结果通常可以应用于更复杂的问题，但也有明显的例外。

研究数值格式稳定性最广泛使用的方法是冯诺依曼方法。我们将在第 1 章中对一个方案进行简要描述。 6.本书中描述的大多数方案都经过了稳定性分析，我们将在描述每个方案时说明重要的结果。然而，当求解复杂的、非线性的、边界条件复杂的耦合方程时，稳定性结果很少，因此我们可能不得不依赖经验和直觉。许多解要求时间步小于某个限制或使用欠松弛。我们将讨论这些问题，并为选择时间步长和欠松弛参数值提供指导。 6、7和8。

\subsection*{2.5.3 收敛}\phantomsection\addcontentsline{toc}{subsection}{2.5.3 收敛}
如果当网格间距趋于零时，离散方程的解趋于微分方程的精确解，则称数值方法是收敛的。对于线性初值问题，Lax 等价定理（Richtmyer 和 Morton 1967，或 Street 1973）指出“给定一个适当提出的线性初值问题和满足一致性条件的有限差分近似，稳定性是收敛的充分必要条件”。显然，除非求解方法收敛，否则一致的方案是没有用的。

对于受边界条件影响较大的非线性问题，方法的稳定性和收敛性很难证明。因此，通常使用数值实验来检查收敛性，即在一系列连续细化的网格上重复计算。如果该方法是稳定的，并且离散化过程中使用的所有近似值都是一致的，我们通常会发现该解确实收敛到与网格无关的解。对于足够小的网格尺寸，收敛速度由主截断误差分量的阶数控制。这使我们能够估计解中的错误。我们将在章节中详细描述这一点。 3和5。

\subsection*{2.5.4 保护}\phantomsection\addcontentsline{toc}{subsection}{2.5.4 保护}
因为要求解的方程是守恒定律，所以数值格式也应该在局部和全局的基础上遵守这些定律。这意味着，在稳态且没有源的情况下，离开封闭体积的守恒量的量等于进入该体积的量。如果使用强守恒形式的方程和有限体积方法，这对于每个单独的控制体积和整个解域来说都是有保证的。如果谨慎选择近似值，其他离散化方法可以变得保守。源或汇项的处理应一致，以便域中的总源或汇等于通过边界的守恒量的净通量。

这是求解方法的一个重要属性，因为它对求解误差施加了约束。如果确保质量、动量和能量守恒，则误差只能将这些量不正确地分布在解域上。非保守方案可以产生人工源和汇，改变局部和全球的平衡。然而，非保守方案可以是一致和稳定的，因此在非常精细的网格的限制下可以得到正确的解。在大多数情况下，由于不守恒而导致的误差仅在相对粗糙的网格上才明显。问题是很难知道这些误差在哪个网格上足够小。因此，保守的方案是首选。

\subsection*{2.5.5 有界性}\phantomsection\addcontentsline{toc}{subsection}{2.5.5 有界性}
数值解也应该在适当的范围内。物理上的非负量（如密度、湍流动能）必须始终为正；其他量（例如浓度）必须介于 0\% 和 100\% 之间。在没有热源的情况下，某些方程（例如，不存在热源时的温度热方程）要求在域的边界上找到变量的最小值和最大值。这些条件应该由数值近似继承。

有界性很难保证。稍后我们将证明只有一些一阶方案才能保证这个性质。所有高阶方案都可以产生无界解；幸运的是，这种情况通常只发生在过于粗糙的网格上，因此具有下冲和过冲的解通常表明解中的误差很大，并且网格需要进行一些细化（至少是局部细化）。问题在于，容易产生无界解的方案可能存在稳定性和收敛性问题。如果可能的话，应该避免这些方法。

\subsection*{2.5.6 可实现性}\phantomsection\addcontentsline{toc}{subsection}{2.5.6 可实现性}
过于复杂而无法直接处理的现象模型（例如湍流、燃烧或多相流）应设计为保证物理上真实的解。这本身不是一个数值问题，但不可实现的模型可能会导致非物理解或导致数值方法发散。我们不会在本书中讨论这些问题，但如果想在 CFD 代码中实现模型，则必须小心这一属性。

\subsection*{2.5.7 准确度}\phantomsection\addcontentsline{toc}{subsection}{2.5.7 准确度}
流体流动和传热问题的数值解只是近似解。除了求解算法开发过程中可能引入的误差外，在编程或设置边界条件时，数值解总是包含三种系统误差：

\begin{itemize}
\item 建模误差，定义为实际流量与数学模型精确解之间的差异；
\item 离散化误差，定义为守恒方程的精确解与通过离散这些方程获得的代数方程组的精确解之间的差异，以及
\item 迭代误差，定义为代数方程组的迭代解和精确解之间的差异。
\end{itemize}
迭代错误通常称为收敛错误（本书早期版本中就是这种情况）。然而，术语“收敛”不仅与迭代求解方法中的误差减少结合使用，而且还（相当恰当地）经常与数值解向网格无关解的收敛相关联，在这种情况下，它与离散化误差密切相关。为了避免混淆，我们应遵守上述错误定义，并且在讨论收敛问题时，始终指出我们正在讨论哪种类型的收敛。

意识到这些错误的存在很重要，更重要的是尝试将其中一个错误与另一个错误区分开来。各种误差可能会相互抵消，因此有时在粗网格上获得的解可能比在更细网格上获得的解更符合实验——根据定义，细网格上的解应该更准确。

建模误差取决于推导变量传输方程时所做的假设。当研究层流时，它们可能被认为可以忽略不计，因为纳维-斯托克斯方程代表了足够精确的流动模型。然而，对于湍流、两相流、燃烧等，建模误差可能非常大——模型方程的精确解可能在质量上是错误的。通过简化解域的几何形状、简化边界条件等也会引入建模误差。这些误差是先验未知的；它们只能通过将离散化和迭代误差可以忽略不计的解与准确的实验数据或更准确的模型获得的数据（例如来自直接模拟湍流的数据等）进行比较来评估。在判断物理现象模型（如湍流模型）之前，控制和估计迭代和离散化误差至关重要。

我们上面提到，离散化近似会引入误差，误差会随着网格的细化而减小，并且近似的阶数是精度的衡量标准。然而，在给定的网格上，相同阶数的方法可能会产生相差一个数量级的解误差。这是因为阶数仅告诉我们误差随着网格间距减小而减小的速率，它没有提供有关单个网格上的误差的信息。我们将在下一章中展示如何估计离散化误差。

迭代求解和舍入产生的误差更容易控制；我们将在第 1 章中看到如何做到这一点。参见图5，其中引入了迭代求解方法。

解有很多种，CFD 代码的开发人员可能很难决定采用哪一种。最终目标是以最少的努力获得所需的精度，或者利用可用资源获得最大的精度。每次我们描述一个特定的方案时，我们都会指出它相对于这些标准的优点或缺点。

\section*{2.6 离散化方法}\phantomsection\addcontentsline{toc}{section}{2.6 离散化方法}
\subsection*{2.6.1 有限差分法}\phantomsection\addcontentsline{toc}{subsection}{2.6.1 有限差分法}
这是最古老的偏微分方程数值求解方法，据信是由欧拉在 18 世纪引入的。这也是用于简单几何形状的最简单的方法。

起点是微分形式的守恒方程。求解域由网格覆盖。在每个网格点，通过用函数节点值的近似值代替偏导数来近似微分方程。结果是每个网格节点都有一个代数方程，其中该节点和一定数量的邻居节点的变量值显示为未知数。

原则上，FD方法可以应用于任何网格类型。然而，在作者已知的 FD 方法的所有应用中，它都被应用于结构化网格。网格线用作局部坐标线。

泰勒级数展开或多项式拟合用于获得变量相对于坐标的一阶和二阶导数的近似值。必要时，这些方法还用于获取网格节点以外位置的变量值（插值）。第 1 章描述了最广泛使用的通过有限差分逼近导数的方法。 3.

在结构化网格上，FD 方法非常简单且有效。在规则网格上特别容易获得高阶方案；有些会在第 1 章中提到。 3. FD 方法的缺点是除非特别小心，否则不会强制执行守恒。此外，对简单几何形状的限制在复杂流动中也是一个显著的缺点。

\subsection*{2.6.2 有限体积法}\phantomsection\addcontentsline{toc}{subsection}{2.6.2 有限体积法}
FV 方法使用守恒方程的积分形式作为起点。将解域细分为有限数量的连续控制体积 (CV)，并将守恒方程应用于每个 CV。每个 CV 的质心是一个计算节点，在该节点处要计算变量值。插值用于根据节点（CV 中心）值来表达 CV 表面的变量值。使用合适的求积公式来近似表面和体积积分。结果，我们获得了每个 CV 的代数方程，其中出现了许多相邻节点值。

FV方法可以适应任何类型的网格，因此它适用于复杂的几何形状。网格仅定义控制体积边界，不需要与坐标系相关。该方法在构造上是保守的，只要共享边界的 CV 的表面积分（表示对流和扩散通量）相同。

FV 方法可能是最容易理解和编程的。所有需要近似的术语都具有物理意义，这就是它受到工程师欢迎的原因。

与 FD 方案相比，FV 方法的缺点是二阶以上的方法更难以在 3D 中开发。这是因为 FV 方法需要三个级别的近似：插值、微分和积分。我们将在第 1 章详细描述 FV 方法。 4 ;这是本书中最常用的方法。

\subsection*{2.6.3 有限元法}\phantomsection\addcontentsline{toc}{subsection}{2.6.3 有限元法}
FE 方法在很多方面与 FV 方法相似。该域被分解为一组通常是非结构化的离散体积或有限元；在 2D 中，它们通常是三角形或四边形，而在 3D 中最常使用四面体或六面体。有限元方法的显著特征是方程在整个域上积分之前先乘以权重函数。在最简单的有限元方法中，解通过每个单元内的线性形状函数来近似，以保证解在单元边界上的连续性。这样的函数可以根据元素角点处的值构造。权重函数通常具有相同的形式。

然后将该近似代入守恒定律的加权积分中，并要求积分对每个节点值的导数为零，从而导出待解方程；这对应于在一组允许的函数中选择最佳解（残差最小的解）。结果是一组非线性代数方程。

有限元方法的一个重要优点是能够处理任意几何形状；有大量文献致力于构建有限元方法的网格。网格很容易细化；每个元素都被简单地细分。有限元方法相对容易进行数学分析，并且可以证明对于某些类型的方程具有最优性。任何使用非结构化网格的方法都存在一个主要缺点，即线性化方程的矩阵不如规则网格的矩阵那么结构化，这使得找到有效的求解方法变得更加困难。有关有限元方法及其在纳维-斯托克斯方程中的应用的更多详细信息，请参阅 Oden (2006)、Zienkiewicz 等人的书籍。 (2005)、Donea 和 Huerta (2003)、Glowinski 和 Pironneau (1992) 或 Fletcher (1991)。

还应该提到一种称为基于控制体积的有限元法（CVFEM）的混合方法。其中，形状函数用于描述元素上变量的变化。通过连接单元的质心，在每个节点周围形成控制体积。积分形式的守恒方程以与有限体积法相同的方式应用于这些 CV。通过 CV 边界的通量和源项按元素计算。我们将在第 1 章中对这种方法进行简短的描述。 9.

\chapter{有限差分法}
\section*{3.1 简介}\phantomsection\addcontentsline{toc}{section}{3.1 简介}
正如第一章中提到的。如图1所示，所有守恒方程都具有相似的结构，并且可以被视为通用输运方程Eq.(1.26)、(1.27)或(1.28)的特例。因此，在本章和后续章节中，我们将仅处理一个通用的守恒方程。它将用于演示所有守恒方程（对流、扩散和源）通用项的离散化方法。纳维-斯托克斯方程的特点和求解耦合非线性问题的技术将在后面介绍。另外，目前不稳定项将被删除，因此我们只考虑与时间无关的问题。

为简单起见，我们此时将仅使用笛卡尔网格。我们要处理的方程是：

\begin{sourceblock}
\begin{tikzpicture}
\node[inner sep=0,anchor=south west] (image) at (0,0) {\includegraphics[page=56,trim={53.00bp 267.28bp 53.87bp 369.75bp},clip,width=0.92\textwidth,max height=0.38\textheight]{Ferziger4th.pdf}};
\begin{scope}[x={(image.south east)},y={(image.north west)}]
\fill[white] (0.00000,0.83889) rectangle (0.06072,1.00000);
\fill[white] (0.06343,0.83889) rectangle (0.11814,1.00000);
\fill[white] (0.12084,0.83889) rectangle (0.17892,1.00000);
\fill[white] (0.18159,0.83889) rectangle (0.21471,1.00000);
\fill[white] (0.29765,0.55239) rectangle (0.36289,0.95878);
\fill[white] (0.36310,0.52422) rectangle (0.40725,0.95878);
\fill[white] (0.46355,0.56132) rectangle (0.48208,0.95843);
\fill[white] (0.41435,0.32120) rectangle (0.44253,0.71865);
\fill[white] (0.32809,0.06905) rectangle (0.36316,0.47544);
\fill[white] (0.36262,0.04122) rectangle (0.37326,0.33562);
\fill[white] (0.44965,0.06905) rectangle (0.48472,0.47544);
\fill[white] (0.48418,0.04122) rectangle (0.49483,0.33562);
\end{scope}
\end{tikzpicture}
\end{sourceblock}
我们假设 ρ 、 u j 和 q φ 是已知的。情况可能并非如此，因为速度可能尚未计算出来，并且流体的属性可能取决于温度，如果使用湍流模型，还可能取决于速度场。正如我们将看到的，用于求解这些方程的迭代方案将 φ 视为唯一的未知数；所有其他变量都固定为前一次迭代确定的值，因此将这些变量视为已知是一种合理的方法。

非正交和非结构化网格的特殊功能将在第 1 章中讨论。 9.此外，在许多可能的离散化技术中，将只描述选定的几种能够说明主要思想的技术；其他的可以在引用的文献中找到。

\section*{3.2 基本概念}\phantomsection\addcontentsline{toc}{section}{3.2 基本概念}
获得数值解的第一步是离散化几何域，即必须定义数值网格。在有限差分（FD）离散化方法中，网格通常是局部结构化的，即每个网格节点可以被认为是局部坐标系的原点，其轴与网格线重合。这也意味着属于同一族的两条网格线（例如 xi 1 ）不相交，并且属于不同族的任何一对网格线（例如 Σ 1 = const ）。且 xi 2 = const.，仅相交一次。在三维空间中，三个网格线在每个节点处相交；这些线均不会在任何其他点相交。图 3.1 显示了 FD 方法中使用的一维 (1D) 和二维 (2D) 笛卡尔网格的示例。

每个节点由一组索引唯一标识，这些索引是与其相交的网格线的索引，2D 中为 ( i , j )，3D 中为 ( i , j , k )。邻居节点是通过统一增加或减少其中一个索引来定义的。

微分形式的通用标量守恒方程 (3.1) 是 FD 方法的起点。由于它在 φ 中是线性的，因此将通过线性代数方程组来近似，其中网格节点处的变量值是未知数。该方程组的解近似于偏微分方程（PDE）的解。

因此，每个节点都有一个与之相关的未知变量值，并且必须提供一个代数方程。后者是该节点的变量值与某些相邻节点的变量值之间的关系。它是通过将特定节点处的偏微分方程的每一项替换为有限差分近似来获得的。当然，方程和未知数的数量必须相等。在边界处

\begin{sourceblock}
\includegraphics[page=57,trim={53.00bp 102.30bp 53.87bp 363.03bp},clip,width=\textwidth,max height=0.38\textheight]{Ferziger4th.pdf}
\par\vspace{0.45\baselineskip}
\small 图 3.1 FD 方法的 1D（上）和 2D（下）笛卡尔网格示例（实心符号表示边界节点，空心符号表示内部计算节点）
\end{sourceblock}
\begin{center}\small 图3.2 关于导数的定义及其近似值\end{center}
\begin{sourceblock}
\begin{tikzpicture}
\node[inner sep=0,anchor=south west] (image) at (0,0) {\includegraphics[page=58,trim={53.00bp 588.12bp 53.87bp 65.25bp},clip,width=0.92\textwidth,max height=0.38\textheight]{Ferziger4th.pdf}};
\begin{scope}[x={(image.south east)},y={(image.north west)}]
\fill[white] (0.00000,0.73062) rectangle (0.04547,1.00000);
\fill[white] (0.04704,0.73062) rectangle (0.08367,1.00000);
\fill[white] (0.09657,0.72357) rectangle (0.13251,1.00000);
\fill[white] (0.13408,0.72357) rectangle (0.17002,1.00000);
\fill[white] (0.17155,0.72357) rectangle (0.27398,1.00000);
\fill[white] (0.27555,0.72357) rectangle (0.30156,1.00000);
\fill[white] (0.00000,0.00000) rectangle (0.01546,0.73218);
\fill[white] (0.01699,0.00000) rectangle (0.12132,0.73218);
\fill[white] (0.12292,0.00000) rectangle (0.16451,0.73218);
\fill[white] (0.16608,0.00000) rectangle (0.19495,0.73218);
\fill[white] (0.54574,0.11433) rectangle (0.61146,0.87392);
\fill[white] (0.75862,0.11433) rectangle (0.87038,0.87392);
\end{scope}
\end{tikzpicture}
\end{sourceblock}
中央

向前

\begin{sourceblock}
\begin{tikzpicture}
\node[inner sep=0,anchor=south west] (image) at (0,0) {\includegraphics[page=58,trim={53.00bp 513.25bp 53.87bp 138.58bp},clip,width=0.92\textwidth,max height=0.38\textheight]{Ferziger4th.pdf}};
\begin{scope}[x={(image.south east)},y={(image.north west)}]
\fill[white] (0.63877,0.19147) rectangle (0.65973,0.91614);
\fill[white] (0.65889,0.17470) rectangle (0.67543,0.89937);
\fill[white] (0.67233,0.08386) rectangle (0.68223,0.61985);
\fill[white] (0.71453,0.15234) rectangle (0.74926,0.89378);
\fill[white] (0.74617,0.08386) rectangle (0.77750,0.61985);
\end{scope}
\end{tikzpicture}
\end{sourceblock}
x 2 − i 1 − i i i+1 i+2
给定变量值的节点（狄利克雷条件），不需要方程。当边界条件涉及导数时（如诺依曼条件），必须对边界条件进行离散化，以便为必须求解的集合提供一个方程。

有限差分近似背后的思想直接借用自导数的定义：

\begin{sourceblock}
\begin{tikzpicture}
\node[inner sep=0,anchor=south west] (image) at (0,0) {\includegraphics[page=58,trim={53.00bp 358.31bp 53.87bp 269.88bp},clip,width=0.92\textwidth,max height=0.38\textheight]{Ferziger4th.pdf}};
\begin{scope}[x={(image.south east)},y={(image.north west)}]
\fill[white] (0.28442,0.47931) rectangle (0.32268,0.78393);
\fill[white] (0.34370,0.03162) rectangle (0.36250,0.27325);
\end{scope}
\end{tikzpicture}
\end{sourceblock}
我们将经常参考图 3.2 中显示的几何解释。某一点的一阶导数 ∂φ/∂ x 是该点处曲线 φ( x ) 的切线斜率，即图中标记为“Exact”的线。其斜率可以通过穿过曲线上两个邻近点的直线的斜率来近似。虚线表示前向差值的近似值； x i 处的导数由穿过点 x i 和 x i + x 处的另一个点的直线的斜率来近似。虚线说明了通过向后差分进行的近似，其中第二个点是 x i − x。标记为“中心”的线表示通过中心差进行近似：它使用穿过位于导数近似点相对侧的两个点的线的斜率。

从图 3.2 中可以明显看出，某些近似值比其他近似值更好。中心差近似线的斜率非常接近精确线的斜率；如果函数 φ( x ) 是二阶多项式并且点在 x 方向上等距分布，则斜率将完全匹配。

从图 3.2 中还可以明显看出，当附加点接近 x i 时，近似的质量会提高，即，随着网格的细化，近似的质量会提高。图 3.2 中所示的近似值是多种可能性中的几种；以下各节概述了推导一阶和二阶导数近似值的主要方法。

在以下两节中，仅考虑一维情况。坐标可以是笛卡尔坐标或曲线坐标，这里的差异并不重要。在多维有限差分中，每个坐标通常单独处理，因此这里开发的方法很容易适应更高的维度。

Fornberg (1988) 提供了差分公式的通用方法以及各种阶导数和精度阶数的导数表达式的有用表。

\section*{3.3 一阶导数的近似}\phantomsection\addcontentsline{toc}{section}{3.3 一阶导数的近似}
方程（3.1）中对流项的离散化需要一阶导数 ∂(ρ u φ)/∂ x 的近似。我们现在将描述一些通用变量 φ 一阶导数的近似方法；这些方法可以应用于任何量的一阶导数。

在上一节中，介绍了一种推导一阶导数近似值的方法。有更系统的方法更适合推导更准确的近似值；其中一些将在稍后描述。

\subsection*{3.3.1 泰勒级数展开}\phantomsection\addcontentsline{toc}{subsection}{3.3.1 泰勒级数展开}
\begin{sourceblock}
\begin{tikzpicture}
\node[inner sep=0,anchor=south west] (image) at (0,0) {\includegraphics[page=59,trim={53.00bp 263.52bp 53.87bp 336.15bp},clip,width=0.92\textwidth,max height=0.38\textheight]{Ferziger4th.pdf}};
\begin{scope}[x={(image.south east)},y={(image.north west)}]
\fill[white] (0.45624,0.70257) rectangle (0.49447,0.87649);
\fill[white] (0.12370,0.59350) rectangle (0.18764,0.77148);
\fill[white] (0.19182,0.59756) rectangle (0.22000,0.77148);
\fill[white] (0.22418,0.58131) rectangle (0.29552,0.77148);
\fill[white] (0.29738,0.59756) rectangle (0.32556,0.77148);
\fill[white] (0.32740,0.59350) rectangle (0.35750,0.77148);
\fill[white] (0.36202,0.58131) rectangle (0.43221,0.77148);
\fill[white] (0.22779,0.26057) rectangle (0.25789,0.43854);
\fill[white] (0.26238,0.24823) rectangle (0.34310,0.45299);
\fill[white] (0.27215,0.04754) rectangle (0.30021,0.22702);
\end{scope}
\end{tikzpicture}
\end{sourceblock}
其中 H 表示“高阶项”。通过将该方程中的 x 替换为 x i + 1 或 x i − 1，我们可以得到这些点处的变量值的表达式，即变量及其在 x i 处的导数。这可以扩展到 x i 附近的任何其他点，例如 x i + 2 和 x i − 2。

使用这些展开式，我们可以根据相邻点的函数值获得点 x i 的一阶和高阶导数的近似表达式。例如，对 x i + 1 处的 φ 使用方程（3.3），我们可以证明：

\begin{sourceblock}
\begin{tikzpicture}
\node[inner sep=0,anchor=south west] (image) at (0,0) {\includegraphics[page=59,trim={53.00bp 131.91bp 53.87bp 496.87bp},clip,width=0.92\textwidth,max height=0.38\textheight]{Ferziger4th.pdf}};
\begin{scope}[x={(image.south east)},y={(image.north west)}]
\fill[white] (0.04583,0.47082) rectangle (0.08406,0.78025);
\fill[white] (0.76126,0.47082) rectangle (0.81146,0.80621);
\fill[white] (0.04683,0.08753) rectangle (0.08189,0.40418);
\fill[white] (0.64818,0.08458) rectangle (0.66797,0.39427);
\end{scope}
\end{tikzpicture}
\end{sourceblock}
可以使用 x i − 1 处的级数表达式 ( 3.3 ) 导出另一个表达式：

\begin{sourceblock}
\begin{tikzpicture}
\node[inner sep=0,anchor=south west] (image) at (0,0) {\includegraphics[page=59,trim={53.00bp 72.11bp 53.87bp 556.67bp},clip,width=0.92\textwidth,max height=0.38\textheight]{Ferziger4th.pdf}};
\begin{scope}[x={(image.south east)},y={(image.north west)}]
\fill[white] (0.04586,0.47109) rectangle (0.08409,0.78051);
\fill[white] (0.76129,0.47082) rectangle (0.81149,0.80648);
\fill[white] (0.04686,0.08753) rectangle (0.08192,0.40418);
\fill[white] (0.64821,0.08485) rectangle (0.66800,0.39427);
\end{scope}
\end{tikzpicture}
\end{sourceblock}
在 x i − 1 和 x i + 1 处使用方程（3.3）可以得到另一个表达式：

\begin{sourceblock}
\begin{tikzpicture}
\node[inner sep=0,anchor=south west] (image) at (0,0) {\includegraphics[page=60,trim={53.00bp 526.53bp 53.87bp 73.14bp},clip,width=0.92\textwidth,max height=0.38\textheight]{Ferziger4th.pdf}};
\begin{scope}[x={(image.south east)},y={(image.north west)}]
\fill[white] (0.12313,0.70272) rectangle (0.16135,0.87664);
\fill[white] (0.12412,0.48714) rectangle (0.15916,0.66511);
\fill[white] (0.23597,0.24718) rectangle (0.30150,0.43854);
\fill[white] (0.30487,0.24823) rectangle (0.38556,0.45314);
\fill[white] (0.38890,0.26448) rectangle (0.41708,0.43854);
\fill[white] (0.41895,0.24823) rectangle (0.45558,0.43854);
\fill[white] (0.46096,0.24718) rectangle (0.56851,0.45314);
\end{scope}
\end{tikzpicture}
\end{sourceblock}
如果保留右侧的所有项，则所有这三个表达式都是精确的。由于高阶导数是未知的，因此这些表达式本身没有多大价值。然而，如果网格点之间的距离，即 x i − x i − 1 和 x i + 1 − x i 之间的距离较小，则高阶项将很小，除非在高阶导数局部非常大的特殊情况下。忽略后一种可能性，一阶导数的近似值是通过截断右侧第一项之后的每个级数得出的：

\begin{sourceblock}
\begin{tikzpicture}
\node[inner sep=0,anchor=south west] (image) at (0,0) {\includegraphics[page=60,trim={53.00bp 321.14bp 53.87bp 233.01bp},clip,width=0.92\textwidth,max height=0.38\textheight]{Ferziger4th.pdf}};
\begin{scope}[x={(image.south east)},y={(image.north west)}]
\fill[white] (0.38710,0.82347) rectangle (0.42532,0.92678);
\fill[white] (0.38710,0.49031) rectangle (0.42532,0.59354);
\fill[white] (0.37489,0.15716) rectangle (0.41311,0.26038);
\end{scope}
\end{tikzpicture}
\end{sourceblock}
这些分别是前面提到的前向（FDS）、后向（BDS）和中心差分（CDS）方案。从右侧删除的项称为截断错误；它们测量近似的准确性，并确定误差随着点间距减小而减小的速率。特别是，第一个截断项通常是误差的主要来源。

截断误差是点之间的间距的幂与点 x = x i 处的高阶导数的乘积之和：

\begin{sourceblock}
\includegraphics[page=60,trim={53.00bp 192.00bp 53.87bp 456.97bp},clip,width=0.92\textwidth,max height=0.38\textheight]{Ferziger4th.pdf}
\end{sourceblock}
其中 x 是点之间的间距（假设当前全部相等），α 是高阶导数乘以常数因子。从方程（3.10）我们看到，对于小间距，包含较高 x 次方的项较小，因此首项（指数最小的项）是主导项。随着 x 减小，上述近似值收敛到精确导数，其误差与 ( x ) m 成比例，其中 m 是前导截断误差项的指数。近似阶数表示网格细化时误差减少的速度；它并不表示误差的绝对大小。因此，对于一阶、二阶、三阶或四阶近似，误差分别减少了两倍、四倍、八倍或十六倍。应该记住，该规则仅对足够小的间距有效； “足够小”的定义取决于函数 φ( x ) 的轮廓。

方程（3.7）和（3.8）表示一阶近似，无论网格间距是均匀的还是非均匀的，因为截断误差中的首项与网格间距成正比（参见方程（3.4）和（3.5））。当网格间距均匀时，式(3.9)的首项消失；剩余的首项与网格间距的平方成正比，因此该方案是二阶精确的。关于网格不均匀性对截断误差的影响的更详细讨论将在第 4 节中给出。 3.3.4.

\subsection*{3.3.2 多项式拟合}\phantomsection\addcontentsline{toc}{subsection}{3.3.2 多项式拟合}
获得导数近似值的另一种方法是将函数拟合到插值曲线并对所得曲线进行微分。例如，如果使用分段线性插值，我们将获得 FDS 或 BDS 近似值，具体取决于第二个点位于点 x i 的左侧还是右侧。

将抛物线拟合到点 x i − 1 、 x i 和 x i + 1 处的数据，并根据插值计算 x i 处的一阶导数，我们得到：

\begin{sourceblock}
\begin{tikzpicture}
\node[inner sep=0,anchor=south west] (image) at (0,0) {\includegraphics[page=61,trim={53.00bp 316.54bp 53.87bp 312.24bp},clip,width=0.92\textwidth,max height=0.38\textheight]{Ferziger4th.pdf}};
\begin{scope}[x={(image.south east)},y={(image.north west)}]
\fill[white] (0.08565,0.47109) rectangle (0.12388,0.78051);
\end{scope}
\end{tikzpicture}
\end{sourceblock}
其中 x i = x i − x i − 1。该近似在任何网格上都存在二阶截断误差。使用泰勒级数方法，通过消除式（3.6）中包含二阶导数的项，可以得到相同的二阶近似；见式(3.26)。对于均匀间距，上述表达式简化为等式（3.9）中给出的 CDS 近似。

其他多项式、样条曲线等可以用作插值，然后逼近导数。一般来说，一阶导数的近似具有与用于近似函数的多项式的次数相同量级的截断误差。我们给出下面两个通过将三次多项式拟合到四个点获得的三阶近似值和通过在均匀网格上拟合四到五次多项式获得的四阶近似值：

\begin{sourceblock}
\begin{tikzpicture}
\node[inner sep=0,anchor=south west] (image) at (0,0) {\includegraphics[page=61,trim={53.00bp 100.66bp 53.87bp 490.82bp},clip,width=0.92\textwidth,max height=0.38\textheight]{Ferziger4th.pdf}};
\begin{scope}[x={(image.south east)},y={(image.north west)}]
\fill[white] (0.18662,0.73533) rectangle (0.22484,0.89017);
\fill[white] (0.18761,0.54340) rectangle (0.22268,0.70185);
\fill[white] (0.17492,0.23547) rectangle (0.21314,0.39044);
\fill[white] (0.17591,0.04366) rectangle (0.21098,0.20212);
\end{scope}
\end{tikzpicture}
\end{sourceblock}
\begin{sourceblock}
\begin{tikzpicture}
\node[inner sep=0,anchor=south west] (image) at (0,0) {\includegraphics[page=62,trim={53.00bp 581.49bp 53.87bp 50.00bp},clip,width=0.92\textwidth,max height=0.38\textheight]{Ferziger4th.pdf}};
\begin{scope}[x={(image.south east)},y={(image.north west)}]
\fill[white] (0.13648,0.50794) rectangle (0.17471,0.84156);
\fill[white] (0.13747,0.09437) rectangle (0.17254,0.43579);
\end{scope}
\end{tikzpicture}
\end{sourceblock}
上述近似分别是三阶BDS、三阶FDS和四阶CDS方案。在非均匀网格上，上述表达式中的系数成为网格扩展比率的函数。

就 FDS 和 BDS 而言，对近似的主要贡献来自于一侧。在对流主导的问题中，当流动局部从节点 x i − 1 到 x i 时，有时使用 BDS；当流动沿负方向时，有时使用 FDS。此类方法称为迎风方案 (UDS)。一阶迎风方案非常不准确；它们的截断误差具有虚假扩散的效果（即，解对应于较大的扩散系数，有时比实际扩散率大得多）。高阶迎风方案更准确，但通常可以用更少的努力实现更高阶的 CDS，因为不需要检查流向（参见上面的表达式）。

我们在这里仅演示了一维多项式拟合；类似的方法可以与任何类型的形状函数或一维、二维或三维插值一起使用。唯一的限制是显而易见的，即用于计算形状函数系数的网格点的数量必须等于可用系数的数量。当使用不规则网格时，这种方法很有吸引力，因为它可以避免使用坐标变换；见第 9.5 节.

\subsection*{3.3.3 紧凑方案}\phantomsection\addcontentsline{toc}{subsection}{3.3.3 紧凑方案}
对于均匀间隔的网格，可以导出许多特殊的方案。其中包括紧凑方案（Lele 1992；Mahesh 1998）和稍后描述的谱方法。这里，仅描述Padé方案。

可以通过使用多项式拟合导出紧凑方案。然而，不是仅使用计算节点处的变量值来导出多项式的系数，而是还使用某些点处的导数值。我们将利用这个想法来推导四阶 Padé 方案。目标是仅使用来自近邻点的信息；这使得结果方程的求解更加简单，并降低了在域边界附近找到近似值的难度。在这里描述的特定方案中，我们将使用节点 i 、 i + 1 和 i − 1 处的变量值以及节点 i + 1 和 i − 1 处的一阶导数，以获得节点 i 处的一阶导数的近似值。为此，在节点 i 附近定义了一个四次多项式：

\begin{sourceblock}
\includegraphics[page=62,trim={53.00bp 105.64bp 53.87bp 540.92bp},clip,width=0.92\textwidth,max height=0.38\textheight]{Ferziger4th.pdf}
\end{sourceblock}
系数a 0 , ..., a 4 可以通过将上述多项式拟合到三个变量和两个导数值来找到。然而，由于我们只对节点 i 处的一阶导数感兴趣，因此我们只需要计算系数 a 1。对式（3.15）微分，可得：

\begin{sourceblock}
\begin{tikzpicture}
\node[inner sep=0,anchor=south west] (image) at (0,0) {\includegraphics[page=63,trim={53.00bp 503.04bp 53.87bp 90.53bp},clip,width=0.92\textwidth,max height=0.38\textheight]{Ferziger4th.pdf}};
\begin{scope}[x={(image.south east)},y={(image.north west)}]
\fill[white] (0.00000,0.35387) rectangle (0.03080,0.51316);
\fill[white] (0.03344,0.35387) rectangle (0.08322,0.51316);
\fill[white] (0.43672,0.24239) rectangle (0.47495,0.40168);
\end{scope}
\end{tikzpicture}
\end{sourceblock}
通过对 x = x i 、 x = x i + 1 和 x = x i − 1 写出方程（ 3.15 ），对 x = x i + 1 和 x = x i − 1 写出方程（ 3.16 ），经过一些重新排列后我们得到：

\begin{sourceblock}
\begin{tikzpicture}
\node[inner sep=0,anchor=south west] (image) at (0,0) {\includegraphics[page=63,trim={53.00bp 431.78bp 53.87bp 196.93bp},clip,width=0.92\textwidth,max height=0.38\textheight]{Ferziger4th.pdf}};
\begin{scope}[x={(image.south east)},y={(image.north west)}]
\fill[white] (0.13338,0.47181) rectangle (0.17164,0.78092);
\fill[white] (0.13441,0.08923) rectangle (0.16944,0.40529);
\fill[white] (0.49931,0.08923) rectangle (0.53435,0.40529);
\end{scope}
\end{tikzpicture}
\end{sourceblock}
如果将节点 i + 2 和 i - 2 处的变量值相加，则可以使用六次多项式；如果还使用这两个节点处的导数，则可以使用八次多项式之一。可以在每个点处写出像方程（3.18）这样的方程。完整的方程组实际上是网格点导数的三对角方程组。为了计算导数，必须求解该系统。

高达六阶的紧致中心近似族可以写成：

\begin{sourceblock}
\begin{tikzpicture}
\node[inner sep=0,anchor=south west] (image) at (0,0) {\includegraphics[page=63,trim={53.00bp 312.43bp 53.87bp 316.29bp},clip,width=0.92\textwidth,max height=0.38\textheight]{Ferziger4th.pdf}};
\begin{scope}[x={(image.south east)},y={(image.north west)}]
\fill[white] (0.06776,0.47167) rectangle (0.10598,0.78087);
\fill[white] (0.02063,0.28514) rectangle (0.04256,0.59407);
\fill[white] (0.06875,0.08899) rectangle (0.10382,0.40513);
\end{scope}
\end{tikzpicture}
\end{sourceblock}
根据参数α、β、γ的选择，得到二阶、四阶CDS、四阶、六阶Padé格式；参数及相应的截断误差如表3.1所示。

显然，对于相同阶的近似，Padé 方案使用更少的计算节点，因此比中心差分近似具有更紧凑的计算分子。如果所有网格点的变量值已知，我们

\begin{sourceblock}
\begin{tikzpicture}
\node[inner sep=0,anchor=south west] (image) at (0,0) {\includegraphics[page=63,trim={53.00bp 97.06bp 53.87bp 456.27bp},clip,width=0.92\textwidth,max height=0.38\textheight]{Ferziger4th.pdf}};
\begin{scope}[x={(image.south east)},y={(image.north west)}]
\fill[white] (0.00000,0.92846) rectangle (0.06406,1.00000);
\fill[white] (0.06529,0.92846) rectangle (0.10192,1.00000);
\fill[white] (0.11483,0.92767) rectangle (0.21161,1.00000);
\fill[white] (0.21287,0.92767) rectangle (0.31104,1.00000);
\fill[white] (0.37441,0.89735) rectangle (0.45985,0.98661);
\fill[white] (0.48740,0.89735) rectangle (0.60304,0.98661);
\fill[white] (0.60460,0.89735) rectangle (0.65889,0.98661);
\fill[white] (0.68529,0.90019) rectangle (0.70466,0.98936);
\fill[white] (0.79230,0.90019) rectangle (0.81149,0.98936);
\fill[white] (0.89934,0.90019) rectangle (0.91693,0.98936);
\fill[white] (0.00000,0.83938) rectangle (0.03528,0.92864);
\fill[white] (0.03681,0.83938) rectangle (0.15335,0.92864);
\fill[white] (0.15489,0.83938) rectangle (0.19648,0.92864);
\fill[white] (0.19805,0.83938) rectangle (0.30614,0.92864);
\fill[white] (0.92280,0.44703) rectangle (0.94033,0.53630);
\end{scope}
\end{tikzpicture}
\end{sourceblock}
可以通过求解三对角系统来计算网格线上所有节点的导数（有关如何完成此操作的详细信息，请参阅第 5 章）。我们将在该节中看到。 5.6，这些方案也可以应用于隐式方法。这个问题将在第二节中再次讨论。 3.7.

这里导出的方案只是几种可能性；扩展到更高阶和多维近似是可能的。还可以导出非均匀网格的方案，但系数是特定于网格的（参见例如 Gamet 等人 1999）。

\subsection*{3.3.4 非均匀网格}\phantomsection\addcontentsline{toc}{subsection}{3.3.4 非均匀网格}
由于截断误差不仅取决于网格间距，还取决于变量的导数，因此我们无法在均匀网格上实现离散化误差的均匀分布。因此我们需要使用非均匀网格。这个想法是在函数导数较大的区域使用较小的 x，在函数平滑的区域使用较大的 x。通过这种方式，应该可以将误差几乎均匀地分布在整个域上，从而对于给定数量的网格点获得更好的解。在本节中，我们将讨论非均匀网格上有限差分近似的准确性。

在某些近似中，当点间距均匀时，截断误差表达式中的首项变为零，即 x i + 1 − x i = x i − x i − 1 = x。 CDS 近似就是这种情况，参见方程（3.6）。尽管对于非均匀间距，不同的近似在形式上具有相同的阶数，但它们不具有相同的截断误差。此外，当 CDS 应用于非均匀网格时，细化网格时误差减小的速率不会恶化，正如我们现在将要展示的那样。

为了证明这一点，文献中存在一些混淆，请注意 CDS 的截断误差为（比较方程（3.9）和（3.6））：

\begin{sourceblock}
\begin{tikzpicture}
\node[inner sep=0,anchor=south west] (image) at (0,0) {\includegraphics[page=64,trim={53.00bp 208.95bp 53.87bp 419.84bp},clip,width=0.92\textwidth,max height=0.38\textheight]{Ferziger4th.pdf}};
\begin{scope}[x={(image.south east)},y={(image.north west)}]
\fill[white] (0.17826,0.44016) rectangle (0.25495,0.80616);
\fill[white] (0.25841,0.47095) rectangle (0.30442,0.78046);
\fill[white] (0.32451,0.44203) rectangle (0.37435,0.80616);
\fill[white] (0.04899,0.25997) rectangle (0.07600,0.59357);
\fill[white] (0.08352,0.28407) rectangle (0.11170,0.59357);
\fill[white] (0.11522,0.28407) rectangle (0.15817,0.78046);
\fill[white] (0.15537,0.08461) rectangle (0.17516,0.39438);
\fill[white] (0.17549,0.09451) rectangle (0.19146,0.40428);
\fill[white] (0.21158,0.06372) rectangle (0.26508,0.39705);
\fill[white] (0.26845,0.09451) rectangle (0.29663,0.40428);
\fill[white] (0.32340,0.06560) rectangle (0.36277,0.40428);
\fill[white] (0.92421,0.00000) rectangle (1.00000,0.11486);
\end{scope}
\end{tikzpicture}
\end{sourceblock}
(3.20) 我们使用了符号（见图 3.2）：

\begin{sourceblock}
\includegraphics[page=64,trim={53.00bp 163.48bp 53.87bp 487.20bp},clip,width=0.92\textwidth,max height=0.38\textheight]{Ferziger4th.pdf}
\end{sourceblock}
首项与 x 成正比，但当 x i + 1 = x i 时变为零。这意味着网格间距越不均匀，误差就越大。

让我们假设网格以常数因子 re 扩展或收缩。这称为复利网格；对于它：

\begin{sourceblock}
\includegraphics[page=64,trim={53.00bp 79.80bp 53.87bp 570.88bp},clip,width=0.92\textwidth,max height=0.38\textheight]{Ferziger4th.pdf}
\end{sourceblock}
在这种情况下，CDS 的前导截断误差项可以重写：

\begin{sourceblock}
\begin{tikzpicture}
\node[inner sep=0,anchor=south west] (image) at (0,0) {\includegraphics[page=65,trim={53.00bp 555.65bp 53.87bp 73.14bp},clip,width=0.92\textwidth,max height=0.38\textheight]{Ferziger4th.pdf}};
\begin{scope}[x={(image.south east)},y={(image.north west)}]
\fill[white] (0.32364,0.25997) rectangle (0.35065,0.59357);
\fill[white] (0.35817,0.28380) rectangle (0.42442,0.78046);
\fill[white] (0.42626,0.44016) rectangle (0.49495,0.78046);
\fill[white] (0.51504,0.44203) rectangle (0.53970,0.77323);
\fill[white] (0.64153,0.03213) rectangle (0.65218,0.26158);
\end{scope}
\end{tikzpicture}
\end{sourceblock}
一阶 FDS 或 BDS 方案的前导误差项为：

\begin{sourceblock}
\begin{tikzpicture}
\node[inner sep=0,anchor=south west] (image) at (0,0) {\includegraphics[page=65,trim={53.00bp 495.86bp 53.87bp 132.93bp},clip,width=0.92\textwidth,max height=0.38\textheight]{Ferziger4th.pdf}};
\begin{scope}[x={(image.south east)},y={(image.north west)}]
\fill[white] (0.46671,0.44177) rectangle (0.49137,0.77323);
\fill[white] (0.37200,0.25971) rectangle (0.39898,0.59331);
\fill[white] (0.40650,0.28380) rectangle (0.43468,0.59331);
\fill[white] (0.59320,0.03213) rectangle (0.60385,0.26158);
\end{scope}
\end{tikzpicture}
\end{sourceblock}
当re接近1时，CDS的一阶截断误差远小于BDS的误差。

现在让我们看看网格细化后会发生什么。我们考虑两种可能性：将两个粗网格点之间的间距减半并插入新点，以便细网格也具有恒定的间距比。

在第一种情况下，新点周围的间距是均匀的，并且旧点处的扩展因子 re 保持与粗网格上的相同。如果重复细化几次，我们会得到一个除了最粗网格点附近之外到处都是均匀的网格。在此阶段，除了属于最粗网格的网格点之外，在所有网格点处，间距是均匀的并且CDS中的前导误差项消失。经过一些细化后，间距不均匀的点的数量将会很少。因此，全局误差的下降速度将比真正的二阶方案慢一些。

在第二种情况下，细网格的扩展因子小于粗网格的扩展因子。简单的算术表明

\begin{sourceblock}
\includegraphics[page=65,trim={53.00bp 283.01bp 53.87bp 361.09bp},clip,width=0.92\textwidth,max height=0.38\textheight]{Ferziger4th.pdf}
\end{sourceblock}
其中 h 表示细化网格，2 h 表示粗化网格。让我们考虑两个网格共有的节点；两个网格上节点 i 处的前导截断误差项之比为（见式（3.22））：

\begin{sourceblock}
\includegraphics[page=65,trim={53.00bp 199.18bp 53.87bp 437.45bp},clip,width=0.92\textwidth,max height=0.38\textheight]{Ferziger4th.pdf}
\end{sourceblock}
两个网格上的网格间距之间存在以下关系（见图 3.3）：

\begin{sourceblock}
\includegraphics[page=65,trim={53.00bp 84.40bp 53.87bp 488.11bp},clip,width=\textwidth,max height=0.38\textheight]{Ferziger4th.pdf}
\par\vspace{0.45\baselineskip}
\small 图 3.3 以常数因子 re 扩展的非均匀网格的细化
\end{sourceblock}
当将它们插入方程（3.24）时，考虑方程（3.23），可以得出CDS的一阶截断误差减小了一个因子

\begin{sourceblock}
\includegraphics[page=66,trim={53.00bp 544.87bp 53.87bp 91.13bp},clip,width=0.92\textwidth,max height=0.38\textheight]{Ferziger4th.pdf}
\end{sourceblock}
当网格细化后。当 r e = 1 时，即网格均匀时，该因子的值为 4。当 r e > 1（扩展网格）或 r e < 1（收缩网格）时，该因子为 r τ > 4，这意味着一阶项引起的误差比二阶误差项减小得更快！因为，在该方法中，随着网格的细化，re → 1，收敛性变为渐近二阶。这将在稍后提供的示例中进行演示。

对于任何方案都可以进行类似的分析，得出相同的结论：非均匀网格的系统细化给出了与均匀网格具有相同阶数的截断误差减少率。

对于给定数量的网格点，通过不均匀的间距几乎总是获得较小的误差。这就是他们的目的。然而，为了使网格完成其工作，用户必须知道哪里需要更小的间距，或者需要使用网格自动适应解的方法。有经验的用户可以识别需要精细网格的区域；见第 12 章 讨论这个问题。那里还将介绍提供自动误差引导网格细化的方法。应该强调的是，随着问题维度的增加，网格生成变得更加困难。事实上，有效网格的生成仍然是计算流体动力学中最困难的问题之一。

通过使用更多的点来消除上述表达式中更多的截断误差项，可以获得一阶导数的高阶近似。例如，使用 φ i − 1 获得 x i 处的二阶导数表达式，并将该表达式代入方程（3.6）中，我们得到以下二阶近似值（对于任何网格，其前导截断误差项（也已给出）与网格间距的平方成正比）：

\begin{sourceblock}
\begin{tikzpicture}
\node[inner sep=0,anchor=south west] (image) at (0,0) {\includegraphics[page=66,trim={53.00bp 180.18bp 53.87bp 419.49bp},clip,width=0.92\textwidth,max height=0.38\textheight]{Ferziger4th.pdf}};
\begin{scope}[x={(image.south east)},y={(image.north west)}]
\fill[white] (0.11197,0.70242) rectangle (0.15020,0.87649);
\fill[white] (0.24974,0.24703) rectangle (0.30325,0.43433);
\fill[white] (0.32487,0.24823) rectangle (0.34953,0.43433);
\fill[white] (0.45125,0.01805) rectangle (0.46189,0.14698);
\end{scope}
\end{tikzpicture}
\end{sourceblock}
\section*{3.4 二阶导数的近似}\phantomsection\addcontentsline{toc}{section}{3.4 二阶导数的近似}
二阶导数出现在扩散项中，见方程（3.1）。为了估计一点的二阶导数，可以使用一阶导数的近似两次。当流体性质可变时，这是唯一可行的方法，因为我们需要扩散系数和一阶导数乘积的导数。接下来我们考虑二阶导数的近似值；稍后将讨论守恒方程中扩散项的应用。

在几何上，二阶导数是与代表一阶导数的曲线相切的线的斜率，见图3.2。通过在位置 x i + 1 和 x i 处插入一阶导数的近似值，可以获得二阶导数的近似值：

\begin{sourceblock}
\begin{tikzpicture}
\node[inner sep=0,anchor=south west] (image) at (0,0) {\includegraphics[page=67,trim={53.00bp 484.32bp 53.87bp 129.98bp},clip,width=0.92\textwidth,max height=0.38\textheight]{Ferziger4th.pdf}};
\begin{scope}[x={(image.south east)},y={(image.north west)}]
\fill[white] (0.00000,0.77488) rectangle (0.11558,0.99788);
\fill[white] (0.46448,0.61863) rectangle (0.50271,0.84182);
\fill[white] (0.31702,0.33912) rectangle (0.36734,0.58102);
\fill[white] (0.31802,0.06289) rectangle (0.36620,0.29552);
\end{scope}
\end{tikzpicture}
\end{sourceblock}
所有这些近似值都涉及至少三个点的数据。

在上式中，外导数是通过FDS估计的。对于内导数，可以使用不同的近似值，例如 BDS；这会产生以下表达式：

\begin{sourceblock}
\begin{tikzpicture}
\node[inner sep=0,anchor=south west] (image) at (0,0) {\includegraphics[page=67,trim={53.00bp 388.66bp 53.87bp 240.12bp},clip,width=0.92\textwidth,max height=0.38\textheight]{Ferziger4th.pdf}};
\begin{scope}[x={(image.south east)},y={(image.north west)}]
\fill[white] (0.08623,0.47056) rectangle (0.13651,0.80621);
\end{scope}
\end{tikzpicture}
\end{sourceblock}
还可以使用 CDS 方法，该方法需要 x i − 1 和 x i + 1 处的一阶导数。更好的选择是在 x i 和 x i + 1 以及 x i 和 x i − 1 中间的点处计算 ∂φ/∂ x。这些一阶导数的 CDS 近似为：

\begin{sourceblock}
\begin{tikzpicture}
\node[inner sep=0,anchor=south west] (image) at (0,0) {\includegraphics[page=67,trim={53.00bp 303.65bp 53.87bp 323.06bp},clip,width=0.92\textwidth,max height=0.38\textheight]{Ferziger4th.pdf}};
\begin{scope}[x={(image.south east)},y={(image.north west)}]
\fill[white] (0.17835,0.49861) rectangle (0.21660,0.79178);
\fill[white] (0.17937,0.13543) rectangle (0.21441,0.43546);
\fill[white] (0.32415,0.11311) rectangle (0.37765,0.42886);
\fill[white] (0.38102,0.11489) rectangle (0.43654,0.43546);
\end{scope}
\end{tikzpicture}
\end{sourceblock}
分别。二阶导数的结果表达式为：

\begin{sourceblock}
\begin{tikzpicture}
\node[inner sep=0,anchor=south west] (image) at (0,0) {\includegraphics[page=67,trim={53.00bp 197.29bp 53.87bp 384.46bp},clip,width=0.92\textwidth,max height=0.38\textheight]{Ferziger4th.pdf}};
\begin{scope}[x={(image.south east)},y={(image.north west)}]
\fill[white] (0.12177,0.57163) rectangle (0.17206,0.72011);
\fill[white] (0.21080,0.48892) rectangle (0.23898,0.62590);
\fill[white] (0.12277,0.40194) rectangle (0.17095,0.54485);
\end{scope}
\end{tikzpicture}
\end{sourceblock}
对于等距的点间距，表达式（3.28）和（3.29）变为：

\begin{sourceblock}
\begin{tikzpicture}
\node[inner sep=0,anchor=south west] (image) at (0,0) {\includegraphics[page=67,trim={53.00bp 137.68bp 53.87bp 491.11bp},clip,width=0.92\textwidth,max height=0.38\textheight]{Ferziger4th.pdf}};
\begin{scope}[x={(image.south east)},y={(image.north west)}]
\fill[white] (0.32719,0.47068) rectangle (0.37750,0.80616);
\end{scope}
\end{tikzpicture}
\end{sourceblock}
泰勒级数展开提供了另一种推导二阶导数近似值的方法。使用 x i − 1 和 x i + 1 处的级数 ( 3.6 )，可以通过错误的显式表达式重新推导方程 ( 3.28 )：

\begin{sourceblock}
\begin{tikzpicture}
\node[inner sep=0,anchor=south west] (image) at (0,0) {\includegraphics[page=68,trim={53.00bp 550.23bp 53.87bp 50.00bp},clip,width=0.92\textwidth,max height=0.38\textheight]{Ferziger4th.pdf}};
\begin{scope}[x={(image.south east)},y={(image.north west)}]
\fill[white] (0.11170,0.72948) rectangle (0.16198,0.91959);
\fill[white] (0.23666,0.24943) rectangle (0.30217,0.44227);
\fill[white] (0.30553,0.25049) rectangle (0.37576,0.44227);
\fill[white] (0.37759,0.25049) rectangle (0.44430,0.44227);
\fill[white] (0.44968,0.24943) rectangle (0.54674,0.44227);
\fill[white] (0.64490,0.01821) rectangle (0.65555,0.14823);
\end{scope}
\end{tikzpicture}
\end{sourceblock}
领先的截断误差项是一阶的，但当点之间的间距均匀时消失，使得近似二阶准确。然而，即使网格不均匀，第 3 节中给出的论点仍然有效。 3.3.4表明，当网格细化时，截断误差以二阶方式减小。当使用复利网格时，误差以与一阶导数的 CDS 近似相同的方式减小，参见方程（3.25）。

二阶导数的高阶近似可以通过包含更多数据点来获得，例如 x i − 2 或 x i + 2。

最后，可以使用插值法来拟合 n + 1 个数据点的 n 次多项式。根据该插值，可以通过微分获得直到第 n 次的所有导数的近似值。对三个点使用二次插值可得出上面给出的公式。像该节中描述的方法一样。 3.3.3也可以推广到二阶导数。

一般来说，二阶导数近似的截断误差是插值多项式减一的次数（抛物线为一阶，三次曲线为二阶等）。当间距均匀且使用偶数阶多项式时，可获得一阶。例如，通过五个点拟合的四次多项式会导致均匀网格上的四阶近似：

\begin{sourceblock}
\begin{tikzpicture}
\node[inner sep=0,anchor=south west] (image) at (0,0) {\includegraphics[page=68,trim={53.00bp 287.20bp 53.87bp 341.59bp},clip,width=0.92\textwidth,max height=0.38\textheight]{Ferziger4th.pdf}};
\begin{scope}[x={(image.south east)},y={(image.north west)}]
\fill[white] (0.06626,0.47095) rectangle (0.11657,0.80643);
\fill[white] (0.06725,0.08728) rectangle (0.11543,0.41017);
\end{scope}
\end{tikzpicture}
\end{sourceblock}
还可以使用二阶导数的近似来提高一阶导数的近似的准确性。例如，使用 FDS 表达式计算一阶导数，方程 ( 3.4 )，仅保留右侧两项，并使用 CDS 表达式 ( 3.29 ) 计算二阶导数，则得到以下一阶导数表达式：

\begin{sourceblock}
\begin{tikzpicture}
\node[inner sep=0,anchor=south west] (image) at (0,0) {\includegraphics[page=68,trim={53.00bp 179.58bp 53.87bp 449.20bp},clip,width=0.92\textwidth,max height=0.38\textheight]{Ferziger4th.pdf}};
\begin{scope}[x={(image.south east)},y={(image.north west)}]
\fill[white] (0.08650,0.47082) rectangle (0.12472,0.78051);
\end{scope}
\end{tikzpicture}
\end{sourceblock}
该表达式在任何网格上都具有二阶截断误差，并简化为均匀网格上一阶导数的标准 CDS 表达式。该近似与方程(3.26)相同。以类似的方式，可以通过消除前导截断误差项中的导数来升级任何近似值。高阶近似总是涉及更多节点，产生需要求解的更复杂的方程以及更复杂的边界条件处理，因此必须进行权衡。二阶近似通常在工程应用中提供了易用性、准确性和成本效益的良好组合。当网格足够细时，三阶和四阶方案对于给定数量的点提供更高的精度，但更难以使用。更高阶的方法仅在特殊情况下使用。

对于扩散项 ( 3.1 ) 的保守形式，必须首先近似内部一阶导数 ∂φ/∂ x，将其相乘并再次对乘积进行微分。如上所示，不必对内导数和外导数使用相同的近似值。

最常用的方法是二阶中心差近似；内导数在节点之间的中间点处进行近似，然后使用网格大小为 x 的中心差。可获得：

\begin{sourceblock}
\begin{tikzpicture}
\node[inner sep=0,anchor=south west] (image) at (0,0) {\includegraphics[page=69,trim={53.00bp 388.81bp 53.87bp 186.97bp},clip,width=0.92\textwidth,max height=0.38\textheight]{Ferziger4th.pdf}};
\begin{scope}[x={(image.south east)},y={(image.north west)}]
\fill[white] (0.19720,0.67751) rectangle (0.21573,0.80544);
\fill[white] (0.27952,0.67751) rectangle (0.31777,0.80544);
\fill[white] (0.37281,0.60015) rectangle (0.40099,0.72820);
\fill[white] (0.18881,0.51892) rectangle (0.22385,0.64985);
\fill[white] (0.35447,0.49624) rectangle (0.36511,0.59108);
\end{scope}
\end{tikzpicture}
\end{sourceblock}
使用内部和外部一阶导数的不同近似可以轻松获得其他近似；可以使用上一节中提出的任何近似值。

\section*{3.5 混合导数的近似}\phantomsection\addcontentsline{toc}{section}{3.5 混合导数的近似}
仅当输运方程用非正交坐标系表示时，才会出现混合导数；见第 9 章 为例。混合导数 ∂ 2 φ/∂ x ∂ y 可以通过组合一维近似来处理，如上面针对二阶导数所述。可以这样写：

\begin{sourceblock}
\begin{tikzpicture}
\node[inner sep=0,anchor=south west] (image) at (0,0) {\includegraphics[page=69,trim={53.00bp 201.92bp 53.87bp 428.93bp},clip,width=0.92\textwidth,max height=0.38\textheight]{Ferziger4th.pdf}};
\begin{scope}[x={(image.south east)},y={(image.north west)}]
\fill[white] (0.37997,0.43978) rectangle (0.43026,0.79456);
\fill[white] (0.49134,0.43978) rectangle (0.50986,0.76736);
\fill[white] (0.44767,0.24199) rectangle (0.47585,0.56957);
\fill[white] (0.37116,0.03400) rectangle (0.43901,0.36923);
\end{scope}
\end{tikzpicture}
\end{sourceblock}
( x i , y j ) 处的混合二阶导数可以使用 CDS 进行估计，方法是首先评估 ( x i + 1 , y j ) 和 ( x i − 1 , y j ) 处关于 y 的一阶导数，然后以上述方式评估该新函数关于 x 的一阶导数。

微分顺序可以改变；数值近似可能取决于阶数。虽然这看起来是一个缺点，但实际上并没有什么问题。所需要的只是数值近似在无限小的网格尺寸的限制下变得精确。两种近似得到的解的差异是由于离散化误差不同造成的。

\section*{3.6 其他项的近似}\phantomsection\addcontentsline{toc}{section}{3.6 其他项的近似}
\subsection*{3.6.1 无差别术语}\phantomsection\addcontentsline{toc}{subsection}{3.6.1 无差别术语}
在标量守恒方程中，可能存在不包含导数的项（我们已将其集中到源项 q φ 中）；这些也必须进行评估。在FD方法中，通常只需要节点处的值。如果非微分项涉及因变量，则可以用变量的节点值来表示。当依赖性是非线性时需要小心。这些术语的处理取决于方程，进一步的讨论被推迟到在第 5 章中、7和8。

\subsection*{3.6.2 边界附近的差异化术语}\phantomsection\addcontentsline{toc}{subsection}{3.6.2 边界附近的差异化术语}
当使用导数的高阶近似时，确实会出现问题；由于它们需要三个以上点的数据，内部节点的近似可能需要边界之外的点的数据。那么可能需要对接近边界的点处的导数使用不同的近似值；通常这些的阶数比内部更深处使用的近似值低，并且可能是单方面的差异。例如，根据对边界值和三个内部点的三次拟合，可以针对紧邻边界点处的一阶导数导出方程(3.13)。通过边界和四个内部点拟合四阶多项式，在第一个内部点 x = x 2 处得到一阶导数的以下近似值：

\begin{sourceblock}
\begin{tikzpicture}
\node[inner sep=0,anchor=south west] (image) at (0,0) {\includegraphics[page=70,trim={53.00bp 247.51bp 53.87bp 381.21bp},clip,width=0.92\textwidth,max height=0.38\textheight]{Ferziger4th.pdf}};
\begin{scope}[x={(image.south east)},y={(image.north west)}]
\fill[white] (0.10776,0.47194) rectangle (0.14598,0.78087);
\fill[white] (0.10875,0.08899) rectangle (0.14382,0.40513);
\end{scope}
\end{tikzpicture}
\end{sourceblock}
使用相同多项式的二阶导数近似可得出：

\begin{sourceblock}
\begin{tikzpicture}
\node[inner sep=0,anchor=south west] (image) at (0,0) {\includegraphics[page=70,trim={53.00bp 187.82bp 53.87bp 440.89bp},clip,width=0.92\textwidth,max height=0.38\textheight]{Ferziger4th.pdf}};
\begin{scope}[x={(image.south east)},y={(image.north west)}]
\fill[white] (0.06177,0.47155) rectangle (0.11209,0.80631);
\fill[white] (0.06277,0.08923) rectangle (0.11095,0.41117);
\end{scope}
\end{tikzpicture}
\end{sourceblock}
以类似的方式，我们可以使用该点、边界点和一定数量的内部网格点处的变量值导出第一内部网格点处的任何阶的单侧导数近似。

\section*{3.7 边界条件的实现}\phantomsection\addcontentsline{toc}{section}{3.7 边界条件的实现}
每个内部网格点都需要偏微分方程的有限差分近似。为了使解唯一，连续问题需要有关域边界处的解的信息。一般来说，给出边界处变量的值（狄利克雷边界条件）或其在特定方向上的梯度（通常垂直于边界——诺依曼边界条件）或这两个量的线性组合。

边界条件可以通过多种方式实现。在一种方法中，传输方程总是仅在内部网格点处求解；如果边界处的边界值未知（诺依曼或罗宾条件），则使用离散边界条件通过内点处的值来表达边界值，从而将其消除为未知。在另一种方法中，在边界之外使用鬼点，并对未知边界值求解输运方程，而鬼点处的值是从离散边界条件获得的。

\subsection*{3.7.1 使用内部网格点实现边界条件}\phantomsection\addcontentsline{toc}{subsection}{3.7.1 使用内部网格点实现边界条件}
如果变量值在某个边界点已知，则无需求解。在包含这些点处的数据的所有 FD 方程中，使用已知值并且不需要更多。

如果梯度是在边界处规定的，则可以使用合适的FD近似（如果仅采用内部网格点，则它必须是单侧近似）来计算变量的边界值。例如，如果规定法线方向的梯度为零，则简单的 FDS 近似可得出（见图 3.1）： ∂φ

\begin{sourceblock}
\begin{tikzpicture}
\node[inner sep=0,anchor=south west] (image) at (0,0) {\includegraphics[page=71,trim={53.00bp 227.58bp 53.87bp 401.13bp},clip,width=0.92\textwidth,max height=0.38\textheight]{Ferziger4th.pdf}};
\begin{scope}[x={(image.south east)},y={(image.north west)}]
\fill[white] (0.00000,0.74566) rectangle (0.05236,1.00000);
\fill[white] (0.05504,0.74566) rectangle (0.16698,1.00000);
\fill[white] (0.31426,0.47181) rectangle (0.35251,0.78066);
\fill[white] (0.31525,0.08923) rectangle (0.35032,0.40529);
\end{scope}
\end{tikzpicture}
\end{sourceblock}
给出 φ 1 = φ 2，允许边界值被边界旁边的节点处的值替换，并作为未知数被消除。这是一阶近似；通过使用更高次数的多项式拟合可以获得更高阶的近似值。

根据对边界和两个内部点的抛物线拟合，对于边界处的一阶导数，可以获得以下在任何网格上都有效的二阶近似：

\begin{sourceblock}
\begin{tikzpicture}
\node[inner sep=0,anchor=south west] (image) at (0,0) {\includegraphics[page=71,trim={53.00bp 96.16bp 53.87bp 532.55bp},clip,width=0.92\textwidth,max height=0.38\textheight]{Ferziger4th.pdf}};
\begin{scope}[x={(image.south east)},y={(image.north west)}]
\fill[white] (0.07185,0.47181) rectangle (0.11008,0.78066);
\end{scope}
\end{tikzpicture}
\end{sourceblock}
在统一网格上，该表达式简化为：

\begin{sourceblock}
\begin{tikzpicture}
\node[inner sep=0,anchor=south west] (image) at (0,0) {\includegraphics[page=72,trim={53.00bp 556.35bp 53.87bp 72.36bp},clip,width=0.92\textwidth,max height=0.38\textheight]{Ferziger4th.pdf}};
\begin{scope}[x={(image.south east)},y={(image.north west)}]
\fill[white] (0.04629,0.47155) rectangle (0.08451,0.78066);
\fill[white] (0.04731,0.08923) rectangle (0.08235,0.40529);
\end{scope}
\end{tikzpicture}
\end{sourceblock}
现在，出现在内部网格节点处的任何离散项中的边界值将被点 2 和 3 处的节点以及指定的边界梯度的组合所替换。这需要调整边界附近节点的系数矩阵元素（如果使用高阶近似，则可能在下一层节点），但要求解的方程组始终保持不变。

\begin{sourceblock}
\begin{tikzpicture}
\node[inner sep=0,anchor=south west] (image) at (0,0) {\includegraphics[page=72,trim={53.00bp 437.66bp 53.87bp 191.06bp},clip,width=0.92\textwidth,max height=0.38\textheight]{Ferziger4th.pdf}};
\begin{scope}[x={(image.south east)},y={(image.north west)}]
\fill[white] (0.00000,0.74559) rectangle (0.05405,1.00000);
\fill[white] (0.05675,0.74559) rectangle (0.14313,1.00000);
\fill[white] (0.27657,0.47167) rectangle (0.31480,0.78060);
\fill[white] (0.27756,0.08899) rectangle (0.31263,0.40513);
\end{scope}
\end{tikzpicture}
\end{sourceblock}
有时需要计算给定变量边界值的点处垂直于边界的一阶导数（例如，计算通过等温表面的热通量）。在这种情况下，上面给出的任何单侧近似都是合适的。结果的准确性不仅取决于所使用的近似值，还取决于内点值的准确性。对于这两个目的，使用相同阶的近似值是明智的。请参阅 Fornberg (1988) 的表 3，了解高达四阶和各种精度的单边导数的表达式列表。

当章节中描述的紧凑方案时。使用3.3.3时，必须提供变量值和边界节点处的导数。通常，其中之一是已知的，另一个必须使用来自内部的信息来计算。例如，当规定变量值时，可以采用边界节点处导数的单侧近似，如式(3.40)。另一方面，如果导数已知，则可以使用多项式插值来计算边界值。通过对四个点的三次拟合，可以获得以下边界值表达式：

\begin{sourceblock}
\begin{tikzpicture}
\node[inner sep=0,anchor=south west] (image) at (0,0) {\includegraphics[page=72,trim={53.00bp 192.90bp 53.87bp 435.81bp},clip,width=0.92\textwidth,max height=0.38\textheight]{Ferziger4th.pdf}};
\begin{scope}[x={(image.south east)},y={(image.north west)}]
\fill[white] (0.30908,0.44082) rectangle (0.37161,0.78066);
\fill[white] (0.37507,0.44082) rectangle (0.45266,0.78066);
\fill[white] (0.45603,0.47181) rectangle (0.48421,0.78066);
\fill[white] (0.48608,0.44082) rectangle (0.53377,0.78066);
\fill[white] (0.23621,0.25434) rectangle (0.26878,0.59391);
\fill[white] (0.27380,0.28507) rectangle (0.30198,0.59418);
\end{scope}
\end{tikzpicture}
\end{sourceblock}
可以用类似的方式获得低阶或高阶的近似值。

\subsection*{3.7.2 使用鬼点实现边界条件}\phantomsection\addcontentsline{toc}{subsection}{3.7.2 使用鬼点实现边界条件}
作为上一节策略的替代方案，可以使用中心差和鬼点（即，点）来实现基于导数的边界条件

\begin{center}\small 图 3.4 一维 FD 网格示例，显示 i = 0 处的幽灵节点（图 3.1 中一维网格的扩展）\end{center}
幽灵节点

内部节点 边界节点
i i+1 i−1 1 0 2
都在问题边界之外。例如，与方程（3.38）相比，直接的结果是准确性的提高。同样，讨论通量边界条件的实现就足够了，因为在狄利克雷边界条件的边界上变量的值是已知的。

例如，如果在法线方向规定第三类边界条件，即 Robin 条件，则简单的 CDS 近似会导致：

\begin{sourceblock}
\begin{tikzpicture}
\node[inner sep=0,anchor=south west] (image) at (0,0) {\includegraphics[page=73,trim={53.00bp 421.67bp 53.87bp 207.04bp},clip,width=0.92\textwidth,max height=0.38\textheight]{Ferziger4th.pdf}};
\begin{scope}[x={(image.south east)},y={(image.north west)}]
\fill[white] (0.22322,0.47208) rectangle (0.26144,0.78092);
\end{scope}
\end{tikzpicture}
\end{sourceblock}
得出 φ 0 = φ 2 + 2 ( x 2 − x 1 ) c φ 1。从图3.4中我们看到节点0在计算域之外；然而，导数项是居中的，因此边界条件的近似是二阶精确的。

有两种方法可以使用此方法。首先，写出节点 1 的差分方程，然后将 φ 0 替换为 φ 2 + 2 ( x 2 − x 1 ) c φ 1。这意味着在施加诺伊曼或罗宾边界条件的情况下，也可以求解边界节点的输运方程，但所有未知数都在计算域内。或者，可以将鬼点处的变量值视为未知数，并且在施加导数边界条件的任何地方将等价于方程(3.42)的方程添加到方程组中。这改变了方程结构，例如，在某些使用直接矩阵求逆的情况下，方程结构不再是三对角的；这种方法当时并不受欢迎。然而，在迭代求解方法中，第二种方法易于编程且高效。

\section*{3.8 代数方程组}\phantomsection\addcontentsline{toc}{section}{3.8 代数方程组}
有限差分近似提供了每个网格节点处的代数方程；它包含该节点处的变量值以及相邻节点处的值。如果微分方程是非线性的，则近似值将包含一些非线性项。数值求解过程将需要线性化；求解这些方程的方法将在第 1 章中讨论。 5.目前，我们仅考虑线性情况。所描述的方法也适用于非线性情况。对于线性情况，离散化的结果是以下形式的线性代数方程组：

\begin{sourceblock}
\begin{tikzpicture}
\node[inner sep=0,anchor=south west] (image) at (0,0) {\includegraphics[page=73,trim={53.00bp 79.20bp 53.87bp 561.47bp},clip,width=0.92\textwidth,max height=0.38\textheight]{Ferziger4th.pdf}};
\begin{scope}[x={(image.south east)},y={(image.north west)}]
\fill[white] (0.00000,0.95367) rectangle (0.02908,1.00000);
\fill[white] (0.03179,0.95367) rectangle (0.07320,1.00000);
\fill[white] (0.07588,0.95367) rectangle (0.14725,1.00000);
\fill[white] (0.34935,0.45387) rectangle (0.41555,0.95289);
\fill[white] (0.41919,0.49902) rectangle (0.44740,0.95289);
\end{scope}
\end{tikzpicture}
\end{sourceblock}
\begin{sourceblock}
\includegraphics[page=74,trim={53.00bp 493.26bp 53.87bp 52.00bp},clip,width=\textwidth,max height=0.38\textheight]{Ferziger4th.pdf}
\par\vspace{0.45\baselineskip}
\small 图 3.5 2D 和 3D 计算分子示例
\end{sourceblock}
其中 P 表示偏微分方程近似的节点，索引 l 遍历有限差分近似中涉及的邻居节点。节点P及其邻居形成所谓的计算分子；由二阶和三阶近似得出的两个例子如图 3.5 所示。 1 系数 A l 取决于几何量、流体特性，对于非线性方程，还取决于变量值本身。 Q P 包含所有不包含未知变量值的项；据推测已知。

方程和未知数的数量必须相等，即每个网格节点必须有一个方程。因此，我们有一大堆必须用数值​​方法求解的线性代数方程。该系统是稀疏的，这意味着每个方程仅包含几个未知数。该系统可以用矩阵表示法写成如下：

\begin{sourceblock}
\begin{tikzpicture}
\node[inner sep=0,anchor=south west] (image) at (0,0) {\includegraphics[page=74,trim={53.00bp 305.15bp 53.87bp 346.66bp},clip,width=0.92\textwidth,max height=0.38\textheight]{Ferziger4th.pdf}};
\begin{scope}[x={(image.south east)},y={(image.north west)}]
\fill[white] (0.00000,0.91835) rectangle (0.10162,1.00000);
\end{scope}
\end{tikzpicture}
\end{sourceblock}
其中A是方稀疏系数矩阵，φ是包含网格节点处的变量值的向量（或列矩阵），Q是包含方程（3.43）右侧项的向量。

矩阵 A 的结构取决于向量 φ 中变量的排序。对于结构化网格，如果变量从角开始标记并以规则方式（字典顺序）逐行遍历，则矩阵具有多对角结构。对于五点计算分子的情况，所有非零系数都位于主对角线、两个相邻对角线以及从主对角线移去 N 个位置的另外两条对角线上，其中 N 是一个方向上的节点数。所有其他系数均为零。这种结构允许使用高效的迭代求解器。

在本书中，为了明确起见，我们将从域的西南角开始对向量 φ 中的条目进行排序，沿着每条网格线向北行进，然后向东穿过域（在三维情况下，我们将从底部计算表面开始，以刚才描述的方式在每个水平面上行进，然后从下到上）。变量是

\footnotetext[1]{例如，诸如 Eq.( 3.43 ) 的方程和图中的分子将通过泊松方程 ∇· ( ∇ ) = f 的离散化生成。}
\begin{center}\small 表 3.2 将网格索引转换为向量或列矩阵的一维存储位置\end{center}
网格位置 指南针符号 存储位置
i , j , k P l = ( k − 1 ) N j N i + ( i − 1 ) N j + j i − 1 , j , k W l − N j i , j − 1 , k S l − 1 i , j + 1 , k N l + 1 i + 1 , j , k El + N j i , j , k − 1 B l − N i N j i , j , k + 1 T l + N i N j
通常以一维数组的形式存储在计算机中。网格位置、罗盘符号和存储位置之间的转换如表 3.2 所示。

由于矩阵 A 是稀疏的，因此将其存储为计算机内存中的二维（2D）或三维（3D）数组是没有意义的（这是完整矩阵的标准做法）。在二维中，将每个非零对角线的元素存储在维度为 1 × N i N j 的单独数组中，其中 N i 和 N j 是两个坐标方向上的网格点数，仅需要 5 N i N j 个字的存储空间；全阵列存储需要 N 2
我N 2
j 存储字数。在 3D 中，数字为 7 N i N j N k 和 N 2 i N 2
jN 2
分别为 k。差异足够大，以至于对角存储方案可以允许将问题保留在主存储器中，而全阵列方案则不允许。

如果使用网格索引引用节点值，例如二维中的 φ i , j，它们看起来像矩阵元素或张量的分量。因为它们实际上是矢量 φ 的分量，所以它们应该只有表 3.2 中指示的单个索引。

二维线性代数方程现在可以写成以下形式：

\begin{sourceblock}
\begin{tikzpicture}
\node[inner sep=0,anchor=south west] (image) at (0,0) {\includegraphics[page=75,trim={53.00bp 232.27bp 53.87bp 405.41bp},clip,width=0.92\textwidth,max height=0.38\textheight]{Ferziger4th.pdf}};
\begin{scope}[x={(image.south east)},y={(image.north west)}]
\fill[white] (0.00000,0.86753) rectangle (0.07074,1.00000);
\end{scope}
\end{tikzpicture}
\end{sourceblock}
如上所述，将矩阵存储为数组没有什么意义。相反，如果对角线保存在单独的数组中，则最好为每个对角线指定单独的名称。 Because each diagonal represents the connection to the variable at a node that lies in a particular direction with respect to the central node, we shall call them A W ,
A S 、 A P 、 A N 和 A E；它们在具有 5 × 5 内部节点的网格的矩阵中的位置如图 3.6 所示。通过这种点排序，每个节点都用索引 l 标识，这也是相对存储位置。在这种符号中，方程（3.45）可以写成

\begin{sourceblock}
\begin{tikzpicture}
\node[inner sep=0,anchor=south west] (image) at (0,0) {\includegraphics[page=75,trim={53.00bp 102.59bp 53.87bp 548.09bp},clip,width=0.92\textwidth,max height=0.38\textheight]{Ferziger4th.pdf}};
\begin{scope}[x={(image.south east)},y={(image.north west)}]
\fill[white] (0.00000,0.92367) rectangle (0.03242,1.00000);
\fill[white] (0.03510,0.92367) rectangle (0.12478,1.00000);
\end{scope}
\end{tikzpicture}
\end{sourceblock}
\begin{sourceblock}
\includegraphics[page=76,trim={53.00bp 394.26bp 53.87bp 52.00bp},clip,width=\textwidth,max height=0.38\textheight]{Ferziger4th.pdf}
\par\vspace{0.45\baselineskip}
\small 图 3.6 五点计算分子的矩阵结构（系数矩阵中五个对角线上的非零项用阴影表示；每一组水平框对应一条网格线）
\end{sourceblock}
其中索引 l 表示方程（3.45）中的行，可以理解，并且指示向量中的列或位置的索引已被相应的字母替换。从现在开始我们将使用这种速记符号。为了清晰起见，必要时将插入索引。类似的处理方法适用于三维问题。

对于块结构和复合网格，该结构保留在每个块内，并且可以使用规则结构网格的求解器。这将在第 1 章中进一步讨论。 5.

对于非结构化网格，系数矩阵仍然稀疏，但不再具有带状结构。对于仅使用四个最近邻节点的四边形和近似的二维网格，任何列或行中只有五个非零系数。主对角线是满的，其他非零系数位于主对角线的一定范围内，但不一定位于确定的对角线上。对于此类矩阵必须使用不同类型的迭代求解器；它们将在第 1 章中讨论。 5.非结构化网格的存储方案将在第 1 章中介绍。如图9所示，因为此类网格主要用于采用FV方法的复杂几何形状。

\section*{3.9 离散化误差}\phantomsection\addcontentsline{toc}{section}{3.9 离散化误差}
因为离散方程表示微分方程的近似，所以后者的精确解（我们用 表示）不满足差分方程。 The imbalance, which is due to truncation of the Taylor series, is called truncation error . For a grid with a reference spacing h , the truncation error τ h is defined as:

\begin{sourceblock}
\begin{tikzpicture}
\node[inner sep=0,anchor=south west] (image) at (0,0) {\includegraphics[page=77,trim={53.00bp 569.98bp 53.87bp 80.71bp},clip,width=0.92\textwidth,max height=0.38\textheight]{Ferziger4th.pdf}};
\begin{scope}[x={(image.south east)},y={(image.north west)}]
\fill[white] (0.00000,0.87767) rectangle (0.02770,1.00000);
\fill[white] (0.03296,0.92427) rectangle (0.05777,1.00000);
\fill[white] (0.06042,0.92427) rectangle (0.15341,1.00000);
\fill[white] (0.15615,0.92427) rectangle (0.19426,1.00000);
\end{scope}
\end{tikzpicture}
\end{sourceblock}
其中L是代表微分方程的符号算子，L h 是代表在网格h上离散化得到的代数方程组的符号算子，由式（3.44）给出。

网格 h 上的离散方程的精确解 φ h 满足以下方程：

\begin{sourceblock}
\begin{tikzpicture}
\node[inner sep=0,anchor=south west] (image) at (0,0) {\includegraphics[page=77,trim={53.00bp 485.85bp 53.87bp 164.39bp},clip,width=0.92\textwidth,max height=0.38\textheight]{Ferziger4th.pdf}};
\begin{scope}[x={(image.south east)},y={(image.north west)}]
\fill[white] (0.00000,0.92579) rectangle (0.11561,1.00000);
\end{scope}
\end{tikzpicture}
\end{sourceblock}
\begin{sourceblock}
\begin{tikzpicture}
\node[inner sep=0,anchor=south west] (image) at (0,0) {\includegraphics[page=77,trim={53.00bp 437.28bp 53.87bp 200.62bp},clip,width=0.92\textwidth,max height=0.38\textheight]{Ferziger4th.pdf}};
\begin{scope}[x={(image.south east)},y={(image.north west)}]
\fill[white] (0.00000,0.97132) rectangle (0.02244,1.00000);
\fill[white] (0.02391,0.97132) rectangle (0.10620,1.00000);
\fill[white] (0.10767,0.97132) rectangle (0.17071,1.00000);
\fill[white] (0.17218,0.97132) rectangle (0.21362,1.00000);
\fill[white] (0.21510,0.97132) rectangle (0.28265,1.00000);
\fill[white] (0.28421,0.97132) rectangle (0.38556,1.00000);
\fill[white] (0.38704,0.97132) rectangle (0.41681,1.00000);
\fill[white] (0.41832,0.97132) rectangle (0.45973,1.00000);
\fill[white] (0.46120,0.97132) rectangle (0.54256,1.00000);
\fill[white] (0.54406,0.97132) rectangle (0.68120,1.00000);
\fill[white] (0.68271,0.97132) rectangle (0.79071,1.00000);
\fill[white] (0.79224,0.97132) rectangle (0.82701,1.00000);
\fill[white] (0.82848,0.97132) rectangle (0.86992,1.00000);
\fill[white] (0.00000,0.55170) rectangle (0.05077,0.96105);
\fill[white] (0.05344,0.54816) rectangle (0.12764,0.96105);
\fill[white] (0.13032,0.56091) rectangle (0.16036,1.00000);
\end{scope}
\end{tikzpicture}
\end{sourceblock}
从式（3.47）和式（3.48）可以看出，对于线性问题，以下关系成立：

\begin{sourceblock}
\begin{tikzpicture}
\node[inner sep=0,anchor=south west] (image) at (0,0) {\includegraphics[page=77,trim={53.00bp 389.88bp 53.87bp 258.57bp},clip,width=0.92\textwidth,max height=0.38\textheight]{Ferziger4th.pdf}};
\begin{scope}[x={(image.south east)},y={(image.north west)}]
\fill[white] (0.00000,0.85076) rectangle (0.12397,1.00000);
\end{scope}
\end{tikzpicture}
\end{sourceblock}
该方程表明截断误差是离散化误差的来源，该误差由算子 L h 进行对流和扩散。 Exact analysis is not possible for non-linear equations, but we expect similar behavior; in any case, if the error is small enough, we can locally linearize about the exact solution and what we will say in this 节 is valid.有关截断误差的大小和分布的信息可以用作网格细化的指导，并有助于实现解域中各处具有相同水平的离散化误差的目标。 However, as the exact solution is not known, the truncation error cannot be calculated exactly. An approximation to it may be obtained by using a solution from another (finer or coarser) grid.由此获得的截断误差的估计并不总是准确的，但它的目的是指出具有较大误差并需要更精细网格的区域。

对于足够细的网格，截断误差（以及离散化误差）与泰勒级数中的首项成正比：

\begin{sourceblock}
\includegraphics[page=77,trim={53.00bp 186.23bp 53.87bp 459.16bp},clip,width=0.92\textwidth,max height=0.38\textheight]{Ferziger4th.pdf}
\end{sourceblock}
其中 H 代表高阶项，α 取决于给定点的导数，但与 h 无关。离散化误差可以根据系统细化（或粗化）网格上获得的解之间的差异来估计。考虑到精确解可表示为（见式（3.49））：

\begin{sourceblock}
\includegraphics[page=77,trim={53.00bp 103.72bp 53.87bp 542.84bp},clip,width=0.92\textwidth,max height=0.38\textheight]{Ferziger4th.pdf}
\end{sourceblock}
指数 p 是该方案的阶数，可以估计如下：

\begin{sourceblock}
\begin{tikzpicture}
\node[inner sep=0,anchor=south west] (image) at (0,0) {\includegraphics[page=78,trim={53.00bp 569.90bp 53.87bp 50.00bp},clip,width=0.92\textwidth,max height=0.38\textheight]{Ferziger4th.pdf}};
\begin{scope}[x={(image.south east)},y={(image.north west)}]
\fill[white] (0.48171,0.60619) rectangle (0.52502,0.88127);
\fill[white] (0.52944,0.60619) rectangle (0.60271,0.88127);
\fill[white] (0.35507,0.18101) rectangle (0.37489,0.43123);
\fill[white] (0.37934,0.18685) rectangle (0.40752,0.43685);
\end{scope}
\end{tikzpicture}
\end{sourceblock}
从式（3.52）还可以看出，网格 h 上的离散化误差可以近似为：

\begin{sourceblock}
\begin{tikzpicture}
\node[inner sep=0,anchor=south west] (image) at (0,0) {\includegraphics[page=78,trim={53.00bp 512.23bp 53.87bp 125.52bp},clip,width=0.92\textwidth,max height=0.38\textheight]{Ferziger4th.pdf}};
\begin{scope}[x={(image.south east)},y={(image.north west)}]
\fill[white] (0.00000,0.88059) rectangle (0.08565,1.00000);
\fill[white] (0.08836,0.88059) rectangle (0.13146,1.00000);
\end{scope}
\end{tikzpicture}
\end{sourceblock}
如果连续网格上的网格大小之比不是 2，则需要用该比率替换最后两个方程中的因子 2（有关网格未系统细化或粗化时的误差估计的详细信息，请参见 Roache 1994）。

当多个网格上的解可用时，通过将误差估计（3.54）与 φ h 相加，可以获得比最细网格上的解 φ h 更精确的近似值；这种方法称为理查森外推法（Richardson 1910）。它很简单，而且当收敛是单调的时，它是准确的。当有多个解可用时，可以重复该过程以进一步提高准确性。

我们在上面已经表明，重要的是网格细化时误差减少的速率，而不是截断误差中首项所定义的方案的正式顺序。方程 (3.53) 考虑了这一点并返回正确的指数 p。这种对方案顺序的估计也是代码验证中的有用工具。如果一个方法应该是二阶准确的，但等式（3.53）发现它只是一阶准确，则代码中可能存在错误。

使用式(3.53)估计的收敛阶仅在收敛是单调的时才有效。仅在足够细的网格上才能预期单调收敛。我们将在示例中表明，当网格较粗时，误差对网格大小的依赖性可能是不规则的。因此，在比较两个网格上的解时应小心；当收敛不是单调时，即使误差不小，两个连续网格上的解也可能相差不大。需要第三个网格来确保解真正收敛。此外，当解不平滑时，使用泰勒级数近似获得的误差估计可能会产生误导。例如，在湍流模拟中，解在很大范围内变化，并且解决方法的顺序可能不能很好地指示解的质量。见第 3.11 节 将表明，对于这些类型的模拟，四阶方案的误差可能不会比二阶方案小很多。

\section*{3.10 有限差分示例}\phantomsection\addcontentsline{toc}{section}{3.10 有限差分示例}
在此示例中，我们求解两端具有狄利克雷边界条件的稳定一维对流扩散方程。目的是针对具有解析解的简单问题展示 FD 离散化技术的特性。

\begin{sourceblock}
\begin{tikzpicture}
\node[inner sep=0,anchor=south west] (image) at (0,0) {\includegraphics[page=79,trim={53.00bp 573.02bp 53.87bp 59.08bp},clip,width=0.92\textwidth,max height=0.38\textheight]{Ferziger4th.pdf}};
\begin{scope}[x={(image.south east)},y={(image.north west)}]
\fill[white] (0.63519,0.66627) rectangle (0.65341,0.96475);
\fill[white] (0.58842,0.11780) rectangle (0.60665,0.41627);
\fill[white] (0.60244,0.03525) rectangle (0.61817,0.25646);
\end{scope}
\end{tikzpicture}
\end{sourceblock}
Pe < 0
佩=0

\begin{sourceblock}
\begin{tikzpicture}
\node[inner sep=0,anchor=south west] (image) at (0,0) {\includegraphics[page=79,trim={53.00bp 482.26bp 53.87bp 148.07bp},clip,width=0.92\textwidth,max height=0.38\textheight]{Ferziger4th.pdf}};
\begin{scope}[x={(image.south east)},y={(image.north west)}]
\fill[white] (0.77089,0.68277) rectangle (0.80144,0.96649);
\fill[white] (0.80307,0.68277) rectangle (0.82241,0.96649);
\fill[white] (0.82403,0.68277) rectangle (0.84171,0.96649);
\fill[white] (0.98141,0.42446) rectangle (0.99765,0.70818);
\fill[white] (0.63053,0.03351) rectangle (0.64821,0.31723);
\fill[white] (0.91591,0.03351) rectangle (0.93504,0.31723);
\end{scope}
\end{tikzpicture}
\end{sourceblock}
待解方程为（见方程（1.28））：

\begin{sourceblock}
\begin{tikzpicture}
\node[inner sep=0,anchor=south west] (image) at (0,0) {\includegraphics[page=79,trim={53.00bp 409.58bp 53.87bp 228.27bp},clip,width=0.92\textwidth,max height=0.38\textheight]{Ferziger4th.pdf}};
\begin{scope}[x={(image.south east)},y={(image.north west)}]
\fill[white] (0.34641,0.53941) rectangle (0.44319,0.95758);
\fill[white] (0.49395,0.54896) rectangle (0.51248,0.95758);
\fill[white] (0.45029,0.30187) rectangle (0.47847,0.71085);
\fill[white] (0.37594,0.04242) rectangle (0.41101,0.46059);
\end{scope}
\end{tikzpicture}
\end{sourceblock}
边界条件： φ = φ 0 at x = 0, φ = φ L at x = L，见图 3.7；在这种情况下，偏导数可以用普通导数代替。假设密度 ρ 和速度 u 恒定。这个问题有准确的解：

\begin{sourceblock}
\begin{tikzpicture}
\node[inner sep=0,anchor=south west] (image) at (0,0) {\includegraphics[page=79,trim={53.00bp 326.10bp 53.87bp 308.04bp},clip,width=0.92\textwidth,max height=0.38\textheight]{Ferziger4th.pdf}};
\begin{scope}[x={(image.south east)},y={(image.north west)}]
\fill[white] (0.42845,0.47656) rectangle (0.50247,0.88781);
\fill[white] (0.50746,0.47656) rectangle (0.55729,0.84969);
\fill[white] (0.29997,0.27000) rectangle (0.32208,0.63125);
\fill[white] (0.32707,0.27000) rectangle (0.35525,0.63125);
\fill[white] (0.35877,0.23375) rectangle (0.39137,0.63125);
\fill[white] (0.39483,0.27000) rectangle (0.42301,0.63125);
\fill[white] (0.44650,0.03750) rectangle (0.48586,0.41750);
\fill[white] (0.48944,0.03750) rectangle (0.53925,0.41031);
\end{scope}
\end{tikzpicture}
\end{sourceblock}
这里 Pe 是 Peclet 数，定义为：

\begin{sourceblock}
\includegraphics[page=79,trim={53.00bp 271.01bp 53.87bp 367.12bp},clip,width=0.92\textwidth,max height=0.38\textheight]{Ferziger4th.pdf}
\end{sourceblock}
因为它非常简单，所以这个问题经常被用作数值方法的测试，包括离散化和求解方案。从物理上讲，它代表了一种对流通过流向扩散来平衡的情况。这种平衡在实际流量中发挥重要作用的情况很少。通常，对流通过压力梯度或垂直于流动方向的扩散来平衡。在文献中，人们发现了许多针对方程（3.55）开发的方法，然后应用于纳维-斯托克斯方程。结果往往很差，最好避免使用这些方法中的大多数。事实上，使用这个问题作为测试用例可能会比该领域的任何其他方法产生更多糟糕的方法选择。尽管存在这些困难，我们仍应考虑这个问题，因为它提出的一些问题值得关注。

让我们考虑 u ≥ 0 且 φ 0 < φ L 的情况；其他情况都很容易处理。在速度较小（u ≈ 0）或扩散系数较大的情况下，佩克莱特数趋于零，对流可以忽略不计；那么解在 x 上是线性的。当佩克莱特数较大时， φ 随着 x 缓慢增长，然后在接近 x = L 的短距离内突然上升至 φ L。 φ梯度的突然变化对离散化方法提出了严峻的考验。

我们将使用使用三点计算分子的FD方案来离散化方程(3.55)。节点 i 处的代数方程如下：

\begin{sourceblock}
\includegraphics[page=80,trim={53.00bp 533.27bp 53.87bp 115.04bp},clip,width=0.92\textwidth,max height=0.38\textheight]{Ferziger4th.pdf}
\end{sourceblock}
通常的做法是使用 CDS 对扩散项进行离散化；因此，对于外导数，我们有：

\begin{sourceblock}
\begin{tikzpicture}
\node[inner sep=0,anchor=south west] (image) at (0,0) {\includegraphics[page=80,trim={53.00bp 443.61bp 53.87bp 174.25bp},clip,width=0.92\textwidth,max height=0.38\textheight]{Ferziger4th.pdf}};
\begin{scope}[x={(image.south east)},y={(image.north west)}]
\fill[white] (0.23952,0.39623) rectangle (0.25808,0.63567);
\fill[white] (0.18117,0.25145) rectangle (0.20935,0.49109);
\fill[white] (0.23113,0.09942) rectangle (0.26620,0.34445);
\fill[white] (0.39179,0.05675) rectangle (0.40244,0.23447);
\end{scope}
\end{tikzpicture}
\end{sourceblock}
内导数的 CDS 近似为：

\begin{sourceblock}
\begin{tikzpicture}
\node[inner sep=0,anchor=south west] (image) at (0,0) {\includegraphics[page=80,trim={53.00bp 382.68bp 53.87bp 251.03bp},clip,width=0.92\textwidth,max height=0.38\textheight]{Ferziger4th.pdf}};
\begin{scope}[x={(image.south east)},y={(image.north west)}]
\fill[white] (0.14406,0.60654) rectangle (0.18229,0.96300);
\fill[white] (0.14505,0.16466) rectangle (0.18012,0.52945);
\end{scope}
\end{tikzpicture}
\end{sourceblock}
因此，扩散项对代数方程 ( 3.58 ) 系数的贡献为：

\begin{sourceblock}
\begin{tikzpicture}
\node[inner sep=0,anchor=south west] (image) at (0,0) {\includegraphics[page=80,trim={53.00bp 264.83bp 53.87bp 311.18bp},clip,width=0.92\textwidth,max height=0.38\textheight]{Ferziger4th.pdf}};
\begin{scope}[x={(image.south east)},y={(image.north west)}]
\fill[white] (0.00000,0.95828) rectangle (0.07651,1.00000);
\fill[white] (0.07922,0.95828) rectangle (0.12060,1.00000);
\fill[white] (0.12328,0.95828) rectangle (0.18638,1.00000);
\end{scope}
\end{tikzpicture}
\end{sourceblock}
如果使用一阶迎风差（UDS—FDS 或 BDS，取决于流向）对对流项进行离散化，我们有：

\begin{sourceblock}
\begin{tikzpicture}
\node[inner sep=0,anchor=south west] (image) at (0,0) {\includegraphics[page=80,trim={53.00bp 168.71bp 53.87bp 435.94bp},clip,width=0.92\textwidth,max height=0.38\textheight]{Ferziger4th.pdf}};
\begin{scope}[x={(image.south east)},y={(image.north west)}]
\fill[white] (0.25874,0.47243) rectangle (0.35552,0.66482);
\end{scope}
\end{tikzpicture}
\end{sourceblock}
这导致对方程（3.58）的系数有以下贡献：

\begin{sourceblock}
\includegraphics[page=80,trim={53.00bp 87.41bp 53.87bp 525.76bp},clip,width=0.92\textwidth,max height=0.38\textheight]{Ferziger4th.pdf}
\end{sourceblock}
Either A c
E or A c
W 为零，具体取决于流动方向。 CDS 近似可得出：

\begin{sourceblock}
\begin{tikzpicture}
\node[inner sep=0,anchor=south west] (image) at (0,0) {\includegraphics[page=81,trim={53.00bp 544.47bp 53.87bp 84.32bp},clip,width=0.92\textwidth,max height=0.38\textheight]{Ferziger4th.pdf}};
\begin{scope}[x={(image.south east)},y={(image.north west)}]
\fill[white] (0.32653,0.46372) rectangle (0.42328,0.78019);
\end{scope}
\end{tikzpicture}
\end{sourceblock}
CDS 对方程（3.58）系数的贡献为：

\begin{sourceblock}
\includegraphics[page=81,trim={53.00bp 465.53bp 53.87bp 147.63bp},clip,width=0.92\textwidth,max height=0.38\textheight]{Ferziger4th.pdf}
\end{sourceblock}
总系数是对流和扩散贡献的总和，A c
和 A d。

边界节点处的 φ 值指定为： φ 1 = φ 0 且 φ N = φ L，其中 N 是包括边界处两个节点在内的节点数。这意味着，对于 i = 2 处的节点，项 A 2
W φ 1 可以计算出来并添加到右边的 Q 2 中，我们设置系数 A 2

\begin{sourceblock}
\begin{tikzpicture}
\node[inner sep=0,anchor=south west] (image) at (0,0) {\includegraphics[page=81,trim={53.00bp 370.13bp 53.87bp 277.75bp},clip,width=0.92\textwidth,max height=0.38\textheight]{Ferziger4th.pdf}};
\begin{scope}[x={(image.south east)},y={(image.north west)}]
\fill[white] (0.00000,0.83406) rectangle (0.05988,1.00000);
\fill[white] (0.06153,0.83406) rectangle (0.10959,1.00000);
\fill[white] (0.11128,0.83406) rectangle (0.15101,1.00000);
\fill[white] (0.15266,0.83406) rectangle (0.19077,1.00000);
\fill[white] (0.19239,0.83406) rectangle (0.23380,1.00000);
\fill[white] (0.23549,0.83406) rectangle (0.36595,1.00000);
\fill[white] (0.37071,0.83954) rectangle (0.40490,1.00000);
\fill[white] (0.38959,0.72015) rectangle (0.41420,1.00000);
\fill[white] (0.41774,0.83406) rectangle (0.44586,1.00000);
\fill[white] (0.44749,0.83406) rectangle (0.49723,1.00000);
\fill[white] (0.49892,0.83406) rectangle (0.60686,1.00000);
\fill[white] (0.60857,0.83406) rectangle (0.63669,1.00000);
\fill[white] (0.63835,0.83406) rectangle (0.70220,1.00000);
\fill[white] (0.70388,0.83406) rectangle (0.86728,1.00000);
\fill[white] (0.86902,0.83406) rectangle (0.90875,1.00000);
\fill[white] (0.91044,0.83406) rectangle (0.95850,1.00000);
\fill[white] (0.96015,0.83406) rectangle (1.00000,1.00000);
\fill[white] (0.00000,0.17908) rectangle (0.09564,0.81271);
\fill[white] (0.10039,0.18456) rectangle (0.16878,0.93428);
\fill[white] (0.88791,0.06572) rectangle (0.90553,0.53505);
\fill[white] (0.94241,0.19934) rectangle (0.97062,0.83242);
\fill[white] (0.97411,0.17908) rectangle (1.00000,0.81271);
\fill[white] (0.03528,0.00000) rectangle (0.08668,0.15772);
\fill[white] (0.08935,0.00000) rectangle (0.19901,0.15772);
\fill[white] (0.20168,0.00000) rectangle (0.33630,0.15772);
\fill[white] (0.33901,0.00000) rectangle (0.42704,0.15772);
\fill[white] (0.42968,0.00000) rectangle (0.45450,0.15772);
\fill[white] (0.45714,0.00000) rectangle (0.53185,0.15772);
\fill[white] (0.53450,0.00000) rectangle (0.62454,0.15772);
\fill[white] (0.62728,0.00000) rectangle (0.67125,0.15772);
\fill[white] (0.67389,0.00000) rectangle (0.73531,0.15772);
\fill[white] (0.73798,0.00000) rectangle (0.79606,0.15772);
\fill[white] (0.79874,0.00000) rectangle (0.89011,0.15772);
\fill[white] (0.89275,0.00000) rectangle (0.93420,0.15772);
\fill[white] (0.93684,0.00000) rectangle (1.00000,0.15772);
\end{scope}
\end{tikzpicture}
\end{sourceblock}
E = 0。由此产生的三对角方程组很容易求解。这里我们只讨论解；用于获得它们的求解器将在第 1 章中介绍。 5.

为了证明与 UDS 相关的虚假扩散以及使用 CDS 时发生振荡的可能性，我们将考虑 Pe = 50 ( L = 1 . 0, ρ = 1 . 0, u = 1 . 0, = 0 . 02, φ 0 = 0 且 φ L = 1 . 0) 的情况。我们从使用具有 11 个节点（10 个相等细分）的统一网格获得的结果开始。使用对流项的 CDS 和 UDS 以及扩散项的 CDS 获得的 φ( x ) 轮廓如图 3.8 所示。

UDS 解显然过度扩散；它对应于 Pe ≈ 18（而不是 50）的精确解。假扩散比真扩散更强！另一方面，CDS 解表现出严重的振荡。振荡是由于最后两点 φ 梯度的突然变化造成的。基于网格间距的局部佩克莱特数 Pe = ρ u x /，在每个节点处等于 5。

如果网格被细化，CDS 振荡就会减少，但当使用 21 个点时它们仍然存在。经过第二次细化（41 个网格节点）后，CDS 解无振荡且非常准确，见图 3.9。 UDSsolution 的精度也通过网格细化得到了提高，但对于 x > 0 仍然存在很大的误差。 8.

CDS 振荡取决于局部佩克莱特数的值。可以证明，如果每个网格节点的局部佩克莱数 Pe ≤ 2，则不会发生振荡（参见 Patankar 1980）。这是 CDS 解有界性的充分条件，但不是必要条件。所谓的混合方案（Spalding 1972）被设计为在Pe≥2的任何节点上从CDS切换到UDS。这限制性太大并且降低了精度。仅当溶液在局部佩克莱特数较高的区域快速变化时，才会出现振荡。

\begin{sourceblock}
\includegraphics[page=82,trim={53.00bp 398.48bp 53.87bp 52.00bp},clip,width=\textwidth,max height=0.38\textheight]{Ferziger4th.pdf}
\par\vspace{0.45\baselineskip}
\small 图 3.8 使用 CDS（左）和 UDS（右）作为对流项以及具有 11 个节点的均匀网格求解 Pe = 50 时的一维对流扩散方程
\end{sourceblock}
\begin{sourceblock}
\includegraphics[page=82,trim={53.00bp 192.09bp 53.87bp 291.40bp},clip,width=\textwidth,max height=0.38\textheight]{Ferziger4th.pdf}
\par\vspace{0.45\baselineskip}
\small 图 3.9 使用 CDS（左）和 UDS（右）作为对流项以及具有 41 个节点的均匀网格求解 Pe = 50 时的一维对流扩散方程
\end{sourceblock}
\begin{sourceblock}
\includegraphics[page=83,trim={53.00bp 433.48bp 53.87bp 52.00bp},clip,width=\textwidth,max height=0.38\textheight]{Ferziger4th.pdf}
\par\vspace{0.45\baselineskip}
\small 图 3.10 使用 CDS（左）和 UDS（右）作为对流项以及具有 11 个节点的非均匀网格（右端附近网格密集）求解 Pe = 50 时的一维对流扩散方程
\end{sourceblock}
为了证明这一点，我们使用具有 11 个节点的非均匀网格重复计算。最小和最大网格间距为 x min = x N − x N − 1 = 0.0125 且 x max = x 2 − x 1 = 0.31，对应于扩展因子re = 0.7，见式(3.21)。因此，最小局部佩克莱特数为 Pe，在右边界附近 min = 0.625，最大值为 Pe，在左边界附近 max = 15.5。因此，在 φ 发生强烈变化的区域中，局部佩克莱特数小于 2，而在 φ 几乎恒定的区域中，局部佩克莱数很大。使用 CDS 和 UDS 方案计算出的网格剖面如图 3.10 所示。 CDS 溶液中没有可见的振荡。此外，它与具有四倍节点数的均匀网格上的解一样准确。 UDS解的精度也通过使用非均匀网格得到了提高，但仍然不可接受。

由于该问题有解析解，方程（3.56），我们可以直接计算数值解中的误差。使用以下平均误差作为衡量标准：

\begin{sourceblock}
\begin{tikzpicture}
\node[inner sep=0,anchor=south west] (image) at (0,0) {\includegraphics[page=83,trim={53.00bp 196.38bp 53.87bp 440.40bp},clip,width=0.92\textwidth,max height=0.38\textheight]{Ferziger4th.pdf}};
\begin{scope}[x={(image.south east)},y={(image.north west)}]
\fill[white] (0.00000,0.82016) rectangle (0.11227,1.00000);
\fill[white] (0.36343,0.29087) rectangle (0.38135,0.68460);
\fill[white] (0.38650,0.29087) rectangle (0.41468,0.68460);
\end{scope}
\end{tikzpicture}
\end{sourceblock}
使用 CDS 和 UDS 以及最多 321 个节点的均匀和非均匀网格解决了该问题。平均误差与平均网格间距的函数关系如图 3.11 所示。 UDS 误差渐近地接近一阶方案的预期斜率。 CDS 显示，从第二个网格开始，二阶方案预期的斜率：当网格间距减小一个数量级时，误差减小两个数量级。

\begin{center}\small 图 3.11 Pe = 50 时一维对流扩散方程解的平均误差与平均网格间距的关系\end{center}
这个例子清楚地表明，非均匀网格上的解以与均匀网格上的解相同的方式收敛，即使截断误差包含一阶项，如第二节中所述。 3.3.4.对于CDS，非均匀网格上的平均误差几乎比具有相同网格节点数的均匀网格上的平均误差小一个数量级。这是因为网格间距较小，误差较大。图 3.11 表明 UDS 在非均匀网格上的误差比均匀网格上的误差大，这是因为均匀网格上少数节点的大误差对平均值的影响较小；均匀网格上的最大节点误差比非均匀网格上的最大节点误差大得多，通过检查图 2-3 可以看出。 3.8 和 3.10。

相关示例请参见下一章最后一节。

\section*{3.11 谱方法简介}\phantomsection\addcontentsline{toc}{section}{3.11 谱方法简介}
频谱方法是一类比 FV 和 FE 方法不太适合通用 CFD 代码的方法，但它们在许多应用中都很重要（例如海洋和大气的全球和气候模拟代码（Washington 和 Parkinson 2005）、大气边界层的高分辨率中尺度建模（Moeng 1984）以及湍流模拟（Moin 和 Kim 1982））。 Some spectral methods are briefly described here. For a more complete description of them, see the books by Canuto et al. ( 2006 and 2007 ), Boyd ( 2001 ), Durran ( 2010 ), and Moin ( 2010 ).

\subsection*{3.11.1 分析工具}\phantomsection\addcontentsline{toc}{subsection}{3.11.1 分析工具}
在光谱方法中，空间导数是借助傅里叶级数或其概括之一来评估的。最简单的谱方法处理由一组均匀间隔的点处的值指定的周期函数。可以用离散傅里叶级数来表示这样的函数：

\begin{sourceblock}
\includegraphics[page=85,trim={53.00bp 462.20bp 53.87bp 167.75bp},clip,width=0.92\textwidth,max height=0.38\textheight]{Ferziger4th.pdf}
\end{sourceblock}
其中 x i = i x , i = 1 , 2 , .。。 N 和 k q = 2 π q / x N。 Equation ( 3.64 ) can be inverted in a surprisingly simple way:

\begin{sourceblock}
\includegraphics[page=85,trim={53.00bp 382.73bp 53.87bp 247.88bp},clip,width=0.92\textwidth,max height=0.38\textheight]{Ferziger4th.pdf}
\end{sourceblock}
as can be proven by using the well-known formula for the summation of geometric series. q 的值集有些任意；将索引从 q 更改为 q ± lN （其中 l 是整数）不会导致网格点处的 e ± i k q x i 值发生变化。此属性称为别名；混叠是非线性微分方程（包括不使用谱方法的方程）数值解中常见且重要的误差源。我们在第二节中再次注意到这一点。 10.3.4.3 与频谱代码在湍流模拟中的应用有关。

这些级数的有用之处在于，方程（3.64）可用于对 f(x) 进行插值。我们简单地将离散变量 x i 替换为连续变量 x；然后为所有 x 定义 f ( x )，而不仅仅是网格点。现在q的范围的选择变得非常重要。不同的 q 集合产生不同的插值；最好的选择是给出最平滑插值的集合，即方程（3.64）中使用的插值。 （集合 − N / 2 + 1 , . . . , N / 2 与所选的一样好。）定义了插值后，我们可以对它进行微分以生成导数的傅立叶级数：

\begin{sourceblock}
\includegraphics[page=85,trim={53.00bp 160.70bp 53.87bp 469.25bp},clip,width=0.92\textwidth,max height=0.38\textheight]{Ferziger4th.pdf}
\end{sourceblock}
这表明 d f / d x 的傅立叶系数是 i k q ˆ f ( k q )。这提供了一种计算导数的方法：

\begin{itemize}
\item 给定f ( x i )，使用方程( 3.65 ) 计算其傅立叶系数ˆ f ( k q )；
\item 计算g = d f / d x 的傅立叶系数； \textasciicircum{} g ( k q ) = i k q \textasciicircum{} f ( k q ) ;
\item 评估序列( 3.66 ) 以获得网格点处的g = d f / d x。
\end{itemize}
有几点需要注意。

\begin{itemize}
\item 该方法很容易推广到更高的导数；例如，d 2 f / d x 2 的傅立叶系数为 − k 2
\end{itemize}
q ˆ f ( k q )。 • 如果f(x) 在x 中是周期性的，则当网格点的数量N 较大时，计算导数的误差随N 呈指数下降。这使得对于大 N 的谱方法比有限差分方法更加准确；然而，对于小 N 来说，情况可能并非如此。 “大”的定义取决于功能。 • 使用方程(3.65) 计算傅立叶系数和/或使用方程(3.64) 计算逆矩阵的成本，如果以最明显的方式完成，则按N 2 缩放。这将是非常昂贵的；该方法因计算傅立叶变换 (FFT) 的快速方法而变得实用，其成本与 N log 2 N 成正比。

为了获得这种特定谱方法的优点，函数必须是周期性的并且网格点均匀间隔。这些条件可以通过使用复指数以外的函数来放松，但几何或边界条件的任何变化都需要方法进行相当大的改变，使得谱方法相对不灵活。对于它们非常适合的问题（例如，几何简单域中的湍流模拟），它们是无与伦比的。

3.11.1.2 离散化误差的另一种观点
谱方法对于提供另一种查看截断误差的方法非常有用，就像它们本身的计算方法一样。只要我们处理周期函数，级数（3.64）就代表该函数，并且我们可以通过我们选择的任何方法来近似其导数。特别是，我们可以使用上面示例的精确谱方法或有限差分近似。这些方法中的任何一种都可以逐项应用于序列，因此考虑 e i kx 的微分就足够了。确切的结果是 i k e i kx。另一方面，如果我们将方程（3.9）的中心差分算子应用于该函数，我们会发现：

\begin{sourceblock}
\includegraphics[page=86,trim={53.00bp 177.12bp 53.87bp 457.01bp},clip,width=0.92\textwidth,max height=0.38\textheight]{Ferziger4th.pdf}
\end{sourceblock}
其中 k eff 称为有效波数 2，因为使用有限差分近似相当于用 k eff 替换精确的波数 k。对于其他方案可以推导出类似的表达式；例如，四阶 CDS，方程（3.14），导致：

\begin{sourceblock}
\begin{tikzpicture}
\node[inner sep=0,anchor=south west] (image) at (0,0) {\includegraphics[page=86,trim={53.00bp 95.86bp 53.87bp 541.89bp},clip,width=0.92\textwidth,max height=0.38\textheight]{Ferziger4th.pdf}};
\begin{scope}[x={(image.south east)},y={(image.north west)}]
\fill[white] (0.00000,0.88975) rectangle (0.12866,1.00000);
\fill[white] (0.13137,0.88975) rectangle (0.19771,1.00000);
\fill[white] (0.20042,0.88975) rectangle (0.23687,1.00000);
\end{scope}
\end{tikzpicture}
\end{sourceblock}
\footnotetext[2]{一些作者将其称为修正波数。}
\begin{center}\small 图 3.12 一阶导数的二阶和四阶 CDS 近似的有效波数，通过 k max = π/ x 归一化\end{center}
1.0

0.8

精确和光谱

0.6

四阶CDS

有效值

*
0.4

0.2

二阶CDS

0.0 0.0 0.2 0.4 0.6 0.8 1.0
*
k

对于低波数（对应于平滑函数），CDS 近似的有效波数可以展开为泰勒级数：

\begin{sourceblock}
\includegraphics[page=87,trim={53.00bp 404.76bp 53.87bp 232.03bp},clip,width=0.92\textwidth,max height=0.38\textheight]{Ferziger4th.pdf}
\end{sourceblock}
这显示了小 k 和小 x 的近似的二阶性质。然而，在任何计算中，可能会遇到高达 k max = π/ x 的波数（参见方程（3.64））。给定傅里叶系数的大小取决于其导数被近似的函数；平滑函数具有较小的高波数分量，但快速变化的函数给出的傅立叶系数随波数缓慢减小。

在图 3.12 中，二阶和四阶 CDS 方案的有效波数由 k max 归一化，显示为归一化波数 k * = k / k max 的函数。如果波数大于最大值的一半，两种方案都会给出较差的近似值。随着网格的细化，会包含更多的波数。在小间距的限制下，函数相对于网格是平滑的，只有小波数才有大系数，并且可以期望得到准确的结果。或者，修改后的波数越接近实际波数，结果越准确。

如果我们解决的问题的解不是很平滑，则离散化方法的阶数可能不再是其准确性的良好指标。人们需要非常小心地声称某一特定方案是准确的，因为所使用的方法是高阶的。仅当解中最高波数的每个波长有足够的节点时，结果才是准确的。

当网格大小的任何功率变为零时，频谱方法产生的误差下降得更快。这经常被认为是该方法的优点。然而，只有当使用足够的点时才会获得这种行为（“足够”的定义取决于函数）。对于少量网格点，谱方法实际上可能比有限差分方法产生更大的误差。

最后，我们注意到迎风差法的有效波数为：

\begin{sourceblock}
\includegraphics[page=88,trim={53.00bp 582.63bp 53.87bp 53.89bp},clip,width=0.92\textwidth,max height=0.38\textheight]{Ferziger4th.pdf}
\end{sourceblock}
并且很复杂；这表明了这种近似的扩散或耗散性质。前者在上一节的 UDS 示例中得到了体现。见第 6.3 节，我们将看到，当用于非定常微分方程时，这种迎风近似是耗散的，例如，导致传播波的振幅随时间不自然地快速衰减。

\subsection*{3.11.2微分方程的解}\phantomsection\addcontentsline{toc}{subsection}{3.11.2微分方程的解}
在这里，我们研究两种著名的求解微分方程的谱方法（参见 Boyd 2001；Durran 2010；或 Moin 2010）。它们基于一系列项中未知解的扩展，例如方程（3.64），以及许多其他谱方法（Canuto et al. 2007），包括谱元素、修补搭配、谱不连续伽辽金等。

第一个公式称为微分方程的弱形式，因为满足方程在解域上的加权积分，但微分方程并非在每个点都满足。伽辽金公式就是这种类型。第二个公式是微分方程的强形式，并且要求域的每个点都满足该方程。该公式的一个重要变体称为光谱搭配或伪光谱方法。在这种情况下，虽然解由截断展开式表示，但微分方程在域中的一组有限网格点处得到满足。该方法用于流体力学，例如全球环流模型和大气边界层的湍流建模（Fox and Orszag 1973；Moeng 1984；Pekurovsky et al. 2006；Sullivan and Patton 2011），因为纳维-斯托克斯或欧拉方程中非线性项的导数可以在物理空间中计算，比在物理空间中计算效率更高且成本更低。在光谱空间中进行评估。

对于这一发展，请考虑污染物在狭窄盆地中扩散的一维方程。污染物在流域的长度上不均匀地流入流域，并在两端被提取，从而使那里的浓度保持为零。如果污染物的扩散率为D x，污染物浓度为 φ，则在长度为 π 的通道上可提出以下边值问题：

\begin{sourceblock}
\includegraphics[page=88,trim={53.00bp 143.07bp 53.87bp 493.82bp},clip,width=0.92\textwidth,max height=0.38\textheight]{Ferziger4th.pdf}
\end{sourceblock}
with

\begin{sourceblock}
\begin{tikzpicture}
\node[inner sep=0,anchor=south west] (image) at (0,0) {\includegraphics[page=88,trim={53.00bp 111.31bp 53.87bp 540.50bp},clip,width=0.92\textwidth,max height=0.38\textheight]{Ferziger4th.pdf}};
\begin{scope}[x={(image.south east)},y={(image.north west)}]
\fill[white] (0.00000,0.91835) rectangle (0.05741,1.00000);
\end{scope}
\end{tikzpicture}
\end{sourceblock}
这里 A ( x ) 表示污染物沿通道的流入，在这个一维问题中，由源项而不是边界上的入口表示（如多维问题中的情况）。

3.11.2.1 使用弱形式和傅里叶级数
这里使用的方法正式称为傅立叶-伽辽金方法（Canuto et al. 2006；Boyd 2001）。该方法基于两个关键原则，即一组完整的基函数能够准确表示任意但合理的函数和正交性概念（Street 1973；Boyd 2001；Moin 2010）。基本思想是用一系列由基函数乘以未知系数组成的序列来表示未知变量[参见方程（3.64）]。然后，通过将给定方程（例如上面的（3.71））乘以一组测试函数并在问题的域上积分来使用基函数的正交性。在此处使用的傅立叶-伽辽金方法中，基函数和测试函数是相同的。

对于方程（3.71）和（3.72）中提出的当前问题，边界条件是齐次狄利克雷条件，而不是周期性条件，因此使用三角正弦函数可以方便地扩展解，作为半范围扩展（即，在 0 < x < π 上）：

\begin{sourceblock}
\begin{tikzpicture}
\node[inner sep=0,anchor=south west] (image) at (0,0) {\includegraphics[page=89,trim={53.00bp 323.64bp 53.87bp 307.02bp},clip,width=0.92\textwidth,max height=0.38\textheight]{Ferziger4th.pdf}};
\begin{scope}[x={(image.south east)},y={(image.north west)}]
\fill[white] (0.00000,0.94335) rectangle (0.01913,1.00000);
\fill[white] (0.02280,0.95378) rectangle (0.05098,1.00000);
\fill[white] (0.05531,0.94617) rectangle (0.07341,1.00000);
\fill[white] (0.07961,0.95378) rectangle (0.10779,1.00000);
\fill[white] (0.11131,0.94335) rectangle (0.15543,1.00000);
\end{scope}
\end{tikzpicture}
\end{sourceblock}
上标 N 只是表示解由 N 个元素的部分（有限）和表示。回想一下 e i kx = cos ( kx ) + i sin ( kx ) 并注意部分和中的每个元素都满足问题的边界条件。一般来说，部分和 φ N ( x ) 将不满足控制方程（3.71）。因此，每个点都会有一个残差：

\begin{sourceblock}
\begin{tikzpicture}
\node[inner sep=0,anchor=south west] (image) at (0,0) {\includegraphics[page=89,trim={53.00bp 226.75bp 53.87bp 409.25bp},clip,width=0.92\textwidth,max height=0.38\textheight]{Ferziger4th.pdf}};
\begin{scope}[x={(image.south east)},y={(image.north west)}]
\fill[white] (0.00000,0.91307) rectangle (0.10120,1.00000);
\fill[white] (0.10391,0.90975) rectangle (0.13035,1.00000);
\fill[white] (0.13305,0.90975) rectangle (0.19272,1.00000);
\fill[white] (0.19546,0.90975) rectangle (0.27017,1.00000);
\end{scope}
\end{tikzpicture}
\end{sourceblock}
该残差等于部分和提供的解估计中的误差。有多种方法可以最小化此残差（例如，参见 Boyd 2001；或 Durran 2010），包括通过最小二乘技术或使用加权残差进行最小化，这是此处说明的方法。在我们的加权残差方法中，我们要求方程的加权平均值为零；这是一个弱解，可以让我们确定未知系数 ˆ φ q。

现在将部分和 ( 3.73 ) 代入方程 ( 3.71 )，将整个方程乘以权重函数 sin ( jx )，并在域上积分得到：

\begin{sourceblock}
\begin{tikzpicture}
\node[inner sep=0,anchor=south west] (image) at (0,0) {\includegraphics[page=89,trim={53.00bp 84.46bp 53.87bp 545.57bp},clip,width=0.92\textwidth,max height=0.38\textheight]{Ferziger4th.pdf}};
\begin{scope}[x={(image.south east)},y={(image.north west)}]
\fill[white] (0.29203,0.57657) rectangle (0.30944,0.81390);
\fill[white] (0.53708,0.24896) rectangle (0.58848,0.57934);
\fill[white] (0.58860,0.25173) rectangle (0.63260,0.57934);
\fill[white] (0.63780,0.24896) rectangle (0.67233,0.57186);
\fill[white] (0.67850,0.25921) rectangle (0.70668,0.57934);
\fill[white] (0.71020,0.24896) rectangle (0.72998,0.56909);
\fill[white] (0.73035,0.25921) rectangle (0.74340,0.57934);
\fill[white] (0.92439,0.24896) rectangle (1.00000,0.56909);
\end{scope}
\end{tikzpicture}
\end{sourceblock}
对于每个 j = 1 , 2 , ... 都必须满足该方程。。。尼。此时，由于源项 A ( x )/ D x 假定已知，但可以采用任意形式，因此有必要将源项展开为正弦级数：

\begin{sourceblock}
\includegraphics[page=90,trim={53.00bp 526.49bp 53.87bp 104.17bp},clip,width=0.92\textwidth,max height=0.38\textheight]{Ferziger4th.pdf}
\end{sourceblock}
正如我们将在下面看到的，源项的形式会影响解，特别是当源存在不连续性时。

将部分和代入方程（3.75）并重新排列产生

\begin{sourceblock}
\includegraphics[page=90,trim={53.00bp 435.19bp 53.87bp 191.90bp},clip,width=0.92\textwidth,max height=0.38\textheight]{Ferziger4th.pdf}
\end{sourceblock}
然而，因为所选择的扩展函数和加权函数是相同的并且是正交的，i。例如，

\begin{sourceblock}
\includegraphics[page=90,trim={53.00bp 362.69bp 53.87bp 266.92bp},clip,width=0.92\textwidth,max height=0.38\textheight]{Ferziger4th.pdf}
\end{sourceblock}
我们得到：

\begin{sourceblock}
\begin{tikzpicture}
\node[inner sep=0,anchor=south west] (image) at (0,0) {\includegraphics[page=90,trim={53.00bp 320.35bp 53.87bp 317.87bp},clip,width=0.92\textwidth,max height=0.38\textheight]{Ferziger4th.pdf}};
\begin{scope}[x={(image.south east)},y={(image.north west)}]
\fill[white] (0.00000,0.81483) rectangle (0.03907,1.00000);
\fill[white] (0.04174,0.81483) rectangle (0.12974,1.00000);
\end{scope}
\end{tikzpicture}
\end{sourceblock}
将方程（3.76）乘以加权函数，在域上积分并再次使用正交性得到

\begin{sourceblock}
\begin{tikzpicture}
\node[inner sep=0,anchor=south west] (image) at (0,0) {\includegraphics[page=90,trim={53.00bp 249.90bp 53.87bp 380.57bp},clip,width=0.92\textwidth,max height=0.38\textheight]{Ferziger4th.pdf}};
\begin{scope}[x={(image.south east)},y={(image.north west)}]
\fill[white] (0.53982,0.25203) rectangle (0.59125,0.58649);
\fill[white] (0.59134,0.25484) rectangle (0.63537,0.58649);
\fill[white] (0.64054,0.25203) rectangle (0.67507,0.57892);
\fill[white] (0.67792,0.26241) rectangle (0.69098,0.58649);
\fill[white] (0.92421,0.25203) rectangle (1.00000,0.57611);
\end{scope}
\end{tikzpicture}
\end{sourceblock}
那么最终的解是

\begin{sourceblock}
\begin{tikzpicture}
\node[inner sep=0,anchor=south west] (image) at (0,0) {\includegraphics[page=90,trim={53.00bp 186.76bp 53.87bp 439.90bp},clip,width=0.92\textwidth,max height=0.38\textheight]{Ferziger4th.pdf}};
\begin{scope}[x={(image.south east)},y={(image.north west)}]
\fill[white] (0.45832,0.61272) rectangle (0.47570,0.82979);
\fill[white] (0.29314,0.48987) rectangle (0.31293,0.78293);
\fill[white] (0.16935,0.31560) rectangle (0.25140,0.65400);
\fill[white] (0.25495,0.32244) rectangle (0.28313,0.61525);
\end{scope}
\end{tikzpicture}
\end{sourceblock}
部分和（3.81）是方程（3.71）的近似（弱）谱解，满足边界条件（3.72）。由于部分和中的 N → 无穷大，读者将认识到该解成为这个简单情况下的标准傅立叶级数解。上面引用的参考文献涉及更复杂的情况以及其他基础和测试函数。

\begin{center}\small 图 3.13 显示了在域中心选择污染输入的结果。窄步长输入近似于 delta 函数，因此解在源附近缓慢收敛，即系数收敛的速率\end{center}
\begin{sourceblock}
\includegraphics[page=91,trim={53.00bp 349.26bp 53.87bp 52.00bp},clip,width=\textwidth,max height=0.38\textheight]{Ferziger4th.pdf}
\par\vspace{0.45\baselineskip}
\small 图 3.13 单位强度源位于 ( 0 . 49 − 0 . 51 )π 之间的污染输入问题的解：显示了源/20（实线）和使用部分展开式的 1 (+)、2 ( −· − )、3 ( • ) 和 512 ( ⋆ ) 项的解
\end{sourceblock}
ˆ φ q 随着 q 的增加而减少，在源项附近快速变化时较小。因此，需要更多的项来达到所需的准确度。在该区域之外，精确解是线性的。同样在这种情况下，部分展开式中的偶数项 (2,4,6, ...) 为零，如图中 1 项和 2 项相同的部分和所示。 100 项展开已经与精确解本质上相同。

3.11.2.2 使用强形式并与傅立叶级数搭配
对于强形式伪谱方法，过程相当简单。考虑到方程（3.71）和（3.72）以及部分展开（3.73）提出的问题，我们只需要在域中的 N − 1 个等距（网格或节点）点处满足方程（3.71）即可。 3 由于部分和已经满足边界条件，因此无需特殊考虑；参见 Boyd (2001) 的讨论

\footnotetext[3]{Boyd (2001) 第 3 章指出，这种伪谱约束可以通过使用 Dirac δ 函数 δ( x − x i ) 作为上述加权残差方法中的测试函数来获得。}
边界条件必须作为方程添加到此处获得的最终集合中的情况。

对于傅立叶级数近似，域中的点可以等距分布为 δ x j = π/ N，其中 j = 1 , 2 , 3 , ...。。。 , N − 1 所以 x j = π j / N。然而，很明显，输入部分展开式的微分方程只能在 N − 1 个点处得到满足，因此将方程（3.73）修改为：

\begin{sourceblock}
\includegraphics[page=92,trim={53.00bp 490.62bp 53.87bp 139.85bp},clip,width=0.92\textwidth,max height=0.38\textheight]{Ferziger4th.pdf}
\end{sourceblock}
所需要的是

\begin{sourceblock}
\begin{tikzpicture}
\node[inner sep=0,anchor=south west] (image) at (0,0) {\includegraphics[page=92,trim={53.00bp 429.68bp 53.87bp 200.35bp},clip,width=0.92\textwidth,max height=0.38\textheight]{Ferziger4th.pdf}};
\begin{scope}[x={(image.south east)},y={(image.north west)}]
\fill[white] (0.07095,0.45223) rectangle (0.16743,0.80642);
\fill[white] (0.92439,0.00000) rectangle (1.00000,0.09665);
\end{scope}
\end{tikzpicture}
\end{sourceblock}
(3.83) 使用方程 (3.82) 得出

\begin{sourceblock}
\includegraphics[page=92,trim={53.00bp 363.84bp 53.87bp 266.63bp},clip,width=0.92\textwidth,max height=0.38\textheight]{Ferziger4th.pdf}
\end{sourceblock}
产生 ( N − 1 ) × ( N − 1 ) 系数矩阵形式的 N − 1 个 ˆ φ q 方程。对于 N = 4，x j = π j / 4 所以

\begin{sourceblock}
\begin{tikzpicture}
\node[inner sep=0,anchor=south west] (image) at (0,0) {\includegraphics[page=92,trim={53.00bp 290.08bp 53.87bp 341.87bp},clip,width=0.92\textwidth,max height=0.38\textheight]{Ferziger4th.pdf}};
\begin{scope}[x={(image.south east)},y={(image.north west)}]
\fill[white] (0.70217,0.04300) rectangle (0.75417,0.44048);
\end{scope}
\end{tikzpicture}
\end{sourceblock}
\begin{sourceblock}
\begin{tikzpicture}
\node[inner sep=0,anchor=south west] (image) at (0,0) {\includegraphics[page=92,trim={53.00bp 176.44bp 53.87bp 373.69bp},clip,width=0.92\textwidth,max height=0.38\textheight]{Ferziger4th.pdf}};
\begin{scope}[x={(image.south east)},y={(image.north west)}]
\fill[white] (0.24974,0.98991) rectangle (0.27284,1.00000);
\fill[white] (0.27696,0.99224) rectangle (0.30514,1.00000);
\fill[white] (0.33708,0.98310) rectangle (0.42671,1.00000);
\fill[white] (0.43991,0.98991) rectangle (0.45970,1.00000);
\fill[white] (0.45994,0.98310) rectangle (0.56472,1.00000);
\fill[white] (0.57786,0.98991) rectangle (0.59765,1.00000);
\fill[white] (0.59783,0.98310) rectangle (0.70250,1.00000);
\fill[white] (0.70217,0.99224) rectangle (0.75417,1.00000);
\fill[white] (0.33714,0.88001) rectangle (0.42668,0.98966);
\fill[white] (0.43994,0.88682) rectangle (0.45973,0.98647);
\fill[white] (0.46000,0.88001) rectangle (0.56466,0.98966);
\fill[white] (0.57792,0.88682) rectangle (0.59771,0.98647);
\fill[white] (0.59786,0.88001) rectangle (0.70244,0.98966);
\end{scope}
\end{tikzpicture}
\end{sourceblock}
通过求解方程组得到解（见5.2节）

\begin{sourceblock}
\includegraphics[page=92,trim={53.00bp 129.98bp 53.87bp 519.84bp},clip,width=0.92\textwidth,max height=0.38\textheight]{Ferziger4th.pdf}
\end{sourceblock}
可以针对任何 N 值生成。

将配置方法应用于上述傅里叶-伽辽金方法解决的问题会产生类似的结果。当平均误差根据方程（3.63）并由域中的最大值归一化时，如图3.14所示

\begin{sourceblock}
\includegraphics[page=93,trim={53.00bp 379.25bp 53.87bp 52.00bp},clip,width=\textwidth,max height=0.38\textheight]{Ferziger4th.pdf}
\par\vspace{0.45\baselineskip}
\small 图 3.14 傅里叶-伽辽金或配置方法的误差与基函数或节点数量以及污染源形状的关系
\end{sourceblock}
对于作为所使用的基函数数量（或者等效地对于配置方法，所使用的节点或网格点的数量）的函数的弱解和强解，揭示了误差的显著差异。

我们在这里使用了三种不同的源形状来说明这对错误的影响。窄盒源如图 3.13 所示，仅在 0 范围内具有恒定的非零值。 49≥x/π≥0.51；宽框有一个从 0 开始的恒定非零源。 45≥x/π≥0.55. 源以及相应的 φ 的二阶导数在这些边缘处是不连续的。傅里叶-伽辽金弱解对这些不连续点进行积分，并且要求微分方程仅在平均值上满足，而不是逐点满足。对解的二阶导数的部分和的检查表明（图 3.15），它在不连续处表现出吉布斯现象（Ferziger 1998；Gibbs 1898，1899），即部分和的值超过和低于正确值（它摆动！），因为级数在那里不收敛。然而，解的一阶导数是连续的，所以那里不会出现这个问题，解本身也没有被感染。事实上，由于该解基于源项的分析积分，因此源材料的质量始终保持不变。如果以数值方式进行积分，则会引入一些误差，该误差取决于积分节点的间距和积分方法。

搭配方法可以明确地看到不连续性，因为我们强制微分方程在跨越不连续性的节点上得到满足，并且源没有准确地表示。这会导致持续出现错误

\begin{sourceblock}
\begin{tikzpicture}
\node[inner sep=0,anchor=south west] (image) at (0,0) {\includegraphics[page=94,trim={53.00bp 597.81bp 53.87bp 57.11bp},clip,width=0.92\textwidth,max height=0.38\textheight]{Ferziger4th.pdf}};
\begin{scope}[x={(image.south east)},y={(image.north west)}]
\fill[white] (0.53320,0.10695) rectangle (0.56821,0.89305);
\end{scope}
\end{tikzpicture}
\par\vspace{0.45\baselineskip}
\small 图 3.15 图 3.13 污染源 Fourier-Galerkin 部分和解的二阶导数吉布斯现象：1024 项
\end{sourceblock}
\begin{sourceblock}
\begin{tikzpicture}
\node[inner sep=0,anchor=south west] (image) at (0,0) {\includegraphics[page=94,trim={53.00bp 489.68bp 53.87bp 100.36bp},clip,width=0.92\textwidth,max height=0.38\textheight]{Ferziger4th.pdf}};
\begin{scope}[x={(image.south east)},y={(image.north west)}]
\fill[white] (0.00000,0.89146) rectangle (0.08620,1.00000);
\fill[white] (0.08776,0.89146) rectangle (0.12226,1.00000);
\fill[white] (0.12385,0.89146) rectangle (0.22063,1.00000);
\fill[white] (0.22220,0.89146) rectangle (0.29344,1.00000);
\fill[white] (0.48445,0.86465) rectangle (0.51386,0.94179);
\fill[white] (0.00000,0.76058) rectangle (0.02535,0.89290);
\fill[white] (0.02692,0.76058) rectangle (0.12695,0.89290);
\fill[white] (0.12857,0.76058) rectangle (0.18430,0.89290);
\fill[white] (0.18589,0.76058) rectangle (0.24728,0.89290);
\fill[white] (0.48556,0.74520) rectangle (0.51498,0.85585);
\fill[white] (0.48556,0.23745) rectangle (0.51498,0.73627);
\fill[white] (0.48556,0.02825) rectangle (0.51498,0.22865);
\fill[white] (0.52051,0.01577) rectangle (0.56821,0.13167);
\end{scope}
\end{tikzpicture}
\end{sourceblock}
\begin{sourceblock}
\begin{tikzpicture}
\node[inner sep=0,anchor=south west] (image) at (0,0) {\includegraphics[page=94,trim={53.00bp 468.07bp 53.87bp 186.85bp},clip,width=0.92\textwidth,max height=0.38\textheight]{Ferziger4th.pdf}};
\begin{scope}[x={(image.south east)},y={(image.north west)}]
\fill[white] (0.53865,0.10695) rectangle (0.56821,0.89305);
\end{scope}
\end{tikzpicture}
\end{sourceblock}
\begin{sourceblock}
\begin{tikzpicture}
\node[inner sep=0,anchor=south west] (image) at (0,0) {\includegraphics[page=94,trim={53.00bp 428.00bp 53.87bp 213.44bp},clip,width=0.92\textwidth,max height=0.38\textheight]{Ferziger4th.pdf}};
\begin{scope}[x={(image.south east)},y={(image.north west)}]
\fill[white] (0.55666,0.59433) rectangle (0.59164,0.95142);
\fill[white] (0.65239,0.59433) rectangle (0.69943,0.95142);
\fill[white] (0.76021,0.59433) rectangle (0.79519,0.95142);
\fill[white] (0.85594,0.59433) rectangle (0.90298,0.95142);
\fill[white] (0.96382,0.59433) rectangle (0.99883,0.95142);
\end{scope}
\end{tikzpicture}
\end{sourceblock}
存在的源材料的量。该误差取决于网格间距及其与规定的非零源面积的关系。标准化源体积以节省总体质量可减少误差。

改变盒形源的宽度会对减少大量节点或基函数的误差产生一些影响，但影响不大。然而，通过使用平滑源 A ( x )/ D x = 1 − cos ( 2 x )，图 3.14 清楚地揭示了源形状的影响。现在，随着项或节点数量的增加，两种方法的误差都会迅速减小。

最后，考虑到弱方法和强方法中确定部分展开系数的不同方法，需要注意的是，即使基函数相同，两种方法的展开系数也可能不同。

搭配方法的应用一般来说很简单，但是非线性项的存在（如纳维-斯托克斯方程中出现的那样）需要特殊处理，这超出了本介绍的范围（参见 Boyd 2001；Moin 2010；Canuto 等人 2006， 2007 ）。

\chapter{有限体积法}
\section*{4.1 简介}\phantomsection\addcontentsline{toc}{section}{4.1 简介}
与前一章一样，我们仅考虑量 φ 的通用守恒方程，并假设速度场和所有流体属性已知。有限体积法以守恒方程的积分形式为起点：

\begin{sourceblock}
\begin{tikzpicture}
\node[inner sep=0,anchor=south west] (image) at (0,0) {\includegraphics[page=95,trim={53.00bp 325.42bp 53.87bp 318.24bp},clip,width=0.92\textwidth,max height=0.38\textheight]{Ferziger4th.pdf}};
\begin{scope}[x={(image.south east)},y={(image.north west)}]
\fill[white] (0.25251,0.42037) rectangle (0.30920,0.94662);
\fill[white] (0.31152,0.43194) rectangle (0.32466,0.94662);
\fill[white] (0.32653,0.42037) rectangle (0.34800,0.93461);
\fill[white] (0.35353,0.41593) rectangle (0.39047,0.93461);
\fill[white] (0.39534,0.43194) rectangle (0.42352,0.94662);
\fill[white] (0.23459,0.05338) rectangle (0.24989,0.43461);
\end{scope}
\end{tikzpicture}
\end{sourceblock}
解域被网格细分为有限数量的小控制体积 (CV)，与有限差分 (FD) 方法相反，网格定义了控制体积边界，而不是计算节点。为了简单起见，我们将使用笛卡尔网格来演示该方法；复杂的几何形状在第 1 章中进行了处理。 9.

通常的方法是通过合适的网格定义 CV，并将计算节点分配给 CV 中心。然而，也可以（对于结构化网格）首先定义节点位置并围绕它们构建 CV，以便 CV 面位于节点之间的中间；见图4.1。应用边界条件的节点在此图中显示为完整符号。

第一种方法的优点是节点值代表 CV 体积的平均值，其精度（二阶）比第二种方法更高，因为节点位于 CV 的质心。第二种方法的优点是，当面位于两个节点之间时，CV 面导数的 CDS 近似更加准确。第一个变体使用得更频繁，本书将采用第一个变体。

FV 类型方法还有其他几种专门的变体（单元顶点方案、双网格方案等）；其中一些将在本章后面和第 1 章中描述。 9.这里我们只描述基本方法。离散化原理

\begin{sourceblock}
\includegraphics[page=96,trim={53.00bp 472.25bp 53.87bp 52.00bp},clip,width=\textwidth,max height=0.38\textheight]{Ferziger4th.pdf}
\par\vspace{0.45\baselineskip}
\small 图 4.1 FV 网格的类型：以 CV 为中心的节点（左）和以节点为中心的 CV 面（右）
\end{sourceblock}
所有变体的原则都是相同的——只需考虑积分体积内不同位置之间的关系。

积分守恒方程 (4.1) 适用于每个 CV，也适用于整个解域。如果我们对所有 CV 的方程求和，我们就得到了全局守恒方程，因为内部 CV 面上的表面积分相互抵消。因此，全球保护被纳入该方法中，这提供了其主要优点之一。

为了获得特定 CV 的代数方程，需要使用求积公式来近似表面和体积积分。根据所使用的近似值，所得方程可能是也可能不是从 FD 方法获得的方程。

\section*{4.2 曲面积分的近似}\phantomsection\addcontentsline{toc}{section}{4.2 曲面积分的近似}
在图中。 4.2 和 4.3 中显示了典型的 2D 和 3D 笛卡尔控制体积以及我们将使用的符号。 CV 曲面由四个（2D 中）或六个（3D 中）平面组成，用与其相对于中心节点 (P) 的方向（e、w、n、s、t 和 b）相对应的小写字母表示。 2D 情况可以被视为 3D 情况的一种特殊情况，其中因变量独立于 z。在本章中，我们将主要讨论二维网格； 3D 问题的扩展很简单。

通过 CV 边界的净通量是四个（2D 中）或六个（3D 中）CV 面的积分之和：

\begin{sourceblock}
\begin{tikzpicture}
\node[inner sep=0,anchor=south west] (image) at (0,0) {\includegraphics[page=96,trim={53.00bp 99.27bp 53.87bp 541.39bp},clip,width=0.92\textwidth,max height=0.38\textheight]{Ferziger4th.pdf}};
\begin{scope}[x={(image.south east)},y={(image.north west)}]
\fill[white] (0.38577,0.48862) rectangle (0.39892,0.94231);
\fill[white] (0.41035,0.48469) rectangle (0.44725,0.94231);
\fill[white] (0.45215,0.49922) rectangle (0.48033,0.95290);
\fill[white] (0.36129,0.16484) rectangle (0.37660,0.50157);
\end{scope}
\end{tikzpicture}
\end{sourceblock}
\begin{sourceblock}
\includegraphics[page=97,trim={53.00bp 429.25bp 53.87bp 52.00bp},clip,width=\textwidth,max height=0.38\textheight]{Ferziger4th.pdf}
\par\vspace{0.45\baselineskip}
\small 图 4.2 典型的 CV 和用于笛卡尔二维网格的符号
\end{sourceblock}
\begin{sourceblock}
\includegraphics[page=97,trim={53.00bp 237.83bp 53.87bp 250.45bp},clip,width=\textwidth,max height=0.38\textheight]{Ferziger4th.pdf}
\par\vspace{0.45\baselineskip}
\small 图 4.3 典型的 CV 和用于笛卡尔 3D 网格的符号
\end{sourceblock}
其中 f 是对流 (ρφ v · n ) 或扩散 ( ∇ φ · n ) 通量矢量在 CV 面法向方向上的分量。由于假设速度场和流体属性已知，唯一的未知数是 φ。如果速度场未知，我们就会遇到涉及非线性耦合方程的更复杂的问题；我们将在章节中处理它。 7和8。

为了保持保护，CV 不重叠很重要；每个 CV 面对其两侧的两个 CV 来说都是唯一的。

在下文中，仅考虑典型的 CV 面，即图 4.2 中标记为“e”的面；通过进行适当的索引替换，可以为所有面导出类似的表达式。

为了准确计算式(4.2)中的曲面积分，需要知道曲面S e 上各处的被积函数f。此信息不可用，因为仅计算 φ 的节点（CV 中心）值，因此必须引入近似值。最好使用两级近似来完成此操作：

\begin{itemize}
\item 根据单元表面上一个或多个位置处的变量值来近似积分；
\item 单元面值根据节点（CV 中心）值进行近似。
\end{itemize}
最简单的积分近似是中点法则：积分近似为单元面中心处的被积函数（其本身是表面平均值的近似值）与单元面面积的乘积：

\begin{sourceblock}
\begin{tikzpicture}
\node[inner sep=0,anchor=south west] (image) at (0,0) {\includegraphics[page=98,trim={53.00bp 441.03bp 53.87bp 201.99bp},clip,width=0.92\textwidth,max height=0.38\textheight]{Ferziger4th.pdf}};
\begin{scope}[x={(image.south east)},y={(image.north west)}]
\fill[white] (0.30692,0.39792) rectangle (0.33789,0.93642);
\fill[white] (0.34313,0.44810) rectangle (0.37131,0.94810);
\end{scope}
\end{tikzpicture}
\end{sourceblock}
如果位置“e”处的 f 值已知，则该积分的近似值具有二阶精度。让我们展示如何在简单的二维问题中确定积分近似的阶数。我们考虑笛卡尔控制体积的东侧，因此沿 y 坐标进行积分。为了简单起见，我们将坐标原点设置在单元面中心“e”；因此 y se = − y / 2 且 y ne = + y / 2。我们首先使用绕面质心的被积函数 f 的泰勒级数展开，假设 f e 已知：

\begin{sourceblock}
\begin{tikzpicture}
\node[inner sep=0,anchor=south west] (image) at (0,0) {\includegraphics[page=98,trim={53.00bp 309.05bp 53.87bp 319.67bp},clip,width=0.92\textwidth,max height=0.38\textheight]{Ferziger4th.pdf}};
\begin{scope}[x={(image.south east)},y={(image.north west)}]
\fill[white] (0.39627,0.46472) rectangle (0.42791,0.78087);
\fill[white] (0.25408,0.27819) rectangle (0.26722,0.58712);
\fill[white] (0.27699,0.28541) rectangle (0.30517,0.59433);
\fill[white] (0.31525,0.25441) rectangle (0.33886,0.58712);
\fill[white] (0.34241,0.28541) rectangle (0.37059,0.59433);
\end{scope}
\end{tikzpicture}
\end{sourceblock}
\begin{sourceblock}
\begin{tikzpicture}
\node[inner sep=0,anchor=south west] (image) at (0,0) {\includegraphics[page=98,trim={53.00bp 250.08bp 53.87bp 379.37bp},clip,width=0.92\textwidth,max height=0.38\textheight]{Ferziger4th.pdf}};
\begin{scope}[x={(image.south east)},y={(image.north west)}]
\fill[white] (0.00000,0.75988) rectangle (0.05071,1.00000);
\fill[white] (0.05341,0.75988) rectangle (0.11531,1.00000);
\fill[white] (0.11798,0.75988) rectangle (0.18105,1.00000);
\end{scope}
\end{tikzpicture}
\end{sourceblock}
将式(4.4)代入式(4.5)可得:

\begin{sourceblock}
\begin{tikzpicture}
\node[inner sep=0,anchor=south west] (image) at (0,0) {\includegraphics[page=98,trim={53.00bp 187.08bp 53.87bp 441.29bp},clip,width=0.92\textwidth,max height=0.38\textheight]{Ferziger4th.pdf}};
\begin{scope}[x={(image.south east)},y={(image.north west)}]
\fill[white] (0.13708,0.58830) rectangle (0.17260,0.82235);
\fill[white] (0.18361,0.28462) rectangle (0.19675,0.59095);
\fill[white] (0.20319,0.28197) rectangle (0.23883,0.59095);
\fill[white] (0.24391,0.29177) rectangle (0.27209,0.59783);
\fill[white] (0.10598,0.07122) rectangle (0.12716,0.29812);
\fill[white] (0.13970,0.06434) rectangle (0.17522,0.29812);
\end{scope}
\end{tikzpicture}
\end{sourceblock}
涉及 f 的一阶导数的项被删除，因此我们得到：

\begin{sourceblock}
\begin{tikzpicture}
\node[inner sep=0,anchor=south west] (image) at (0,0) {\includegraphics[page=98,trim={53.00bp 125.63bp 53.87bp 503.09bp},clip,width=0.92\textwidth,max height=0.38\textheight]{Ferziger4th.pdf}};
\begin{scope}[x={(image.south east)},y={(image.north west)}]
\fill[white] (0.26833,0.58471) rectangle (0.30385,0.82068);
\end{scope}
\end{tikzpicture}
\end{sourceblock}
右侧第一项是中点法则近似；第二项是截断误差的首项，通常是最大的。局部误差与 ( y ) 3 成正比，但我们必须考虑积分需要在包含 n 个段的有限区域上执行，其中 n = Y / y，其中
Y 是总积分距离。因此，总误差与 (y) 2 成正比，因为局部误差发生了 n 次。

由于f的值在单元面中心‘e’处不可用，因此必须通过插值来获得。为了保持表面积分的中点法则近似的二阶精度，f e 的值必须至少以二阶精度计算。我们将在第 4 节中介绍一些广泛使用的近似值。 4.4.

二维表面积分的另一个二阶近似是梯形规则，它导致：

\begin{sourceblock}
\begin{tikzpicture}
\node[inner sep=0,anchor=south west] (image) at (0,0) {\includegraphics[page=99,trim={53.00bp 474.53bp 53.87bp 161.77bp},clip,width=0.92\textwidth,max height=0.38\textheight]{Ferziger4th.pdf}};
\begin{scope}[x={(image.south east)},y={(image.north west)}]
\fill[white] (0.00000,0.87265) rectangle (0.05820,1.00000);
\fill[white] (0.06090,0.87265) rectangle (0.13892,1.00000);
\fill[white] (0.14162,0.87265) rectangle (0.20797,1.00000);
\fill[white] (0.21068,0.87265) rectangle (0.24713,1.00000);
\fill[white] (0.29558,0.30831) rectangle (0.32659,0.72554);
\fill[white] (0.33179,0.34718) rectangle (0.35997,0.73458);
\end{scope}
\end{tikzpicture}
\end{sourceblock}
在这种情况下，我们需要评估 CV 角处的通量矢量分量。使用与上述中点规则近似相同的方法，但在位置“se”周围执行泰勒级数展开并假设 f ne 和 f se 已知，可以表明梯形规则的局部截断误差是中点规则的两倍大，但符号相反：

\begin{sourceblock}
\begin{tikzpicture}
\node[inner sep=0,anchor=south west] (image) at (0,0) {\includegraphics[page=99,trim={53.00bp 366.27bp 53.87bp 262.44bp},clip,width=0.92\textwidth,max height=0.38\textheight]{Ferziger4th.pdf}};
\begin{scope}[x={(image.south east)},y={(image.north west)}]
\fill[white] (0.21943,0.58456) rectangle (0.25498,0.82046);
\end{scope}
\end{tikzpicture}
\end{sourceblock}
但请注意，参考位置现在是面角，因为未使用中点值。同样，从 CV 质心到面角的插值必须至少为二阶精度，以保留梯形规则积分近似的二阶。

对于表面积分的高阶近似，必须在两个以上位置评估通量矢量分量。四阶近似是辛普森法则，它将 S e 上的积分估计为：

\begin{sourceblock}
\begin{tikzpicture}
\node[inner sep=0,anchor=south west] (image) at (0,0) {\includegraphics[page=99,trim={53.00bp 236.21bp 53.87bp 400.10bp},clip,width=0.92\textwidth,max height=0.38\textheight]{Ferziger4th.pdf}};
\begin{scope}[x={(image.south east)},y={(image.north west)}]
\fill[white] (0.25368,0.30808) rectangle (0.28466,0.72544);
\fill[white] (0.28986,0.34697) rectangle (0.31805,0.73450);
\end{scope}
\end{tikzpicture}
\end{sourceblock}
这里三个位置需要 f 的值：单元面中心“e”和两个角“ne”和“se”。为了保持四阶精度，这些值必须通过节点值插值来获得，至少与辛普森规则一样准确。三次多项式是合适的，如下所示。

在 3D 中，中点规则又是最简单的二阶近似。高阶近似，要求被积函数位于单元面中心以外的位置（例如，角和边缘的中心）是可能的，但它们更难以实现。下一节提到了一种可能性。

如果假设 f 的变化具有某种特定的简单形状（例如插值多项式），则积分很容易。近似的精度取决于形状函数的阶数。

\section*{4.3 体积积分的近似}\phantomsection\addcontentsline{toc}{section}{4.3 体积积分的近似}
传输方程中的某些项需要对 CV 体积进行积分。最简单的二阶精确逼近是将体积积分替换为被积函数平均值与 CV 体积的乘积，并将前者近似为 CV 中心处的值：

\begin{sourceblock}
\begin{tikzpicture}
\node[inner sep=0,anchor=south west] (image) at (0,0) {\includegraphics[page=100,trim={53.00bp 496.60bp 53.87bp 147.06bp},clip,width=0.92\textwidth,max height=0.38\textheight]{Ferziger4th.pdf}};
\begin{scope}[x={(image.south east)},y={(image.north west)}]
\fill[white] (0.28337,0.38078) rectangle (0.32313,0.93461);
\fill[white] (0.32848,0.43194) rectangle (0.35666,0.94662);
\end{scope}
\end{tikzpicture}
\end{sourceblock}
其中 q P 代表 CV 中心处的 q 值。这个数量很容易计算；因为所有变量在节点 P 处都可用，因此不需要插值。如果 q 是常数或在 CV 内线性变化，则上述近似变得精确；否则，它包含二阶误差，这一点很容易看出。

高阶近似需要 q 值位于更多位置而不仅仅是中心。这些值必须通过插值节点值或等效地使用形状函数来获得。

在二维中，体积积分变成面积积分。使用二二次形函数可以得到四阶近似：

q ( x , y ) = a 0 + a 1 x + a 2 y + a 3 x 2 + a 4 y 2 + a 5 xy + a 6 x 2 y + a 7 xy 2 + a 8 x 2 y 2。

(4.12) 九个系数是通过将函数拟合到九个位置（‘nw’、‘w’、‘sw’、‘n’、P、‘s’、‘ne’、‘e’和‘se’处的 q 值来获得的，见图 4.2。然后可以评估积分。在二维中，积分给出（对于笛卡尔网格）：

\begin{sourceblock}
\begin{tikzpicture}
\node[inner sep=0,anchor=south west] (image) at (0,0) {\includegraphics[page=100,trim={53.00bp 271.98bp 53.87bp 360.69bp},clip,width=0.92\textwidth,max height=0.38\textheight]{Ferziger4th.pdf}};
\begin{scope}[x={(image.south east)},y={(image.north west)}]
\fill[white] (0.02911,0.25575) rectangle (0.06884,0.62773);
\fill[white] (0.07417,0.29011) rectangle (0.10235,0.63579);
\fill[white] (0.92439,0.00000) rectangle (1.00000,0.12817);
\end{scope}
\end{tikzpicture}
\end{sourceblock}
(4.13) 只需要确定四个系数，但它们取决于上面列出的所有九个位置处的 q 值。在统一的笛卡尔网格上我们得到：

\begin{sourceblock}
\begin{tikzpicture}
\node[inner sep=0,anchor=south west] (image) at (0,0) {\includegraphics[page=100,trim={53.00bp 203.84bp 53.87bp 433.91bp},clip,width=0.92\textwidth,max height=0.38\textheight]{Ferziger4th.pdf}};
\begin{scope}[x={(image.south east)},y={(image.north west)}]
\fill[white] (0.92433,0.00000) rectangle (1.00000,0.18246);
\end{scope}
\end{tikzpicture}
\end{sourceblock}
(4.14) 由于只有 P 处的值可用，因此必须使用插值来获得其他位置处的 q。它必须至少达到四阶精度才能保持积分近似的精度。下一节将描述一些可能性。

上述 2D 体积积分的四阶近似可用于近似 3D 表面积分。 3D 体积积分的高阶近似更复杂，但可以使用相同的技术找到。

\section*{4.4 插值和微分实践}\phantomsection\addcontentsline{toc}{section}{4.4 插值和微分实践}
积分的近似值需要计算节点（CV 中心）以外的位置处的变量值。被积函数（在前面的部分中用 f 表示）涉及这些位置处的多个变量和/或变量梯度的乘积：对于对流通量，f c = ρφ v · n；对于扩散通量，f d = ∇ φ · n。我们假设速度场和流体特性 ρ 和 ρ 在所有位置都是已知的。为了计算对流和扩散通量，需要 CV 表面上一个或多个位置处的 φ 值及其垂直于单元面的梯度。源项的体积积分也可能需要这些值。它们必须通过插值以节点值表示。

有多种可能性可用于计算 φ 及其在单元表面的梯度。事实上，Waterson 和 Deconinck（2007）已经审查了有限体积环境中数十种方案的性能。商业代码通常实施许多不同的方案，并就其使用和准确性提供建议。在这里，我们的策略是详细探索一些常用的方案，然后总结并将我们的观点扩展到更通用的方案，包括 κ 方法、通量限制器和总变差递减 (TVD) 方案。 1

\subsection*{4.4.1 迎风插值（UDS）}\phantomsection\addcontentsline{toc}{subsection}{4.4.1 迎风插值（UDS）}
通过“e”上游节点处的值来近似 φ e 相当于对一阶导数使用后向或前向差分近似（取决于流向），因此该近似被称为迎风差分方案 (UDS)。在 UDS 中，φ e 近似为：

\begin{sourceblock}
\includegraphics[page=101,trim={53.00bp 243.26bp 53.87bp 387.18bp},clip,width=0.92\textwidth,max height=0.38\textheight]{Ferziger4th.pdf}
\end{sourceblock}
这是唯一无条件满足有界准则的近似，即它永远不会产生振荡解。然而，它通过数值扩散来实现这一点。这在前一章中已经介绍过，下面将再次介绍。

关于 P 的泰勒级数展开给出（对于笛卡尔网格且 ( v · n ) e > 0）：

\begin{sourceblock}
\begin{tikzpicture}
\node[inner sep=0,anchor=south west] (image) at (0,0) {\includegraphics[page=101,trim={53.00bp 134.37bp 53.87bp 494.35bp},clip,width=0.92\textwidth,max height=0.38\textheight]{Ferziger4th.pdf}};
\begin{scope}[x={(image.south east)},y={(image.north west)}]
\fill[white] (0.38244,0.47194) rectangle (0.42069,0.78087);
\fill[white] (0.10433,0.25441) rectangle (0.13576,0.59433);
\fill[white] (0.14096,0.28541) rectangle (0.16914,0.59433);
\fill[white] (0.17263,0.25441) rectangle (0.20641,0.59433);
\fill[white] (0.21008,0.28541) rectangle (0.23826,0.59433);
\fill[white] (0.24012,0.25441) rectangle (0.28039,0.59433);
\fill[white] (0.28394,0.25441) rectangle (0.35847,0.59433);
\end{scope}
\end{tikzpicture}
\end{sourceblock}
\footnotetext[1]{在本节中，我们将研究定义 φ 和/或其导数的方案。见第 11.3 节，我们将研究所谓的通量校正传输（FCT）。}
其中 H 表示高阶项。 UDS 近似仅保留右侧第一项，因此它是一阶方案。它的主要截断误差项是扩散的，即它类似于扩散通量：

\begin{sourceblock}
\begin{tikzpicture}
\node[inner sep=0,anchor=south west] (image) at (0,0) {\includegraphics[page=102,trim={53.00bp 533.20bp 53.87bp 95.52bp},clip,width=0.92\textwidth,max height=0.38\textheight]{Ferziger4th.pdf}};
\begin{scope}[x={(image.south east)},y={(image.north west)}]
\fill[white] (0.40722,0.40406) rectangle (0.42250,0.63335);
\fill[white] (0.39263,0.27793) rectangle (0.40577,0.58712);
\end{scope}
\end{tikzpicture}
\end{sourceblock}
数值、人工或虚假扩散的系数（它有各种不恭维的名称！）是 num
e = (ρ u ) e x / 2。如果流动与网格倾斜，则在多维问题中这种数值扩散会被放大；然后，截断误差会在垂直于流动的方向以及流向方向上产生扩散，这是一种特别严重的误差类型。变量的峰值或快速变化将被抹掉，并且由于误差减少率只是一阶的，因此需要非常精细的网格才能获得准确的解。

\subsection*{4.4.2 线性插值（CDS）}\phantomsection\addcontentsline{toc}{subsection}{4.4.2 线性插值（CDS）}
CV 面中心值的另一个直接近似是两个最近节点之间的线性插值。在笛卡尔网格上的位置“e”处，我们有（见图 4.2 和 4.3）：

\begin{sourceblock}
\includegraphics[page=102,trim={53.00bp 326.51bp 53.87bp 324.17bp},clip,width=0.92\textwidth,max height=0.38\textheight]{Ferziger4th.pdf}
\end{sourceblock}
其中线性插值因子 λ e 定义为：

\begin{sourceblock}
\includegraphics[page=102,trim={53.00bp 267.63bp 53.87bp 368.99bp},clip,width=0.92\textwidth,max height=0.38\textheight]{Ferziger4th.pdf}
\end{sourceblock}
方程（4.18）是二阶精确的，可以通过使用 φ E 关于点 x P 的泰勒级数展开来消除方程（4.16）中的一阶导数。结果是：

\begin{sourceblock}
\begin{tikzpicture}
\node[inner sep=0,anchor=south west] (image) at (0,0) {\includegraphics[page=102,trim={53.00bp 186.53bp 53.87bp 442.18bp},clip,width=0.92\textwidth,max height=0.38\textheight]{Ferziger4th.pdf}};
\begin{scope}[x={(image.south east)},y={(image.north west)}]
\fill[white] (0.46138,0.44109) rectangle (0.57483,0.78092);
\fill[white] (0.57850,0.44109) rectangle (0.65053,0.78092);
\fill[white] (0.08066,0.25461) rectangle (0.11209,0.59418);
\fill[white] (0.11729,0.28533) rectangle (0.14547,0.59418);
\fill[white] (0.14899,0.25461) rectangle (0.21146,0.59418);
\fill[white] (0.21501,0.28533) rectangle (0.24319,0.59418);
\fill[white] (0.24502,0.25461) rectangle (0.30677,0.59418);
\fill[white] (0.30863,0.25461) rectangle (0.38205,0.59418);
\fill[white] (0.38391,0.28533) rectangle (0.45780,0.78092);
\fill[white] (0.74893,0.03206) rectangle (0.76541,0.26102);
\end{scope}
\end{tikzpicture}
\end{sourceblock}
在均匀或非均匀网格上，前导截断误差项与网格间距的平方成正比。

与所有高于一阶的近似一样，该方案可能会产生振荡解。这是最简单的二阶方案，也是应用最广泛的一种。它对应于FD方法中一阶导数的中心差分近似；因此缩写为 CDS。

P 和 E 节点之间的线性分布假设也提供了评估扩散通量所需的最简单的梯度近似：

\begin{sourceblock}
\begin{tikzpicture}
\node[inner sep=0,anchor=south west] (image) at (0,0) {\includegraphics[page=103,trim={53.00bp 581.42bp 53.87bp 50.00bp},clip,width=0.92\textwidth,max height=0.38\textheight]{Ferziger4th.pdf}};
\begin{scope}[x={(image.south east)},y={(image.north west)}]
\fill[white] (0.39600,0.50893) rectangle (0.43423,0.84188);
\end{scope}
\end{tikzpicture}
\end{sourceblock}
通过使用围绕 φ e 的泰勒级数展开，可以看出上述近似的截断误差为：

\begin{sourceblock}
\begin{tikzpicture}
\node[inner sep=0,anchor=south west] (image) at (0,0) {\includegraphics[page=103,trim={53.00bp 484.75bp 53.87bp 118.94bp},clip,width=0.92\textwidth,max height=0.38\textheight]{Ferziger4th.pdf}};
\begin{scope}[x={(image.south east)},y={(image.north west)}]
\fill[white] (0.31711,0.66501) rectangle (0.35738,0.86869);
\fill[white] (0.36096,0.66501) rectangle (0.44592,0.88423);
\fill[white] (0.44938,0.66501) rectangle (0.52319,0.86869);
\fill[white] (0.52689,0.66501) rectangle (0.60938,0.88423);
\fill[white] (0.24734,0.55741) rectangle (0.27432,0.75693);
\fill[white] (0.28183,0.57166) rectangle (0.31002,0.75693);
\fill[white] (0.26538,0.26421) rectangle (0.30565,0.46789);
\fill[white] (0.30920,0.26421) rectangle (0.39420,0.48343);
\fill[white] (0.39753,0.28279) rectangle (0.42571,0.46789);
\fill[white] (0.42758,0.26421) rectangle (0.47134,0.46789);
\fill[white] (0.47501,0.26421) rectangle (0.55753,0.48343);
\end{scope}
\end{tikzpicture}
\end{sourceblock}
当位置“e”位于 P 和 E 之间（例如在均匀网格上）时，近似值为二阶精度，因为右侧第一项消失，并且主误差项与 ( x ) 2 成比例。当网格不均匀时，前导误差项与 x 和网格扩展因子减去 1 的乘积成正比。尽管具有形式上的一阶精度，但即使在非均匀网格上，细化网格时的误差减少也与二阶近似相似。见第 3.3.4 节 对此行为的详细解释。

\subsection*{4.4.3 二次迎风插值（QUICK）}\phantomsection\addcontentsline{toc}{subsection}{4.4.3 二次迎风插值（QUICK）}
下一个逻辑改进是用抛物线而不是直线来近似 P 和 E 之间的可变轮廓。为了构造抛物线，我们需要再使用一点数据；根据对流的性质，第三点取在上游侧，即，如果流动是从P到E（即u x > 0），则取W；如果u x < 0，则取EE，见图4.2。由此我们得到：

\begin{sourceblock}
\includegraphics[page=103,trim={53.00bp 248.17bp 53.87bp 402.51bp},clip,width=0.92\textwidth,max height=0.38\textheight]{Ferziger4th.pdf}
\end{sourceblock}
其中D、U和UU分别表示下游、第一上游和第二上游节点（E、P和W或P、E和EE，取决于流向）。系数 g 1 和 g 2 可以用节点坐标表示：

\begin{sourceblock}
\includegraphics[page=103,trim={53.00bp 152.39bp 53.87bp 484.24bp},clip,width=0.92\textwidth,max height=0.38\textheight]{Ferziger4th.pdf}
\end{sourceblock}
对于均匀网格，插值涉及的三个节点值的系数为： 3
8 表示下游点，6 8 表示第一个上游节点，− 1
8为第二上游节点。该方案比 CDS 方案稍微复杂一些：它在每个方向上将计算分子多扩展一个节点（在 2D 中，包括节点 EE、WW、NN 和 SS），并且在非正交和/或非均匀网格上，系数 g i 的表达式并不简单。伦纳德
(1979) 使这个方案变得流行，并将其命名为 QUICK（对流运动学的二次迎风插值）。

这种二次插值方案在均匀和非均匀网格上都具有三阶截断误差。这可以通过使用 φ W 消除方程 ( 4.20 ) 的二阶导数来证明，在 u x > 0 的均匀笛卡尔网格上，导致：

\begin{sourceblock}
\begin{tikzpicture}
\node[inner sep=0,anchor=south west] (image) at (0,0) {\includegraphics[page=104,trim={53.00bp 507.76bp 53.87bp 120.95bp},clip,width=0.92\textwidth,max height=0.38\textheight]{Ferziger4th.pdf}};
\begin{scope}[x={(image.south east)},y={(image.north west)}]
\fill[white] (0.24379,0.46166) rectangle (0.26358,0.77077);
\fill[white] (0.17185,0.25407) rectangle (0.20328,0.59391);
\fill[white] (0.20848,0.28507) rectangle (0.23666,0.59391);
\fill[white] (0.73032,0.03206) rectangle (0.74680,0.26102);
\end{scope}
\end{tikzpicture}
\end{sourceblock}
右侧的前三项代表 QUICK 近似，而最后一项是主要截断误差。然而，当该插值方案与曲面积分的中点法则近似结合使用时，整体近似仍然是二阶精度（求积近似的精度）。尽管 QUICK 近似比 CDS 稍微准确，但这两种方案都以二阶方式渐近收敛，并且差异很少很大。

\subsection*{4.4.4 高阶方案}\phantomsection\addcontentsline{toc}{subsection}{4.4.4 高阶方案}
仅当使用高阶公式近似积分时，高于三阶的插值才有意义。如果在二维表面积分中使用辛普森法则，则必须使用至少三阶多项式进行插值，以保留求积近似的四阶精度，这会导致四阶插值误差。例如，通过拟合多项式

\begin{sourceblock}
\includegraphics[page=104,trim={53.00bp 278.88bp 53.87bp 367.41bp},clip,width=0.92\textwidth,max height=0.38\textheight]{Ferziger4th.pdf}
\end{sourceblock}
通过四个节点（“e”两侧各两个：W、P、E 和 EE）处的 φ 值，可以确定四个系数 a i 并找到 φ e 作为节点值的函数。对于均匀笛卡尔网格，可得到以下表达式：

\begin{sourceblock}
\includegraphics[page=104,trim={53.00bp 196.69bp 53.87bp 441.06bp},clip,width=0.92\textwidth,max height=0.38\textheight]{Ferziger4th.pdf}
\end{sourceblock}
可以使用相同的多项式来确定导数；我们只需对其微分一次即可得到：

\begin{sourceblock}
\begin{tikzpicture}
\node[inner sep=0,anchor=south west] (image) at (0,0) {\includegraphics[page=104,trim={53.00bp 128.69bp 53.87bp 500.02bp},clip,width=0.92\textwidth,max height=0.38\textheight]{Ferziger4th.pdf}};
\begin{scope}[x={(image.south east)},y={(image.north west)}]
\fill[white] (0.32340,0.47181) rectangle (0.36165,0.78066);
\end{scope}
\end{tikzpicture}
\end{sourceblock}
在统一的笛卡尔网格上，产生：

\begin{sourceblock}
\begin{tikzpicture}
\node[inner sep=0,anchor=south west] (image) at (0,0) {\includegraphics[page=104,trim={53.00bp 69.77bp 53.87bp 558.95bp},clip,width=0.92\textwidth,max height=0.38\textheight]{Ferziger4th.pdf}};
\begin{scope}[x={(image.south east)},y={(image.north west)}]
\fill[white] (0.28250,0.47194) rectangle (0.32072,0.78087);
\end{scope}
\end{tikzpicture}
\end{sourceblock}
上述近似有时称为四阶 CDS。当然，可以使用更高次数的多项式和/或多维多项式。也可以使用三次样条，它确保插值函数及其前两个导数在整个解域中的连续性（但成本会有所增加）。

一旦在单元面中心获得变量值及其导数，就可以在单元面上进行插值以获得 CV 角处的值。这对于显式方法来说并不困难，但基于辛普森规则和多项式插值的四阶方案会产生太大的计算分子，无法进行隐式处理。人们可以通过使用第 1 节中描述的延迟校正方法来避免这种复杂性。 5.6.

另一种方法是使用 FD 方法中导出紧凑 (Padé) 方案的技术。例如，可以通过将多项式 (4.25) 拟合到单元面两侧两个节点处的变量值和一阶导数来获得多项式的系数。对于均匀笛卡尔网格，可得出以下 φ e 表达式：

\begin{sourceblock}
\begin{tikzpicture}
\node[inner sep=0,anchor=south west] (image) at (0,0) {\includegraphics[page=105,trim={53.00bp 401.18bp 53.87bp 227.54bp},clip,width=0.92\textwidth,max height=0.38\textheight]{Ferziger4th.pdf}};
\begin{scope}[x={(image.south east)},y={(image.north west)}]
\fill[white] (0.00000,0.73918) rectangle (0.03907,1.00000);
\fill[white] (0.04174,0.71833) rectangle (0.07317,1.00000);
\fill[white] (0.07756,0.73918) rectangle (0.16896,1.00000);
\fill[white] (0.29555,0.44094) rectangle (0.32932,0.78087);
\fill[white] (0.33302,0.47167) rectangle (0.36120,0.78087);
\fill[white] (0.36304,0.44094) rectangle (0.39795,0.78087);
\fill[white] (0.46376,0.46472) rectangle (0.48186,0.77365);
\fill[white] (0.22364,0.25441) rectangle (0.25507,0.59407);
\fill[white] (0.26030,0.28514) rectangle (0.28848,0.59407);
\fill[white] (0.40523,0.28514) rectangle (0.43341,0.59407);
\fill[white] (0.33771,0.08632) rectangle (0.35750,0.39524);
\fill[white] (0.45179,0.08632) rectangle (0.47158,0.39524);
\fill[white] (0.71991,0.03207) rectangle (0.73756,0.26109);
\end{scope}
\end{tikzpicture}
\end{sourceblock}
上式右边第一项表示线性插值的二阶近似；第二项表示线性插值的前导截断误差项的近似值，见方程(4.20)，其中二阶导数由CDS近似。

问题是节点 P 和 E 处的导数是未知的并且必须对其本身进行近似。然而，即使我们用二阶 CDS 来近似一阶导数，即：

\begin{sourceblock}
\begin{tikzpicture}
\node[inner sep=0,anchor=south west] (image) at (0,0) {\includegraphics[page=105,trim={53.00bp 270.53bp 53.87bp 358.18bp},clip,width=0.92\textwidth,max height=0.38\textheight]{Ferziger4th.pdf}};
\begin{scope}[x={(image.south east)},y={(image.north west)}]
\fill[white] (0.23537,0.47155) rectangle (0.27359,0.78039);
\fill[white] (0.23636,0.08897) rectangle (0.27143,0.40502);
\end{scope}
\end{tikzpicture}
\end{sourceblock}
所得到的单元面值的近似值保留了多项式的四阶精度：

\begin{sourceblock}
\includegraphics[page=105,trim={53.00bp 202.57bp 53.87bp 435.18bp},clip,width=0.92\textwidth,max height=0.38\textheight]{Ferziger4th.pdf}
\end{sourceblock}
该表达式与式（4.26）相同。

如果我们使用单元面两侧的变量值和上游侧的导数作为数据，我们可以拟合抛物线。这导致了与上述 QUICK 方案等效的近似：

\begin{sourceblock}
\begin{tikzpicture}
\node[inner sep=0,anchor=south west] (image) at (0,0) {\includegraphics[page=105,trim={53.00bp 108.89bp 53.87bp 519.83bp},clip,width=0.92\textwidth,max height=0.38\textheight]{Ferziger4th.pdf}};
\begin{scope}[x={(image.south east)},y={(image.north west)}]
\fill[white] (0.35582,0.46179) rectangle (0.37561,0.77098);
\fill[white] (0.28388,0.25414) rectangle (0.31528,0.59407);
\fill[white] (0.32048,0.28514) rectangle (0.34866,0.59407);
\end{scope}
\end{tikzpicture}
\end{sourceblock}
可以使用相同的方法来获得单元面中心处的导数的近似值；根据多项式 ( 4.25 ) 的导数，我们得到：

\begin{sourceblock}
\begin{tikzpicture}
\node[inner sep=0,anchor=south west] (image) at (0,0) {\includegraphics[page=106,trim={53.00bp 581.42bp 53.87bp 50.00bp},clip,width=0.92\textwidth,max height=0.38\textheight]{Ferziger4th.pdf}};
\begin{scope}[x={(image.south east)},y={(image.north west)}]
\fill[white] (0.12232,0.50893) rectangle (0.16057,0.84188);
\fill[white] (0.12334,0.09620) rectangle (0.15838,0.43692);
\fill[white] (0.18105,0.03456) rectangle (0.19519,0.28168);
\end{scope}
\end{tikzpicture}
\end{sourceblock}
显然，右侧第一项是二阶 CDS 近似。其余项表示提高准确性的修正。

近似值 ( 4.29 )、( 4.31 ) 和 ( 4.32 ) 的问题在于它们包含 CV 中心处的一阶导数，而这些一阶导数是未知的。尽管我们可以用以节点变量值表示的二阶近似值来替换它们，而不破坏它们的精度顺序，但所得的计算分子将比我们希望的大得多。例如，在二维中，使用辛普森法则和四阶多项式插值，我们发现每个通量取决于 15 个节点值，并且一个 CV 的代数方程涉及 25 个值。由此产生的方程组的解将非常昂贵（参见第 5 章）。

解决这个问题的方法在于延迟校正方法，这将在第 1 节中描述。 5.6.

人们应该记住，高阶近似并不一定能保证在任何单个网格上都能得到更准确的解；只有当网格足够细以捕获解的所有基本细节时，才能实现高精度；这种情况发生的网格大小只能通过系统的网格细化来确定。

\subsection*{4.4.5 其他方案}\phantomsection\addcontentsline{toc}{subsection}{4.4.5 其他方案}
人们已经提出了大量对流通量的近似值；讨论所有这些问题超出了本书的范围。上面使用的方法几乎可以用来推导所有这些。我们在这里和下一节中描述其中的一些。

人们可以通过从两个上游节点线性外推来近似 φ e，从而产生所谓的线性迎风方案（LUDS）。该方案具有二阶精度，但由于它比 CDS 更复杂并且可以产生无界解，因此后者是更好的选择。

Raithby (1976) 提出的另一种方法是从迎风侧进行推断，但沿着流线而不是网格线（倾斜迎风方案）。提出了对应于迎风和线性迎风方案的一阶和二阶方案。它们比基于沿网格线外推的方案具有更好的准确性。然而，这些方案非常复杂（有许多可能的流动方向）并且需要大量插值。由于当网格不够精细且难以编程时，这些方案可能会产生振荡解，因此尚未得到广泛使用。

也可以混合两个或多个不同的近似值（参见例如第 4.4.6 节）。在 20 世纪 70 年代和 1980 年代大量使用（并且仍在一些商业代码中使用）的一个例子是 Spalding (1972) 的混合方案，它根据 Peclet 数的本地值在 UDS 和 CDS 之间切换（例如，当本地 Pe > 2 时切换到 UDS）。其他研究人员提出了低阶和高阶方案的混合，以避免非物理振荡，特别是对于带有冲击的可压缩流。其中一些想法将在下一节中介绍，并将在第 1 章中提到。 11 也是如此。混合可用于提高某些迭代求解器的收敛速度，如下所示。

\subsection*{4.4.6 总体策略、TVD 方案和通量限制器}\phantomsection\addcontentsline{toc}{subsection}{4.4.6 总体策略、TVD 方案和通量限制器}
上述方案的准确性和生成良好（即有界）结果的能力范围广泛。特别是，如上所述，只有 UDS 方案才能保证产生没有振荡的解。 Waterson 和 Deconinck (2007) 对有界高阶对流方案进行了全面回顾。在这里，我们简要概述了用于分类线性模型的 κ 方案、总变差递减 (TVD) 方案以及通量限制器方案的概念。尽管我们直到第 1 章才明确处理不稳态问题。如图6所示，这里考虑时间对空间离散化方案的影响是很方便的，因为在第6章中。如图6所示，重点在于时间步进方案本身。在下文中，假设网格和面间距均匀；导出非均匀网格的表达式很简单。

上述许多方案可以被纳入一个通用框架中，称为 κ 方案（Van Leer 1985；Waterson 和 Deconinck 2007）。我们遵循后者列出的模式，部分是为了方便读者使用该评论作为选择适当方案的来源。在图的上下文中。 4.2 和 4.3 中，对于单元面变量 φ e，在速度为正方向的均匀网格上，可以写出 κ 格式：

\begin{sourceblock}
\begin{tikzpicture}
\node[inner sep=0,anchor=south west] (image) at (0,0) {\includegraphics[page=107,trim={53.00bp 259.51bp 53.87bp 371.24bp},clip,width=0.92\textwidth,max height=0.38\textheight]{Ferziger4th.pdf}};
\begin{scope}[x={(image.south east)},y={(image.north west)}]
\fill[white] (0.17708,0.21136) rectangle (0.20851,0.57078);
\fill[white] (0.21371,0.24414) rectangle (0.24189,0.57078);
\fill[white] (0.24541,0.21136) rectangle (0.27916,0.57078);
\fill[white] (0.28283,0.24414) rectangle (0.31101,0.57078);
\end{scope}
\end{tikzpicture}
\end{sourceblock}
其中“主导”方案是迎风差分方案（4.15），添加项是“反扩散”项，用于抵消扩散和迎风偏向 UDS。 κ 分配每个活动方案的比例，如表 4.1 所示。

\begin{center}\small 表 4.1 FVM 中线性对流方程的 κ 方案示例\end{center}
\begin{sourceblock}
\includegraphics[page=107,trim={53.00bp 164.70bp 53.87bp 487.94bp},clip,width=0.92\textwidth,max height=0.38\textheight]{Ferziger4th.pdf}
\end{sourceblock}
CDS 1 1 2 (φ P + φ E ) 二阶精确；参见方程( 4.18 ) 对于 λ e = 1 2 QUICK 1 2 6 8 φ P + 3 8 φ E − 1 8 φ W 二阶精确；参见等式(4.24)

LUI − 1 3 2 φ P − 1 2 φ W 二阶精确；完全迎风
CUI 1 3 5 6 Φ P + 2 6 Φ E − 1 6 Φ W 二阶精确； Waterson-Deconinck 2007 测试中的最佳 κ 方案
在计算机代码中，只需编写通用方程（4.33）即可通过指定 κ 来利用它们中的任何一个。 Waterson 和 Deconinck (2007) 总结了线性标量对流方程的测试，并表明稳定一维对流的修正微分方程可以用 κ 格式写成：

\begin{sourceblock}
\begin{tikzpicture}
\node[inner sep=0,anchor=south west] (image) at (0,0) {\includegraphics[page=108,trim={53.00bp 484.76bp 53.87bp 118.93bp},clip,width=0.92\textwidth,max height=0.38\textheight]{Ferziger4th.pdf}};
\begin{scope}[x={(image.south east)},y={(image.north west)}]
\fill[white] (0.26659,0.68359) rectangle (0.30484,0.86869);
\fill[white] (0.24644,0.56749) rectangle (0.26623,0.75260);
\fill[white] (0.26761,0.45412) rectangle (0.30265,0.64355);
\fill[white] (0.32313,0.24083) rectangle (0.33841,0.37806);
\fill[white] (0.54893,0.28279) rectangle (0.59919,0.48343);
\fill[white] (0.32313,0.11978) rectangle (0.36791,0.35612);
\fill[white] (0.37215,0.16509) rectangle (0.46090,0.35612);
\fill[white] (0.48099,0.16669) rectangle (0.52343,0.37166);
\end{scope}
\end{tikzpicture}
\end{sourceblock}
对于以点 P 为中心的 CV。请注意，上式左侧的对流项是使用动量方程的非保守微分形式编写的，请参见方程 (1.20)，以与 Waterson 和 Deconinck (2007) 使用的符号保持一致。它代表“e”面和“w”面的对流通量之间的差异。

修正方程是通过使用适当的插值公式来表达微分方程的FV形式，然后将得到的表达式展开为关于点P的泰勒级数而获得的（Warming和Hyett 1974；Fletcher 1991（见第 9.2 节）；Ferziger 1998）。结果是原始方程加上隐含在 FV 计算公式中的其他项。这种修正方程方法很有用，因为（1）它也可以应用于非线性方案，（2）从结果中我们可以理解采用特定方案的后果。例如，在这种情况下，可能同时存在三阶和四阶导数项，具体取决于 κ 的值。三阶项是色散的，这意味着在不稳定的情况下，溶液的各个成分以不同的速度移动，从而错误地色散了初始波形。 2 四阶项具有扩散性，会导致解的耗散或退化。因此，对于 CDS 方案，其中 κ = 1，该方法是色散的，但不是扩散的。另一方面，对于立方迎风方案 (CUI)，κ = 1
3、该方法比较准确，有扩散性，但不分散。

虽然一般来说，上述方法通常是足够的，但也存在一些相关的具体问题，例如，与标量对流相关的特定问题，其中产生负密度或盐度的振荡是非物理的，或者与快速变化区域附近的速度有关，即在极值附近，例如冲击。因此，创建其结果以某种方式受到限制的技术是有用的。在上述技术中，只有 UDS 能够保证产生有界和/或单调行为。全变分递减 (TVD) 方案和通量限制器的作用是提供有界解。

在没有任何材料源或汇的情况下，考虑一维标量 φ 的对流作为时间的函数（例如，通量的时间变化率）

\begin{sourceblock}
\begin{tikzpicture}
\node[inner sep=0,anchor=south west] (image) at (0,0) {\includegraphics[page=108,trim={53.00bp 100.44bp 53.87bp 546.49bp},clip,width=0.92\textwidth,max height=0.38\textheight]{Ferziger4th.pdf}};
\begin{scope}[x={(image.south east)},y={(image.north west)}]
\fill[white] (0.00000,0.20562) rectangle (0.02908,0.80791);
\fill[white] (0.03564,0.20562) rectangle (0.07705,0.80791);
\fill[white] (0.08355,0.20562) rectangle (0.13498,0.80791);
\fill[white] (0.17874,0.06247) rectangle (0.20472,0.51848);
\fill[white] (0.24722,0.20562) rectangle (0.34520,0.80791);
\fill[white] (0.35176,0.20562) rectangle (0.39317,0.80791);
\fill[white] (0.39970,0.20562) rectangle (0.46358,0.80791);
\fill[white] (0.47011,0.20562) rectangle (0.54689,0.80791);
\fill[white] (0.55347,0.20562) rectangle (0.60322,0.80791);
\fill[white] (0.60974,0.20562) rectangle (0.65116,0.80791);
\fill[white] (0.65771,0.20562) rectangle (0.82553,0.80791);
\fill[white] (0.83215,0.22488) rectangle (0.85426,0.82665);
\fill[white] (0.86229,0.20562) rectangle (0.88707,0.80791);
\fill[white] (0.89359,0.20562) rectangle (1.00000,0.80791);
\fill[white] (0.22238,0.00000) rectangle (0.23883,0.33368);
\end{scope}
\end{tikzpicture}
\end{sourceblock}
最初，以后的值也必须受到限制，并且不应发生可能导致非物理条件的波动。 Harten (1983)（参见 Hirsch 2007 或 Durran 2010）是第一个量化方案中实现所需有界性的方法的人：标量的总变化必须不随时间增加（有关该主题的文献，因为时间很丰富，仔细阅读上述参考文献即可证实）。因此，如果数量 φ 的总变化在时间 n 定义为：

\begin{sourceblock}
\begin{tikzpicture}
\node[inner sep=0,anchor=south west] (image) at (0,0) {\includegraphics[page=109,trim={53.00bp 498.52bp 53.87bp 140.60bp},clip,width=0.92\textwidth,max height=0.38\textheight]{Ferziger4th.pdf}};
\begin{scope}[x={(image.south east)},y={(image.north west)}]
\fill[white] (0.32520,0.46077) rectangle (0.34668,0.88860);
\fill[white] (0.34725,0.46077) rectangle (0.37038,0.88860);
\fill[white] (0.37011,0.47076) rectangle (0.42899,0.95559);
\fill[white] (0.43248,0.47076) rectangle (0.46066,0.89859);
\end{scope}
\end{tikzpicture}
\end{sourceblock}
其中 k 是网格点索引，标准为

\begin{sourceblock}
\includegraphics[page=109,trim={53.00bp 453.00bp 53.87bp 197.09bp},clip,width=0.92\textwidth,max height=0.38\textheight]{Ferziger4th.pdf}
\end{sourceblock}
即，时间 n + 1 处的总变化应小于或等于时间 n 处的总变化。其效果是将进入控制体积的守恒量的通量限制到不会在该控制体积中产生该量的分布的局部最大值或最小值的水平。尽管有定义，这种约束在文献中已被称为总变异递减或 TVD（Durran 2010）。那么问题就变成了如何满足式（4.36）的检验？前面已经指出，并由 Godunov 在 1959 年证明（参见 Roe 1986 或 Hirsch 2007），在线性方案中，只有一阶方案才能保证满足测试。这激发了对非线性方案的探索。其中最成功的是通量限制器；它们是“基于解场中局部梯度比率定义对流方案的简单函数”（Waterson 和 Deconinck 2007）。 Roe (1986) 指出，一种策略是从满足测试的方案（即 UDS）开始，然后添加非线性项以提高精度。我们在这里遵循[与 κ 方案（方程（4.33））中的方法并行，其中“主导”项也是 UDS]，写道：

\begin{sourceblock}
\begin{tikzpicture}
\node[inner sep=0,anchor=south west] (image) at (0,0) {\includegraphics[page=109,trim={53.00bp 225.94bp 53.87bp 412.17bp},clip,width=0.92\textwidth,max height=0.38\textheight]{Ferziger4th.pdf}};
\begin{scope}[x={(image.south east)},y={(image.north west)}]
\fill[white] (0.00000,0.89262) rectangle (0.05242,1.00000);
\fill[white] (0.05510,0.89262) rectangle (0.13729,1.00000);
\fill[white] (0.14000,0.89262) rectangle (0.23970,1.00000);
\end{scope}
\end{tikzpicture}
\end{sourceblock}
where

\begin{sourceblock}
\begin{tikzpicture}
\node[inner sep=0,anchor=south west] (image) at (0,0) {\includegraphics[page=109,trim={53.00bp 178.10bp 53.87bp 453.67bp},clip,width=0.92\textwidth,max height=0.38\textheight]{Ferziger4th.pdf}};
\begin{scope}[x={(image.south east)},y={(image.north west)}]
\fill[white] (0.00000,0.85831) rectangle (0.07732,1.00000);
\end{scope}
\end{tikzpicture}
\end{sourceblock}
请注意，通过将函数 r 写成式（4.38）所示，我们可以看到它实际上是中心导数与上游导数的比率。可以选择权重函数 ( r )，以便总方案是 TVD，特别是在变量快速变化附近无振荡。这种方案的一个特点是，虽然当变量变化平滑时它是二阶精确的，但它在局部极值处降低到一阶。有大量文献报道了实现这种行为的各种方法，例如 Waterson 和 Deconinck (2007)、Hirsch (2007)、Sweby (1984, 1985)、Yang 和 Przekwas (1992) 以及 Jakobsen (2003)。基本的

\begin{center}\small 表 4.2 ( r ) 的各种方案 ( r ) 的方案表达式\end{center}
迎风 0 CDS r

\begin{sourceblock}
\begin{tikzpicture}
\node[inner sep=0,anchor=south west] (image) at (0,0) {\includegraphics[page=110,trim={53.00bp 488.37bp 53.87bp 88.55bp},clip,width=0.92\textwidth,max height=0.38\textheight]{Ferziger4th.pdf}};
\begin{scope}[x={(image.south east)},y={(image.north west)}]
\fill[white] (0.37441,0.93421) rectangle (0.42875,1.00000);
\fill[white] (0.66770,0.93522) rectangle (0.68244,1.00000);
\fill[white] (0.37441,0.73694) rectangle (0.46695,0.84981);
\fill[white] (0.66824,0.73694) rectangle (0.71693,0.84981);
\fill[white] (0.73197,0.73694) rectangle (0.75811,0.85328);
\fill[white] (0.75753,0.73694) rectangle (0.80198,0.84981);
\fill[white] (0.69314,0.49025) rectangle (0.74147,0.60659);
\fill[white] (0.74490,0.49384) rectangle (0.76959,0.60659);
\fill[white] (0.77044,0.49025) rectangle (0.79747,0.60659);
\fill[white] (0.67185,0.35631) rectangle (0.72063,0.47926);
\fill[white] (0.72310,0.35978) rectangle (0.74776,0.47265);
\fill[white] (0.74809,0.35732) rectangle (0.76280,0.47007);
\fill[white] (0.76623,0.35978) rectangle (0.79089,0.47265);
\fill[white] (0.79173,0.35631) rectangle (0.81877,0.47265);
\end{scope}
\end{tikzpicture}
\end{sourceblock}
\begin{sourceblock}
\begin{tikzpicture}
\node[inner sep=0,anchor=south west] (image) at (0,0) {\includegraphics[page=110,trim={53.00bp 461.15bp 53.87bp 180.25bp},clip,width=0.92\textwidth,max height=0.38\textheight]{Ferziger4th.pdf}};
\begin{scope}[x={(image.south east)},y={(image.north west)}]
\fill[white] (0.37441,0.29184) rectangle (0.41883,0.69846);
\fill[white] (0.42036,0.29184) rectangle (0.47185,0.69846);
\fill[white] (0.47338,0.29184) rectangle (0.58003,0.69846);
\end{scope}
\end{tikzpicture}
\end{sourceblock}
MINMOD 最大值 [0 , 最小值 ( r , 1 ) ]

超级蜜蜂 max [ 0 , min ( 2 r , 1 ), min ( r , 2 ) ]

选择限制器 (r) 的想法是在 TVD 约束内最大化“反扩散”高阶项的影响。 Hirsch（2007，第 8.3.4 节）给出了如何选择适当的限制器区域并将其显示在 Sweby 通量限制器图上的示例。 TVD 方案 ( 4.37 ) 的一般结果是

\begin{sourceblock}
\includegraphics[page=110,trim={53.00bp 317.61bp 53.87bp 334.20bp},clip,width=0.92\textwidth,max height=0.38\textheight]{Ferziger4th.pdf}
\end{sourceblock}
Waterson 和 Deconinck (2007) 报道了大约 20 个版本的通量限制器的测试，其中 MUSCL (Van Leer 1977) 给出了最好的性能。表 4.2 列出了一小部分限制器功能以供参考。

如果( r ) = 0，则该方案简化为迎风方案，该方案是一阶精确且TVD的；如果( r ) = r，则该方案是CDS方案，该方案是二阶精确的并且不是TVD。在所示的 TVD 方案中，前四个（MUSCL、OSPRE、H-CUI 和 Van Leer Harmonic）在 Waterson 和 Deconinck（2007）测试报告中得分最高（按得分降序排列），并且基本上是二阶准确的。底部的两个是常用的，但测试分数要低得多，并且明显低于二阶准确度。弗林格等人。 (2005)，在使用表中所示的一些 TVD 限制器时，证明了在计算过程中更改限制器以保留特定行为可能很有用；这种变化是通过对他们案例中流动势能行为的测试而促成的。

顺便说一句，我们注意到方程（4.37）和（4.33）的相似性，并且可以推断出 κ 方案（Waterson 和 Deconinck 2007）：

\begin{sourceblock}
\includegraphics[page=110,trim={53.00bp 78.75bp 53.87bp 559.00bp},clip,width=0.92\textwidth,max height=0.38\textheight]{Ferziger4th.pdf}
\end{sourceblock}
本节中描述的用于生成变量面值的插值方案对于稳定情况和不稳定情况都很有用（无论使用何种时间步进方案）。尽管可以导出二维和三维空间的方案，但是通常通过在每个方向上使用一维方案将一维方案扩展到多个维度。如果对流项中的传输速度在空间中变化，则必须使用每项中的实际面通量重新推导这些方案，因为传输速度可能不同，例如，在图 1 和 2 中 CV 的 e 和 w 、n 和 s 或 t 和 b 面上。 4.2 和 4.3。

此外，如上所述，许多流场包括振荡和/或急剧梯度，这需要采取一些措施来防止数值解中出现问题。我们已经讨论了处理这个问题的一些方法，但是还存在我们尚未讨论的其他方案。一类方法称为“本质无振荡”（ENO）和“加权本质无振荡”（WENO）方案，请参见第 1 节。 11.3. Durran（2010）对这些方法进行了深入讨论，并引用了“它们在平滑最大值和最小值附近保持真正高阶精度的能力”。由于它们用于地球物理应用，因此当前有活跃的文献，例如 Gottlieb 等人。 （2006），王等人。 (2016)，以及 Li 和 Xing (1967)，最后一个包括许多说明性例子。

\section*{4.5 边界条件的实现}\phantomsection\addcontentsline{toc}{section}{4.5 边界条件的实现}
每个 CV 提供一个代数方程。对于每个 CV，体积积分的计算方式相同，但通过与域边界一致的 CV 面的通量需要特殊处理。这些边界通量必须是已知的，或者表示为内部值和边界数据的组合。它们不应该引入额外的未知数，因为方程的数量等于单元的数量，因此只有单元中心值可以被视为未知数。由于边界之外没有节点，因此这些近似通常基于单方面差异或外推。

然而，正如我们在该节中所展示的那样。 3.7.2，可以实现具有中心差异的基于导数的边界条件并在域外使用鬼点。然后，在诺依曼边界条件的情况下，通常使用鬼点和第一内部点之间的边界条件关系，将外部点合并到在边界处应用的微分方程中。在本书中我们不会这样做。

通常，对流通量是在流入边界处规定的。对流通量在不渗透壁和对称平面处为零，并且通常假设与流出边界的法向坐标无关；在这种情况下，可以使用迎风近似。有时在壁上指定扩散通量，例如规定指定的热通量（包括零热通量的绝热表面的特殊情况）或变量的边界值。在这种情况下，使用法向梯度的单侧近似来评估扩散通量，如第 1 节中所述。 3.7.如果指定了梯度本身，则将其用于计算通量，并且可以使用以节点值表示的通量近似值来计算变量的边界值。这将在下面的示例中进行演示。

\section*{4.6 代数方程组}\phantomsection\addcontentsline{toc}{section}{4.6 代数方程组}
通过对所有通量近似值和源项求和，我们生成一个代数方程，该方程将 CV 中心的变量值与几个相邻 CV 的值联系起来。方程和未知数的数量都等于 CV 的数量，因此系统是适定的。特定 CV 的代数方程具有 ( 3.43 ) 形式，整个解域的方程组具有方程 ( 3.44 ) 给出的矩阵形式。当节的订购方案时。采用图3.8，矩阵A具有图3.6所示的形式。这仅适用于具有四边形或六面体 CV 的结构化网格；对于其他几何形状，矩阵结构会更复杂（详细信息请参见第 9 章），但它总是稀疏的。任何行中元素的最大数量等于二阶近似的近邻元素数量。对于高阶近似，它取决于方案中使用的邻居的数量。

\section*{4.7 示例}\phantomsection\addcontentsline{toc}{section}{4.7 示例}
为了演示 FV 方法并显示上述离散化方法的一些属性，我们提供了三个示例。

\subsection*{4.7.1 测试 FV 近似值的阶数}\phantomsection\addcontentsline{toc}{subsection}{4.7.1 测试 FV 近似值的阶数}
由于文献中对 FV 近似顺序存在误解，因此此处列出了一些代表性测试的结果。考虑采用 6 级细化的统一二维笛卡尔网格（见图 4.4，其中显示了三个最粗糙的网格）。为了简单起见，考虑定义变量 φ 在二维上的变化的两个解析函数：

\begin{sourceblock}
\includegraphics[page=112,trim={53.00bp 140.71bp 53.87bp 506.97bp},clip,width=0.92\textwidth,max height=0.38\textheight]{Ferziger4th.pdf}
\end{sourceblock}
最粗网格的间距 x = y = 1；然后将其细化五倍，以便最细网格上的 32 个面对应于最粗网格上的单个面。我们测试了计算面质心处变量值的各种插值近似的精度，以及对流通量的表面积分近似的精度。仅针对单个值评估插值近似值

\begin{sourceblock}
\includegraphics[page=113,trim={53.00bp 473.25bp 53.87bp 52.00bp},clip,width=\textwidth,max height=0.38\textheight]{Ferziger4th.pdf}
\par\vspace{0.45\baselineskip}
\small 图4.4 用于测试插值和积分阶数的三个最粗网格的示意图
\end{sourceblock}
位置，同时针对与一个粗网格面相对应的表面评估积分近似。

我们首先考虑二阶 CDS（标记为 CDS2）和四阶 CDS（标记为 CDS4）插值的精度，如方程（4.18）和（4.26）中所定义。变量的节点值是将节点坐标插入式（4.40）的表达式中得到的精确函数值。 x = 0 和 y = 0 处的面质心的插值。将每个网格中的 5 与方程（4.40）中的精确值进行比较。图 4.5 显示了针对多项式和余弦函数细化网格时插值和误差的变化。请注意，在多项式函数的情况下，CDS2 高估了精确值，而 CDS4 低估了精确值。当网格间距连续减半时，两种近似值都以预期的顺序收敛到指定位置处的精确函数值。 CDS4 近似在所有网格上都更准确；最精细网格上的误差要低 4 个数量级。

接下来考虑中点积分近似和辛普森规则的精度，如方程（4.3）和（4.10）中所定义。在这两种情况下，都使用积分点（面质心和面角）处的变量值的三种变体：CDS2 和 CDS4 的近似值，以及方程（4.40）给出的解析表达式的精确值。中点规则和两个函数（多项式和余弦）的结果如图 1 和 2 所示。 4.6.结果非常令人惊讶：在多项式函数的情况下，CDS2 高估了积分，而 CDS4 低估了积分，但 CDS2 的积分误差比 CDS4 或精确的中点变量值低一个数量级。在余弦函数的情况下，CDS2 低估积分，而 CDS4（以及精确的面质心值）高估积分，并且 CDS2 的误差仅比 CDS4 大 3 倍左右。

无论如何，在所有情况下都获得二阶收敛。这些测试证实，使用更准确的插值方案并不一定意味着积分近似也会更准确：这里细网格上的 CDS4 产生的插值误差低了 4 个数量级，但积分近似的误差对于余弦函数来说仅低约 3 倍，而对于多项式函数则高 10 倍！

\begin{sourceblock}
\includegraphics[page=114,trim={53.00bp 288.25bp 53.87bp 52.00bp},clip,width=\textwidth,max height=0.38\textheight]{Ferziger4th.pdf}
\par\vspace{0.45\baselineskip}
\small 图 4.5 对于多项式（上）和余弦函数（下），当网格间距减半时，二阶 CDS 和四阶 CDS 的插值面值（左）以及相关误差（右）的变化
\end{sourceblock}
特别明显的是，插值的阶数不会影响中点规则近似的阶数，只要它是二阶近似或更好（即使面质心处的精确变量值也没有太大区别）。因此，使用QUICK或CDS4等方案的求解方法仅以高于二阶的阶数进行插值，但只要使用中点法则进行积分逼近，整体逼近阶数就不可能优于二阶。

最后，研究了辛普森规则的积分近似。同样，三个积分点（面质心和两个角）处的变量值要么完全从指定的分析函数中获得，要么从 CDS2 或 CDS4 插值中获得。图 4.7 显示了多项式和余弦函数网格细化后积分值和误差的变化。对于这两个函数，CDS4 会导致积分被低估，而 CDS2 会导致多项式函数被高估而余弦函数被低估。这里使用

\begin{sourceblock}
\includegraphics[page=115,trim={53.00bp 295.26bp 53.87bp 52.00bp},clip,width=\textwidth,max height=0.38\textheight]{Ferziger4th.pdf}
\par\vspace{0.45\baselineskip}
\small 图 4.6 对于多项式（上）和余弦函数（下），当面质心处的变量值是精确的或通过 CDS2 或 CDS4 插值获得时，中点法则近似中表面积分的变化（左）以及当网格间距减半时相关误差（右）
\end{sourceblock}
积分点处的精确函数值会导致最低的误差（在所有网格上都超过一个数量级）。 CDS4 总是比 CDS2 产生更低的积分误差；在最粗的网格上，差异并没有那么大，但在最细的网格上，差异超过3个数量级。

这些结果还表明积分计算的阶数等于所涉及的近似（插值和积分近似）的最低阶。使用 CDS2 进行插值导致辛普森规则的积分近似仅变为二阶近似。如果采用一阶迎风格式插值，则总阶优先。然而，还应该注意的是，当使用 CDS2 进行插值时，使用辛普森规则来近似积分时的误差比中点规则的情况下低一个数量级以上，尽管积分的两个近似值都是二阶的（参见图 4.6 和 4.7 中标记为 CD2 的线）。

\begin{sourceblock}
\includegraphics[page=116,trim={53.00bp 290.25bp 53.87bp 52.00bp},clip,width=\textwidth,max height=0.38\textheight]{Ferziger4th.pdf}
\par\vspace{0.45\baselineskip}
\small 图 4.7 对于多项式（上）和余弦函数（下），当面质心处的变量值是精确的或通过 CDS2 或 CDS4 插值获得时，辛普森规则近似中表面积分的变化（左）以及当网格间距减半时相关误差（右）
\end{sourceblock}
如果我们比较计算扩散通量所需的单元表面一阶导数的二阶和四阶近似值以及不同的表面积分近似值，将会得到相同的结论。

本练习的目的是证明 FV 方法的准确性取决于三个因素：(i) 插值到细胞质心以外的位置，(ii) 导数的近似值（对于扩散通量或源项），以及 (iii) 积分的近似值。为了获得最佳结果，需要在这三种近似值之间取得良好的平衡；使一个比其他更准确并不一定会改善整体结果。

\begin{sourceblock}
\begin{tikzpicture}
\node[inner sep=0,anchor=south west] (image) at (0,0) {\includegraphics[page=117,trim={53.00bp 579.30bp 53.87bp 58.59bp},clip,width=0.92\textwidth,max height=0.38\textheight]{Ferziger4th.pdf}};
\begin{scope}[x={(image.south east)},y={(image.north west)}]
\fill[white] (0.67065,0.44743) rectangle (0.72397,0.76743);
\fill[white] (0.72253,0.43221) rectangle (0.74006,0.84035);
\fill[white] (0.85802,0.04248) rectangle (0.90568,0.39044);
\end{scope}
\end{tikzpicture}
\end{sourceblock}
\begin{sourceblock}
\begin{tikzpicture}
\node[inner sep=0,anchor=south west] (image) at (0,0) {\includegraphics[page=117,trim={53.00bp 471.36bp 53.87bp 115.25bp},clip,width=0.92\textwidth,max height=0.38\textheight]{Ferziger4th.pdf}};
\begin{scope}[x={(image.south east)},y={(image.north west)}]
\fill[white] (0.76577,0.54407) rectangle (0.88869,0.65761);
\fill[white] (0.76770,0.46247) rectangle (0.79690,0.58456);
\fill[white] (0.79889,0.46247) rectangle (0.81850,0.58456);
\fill[white] (0.82054,0.47240) rectangle (0.87817,0.59613);
\fill[white] (0.62836,0.39545) rectangle (0.67666,0.50899);
\fill[white] (0.50592,0.12159) rectangle (0.52295,0.25475);
\fill[white] (0.60644,0.12159) rectangle (0.62346,0.25475);
\end{scope}
\end{tikzpicture}
\end{sourceblock}
\subsection*{4.7.2 已知速度场中的标量输运}\phantomsection\addcontentsline{toc}{subsection}{4.7.2 已知速度场中的标量输运}
考虑如图 4.8 所示的已知速度场中标量的传输问题。后者由 u x = x 和 u y = − y 给出，表示驻点附近的无粘流。流线是 xy = const 线。并改变相对于笛卡尔网格的方向。另一方面，在任何单元面上，法向速度分量都是恒定的，因此对流通量近似值中的误差仅取决于用于 φ e 的近似值。这有助于分析准确性。

需要求解的标量传输方程为：

\begin{sourceblock}
\begin{tikzpicture}
\node[inner sep=0,anchor=south west] (image) at (0,0) {\includegraphics[page=117,trim={53.00bp 286.72bp 53.87bp 356.95bp},clip,width=0.92\textwidth,max height=0.38\textheight]{Ferziger4th.pdf}};
\begin{scope}[x={(image.south east)},y={(image.north west)}]
\fill[white] (0.31603,0.05340) rectangle (0.33131,0.43480);
\end{scope}
\end{tikzpicture}
\end{sourceblock}
并应用以下边界条件：

\begin{itemize}
\item 沿北（入口）边界 φ = 0；
\item 沿着西边界，从 y = 1 时的 φ = 0 到 y = 0 时的 φ = 1，φ 的线性变化；
\item 南边界的对称条件（垂直于边界的零梯度）；
\item 出口（东）边界处流动方向的梯度为零。
\end{itemize}
几何形状和流场如图4.8所示。我们给出了“e”面的离散化细节。

对流通量将使用中点规则和 UDS 或 CDS 插值进行评估。我们将对流通量表示为质量通量与 φ 平均面值的乘积：

\begin{sourceblock}
\begin{tikzpicture}
\node[inner sep=0,anchor=south west] (image) at (0,0) {\includegraphics[page=117,trim={53.00bp 97.56bp 53.87bp 543.99bp},clip,width=0.92\textwidth,max height=0.38\textheight]{Ferziger4th.pdf}};
\begin{scope}[x={(image.south east)},y={(image.north west)}]
\fill[white] (0.31450,0.41033) rectangle (0.34965,0.95120);
\end{scope}
\end{tikzpicture}
\end{sourceblock}
其中 ̇ me e 是通过“e”面的质量通量：

\begin{sourceblock}
\begin{tikzpicture}
\node[inner sep=0,anchor=south west] (image) at (0,0) {\includegraphics[page=118,trim={53.00bp 557.55bp 53.87bp 85.47bp},clip,width=0.92\textwidth,max height=0.38\textheight]{Ferziger4th.pdf}};
\begin{scope}[x={(image.south east)},y={(image.north west)}]
\fill[white] (0.29326,0.39836) rectangle (0.33080,0.94810);
\fill[white] (0.33603,0.44810) rectangle (0.36421,0.94810);
\end{scope}
\end{tikzpicture}
\end{sourceblock}
表达式 ( 4.43 ) 在任何网格上都是精确的，因为速度 u x , e 沿表面是恒定的。则通量近似为：

\begin{sourceblock}
\includegraphics[page=118,trim={53.00bp 488.30bp 53.87bp 142.51bp},clip,width=0.92\textwidth,max height=0.38\textheight]{Ferziger4th.pdf}
\end{sourceblock}
线性插值系数 λ e 由方程 ( 4.19 ) 定义。通过简单地绕 P 旋转面“e”直至其折叠到特定面上并替换索引，即可获得通过其他 CV 面的通量的类似表达式。请注意，每个面上的 n 都指向外，即从单元中心 P 到相邻单元的中心，参见。图4.2。例如，在面“w”处，我们得到：

\begin{sourceblock}
\includegraphics[page=118,trim={53.00bp 381.91bp 53.87bp 248.91bp},clip,width=0.92\textwidth,max height=0.38\textheight]{Ferziger4th.pdf}
\end{sourceblock}
\begin{sourceblock}
\begin{tikzpicture}
\node[inner sep=0,anchor=south west] (image) at (0,0) {\includegraphics[page=118,trim={53.00bp 337.93bp 53.87bp 298.69bp},clip,width=0.92\textwidth,max height=0.38\textheight]{Ferziger4th.pdf}};
\begin{scope}[x={(image.south east)},y={(image.north west)}]
\fill[white] (0.00000,0.78523) rectangle (0.04740,1.00000);
\end{scope}
\end{tikzpicture}
\end{sourceblock}
对于 UDS 情况，获得了对流通量对代数方程中的系数的以下贡献：

\begin{sourceblock}
\includegraphics[page=118,trim={53.00bp 256.25bp 53.87bp 368.21bp},clip,width=0.92\textwidth,max height=0.38\textheight]{Ferziger4th.pdf}
\end{sourceblock}
对于 CDS 情况，系数为：

\begin{sourceblock}
\includegraphics[page=118,trim={53.00bp 186.48bp 53.87bp 437.98bp},clip,width=0.92\textwidth,max height=0.38\textheight]{Ferziger4th.pdf}
\end{sourceblock}
A c 的表达式
P 由连续性条件得出：

\begin{sourceblock}
\includegraphics[page=118,trim={53.00bp 140.52bp 53.87bp 510.17bp},clip,width=0.92\textwidth,max height=0.38\textheight]{Ferziger4th.pdf}
\end{sourceblock}
速度场满足。请注意，以节点 P 为中心的 CV 的 ̇ m w 和 λ w 分别等于以节点 W 为中心的 CV 的 - ̇ m e 和 1 - λ e。因此，在计算机代码中，质量通量和插值因子计算一次并存储为每个 CV 的 ̇ me 、 ̇ m n 和 λ e 、 λ n。

使用中点法则和法向导数的 CDS 近似来评估扩散通量积分；这是最简单且使用最广泛的近似：

\begin{sourceblock}
\begin{tikzpicture}
\node[inner sep=0,anchor=south west] (image) at (0,0) {\includegraphics[page=119,trim={53.00bp 534.61bp 53.87bp 101.12bp},clip,width=0.92\textwidth,max height=0.38\textheight]{Ferziger4th.pdf}};
\begin{scope}[x={(image.south east)},y={(image.north west)}]
\fill[white] (0.08800,0.34199) rectangle (0.12436,0.77902);
\fill[white] (0.10487,0.27359) rectangle (0.11898,0.55574);
\fill[white] (0.12980,0.35087) rectangle (0.15798,0.73101);
\end{scope}
\end{tikzpicture}
\end{sourceblock}
请注意 xE = 1
2 ( x i + 1 + x i ) 且 x P = 1 2 ( x i + x i − 1 )，见图 4.2。扩散系数假设为常数；如果不是，则可以在 P 和 E 处的节点值之间进行线性插值。扩散项对代数方程系数的贡献为：

\begin{sourceblock}
\includegraphics[page=119,trim={53.00bp 399.63bp 53.87bp 197.22bp},clip,width=0.92\textwidth,max height=0.38\textheight]{Ferziger4th.pdf}
\end{sourceblock}
对其他 CV 面应用相同的近似，积分方程变为：

\begin{sourceblock}
\begin{tikzpicture}
\node[inner sep=0,anchor=south west] (image) at (0,0) {\includegraphics[page=119,trim={53.00bp 353.88bp 53.87bp 296.80bp},clip,width=0.92\textwidth,max height=0.38\textheight]{Ferziger4th.pdf}};
\begin{scope}[x={(image.south east)},y={(image.north west)}]
\fill[white] (0.00000,0.92367) rectangle (0.11726,1.00000);
\end{scope}
\end{tikzpicture}
\end{sourceblock}
它表示通用节点 P 的方程。系数 A l 是通过对流和扩散贡献求和获得的，参见方程（4.47）、（4.48）和（4.50）：

\begin{sourceblock}
\begin{tikzpicture}
\node[inner sep=0,anchor=south west] (image) at (0,0) {\includegraphics[page=119,trim={53.00bp 292.91bp 53.87bp 355.12bp},clip,width=0.92\textwidth,max height=0.38\textheight]{Ferziger4th.pdf}};
\begin{scope}[x={(image.south east)},y={(image.north west)}]
\fill[white] (0.00000,0.85422) rectangle (0.08484,1.00000);
\end{scope}
\end{tikzpicture}
\end{sourceblock}
其中 l 表示索引 P、E、W、N、S 中的任何一个。A P 等于所有相邻系数的负和是所有保守方案的一个特征，并确保均匀场是离散方程的解。

上述表达式对所有内部 CV 均有效。对于靠近边界的 CV，边界条件要求对方程进行一些修改。在规定了 的北边界和西边界处，法线方向上的梯度使用单侧差分来近似，例如，在西边界处：

\begin{sourceblock}
\begin{tikzpicture}
\node[inner sep=0,anchor=south west] (image) at (0,0) {\includegraphics[page=119,trim={53.00bp 160.92bp 53.87bp 467.80bp},clip,width=0.92\textwidth,max height=0.38\textheight]{Ferziger4th.pdf}};
\begin{scope}[x={(image.south east)},y={(image.north west)}]
\fill[white] (0.38881,0.47194) rectangle (0.42704,0.78087);
\end{scope}
\end{tikzpicture}
\end{sourceblock}
其中 W 表示边界节点，其位置与单元面中心“w”重合。该近似具有一阶精度，但适用于半宽 CV。系数和边界值的乘积被添加到源项中。例如，沿着西边界（索引 i = 2 的 CV），将 A W φ W 添加到 Q P 并将系数 A W 设置为零。这同样适用于北边界处的系数A N。

\begin{sourceblock}
\includegraphics[page=120,trim={53.00bp 439.48bp 53.87bp 52.00bp},clip,width=\textwidth,max height=0.38\textheight]{Ferziger4th.pdf}
\par\vspace{0.45\baselineskip}
\small 图 4.9 φ 的等值线，从 0.05 到 0.95，步长为 0.1（从上到下），对于 = 0.01（左）和 = 0.001（右）
\end{sourceblock}
在南边界， 的法向梯度为零，当应用上述近似时，这意味着边界值等于 CV 中心处的值。因此，对于索引 j = 2 的单元格， φ S = φ P 并且这些 CV 的代数方程修改为：

\begin{sourceblock}
\includegraphics[page=120,trim={53.00bp 319.77bp 53.87bp 330.91bp},clip,width=0.92\textwidth,max height=0.38\textheight]{Ferziger4th.pdf}
\end{sourceblock}
这需要将 A S 添加到 A P，然后设置 A S = 0。出口（东）边界处的零梯度条件以类似的方式实现。

我们现在转向结果。图 4.9 显示了使用 CDS 在 40 × 40 CV 均匀网格上计算出的 φ 等值线，其中对流通量具有两个值：0.001 和 0.01 (ρ = 1. 0)。我们看到，正如预期的那样，随着 的增加，通过流的扩散进行的传输要强得多。

为了评估预测的准确性，我们监测通过西（左）边界的总通量，在该边界处，对流通量为零。由于该边界上的对流通量为零，因此该量是通过将沿该边界的所有 CV 面上的扩散通量相加而获得的，其近似公式为（4.49）和（4.53）。图 4.10 显示了当网格针对对流通量的 UDS 和 CDS 离散化而细化时通量的变化；扩散通量始终使用 CDS 进行离散化。网格从 10 × 10 CV 细化至 320 × 320 CV。在最粗糙的网格上，CDS 不会为 = 0 生成有意义的解。 001;对流占主导地位，在这样的粗网格上，西边界附近短距离内 φ 的快速变化（见图 4.9）会导致振荡如此强烈，以至于大多数迭代求解器无法收敛（局部单元佩克莱特数，Pe = ρ u x x /，在此网格上的范围在 10 到 100 之间）。 （可以借助延迟校正获得收敛解，但这会非常不准确。）随着网格的细化，CDS 结果单调收敛到独立于网格的解。在

\begin{sourceblock}
\includegraphics[page=121,trim={53.00bp 433.48bp 53.87bp 52.00bp},clip,width=\textwidth,max height=0.38\textheight]{Ferziger4th.pdf}
\par\vspace{0.45\baselineskip}
\small 图 4.10 通过西墙的总通量 φ 的收敛（左）以及计算通量的误差作为网格间距的函数，对于 = 0.001
\end{sourceblock}
在 40 × 40 CV 网格中，局部 Peclet 数范围为 2.5 到 25，但解中没有振荡，如图 4.9 所示。

正如预期的那样，UDS 解不会在任何网格上振荡。然而，收敛不是单调的：两个最粗糙网格上的通量低于收敛值；它在下一个网格上太高，然后单调地接近正确的结果。通过假设 CDS 方案的二阶收敛，我们通过 Richardson 外推法估计了与网格无关的解（详细信息请参见第 3.9 节），并且能够确定每个解中的误差。图 4.10 中绘制了 UDS 和 CDS 的误差与归一化网格大小（对于最粗网格，x = 1）的关系。还显示了一阶和二阶方案的预期斜率。 CDS 误差曲线具有二阶方案预期的斜率。 UDS 错误显示前三个网格上的不规则行为。从第四个网格开始，误差曲线接近预期斜率。 320×320 CV网格上的UDS解仍然误差超过1\%； CDS 在 80 × 80 网格上产生更准确的结果！

\subsection*{4.7.3 测试数值扩散}\phantomsection\addcontentsline{toc}{subsection}{4.7.3 测试数值扩散}
另一个流行的测试用例是倾斜于网格线的均匀流动中的阶梯轮廓的对流；见图4.11。可以用上述方法通过调整边界条件（西边界和南边界的φ规定值，北边界和东边界的流出条件）来求解。我们在下面显示了使用 UDS 和 CDS 离散化获得的结果。

\begin{center}\small 图 4.11 与网格线倾斜的均匀流中阶梯轮廓对流的几何形状和边界条件\end{center}
(1,1)

y

\begin{sourceblock}
\begin{tikzpicture}
\node[inner sep=0,anchor=south west] (image) at (0,0) {\includegraphics[page=122,trim={53.00bp 515.89bp 53.87bp 107.93bp},clip,width=0.92\textwidth,max height=0.38\textheight]{Ferziger4th.pdf}};
\begin{scope}[x={(image.south east)},y={(image.north west)}]
\fill[white] (0.57035,0.58837) rectangle (0.60057,0.97164);
\fill[white] (0.72653,0.67698) rectangle (0.74502,0.92982);
\fill[white] (0.75047,0.49953) rectangle (0.77802,0.76063);
\fill[white] (0.72283,0.26465) rectangle (0.74087,0.50236);
\fill[white] (0.74123,0.02836) rectangle (0.76872,0.28875);
\end{scope}
\end{tikzpicture}
\end{sourceblock}
x

(0,0)

由于这种情况下不存在扩散，因此要求解的方程为（微分形式）：

\begin{sourceblock}
\begin{tikzpicture}
\node[inner sep=0,anchor=south west] (image) at (0,0) {\includegraphics[page=122,trim={53.00bp 425.76bp 53.87bp 212.10bp},clip,width=0.92\textwidth,max height=0.38\textheight]{Ferziger4th.pdf}};
\begin{scope}[x={(image.south east)},y={(image.north west)}]
\fill[white] (0.00000,0.86987) rectangle (0.08072,1.00000);
\end{scope}
\end{tikzpicture}
\end{sourceblock}
对于这种情况，两个方向上均匀网格上的 UDS 给出了非常简单的方程：

\begin{sourceblock}
\begin{tikzpicture}
\node[inner sep=0,anchor=south west] (image) at (0,0) {\includegraphics[page=122,trim={53.00bp 368.45bp 53.87bp 269.40bp},clip,width=0.92\textwidth,max height=0.38\textheight]{Ferziger4th.pdf}};
\begin{scope}[x={(image.south east)},y={(image.north west)}]
\fill[white] (0.00000,0.88017) rectangle (0.11558,1.00000);
\end{scope}
\end{tikzpicture}
\end{sourceblock}
这很容易以顺序方式解决，无需迭代。另一方面，CDS 给出主对角线上的系数 A P 为零值，使得求解变得困难。对于这个问题，大多数迭代求解器都无法收敛；然而，通过使用上面提到的和第 3 节中描述的延迟校正方法。 5.6 即可求得解。

如果流动平行于 x 坐标，则两种方案都会给出正确的结果：轮廓只是在下游对流。当流动与网格线倾斜时，UDS 在任何下游横截面产生涂抹的阶梯轮廓，而 CDS 产生振荡。在图 4.12 中，我们显示了 x = 0 时 φ 的轮廓。 5 为流量与网格成 45 ° 的情况（ u x = u y = 1；ρ = 1；在西边界 φ = 0（y < 0 . 1 时）和 φ = 1（0 . 1 < y < 1 时）；在南边界， φ = 0），在具有 20 × 20、40 × 40 和 80 × 的三个均匀网格序列上获得使用三种不同离散化的 80 个 CV：UDS、CDS 以及 95\% CDS 和 5\% UDS 的混合。在 UDS 解中可以清楚地看到数值扩散的效果；步骤高度模糊，并且连续网格上的解之间的差异几乎相同，表明尚未达到一阶渐近收敛 - 当网格间距减半时，需要多次细化网格才能使差异减半。另一方面，CDS 产生的轮廓在台阶处具有适当的陡度，但它会在台阶两侧振荡并产生上冲和下冲。摆动的幅度不会随着网格的细化而减小，只是波长会减小。当 5\% 的 UDS 与 95\% 的 CDS 混合时，摆动幅度显著减小，并且随着网格的细化，摆动幅度明显变小。在这个例子中，不存在物理扩散；在实际流动中，粘度和扩散率始终存在并且处于足够细的网格中，

\begin{sourceblock}
\begin{tikzpicture}
\node[inner sep=0,anchor=south west] (image) at (0,0) {\includegraphics[page=123,trim={53.00bp 518.59bp 53.87bp 57.20bp},clip,width=0.92\textwidth,max height=0.38\textheight]{Ferziger4th.pdf}};
\begin{scope}[x={(image.south east)},y={(image.north west)}]
\fill[white] (0.46740,0.83221) rectangle (0.49335,0.91987);
\fill[white] (0.49383,0.83221) rectangle (0.50815,0.91987);
\fill[white] (0.50863,0.83221) rectangle (0.53456,0.91987);
\fill[white] (0.53504,0.83221) rectangle (0.56626,0.91987);
\fill[white] (0.00000,0.76259) rectangle (0.01612,0.87393);
\fill[white] (0.02066,0.76502) rectangle (0.04532,0.87648);
\fill[white] (0.04758,0.76159) rectangle (0.08502,0.87648);
\fill[white] (0.08659,0.76159) rectangle (0.19468,0.87305);
\fill[white] (0.19618,0.76159) rectangle (0.22647,0.87305);
\fill[white] (0.22803,0.76159) rectangle (0.28376,0.87305);
\fill[white] (0.46740,0.76226) rectangle (0.49335,0.84992);
\fill[white] (0.49383,0.76226) rectangle (0.50815,0.84992);
\fill[white] (0.50863,0.76226) rectangle (0.53456,0.84992);
\fill[white] (0.53504,0.76226) rectangle (0.56626,0.84992);
\fill[white] (0.00000,0.65136) rectangle (0.05507,0.76281);
\fill[white] (0.05663,0.65136) rectangle (0.11660,0.76281);
\fill[white] (0.11817,0.65136) rectangle (0.17389,0.76281);
\fill[white] (0.17549,0.65136) rectangle (0.23615,0.76281);
\fill[white] (0.23774,0.65136) rectangle (0.25386,0.76281);
\fill[white] (0.46740,0.69231) rectangle (0.49335,0.77997);
\fill[white] (0.49383,0.69231) rectangle (0.50815,0.77997);
\fill[white] (0.50860,0.69231) rectangle (0.53456,0.77997);
\fill[white] (0.51389,0.62236) rectangle (0.56626,0.71002);
\fill[white] (0.53504,0.69231) rectangle (0.56626,0.77997);
\fill[white] (0.00000,0.54112) rectangle (0.06069,0.65257);
\fill[white] (0.06226,0.54112) rectangle (0.08827,0.65257);
\fill[white] (0.08983,0.54112) rectangle (0.14132,0.65257);
\fill[white] (0.14292,0.54112) rectangle (0.19723,0.65257);
\fill[white] (0.19880,0.54112) rectangle (0.24036,0.65257);
\fill[white] (0.24192,0.54112) rectangle (0.28069,0.65257);
\fill[white] (0.00000,0.43088) rectangle (0.05507,0.54234);
\fill[white] (0.05663,0.43088) rectangle (0.14914,0.54234);
\fill[white] (0.15071,0.43088) rectangle (0.19227,0.54234);
\fill[white] (0.19380,0.43088) rectangle (0.24818,0.54234);
\fill[white] (0.00000,0.32064) rectangle (0.09326,0.43210);
\fill[white] (0.35465,0.16624) rectangle (0.38764,0.22557);
\end{scope}
\end{tikzpicture}
\end{sourceblock}
\begin{sourceblock}
\begin{tikzpicture}
\node[inner sep=0,anchor=south west] (image) at (0,0) {\includegraphics[page=123,trim={53.00bp 492.29bp 53.87bp 162.37bp},clip,width=0.92\textwidth,max height=0.38\textheight]{Ferziger4th.pdf}};
\begin{scope}[x={(image.south east)},y={(image.north west)}]
\fill[white] (0.39771,0.10453) rectangle (0.43377,0.89547);
\end{scope}
\end{tikzpicture}
\end{sourceblock}
\begin{sourceblock}
\begin{tikzpicture}
\node[inner sep=0,anchor=south west] (image) at (0,0) {\includegraphics[page=123,trim={53.00bp 408.49bp 53.87bp 188.67bp},clip,width=0.92\textwidth,max height=0.38\textheight]{Ferziger4th.pdf}};
\begin{scope}[x={(image.south east)},y={(image.north west)}]
\fill[white] (0.41645,0.85097) rectangle (0.43377,0.98260);
\fill[white] (0.43438,0.76341) rectangle (0.47789,0.89504);
\fill[white] (0.52863,0.76341) rectangle (0.54592,0.89504);
\fill[white] (0.60027,0.76341) rectangle (0.63633,0.89504);
\fill[white] (0.67982,0.76341) rectangle (0.71588,0.89504);
\fill[white] (0.75940,0.76341) rectangle (0.79546,0.89504);
\fill[white] (0.83898,0.76341) rectangle (0.87504,0.89504);
\fill[white] (0.92647,0.76341) rectangle (0.94376,0.89504);
\fill[white] (0.41645,0.39751) rectangle (0.43377,0.52914);
\fill[white] (0.46920,0.32705) rectangle (0.49510,0.44172);
\fill[white] (0.49555,0.32705) rectangle (0.50986,0.44172);
\fill[white] (0.51032,0.32705) rectangle (0.53621,0.44172);
\fill[white] (0.53666,0.32705) rectangle (0.56782,0.44172);
\fill[white] (0.46920,0.23572) rectangle (0.49510,0.35039);
\fill[white] (0.49555,0.23572) rectangle (0.50986,0.35039);
\fill[white] (0.51032,0.23572) rectangle (0.53621,0.35039);
\fill[white] (0.53666,0.23572) rectangle (0.56782,0.35039);
\fill[white] (0.46920,0.14439) rectangle (0.49510,0.25892);
\fill[white] (0.49555,0.14439) rectangle (0.50986,0.25892);
\fill[white] (0.51032,0.14439) rectangle (0.53621,0.25892);
\fill[white] (0.53666,0.14439) rectangle (0.56782,0.25892);
\fill[white] (0.39771,0.01740) rectangle (0.43377,0.14888);
\fill[white] (0.51558,0.05291) rectangle (0.56782,0.16758);
\end{scope}
\end{tikzpicture}
\end{sourceblock}
\begin{sourceblock}
\begin{tikzpicture}
\node[inner sep=0,anchor=south west] (image) at (0,0) {\includegraphics[page=123,trim={53.00bp 356.06bp 53.87bp 272.41bp},clip,width=0.92\textwidth,max height=0.38\textheight]{Ferziger4th.pdf}};
\begin{scope}[x={(image.south east)},y={(image.north west)}]
\fill[white] (0.35465,0.39766) rectangle (0.38764,0.53995);
\fill[white] (0.39771,0.03186) rectangle (0.43377,0.27263);
\end{scope}
\end{tikzpicture}
\end{sourceblock}
\begin{sourceblock}
\begin{tikzpicture}
\node[inner sep=0,anchor=south west] (image) at (0,0) {\includegraphics[page=123,trim={53.00bp 329.83bp 53.87bp 324.83bp},clip,width=0.92\textwidth,max height=0.38\textheight]{Ferziger4th.pdf}};
\begin{scope}[x={(image.south east)},y={(image.north west)}]
\fill[white] (0.39771,0.10453) rectangle (0.43377,0.89547);
\end{scope}
\end{tikzpicture}
\end{sourceblock}
\begin{sourceblock}
\begin{tikzpicture}
\node[inner sep=0,anchor=south west] (image) at (0,0) {\includegraphics[page=123,trim={53.00bp 245.85bp 53.87bp 351.06bp},clip,width=0.92\textwidth,max height=0.38\textheight]{Ferziger4th.pdf}};
\begin{scope}[x={(image.south east)},y={(image.north west)}]
\fill[white] (0.41645,0.85165) rectangle (0.43377,0.98267);
\fill[white] (0.43624,0.76455) rectangle (0.47979,0.89571);
\fill[white] (0.53029,0.76455) rectangle (0.54758,0.89571);
\fill[white] (0.60171,0.76455) rectangle (0.63777,0.89571);
\fill[white] (0.68108,0.76455) rectangle (0.71714,0.89571);
\fill[white] (0.76042,0.76455) rectangle (0.79648,0.89571);
\fill[white] (0.83979,0.76455) rectangle (0.87585,0.89571);
\fill[white] (0.92704,0.76455) rectangle (0.94436,0.89571);
\fill[white] (0.41645,0.39867) rectangle (0.43377,0.52968);
\fill[white] (0.46824,0.32804) rectangle (0.49429,0.44287);
\fill[white] (0.49477,0.32804) rectangle (0.50914,0.44287);
\fill[white] (0.50962,0.32804) rectangle (0.53567,0.44287);
\fill[white] (0.53615,0.32804) rectangle (0.56749,0.44287);
\fill[white] (0.46824,0.23631) rectangle (0.49429,0.35115);
\fill[white] (0.49477,0.23631) rectangle (0.50914,0.35115);
\fill[white] (0.50962,0.23631) rectangle (0.53567,0.35115);
\fill[white] (0.53615,0.23631) rectangle (0.56749,0.35115);
\fill[white] (0.46824,0.14473) rectangle (0.49429,0.25943);
\fill[white] (0.49477,0.14473) rectangle (0.50914,0.25943);
\fill[white] (0.50962,0.14473) rectangle (0.53567,0.25943);
\fill[white] (0.53615,0.14473) rectangle (0.56749,0.25943);
\fill[white] (0.39771,0.01733) rectangle (0.43377,0.14835);
\fill[white] (0.51492,0.05301) rectangle (0.56749,0.16785);
\end{scope}
\end{tikzpicture}
\end{sourceblock}
\begin{sourceblock}
\begin{tikzpicture}
\node[inner sep=0,anchor=south west] (image) at (0,0) {\includegraphics[page=123,trim={53.00bp 193.07bp 53.87bp 435.22bp},clip,width=0.92\textwidth,max height=0.38\textheight]{Ferziger4th.pdf}};
\begin{scope}[x={(image.south east)},y={(image.north west)}]
\fill[white] (0.35465,0.39868) rectangle (0.38764,0.54029);
\fill[white] (0.39771,0.03170) rectangle (0.43377,0.27160);
\end{scope}
\end{tikzpicture}
\end{sourceblock}
\begin{sourceblock}
\begin{tikzpicture}
\node[inner sep=0,anchor=south west] (image) at (0,0) {\includegraphics[page=123,trim={53.00bp 166.67bp 53.87bp 488.00bp},clip,width=0.92\textwidth,max height=0.38\textheight]{Ferziger4th.pdf}};
\begin{scope}[x={(image.south east)},y={(image.north west)}]
\fill[white] (0.39771,0.10462) rectangle (0.43377,0.89538);
\end{scope}
\end{tikzpicture}
\end{sourceblock}
\begin{sourceblock}
\begin{tikzpicture}
\node[inner sep=0,anchor=south west] (image) at (0,0) {\includegraphics[page=123,trim={53.00bp 122.36bp 53.87bp 514.40bp},clip,width=0.92\textwidth,max height=0.38\textheight]{Ferziger4th.pdf}};
\begin{scope}[x={(image.south east)},y={(image.north west)}]
\fill[white] (0.41645,0.65044) rectangle (0.43377,0.95916);
\fill[white] (0.43510,0.44384) rectangle (0.47865,0.75289);
\fill[white] (0.52971,0.44384) rectangle (0.54701,0.75289);
\fill[white] (0.60162,0.44384) rectangle (0.63768,0.75289);
\fill[white] (0.68153,0.44384) rectangle (0.71756,0.75289);
\fill[white] (0.76138,0.44384) rectangle (0.79744,0.75289);
\fill[white] (0.84129,0.44384) rectangle (0.87735,0.75289);
\fill[white] (0.92914,0.44384) rectangle (0.94644,0.75289);
\end{scope}
\end{tikzpicture}
\end{sourceblock}
摆动会在 CDS 溶液中消失，如第 1 节所示。 3.10.局部网格细化将有助于定位，甚至可能消除振荡，这将在第 1 章中讨论。 12.还可以通过局部引入数值扩散来消除振荡（例如，仅在必要时将 CDS 与 UDS 混合，而不是像上面那样在整个解域中均匀混合）。有时，这会在激波附近的可压缩流中完成（有关此类混合的系统方法，请参见第 4.4.6 节）。

通过使用第 3 节中描述的修正方程方法。由图 4.4.6 可以看出（使用关于单元中心的差分方程的泰勒级数展开）方程（4.56）中体现的 UDS 方法更接近于解决对流扩散问题：

\begin{sourceblock}
\includegraphics[page=124,trim={53.00bp 438.16bp 53.87bp 198.73bp},clip,width=0.92\textwidth,max height=0.38\textheight]{Ferziger4th.pdf}
\end{sourceblock}
比原来的方程（4.55）。方程(4.57)是该问题的修正方程。通过将该方程变换为与流动平行和垂直的坐标，可以得出法线方向的有效扩散率为：

\begin{sourceblock}
\includegraphics[page=124,trim={53.00bp 367.34bp 53.87bp 283.34bp},clip,width=0.92\textwidth,max height=0.38\textheight]{Ferziger4th.pdf}
\end{sourceblock}
其中 U 是速度大小，θ 是流相对于 x 方向的角度。 de Vahl Davis (1972) 得出了类似且被广泛引用的结果。正如在该节中一样。 4.4.6 中，我们看到，由于修正方程所暴露的截断误差，FV 离散方程不能精确地求解我们寻求其解的微分方程。

总结我们的发现，我们已经表明：

\begin{itemize}
\item 高阶方案在粗网格上振荡，但随着网格的细化，它比低阶方案更快地收敛到精确解。
\item 一阶UDS 不准确，不应使用。提到这个方案是因为它仍然在一些商业代码中使用。用户应该意识到，使用这种方法无法在经济实惠的网格上获得高精度，尤其是在 3D 中。它在流向和法向方向上引入了很大的扩散误差。 3
\end{itemize}
\begin{itemize}
\item CDS 是最简单的二阶精度方案，在精度、简单性和效率之间提供了良好的折衷。然而，在对流主导的问题中需要小心，如果网格不够精细，可能需要某种 TVD 方法来避免摆动。
\end{itemize}
3 Spalding 混合方案（第 4.4.5 节）也可以产生这种效果，如 Freitas 等人所见。 ( 1985 )，在三维非定常代码的模拟中用 QUICK（第 4.4.3 节）取代该方案揭示了被混合方案的虚假数值扩散所隐藏的涡流和其他三维效应。

\chapter{线性方程组的求解}
\section*{5.1 简介}\phantomsection\addcontentsline{toc}{section}{5.1 简介}
在前两章中，我们展示了如何使用 FD 和 FV 方法离散对流扩散方程。无论哪种情况，离散化过程的结果都是代数方程组，根据推导它们的偏微分方程的性质，代数方程组是线性的或非线性的。在非线性情况下，必须通过迭代技术来求解离散方程，其中包括猜测解、对该解的方程进行线性化以及改进解；重复该过程直到获得收敛结果。因此，无论方程是否是线性的，都需要求解线性代数方程组的有效方法。

从偏微分方程导出的矩阵总是稀疏的，即它们的大部分元素为零。下面将描述一些用于求解使用结构化网格时出现的方程的方法；矩阵的所有非零元素都位于少量明确定义的对角线上；我们可以利用这个结构。一些方法也适用于由非结构化网格产生的矩阵。

用五点近似（迎风或中心差）离散的二维问题的系数矩阵结构如图 3.6 所示。一个 CV 或网格节点的代数方程由方程 1 给出。 ( 3.43 )，完整问题的矩阵版本由式(3.43)给出。 ( 3.44 )，见第 3 节。 3.8，这里重复一下：

\begin{sourceblock}
\begin{tikzpicture}
\node[inner sep=0,anchor=south west] (image) at (0,0) {\includegraphics[page=125,trim={53.00bp 159.00bp 53.87bp 492.81bp},clip,width=0.92\textwidth,max height=0.38\textheight]{Ferziger4th.pdf}};
\begin{scope}[x={(image.south east)},y={(image.north west)}]
\fill[white] (0.44054,0.09072) rectangle (0.48198,0.91626);
\fill[white] (0.48698,0.10956) rectangle (0.51516,0.91626);
\fill[white] (0.51868,0.09072) rectangle (0.54680,0.89742);
\fill[white] (0.55032,0.10956) rectangle (0.56337,0.91626);
\fill[white] (0.93937,0.08374) rectangle (1.00000,0.89044);
\end{scope}
\end{tikzpicture}
\end{sourceblock}
除了描述表示离散偏微分方程的线性代数系统的一些更好的求解方法之外，我们还将在本章中讨论非线性方程组的求解。然而，我们从线性方程开始。假设读者已经接触过线性系统的求解方法，因此描述很简短。

\section*{5.2 直接法}\phantomsection\addcontentsline{toc}{section}{5.2 直接法}
假设矩阵 A 非常稀疏。事实上，我们将遇到的最复杂的矩阵是块类型的带状矩阵；这极大地简化了求解任务，但我们将简要回顾一下一般矩阵的方法，因为稀疏矩阵的方法与它们密切相关。对于设计用于处理全矩阵的方法的描述，使用全矩阵表示法（相对于之前介绍的对角线表示法）更明智，并将被采用。

\subsection*{5.2.1 高斯消除}\phantomsection\addcontentsline{toc}{subsection}{5.2.1 高斯消除}
求解线性代数方程组的基本方法是高斯消去法。其基础是将大型方程组系统地简化为较小的方程组。在此过程中，矩阵的元素被修改，但由于因变量名称没有改变，因此可以方便地仅用矩阵来描述该方法：

\begin{sourceblock}
\begin{tikzpicture}
\node[inner sep=0,anchor=south west] (image) at (0,0) {\includegraphics[page=126,trim={53.00bp 334.59bp 53.87bp 265.77bp},clip,width=0.92\textwidth,max height=0.38\textheight]{Ferziger4th.pdf}};
\begin{scope}[x={(image.south east)},y={(image.north west)}]
\fill[white] (0.00000,0.82700) rectangle (0.08235,1.00000);
\fill[white] (0.08502,0.82700) rectangle (0.16307,1.00000);
\end{scope}
\end{tikzpicture}
\end{sourceblock}
该算法的核心是消除 A 21 的技术，即用零替换它。这是通过将第一个方程（矩阵的第一行）乘以 A 21 / A 11 并从第二行或方程中减去它来完成的；在此过程中，矩阵第二行中的所有元素都将被修改，方程右侧的力向量的第二个元素也是如此。矩阵第一列的其他元素 A 31 、A 41 、...。。。 , A n 1 类似处理；例如，要消除 A i 1，请将矩阵的第一行乘以 A i 1 / A 11 并从第 i 行中减去。通过系统地向下推进矩阵的第一列，A 11 以下的所有元素都被消除。当这个过程完成时，方程 2 、 3 、 .。。 , n 包含变量 φ 1；它们是变量 φ 2 、 φ 3 、...的一组 n − 1 个方程。。。， φ n。然后将相同的过程应用于这个较小的方程组 - 第二列中 A 22 以下的所有元素都被消除。

此过程针对第 1、2、3 列执行。。。 , n − 1。此过程完成后，原始矩阵已被替换为上三角矩阵：

\begin{sourceblock}
\includegraphics[page=126,trim={53.00bp 90.72bp 53.87bp 510.32bp},clip,width=0.92\textwidth,max height=0.38\textheight]{Ferziger4th.pdf}
\end{sourceblock}
除第一行中的元素外，所有元素都与原始矩阵 A 中的元素不同。由于不再需要原始矩阵的元素，因此存储修改后的元素来代替原始元素是有效的。 （在极少数情况下需要保存原始矩阵，可以在开始消除过程之前创建一个副本。）

刚刚描述的算法的这一部分称为前向消除。方程右侧的元素 Q i 也在此过程中进行了修改。

由前向消除得到的上三角方程组很容易求解。最后一个方程仅包含一个变量 φ n 并且很容易求解：

\begin{sourceblock}
\includegraphics[page=127,trim={53.00bp 450.86bp 53.87bp 186.44bp},clip,width=0.92\textwidth,max height=0.38\textheight]{Ferziger4th.pdf}
\end{sourceblock}
倒数第二个方程仅包含 φ n − 1 和 φ n，一旦已知 φ n，就可以求解 φ n − 1。如此往上，依次求解各个方程；第 i 个方程得出 φ i：

\begin{sourceblock}
\includegraphics[page=127,trim={53.00bp 349.07bp 53.87bp 267.26bp},clip,width=0.92\textwidth,max height=0.38\textheight]{Ferziger4th.pdf}
\end{sourceblock}
右边是可计算的，因为总和中出现的所有 φ k 都已被评估。以这种方式，可以计算所有变量。高斯消除算法中从三角矩阵开始计算未知数的部分称为回代。

不难看出，对于较大的 n，通过高斯消元法求解 n 个方程组的线性系统所需的运算次数与 n 3 / 3 成正比。这项工作的大部分是在前向消元阶段；后向替换仅需要 n 2 / 2 算术运算，并且比前向消除的成本要低得多。因此，高斯消除法的成本很高，但对于满矩阵，它与任何可用的方法一样好。高斯消元法的高成本激励人们寻找更有效的特殊矩阵求解器，例如微分方程离散化产生的稀疏矩阵。

对于非稀疏的大型系统，高斯消除法很容易受到错误累积的影响（参见 Golub 和 van Loan 1996；Watkins 2010 ），如果不进行修改，这将使其变得不可靠。添加枢轴或交换行以使枢轴元素（分母中出现的对角线元素）尽可能大，可以控制误差的增长。幸运的是，对于稀疏矩阵，误差累积很少成为问题，因此这个问题在这里并不重要。

高斯消除不能很好地矢量化或并行化，并且在 CFD 问题中很少未经修改而使用。

\subsection*{5.2.2 LU分解}\phantomsection\addcontentsline{toc}{subsection}{5.2.2 LU分解}
已经提出了高斯消除的多种变体。大多数人在这里没什么兴趣。 CFD 价值的一种变体是 LU 分解。它是在没有推导的情况下呈现的。

我们已经看到，在高斯消除中，前向消除将全矩阵简化为上三角矩阵。这个过程可以通过将原始矩阵 A 乘以下三角矩阵来以更正式的方式执行。就其本身而言，这没什么意义，但是，由于下三角矩阵的逆也是下三角矩阵，因此该结果表明，任何矩阵 A 都可以分解为下三角矩阵 ( L ) 和上三角矩阵 ( U ) 的乘积，但受到此处可以忽略的一些限制：

\begin{sourceblock}
\includegraphics[page=128,trim={53.00bp 447.57bp 53.87bp 204.24bp},clip,width=0.92\textwidth,max height=0.38\textheight]{Ferziger4th.pdf}
\end{sourceblock}
为了使分解唯一，我们要求 L 、 L ii 的对角线元素全部为 1；或者，可以要求 U 的对角线元素是统一的。

这种因式分解的有用之处在于它很容易构造。上三角矩阵U正是高斯消元法前向阶段产生的矩阵。此外，L 的元素是消除过程中使用的乘法因子（例如 A ji / A ii ）。这允许通过对高斯消元法进行微小修改来构建因式分解。此外，L和U的元素可以存储在A的元素所在的位置。

这种因式分解的存在允许方程组（5.1）分两个阶段求解。根据定义：

\begin{sourceblock}
\includegraphics[page=128,trim={53.00bp 292.15bp 53.87bp 359.66bp},clip,width=0.92\textwidth,max height=0.38\textheight]{Ferziger4th.pdf}
\end{sourceblock}
方程组 ( 5.1 ) 变为：

\begin{sourceblock}
\includegraphics[page=128,trim={53.00bp 244.33bp 53.87bp 407.48bp},clip,width=0.92\textwidth,max height=0.38\textheight]{Ferziger4th.pdf}
\end{sourceblock}
后一组方程可以通过高斯消去法的向后代入阶段所使用的方法的变体来求解，其中从系统的顶部而不是底部开始。一旦方程。 ( 5.8 ) 已解出 Y 方程。 (5.7) 与高斯消元法后代入阶段求解的三角方程组相同，可解出 φ。

LU 分解相对于高斯消除的优点是可以在不知道向量 Q 的情况下执行分解。因此，如果要求解涉及相同矩阵的多个系统，首先进行因式分解可以节省大量成本；然后可以根据需要对系统进行求解。正如我们将在下面看到的，LU 分解的变体是一些求解线性方程组的更好迭代方法的基础；这是在这里介绍它的主要原因。

\subsection*{5.2.3 三对角系}\phantomsection\addcontentsline{toc}{subsection}{5.2.3 三对角系}
当常微分方程（一维问题）进行有限差分时，例如使用 CDS 近似，所得代数方程具有特别简单的结构。每个方程仅包含其自身节点及其左、右邻居的变量：

\begin{sourceblock}
\includegraphics[page=129,trim={53.00bp 505.37bp 53.87bp 142.93bp},clip,width=0.92\textwidth,max height=0.38\textheight]{Ferziger4th.pdf}
\end{sourceblock}
相应的矩阵 A 仅在其主对角线（由 A P 表示）以及紧邻其上方和下方的对角线（由
分别为 A E 和 A W ）。这样的矩阵称为三对角矩阵；包含三对角矩阵的系统特别容易求解。矩阵元素最好存储为三个 n × 1 数组。

对于三对角系统，高斯消去尤其容易：在前向消去过程中，每行只需消去一个元素。当算法进行到第 i 行时，只有 A i
P需要修改；新值为：

\begin{sourceblock}
\begin{tikzpicture}
\node[inner sep=0,anchor=south west] (image) at (0,0) {\includegraphics[page=129,trim={53.00bp 360.79bp 53.87bp 272.08bp},clip,width=0.92\textwidth,max height=0.38\textheight]{Ferziger4th.pdf}};
\begin{scope}[x={(image.south east)},y={(image.north west)}]
\fill[white] (0.00000,0.84761) rectangle (0.03248,1.00000);
\end{scope}
\end{tikzpicture}
\end{sourceblock}
其中该方程应从程序员的角度理解，即结果存储在原始 A i 的位置
P。强迫术语也被修改：

\begin{sourceblock}
\includegraphics[page=129,trim={53.00bp 286.78bp 53.87bp 346.13bp},clip,width=0.92\textwidth,max height=0.38\textheight]{Ferziger4th.pdf}
\end{sourceblock}
该方法的回代部分也很简单。第 i 个变量的计算公式为：

\begin{sourceblock}
\begin{tikzpicture}
\node[inner sep=0,anchor=south west] (image) at (0,0) {\includegraphics[page=129,trim={53.00bp 228.38bp 53.87bp 405.49bp},clip,width=0.92\textwidth,max height=0.38\textheight]{Ferziger4th.pdf}};
\begin{scope}[x={(image.south east)},y={(image.north west)}]
\fill[white] (0.00000,0.84630) rectangle (0.07071,1.00000);
\end{scope}
\end{tikzpicture}
\end{sourceblock}
这种三对角求解方法有时称为托马斯算法或三对角矩阵算法 (TDMA)。它很容易编程（FORTRAN 代码只需要八个可执行行），更重要的是，运算次数与未知数的数量 n 成正比，而不是与全矩阵高斯消元法的 n 3 成正比。换句话说，每个未知数的成本与未知数的数量无关，这几乎是人们所期望的良好缩放。低成本表明只要有可能就采用该算法。许多求解方法都利用该方法的低成本优势，减少涉及三对角矩阵的问题。

\subsection*{5.2.4 循环归约}\phantomsection\addcontentsline{toc}{subsection}{5.2.4 循环归约}
有些情况甚至更特殊，可以比 TDMA 提供更大的成本降低。系统提供了一个有趣的例子，其中矩阵不仅是三对角的，而且每条对角线上的所有元素都是相同的。循环约简法可用于求解此类系统，每个变量的成本实际上随着系统变大而降低。让我们看看这怎么可能。

假设在系统( 5.9 )中，系数 A i
W , A i
P and A i
E 与索引 i 无关；然后我们可以删除索引。然后，对于 i 的偶数值，我们将行 i − 1 乘以 A W / A P 并从行 i 中减去它。然后我们将第 i + 1 行乘以
A E / A P 并从第 i 行中减去它。这消除了偶数行中主对角线左侧和右侧的元素，但将主对角线左侧两列的零元素替换为 - A 2
W / A P 和零元素在主对角线右侧两列 - A 2
E / A P;对角线元素变为 A P − 2 A W A E / A P。由于每偶数行的元素相同，因此新元素的计算只需进行一次；这就是节省的来源。

完成这些操作后，偶数方程仅包含偶数索引变量，并构成这些变量的一组 n/2 方程；被视为一个单独的系统，这些方程是三对角的，并且简化矩阵的每个对角线上的元素再次相等。换句话说，简化后的方程组与原始方程组具有相同的形式，但大小只有原来的一半。可以用同样的方法进一步减少。如果原始集合中方程的数量是2的幂（或某些其他方便的数字），则可以继续该方法，直到只剩下一个方程；后者直接解决。然后可以通过回代的变体找到剩余的变量。

可以看出，这种方法的成本与 log 2 n 成正比，因此每个变量的成本随着变量数量的增加而降低。尽管该方法可能看起来相当专业，但它在 CFD 应用中发挥着作用。这些是非常规则的几何形状的流动，例如在湍流的间接大涡流模拟和一些气象应用中使用的矩形框。

在这些应用中，循环约简和相关方法为直接（即非迭代）求解椭圆方程（例如拉普拉斯方程和泊松方程）的方法提供了基础。由于解在不包含迭代误差的意义上也是精确的，因此无论何时使用该方法都是无价的。 1
循环约简与快速傅里叶变换密切相关，快速傅里叶变换也用于求解简单几何中的椭圆方程。傅里叶方法也可用于评估导数，如第 1 节所示。 3.11.

\footnotetext[1]{比尼等人。 ( 2009 ) 回顾循环约简的历史、扩展以及新的证明和公式。它被用作高度并行的多重网格应用程序的平滑器，图形处理器（GPU）现在用于流体流动计算（Göddeke 和 Strzodka 2011）。}
\section*{5.3 迭代方法}\phantomsection\addcontentsline{toc}{section}{5.3 迭代方法}
\subsection*{5.3.1 基本概念}\phantomsection\addcontentsline{toc}{subsection}{5.3.1 基本概念}
任何方程组都可以通过高斯消元法或 LU 分解来求解。不幸的是，稀疏矩阵的三角因子并不稀疏，因此这些方法的成本相当高。此外，离散化误差通常远大于计算机算术的精度，因此没有理由如此准确地求解系统。比离散化方案更精确的解就足够了。

这为迭代方法留下了空间。它们出于非线性问题的需要而被使用，但它们对于稀疏线性系统同样有价值。在迭代方法中，人们猜测一个解，并使用方程来系统地改进它。如果每次迭代都很便宜并且迭代次数很少，则迭代求解器的成本可能低于直接方法。在 CFD 问题中通常是这种情况。

考虑方程表示的矩阵问题。 ( 5.1 ) 这可能是由流动问题的 FD 或 FV 近似产生的。经过 n 次迭代后，我们得到了一个近似解 φ n，它并不完全满足这些方程。相反，有一个非零残差 ρ n：

\begin{sourceblock}
\begin{tikzpicture}
\node[inner sep=0,anchor=south west] (image) at (0,0) {\includegraphics[page=131,trim={53.00bp 357.90bp 53.87bp 289.78bp},clip,width=0.92\textwidth,max height=0.38\textheight]{Ferziger4th.pdf}};
\begin{scope}[x={(image.south east)},y={(image.north west)}]
\fill[white] (0.00000,0.71289) rectangle (0.09898,1.00000);
\fill[white] (0.10168,0.71289) rectangle (0.14586,1.00000);
\end{scope}
\end{tikzpicture}
\end{sourceblock}
通过从方程中减去该方程。 ( 5.1 )，我们得到迭代误差之间的关系，定义为：

\begin{sourceblock}
\begin{tikzpicture}
\node[inner sep=0,anchor=south west] (image) at (0,0) {\includegraphics[page=131,trim={53.00bp 310.09bp 53.87bp 337.59bp},clip,width=0.92\textwidth,max height=0.38\textheight]{Ferziger4th.pdf}};
\begin{scope}[x={(image.south east)},y={(image.north west)}]
\fill[white] (0.00000,0.71235) rectangle (0.06232,1.00000);
\fill[white] (0.06499,0.71235) rectangle (0.15798,1.00000);
\fill[white] (0.16072,0.71235) rectangle (0.20382,1.00000);
\end{scope}
\end{tikzpicture}
\end{sourceblock}
其中 φ 是收敛解，残差：

\begin{sourceblock}
\includegraphics[page=131,trim={53.00bp 262.26bp 53.87bp 385.43bp},clip,width=0.92\textwidth,max height=0.38\textheight]{Ferziger4th.pdf}
\end{sourceblock}
迭代过程的目的是使残差为零；在此过程中，ε也变为零。要了解如何做到这一点，请考虑线性系统的迭代方案；这样的方案可以写成：

\begin{sourceblock}
\includegraphics[page=131,trim={53.00bp 190.53bp 53.87bp 457.15bp},clip,width=0.92\textwidth,max height=0.38\textheight]{Ferziger4th.pdf}
\end{sourceblock}
迭代方法必须要求的一个明显性质是收敛结果满足式(1)。 （5.1）。根据定义，收敛时 φ n + 1 = φ n = φ，我们必须有：

\begin{sourceblock}
\begin{tikzpicture}
\node[inner sep=0,anchor=south west] (image) at (0,0) {\includegraphics[page=131,trim={53.00bp 130.75bp 53.87bp 521.06bp},clip,width=0.92\textwidth,max height=0.38\textheight]{Ferziger4th.pdf}};
\begin{scope}[x={(image.south east)},y={(image.north west)}]
\fill[white] (0.00000,0.91835) rectangle (0.03907,1.00000);
\fill[white] (0.04174,0.91835) rectangle (0.10484,1.00000);
\fill[white] (0.10752,0.91835) rectangle (0.17615,1.00000);
\end{scope}
\end{tikzpicture}
\end{sourceblock}
或者，我们可以更一般地写，

\begin{sourceblock}
\includegraphics[page=131,trim={53.00bp 82.93bp 53.87bp 568.88bp},clip,width=0.92\textwidth,max height=0.38\textheight]{Ferziger4th.pdf}
\end{sourceblock}
其中 P 是非奇异的，即所谓的预处理矩阵。预处理矩阵可以显著加速迭代的收敛，在第 4 节中讨论。见下文 5.3.6.1。该迭代方法的替代版本可以通过从方程的每一侧减去 M φ n 来获得。 （5.16）。我们得到：

\begin{sourceblock}
\includegraphics[page=132,trim={53.00bp 523.29bp 53.87bp 124.28bp},clip,width=0.92\textwidth,max height=0.38\textheight]{Ferziger4th.pdf}
\end{sourceblock}
其中 δ n = φ n + 1 − φ n 称为校正或更新，是迭代误差的近似值。

为了使迭代方法有效，求解系统（5.16）必须便宜并且该方法必须快速收敛。廉价的迭代需要计算
N φ n 和系统的解必须易于执行。第一个要求很容易满足；因为A是稀疏的，所以N也是稀疏的，并且N φ n 的计算很简单。第二个要求意味着迭代矩阵M必须易于求逆；从实用的角度来看，M应该是对角线、三对角线、三角形，或者可能是块三对角线或三角形；另一种可能性描述如下。为了快速收敛，M 应该很好地逼近 A，使得 N φ 在某种意义上很小。这将在下一节中讨论。

\subsection*{5.3.2 收敛}\phantomsection\addcontentsline{toc}{subsection}{5.3.2 收敛}
正如我们所指出的，迭代方法的快速收敛是其有效性的关键。在这里，我们给出了一个简单的分析，有助于理解决定收敛速度的因素，并提供如何改进收敛速度的见解。

首先，我们推导确定迭代误差行为的方程。为了找到它，我们回想一下，在收敛时， φ n + 1 = φ n = φ，因此收敛解遵循以下方程：

\begin{sourceblock}
\begin{tikzpicture}
\node[inner sep=0,anchor=south west] (image) at (0,0) {\includegraphics[page=132,trim={53.00bp 236.36bp 53.87bp 415.45bp},clip,width=0.92\textwidth,max height=0.38\textheight]{Ferziger4th.pdf}};
\begin{scope}[x={(image.south east)},y={(image.north west)}]
\fill[white] (0.00000,0.91835) rectangle (0.10066,1.00000);
\fill[white] (0.10334,0.91835) rectangle (0.17756,1.00000);
\fill[white] (0.18033,0.91835) rectangle (0.22174,1.00000);
\fill[white] (0.22439,0.91835) rectangle (0.34069,1.00000);
\end{scope}
\end{tikzpicture}
\end{sourceblock}
从方程中减去这个方程。 ( 5.16 ) 并使用迭代误差的定义( 5.14 )，我们发现：

\begin{sourceblock}
\begin{tikzpicture}
\node[inner sep=0,anchor=south west] (image) at (0,0) {\includegraphics[page=132,trim={53.00bp 188.54bp 53.87bp 459.16bp},clip,width=0.92\textwidth,max height=0.38\textheight]{Ferziger4th.pdf}};
\begin{scope}[x={(image.south east)},y={(image.north west)}]
\fill[white] (0.00000,0.71312) rectangle (0.06403,1.00000);
\fill[white] (0.06671,0.71312) rectangle (0.13603,1.00000);
\fill[white] (0.13868,0.71312) rectangle (0.17844,1.00000);
\fill[white] (0.18111,0.71312) rectangle (0.24087,1.00000);
\end{scope}
\end{tikzpicture}
\end{sourceblock}
\begin{sourceblock}
\includegraphics[page=132,trim={53.00bp 152.68bp 53.87bp 487.55bp},clip,width=0.92\textwidth,max height=0.38\textheight]{Ferziger4th.pdf}
\end{sourceblock}
n →∞ ε n = 0。关键作用是由
如果 lim 则迭代方法收敛
迭代矩阵 M − 1 N 的特征值 λ k 和特征向量 ψ k 定义为：

\begin{sourceblock}
\begin{tikzpicture}
\node[inner sep=0,anchor=south west] (image) at (0,0) {\includegraphics[page=132,trim={53.00bp 87.04bp 53.87bp 559.81bp},clip,width=0.92\textwidth,max height=0.38\textheight]{Ferziger4th.pdf}};
\begin{scope}[x={(image.south east)},y={(image.north west)}]
\fill[white] (0.00000,0.71903) rectangle (0.04241,1.00000);
\end{scope}
\end{tikzpicture}
\end{sourceblock}
其中 K 是方程（网格点）的数量。我们假设特征向量形成一个完整的集合，即 R n 的基础，即所有 n 分量向量的向量空间。如果是这样，则初始误差可以用以下形式表示：

\begin{sourceblock}
\includegraphics[page=133,trim={53.00bp 526.43bp 53.87bp 104.17bp},clip,width=0.92\textwidth,max height=0.38\textheight]{Ferziger4th.pdf}
\end{sourceblock}
其中 a k 是常数。然后迭代过程（5.22）产生：

\begin{sourceblock}
\includegraphics[page=133,trim={53.00bp 460.09bp 53.87bp 170.52bp},clip,width=0.92\textwidth,max height=0.38\textheight]{Ferziger4th.pdf}
\end{sourceblock}
并且，通过归纳法，不难证明

\begin{sourceblock}
\includegraphics[page=133,trim={53.00bp 393.94bp 53.87bp 236.66bp},clip,width=0.92\textwidth,max height=0.38\textheight]{Ferziger4th.pdf}
\end{sourceblock}
显然，当n很大时，若要使ε n 为零，充分必要条件是所有特征值的大小必须小于1。特别是，对于最大特征值来说一定是这样，其大小称为矩阵 M − 1 N 的谱半径。事实上，经过多次迭代后，方程中的项。 ( 5.26 ) 包含小幅度特征值的项变得非常小，并且仅包含最大特征值的项（我们可以将其视为 λ 1 并假设是唯一的）仍然存在：

\begin{sourceblock}
\begin{tikzpicture}
\node[inner sep=0,anchor=south west] (image) at (0,0) {\includegraphics[page=133,trim={53.00bp 287.82bp 53.87bp 358.97bp},clip,width=0.92\textwidth,max height=0.38\textheight]{Ferziger4th.pdf}};
\begin{scope}[x={(image.south east)},y={(image.north west)}]
\fill[white] (0.00000,0.71990) rectangle (0.02746,1.00000);
\fill[white] (0.03014,0.71990) rectangle (0.06322,1.00000);
\fill[white] (0.06589,0.71990) rectangle (0.16223,1.00000);
\fill[white] (0.16496,0.71990) rectangle (0.27296,1.00000);
\end{scope}
\end{tikzpicture}
\end{sourceblock}
如果收敛定义为将迭代误差减少到某个容差 δ 以下，我们需要：

\begin{sourceblock}
\begin{tikzpicture}
\node[inner sep=0,anchor=south west] (image) at (0,0) {\includegraphics[page=133,trim={53.00bp 239.99bp 53.87bp 406.91bp},clip,width=0.92\textwidth,max height=0.38\textheight]{Ferziger4th.pdf}};
\begin{scope}[x={(image.south east)},y={(image.north west)}]
\fill[white] (0.00000,0.72505) rectangle (0.05901,1.00000);
\fill[white] (0.06168,0.72505) rectangle (0.08806,1.00000);
\fill[white] (0.09077,0.72505) rectangle (0.13050,1.00000);
\fill[white] (0.13320,0.72505) rectangle (0.23119,1.00000);
\end{scope}
\end{tikzpicture}
\end{sourceblock}
对该方程两边取对数，我们找到所需迭代次数的表达式：

\begin{sourceblock}
\begin{tikzpicture}
\node[inner sep=0,anchor=south west] (image) at (0,0) {\includegraphics[page=133,trim={53.00bp 165.89bp 53.87bp 450.51bp},clip,width=0.92\textwidth,max height=0.38\textheight]{Ferziger4th.pdf}};
\begin{scope}[x={(image.south east)},y={(image.north west)}]
\fill[white] (0.00000,0.80860) rectangle (0.10394,1.00000);
\fill[white] (0.10665,0.80860) rectangle (0.20295,1.00000);
\fill[white] (0.20568,0.80860) rectangle (0.23543,1.00000);
\fill[white] (0.23814,0.80860) rectangle (0.36274,1.00000);
\fill[white] (0.40632,0.18416) rectangle (0.42611,0.41677);
\fill[white] (0.43098,0.18959) rectangle (0.45916,0.42199);
\end{scope}
\end{tikzpicture}
\end{sourceblock}
我们看到，如果谱半径非常接近1，迭代过程将收敛得非常慢。

作为一个简单（琐碎可能更具描述性）的例子，考虑单个方程的情况（人们永远不会梦想使用迭代方法）。假设我们要解决：

\begin{sourceblock}
\begin{tikzpicture}
\node[inner sep=0,anchor=south west] (image) at (0,0) {\includegraphics[page=133,trim={53.00bp 85.86bp 53.87bp 565.95bp},clip,width=0.92\textwidth,max height=0.38\textheight]{Ferziger4th.pdf}};
\begin{scope}[x={(image.south east)},y={(image.north west)}]
\fill[white] (0.00000,0.91835) rectangle (0.10565,1.00000);
\fill[white] (0.10839,0.91835) rectangle (0.14812,1.00000);
\fill[white] (0.15080,0.91835) rectangle (0.21353,1.00000);
\fill[white] (0.21627,0.91835) rectangle (0.24436,1.00000);
\fill[white] (0.24704,0.91835) rectangle (0.32295,1.00000);
\end{scope}
\end{tikzpicture}
\end{sourceblock}
我们使用迭代方法（注意 m = a + n 且 p 是迭代计数器）：

\begin{sourceblock}
\includegraphics[page=134,trim={53.00bp 571.10bp 53.87bp 76.59bp},clip,width=0.92\textwidth,max height=0.38\textheight]{Ferziger4th.pdf}
\end{sourceblock}
那么误差服从等式的标量。 （5.22）：

\begin{sourceblock}
\includegraphics[page=134,trim={53.00bp 513.60bp 53.87bp 124.53bp},clip,width=0.92\textwidth,max height=0.38\textheight]{Ferziger4th.pdf}
\end{sourceblock}
我们看到，如果 n / m 很小，即如果 n 很小，这意味着 m ≈ a，误差会很快减小。在构建系统的迭代方法时，我们会发现类似的结果成立：M越接近A，收敛越快。

在迭代方法中，能够估计迭代误差以便决定何时停止迭代非常重要。计算迭代矩阵的特征值很困难（通常不明确已知），因此必须使用近似值。我们将在本章后面描述一些估计迭代误差的方法和停止迭代的标准。

\subsection*{5.3.3 一些基本方法}\phantomsection\addcontentsline{toc}{subsection}{5.3.3 一些基本方法}
在最简单的方法（雅可比方法）中， M 是一个对角矩阵，其元素是 A 的对角元素。对于拉普拉斯方程的五点离散，如果每次迭代都从域的左下（西南）角开始，并且我们使用上面介绍的地理符号，则方法为：

\begin{sourceblock}
\includegraphics[page=134,trim={53.00bp 253.65bp 53.87bp 382.12bp},clip,width=0.92\textwidth,max height=0.38\textheight]{Ferziger4th.pdf}
\end{sourceblock}
可以看出，为了收敛，该方法需要与一个方向上的网格点数的平方成正比的迭代次数。这意味着它比直接求解器更昂贵，因此没有理由使用它。

在高斯-赛德尔方法中，M 是 A 的下三角部分。由于它是下面给出的SOR方法的特例，我们不会单独给出方程。它的收敛速度是雅可比方法的两倍，但这还不足以带来有用的改进。

更好的方法之一是高斯-赛德尔方法的加速版本，称为连续过度松弛或 SOR，我们将在下面描述。有关 Jacobi 和 Gauss-Seidel 方法的介绍和分析，请参阅数值方法的介绍性文本，例如 Ferziger (1998) 或 Press 等人。 （2007）

如果每次迭代都从域的左下（西南）角开始，并且我们再次使用地理符号，则 SOR 方法可以编写为：

\begin{sourceblock}
\includegraphics[page=135,trim={53.00bp 581.27bp 53.87bp 53.84bp},clip,width=0.92\textwidth,max height=0.38\textheight]{Ferziger4th.pdf}
\end{sourceblock}
其中 ω 是过度松弛因子，必须大于 1 才能加速，n 是迭代计数器。对于矩形域中的拉普拉斯方程等简单问题，有理论指导最佳过松弛因子的选择，但很难将该理论应用于更复杂的问题；幸运的是，该方法的行为通常与简单情况中的行为类似。一般来说，网格点数越多，最佳过松弛因子就越大（见5.7节）。通常，1.6≤ω≤1.9 会工作得很好；对于 ω = 2.0 方案存在分歧。对于小于最佳值的 ω 值，收敛是单调的，并且收敛速度随着 ω 的增加而增加。当超过最佳 ω 时，收敛速度会变差并且收敛会振荡。该知识可用于搜索最佳的过松弛因子。当采用最佳过松弛因子时，迭代次数与一个方向上的网格点数量成正比，比上述方法有了实质性的改进。当 ω = 1 时，SOR 简化为 Gauss-Seidel 方法。

\subsection*{5.3.4 不完全LU分解：Stone法}\phantomsection\addcontentsline{toc}{subsection}{5.3.4 不完全LU分解：Stone法}
我们提出两个观察结果。 LU 分解是一种优秀的通用线性系统求解器，但它无法利用矩阵的稀疏性。在迭代方法中，如果 M 是 A 的良好近似，则可以快速收敛。这些观察结果引发了使用 A 的近似 LU 分解作为迭代矩阵 M 的想法，即：

\begin{sourceblock}
\includegraphics[page=135,trim={53.00bp 262.14bp 53.87bp 389.67bp},clip,width=0.92\textwidth,max height=0.38\textheight]{Ferziger4th.pdf}
\end{sourceblock}
其中 L 和 U 都很稀疏，N 很小。

这种对称矩阵方法的一个版本称为不完全 Cholesky 分解，通常与共轭梯度方法结合使用。由于离散对流扩散问题或纳维-斯托克斯方程产生的矩阵不是对称的，因此该方法不能应用于它们。这种方法的非对称版本，称为不完全 LU 分解或 ILU，是可能的，但尚未得到广泛使用。在 ILU 方法中，与 LU 分解一样进行，但是对于原始矩阵 A 中为零的每个元素，L 或 U 的相应元素被设置为零。这种因式分解并不精确，但是这些因式分解的乘积可以用作迭代方法的矩阵M。这种方法收敛得相当慢。

Stone (1968) 提出了另一种不完全下上分解方法，该方法已在 CFD 中得到应用。该方法也称为强隐式过程 (SIP)，专门针对偏微分方程离散化的代数方程而设计，不适用于一般方程组。

我们将描述五点计算分子的SIP方法，即具有图3.6所示结构的矩阵。相同的原理可用于构建 7 点（3D 中）和 9 点（2D 非正交网格）计算分子的求解器。

与 ILU 中一样，L 和 U 矩阵仅在 A 具有非零元素的对角线上具有非零元素。具有这些结构的下三角矩阵和上三角矩阵的乘积具有比 A 更多的非零对角线。对于标准五点分子，还有两条对角线（对应于节点 NW 和 SE 或 NE 和 SW，具体取决于向量中节点的顺序），对于 3D 中的七点分子，还有 6 条对角线。对于本书中用于二维问题的节点排序，额外的两条对角线对应于节点NW和SE（网格索引（i，j）和一维存储位置索引l的对应关系见表3.2）。

为了使这些矩阵唯一，U 主对角线上的每个元素都设置为 1。因此需要确定五组元素（三组在 L 中，两组在 U 中）。对于图 5.1 所示形式的矩阵，矩阵乘法规则给出 L 和 U 的乘积的元素，M = LU：

\begin{sourceblock}
\includegraphics[page=136,trim={53.00bp 124.57bp 53.87bp 287.32bp},clip,width=\textwidth,max height=0.38\textheight]{Ferziger4th.pdf}
\par\vspace{0.45\baselineskip}
\small 图 5.1 矩阵 L 和 U 以及乘积矩阵 M 的示意图；对角线
\end{sourceblock}
A 中未找到的 M 用虚线表示
我们希望选择 L 和 U，使得 M 尽可能好地逼近 A。至少， N 包含 M 的两个非零对角线，它们对应于 A 的零对角线；参见等式。 （5.36）。一个明显的选择是让 N 仅在这两条对角线上具有非零元素，并强制 M 的其他对角线等于 A 的相应对角线。这是可能的；事实上，这就是前面提到的标准ILU方法。不幸的是，这种方法收敛速度很慢。

Stone (1968) 认识到，通过允许 N 在与 LU 的所有七个非零对角线相对应的对角线上具有非零元素，可以提高收敛性。通过考虑向量 M φ 最容易推导该方法：

\begin{sourceblock}
\includegraphics[page=137,trim={53.00bp 459.39bp 53.87bp 176.35bp},clip,width=0.92\textwidth,max height=0.38\textheight]{Ferziger4th.pdf}
\end{sourceblock}
最后两项是“额外”的。该方程中的每一项对应于 M = LU 的对角线。

矩阵 N 必须包含 M 的两条“额外”对角线，并且我们要选择剩余对角线上的元素，以便 N φ ≈ 0，或者换句话说，

\begin{sourceblock}
\includegraphics[page=137,trim={53.00bp 375.70bp 53.87bp 274.98bp},clip,width=0.92\textwidth,max height=0.38\textheight]{Ferziger4th.pdf}
\end{sourceblock}
这要求上式中两个“额外”项的贡献几乎被其他对角线的贡献抵消。因此，等式。 ( 5.38 ) 应简化为以下表达式：

\begin{sourceblock}
\begin{tikzpicture}
\node[inner sep=0,anchor=south west] (image) at (0,0) {\includegraphics[page=137,trim={53.00bp 278.41bp 53.87bp 344.99bp},clip,width=0.92\textwidth,max height=0.38\textheight]{Ferziger4th.pdf}};
\begin{scope}[x={(image.south east)},y={(image.north west)}]
\fill[white] (0.20968,0.04445) rectangle (0.23898,0.24520);
\fill[white] (0.24346,0.09336) rectangle (0.28484,0.36383);
\fill[white] (0.28752,0.09336) rectangle (0.47871,0.36383);
\fill[white] (0.48150,0.09336) rectangle (0.50962,0.36383);
\fill[white] (0.51227,0.07487) rectangle (0.56929,0.37248);
\fill[white] (0.57383,0.09336) rectangle (0.62189,0.36383);
\fill[white] (0.62463,0.07487) rectangle (0.68051,0.37248);
\fill[white] (0.03528,0.00000) rectangle (0.12830,0.08400);
\fill[white] (0.13483,0.00000) rectangle (0.18211,0.08400);
\fill[white] (0.18863,0.00000) rectangle (0.24334,0.08400);
\fill[white] (0.24986,0.00000) rectangle (0.27465,0.08400);
\fill[white] (0.28114,0.00000) rectangle (0.33835,0.08400);
\fill[white] (0.34484,0.00000) rectangle (0.44445,0.08400);
\fill[white] (0.45101,0.00000) rectangle (0.49242,0.08400);
\fill[white] (0.49892,0.00000) rectangle (0.61853,0.08400);
\fill[white] (0.62505,0.00000) rectangle (0.77961,0.08400);
\fill[white] (0.78614,0.00000) rectangle (0.81922,0.08400);
\fill[white] (0.82571,0.00000) rectangle (0.91371,0.08400);
\fill[white] (0.92021,0.00000) rectangle (1.00000,0.08400);
\end{scope}
\end{tikzpicture}
\end{sourceblock}
SE 是 φ NW 和 φ SE 的近似值。斯通的关键思想是，由于方程近似于椭圆偏微分方程，因此可以预期解是平滑的。既然这样，则 φ ＊
西北和 φ *
SE 可以根据与 A 对角线对应的节点处的 φ 值来近似。 Stone 提出了以下近似值（其他近似值也是可能的；参见 Schneider 和 Zedan (1981)，例如）：

\begin{sourceblock}
\includegraphics[page=137,trim={53.00bp 183.56bp 53.87bp 452.60bp},clip,width=0.92\textwidth,max height=0.38\textheight]{Ferziger4th.pdf}
\end{sourceblock}
如果 α = 1，则这些是二阶精确插值，但斯通发现稳定性需要 α < 1。这些近似基于与偏微分方程的联系，对于一般代数方程没有什么意义。

如果将这些近似值代入方程。 ( 5.39 ) 结果等于式(5.39) (5.38)，我们获得 N 的所有元素作为 M NW 和 M SE 的线性组合。 M 的元素，方程。 ( 5.36 )，现在可以设置为等于 A 和
尼。所得方程不仅足以确定 L 和 U 的所有元素，而且可以从网格的西南角开始按顺序求解：

\begin{sourceblock}
\begin{tikzpicture}
\node[inner sep=0,anchor=south west] (image) at (0,0) {\includegraphics[page=138,trim={53.00bp 487.91bp 53.87bp 53.89bp},clip,width=0.92\textwidth,max height=0.38\textheight]{Ferziger4th.pdf}};
\begin{scope}[x={(image.south east)},y={(image.north west)}]
\fill[white] (0.33621,0.86843) rectangle (0.36541,0.97587);
\fill[white] (0.35522,0.85170) rectangle (0.37982,0.92062);
\fill[white] (0.38520,0.87052) rectangle (0.41338,0.96357);
\fill[white] (0.41988,0.86843) rectangle (0.44899,0.97587);
\fill[white] (0.35215,0.66447) rectangle (0.38135,0.77200);
\fill[white] (0.37116,0.64782) rectangle (0.38764,0.71674);
\fill[white] (0.39266,0.66664) rectangle (0.42084,0.75961);
\fill[white] (0.42731,0.66447) rectangle (0.45642,0.77200);
\fill[white] (0.09441,0.44869) rectangle (0.12364,0.55622);
\fill[white] (0.11341,0.43204) rectangle (0.12989,0.50097);
\fill[white] (0.13525,0.45086) rectangle (0.16343,0.54383);
\fill[white] (0.16989,0.44869) rectangle (0.19898,0.55622);
\fill[white] (0.30770,0.23291) rectangle (0.34298,0.34044);
\fill[white] (0.32947,0.21626) rectangle (0.34944,0.28519);
\fill[white] (0.35477,0.23508) rectangle (0.38295,0.32805);
\fill[white] (0.32256,0.02903) rectangle (0.35780,0.13656);
\fill[white] (0.34430,0.01239) rectangle (0.36192,0.08131);
\fill[white] (0.36725,0.03120) rectangle (0.39543,0.12418);
\end{scope}
\end{tikzpicture}
\end{sourceblock}
系数必须按此顺序计算。对于边界旁边的节点，携带边界节点索引的任何矩阵元素都被理解为零。因此，沿着西边界 ( i = 2)，索引为 l − N j 的元素为零；沿着南边界 ( j = 2)，索引为 l − 1 的元素为零；沿着北边界 ( j = N j − 1)，索引为 l + 1 的元素为零；最后，沿着东边界 ( i = N i − 1)，索引为 l + N j 的元素为零。

我们现在转向借助这种近似因式分解来求解方程组。与残差更新相关的方程为（见方程（5.19））：

\begin{sourceblock}
\includegraphics[page=138,trim={53.00bp 358.71bp 53.87bp 288.87bp},clip,width=0.92\textwidth,max height=0.38\textheight]{Ferziger4th.pdf}
\end{sourceblock}
方程的求解与一般 LU 分解一样。将上式乘以 L − 1 可得出：

\begin{sourceblock}
\includegraphics[page=138,trim={53.00bp 298.93bp 53.87bp 348.64bp},clip,width=0.92\textwidth,max height=0.38\textheight]{Ferziger4th.pdf}
\end{sourceblock}
R n 很容易计算：

\begin{sourceblock}
\begin{tikzpicture}
\node[inner sep=0,anchor=south west] (image) at (0,0) {\includegraphics[page=138,trim={53.00bp 248.80bp 53.87bp 392.69bp},clip,width=0.92\textwidth,max height=0.38\textheight]{Ferziger4th.pdf}};
\begin{scope}[x={(image.south east)},y={(image.north west)}]
\fill[white] (0.26695,0.14645) rectangle (0.29690,0.68884);
\fill[white] (0.30340,0.15740) rectangle (0.33158,0.62637);
\end{scope}
\end{tikzpicture}
\end{sourceblock}
该方程通过按 l 递增的顺序进行求解。当 R 的计算完成后，我们需要求解式（1）。 （5.43）：

\begin{sourceblock}
\includegraphics[page=138,trim={53.00bp 189.37bp 53.87bp 456.36bp},clip,width=0.92\textwidth,max height=0.38\textheight]{Ferziger4th.pdf}
\end{sourceblock}
按索引 l 递减的顺序。

在SIP方法中，矩阵L和U的元素仅需要在第一次迭代之前计算一次。在后续迭代中，我们只需通过求解两个三角系统来计算残差，然后计算 R，最后计算 δ。

Stone 的方法可以推广为针对在二维中应用紧致差分近似时出现的九对角矩阵以及在三维中使用中心差值时出现的七对角矩阵产生有效的求解器。 Leister 和 Peri c (1994)) 给出了 3D（7 点）矢量化版本； Schneider 和 Zedan (1981) 以及 Peri ́c (1987) 描述了 2D 问题的两个 9 点版本。五对角 (2D) 和七对角 (3D) 矩阵的计算机代码可通过互联网获得；详细内容请参见附录。 SIP 对于模型问题的性能将在第 4 节中介绍。 5.8.

与其他方法不同，Stone 方法本身就是一种很好的迭代技术，也是共轭梯度方法（称为预处理器）和多重网格方法（用作平滑器）的良好基础。下面将描述这些方法。

\subsection*{5.3.5 ADI和其他分割方法}\phantomsection\addcontentsline{toc}{subsection}{5.3.5 ADI和其他分割方法}
求解椭圆问题的常用方法是在方程中添加包含一阶时间导数的项，并求解所得的抛物线问题，直到达到稳态。此时，时间导数为零，解满足原椭圆方程。许多求解椭圆方程的迭代方法，包括已经描述的大多数方法，都可以用这种方式解释。在本节中，我们提出一种方法，它与抛物线方程的联系非常密切，如果只考虑椭圆方程，可能不会发现它。

考虑到稳定性，要求抛物线方程的方法在时间上是隐式的。在二维或三维中，这需要在每个时间步求解 2D 或 3D 椭圆问题；成本可能会很大，但通过使用交替方向隐式或 ADI 方法可以大大降低。我们只给出最简单的二维方法及其变体。 ADI 是许多其他方法的基础；有关其中一些方法的更多详细信息，请参见 Hageman 和 Young (2004)。

假设我们要求解二维拉普拉斯方程。添加时间导数可将其转换为二维热方程：

\begin{sourceblock}
\begin{tikzpicture}
\node[inner sep=0,anchor=south west] (image) at (0,0) {\includegraphics[page=139,trim={53.00bp 187.10bp 53.87bp 443.75bp},clip,width=0.92\textwidth,max height=0.38\textheight]{Ferziger4th.pdf}};
\begin{scope}[x={(image.south east)},y={(image.north west)}]
\fill[white] (0.34433,0.44007) rectangle (0.38256,0.76764);
\end{scope}
\end{tikzpicture}
\end{sourceblock}
如果使用时间梯形规则对该方程进行离散化（应用于偏微分方程时称为 Crank-Nicolson 方法；请参阅下一章），并使用中心差来近似均匀网格上的空间导数，我们得到：

\begin{sourceblock}
\begin{tikzpicture}
\node[inner sep=0,anchor=south west] (image) at (0,0) {\includegraphics[page=139,trim={53.00bp 91.89bp 53.87bp 538.95bp},clip,width=0.92\textwidth,max height=0.38\textheight]{Ferziger4th.pdf}};
\begin{scope}[x={(image.south east)},y={(image.north west)}]
\fill[white] (0.08162,0.43994) rectangle (0.14346,0.80227);
\fill[white] (0.14686,0.43994) rectangle (0.21095,0.79688);
\fill[white] (0.22051,0.24193) rectangle (0.24869,0.56969);
\fill[white] (0.83576,0.24193) rectangle (0.85050,0.56969);
\fill[white] (0.92439,0.23173) rectangle (1.00000,0.55921);
\fill[white] (0.15164,0.03399) rectangle (0.16478,0.36147);
\end{scope}
\end{tikzpicture}
\end{sourceblock}
我们使用简写符号：

\begin{sourceblock}
\begin{tikzpicture}
\node[inner sep=0,anchor=south west] (image) at (0,0) {\includegraphics[page=140,trim={53.00bp 541.34bp 53.87bp 50.00bp},clip,width=0.92\textwidth,max height=0.38\textheight]{Ferziger4th.pdf}};
\begin{scope}[x={(image.south east)},y={(image.north west)}]
\fill[white] (0.30111,0.76163) rectangle (0.34938,0.92914);
\fill[white] (0.30229,0.23503) rectangle (0.35056,0.40254);
\end{scope}
\end{tikzpicture}
\end{sourceblock}
为空间有限差分。重新排列方程。 ( 5.47 )，我们发现，在时间步 n + 1 处，我们必须求解方程组：

\begin{sourceblock}
\begin{tikzpicture}
\node[inner sep=0,anchor=south west] (image) at (0,0) {\includegraphics[page=140,trim={53.00bp 394.77bp 53.87bp 166.75bp},clip,width=0.92\textwidth,max height=0.38\textheight]{Ferziger4th.pdf}};
\begin{scope}[x={(image.south east)},y={(image.north west)}]
\fill[white] (0.36355,0.86618) rectangle (0.37669,0.97677);
\fill[white] (0.39014,0.86876) rectangle (0.41952,0.98853);
\fill[white] (0.26376,0.79851) rectangle (0.28355,0.90900);
\fill[white] (0.28538,0.80205) rectangle (0.31356,0.91254);
\fill[white] (0.33922,0.73084) rectangle (0.35901,0.84143);
\fill[white] (0.38177,0.73179) rectangle (0.42791,0.84716);
\fill[white] (0.38361,0.50650) rectangle (0.39675,0.61709);
\fill[white] (0.41023,0.50908) rectangle (0.43958,0.62875);
\fill[white] (0.28385,0.43883) rectangle (0.30364,0.54932);
\fill[white] (0.30547,0.44236) rectangle (0.33365,0.55286);
\fill[white] (0.35931,0.37115) rectangle (0.37910,0.48174);
\fill[white] (0.40183,0.37220) rectangle (0.44794,0.48748);
\fill[white] (0.30454,0.14930) rectangle (0.32054,0.25989);
\fill[white] (0.36024,0.14682) rectangle (0.39777,0.26907);
\end{scope}
\end{tikzpicture}
\end{sourceblock}
由于 φ n + 1 − φ n ≈ t ∂φ/∂ t，对于较小的 t，最后一项与 (t) 3 成比例。由于 FD 近似是二阶的，因此对于较小的 t，最后一项与离散化误差相比很小，可以忽略不计。剩下的方程可以分解为两个更简单的方程：

\begin{sourceblock}
\begin{tikzpicture}
\node[inner sep=0,anchor=south west] (image) at (0,0) {\includegraphics[page=140,trim={53.00bp 260.29bp 53.87bp 338.86bp},clip,width=0.92\textwidth,max height=0.38\textheight]{Ferziger4th.pdf}};
\begin{scope}[x={(image.south east)},y={(image.north west)}]
\fill[white] (0.35056,0.79116) rectangle (0.36370,0.96373);
\fill[white] (0.37717,0.79504) rectangle (0.40653,0.98209);
\fill[white] (0.25080,0.68533) rectangle (0.27059,0.85789);
\fill[white] (0.27242,0.69085) rectangle (0.30060,0.86341);
\fill[white] (0.32626,0.57979) rectangle (0.34605,0.75235);
\fill[white] (0.36878,0.58128) rectangle (0.41492,0.76131);
\fill[white] (0.33802,0.22929) rectangle (0.35116,0.40185);
\fill[white] (0.36400,0.23317) rectangle (0.39335,0.42036);
\fill[white] (0.23820,0.12345) rectangle (0.25798,0.29616);
\fill[white] (0.25982,0.12897) rectangle (0.28800,0.30154);
\fill[white] (0.31365,0.01791) rectangle (0.33347,0.19048);
\fill[white] (0.35621,0.01941) rectangle (0.40114,0.19943);
\end{scope}
\end{tikzpicture}
\end{sourceblock}
每个方程组都是一组三对角方程，可以使用高效的 TDMA 方法求解；这不需要迭代，并且比求解方程便宜得多。 （5.47）。任一方程。 ( 5.49 ) 或 ( 5.50 ) 本身作为一种方法，仅在时间上是一阶精确且有条件稳定的，但组合方法是二阶精确且无条件稳定的！基于这些思想的方法族被称为分裂或近似因式分解方法；已经开发出各种各样的产品。

仅当时间步长较小时，忽略对因式分解至关重要的三阶项才是合理的。因此，虽然该方法无条件稳定，但如果时间步长较大，则可能不及时准确。对于椭圆方程，目标是尽快获得稳态解；这最好通过尽可能大的时间步长来完成。然而，当时间步长较大时，因式分解误差会变大，因此该方法失去了部分有效性。事实上，存在一个最佳时间步长，可以实现最快的收敛。当使用这一时间步长时，ADI 方法非常高效——它在与一个方向上的点数成正比的迭代次数中收敛。

更好的策略以循环方式使用不同的时间步进行多次迭代。这种方法可以使收敛的迭代次数与一个方向上网格点数的平方根成正比，使ADI成为一种优秀的方法。

涉及对流项和源项的方程需要对此方法进行一些推广。在 CFD 中，压力或压力校正方程属于上述类型，并且经常使用刚刚描述的方法的变体来求解它。 ADI 方法在解决可压缩流动问题时非常常用。它们也非常适合并行计算。

本节中描述的方法利用了矩阵的结构，而矩阵的结构又归因于结构化网格的使用。然而，对开发的仔细检查表明该方法的基础是矩阵的加法分解：

\begin{sourceblock}
\begin{tikzpicture}
\node[inner sep=0,anchor=south west] (image) at (0,0) {\includegraphics[page=141,trim={53.00bp 415.69bp 53.87bp 236.12bp},clip,width=0.92\textwidth,max height=0.38\textheight]{Ferziger4th.pdf}};
\begin{scope}[x={(image.south east)},y={(image.north west)}]
\fill[white] (0.00000,0.91766) rectangle (0.04075,1.00000);
\fill[white] (0.04343,0.91766) rectangle (0.13477,1.00000);
\end{scope}
\end{tikzpicture}
\end{sourceblock}
其中 H 是表示二阶导数对 x 和 V 贡献的项的矩阵，这些项来自 y 方向的二阶导数。

没有理由不能使用其他加法分解。一项有用的建议是考虑加法 LU 分解：

\begin{sourceblock}
\includegraphics[page=141,trim={53.00bp 320.05bp 53.87bp 331.76bp},clip,width=0.92\textwidth,max height=0.38\textheight]{Ferziger4th.pdf}
\end{sourceblock}
这与节的乘法LU分解不同。 5.2.2.通过这种分解，方程。 ( 5.49 ) 和 ( 5.50 ) 替换为：

\begin{sourceblock}
\includegraphics[page=141,trim={53.00bp 248.84bp 53.87bp 387.50bp},clip,width=0.92\textwidth,max height=0.38\textheight]{Ferziger4th.pdf}
\end{sourceblock}
每个步骤本质上都是高斯-赛德尔迭代。该方法的收敛速度与上面给出的ADI方法的收敛速度相似。它还具有非常重要的优点，即它不依赖于网格或矩阵的结构，因此可以应用于非结构化网格和结构化网格上的问题。然而，它的并行性不如 ADI 的 HV 版本。

\subsection*{5.3.6 克雷洛夫方法}\phantomsection\addcontentsline{toc}{subsection}{5.3.6 克雷洛夫方法}
在本节中，我们将介绍基于求解非线性方程的技术的方法。非线性求解器可以分为两大类：类牛顿方法（第 5.5.1 节）和全局方法。如果可以精确估计解，则前者收敛得非常快，但如果初始猜测与精确解相差甚远，则可能会灾难性地失败。 “远”是一个相对词；每个方程都是不同的。除非通过反复试验，否则无法确定估计值是否“足够接近”。全局方法保证能找到解（如果存在），但速度不是很快。经常使用两种方法的组合；最初使用全局方法，然后在接近收敛时使用类牛顿方法。

本节方法的目的本质上是通过创建大矩阵的近似值然后在迭代过程中使用这些近似值来创建将大型线性方程组简化为较小问题的方案；这称为投影方法，因为我们将问题从 N × N 投影到更小的维度。 van der Vorst (2002) 为这种方法提供了如下背景。给定我们的方程组 ( 5.1 ) 和矩阵 A，第 k 步的迭代方案如下所示

\begin{sourceblock}
\includegraphics[page=142,trim={53.00bp 425.66bp 53.87bp 224.17bp},clip,width=0.92\textwidth,max height=0.38\textheight]{Ferziger4th.pdf}
\end{sourceblock}
yields 2

\begin{sourceblock}
\begin{tikzpicture}
\node[inner sep=0,anchor=south west] (image) at (0,0) {\includegraphics[page=142,trim={53.00bp 389.80bp 53.87bp 258.31bp},clip,width=0.92\textwidth,max height=0.38\textheight]{Ferziger4th.pdf}};
\begin{scope}[x={(image.south east)},y={(image.north west)}]
\fill[white] (0.00000,0.83916) rectangle (0.08620,1.00000);
\end{scope}
\end{tikzpicture}
\end{sourceblock}
其中，如果 φ 0 是初始猜测解，则初始残差为 ρ 0 = Q − A φ 0。可以看出迭代过程在所谓的移位 Krylov 子空间 K k(A;ρ 0) 中生成近似解。这种特殊的方案（理查森迭代）虽然简单，但既不高效也不是最优的。然后，克雷洛夫子空间投影方法尝试通过构建更好的近似解来改进它。

van der Vorst (2002) 简洁地描述了三类实现此目的的方法：

1. Ritz-Galerkin 方法 (R-G)：构造 φ k，其残差与当前子空间 Q − A φ k ⊥ K k ( A ; ρ 0 ) 正交。 2. 最小残差方法：确定 φ k，其中欧几里得范数 ∥ Q − A φ k ∥ 2 在 K k ( A ; ρ 0 ) 上最小。 3. Petrov-Galerkin 方法 (P-G)：找到 φ k，使得残差 Q − A φ k 与其他合适的 k 维空间正交。

下面介绍的方法来自这些类；作为参考，另见 Saad (2003)。共轭梯度法 (CG) 是 R-G 方法，双共轭梯度和 CGSTAB 是 P-G 方法，GMRES 是最小残差方法。

\begin{sourceblock}
\includegraphics[page=142,trim={53.00bp 67.64bp 53.87bp 580.94bp},clip,width=0.92\textwidth,max height=0.38\textheight]{Ferziger4th.pdf}
\end{sourceblock}
5.3.6.1 共轭梯度法
许多全局方法都是下降方法。 3 这些方法首先将原始方程组转换为最小化问题。假设要求解的方程组由方程给出。 ( 5.1 ) 并且矩阵A是对称的且其特征值为正；这样的矩阵称为正定矩阵。 （大多数与流体力学问题相关的矩阵都不是对称的或正定的，因此我们稍后需要推广此方法。）对于正定矩阵，求解方程组 ( 5.1 ) 相当于找到以下最小值的问题

\begin{sourceblock}
\includegraphics[page=143,trim={53.00bp 456.20bp 53.87bp 174.41bp},clip,width=0.92\textwidth,max height=0.38\textheight]{Ferziger4th.pdf}
\end{sourceblock}
相对于所有 φ i；这可以通过对每个变量求 F 的导数并将其设置为零来验证。将原始系统转换为不需要正定性的最小化问题的一种方法是取所有方程的平方和，但这会带来额外的困难。

寻找函数最小值的最古老且最著名的方法是最速下降法。函数 F 可以被认为是（超）空间中的表面。假设我们有一些起始猜测，可以将其表示为该（超）空间中的一个点。在那一点上，我们找到了从表面开始最陡的向下路径；它位于与函数梯度相反的方向。然后我们搜索该线上的最低点。通过构造，它的F值比起始点低；从这个意义上说，新的估计更接近解。然后将新值用作下一次迭代的起点，并继续该过程直到收敛。不幸的是，虽然保证收敛，但最速下降法通常收敛得非常慢。

如果函数 F 的幅度的等高线图有一个狭窄的谷，则该方法往往会在该谷上来回振荡，并且可能需要许多步骤才能找到解。换句话说，该方法倾向于一遍又一遍地使用相同的搜索方向。

已提出许多改进建议。最简单的要求新的搜索方向与旧的搜索方向尽可能不同。其中有共轭梯度法。这里我们只给出算法的总体思路和描述；更完整的演示可以在 Shewchuk (1994)、Watkins (2010) 或 Golub 和 van Loan (1996) 中找到。

共轭梯度法基于一项显著的发现：在一次搜索一个方向的同时，可以同时在多个方向上最小化一个函数。这是通过巧妙选择方向而实现的。我们将针对两个方向的情况来描述这一点；假设我们希望找到 α 1 和 α 2 的值

\begin{sourceblock}
\includegraphics[page=144,trim={53.00bp 593.62bp 53.87bp 52.66bp},clip,width=0.92\textwidth,max height=0.38\textheight]{Ferziger4th.pdf}
\end{sourceblock}
最小化 F；也就是说，我们尝试最小化 p 1 − p 2 平面中的 F。这个问题可以简化为关于 p 1 和 p 2 分别最小化的问题，前提是这两个方向在以下意义上是共轭的：

\begin{sourceblock}
\includegraphics[page=144,trim={53.00bp 523.29bp 53.87bp 124.41bp},clip,width=0.92\textwidth,max height=0.38\textheight]{Ferziger4th.pdf}
\end{sourceblock}
该属性类似于正交性；向量 p 1 和 p 2 被认为与矩阵 A 共轭，该方法因此得名。该陈述和下面引用的其他陈述的详细证明可以在 Golub 和 van Loan (1996) 的书中找到。

该属性可以扩展到任意多个方向。在共轭梯度法中，每个新的搜索方向都需要与之前的所有搜索方向共轭。如果矩阵是非奇异的，就像几乎所有工程问题的情况一样，则方向保证是线性无关的。因此，如果采用精确（无舍入误差）算法，则当迭代次数等于矩阵大小时，共轭梯度法将精确收敛。这个数字可能相当大，并且在实践中，由于算术错误而无法实现精确收敛。因此，将共轭梯度法视为迭代方法更为明智。

虽然共轭梯度法保证每次迭代时误差都会减少，但减少的大小取决于搜索方向。对于这种方法来说，在多次迭代中仅稍微减少误差，然后找到在一次迭代中将误差减少一个数量级或更多的方向并不罕见。事实上，该方法（实际上也是一般迭代方法）的收敛速度取决于系数矩阵 A（方程（5.1））。未知数 φ 相对于方程右侧 Q 的扰动的变化可以显示 (Watkins 2010) 与

\begin{sourceblock}
\includegraphics[page=144,trim={53.00bp 236.36bp 53.87bp 411.34bp},clip,width=0.92\textwidth,max height=0.38\textheight]{Ferziger4th.pdf}
\end{sourceblock}
其中 ∥ A ∥ 表示 A 的范数。 Saad (2003) 表明该分析可以扩展到迭代过程，如下所示。首先定义残差范数，

\begin{sourceblock}
\includegraphics[page=144,trim={53.00bp 176.69bp 53.87bp 471.09bp},clip,width=0.92\textwidth,max height=0.38\textheight]{Ferziger4th.pdf}
\end{sourceblock}
其中 φ k 是 k 次迭代后解的估计。然后，它将遵循等式[经过一些工作！] ( 5.58 ) 迭代误差范数之比 ε k = φ − φ k
解向量 φ 的范数与残差范数与 Q 范数的比值有关：

\begin{sourceblock}
\begin{tikzpicture}
\node[inner sep=0,anchor=south west] (image) at (0,0) {\includegraphics[page=144,trim={53.00bp 91.81bp 53.87bp 542.42bp},clip,width=0.92\textwidth,max height=0.38\textheight]{Ferziger4th.pdf}};
\begin{scope}[x={(image.south east)},y={(image.north west)}]
\fill[white] (0.00000,0.82106) rectangle (0.06737,1.00000);
\fill[white] (0.07008,0.82106) rectangle (0.09982,1.00000);
\fill[white] (0.10250,0.82419) rectangle (0.13065,1.00000);
\fill[white] (0.13332,0.82106) rectangle (0.16310,1.00000);
\fill[white] (0.16577,0.82106) rectangle (0.26803,1.00000);
\end{scope}
\end{tikzpicture}
\end{sourceblock}
它给出了迭代误差的上限。显然，κ 的大小对迭代有显著影响。现在，从基本线性代数（Golub 和 van Loan 1996）中，我们回想起矩阵 A 的特征值是允许非零解的值

\begin{sourceblock}
\begin{tikzpicture}
\node[inner sep=0,anchor=south west] (image) at (0,0) {\includegraphics[page=145,trim={53.00bp 547.29bp 53.87bp 104.62bp},clip,width=0.92\textwidth,max height=0.38\textheight]{Ferziger4th.pdf}};
\begin{scope}[x={(image.south east)},y={(image.north west)}]
\fill[white] (0.00000,0.91778) rectangle (0.11062,1.00000);
\fill[white] (0.11332,0.91778) rectangle (0.22638,1.00000);
\fill[white] (0.22905,0.91778) rectangle (0.25717,1.00000);
\end{scope}
\end{tikzpicture}
\end{sourceblock}
其中 x 是 A 的特征向量。从这个结果和方程。 ( 5.58 ) 可以证明

\begin{sourceblock}
\includegraphics[page=145,trim={53.00bp 486.75bp 53.87bp 150.21bp},clip,width=0.92\textwidth,max height=0.38\textheight]{Ferziger4th.pdf}
\end{sourceblock}
当 λ max 和 λ min 是矩阵 A 的最大和最小特征值时。因此，收敛速度取决于所谓的条件数 κ，它只是系数矩阵的属性。探索条件数的一个有用方法是将 MATLAB TM 中的“cond”函数应用于各种代表性矩阵。

CFD 问题中出现的矩阵条件数通常近似于任意方向上最大网格点数的平方。每个方向有100个网格点，条件数应约为10 4，标准共轭梯度法收敛速度较慢。尽管对于给定的条件数，共轭梯度法明显快于最速下降法，但这种基本方法并不是很有用。

可以通过将我们寻求解的问题替换为另一个具有相同解但条件数更小的问题来改进该方法。这是通过预处理来完成的。预处理问题的一种方法是将方程预先乘以另一个（精心选择的）矩阵。由于这会破坏矩阵的对称性，因此预处理必须采用以下形式：

\begin{sourceblock}
\includegraphics[page=145,trim={53.00bp 275.14bp 53.87bp 374.95bp},clip,width=0.92\textwidth,max height=0.38\textheight]{Ferziger4th.pdf}
\end{sourceblock}
将共轭梯度法应用于矩阵 C − 1 AC − 1，即修改问题（5.61）。如果这样做并使用迭代方法的残差形式，则将得到以下算法结果（有关详细推导，请参阅 Golub 和 van Loan 1996 或 Shewchuk 1994 ）。在本描述中，ρ k 是第 k 次迭代的残差，p k 是第 k 个搜索方向，z k 是辅助向量，α k 和 β k 是用于构造新解、残差和搜索方向的参数。该算法可以总结如下：

\begin{itemize}
\item 通过设置进行初始化： k = 0, φ 0 = φ in , ρ 0 = Q − A φ in , p 0 = 0 , s 0 = 10 30
\end{itemize}
\begin{sourceblock}
\begin{tikzpicture}
\node[inner sep=0,anchor=south west] (image) at (0,0) {\includegraphics[page=145,trim={53.00bp 89.93bp 53.87bp 502.20bp},clip,width=0.92\textwidth,max height=0.38\textheight]{Ferziger4th.pdf}};
\begin{scope}[x={(image.south east)},y={(image.north west)}]
\fill[white] (0.00000,0.98906) rectangle (0.01913,1.00000);
\fill[white] (0.02779,0.98406) rectangle (0.13910,1.00000);
\fill[white] (0.14174,0.98406) rectangle (0.17651,1.00000);
\fill[white] (0.17922,0.98406) rectangle (0.27227,1.00000);
\fill[white] (0.27501,0.98541) rectangle (0.29311,1.00000);
\fill[white] (0.29841,0.98906) rectangle (0.32659,1.00000);
\fill[white] (0.33011,0.98406) rectangle (0.35756,1.00000);
\fill[white] (0.36027,0.98906) rectangle (0.39477,1.00000);
\fill[white] (0.39991,0.98906) rectangle (0.42809,1.00000);
\fill[white] (0.43161,0.96203) rectangle (0.48129,1.00000);
\fill[white] (0.48394,0.98906) rectangle (0.51738,1.00000);
\fill[white] (0.52253,0.98906) rectangle (0.55071,1.00000);
\fill[white] (0.55423,0.98541) rectangle (0.58235,1.00000);
\fill[white] (0.58418,0.96203) rectangle (0.68574,1.00000);
\fill[white] (0.68842,0.98541) rectangle (0.72039,1.00000);
\fill[white] (0.72553,0.98906) rectangle (0.75371,1.00000);
\fill[white] (0.75723,0.98406) rectangle (0.78469,1.00000);
\fill[white] (0.78734,0.98541) rectangle (0.81609,1.00000);
\fill[white] (0.82123,0.98906) rectangle (0.84941,1.00000);
\fill[white] (0.85293,0.98406) rectangle (0.90884,1.00000);
\end{scope}
\end{tikzpicture}
\end{sourceblock}
\begin{sourceblock}
\begin{tikzpicture}
\node[inner sep=0,anchor=south west] (image) at (0,0) {\includegraphics[page=146,trim={53.00bp 583.16bp 53.87bp 52.66bp},clip,width=0.92\textwidth,max height=0.38\textheight]{Ferziger4th.pdf}};
\begin{scope}[x={(image.south east)},y={(image.north west)}]
\fill[white] (0.00000,0.00000) rectangle (0.01913,0.03529);
\end{scope}
\end{tikzpicture}
\end{sourceblock}
\begin{itemize}
\item 重复直到收敛。
\end{itemize}
该算法涉及第一步求解线性方程组。涉及的矩阵是 M = C − 1，其中 C 是预处理矩阵，实际上从未实际构造过。为了使该方法有效，M 必须易于求逆。最常用的 M 选择是 A 的不完全 Cholesky 分解，但在测试中发现，如果 M = LU （其中 L 和 U 是 Stone 的 SIP 方法中使用的因子），可以获得更快的收敛。下面将给出示例。 Saad (2003) 对串行和并行计算的预处理器进行了广泛的讨论。

5.3.6.2 双共轭梯度和 CGSTAB
上面给出的共轭梯度法仅适用于对称系统；将泊松方程离散化得到的矩阵通常是对称的（例如第7章将介绍的热传导方程和压力或压力校正方程）。为了将该方法应用于不一定对称的方程组（例如任何对流扩散方程），我们需要将不对称问题转换为对称问题。有几种方法可以做到这一点，其中以下可能是最简单的。考虑系统：0 A A T 0

\begin{sourceblock}
\begin{tikzpicture}
\node[inner sep=0,anchor=south west] (image) at (0,0) {\includegraphics[page=146,trim={53.00bp 301.71bp 53.87bp 329.40bp},clip,width=0.92\textwidth,max height=0.38\textheight]{Ferziger4th.pdf}};
\begin{scope}[x={(image.south east)},y={(image.north west)}]
\fill[white] (0.00000,0.72823) rectangle (0.09570,1.00000);
\fill[white] (0.35777,0.42506) rectangle (0.37756,0.75535);
\fill[white] (0.41335,0.42792) rectangle (0.43648,0.75821);
\fill[white] (0.49901,0.43563) rectangle (0.52478,0.76591);
\fill[white] (0.46379,0.23637) rectangle (0.47693,0.56637);
\fill[white] (0.65501,0.23637) rectangle (0.66806,0.56637);
\fill[white] (0.92424,0.22581) rectangle (1.00000,0.55581);
\fill[white] (0.35609,0.02683) rectangle (0.39194,0.37311);
\fill[white] (0.42174,0.02398) rectangle (0.44153,0.35398);
\end{scope}
\end{tikzpicture}
\end{sourceblock}
该系统可以分解为两个子系统。首先是原来的系统；第二个涉及转置矩阵并且是不相关的。 （如果需要这样做，可以以很少的额外成本求解涉及转置矩阵的方程组。）当将预条件共轭梯度法应用于该系统时，将产生以下称为双共轭梯度的方法：

\begin{itemize}
\item 通过设置进行初始化： k = 0, φ 0 = φ in , ρ 0 = Q − A φ in , ρ 0 = Q − A T φ in , p 0 = p 0 = 0 , s 0 = 10 30
\end{itemize}
\begin{sourceblock}
\begin{tikzpicture}
\node[inner sep=0,anchor=south west] (image) at (0,0) {\includegraphics[page=146,trim={53.00bp 82.25bp 53.87bp 462.06bp},clip,width=0.92\textwidth,max height=0.38\textheight]{Ferziger4th.pdf}};
\begin{scope}[x={(image.south east)},y={(image.north west)}]
\fill[white] (0.02779,0.99114) rectangle (0.05976,1.00000);
\fill[white] (0.06490,0.99335) rectangle (0.09308,1.00000);
\fill[white] (0.09657,0.99031) rectangle (0.12406,1.00000);
\fill[white] (0.12671,0.99114) rectangle (0.15546,1.00000);
\fill[white] (0.16060,0.99335) rectangle (0.18878,1.00000);
\fill[white] (0.19230,0.99031) rectangle (0.24821,1.00000);
\fill[white] (0.00000,0.00000) rectangle (0.01913,0.00878);
\end{scope}
\end{tikzpicture}
\end{sourceblock}
\begin{itemize}
\item 重复直到收敛。
\end{itemize}
上述算法由Fletcher (1976)发表。它每次迭代所需的工作量几乎是标准共轭梯度法的两倍，但收敛的迭代次数大约相同。它尚未在 CFD 应用中广泛使用，但它似乎非常强大（这意味着它可以毫无困难地处理各种问题）。

已经开发出更稳定、更稳健的双共轭梯度法类型的其他变体。这里我们提到CGS（共轭梯度平方）算法，由Sonneveld（1989）提出； CGSTAB（CGS 稳定），由 van der Vorst 和 Sonneveld (1990) 提出，另一个版本由 van der Vorst (1992) 提出；和 GMRES（见下文第 5.3.6.3 节）。所有这些都可以应用于非对称矩阵以及结构化和非结构化网格。下面我们给出没有正式推导的 CGSTAB 算法：

\begin{itemize}
\item 通过设置进行初始化： k = 0, φ 0 = φ in , ρ 0 = Q − A φ in , u 0 = p 0 = 0
\item 使计数器前进 k = k + 1 并计算： β k = ρ 0
\item ρ k − 1
\end{itemize}
\begin{sourceblock}
\begin{tikzpicture}
\node[inner sep=0,anchor=south west] (image) at (0,0) {\includegraphics[page=147,trim={53.00bp 350.04bp 53.87bp 238.48bp},clip,width=0.92\textwidth,max height=0.38\textheight]{Ferziger4th.pdf}};
\begin{scope}[x={(image.south east)},y={(image.north west)}]
\fill[white] (0.02779,0.94306) rectangle (0.06087,1.00000);
\fill[white] (0.06719,0.94306) rectangle (0.09537,1.00000);
\fill[white] (0.09886,0.94306) rectangle (0.13230,1.00000);
\fill[white] (0.13579,0.94306) rectangle (0.14893,1.00000);
\fill[white] (0.15074,0.94306) rectangle (0.21110,1.00000);
\end{scope}
\end{tikzpicture}
\end{sourceblock}
\begin{sourceblock}
\includegraphics[page=147,trim={53.00bp 266.25bp 53.87bp 325.78bp},clip,width=0.92\textwidth,max height=0.38\textheight]{Ferziger4th.pdf}
\end{sourceblock}
请注意， u 、 v 、 w 、 y 和 z 是辅助向量，与速度向量或坐标 y 和 z 无关。该算法可以按照上面给出的方式进行编程；具有不完全 Cholesky 预处理的共轭梯度法（ICCG，用于对称矩阵，2D 和 3D 版本）和 3D CGSTAB 求解器的计算机代码可通过互联网获得；详情请参见附录。

5.3.6.3 广义最小残差法
Saad 和 Schultz (1986) 提出的广义最小残差 (GMRES) 方法处理非对称 A。 Saad (2003) 详细描述了该方法。 GMRES 有缺点，但因其强大而变得流行。两个问题是：

1. 迭代使用所有先前的搜索方向向量来计算下一个搜索方向，因此存储空间和操作次数线性增长，如果 A 很大，这可能是一个严重的问题。

2. 当矩阵不是正定时，迭代可能会停止。

通过 (1) 在预定数量的步骤后重新启动迭代和 (2) 预处理矩阵（参见第 5.3.6.1 节）以减少收敛所需的迭代次数来实现改进。在一系列论文中（例如，参见 Armfield 和 Street 2004），我们成功地使用了重新启动的预处理 GMRES 方法，发现与其他求解器（包括 CG）相比，它对于求解泊松和压力校正方程（参见第 7.1.5 节）最有效。

求解 A φ = Q 的基本 GMRES 算法为（改编自 Golub 和 van Loan 1996 以及 Saad 2003）：

\begin{itemize}
\item 开始：设置误差容限并选择m 来限制迭代次数。初始化设置： k = 0, φ 0 = φ in , ρ 0 = Q − A φ in , h 10 =∥ ρ 0 ∥ 2
\item 如果 h k + 1 , k > 0 , 计算 β k + 1 = ρ k / k k + 1 , k k = k + 1 ρ k = A β k
\end{itemize}
\begin{sourceblock}
\begin{tikzpicture}
\node[inner sep=0,anchor=south west] (image) at (0,0) {\includegraphics[page=148,trim={53.00bp 342.46bp 53.87bp 248.08bp},clip,width=0.92\textwidth,max height=0.38\textheight]{Ferziger4th.pdf}};
\begin{scope}[x={(image.south east)},y={(image.north west)}]
\fill[white] (0.07444,0.98942) rectangle (0.10671,1.00000);
\fill[white] (0.11296,0.98942) rectangle (0.14114,1.00000);
\fill[white] (0.14764,0.98585) rectangle (0.19874,1.00000);
\fill[white] (0.11191,0.00000) rectangle (0.14165,0.03042);
\fill[white] (0.14433,0.00000) rectangle (0.18574,0.03042);
\fill[white] (0.18845,0.00000) rectangle (0.21774,0.03532);
\fill[white] (0.22138,0.00000) rectangle (0.24956,0.03532);
\fill[white] (0.25140,0.00000) rectangle (0.28238,0.03532);
\fill[white] (0.28421,0.00000) rectangle (0.31251,0.03532);
\fill[white] (0.31438,0.00000) rectangle (0.33248,0.03175);
\fill[white] (0.33696,0.00000) rectangle (0.39669,0.03042);
\fill[white] (0.39937,0.00000) rectangle (0.49402,0.03042);
\fill[white] (0.49672,0.00000) rectangle (0.60968,0.03042);
\end{scope}
\end{tikzpicture}
\end{sourceblock}
( k + 1 ) × k 最小二乘问题的：

∥ h 10 e 1 − ̃ H k y k ∥ 2 = min 如果满足公差，则设置 φ = φ k 并停止 否则，如果 k ≥ m，则设置 φ 0 = φ k 并返回： 开始 返回： 如果 h k + 1，k > 0
上述算法中共有三个辅助矩阵：

1. Q k 的列是正交的 Arnoldi 向量。 2. H k 是包含 hi j 的上 Hessenberg 矩阵。 3. e 1 是单位矩阵 I n 的第一列，所以

\begin{sourceblock}
\includegraphics[page=148,trim={53.00bp 199.37bp 53.87bp 449.78bp},clip,width=0.92\textwidth,max height=0.38\textheight]{Ferziger4th.pdf}
\end{sourceblock}
最后，我们注意到∥·∥ 2 是矩阵范数。

\subsection*{5.3.7 多重网格方法}\phantomsection\addcontentsline{toc}{subsection}{5.3.7 多重网格方法}
这里要讨论的最后一种求解线性系统的方法是多重网格方法；我们将应用节中的流动计算方法。 12.4.多重网格概念的基础是对迭代方法的观察。它们的收敛速度取决于与该方法相关的迭代矩阵的特征值。特别是，具有最大幅度的特征值（矩阵的谱半径）决定了达到解的速度；见第 5.3.2 节.与该特征值关联的特征向量决定了迭代误差的空间分布，并且不同方法的差异很大。让我们简要回顾一下上述一些方法中这些实体的行为。给出了拉普拉斯方程的性质；它们中的大多数都推广到其他椭圆偏微分方程。

对于拉普拉斯方程，雅可比方法的两个最大特征值是实数且符号相反。一个特征向量表示空间坐标的平滑函数，另一个特征向量表示快速振荡函数。因此，雅可比方法的迭代误差是非常平滑和非常粗糙分量的混合；这使得加速变得困难。另一方面，高斯-塞德尔方法具有单个实数正最大特征值，其特征向量使迭代误差成为空间坐标的平滑函数。

最佳过松弛因子的SOR方法的最大特征值位于复平面上的一个圆上，且有多个；因此，错误的表现方式非常复杂。在 ADI 中，错误的性质取决于参数，但往往相当复杂。最后，SIP具有相对平滑的迭代误差。

其中一些方法产生的误差是空间坐标的平滑函数。让我们考虑其中一种方法。第 n 次迭代后的迭代误差 ε n 和残差 ρ n 之间的关系如下： （5.15）。在 Gauss-Seidel 和 SIP 方法中，经过几次迭代后，误差的快速变化分量已被消除，并且误差成为空间坐标的平滑函数。如果误差是平滑的，则可以在较粗糙的网格上计算更新（迭代误差的近似值）。在二维网格上，网格的粗度是原始网格的两倍，迭代成本是原始网格的 1 / 4；在三个维度上，成本是细网格成本的八分之一。此外，迭代方法在较粗糙的网格上收敛得更快。 Gauss-Seidel 在粗略两倍的网格上收敛速度是原来的四倍；对于 SIP 来说，这一比率不太有利，但仍然很大。

这表明大部分工作可以在较粗糙的网格上完成。为此，我们需要定义：两个网格之间的关系、粗网格上的有限差分算子、将残差从细网格平滑（限制）到粗网格的方法以及从粗网格到细网格插值（延长）更新或校正的方法；括号中的单词是多重网格文献中常用的特殊术语。每个项目都有多种选择；它们会影响方法的行为，但在良好的选择范围内，差异并不大。因此，我们将为每一项只提供一个好的选择。

在有限差分方案中，粗网格通常由细网格的每隔一行组成。在有限体积方法中，通常将粗网格CV由2个（二维为4个，三维为8个）细网格CV组成；粗网格节点则位于细网格节点之间。

虽然没有理由在一维中使用多重网格方法（因为 TDMA 算法非常有效），但很容易说明多重网格方法的原理并导出一般情况下使用的一些程序。因此考虑这个问题：

\begin{sourceblock}
\begin{tikzpicture}
\node[inner sep=0,anchor=south west] (image) at (0,0) {\includegraphics[page=150,trim={53.00bp 557.61bp 53.87bp 79.17bp},clip,width=0.92\textwidth,max height=0.38\textheight]{Ferziger4th.pdf}};
\begin{scope}[x={(image.south east)},y={(image.north west)}]
\fill[white] (0.00000,0.90770) rectangle (0.10562,1.00000);
\fill[white] (0.10833,0.90770) rectangle (0.14974,1.00000);
\fill[white] (0.15242,0.90770) rectangle (0.26538,1.00000);
\end{scope}
\end{tikzpicture}
\end{sourceblock}
均匀网格上的标准 FD 近似为：

\begin{sourceblock}
\begin{tikzpicture}
\node[inner sep=0,anchor=south west] (image) at (0,0) {\includegraphics[page=150,trim={53.00bp 502.25bp 53.87bp 135.97bp},clip,width=0.92\textwidth,max height=0.38\textheight]{Ferziger4th.pdf}};
\begin{scope}[x={(image.south east)},y={(image.north west)}]
\fill[white] (0.32668,0.54262) rectangle (0.34647,0.95702);
\fill[white] (0.29654,0.05229) rectangle (0.31251,0.46669);
\fill[white] (0.33260,0.04298) rectangle (0.37504,0.47493);
\end{scope}
\end{tikzpicture}
\end{sourceblock}
在x间距的网格上进行n次迭代后，得到一个近似解 φ n，并且在残差 ρ n 范围内满足上式：

\begin{sourceblock}
\begin{tikzpicture}
\node[inner sep=0,anchor=south west] (image) at (0,0) {\includegraphics[page=150,trim={53.00bp 433.25bp 53.87bp 204.97bp},clip,width=0.92\textwidth,max height=0.38\textheight]{Ferziger4th.pdf}};
\begin{scope}[x={(image.south east)},y={(image.north west)}]
\fill[white] (0.29023,0.54298) rectangle (0.31002,0.95702);
\fill[white] (0.26003,0.05265) rectangle (0.27603,0.46669);
\fill[white] (0.29612,0.04298) rectangle (0.33856,0.47493);
\end{scope}
\end{tikzpicture}
\end{sourceblock}
从方程中减去这个方程。 ( 5.64 ) 给出

\begin{sourceblock}
\begin{tikzpicture}
\node[inner sep=0,anchor=south west] (image) at (0,0) {\includegraphics[page=150,trim={53.00bp 375.99bp 53.87bp 262.22bp},clip,width=0.92\textwidth,max height=0.38\textheight]{Ferziger4th.pdf}};
\begin{scope}[x={(image.south east)},y={(image.north west)}]
\fill[white] (0.32659,0.54279) rectangle (0.34638,0.95704);
\fill[white] (0.29642,0.05263) rectangle (0.31242,0.46652);
\fill[white] (0.33251,0.04296) rectangle (0.37492,0.47512);
\end{scope}
\end{tikzpicture}
\end{sourceblock}
这是方程。 ( 5.15 ) 对于节点 i。这是我们要在粗网格上迭代的方程。

为了推导粗网格上的离散方程，我们注意到粗网格节点 I 周围的控制体积由节点 i 周围的整个控制体积加上细网格控制体积 i − 1 和 i + 1 的一半组成（见图 5.2）。这表明我们将指数为 i − 1 和 i + 1 的方程 ( 5.66 ) 的一半添加到指数为 i 的完整方程中，这导致（省略上标 n ）：

\begin{sourceblock}
\begin{tikzpicture}
\node[inner sep=0,anchor=south west] (image) at (0,0) {\includegraphics[page=150,trim={53.00bp 246.91bp 53.87bp 389.82bp},clip,width=0.92\textwidth,max height=0.38\textheight]{Ferziger4th.pdf}};
\begin{scope}[x={(image.south east)},y={(image.north west)}]
\fill[white] (0.21293,0.51853) rectangle (0.23272,0.91193);
\fill[white] (0.28376,0.25434) rectangle (0.33621,0.68684);
\fill[white] (0.33967,0.25672) rectangle (0.40842,0.68684);
\fill[white] (0.41377,0.29378) rectangle (0.44198,0.68684);
\fill[white] (0.44385,0.25434) rectangle (0.49630,0.68684);
\fill[white] (0.81871,0.29378) rectangle (0.83176,0.68684);
\fill[white] (0.92439,0.28120) rectangle (1.00000,0.67426);
\fill[white] (0.17272,0.04080) rectangle (0.19251,0.43387);
\fill[white] (0.19275,0.05304) rectangle (0.20875,0.44645);
\fill[white] (0.22884,0.04420) rectangle (0.27125,0.45427);
\end{scope}
\end{tikzpicture}
\end{sourceblock}
利用两个网格之间的关系（X = 2 x，见图5.2），这相当于粗网格上的以下方程：

\begin{sourceblock}
\includegraphics[page=150,trim={53.00bp 155.79bp 53.87bp 453.68bp},clip,width=\textwidth,max height=0.38\textheight]{Ferziger4th.pdf}
\par\vspace{0.45\baselineskip}
\small 图 5.2 一维多重网格技术中使用的网格 i i+1 i+2 i+3 i−1 i−2 i−3
\end{sourceblock}
细格子

x

I−1 I+1 I
X

粗网格也用于定义 ρ I。该方程的左边是粗网格上二阶导数的标准近似，表明粗网格上的明显离散化是合理的。右侧是细网格强制项的平滑或过滤，并提供平滑或限制操作的自然定义。

从粗网格到细网格的最简单的量延长或插值是线性插值。在两个网格的重合点处，粗网格点处的值被简单地注入到相应的细网格点上。在位于粗网格点之间的细网格点处，注入的值是相邻粗网格值的平均值。

因此，两网格迭代方法是：

\begin{itemize}
\item 在精细网格上，使用给出平滑误差的方法执行迭代（限制）；
\item 计算细网格上的残差；
\item 将残差限制（平滑）到粗网格；
\item 在粗网格上执行校正方程的迭代；
\item 延长（插值）对细网格的修正；
\item 更新细网格上的解；
\item 重复整个过程，直至残留物减少至所需水平。
\end{itemize}
人们自然会问：为什么不使用更粗的网格来进一步提高收敛速度呢？这是个好主意。事实上，我们应该继续这一过程，直到无法定义更粗糙的网格为止；在最粗糙的网格上，未知数的数量非常少，以至于可以以可忽略的成本精确求解方程。

多重网格与其说是一种特定方法，不如说是一种策略。在刚刚描述的框架内，有许多参数可以或多或少地任意选择：粗网格结构、限制器（平滑器）、每个网格上的迭代次数、访问各个网格的顺序以及限制和延长（插值）方案是其中最重要的。当然，收敛速度确实取决于所做的选择，但最差和最好方法之间的性能范围可能小于两倍。

多重网格方法最重要的特性是达到给定收敛水平所需的最精细网格上的迭代次数大致与网格节点的数量无关。这正如人们所期望的那样好——计算成本与网格节点的数量成正比。在每个方向大约有 100 个节点的 2D 和 3D 问题中，多重网格方法可以在基本方法所需时间的十分之一到百分之一内收敛。一个例子将在第 4 节中介绍。 5.8.

多重网格方法所基于的迭代方法必须具有良好的平滑性；作为一种独立的方法，它的收敛特性不太重要。 Gauss-Seidel 和 SIP 是两个不错的选择，但还有其他可能性。

在二维中，限制（平滑）算子有多种可能性。如果在每个方向上都使用上述方法，则结果将是九点方案。一个更简单但几乎同样有效的限制是五点方案：

\begin{sourceblock}
\begin{tikzpicture}
\node[inner sep=0,anchor=south west] (image) at (0,0) {\includegraphics[page=152,trim={53.00bp 589.53bp 53.87bp 53.79bp},clip,width=0.92\textwidth,max height=0.38\textheight]{Ferziger4th.pdf}};
\begin{scope}[x={(image.south east)},y={(image.north west)}]
\fill[white] (0.27059,0.44040) rectangle (0.29041,0.94741);
\fill[white] (0.17642,0.15074) rectangle (0.19756,0.65732);
\fill[white] (0.19606,0.07099) rectangle (0.22779,0.45486);
\fill[white] (0.23531,0.15074) rectangle (0.26349,0.65732);
\fill[white] (0.81140,0.15074) rectangle (0.82445,0.65732);
\fill[white] (0.92439,0.13453) rectangle (1.00000,0.64110);
\fill[white] (0.27059,0.00000) rectangle (0.29041,0.33129);
\end{scope}
\end{tikzpicture}
\end{sourceblock}
类似地，有效的延长器是双线性插值。在二维中，细网格上存在三种点。与粗网格点对应的那些被赋予对应点处的值。位于连接两个粗网格点的线上的那些接收两个粗网格值的平均值。最后，粗网格体中心的点取四个相邻值的平均值。对于 FV 方法和 3D 问题，可以导出类似的方案。

迭代求解方法中的初始猜测通常与收敛解相距甚远（通常使用零域）。因此，首先在非常粗糙的网格上求解方程（这很便宜）并使用该解为下一个更精细的网格上的初始场提供更好的猜测是有意义的。当我们到达最精细的网格时，我们已经有了一个相当好的起始解。这种类型的多重网格方法称为全多重网格（FMG）方法。获得最精细网格初始解的成本通常可以通过精细网格迭代节省的成本来补偿。

最后，我们指出，可以构造一种方法，在该方法中求解方程以获得解的近似值，而不是在每个网格处进行校正。这称为完全近似方案（FAS），通常用于解决非线性问题。需要注意的是，FAS 中每个网格上获得的解并不是单独使用该网格时获得的解，而是精细网格解的平滑版本；这是通过将每个网格的校正传递到下一个较粗网格来实现的。纳维-斯托克斯方程的 FAS 方案的一种变体将在第 1 章中介绍。 12.

有关多重网格方法的详细分析，请参阅 Briggs 等人的书籍。 (2000)、Hackbusch (2003) 和 Brandt (1984)。使用 Gauss-Seidel、SIP 或 ICCG 方法作为平滑器的 2D 多重网格求解器可通过互联网获得；详情请参见附录。

\subsection*{5.3.8 其他迭代求解器}\phantomsection\addcontentsline{toc}{subsection}{5.3.8 其他迭代求解器}
还有许多其他迭代求解器，这里无法详细描述。我们提到高斯-赛德尔求解器的“红黑”变体，它通常与多重网格方法结合使用。在结构化网格上，节点被想象为以与棋盘相同的方式“着色”。该方法由两个雅可比步骤组成：首先更新黑色节点，然后更新红色节点。当黑色节点的值更新时，仅使用“旧”红色值，请参见等式： （5.33）。在下一步中，使用更新的黑色值重新计算红色值。雅可比方法对两组节点的交替应用给出了具有与高斯-赛德尔方法相同的收敛特性的整体方法。红黑高斯-赛德尔求解器的一个优点是它可以很好地进行矢量化和并行化，因为任一步骤中都没有数据依赖性。

经常应用于多维问题的另一种做法是使用对应于低维问题的迭代矩阵。其中一个版本是上述 ADI 方法，它将 2D 问题简化为一系列 1D 问题。由此产生的三对角问题逐行解决。迭代过程中改变求解方向以提高收敛速度。该方法通常以高斯-塞德尔方式使用，即使用已访问过的行中的新变量值。

红黑高斯-赛德尔方法的对应物是“斑马”逐行求解器：首先在偶数行上找到解，然后处理奇数行。这提供了更好的并行化和矢量化可能性，而不会牺牲收敛特性。

还可以使用 2D SIP 方法来解决 3D 问题，逐个平面地应用该方法，并将相邻平面的贡献归入方程的右侧。然而，这种方法既不比 3D 版本的 SIP 更便宜，也不更快，因此不经常使用。

\section*{5.4 耦合方程及其解}\phantomsection\addcontentsline{toc}{section}{5.4 耦合方程及其解}
流体动力学和传热中的大多数问题都需要耦合方程组的解决，即每个方程的主导变量出现在其他一些方程中。解决此类问题有两种方法。第一种方法是同时求解所有变量。另一种方法是求解每个方程的主变量，将其他变量视为已知，然后迭代方程直到获得耦合系统的解。这两种方法也可以混合使用。我们分别将这些方法称为同时方法和顺序方法，下面将更详细地描述它们。

\subsection*{5.4.1 同时解}\phantomsection\addcontentsline{toc}{subsection}{5.4.1 同时解}
在联立方法中，所有方程都被视为单个系统的一部分。流体力学的离散方程在线性化后具有块带结构。直接求解方程的成本非常高，特别是当方程是非线性的并且问题是三维的时。耦合系统的迭代求解技术是单个方程方法的推广。选择上述方法是因为它们适用于耦合系统。多位作者开发了基于迭代求解器的同步求解方法；例如，参见 Galpin 和 Raithby (1986)、Deng 等人的论文。 （1994）和韦斯等人。 （1999）。

\subsection*{5.4.2 顺序求解}\phantomsection\addcontentsline{toc}{subsection}{5.4.2 顺序求解}
当方程是线性且紧密耦合时，联立方法是最好的。然而，方程可能非常复杂和非线性，以至于耦合方法使用起来很困难且昂贵。然后，最好将每个方程视为只有一个未知数，暂时将其他变量视为已知，并使用当前可用的最佳值。然后依次求解方程，重复循环，直到满足所有方程。使用此类方法时，需要注意两点：

由于某些项（例如依赖于其他变量的系数和源项）会随着计算的进行而变化，因此在每次迭代时准确求解方程的效率很低。在这种情况下，直接求解器是不必要的，而迭代求解器是首选。对每个方程执行的迭代称为内迭代。 • 为了获得满足所有方程的解，必须在每次循环后更新系数矩阵和源向量，并重复该过程。这些循环称为外迭代。

此类求解方法的优化需要仔细选择每个外部迭代的内部迭代次数。还需要限制从一次外部迭代到下一次迭代（欠松弛）的每个变量的变化，因为一个变量的变化会改变其他方程中的系数，这可能会减慢或阻止收敛。不幸的是，很难分析这些方法的收敛性，因此欠松弛因子的选择很大程度上是凭经验的。

上面描述的多重网格方法作为内部迭代（线性问题）的收敛加速器，可以应用于耦合问题。它也可以用于加速外部迭代，如第 1 章中所述。 12.

\subsection*{5.4.3 欠松弛}\phantomsection\addcontentsline{toc}{subsection}{5.4.3 欠松弛}
我们将介绍一种广泛使用的欠松弛技术。在第 n 次外部迭代中，一般变量 φ 在典型点 P 处的代数方程可写为：

\begin{sourceblock}
\begin{tikzpicture}
\node[inner sep=0,anchor=south west] (image) at (0,0) {\includegraphics[page=154,trim={53.00bp 161.46bp 53.87bp 477.67bp},clip,width=0.92\textwidth,max height=0.38\textheight]{Ferziger4th.pdf}};
\begin{scope}[x={(image.south east)},y={(image.north west)}]
\fill[white] (0.00000,0.89930) rectangle (0.09732,1.00000);
\fill[white] (0.34586,0.42799) rectangle (0.41236,0.95557);
\end{scope}
\end{tikzpicture}
\end{sourceblock}
其中 Q 包含所有不显式依赖于 φ n 的项；系数A l 和源Q 可能涉及φ n − 1。离散化方案在这里并不重要。该方程是线性的，整个解域的方程组通常通过迭代（内部迭代）求解。

在早期的外部迭代中，允许 φ 改变与方程 1 一样多的值。 ( 5.70 ) 要求可能会导致不稳定，因此我们允许 φ n 仅改变可能差异的一小部分 α φ：

\begin{sourceblock}
\begin{tikzpicture}
\node[inner sep=0,anchor=south west] (image) at (0,0) {\includegraphics[page=155,trim={53.00bp 558.62bp 53.87bp 88.54bp},clip,width=0.92\textwidth,max height=0.38\textheight]{Ferziger4th.pdf}};
\begin{scope}[x={(image.south east)},y={(image.north west)}]
\fill[white] (0.00000,0.72076) rectangle (0.13311,1.00000);
\end{scope}
\end{tikzpicture}
\end{sourceblock}
其中 φ new 是方程的结果。 ( 5.70 ) 且欠松弛因子满足 0 < α φ < 1。

因为更新系数矩阵和源向量后通常不再需要旧的迭代，因此可以将新的解重写。将方程式中的 φ 替换为新值。 ( 5.71 ) 由

\begin{sourceblock}
\begin{tikzpicture}
\node[inner sep=0,anchor=south west] (image) at (0,0) {\includegraphics[page=155,trim={53.00bp 461.89bp 53.87bp 167.55bp},clip,width=0.92\textwidth,max height=0.38\textheight]{Ferziger4th.pdf}};
\begin{scope}[x={(image.south east)},y={(image.north west)}]
\fill[white] (0.00000,0.73842) rectangle (0.05735,1.00000);
\fill[white] (0.06156,0.72834) rectangle (0.08968,1.00000);
\fill[white] (0.09236,0.72834) rectangle (0.13792,1.00000);
\fill[white] (0.14063,0.72834) rectangle (0.21780,1.00000);
\fill[white] (0.22051,0.72834) rectangle (0.25525,1.00000);
\end{scope}
\end{tikzpicture}
\end{sourceblock}
由方程式得出( 5.70 ) 得出节点 P 处的修正方程：

\begin{sourceblock}
\begin{tikzpicture}
\node[inner sep=0,anchor=south west] (image) at (0,0) {\includegraphics[page=155,trim={53.00bp 388.25bp 53.87bp 232.18bp},clip,width=0.92\textwidth,max height=0.38\textheight]{Ferziger4th.pdf}};
\begin{scope}[x={(image.south east)},y={(image.north west)}]
\fill[white] (0.23528,0.69569) rectangle (0.27065,0.96806);
\fill[white] (0.23411,0.39401) rectangle (0.26812,0.66659);
\end{scope}
\end{tikzpicture}
\end{sourceblock}
where A ∗
P and Q ∗
P 是修改后的主对角矩阵元素和源向量分量。该修改后的方程在内部迭代中求解。当外层迭代收敛时，涉及 α 的项相互抵消，我们得到了原问题的解。

这种欠松弛是由 Patankar (1980) 提出的。它对许多迭代求解方法有积极的影响，因为矩阵 A 的对角优势增加了（元素 A *
P 大于 A P，而 A l 保持不变）。它比显式应用表达式（5.71）更有效。

最佳欠松弛因子取决于问题。一个好的策略是在早期迭代中使用小的欠松弛因子，并在接近收敛时将其增加到统一。第 1 章将给出一些选择欠松弛因子来求解纳维-斯托克斯方程的指导。 7、8、9和12。欠松弛不仅可以应用于因变量，还可以应用于方程中的各个项。当流体属性（粘度、密度、普朗特数等）取决于解并且需要更新时，通常有必要这样做。

我们上面提到，迭代求解方法通常可以被视为求解非稳态问题，直到达到稳态。时间步长的控制对于控制解的演化非常重要。在下一章中，我们将展示时间步长可以被解释为欠松弛因素。上述欠松弛方案可以被解释为在不同节点处使用不同的时间步长。

\section*{5.5 非线性方程及其解}\phantomsection\addcontentsline{toc}{section}{5.5 非线性方程及其解}
如上所述，求解非线性方程的技术有两种：类牛顿技术和全局技术。当可以很好地估计解时，前者要快得多，但后者保证不会发散；速度和安全性之间需要权衡。经常使用这两种方法的组合。有大量文献致力于解决非线性方程的方法，并且最先进的技术仍在不断发展。我们甚至无法在这里涵盖大部分方法，并且只能对某些方法进行简短概述。

\subsection*{5.5.1 类牛顿技术}\phantomsection\addcontentsline{toc}{subsection}{5.5.1 类牛顿技术}
求解非线性方程的主要方法是牛顿法。假设需要找到单个代数方程 f ( x ) = 0 的根。牛顿方法使用泰勒级数的前两项对有关 x 估计值的函数进行线性化：

\begin{sourceblock}
\begin{tikzpicture}
\node[inner sep=0,anchor=south west] (image) at (0,0) {\includegraphics[page=156,trim={53.00bp 375.05bp 53.87bp 274.25bp},clip,width=0.92\textwidth,max height=0.38\textheight]{Ferziger4th.pdf}};
\begin{scope}[x={(image.south east)},y={(image.north west)}]
\fill[white] (0.00000,0.82779) rectangle (0.02908,1.00000);
\fill[white] (0.03179,0.82779) rectangle (0.07320,1.00000);
\fill[white] (0.07588,0.82779) rectangle (0.15814,1.00000);
\fill[white] (0.16087,0.82779) rectangle (0.24223,1.00000);
\end{scope}
\end{tikzpicture}
\end{sourceblock}
将线性化函数设置为零可提供根的新估计：

\begin{sourceblock}
\includegraphics[page=156,trim={53.00bp 315.12bp 53.87bp 321.84bp},clip,width=0.92\textwidth,max height=0.38\textheight]{Ferziger4th.pdf}
\end{sourceblock}
我们继续下去，直到根 x k − x k − 1 的变化达到期望的程度。该方法相当于通过 x k 处的切线来近似表示函数的曲线。当估计值足够接近根时，该方法二次收敛，即迭代 k + 1 处的误差与迭代 k 处的误差的平方成正比。这意味着一旦解估计接近根，只需要几次迭代。因此，只要可行，就会采用它。

牛顿方法很容易推广到方程组。一个通用的非线性方程组可以写成：

\begin{sourceblock}
\includegraphics[page=156,trim={53.00bp 173.96bp 53.87bp 477.07bp},clip,width=0.92\textwidth,max height=0.38\textheight]{Ferziger4th.pdf}
\end{sourceblock}
该系统可以按照与单个方程完全相同的方式进行线性化。唯一的区别是现在我们需要使用多变量泰勒级数：

\begin{sourceblock}
\includegraphics[page=156,trim={53.00bp 84.68bp 53.87bp 534.90bp},clip,width=0.92\textwidth,max height=0.38\textheight]{Ferziger4th.pdf}
\end{sourceblock}
对于 i = 1 , 2 , .。。，n。当它设置为零时，我们就有了一个线性代数方程组，可以通过高斯消去法或其他一些技术来求解。系统的矩阵是偏导数的集合：

\begin{sourceblock}
\includegraphics[page=157,trim={53.00bp 532.78bp 53.87bp 103.07bp},clip,width=0.92\textwidth,max height=0.38\textheight]{Ferziger4th.pdf}
\end{sourceblock}
称为系统的雅可比行列式。方程组为：

\begin{sourceblock}
\includegraphics[page=157,trim={53.00bp 467.66bp 53.87bp 162.95bp},clip,width=0.92\textwidth,max height=0.38\textheight]{Ferziger4th.pdf}
\end{sourceblock}
对于接近正确根的估计，牛顿系统方法的收敛速度与单个方程的方法一样快。然而，对于大型系统，其主要缺点足以抵消快速收敛的影响。为了使该方法有效，必须在每次迭代时评估雅可比行列式。这带来了两个困难。首先，在一般情况下，雅可比行列式有 n 2 个元素，它们的评估成为该方法中最昂贵的部分。第二是可能不存在评估雅可比行列式的直接方法；许多系统的方程是隐式的，或者方程可能非常复杂，以致微分几乎不可能。

对于一般的非线性方程组，割线法更为有效。对于单个方程，割线法通过在曲线上两点之间绘制的割线来近似函数的导数（参见 Ferziger 1998 或 Moin 2010）。该方法比牛顿法收敛得慢，但由于不需要求导数，因此可以以较低的总成本找到解，并且可以应用于无法直接求导数的问题。割线方法对系统有许多推广，其中大多数都非常有效，但由于它们尚未在 CFD 中应用，因此我们在此不对其进行回顾。

\subsection*{5.5.2 其他技术}\phantomsection\addcontentsline{toc}{subsection}{5.5.2 其他技术}
求解耦合非线性方程的常用方法是上一节中描述的顺序解耦方法。非线性项（对流通量、源项）通常使用皮卡迭代方法进行线性化。对于对流项，这意味着质量通量被视为已知，因此 u i 动量分量方程中的非线性对流项近似为：

\begin{sourceblock}
\begin{tikzpicture}
\node[inner sep=0,anchor=south west] (image) at (0,0) {\includegraphics[page=157,trim={53.00bp 97.06bp 53.87bp 552.08bp},clip,width=0.92\textwidth,max height=0.38\textheight]{Ferziger4th.pdf}};
\begin{scope}[x={(image.south east)},y={(image.north west)}]
\fill[white] (0.00000,0.84000) rectangle (0.04241,1.00000);
\end{scope}
\end{tikzpicture}
\end{sourceblock}
其中索引 o 表示这些值取自先前外部迭代的结果。类似地，源术语被分解为两部分：

\begin{sourceblock}
\includegraphics[page=158,trim={53.00bp 558.03bp 53.87bp 92.66bp},clip,width=0.92\textwidth,max height=0.38\textheight]{Ferziger4th.pdf}
\end{sourceblock}
部分b 0 被吸收到代数方程的右侧，而b 1 则贡献于系数矩阵A。类似的方法可用于涉及多个变量的非线性项。

这种线性化比使用类牛顿线性化的耦合技术需要更多的迭代。然而，每次迭代的计算量要小得多，并且使用多重网格技术可以减少外部迭代的数量，这使得这种方法很有吸引力。

牛顿法有时用于对非线性项进行线性化；例如，动量分量 u i 的方程中的对流项可以表示为：

\begin{sourceblock}
\begin{tikzpicture}
\node[inner sep=0,anchor=south west] (image) at (0,0) {\includegraphics[page=158,trim={53.00bp 413.11bp 53.87bp 234.58bp},clip,width=0.92\textwidth,max height=0.38\textheight]{Ferziger4th.pdf}};
\begin{scope}[x={(image.south east)},y={(image.north west)}]
\fill[white] (0.00000,0.85257) rectangle (0.03744,1.00000);
\end{scope}
\end{tikzpicture}
\end{sourceblock}
非线性源项可以用相同的方式处理。这导致难以求解的耦合线性方程组，并且除非使用完整的牛顿技术，否则收敛性不是二次的。然而，可以开发受益于这种线性化技术的特殊耦合迭代技术，如 Galpin 和 Raithby (1986) 所示。

\section*{5.6 延迟修正方法}\phantomsection\addcontentsline{toc}{section}{5.6 延迟修正方法}
如果包含未知变量的节点值的所有项都保留在方程的左侧。 （3.43），计算分子可能会变得非常大。由于计算分子的大小会影响存储要求和求解线性方程组所需的工作量，因此我们希望使其尽可能小；通常，只有节点 P 的最近邻居保留在方程的左侧。然而，产生如此简单的计算分子的近似通常不够准确，因此我们经常被迫使用引用更多节点而不仅仅是最近邻居的近似。

解决这个问题的一种方法是仅将包含最近邻的项保留在方程的左侧。 ( 3.43 ) 并将所有其他项移至右侧。这要求使用先前迭代中的值来评估这些项。然而，这不是一个好的实践，并且可能导致迭代的发散，因为明确处理的术语可能是大量的。为了防止发散，需要对从一次迭代到下一次迭代的变化进行强烈的欠松弛（参见第 5.4.3 节），从而导致收敛缓慢。

更好的方法是显式计算用高阶近似来近似的项，并将它们放在方程的右侧。然后对这些项进行更简单的近似（给出一个小的计算分子）并将其放在方程的左侧（具有未知变量值）和右侧（使用现有值显式计算它）。右侧现在是同一项的两个近似值之间的差异，并且可能很小。因此，它不会在迭代求解过程中造成任何问题。一旦迭代收敛，低阶近似项就会消失，获得的解对应于高阶近似项。

由于要求解的方程是非线性的，通常需要迭代方法，因此向显式处理的部分添加一个小项仅会增加少量的计算工作量。另一方面，当隐式处理的方程部分中的计算分子的大小较小时，所需的存储器和计算时间都大大减少。

我们会经常提到这种技术。它用于处理高阶近似、网格非正交性以及避免解中振荡等不良影响所需的校正。因为等式右边可以看作是“修正”，所以这种方法称为延迟修正。这里将结合 FD 中的 Padé 方案（参见第 3.3.3 节）和 FV 方法中的高阶插值（参见第 4.4.4 节）来演示其使用。

如果将 Padé 方案用于隐式 FD 方法，则必须采用延迟校正，因为一个节点处导数的逼近涉及相邻节点处的导数。一种方法是使用相邻节点处导数的“旧值”和远处节点处的变量值。这些通常取自先前迭代的结果，我们有：

\begin{sourceblock}
\begin{tikzpicture}
\node[inner sep=0,anchor=south west] (image) at (0,0) {\includegraphics[page=159,trim={53.00bp 291.42bp 53.87bp 337.30bp},clip,width=0.92\textwidth,max height=0.38\textheight]{Ferziger4th.pdf}};
\begin{scope}[x={(image.south east)},y={(image.north west)}]
\fill[white] (0.04355,0.47167) rectangle (0.08180,0.78060);
\fill[white] (0.92439,0.00000) rectangle (1.00000,0.09781);
\end{scope}
\end{tikzpicture}
\end{sourceblock}
(5.83) 在这种情况下，只有该方程右侧的第一项将被移动到方程的左侧，以便在新的外部迭代中求解。

然而，这种方法可能会对收敛速度产生不利影响，因为隐式处理的部分不是导数的近似值，而是导数的某个倍数。以下版本的延迟修正更有效：

\begin{sourceblock}
\begin{tikzpicture}
\node[inner sep=0,anchor=south west] (image) at (0,0) {\includegraphics[page=159,trim={53.00bp 179.04bp 53.87bp 446.75bp},clip,width=0.92\textwidth,max height=0.38\textheight]{Ferziger4th.pdf}};
\begin{scope}[x={(image.south east)},y={(image.north west)}]
\fill[white] (0.12965,0.08079) rectangle (0.16472,0.37398);
\end{scope}
\end{tikzpicture}
\end{sourceblock}
这里，左侧使用完整的二阶 CDS 近似。在右侧，我们可以看到显式计算的 Padé 方案导数和显式计算的 CDS 近似之间的差异。这给出了更平衡的表达式，因为在二阶 CDS 足够准确的情况下，方括号中的项可以忽略不计。我们可以使用 UDS 作为隐式部分，而不是 CDS；在这种情况下，方程两边都应使用 UDS 近似。

当使用高阶方案时，延迟校正在 FV 方法中也很有用（参见第 4.4.4 节）。高阶通量近似值是显式计算的，然后该近似值与隐式低阶近似值（仅使用最近邻的变量值）相结合（首先由 Khosla 和 Rubin 1974 年建议）：

\begin{sourceblock}
\begin{tikzpicture}
\node[inner sep=0,anchor=south west] (image) at (0,0) {\includegraphics[page=160,trim={53.00bp 520.97bp 53.87bp 120.51bp},clip,width=0.92\textwidth,max height=0.38\textheight]{Ferziger4th.pdf}};
\begin{scope}[x={(image.south east)},y={(image.north west)}]
\fill[white] (0.33344,0.11071) rectangle (0.36445,0.61557);
\fill[white] (0.36965,0.15734) rectangle (0.39783,0.62652);
\fill[white] (0.40343,0.14680) rectangle (0.44211,0.68573);
\end{scope}
\end{tikzpicture}
\end{sourceblock}
F L
e 代表一些低阶方案的近似（UDS 通常用于对流，CDS 用于扩散通量）和 F H
e 是高阶近似值。括号中的术语使用上一次迭代的值进行评估，如上标“旧”所示。与隐式部分相比，它通常较小，因此其显式处理不会显著影响收敛性。

相同的方法可以应用于所有高阶近似，包括谱方法。尽管相对于纯低阶方案，延迟校正增加了每次迭代的计算时间，但额外的工作量比隐式处理整个高阶近似所需的工作量小得多。人们还可以将“旧”项与零和单位之间的混合因子相乘，以产生纯低阶和纯高阶方案的混合。有时，这用于避免在不够精细的网格上使用高阶方案时导致的振荡。例如，当计算物体周围的流动时，人们希望在物体附近使用精细网格，而在远离物体的地方使用较粗糙的网格。高阶方案可能会在粗网格区域产生振荡，从而破坏整个解。由于粗网格区域中变量变化缓慢，我们可以在不影响细网格区域解的情况下降低粗网格区域的近似阶数；这可以通过仅在外部区域使用混合因子来实现。

关于延迟修正方法的其他用途的更多细节将在后续章节中给出。

\section*{5.7 收敛准则和迭代误差}\phantomsection\addcontentsline{toc}{section}{5.7 收敛准则和迭代误差}
使用迭代求解器时，知道何时退出非常重要。最常见的过程是基于两次连续迭代之间的差异；当通过某种标准测量的该差异小于预先选择的值时，该过程停止。不幸的是，当误差不小并且适当的标准化是必要的时，这种差异可能很小。

从《节》中提出的分析来看。 5.3.2，我们发现（见式（5.14）和（5.27））：

\begin{sourceblock}
\includegraphics[page=160,trim={53.00bp 102.84bp 53.87bp 542.72bp},clip,width=0.92\textwidth,max height=0.38\textheight]{Ferziger4th.pdf}
\end{sourceblock}
其中 δ n 是迭代 n + 1 和 n 处的解之间的差，λ 1 是迭代矩阵的最大特征值或谱半径。可以从（Ferziger 1998）估计：

\begin{sourceblock}
\includegraphics[page=161,trim={53.00bp 536.16bp 53.87bp 99.40bp},clip,width=0.92\textwidth,max height=0.38\textheight]{Ferziger4th.pdf}
\end{sourceblock}
对于足够大的 n 且其中 ∥ a ∥ 表示范数（例如，均方根或
L 2 范数) a。

一旦获得了特征值的估计，估计迭代误差就不难了。事实上，通过重新排列方程。 ( 5.86 )，我们发现：

\begin{sourceblock}
\includegraphics[page=161,trim={53.00bp 440.73bp 53.87bp 194.24bp},clip,width=0.92\textwidth,max height=0.38\textheight]{Ferziger4th.pdf}
\end{sourceblock}
因此，迭代误差的良好估计是：

\begin{sourceblock}
\includegraphics[page=161,trim={53.00bp 383.11bp 53.87bp 251.86bp},clip,width=0.92\textwidth,max height=0.38\textheight]{Ferziger4th.pdf}
\end{sourceblock}
该误差估计可以根据解的两个连续迭代来计算。虽然该方法是为线性系统设计的，但所有系统本质上都是线性接近收敛的；由于此时最需要误差估计，因此该方法也可以应用于非线性系统。

不幸的是，迭代方法通常具有复杂的特征值。在这种情况下，误差减少不是指数的，也可能不是单调的。由于方程是实数，因此复数特征值必须以共轭对的形式出现。他们的估计需要扩展上述过程。特别是，需要来自更多迭代的数据。下面使用的一些想法可以在 Golub 和 van Loan (1996) 中找到。

如果最大幅值的特征值是复数，则至少有两个特征值，方程： ( 5.27 ) 必须替换为

\begin{sourceblock}
\includegraphics[page=161,trim={53.00bp 205.47bp 53.87bp 439.92bp},clip,width=0.92\textwidth,max height=0.38\textheight]{Ferziger4th.pdf}
\end{sourceblock}
其中 * 表示复数的共轭。和之前一样，我们减去两次连续的迭代以获得 δ n；参见等式。 （5.86）。如果我们进一步让：

\begin{sourceblock}
\includegraphics[page=161,trim={53.00bp 145.99bp 53.87bp 503.80bp},clip,width=0.92\textwidth,max height=0.38\textheight]{Ferziger4th.pdf}
\end{sourceblock}
则得到如下表达式：

\begin{sourceblock}
\includegraphics[page=161,trim={53.00bp 98.20bp 53.87bp 547.39bp},clip,width=0.92\textwidth,max height=0.38\textheight]{Ferziger4th.pdf}
\end{sourceblock}
因为特征值 λ 1 的大小是最感兴趣的量，所以我们写：

\begin{sourceblock}
\begin{tikzpicture}
\node[inner sep=0,anchor=south west] (image) at (0,0) {\includegraphics[page=162,trim={53.00bp 569.98bp 53.87bp 76.59bp},clip,width=0.92\textwidth,max height=0.38\textheight]{Ferziger4th.pdf}};
\begin{scope}[x={(image.south east)},y={(image.north west)}]
\fill[white] (0.00000,0.72969) rectangle (0.07405,1.00000);
\end{scope}
\end{tikzpicture}
\end{sourceblock}
简单的计算表明：

\begin{sourceblock}
\includegraphics[page=162,trim={53.00bp 523.29bp 53.87bp 124.28bp},clip,width=0.92\textwidth,max height=0.38\textheight]{Ferziger4th.pdf}
\end{sourceblock}
从中不难看出：

\begin{sourceblock}
\begin{tikzpicture}
\node[inner sep=0,anchor=south west] (image) at (0,0) {\includegraphics[page=162,trim={53.00bp 461.27bp 53.87bp 176.29bp},clip,width=0.92\textwidth,max height=0.38\textheight]{Ferziger4th.pdf}};
\begin{scope}[x={(image.south east)},y={(image.north west)}]
\fill[white] (0.42830,0.29846) rectangle (0.47883,0.70294);
\end{scope}
\end{tikzpicture}
\end{sourceblock}
是特征值大小的估计。

误差的估计需要进一步的近似。复杂的特征值会导致误差振荡，并且误差的形状并不独立于迭代次数，即使对于较大的 n 也是如此。为了估计误差，我们根据上面给出的表达式计算 δ n 和 l。由于复特征值和特征向量，结果包含与相位角余弦成比例的项。由于我们只对幅度感兴趣，因此我们假设这些项在平均意义上为零并丢弃它们。这使我们能够找到两个数量之间的简单关系：

\begin{sourceblock}
\includegraphics[page=162,trim={53.00bp 318.73bp 53.87bp 316.01bp},clip,width=0.92\textwidth,max height=0.38\textheight]{Ferziger4th.pdf}
\end{sourceblock}
这是所需的误差估计。由于解中的振荡，任何特定迭代的估计可能不准确，但正如我们将在下面展示的，平均而言它相当不错。

为了消除振荡的一些影响，应在一系列迭代中对特征值估计进行平均。根据问题和预期迭代次数，平均范围可能在 1 到 50 之间变化（通常为预期迭代次数的 1\%）。最后，我们需要一种可以处理实数和复数特征值的方法。如果主特征值 ( λ 1 ) 为实数，则复杂情况 ( 5.96 ) 的误差估计器会给出较低的估计值。此外，在这种情况下，λ 1 对 z n 的贡献消失，因此特征值估计非常差。然而，这个事实可以用来确定 λ 1 是实数还是复数。如果比例

\begin{sourceblock}
\includegraphics[page=162,trim={53.00bp 130.96bp 53.87bp 506.12bp},clip,width=0.92\textwidth,max height=0.38\textheight]{Ferziger4th.pdf}
\end{sourceblock}
很小，特征值可能是实数；如果 r 很大，则特征值可能很复数。对于实数特征值，r 往往小于 10 − 2，对于复数特征值，r ≈ 1。因此可以采用 r = 0 的值。 1 作为特征值类型的指示符，并使用适当的误差估计表达式。

一种折衷方案是使用残差的减少作为停止标准。当残差范数减少到其原始大小的一部分（通常减少三或四个数量级）时，迭代停止。正如我们所表明的，迭代误差通过方程与残差相关。 （5.15）因此残差的减少伴随着迭代误差的减少。如果迭代从零初始值开始，则初始误差等于解本身。当残留水平下降到比初始水平低三到四个数量级时，误差可能下降了相当的量，即，约为溶液的 0.1\%。在迭代过程开始时，残差和误差通常不会以相同的方式下降；还需要小心，因为如果矩阵条件不好，即使残差很小，误差也可能很大。

许多迭代求解器需要计算残差。上述方法很有吸引力，因为它不需要额外的计算。第一次内迭代之前的残差范数为检查内迭代的收敛性提供了参考。同时它提供了外部迭代收敛的度量。经验表明，当残差下降一到两个数量级时，可以停止内部迭代。在残差减少三到五个数量级之前，不应停止外部迭代，具体取决于所需的精度。可以使用绝对残差之和（L 1 范数）代替均方根。 (L 2 )范数。细化网格的收敛准则应该更加严格，因为细化网格上的离散化误差比粗网格上的要小。

如果初始误差的阶数已知，则可以监视两次迭代之间差异的范数，并将其与迭代过程开始时的相同量进行比较。当差异范数下降三到四个数量级时，误差通常也会下降相当的量。

这两种方法都只是近似方法；然而，它们优于基于两个连续迭代之间的非标准化差异的标准。

为了测试估计迭代误差的方法，我们首先使用 SOR 求解器给出线性二维问题的解。线性问题是平方域中的拉普拉斯方程 \{ 0 < x < 1 ; 0 < y < 1 \}，选择狄利克雷边界条件来对应解 φ( x , y ) = 100 xy。这种选择的优点是，收敛解的二阶中心差分近似在任何网格上都是精确的，因此可以轻松计算当前迭代和收敛解之间的实际差异。解的初始猜测在域内的任何地方都为零。我们选择SOR方法作为迭代技术，因为如果松弛参数大于最优值，则特征值是复数。

\begin{center}\small 图 2 和 3 中显示了具有 20 × 20 和 80 × 80 CV 的均匀网格的结果。 5.3 和 5.4。在每种情况下，都会显示精确迭代误差的范数、使用上述技术的误差估计、两次迭代之间的差异以及残差。对于这两种情况，显示了松弛参数的两个值的计算结果：一个低于最佳值，具有实特征值，\end{center}
\begin{sourceblock}
\includegraphics[page=164,trim={53.00bp 392.25bp 53.87bp 52.00bp},clip,width=\textwidth,max height=0.38\textheight]{Ferziger4th.pdf}
\par\vspace{0.45\baselineskip}
\small 图 5.3 在 20 × 20 CV 网格上使用 SOR 求解器解决拉普拉斯问题的精确迭代误差、误差估计、残差和两次迭代之间的差的范数的变化：松弛参数比最优值小（左）和大（右），约为 1.73
\end{sourceblock}
\begin{sourceblock}
\includegraphics[page=164,trim={53.00bp 164.90bp 53.87bp 307.37bp},clip,width=\textwidth,max height=0.38\textheight]{Ferziger4th.pdf}
\par\vspace{0.45\baselineskip}
\small 图 5.4 在 80 × 80 CV 网格上使用 SOR 求解器解决拉普拉斯问题的精确迭代误差、误差估计、残差和两次迭代之间的差的范数的变化：松弛参数比最佳值小（左）和大（右），约为 1.92，比最佳值高 1，导致复杂的特征值。 4 对于实特征值的情况，指数收敛结果平滑。在这种情况下，误差估计几乎是准确的（除了初始阶段）。然而，两次迭代之间的残差和差值的范数最初下降得太快，并且不跟随迭代误差的下降。随着网格的细化，这种效果更加明显。在80×80 CV网格上，残差范数迅速降低了两个数量级，而误差仅略有降低。一旦达到渐近缩减率，所有四条曲线的斜率都相同。
\end{sourceblock}
在迭代矩阵的特征值是复数的情况下，收敛不是单调的——误差存在振荡。在这种情况下，预测误差和精确误差的比较也非常令人满意。在这种情况下，所有上述收敛标准似乎都同样好。

稍后将介绍迭代误差估计的进一步示例，特别是在解决耦合流问题的情况下的外部迭代。

\section*{5.8 例子}\phantomsection\addcontentsline{toc}{section}{5.8 例子}
在上一章中，我们提出了一些二维问题的解，但没有讨论解决方法。现在，我们将展示针对驻点流中标量传输情况的各种求解器的性能；见第 4.7 节 用于描述问题和用于推导线性代数方程的离散化技术。

我们考虑 = 0 的情况。 01 和具有 20 × 20、40 × 40 和 80 × 80 CV 的均匀网格。方程矩阵 A 不是对称的，并且在 CDS 离散化的情况下，它不是对角占优的。在对角主导矩阵中，主对角线上的元素满足以下条件：

\begin{sourceblock}
\begin{tikzpicture}
\node[inner sep=0,anchor=south west] (image) at (0,0) {\includegraphics[page=165,trim={53.00bp 235.51bp 53.87bp 405.16bp},clip,width=0.92\textwidth,max height=0.38\textheight]{Ferziger4th.pdf}};
\begin{scope}[x={(image.south east)},y={(image.north west)}]
\fill[white] (0.40806,0.45387) rectangle (0.44343,0.94268);
\fill[white] (0.44875,0.49902) rectangle (0.47615,0.95289);
\end{scope}
\end{tikzpicture}
\end{sourceblock}
可以看出，迭代求解方法收敛的充分条件是满足上述关系，并且不等式必须至少适用于一个节点。只有对流项的 UDS 离散化才能满足此条件。虽然像 Jacobi 和 Gauss-Seidel 这样的简单求解器在违反上述条件时通常会发散，但 ILU、SIP 和共轭梯度求解器对矩阵的对角优势不太敏感。

我们考虑了五个求解器：

\begin{itemize}
\item 高斯-赛德尔，用GS 表示；
\item 在 x = const 线上使用 TDMA 的高斯-赛德尔线，用 LGS-X 表示；
\item 在y = const 线上使用TDMA 的高斯-塞德尔线，用LGS-Y 表示；
\end{itemize}
\footnotetext[4]{对于简单的矩形几何和狄利克雷边界条件 (Brazier 1974)： ω = 2 /( 1 + sin (π/ N CV ))。}
\begin{center}\small 表 5.1 求解驻点流中的二维标量输运问题时，各种求解器将 L 1 残差范数降低 4 个数量级所需的迭代次数\end{center}
方案 网格 GS LGS-X LGS-Y LGS-ADI SIP
UDS 20 × 20 68 40 35 18 14 40 × 40 211 114 110 52 21 80 × 80 720 381 384 175 44 CDS 20 × 20 – – – 12 19 40 × 40 163 95 77 39 19 80×80 633 349 320 153 40

\begin{itemize}
\item 线高斯-赛德尔在x = const 线上交替使用TDMA。 y = const.，用 LGS-ADI 表示；
\item Stone 的 ILU 方法，用 SIP 表示。
\end{itemize}
\begin{center}\small 表 5.1 显示了上述求解器将绝对残差和减少四个数量级所需的迭代次数。\end{center}
从表中我们可以看出 LGS-X 和 LGS-Y 求解器的速度大约是 GS 的两倍； LGS-ADI 的速度约为 LGS-X 的两倍，而在更精细的网格上，SIP 的速度约为 LGS-ADI 的四倍。对于 GS 和 LGS 求解器，每次细化网格时，迭代次数都会增加约四倍；在 SIP 和 LGS-ADI 的情况下，该因子较小，但是，正如我们将在下一个示例中看到的，在非常精细的网格的限制下，该因子会增加并渐近地接近 4。

另一个有趣的观察结果是，GS 和 LGS 求解器不会在 CDS 离散化的 20 × 20 CV 网格上收敛。这是因为在这种情况下矩阵不是对角占优的。即使在 40 × 40 CV 网格上，矩阵也不是完全对角占优，但违规发生在变量均匀分布的区域（低梯度），因此对求解器的影响并不那么严重。 LGS-ADI 和 SIP 求解器不受影响。

我们现在转向一个测试用例，该用例存在解析解，并且 CDS 近似在任何网格上都会产生精确解。这有助于评估迭代误差，但求解器的行为不会受益，因此这种情况非常适合评估求解器性能。我们正在用狄利克雷边界条件求解拉普拉斯方程，其精确解为 φ = xy。解域是一个矩形，解在所有边界处都有规定，并且内部的初始值全部为零。因此，初始误差等于解，并且是空间坐标的平滑函数。使用前一章中描述的FV方法和CDS方案进行离散化。因为不存在对流，所以问题是完全椭圆的。

考虑的求解器是：

\begin{itemize}
\item 高斯-赛德尔求解器，用GS 表示；
\item 使用TDMA 的线高斯-赛德尔解算器交替沿线x = const。 y = const.，用 LGS-ADI 表示；
\item 第1 节中描述的ADI 求解器。 5.3.5；
\end{itemize}
\begin{center}\small 表 5.2 对于在方形域 X × Y = 1 × 1 上具有狄利克雷边界条件的二维拉普拉斯方程，将归一化 L 1 误差范数降低到 10 − 5 以下所需的迭代次数，并且在两个方向上都有均匀网格\end{center}
网格 GS LGS-ADI ADI SIP ICCG MG-GS MG-SIP
8 × 8 74 22 16 8 7 12 7 16 × 16 292 77 31 20 13 10 6 32 × 32 1160 294 64 67 23 10 6 64 × 64 4622 1160 132 254 46 10 6 128 × 128 – – 274 1001 91 10 6 256 × 256 – – – – 181 10 6

\begin{center}\small 表 5.3 ADI 求解器所需的迭代次数与时间步长的函数关系（双向均匀网格，64 × 64 CV）\end{center}
1 / t 80 68 64 60 32 16 8
不，Iter。 152 134 132 134 234 468 936

\begin{itemize}
\item Stone 的 ILU 方法，用 SIP 表示；
\item 共轭梯度法，使用不完全Cholesky 分解进行预处理，表示为ICCG；
\item 使用GS 作为平滑器的多重网格方法，表示为MG-GS；
\item 使用SIP 作为平滑器的多重网格方法，表示为MG-SIP。
\end{itemize}
\begin{center}\small 表 5.2 显示了在方形解域上的均匀网格上获得的结果。 LGS-ADI 的速度大约是 GS 的四倍，SIP 的速度大约是 LGS-ADI 的四倍。在粗网格上，ADI 的效率低于 SIP，但是，当选择最佳时间步长时，只有当一个方向上的网格点数加倍时，迭代次数才会加倍，因此对于细网格，ADI 相当有效。当时间步长以循环方式变化时，求解器变得更加高效。 SIP 也是如此，但参数的循环变化会增加每次迭代的成本。表 5.3 给出了 ADI 在不同时间步长下达到收敛所需的迭代次数。当网格细化时，最佳时间步长会减少两倍。\end{center}
ICCG 比 SIP 快得多；当网格细化时，迭代次数加倍，因此在细网格上其优势更大。多重网格求解器非常高效；使用 SIP 更平滑，只需要在最精细的网格上进行 6 次迭代。在 MG 求解器中，最粗级别为 2 × 2 CV，因此 8 × 8 CV 网格上有 3 个级别，256 × 256 CV 网格上有 8 个级别。在最细网格和延长后的所有网格上进行 1 次迭代，而在限制阶段进行 4 次迭代。在SIP中，参数α设置为0.92。没有尝试寻找参数的最佳值；获得的结果具有足够的代表性，足以显示各种求解器的趋势和相对性能。另请注意，每次迭代的计算量对于每个迭代都是不同的

\begin{sourceblock}
\includegraphics[page=168,trim={53.00bp 424.25bp 53.87bp 52.00bp},clip,width=\textwidth,max height=0.38\textheight]{Ferziger4th.pdf}
\par\vspace{0.45\baselineskip}
\small 图 5.5 残差（左）和迭代误差（右）的 L 1 范数随各种求解器和 64 × 64 CV 网格执行迭代次数的变化
\end{sourceblock}
求解器。以 GS 迭代的成本为参考，我们发现以下相对成本：LGS-ADI — 2.5，ADI — 3.0，SIP — 第一次迭代为 4.0，之后为 2.0，ICCG — 第一次迭代为 4.5，之后为 3.0。对于 MG 方法，需要将最细网格上的迭代次数乘以大约 1.5，以考虑粗网格上的迭代成本。因此，在这种情况下，MG-GS 是计算效率最高的求解器。

由于每个求解器的收敛速度不同，因此相对成本取决于我们想要求解方程的精确程度。为了分析这个问题，我们在图 5.5 中绘制了绝对残差之和的变化以及迭代误差随迭代的变化。可以得出两个结论：

\begin{itemize}
\item 残差和的下降最初是不规则的，但经过一定次数的迭代后，收敛率变得恒定。ICCG 求解器是一个例外，它随着迭代的进行而变得更快。当需要非常精确的解时，MG 求解器和 ICCG 是最佳选择。如果（就像解决非线性问题时的情况一样）需要中等精度，SIP 就会变得具有竞争力，在这种情况下甚至 ADI 也可能足够好。
\item GS、SIP、ICCG 和ADI 求解器初始残差范数的减少并未伴随迭代误差的同等减少。只有 MG 求解器才能以相同的速度减少误差和残差范数。
\end{itemize}
这些结论非常普遍，尽管存在与问题相关的特征。我们将在第 1 章中展示纳维-斯托克斯方程的类似结果。 8.

由于 SIP 求解器用于 CFDon 结构化网格，因此我们在图 5.6 中显示了达到收敛所需的迭代次数对参数 α 的依赖性。当 α = 0 时，SIP 简化为标准 ILU 求解器。在 α 的最佳值下，SIP 的速度约为 ILU 的六倍。 SIP 的问题在于最优

\begin{center}\small 图 5.6 使用 SIP 求解器将上述 2D 拉普拉斯问题中的 L 1 残差范数降低到 10 − 4 以下所需的迭代次数，作为参数 α 的函数\end{center}
\begin{center}\small 表 5.4 对于在矩形域 X × Y = 10 × 1 上具有狄利克雷边界条件的二维拉普拉斯方程，将归一化 L 1 残差范数降低到 10 − 5 以下所需的迭代次数，在两个方向上均具有均匀网格\end{center}
网格 GS LGS-ADI SIP ICCG MG-GS MG-SIP
8 × 8 74 5 4 4 54 3 16 × 16 293 8 6 6 140 4 32 × 32 1164 18 13 11 242 5 64 × 64 4639 53 38 21 288 6 128 × 128 – 189 139 41 283 6 256 × 256 – – – 82 270 6
α 的值位于可用值范围的末端：当 α 略大于最佳值时，该方法不会收敛。最佳值通常在 0.92 到 0.96 之间。使用 α = 0 是安全的。 92，通常不是最佳的，但它提供的速度大约是标准 ILU 方法的五倍。

一些求解器会受到单元纵横比的影响，因为系数的大小变得不均匀。使用 x = 10 y 的网格使得 A N 和 A S 比 A W 和 A E 大 100 倍（请参阅上一章的示例部分）。为了研究这种效应，我们在矩形区域 X × Y = 10 × 1 上解决了上述拉普拉斯方程问题，在每个方向上使用相同数量的网格节点。表 5.4 显示了减少标准化所需的迭代次数
对于各种求解器，L 1 残差范数低于 10 − 5。 GS 求解器不受影响，但它不再适合 MG 方法的平滑器。与方形网格问题相比，LGS-ADI 和 SIP 求解器的速度要快得多。 ICCG 的表现也稍好一些。 MG-SIP 不受影响，但 MG-GS 显著恶化。

这种行为很典型，也出现在对流扩散问题中（尽管效果不太明显）以及具有小和大单元纵横比的不均匀网格。对恶化的数学解释

\begin{center}\small 表 5.5 对于具有 Neumann 边界条件的 3D Poisson 方程，各种求解器将 L 1 残差范数降低到 10 − 4 以下所需的迭代次数 Grid GS SIP ICCG CGSTAB FMG-GS FMG-SIP\end{center}
8 3 66 27 10 7 10 6 16 3 230 81 19 12 10 6 32 3 882 316 34 21 9 6 64 3 – 1288 54 41 7 6
Brandt (1984) 给出了 GS 和 ILU 性能随着长径比的增加而提高。

最后，我们给出了在 3D 环境中求解具有诺伊曼边界条件的泊松方程的一些结果。 CFD 中的压力和压力校正方程属于这种类型。解出的方程为：

\begin{sourceblock}
\includegraphics[page=170,trim={53.00bp 389.49bp 53.87bp 247.39bp},clip,width=0.92\textwidth,max height=0.38\textheight]{Ferziger4th.pdf}
\end{sourceblock}
其中 x * = x / X 、y * = y / Y 、z * = z / Z 和 X 、 Y 、 Z 是解域的维度。使用 FV 方法对方程进行离散化。域上的源项之和为零，并且在所有边界处指定了诺依曼边界条件（垂直于边界的零梯度）。除了上面介绍的 GS、SIP 和 ICCG 求解器外，我们还使用了不完全 Cholesky 预处理的 CGSTAB 方法。初始解为零。表 5.5 列出了将绝对残差归一化总和减少四个数量级所需的迭代次数。

从这个练习中得出的结论与从具有狄利克雷边界条件的二维问题中得出的结论类似。当需要精确的解时，GS 和 SIP 在精细网格上变得低效；共轭梯度求解器是更好的选择，多重网格方法是最好的。 FMG 策略比直接多重网格更好，其中粗网格上的解为下一个更精细网格提供初始解。使用 ICCG 或 CGSTAB 作为平滑器的 FMG 需要更少的迭代（在最精细的网格上三到四次），但计算时间比 MG-SIP 长。 FMG 原理也可以应用于其他求解器。

\chapter{非稳态问题的求解方法}
\section*{6.1 简介}\phantomsection\addcontentsline{toc}{section}{6.1 简介}
在计算非定常流时，我们需要考虑第四个坐标方向：时间。正如空间坐标一样，时间也必须离散化。我们可以将有限差分精神中的时间“网格”视为时间上的离散点，或者将有限体积视图中的时间“网格”视为“时间体积”。空间和时间坐标之间的主要区别在于影响的方向：而任何空间位置的力都可能（在椭圆问题中）影响其他任何地方的流动，而给定时刻的力只会影响未来的流动 - 没有向后影响。因此，非稳态流动在时间上呈抛物线状。这意味着计算开始后的任何时间都不能对解施加任何条件（边界处除外），这对解策略的选择有很大影响。为了忠实于时间的本质，本质上所有的求解方法都是以逐步或“行进”的方式在时间上前进的。这些方法与应用于常微分方程 (ODE) 初值问题的方法非常相似，因此我们将在下一节中简要回顾这些方法。

\section*{6.2 常微分方程初值问题的求解方法}\phantomsection\addcontentsline{toc}{section}{6.2 常微分方程初值问题的求解方法}
\subsection*{6.2.1 两水平方法}\phantomsection\addcontentsline{toc}{subsection}{6.2.1 两水平方法}
\begin{sourceblock}
\begin{tikzpicture}
\node[inner sep=0,anchor=south west] (image) at (0,0) {\includegraphics[page=171,trim={53.00bp 93.10bp 53.87bp 544.66bp},clip,width=0.92\textwidth,max height=0.38\textheight]{Ferziger4th.pdf}};
\begin{scope}[x={(image.south east)},y={(image.north west)}]
\fill[white] (0.28352,0.53735) rectangle (0.35732,0.95772);
\end{scope}
\end{tikzpicture}
\end{sourceblock}
\begin{sourceblock}
\includegraphics[page=172,trim={53.00bp 522.48bp 53.87bp 52.00bp},clip,width=\textwidth,max height=0.38\textheight]{Ferziger4th.pdf}
\par\vspace{0.45\baselineskip}
\small 图 6.1 f ( t ) 在区间 t 上的时间积分近似（从左到右分别为显式欧拉、隐式欧拉、梯形法则和中点法则）
\end{sourceblock}
基本问题是在初始点后短时间 t 内找到解 φ。 t 1 = t 0 + t , φ 1 处的解可被视为新的初始条件，并且该解可以推进至t 2 = t 1 + t , t 3 = t 2 + t ,。。。等等
最简单的方法可以通过将方程（6.1）从 t n 积分到 t n + 1 = t n + t 来构造：

\begin{sourceblock}
\begin{tikzpicture}
\node[inner sep=0,anchor=south west] (image) at (0,0) {\includegraphics[page=172,trim={53.00bp 377.72bp 53.87bp 246.12bp},clip,width=0.92\textwidth,max height=0.38\textheight]{Ferziger4th.pdf}};
\begin{scope}[x={(image.south east)},y={(image.north west)}]
\fill[white] (0.00000,0.91300) rectangle (0.02298,1.00000);
\fill[white] (0.02725,0.93853) rectangle (0.05543,1.00000);
\fill[white] (0.08138,0.92979) rectangle (0.10556,1.00000);
\end{scope}
\end{tikzpicture}
\end{sourceblock}
其中我们使用简写符号 φ n + 1 = φ( t n + 1 )。这个方程是精确的。然而，在不知道解的情况下无法评估右侧；所以一些近似是必要的。微积分的中值定理保证，如果被积函数在 t n 和 t n + 1 之间的适当点 t = τ 处求值，则积分等于 f (τ, φ(τ)) t 但这没什么用处，因为 τ 是未知的。因此，我们使用一些近似数值求积来评估积分。

下面给出四个相对简单的过程；几何图如图6.1所示。

如果使用初始点处被积函数的值来估计式（6.2）右侧的积分，我们有：

\begin{sourceblock}
\includegraphics[page=172,trim={53.00bp 221.90bp 53.87bp 425.08bp},clip,width=0.92\textwidth,max height=0.38\textheight]{Ferziger4th.pdf}
\end{sourceblock}
这被称为显式或前向欧拉方法。

相反，如果我们使用最终点来估计积分，我们将获得隐式或后向欧拉方法：

\begin{sourceblock}
\includegraphics[page=172,trim={53.00bp 150.10bp 53.87bp 496.81bp},clip,width=0.92\textwidth,max height=0.38\textheight]{Ferziger4th.pdf}
\end{sourceblock}
还可以通过使用区间的中点来获得另一种方法：

\begin{sourceblock}
\begin{tikzpicture}
\node[inner sep=0,anchor=south west] (image) at (0,0) {\includegraphics[page=172,trim={53.00bp 95.85bp 53.87bp 540.41bp},clip,width=0.92\textwidth,max height=0.38\textheight]{Ferziger4th.pdf}};
\begin{scope}[x={(image.south east)},y={(image.north west)}]
\fill[white] (0.28842,0.20482) rectangle (0.35026,0.64960);
\fill[white] (0.35531,0.20482) rectangle (0.38349,0.59170);
\fill[white] (0.38701,0.20482) rectangle (0.42108,0.64324);
\fill[white] (0.42538,0.20482) rectangle (0.45356,0.59170);
\fill[white] (0.46198,0.19578) rectangle (0.47510,0.58300);
\end{scope}
\end{tikzpicture}
\end{sourceblock}
这被称为中点法则，可以被视为广泛使用的求解偏微分方程方法——蛙跳法的基础。

最后，可以使用初始点和最终点之间的直线插值来构造近似值：

\begin{sourceblock}
\begin{tikzpicture}
\node[inner sep=0,anchor=south west] (image) at (0,0) {\includegraphics[page=173,trim={53.00bp 529.67bp 53.87bp 114.73bp},clip,width=0.92\textwidth,max height=0.38\textheight]{Ferziger4th.pdf}};
\begin{scope}[x={(image.south east)},y={(image.north west)}]
\fill[white] (0.37089,0.41260) rectangle (0.39068,0.94480);
\fill[white] (0.20024,0.10856) rectangle (0.26211,0.71941);
\fill[white] (0.26713,0.10856) rectangle (0.29531,0.64029);
\fill[white] (0.29883,0.10856) rectangle (0.33290,0.71113);
\fill[white] (0.33717,0.10856) rectangle (0.36538,0.64029);
\fill[white] (0.76659,0.09614) rectangle (0.77973,0.62787);
\fill[white] (0.78595,0.10856) rectangle (0.80066,0.64029);
\fill[white] (0.93937,0.09154) rectangle (1.00000,0.62328);
\fill[white] (0.37083,0.00000) rectangle (0.39062,0.29807);
\end{scope}
\end{tikzpicture}
\end{sourceblock}
这称为梯形法则，是求解偏微分方程的流行方法——克兰克-尼科尔森法的基础。

总的来说，这些方法被称为两级方法，因为它们仅涉及两次未知值（中点规则可能会或可能不会被视为两级方法，具体取决于采用的进一步近似）。这些方法的分析可以在常微分方程数值解的文本中找到（例如参见Ferziger 1998，或Moin 2010），这里不再重复。

我们将简单回顾一下它们的一些最重要的属性。首先我们注意到，除了第一个方法之外，所有方法都需要除 t = t n 之外的某个点（这是已知解的积分区间的初始点）处的 φ( t ) 值。因此，对于这些方法，如果不进一步近似或迭代就无法计算右侧。因此，第一个方法属于称为显式方法的类，而所有其他方法都是隐式方法。

如果 t 很小，所有方法都会产生良好的解。然而，大步长方法的行为很重要，因为在时间尺度变化很大的问题（包括流体力学中的许多问题）中，目标通常是计算解的缓慢、长期行为，而短时间尺度只是一个麻烦。具有宽时间尺度的问题称为刚性问题，是求解常微分方程时面临的最大困难。因此，当步长很大时，了解方法的行为非常重要。这就提出了稳定性问题。

文献中有许多稳定性的定义。我们将使用一个粗略的定义，如果当基础微分方程的解也是有界时，该方法产生有界解，则该方法称为稳定方法。对于显式欧拉方法，稳定性要求：

\begin{sourceblock}
\includegraphics[page=173,trim={53.00bp 168.49bp 53.87bp 462.17bp},clip,width=0.92\textwidth,max height=0.38\textheight]{Ferziger4th.pdf}
\end{sourceblock}
其中，如果允许 f ( t , φ) 具有复数值，则要求 t ∂ f ( t , φ)/∂φ 限制在以 − 1 为圆心的单位圆内。（必须考虑复数值，因为高阶系统可能具有复特征值。只有实部为零或负数的值才有意义，因为它们会导致有界解。）具有此属性的方法称为条件稳定；对于 f 的实际值，方程（6.7）简化为（参见方程（6.1））：

\begin{sourceblock}
\begin{tikzpicture}
\node[inner sep=0,anchor=south west] (image) at (0,0) {\includegraphics[page=174,trim={53.00bp 583.83bp 53.87bp 50.00bp},clip,width=0.92\textwidth,max height=0.38\textheight]{Ferziger4th.pdf}};
\begin{scope}[x={(image.south east)},y={(image.north west)}]
\fill[white] (0.43233,0.46394) rectangle (0.46397,0.83008);
\fill[white] (0.46541,0.46394) rectangle (0.50205,0.83008);
\fill[white] (0.50223,0.47230) rectangle (0.53699,0.83008);
\fill[white] (0.41768,0.24791) rectangle (0.43083,0.60569);
\fill[white] (0.46478,0.03714) rectangle (0.50301,0.39492);
\end{scope}
\end{tikzpicture}
\end{sourceblock}
上面定义的所有其他方法都是无条件稳定的，即，如果 ∂ f ( t , φ)/∂φ < 0，它们会为任何时间步产生有界解。然而，即使 t 非常大，隐式欧拉方法也往往会产生平滑解，而梯形规则经常会产生小阻尼振荡的解。因此，即使对于非线性方程，隐式欧拉方法也往往表现良好，而梯形规则对于非线性问题可能不稳定。

最后，需要考虑准确性问题。在全球范围内，由于人们可能需要处理的方程种类繁多，因此很难对这个问题说太多。对于单个小步骤，可以使用泰勒级数来表明显式欧拉方法从 t n 处的已知解开始，产生时间 t n + t 处的解，其误差与 (t) 2 成比例。然而，由于计算到某个有限最终时间 t = t 0 + T 所需的步骤数与 t 成反比，并且每一步都会产生误差，因此误差最终将与 t 本身成正比。因此，显式欧拉方法是一阶方法。隐式欧拉方法也是一阶方法，而梯形和中点法则方法的误差与( t ) 2 成比例，因此是二阶方法。可以证明二阶是二级方案可实现的最高阶。

值得注意的是，方法的顺序并不是其准确性的唯一指标。虽然对于足够小的步长，高阶方法确实会比低阶方法具有更小的误差，但同样阶数的两种方法也可能具有相差一个数量级的误差。该顺序仅决定当步长变为零时误差变为零的速率，并且仅在步长变得足够小时之后才确定。 “足够小”既取决于方法又取决于问题，并且不能先验地确定。

当步长足够小时，可以通过比较使用两种不同步长获得的解来估计解中的离散化误差。这种方法称为理查森外推法，已在第 1 章中进行了描述。 3，并且适用于空间和时间离散化误差。还可以通过分析不同阶数的两个方案产生的解的差异来估计误差；这将在第 1 章中讨论。 12.

\subsection*{6.2.2 预测校正器和多点方法}\phantomsection\addcontentsline{toc}{subsection}{6.2.2 预测校正器和多点方法}
我们发现的两级方法的属性非常通用。显式方法非常容易编程，并且每步使用很少的计算机内存和计算时间，但如果时间步长很大，则不稳定。另一方面，隐式方法需要迭代求解才能获得新时间步的值。这使得它们更难编程，并且每个时间步使用更多的计算机内存和时间，但它们更加稳定。 （上述隐式方法是无条件稳定的；并非所有隐式方法都是如此，但它们通常比显式方法更稳定。）有人可能会问是否有可能结合这两种方法中的最佳方法。预测校正方法是实现这一目标的尝试。

已经开发出多种预测校正方法；目前，我们只给出一种众所周知的方法，它通常被称为预测校正方法。在此方法中，使用显式欧拉方法预测新时间步的解：

\begin{sourceblock}
\includegraphics[page=175,trim={53.00bp 485.11bp 53.87bp 160.27bp},clip,width=0.92\textwidth,max height=0.38\textheight]{Ferziger4th.pdf}
\end{sourceblock}
其中 * 表示这不是 t n + 1 时解的最终值。相反，通过应用梯形规则使用 φ ∗ 来修正解
n + 1 计算导数：

\begin{sourceblock}
\begin{tikzpicture}
\node[inner sep=0,anchor=south west] (image) at (0,0) {\includegraphics[page=175,trim={53.00bp 423.05bp 53.87bp 220.17bp},clip,width=0.92\textwidth,max height=0.38\textheight]{Ferziger4th.pdf}};
\begin{scope}[x={(image.south east)},y={(image.north west)}]
\fill[white] (0.00000,0.77443) rectangle (0.12950,1.00000);
\fill[white] (0.37585,0.44328) rectangle (0.39564,0.94764);
\fill[white] (0.20517,0.15445) rectangle (0.26704,0.73429);
\fill[white] (0.27206,0.15445) rectangle (0.30024,0.65925);
\fill[white] (0.30376,0.15445) rectangle (0.33783,0.72600);
\fill[white] (0.34214,0.15445) rectangle (0.37032,0.65925);
\fill[white] (0.76334,0.14311) rectangle (0.77648,0.64747);
\fill[white] (0.78271,0.15445) rectangle (0.79576,0.65925);
\fill[white] (0.92421,0.13831) rectangle (1.00000,0.64311);
\fill[white] (0.37576,0.00000) rectangle (0.39555,0.33421);
\end{scope}
\end{tikzpicture}
\end{sourceblock}
该方法可以证明是二阶精度（梯形规则的精度），但大致具有显式欧拉方法的稳定性。人们可能认为通过迭代校正器可以提高稳定性，但事实证明并非如此，因为只有当 t 足够小时，该迭代过程才会收敛到梯形规则解。

这种预测校正方法属于双级系列，其可能的最高精度是二阶的。对于高阶近似，必须使用更多点的信息。附加点可以是已经计算出数据的点，也可以是t n 和t n + 1 之间的点，严格使用这些点是为了计算方便；前者称为多点法，后者称为龙格-库塔法。在这里，我们将介绍多点方法；龙格-库塔方法将在下一节中介绍。

最著名的多点方法 Adams 方法是通过将多项式与多个时间点的导数拟合而得出的。如果拉格朗日多项式（Ferziger 1998 或 Moin 2010）适合 f ( t n − m , φ n − m ) , f ( t n − m + 1 , φ n − m + 1 ), .。。 ,
f ( t n , φ n )，并将结果用于计算式(6.2)中的积分，得到m+1阶的显式方法；这种类型的方法称为 Adams-Bashforth 方法。对于偏微分方程的求解，仅使用低阶方法。一阶方法是显式欧拉方法，而二阶和三阶方法是：

\begin{sourceblock}
\begin{tikzpicture}
\node[inner sep=0,anchor=south west] (image) at (0,0) {\includegraphics[page=175,trim={53.00bp 143.96bp 53.87bp 496.07bp},clip,width=0.92\textwidth,max height=0.38\textheight]{Ferziger4th.pdf}};
\begin{scope}[x={(image.south east)},y={(image.north west)}]
\fill[white] (0.40746,0.34738) rectangle (0.42060,0.79012);
\fill[white] (0.21278,0.09000) rectangle (0.27465,0.59900);
\fill[white] (0.27967,0.09000) rectangle (0.30785,0.53313);
\fill[white] (0.31137,0.09000) rectangle (0.34544,0.59173);
\fill[white] (0.34974,0.09000) rectangle (0.37792,0.53313);
\fill[white] (0.92439,0.07622) rectangle (1.00000,0.51896);
\end{scope}
\end{tikzpicture}
\end{sourceblock}
and

\begin{sourceblock}
\begin{tikzpicture}
\node[inner sep=0,anchor=south west] (image) at (0,0) {\includegraphics[page=175,trim={53.00bp 90.82bp 53.87bp 549.20bp},clip,width=0.92\textwidth,max height=0.38\textheight]{Ferziger4th.pdf}};
\begin{scope}[x={(image.south east)},y={(image.north west)}]
\fill[white] (0.23110,0.34763) rectangle (0.24424,0.79020);
\fill[white] (0.03645,0.09035) rectangle (0.09829,0.59877);
\fill[white] (0.10331,0.09035) rectangle (0.13152,0.53292);
\fill[white] (0.13501,0.09035) rectangle (0.16911,0.59188);
\fill[white] (0.17341,0.09035) rectangle (0.20159,0.53292);
\fill[white] (0.87886,0.09035) rectangle (0.89191,0.53292);
\fill[white] (0.92424,0.07619) rectangle (1.00000,0.51876);
\end{scope}
\end{tikzpicture}
\end{sourceblock}
如果 t n + 1 处的数据包含在插值多项式中，则可以获得隐式方法，称为 Adams-Moulton 方法。一阶方法为隐式欧拉法，二阶方法为梯形法则，三阶方法为：

\begin{sourceblock}
\begin{tikzpicture}
\node[inner sep=0,anchor=south west] (image) at (0,0) {\includegraphics[page=176,trim={53.00bp 541.64bp 53.87bp 98.38bp},clip,width=0.92\textwidth,max height=0.38\textheight]{Ferziger4th.pdf}};
\begin{scope}[x={(image.south east)},y={(image.north west)}]
\fill[white] (0.25212,0.34724) rectangle (0.26526,0.79020);
\fill[white] (0.05744,0.09035) rectangle (0.11931,0.59877);
\fill[white] (0.12433,0.09035) rectangle (0.15251,0.53292);
\fill[white] (0.15603,0.09035) rectangle (0.19011,0.59188);
\fill[white] (0.19438,0.09035) rectangle (0.22256,0.53292);
\fill[white] (0.85802,0.09035) rectangle (0.87107,0.53292);
\fill[white] (0.92439,0.07619) rectangle (1.00000,0.51876);
\end{scope}
\end{tikzpicture}
\end{sourceblock}
常见的方法使用 ( m − 1 ) 阶 Adams-Bashforth 方法作为预测器，使用 m 阶 Adams-Moulton 方法作为校正器。因此，可以获得任何阶的预测校正方法。

多点方法的优点是相对容易构造任意阶的方法。这些方法也易于使用和编程。最后一个优点是，它们每个时间步只需要对导数进行一次评估（这可能非常复杂，特别是在涉及偏微分方程的应用中），这使得它们相对便宜。 （f ( t , φ( t )) 的值需要多次，但一旦计算出来就可以存储，因此每个时间步只需要一次评估。）它们的主要缺点是，因为它们需要当前点之前的许多点的数据，所以不能仅使用初始时间点的数据开始。必须使用其他方法来开始计算。一种方法是使用较小的步长和较低阶的方法（以便达到所需的精度），并随着更多数据的可用而缓慢增加阶数。

这些方法是许多精确的常微分方程求解器的基础。在这些求解器中，误差估计器用于确定每一步解的准确性。如果解不够准确，则方法的阶数会增加到程序允许的最大阶数。另一方面，如果解比所需的准确得多，则可以减少方法的阶数以节省计算时间。由于多点方法中的步长很难改变，因此仅当已经达到该方法的最大阶数时才执行此操作。 1
由于多点方法使用来自多个时间步的数据，因此它们可能会产生非物理解。由于篇幅不允许在此进行分析，但值得注意的是，多点方法的不稳定性通常是由于非物理解造成的。通过仔细选择起始方法，这些可能会被抑制，但不能完全抑制。这些方法通常会在一段时间内提供准确的解，然后随着解的非物理部分的增长而开始表现不佳。此问题的常见补救措施是经常重新启动该方法，这是一种有效的技巧，但可能会降低方案的准确性和/或效率。

值得一提的一种特殊的多点（三级）方案是蛙跳法。它本质上是将中点规则积分应用于大小为 2 t 的时间间隔：

\footnotetext[1]{上面使用的正交近似假设时间步长一致；如果允许时间步长变化，则与不同时间级别的函数值相乘的系数将成为步长大小的复杂函数，正如我们在第 1 章中看到的那样。 3 空间有限差分。}
\begin{sourceblock}
\includegraphics[page=177,trim={53.00bp 592.95bp 53.87bp 54.94bp},clip,width=0.92\textwidth,max height=0.38\textheight]{Ferziger4th.pdf}
\end{sourceblock}
尽管在 2 t 宽的间隔上进行积分，但行进步骤为 t，因此积分间隔重叠。该方案在过去已被广泛使用，其一些功能将在第 2 节中描述。 6.3.1.2.

\subsection*{6.2.3 龙格-库塔方法}\phantomsection\addcontentsline{toc}{subsection}{6.2.3 龙格-库塔方法}
通过使用 t n 和 t n + 1 之间的点而不是更早的点可以克服启动多点方法的困难。这种方法称为龙格-库塔方法，在流体力学应用中非常流行。这些方法可以以系统的方式推导出任何精度级别的方法；我们在这里展示了三个级别。在每种情况下，方案的水平也反映了其准确性。

一阶方法实际上就是上面描述的显式欧拉方法，但二阶龙格-库塔方法（RK2）由两个步骤组成。第一个可以被视为基于显式欧拉方法的半步预测器；它后面是一个中点规则校正器，使该方法成为二阶：

\begin{sourceblock}
\begin{tikzpicture}
\node[inner sep=0,anchor=south west] (image) at (0,0) {\includegraphics[page=177,trim={53.00bp 315.12bp 53.87bp 293.67bp},clip,width=0.92\textwidth,max height=0.38\textheight]{Ferziger4th.pdf}};
\begin{scope}[x={(image.south east)},y={(image.north west)}]
\fill[white] (0.29002,0.10671) rectangle (0.35185,0.33827);
\fill[white] (0.35690,0.10671) rectangle (0.38508,0.30828);
\fill[white] (0.38860,0.10671) rectangle (0.42268,0.33514);
\fill[white] (0.42698,0.10671) rectangle (0.45516,0.30828);
\fill[white] (0.48108,0.10201) rectangle (0.49423,0.30357);
\fill[white] (0.50367,0.10201) rectangle (0.51678,0.30357);
\fill[white] (0.69783,0.10671) rectangle (0.71089,0.30828);
\fill[white] (0.92439,0.10026) rectangle (1.00000,0.30183);
\end{scope}
\end{tikzpicture}
\end{sourceblock}
该方法使用方便，并且是自启动的，即除了微分方程本身所需的初始条件外，不需要任何数据。事实上，它在很多方面与上述预测校正方法非常相似。然而，Durran (2010) 表明 RK2 方案的幅度因子大于 1，因此会放大误差。

6.2.3.1 三阶龙格-库塔方法
三阶龙格-库塔方法 (RK3) 正在气象模拟代码中得到应用，它们取代了上一节中描述的蛙跳方法。与其他高阶龙格-库塔格式一样，三阶格式并不是唯一的。 Heun 方法是最著名的版本之一。它由三个步骤组成；第一个是超过 1/3 时间步长的显式欧拉预测器：

\begin{sourceblock}
\includegraphics[page=177,trim={53.00bp 102.48bp 53.87bp 535.27bp},clip,width=0.92\textwidth,max height=0.38\textheight]{Ferziger4th.pdf}
\end{sourceblock}
接下来是应用于 2/3 时间步长的中点规则近似（即，以 t + t 为中心）

\footnotetext[3]{):}
\begin{sourceblock}
\begin{tikzpicture}
\node[inner sep=0,anchor=south west] (image) at (0,0) {\includegraphics[page=178,trim={53.00bp 584.32bp 53.87bp 53.43bp},clip,width=0.92\textwidth,max height=0.38\textheight]{Ferziger4th.pdf}};
\begin{scope}[x={(image.south east)},y={(image.north west)}]
\fill[white] (0.27753,0.30468) rectangle (0.31206,0.77246);
\fill[white] (0.29483,0.18387) rectangle (0.32743,0.49243);
\fill[white] (0.32623,0.29729) rectangle (0.33853,0.52906);
\end{scope}
\end{tikzpicture}
\end{sourceblock}
最后一步可以被视为整个时间步长的梯形规则近似，其中时间步长结束时的被积函数值是通过 t n 和 t n + 2 处的值进行线性外推获得的

\begin{sourceblock}
\begin{tikzpicture}
\node[inner sep=0,anchor=south west] (image) at (0,0) {\includegraphics[page=178,trim={53.00bp 502.34bp 53.87bp 111.30bp},clip,width=0.92\textwidth,max height=0.38\textheight]{Ferziger4th.pdf}};
\begin{scope}[x={(image.south east)},y={(image.north west)}]
\fill[white] (0.00000,0.98438) rectangle (0.05988,1.00000);
\fill[white] (0.06418,0.98438) rectangle (0.17212,1.00000);
\fill[white] (0.17645,0.98438) rectangle (0.21786,1.00000);
\fill[white] (0.22217,0.98438) rectangle (0.29110,1.00000);
\fill[white] (0.29543,0.98438) rectangle (0.32517,1.00000);
\fill[white] (0.32950,0.98438) rectangle (0.37092,1.00000);
\fill[white] (0.37522,0.98438) rectangle (0.49269,1.00000);
\fill[white] (0.49702,0.98438) rectangle (0.52349,1.00000);
\fill[white] (0.52779,0.98438) rectangle (0.56920,1.00000);
\fill[white] (0.57350,0.98438) rectangle (0.62156,1.00000);
\fill[white] (0.62589,0.98438) rectangle (0.65567,1.00000);
\fill[white] (0.65994,0.98438) rectangle (0.70135,1.00000);
\fill[white] (0.70565,0.98438) rectangle (0.76373,1.00000);
\fill[white] (0.76803,0.98438) rectangle (0.82111,1.00000);
\fill[white] (0.82541,0.98438) rectangle (0.85020,1.00000);
\fill[white] (0.85447,0.98438) rectangle (0.96244,1.00000);
\fill[white] (0.96680,0.98438) rectangle (1.00000,1.00000);
\fill[white] (0.00000,0.75676) rectangle (0.07233,0.97714);
\fill[white] (0.07504,0.75676) rectangle (0.23741,0.97714);
\fill[white] (0.24018,0.75676) rectangle (0.33817,0.97714);
\fill[white] (0.34087,0.75676) rectangle (0.42144,0.97714);
\fill[white] (0.42415,0.75676) rectangle (0.45059,0.97714);
\fill[white] (0.45332,0.74324) rectangle (0.47693,0.97905);
\fill[white] (0.48208,0.75676) rectangle (0.53014,0.97714);
\fill[white] (0.53284,0.72476) rectangle (0.57377,0.97905);
\fill[white] (0.57254,0.78610) rectangle (0.58484,0.91143);
\fill[white] (0.20289,0.15962) rectangle (0.21603,0.37981);
\fill[white] (0.21747,0.14267) rectangle (0.29149,0.38495);
\fill[white] (0.29164,0.16476) rectangle (0.36620,0.41771);
\fill[white] (0.36971,0.16476) rectangle (0.42487,0.51086);
\end{scope}
\end{tikzpicture}
\end{sourceblock}
导致：

\begin{sourceblock}
\begin{tikzpicture}
\node[inner sep=0,anchor=south west] (image) at (0,0) {\includegraphics[page=178,trim={53.00bp 461.52bp 53.87bp 178.75bp},clip,width=0.92\textwidth,max height=0.38\textheight]{Ferziger4th.pdf}};
\begin{scope}[x={(image.south east)},y={(image.north west)}]
\fill[white] (0.00000,0.86896) rectangle (0.09233,1.00000);
\fill[white] (0.09501,0.86896) rectangle (0.13146,1.00000);
\fill[white] (0.38800,0.49633) rectangle (0.40114,0.94318);
\fill[white] (0.19332,0.23657) rectangle (0.25519,0.75029);
\fill[white] (0.26021,0.23657) rectangle (0.28839,0.68342);
\fill[white] (0.29191,0.23657) rectangle (0.32598,0.74295);
\fill[white] (0.33029,0.23657) rectangle (0.35847,0.68342);
\fill[white] (0.37395,0.00000) rectangle (0.39374,0.39583);
\fill[white] (0.79453,0.23657) rectangle (0.80758,0.68342);
\fill[white] (0.92439,0.22227) rectangle (1.00000,0.66950);
\end{scope}
\end{tikzpicture}
\end{sourceblock}
如上所述，三阶龙格-库塔格式并不是唯一的。接下来描述地球物理应用中使用的另一个版本。

Wicker 和 Skamarock (2002) 提出了一种用于使用时间分割的代码的新颖方案。 2 他们的方案已在天气研究和预报 (WRF) 模型的高级研究 (ARF) 版本中实施（Skamarock 和 Klemp 2008）。对于“慢速或低频”（气象学上重要的）运动模式，它已被证明可以轻松适应稳定的时间分割，并且对于振荡和阻尼运动模式都具有出色的稳定性。其他方案用于声波和重力波运动。

Wicker 和 Skamarock 积分需要三个步骤才能从 t n 推进到 t n + 1 = t n + t：

\begin{sourceblock}
\begin{tikzpicture}
\node[inner sep=0,anchor=south west] (image) at (0,0) {\includegraphics[page=178,trim={53.00bp 233.66bp 53.87bp 341.55bp},clip,width=0.92\textwidth,max height=0.38\textheight]{Ferziger4th.pdf}};
\begin{scope}[x={(image.south east)},y={(image.north west)}]
\fill[white] (0.00000,0.95304) rectangle (0.05074,1.00000);
\fill[white] (0.05573,0.96569) rectangle (0.08391,1.00000);
\fill[white] (0.08743,0.95381) rectangle (0.11107,1.00000);
\fill[white] (0.11537,0.96569) rectangle (0.14355,1.00000);
\fill[white] (0.16947,0.96162) rectangle (0.19365,1.00000);
\fill[white] (0.28502,0.41350) rectangle (0.33050,0.55966);
\fill[white] (0.30232,0.37589) rectangle (0.33492,0.47223);
\fill[white] (0.33371,0.41131) rectangle (0.34602,0.48367);
\fill[white] (0.29149,0.05466) rectangle (0.35185,0.19444);
\fill[white] (0.35690,0.06730) rectangle (0.38508,0.19444);
\fill[white] (0.38860,0.05543) rectangle (0.42117,0.19444);
\fill[white] (0.42544,0.06730) rectangle (0.45362,0.19444);
\fill[white] (0.47958,0.06434) rectangle (0.49272,0.19158);
\fill[white] (0.50217,0.06434) rectangle (0.51531,0.19158);
\fill[white] (0.69636,0.06730) rectangle (0.70941,0.19444);
\end{scope}
\end{tikzpicture}
\end{sourceblock}
Purser (2007) 表明，这不是真正的 Runge-Kutta 方案，虽然它对于线性方程具有三阶精度，但对于非线性方程仅具有二阶精度（尽管与 RK2 方案相比，测试中的误差要小得多）。

2 在许多地球物理应用中，某些运动（例如声波）相对于其他运动（例如风或洋流）移动得非常快。那么，实际分割计算并对系统的快速部分和慢速部分进行单独的计算是有用的（Klemp 等人，2007 年，或 Blumberg 和 Mellor，1987 年）。

6.2.3.2 四阶龙格-库塔方法
从历史上看，最流行的龙格-库塔方法是四阶版本（RK4）。该方法的前两个步骤在 t n + 1 处使用显式欧拉预测器和隐式欧拉校正器
2.接下来是整个步骤的中点规则预测器和赋予该方法第四阶的辛普森规则最终校正器。方法是：

\begin{sourceblock}
\begin{tikzpicture}
\node[inner sep=0,anchor=south west] (image) at (0,0) {\includegraphics[page=179,trim={53.00bp 387.59bp 53.87bp 136.10bp},clip,width=0.92\textwidth,max height=0.38\textheight]{Ferziger4th.pdf}};
\begin{scope}[x={(image.south east)},y={(image.north west)}]
\fill[white] (0.00000,0.96104) rectangle (0.05074,1.00000);
\fill[white] (0.05341,0.96104) rectangle (0.14809,1.00000);
\fill[white] (0.15080,0.96104) rectangle (0.18391,1.00000);
\fill[white] (0.28418,0.62562) rectangle (0.32965,0.71892);
\fill[white] (0.30147,0.60161) rectangle (0.33408,0.66311);
\fill[white] (0.33290,0.62422) rectangle (0.34520,0.67041);
\fill[white] (0.28992,0.40463) rectangle (0.32445,0.49793);
\fill[white] (0.30722,0.38821) rectangle (0.35029,0.45019);
\fill[white] (0.35531,0.40463) rectangle (0.38349,0.48578);
\fill[white] (0.38701,0.40463) rectangle (0.42108,0.49660);
\fill[white] (0.42538,0.40463) rectangle (0.45356,0.48578);
\fill[white] (0.47952,0.40274) rectangle (0.49266,0.48396);
\fill[white] (0.50208,0.40274) rectangle (0.51522,0.48396);
\fill[white] (0.69627,0.40463) rectangle (0.71101,0.48578);
\fill[white] (0.92439,0.40204) rectangle (1.00000,0.48326);
\fill[white] (0.38905,0.22022) rectangle (0.40220,0.30144);
\fill[white] (0.19438,0.17311) rectangle (0.25624,0.26290);
\fill[white] (0.26126,0.17311) rectangle (0.28944,0.25426);
\fill[white] (0.29296,0.17311) rectangle (0.32704,0.26157);
\fill[white] (0.33134,0.17311) rectangle (0.35952,0.25426);
\fill[white] (0.37504,0.12081) rectangle (0.39483,0.20204);
\fill[white] (0.77835,0.17311) rectangle (0.80653,0.25426);
\fill[white] (0.92433,0.11183) rectangle (1.00000,0.19298);
\fill[white] (0.36328,0.04044) rectangle (0.38307,0.12159);
\fill[white] (0.38998,0.04114) rectangle (0.40310,0.12229);
\fill[white] (0.42589,0.02871) rectangle (0.46683,0.12229);
\fill[white] (0.46562,0.05132) rectangle (0.47792,0.09751);
\end{scope}
\end{tikzpicture}
\end{sourceblock}
已经开发了该方法的许多变体。特别是，有几种方法添加第四或第五阶的第五步骤，以允许估计误差，从而实现自动误差控制的可能性。

从上面给出的方法中很容易看出，n 阶龙格-库塔方法要求每个时间步对导数求值 n 次，使得这些方法比可比阶数的多点方法更昂贵。在部分补偿中，给定阶数的龙格-库塔法比相同阶数的多点法更准确（即误差项的系数更小）并且更稳定。 Purser (2007) 回顾了生成 RK 方法的系统策略，其中包括用于显示系数的 Butcher 数组以及使每种方法一致所需的必要约束（另请参见 \url{http://en.wikipedia.org/wiki/Butcher_tableau} 或 Butcher 2008）。

\subsection*{6.2.4 其他方法}\phantomsection\addcontentsline{toc}{subsection}{6.2.4 其他方法}
还有许多其他可能性可以及时推进解。我们仅描述一种在工程 CFD 中广泛使用的多点全隐式方案。

通过在以 t n + 1 为中心的时间间隔 t（即从 t n + 1 − t / 2 到 t n + 1 + t / 2 ）上积分并将中点规则应用到方程（ 6.2 ）的左侧和右侧，可以构造完全隐式的三级二阶方案。 t n + 1 处的时间导数通过对强制通过三个时间级别 t n − 1 、t n 和 t n + 1 处的解的抛物线求微分来近似，从而得出：

\begin{sourceblock}
\begin{tikzpicture}
\node[inner sep=0,anchor=south west] (image) at (0,0) {\includegraphics[page=180,trim={53.00bp 580.66bp 53.87bp 50.00bp},clip,width=0.92\textwidth,max height=0.38\textheight]{Ferziger4th.pdf}};
\begin{scope}[x={(image.south east)},y={(image.north west)}]
\fill[white] (0.30045,0.48732) rectangle (0.33817,0.82356);
\end{scope}
\end{tikzpicture}
\end{sourceblock}
右侧仅在 t n + 1 时计算。由于方程两边都是积分区间中心处被积函数的二阶近似，因此应用中点法则来近似积分意味着只需将两边乘以 t — 这一步骤可以省略。这导致了以下时间提前方法：

\begin{sourceblock}
\includegraphics[page=180,trim={53.00bp 476.23bp 53.87bp 161.88bp},clip,width=0.92\textwidth,max height=0.38\textheight]{Ferziger4th.pdf}
\end{sourceblock}
该方案是完全隐式的，因为 f 仅在新的时间级别进行评估。它是二阶的，非常容易实现，但由于它是隐式的，因此需要在每个时间步求解代数方程组。

该方案在CFD工程应用中得到广泛应用；所有主要的商业 CFD 包以及公共代码都提供了这种时间推进方案。我们将在第 1 节中更详细地讨论其原因。 6.3.2.4，但由于方案的流行，我们在这里也给出了它的可变时间步长的版本。与其他方法相比，存在概念上的差异，我们需要强调这一点。

按照图 6.2 中非均匀时间步长的表示法，通过在三个时间水平上通过 φ 值拟合抛物线，可以获得以下 t n + 1 处一阶导数的表达式：

\begin{sourceblock}
\includegraphics[page=180,trim={53.00bp 253.71bp 53.87bp 331.58bp},clip,width=\textwidth,max height=0.38\textheight]{Ferziger4th.pdf}
\par\vspace{0.45\baselineskip}
\small 图 6.2 使用中点法则和非均匀时间步长逼近 φ 在 t n + 1 处的时间导数和 f ( t ) 的时间积分
\end{sourceblock}
\begin{sourceblock}
\includegraphics[page=181,trim={53.00bp 583.22bp 53.87bp 53.41bp},clip,width=0.92\textwidth,max height=0.38\textheight]{Ferziger4th.pdf}
\end{sourceblock}
值得注意的是，计算解的时间级别的位置（ t n − 1 、 t n 和 t n + 1 ）现在不同了：它不是在用户指定的积分间隔的末尾（与所有其他方案的情况一样），而是在它们的中心。因此，

\begin{sourceblock}
\includegraphics[page=181,trim={53.00bp 494.20bp 53.87bp 136.22bp},clip,width=0.92\textwidth,max height=0.38\textheight]{Ferziger4th.pdf}
\end{sourceblock}
另请注意，该方案是多级类型，不能在计算开始时的第一个时间步中使用；人们必须使用不同的方案才能开始。通常从隐式欧拉方案和较小的时间步长开始。

\section*{6.3 通用输运方程的应用}\phantomsection\addcontentsline{toc}{section}{6.3 通用输运方程的应用}
接下来我们考虑将上面给出的一些方法应用于通用传输方程（1.28）。在章节中。参照图3和图4，讨论了稳态问题的对流和扩散通量以及源项的离散化。对于非定常流，这些项可以以相同的方式处理；然而，必须回答何时评估通量和源的问题。

如果将守恒方程重写为类似于常微分方程（6.1）的形式，例如：

\begin{sourceblock}
\includegraphics[page=181,trim={53.00bp 269.44bp 53.87bp 368.41bp},clip,width=0.92\textwidth,max height=0.38\textheight]{Ferziger4th.pdf}
\end{sourceblock}
可以使用任何时间积分方法。函数 f ( t , φ) 表示对流项、扩散项和源项之和，所有这些项现在都出现在方程的右侧。因为这些项是未知的，所以我们必须使用上面介绍的正交近似之一。使用第 1 章中介绍的方法之一对对流、扩散和源项进行离散化。 3或4，在一个或多个时间级别。如果使用显式方法进行时间积分，则只需在解已知的时间评估这些项，以便可以计算它们。对于隐式方法，需要在新的时间水平上对上式右侧进行离散化，但其解尚未知。因此，必须求解与稳态问题所获得的代数方程组不同的代数方程组。我们将分析一些常见方案应用于下面的一维问题时的属性；一维和二维问题的解将在示例部分讨论。

\subsection*{6.3.1 显式方法}\phantomsection\addcontentsline{toc}{subsection}{6.3.1 显式方法}
最简单的方法是显式欧拉法，其中使用 t n 处的已知值来评估所有通量和源。在 CV 或网格点的方程中，新时间级别的唯一未知数是该节点的值；邻居值都是在较早的时间级别评估的。因此，我们可以明确地计算每个节点处的未知数的新值。

为了研究显式欧拉和其他简单格式的性质，我们考虑具有恒定速度、恒定流体性质和无源项的方程（6.29）的一维版本：

\begin{sourceblock}
\begin{tikzpicture}
\node[inner sep=0,anchor=south west] (image) at (0,0) {\includegraphics[page=182,trim={53.00bp 436.34bp 53.87bp 200.54bp},clip,width=0.92\textwidth,max height=0.38\textheight]{Ferziger4th.pdf}};
\begin{scope}[x={(image.south east)},y={(image.north west)}]
\fill[white] (0.00000,0.83322) rectangle (0.04740,1.00000);
\fill[white] (0.05011,0.83322) rectangle (0.08487,1.00000);
\fill[white] (0.08758,0.83322) rectangle (0.17056,1.00000);
\fill[white] (0.17326,0.83322) rectangle (0.25296,1.00000);
\end{scope}
\end{tikzpicture}
\end{sourceblock}
该方程经常在文献中用作纳维-斯托克斯方程的模型方程。它是方程（3.55）的时间相关版本，用于说明稳态问题的方法。与该方程一样，它假设重要的平衡存在于平流和流向扩散之间，这种平衡在实际流动中很少出现。因此，在将从该方程中学到的知识扩展到纳维-斯托克斯方程时必须小心。尽管有这个重要的缺点，我们还是可以通过考虑方程（6.30）学到一些东西。

我们首先假设空间导数是使用 CDS 近似的，并且网格在 x 方向上是均匀的。在这种情况下，FD 和 FV 离散化会产生相同的代数方程。新变量值 φ n + 1
i , is:

\begin{sourceblock}
\begin{tikzpicture}
\node[inner sep=0,anchor=south west] (image) at (0,0) {\includegraphics[page=182,trim={53.00bp 269.78bp 53.87bp 366.24bp},clip,width=0.92\textwidth,max height=0.38\textheight]{Ferziger4th.pdf}};
\begin{scope}[x={(image.south east)},y={(image.north west)}]
\fill[white] (0.08087,0.28718) rectangle (0.14274,0.73240);
\fill[white] (0.09805,0.20651) rectangle (0.10869,0.49137);
\fill[white] (0.14776,0.28718) rectangle (0.17594,0.67098);
\fill[white] (0.17946,0.28718) rectangle (0.21353,0.72178);
\end{scope}
\end{tikzpicture}
\end{sourceblock}
可以重写：

\begin{sourceblock}
\begin{tikzpicture}
\node[inner sep=0,anchor=south west] (image) at (0,0) {\includegraphics[page=182,trim={53.00bp 221.62bp 53.87bp 421.15bp},clip,width=0.92\textwidth,max height=0.38\textheight]{Ferziger4th.pdf}};
\begin{scope}[x={(image.south east)},y={(image.north west)}]
\fill[white] (0.16418,0.16645) rectangle (0.22602,0.74069);
\fill[white] (0.18132,0.06205) rectangle (0.19197,0.42918);
\fill[white] (0.23107,0.16645) rectangle (0.25925,0.66110);
\fill[white] (0.26277,0.15062) rectangle (0.29371,0.66110);
\fill[white] (0.29555,0.15062) rectangle (0.37392,0.66110);
\fill[white] (0.37408,0.16645) rectangle (0.40815,0.72657);
\fill[white] (0.39125,0.06889) rectangle (0.40189,0.43603);
\fill[white] (0.41248,0.16645) rectangle (0.44066,0.66110);
\fill[white] (0.76033,0.16645) rectangle (0.79441,0.72657);
\fill[white] (0.77750,0.06590) rectangle (0.81699,0.44373);
\fill[white] (0.82202,0.16645) rectangle (0.83675,0.66110);
\fill[white] (0.92424,0.15062) rectangle (1.00000,0.64527);
\end{scope}
\end{tikzpicture}
\end{sourceblock}
我们引入了无量纲参数：

\begin{sourceblock}
\includegraphics[page=182,trim={53.00bp 161.67bp 53.87bp 476.18bp},clip,width=0.92\textwidth,max height=0.38\textheight]{Ferziger4th.pdf}
\end{sourceblock}
参数 d 是时间步长 t 与特征扩散时间 ρ( x ) 2 / 的比值，其大致是扰动扩散传播距离 x 所需的时间。第二个量 (c) 是时间步长 t 与特征对流时间 x / u 的比率，即扰动在距离 x 上对流所需的时间。这个比率称为库朗数，是计算流体动力学的关键参数之一。 3
如果 φ 是温度（一种可能性），则方程（6.32）需要满足几个条件。借助扩散，旧时间水平下三个点 x i - 1 、x i 和 x i + 1 中任意一个的温度升高都会导致新时间水平下 x i 点的温度升高。对于对流的点 x i − 1 和 x i 也可以这样说，假设 u > 0。如果 φ 表示物质的浓度，则它不应为负。

\begin{sourceblock}
\begin{tikzpicture}
\node[inner sep=0,anchor=south west] (image) at (0,0) {\includegraphics[page=183,trim={53.00bp 497.32bp 53.87bp 150.58bp},clip,width=0.92\textwidth,max height=0.38\textheight]{Ferziger4th.pdf}};
\begin{scope}[x={(image.south east)},y={(image.north west)}]
\fill[white] (0.00000,0.83388) rectangle (0.08403,1.00000);
\fill[white] (0.08674,0.83388) rectangle (0.12980,1.00000);
\fill[white] (0.13251,0.83388) rectangle (0.16559,1.00000);
\fill[white] (0.16830,0.83388) rectangle (0.28021,1.00000);
\fill[white] (0.03528,0.17818) rectangle (0.08668,0.81250);
\fill[white] (0.09080,0.17818) rectangle (0.22045,0.81250);
\fill[white] (0.22451,0.17818) rectangle (0.27426,0.81250);
\fill[white] (0.27835,0.17818) rectangle (0.34641,0.81250);
\fill[white] (0.35050,0.17818) rectangle (0.38027,0.81250);
\fill[white] (0.38436,0.17818) rectangle (0.42577,0.81250);
\fill[white] (0.42986,0.17818) rectangle (0.57203,0.81250);
\fill[white] (0.57618,0.17818) rectangle (0.60595,0.81250);
\fill[white] (0.60998,0.19846) rectangle (0.67182,0.93421);
\fill[white] (0.74674,0.06195) rectangle (0.78626,0.54605);
\fill[white] (0.79705,0.17818) rectangle (0.82517,0.81250);
\fill[white] (0.82923,0.17818) rectangle (0.95110,0.81250);
\fill[white] (0.95519,0.17818) rectangle (1.00000,0.81250);
\fill[white] (0.00000,0.00000) rectangle (0.09729,0.15680);
\fill[white] (0.10244,0.00000) rectangle (0.20686,0.15680);
\fill[white] (0.21206,0.00000) rectangle (0.29675,0.15680);
\fill[white] (0.30186,0.00000) rectangle (0.35991,0.15680);
\fill[white] (0.36499,0.00000) rectangle (0.39645,0.15680);
\fill[white] (0.40153,0.00000) rectangle (0.42965,0.15680);
\fill[white] (0.43474,0.00000) rectangle (0.53774,0.15680);
\fill[white] (0.54283,0.00000) rectangle (0.63248,0.15680);
\fill[white] (0.63759,0.00000) rectangle (0.68565,0.15680);
\fill[white] (0.69080,0.00000) rectangle (0.80211,0.15680);
\fill[white] (0.80725,0.00000) rectangle (0.82535,0.15680);
\fill[white] (0.83044,0.00000) rectangle (0.89681,0.15680);
\fill[white] (0.90192,0.00000) rectangle (1.00000,0.15680);
\end{scope}
\end{tikzpicture}
\end{sourceblock}
方程（6.32）中的i+1可能变为负值，这应该提醒我们可能出现的问题并需要进行更详细的分析。在可能的范围内，我们希望这种分析能够模仿用于普通微分方程的分析。一种简单的方法是由冯·诺依曼发明的，该方法就是以他的名字命名的。他认为边界条件很少是问题的原因（也有与这里无关的例外），所以为什么不完全忽略它们呢？如果这样做了，分析就会简化。由于这种分析方法基本上可以应用于本章讨论的所有方法，因此我们将对其进行详细描述；有关更多详细信息，请参阅 Moin (2010) 或 Fletcher (1991)。

本质上，这个想法可以如下得出。式(6.32)的集合可以写成矩阵形式：

\begin{sourceblock}
\begin{tikzpicture}
\node[inner sep=0,anchor=south west] (image) at (0,0) {\includegraphics[page=183,trim={53.00bp 355.91bp 53.87bp 291.76bp},clip,width=0.92\textwidth,max height=0.38\textheight]{Ferziger4th.pdf}};
\begin{scope}[x={(image.south east)},y={(image.north west)}]
\fill[white] (0.00000,0.71197) rectangle (0.02746,1.00000);
\fill[white] (0.03014,0.71197) rectangle (0.11320,1.00000);
\fill[white] (0.11585,0.71197) rectangle (0.18728,1.00000);
\end{scope}
\end{tikzpicture}
\end{sourceblock}
式中，三对角矩阵A的元素可以通过式(6.32)求得。该方程根据前一步的解给出了新步骤的解。因此，可以通过将初始解 φ 0 与矩阵 A 重复相乘来获得 t n + 1 处的解。问题是：以任何方便的方式测量的连续时间步长（对于不变边界条件）的解之间的差异是否随着 n 的增加而增加、减少或保持不变？例如，一种衡量标准是：

\begin{sourceblock}
\begin{tikzpicture}
\node[inner sep=0,anchor=south west] (image) at (0,0) {\includegraphics[page=183,trim={53.00bp 218.92bp 53.87bp 417.63bp},clip,width=0.92\textwidth,max height=0.38\textheight]{Ferziger4th.pdf}};
\begin{scope}[x={(image.south east)},y={(image.north west)}]
\fill[white] (0.24743,0.42954) rectangle (0.26532,0.82055);
\fill[white] (0.27050,0.42954) rectangle (0.29868,0.82055);
\fill[white] (0.30220,0.42954) rectangle (0.35344,0.87225);
\fill[white] (0.35771,0.42954) rectangle (0.46827,0.87868);
\fill[white] (0.47176,0.42954) rectangle (0.49994,0.82055);
\end{scope}
\end{tikzpicture}
\end{sourceblock}
微分方程要求这个量通过扩散作用随时间减小。如果边界条件不变，最终将得到稳态解。当然，我们希望数值方法能够保留精确方程的这一性质。

这个问题与矩阵 A 的特征值密切相关。如果其中一些大于1，不难看出ε会随着n的增长而增长，而如果全部小于1，则ε会衰减。通常，矩阵的特征值很难估计，对于比这更复杂的问题，我们会遇到严重的麻烦。该问题的优点是，由于矩阵的每条对角线都是常数，因此很容易找到特征向量。它们可以用正弦和余弦表示，但使用复指数形式更简单：

\begin{sourceblock}
\includegraphics[page=184,trim={53.00bp 556.84bp 53.87bp 88.54bp},clip,width=0.92\textwidth,max height=0.38\textheight]{Ferziger4th.pdf}
\end{sourceblock}
其中 i = √ − 1 且 α 是可以任意选择的波数；下面将讨论这个选择。如果将式（ 6.36 ）代入式（ 6.32 ），则复指数项 e i α j 对于每一项都是公共的，可以将其删除，得到特征值 σ 的显式表达式：

\begin{sourceblock}
\includegraphics[page=184,trim={53.00bp 473.48bp 53.87bp 178.33bp},clip,width=0.92\textwidth,max height=0.38\textheight]{Ferziger4th.pdf}
\end{sourceblock}
这个数量的大小才是重要的。因为复数的大小是实部和虚部的平方和，所以我们有：

\begin{sourceblock}
\includegraphics[page=184,trim={53.00bp 413.70bp 53.87bp 233.99bp},clip,width=0.92\textwidth,max height=0.38\textheight]{Ferziger4th.pdf}
\end{sourceblock}
我们现在研究 σ 2 小于 1 的条件。

由于 σ 的表达式中有两个独立参数，因此首先考虑特殊情况是最简单的。当没有扩散时（因此扩散参数 d = 0），对于任何 α σ > 1 并且该方法对于任何 c 值都是不稳定的，即该方法无条件不稳定，使其无用。另一方面，当没有对流时（因此库朗数 c = 0），我们发现当 cos α = − 1 时 σ 最大，因此如果 d < 1，该方法是稳定的
2，即条件稳定。

直接诉诸我们对方程（6.32）的直觉，即所有旧节点值的系数应该为正，得出类似的结论：d < 0.5 且 c < 2 d。第一个条件导致 t 的限制：

\begin{sourceblock}
\includegraphics[page=184,trim={53.00bp 234.85bp 53.87bp 401.94bp},clip,width=0.92\textwidth,max height=0.38\textheight]{Ferziger4th.pdf}
\end{sourceblock}
第二个要求对时间步长没有限制，但给出了对流系数和扩散系数之间的关系：

\begin{sourceblock}
\includegraphics[page=184,trim={53.00bp 170.13bp 53.87bp 467.99bp},clip,width=0.92\textwidth,max height=0.38\textheight]{Ferziger4th.pdf}
\end{sourceblock}
即，细胞佩克莱特数应小于二。这已经被提及为使用对流通量的 CDS 获得的解有界性的充分（但不是必要）条件。

由于该方法基于常微分方程的显式欧拉方法和空间导数的中心差分近似的组合，因此它继承了两者的精度。因此，该方法在时间上是一阶的，在空间上是二阶的。要求 d < 0.5意味着，每次空间网格减半，时间步长必须减少四倍。这使得该方案不适合不需要高时间分辨率的问题（缓慢变化的解、接近稳态的解）；这些通常是人们想要及时使用一阶方法的情况。下面将显示一个说明稳定性问题的示例。

Courant 和 Friedrichs 在 1920 年代认识到与方程（6.40）条件相关的不稳态问题，并提出了一种至今仍在使用的解。他们指出，对于以对流为主的问题，系数 φ n − 1 是可能的
方程(6.32)中的i+1为负值，表明可以利用迎风差异来解决该问题。我们没有像上面那样通过 CDS 近似来评估对流项，而是使用 UDS（参见第 3 章）。然后我们可以代替方程（6.31）得到：

\begin{sourceblock}
\begin{tikzpicture}
\node[inner sep=0,anchor=south west] (image) at (0,0) {\includegraphics[page=185,trim={53.00bp 426.96bp 53.87bp 209.17bp},clip,width=0.92\textwidth,max height=0.38\textheight]{Ferziger4th.pdf}};
\begin{scope}[x={(image.south east)},y={(image.north west)}]
\fill[white] (0.09254,0.28457) rectangle (0.15441,0.73176);
\fill[white] (0.10971,0.20393) rectangle (0.12036,0.48950);
\fill[white] (0.15943,0.28457) rectangle (0.18761,0.67011);
\fill[white] (0.19113,0.28457) rectangle (0.22520,0.72143);
\end{scope}
\end{tikzpicture}
\end{sourceblock}
代替方程（6.32），得出：

\begin{sourceblock}
\includegraphics[page=185,trim={53.00bp 378.70bp 53.87bp 268.94bp},clip,width=0.92\textwidth,max height=0.38\textheight]{Ferziger4th.pdf}
\end{sourceblock}
由于相邻值的系数始终为正，因此它们不会导致不稳定的非物理行为。然而，系数 φ n
我可能会持消极态度，从而产生潜在的问题。为了使该系数为正，时间步长应满足以下条件：

\begin{sourceblock}
\includegraphics[page=185,trim={53.00bp 272.88bp 53.87bp 352.65bp},clip,width=0.92\textwidth,max height=0.38\textheight]{Ferziger4th.pdf}
\end{sourceblock}
当对流可以忽略不计时，稳定所需时间步长的限制与式（6.39）相同。对于可忽略不计的扩散，要满足的标准是：

\begin{sourceblock}
\includegraphics[page=185,trim={53.00bp 203.75bp 53.87bp 434.10bp},clip,width=0.92\textwidth,max height=0.38\textheight]{Ferziger4th.pdf}
\end{sourceblock}
即库朗数应小于 1。

也可以应用冯诺依曼稳定性分析来解决这个问题，分析结果与刚才得出的结论是一致的。因此，与中心差分法不同，迎风近似提供了一定的稳定性，加上其易用性，使得此类方法多年来非常流行。它至今仍在使用。当对流和扩散同时存在时，稳定性判据更加复杂。大多数人并没有处理这种复杂性，而是要求满足每个单独的标准，这个条件可能比必要的限制更多一点，但很安全。

该方法在空间和时间上都存在一阶截断误差，并且如果要保持较小的误差，则需要在两个变量中使用非常小的步长。由于这个原因，它很少被使用。

对库朗数的限制还解释为流体粒子在单个时间步中不应移动超过一个网格长度。这种对信息传播速率的限制看起来非常合理，但当采用此类方法解决稳态问题或使用局部高度精细的网格（例如，靠近实体墙）时，可能会限制收敛速率。

其他显式方案可以基于常微分方程的其他方法。事实上，前面描述的所有方法都曾在 CFD 中使用过。接下来将描述基于中心差的三时间级别方法，称为蛙跳方法。

6.3.1.2 蛙跳法
当将蛙跳法（定义见 6.2.2 节）应用于通用输运方程 Eq.(6.30)，其中使用 CDS 进行空间离散时，我们得到：

\begin{sourceblock}
\begin{tikzpicture}
\node[inner sep=0,anchor=south west] (image) at (0,0) {\includegraphics[page=186,trim={53.00bp 345.34bp 53.87bp 290.69bp},clip,width=0.92\textwidth,max height=0.38\textheight]{Ferziger4th.pdf}};
\begin{scope}[x={(image.south east)},y={(image.north west)}]
\fill[white] (0.06075,0.28695) rectangle (0.12259,0.73265);
\fill[white] (0.07789,0.20658) rectangle (0.08854,0.49120);
\fill[white] (0.12764,0.28695) rectangle (0.15582,0.67087);
\fill[white] (0.15934,0.28695) rectangle (0.22117,0.73265);
\end{scope}
\end{tikzpicture}
\end{sourceblock}
就无量纲系数 d 和 c 而言，上式为：

\begin{sourceblock}
\includegraphics[page=186,trim={53.00bp 297.10bp 53.87bp 350.54bp},clip,width=0.92\textwidth,max height=0.38\textheight]{Ferziger4th.pdf}
\end{sourceblock}
在这种方法中，永远无法满足物理真实模拟热传导的启发式要求，因为 φ n 的系数
我无条件否定！事实上，该方案是无条件不稳定的，并且对于非稳态问题的数值求解似乎没有用处。然而，我们需要仔细研究扩散在传输方程中的作用。 Mesinger和Arakawa（1976）给出了清晰的描述。

不幸的是，由于蛙跳方法的构造，它既有“正确”的物理模式，也有错误的“计算”模式（所有三级显式方法都表现出这些模式）。对式（6.46）进行冯诺依曼稳定性分析表明，该方法始终不稳定，但当 d = 0，即不存在扩散时，稳定性分析表明，所得线性对流方程的蛙跳格式对于库朗数 c ≤ 1 是稳定的，并且具有一定的数值​​色散，但没有数值扩散。这两种模式是解耦的，并且计算模式随着时间的推移而增长。另一方面，如果时间步长较小，则不稳定性非常弱，并且与其他方法相比，波状解的阻尼非常弱。由于这些原因，这种方法实际上在许多应用中使用（通过一些技巧来稳定它），特别是在对流占主导地位且扩散很小的气象学和海洋学中。

Williams（2009）列出了地球物理流体力学中使用蛙跳法的大量数值代码，并指出它们主要使用 Robert-Asselin（RA）时间滤波器（Asselin 1972）来抑制影响解的方法的计算模式。 RA 方案修改 φ n

\begin{sourceblock}
\begin{tikzpicture}
\node[inner sep=0,anchor=south west] (image) at (0,0) {\includegraphics[page=187,trim={53.00bp 529.58bp 53.87bp 104.99bp},clip,width=0.92\textwidth,max height=0.38\textheight]{Ferziger4th.pdf}};
\begin{scope}[x={(image.south east)},y={(image.north west)}]
\fill[white] (0.00000,0.97466) rectangle (0.07236,1.00000);
\fill[white] (0.07296,0.97466) rectangle (0.14433,1.00000);
\fill[white] (0.14499,0.97466) rectangle (0.22301,1.00000);
\fill[white] (0.22361,0.97466) rectangle (0.30830,1.00000);
\fill[white] (0.30890,0.97466) rectangle (0.35032,1.00000);
\fill[white] (0.35092,0.97466) rectangle (0.45973,1.00000);
\fill[white] (0.46036,0.97466) rectangle (0.51173,1.00000);
\fill[white] (0.51239,0.97466) rectangle (0.65859,1.00000);
\fill[white] (0.65922,0.97466) rectangle (0.76722,1.00000);
\fill[white] (0.78505,0.90941) rectangle (0.79570,1.00000);
\fill[white] (0.80499,0.97466) rectangle (0.86469,1.00000);
\fill[white] (0.86529,0.97466) rectangle (1.00000,1.00000);
\fill[white] (0.00000,0.59582) rectangle (0.02908,0.96199);
\fill[white] (0.03179,0.60754) rectangle (0.09362,1.00000);
\fill[white] (0.20024,0.12290) rectangle (0.23435,0.53785);
\fill[white] (0.21741,0.05100) rectangle (0.22806,0.32246);
\fill[white] (0.24030,0.12290) rectangle (0.31561,0.53785);
\fill[white] (0.69675,0.05100) rectangle (0.70740,0.32246);
\fill[white] (0.71798,0.12290) rectangle (0.74617,0.48907);
\fill[white] (0.74800,0.11435) rectangle (0.76779,0.48052);
\fill[white] (0.76758,0.08869) rectangle (0.77820,0.36015);
\fill[white] (0.78758,0.12290) rectangle (0.80063,0.48907);
\fill[white] (0.92439,0.11118) rectangle (1.00000,0.47735);
\end{scope}
\end{tikzpicture}
\end{sourceblock}
该校正 d f 是中央二阶导数时间滤波器。滤波器参数 ν 尽可能小（0.01-0.2），与消除杂散计算模式一致。这种滤波将蛙跳方案的精度从二阶降低到一阶，并抑制解的幅度。 Williams (2009)对该方案的修改使用

\begin{sourceblock}
\includegraphics[page=187,trim={53.00bp 416.45bp 53.87bp 216.82bp},clip,width=0.92\textwidth,max height=0.38\textheight]{Ferziger4th.pdf}
\end{sourceblock}
这减少了不需要的数值阻尼并提高了求解精度（Amezcua 等人，2011）。威廉姆斯证明 α ≥ 0.5 是稳定性所必需的，且 α = 0.53 且 ν = 0.2 比 RA 方案有显著改进；阿梅兹夸等人。使用 ν = 0.1. Williams 的 RAW 方案（由 Amezcua 等人在 2011 年命名）正在气象规范中采用，选择保持跨越式时间提前。一些代码管理员选择实施替代方案，例如第 1 节中描述的三阶龙格-库塔方案。 6.2.3.

\subsection*{6.3.2 隐式方法}\phantomsection\addcontentsline{toc}{subsection}{6.3.2 隐式方法}
如果稳定性是首要要求，则常微分方程方法的分析建议使用向后或隐式欧拉方法。应用于通用传输方程（6.30），利用空间导数的 CDS 近似，它给出：

\begin{sourceblock}
\begin{tikzpicture}
\node[inner sep=0,anchor=south west] (image) at (0,0) {\includegraphics[page=187,trim={53.00bp 151.74bp 53.87bp 483.62bp},clip,width=0.92\textwidth,max height=0.38\textheight]{Ferziger4th.pdf}};
\begin{scope}[x={(image.south east)},y={(image.north west)}]
\fill[white] (0.05468,0.28070) rectangle (0.11654,0.71670);
\fill[white] (0.07185,0.20208) rectangle (0.08250,0.48051);
\fill[white] (0.12156,0.28070) rectangle (0.14974,0.65627);
\fill[white] (0.15326,0.28070) rectangle (0.18734,0.70630);
\end{scope}
\end{tikzpicture}
\end{sourceblock}
或者，根据式(6.33)中引入的无量纲系数c和d重新排列：

\begin{sourceblock}
\begin{tikzpicture}
\node[inner sep=0,anchor=south west] (image) at (0,0) {\includegraphics[page=187,trim={53.00bp 81.08bp 53.87bp 552.66bp},clip,width=0.92\textwidth,max height=0.38\textheight]{Ferziger4th.pdf}};
\begin{scope}[x={(image.south east)},y={(image.north west)}]
\fill[white] (0.14812,0.25525) rectangle (0.17910,0.62377);
\fill[white] (0.18093,0.26667) rectangle (0.20911,0.62377);
\fill[white] (0.21095,0.25525) rectangle (0.25928,0.62377);
\fill[white] (0.25946,0.26667) rectangle (0.32129,0.68056);
\end{scope}
\end{tikzpicture}
\end{sourceblock}
上式可以写成：

\begin{sourceblock}
\includegraphics[page=188,trim={53.00bp 568.64bp 53.87bp 78.85bp},clip,width=0.92\textwidth,max height=0.38\textheight]{Ferziger4th.pdf}
\end{sourceblock}
其中矩阵系数为（见方程（6.33））：

\begin{sourceblock}
\includegraphics[page=188,trim={53.00bp 503.06bp 53.87bp 124.53bp},clip,width=0.92\textwidth,max height=0.38\textheight]{Ferziger4th.pdf}
\end{sourceblock}
在此方法中，所有通量和源项均根据新时间级别的未知变量值进行评估。结果是一个代数方程组，与稳态问题所获得的方程组非常相似；实际上，唯一的区别在于对系数 A P 和源项 Q P 的额外贡献，这源于不稳定项。

与常微分方程的情况一样，使用隐式欧拉方法允许采用任意大的时间步长；该属性对于研究缓慢瞬变或稳态流动的流动很有用。当 CDS 用于粗网格时可能会出现问题（如果在变量梯度变化强烈的区域佩克莱特数太大）；产生了振荡解，但方案保持稳定。

该方法的缺点是其时间上的一阶截断误差以及需要在每个时间步求解大量耦合方程组。它还需要比显式方案更多的存储空间，因为必须存储整个系数矩阵 A 和源向量。优点是可以使用大时间步长，尽管有缺点，但可能会导致更有效的过程，特别是在迈向稳态解时。

6.3.2.2 曲柄-尼科尔森法
梯形法则方法的二阶精度及其相对简单性表明，当时间精度很重要时，它可以应用于偏微分方程。它后来被称为克兰克-尼科尔森方法。特别是，当应用于具有空间导数 CDS 离散化的一维通用输运方程时，有：

\begin{sourceblock}
\begin{tikzpicture}
\node[inner sep=0,anchor=south west] (image) at (0,0) {\includegraphics[page=188,trim={53.00bp 126.71bp 53.87bp 477.76bp},clip,width=0.92\textwidth,max height=0.38\textheight]{Ferziger4th.pdf}};
\begin{scope}[x={(image.south east)},y={(image.north west)}]
\fill[white] (0.08000,0.64115) rectangle (0.14183,0.85860);
\fill[white] (0.09714,0.60175) rectangle (0.10779,0.74072);
\fill[white] (0.14752,0.64115) rectangle (0.17570,0.82860);
\fill[white] (0.17991,0.64115) rectangle (0.21398,0.85358);
\fill[white] (0.20758,0.24907) rectangle (0.22072,0.43652);
\fill[white] (0.19347,0.01946) rectangle (0.21329,0.20691);
\end{scope}
\end{tikzpicture}
\end{sourceblock}
该方案是隐式的；新时间水平上通量和源的贡献产生了一组类似于隐式欧拉方案的耦合方程。上式可以重写为：

\begin{sourceblock}
\includegraphics[page=189,trim={53.00bp 592.54bp 53.87bp 54.94bp},clip,width=0.92\textwidth,max height=0.38\textheight]{Ferziger4th.pdf}
\end{sourceblock}
其中（以无量纲系数 c 和 d 表示，在方程（6.33）中引入）：

\begin{sourceblock}
\includegraphics[page=189,trim={53.00bp 512.47bp 53.87bp 102.76bp},clip,width=0.92\textwidth,max height=0.38\textheight]{Ferziger4th.pdf}
\end{sourceblock}
Qt 项
i 代表“附加”源项，其中包含前一个时间级别的贡献；它在新时间级别的迭代过程中保持不变。该方程还可能包含依赖于新解的源项，因此上述项需要单独存储。

与一阶隐式欧拉方案相比，该方案每步只需要很少的计算量。冯·诺依曼稳定性分析表明，该方案是无条件稳定的，但对于大时间步长，振荡解（甚至不稳定）是可能的。这可能归因于 φ n 的系数的可能性
i 在大 t 处变为负值，但如果 1 − d > 0 或 t < ρ( x ) 2 /，则保证为正值，这是显式欧拉方法允许的最大步长的两倍。实际上，可以使用更大的时间步长而不会产生振荡；该限制取决于问题。

如果将修正方程方法 4（第 4.4.6 节）应用于方程（6.54），则结果为（仅保留前序截断项时；参见 Fletcher 1991，表 9.3）

\begin{sourceblock}
\begin{tikzpicture}
\node[inner sep=0,anchor=south west] (image) at (0,0) {\includegraphics[page=189,trim={53.00bp 297.93bp 53.87bp 332.82bp},clip,width=0.92\textwidth,max height=0.38\textheight]{Ferziger4th.pdf}};
\begin{scope}[x={(image.south east)},y={(image.north west)}]
\fill[white] (0.04638,0.44137) rectangle (0.08460,0.76801);
\fill[white] (0.73600,0.43402) rectangle (0.76725,0.79542);
\fill[white] (0.67964,0.04408) rectangle (0.70078,0.37101);
\fill[white] (0.92439,0.00000) rectangle (1.00000,0.07262);
\end{scope}
\end{tikzpicture}
\end{sourceblock}
(6.57) 有两件事是显而易见的。首先，截断误差中剩余的三阶导数表明该方案是分散的，这意味着格式良好的初始条件将随着时间的推移而失真，因为其分量在解中以不同的速度移动。 Fletcher（1991）表明这种效应可能是显著的！其次，该方案具有正系数的四阶导数，这将稍微削弱物理扩散。 5 沃宁和海耶特 (1974) 和多尼亚
4 回想一下，该过程是 (1) 将差分格式中的项展开关于单点（在本例中为 x i , t n ）的二维 ( x , t ) 泰勒级数中，以获得扩展偏微分方程，(2) 使用该偏微分方程本身（及其导数）将所有高阶和混合时间项替换为扩展 PDE 中的空间导数项，以及 (3) 重新排列，使原始 PDE 位于左侧，右侧的其余项。 Warming 和 Hyett (1974) 的程序表在手动执行该方法时非常有用。其余项是截断误差，即差分方程求解结果与原始方程求解结果之间的差异。保留最低阶的以查看最重要的物理效果。 5 为了发生扩散（耗散），修正方程中偶数阶导数的系数必须具有交替符号，二阶项为正，四阶项为负，依此类推。请注意，此处的原始方程中存在二阶物理扩散。

等人。 (1987)表明格式稳定的必要条件是修正方程中的首偶次项是扩散的；然而，对于完整的稳定性分析，可能需要应用冯诺依曼方法。

克兰克-尼科尔森格式可以被视为一阶显式和隐式欧拉格式的平等混合。只有等量混合才能获得二阶精度；对于可能随空间和时间变化的其他混合因素，该方法保持一阶准确。如果隐式贡献增加，稳定性会增加，但准确度会降低，如下一节所述。

6.3.2.3 θ –Method
显式和隐式欧拉方案以及克兰克-尼科尔森方案可以组合成一个更通用的（加权）方案，允许在稳定性和准确性之间进行选择，这在非线性问题中非常重要。对于一般输运方程，为了从 t n 推进到 t n + 1 = t n + t，我们代替方程（6.54）写道：

\begin{sourceblock}
\begin{tikzpicture}
\node[inner sep=0,anchor=south west] (image) at (0,0) {\includegraphics[page=190,trim={53.00bp 345.17bp 53.87bp 259.30bp},clip,width=0.92\textwidth,max height=0.38\textheight]{Ferziger4th.pdf}};
\begin{scope}[x={(image.south east)},y={(image.north west)}]
\fill[white] (0.07546,0.64115) rectangle (0.13732,0.85876);
\fill[white] (0.09263,0.60175) rectangle (0.10328,0.74088);
\fill[white] (0.14304,0.64115) rectangle (0.17122,0.82860);
\fill[white] (0.17540,0.64115) rectangle (0.20947,0.85341);
\fill[white] (0.17540,0.13410) rectangle (0.20638,0.32755);
\fill[white] (0.20821,0.14010) rectangle (0.27041,0.32755);
\fill[white] (0.28968,0.13572) rectangle (0.30283,0.32333);
\end{scope}
\end{tikzpicture}
\end{sourceblock}
该方案是隐式的，除非 θ = 0，在这种情况下，它恢复为前向欧拉方案。当 θ = 1 时，检索向后欧拉方案，而当 θ = 0 时。 50，上面的克兰克-尼科尔森方案就出现了。由此可见，可以应用式（6.55）和（6.56）中描述的求解方法。当θ<1时
如图2所示，根据冯诺依曼方法，该方法是不稳定的。该方法在 1 内稳定
2 ≤ θ ≤ 1，并且正如预期的那样，对于 θ = 1 2，即二阶，最准确。对于 θ > 1
如图2所示，该方法是一阶精确且耗散的。虽然使用此方法的代码可以运行用于 θ = 1 的测试
2 , 0 . 5 处于稳定性极限，因此现实世界模拟运行的代码通常采用 0.52≤θ≤0.60.

Casulli 和 Cattani (1994) 中描述了这种 θ 方法，随后 Casulli 在一系列与自由表面和非静流动相关的令人印象深刻的论文中使用了该方法。弗林格等人。 (2006)在非静水力沿海海洋模拟的代码中采用了该方法。

6.3.2.4 三时间水平法
完全隐式的二阶精度方案可以通过使用时间上的二次向后近似来获得，如第 1 节中所述。 6.2.4.对于一维通用输运方程和空间中的 CDS 离散化，我们得到：

\begin{sourceblock}
\begin{tikzpicture}
\node[inner sep=0,anchor=south west] (image) at (0,0) {\includegraphics[page=191,trim={53.00bp 552.10bp 53.87bp 53.84bp},clip,width=0.92\textwidth,max height=0.38\textheight]{Ferziger4th.pdf}};
\begin{scope}[x={(image.south east)},y={(image.north west)}]
\fill[white] (0.25759,0.26844) rectangle (0.31943,0.49120);
\fill[white] (0.19095,0.13920) rectangle (0.25221,0.33571);
\end{scope}
\end{tikzpicture}
\end{sourceblock}
由此产生的代数方程可以写成：

\begin{sourceblock}
\includegraphics[page=191,trim={53.00bp 491.41bp 53.87bp 146.71bp},clip,width=0.92\textwidth,max height=0.38\textheight]{Ferziger4th.pdf}
\end{sourceblock}
系数 A E 和 A W 与隐式欧拉方案的系数相同，见式（1）。 6.53。现在，中心系数来自时间导数的贡献更大：

\begin{sourceblock}
\begin{tikzpicture}
\node[inner sep=0,anchor=south west] (image) at (0,0) {\includegraphics[page=191,trim={53.00bp 426.22bp 53.87bp 211.90bp},clip,width=0.92\textwidth,max height=0.38\textheight]{Ferziger4th.pdf}};
\begin{scope}[x={(image.south east)},y={(image.north west)}]
\fill[white] (0.00000,0.81585) rectangle (0.12950,1.00000);
\end{scope}
\end{tikzpicture}
\end{sourceblock}
源项包含时间 t n − 1 的贡献；参见等式。 6.60。

这种三时间电平（TTL）方案比 Crank-Nicolson（CN）方案更容易实现；它也不太容易产生振荡解，尽管 t 值较大时可能会发生这种情况。必须存储三个时间级别的变量值，但内存要求与 Crank-Nicolson 方案相同。该格式在时间上是二阶精确的，可以证明该格式是无条件稳定的。从式（6.60）中我们还可以看出，节点i处旧值的系数始终为正；然而，t n − 1 处的值的系数始终为负，这就是为什么如果步长很大，该方案可能会产生振荡解。

如果将修正方程法（第 4.4.6 节）应用于方程（6.59），则结果为（保留前序截断项；参见 Fletcher 1991，表 9.3）

\begin{sourceblock}
\begin{tikzpicture}
\node[inner sep=0,anchor=south west] (image) at (0,0) {\includegraphics[page=191,trim={53.00bp 227.98bp 53.87bp 408.81bp},clip,width=0.92\textwidth,max height=0.38\textheight]{Ferziger4th.pdf}};
\begin{scope}[x={(image.south east)},y={(image.north west)}]
\fill[white] (0.92424,0.00000) rectangle (1.00000,0.11175);
\end{scope}
\end{tikzpicture}
\end{sourceblock}
(6.62) 由式(6.62)可知，对于CN方法，截断误差为
O(x2)。有趣的是，在小 c 2 的限制下，CN 和 TTL 误差是相同的，但对于 c 2 = u t
x = O ( 1 )，对于色散，TTL 误差大 2 倍；对于“反”扩散/耗散，TTL 误差大 3.25 倍。

该方案可以与一阶隐式欧拉方案混合。只有对中心系数和源项的贡献需要以第 1 节中描述的延迟校正方法的方式进行修改。 5.6.这在开始计算时很有用，因为只有一个旧级别的解可用。此外，如果需要稳态解，则切换到隐式欧拉方案可确保稳定性并允许使用大时间步长。混合少量的一阶方案有助于防止振荡，从而有助于解的美观（没有振荡，精度不会更好，但图形看起来更好）。

如果确实发生振荡，则必须减小时间步长，因为振荡表明存在较大的时间离散化误差。此评论不适用于仅有条件稳定的方案。

该方案适用于大多数商业和公共 CFD 代码。其受欢迎的一个原因是它是完全隐式的，即对流、扩散和源项仅在新的时间水平上计算，就像一阶隐式欧拉方案一样。这意味着可以在两个时间步之间更改网格；先前时间级别所需的唯一信息是网格点处的变量值，因此需要将旧解插值到新网格。然而，我们不需要先前时间水平的通量或源项，就像克兰克-尼科尔森方案的情况一样。

\subsection*{6.3.3 其他方法}\phantomsection\addcontentsline{toc}{subsection}{6.3.3 其他方法}
上述方案是通用 CFD 代码中最常用的方案。对于特殊目的，例如在大涡流和湍流的直接模拟中，经常使用高阶方案，例如三阶或四阶龙格-库塔或亚当斯方法。通常，当空间离散化也是高阶时，使用高阶时间离散化，这是当解域具有规则形状时的情况，以便易于应用空间中的高阶方法。将常微分方程的高阶方法应用于 CFD 问题非常简单。

人们可以以与在 θ 方法中混合显式和隐式欧拉方案类似的方式混合任意两种方案。通常，一种方案阶数较高，但在某些条件下容易出现稳定性问题；另一个通常是无条件稳定的，但不太准确（通常是隐式欧拉方案）。当高阶方案遇到问题时，可以在本地应用这种混合。然而，识别此类条件并确定最佳混合因子以最大限度地提高准确性和稳定性并非易事。尽管我们没有提供这种方法的示例，但我们想指出这种可能性；类似的空间离散化方法在第 3 节中进行了解释。 4.4.6.

\section*{6.4 示例}\phantomsection\addcontentsline{toc}{section}{6.4 示例}
为了演示上述一些方法的性能，我们首先看一下第 1 章示例问题的不稳定版本。 3.待解问题由式（6.30）给出，初始条件和边界条件如下：t=0时，φ0=0；对于所有后续时间，在 x = 0 时 φ = 0，在 x = L = 1、ρ = 1、u = 1 和 = 0 时，φ = 1.1. 由于边界条件不随时间变化，解从初始零场发展到第 1 节中给出的稳态解。 3.10.对于空间离散化，使用二阶 CDS。对于时间

\begin{sourceblock}
\includegraphics[page=193,trim={53.00bp 431.48bp 53.87bp 52.00bp},clip,width=\textwidth,max height=0.38\textheight]{Ferziger4th.pdf}
\par\vspace{0.45\baselineskip}
\small 图 6.3 使用时间步 t = 0 通过显式欧拉方法求解的时间演化。 003（左；d < 0 . 5）且 t = 0 . 00325（右，d > 0 . 5）
\end{sourceblock}
离散化方面，我们使用显式和隐式一阶欧拉方法、Crank-Nicolson 方法和具有三个时间级别的完全隐式方法。

我们首先演示当显式欧拉方案（仅有条件稳定）与违反稳定性条件的时间步一起使用时会发生什么。图 6.3 显示了解在一小段时间内的演变，使用略低于方程 (6.39) 给出的临界值和稍高于方程 (6.39) 给出的临界值的时间步进行计算。当时间步长大于临界值时，就会产生振荡，并且振荡会随着时间无限增长。比图 6.3 所示的最后一个解晚了几个时间步，数字变得太大而无法由计算机处理。对于隐式方案，即使使用非常大的时间步长也不会出现问题。

为了研究时间离散化的准确性，我们查看 x = 0 处节点的解。 95 个均匀网格，在时间 t = 0 时具有 41 个节点 ( x = 0 . 025)。 01. 我们使用上述四种方案使用 5、10、20 和 40 个时间步长进行了计算。 φ 在 x = 0 处收敛。 95（t = 0）。 01 随着时间步长减小如图 6.4 所示。隐式欧拉和三水平方法低估了正确值，而显式欧拉和克兰克-尼科尔森方法则高估了正确值。所有方案都显示出向时间步独立解的单调收敛性。

因为我们没有精确的解来比较，所以我们在时间 t = 0 时获得了精确的参考解。 01 使用 Crank-Nicolson 方案（最准确的方法），t = 0.0001（100 个时间步长）。该解比上述任何解都要准确得多，因此可以将其视为用于误差估计目的的精确解。通过从该参考解中减去上述解，我们获得了每个方案和时间步长的时间离散化误差的估计。所有空间离散化误差都是相同的

\begin{sourceblock}
\includegraphics[page=194,trim={53.00bp 422.48bp 53.87bp 52.00bp},clip,width=\textwidth,max height=0.38\textheight]{Ferziger4th.pdf}
\par\vspace{0.45\baselineskip}
\small 图 6.4 φ 在 x = 0 处的收敛。 95 且 t = 0.01 各种时间积分方案的时间步长减小（左）和时间离散化误差（右）
\end{sourceblock}
例，并且在这里不起作用。由此获得的误差相对于时间步长大小绘制在图 6.4 中。

两种欧拉方法显示了预期的一阶行为：当时间步减少相同量时，误差减少一个数量级。两个二阶方案还显示了预期的误差减少率，该率非常接近理想的斜率。然而，Crank-Nicolson 方法给出了更准确的解，因为它的初始误差要小得多。三层方案采用隐式欧拉法启动，初始误差较大。因为在这个问题中，时间变化从初始状态到稳态是单调的，所以初始误差在整个解中仍然很重要。两种方法的错误减少率相同，但在这种情况下，错误级别由其初始值决定。

接下来我们检查第 1 章的 2D 测试用例。 4，涉及从具有规定温度、暴露于驻点流的壁的传热；见第 4.7 节.初始解还是 φ 0 = 0；边界条件不随时间变化，与第 2 节研究的稳态问题相同。 4.7，其中 ρ = 1.2 和 = 0.1. 空间离散采用CDS，使用20×20 CV的均匀网格。隐式方案中的线性方程系统使用 SIP 求解器（在第 5 章中描述）进行求解，并且迭代误差降低到 10 − 5 以下。我们计算解向稳态的时间演化。图 6.5 显示了四个时刻的等温线。

为了研究这种情况下各种方案的准确性，我们观察时间 t = 0 时通过等温壁的热通量。 12. 热通量 Q 随时间步长的变化如图 6.6 所示。与前面的测试用例一样，使用二阶方案获得的结果随着数量的变化很小

\begin{sourceblock}
\includegraphics[page=195,trim={53.00bp 344.48bp 53.87bp 52.00bp},clip,width=\textwidth,max height=0.38\textheight]{Ferziger4th.pdf}
\par\vspace{0.45\baselineskip}
\small 图 6.5 t = 0 时刻非定常二维问题的等温线。 2（左上），t = 0.5（右上），t = 1.0（左下）和 t = 2.0（右下），在统一的 20 × 20 CV 网格上计算，使用 CDS 进行空间离散化，使用 Crank-Nicolson 方法进行时间离散化
\end{sourceblock}
\begin{sourceblock}
\includegraphics[page=195,trim={53.00bp 114.12bp 53.87bp 355.37bp},clip,width=\textwidth,max height=0.38\textheight]{Ferziger4th.pdf}
\par\vspace{0.45\baselineskip}
\small 图 6.6 t = 0 时通过等温壁的热通量。 12（左）和计算的壁热通量中的时间离散误差（右）作为各种方案的时间步长的函数（通过 CDS 进行空间离散，20 × 20 CV 均匀网格）
\end{sourceblock}
时间步长增加，而一阶方案的准确性却大大降低。显式欧拉方案不是单调收敛的；使用最大时间步长获得的解位于与时间步长无关的值与使用较小时间步长获得的值的相反侧。

为了误差估计，我们通过 Crank-Nicolson 方案使用非常小的时间步长（t = 0 . 0003，400 步到 t = 0 . 12）获得了准确的参考解。所有情况下的空间离散化都是相同的，因此空间离散化误差再次抵消。通过从参考解中减去使用不同方案和时间步计算的热通量值，我们获得时间离散误差的估计。这些是根据标准化时间步长（用最大时间步长标准化）绘制的，如图 6.6 所示。

再次获得一阶和二阶格式的预期渐近收敛率。然而，现在通过具有三个时间级别的二阶方案获得了最低的误差。与前面的示例一样，这是由于初始误差的主导效应所致，当使用隐式欧拉方法启动三级方案时，初始误差比在 Crank-Nicolson 方案中更小。两种二阶方案的误差都比一阶方案小得多：使用最大时间步的二阶方案的结果比使用小八倍时间步的一阶方案更准确！

尽管我们正在处理从初始状态到稳态平稳过渡的简单非稳态问题，但我们发现一阶欧拉方案非常不准确。在这两个测试用例中，二阶方案的误差比欧拉方案（其误差在最小时间步长下为 1\% 左右）低两个数量级以上。可以预见，在瞬态流动中，可能会出现更大的差异。一阶格式中，求稳态解时只能使用隐式欧拉法；对于瞬态流动问题，建议使用二阶（或更高）阶方案。

非定常流动问题将在各节中讨论。 8.4.2 和 9.12。

\chapter{Navier--Stokes 方程的求解：第一部分}
纳维-斯托克斯方程的解是本书的主要主题；因此，我们对材料给予了极大的关注，并将其分为两个主要部分，每个部分组成一章。本章我们将介绍方程的基本问题和特征以及求解方法；一些常用方法的实现细节及其应用示例将在下一章中介绍。这种布局旨在为新手和/或专业人士提供一个方便的结构，允许从一开始就进行系统学习，或者直接访问特定问题、方法或工具，例如，(1) 使用商业代码并希望了解这些代码中的可用方法，或 (2) 需要在替代方案中进行选择的代码开发人员。

\section*{7.1 基础知识}\phantomsection\addcontentsline{toc}{section}{7.1 基础知识}
在章节中。在图 3、4 和 6 中，我们处理了一般守恒方程的离散化。其中描述的离散化原理适用于动量方程和连续性方程（我们统称为纳维-斯托克斯方程）。我们在这里重述这些方程的积分形式：

\begin{sourceblock}
\begin{tikzpicture}
\node[inner sep=0,anchor=south west] (image) at (0,0) {\includegraphics[page=197,trim={53.00bp 140.23bp 53.87bp 460.80bp},clip,width=0.92\textwidth,max height=0.38\textheight]{Ferziger4th.pdf}};
\begin{scope}[x={(image.south east)},y={(image.north west)}]
\fill[white] (0.13368,0.80387) rectangle (0.15224,0.98157);
\fill[white] (0.20283,0.68008) rectangle (0.24800,0.87421);
\fill[white] (0.25672,0.69098) rectangle (0.29630,0.87022);
\fill[white] (0.30268,0.69667) rectangle (0.33086,0.87421);
\fill[white] (0.12818,0.58393) rectangle (0.15744,0.76563);
\fill[white] (0.18033,0.56581) rectangle (0.19795,0.69759);
\fill[white] (0.35065,0.01843) rectangle (0.36827,0.15021);
\end{scope}
\end{tikzpicture}
\end{sourceblock}
并且，回顾方程（1.18），t i = τ i j i j − p i i，我们给出一个微分形式：

\begin{sourceblock}
\begin{tikzpicture}
\node[inner sep=0,anchor=south west] (image) at (0,0) {\includegraphics[page=197,trim={53.00bp 82.18bp 53.87bp 554.84bp},clip,width=0.92\textwidth,max height=0.38\textheight]{Ferziger4th.pdf}};
\begin{scope}[x={(image.south east)},y={(image.north west)}]
\fill[white] (0.24093,0.52404) rectangle (0.32815,0.95845);
\end{scope}
\end{tikzpicture}
\end{sourceblock}
\begin{sourceblock}
\includegraphics[page=198,trim={53.00bp 584.34bp 53.87bp 53.51bp},clip,width=0.92\textwidth,max height=0.38\textheight]{Ferziger4th.pdf}
\end{sourceblock}
由式（1.17）、（1.21）、（1.5）和（1.6）得出；见第 1.4 节 应力和强迫术语的定义。接下来，我们将描述如何处理动量方程中与一般守恒方程中不同的项。

动量方程中的非定常项和平流项与一般守恒方程具有相同的形式。扩散（粘性）项与通用方程中的对应项类似，但由于动量方程是矢量方程，这些贡献变得更加复杂，需要更详细地考虑它们的处理。动量方程还包含压力的贡献，这在通用方程中没有类似物。它可以被视为源项（将压力梯度视为体积力——非保守地）或表面力（保守地处理），但由于压力和连续性方程的紧密联系，需要特别注意。最后，主变量是向量这一事实使得网格的选择具有更大的自由度。

\subsection*{7.1.1 对流项和粘性项的离散化}\phantomsection\addcontentsline{toc}{subsection}{7.1.1 对流项和粘性项的离散化}
\begin{sourceblock}
\begin{tikzpicture}
\node[inner sep=0,anchor=south west] (image) at (0,0) {\includegraphics[page=198,trim={53.00bp 294.12bp 53.87bp 342.55bp},clip,width=0.92\textwidth,max height=0.38\textheight]{Ferziger4th.pdf}};
\begin{scope}[x={(image.south east)},y={(image.north west)}]
\fill[white] (0.27711,0.53003) rectangle (0.36812,0.95928);
\fill[white] (0.36833,0.53003) rectangle (0.39371,0.95928);
\fill[white] (0.43916,0.31693) rectangle (0.48722,0.70920);
\end{scope}
\end{tikzpicture}
\end{sourceblock}
动量方程中对流项的处理遵循通用方程中对流项的处理；章节中描述的任何方法。可以使用3和4。

动量方程中的粘性项对应于通用方程中的扩散项；它们的微分和积分形式是：

\begin{sourceblock}
\begin{tikzpicture}
\node[inner sep=0,anchor=south west] (image) at (0,0) {\includegraphics[page=198,trim={53.00bp 187.98bp 53.87bp 448.68bp},clip,width=0.92\textwidth,max height=0.38\textheight]{Ferziger4th.pdf}};
\begin{scope}[x={(image.south east)},y={(image.north west)}]
\fill[white] (0.30301,0.53019) rectangle (0.34262,0.95929);
\fill[white] (0.34117,0.53019) rectangle (0.35182,0.82123);
\fill[white] (0.40081,0.31716) rectangle (0.44884,0.70929);
\end{scope}
\end{tikzpicture}
\end{sourceblock}
其中，对于牛顿流体和不可压缩流：

\begin{sourceblock}
\begin{tikzpicture}
\node[inner sep=0,anchor=south west] (image) at (0,0) {\includegraphics[page=198,trim={53.00bp 129.26bp 53.87bp 500.77bp},clip,width=0.92\textwidth,max height=0.38\textheight]{Ferziger4th.pdf}};
\begin{scope}[x={(image.south east)},y={(image.north west)}]
\fill[white] (0.34830,0.22902) rectangle (0.37179,0.57906);
\fill[white] (0.37035,0.22902) rectangle (0.38099,0.46635);
\fill[white] (0.38806,0.25893) rectangle (0.41624,0.57906);
\fill[white] (0.41976,0.25893) rectangle (0.44511,0.57906);
\end{scope}
\end{tikzpicture}
\end{sourceblock}
由于动量方程是矢量方程，因此粘性项比一般扩散项更复杂。动量方程中的粘性项与一般守恒方程中的扩散项相对应的部分是

\begin{sourceblock}
\begin{tikzpicture}
\node[inner sep=0,anchor=south west] (image) at (0,0) {\includegraphics[page=199,trim={53.00bp 582.39bp 53.87bp 54.29bp},clip,width=0.92\textwidth,max height=0.38\textheight]{Ferziger4th.pdf}};
\begin{scope}[x={(image.south east)},y={(image.north west)}]
\fill[white] (0.25585,0.56687) rectangle (0.27438,0.95927);
\fill[white] (0.33985,0.52987) rectangle (0.38304,0.95927);
\fill[white] (0.31471,0.32960) rectangle (0.34006,0.72234);
\fill[white] (0.24195,0.08045) rectangle (0.27702,0.48201);
\fill[white] (0.27651,0.05295) rectangle (0.28716,0.34386);
\end{scope}
\end{tikzpicture}
\end{sourceblock}
该项可以使用第 1 章中针对通用方程的相应项描述的任何方法进行离散化。 3 和 4 类似，但这只是粘性效应对动量第 i 个分量的贡献之一。

方程 (1.28)、(1.26)、(1.21) 和 (1.17) 使我们能够将其他粘性效应识别为体积粘度（仅在可压缩流中非零）的贡献以及由于粘度的空间变化而产生的进一步贡献。对于具有恒定流体特性的不可压缩流，这些贡献消失（由于连续性方程）。

当不可压缩流中的粘度在空间上变化时，非零的额外项可以以与项（7.8）相同的方式处理：

\begin{sourceblock}
\begin{tikzpicture}
\node[inner sep=0,anchor=south west] (image) at (0,0) {\includegraphics[page=199,trim={53.00bp 415.71bp 53.87bp 220.96bp},clip,width=0.92\textwidth,max height=0.38\textheight]{Ferziger4th.pdf}};
\begin{scope}[x={(image.south east)},y={(image.north west)}]
\fill[white] (0.22189,0.56668) rectangle (0.24045,0.95894);
\fill[white] (0.28075,0.32949) rectangle (0.34081,0.95928);
\fill[white] (0.34102,0.53003) rectangle (0.35167,0.82117);
\fill[white] (0.20800,0.08042) rectangle (0.24307,0.48185);
\fill[white] (0.24253,0.05294) rectangle (0.25317,0.34374);
\fill[white] (0.30749,0.05294) rectangle (0.34911,0.48185);
\fill[white] (0.53832,0.04072) rectangle (0.55362,0.33152);
\end{scope}
\end{tikzpicture}
\end{sourceblock}
其中 n 是垂直于控制体积表面的向外单位，并且适用 j 的求和。如上所述，对于常数 μ，此项消失。因此，即使使用隐式求解方法，该术语也经常被明确处理。有人认为，即使当粘度变化时，该项与项（7.8）相比也很小，因此其处理对收敛速度仅产生轻微影响。然而，这一论点严格地仅适用于整体意义上。任何一个 CV 面上的额外项都可能相当大。

\subsection*{7.1.2 压力项和体积力的离散化}\phantomsection\addcontentsline{toc}{subsection}{7.1.2 压力项和体积力的离散化}
\begin{sourceblock}
\begin{tikzpicture}
\node[inner sep=0,anchor=south west] (image) at (0,0) {\includegraphics[page=199,trim={53.00bp 238.13bp 53.87bp 411.52bp},clip,width=0.92\textwidth,max height=0.38\textheight]{Ferziger4th.pdf}};
\begin{scope}[x={(image.south east)},y={(image.north west)}]
\fill[white] (0.00000,0.84536) rectangle (0.03741,1.00000);
\fill[white] (0.03889,0.84536) rectangle (0.11026,1.00000);
\fill[white] (0.11179,0.84536) rectangle (0.13988,1.00000);
\fill[white] (0.14138,0.84536) rectangle (0.24325,1.00000);
\fill[white] (0.24478,0.84536) rectangle (0.28451,1.00000);
\fill[white] (0.28605,0.84536) rectangle (0.37738,1.00000);
\fill[white] (0.37886,0.84536) rectangle (0.43359,1.00000);
\fill[white] (0.43510,0.84536) rectangle (0.49317,1.00000);
\fill[white] (0.49468,0.84536) rectangle (0.53609,1.00000);
\fill[white] (0.53756,0.84536) rectangle (0.66881,1.00000);
\fill[white] (0.67032,0.84536) rectangle (0.69844,1.00000);
\fill[white] (0.69991,0.84536) rectangle (0.77131,1.00000);
\fill[white] (0.77278,0.84536) rectangle (0.80256,1.00000);
\fill[white] (0.80406,0.84536) rectangle (0.84547,1.00000);
\fill[white] (0.84695,0.84536) rectangle (1.00000,1.00000);
\fill[white] (0.19564,0.00000) rectangle (0.25365,0.84354);
\fill[white] (0.25549,0.12007) rectangle (0.28322,0.82777);
\fill[white] (0.28704,0.12007) rectangle (0.31681,0.82171);
\fill[white] (0.32066,0.12007) rectangle (0.50686,0.82171);
\fill[white] (0.51068,0.12007) rectangle (0.58713,0.82171);
\fill[white] (0.59098,0.12007) rectangle (0.63242,0.82171);
\fill[white] (0.63624,0.12007) rectangle (0.68268,0.82171);
\fill[white] (0.68650,0.12007) rectangle (0.74623,0.82171);
\fill[white] (0.75005,0.12007) rectangle (0.77486,0.82171);
\fill[white] (0.77865,0.12007) rectangle (0.84253,0.82171);
\fill[white] (0.84638,0.12007) rectangle (0.90111,0.82171);
\fill[white] (0.90493,0.12007) rectangle (0.96800,0.82171);
\fill[white] (0.97182,0.12007) rectangle (1.00000,0.82171);
\fill[white] (0.00000,0.00000) rectangle (0.04075,0.09642);
\fill[white] (0.04274,0.00000) rectangle (0.18403,0.09642);
\fill[white] (0.18605,0.00000) rectangle (0.30565,0.09642);
\fill[white] (0.30773,0.00000) rectangle (0.36078,0.09642);
\fill[white] (0.36274,0.00000) rectangle (0.47958,0.09642);
\fill[white] (0.48156,0.00000) rectangle (0.58623,0.09642);
\fill[white] (0.58827,0.00000) rectangle (0.62971,0.09642);
\fill[white] (0.63170,0.00000) rectangle (0.73468,0.09642);
\fill[white] (0.73672,0.00000) rectangle (0.76647,0.09642);
\fill[white] (0.76848,0.00000) rectangle (0.81657,0.09642);
\fill[white] (0.81856,0.00000) rectangle (0.92156,0.09642);
\fill[white] (0.92361,0.00000) rectangle (1.00000,0.09642);
\end{scope}
\end{tikzpicture}
\end{sourceblock}
3∇·v。在不可压缩流动中，最后一项为零。动量方程的一种形式（参见方程（7.3））包含该量的梯度，可以通过第 1 章中描述的 FD 方法来近似。 3.然而，由于网格上的压力和速度节点可能不一致，因此用于其导数的近似值可能不同。

在 FV 方法中，压力项通常被视为表面力（保守方法），即，在 u i 的方程中，积分

\begin{sourceblock}
\begin{tikzpicture}
\node[inner sep=0,anchor=south west] (image) at (0,0) {\includegraphics[page=199,trim={53.00bp 130.36bp 53.87bp 513.30bp},clip,width=0.92\textwidth,max height=0.38\textheight]{Ferziger4th.pdf}};
\begin{scope}[x={(image.south east)},y={(image.north west)}]
\fill[white] (0.40932,0.43238) rectangle (0.43750,0.94662);
\end{scope}
\end{tikzpicture}
\end{sourceblock}
是必需的。第 1 章中描述的方法。然后可以使用图4来近似表面积分。正如我们下面将要说明的，此项的处理以及网格上变量的排列对于保证数值求解方法的计算效率和精度起着重要作用。

或者，可以通过将上述积分保留为其体积形式来非保守地处理压力：

\begin{sourceblock}
\begin{tikzpicture}
\node[inner sep=0,anchor=south west] (image) at (0,0) {\includegraphics[page=200,trim={53.00bp 546.24bp 53.87bp 97.42bp},clip,width=0.92\textwidth,max height=0.38\textheight]{Ferziger4th.pdf}};
\begin{scope}[x={(image.south east)},y={(image.north west)}]
\fill[white] (0.39507,0.43238) rectangle (0.42325,0.94662);
\end{scope}
\end{tikzpicture}
\end{sourceblock}
在这种情况下，导数（或者，对于非正交网格，所有三个导数）需要在 CV 内的一个或多个位置进行近似。非保守方法引入了全局非保守误差；尽管随着网格大小变为零，此误差趋于零，但对于有限的网格大小来说，它可能很重要。

两种方法之间的差异仅在 FV 方法中显著。在 FD 方法中，尽管可以产生保守近似值和非保守近似值，但两种版本之间没有区别。

其他体积力，例如在非笛卡尔坐标系中使用协变或逆变速度时产生的非保守力，很容易在有限差分格式中处理：它们通常是一个或多个变量的简单函数，可以使用第 1 章中描述的技术进行评估。 3.如果这些项涉及未知数，例如圆柱坐标中粘性项的分量：

\begin{sourceblock}
\begin{tikzpicture}
\node[inner sep=0,anchor=south west] (image) at (0,0) {\includegraphics[page=200,trim={53.00bp 363.42bp 53.87bp 274.70bp},clip,width=0.92\textwidth,max height=0.38\textheight]{Ferziger4th.pdf}};
\begin{scope}[x={(image.south east)},y={(image.north west)}]
\fill[white] (0.00000,0.71842) rectangle (0.15056,1.00000);
\end{scope}
\end{tikzpicture}
\end{sourceblock}
他们可能会被隐含地对待。这通常是在该项对离散方程中的中心系数 A P 的贡献为正时进行的，以避免通过减少矩阵的对角优势来避免迭代求解方案的不稳定。否则，将明确处理额外的术语。

在 FV 方法中，这些项在 CV 体积上积分。通常，使用中点规则近似，以便 CV 中心的值乘以细胞体积。更复杂的方案是可能的，但很少使用。

在某些情况下，被视为体积力的非保守项主导着输运方程（例如，当在极坐标中计算旋流流时或在旋转坐标系中处理流时，例如在涡轮机械流中）。非线性源项和变量耦合的处理可能变得非常重要。

\subsection*{7.1.3 保护特性}\phantomsection\addcontentsline{toc}{subsection}{7.1.3 保护特性}
纳维-斯托克斯方程具有这样的属性：任何控制体积（微观或宏观）中的动量仅通过穿过表面的流动、作用在表面上的力和体积体力来改变。如果使用 FV 方法并且相邻控制体积的表面通量相同，则离散方程将继承这一重要属性。如果这样做，则整个域上的积分（即微观控制体积上的积分之和）会减少到域表面上的总和。总体质量守恒以相同的方式从连续性方程得出。

节能是一个更为复杂的问题。在不可压缩等温流中，唯一重要的能量是动能。 1 当传热很重要时，动能通常比热能小，因此为解释能量传输而引入的方程是热能守恒方程。只要流体性质对温度的依赖性不显著，热能方程就可以在动量方程求解完成后求解。这样，耦合就完全是单向的，并且能量方程变成了无源标量传输的方程，这种情况在第 1 章中进行了处理。 3 – 6。

动能方程可以通过动量方程与速度的标量积来导出，该过程模拟了经典力学中能量方程的推导。请注意，与可压缩流相比，总能量有一个单独的守恒方程，在不可压缩等温流中，动量和能量守恒都是同一个方程的结果；这提出了本节主题的问题。

我们主要对宏观控制体积的动能守恒方程感兴趣，该控制体积可以是整个考虑的域，也可以是有限体积方法中使用的小 CV 之一。如果将按上述方式获得的局部动能方程在控制体积上积分，则利用高斯定理我们可以得到：

\begin{sourceblock}
\begin{tikzpicture}
\node[inner sep=0,anchor=south west] (image) at (0,0) {\includegraphics[page=201,trim={53.00bp 278.91bp 53.87bp 324.62bp},clip,width=0.92\textwidth,max height=0.38\textheight]{Ferziger4th.pdf}};
\begin{scope}[x={(image.south east)},y={(image.north west)}]
\fill[white] (0.07414,0.69094) rectangle (0.09266,0.87558);
\fill[white] (0.06860,0.46223) rectangle (0.09786,0.65117);
\end{scope}
\end{tikzpicture}
\end{sourceblock}
这里S代表应力张量的粘性部分，其分量为式(1.13)中定义的τ i j，即S = T + p I。如果流动是无粘性的，则右侧体积积分中的第一项消失；如果流量不可压缩，则第二个为零；在没有体积力的情况下，第三个为零。

与该方程相关的几点值得一提。

\begin{itemize}
\item 其右侧的前三项是控制体积表面上的积分。这意味着控制体积内的动能不会因控制体积内的对流和/或压力的作用而改变。在没有粘度的情况下，只有通过表面的能量流或作用在控制体积表面的力所做的功才能影响其内部的动能；
\end{itemize}
1 在重力场中流体密度随高度变化的流动中，流动被称为分层流动，并且流动可能会携带重流体向上或携带轻流体向下，因此它现在具有与其周围环境不同的密度；这样，流体不仅具有动能，还具有由于其位置而产生的能量，称为势能。其结果是浮力，本章稍后将谈到，并且在气象学和海洋学中都非常重要。

从这个意义上来说，动能是全局守恒的。这是我们希望在数值方法中保留的属性。 • 用数值方法保证全球节能是一个值得实现的目标，但并不容易实现。由于动能方程是动量方程的结果，而不是独特的守恒定律，因此不能单独执行。 • 如果数值方法是能量守恒的并且通过表面的净能量通量为零，则域中的总动能不会随时间增长。如果使用这种方法，域中每个点的速度必须保持有界，从而提供一种重要的数值稳定性。事实上，能量方法（有时与物理无关）经常被用来证明数值方法的稳定性。能量守恒并没有说明方法的收敛性或准确性。通过动能不保守的方法可以获得精确的解。然而，动能守恒在计算非定常流动时尤其重要。 • 因为动能方程是动量方程的结果，并且不能在数值方法中独立执行，所以全局动能守恒必须是离散动量方程的结果。因此，它是离散化方法的一个属性，但不是一个明显的属性。为了了解它是如何出现的，我们通过取离散动量方程与速度的标量积并对所有控制体积求和，形成与离散动量方程相对应的动能方程。我们将逐项考虑结果。 • 压力梯度项特别重要，因此让我们进一步研究它们。为了将压力梯度项转化为方程（7.12）所示的形式，我们使用了以下等式：

\begin{sourceblock}
\includegraphics[page=202,trim={53.00bp 284.18bp 53.87bp 367.63bp},clip,width=0.92\textwidth,max height=0.38\textheight]{Ferziger4th.pdf}
\end{sourceblock}
对于不可压缩流，p ∇ · v = 0，因此仅保留右侧第一项。由于它是散度，因此其体积积分可以转换为表面积分。如前所述，这意味着压力仅通过其在表面的作用影响总动能收支。我们希望离散化能够保留这个属性。让我们看看这会如何发生。如果 G i p 表示压力梯度第 i 个分量的数值近似，则当离散化的 u i -动量方程乘以 u i 时，压力梯度项给出 u i G i p V 的贡献。能量守恒要求该贡献等于（参见方程（7.13））：

\begin{sourceblock}
\begin{tikzpicture}
\node[inner sep=0,anchor=south west] (image) at (0,0) {\includegraphics[page=202,trim={53.00bp 119.44bp 53.87bp 510.65bp},clip,width=0.92\textwidth,max height=0.38\textheight]{Ferziger4th.pdf}};
\begin{scope}[x={(image.south east)},y={(image.north west)}]
\fill[white] (0.26114,0.33897) rectangle (0.28824,0.68239);
\fill[white] (0.29251,0.33897) rectangle (0.32659,0.68239);
\fill[white] (0.32875,0.36172) rectangle (0.34854,0.68239);
\fill[white] (0.37456,0.36172) rectangle (0.39768,0.68239);
\fill[white] (0.40574,0.36893) rectangle (0.43392,0.68988);
\fill[white] (0.21756,0.04771) rectangle (0.25705,0.29265);
\end{scope}
\end{tikzpicture}
\end{sourceblock}
其中上标 N 表示总和是在所有 CV（网格节点）上进行的，S b 是解域的边界，v n 是垂直于边界的速度分量，D i u i 是连续性方程中使用的离散速度散度。如果是这样，则每个节点处的 D i u i = 0，因此上式中的第二项为零。只有当 G i 和 D i 在以下意义上兼容时，才能确保左侧和右侧的相等性：

\begin{sourceblock}
\begin{tikzpicture}
\node[inner sep=0,anchor=south west] (image) at (0,0) {\includegraphics[page=203,trim={53.00bp 526.43bp 53.87bp 104.17bp},clip,width=0.92\textwidth,max height=0.38\textheight]{Ferziger4th.pdf}};
\begin{scope}[x={(image.south east)},y={(image.north west)}]
\fill[white] (0.70448,0.34947) rectangle (0.77585,0.67473);
\end{scope}
\end{tikzpicture}
\end{sourceblock}
这表明，如果动能守恒成立，则压力梯度的近似值和速度的散度必须兼容。一旦选择了其中一个近似值，就失去了选择另一个近似值的自由。为了使这一点更具体，假设压力梯度是用后向差分来近似的，散度算子是用前向差分来近似的（交错网格上的通常选择）。均匀网格上方程（7.15）的一维版本如下：

\begin{sourceblock}
\includegraphics[page=203,trim={53.00bp 388.56bp 53.87bp 242.04bp},clip,width=0.92\textwidth,max height=0.38\textheight]{Ferziger4th.pdf}
\end{sourceblock}
求和时剩下的唯一两项是右侧的“表面项”。因此，这两个运算符在上述意义上是兼容的。相反，如果对压力梯度使用前向差分，则连续性方程将需要使用后向差分。如果其中一个使用中心差异，则另一个也需要它们。当对所有 CV（网格节点）求和时仅保留边界项的要求适用于其他两个保守项，即对流应力项和粘性应力项。在任何情况下，满足这一要求都不容易，并且对于任意和非结构化网格来说尤其困难（例如，参见 Mahesh 等人 2004，对于交错和共置的非结构化网格）。如果一种方法在均匀规则网格上不是能量守恒的，那么在更复杂的网格上也肯定不会如此。另一方面，在均匀网格上保守的方法在复杂网格上可能几乎如此。 • 泊松方程通常用于计算压力。正如我们将看到的，它是通过动量方程的散度导出的。因此，泊松方程中的拉普拉斯算子是连续性方程中的散度算子和动量方程中的梯度算子的乘积，即 L = D ( G ( ))。泊松方程的近似值不能独立选择；如果要获得质量守恒，它必须与散度和梯度算子一致。能量守恒进一步要求散度和梯度近似在上面定义的意义上保持一致。 • 对于没有体积力的不可压缩流动，唯一剩余的体积积分是粘性项。对于牛顿流体，该术语变为：

\begin{sourceblock}
\begin{tikzpicture}
\node[inner sep=0,anchor=south west] (image) at (0,0) {\includegraphics[page=204,trim={53.00bp 583.16bp 53.87bp 53.51bp},clip,width=0.92\textwidth,max height=0.38\textheight]{Ferziger4th.pdf}};
\begin{scope}[x={(image.south east)},y={(image.north west)}]
\fill[white] (0.40617,0.32983) rectangle (0.43435,0.72209);
\end{scope}
\end{tikzpicture}
\end{sourceblock}
检查发现，被积函数是平方和（参见方程（7.7）中τ i j 的定义），所以总是负（或零）。它代表流动的动能到流体内能的可逆（在热力学意义上）转换，称为粘性耗散。由于不可压缩流动通常是低速流动，因此内能的增加很少显著，但动能的损失通常对流动非常重要。在可压缩流中，能量传递通常对双方都很重要。 • 时差法会破坏能量守恒性质。除了上述空间离散化的要求外，还应适当选择时间导数的近似值。克兰克-尼科尔森方案是一个特别好的选择。初始，动量方程中的时间导数近似为：

\begin{sourceblock}
\begin{tikzpicture}
\node[inner sep=0,anchor=south west] (image) at (0,0) {\includegraphics[page=204,trim={53.00bp 395.48bp 53.87bp 242.37bp},clip,width=0.92\textwidth,max height=0.38\textheight]{Ferziger4th.pdf}};
\begin{scope}[x={(image.south east)},y={(image.north west)}]
\fill[white] (0.39519,0.54860) rectangle (0.41633,0.95758);
\fill[white] (0.44328,0.53941) rectangle (0.46638,0.94804);
\fill[white] (0.43720,0.04242) rectangle (0.45035,0.45104);
\end{scope}
\end{tikzpicture}
\end{sourceblock}
如果我们将此项与 u n + 1 / 2 进行标量积

\begin{sourceblock}
\begin{tikzpicture}
\node[inner sep=0,anchor=south west] (image) at (0,0) {\includegraphics[page=204,trim={53.00bp 321.10bp 53.87bp 289.99bp},clip,width=0.92\textwidth,max height=0.38\textheight]{Ferziger4th.pdf}};
\begin{scope}[x={(image.south east)},y={(image.north west)}]
\fill[white] (0.02779,0.95695) rectangle (0.05254,1.00000);
\fill[white] (0.05317,0.95695) rectangle (0.09293,1.00000);
\fill[white] (0.09353,0.95695) rectangle (0.14794,1.00000);
\fill[white] (0.14860,0.95695) rectangle (0.19002,1.00000);
\fill[white] (0.19065,0.95695) rectangle (0.26535,1.00000);
\fill[white] (0.26598,0.95695) rectangle (0.36232,1.00000);
\fill[white] (0.36301,0.95695) rectangle (0.39278,1.00000);
\fill[white] (0.39341,0.95695) rectangle (0.44150,1.00000);
\fill[white] (0.44214,0.95695) rectangle (0.50186,1.00000);
\fill[white] (0.50250,0.95695) rectangle (0.56057,1.00000);
\fill[white] (0.56117,0.95876) rectangle (0.64114,1.00000);
\fill[white] (0.57759,0.91971) rectangle (0.58824,1.00000);
\fill[white] (0.63795,0.95695) rectangle (0.65026,1.00000);
\fill[white] (0.65089,0.95695) rectangle (0.72893,1.00000);
\fill[white] (0.72956,0.95695) rectangle (0.75768,1.00000);
\fill[white] (0.75832,0.95695) rectangle (0.79973,1.00000);
\fill[white] (0.80036,0.95695) rectangle (1.00000,1.00000);
\fill[white] (0.02779,0.73987) rectangle (0.12241,0.94986);
\fill[white] (0.12571,0.73987) rectangle (0.15053,0.94986);
\fill[white] (0.15374,0.73987) rectangle (0.32328,0.94986);
\fill[white] (0.32659,0.73987) rectangle (0.36135,0.94986);
\fill[white] (0.36457,0.74169) rectangle (0.43540,0.99019);
\fill[white] (0.48523,0.70227) rectangle (0.49585,0.85813);
\fill[white] (0.49829,0.73987) rectangle (0.55098,0.95658);
\fill[white] (0.55426,0.73987) rectangle (0.59567,0.94986);
\fill[white] (0.59892,0.73987) rectangle (0.67029,0.94986);
\fill[white] (0.67356,0.73987) rectangle (0.69835,0.94986);
\fill[white] (0.70159,0.73987) rectangle (0.74304,0.94986);
\fill[white] (0.74626,0.73987) rectangle (0.83591,0.94986);
\fill[white] (0.83925,0.73987) rectangle (0.86737,0.94986);
\fill[white] (0.87062,0.73987) rectangle (0.91203,0.94986);
\fill[white] (0.91528,0.73987) rectangle (1.00000,0.94986);
\fill[white] (0.02779,0.52262) rectangle (0.12186,0.73261);
\fill[white] (0.64791,0.39019) rectangle (0.66322,0.54605);
\fill[white] (0.69278,0.15513) rectangle (0.70752,0.36530);
\end{scope}
\end{tikzpicture}
\end{sourceblock}
其中 v 2 = u i u i （隐含求和）。通过正确选择其他项的近似值，克兰克-尼科尔森方案是能量守恒的。

动量守恒和能量守恒均受同一个方程控制，这一事实使得构建保持这两种性质的数值近似变得困难。正如已经指出的，动能守恒不能独立执行。如果动量方程以强守恒形式编写并使用有限体积方法，则通常可以保证全局动量守恒。能量守恒方法的构建是一件偶然的事情。选择一种方法并确定它是否保守；如果没有，则进行调整，直到实现保护。

保证动能守恒的另一种方法是使用不同形式的动量方程。例如，对于不可压缩流，可以使用以下方程：

\begin{sourceblock}
\begin{tikzpicture}
\node[inner sep=0,anchor=south west] (image) at (0,0) {\includegraphics[page=204,trim={53.00bp 94.09bp 53.87bp 522.38bp},clip,width=0.92\textwidth,max height=0.38\textheight]{Ferziger4th.pdf}};
\begin{scope}[x={(image.south east)},y={(image.north west)}]
\fill[white] (0.19765,0.30723) rectangle (0.24087,0.56171);
\end{scope}
\end{tikzpicture}
\end{sourceblock}
其中 ε i jk 是 Levi-Civita 符号（如果 \{ i jk \} = \{ 123 \} 或其偶数排列，则为 + 1；如果 \{ i jk \} 是 \{ 123 \} 的奇数排列（例如 \{ 321 \}），则为 − 1，否则为零）。 ω 是方程（7.105）定义的涡度。能量守恒由这种形式的动量方程通过对称性得出；当方程乘以 u i 时，由于 ε i jk 的反​​对称性，左侧第二项同样为零。然而，由于这不是动量方程的保守形式，因此构建动量守恒方法需要小心。

动能守恒在计算复杂的非定常流动时特别重要。例子包括全球天气模式的模拟和湍流的模拟。这些模拟中缺乏有保证的能量守恒通常会导致动能增长和不稳定。对于稳态流，能量守恒不太重要，但它确实可以通过迭代求解方法防止某些类型的不当行为。

动能并不是唯一一个需要守恒但不能独立执行的量：角动量是另一个这样的量。旋转机械、内燃机和许多其他设备中的流动表现出明显的旋转或漩涡。如果数值方案不守恒全局角动量，计算可能会遇到麻烦。就角动量守恒而言，中心差分方案通常比迎风方案好得多。

\subsection*{7.1.4 网格上变量排列的选择}\phantomsection\addcontentsline{toc}{subsection}{7.1.4 网格上变量排列的选择}
现在让我们转向离散化。第一个问题是选择域中要计算未知因变量值的点。这比人们想象的要复杂得多。第 1 章概述了数值网格的基本特征。 2.然而，解域内计算点的分布有许多变体。与 FD 和 FV 离散化方法相关的基本布置如图 1 和 2 所示。 3.1 和 4.1。当求解矢量场的耦合方程（如纳维-斯托克斯方程）时，这些排列可能会变得更加复杂。这些问题将在下面讨论。

7.1.4.1 并置安排
显而易见的选择是将所有变量存储在同一组网格点处，并对所有变量使用相同的控制体积；这种安排称为并置，见图 7.1。由于每个方程中的许多项本质上是相同的，因此必须计算和存储的系数数量被最小化，并且通过这种选择简化了编程。此外，当使用多重网格过程时，用于在各个网格之间传输信息的相同限制和延长算子可以用于所有变量。

\begin{sourceblock}
\includegraphics[page=206,trim={53.00bp 506.25bp 53.87bp 52.00bp},clip,width=\textwidth,max height=0.38\textheight]{Ferziger4th.pdf}
\par\vspace{0.45\baselineskip}
\small 图 7.1 FD（左）和 FV（右）网格上速度分量和压力的同位排列
\end{sourceblock}
共置布置在复杂的解域中也具有显著的优势，特别是当边界具有坡度不连续性或边界条件不连续时。可以设计一组控制体积来拟合包括不连续性的边界。变量的其他排列导致一些变量位于网格的奇点处，这可能导致离散方程中的奇点。

由于压力-速度耦合的困难以及压力振荡的发生，同位布置在不可压缩流动计算中长期以来不受欢迎。从20世纪60年代中期引入交错电网到1980年代初，共置布置几乎没有被使用。然后，随着复杂几何形状的问题开始得到解决，非正交网格的使用变得更加普遍。仅当向量和张量的逆变（或其他面向网格）分量为工作变量时，交错方法才能在广义坐标中使用。这通过引入难以数值处理的曲率项使方程变得复杂，并且当网格不平滑时可能会产生非保守误差，如第 1 章所示。 9.当改进的压力-速度耦合算法在 20 世纪 80 年代开发出来时，并置布置的受欢迎程度开始上升。如今，所有主要商业和公共领域 CFD 代码都使用变量的并置排列。下面将展示其优点。

7.1.4.2 交错排列
不需要所有变量共享同一个网格；不同的安排可能是有利的。在笛卡尔坐标系中，Harlow 和 Welsh (1965) 引入的交错布置比共置布置具有多个优点。这种布置如图7.2所示。需要使用同位布置进行插值的几个项可以在不进行插值的情况下计算（到二阶近似）。这可以从图8.1所示的x动量CV看出。压力项和扩散项都很自然地通过中心差近似来近似，无需插值，因为压力项

\begin{sourceblock}
\includegraphics[page=207,trim={53.00bp 505.25bp 53.87bp 52.00bp},clip,width=\textwidth,max height=0.38\textheight]{Ferziger4th.pdf}
\par\vspace{0.45\baselineskip}
\small 图 7.2 FV 网格上速度分量和压力的完全交错排列（右）和部分交错排列（左）
\end{sourceblock}
某些节点位于 CV 面中心，并且扩散项所需的速度导数很容易在 CV 面上计算。此外，压力 CV 面上连续性方程中质量通量的评估也很简单。详细信息如下。

也许交错布置的最大优点是速度和压力之间的强耦合。这有助于避免某些类型的收敛问题以及压力场和速度场的振荡，这个问题将在稍后进一步讨论。

交错网格上的数值近似也是动能守恒的，这具有前面讨论过的优点。这个命题的证明很简单，但是很冗长，这里不再给出。

其他交错方法也被提出，例如，部分交错 ALE（任意拉格朗日-欧拉）方法（Hirt 等人，1997 年；Donea 等人，2004 年），其中两个速度分量都存储在压力 CV 的拐角处，见图 7.2。当网格非正交时，这种变体具有一些优点，其中一个重要的优点是不需要指定边界处的压力。然而，它也有缺点，特别是可能产生振荡压力或速度场。

其他安排还没有得到广泛普及，这里不再进一步讨论。

\subsection*{7.1.5 压力的计算}\phantomsection\addcontentsline{toc}{subsection}{7.1.5 压力的计算}
纳维-斯托克斯方程的解因缺乏独立的压力方程而变得复杂，压力的梯度对三个动量方程都有贡献。此外，虽然可压缩流的质量守恒方程以密度作为主导变量，但连续性方程在不可压缩流中没有主导变量，质量守恒成为速度场的运动学约束而不是动力学方程。解决这一困难的方法之一是构造压力场以保证连续性方程的满足。乍一看这可能有点奇怪，但我们将在下面证明这是可能的。请注意，绝对压力在不可压缩流中并不重要。只有压力梯度（压差）影响流量。

在可压缩流中，连续性方程可用于确定密度，压力可根据状态方程计算。这种方法不适用于不可压缩或低马赫数流。

在本节中，我们将介绍一些最流行的压力-速度耦合方法背后的基本原理。第 8 章介绍了离散方程组和其他指导，它们构成了编写计算机代码的基础。

7.1.5.1 压力方程及其解
动量方程清楚地确定了各自的速度分量，因此它们的作用得到了明确的定义。这就留下了不包含压力的连续性方程来确定压力。这怎么能做到呢？最常见的方法是将两个方程结合起来。

连续性方程的形式表明我们取动量散度方程(1.15)。连续性方程可用于简化所得方程，留下压力的泊松方程：

\begin{sourceblock}
\begin{tikzpicture}
\node[inner sep=0,anchor=south west] (image) at (0,0) {\includegraphics[page=208,trim={53.00bp 319.38bp 53.87bp 318.48bp},clip,width=0.92\textwidth,max height=0.38\textheight]{Ferziger4th.pdf}};
\begin{scope}[x={(image.south east)},y={(image.north west)}]
\fill[white] (0.17423,0.30163) rectangle (0.21811,0.71075);
\fill[white] (0.21997,0.29243) rectangle (0.29056,0.71075);
\fill[white] (0.29408,0.30163) rectangle (0.32226,0.71075);
\fill[white] (0.32577,0.30163) rectangle (0.39305,0.71075);
\end{scope}
\end{tikzpicture}
\end{sourceblock}
在笛卡尔坐标中，该方程为：

\begin{sourceblock}
\begin{tikzpicture}
\node[inner sep=0,anchor=south west] (image) at (0,0) {\includegraphics[page=208,trim={53.00bp 259.22bp 53.87bp 370.81bp},clip,width=0.92\textwidth,max height=0.38\textheight]{Ferziger4th.pdf}};
\begin{scope}[x={(image.south east)},y={(image.north west)}]
\fill[white] (0.10650,0.45251) rectangle (0.12502,0.77292);
\fill[white] (0.17089,0.44531) rectangle (0.20755,0.77292);
\fill[white] (0.09438,0.03323) rectangle (0.13600,0.38355);
\fill[white] (0.16713,0.03323) rectangle (0.20872,0.38355);
\end{scope}
\end{tikzpicture}
\end{sourceblock}
对于密度、粘度和体积力恒定的情况，该方程进一步简化；粘性项和不稳定项因连续性方程而消失，留下：

\begin{sourceblock}
\begin{tikzpicture}
\node[inner sep=0,anchor=south west] (image) at (0,0) {\includegraphics[page=208,trim={53.00bp 188.19bp 53.87bp 441.85bp},clip,width=0.92\textwidth,max height=0.38\textheight]{Ferziger4th.pdf}};
\begin{scope}[x={(image.south east)},y={(image.north west)}]
\fill[white] (0.00000,0.73657) rectangle (0.10003,1.00000);
\fill[white] (0.29898,0.45263) rectangle (0.31750,0.77285);
\fill[white] (0.36337,0.44515) rectangle (0.40003,0.77285);
\fill[white] (0.28686,0.03324) rectangle (0.32848,0.38338);
\end{scope}
\end{tikzpicture}
\end{sourceblock}
压力方程可以通过第 1 章中描述的椭圆方程数值方法之一来求解。 3和4。值得注意的是，压力方程的右侧是动量方程项的导数之和；这些必须以与导出它们的方程中的处理方式一致的方式进行近似。下面有两种获取压力的方法。在第一种情况下，我们寻求在每个方案的下一次迭代或下一个时间步骤中找到压力本身。或者，我们可以根据旧值（在先前的迭代级别或时间步）加上压力校正来定义新压力，即 p new = p old + p ′。然后我们寻求找到更小的压力修正 p ′。正如预期的那样，这种方法通常有优点。

压力方程中的拉普拉斯算子是源自连续性方程的散度算子和源自动量方程的梯度算子的乘积。在数值近似中，保持这些算子的一致性非常重要，即，泊松方程的近似被定义为基本方程中使用的散度和梯度近似的乘积。为了强调这个问题的重要性，将上述方程中压力的两个导数分开：外导数来自连续性方程，内导数来自动量方程。外导数和内导数可以使用不同的方案进行离散化——它们必须是动量方程和连续性方程中使用的方案。违反此约束会导致连续性方程不满足（不幸的是，我们稍后将看到（第 8.2 节）这是共置网格方法的一个问题）。

这种压力方程用于以显式和隐式求解方法计算压力。为了保持所使用的近似值之间的一致性，最好从离散动量和连续性方程导出压力方程，而不是通过近似上述泊松方程。压力方程还可用于计算通过求解涡流函数方程获得的速度场的压力，请参见第 1 节。 7.2.4.

7.1.5.2 关于压力和不可压缩性的注释
假设我们有一个速度场v*，它不满足连续性条件；例如，v*可能是通过对纳维-斯托克斯方程进行时间推进而无需调用连续性而获得的。我们希望创建一个新的速度场 v，其中：

\begin{itemize}
\item 满足连续性，
\item 尽可能接近原始场v *。
\end{itemize}
从数学上讲，我们可以将此问题视为最小化问题之一：

\begin{sourceblock}
\begin{tikzpicture}
\node[inner sep=0,anchor=south west] (image) at (0,0) {\includegraphics[page=209,trim={53.00bp 171.63bp 53.87bp 465.42bp},clip,width=0.92\textwidth,max height=0.38\textheight]{Ferziger4th.pdf}};
\begin{scope}[x={(image.south east)},y={(image.north west)}]
\fill[white] (0.37868,0.56136) rectangle (0.39847,0.95875);
\fill[white] (0.31525,0.32485) rectangle (0.33835,0.80303);
\fill[white] (0.34331,0.33379) rectangle (0.37149,0.73118);
\fill[white] (0.41874,0.04125) rectangle (0.43636,0.33585);
\end{scope}
\end{tikzpicture}
\end{sourceblock}
其中 r 是位置向量，V 是定义速度场的域，受连续性约束

\begin{sourceblock}
\includegraphics[page=209,trim={53.00bp 115.25bp 53.87bp 536.56bp},clip,width=0.92\textwidth,max height=0.38\textheight]{Ferziger4th.pdf}
\end{sourceblock}
在该领域的任何地方都感到满意。下面将讨论边界条件的问题。

这是变分法的标准类型问题。处理这个问题的一个有用方法是引入拉格朗日乘子。原问题（7.21）被最小化问题替代：

\begin{sourceblock}
\begin{tikzpicture}
\node[inner sep=0,anchor=south west] (image) at (0,0) {\includegraphics[page=210,trim={53.00bp 534.26bp 53.87bp 102.78bp},clip,width=0.92\textwidth,max height=0.38\textheight]{Ferziger4th.pdf}};
\begin{scope}[x={(image.south east)},y={(image.north west)}]
\fill[white] (0.24120,0.56151) rectangle (0.26099,0.95876);
\fill[white] (0.17774,0.32474) rectangle (0.20087,0.72199);
\fill[white] (0.20583,0.33402) rectangle (0.23402,0.73127);
\fill[white] (0.24111,0.07835) rectangle (0.26093,0.47560);
\end{scope}
\end{tikzpicture}
\end{sourceblock}
其中 λ 是拉格朗日乘子。包含拉格朗日乘数项不会影响最小值，因为约束 ( 7.22 ) 要求该项为零。

假设最小化函数 R 的函数是 v +；当然，v+
也满足(7.22)。因此：

\begin{sourceblock}
\begin{tikzpicture}
\node[inner sep=0,anchor=south west] (image) at (0,0) {\includegraphics[page=210,trim={53.00bp 429.57bp 53.87bp 207.49bp},clip,width=0.92\textwidth,max height=0.38\textheight]{Ferziger4th.pdf}};
\begin{scope}[x={(image.south east)},y={(image.north west)}]
\fill[white] (0.38722,0.56121) rectangle (0.40701,0.95873);
\fill[white] (0.29053,0.29402) rectangle (0.34653,0.72215);
\fill[white] (0.35185,0.33356) rectangle (0.38003,0.73143);
\fill[white] (0.42728,0.04127) rectangle (0.44490,0.33597);
\end{scope}
\end{tikzpicture}
\end{sourceblock}
如果 R min 是真正的最小值，则任何与 v + 的偏差都必定会导致 R 发生二阶变化。因此假设：

\begin{sourceblock}
\includegraphics[page=210,trim={53.00bp 372.19bp 53.87bp 275.51bp},clip,width=0.92\textwidth,max height=0.38\textheight]{Ferziger4th.pdf}
\end{sourceblock}
其中 δ v 很小但任意。将 v 代入表达式（7.23）时，结果为 R min + δ R 其中：

\begin{sourceblock}
\begin{tikzpicture}
\node[inner sep=0,anchor=south west] (image) at (0,0) {\includegraphics[page=210,trim={53.00bp 300.17bp 53.87bp 341.77bp},clip,width=0.92\textwidth,max height=0.38\textheight]{Ferziger4th.pdf}};
\begin{scope}[x={(image.south east)},y={(image.north west)}]
\fill[white] (0.09293,0.39008) rectangle (0.13257,0.87893);
\fill[white] (0.13753,0.40124) rectangle (0.16571,0.87893);
\end{scope}
\end{tikzpicture}
\end{sourceblock}
我们删除了与 (δ v ) 2 成比例的项，因为它是二阶的。现在，对最后一项进行分部积分并应用高斯定理，我们得到：

\begin{sourceblock}
\begin{tikzpicture}
\node[inner sep=0,anchor=south west] (image) at (0,0) {\includegraphics[page=210,trim={53.00bp 229.32bp 53.87bp 412.61bp},clip,width=0.92\textwidth,max height=0.38\textheight]{Ferziger4th.pdf}};
\begin{scope}[x={(image.south east)},y={(image.north west)}]
\fill[white] (0.04223,0.39033) rectangle (0.08186,0.87898);
\fill[white] (0.08683,0.40149) rectangle (0.11501,0.87898);
\end{scope}
\end{tikzpicture}
\end{sourceblock}
在给定 v 边界条件的域表面部分（壁、流入），假定 v 和 v + 都满足给定条件，因此 δ v 在那里为零。边界的这些部分对式（7.27）中的表面积分没有贡献，因此不需要对它们设置 λ 条件；然而，下面将提出一个条件。在给定其他类型边界条件的边界部分（对称平面、流出），δ v 不一定为零；为了使表面积分消失，我们需要要求边界的这些部分上 λ = 0。

如果 δ R 对于任意 δ v 都消失，我们必须要求方程（7.27）中的体积积分也消失，即：

\begin{sourceblock}
\includegraphics[page=210,trim={53.00bp 77.39bp 53.87bp 572.70bp},clip,width=0.92\textwidth,max height=0.38\textheight]{Ferziger4th.pdf}
\end{sourceblock}
最后，我们记得 v + ( r ) 必须满足连续性方程 ( 7.22 )。取方程（7.28）的散度并应用这个条件，我们发现：

\begin{sourceblock}
\includegraphics[page=211,trim={53.00bp 559.15bp 53.87bp 90.94bp},clip,width=0.92\textwidth,max height=0.38\textheight]{Ferziger4th.pdf}
\end{sourceblock}
这是 λ( r ) 的泊松方程。在 v 上给出边界条件的边界部分上，v + = v *。方程 (7.28) 表明，在这种情况下 ∇ λ( r ) = 0 并且我们在 λ 上有一个边界条件。

如果满足式(7.29)和边界条件，则速度场将是无散度的。值得注意的是，可以重复整个练习，并将连续运算符替换为离散运算符。

一旦求解了泊松方程，就可以从方程（7.28）中获得修正后的速度场，其形式为：

\begin{sourceblock}
\includegraphics[page=211,trim={53.00bp 427.64bp 53.87bp 222.44bp},clip,width=0.92\textwidth,max height=0.38\textheight]{Ferziger4th.pdf}
\end{sourceblock}
这表明拉格朗日乘子 λ(r) 本质上起着压力的作用，并再次表明，在不可压缩流动中，压力的作用是满足连续性。

\subsection*{7.1.6 纳维-斯托克斯方程的初始条件和边界条件}\phantomsection\addcontentsline{toc}{subsection}{7.1.6 纳维-斯托克斯方程的初始条件和边界条件}
在迭代求解过程开始之前，需要初始化所有变量的值。在稳态流的情况下，初始值对于最终解而言并不重要；然而，它们会影响收敛速度和整体计算工作。因此，需要初始化解，使得与最终解的差异尽可能小。然而，在许多实际情况下，很难提供初始字段来很好地估计最终解；在大多数情况下，我们从一些琐碎的初始化开始（例如，零或恒定的非零速度、恒定的压力和温度）。

当计算非定常流动时，对初始条件的要求更加严格。我们在 t = 0 时指定的速度场和压力场必须满足纳维-斯托克斯方程。由于 t = t 时的解强烈依赖于 t = 0 时的解，因此初始解中的任何错误都将被带入未来的时间步。仅当计算周期性流动（或具有随机性质的流动，例如湍流的直接或大涡流模拟，参见第 10 章）时，初始条件在一段时间后就会消失。无论如何，不​​适当的初始条件都会造成相当大的麻烦，因此需要小心。

每个时间步都需要应用边界条件；它们可能是恒定的，也可能随时间变化。第章中关于边界条件的所有内容。 3

\begin{center}\small 图7.3 壁面和对称面处速度的边界条件\end{center}
y y
v v u u
墙对称平面 x x
4 表示通用守恒方程，适用于动量方程。本节将讨论一些特殊功能。

在壁上，应用无滑移边界条件，即流体的速度等于壁速度，即狄利克雷边界条件。然而，还有另一个条件可以直接施加在 FV 方法中：壁上的法向粘性应力为零。这是从连续性方程得出的，例如，对于 y = 0 处的墙（见图 7.3）：

\begin{sourceblock}
\begin{tikzpicture}
\node[inner sep=0,anchor=south west] (image) at (0,0) {\includegraphics[page=212,trim={53.00bp 393.85bp 53.87bp 237.00bp},clip,width=0.92\textwidth,max height=0.38\textheight]{Ferziger4th.pdf}};
\begin{scope}[x={(image.south east)},y={(image.north west)}]
\fill[white] (0.08944,0.43242) rectangle (0.12535,0.76764);
\fill[white] (0.79976,0.24199) rectangle (0.82794,0.56957);
\fill[white] (0.83146,0.23151) rectangle (0.85125,0.55937);
\fill[white] (0.85495,0.24199) rectangle (0.86800,0.56957);
\fill[white] (0.92439,0.23151) rectangle (1.00000,0.55937);
\fill[white] (0.08935,0.03400) rectangle (0.12442,0.36923);
\fill[white] (0.75383,0.00000) rectangle (0.79453,0.21678);
\end{scope}
\end{tikzpicture}
\end{sourceblock}
因此，南边界v方程中的扩散通量为：

\begin{sourceblock}
\begin{tikzpicture}
\node[inner sep=0,anchor=south west] (image) at (0,0) {\includegraphics[page=212,trim={53.00bp 334.00bp 53.87bp 307.55bp},clip,width=0.92\textwidth,max height=0.38\textheight]{Ferziger4th.pdf}};
\begin{scope}[x={(image.south east)},y={(image.north west)}]
\fill[white] (0.36824,0.41033) rectangle (0.40460,0.95120);
\end{scope}
\end{tikzpicture}
\end{sourceblock}
这应该直接实现，而不是仅使用墙处 v = 0 的条件。因为在单元中心 v P ̸ = 0，如果不这样做，我们将在离散通量表达式中获得非零导数； v = 0 用作连续性方程中的边界条件。剪应力可以通过使用导数 ∂ u /∂ y 的单边近似来计算；一种可能的近似是（对于 u 方程和图 7.4 中的情况）：

\begin{sourceblock}
\begin{tikzpicture}
\node[inner sep=0,anchor=south west] (image) at (0,0) {\includegraphics[page=212,trim={53.00bp 214.84bp 53.87bp 413.89bp},clip,width=0.92\textwidth,max height=0.38\textheight]{Ferziger4th.pdf}};
\begin{scope}[x={(image.south east)},y={(image.north west)}]
\fill[white] (0.20797,0.27800) rectangle (0.24433,0.63325);
\fill[white] (0.22484,0.22240) rectangle (0.23780,0.45175);
\fill[white] (0.24977,0.28522) rectangle (0.27795,0.59423);
\end{scope}
\end{tikzpicture}
\end{sourceblock}
壁面处 u 对 y 的导数可以通过多种方式进行近似。一种方法是将节点S的位置设置为与面质心“s”相同，见图7.4的左侧。另一种方法是将节点S设置在计算域之外，就好像近边界单元在边界处进行镜像一样，见图7.4的右侧。

在第一个例子中，假设 u 随 y 线性变化的以下简单近似为：

\begin{sourceblock}
\begin{tikzpicture}
\node[inner sep=0,anchor=south west] (image) at (0,0) {\includegraphics[page=212,trim={53.00bp 84.17bp 53.87bp 544.54bp},clip,width=0.92\textwidth,max height=0.38\textheight]{Ferziger4th.pdf}};
\begin{scope}[x={(image.south east)},y={(image.north west)}]
\fill[white] (0.40099,0.46460) rectangle (0.43690,0.78066);
\end{scope}
\end{tikzpicture}
\par\vspace{0.45\baselineskip}
\small 图7.4 关于速度边界条件在笛卡尔网格上的实现
\end{sourceblock}
N

N

近界 CV
n

n

P e w
E

W

P e w
E

W

s

s

S

边界

S

因为这是单方面的近似，所以它只是一阶精确的。然而，如果在边界处使用半电池宽度上的这种近似以及内部的中心差近似，则解将以二阶收敛。基于 u 随 y 二次变化的假设（抛物线拟合边界位置“s”以及单元中心 P 和 N 处的值，参见图 7.4 的左侧），可以通过使用二阶单边近似来提高精度，对于均匀网格，其读数为：

\begin{sourceblock}
\begin{tikzpicture}
\node[inner sep=0,anchor=south west] (image) at (0,0) {\includegraphics[page=213,trim={53.00bp 371.72bp 53.87bp 257.00bp},clip,width=0.92\textwidth,max height=0.38\textheight]{Ferziger4th.pdf}};
\begin{scope}[x={(image.south east)},y={(image.north west)}]
\fill[white] (0.23074,0.46472) rectangle (0.26665,0.78087);
\end{scope}
\end{tikzpicture}
\end{sourceblock}
如果“边界”节点放置在解域之外，则可以在边界面“s”和内面“n”处使用 u 对 y 的导数的相同中心差分近似，即：

\begin{sourceblock}
\begin{tikzpicture}
\node[inner sep=0,anchor=south west] (image) at (0,0) {\includegraphics[page=213,trim={53.00bp 288.89bp 53.87bp 339.82bp},clip,width=0.92\textwidth,max height=0.38\textheight]{Ferziger4th.pdf}};
\begin{scope}[x={(image.south east)},y={(image.north west)}]
\fill[white] (0.32743,0.46460) rectangle (0.36334,0.78066);
\end{scope}
\end{tikzpicture}
\end{sourceblock}
然而，边界条件指定 u s = u wall，因此需要使用指定的边界值和内部单元中心的一个或多个值通过外推法获得辅助节点 S 处的 u 值。基于线性外推法的最简单近似可得出：

\begin{sourceblock}
\includegraphics[page=213,trim={53.00bp 207.74bp 53.87bp 442.95bp},clip,width=0.92\textwidth,max height=0.38\textheight]{Ferziger4th.pdf}
\end{sourceblock}
使用二次外推法可以得到更准确的近似值（与式（7.35）中导数的近似值相同）；对于均匀网格，可以得到：

\begin{sourceblock}
\includegraphics[page=213,trim={53.00bp 125.27bp 53.87bp 512.85bp},clip,width=0.92\textwidth,max height=0.38\textheight]{Ferziger4th.pdf}
\end{sourceblock}
在容易导出的非均匀网格的相应表达式中，节点变量值的乘数变成网格间距的函数。

在对称平面上，我们有相反的情况：剪应力为零，但法向应力不是，因为（对于图 7.3 中的情况）：

\begin{sourceblock}
\begin{tikzpicture}
\node[inner sep=0,anchor=south west] (image) at (0,0) {\includegraphics[page=214,trim={53.00bp 583.56bp 53.87bp 50.00bp},clip,width=0.92\textwidth,max height=0.38\textheight]{Ferziger4th.pdf}};
\begin{scope}[x={(image.south east)},y={(image.north west)}]
\fill[white] (0.31991,0.46839) rectangle (0.35585,0.83149);
\fill[white] (0.64156,0.26212) rectangle (0.66974,0.61725);
\fill[white] (0.67326,0.25077) rectangle (0.69305,0.60589);
\fill[white] (0.69675,0.26212) rectangle (0.70980,0.61725);
\fill[white] (0.92439,0.25077) rectangle (1.00000,0.60589);
\fill[white] (0.32042,0.03683) rectangle (0.35537,0.39994);
\fill[white] (0.59534,0.00000) rectangle (0.63510,0.31584);
\end{scope}
\end{tikzpicture}
\end{sourceblock}
u 方程中的扩散通量为零，v 方程中的扩散通量需要 v 对 y 的导数的近似：

\begin{sourceblock}
\begin{tikzpicture}
\node[inner sep=0,anchor=south west] (image) at (0,0) {\includegraphics[page=214,trim={53.00bp 509.04bp 53.87bp 119.68bp},clip,width=0.92\textwidth,max height=0.38\textheight]{Ferziger4th.pdf}};
\begin{scope}[x={(image.south east)},y={(image.north west)}]
\fill[white] (0.21426,0.27819) rectangle (0.25062,0.63335);
\fill[white] (0.23110,0.22261) rectangle (0.24409,0.45190);
\fill[white] (0.25606,0.28541) rectangle (0.28424,0.59433);
\end{scope}
\end{tikzpicture}
\end{sourceblock}
边界条件指定 v s = 0（没有流过对称平面）。上面针对壁边界处的 (∂ u /∂ y ) s 引入的相同近似可以应用于对称边界处的 v。

在使用交错网格的 FV 方法中，边界处不需要压力（除非指定了边界处的压力；这将在第 11 章中处理）。这是因为垂直于边界的速度分量的最近 CV 仅延伸到标量 CV 的中心，在此计算压力。当使用并置布置时，所有 CV 都会延伸到边界，我们需要边界压力来计算动量方程中的压力。我们必须使用内部外推法来获得边界处的压力。在大多数情况下，线性外推法对于二阶方法来说足够准确，但二次外推法甚至更好。然而，在某些情况下，法向速度分量的方程中需要靠近壁面的大压力梯度来平衡体积力（浮力、离心力等）。如果压力外推不准确，则可能无法满足该条件，并且可能会导致边界附近出现较大的法向速度。通过根据连续性方程计算第一个 CV 的法向速度分量、调整压力外推法或通过局部网格细化，可以避免这种情况。

压力校正方程的边界条件（在 7.2.2 节中定义）也值得关注。当规定通过边界的质量通量时，压力校正方程中的质量通量校正为零。在推导压力校正方程时，应直接在连续性方程中应用该条件。它相当于为压力校正指定诺伊曼边界条件（零梯度）。第 2 节讨论了分数阶方法中压力校正的边界条件。 8.3.4.

在出口处，如果给定入口质量通量，当流出边界远离感兴趣区域且雷诺数很大时，通常可以使用边界速度外推（零梯度，例如 u E = u P ）来获得稳态流。然后校正外推速度以给出与入口处完全相同的总质量通量（这不能通过任何外推来保证）。然后，将校正后的速度视为以下外部迭代的规定，并且在连续性方程中将流出边界处的质量通量校正设置为零。这导致压力校正方程在所有边界上都具有诺依曼条件，并使其在数学上是奇异的。为了使解唯一，可以将某一点的压力固定，从而从所有校正压力中减去该点计算的压力校正。另一种选择是将平均压力设置为某个值，例如零。实际上，没有必要采取这些步骤，因为大多数方法都以起始压力开始每个步骤，并且实际压力本身通常在不可压缩流中并不重要，因为只有梯度很重要。在已知或需要实际压力的情况下，它通常由物理情况定义，例如，表面压力已知的自由表面流。常识表明无论如何都要小心。如果与梯度中网格点之间的压力差相比，平均压力变得非常大，则数值精度可能会下降，并且在某个点分配固定压力将防止这种情况发生。

当指定入口和出口边界之间的压力差时，会得到另一种情况。那么这些边界处的速度就无法指定——必须计算它们以使压力损失达到指定值。这可以通过多种方式实现。在任何情况下，必须从内部节点推断边界速度（以类似于并置布置中单元面的插值的方式），然后进行校正。第 1 章给出了如何处理指定静压的示例。 11.

\subsection*{7.1.7 示例性简单方案}\phantomsection\addcontentsline{toc}{subsection}{7.1.7 示例性简单方案}
为了开始纳维-斯托克斯方程的实际求解，我们描述了两个简单的方案。第一种是针对非定常流的显式时间步进方法。第二种方案向我们介绍了隐式方案的附加功能。然后在第 4 节中开始对常见方案的详细检查。 7.2.

7.1.7.1 一种简单的显式时间推进方法
让我们看一下非定常方程的方法，该方法说明了如何构造压力的数值泊松方程以及它在增强连续性方面所起的作用。空间导数的近似值的选择在这里并不重要，因此半离散（空间上离散，但时间上不离散）动量方程可以象征性地写为：

\begin{sourceblock}
\includegraphics[page=215,trim={53.00bp 151.80bp 53.87bp 485.23bp},clip,width=0.92\textwidth,max height=0.38\textheight]{Ferziger4th.pdf}
\end{sourceblock}
其中 δ/δ x 表示离散空间导数（可以表示每个项中的不同近似值），H i 是平流项和粘性项的简写符号，其处理在这里并不重要。

为简单起见，假设我们希望使用时间推进的显式欧拉方法求解方程（7.41）。然后我们有：

\begin{sourceblock}
\begin{tikzpicture}
\node[inner sep=0,anchor=south west] (image) at (0,0) {\includegraphics[page=216,trim={53.00bp 582.75bp 53.87bp 53.25bp},clip,width=0.92\textwidth,max height=0.38\textheight]{Ferziger4th.pdf}};
\begin{scope}[x={(image.south east)},y={(image.north west)}]
\fill[white] (0.25218,0.27439) rectangle (0.36153,0.75116);
\fill[white] (0.36490,0.27439) rectangle (0.47648,0.74486);
\fill[white] (0.48244,0.31022) rectangle (0.51062,0.69409);
\fill[white] (0.53823,0.30159) rectangle (0.55137,0.68514);
\end{scope}
\end{tikzpicture}
\end{sourceblock}
为了应用此方法，时间步 n 处的速度用于计算 H n
i，如果压力可用，也可以计算 δ p n /δ x i。这给出了新时间步 n + 1 处的 ρ u i 估计。一般来说，该速度场不满足我们想要强制执行的连续性方程：

\begin{sourceblock}
\includegraphics[page=216,trim={53.00bp 487.18bp 53.87bp 148.63bp},clip,width=0.92\textwidth,max height=0.38\textheight]{Ferziger4th.pdf}
\end{sourceblock}
我们已经表达了对不可压缩流的兴趣，但其中包括具有变密度的流；通过包括密度来强调这一点。为了了解如何强制连续性，让我们采用方程（7.42）的数值散度（使用用于近似连续性方程的数值运算符）。结果是：

\begin{sourceblock}
\begin{tikzpicture}
\node[inner sep=0,anchor=south west] (image) at (0,0) {\includegraphics[page=216,trim={53.00bp 392.77bp 53.87bp 237.26bp},clip,width=0.92\textwidth,max height=0.38\textheight]{Ferziger4th.pdf}};
\begin{scope}[x={(image.south east)},y={(image.north west)}]
\fill[white] (0.19254,0.42260) rectangle (0.31639,0.80670);
\fill[white] (0.23365,0.03323) rectangle (0.27320,0.38355);
\end{scope}
\end{tikzpicture}
\end{sourceblock}
第一项是新速度场的散度，我们希望其为零。如果在时间步 n 强制执行连续性，则第二项为零；我们假设情况确实如此，但如果情况并非如此，则该项应保留在方程中。当采用迭代方法求解压力泊松方程且迭代过程未完全收敛时，需要保留此项。类似地，对于常数 ρ，H i 的粘性分量的散度应该为零，但很容易考虑非零值。考虑到所有这些，结果是压力 p n 的离散泊松方程：

\begin{sourceblock}
\begin{tikzpicture}
\node[inner sep=0,anchor=south west] (image) at (0,0) {\includegraphics[page=216,trim={53.00bp 250.52bp 53.87bp 379.51bp},clip,width=0.92\textwidth,max height=0.38\textheight]{Ferziger4th.pdf}};
\begin{scope}[x={(image.south east)},y={(image.north west)}]
\fill[white] (0.38388,0.45251) rectangle (0.40156,0.77264);
\fill[white] (0.37179,0.03323) rectangle (0.41134,0.38355);
\end{scope}
\end{tikzpicture}
\end{sourceblock}
请注意，括号外的算子 δ/δ x i 是从连续性方程继承的散度算子，而 δ p /δ x i 是动量方程的压力梯度。如果压力 p n 满足离散泊松方程，则时间步 n + 1 处的速度场将是无散的（就离散散度算子而言）。请注意，该压力所属的时间步长是任意的。如果隐式处理压力梯度项，我们将得到 p n + 1
代替 p n 但其他一切都保持不变。

这为纳维-斯托克斯方程的时移提供了以下算法：

\begin{itemize}
\item 从时间t n 处的速度场u n i 开始，假设无散度。 （如前所述，如果它不是无散度的，则可以对其进行纠正。）
\item 计算平流项和粘性项的组合 H n i 及其散度（两者都需要保留以供以后使用）。
\end{itemize}
\begin{itemize}
\item 求解压力p n 的泊松方程。
\item 计算新时间步的速度场。它将是无分歧的。
\item 现在已为下一个时间步骤做好准备。
\end{itemize}
当需要精确的流动时程时，通常使用与此类似的方法来求解纳维-斯托克斯方程。实践中的主要区别在于，使用了比一阶欧拉方法更准确的时间提前方法，并且某些术语可能会被隐式处理。稍后将描述其中一些方法。

我们已经展示了求解压力的泊松方程如何确保速度场满足连续性方程，即它是无散度的。这个想法贯穿于许多用于求解稳态和非稳态纳维-斯托克斯方程的方法中。

7.1.7.2 一种简单的隐式时间推进方法
为了了解使用隐式方法求解纳维-斯托克斯方程时会出现哪些额外困难，让我们构建这样一个方法。因为我们有兴趣阐明某些问题，所以让我们使用基于最简单的隐式方法（向后或隐式欧拉方法）的方案。如果我们将此方法应用于方程（7.41），我们有：

\begin{sourceblock}
\begin{tikzpicture}
\node[inner sep=0,anchor=south west] (image) at (0,0) {\includegraphics[page=217,trim={53.00bp 310.83bp 53.87bp 319.21bp},clip,width=0.92\textwidth,max height=0.38\textheight]{Ferziger4th.pdf}};
\begin{scope}[x={(image.south east)},y={(image.north west)}]
\fill[white] (0.10382,0.22909) rectangle (0.21314,0.62715);
\fill[white] (0.21651,0.22909) rectangle (0.32809,0.62188);
\fill[white] (0.33405,0.25928) rectangle (0.36223,0.57950);
\fill[white] (0.38983,0.25180) rectangle (0.40298,0.57202);
\end{scope}
\end{tikzpicture}
\end{sourceblock}
我们立即发现存在上一节中描述的显式方法中不存在的困难。让我们一次考虑这些。

首先，压力有问题。新时间步的速度场散度必须为零。这可以通过与显式方法大致相同的方式来完成。我们取方程（7.46）的散度，假设时间步n处的速度场是无散度的（如果需要的话可以进行修正），并要求新时间步n+1处的散度也为零。由此得出压力的泊松方程：

\begin{sourceblock}
\includegraphics[page=217,trim={53.00bp 166.14bp 53.87bp 463.90bp},clip,width=0.92\textwidth,max height=0.38\textheight]{Ferziger4th.pdf}
\end{sourceblock}
问题是，只有在 n + 1 时刻的速度场计算完成后，才能计算右侧的项，反之亦然。因此，必须同时求解泊松方程和动量方程。这只能通过某种类型的迭代过程来完成。

接下来，即使压力已知，方程（7.46）也是一个大型非线性方程组，必须求解速度场。该方程组的结构类似于压力方程的有限差分拉普拉斯算子的矩阵结构。然而，由于动量方程包含平流项的贡献，并且狄利克雷边界条件应用于某些边界（例如入口、壁、对称面），因此它们的求解通常比压力或压力校正方程的求解更容易。这将在本章末尾的示例中进行演示。

如果希望准确地求解方程（7.46）（下一节将介绍非迭代分步法的情况），最好的程序是采用前一时间步的收敛结果作为新速度场的初始猜测，然后使用牛顿迭代法（第5.5.1节）或为系统设计的割线法（Ferziger 1998；莫因 2010）。

在了解了如何构造显式和隐式方案来求解纳维-斯托克斯方程之后，我们现在将研究一些更常用的求解它们的方法。我们的报道并不详尽，其他方法可以在文献中找到，例如 Chang 等人的精确投影方法。 （2002）。

\section*{7.2 稳态流和非稳态流的计算策略}\phantomsection\addcontentsline{toc}{section}{7.2 稳态流和非稳态流的计算策略}
本节介绍一组常用的解决方法。给出了适用于稳态流和非稳态流的方法，并使用迭代和非迭代方法。皆与上一宗中所阐述的元素相似。 7.1.7.特别是，对于不可压缩流，稳态和非稳态问题均以行进方案（及时或通过一组迭代）的形式求解，从而需要求解泊松方程（对于压力、压力校正或流函数）。下面将类似的方法分组在一起。

\subsection*{7.2.1 分数步法}\phantomsection\addcontentsline{toc}{subsection}{7.2.1 分数步法}
Harlow 和 Welsh (1965) 以及 Chorin (1968) 的开创性论文记录了以分离方式对非定常纳维-斯托克斯方程进行时间前向积分的想法。本文描述的许多方法都源自或建立在这些早期工作的基础上。事实上，Patankar 和 Spalding (1972) 注意到这些论文对 SIMPLE 算法发展的影响（在下一节中描述）。虽然 SIMPLE 是一种迭代方法，但这种方法不是必需的，实际上甚至可以忽略预测器步骤中的压力。然而，所有这些方法以及一般而言后续的方法都使用不可压缩流中的压力来强制连续性。 Armfield (1991, 1994) 以及 Armfield 和 Street (2002) 为许多分数步方法提供了背景和背景。

分数步法的基本主题是（i）使用可用压力信息（如果需要）和当前速度场找到下一时间步速度场的准确估计，（ii）求解新压力的泊松方程或修正旧压力的方程，以及（iii）使用新压力或校正压力将下一时间步的估计速度更新为新的质量守恒速度。没有使用迭代，但是可以执行，我们稍后会讨论。

Kim 和 Moin (1985) 使用分数步骤创建了交错网格上不可压缩流动的方案，其中压力未在预测器中使用，而是计算出来，然后在校正器步骤中使用。臧等人。 ( 1994 ) 将该方法扩展到并置网格和曲线坐标。 Kim、Moin 和 Zang 等人。在预测器步骤中使用近似因式分解 ADI 方案（参见例如 Beam and Warming 1976 或 5.3.5 节），以允许通过每个坐标方向上的三对角矩阵的解来估计速度场；因式分解方法本身的误差为 O (( t ) 3 )，即比整个方案期望的 O (( t ) 2 ) 小一阶。这些公式中重要的是对粘性（或湍流）项的隐式处理；这避免了对类似扩散项的时间步长限制，这对于显式方案来说非常严格（参见 Ferziger 1998）。这些方法已被证明很受欢迎并已被广泛使用，并且近似因式分解 ADI 方案在其他分数步方法中很常见。此外，臧等人。 (1994) 使用紧凑的压力模板来进行 FV 公式计算，并将同位网格上以 CV 为中心的速度插值到体积面上，以帮助应用连续性约束。叶等人。 (1999) 讨论了这种方法的价值，Kim 和 Choi (2000) 将其扩展到非结构化网格。

Gresho (1990) 为分数步方案创造了名称，它们可用作速记。 P1 方法将用于估计新速度场的动量方程中的压力场设置为零，然后求解压力泊松方程以获得新压力。 Kim 和 Moin (1985) 以及 Zang 等人。 (1994) 方案是 P1，并求解与很少需要或计算的实际压力相关的伪压力。 P2 方法将动量方程中的压力设置为等于前一时间步的值，然后求解压力泊松方程以找到压力校正。在 P3 方法中，动量方程中使用的压力是通过二阶精确方案从前面的两个压力场推断出来的。我们稍后将评论最后一个方案。下面，我们将在本节中对P2方案进行一般性描述；见第 8.3 节 我们给出了交错电网和共置电网方案实施的一些细节；最后见第 8.4 节 我们检查了盖子中的稳态和非稳态流动以及矩形外壳中的浮力驱动流动所获得的结果，包括与替代解决方法的直接比较。

在这里，我们紧随 Armfield 和 Street（2000、2003、2004）的论文描述了分数步方法。 Armfield 和 Street (2003) 的一个重要结果（图 7.5）是，在 P2 方案中，压力在时间上是二阶精确的。 2 此外，我们以向量和算子的形式编写方程，因为这样编写时特别容易理解方法的设置和流程。

控制方程是不可压缩流体的三维纳维-斯托克斯方程（参见第 1.3 节和第 1.4 节以及参见方程（7.3）和（7.4））：

\footnotetext[2]{在这种情况下，使用二阶外推 p n + 1 = ( 3 / 2 ) p n + 1 / 2 − ( 1 / 2 ) p n − 1 / 2。}
\begin{sourceblock}
\includegraphics[page=220,trim={53.00bp 473.25bp 53.87bp 52.00bp},clip,width=\textwidth,max height=0.38\textheight]{Ferziger4th.pdf}
\par\vspace{0.45\baselineskip}
\small 图 7.5 腔体内自然对流的各种方法的压力精度（来自 Armfield 和 Street 2003；经许可转载）
\end{sourceblock}
\begin{sourceblock}
\begin{tikzpicture}
\node[inner sep=0,anchor=south west] (image) at (0,0) {\includegraphics[page=220,trim={53.00bp 372.83bp 53.87bp 244.98bp},clip,width=0.92\textwidth,max height=0.38\textheight]{Ferziger4th.pdf}};
\begin{scope}[x={(image.south east)},y={(image.north west)}]
\fill[white] (0.28310,0.73040) rectangle (0.35994,0.97517);
\end{scope}
\end{tikzpicture}
\end{sourceblock}
其中 v = ( u i ) 是速度，p 是压力，并且与之前一样，S = T + p I 表示应力张量的粘性部分（参见方程（1.9）和（1.13））。请注意，我们在这里忽略了身体力，因为它们对方法没有影响。

该方法的目标是空间和时间上的二阶精度。为此，对流项使用二阶 Adams-Bashforth 方法（第 6.2.2 节），对粘性项使用二阶 Adams-Bashforth 方法（第 6.3.2.2 节），对方程（7.48）和（7.49）进行时间离散。因此，该方法是半隐式的。离散化方程可以写为：

\begin{sourceblock}
\begin{tikzpicture}
\node[inner sep=0,anchor=south west] (image) at (0,0) {\includegraphics[page=220,trim={53.00bp 218.03bp 53.87bp 416.10bp},clip,width=0.92\textwidth,max height=0.38\textheight]{Ferziger4th.pdf}};
\begin{scope}[x={(image.south east)},y={(image.north west)}]
\fill[white] (0.05362,0.47985) rectangle (0.15263,0.88785);
\fill[white] (0.15597,0.47985) rectangle (0.25720,0.88191);
\end{scope}
\end{tikzpicture}
\end{sourceblock}
\begin{sourceblock}
\begin{tikzpicture}
\node[inner sep=0,anchor=south west] (image) at (0,0) {\includegraphics[page=220,trim={53.00bp 189.24bp 53.87bp 458.45bp},clip,width=0.92\textwidth,max height=0.38\textheight]{Ferziger4th.pdf}};
\begin{scope}[x={(image.south east)},y={(image.north west)}]
\fill[white] (0.40319,0.07046) rectangle (0.52508,0.80542);
\fill[white] (0.53011,0.08509) rectangle (0.55829,0.71165);
\fill[white] (0.56180,0.06504) rectangle (0.58159,0.69214);
\fill[white] (0.58526,0.08509) rectangle (0.60000,0.71165);
\fill[white] (0.92439,0.06504) rectangle (1.00000,0.69214);
\end{scope}
\end{tikzpicture}
\end{sourceblock}
其中离散项包括速度 v 、压力 p 、对流算子 C ( v ) = ∇ · (ρ vv ) 、梯度算子 G () = G i () = ∇ () (Eq.( 1.15 ))、散度算子 D () = Di () = ∇ · () (Eq.( 1.6 )) 以及粘性项 L ( v ) = ∇ · S。如果使用标准二阶中心差分 (CDS)，则当以 n + 1 / 2 时间步为中心时，该方案在时间和空间上都是二阶。请注意 Adams-Bashforth 表示得出

\begin{sourceblock}
\begin{tikzpicture}
\node[inner sep=0,anchor=south west] (image) at (0,0) {\includegraphics[page=220,trim={53.00bp 69.79bp 53.87bp 568.33bp},clip,width=0.92\textwidth,max height=0.38\textheight]{Ferziger4th.pdf}};
\begin{scope}[x={(image.south east)},y={(image.north west)}]
\fill[white] (0.20493,0.29872) rectangle (0.33035,0.78266);
\fill[white] (0.33383,0.30835) rectangle (0.38899,0.95717);
\end{scope}
\end{tikzpicture}
\end{sourceblock}
我们的分数步方法现在通过以下方式求解方程（7.50）和（7.51）

1. 通过求解以下方程，使用旧压力值 p n − 1 / 2 找到新速度 v * 的估计值：

\begin{sourceblock}
\includegraphics[page=221,trim={53.00bp 529.08bp 53.87bp 107.44bp},clip,width=0.92\textwidth,max height=0.38\textheight]{Ferziger4th.pdf}
\end{sourceblock}
2. 根据下式定义压力修正 p '

\begin{sourceblock}
\includegraphics[page=221,trim={53.00bp 484.73bp 53.87bp 162.97bp},clip,width=0.92\textwidth,max height=0.38\textheight]{Ferziger4th.pdf}
\end{sourceblock}
3. 接下来定义对 v * 的修正，要求修正后的速度满足以下近似版本的方程（7.50）（其中 L 算子中的 v n + 1 被 v * 取代，原因很快就会变得显而易见）：

\begin{sourceblock}
\includegraphics[page=221,trim={53.00bp 399.51bp 53.87bp 234.63bp},clip,width=0.92\textwidth,max height=0.38\textheight]{Ferziger4th.pdf}
\end{sourceblock}
将式（7.55）减去式（7.53），得到修正后的速度表达式如下：

\begin{sourceblock}
\includegraphics[page=221,trim={53.00bp 343.74bp 53.87bp 303.96bp},clip,width=0.92\textwidth,max height=0.38\textheight]{Ferziger4th.pdf}
\end{sourceblock}
最后 4. 将式（ 7.56 ）代入连续性方程（ 7.51 ）求 p ′，得

\begin{sourceblock}
\includegraphics[page=221,trim={53.00bp 271.93bp 53.87bp 364.70bp},clip,width=0.92\textwidth,max height=0.38\textheight]{Ferziger4th.pdf}
\end{sourceblock}
已知 p ′，最终速度 v n + 1 和压力 p n + 1 / 2 分别由方程（7.56）和（7.54）给出。在典型情况下，对流项使用QUICK方案在空间中离散化（第4.4.3节），其他项使用CDS离散化，使用ADI方案（第5.3.5或8.3.2节）求解方程（7.53），并且通过泊松求解器求解方程（7.57），例如，预处理重新启动 GMRES（第 5.3.6.3 节）。 Armfield 和 Street（1999）报告称，ADI 方案的四次扫描足以为他们测试的情况获得准确的解。

一旦解前进了一个时间步长，就可以确定是否犯了任何错误。我们首先将方程（7.56）改写为：

\begin{sourceblock}
\includegraphics[page=221,trim={53.00bp 130.94bp 53.87bp 516.76bp},clip,width=0.92\textwidth,max height=0.38\textheight]{Ferziger4th.pdf}
\end{sourceblock}
现在将其插入修正后的速度和压力满足的差值方程（7.55）中（使用方程（7.54）中的定义）以获得

\begin{sourceblock}
\begin{tikzpicture}
\node[inner sep=0,anchor=south west] (image) at (0,0) {\includegraphics[page=222,trim={53.00bp 557.04bp 53.87bp 51.45bp},clip,width=0.92\textwidth,max height=0.38\textheight]{Ferziger4th.pdf}};
\begin{scope}[x={(image.south east)},y={(image.north west)}]
\fill[white] (0.77735,0.46574) rectangle (0.79714,0.66626);
\fill[white] (0.61985,0.14987) rectangle (0.66671,0.46522);
\end{scope}
\end{tikzpicture}
\end{sourceblock}
将该方程（我们已经实际求解）与方程（7.50）（我们想要求解）进行比较表明，求解的方程包含一个额外的项——上式中的最后一项——它表示由于不校正 L 算子中的 v ∗ 而引入的误差；回想一下，本质上，上面最后一项中的 L = D ( G ())。通过认识到压力修正可以表示为

\begin{sourceblock}
\includegraphics[page=222,trim={53.00bp 451.32bp 53.87bp 186.53bp},clip,width=0.92\textwidth,max height=0.38\textheight]{Ferziger4th.pdf}
\end{sourceblock}
我们看到这个附加项与 ( t ) 2 成正比，因此与基本离散化一致（Armfield 和 Street 2002）。总之，无论网格是交错的还是非交错的，这里描述的 P2 方法都不能提供离散方程的精确解。虽然速度场将是无散的（达到压力泊松解的精度），但此误差意味着更新的速度场加上更新的压力并不完全满足离散动量方程。虽然误差在时间上是二阶的（意味着当时间步长减半时误差会减小 4 倍），但如果所选时间步长不够小，误差可能仍然很大。不幸的是，什么足够小取决于问题。

上述方法仅代表众多可能选择中的一种。显然，人们可以轻松地为对流项实现一种不同的显式时间提前方案，而无需改变压力校正方程。请注意，在上述算法中，压力在 t n + 1 / 2 时计算，而速度在 t n + 1 时计算；如果在同一时间水平需要这两个量，则必须对其中一个进行插值。然而，我们可以将曲柄尼科尔森时间推进方案应用于压力项，就像对粘性项所做的那样；唯一的区别是，现在一半的压力项（来自时间水平 t n ）将被固定，并且只有隐含的一半（来自时间水平 t n + 1 ）将被纠正。

人们还可以使用显式方案及时推进粘性项；如果解域中没有极小的单元，则这可能是可以接受的，因此时间步长的稳定性限制不会太严格。在这种情况下，可以显式求解动量方程（7.50），而无需求解方程组。

如果对流项和扩散项都选择隐式时间提前，事情会变得有点复杂。现在非线性对流项需要线性化。皮卡德迭代方案是通常的选择，但是仍然有很多可能性来近似假设已知的部分。这很大程度上取决于是否要在一个时间步长内迭代，或者动量方程是否只求解一次。在前一种情况下，选择是显而易见的：使用上一次迭代的值：

\begin{sourceblock}
\includegraphics[page=223,trim={53.00bp 595.02bp 53.87bp 52.68bp},clip,width=0.92\textwidth,max height=0.38\textheight]{Ferziger4th.pdf}
\end{sourceblock}
这里 m 是迭代计数器。问题只是第一次迭代要做什么？如果要进行更多迭代，则选择并不重要 - 随着迭代收敛，后续两次迭代的值实际上将相同。如果想要一种非迭代方案，那么显式部分的选择就变得至关重要。如果要保留整个方案的二阶精度，它应该是新时间水平上解的二阶近似。为此，可以使用任何二阶或更高阶的 Adams-Bashforth 方案，例如上述算法中用于完全对流的方案。实际上，在隐式方法中，所有变量值都应该使用二阶显式时间推进方案在新的时间级别进行初始化，特别是当网格是非正交的并且离散化涉及使用主要变量值计算的延迟校正时。否则，如果仅使用前一时间步长的值，则整体方案将是一阶的。

使用隐式迭代方法，甚至可以做出更多选择：可以首先单独迭代动量方程，以更新非线性和延迟校正，然后进行单个压力校正步骤，其中必须将压力校正方程求解到相对严格的公差，以确保充分执行连续性要求。另一种选择是扩展迭代循环来求解动量和压力校正方程，以便更新非线性和压力速度耦合。以下是此类迭代方案的示例算法：

1. 在新时间步内的第 m 次迭代中，使用二阶全隐式三时间级时间积分方案求解动量方程，以估计以下形式的新时间级解（参见第 6.3.2.4 节）：

\begin{sourceblock}
\includegraphics[page=223,trim={53.00bp 239.19bp 53.87bp 394.94bp},clip,width=0.92\textwidth,max height=0.38\textheight]{Ferziger4th.pdf}
\end{sourceblock}
其中 v ∗ 是 v m 的预测值；需要对其进行校正以强制连续性。 2. 要求修正后的速度和压力满足这种形式的动量方程：

\begin{sourceblock}
\includegraphics[page=223,trim={53.00bp 157.90bp 53.87bp 476.23bp},clip,width=0.92\textwidth,max height=0.38\textheight]{Ferziger4th.pdf}
\end{sourceblock}
将式(7.63)减去式(7.62)可得速度与压力修正关系式：

\begin{sourceblock}
\begin{tikzpicture}
\node[inner sep=0,anchor=south west] (image) at (0,0) {\includegraphics[page=223,trim={53.00bp 90.49bp 53.87bp 546.24bp},clip,width=0.92\textwidth,max height=0.38\textheight]{Ferziger4th.pdf}};
\begin{scope}[x={(image.south east)},y={(image.north west)}]
\fill[white] (0.12388,0.51853) rectangle (0.14367,0.91193);
\fill[white] (0.10626,0.04080) rectangle (0.12605,0.43421);
\fill[white] (0.14547,0.04454) rectangle (0.15862,0.43761);
\end{scope}
\end{tikzpicture}
\end{sourceblock}
3. 对速度 v m 施加连续性要求，

\begin{sourceblock}
\includegraphics[page=224,trim={53.00bp 559.25bp 53.87bp 78.87bp},clip,width=0.92\textwidth,max height=0.38\textheight]{Ferziger4th.pdf}
\end{sourceblock}
并求解所得的压力校正方程。 4. 增加迭代计数器并重复步骤 1-3，直到残差变得足够小。然后设置 v n + 1 = v m，p n + 1 = p m 并进入下一个时间级别。

请注意，上面的压力校正方程看起来与之前版本的方程（7.57）完全相同，除了 3/2 乘数之外；然而，右边是不同的，因为 v ∗ 来自动量方程的不同形式。在这种情况下，不必以非常严格的公差求解线性动量方程或压力校正方程，因为该求解将在下一次迭代中继续；通常，如果每个时间步执行三次或多次迭代，则每次迭代中将残差减少一个数量级就足够了。

上述算法的非迭代版本很容易导出；它使用类似于方程（7.52）的 Adams-Bashforth 方案对 t n + 1 处的对流通量进行显式估计，只是我们现在需要 t n + 1 处的对流通量而不是 t n + 1 / 2 处的估计：

\begin{sourceblock}
\includegraphics[page=224,trim={53.00bp 352.26bp 53.87bp 297.83bp},clip,width=0.92\textwidth,max height=0.38\textheight]{Ferziger4th.pdf}
\end{sourceblock}
外部迭代循环被删除，动量方程和压力校正方程每个时间步仅求解一次（但现在具有更严格的容差）。

上述迭代隐式分数步法（我们经常使用缩写词 IFSM 来引用）与下一节中描述的 SIMPLE 算法非常相似。细微的差别将在下一节的末尾讨论。包含迭代和非迭代算法的计算机代码是可用的（参见附录），其应用的一些结果将在下一章末尾的示例部分中显示。

\subsection*{7.2.2 简单、更简单、简单和 PISO}\phantomsection\addcontentsline{toc}{subsection}{7.2.2 简单、更简单、简单和 PISO}
正如第 1 章中所指出的。如图6所示，许多稳态问题的方法可以被视为解决非稳态问题直到达到稳态。主要区别在于，在解决不稳态问题时，选择时间步长以获得准确的历史记录，而在寻求稳态解时，则使用较大的时间步长来尝试快速达到稳态。对于稳定和缓慢的瞬态流，隐式方法是首选，因为它们的时间步长限制不如显式方案严格（实际上，它们可能没有任何限制！）。它们还经常用于瞬态问题的时间精确求解（特别是在使用商业 CFD 代码时，其通常不提供显式版本），特别是当网格被局部细化时，由于稳定​​性原因，显式方法需要的时间步长比足够解决精度所需的时间步长要小。

许多为稳定不可压缩流开发的求解方法都是隐式的；一些最流行的方法可以被视为前一节方法的变体。他们使用压力（或压力校正）方程在每个时间步上强制质量守恒，或者用稳定求解器首选的语言在每个外部迭代中强制质量守恒。我们现在看看其中一些方法。在本节中，我们对SIMPLE和相关方案进行一般描述；见第 8.1 和 8.2 节 我们详细介绍了交错和共置电网的 SIMPLE 和 IFSM 方案的实现；最后见第 8.4.1 节 我们检查两种网格类型上的 SIMPLE 和 IFSM 模拟结果。

我们使用与前一节类似的符号，以证明分数步方法和 SIMPLE 类型方法之间的相似性和差异。

起点是上一节中给出的方程（7.48）和（7.49）。 SIMPLE 型方法（包括 SIMPLEC 和 PISO，大多数商业 CFD 代码中都有）和分数步方法之间的主要区别在于，前者通常是完全隐式的（但也可以使用 Crank-Nicolson 方案）。这意味着所有通量和源项都是在新的时间级别计算的；先前时间级别的值仅出现在离散时间导数中。大多数代码提供隐式欧拉方案（一阶，不适合时间精确模拟）和二次后向方案（也称为三级方案；时间二阶，适合时间精确模拟）之间的选项。如果我们使用上一节中相同的运算符符号编写离散动量方程和连续性方程，则离散动量方程的两个版本如下：

\begin{sourceblock}
\includegraphics[page=225,trim={53.00bp 237.46bp 53.87bp 361.29bp},clip,width=0.92\textwidth,max height=0.38\textheight]{Ferziger4th.pdf}
\end{sourceblock}
由于第一个方程使用隐式欧拉格式，因此它是时间上的一阶格式（时间导数的近似值是相对于计算所有其他项的时间级别的一阶后向格式）。第二个方程使用时间导数近似，该近似在计算所有其他项的时间水平上是二阶精确的；它是通过在时间上对涉及新时间级别的二次插值进行微分而获得的，如第 1 节中所述。 6.3.2.4.由于对流、扩散和源项始终在 t n + 1 处计算，因此很容易在两种方案之间切换甚至混合它们。

这里我们假设流动是不可压缩的并且密度是恒定的；我们在第 1 章中展示。 11 该方法如何扩展到可压缩流。在这种假设下，连续性方程中没有时间导数，因此对于两种时间推进方案，它与上一节中的相同，见方程（7.51）。

由于所有项均使用完全隐式方案进行离散化，因此我们必须对非线性项进行线性化并迭代求解 v n + 1 和 p n + 1 的方程。如果我们计算非定常流并且需要时间精度，则必须在每个时间步内继续迭代，直到整个非线性方程组满足在窄公差范围内。对于稳定的流量，容差可以更大；然后，我们可以采取无限时间步长并迭代，直到满足稳定的非线性方程，或者及时前进，而不需要在每个时间步完全满足非线性方程（在这种情况下，通常每个时间步只执行一次迭代）。

在一个时间步内更新非线性项和耦合项的迭代称为外迭代，以区别于在具有固定系数的线性系统上执行的内迭代。

现在，我们删除上标（ n + 1 ）并使用外部迭代计数器 m 来表示新解的当前估计；当这些迭代收敛时，我们得到 v n + 1 ≈ v m。我们假设使用皮卡德迭代执行线性化，如上一节中的方程（7.61）所示。对于稳态问题和迭代时间推进方案（SIMPLE、SIMPLEC 或 PISO），这几乎是专门使用的。如果要有效地获得时间精确的解，只有新时间级别上的第一次迭代至关重要。如果简单地采用前一时间步长的值来表示 ( m = 0) 值，则时间步长内所需的迭代次数将高于使用更好的估计（例如，使用二阶显式时间推进方案）的次数。这个问题对于 PISO 算法尤其重要，该算法每个时间步仅求解一次动量方程；因此，无法通过迭代来改善初始误差。这个问题将在本节末尾再次讨论。源项和可变流体属性以类似的方式处理，即这些项的部分内容在前一次迭代中进行评估。

线性动量方程依次求解。 3 当所有隐式离散项组合在一起时，我们为每个速度分量获得以下形式的矩阵方程：

\begin{sourceblock}
\includegraphics[page=226,trim={53.00bp 234.05bp 53.87bp 413.59bp},clip,width=0.92\textwidth,max height=0.38\textheight]{Ferziger4th.pdf}
\end{sourceblock}
其中 G i 是梯度算子的 i 分量的简写符号。源项 Q 包含可以用 u m − 1 显式计算的所有项
i 以及任何体积力或其他可能取决于 u n + 1 的线性化或延迟校正项
i 或新时间级别的其他变量（例如温度）。它还包含涉及先前时间步解的不稳定项的部分，请参见方程（7.67）和（7.68）。请注意，矩阵 A 对于所有速度分量可能并不相同，但我们暂时忽略这一点。这里我们假设 FD 方法用于离散方程，但相同的过程也适用于 FV 方法 - 上述方程只需乘以细胞体积。

\footnotetext[3]{耦合（或整体）求解器也可用，并且大多数商业软件包现在都提供它们；关于求解器替代方案的讨论可以在代码文档或文献中找到，例如 Heil 等人。 （2008）或马林（2012）。另见第 1 章。参见图11，其中简要描述了一种此类方法。}
为了清楚起见，我们将删除矩阵 A 和源项 Q 的上标 ( m − 1 )，现在我们已经认识到它们确实取决于解，但使用先前外部迭代的值进行计算。

上述方程的一行内容为：

\begin{sourceblock}
\begin{tikzpicture}
\node[inner sep=0,anchor=south west] (image) at (0,0) {\includegraphics[page=227,trim={53.00bp 518.45bp 53.87bp 108.23bp},clip,width=0.92\textwidth,max height=0.38\textheight]{Ferziger4th.pdf}};
\begin{scope}[x={(image.south east)},y={(image.north west)}]
\fill[white] (0.25802,0.29270) rectangle (0.32695,0.65408);
\fill[white] (0.30686,0.26280) rectangle (0.33696,0.48657);
\fill[white] (0.34063,0.32184) rectangle (0.36881,0.61505);
\end{scope}
\end{tikzpicture}
\end{sourceblock}
压力项以符号差分形式书写，以强调求解方法与空间导数的离散化近似的独立性。空间导数的离散化可以是第 1 章中描述的任何阶数或任何类型。 3和4。请注意，系数 A k 包含离散对流和扩散项的贡献，而对角线系数
此外，P 还包含不稳定项的贡献（参见方程（7.67）和（7.68））。

现在我们回到更简单的矩阵表示法（7.69），但将矩阵 A 分为对角部分 A D 和非对角部分 A OD，原因很快就会变得显而易见。我们还用一个或多个星号替换表示当前外部迭代的值的上标 m，具体取决于我们所处的近似级别。

在外部迭代 m 中，我们首先使用上一次迭代的压力求解该方程：

\begin{sourceblock}
\begin{tikzpicture}
\node[inner sep=0,anchor=south west] (image) at (0,0) {\includegraphics[page=227,trim={53.00bp 327.22bp 53.87bp 320.55bp},clip,width=0.92\textwidth,max height=0.38\textheight]{Ferziger4th.pdf}};
\begin{scope}[x={(image.south east)},y={(image.north west)}]
\fill[white] (0.00000,0.84268) rectangle (0.10653,1.00000);
\fill[white] (0.10926,0.84268) rectangle (0.22226,1.00000);
\end{scope}
\end{tikzpicture}
\end{sourceblock}
通过求解该方程 v m * 获得的速度场通常不满足连续性方程。这就是我们想要通过纠正压力和速度来强制执行的，

\begin{sourceblock}
\begin{tikzpicture}
\node[inner sep=0,anchor=south west] (image) at (0,0) {\includegraphics[page=227,trim={53.00bp 267.45bp 53.87bp 377.93bp},clip,width=0.92\textwidth,max height=0.38\textheight]{Ferziger4th.pdf}};
\begin{scope}[x={(image.south east)},y={(image.north west)}]
\fill[white] (0.00000,0.74518) rectangle (0.04740,1.00000);
\fill[white] (0.05011,0.74518) rectangle (0.17675,1.00000);
\end{scope}
\end{tikzpicture}
\end{sourceblock}
要求修正后的速度和压力满足以下简化版式(7.69)，得到速度和压力修正的关系：

\begin{sourceblock}
\includegraphics[page=227,trim={53.00bp 207.67bp 53.87bp 440.10bp},clip,width=0.92\textwidth,max height=0.38\textheight]{Ferziger4th.pdf}
\end{sourceblock}
现在，通过从方程（7.73）减去方程（7.71），我们得到速度和压力修正之间的以下关系：

\begin{sourceblock}
\includegraphics[page=227,trim={53.00bp 147.90bp 53.87bp 499.88bp},clip,width=0.92\textwidth,max height=0.38\textheight]{Ferziger4th.pdf}
\end{sourceblock}
这个关系很简单，因为对角矩阵很容易求逆；如果你**
i 也应用于方程（7.73）中的非对角矩阵，对于推导压力校正方程而言，该关系过于复杂。这种简化的合理性在于，当外部迭代收敛时，所有校正都趋于零，因此最终的解不受影响。然而，这种简化确实会影响收敛速度，稍后我们将描述如何通过正确选择欠松弛因子来改进收敛速度。

我们现在要求修正后的速度 u **
i 满足离散连续性方程，因此：

\begin{sourceblock}
\begin{tikzpicture}
\node[inner sep=0,anchor=south west] (image) at (0,0) {\includegraphics[page=228,trim={53.00bp 547.19bp 53.87bp 102.89bp},clip,width=0.92\textwidth,max height=0.38\textheight]{Ferziger4th.pdf}};
\begin{scope}[x={(image.south east)},y={(image.north west)}]
\fill[white] (0.00000,0.81943) rectangle (0.11477,1.00000);
\fill[white] (0.11753,0.81943) rectangle (0.18060,1.00000);
\end{scope}
\end{tikzpicture}
\end{sourceblock}
i 通过 p ′ 借助表达式（7.74），我们得到压力修正方程：

通过表达 u ′

\begin{sourceblock}
\begin{tikzpicture}
\node[inner sep=0,anchor=south west] (image) at (0,0) {\includegraphics[page=228,trim={53.00bp 498.59bp 53.87bp 150.71bp},clip,width=0.92\textwidth,max height=0.38\textheight]{Ferziger4th.pdf}};
\begin{scope}[x={(image.south east)},y={(image.north west)}]
\fill[white] (0.00000,0.82779) rectangle (0.12556,1.00000);
\fill[white] (0.12827,0.82779) rectangle (0.24454,1.00000);
\end{scope}
\end{tikzpicture}
\end{sourceblock}
请注意，该方程与上一节中的方程（7.57）具有相同的形式；我们将在下面进一步讨论相似之处和不同之处。

这种方法本质上是上一节中介绍的方法的变体。这种方法首先构造一个不满足连续性方程的速度场，然后通过减去一些东西（通常是压力梯度）来校正它，称为投影方法。从矢量数学的角度来看，压力通过连续性约束作为算子，将发散的速度矢量场投影到非发散的矢量场（参见 Kim 和 Moin 1985）。

一旦压力校正得到解决，速度和压力就可以使用方程（7.74）和（7.72）进行更新。这些被用来表示外部迭代 m 处的解，并且新的迭代可以开始。这被称为 SIMPLE 算法（Caretto 等人，1972）； SIMPLE 是“压力关联方程的半隐式方法”的缩写。我们将在下面讨论它的属性。

我们可以近似估计其效果，而不是忽略非对角线项中的速度校正。这可以通过近似速度修正 u ′ 来实现
i 在任何节点处通过邻居值的加权平均值计算，例如，

\begin{sourceblock}
\begin{tikzpicture}
\node[inner sep=0,anchor=south west] (image) at (0,0) {\includegraphics[page=228,trim={53.00bp 241.74bp 53.87bp 389.72bp},clip,width=0.92\textwidth,max height=0.38\textheight]{Ferziger4th.pdf}};
\begin{scope}[x={(image.south east)},y={(image.north west)}]
\fill[white] (0.18818,0.35900) rectangle (0.21522,0.74971);
\fill[white] (0.20460,0.29931) rectangle (0.23474,0.55392);
\fill[white] (0.24006,0.36678) rectangle (0.26824,0.70012);
\end{scope}
\end{tikzpicture}
\end{sourceblock}
这使我们能够近似 A OD u ′
当从方程（7.69）而不是（7.73）减去方程（7.71）时，i出现在速度和压力修正之间的关系中：

\begin{sourceblock}
\begin{tikzpicture}
\node[inner sep=0,anchor=south west] (image) at (0,0) {\includegraphics[page=228,trim={53.00bp 181.23bp 53.87bp 466.54bp},clip,width=0.92\textwidth,max height=0.38\textheight]{Ferziger4th.pdf}};
\begin{scope}[x={(image.south east)},y={(image.north west)}]
\fill[white] (0.00000,0.84268) rectangle (0.08484,1.00000);
\end{scope}
\end{tikzpicture}
\end{sourceblock}
通过使用表达式（7.77），我们可以写出上述方程的简化版本，如下所示（该简化比忽略左侧第二项简单得多，如 SIMPLE 中所做的那样）：

\begin{sourceblock}
\includegraphics[page=228,trim={53.00bp 109.50bp 53.87bp 537.90bp},clip,width=0.92\textwidth,max height=0.38\textheight]{Ferziger4th.pdf}
\end{sourceblock}
其中̃ A D 表示非对角矩阵元素之和，见式(7.77)。

值得注意的是，当使用FV方法时，对流和扩散通量对对角矩阵元素的贡献等于对非对角元素的贡献的负和；见第 8.1 节 了解更多详情。在没有其他贡献的情况下，这将使上式中的 A D + ̃ A D = 0。幸运的是，当解决瞬态问题时，非稳态项的贡献出现在 A D 中，或者当解决稳态问题时，它除以欠松弛因子 α u < 1；见节。有关这种欠松弛的详细信息，请参见 5.4.3。因此，A D 为正，并且其大小始终大于 ̃ A D 中非对角线元素之和（通常为负）。然而，总和
式（7.79）中出现的 A D + ̃ A D 远小于 SIMPLE 相应表达式中出现的 A D，见式（7.74）。因为压力校正的梯度乘以此项的倒数值，所以对于两种方法中相同的速度校正，来自该校正方法（称为 SIMPLEC；见下文）的压力校正必须比来自 SIMPLE 的压力校正小得多。这就是为什么 SIMPLE 需要压力校正的欠松弛（由于其推导过程的简化而导致过度预测），如下所述。

下一步是要求修正后的速度满足连续性方程，这导致压力修正方程的形式与方程（7.76）相同，只是将 A D 替换为总和 A D + ̃ A D。这称为 SIMPLEC 算法（SIMPLE-Corrected）（Van Doormal 和 Raithby 1984）。

这种通用类型的又一种方法是通过将简单步骤视为预测器以及随后的一系列校正器步骤来导出的。我们只需再添加一个星号来继续修正速度和压力的过程即可；在 SIMPLE 之后的第一次修正中，我们希望速度和压力满足动量方程的这种形式：

\begin{sourceblock}
\includegraphics[page=229,trim={53.00bp 270.33bp 53.87bp 377.86bp},clip,width=0.92\textwidth,max height=0.38\textheight]{Ferziger4th.pdf}
\end{sourceblock}
现在，通过从方程（7.80）减去方程（7.73），我们得到第二速度和压力修正之间的以下关系：

\begin{sourceblock}
\includegraphics[page=229,trim={53.00bp 210.14bp 53.87bp 437.63bp},clip,width=0.92\textwidth,max height=0.38\textheight]{Ferziger4th.pdf}
\end{sourceblock}
Here u ′
i 由上一步可知。我们现在要求 u ******
i 也满足连续性方程：

\begin{sourceblock}
\begin{tikzpicture}
\node[inner sep=0,anchor=south west] (image) at (0,0) {\includegraphics[page=229,trim={53.00bp 164.63bp 53.87bp 485.46bp},clip,width=0.92\textwidth,max height=0.38\textheight]{Ferziger4th.pdf}};
\begin{scope}[x={(image.south east)},y={(image.north west)}]
\fill[white] (0.00000,0.81994) rectangle (0.12562,1.00000);
\fill[white] (0.12833,0.81994) rectangle (0.24460,1.00000);
\end{scope}
\end{tikzpicture}
\end{sourceblock}
i 借助表达式 ( 7.81 ) 通过 p ′′ 并认识到 u **
通过表达 u ′′
i 已经满足连续性方程，我们得到第二个压力修正方程：

\begin{sourceblock}
\begin{tikzpicture}
\node[inner sep=0,anchor=south west] (image) at (0,0) {\includegraphics[page=229,trim={53.00bp 102.96bp 53.87bp 545.23bp},clip,width=0.92\textwidth,max height=0.38\textheight]{Ferziger4th.pdf}};
\begin{scope}[x={(image.south east)},y={(image.north west)}]
\fill[white] (0.00000,0.83844) rectangle (0.11558,1.00000);
\end{scope}
\end{tikzpicture}
\end{sourceblock}
这个过程可以通过在方程（7.80）中的每一项上再添加一个星号来继续。

请注意，每个步骤的压力校正方程都具有相同的系数矩阵，可以在某些求解器中利用该系数矩阵（可以存储并重复使用矩阵的因式分解）。这个过程被称为 PISO 算法（带有算子分裂的隐式压力）（Issa 1986）。除了 SIMPLE 或 SIMPLEC 之外，一些商业和开放式 CFD 软件包还提供此算法。通常，执行 3-5 个校正步骤。

最后，Patankar（1980）提出了另一种类似的方法，称为SIMPLER（SIMPLE-Revised）。其中，首先求解压力校正方程（7.76），并像 SIMPLE 中一样校正速度。新的压力场是根据对式(7.69)求散度得到的压力方程计算得到的。该算法尚未得到广泛应用，因为与迄今为止讨论的其他替代方案相比，它没有任何优势。

如前所述，由于忽略了非对角项中速度修正的影响，SIMPLE 算法收敛速度不快。实际上，如果使用计算的压力校正和上面给出的方程来校正压力和速度，则它可能根本不会收敛，除非时间步长非常小。特别是对于使用无限时间步长计算的稳态流，其性能很大程度上取决于动量方程中使用的欠松弛参数的值。通过反复试验首先发现，与式(7.72)中的说法相反，在求解压力校正方程后，如果仅将p′的一部分添加到p m − 1 中，即更新压力，则可以提高收敛性：

\begin{sourceblock}
\begin{tikzpicture}
\node[inner sep=0,anchor=south west] (image) at (0,0) {\includegraphics[page=230,trim={53.00bp 330.61bp 53.87bp 315.70bp},clip,width=0.92\textwidth,max height=0.38\textheight]{Ferziger4th.pdf}};
\begin{scope}[x={(image.south east)},y={(image.north west)}]
\fill[white] (0.00000,0.73323) rectangle (0.02911,1.00000);
\fill[white] (0.03179,0.73323) rectangle (0.13405,1.00000);
\end{scope}
\end{tikzpicture}
\end{sourceblock}
其中 0 ≤ α p ≤ 1。SIMPLEC、SIMPLER 和 PISO 不需要对笛卡尔网格上的压力校正进行欠松弛。我们将在下一章中展示，由于应用于（或忽略）与网格非正交性相关的项的延迟校正方法，非正交网格上压力校正方程的离散化和求解可能需要欠松弛。

通过要求 SIMPLE 和 SIMPLEC 产生相同的速度修正，我们可以得出速度和压力的欠松弛因子之间的最佳关系，因为在后者中，邻近速度修正的影响是近似的而不是被忽略（见方程（7.79）和（7.74））：

\begin{sourceblock}
\begin{tikzpicture}
\node[inner sep=0,anchor=south west] (image) at (0,0) {\includegraphics[page=230,trim={53.00bp 173.31bp 53.87bp 454.69bp},clip,width=0.92\textwidth,max height=0.38\textheight]{Ferziger4th.pdf}};
\begin{scope}[x={(image.south east)},y={(image.north west)}]
\fill[white] (0.20427,0.08285) rectangle (0.23964,0.40928);
\end{scope}
\end{tikzpicture}
\end{sourceblock}
该表达式可以重写如下：

\begin{sourceblock}
\includegraphics[page=230,trim={53.00bp 112.97bp 53.87bp 515.22bp},clip,width=0.92\textwidth,max height=0.38\textheight]{Ferziger4th.pdf}
\end{sourceblock}
最佳欠松弛对于稳态问题至关重要；在这种情况下（没有不稳定项的贡献）并且在没有源项的贡献的情况下，我们可以通过查看等式来编写。 7.86（参见第 5.4.3 节）：

\begin{sourceblock}
\begin{tikzpicture}
\node[inner sep=0,anchor=south west] (image) at (0,0) {\includegraphics[page=231,trim={53.00bp 582.68bp 53.87bp 50.00bp},clip,width=0.92\textwidth,max height=0.38\textheight]{Ferziger4th.pdf}};
\begin{scope}[x={(image.south east)},y={(image.north west)}]
\fill[white] (0.51236,0.03586) rectangle (0.54478,0.41393);
\end{scope}
\end{tikzpicture}
\end{sourceblock}
将这个关系代入式（7.86），可得：

\begin{sourceblock}
\includegraphics[page=231,trim={53.00bp 521.66bp 53.87bp 107.06bp},clip,width=0.92\textwidth,max height=0.38\textheight]{Ferziger4th.pdf}
\end{sourceblock}
这表明由 SIMPLE 产生的压力修正应乘以

\begin{sourceblock}
\begin{tikzpicture}
\node[inner sep=0,anchor=south west] (image) at (0,0) {\includegraphics[page=231,trim={53.00bp 476.37bp 53.87bp 174.31bp},clip,width=0.92\textwidth,max height=0.38\textheight]{Ferziger4th.pdf}};
\begin{scope}[x={(image.south east)},y={(image.north west)}]
\fill[white] (0.00000,0.92367) rectangle (0.03411,1.00000);
\end{scope}
\end{tikzpicture}
\end{sourceblock}
为了获得与 SIMPLEC 相同的速度修正。 4
我们将见第 8.4 节 示例中速度和压力的欠松弛因素如何影响 SIMPLE 的效率。

此类方法的求解算法可概括如下（见图7.6）：

1. 使用最新解 u n i 和 p n 作为 u n + 1 的起始估计，开始计算新时间 t n + 1 处的场
我和 p n + 1.2. 使用迭代计数器 m 启动外迭代循环。 3. 组装并求解速度分量的线性代数方程组（动量方程）以获得 u ∗
我。 4. 建立并求解压力校正方程以获得 p ′。 5. 修正速度和压力，得到满足连续性方程的速度场 u ＊＊ i 和新的压力 p ＊。对于 SIMPLE 和 SIMPLEC，这些是迭代 m 的最终值；对于 SIMPLER，只有速度是最终的。对于 PISO 算法，求解第二个压力校正方程并再次校正速度和压力。重复此操作，直到修正变得足够小并继续下一个时间步骤。对于更简单的情况，求解 u m 之后的 p m 的压力方程
i 是上面得到的。 6. 如果需要求解其他传输方程（例如，温度、物种、湍流量等），请在此处完成；在需要时使用校正的质量通量、速度和压力。 7. 如果流体属性是可变的，现在可以使用单元中心更新的变量值重新计算它们。 8. 返回步骤 2 并重复，使用 u m i 和 p m 作为 u n + 1 的改进估计
i 和 p n + 1，直到所有修正都小到可以忽略不计。 9. 前进到下一个时间步骤。

这种方法对于解决稳态问题相当有效。它们的融合可以通过多重网格策略进一步提高，这将被证明

\footnotetext[4]{这种关系首先由 Raithby 和 Schneider (1979) 导出，后来由 Peri ́c (1985) 使用不同的论证重新发现。}
\begin{center}\small 图7.6 类SIMPLE算法流程图\end{center}
提前时间

Moment u m
时间步进循环

线性压力方程

外迭代

s olver
S c a l a r s
内迭代

Propertie s
No

Converged?

Ye s
No

时间限制？

Ye s
章中阐述。 12.上述方法还有很多衍生方法，名称各异，但都源于上述思想，这里不再一一列举。下面我们将证明人工压缩方法也可以用类似的方式解释。

类似 SIMPLE 的方法的一个很好的功能是，它们可以轻松扩展以求解其他传输方程，如步骤 6 中所述。此外，可以轻松处理可变流体属性：在一次外部迭代期间简单地假设它们为常数，并在循环结束时更新所有变量时重新计算。如图 7.6 所示；外部迭代循环可以轻松扩展以更新任何非线性或延迟校正，或求解其他方程。对于上一节中描述的隐式迭代分数步方法也是如此，其基本上遵循相同的流程图。

通过将上述一类方法与上一节中介绍的分数步方法进行比较，我们可以看到它们与后者的隐式版本非常相似。唯一的区别是，类 SIMPLE 方法在推导压力校正方程时，不仅更新非稳态项中的新速度（如分数阶方法所做的那样），而且还更新离散对流和扩散通量的对角线部分。当使用小时间步长计算非稳态流动时，这种差异并不那么重要，但如果使用大时间步长来迈向稳态解，这种差异就会变得很重要。我们将在下一节中更详细地研究这个问题。

7.2.2.1 简单法与迭代隐式分步法
仔细观察方程（7.57）和（7.65）可以发现，随着时间步长趋于无穷大，用于压力校正的泊松方程中的源项趋于零。这清楚地表明，给定形式的分数步法不能用于解决稳态问题，即当不稳定项缺失或时间步长太大时。

如果我们检查类似 SIMPLE 方法的压力校正方程，我们会发现它们不会遇到同样的问题，参见方程（7.76）。如果我们比较两种方法的单个网格点的压力校正方程，可以更好地分析这一点。对于SIMPLE，我们从方程（7.76）得到：

\begin{sourceblock}
\begin{tikzpicture}
\node[inner sep=0,anchor=south west] (image) at (0,0) {\includegraphics[page=233,trim={53.00bp 497.82bp 53.87bp 132.14bp},clip,width=0.92\textwidth,max height=0.38\textheight]{Ferziger4th.pdf}};
\begin{scope}[x={(image.south east)},y={(image.north west)}]
\fill[white] (0.35041,0.45357) rectangle (0.36809,0.77336);
\fill[white] (0.33829,0.03510) rectangle (0.37783,0.38474);
\end{scope}
\end{tikzpicture}
\end{sourceblock}
如果我们假设 ρ/ A P 在网格点附近不会发生空间变化（这通常是不正确的，但对于目前的目的来说也不是一个不合理的假设），我们可以将上面的方程重写如下：

\begin{sourceblock}
\begin{tikzpicture}
\node[inner sep=0,anchor=south west] (image) at (0,0) {\includegraphics[page=233,trim={53.00bp 415.33bp 53.87bp 214.70bp},clip,width=0.92\textwidth,max height=0.38\textheight]{Ferziger4th.pdf}};
\begin{scope}[x={(image.south east)},y={(image.north west)}]
\fill[white] (0.35041,0.45251) rectangle (0.36809,0.77264);
\fill[white] (0.33829,0.03323) rectangle (0.37783,0.38327);
\end{scope}
\end{tikzpicture}
\end{sourceblock}
对于分数步法，方程（7.65）给出：

\begin{sourceblock}
\begin{tikzpicture}
\node[inner sep=0,anchor=south west] (image) at (0,0) {\includegraphics[page=233,trim={53.00bp 356.55bp 53.87bp 273.49bp},clip,width=0.92\textwidth,max height=0.38\textheight]{Ferziger4th.pdf}};
\begin{scope}[x={(image.south east)},y={(image.north west)}]
\fill[white] (0.34295,0.45263) rectangle (0.36066,0.77285);
\fill[white] (0.33086,0.03324) rectangle (0.37041,0.38338);
\end{scope}
\end{tikzpicture}
\end{sourceblock}
现在我们看到，如果 3 /( 2 t ) = A P /ρ，这两个方程本质上是相同的。

现在让我们看看 A P 背后隐藏着什么。从方程（7.68）我们看到它肯定包含 3 ρ/( 2 t )，加上离散对流和粘性项的贡献，其量级分别为 ρ u i / x 和 μ/( x ) 2 （取决于使用哪种近似来离散化这些项）。人们可以证明，对于不可压缩流和所有保守方案，这一贡献等于 -
k A k，其中 A k 是离散和线性动量方程矩阵的非对角元素（参见方程（8.20））。因此，对于三时间级别方案：

\begin{sourceblock}
\includegraphics[page=233,trim={53.00bp 199.61bp 53.87bp 431.60bp},clip,width=0.92\textwidth,max height=0.38\textheight]{Ferziger4th.pdf}
\end{sourceblock}
其中求和涉及节点 P 的计算分子中包含的所有相邻网格点。显然，如果在 SIMPLE 中省略对流和扩散的贡献，则它变成上一节中描述的隐式分数阶方法，因为速度校正也将是相同的（参见方程（7.56）和（7.81））：

\begin{sourceblock}
\begin{tikzpicture}
\node[inner sep=0,anchor=south west] (image) at (0,0) {\includegraphics[page=233,trim={53.00bp 68.67bp 53.87bp 530.84bp},clip,width=0.92\textwidth,max height=0.38\textheight]{Ferziger4th.pdf}};
\begin{scope}[x={(image.south east)},y={(image.north west)}]
\fill[white] (0.00000,0.94222) rectangle (0.13284,1.00000);
\fill[white] (0.13555,0.94222) rectangle (0.18361,1.00000);
\fill[white] (0.18629,0.94222) rectangle (0.28180,1.00000);
\end{scope}
\end{tikzpicture}
\end{sourceblock}
当时间步长非常小时，与 ρ/ t 相比，邻近系数对 A P 的贡献变小，因此可以预期瞬态 SIMPLE 和迭代隐式分数步法的行为类似。然而，对于较大的时间步长，由于欠松弛以及对流和扩散通量对 SIMPLE 中压力校正方程矩阵对角线元素的贡献而产生差异，而分数步法中不存在这些差异。我们在节 中比较了两种方法的性能。 8.4 适用于稳态流和非稳态流。

7.2.2.2 SIMPLE 中的欠松弛和时间推进
当使用类似 SIMPLE 的方法来计算稳态流时，它们通常会丢弃不稳定项，从而假设无限大的时间步长。然而，如第 1 节所述，如果不调用欠松弛，该方法将无法工作。 5.4.3.解决稳态问题时使用欠松弛所产生的代数方程与将隐式欧拉格式应用于非定常方程所产生的代数方程之间存在很强的相似性。欠松弛和隐式时间离散都会产生额外的源项以及对中心系数 A P 的贡献。

如果我们要求方程（7.93）中给出的时间推进方案的对角线系数等于欠松弛稳态计算中使用的对角线系数，我们得到以下关系（注意非对角线系数通常都是负数）：

\begin{sourceblock}
\begin{tikzpicture}
\node[inner sep=0,anchor=south west] (image) at (0,0) {\includegraphics[page=234,trim={53.00bp 291.24bp 53.87bp 342.44bp},clip,width=0.92\textwidth,max height=0.38\textheight]{Ferziger4th.pdf}};
\begin{scope}[x={(image.south east)},y={(image.north west)}]
\fill[white] (0.00000,0.92452) rectangle (0.12208,1.00000);
\fill[white] (0.34442,0.60659) rectangle (0.36556,0.96303);
\fill[white] (0.38129,0.39156) rectangle (0.40947,0.74769);
\fill[white] (0.36000,0.16543) rectangle (0.37314,0.52157);
\end{scope}
\end{tikzpicture}
\end{sourceblock}
其中 k 在计算分子的相邻细胞中心上运行。从这个方程中，我们可以导出等效时间步长的表达式，该表达式表示为欠松弛因子 α u 的函数，反之亦然，这将导致两种计算方法具有相同的行为：

\begin{sourceblock}
\includegraphics[page=234,trim={53.00bp 194.31bp 53.87bp 433.17bp},clip,width=0.92\textwidth,max height=0.38\textheight]{Ferziger4th.pdf}
\end{sourceblock}
这些表达式是在假设使用 FD 方法的情况下推导出来的；对于 FV 方法，只需将 ρ 替换为 ρ V，其中 V 是细胞体积。

在新时间步的迭代中，通常的初始猜测是前一步的解。如果最终稳态是唯一感兴趣的结果，并且从初始猜测到最终阶段的发展细节并不重要，则每个时间步仅执行一次迭代就足够了。这样就不必存储旧的解——只需要组合矩阵和源项。

使用时间推进和欠松弛之间的主要区别在于，在时间推进中对所有网格点使用相同的时间步长相当于使用可变的欠松弛因子；相反，使用恒定的欠松弛因子相当于对每个网格点应用不同的时间步长。

需要注意的是，如果每个时间步仅执行一次迭代，则该方案不会保留隐式欧拉方法的所有稳定性。因此，在这种时间推进方案中可以采用的时间步长是有限的。另一方面，当在稳态计算中将欠松弛与外部迭代一起使用时，参数 α u 的选择也受到限制，并且肯定必须小于 1。典型值范围在 0.7 到 0.9 之间，具体取决于问题和网格质量（网格质量差和僵硬问题的值较低）。

我们在这里比较了在时间推进和松弛模式下运行的 SIMPLE 算法。同样的分析适用于第 4 节中提出的迭代隐式分数步方法的比较。 7.2.1 和 SIMPLE 处于欠松弛模式。通过适当选择一个时间步长和另一个欠松弛参数，可以使它们或多或少以完全相同的方式执行。

请注意，类似 SIMPLE 的方法通常会保留欠松弛，即使它们用于计算不稳态流，即对角矩阵元素始终具有以下形式：

\begin{sourceblock}
\begin{tikzpicture}
\node[inner sep=0,anchor=south west] (image) at (0,0) {\includegraphics[page=235,trim={53.00bp 366.35bp 53.87bp 253.46bp},clip,width=0.92\textwidth,max height=0.38\textheight]{Ferziger4th.pdf}};
\begin{scope}[x={(image.south east)},y={(image.north west)}]
\fill[white] (0.00000,0.83272) rectangle (0.04406,1.00000);
\fill[white] (0.04677,0.83272) rectangle (0.08818,1.00000);
\fill[white] (0.09089,0.83272) rectangle (0.16226,1.00000);
\fill[white] (0.37131,0.17699) rectangle (0.40665,0.44572);
\fill[white] (0.41200,0.20203) rectangle (0.44018,0.45154);
\end{scope}
\end{tikzpicture}
\end{sourceblock}
如果没有欠松弛，当非线性和耦合效应很强时（例如，使用雷诺平均纳维-斯托克斯方程与湍流模型相结合计算湍流），该方法将受到可使用的时间步长的限制。因为外部迭代无论如何都是在一个时间步内执行的，所以在任何情况下允许少量的欠松弛都是更安全的；仅对于非常小的时间步长，例如在湍流的直接或大涡模拟中，可以将欠松弛因子设置为1。

\subsection*{7.2.3 人工压缩方法}\phantomsection\addcontentsline{toc}{subsection}{7.2.3 人工压缩方法}
可压缩流动是流体力学的一个非常重要的领域。它的应用，特别是在空气动力学和涡轮发动机设计中，引起了人们对可压缩流动方程数值求解方法开发的极大关注。已经开发了许多这样的方法。一个明显的问题是它们是否可以适应不可压缩流动的解。我们在这里展示了这些方法如何适用于不可压缩流，并描述了人工可压缩性方法的一些关键属性。有关此方法的深入讨论，请参阅 Kwak 和 Kiris（2011 年，第 4 章），以及 Louda 等人。 ( 2008 ) 最近使用雷诺平均纳维-斯托克斯方程在湍流中的应用。

可压缩流方程和不可压缩流方程的主要区别在于它们的数学特性。可压缩流方程是双曲线的，这意味着它们具有信号以有限传播速度传播的实数特征；这反映了可压缩流体支持声波的能力。相比之下，我们看到不可压缩方程具有混合抛物线-椭圆特征。如果要使用可压缩流动的方法来计算不可压缩流动，则需要修改方程的特征。

特征上的差异可以追溯到不可压缩连续性方程中缺少时间导数项。可压缩版本包含密度的时间导数。因此，赋予不可压缩方程双曲特征的最直接方法是在连续性方程上附加时间导数。由于密度是常数，因此不可能添加 ∂ρ/∂ t，即使用可压缩方程。速度分量的时间导数出现在动量方程中，因此它们不是逻辑选择。这使得压力的时间导数成为明确的选择。

将压力的时间导数添加到连续性方程中意味着我们不再求解真正不可压缩流的方程。因此，生成的时间历程不可能准确，并且人工压缩方法对非稳态不可压缩流的适用性是一个值得怀疑的事业，尽管已经进行了尝试。另一方面，在收敛时，时间导数为零并且解满足不可压缩方程。这种方法首先由 Chorin (1967) 提出，并且文献中已经提出了许多版本，主要区别在于所使用的基础可压缩流方法。如上所述，基本思想是将压力的时间导数添加到连续性方程中：

\begin{sourceblock}
\begin{tikzpicture}
\node[inner sep=0,anchor=south west] (image) at (0,0) {\includegraphics[page=236,trim={53.00bp 308.22bp 53.87bp 329.64bp},clip,width=0.92\textwidth,max height=0.38\textheight]{Ferziger4th.pdf}};
\begin{scope}[x={(image.south east)},y={(image.north west)}]
\fill[white] (0.00000,0.88331) rectangle (0.11558,1.00000);
\end{scope}
\end{tikzpicture}
\end{sourceblock}
其中 β 是人工压缩性参数，其值对该方法的性能至关重要。显然，β 的值越大，方程就越“不可压缩”；然而，大的 β 使得方程在数值上非常僵硬。我们只考虑恒定密度的情况，但该方法可以应用于变密度流。

为了求解这些方程，有许多方法可用。事实上，因为现在每个方程都包含一个时间导数，所以用于求解它们的方法可以按照第 1 章中介绍的用于求解常微分方程的方法进行建模。 6.由于人工可压缩性方法主要针对稳态流，因此应优先考虑隐式方法。还有一点很重要，即可压缩流面临的主要困难，即亚音速流向超音速流过渡的可能性，特别是可以避免激波的可能存在。二维或三维问题的求解方法的最佳选择是隐式方法，它不需要在每个时间步求解完整的二维或三维问题，这意味着交替方向隐式或近似因式分解方法是最佳选择。下面介绍了使用人工压缩性推导压力方程的方案示例。

最简单的方案使用时间上的一阶显式离散化；它可以逐点计算压力，但对时间步长有严格的限制。因为压力的时间发展并不重要，并且我们感兴趣的是尽快获得稳态解，所以完全隐式欧拉方案是更好的选择。为了将该方法与前面章节中描述的方法联系起来，我们注意到，通过使用旧压力求解动量方程获得的中间速度场 v * 不满足不可压缩连续性方程，因此需要对其进行修正。速度修正需要与压力修正联系起来。修正定义为：

\begin{sourceblock}
\includegraphics[page=237,trim={53.00bp 499.37bp 53.87bp 148.32bp},clip,width=0.92\textwidth,max height=0.38\textheight]{Ferziger4th.pdf}
\end{sourceblock}
从动量方程可以看出，速度修正必须与压力梯度成正比；分数步法和类SIMPLE方法都已经导出了合适的关系，让我们采用SIMPLE方法的定义，见方程（7.74）：

\begin{sourceblock}
\begin{tikzpicture}
\node[inner sep=0,anchor=south west] (image) at (0,0) {\includegraphics[page=237,trim={53.00bp 426.86bp 53.87bp 220.05bp},clip,width=0.92\textwidth,max height=0.38\textheight]{Ferziger4th.pdf}};
\begin{scope}[x={(image.south east)},y={(image.north west)}]
\fill[white] (0.00000,0.72491) rectangle (0.10147,1.00000);
\fill[white] (0.10418,0.72491) rectangle (0.14725,1.00000);
\fill[white] (0.14992,0.72491) rectangle (0.28012,1.00000);
\end{scope}
\end{tikzpicture}
\end{sourceblock}
其中 A D 是离散动量方程系数矩阵的对角部分，G 是离散梯度算子。利用上述两个方程中引入的定义，修正连续性方程（7.99）的离散化版本可以写为：

\begin{sourceblock}
\begin{tikzpicture}
\node[inner sep=0,anchor=south west] (image) at (0,0) {\includegraphics[page=237,trim={53.00bp 331.85bp 53.87bp 304.78bp},clip,width=0.92\textwidth,max height=0.38\textheight]{Ferziger4th.pdf}};
\begin{scope}[x={(image.south east)},y={(image.north west)}]
\fill[white] (0.26493,0.51677) rectangle (0.29134,0.95934);
\fill[white] (0.31456,0.27991) rectangle (0.37332,0.68079);
\fill[white] (0.24520,0.04981) rectangle (0.26692,0.44155);
\fill[white] (0.29329,0.04066) rectangle (0.30644,0.43240);
\end{scope}
\end{tikzpicture}
\end{sourceblock}
其中 D 表示离散散度算子。

除了左侧第一项之外，该方程与 SIMPLE 方法的压力校正方程 Eq.(7.76) 相同。它引入了对压力校正方程的对角矩阵元素的额外贡献，因此起到欠松弛的作用（有关这种欠松弛的更多详细信息，请参见第 5.4.3 节）。由于这个附加项，新时间水平下的修正速度 v n + 1 也将不满足不可压缩连续性方程；然而，当我们接近稳态时，所有校正都趋于零，因此将满足正确的方程。

因此，到目前为止提出的所有压力计算方法，尽管是通过略有不同的途径得出的，但都归结为相同的基本方法。同样，重要的是括号内的压力导数以与动量方程相同的方式近似，而外部导数是来自连续性方程的导数。

基于人工可压缩性的方法收敛的关键因素是参数β的选择。最佳值取决于问题，尽管一些作者建议使用自动程序来选择它。非常大的值将需要校正的速度场以满足不可压缩连续性方程。在该方法的上述版本中，这对应于没有压力校正欠松弛的SIMPLE方案；该过程将仅在小 t 时收敛。然而，如果像 SIMPLE 中那样仅将 p ' 的一部分添加到压力中，则可以使用无限大的 β。事实上，SIMPLE 可以被视为使用无限 β 的人工压缩方法的特殊版本。

允许的 β 最低值可以通过查看压力波的传播速度来确定。伪声速为：

\begin{sourceblock}
\begin{tikzpicture}
\node[inner sep=0,anchor=south west] (image) at (0,0) {\includegraphics[page=238,trim={53.00bp 534.03bp 53.87bp 115.00bp},clip,width=0.92\textwidth,max height=0.38\textheight]{Ferziger4th.pdf}};
\begin{scope}[x={(image.south east)},y={(image.north west)}]
\fill[white] (0.40941,0.07013) rectangle (0.42752,0.74635);
\fill[white] (0.43164,0.08591) rectangle (0.45982,0.76154);
\end{scope}
\end{tikzpicture}
\end{sourceblock}
通过要求压力波的传播速度比涡量传播的速度快得多，可以针对简单的通道流得出以下标准（参见 Kwak 等人，1986）：

\begin{sourceblock}
\begin{tikzpicture}
\node[inner sep=0,anchor=south west] (image) at (0,0) {\includegraphics[page=238,trim={53.00bp 456.00bp 53.87bp 171.24bp},clip,width=0.92\textwidth,max height=0.38\textheight]{Ferziger4th.pdf}};
\begin{scope}[x={(image.south east)},y={(image.north west)}]
\fill[white] (0.25955,0.23548) rectangle (0.28123,0.53265);
\fill[white] (0.28689,0.23548) rectangle (0.32117,0.53265);
\fill[white] (0.67332,0.22596) rectangle (0.72313,0.53265);
\fill[white] (0.72665,0.23548) rectangle (0.74138,0.53265);
\end{scope}
\end{tikzpicture}
\end{sourceblock}
其中 x L 是入口和出口之间的距离，x δ 是两壁之间距离的一半，x ref 是参考长度。基于人工压缩性的各种方法中使用的 β 典型值在 0.1 到 10 之间。

显然，如果修正后的速度场要严格满足连续性方程，则 1 /(β t ) 应该比式(7.102)第二项产生的系数小。如果要获得快速收敛，这是必要的。对于某些迭代求解方法（例如，在并行处理中使用域分解技术或在复杂几何形状中使用块结构网格），我们发现将 SIMPLE 中压力校正方程的 A P 系数除以小于 1 的因子（0.95 到 0.99）很有用。这相当于 1 /(β t ) ≈ (0.01 to 0.05) A P 的人工压缩方法。

\subsection*{7.2.4 流函数-涡度法}\phantomsection\addcontentsline{toc}{subsection}{7.2.4 流函数-涡度法}
对于具有恒定流体性质的不可压缩二维流，可以通过引入流函数 ψ 和涡度 ω 作为因变量来简化纳维-斯托克斯方程。这两个量（在二维中）根据笛卡尔速度分量定义为：

\begin{sourceblock}
\includegraphics[page=238,trim={53.00bp 168.70bp 53.87bp 469.16bp},clip,width=0.92\textwidth,max height=0.38\textheight]{Ferziger4th.pdf}
\end{sourceblock}
\begin{sourceblock}
\begin{tikzpicture}
\node[inner sep=0,anchor=south west] (image) at (0,0) {\includegraphics[page=238,trim={53.00bp 125.40bp 53.87bp 512.44bp},clip,width=0.92\textwidth,max height=0.38\textheight]{Ferziger4th.pdf}};
\begin{scope}[x={(image.south east)},y={(image.north west)}]
\fill[white] (0.00000,0.80707) rectangle (0.04740,1.00000);
\end{scope}
\end{tikzpicture}
\end{sourceblock}
恒定 ψ 的线是流线（到处都与流动平行的线），该变量因此得名。涡度与旋转运动相关；方程（7.104）是也适用于 3D 的更一般定义的一个特例：

\begin{sourceblock}
\includegraphics[page=239,trim={53.00bp 595.02bp 53.87bp 56.79bp},clip,width=0.92\textwidth,max height=0.38\textheight]{Ferziger4th.pdf}
\end{sourceblock}
在二维流中，涡量矢量与流平面正交，方程（7.105）简化为方程（7.104）。引入流函数的主要原因是，对于ρ、μ、g为常数的流，连续性方程同样满足，无需明确处理。将方程（7.103）代入涡度的定义（7.104）得到连接流函数和涡度的运动学方程：

\begin{sourceblock}
\includegraphics[page=239,trim={53.00bp 474.02bp 53.87bp 162.86bp},clip,width=0.92\textwidth,max height=0.38\textheight]{Ferziger4th.pdf}
\end{sourceblock}
最后，通过分别对 y 和 x 求导 x 和 y 动量方程，并将结果相减，我们得到涡量的动态方程：

\begin{sourceblock}
\begin{tikzpicture}
\node[inner sep=0,anchor=south west] (image) at (0,0) {\includegraphics[page=239,trim={53.00bp 391.15bp 53.87bp 239.70bp},clip,width=0.92\textwidth,max height=0.38\textheight]{Ferziger4th.pdf}};
\begin{scope}[x={(image.south east)},y={(image.north west)}]
\fill[white] (0.22301,0.44007) rectangle (0.26358,0.76764);
\fill[white] (0.20132,0.24199) rectangle (0.22247,0.56957);
\end{scope}
\end{tikzpicture}
\end{sourceblock}
压力没有出现在这两个方程中，即它已作为因变量被消除。因此，纳维-斯托克斯方程被一组仅两个偏微分方程所取代，代替了速度分量和压力的三个偏微分方程。因变量和方程数量的减少使得这种方法具有吸引力。

这两个方程通过涡度方程中 u x 和 u y （它们是 ψ 的导数）的出现以及涡度 ω 充当 ψ 泊松方程中的源项进行耦合。速度分量是通过对流函数求导而获得的。如果需要的话，压力可以通过求解泊松方程来获得，如第1节所述。 7.1.5.1.

这些方程的求解方法如下。给定初始速度场，通过微分计算涡度。然后使用动态涡度方程计算新时间步的涡度；任何标准的时间提前方法都可以用于此目的。有了涡度，就可以通过求解泊松方程来计算新时间步的流函数；可以使用椭圆方程的任何迭代方案。最后，有了流函数，通过微分很容易获得速度分量，我们准备开始下一个时间步的计算。

这种方法的问题在于边界条件，特别是在复杂的几何形状中。因为流动与它们平行，所以实体边界和对称平面是恒定流函数的表面。然而，只有速度已知时才能计算这些边界处的流函数值。一个更困难的问题是，涡度及其在边界处的导数通常都是事先未知的。例如，壁面处的涡度等于 ω wall = − τ wall /μ，其中 τ wall 是壁面剪应力，这通常是我们寻求确定的量。涡度的边界值可以通过流函数计算

\begin{center}\small 图 7.7 用于计算肋上流动的有限差分网格，显示了需要特殊处理以确定涡度边界值的突出角 A 和 B\end{center}
A

B

C

D

通过在垂直于边界的方向上使用单侧有限差分进行微分，参见 Spotz 和 Carey (1995) 或 Spotz (1998)。这种方法通常会减慢收敛速度。

Calhoun (2002) 针对高度不规则边界的问题提出了一种新颖的分数步方法。她采用了嵌入式（又称浸入式）边界方法。虽然卡尔霍恩在她的领域内处理了物体的一般形状，但涡量在边界的尖角处是奇异的，在处理它们时需要特别小心。例如，在图 7.7 中标记为 A 和 B 的角点处，∂ u y /∂ x 和 ∂ u x /∂ y 的导数不连续，这意味着那里的涡度 ω 也不连续，无法使用上述方法计算。一些作者将涡度从内部推断到边界，但这也没有在 A 和 B 处提供唯一的结果。可以导出角点附近涡量的分析行为并用它来校正解，但这很困难，因为必须单独处理每个特殊情况。避免大误差（可能向下游对流）的一种更简单但有效的方法是围绕奇点局部细化网格。

涡流函数方法在二维不可压缩流中得到了广泛的应用。近年来它已经不再流行，因为它扩展到三维流是困难的（但可行；参见 Weinan 和 Liu 1997 的非交错网格上的涡量矢量势方案；Wakashima 和 Saitoh 2004 ）。涡度和流函数都成为三维的三分量矢量，因此可以使用六个偏微分方程组来代替速度-压力公式中所需的四个偏微分方程组。三维方案还继承了处理上述二维流动的可变流体属性、可压缩性和边界条件的困难。

\chapter{Navier--Stokes 方程的求解：第二部分}
我们已经描述了输运方程中各项的离散化方法。论证了不可压缩流中压力和速度分量的联系，并给出了一些求解方法。可以设计许多其他求解纳维-斯托克斯方程的方法。在这里不可能描述所有这些。然而，其中许多与已经描述的方法具有共同的元素。熟悉这些方法应该能让读者理解其他方法。

我们下面详细描述一些代表更大组方法的方法。使用压力校正方程和交错网格的第一个隐式方法（SIMPLE 和迭代隐式分数阶方法）被详细描述，以便可以直接转换为计算机代码。相应的代码和许多注释指向书中相应的方程，可以通过互联网获得；详情请参见附录。然后，我们以类似的模式工作，描述共置网格的这些方法，最后，我们讨论在交错网格和共置网格上实现分数步方法的一些其他问题。本章最后提供了稳态和非稳态流动模拟的示例。

\section*{8.1 交错网格上的隐式迭代方法}\phantomsection\addcontentsline{toc}{section}{8.1 交错网格上的隐式迭代方法}
在本节中，我们提出两种隐式有限体积方案，它们在交错二维笛卡尔网格上使用压力校正方法。一种基于 SIMPLE 算法，另一种基于分数步法；我们将后者称为“隐式分数步法”（IFSM）。它们可用于计算稳态和非稳态流动。复杂几何形状的求解方法将在下一章中介绍。

\begin{sourceblock}
\includegraphics[page=242,trim={53.00bp 484.26bp 53.87bp 52.00bp},clip,width=\textwidth,max height=0.38\textheight]{Ferziger4th.pdf}
\par\vspace{0.45\baselineskip}
\small 图 8.1 交错网格的控制体积：质量守恒和标量（左）、x 动量（中）和 y 动量（右）
\end{sourceblock}
积分形式的纳维-斯托克斯方程如下：

\begin{sourceblock}
\begin{tikzpicture}
\node[inner sep=0,anchor=south west] (image) at (0,0) {\includegraphics[page=242,trim={53.00bp 349.89bp 53.87bp 254.81bp},clip,width=0.92\textwidth,max height=0.38\textheight]{Ferziger4th.pdf}};
\begin{scope}[x={(image.south east)},y={(image.north west)}]
\fill[white] (0.03023,0.40592) rectangle (0.04737,0.57780);
\fill[white] (0.09600,0.27604) rectangle (0.13823,0.47559);
\fill[white] (0.14162,0.29818) rectangle (0.17771,0.47168);
\fill[white] (0.18298,0.30355) rectangle (0.20884,0.47559);
\fill[white] (0.02526,0.19596) rectangle (0.05209,0.37191);
\fill[white] (0.07266,0.17301) rectangle (0.09137,0.32227);
\end{scope}
\end{tikzpicture}
\end{sourceblock}
为了方便起见，假设唯一的体积力是浮力。宏观动量通量向量 t i，见方程（1.18），被分成粘性贡献 τ i j i j 和压力贡献 p i i。我们假设除浮力项外的密度常数，即我们使用 Boussinesq 近似。平均重力被纳入压力项中，如第 1 节所示。 1.4.

典型的交错控制体积如图 8.1 所示。 u x 和 u y 的控制体积相对于连续性方程的控制体积发生位移。对于非均匀网格，速度节点不在其控制体积的中心。单元面向 u x 的“e”和“w”，以及 u y 的“n”和“s”位于节点之间的中间。为了方便起见，我们有时会使用 u 代替 u x，使用 v 代替 u y。

SIMPLE 和 IFSM 的大部分离散化和求解算法是相同的；这些差异仅限于压力校正方程，在此将予以强调。我们将使用第 1 节中描述的二阶隐式三时间级方案。 6.2.4及时整合。这导致不稳定项的近似如下：

\begin{sourceblock}
\begin{tikzpicture}
\node[inner sep=0,anchor=south west] (image) at (0,0) {\includegraphics[page=242,trim={53.00bp 109.07bp 53.87bp 519.65bp},clip,width=0.92\textwidth,max height=0.38\textheight]{Ferziger4th.pdf}};
\begin{scope}[x={(image.south east)},y={(image.north west)}]
\fill[white] (0.09792,0.47167) rectangle (0.11648,0.78060);
\fill[white] (0.09242,0.08899) rectangle (0.12168,0.40513);
\fill[white] (0.14457,0.05746) rectangle (0.16220,0.28648);
\end{scope}
\end{tikzpicture}
\end{sourceblock}
\begin{sourceblock}
\begin{tikzpicture}
\node[inner sep=0,anchor=south west] (image) at (0,0) {\includegraphics[page=243,trim={53.00bp 573.29bp 53.87bp 64.47bp},clip,width=0.92\textwidth,max height=0.38\textheight]{Ferziger4th.pdf}};
\begin{scope}[x={(image.south east)},y={(image.north west)}]
\fill[white] (0.00000,0.80796) rectangle (0.07732,1.00000);
\fill[white] (0.26481,0.53735) rectangle (0.30093,0.95772);
\fill[white] (0.32785,0.54087) rectangle (0.35098,0.94820);
\fill[white] (0.18884,0.29493) rectangle (0.21838,0.76357);
\fill[white] (0.20773,0.22199) rectangle (0.22421,0.52431);
\fill[white] (0.22953,0.30409) rectangle (0.25771,0.71177);
\end{scope}
\end{tikzpicture}
\end{sourceblock}
从现在开始，我们将删除上标 n + 1；除非另有说明，所有项均在 t n + 1 处求值。由于该方案很简单，因此方程需要迭代解。如果时间步长与显式方案中使用的时间步长一样小，则每个时间步长迭代一到两次就足够了。对于具有缓慢瞬态的流，我们可以使用更大的时间步长，并且每个时间步长需要更多的迭代。如前所述，这些迭代称为外迭代，以区别于用于求解线性方程（例如压力校正方程）的内迭代。我们假设第 1 章中描述的求解器之一。 5 用于后者并且应集中于外部迭代。

我们现在考虑对流和扩散通量以及源项的近似。表面积分可以分为四个 CV 面积分。让我们把注意力集中在CV面“e”上；其他面同样处理，通过索引代换即可得到结果。我们采用第二章中提出的二阶中心差分近似。 4.通过假设 CV 面中心处的量值表示该面上的平均值（中点规则近似）来近似通量。在第 m 次外迭代中，所有非线性项均由“旧”值（来自前一外迭代）和“新”值的乘积来近似。因此，在离散动量方程时，使用现有速度场评估通过每个 CV 面的质量通量，并假设已知：

\begin{sourceblock}
\begin{tikzpicture}
\node[inner sep=0,anchor=south west] (image) at (0,0) {\includegraphics[page=243,trim={53.00bp 291.14bp 53.87bp 350.15bp},clip,width=0.92\textwidth,max height=0.38\textheight]{Ferziger4th.pdf}};
\begin{scope}[x={(image.south east)},y={(image.north west)}]
\fill[white] (0.28842,0.40604) rectangle (0.33179,0.94406);
\end{scope}
\end{tikzpicture}
\end{sourceblock}
这种线性化本质上是皮卡德迭代的第一步；隐式方案的其他线性化在第 1 章中进行了描述。 5.除非另有明确说明，本节其余部分中的所有变量都属于第 m 个外迭代。质量通量 (8.4) 满足“标量”CV 上的连续性方程，见图 8.1。动量 CV 表面的质量通量必须通过插值获得；理想情况下，这些通量将为动量 CV 提供质量守恒，但这只能保证插值的准确性。另一种可能性是使用标量 CV 面的质量通量。由于 u -CV 的东面和西面位于标量 CV 面的中间，因此质量通量可计算为：

\begin{sourceblock}
\begin{tikzpicture}
\node[inner sep=0,anchor=south west] (image) at (0,0) {\includegraphics[page=243,trim={53.00bp 141.16bp 53.87bp 496.95bp},clip,width=0.92\textwidth,max height=0.38\textheight]{Ferziger4th.pdf}};
\begin{scope}[x={(image.south east)},y={(image.north west)}]
\fill[white] (0.00000,0.81555) rectangle (0.12556,1.00000);
\fill[white] (0.12830,0.81555) rectangle (0.16638,1.00000);
\end{scope}
\end{tikzpicture}
\end{sourceblock}
通过 u-CV 北面和南面的质量通量可近似为两个标量 CV 面质量通量之和的一半：

\begin{sourceblock}
\includegraphics[page=243,trim={53.00bp 76.06bp 53.87bp 562.06bp},clip,width=0.92\textwidth,max height=0.38\textheight]{Ferziger4th.pdf}
\end{sourceblock}
上标u表示索引指的是u-CV，见图8.1。因此，u -CV 的四个质量通量之和是进入两个相邻标量 CV 的质量通量之和的一半。因此，它们满足双标量 CV 的连续性方程，因此通过 u-CV 面的质量通量也守恒质量。该结果也适用于 v 动量 CV。必须保证通过动量CV的质量通量满足连续性方程；否则，动量将不守恒。

则 u i -动量通过 u i -CV 的‘e’面的对流通量为（见 4.2 节和方程（8.4））：

\begin{sourceblock}
\begin{tikzpicture}
\node[inner sep=0,anchor=south west] (image) at (0,0) {\includegraphics[page=244,trim={53.00bp 462.48bp 53.87bp 179.65bp},clip,width=0.92\textwidth,max height=0.38\textheight]{Ferziger4th.pdf}};
\begin{scope}[x={(image.south east)},y={(image.north west)}]
\fill[white] (0.30713,0.39650) rectangle (0.34229,0.95002);
\end{scope}
\end{tikzpicture}
\end{sourceblock}
尽管需要具有相同精度的近似值，但该表达式中使用的 u i 的 CV 面值不必是用于计算质量通量的值。线性插值是最简单的二阶近似；见第 4.4.2 节 详情。我们将其称为中心差分方案 (CDS)，尽管不涉及差分。这是因为，在均匀网格上，它会产生与 CDS 有限差分法相同的代数方程。

当应用于从对流通量的中心差近似导出的代数方程组时，一些迭代求解器无法收敛。这是因为矩阵可能不是对角占优的；这些方程最好使用第 1 节中描述的延迟修正方法来求解。 5.6.在该方法中，通量表示为：

\begin{sourceblock}
\includegraphics[page=244,trim={53.00bp 296.45bp 53.87bp 348.93bp},clip,width=0.92\textwidth,max height=0.38\textheight]{Ferziger4th.pdf}
\end{sourceblock}
其中上标 CDS 和 UDS 分别表示中心差和迎风差的近似值（参见第 4.4 节）。括号中的项使用前一次迭代的值进行计算，而矩阵则使用 UDSapproximation 计算。在收敛时，UDS 贡献相互抵消，留下 CDS 解。该过程通常以大约纯迎风近似所获得的速率收敛。

这两种方案也可以混合使用；这是通过将显式部分（方程（8.8）中括号内的项）乘以因子0≤γ≤1来实现的。这种做法可以消除粗网格上中心差异所获得的振荡。然而，它以降低准确性为代价提高了结果的美观性。混合可以在本地使用，例如，允许使用 CDS 计算具有冲击的流量；这比到处应用它更好。

计算扩散通量需要评估 CV 面“e”处的应力 τ xx 和 τ yx。因为该 CV 面上的向外单位法向量是 i，所以我们有：

\begin{sourceblock}
\begin{tikzpicture}
\node[inner sep=0,anchor=south west] (image) at (0,0) {\includegraphics[page=244,trim={53.00bp 84.14bp 53.87bp 557.41bp},clip,width=0.92\textwidth,max height=0.38\textheight]{Ferziger4th.pdf}};
\begin{scope}[x={(image.south east)},y={(image.north west)}]
\fill[white] (0.32496,0.41033) rectangle (0.36132,0.95120);
\end{scope}
\end{tikzpicture}
\end{sourceblock}
\begin{sourceblock}
\begin{tikzpicture}
\node[inner sep=0,anchor=south west] (image) at (0,0) {\includegraphics[page=245,trim={53.00bp 593.89bp 53.87bp 55.42bp},clip,width=0.92\textwidth,max height=0.38\textheight]{Ferziger4th.pdf}};
\begin{scope}[x={(image.south east)},y={(image.north west)}]
\fill[white] (0.64623,0.00000) rectangle (0.69161,0.84730);
\fill[white] (0.68998,0.09150) rectangle (0.72950,0.61616);
\fill[white] (0.73290,0.14439) rectangle (0.78298,0.84730);
\fill[white] (0.78135,0.09150) rectangle (0.83356,0.84730);
\fill[white] (0.83777,0.13844) rectangle (0.87750,0.82531);
\fill[white] (0.88171,0.13844) rectangle (0.92316,0.82531);
\fill[white] (0.92734,0.13844) rectangle (1.00000,0.83125);
\fill[white] (0.00000,0.00000) rectangle (0.05074,0.11468);
\fill[white] (0.05374,0.00000) rectangle (0.15011,0.11468);
\fill[white] (0.15305,0.00000) rectangle (0.17949,0.11468);
\fill[white] (0.18253,0.00000) rectangle (0.22893,0.11468);
\fill[white] (0.23194,0.00000) rectangle (0.28635,0.11468);
\fill[white] (0.28938,0.00000) rectangle (0.37901,0.11468);
\fill[white] (0.38205,0.00000) rectangle (0.56159,0.11468);
\fill[white] (0.56463,0.00000) rectangle (0.59441,0.11468);
\fill[white] (0.59741,0.00000) rectangle (0.63886,0.11468);
\fill[white] (0.64183,0.00000) rectangle (0.78370,0.11468);
\fill[white] (0.78680,0.00000) rectangle (0.87311,0.11468);
\fill[white] (0.87615,0.00000) rectangle (1.00000,0.11468);
\end{scope}
\end{tikzpicture}
\end{sourceblock}
2 ( y j + 1 − y j − 1 ) 对于 v -CV。 CV 面的应力需要导数的近似；中心差近似导致：

\begin{sourceblock}
\begin{tikzpicture}
\node[inner sep=0,anchor=south west] (image) at (0,0) {\includegraphics[page=245,trim={53.00bp 495.23bp 53.87bp 103.28bp},clip,width=0.92\textwidth,max height=0.38\textheight]{Ferziger4th.pdf}};
\begin{scope}[x={(image.south east)},y={(image.north west)}]
\fill[white] (0.27669,0.69097) rectangle (0.35116,0.87905);
\fill[white] (0.35636,0.70812) rectangle (0.38454,0.87905);
\fill[white] (0.38806,0.70265) rectangle (0.40785,0.87358);
\fill[white] (0.16851,0.14077) rectangle (0.24295,0.32870);
\fill[white] (0.24815,0.15777) rectangle (0.27633,0.32870);
\fill[white] (0.27985,0.15777) rectangle (0.30523,0.32870);
\end{scope}
\end{tikzpicture}
\end{sourceblock}
请注意，τ xx 在 u -CV 的“e”面评估，τ yx 在 v -CV 的“e”面评估，因此索引指的是相应 CV 上的位置，见图 8.1。因此，v -CV 上的 u ne 和 use 实际上是 u 速度的节点值，不需要插值。

在其他 CV 面上，我们获得了类似的表达式。对于 u -CV，我们需要在面“e”和“w”处近似 τ xx，在面“n”和“s”处近似 τ xy。对于 v -CV，在面“e”和“w”处需要 τ yx，在面“n”和“s”处需要 τ yy。

压力项近似为：

\begin{sourceblock}
\begin{tikzpicture}
\node[inner sep=0,anchor=south west] (image) at (0,0) {\includegraphics[page=245,trim={53.00bp 354.47bp 53.87bp 287.47bp},clip,width=0.92\textwidth,max height=0.38\textheight]{Ferziger4th.pdf}};
\begin{scope}[x={(image.south east)},y={(image.north west)}]
\fill[white] (0.22851,0.39050) rectangle (0.26710,0.93967);
\end{scope}
\end{tikzpicture}
\end{sourceblock}
对于 u 方程和

\begin{sourceblock}
\begin{tikzpicture}
\node[inner sep=0,anchor=south west] (image) at (0,0) {\includegraphics[page=245,trim={53.00bp 297.22bp 53.87bp 344.72bp},clip,width=0.92\textwidth,max height=0.38\textheight]{Ferziger4th.pdf}};
\begin{scope}[x={(image.south east)},y={(image.north west)}]
\fill[white] (0.23411,0.39050) rectangle (0.27269,0.93967);
\end{scope}
\end{tikzpicture}
\end{sourceblock}
为 v 方程。在笛卡尔网格上，“n”和“s”CV 面对 u 方程或“e”和“w”面对 v 方程没有压力贡献。

如果存在浮力，则浮力近似为：

\begin{sourceblock}
\begin{tikzpicture}
\node[inner sep=0,anchor=south west] (image) at (0,0) {\includegraphics[page=245,trim={53.00bp 204.11bp 53.87bp 437.71bp},clip,width=0.92\textwidth,max height=0.38\textheight]{Ferziger4th.pdf}};
\begin{scope}[x={(image.south east)},y={(image.north west)}]
\fill[white] (0.21005,0.38816) rectangle (0.24860,0.93503);
\end{scope}
\end{tikzpicture}
\end{sourceblock}
\begin{sourceblock}
\begin{tikzpicture}
\node[inner sep=0,anchor=south west] (image) at (0,0) {\includegraphics[page=245,trim={53.00bp 168.55bp 53.87bp 468.81bp},clip,width=0.92\textwidth,max height=0.38\textheight]{Ferziger4th.pdf}};
\begin{scope}[x={(image.south east)},y={(image.north west)}]
\fill[white] (0.09074,0.00000) rectangle (0.14156,0.45274);
\fill[white] (0.14695,0.01112) rectangle (0.27041,0.45274);
\fill[white] (0.26878,0.01112) rectangle (0.30830,0.31793);
\fill[white] (0.31170,0.04170) rectangle (0.36174,0.45274);
\fill[white] (0.36012,0.01112) rectangle (0.41236,0.45274);
\fill[white] (0.41266,0.03822) rectangle (0.45239,0.44024);
\fill[white] (0.45266,0.03822) rectangle (0.49411,0.44024);
\fill[white] (0.49441,0.03822) rectangle (0.56857,0.44371);
\fill[white] (0.56890,0.03822) rectangle (0.62481,0.44024);
\fill[white] (0.62514,0.03822) rectangle (0.69149,0.44024);
\fill[white] (0.69182,0.03822) rectangle (0.75654,0.44024);
\fill[white] (0.75684,0.03822) rectangle (0.82319,0.44024);
\fill[white] (0.82346,0.03822) rectangle (0.86983,0.44024);
\fill[white] (0.87017,0.03822) rectangle (0.90325,0.44024);
\fill[white] (0.90358,0.03822) rectangle (1.00000,0.44024);
\end{scope}
\end{tikzpicture}
\end{sourceblock}
完整的 u i -动量方程的近似值为：

\begin{sourceblock}
\includegraphics[page=245,trim={53.00bp 118.16bp 53.87bp 528.65bp},clip,width=0.92\textwidth,max height=0.38\textheight]{Ferziger4th.pdf}
\end{sourceblock}
where

\begin{sourceblock}
\begin{tikzpicture}
\node[inner sep=0,anchor=south west] (image) at (0,0) {\includegraphics[page=245,trim={53.00bp 70.50bp 53.87bp 574.89bp},clip,width=0.92\textwidth,max height=0.38\textheight]{Ferziger4th.pdf}};
\begin{scope}[x={(image.south east)},y={(image.north west)}]
\fill[white] (0.17630,0.07422) rectangle (0.19044,0.48723);
\fill[white] (0.19795,0.18699) rectangle (0.22614,0.74410);
\fill[white] (0.23005,0.17398) rectangle (0.26523,0.81446);
\end{scope}
\end{tikzpicture}
\end{sourceblock}
如果 ρ 和 μ 恒定，则部分扩散通量项会根据连续性方程被抵消，请参见第 1 节。 7.1. （它可能不会在数值近似中完全抵消，但可以通过在离散化之前删除这些项来简化方程）。例如，在 u 方程中，“e”和“w”面上的 τ xx 项将减少一半，而在“n”和“s”面上的 τ yx 项中，∂ v /∂ x 贡献将被删除。即使当 ρ 和 μ 不是常数时，这些项的总和对 F d 的贡献也很小。这就是为什么有一个明确的“扩散源项”，例如，对于 u：

\begin{sourceblock}
\begin{tikzpicture}
\node[inner sep=0,anchor=south west] (image) at (0,0) {\includegraphics[page=246,trim={53.00bp 443.75bp 53.87bp 163.06bp},clip,width=0.92\textwidth,max height=0.38\textheight]{Ferziger4th.pdf}};
\begin{scope}[x={(image.south east)},y={(image.north west)}]
\fill[white] (0.23083,0.66273) rectangle (0.26941,0.88674);
\end{scope}
\end{tikzpicture}
\end{sourceblock}
通常是根据先前的外部迭代 m − 1 计算并明确处理的。只有 F d − Q d
u 被隐式对待。这种近似的结果是，在同位网格上，方程（8.15）隐含的矩阵对于所有三个速度分量都是相同的。

当所有通量和源项的近似值代入方程（8.15）时，我们得到以下形式的代数方程：

\begin{sourceblock}
\begin{tikzpicture}
\node[inner sep=0,anchor=south west] (image) at (0,0) {\includegraphics[page=246,trim={53.00bp 326.44bp 53.87bp 312.68bp},clip,width=0.92\textwidth,max height=0.38\textheight]{Ferziger4th.pdf}};
\begin{scope}[x={(image.south east)},y={(image.north west)}]
\fill[white] (0.23035,0.46077) rectangle (0.26454,0.95559);
\end{scope}
\end{tikzpicture}
\end{sourceblock}
v 的方程具有相同的形式。系数取决于所使用的近似值；对于上面应用的 UDS 近似，u 方程的系数为：

\begin{sourceblock}
\begin{tikzpicture}
\node[inner sep=0,anchor=south west] (image) at (0,0) {\includegraphics[page=246,trim={53.00bp 173.41bp 53.87bp 392.60bp},clip,width=0.92\textwidth,max height=0.38\textheight]{Ferziger4th.pdf}};
\begin{scope}[x={(image.south east)},y={(image.north west)}]
\fill[white] (0.26144,0.12424) rectangle (0.29564,0.25776);
\end{scope}
\end{tikzpicture}
\end{sourceblock}
请注意，对流和扩散通量对中心系数 A P 的贡献可以表示为相邻系数的负和；实际上，当从离散通量近似中提取这一贡献时，可以得到：

\begin{sourceblock}
\begin{tikzpicture}
\node[inner sep=0,anchor=south west] (image) at (0,0) {\includegraphics[page=246,trim={53.00bp 94.23bp 53.87bp 544.90bp},clip,width=0.92\textwidth,max height=0.38\textheight]{Ferziger4th.pdf}};
\begin{scope}[x={(image.south east)},y={(image.north west)}]
\fill[white] (0.31808,0.46057) rectangle (0.35227,0.95557);
\fill[white] (0.33696,0.38393) rectangle (0.35344,0.70122);
\fill[white] (0.35880,0.47057) rectangle (0.38698,0.89856);
\fill[white] (0.39347,0.46057) rectangle (0.42301,0.95298);
\end{scope}
\end{tikzpicture}
\end{sourceblock}
右侧最后一项表示离散连续性方程，对于不可压缩流，该方程为零，因此在方程（8.19）中被省略。对于所有保守方案都是如此。

另请注意，以节点 P 为中心的 CV 的 ̇ m w 等于以节点 W 为中心的 CV 的 − ̇ m e。源项 Qu
P 不仅包含压力和浮力项，还包含延迟校正产生的对流和扩散通量部分以及非稳态项的贡献，即：

\begin{sourceblock}
\includegraphics[page=247,trim={53.00bp 497.06bp 53.87bp 150.90bp},clip,width=0.92\textwidth,max height=0.38\textheight]{Ferziger4th.pdf}
\end{sourceblock}
\begin{sourceblock}
\begin{tikzpicture}
\node[inner sep=0,anchor=south west] (image) at (0,0) {\includegraphics[page=247,trim={53.00bp 461.19bp 53.87bp 180.29bp},clip,width=0.92\textwidth,max height=0.38\textheight]{Ferziger4th.pdf}};
\begin{scope}[x={(image.south east)},y={(image.north west)}]
\fill[white] (0.00000,0.62733) rectangle (0.07732,1.00000);
\fill[white] (0.30096,0.14680) rectangle (0.33838,0.68573);
\end{scope}
\end{tikzpicture}
\end{sourceblock}
该“对流源”是使用先前外部迭代 m − 1 的速度计算的。

v方程中的系数以相同的方式获得并且具有相同的形式；然而，网格位置“e”、“n”等具有不同的坐标，见图8.1。

线性动量方程通过序贯求解方法（参见第 5.4 节）求解，使用“旧”质量通量和来自先前外部迭代的压力。这会产生新的速度 u * 和 v *，它们不一定满足连续性方程，因此：

\begin{sourceblock}
\includegraphics[page=247,trim={53.00bp 329.69bp 53.87bp 318.09bp},clip,width=0.92\textwidth,max height=0.38\textheight]{Ferziger4th.pdf}
\end{sourceblock}
其中质量通量是根据方程（8.4）使用u*和v*计算的。由于变量的排列是交错的，因此质量 CV 上的单元面速度是节点值。除非另有说明，以下索引均指此 CV，见图 8.1。

SIMPLE 和 IFSM 中离散压力校正方程的推导遵循略有不同的路线；我们首先使用 SIMPLE 算法。

\subsection*{8.1.1 交错网格的简单}\phantomsection\addcontentsline{toc}{subsection}{8.1.1 交错网格的简单}
从动量方程计算出的速度分量 u ＊ 和 v ＊ 可以表示如下（通过将方程（8.18）除以 A P 并显式地写出压力项；注意质量 CV 上的索引“e”表示 u -CV 上的索引 P）：

\begin{sourceblock}
\includegraphics[page=247,trim={53.00bp 114.76bp 53.87bp 521.19bp},clip,width=0.92\textwidth,max height=0.38\textheight]{Ferziger4th.pdf}
\end{sourceblock}
\begin{sourceblock}
\begin{tikzpicture}
\node[inner sep=0,anchor=south west] (image) at (0,0) {\includegraphics[page=248,trim={53.00bp 556.17bp 53.87bp 55.57bp},clip,width=0.92\textwidth,max height=0.38\textheight]{Ferziger4th.pdf}};
\begin{scope}[x={(image.south east)},y={(image.north west)}]
\fill[white] (0.09657,0.69982) rectangle (0.11071,0.85754);
\fill[white] (0.11651,0.73621) rectangle (0.14129,0.94871);
\fill[white] (0.14397,0.73621) rectangle (0.26692,0.94871);
\fill[white] (0.26962,0.73621) rectangle (0.37263,0.94871);
\fill[white] (0.37537,0.73621) rectangle (0.41513,0.94871);
\end{scope}
\end{tikzpicture}
\end{sourceblock}
Analogously, v ∗
n 可以表示为：

\begin{sourceblock}
\includegraphics[page=248,trim={53.00bp 497.19bp 53.87bp 138.74bp},clip,width=0.92\textwidth,max height=0.38\textheight]{Ferziger4th.pdf}
\end{sourceblock}
需要修正速度 u * 和 v * 以强制质量守恒。这是按照第 2 节中概述的那样完成的。 7.2.2、通过校正压力，使̃ u 和̃ v 保持不变。修正后的速度 - 第 m 次外迭代的最终值 − u m = u ∗ + u ′
和 v m = v ∗ + v ′ 需要满足线性动量方程，这只有在压力被校正的情况下才有可能。所以我们写：

\begin{sourceblock}
\includegraphics[page=248,trim={53.00bp 390.92bp 53.87bp 245.02bp},clip,width=0.92\textwidth,max height=0.38\textheight]{Ferziger4th.pdf}
\end{sourceblock}
and

\begin{sourceblock}
\begin{tikzpicture}
\node[inner sep=0,anchor=south west] (image) at (0,0) {\includegraphics[page=248,trim={53.00bp 346.87bp 53.87bp 289.08bp},clip,width=0.92\textwidth,max height=0.38\textheight]{Ferziger4th.pdf}};
\begin{scope}[x={(image.south east)},y={(image.north west)}]
\fill[white] (0.00000,0.82809) rectangle (0.04740,1.00000);
\end{scope}
\end{tikzpicture}
\end{sourceblock}
其中 pm = pm − 1 + p ′ 是新压力。这些指标指的是质量CV。由式（8.27）减去式（8.24）得到速度与压力修正值的关系：

\begin{sourceblock}
\begin{tikzpicture}
\node[inner sep=0,anchor=south west] (image) at (0,0) {\includegraphics[page=248,trim={53.00bp 262.00bp 53.87bp 366.32bp},clip,width=0.92\textwidth,max height=0.38\textheight]{Ferziger4th.pdf}};
\begin{scope}[x={(image.south east)},y={(image.north west)}]
\fill[white] (0.24848,0.28556) rectangle (0.27552,0.64410);
\end{scope}
\end{tikzpicture}
\end{sourceblock}
类似地，可得：

\begin{sourceblock}
\begin{tikzpicture}
\node[inner sep=0,anchor=south west] (image) at (0,0) {\includegraphics[page=248,trim={53.00bp 202.69bp 53.87bp 425.63bp},clip,width=0.92\textwidth,max height=0.38\textheight]{Ferziger4th.pdf}};
\begin{scope}[x={(image.south east)},y={(image.north west)}]
\fill[white] (0.24770,0.28556) rectangle (0.27236,0.64410);
\end{scope}
\end{tikzpicture}
\end{sourceblock}
修正后的速度需要满足连续性方程，因此我们将 u m 和 v m 代入质量通量表达式（8.4），并使用式（8.23）：

\begin{sourceblock}
\includegraphics[page=248,trim={53.00bp 144.26bp 53.87bp 503.51bp},clip,width=0.92\textwidth,max height=0.38\textheight]{Ferziger4th.pdf}
\end{sourceblock}
最后，将上述表达式（8.29）和（8.30）中的u′和v′代入连续性方程，得到压力校正方程：

\begin{sourceblock}
\begin{tikzpicture}
\node[inner sep=0,anchor=south west] (image) at (0,0) {\includegraphics[page=248,trim={53.00bp 73.98bp 53.87bp 564.18bp},clip,width=0.92\textwidth,max height=0.38\textheight]{Ferziger4th.pdf}};
\begin{scope}[x={(image.south east)},y={(image.north west)}]
\fill[white] (0.31952,0.44496) rectangle (0.35612,0.95711);
\end{scope}
\end{tikzpicture}
\end{sourceblock}
其中系数为：

\begin{sourceblock}
\begin{tikzpicture}
\node[inner sep=0,anchor=south west] (image) at (0,0) {\includegraphics[page=249,trim={53.00bp 518.97bp 53.87bp 70.98bp},clip,width=0.92\textwidth,max height=0.38\textheight]{Ferziger4th.pdf}};
\begin{scope}[x={(image.south east)},y={(image.north west)}]
\fill[white] (0.24586,0.64549) rectangle (0.28247,0.83357);
\fill[white] (0.26475,0.61688) rectangle (0.28238,0.72949);
\fill[white] (0.28815,0.64890) rectangle (0.31633,0.80076);
\fill[white] (0.31985,0.64890) rectangle (0.34803,0.80076);
\fill[white] (0.25110,0.14175) rectangle (0.28770,0.32983);
\fill[white] (0.26998,0.11327) rectangle (0.28992,0.22575);
\fill[white] (0.29525,0.14529) rectangle (0.32343,0.29702);
\fill[white] (0.32695,0.14529) rectangle (0.35513,0.29702);
\end{scope}
\end{tikzpicture}
\end{sourceblock}
\begin{sourceblock}
\begin{tikzpicture}
\node[inner sep=0,anchor=south west] (image) at (0,0) {\includegraphics[page=249,trim={53.00bp 483.71bp 53.87bp 154.46bp},clip,width=0.92\textwidth,max height=0.38\textheight]{Ferziger4th.pdf}};
\begin{scope}[x={(image.south east)},y={(image.north west)}]
\fill[white] (0.28198,0.44512) rectangle (0.31856,0.95710);
\end{scope}
\end{tikzpicture}
\end{sourceblock}
求解压力校正方程后，对速度和压力进行校正。正如第 1 节中所指出的。 7.2.2 中，如果尝试使用非常大的时间步长来计算稳态流，则动量方程必须是欠松弛的，如第 7.2.2 节中所述。 5.4.2，因此仅将部分压力修正p ' 添加到pm − 1 中。在具有大时间步长的不稳态计算中也可能需要欠松弛。

修正后的速度满足连续性方程，其精度可用于求解压力修正方程。然而，它们不满足非线性动量方程，因为方程（8.27）和（8.28）中的̃ u 和 ̃ v 没有得到修正，所以我们必须开始另一次外迭代。当连续性方程和动量方程都满足所需的公差时，u ′
i 和 p ′ 将可以忽略不计，我们可以进入下一个时间级别。为了在新的时间步开始迭代，前一个时间步的解提供了初始猜测。这可以通过使用外推法来改进；任何显式时间推进方案都可以作为预测器。对于小时间步长，外推法相当准确，并且可以节省一些迭代次数。

如果计算稳态流，则可以采用单个无限时间步长；由不稳定项产生的所有贡献都会消失，并继续进行外部迭代，直到所有修正都可以忽略不计。计算非定常流时，每个时间步长的外部迭代次数通常在 3（对于小时间步长和高欠松弛因子）和 10（对于较大时间步长和中等欠松弛因子）之间。一些例子将在第三节中介绍。 8.4.使用该算法的计算机代码可通过互联网获得；详细内容请参见附录。

请注意，压力校正方程 (8.33) 中的系数与面部面积的平方成正比。因此，比率 A E / A N 与纵横比 a r = y / x 的平方成正比。如果细胞高度拉伸（例如，当要解析粘性子层时，靠近墙壁），这可能会使压力校正方程变得僵硬并且更难以求解。如果可能，应避免纵横比大于 100，因为一个方向（例如，上述示例中垂直于墙壁）上的系数比其他方向上的系数大 10 4 倍以上。

上述算法很容易修改为给出第 1 节中描述的 SIMPLEC 方法。 7.2.2.压力校正方程的形式为 ( 8.32 )，但表达式为

\begin{sourceblock}
\begin{tikzpicture}
\node[inner sep=0,anchor=south west] (image) at (0,0) {\includegraphics[page=250,trim={53.00bp 592.27bp 53.87bp 50.00bp},clip,width=0.92\textwidth,max height=0.38\textheight]{Ferziger4th.pdf}};
\begin{scope}[x={(image.south east)},y={(image.north west)}]
\fill[white] (0.00000,0.16548) rectangle (0.03907,0.64977);
\fill[white] (0.03886,0.16548) rectangle (0.18848,0.64977);
\fill[white] (0.18830,0.16548) rectangle (0.31762,0.64977);
\fill[white] (0.32039,0.16967) rectangle (0.35459,0.71512);
\fill[white] (0.76935,0.07876) rectangle (0.78349,0.43821);
\fill[white] (0.78704,0.16548) rectangle (0.83510,0.64977);
\fill[white] (0.83786,0.16967) rectangle (0.87089,0.71052);
\fill[white] (0.85675,0.07834) rectangle (0.87323,0.43779);
\fill[white] (0.87693,0.18056) rectangle (0.90511,0.66527);
\fill[white] (0.93829,0.07876) rectangle (0.95239,0.43821);
\fill[white] (0.95832,0.16967) rectangle (0.99134,0.71052);
\fill[white] (0.97720,0.07876) rectangle (1.00000,0.64977);
\fill[white] (0.00000,0.00000) rectangle (0.15320,0.14872);
\fill[white] (0.15353,0.00000) rectangle (0.20493,0.14872);
\fill[white] (0.20526,0.00000) rectangle (0.32442,0.14872);
\fill[white] (0.32475,0.00000) rectangle (0.35287,0.14872);
\fill[white] (0.35314,0.00000) rectangle (0.39456,0.14872);
\fill[white] (0.39486,0.00000) rectangle (0.46457,0.14872);
\fill[white] (0.46487,0.00000) rectangle (0.58620,0.14872);
\fill[white] (0.58650,0.00000) rectangle (0.61131,0.14872);
\fill[white] (0.61155,0.00000) rectangle (0.66463,0.14872);
\fill[white] (0.66490,0.00000) rectangle (0.86159,0.14872);
\fill[white] (0.86189,0.00000) rectangle (0.91329,0.14872);
\fill[white] (0.91356,0.00000) rectangle (1.00000,0.14872);
\end{scope}
\end{tikzpicture}
\end{sourceblock}
分别为 k。 PISO 算法的扩展也很简单。第二个压力校正方程与第一个压力校正方程具有相同的系数矩阵，但源项现在基于 ̃ u ′
我。在第一个压力修正方程中忽略了这一贡献，但现在可以使用第一个速度修正 u ′ 来计算它
我。高阶离散化很容易纳入上述解策略中。边界条件的实现在第 2 节中进行了讨论。 7.1.6.

\subsection*{8.1.2 交错网格的 IFSM}\phantomsection\addcontentsline{toc}{subsection}{8.1.2 交错网格的 IFSM}
SIMPLE 和 IFSM 之间的主要区别在于后者始终使用有限时间步长，即使在计算稳态流时也是如此。然而，它不需要欠松弛，因此只需为每个应用选择合适的时间步长。

速度与压力修正之间的关系由式（7.64）推导出来，即速度仅在非定常项进行修正，留下u*
i 在对流和扩散通量中：

\begin{sourceblock}
\begin{tikzpicture}
\node[inner sep=0,anchor=south west] (image) at (0,0) {\includegraphics[page=250,trim={53.00bp 285.94bp 53.87bp 291.61bp},clip,width=0.92\textwidth,max height=0.38\textheight]{Ferziger4th.pdf}};
\begin{scope}[x={(image.south east)},y={(image.north west)}]
\fill[white] (0.19612,0.11751) rectangle (0.22316,0.27057);
\end{scope}
\end{tikzpicture}
\end{sourceblock}
类似地：

\begin{sourceblock}
\begin{tikzpicture}
\node[inner sep=0,anchor=south west] (image) at (0,0) {\includegraphics[page=250,trim={53.00bp 227.02bp 53.87bp 401.69bp},clip,width=0.92\textwidth,max height=0.38\textheight]{Ferziger4th.pdf}};
\begin{scope}[x={(image.south east)},y={(image.north west)}]
\fill[white] (0.18755,0.27812) rectangle (0.21221,0.64040);
\end{scope}
\end{tikzpicture}
\end{sourceblock}
将这些表达式代入连续性方程（8.31），得到与SIMPLE情况下的压力校正方程相同的形式，方程（8.32），只是系数不同；现在我们有：而不是等式（8.33）：

\begin{sourceblock}
\includegraphics[page=250,trim={53.00bp 110.47bp 53.87bp 490.74bp},clip,width=0.92\textwidth,max height=0.38\textheight]{Ferziger4th.pdf}
\end{sourceblock}
\begin{sourceblock}
\begin{tikzpicture}
\node[inner sep=0,anchor=south west] (image) at (0,0) {\includegraphics[page=250,trim={53.00bp 75.03bp 53.87bp 533.14bp},clip,width=0.92\textwidth,max height=0.38\textheight]{Ferziger4th.pdf}};
\begin{scope}[x={(image.south east)},y={(image.north west)}]
\fill[white] (0.35880,0.89408) rectangle (0.37859,1.00000);
\fill[white] (0.39802,0.88046) rectangle (0.44039,1.00000);
\fill[white] (0.64589,0.89408) rectangle (0.66568,1.00000);
\fill[white] (0.68511,0.88046) rectangle (0.72517,1.00000);
\fill[white] (0.22806,0.77540) rectangle (0.26466,1.00000);
\fill[white] (0.24695,0.73762) rectangle (0.26689,0.88563);
\fill[white] (0.27224,0.77989) rectangle (0.30042,0.97930);
\fill[white] (0.30391,0.77989) rectangle (0.33212,0.97930);
\fill[white] (0.51762,0.77540) rectangle (0.55423,1.00000);
\fill[white] (0.53651,0.73762) rectangle (0.55299,0.88563);
\fill[white] (0.55991,0.77989) rectangle (0.58809,0.97930);
\fill[white] (0.59158,0.77989) rectangle (0.61976,0.97930);
\fill[white] (0.92439,0.77350) rectangle (1.00000,0.97309);
\fill[white] (0.28198,0.21477) rectangle (0.31856,0.46179);
\end{scope}
\end{tikzpicture}
\end{sourceblock}
计算非稳态流动时，选择时间步长，以便充分解决流动变化（例如，如果流动是周期性的，则每个周期 100 个时间步长）。对于 SIMPLE 方法，每个时间步所需的外部迭代次数取决于时间步大小。

当迈向稳态时，仍然需要使用有限的时间步长，但时间步长可能相对较大。当解的变化相对较大时（特别是当初始化不是最终解的良好近似时），最初使用较小的时间步通常是有利的；随着接近稳态，时间步长可以逐渐增加。最大允许时间步长取决于问题，但通常与 CFL 数相关：

\begin{sourceblock}
\includegraphics[page=251,trim={53.00bp 451.77bp 53.87bp 186.08bp},clip,width=0.92\textwidth,max height=0.38\textheight]{Ferziger4th.pdf}
\end{sourceblock}
1 到 100 之间的值通常是合适的；见第 8.4 节 一些例子。

\section*{8.2 共置网格的隐式迭代方法}\phantomsection\addcontentsline{toc}{section}{8.2 共置网格的隐式迭代方法}
前面提到，数值网格上变量的共置排列会产生问题，导致它在一段时间内失宠。在这里，我们将首先说明出现问题的原因，然后提出解决方法。

\subsection*{8.2.1 共置变量的压力处理}\phantomsection\addcontentsline{toc}{subsection}{8.2.1 共置变量的压力处理}
我们首先研究有限差分格式和第 4 节中提出的简单时间提前方法。 7.1.7.1.我们推导出压力的离散泊松方程，可以写成：

\begin{sourceblock}
\begin{tikzpicture}
\node[inner sep=0,anchor=south west] (image) at (0,0) {\includegraphics[page=251,trim={53.00bp 201.49bp 53.87bp 428.55bp},clip,width=0.92\textwidth,max height=0.38\textheight]{Ferziger4th.pdf}};
\begin{scope}[x={(image.south east)},y={(image.north west)}]
\fill[white] (0.38304,0.45235) rectangle (0.40072,0.77258);
\fill[white] (0.37095,0.03324) rectangle (0.41050,0.38338);
\end{scope}
\end{tikzpicture}
\end{sourceblock}
where H n
i 是平流项和粘性项之和的简写符号：

\begin{sourceblock}
\includegraphics[page=251,trim={53.00bp 140.94bp 53.87bp 493.54bp},clip,width=0.92\textwidth,max height=0.38\textheight]{Ferziger4th.pdf}
\end{sourceblock}
（暗示对 j 求和）。用于近似导数的离散化方案在方程（8.39）中并不重要；这就是使用符号表示法的原因。此外，该方程并不特定于任何网格排列。

现在让我们看看图 8.2 中所示的并置布置以及动量方程中压力梯度项和

\begin{center}\small 图 8.2 并置网格中的控制量和使用的符号\end{center}
x

N

j y
n

西北北

y

W E
P

e

w

西南

s

j−1 y
S

x x i
i−1
连续性方程的发散。我们首先考虑压力项的前向差分格式和连续性方程的后向差分格式。第 7.1.3 节表明这种组合是节能的。为了简单起见，我们假设网格是均匀的，间距为 x 和 y。

用后向差分格式逼近压力方程中的外差分算子 δ/δ x i，可得：

\begin{sourceblock}
\begin{tikzpicture}
\node[inner sep=0,anchor=south west] (image) at (0,0) {\includegraphics[page=252,trim={53.00bp 350.66bp 53.87bp 266.29bp},clip,width=0.92\textwidth,max height=0.38\textheight]{Ferziger4th.pdf}};
\begin{scope}[x={(image.south east)},y={(image.north west)}]
\fill[white] (0.04433,0.59280) rectangle (0.09038,0.85444);
\fill[white] (0.05077,0.30697) rectangle (0.08376,0.54747);
\fill[white] (0.11185,0.26387) rectangle (0.12833,0.43810);
\fill[white] (0.92439,0.00000) rectangle (1.00000,0.06221);
\end{scope}
\end{tikzpicture}
\end{sourceblock}
(8.41) 将右侧表示为 Q H
P 并使用压力导数的前向差分近似，我们得出：

\begin{sourceblock}
\includegraphics[page=252,trim={53.00bp 267.85bp 53.87bp 355.63bp},clip,width=0.92\textwidth,max height=0.38\textheight]{Ferziger4th.pdf}
\end{sourceblock}
压力的代数方程组采用以下形式：

\begin{sourceblock}
\begin{tikzpicture}
\node[inner sep=0,anchor=south west] (image) at (0,0) {\includegraphics[page=252,trim={53.00bp 209.47bp 53.87bp 428.69bp},clip,width=0.92\textwidth,max height=0.38\textheight]{Ferziger4th.pdf}};
\begin{scope}[x={(image.south east)},y={(image.north west)}]
\fill[white] (0.20962,0.44496) rectangle (0.24623,0.95711);
\end{scope}
\end{tikzpicture}
\end{sourceblock}
其中系数为：

\begin{sourceblock}
\begin{tikzpicture}
\node[inner sep=0,anchor=south west] (image) at (0,0) {\includegraphics[page=252,trim={53.00bp 150.26bp 53.87bp 483.79bp},clip,width=0.92\textwidth,max height=0.38\textheight]{Ferziger4th.pdf}};
\begin{scope}[x={(image.south east)},y={(image.north west)}]
\fill[white] (0.05173,0.38797) rectangle (0.08833,0.83422);
\end{scope}
\end{tikzpicture}
\end{sourceblock}
如果图 8.2 中所示的 CV 用于动量方程和连续性方程，并且使用以下近似值，则可以验证 FV 方法将重现方程（8.42）：ue = u P , p e = p E； v n = v P , p n = p N ; u w = u W , p w = p P ; v s = v S， p s = p P。

压力或压力校正方程与在具有中心差近似的交错网格上获得的方程具有相同的形式；这是因为通过一阶导数的前向和后向差分近似值的乘积来近似二阶导数可以得到中心差分近似值。然而，动量方程现在由于使用主要驱动力项（压力梯度）的一阶近似而受到影响。最好使用高阶近似。

现在考虑一下，如果我们为动量方程中的压力梯度和连续性方程中的散度选择中心差分近似，会发生什么。用中心差分来近似式（8.39）中的外差分算子，可得：

\begin{sourceblock}
\begin{tikzpicture}
\node[inner sep=0,anchor=south west] (image) at (0,0) {\includegraphics[page=253,trim={53.00bp 425.68bp 53.87bp 191.16bp},clip,width=0.92\textwidth,max height=0.38\textheight]{Ferziger4th.pdf}};
\begin{scope}[x={(image.south east)},y={(image.north west)}]
\fill[white] (0.03973,0.59351) rectangle (0.08577,0.85477);
\fill[white] (0.04614,0.30852) rectangle (0.07913,0.54848);
\fill[white] (0.10722,0.26531) rectangle (0.12484,0.43915);
\fill[white] (0.92439,0.00000) rectangle (1.00000,0.05984);
\end{scope}
\end{tikzpicture}
\end{sourceblock}
(8.45) 我们再次将右侧表示为 Q H
;然而，这个数量不是之前获得的数量。插入压力导数的中心差近似，我们发现：

\begin{sourceblock}
\includegraphics[page=253,trim={53.00bp 331.46bp 53.87bp 291.91bp},clip,width=0.92\textwidth,max height=0.38\textheight]{Ferziger4th.pdf}
\end{sourceblock}
压力代数方程组的形式为：

\begin{sourceblock}
\begin{tikzpicture}
\node[inner sep=0,anchor=south west] (image) at (0,0) {\includegraphics[page=253,trim={53.00bp 273.19bp 53.87bp 364.98bp},clip,width=0.92\textwidth,max height=0.38\textheight]{Ferziger4th.pdf}};
\begin{scope}[x={(image.south east)},y={(image.north west)}]
\fill[white] (0.17212,0.44476) rectangle (0.20872,0.95710);
\end{scope}
\end{tikzpicture}
\end{sourceblock}
其中系数为：

\begin{sourceblock}
\begin{tikzpicture}
\node[inner sep=0,anchor=south west] (image) at (0,0) {\includegraphics[page=253,trim={53.00bp 213.97bp 53.87bp 420.08bp},clip,width=0.92\textwidth,max height=0.38\textheight]{Ferziger4th.pdf}};
\begin{scope}[x={(image.south east)},y={(image.north west)}]
\fill[white] (0.04427,0.38797) rectangle (0.08087,0.83453);
\fill[white] (0.92439,0.00000) rectangle (1.00000,0.09910);
\end{scope}
\end{tikzpicture}
\end{sourceblock}
（8.48） 该方程与方程（8.43）具有相同的形式，但它涉及相距 2 x 或 2 y 的节点！它是一个离散的泊松方程，网格的粗度是基本方程的两倍，但方程分为四个不相连的系统，一个 i 和 j 均为偶数，一个 i 偶数和 j 奇数，一个 i 奇数和 j 偶数，以及一个 i 和 j 均为奇数。这些系统中的每一个都给出了不同的解。对于具有均匀压力场的流动，可以产生图8.3所示的棋盘压力分布满足这些方程。然而，压力梯度不受影响并且速度场可能是平滑的。还有一种可能性是无法获得收敛的稳态解。

如果通过两个相邻节点的线性插值来计算通量的 CV 面值，则使用有限体积方法可以获得类似的结果。

\begin{center}\small 图 8.3 棋盘压力场，由 2 个间距上的四个均匀场叠加组成，CDS 将其解释为均匀场\end{center}
−1 −1
0 2
−1
−1 1
0 2
1 −1
−1
上述问题的根源可以追溯到使用一阶导数的 2 x 近似值。已经提出了多种治疗方法。在不可压缩流动中，绝对压力水平并不重要——只有差异才重要。除非在某处指定了压力的绝对值，否则压力方程是奇异的，并且具有无限多个解，所有解都相差一个常数。这使得一个简单的解决方法成为可能：过滤掉振荡，正如 van der Wijngaart (1990) 所做的那样。

我们将提出一种处理并置网格上的压力-速度耦合的方法，该方法已在复杂的几何形状中广泛使用，并且简单而有效。它几乎用于所有商业和公共 CFD 代码。

在交错网格上，中心差近似基于 x 差。压力方程（8.39）中的外一阶导数的 x 近似具有以下形式：

\begin{sourceblock}
\begin{tikzpicture}
\node[inner sep=0,anchor=south west] (image) at (0,0) {\includegraphics[page=254,trim={53.00bp 279.22bp 53.87bp 337.73bp},clip,width=0.92\textwidth,max height=0.38\textheight]{Ferziger4th.pdf}};
\begin{scope}[x={(image.south east)},y={(image.north west)}]
\fill[white] (0.05651,0.59280) rectangle (0.10256,0.85444);
\fill[white] (0.06292,0.30697) rectangle (0.09591,0.54747);
\fill[white] (0.12400,0.26367) rectangle (0.13814,0.43810);
\fill[white] (0.92439,0.00000) rectangle (1.00000,0.06221);
\end{scope}
\end{tikzpicture}
\end{sourceblock}
(8.49) 问题是压力导数和量 H 的值在单元面位置不可用，因此我们必须使用插值。让我们选择线性插值，它与导数的 CDS 近似具有相同的精度。还让式(8.39)中压力的内导数用中心差来近似。细胞中心导数的线性插值导致：

\begin{sourceblock}
\begin{tikzpicture}
\node[inner sep=0,anchor=south west] (image) at (0,0) {\includegraphics[page=254,trim={53.00bp 171.07bp 53.87bp 457.64bp},clip,width=0.92\textwidth,max height=0.38\textheight]{Ferziger4th.pdf}};
\begin{scope}[x={(image.south east)},y={(image.north west)}]
\fill[white] (0.27077,0.46460) rectangle (0.31681,0.80844);
\fill[white] (0.59765,0.44082) rectangle (0.64310,0.77344);
\fill[white] (0.64683,0.44082) rectangle (0.71176,0.78066);
\fill[white] (0.56060,0.28507) rectangle (0.58878,0.59391);
\fill[white] (0.27714,0.08923) rectangle (0.31017,0.40529);
\fill[white] (0.46490,0.08656) rectangle (0.48469,0.39540);
\fill[white] (0.50989,0.08923) rectangle (0.52803,0.39808);
\end{scope}
\end{tikzpicture}
\end{sourceblock}
通过这种插值，压力方程（8.46）被恢复。

我们可以使用中心差和 x 间距来评估单元表面的压力导数，如下所示：

\begin{sourceblock}
\begin{tikzpicture}
\node[inner sep=0,anchor=south west] (image) at (0,0) {\includegraphics[page=254,trim={53.00bp 88.24bp 53.87bp 540.48bp},clip,width=0.92\textwidth,max height=0.38\textheight]{Ferziger4th.pdf}};
\begin{scope}[x={(image.south east)},y={(image.north west)}]
\fill[white] (0.39044,0.46472) rectangle (0.43648,0.80866);
\end{scope}
\end{tikzpicture}
\end{sourceblock}
如果将此近似应用于所有单元面，我们将得到以下压力方程（这在非均匀网格上也有效）：

\begin{sourceblock}
\includegraphics[page=255,trim={53.00bp 533.86bp 53.87bp 89.61bp},clip,width=0.92\textwidth,max height=0.38\textheight]{Ferziger4th.pdf}
\end{sourceblock}
与式（8.42）相同，只是右边现在是通过插值获得的：

\begin{sourceblock}
\begin{tikzpicture}
\node[inner sep=0,anchor=south west] (image) at (0,0) {\includegraphics[page=255,trim={53.00bp 472.46bp 53.87bp 163.55bp},clip,width=0.92\textwidth,max height=0.38\textheight]{Ferziger4th.pdf}};
\begin{scope}[x={(image.south east)},y={(image.north west)}]
\fill[white] (0.00000,0.96216) rectangle (0.16556,1.00000);
\end{scope}
\end{tikzpicture}
\end{sourceblock}
使用这种近似消除了压力场中的振荡，但是为了实现这一点，我们在动量和压力方程中的压力梯度处理中引入了不一致。让我们比较一下这两个近似值。很容易证明方程（8.52）和（8.46）左边的区别在于：

\begin{sourceblock}
\begin{tikzpicture}
\node[inner sep=0,anchor=south west] (image) at (0,0) {\includegraphics[page=255,trim={53.00bp 367.21bp 53.87bp 270.54bp},clip,width=0.92\textwidth,max height=0.38\textheight]{Ferziger4th.pdf}};
\begin{scope}[x={(image.south east)},y={(image.north west)}]
\fill[white] (0.92433,0.00000) rectangle (1.00000,0.10356);
\end{scope}
\end{tikzpicture}
\end{sourceblock}
(8.54) 表示四阶压力导数的中心差分近似：

\begin{sourceblock}
\begin{tikzpicture}
\node[inner sep=0,anchor=south west] (image) at (0,0) {\includegraphics[page=255,trim={53.00bp 308.17bp 53.87bp 320.56bp},clip,width=0.92\textwidth,max height=0.38\textheight]{Ferziger4th.pdf}};
\begin{scope}[x={(image.south east)},y={(image.north west)}]
\fill[white] (0.00000,0.70890) rectangle (0.14117,1.00000);
\fill[white] (0.24767,0.43170) rectangle (0.26298,0.66105);
\fill[white] (0.36340,0.46458) rectangle (0.40580,0.80674);
\fill[white] (0.22550,0.27800) rectangle (0.24863,0.58701);
\fill[white] (0.24409,0.21999) rectangle (0.26054,0.44908);
\fill[white] (0.26866,0.28522) rectangle (0.29684,0.59423);
\fill[white] (0.30036,0.28522) rectangle (0.34331,0.78081);
\fill[white] (0.66854,0.08875) rectangle (0.71552,0.41139);
\end{scope}
\end{tikzpicture}
\end{sourceblock}
通过应用二阶导数的标准 CDS 近似两次可以轻松获得表达式 (8.54)，请参见第 1 节。 3.4.

随着网格的细化，这种差异趋于零，并且引入的误差与基本离散化中的误差具有相同的量级，因此不会显著增加后者。然而，该方案的节能特性在此过程中被破坏。事实上，Ham 和 Iaccarino (2004) 已经表明（在非迭代（分数步）方法的背景下）对动能的净效应是耗散的。

上述结果是通过二阶 CDS 离散化和线性插值得出的。可以为任何离散化方案和插值构建类似的推导。让我们看看如何将上述想法转化为使用 FV 离散化的隐式压力校正方法。

\subsection*{8.2.2 共置网格的简单}\phantomsection\addcontentsline{toc}{subsection}{8.2.2 共置网格的简单}
请记住，所有变量的 CV 现在都是相同的。单元面中心处的压力（不是节点位置）必须通过插值获得；线性插值是合适的二阶近似，但也可以使用更高阶的方法。计算细胞面速度所需的 CV 中心的梯度可以使用高斯定理获得。将所有面上的 x 和 y 方向的压力相加，然后除以单元体积，得到相应的平均压力导数，例如：

\begin{sourceblock}
\begin{tikzpicture}
\node[inner sep=0,anchor=south west] (image) at (0,0) {\includegraphics[page=256,trim={53.00bp 483.84bp 53.87bp 144.89bp},clip,width=0.92\textwidth,max height=0.38\textheight]{Ferziger4th.pdf}};
\begin{scope}[x={(image.south east)},y={(image.north west)}]
\fill[white] (0.42180,0.46458) rectangle (0.45642,0.78081);
\end{scope}
\end{tikzpicture}
\end{sourceblock}
where Q p
u 表示所有 CV 面上 x 方向的压力之和，见式（8.12）。在笛卡尔网格上，这可简化为标准 CDS 近似。

线性动量方程的解产生 u ＊ 和 v ＊。对于离散连续性方程，我们需要单元面速度，该速度必须通过插值计算；线性插值是显而易见的选择。 SIMPLE算法的压力修正方程可以按照节的思路推导出来。 7.2.2 和 8.1。连续性方程中所需的插值单元面速度涉及插值压力梯度，因此它们的校正与插值压力校正梯度成正比（见方程（8.29））：

\begin{sourceblock}
\begin{tikzpicture}
\node[inner sep=0,anchor=south west] (image) at (0,0) {\includegraphics[page=256,trim={53.00bp 324.93bp 53.87bp 303.38bp},clip,width=0.92\textwidth,max height=0.38\textheight]{Ferziger4th.pdf}};
\begin{scope}[x={(image.south east)},y={(image.north west)}]
\fill[white] (0.36851,0.28549) rectangle (0.39555,0.64393);
\end{scope}
\end{tikzpicture}
\end{sourceblock}
在均匀网格上，使用该单元面速度校正表达式导出的压力校正方程对应于方程（8.46）。在非均匀网格上，压力校正方程的计算分子涉及节点 P、E、W、N、S、EE、WW、NN 和 SS。如上一节所示，该方程可能有振荡解。尽管可以滤除振荡（参见 van der Wijngaart 1990），但压力校正方程在任意网格上变得复杂，并且求解算法的收敛可能很慢。使用上一节中讨论的方法可以获得类似于交错网格方程的紧凑压力校正方程。下面进行描述。

上一节表明，插值压力梯度可以用单元表面的紧凑中心差近似代替。因此，插值单元面速度由插值压力梯度与单元表面计算的梯度之间的差异进行修改（参见方程（8.27）和（8.28））：

\begin{sourceblock}
\begin{tikzpicture}
\node[inner sep=0,anchor=south west] (image) at (0,0) {\includegraphics[page=256,trim={53.00bp 115.04bp 53.87bp 507.69bp},clip,width=0.92\textwidth,max height=0.38\textheight]{Ferziger4th.pdf}};
\begin{scope}[x={(image.south east)},y={(image.north west)}]
\fill[white] (0.00000,0.81594) rectangle (0.04740,1.00000);
\fill[white] (0.05008,0.81594) rectangle (0.14559,1.00000);
\fill[white] (0.18920,0.30845) rectangle (0.22150,0.62082);
\fill[white] (0.20577,0.26077) rectangle (0.21991,0.45819);
\fill[white] (0.22650,0.31467) rectangle (0.25468,0.58097);
\fill[white] (0.25820,0.27459) rectangle (0.32367,0.59226);
\fill[white] (0.32722,0.31467) rectangle (0.35540,0.58097);
\fill[white] (0.38220,0.28818) rectangle (0.41377,0.57498);
\end{scope}
\end{tikzpicture}
\end{sourceblock}
这种对插值单元面速度的校正被称为 Rhie-Chow 校正（Rhie 和 Chow 1983）。上横线表示插值，以单元面为中心的体积由下式定义：

\begin{sourceblock}
\includegraphics[page=257,trim={53.00bp 593.89bp 53.87bp 56.79bp},clip,width=0.92\textwidth,max height=0.38\textheight]{Ferziger4th.pdf}
\end{sourceblock}
对于笛卡尔网格。

此过程添加了对插值速度的校正，该校正与压力乘以 ( x ) 2 / 4 的三阶导数成正比；单元中心的四阶导数是通过应用散度算子得出的。在二阶方案中，单元表面的压力导数使用 CDS 计算，见方程（8.51）。如果在非均匀网格上使用 CDS 近似（8.51），则单元中心压力梯度应使用权重 1/2 进行插值，因为这种近似不会“看到”网格的非均匀性。

如果压力波动较快，则修正量较大；那么三阶导数就很大并且将激活压力校正并平滑压力。

SIMPLE 方法中对单元面速度的修正现在为：

\begin{sourceblock}
\begin{tikzpicture}
\node[inner sep=0,anchor=south west] (image) at (0,0) {\includegraphics[page=257,trim={53.00bp 410.21bp 53.87bp 218.12bp},clip,width=0.92\textwidth,max height=0.38\textheight]{Ferziger4th.pdf}};
\begin{scope}[x={(image.south east)},y={(image.north west)}]
\fill[white] (0.34307,0.46760) rectangle (0.36286,0.77334);
\fill[white] (0.16463,0.28564) rectangle (0.19167,0.64401);
\fill[white] (0.18120,0.23063) rectangle (0.19531,0.45755);
\fill[white] (0.20054,0.29252) rectangle (0.22872,0.59852);
\fill[white] (0.23224,0.29252) rectangle (0.26042,0.59852);
\fill[white] (0.28054,0.26210) rectangle (0.31212,0.59138);
\end{scope}
\end{tikzpicture}
\end{sourceblock}
与其他细胞表面的相应表达。当将它们插入离散连续性方程时，结果又是压力校正方程（8.32）。唯一的区别是系数 1 / A u
P and 1 / A v
单元面上的 P 不是交错排列中的节点值，而是插值单元中心值。

因为式(8.58)中的修正项乘以1 / A u
P，其中包含的欠松弛参数的值可能会影响收敛的细胞面速度。然而，没有什么理由担心，因为使用不同欠松弛参数获得的两个解的差异远小于离散化误差，如下面的示例所示。我们还表明，使用共置网格的隐式算法与交错网格算法具有相同的收敛速度、对欠松弛​​因素的依赖性以及计算成本。此外，不同变量排列得到的解之间的差异也远小于离散化误差。

我们推导了二阶近似的同位网格上的压力校正方程。该方法可以适应更高阶的近似；微分和插值的阶数相同很重要。有关四阶方法的描述，请参见 Lilek 和 Peri ́c (1995)。

\subsection*{8.2.3 共置电网的 IFSM}\phantomsection\addcontentsline{toc}{subsection}{8.2.3 共置电网的 IFSM}
前面描述的用于交错网格的 IFSM 方法很容易扩展到并置网格。实际上，两种情况下的压力校正方程是相同的，参见方程（8.32）和（8.37）。唯一的区别在于计算质量通量所需的细胞面速度的计算。在交错网格的情况下，速度存储在连续性 CV 的面上，因此不需要插值。在共置网格上，必须使用插值根据节点值计算面速度。当使用 IFSM 计算非定常流时，无需 Rhie-Chow 校正的线性插值通常效果很好。然而，在节中提出的应用程序中。 8.4 我们对单元表面的线性插值速度应用了以下校正：

\begin{sourceblock}
\begin{tikzpicture}
\node[inner sep=0,anchor=south west] (image) at (0,0) {\includegraphics[page=258,trim={53.00bp 514.74bp 53.87bp 107.99bp},clip,width=0.92\textwidth,max height=0.38\textheight]{Ferziger4th.pdf}};
\begin{scope}[x={(image.south east)},y={(image.north west)}]
\fill[white] (0.00000,0.77517) rectangle (0.02577,1.00000);
\fill[white] (0.02845,0.77517) rectangle (0.07651,1.00000);
\fill[white] (0.07922,0.77517) rectangle (0.15359,1.00000);
\fill[white] (0.24177,0.30868) rectangle (0.27408,0.62082);
\end{scope}
\end{tikzpicture}
\end{sourceblock}
该表达式与 SIMPLE 中使用的表达式相同，除了方括号中的项的乘数不同之外，请参见。方程(8.58)和(8.60)。按照与 SIMPLE 中相同的方法，为 IFSM 获得速度和压力校正之间的以下关系：

\begin{sourceblock}
\begin{tikzpicture}
\node[inner sep=0,anchor=south west] (image) at (0,0) {\includegraphics[page=258,trim={53.00bp 419.96bp 53.87bp 208.76bp},clip,width=0.92\textwidth,max height=0.38\textheight]{Ferziger4th.pdf}};
\begin{scope}[x={(image.south east)},y={(image.north west)}]
\fill[white] (0.20704,0.27819) rectangle (0.23408,0.64030);
\end{scope}
\end{tikzpicture}
\end{sourceblock}
v * 的类似表达式如下
n and v ′
名词互联网上提供了包含该算法的计算机代码；详细内容请参见附录。

\section*{8.3 非定常流的非迭代隐式方法}\phantomsection\addcontentsline{toc}{section}{8.3 非定常流的非迭代隐式方法}
与完全显式方法相比，非迭代隐式方法允许使用更大的时间步长。由于在不可压缩流动的情况下，无论如何都必须求解用于压力或压力校正的泊松型方程，因此通常优选隐式方法。当局部创建非常精细的网格（靠近壁、锐边周围的边界层等）时，这一点尤其重要，因为如果方法完全显式，则扩散项会对时间步长大小产生过于严格的稳定性限制。

术语“非迭代”在这里意味着缺少外部迭代循环：每个时间步仅求解动量和压力校正方程一次。分数步法 (FSM) 的许多版本都属于这一类；从这个意义上说，PISO 可能被认为是非迭代的，也可能不是非迭代的（它只求解一次动量方程，但外部迭代循环包含压力校正方程和显式速度校正）。

前面提出的 IFSM 和 FSM 的非迭代版本之间的主要区别是：（i）后者对全部或部分对流通量使用显式时间超前方案，以及（ii）线性方程的求解比迭代方法具有更严格的容差，以确保迭代误差足够小，并且充分满足动量方程和连续性方程。非迭代方法会引入分裂误差，因为它们仅修正非稳态项中的速度；该误差通常与 (t) 2 成正比。 SIMPLE 和 IFSM 都通过外部迭代消除了分裂误差。通常，非迭代 FSM 中的一个时间步或
PISO 的成本（就计算量而言）相当于 SIMPLE 或 IFSM 中 2-3 次外部迭代。当我们在第三节中介绍示例计算时，我们将再次指出这一事实。 8.4.

前面关于 SIMPLE 和相关方法的部分中描述的许多工具和算法都可以在分数步方法中使用。因此，本节不像前面的那么详细，而是重点关注非迭代分数步方法的一些典型项目，以突出特定步骤。接下来，我们回顾该方法的主要步骤，从 Adams-Bashforth 离散化开始。接下来，我们按顺序处理：近似因式分解 ADI、压力泊松方程以及初始条件和边界条件。最后，我们评论单步和迭代分数步方法之间的效率和准确性的差异。

如前几节所示，同时考虑了并置网格和交错网格。然而，在手边没有图片的情况下，很容易忘记变量的放置位置以及并置和交错网格的适当控制量是什么。例如，假设速度分量位于交错网格的单元面上，人们可能会天真地想知道为什么需要在交错网格的对流项中使用 QUICK。因此，我们在图 8.4 中重现了网格和相关控制体积。对于并置网格，我们现在回想一下

\begin{sourceblock}
\includegraphics[page=259,trim={53.00bp 115.02bp 53.87bp 287.32bp},clip,width=\textwidth,max height=0.38\textheight]{Ferziger4th.pdf}
\par\vspace{0.45\baselineskip}
\small 图 8.4 控制体积和符号：对于并置网格上的所有变量（上）、标量变量（左下）、交错网格上的 x 动量（中下）和 y 动量（右下），只有一个速度和压力控制体积，并且所有变量都在该网格上的 P 点定义，以用于我们在此使用的 FV 方法。因此，需要插值来找到 CV 面上的变量值，因为在 FV 方法中，在每种情况下对通过 CV 面的通量进行求和。在连续性方程中，我们需要 CV 面上的速度分量来计算质量通量；在所有其他传输方程中，我们需要对流质量通量和对流变量（速度分量、温度等）来计算对流通量。还需要对压力进行插值以获得每个单元表面的压力。
\end{sourceblock}
另一方面，在交错网格上，每个速度分量以及连续性方程和标量变量（例如，温度、压力和流体属性也存储在该 CV 中心的节点处）都有单独的控制体积。在图 8.4 的下半部分，箭头所在的位置速度已知，而点所在的位置压力已知。对于连续性方程，不需要插值：因为速度存储在 CV 面上，所以可以直接计算质量通量（使用中点规则在二阶近似内）。然而，为了计算所有其他传输方程中的对流通量，需要插值：对于标量变量（如温度），我们需要对存储在单元中心的变量进行插值以获得面质心处的对流值，而对于速度 CV，需要插值来计算单元面处的对流（质量通量）和对流速度。例如，为了获得 x 动量方程中“e”点的对流通量，我们需要对速度进行插值，而东面和西面的相关压力可以直接从存储在那里的压力值计算出来（在笛卡尔网格上，x 动量 CV 的“n”和“s”面上的压力不相关，因为它们在 x 方向上没有分量；因此，y 动量的“e”和“w”面上的压力CV 不相关，而“n”和“s”面上的力可以直接计算）。

请注意，对于除中点规则之外的积分近似，在任何情况下，对于任何网格类型和变量排列，都需要将节点变量值插值到单元面内的积分点。

\subsection*{8.3.1 Adams-Bashforth 对流项的空间离散化}\phantomsection\addcontentsline{toc}{subsection}{8.3.1 Adams-Bashforth 对流项的空间离散化}
方程（7.50）中使用的 Adams-Bashforth 离散化是典型的分数阶方法，并且至关重要，因为它使方程的该部分变得明确，同时保持时间上的二阶精度。共置网格方案和交错网格方案都需要插值来定义对流项。例如，使用对流项的保守形式需要控制体积每个面上的动量通量。对于均匀笛卡尔网格，图 8.4 中 x 动量 CV 的 x 分量为 (ρ uu ) e。由式（4.24）得出，假设u e 为正：

\begin{sourceblock}
\begin{tikzpicture}
\node[inner sep=0,anchor=south west] (image) at (0,0) {\includegraphics[page=261,trim={53.00bp 555.56bp 53.87bp 50.00bp},clip,width=0.92\textwidth,max height=0.38\textheight]{Ferziger4th.pdf}};
\begin{scope}[x={(image.south east)},y={(image.north west)}]
\fill[white] (0.11708,0.58419) rectangle (0.20478,0.79416);
\fill[white] (0.21005,0.60317) rectangle (0.23823,0.79416);
\fill[white] (0.24174,0.58419) rectangle (0.43074,0.79416);
\fill[white] (0.43639,0.60317) rectangle (0.46457,0.79416);
\fill[white] (0.46809,0.58419) rectangle (0.53925,0.79416);
\fill[white] (0.37212,0.23886) rectangle (0.44677,0.44883);
\fill[white] (0.45044,0.25801) rectangle (0.47862,0.44883);
\fill[white] (0.48048,0.23886) rectangle (0.55398,0.44883);
\fill[white] (0.45398,0.01981) rectangle (0.47377,0.21063);
\end{scope}
\end{tikzpicture}
\end{sourceblock}
这个结果对于交错网格或共置网格都有效，可以与隐式迭代方法中的方程（8.7）进行比较；请注意，在这两种情况下，QUICK 插值均应用于对流量，而不是应用于定义质量通量的对流速度。后者通常在需要时进行线性插值（在共置网格上和交错网格上的动量 CV）。

除节点之外的其他点所需的变量值通常通过线性插值获得。然而，其他插值有时也是有用的，例如，质量通量中的面速度可以通过压力泊松方程的形成中的高阶插值和/或压力加权项来近似（参见例如Rhie和Chow 1983；Armfield 1991；Zang等人1994；Ye等人1999）。同样，QUICK 项可以被替代的离散化所取代，例如第 1 节中给出的离散化。 4.4，例如 UDS、CDS、CUI 或 TVD。

\subsection*{8.3.2 交替方向隐式方案}\phantomsection\addcontentsline{toc}{subsection}{8.3.2 交替方向隐式方案}
ADI 方案通常称为近似因式分解 ADI 方案，可以按照第 1 节中的概述构建。 5.3.5.我们将方程（7.53）（对于笛卡尔坐标）变换为

\begin{sourceblock}
\begin{tikzpicture}
\node[inner sep=0,anchor=south west] (image) at (0,0) {\includegraphics[page=261,trim={53.00bp 270.03bp 53.87bp 360.72bp},clip,width=0.92\textwidth,max height=0.38\textheight]{Ferziger4th.pdf}};
\begin{scope}[x={(image.south east)},y={(image.north west)}]
\fill[white] (0.30929,0.43402) rectangle (0.32241,0.76067);
\fill[white] (0.22677,0.24414) rectangle (0.24791,0.57078);
\fill[white] (0.25158,0.24414) rectangle (0.27976,0.57078);
\end{scope}
\end{tikzpicture}
\end{sourceblock}
其中 AP 包括先前时间步骤中已知的平流和压力项。现在，如果我们假设密度和粘度恒定，我们可以将该方程重新排列为以下形式：

\begin{sourceblock}
\includegraphics[page=261,trim={53.00bp 188.35bp 53.87bp 449.40bp},clip,width=0.92\textwidth,max height=0.38\textheight]{Ferziger4th.pdf}
\end{sourceblock}
where

\begin{sourceblock}
\begin{tikzpicture}
\node[inner sep=0,anchor=south west] (image) at (0,0) {\includegraphics[page=261,trim={53.00bp 141.23bp 53.87bp 493.15bp},clip,width=0.92\textwidth,max height=0.38\textheight]{Ferziger4th.pdf}};
\begin{scope}[x={(image.south east)},y={(image.north west)}]
\fill[white] (0.00000,0.84666) rectangle (0.07732,1.00000);
\end{scope}
\end{tikzpicture}
\end{sourceblock}
因式分解( 8.64 ) 得出：

\begin{sourceblock}
\includegraphics[page=261,trim={53.00bp 70.72bp 53.87bp 550.62bp},clip,width=0.92\textwidth,max height=0.38\textheight]{Ferziger4th.pdf}
\end{sourceblock}
\begin{sourceblock}
\begin{tikzpicture}
\node[inner sep=0,anchor=south west] (image) at (0,0) {\includegraphics[page=262,trim={53.00bp 595.02bp 53.87bp 53.17bp},clip,width=0.92\textwidth,max height=0.38\textheight]{Ferziger4th.pdf}};
\begin{scope}[x={(image.south east)},y={(image.north west)}]
\fill[white] (0.29750,0.00000) rectangle (0.31943,0.40111);
\fill[white] (0.33251,0.06685) rectangle (0.36394,0.71086);
\fill[white] (0.37002,0.06685) rectangle (0.41143,0.71086);
\fill[white] (0.41750,0.06685) rectangle (0.46394,0.71086);
\fill[white] (0.46995,0.06685) rectangle (0.52968,0.71086);
\fill[white] (0.53576,0.06685) rectangle (0.56057,0.71086);
\fill[white] (0.56659,0.06685) rectangle (0.71952,0.71086);
\fill[white] (0.72562,0.06685) rectangle (0.75374,0.71086);
\fill[white] (0.75976,0.08691) rectangle (0.77576,0.73148);
\fill[white] (0.79504,0.07242) rectangle (0.83257,0.78552);
\fill[white] (0.84009,0.06685) rectangle (0.87982,0.71086);
\fill[white] (0.88586,0.06685) rectangle (0.95564,0.71086);
\fill[white] (0.98574,0.07242) rectangle (0.99889,0.71643);
\fill[white] (0.00000,0.00000) rectangle (0.04740,0.04457);
\fill[white] (0.04977,0.00000) rectangle (0.09615,0.04457);
\fill[white] (0.09853,0.00000) rectangle (0.13161,0.04457);
\fill[white] (0.13395,0.00000) rectangle (0.26054,0.04457);
\fill[white] (0.26295,0.00000) rectangle (0.37591,0.04457);
\fill[white] (0.37832,0.00000) rectangle (0.45549,0.04457);
\fill[white] (0.45786,0.00000) rectangle (0.55414,0.04457);
\fill[white] (0.55654,0.00000) rectangle (0.59795,0.04457);
\fill[white] (0.60030,0.00000) rectangle (0.70168,0.04457);
\fill[white] (0.70403,0.00000) rectangle (0.79865,0.04457);
\fill[white] (0.80105,0.00000) rectangle (0.82917,0.04457);
\fill[white] (0.83149,0.00000) rectangle (0.87293,0.04457);
\fill[white] (0.87528,0.00000) rectangle (1.00000,0.04457);
\end{scope}
\end{tikzpicture}
\end{sourceblock}
∂ t 因此，对于较小的 t，最后一项与 ( t ) 3 成正比，可以忽略不计。方程(8.65)将求解过程简化为三对角矩阵在交替方向上的顺序求逆。因此，使用速度的分量符号（参见方程（5.49）和（5.50））：

\begin{sourceblock}
\begin{tikzpicture}
\node[inner sep=0,anchor=south west] (image) at (0,0) {\includegraphics[page=262,trim={53.00bp 448.64bp 53.87bp 112.87bp},clip,width=0.92\textwidth,max height=0.38\textheight]{Ferziger4th.pdf}};
\begin{scope}[x={(image.south east)},y={(image.north west)}]
\fill[white] (0.17296,0.79843) rectangle (0.19275,0.90892);
\fill[white] (0.19462,0.80197) rectangle (0.25362,0.97926);
\fill[white] (0.27287,0.86620) rectangle (0.28602,0.97668);
\fill[white] (0.29949,0.86868) rectangle (0.32884,0.98853);
\fill[white] (0.23946,0.73086) rectangle (0.27573,0.84488);
\fill[white] (0.29110,0.73182) rectangle (0.33723,0.84708);
\fill[white] (0.34770,0.50655) rectangle (0.36084,0.61703);
\fill[white] (0.37371,0.50903) rectangle (0.40310,0.62879);
\fill[white] (0.24782,0.43878) rectangle (0.26761,0.54936);
\fill[white] (0.26944,0.44232) rectangle (0.32842,0.61961);
\fill[white] (0.31426,0.37121) rectangle (0.35053,0.48523);
\fill[white] (0.36592,0.37217) rectangle (0.41086,0.48743);
\fill[white] (0.25432,0.07914) rectangle (0.27411,0.18962);
\fill[white] (0.27597,0.08258) rectangle (0.33495,0.25987);
\fill[white] (0.35423,0.14680) rectangle (0.36737,0.25729);
\fill[white] (0.37925,0.14938) rectangle (0.40863,0.26914);
\fill[white] (0.32078,0.01147) rectangle (0.35705,0.12549);
\fill[white] (0.37242,0.01242) rectangle (0.41546,0.12769);
\end{scope}
\end{tikzpicture}
\end{sourceblock}
每个方程组都有一个三对角系数矩阵，因此可以用高效的 TDMA 方法求解；这不需要迭代。然而，这种求解方法仅适用于结构化网格（笛卡尔网格和非正交网格）；它无法如此简单和高效地应用于非结构化网格。

\subsection*{8.3.3 压力的泊松方程}\phantomsection\addcontentsline{toc}{subsection}{8.3.3 压力的泊松方程}
压力的泊松方程可以通过任何椭圆方程求解器求解，例如多重网格、ADI、共轭梯度、GMRES 等（参见第 5 章）。在这里，我们解决与泊松方程离散化相关的一些问题。

1. 对于非迭代分步方法，Hirt 和 Harlow (1967) 认为，鉴于泊松方程在每一步都迭代求解以强制连续性，因此应注意避免不可压缩误差的累积。他们的见解是，人们可以（1）将方程的迭代解提高到高精度，或者（2）使用一些自校正程序。有趣的是，对于方程（7.44），值得注意的是，压力泊松方程的公式首先包含当前（n）和未来（n+1）时间步长的速度散度；我们可以证明节中列出的公式。 7.2.1 实际上包括解释发散 v n 的能力。我们首先将不带粘性项的方程（7.53）写成（它们在此图中不起任何作用）：

\begin{sourceblock}
\includegraphics[page=262,trim={53.00bp 119.09bp 53.87bp 511.66bp},clip,width=0.92\textwidth,max height=0.38\textheight]{Ferziger4th.pdf}
\end{sourceblock}
并将其引入压力泊松方程（7.57）中，导出压力泊松方程（7.57），迫使v n + 1 的散度为零。结果是

\begin{sourceblock}
\begin{tikzpicture}
\node[inner sep=0,anchor=south west] (image) at (0,0) {\includegraphics[page=263,trim={53.00bp 531.22bp 53.87bp 52.48bp},clip,width=0.92\textwidth,max height=0.38\textheight]{Ferziger4th.pdf}};
\begin{scope}[x={(image.south east)},y={(image.north west)}]
\fill[white] (0.25768,0.43122) rectangle (0.31645,0.57472);
\fill[white] (0.20698,0.34546) rectangle (0.22012,0.48569);
\end{scope}
\end{tikzpicture}
\end{sourceblock}
因此，上一步中的散度中的任何错误都会被反馈作为下一步的修正。 Hirt 和 Harlow（1967）表明，这可以减少不可压缩误差的累积，例如通过以较少的迭代次数终止迭代来节省计算时间，从而降低精度。 2. 交错网格的设置和策略与第 2 节中描述的基本相同。 8.1 所以我们这里不关注交错网格；相反，我们把时间花在同地网格上。在第 4 节中，针对 SIMPLE 和相关方案讨论了与并置网格上的压力计算相关的问题和解。 8.2.1.论证了稀疏拉普拉斯形式的棋盘压力模式，并导出了替代的紧凑拉普拉斯形式。然而，那里和宗门之中。 7.1.3和7.1.5.1，注意到使用紧凑形式时能量守恒的损失。我们在这里描述了两种不会经历压力场解耦的替代公式；一种使用紧致拉普拉斯算子及其伴随的连续性误差，另一种使用稀疏拉普拉斯算子的修改方法。 3. 在同位网格上使用第 2 节中描述的方程。 7.2.1，压力泊松方程（7.57）可以用两种方式离散化。 1 我们首先直接在拉普拉斯算子上使用CDS对方程进行二维离散；得出（参见方程（8.52））：

\begin{sourceblock}
\begin{tikzpicture}
\node[inner sep=0,anchor=south west] (image) at (0,0) {\includegraphics[page=263,trim={53.00bp 219.28bp 53.87bp 383.11bp},clip,width=0.92\textwidth,max height=0.38\textheight]{Ferziger4th.pdf}};
\begin{scope}[x={(image.south east)},y={(image.north west)}]
\fill[white] (0.27835,0.68565) rectangle (0.30472,0.89129);
\fill[white] (0.29335,0.65161) rectangle (0.31098,0.78604);
\fill[white] (0.31468,0.68408) rectangle (0.38968,0.89129);
\end{scope}
\end{tikzpicture}
\end{sourceblock}
另一方面，通过使用离散形式的散度和梯度算子构建方程，如第 4 节所示。 7.1.5.1 产量（2×等网格点见图3.5，即EE、WW等）

\begin{sourceblock}
\begin{tikzpicture}
\node[inner sep=0,anchor=south west] (image) at (0,0) {\includegraphics[page=263,trim={53.00bp 108.33bp 53.87bp 494.05bp},clip,width=0.92\textwidth,max height=0.38\textheight]{Ferziger4th.pdf}};
\begin{scope}[x={(image.south east)},y={(image.north west)}]
\fill[white] (0.24866,0.68585) rectangle (0.27504,0.89131);
\fill[white] (0.26367,0.65166) rectangle (0.29411,0.78607);
\fill[white] (0.29777,0.68428) rectangle (0.37278,0.89131);
\end{scope}
\end{tikzpicture}
\end{sourceblock}
使用方程（8.71）的紧凑拉普拉斯算子实际上产生的连续性误差等于

\begin{sourceblock}
\begin{tikzpicture}
\node[inner sep=0,anchor=south west] (image) at (0,0) {\includegraphics[page=264,trim={53.00bp 528.90bp 53.87bp 76.79bp},clip,width=0.92\textwidth,max height=0.38\textheight]{Ferziger4th.pdf}};
\begin{scope}[x={(image.south east)},y={(image.north west)}]
\fill[white] (0.03829,0.91580) rectangle (0.10797,1.00000);
\fill[white] (0.11068,0.91580) rectangle (0.13880,1.00000);
\fill[white] (0.28589,0.61787) rectangle (0.31408,0.80926);
\fill[white] (0.33335,0.61356) rectangle (0.34650,0.80480);
\fill[white] (0.28589,0.14127) rectangle (0.32523,0.33251);
\fill[white] (0.34451,0.13681) rectangle (0.38205,0.34855);
\end{scope}
\end{tikzpicture}
\end{sourceblock}
因为 p ′ ∼ t (∂ p /∂ t )。在节中获得了相同的结果。 8.2.1.相反，对于交错网格，使用离散散度和梯度形式的 FV 公式给出了紧凑的拉普拉斯算子并且没有连续性误差。方程（8.71）在Armfield（2000）的非迭代模拟中表现良好； Armfield 和 Street (2005) 以及 Armfield 等人。 （2010）。在那里，在每个时间步长，使用当前的 p ' 值，校正 CV 速度，例如，

\begin{sourceblock}
\begin{tikzpicture}
\node[inner sep=0,anchor=south west] (image) at (0,0) {\includegraphics[page=264,trim={53.00bp 410.60bp 53.87bp 227.16bp},clip,width=0.92\textwidth,max height=0.38\textheight]{Ferziger4th.pdf}};
\begin{scope}[x={(image.south east)},y={(image.north west)}]
\fill[white] (0.03829,0.88020) rectangle (0.07802,1.00000);
\fill[white] (0.08069,0.88020) rectangle (0.19398,1.00000);
\end{scope}
\end{tikzpicture}
\end{sourceblock}
并且每个 CV 面速度也根据例如以下公式进行校正：

\begin{sourceblock}
\includegraphics[page=264,trim={53.00bp 353.52bp 53.87bp 284.33bp},clip,width=0.92\textwidth,max height=0.38\textheight]{Ferziger4th.pdf}
\end{sourceblock}
这会产生这些对流速度的无散度速度场。据观察，压力中的高阶误差似乎限制了压力场中网格尺度误差的增长。此外，将连续性误差反馈到下一个时间步（方程（8.70））减少了连续性误差的累积。 4. 使用稀疏形式（8.72）效率低下，并且会导致迭代共置方案中的压力振荡（参见第 8.2.1 节）。阿姆菲尔德等人。 (2010) 对非迭代情况的稀疏方案进行了改进（参见 Choi 和 Moin 1994 的交错网格分数步法）。修改后的方案如下（原始分步过程在第 7.2.1 节中描述）： 求解方程（7.50）和（7.51）：

a.使用旧压力值求出新速度 v * 的估计值

\begin{sourceblock}
\begin{tikzpicture}
\node[inner sep=0,anchor=south west] (image) at (0,0) {\includegraphics[page=264,trim={53.00bp 142.41bp 53.87bp 456.34bp},clip,width=0.92\textwidth,max height=0.38\textheight]{Ferziger4th.pdf}};
\begin{scope}[x={(image.south east)},y={(image.north west)}]
\fill[white] (0.06544,0.96439) rectangle (0.09104,1.00000);
\fill[white] (0.10120,0.96439) rectangle (0.19089,1.00000);
\fill[white] (0.19203,0.96439) rectangle (0.22511,1.00000);
\fill[white] (0.22623,0.96439) rectangle (0.33086,1.00000);
\fill[white] (0.33194,0.96439) rectangle (0.36168,1.00000);
\fill[white] (0.36277,0.96439) rectangle (0.40418,1.00000);
\fill[white] (0.40526,0.96439) rectangle (0.45922,1.00000);
\fill[white] (0.46033,0.96439) rectangle (0.56117,1.00000);
\fill[white] (0.56232,0.96439) rectangle (0.66580,1.00000);
\fill[white] (0.66686,0.96439) rectangle (0.70830,1.00000);
\fill[white] (0.70938,0.96439) rectangle (0.75248,1.00000);
\fill[white] (0.75353,0.96439) rectangle (0.82247,1.00000);
\fill[white] (0.82361,0.96439) rectangle (0.85338,1.00000);
\fill[white] (0.85447,0.96439) rectangle (0.89591,1.00000);
\fill[white] (0.89696,0.96439) rectangle (1.00000,1.00000);
\fill[white] (0.46839,0.53331) rectangle (0.49657,0.70500);
\end{scope}
\end{tikzpicture}
\end{sourceblock}
b.将旧的压力梯度添加到估计的速度场中

\begin{sourceblock}
\includegraphics[page=264,trim={53.00bp 98.95bp 53.87bp 551.13bp},clip,width=0.92\textwidth,max height=0.38\textheight]{Ferziger4th.pdf}
\end{sourceblock}
因此近似抵消了前面方程中的压力梯度。在没有隐含粘性项的情况下，抵消将是精确的。 c.定义对以下形式的 \textasciicircum{} v ∗ 的修正

\begin{sourceblock}
\includegraphics[page=265,trim={53.00bp 535.24bp 53.87bp 112.46bp},clip,width=0.92\textwidth,max height=0.38\textheight]{Ferziger4th.pdf}
\end{sourceblock}
最后是 d。将方程(8.78)代入连续性方程(7.51)求出p n + 1 / 2

\begin{sourceblock}
\begin{tikzpicture}
\node[inner sep=0,anchor=south west] (image) at (0,0) {\includegraphics[page=265,trim={53.00bp 464.18bp 53.87bp 171.04bp},clip,width=0.92\textwidth,max height=0.38\textheight]{Ferziger4th.pdf}};
\begin{scope}[x={(image.south east)},y={(image.north west)}]
\fill[white] (0.10117,0.79043) rectangle (0.12929,1.00000);
\fill[white] (0.13197,0.79043) rectangle (0.21167,1.00000);
\end{scope}
\end{tikzpicture}
\end{sourceblock}
使用稀疏拉普拉斯算子 (8.72) 产生

\begin{sourceblock}
\begin{tikzpicture}
\node[inner sep=0,anchor=south west] (image) at (0,0) {\includegraphics[page=265,trim={53.00bp 378.03bp 53.87bp 224.92bp},clip,width=0.92\textwidth,max height=0.38\textheight]{Ferziger4th.pdf}};
\begin{scope}[x={(image.south east)},y={(image.north west)}]
\fill[white] (0.54174,0.23991) rectangle (0.56153,0.42285);
\fill[white] (0.55576,0.01899) rectangle (0.56890,0.20193);
\end{scope}
\end{tikzpicture}
\end{sourceblock}
泊松方程求解器中每个时间步使用的初始压力应该是前一个时间步的压力。

在节。由图7.2.1可知，我们用分数步法的应用结果可以看出，式(7.59)产生了O(t)2误差(1/2)tL(G(p′))，但与基本离散化是一致的。此处遵循相同的过程会显示完全相同的错误，即

\begin{sourceblock}
\includegraphics[page=265,trim={53.00bp 252.48bp 53.87bp 385.64bp},clip,width=0.92\textwidth,max height=0.38\textheight]{Ferziger4th.pdf}
\end{sourceblock}
这意味着与 P1 方法相比，该过程中的额外压力消除步骤将分数阶误差降低了一个数量级；再次回忆一下 L = D ( G ( ))。 Armfield 等人提出的结果。 (2010)表明这种稀疏方案基本上没有散度误差。然而，它们的紧凑压力校正方案（通过直接微分泊松方程导出；请参阅第 7.1.3 节或第 7.1.5.1 节中的警告）具有大致相同的发散误差，无论压力泊松方程的收敛程度如何。与方法误差的结果一致，两种方案具有大致相同的精度。两种方案均未观察到电网规模振荡，但在修改方案中求解新的全压力可防止电网规模振荡在任何情况下累积。

\subsection*{8.3.4 初始条件和边界条件}\phantomsection\addcontentsline{toc}{subsection}{8.3.4 初始条件和边界条件}
速度的初始条件应该是无散度的。此外，大多数分数步方案都采用 Adams-Bashforth 方案进行对流。由于 Adams-Bashforth 方案是多级方法，因此无法仅使用初始时间点的数据来启动解。必须使用其他方法来开始计算，例如带有压力迭代的 Crank-Nicolson 方法；见第 6.2.2 节.通常，第一步计算的压力校正场可能包括一阶时间误差，因为初始压力到处都是零，并且是在错误的时间指定的，即不是在 1 / 2 时间水平；弗林格等人。 （2003）举个例子。然而，如果使用多级启动，则不会出现此问题，因为在正确的时间（即时间步长的中间）计算全压力，直到收集到足够的时间数据以使用常规方法。

边界条件是一个更复杂的问题。对于原始的 P1 方案，在估计新速度的第一步之后需要特殊的中间边界条件（Kim 和 Moin 1985；Zang 等人 1994）。然而，一般来说，物理边界条件可用于速度和标量，而对于压力，所需的条件取决于方法，如下所示：

1. 对于交错网格，不需要压力边界条件，但将垂直于表面的压力校正（或 P1 方案中的伪压力）的梯度设置为零是合适的。 2
2. 对于使用边界外部辅助节点的并置方案，直接内部节点处的法向动量方程需要直接外部节点处的压力，因此使用从内部进行的高阶外推。同样，将垂直于表面的压力校正（或 P1 方案中的伪压力）的梯度设置为零是合适的。 3
3. 对于壁面处的切向速度，值得注意的是，虽然该方法的速度估计步骤中可以为估计速度设置边界条件，但投影步骤本质上是一个无旋步骤，无法约束切向速度修正；因此，会犯一个小错误（Armfield and Street 2002）。对于具有较小分数步误差的方案（参见方程（7.59）和（8.81）），即当发散约束的修正较小时，它较小，因为原始速度估计更好。

2 当存在狄利克雷速度边界条件时，这是正确的；在速度具有诺依曼边界条件的流出处，根据流量，压力边界条件有多种选项。专业代码文档通常描述这些选项。另请参见 Sani 等人。 （2006）。 3 同上。

\subsection*{8.3.5 迭代与非迭代方案}\phantomsection\addcontentsline{toc}{subsection}{8.3.5 迭代与非迭代方案}
总是出现这样的问题：在分数步法中使用迭代是否有好处，即在压力校正和速度更新完成后返回，以在第二次（或第三次……）过程中使用该信息。出现两个问题：准确性和效率。请注意，此讨论适用于经典的分数步方法，其中对流项使用显式 Adams-Bashforth 方案进行近似；上一节描述了迭代完全隐式方案。

ArmfieldandStreet（2000、2002、2003、2004）在交错和共置网格的背景下对迭代方案进行了检查。结果不会因网格类型而有显著差异，但会根据方案而有所不同。对于所有测试，迭代方案的效率较低，因为它们需要更多的 CPU 时间才能达到规定的精度水平，但由于它们通过将连续迭代中的压力校正或压力差驱动为零来消除分割误差 ( 1 / 2 ) t L ( G ( p ′ ))，因此迭代方法对于给定的网格和时间步长更准确。

Armfield 和 Street (2004) 使用基本方程组 (7.50) 和 (7.51) 检查了迭代、P2 和 P3 方案以及交错网格上的新“压力”方案。

P2 方法：P2 方法是由式（7.53 – 7.57）描述的，以获得 v n + 1
和 p n + 1 / 2 以及上面提到的固定投影误差并在方程（7.59）中给出。

迭代方法：迭代方法只是循环方程，直到压力修正为零并且解满足动量方程和连续性方程。请注意，与同一算法的非迭代版本相比，迭代方法的稳定性限制没有显著改善，因为该限制是由对流的显式 Adams-Bashforth 处理施加的。

P3 方法：P3 方法不是迭代的。 Gresho (1990) 提出了类似的方法，但由于稳定性问题显然没有实施。基本思想是通过使用二阶外推法更好地估计该方案第一步中的压力

\begin{sourceblock}
\includegraphics[page=267,trim={53.00bp 220.42bp 53.87bp 427.28bp},clip,width=0.92\textwidth,max height=0.38\textheight]{Ferziger4th.pdf}
\end{sourceblock}
所以方程(7.53)变为

\begin{sourceblock}
\begin{tikzpicture}
\node[inner sep=0,anchor=south west] (image) at (0,0) {\includegraphics[page=267,trim={53.00bp 159.88bp 53.87bp 470.87bp},clip,width=0.92\textwidth,max height=0.38\textheight]{Ferziger4th.pdf}};
\begin{scope}[x={(image.south east)},y={(image.north west)}]
\fill[white] (0.21802,0.24414) rectangle (0.24620,0.57078);
\fill[white] (0.92439,0.00000) rectangle (1.00000,0.07855);
\end{scope}
\end{tikzpicture}
\end{sourceblock}
(8.83) 压力修正由下式给出

\begin{sourceblock}
\includegraphics[page=267,trim={53.00bp 114.03bp 53.87bp 533.67bp},clip,width=0.92\textwidth,max height=0.38\textheight]{Ferziger4th.pdf}
\end{sourceblock}
动量方程中的压力现在在时间上近似为二阶，固定投影误差为三阶，因此解将比 P2 给出的解更准确。由于对压力和对流进行二阶外推，该方案的前两步需要使用不同的方法； Kirkpatrick 和 Armfield (2008) 对动量方程和压力迭代中的所有项使用 Crank-Nicolson。

压力法：这里的想法是首先直接求解压力修正，给定 p n − 1 / 2，通过动量方程的散度为

\begin{sourceblock}
\includegraphics[page=268,trim={53.00bp 498.31bp 53.87bp 125.53bp},clip,width=0.92\textwidth,max height=0.38\textheight]{Ferziger4th.pdf}
\end{sourceblock}
然后，在动量方程中使用压力 p n + 1 / 2 = p n − 1 / 2 + p ′ 来求出新的速度。这样就完成了时间步。

上述每个方案均使用上述工具求解，包括 QUICK、ADI（使用完整系统的四次扫描）和 GMRES。这里给出的结果来自 Ra = 6 × 10 5 和 Pr = 7 的二维方形空腔中自然对流的测试用例。 5（见第 8.4 节）；模拟了从初始条件（恒定温度下静止的流体）到 t = 2 状态的流动发展。改变时间步长来检查收敛性，但网格保持固定在 50 × 50 均匀网格。该代码从无量纲时间 t = 0 到 2 运行，步长大小从 0 变化。 003125 至 0.1. 时间步长 = 7 的基准。运行 8125 × 10 − 4 是为了将误差评估为运行与基准之间差异的 L 2 范数。虽然 ADI 扫描是有限的，但 GMRES 扫描的数量取决于具体情况；对于非迭代情况下最严格的收敛标准，最多可达一百次，每个迭代步骤仅进行五次扫描。更多细节可以在论文中找到。自然对流的基本物理原理将在下面的示例部分中进行解释。

\begin{sourceblock}
\includegraphics[page=268,trim={53.00bp 224.93bp 53.87bp 372.07bp},clip,width=\textwidth,max height=0.38\textheight]{Ferziger4th.pdf}
\end{sourceblock}
\begin{sourceblock}
\begin{tikzpicture}
\node[inner sep=0,anchor=south west] (image) at (0,0) {\includegraphics[page=268,trim={53.00bp 137.41bp 53.87bp 441.84bp},clip,width=0.92\textwidth,max height=0.38\textheight]{Ferziger4th.pdf}};
\begin{scope}[x={(image.south east)},y={(image.north west)}]
\fill[white] (0.53104,0.71976) rectangle (0.61065,0.83324);
\fill[white] (0.53104,0.63851) rectangle (0.55955,0.75199);
\fill[white] (0.53104,0.56082) rectangle (0.61188,0.67430);
\fill[white] (0.53104,0.48130) rectangle (0.55955,0.59478);
\fill[white] (0.35447,0.05582) rectangle (0.39074,0.30372);
\fill[white] (0.39934,0.01381) rectangle (0.45525,0.12729);
\end{scope}
\end{tikzpicture}
\end{sourceblock}
\begin{sourceblock}
\begin{tikzpicture}
\node[inner sep=0,anchor=south west] (image) at (0,0) {\includegraphics[page=268,trim={53.00bp 112.54bp 53.87bp 541.34bp},clip,width=0.92\textwidth,max height=0.38\textheight]{Ferziger4th.pdf}};
\begin{scope}[x={(image.south east)},y={(image.north west)}]
\fill[white] (0.39976,0.09788) rectangle (0.45567,0.90212);
\end{scope}
\end{tikzpicture}
\end{sourceblock}
\begin{sourceblock}
\begin{tikzpicture}
\node[inner sep=0,anchor=south west] (image) at (0,0) {\includegraphics[page=268,trim={53.00bp 81.49bp 53.87bp 564.45bp},clip,width=0.92\textwidth,max height=0.38\textheight]{Ferziger4th.pdf}};
\begin{scope}[x={(image.south east)},y={(image.north west)}]
\fill[white] (0.39976,0.45248) rectangle (0.45567,0.94059);
\fill[white] (0.43690,0.05941) rectangle (0.49221,0.54752);
\fill[white] (0.67480,0.05941) rectangle (0.71889,0.54752);
\fill[white] (0.91269,0.05941) rectangle (0.94553,0.54752);
\end{scope}
\end{tikzpicture}
\end{sourceblock}
\begin{sourceblock}
\begin{tikzpicture}
\node[inner sep=0,anchor=south west] (image) at (0,0) {\includegraphics[page=269,trim={53.00bp 524.49bp 53.87bp 58.02bp},clip,width=0.92\textwidth,max height=0.38\textheight]{Ferziger4th.pdf}};
\begin{scope}[x={(image.south east)},y={(image.north west)}]
\fill[white] (0.40364,0.87158) rectangle (0.45696,0.98565);
\fill[white] (0.78433,0.71888) rectangle (0.86096,0.83295);
\fill[white] (0.78433,0.63793) rectangle (0.81191,0.75188);
\fill[white] (0.78433,0.56033) rectangle (0.86220,0.67440);
\fill[white] (0.78433,0.48093) rectangle (0.81191,0.59500);
\fill[white] (0.35453,0.07653) rectangle (0.38956,0.32476);
\fill[white] (0.40195,0.01435) rectangle (0.45585,0.12842);
\end{scope}
\end{tikzpicture}
\end{sourceblock}
\begin{sourceblock}
\begin{tikzpicture}
\node[inner sep=0,anchor=south west] (image) at (0,0) {\includegraphics[page=269,trim={53.00bp 500.60bp 53.87bp 153.61bp},clip,width=0.92\textwidth,max height=0.38\textheight]{Ferziger4th.pdf}};
\begin{scope}[x={(image.south east)},y={(image.north west)}]
\fill[white] (0.40235,0.10059) rectangle (0.45627,0.89941);
\end{scope}
\end{tikzpicture}
\end{sourceblock}
\begin{sourceblock}
\begin{tikzpicture}
\node[inner sep=0,anchor=south west] (image) at (0,0) {\includegraphics[page=269,trim={53.00bp 458.17bp 53.87bp 176.37bp},clip,width=0.92\textwidth,max height=0.38\textheight]{Ferziger4th.pdf}};
\begin{scope}[x={(image.south east)},y={(image.north west)}]
\fill[white] (0.40235,0.66013) rectangle (0.45627,0.96203);
\fill[white] (0.45615,0.43671) rectangle (0.47173,0.73829);
\fill[white] (0.62093,0.43671) rectangle (0.64731,0.73829);
\fill[white] (0.78659,0.43671) rectangle (0.82373,0.73829);
\fill[white] (0.95242,0.43671) rectangle (1.00000,0.73829);
\fill[white] (0.61714,0.03797) rectangle (0.67507,0.40665);
\fill[white] (0.67708,0.03797) rectangle (0.73044,0.40665);
\fill[white] (0.73245,0.03797) rectangle (0.84192,0.40665);
\end{scope}
\end{tikzpicture}
\end{sourceblock}
显示迭代方案是四种方案中效率最低的，而 P3 和压力方案效率最高。在最后这些情况下，虽然这些方案不是迭代的，但它们的设置意味着实现一定精度水平的工作量显著减少，即，对于给定的精度水平，它们只需要迭代方案的 50\% 的 CPU 时间和 P2 方案的 60\% 的时间。 Shen (1993) 指出，类似 P3 的方法可以产生无限时间的解；在任何这些测试或柯克帕特里克和阿姆菲尔德（2008）更广泛的测试中都没有观察到这种行为。

该测试用例的结论（从初始条件发展到稳态的非定常流）可能无法代表所有非定常流。然而，可以公平地说，对于相同的时间步长，分数步方法的非迭代版本通常需要迭代版本的一半计算时间。与显式处理对流项的方法（无论是迭代还是非迭代）相比，上一节中描述的完全隐式迭代方法允许使用更大的时间步长；我们将在下一节中介绍完全隐式迭代方法和 P2 非迭代方法版本的一些结果。

\section*{8.4 例子}\phantomsection\addcontentsline{toc}{section}{8.4 例子}
在本节中，我们将给出使用三种不同求解技术的示例：用于稳态流的 SIMPLE 和隐式迭代分步法 (IFSM)，以及用于非稳态流的 SIMPLE、IFSM 和 P2 非迭代分步法的版本。首先，使用交错和共置网格，我们在具有两种配置的方形外壳中处理稳态流。我们从盖子驱动的空腔开始，然后研究浮力作为驱动力的情况。然后，我们使用同位网格来检查由振荡盖或振荡热壁温度驱动的不稳定周期性流动的特征。我们评估迭代和离散化误差，以及各种其他参数对精度和效率的影响（例如，SIMPLE 中的欠松弛因子和分数步法中的时间步长、交错排列与共置排列等）。

\subsection*{8.4.1 方形外壳中的稳态流动}\phantomsection\addcontentsline{toc}{subsection}{8.4.1 方形外壳中的稳态流动}
在本节中，我们将演示隐式迭代求解方法在层流稳态流计算中的应用。用于获得所提出解的计算机代码以及必要的输入数据可供下载；详细内容请参见附录。作为测试用例，我们选择了方形外壳中的两个流；一种流动由移动的盖子驱动，另一种流动由浮力驱动。几何形状和边界条件如图8.7 所示。这两个测试用例都被许多作者使用过，并且文献中提供了准确的解；例如，参见 Ghia 等人。 （1982）和霍特曼等人。 （1990）。我们比较了 SIMPLE 和迭代分数步法 (IFSM) 的性能，并特别演示了如何估计迭代和离散化误差。使用交错网格和共置网格。

我们首先考虑盖子驱动的型腔流动，这是一个流行的基准问题。请参阅 Erturk (2009) 进行富有洞察力的评论。移动的盖子在下面的两个角产生一个强涡流和一系列较弱的涡流（在较高雷诺数下的模拟显示在左上角也形成了第三个涡流）。基于腔高度 H 和盖速度 U L 的非均匀网格样本和雷诺数流线，Re = U L H /ν = 1000，如图 8.8 所示；如图所示，计算是在更精细的非均匀网格上进行的。

现在我们可以看看迭代误差的估计。 节 中提出了几种方法。 5.7.首先，通过迭代直到残差范数变得可以忽略不计（双精度舍入误差的量级）来获得准确的解。然后重复计算，并将迭代误差计算为先前获得的收敛解与中间解之间的差。

\begin{sourceblock}
\includegraphics[page=270,trim={53.00bp 90.72bp 53.87bp 430.78bp},clip,width=\textwidth,max height=0.38\textheight]{Ferziger4th.pdf}
\par\vspace{0.45\baselineskip}
\small 图 8.7 二维稳态流测试用例的几何形状和边界条件：盖驱动（左）和浮力驱动（右）空腔流
\end{sourceblock}
\begin{sourceblock}
\includegraphics[page=271,trim={53.00bp 442.48bp 53.87bp 52.00bp},clip,width=\textwidth,max height=0.38\textheight]{Ferziger4th.pdf}
\par\vspace{0.45\baselineskip}
\small 图 8.8 用于解决空腔流动问题的 64 × 64 CV 非均匀网格（左）和 Re = 1000 时盖驱动空腔流动的流线（右），在 256 × 256 CV 非均匀网格上计算（一个涡流中任意两条相邻流线之间的质量流量恒定）
\end{sourceblock}
\begin{center}\small 图 8.9 显示了迭代误差的范数、从方程（5.89）或（5.96）获得的估计、两次迭代之间的差值以及残差。计算在 32 × 32 CV 网格上进行，速度欠松弛因子为 0.7，压力欠松弛因子为 0.3（未针对效率进行优化）。由于该算法需要多次迭代才能收敛，因此误差估计器所需的特征值是最近 50 次迭代的平均值。这些场是通过对下一个较粗网格的解进行插值来启动的，这就是初始误差相对较低的原因。\end{center}
该图表明误差估计技术对于非线性流动问题给出了良好的结果。在求解过程的开始阶段估计不好，误差很大。使用两次迭代之间的差异或残差的绝对水平并不是迭代误差的可靠度量。这些量确实以与误差相同的速率减少，但需要适当地归一化以定量地表示迭代误差。而且，它们最初下降得非常快，而误差减少则慢得多。然而，一段时间后，所有曲线都变得几乎平行，并且如果人们粗略地知道初始误差的阶数（如果从零初始场开始，它就是解本身），那么停止迭代的可靠标准是将两次迭代之间的差值或残差的范数减少某个因子，例如三或四个数量级。在其他网格和其他流动问题上获得了与图 8.9 中所示类似的结果。

接下来我们转向离散化误差的估计。我们使用CDS和UDS离散化对五个网格进行计算；最粗的有 16 × 16 CV，最细的有 256 × 256 CV。使用均匀和非均匀共置网格，并应用 SIMPLE 和 IFSM 方法。主涡流的强度 ψ min，表示涡流中心和涡流中心之间的质量流量

\begin{center}\small 图 8.9 SIMPLE 方法的精确迭代误差和估计迭代误差范数的比较、两次迭代之间的差异以及在 32 × 32 CV 网格上 Re = 10 3 时盖驱动腔流解的残差，对流和扩散均采用 CDS 离散化\end{center}
在所有网格上比较了边界和较大次级涡流 ψ max 的强度。这些值是通过查找单元角中流函数的最大值和最小值来估计的。流函数值是通过对通过单元面的体积通量求和来计算的：我们将左下角的值设置为零，并通过添加通过连接两个顶点的面的体积通量来计算下一个顶点的值。图 8.10 显示了计算出的涡强度随着网格的细化而变化。四个最精细网格的结果显示两个量向独立于网格的解的单调收敛。非均匀网格的结果显然比均匀网格的结果更准确。细胞从壁向空腔中心生长的膨胀比在最粗网格上为1.17166，在最细网格上为1.01；下一个较细网格上的值等于下一个较粗网格上的值的平方根，以确保来自较粗网格的网格线保留在较细网格中（即每个粗网格CV恰好包含4个细网格CV）。

为了实现定量误差估计，使用在两个最精细网格上获得的结果和理查森外推法来估计与网格无关的解（参见第 3.9 节）。这些值为：ψ min = − 0.11893 且 ψ max = 0.00173. 请注意，理查森外推应用于均匀或非均匀网格上获得的解，使用 SIMPLE 或 IFSM，提供与四个有效数字相同的估计值；主涡流相差 0.007\%，次要涡流相差 0.027\%（这适用于 CDS 离散化）。通过从参考解中减去给定网格上的结果，可以获得误差估计。图 8.11 中根据平均网格大小绘制了误差。对于均匀网格和非均匀网格上的两个量，二阶误差的预期减少

\begin{sourceblock}
\includegraphics[page=273,trim={53.00bp 457.48bp 53.87bp 52.00bp},clip,width=\textwidth,max height=0.38\textheight]{Ferziger4th.pdf}
\par\vspace{0.45\baselineskip}
\small 图 8.10 Re = 1000 时盖驱动空腔流中初级（ψ min；左）和次级（ψ max；右）涡流强度的变化，使用 SIMPLE 和 IFSM 求解方法以及均匀和非均匀网格计算
\end{sourceblock}
使用CDS离散化时得到的方法；对于UDS，误差减小只是接近渐近一阶收敛，因为误差太大。非均匀网格上的误差较低，特别是对于 ψ max；由于次级涡流被限制在一个角落，因此那里的非均匀网格更精细，从而产生更高的精度。然而，由于该数量比 ψ min 小两个数量级，因此误差估计不太可靠（必须迭代到更严格的容差才能获得 ψ max 的准确值）。

请注意，使用对流一阶迎风离散化产生的误差远高于使用 CDS 获得的解中的误差。即使在 256 × 256 CV 的网格上（可以认为非常精细），使用 UDS 时 ψ min 的误差也在 10\% 左右，这几乎比使用 CDS 产生的误差大两个数量级；见图8.11。

使用SIMPLE和IFSM获得的结果只能在两个最粗的网格上相互区分；对于所有更精细的网格，差异太小而无法在图表中看到。这是预期的，因为在这两种情况下都使用相同的离散化；这两种方法的区别仅在于压力修正方程，预计压力修正方程只会影响迭代求解过程的收敛速度，而不影响解本身。因此，我们在图 8.12 中仅显示使用 IFSM 获得的中心线速度剖面。使用二阶CDS导致单调收敛；连续网格上的错误比率是四倍。两个最细网格的速度分布很难彼此区分。

接下来我们研究使用 CDS 离散化在均匀共置网格和交错网格上获得的解之间的差异。由于速度节点位于两个网格上的不同位置，因此交错值被线性插值到单元中心（高阶插值会更好，但线性插值就足够了）。每个变量 ( φ = ( u x , u y , p ) ) 的平均差异确定为：

\begin{sourceblock}
\includegraphics[page=274,trim={53.00bp 454.48bp 53.87bp 52.00bp},clip,width=\textwidth,max height=0.38\textheight]{Ferziger4th.pdf}
\par\vspace{0.45\baselineskip}
\small 图 8.11 Re = 1000 时盖驱动空腔流中主涡流（ψ min；左）和次涡流（ψ max；右）强度解的估计离散化误差，使用 SIMPLE 和 IFSM 求解方法以及均匀和非均匀网格计算
\end{sourceblock}
\begin{sourceblock}
\includegraphics[page=274,trim={53.00bp 258.90bp 53.87bp 245.37bp},clip,width=\textwidth,max height=0.38\textheight]{Ferziger4th.pdf}
\par\vspace{0.45\baselineskip}
\small 图 8.12 Re = 1000 时盖驱动腔流中 u x 速度沿垂直中心线（左）和 u y 速度沿水平中心线（右）的分布图，使用 IFSM 在五个非均匀网格上计算
\end{sourceblock}
\begin{sourceblock}
\includegraphics[page=274,trim={53.00bp 176.85bp 53.87bp 453.47bp},clip,width=0.92\textwidth,max height=0.38\textheight]{Ferziger4th.pdf}
\end{sourceblock}
其中 N 是 CV 的数量。对于 u x 和 u y，任何网格上的 ε 都比同一网格上的离散化误差小一个数量级。压力差异稍小（无需插值）。

接下来研究使用 CDS 和共置非均匀网格的 SIMPLE 和 IFSM 方法的收敛特性。 SIMPLE 算法有两个可调参数：速度欠松弛和压力欠松弛。 IFSM 不需要任何变量欠松弛，并且只有一个参数——时间步长（通常以标准化形式表示为 CFL 数，

\begin{sourceblock}
\includegraphics[page=275,trim={53.00bp 414.25bp 53.87bp 52.00bp},clip,width=\textwidth,max height=0.38\textheight]{Ferziger4th.pdf}
\par\vspace{0.45\baselineskip}
\small 图 8.13 使用具有欠松弛参数的各种组合的 SIMPLE 和具有交错（左）和共置（右）变量排列的 32 × 32 CV 均匀网格（Re = 1000 时的盖驱动腔流），将所有方程中的残差水平降低三个数量级所需的外部迭代次数
\end{sourceblock}
CFL x = u x t / x 或 CFL y = u y t / y )。我们首先使用不同的速度欠松弛参数来检查压力欠松弛参数 α p （参见方程（7.84））对 SIMPLE 收敛的影响。图 8.13 显示了使用欠松弛参数和 32 × 32 CV 均匀网格的各种组合将所有方程中的残差水平降低三个数量级所需的外部迭代次数。

该图显示，对于两种类型的变量排列，对欠松弛参数 α p 的依赖性几乎相同，尽管对于共置网格来说，良好值的范围稍宽一些。当速度更加欠松弛时，我们可以使用 0.1 到 1.0 之间的任何 α p 值，但该方法收敛缓慢。对于较大的α u 值，收敛速度较快，但α p 的有用范围受到限制。 α p 的值由方程式建议。 7.89 接近最佳； α p = 1.1 − α u 给出了该流程的最佳结果。通常，改变一个参数并根据上述关系确定另一个参数。

在图 8.14 中，我们展示了在 Re = 1000 的盖驱动腔流的情况下，速度欠松弛因子 α u 对 SIMPLE 收敛速度的影响，以及 CFL 数对变量同位排列和非均匀网格的 IFSM 收敛的影响，在 Re = 1000 的盖驱动腔流的情况下。对于交错变量排列和均匀网格也获得了类似的图表。对于这两种求解方法，所需的迭代次数随着网格的细化而增加。在SIMPLE情况下，欠松弛的最佳值接近但小于1；所有网格上的迭代均出现 α u = 0.99 的偏差，两个最细网格上的 α u = 0.98 的迭代也出现偏差。网格越细，随着 α u 的减小，所需迭代次数的增加就越陡峭

\begin{sourceblock}
\includegraphics[page=276,trim={53.00bp 456.25bp 53.87bp 52.00bp},clip,width=\textwidth,max height=0.38\textheight]{Ferziger4th.pdf}
\par\vspace{0.45\baselineskip}
\small 图 8.14 将所有方程中的残差水平降低四个数量级所需的外部迭代次数，作为 SIMPLE 中的欠松弛因子 α u（左）或 IFSM 中的 CFL 数（右）的函数； Re = 1000 时盖子驱动的腔体流动
\end{sourceblock}
低于最佳值。在合理的范围内，IFSM 似乎对 CFLnumber 的值不太敏感：如果 CFL 低于 10，则需要多次迭代，但在 20 到 200 之间，增加是适度的，特别是对于较粗的网格。我们在第 1 章中展示。 12 如何使用多重网格方法提高精细网格的计算效率。

接下来我们研究 SIMPLE 中的速度欠松弛参数和 IFSM 中的 CFL 数对同位变量排列和非均匀网格解的影响。我们比较使用 α u = 0 获得的解。 95 且α u = 0。 SIMPLE 中为 6，IFSM 中 CFL = 6.4 和 CFL = 102.4（CFL 数是使用盖速度和平均网格间距确定的，就像在均匀网格上一样；对于使用的非均匀网格，整个网格中发现的最大 CFL 数大约高出 50\%）。使用表达式 (8.86) 评估解的差异（在残余水平降低四个数量级之后，以最小化迭代误差的影响）。对于 SIMPLE，我们在 32 × 32 CV 的网格上获得约 5 × 10 − 4 的速度 ε 值，在 64 × 64 CV 的网格上获得约 8 × 10 − 5 的速度 ε 值，在 128 × 128 CV 的网格上获得约 1.5 × 10 − 5 的速度 ε 值。对于IFSM，解的差异较大；我们在 32 × 32 CV 的网格上获得了大约 2 × 10 − 2 的速度 ε 值，在 64 × 64 CV 的网格上获得了 6 × 10 − 3 的速度，在 128 × 128 CV 的网格上获得了 1.4 × 10 − 3 的速度 ε 值。这些差异远小于相应网格上的离散化误差（分别约为 10\%、2.5\% 和 0.6\%），并且随着网格细化以与离散化误差减小相同的速率减小；因此，它们可以被忽略。

稳态解对欠松弛参数或时间步长的依赖性源于以下事实：当使用变量的同位排列时，这些参数影响用于计算通过单元面的质量通量的插值速度的校正。 Rhie 和 Chow（1983）之后的修正或多或少是避免压力-速度解耦的标准方法，并且它是

\begin{sourceblock}
\includegraphics[page=277,trim={53.00bp 441.25bp 53.87bp 52.00bp},clip,width=\textwidth,max height=0.38\textheight]{Ferziger4th.pdf}
\par\vspace{0.45\baselineskip}
\small 图 8.15 瑞利数 Ra = 10 5 和普朗特数 Pr = 0.1 处的浮力驱动空腔流等温线（左）和速度矢量（右）（任意两条相邻等温线之间的温度差相同）
\end{sourceblock}
也用于上述计算。然而，还有其他方法可以避免并置网格上的压力或速度振荡（参见 Armfield 和 Street 2005）。此外，还可以将插值固定为始终对应于无欠松弛且与时间步长无关；请参阅 Pascau (2011) 了解详细讨论和一种解，以及 Tukovic 等人。 （2018）扩展到移动网格。在大多数应用中，这种相关性不会造成任何问题，因为它比离散化误差小。然而，当使用非常小的时间步长而流量没有及时改变时，可能会出现问题，并且可能需要引用的参考文献中提出的解决方法之一。

接下来考虑方腔中的 2D 浮力驱动流动，如图 1 所示。 8.15。我们现在需要求解与纳维-斯托克斯方程耦合的能量方程，因为速度场取决于解域内的温度分布。这里考虑的不可压缩流体的能量方程由我们到目前为止考虑的通用标量输运方程表示；我们只需将扩散率设置为粘度与普朗特数之比即可。重力方向速度分量的动量方程获得了一个额外的源项；使用本示例中使用的 Bousinesq 近似，该源项为（请参阅第 1.4 节末尾的解释）：

\begin{sourceblock}
\includegraphics[page=277,trim={53.00bp 144.76bp 53.87bp 505.92bp},clip,width=0.92\textwidth,max height=0.38\textheight]{Ferziger4th.pdf}
\end{sourceblock}
其中 β 是体积膨胀系数，g i 是 i 方向上的重力分量，T ref 是参考温度，在该温度下计算密度 ρ ref 作为温度的函数。然后密度被认为是常数，而动量方程中的上述源项表示围绕 ρ ref 的（线性化）密度变化的影响。请注意，只有当密度变化接近线性时，这种近似才可行。当模拟液体流动时，只能在相对较小的温度范围内假设这一点。

在此测试用例中，重复使用了用于盖驱动腔流计算的相同五个网格。冷壁和热壁是等温的。受热的流体沿着热壁上升，而冷却的流体沿着冷壁下降。普朗特数为 0.1（意味着热扩散率占主导地位），选择温差和其他流体特性，使得瑞利数为

\begin{sourceblock}
\includegraphics[page=278,trim={53.00bp 488.41bp 53.87bp 148.74bp},clip,width=0.92\textwidth,max height=0.38\textheight]{Ferziger4th.pdf}
\end{sourceblock}
此处使用以下值：T hot = 10、T Cold = 0、ρ = 1、β = 0.01、g = 10、μ = 0.001 和 H = 1（通篇使用 SI 单位）。对于每个网格，计算从初始字段 u x = 0、u y = 0 和 T = 6 开始。

预测的速度矢量和等温线如图 8.15 所示。流动结构很大程度上取决于普朗特数。在空腔的中心区域形成一个几乎停滞、稳定分层的大流体核心。请注意，等温线以直角接近绝热（顶部和底部）壁；在热通量为零的任何边界（通常是绝热壁和对称平面）上都必须如此。在检查数值解的合理性时，这是要检查的特征之一。我们已经看到（甚至在期刊出版物中）违反此条件的结果。当等温线更密集时，垂直于等温线方向的温度梯度更高，因此传导的热通量也更高。这在图 8.16 中可见，它显示了沿冷壁每单位面积的局部热通量的变化（热通量为负，因为热量“进入”解域）：热流体在顶部附近与冷壁相遇的地方（另请参见图 8.15 中的等温线和速度矢量）比在底部附近冷却流体离开冷壁的地方传热要强烈得多。

\begin{center}\small 图 8.16 显示了解对网格精细度的依赖性。代表两个最细网格上的解的线在视觉上无法彼此区分，这表明精度很高。接下来提出离散化误差的估计。\end{center}
预计非均匀网格将比均匀网格给出更准确的结果。确实如此。图 8.17 显示了通过等温壁的总热通量以及使用理查森外推法估计的离散化误差作为均匀和非均匀网格的网格细度的函数，使用 SIMPLE 或 IFSM 求解方法进行计算。当应用于网格类型和求解方法的两个最精细水平的结果时，理查森外推法会产生与网格无关的值的相同估计，直到五个有效数字。该估计为 Q = 0.39248，当通过纯热传导的热通量标准化时，
Q 条件 = 0.1，给出努塞尔数 Nu = 3.9248。通过从估计的与网格无关的解中减去所有网格上的解，我们获得了离散化误差的估计。预测总热通量的误差如图 8.17 所示。所有误差均渐近趋于二阶方案预期的斜率（当网格间距减小一个数量级时，误差减小两个数量级）

\begin{sourceblock}
\includegraphics[page=279,trim={53.00bp 455.25bp 53.87bp 52.00bp},clip,width=\textwidth,max height=0.38\textheight]{Ferziger4th.pdf}
\par\vspace{0.45\baselineskip}
\small 图 8.16 Ra = 10 5 时浮力驱动腔流中沿着冷壁的单位面积局部热通量（左）和沿着垂直中心线的 u x 速度（右）的预测变化，使用 IFSM 在五个非均匀网格上计算
\end{sourceblock}
\begin{sourceblock}
\includegraphics[page=279,trim={53.00bp 256.90bp 53.87bp 244.37bp},clip,width=\textwidth,max height=0.38\textheight]{Ferziger4th.pdf}
\par\vspace{0.45\baselineskip}
\small 图 8.17 在瑞利数 Ra = 10 5 和普朗特数 Pr = 0.1 时，通过等温壁的热通量 Q （左）以及浮力驱动腔流中相关的离散化误差（右），作为网格细度的函数（在均匀和非均匀网格上使用 CDS 进行计算，采用 SIMPLE 或 IFSM 求解方法）
\end{sourceblock}
幅度）。非均匀网格上的热通量误差远小于均匀网格上的热通量误差，因为这些网格可以更好地解析边界层。

请注意，对于任何给定的网格，IFSM 的盖驱动腔流中主涡流强度的误差大于 SIMPLE（见图 8.11），而在浮力驱动流的情况下，IFSM 的总热通量误差小于 SIMPLE（见图 8.17）。然而，在这两种情况下，使用任一方法获得的解之间的差异都相对较小。

我们再次检查两种求解方法的效率及其对 SIMPLE 中欠松弛因子和 IFSM 中 CFL 数的依赖性。对于盖驱动腔，SIMPLE 中速度的欠松弛因子在

\begin{sourceblock}
\includegraphics[page=280,trim={53.00bp 456.25bp 53.87bp 52.00bp},clip,width=\textwidth,max height=0.38\textheight]{Ferziger4th.pdf}
\par\vspace{0.45\baselineskip}
\small 图 8.18 将所有方程中的残差水平降低四个数量级所需的外部迭代次数，作为 SIMPLE 中的欠松弛因子 α u（左）或 IFSM 中的 CFL 数（右）的函数； Ra = 10 5 时的浮力驱动腔流
\end{sourceblock}
0.5 和 0.98。对于 IFSM，时间步长的选择不太明显，因为我们事先不知道速度有多高以及会产生哪些 CFL 数。有趣的是，对于盖驱动腔流使用相同的时间步长，并且在相同的时间步长下实现了最高效率（即所需迭代次数最少），尽管计算出的 CFL 数大约小了 3 倍。这是因为在两种情况下都使用了相同的流体属性（密度和粘度）和相同大小的解域。从式（7.94）中我们可以看出，速度修正量与时间步长和密度之比乘以压力修正梯度成正比；因此，可以预期，对于相同的时间步长、密度和网格间距，该方法将以类似的方式表现。一个例外是粗网格的最大时间步长：我们无法在浮力驱动空腔中获得 6.4 s 时间步长的解，但对于盖子驱动空腔流则可以获得解。

\begin{center}\small 图 8.18 显示了效率分析的结果。由于 IFSM 不使用欠松弛，因此时间步长（CFL 数）仍然是唯一的参数；另一方面，在 SIMPLE 中出现了另一个因素——温度的欠松弛因素。如果将其设置为等于速度的欠松弛因子，收敛会变得非常慢。然而，我们从经验中知道，能量方程在层流中表现良好，通常需要很少或不需要欠松弛。因此，这里将能量方程的欠松弛因子设置为0.99；以此值来看，SIMPLE的效率与IFSM相当。\end{center}
使用 SIMPLE 或 IFSM 等隐式方法计算稳定不可压缩流的效率可以通过使用多重网格方法进行外部迭代以及通过使用从较粗网格插值的解在较细网格上开始计算来显著提高，而不是像此处所做的那样从猜测的初始场开始。这将在第 1 章中进行演示。此处研究的两个测试用例为 12。

\subsection*{8.4.2 方形外壳中的非定常流}\phantomsection\addcontentsline{toc}{subsection}{8.4.2 方形外壳中的非定常流}
在本节中，我们将注意力转向非定常流动的预测。当迈向稳态流时，我们可以使用较大的时间步长，并且每个时间步长仅执行一次迭代，因为只要迭代过程收敛，中间解就无关紧要。然而，在计算非定常流动时，我们必须确保在进入下一个时间步骤之前所有方程都得到足够的满足。这意味着，一方面，我们必须选择合适的时间步长来充分解决变量的时间变化，另一方面，我们需要在每个时间步长执行足够数量的迭代，以确保迭代误差足够小。

对于时间精确的流动模拟，时间积分方案应该至少是二阶的。一阶隐式欧拉方案对于迈向稳态非常有用（例如，在 IFSM 中），但在周期流中，它需要太小的时间步长才能获得令人满意的精度。通常的选择是 Crank-Nicolson 方案（几乎专门用于分数步方法）和完全隐式三时间级别方案（几乎专门用于 SIMPLE 类型方法）。中心差用于对流项和扩散项。

使用的非迭代分数步方法与第 2 节中描述的方法略有不同。 7.2.1，因为它基于完全隐式的三时间电平方案而不是通常的克兰克-尼科尔森方案（计算机代码可在互联网上找到；详细信息请参阅附录）：

\begin{sourceblock}
\includegraphics[page=281,trim={53.00bp 302.57bp 53.87bp 331.57bp},clip,width=0.92\textwidth,max height=0.38\textheight]{Ferziger4th.pdf}
\end{sourceblock}
为了使方法非迭代，使用 Adams-Bashforth 方案估计时间水平 t n + 1 处的非线性对流通量；该表达式与式（7.52）中给出的表达式略有不同：

\begin{sourceblock}
\includegraphics[page=281,trim={53.00bp 235.05bp 53.87bp 415.03bp},clip,width=0.92\textwidth,max height=0.38\textheight]{Ferziger4th.pdf}
\end{sourceblock}
算法的其余部分遵循与第 2 节中描述的相同的路线。 7.2.1.

第一个流动问题是带有振荡盖子的盖子驱动的腔体流动：它以一定速度以正弦曲线移动

\begin{sourceblock}
\begin{tikzpicture}
\node[inner sep=0,anchor=south west] (image) at (0,0) {\includegraphics[page=281,trim={53.00bp 173.78bp 53.87bp 476.53bp},clip,width=0.92\textwidth,max height=0.38\textheight]{Ferziger4th.pdf}};
\begin{scope}[x={(image.south east)},y={(image.north west)}]
\fill[white] (0.00000,0.92546) rectangle (0.14899,1.00000);
\fill[white] (0.15170,0.92546) rectangle (0.20977,1.00000);
\fill[white] (0.21242,0.92546) rectangle (0.23053,1.00000);
\fill[white] (0.23323,0.92546) rectangle (0.33411,1.00000);
\end{scope}
\end{tikzpicture}
\end{sourceblock}
其中 ω = 2 π/ P，其中 P 是振荡周期。我们在这里设置 u max = 1 和 P = 10 s；腔高度为 1，粘度设置为 μ = 0.001（自始至终使用的 SI 单位），导致雷诺数在 0 到 1000 之间变化。解使用速度和压力（静止流体）的零值进行初始化。使用具有 128 × 128 CV 的非均匀网格（上一节中使用相同的网格来分析稳态流）。

我们首先通过计算不同时间步长的第一四分之一的 lid 振荡（从 t = 0 到 t = 2 . 5 s，对应于 ω t = π/ 2）来分析隐式欧拉和三时间级方案的时间精度。完全隐式方法（SIMPLE 和 IFSM）允许使用较大的时间步长，因此我们从 t = 0 开始。 125 s（相当于每个振荡周期 10 s 80 个时间步长）并减半 4 次。最精细的时间步长 ( t = 0 . 0078125 s) 对应于每个振荡周期 1280 个时间步长。非迭代分数步法（我们在此标记为 EFSM，因为它对于对流通量是显式的）无法使用大于 0.0125 s 的时间步长，因此在该方法中，我们使用时间步长 0.0125 s、0.00625 s 和 0.003125 s。为了评估离散化误差，使用积分很方便。我们选择在这里评估 t = 2 时移动盖子产生的主涡流强度的差异。 5秒。它是通过对质量通量进行积分来计算的，以获得单元顶点处的流函数值（从左下角的值 0 开始）；解域中的最小值表示给定盖子运动方向的涡流强度。使用理查森外推法和在两个最细网格上获得的值来估计与网格无关的值；所有方法都预测该值为 ψ min = − 0.043682.通过从该参考值中减去用不同方法和时间步长计算的值，估计离散化误差。

\begin{center}\small 图 8.19 显示了预测涡流强度的变化以及相关的时间离散误差作为时间步长和所使用方法的函数。在此示例中，对于每个时间步长，使用 SIMPLE 获得的结果似乎比使用 IFSM 获得的结果稍微准确一些；然而，情况并非总是如此，当我们对高 10 倍的粘度（雷诺数小 10 倍）重复此练习时，IFSM 结果的误差更小。两种方法均采用两种时间积分方案，表现出相同的行为。显然，一阶隐式欧拉方案（此处标记为 IE）导致的误差比使用二阶三时间级方案（此处标记为 TTL）时大一个数量级。当 IE 的时间步长减半时，错误会减半；而使用 TTL 方案时，错误会减少四倍。\end{center}
我们强制代码在每个时间步将残差减少四到五个数量级，以确保迭代误差可以忽略不计；这远远超过了通常情况下的情况，特别是当时间步长很小时，因此我们没有过多关注计算效率。然而，可以说，为了达到相同的残差水平，IFSM 每个时间步所需的计算时间比 SIMPLE 略少。对于相同的时间步长，EFSM比迭代方法需要更少的计算时间；然而，IFSM 达到规定精度所需的努力最少。请注意，IFSM 和 SIMPLE 已经达到了 0.1\% 左右的离散化误差水平，其中 EFSM 由于稳定性原因而无法工作。

\begin{center}\small 图 8.20 显示了在仿真开始后第五个振荡周期内创建的与 ω t = π/ 2 ( u L = 1)、ω t = π ( u L = 0)、ω t = 3 π/ 2 ( u L = − 1) 和 ω t = 2 π ( u L = 0) 对应的时间的速度矢量。此时的流动是完全周期性的，这可以通过解的完美对称性来识别，\end{center}
4.36 4.37

\begin{center}\small 图 8.19 t = 2 时主涡的强度。盖驱动腔流中的 5 s（上）和时间离散误差（下），其中盖振荡运动是时间步长的函数\end{center}
4.35

4.34

4.33

4.32

min

4.31

4.3

简单，TTL
4.29

IFSM、TTL 简单、IE
4.28

IFSM、IE EFSM、TTL
4.27

EFSM、IE
4.26

0 0.02 0.04 0.06 0.08 0.1 0.12 0.14
t

\begin{sourceblock}
\begin{tikzpicture}
\node[inner sep=0,anchor=south west] (image) at (0,0) {\includegraphics[page=283,trim={53.00bp 365.29bp 53.87bp 267.38bp},clip,width=0.92\textwidth,max height=0.38\textheight]{Ferziger4th.pdf}};
\begin{scope}[x={(image.south east)},y={(image.north west)}]
\fill[white] (0.35820,0.57873) rectangle (0.38487,0.96415);
\fill[white] (0.35820,0.03585) rectangle (0.38487,0.56618);
\end{scope}
\end{tikzpicture}
\end{sourceblock}
\begin{sourceblock}
\begin{tikzpicture}
\node[inner sep=0,anchor=south west] (image) at (0,0) {\includegraphics[page=283,trim={53.00bp 295.26bp 53.87bp 318.54bp},clip,width=0.92\textwidth,max height=0.38\textheight]{Ferziger4th.pdf}};
\begin{scope}[x={(image.south east)},y={(image.north west)}]
\fill[white] (0.39206,0.80760) rectangle (0.43940,0.97707);
\end{scope}
\end{tikzpicture}
\end{sourceblock}
相位相差 π。显然，在一个时期内会发生非常复杂的流型变化，数值方法需要准确地捕捉这些变化。

\begin{center}\small 图 8.19 中的误差分析表明，精确的解是通过隐式迭代方案（SIMPLE 和 IFSM）产生的，时间约为 100 秒。每个振荡周期 100 个时间步长。事实上，如图 8.21 所示，随着时间步长的减小，解仅在尖峰处略有变化。正如预期的那样，使用一阶隐式欧拉方案获得的解与使用二阶三时间级方案获得的解存在很大偏差。然而，值得注意的是，二阶方案在解中再现了如此复杂的时间变化。每个振荡周期 60 个时间步长。\end{center}
我们简要讨论的第二个非定常流是上一节中的浮力驱动的空腔流，但具有正弦变化的热壁温度：

\begin{sourceblock}
\includegraphics[page=283,trim={53.00bp 90.25bp 53.87bp 560.06bp},clip,width=0.92\textwidth,max height=0.38\textheight]{Ferziger4th.pdf}
\end{sourceblock}
\begin{sourceblock}
\includegraphics[page=284,trim={53.00bp 272.26bp 53.87bp 52.00bp},clip,width=\textwidth,max height=0.38\textheight]{Ferziger4th.pdf}
\par\vspace{0.45\baselineskip}
\small 图 8.20 四分之一周期、二分之一周期、四分之三周期和全周期时盖子驱动腔流中盖子振动运动的速度矢量（分别从左到右、从上到下）
\end{sourceblock}
其中，T 0 = 10 为平均热壁温度，Ta = 5 为温度波动幅度，ω = 2 π/ P，其中 P = 10 s 为振荡周期。流体最初处于静止状态，温度设置为 5.1，冷壁温度 T C = 0，热壁温度 T H = 10。平均瑞利数与上一节中介绍的稳态流动相同，但由于热壁温度周期性变化，流型随时间发生显著变化。经过几个周期后，流程会“忘记”初始化的影响并变得完全周期性。

\begin{center}\small 图 8.22 显示了速度矢量和图 8.23 等温线，分别对应于 ω t = π/ 2 ( T H = 15)、 ω t = π ( T H = 10)、 ω t = 3 π/ 2 ( T H = 5) 和 ω t = 2 π ( T H = 10)，当流动完全周期性时创建；它不是对称的（就像在盖驱动腔流的情况下一样），因为这里只有热壁温度\end{center}
\begin{sourceblock}
\begin{tikzpicture}
\node[inner sep=0,anchor=south west] (image) at (0,0) {\includegraphics[page=285,trim={53.00bp 571.75bp 53.87bp 57.63bp},clip,width=0.92\textwidth,max height=0.38\textheight]{Ferziger4th.pdf}};
\begin{scope}[x={(image.south east)},y={(image.north west)}]
\fill[white] (0.41194,0.73069) rectangle (0.44644,0.96736);
\fill[white] (0.65862,0.60773) rectangle (0.70150,0.82807);
\fill[white] (0.71411,0.56393) rectangle (0.74385,0.80060);
\fill[white] (0.66403,0.43063) rectangle (0.70150,0.65125);
\fill[white] (0.71411,0.38683) rectangle (0.74385,0.62350);
\fill[white] (0.40006,0.04761) rectangle (0.44644,0.28428);
\fill[white] (0.66403,0.07644) rectangle (0.70150,0.29706);
\fill[white] (0.66403,0.25354) rectangle (0.70150,0.47416);
\fill[white] (0.71411,0.03264) rectangle (0.74385,0.26931);
\fill[white] (0.71411,0.20974) rectangle (0.74385,0.44641);
\fill[white] (0.00096,0.00000) rectangle (0.01708,0.16023);
\fill[white] (0.02066,0.00000) rectangle (0.04532,0.16621);
\fill[white] (0.04758,0.00000) rectangle (0.14439,0.16621);
\fill[white] (0.14602,0.00000) rectangle (0.21729,0.15778);
\fill[white] (0.21886,0.00000) rectangle (0.26165,0.15778);
\end{scope}
\end{tikzpicture}
\end{sourceblock}
y = 0.90763）在两个振荡周期期间：解对时间步长大小（上；IFSM）和最大时间步使用的方案（下；每个振荡周期 62.5 个时间步）的依赖性

\begin{sourceblock}
\begin{tikzpicture}
\node[inner sep=0,anchor=south west] (image) at (0,0) {\includegraphics[page=285,trim={53.00bp 521.91bp 53.87bp 133.00bp},clip,width=0.92\textwidth,max height=0.38\textheight]{Ferziger4th.pdf}};
\begin{scope}[x={(image.south east)},y={(image.north west)}]
\fill[white] (0.00000,0.50935) rectangle (0.04093,1.00000);
\fill[white] (0.04250,0.50935) rectangle (0.07275,1.00000);
\fill[white] (0.07432,0.50935) rectangle (0.11026,1.00000);
\fill[white] (0.11182,0.50935) rectangle (0.19299,1.00000);
\fill[white] (0.19456,0.50935) rectangle (0.24605,1.00000);
\fill[white] (0.24761,0.50935) rectangle (0.29771,1.00000);
\fill[white] (0.35750,0.10686) rectangle (0.38722,0.78183);
\fill[white] (0.39296,0.11843) rectangle (0.44644,0.89314);
\fill[white] (0.00000,0.00000) rectangle (0.03528,0.51915);
\fill[white] (0.03681,0.00000) rectangle (0.10908,0.51915);
\fill[white] (0.11065,0.00000) rectangle (0.16075,0.51915);
\fill[white] (0.16229,0.00000) rectangle (0.20812,0.51915);
\fill[white] (0.20965,0.00000) rectangle (0.28737,0.51915);
\end{scope}
\end{tikzpicture}
\end{sourceblock}
\begin{sourceblock}
\begin{tikzpicture}
\node[inner sep=0,anchor=south west] (image) at (0,0) {\includegraphics[page=285,trim={53.00bp 496.90bp 53.87bp 158.14bp},clip,width=0.92\textwidth,max height=0.38\textheight]{Ferziger4th.pdf}};
\begin{scope}[x={(image.south east)},y={(image.north west)}]
\fill[white] (0.00000,0.07568) rectangle (0.11026,0.98288);
\fill[white] (0.11179,0.07568) rectangle (0.19014,0.98288);
\fill[white] (0.40481,0.10811) rectangle (0.44644,0.89189);
\end{scope}
\end{tikzpicture}
\end{sourceblock}
\begin{sourceblock}
\begin{tikzpicture}
\node[inner sep=0,anchor=south west] (image) at (0,0) {\includegraphics[page=285,trim={53.00bp 418.24bp 53.87bp 183.26bp},clip,width=0.92\textwidth,max height=0.38\textheight]{Ferziger4th.pdf}};
\begin{scope}[x={(image.south east)},y={(image.north west)}]
\fill[white] (0.39296,0.84684) rectangle (0.44644,0.98144);
\fill[white] (0.43639,0.73004) rectangle (0.47681,0.86463);
\fill[white] (0.56725,0.73004) rectangle (0.60767,0.86463);
\fill[white] (0.69820,0.73004) rectangle (0.73865,0.86463);
\fill[white] (0.82908,0.73004) rectangle (0.86950,0.86463);
\fill[white] (0.95994,0.73004) rectangle (1.00000,0.86463);
\fill[white] (0.41131,0.40563) rectangle (0.44565,0.53976);
\fill[white] (0.65546,0.33586) rectangle (0.73720,0.46086);
\fill[white] (0.73780,0.33586) rectangle (0.76102,0.46086);
\fill[white] (0.63380,0.23561) rectangle (0.71552,0.36061);
\fill[white] (0.71615,0.23561) rectangle (0.76102,0.36061);
\fill[white] (0.65874,0.13537) rectangle (0.71552,0.26037);
\fill[white] (0.71615,0.13537) rectangle (0.76102,0.26037);
\fill[white] (0.39946,0.01856) rectangle (0.44565,0.15285);
\end{scope}
\end{tikzpicture}
\end{sourceblock}
\begin{sourceblock}
\begin{tikzpicture}
\node[inner sep=0,anchor=south west] (image) at (0,0) {\includegraphics[page=285,trim={53.00bp 368.19bp 53.87bp 261.87bp},clip,width=0.92\textwidth,max height=0.38\textheight]{Ferziger4th.pdf}};
\begin{scope}[x={(image.south east)},y={(image.north west)}]
\fill[white] (0.42902,0.72644) rectangle (0.44565,0.96674);
\fill[white] (0.35871,0.41048) rectangle (0.38833,0.61973);
\fill[white] (0.39239,0.03326) rectangle (0.44565,0.27356);
\end{scope}
\end{tikzpicture}
\end{sourceblock}
\begin{sourceblock}
\begin{tikzpicture}
\node[inner sep=0,anchor=south west] (image) at (0,0) {\includegraphics[page=285,trim={53.00bp 343.14bp 53.87bp 311.93bp},clip,width=0.92\textwidth,max height=0.38\textheight]{Ferziger4th.pdf}};
\begin{scope}[x={(image.south east)},y={(image.north west)}]
\fill[white] (0.40421,0.10840) rectangle (0.44565,0.89160);
\end{scope}
\end{tikzpicture}
\end{sourceblock}
\begin{sourceblock}
\begin{tikzpicture}
\node[inner sep=0,anchor=south west] (image) at (0,0) {\includegraphics[page=285,trim={53.00bp 297.27bp 53.87bp 336.94bp},clip,width=0.92\textwidth,max height=0.38\textheight]{Ferziger4th.pdf}};
\begin{scope}[x={(image.south east)},y={(image.north west)}]
\fill[white] (0.39239,0.69089) rectangle (0.44565,0.96242);
\fill[white] (0.43564,0.43783) rectangle (0.47594,0.70936);
\fill[white] (0.56598,0.43783) rectangle (0.60629,0.70936);
\fill[white] (0.69642,0.43783) rectangle (0.73669,0.70936);
\fill[white] (0.82677,0.43783) rectangle (0.86707,0.70936);
\fill[white] (0.95711,0.43783) rectangle (0.99741,0.70936);
\end{scope}
\end{tikzpicture}
\end{sourceblock}
振荡，但 t 时刻和 t + P 时刻的解是相同的。在一个时期内，流动模式会发生复杂的变化，这使得其预测成为一项艰巨的任务。

请注意，根据零热通量边界条件的要求，等温线始终与绝热顶部和底部边界保持正交。等温线之间间距的变化显示了热通量的局部变化：密集的等温线表示垂直于等温线的方向上存在高温度梯度。在热壁温度和冷壁温度恒定的稳态流动的情况下，两个等温壁的净热通量相同。这里的热通量随时间变化，并且在两个壁处不同，因为流体在一段时间内积聚热量，并在剩余部分中散发热量。通常情况下，热量通过热壁进入，但由于热壁温度的振荡变化，有一个短时期，沿着热壁流动的流体实际上比壁本身更热，因此净热通量改变符号——腔体通过两个壁散失热量。从图 8.24 中可以看出这一点，该图显示了两个时期内通过热壁的总热通量的变化：在短时间内，热通量为正，表明流出通量。还要注意，通过热壁的最大热通量发生在温度达到最大值之前（发生在四分之一的时间）

\begin{sourceblock}
\includegraphics[page=286,trim={53.00bp 266.25bp 53.87bp 52.00bp},clip,width=\textwidth,max height=0.38\textheight]{Ferziger4th.pdf}
\par\vspace{0.45\baselineskip}
\small 图8.22 四分之一周期、二分之一周期、四分之三周期和全周期时具有振荡热壁温度的浮力驱动腔流中的速度矢量（分别从左到右、从上到下）
\end{sourceblock}
振荡周期，此处为 t = 182.5 且 t = 192.5 s)：速度场和温度场变化之间存在相移。

\begin{center}\small 图 8.24 还表明，二阶时间积分方案已经产生了一个非常精确的解，每个振荡周期有 62.5 个时间步长：图中无法区分对应于三个不同时间步长的三条曲线。这在图 8.25 中得到了证实，图 8.25 显示了一个监测位置在两个周期内的 u x 速度，使用三种方法的最大时间步长进行计算。可以看出，一阶隐式欧拉格式与二阶三时间级格式得到的解存在显著差异：前者由于一阶误差导致数值在时间上扩散，出现了峰值，而没有达到最低值。 SIMPLE 和使用二阶方案的 IFSM 之间的差异可以忽略不计；对于较小的时间步长，它会变得更小。\end{center}
\begin{sourceblock}
\includegraphics[page=287,trim={53.00bp 276.25bp 53.87bp 52.00bp},clip,width=\textwidth,max height=0.38\textheight]{Ferziger4th.pdf}
\par\vspace{0.45\baselineskip}
\small 图8.23 四分之一周期、二分之一周期、四分之三周期和全周期时具有振荡热壁温度的浮力驱动腔流等温线（分别从左到右、从上到下）
\end{sourceblock}
\begin{center}\small 图 8.24 在两个振荡周期内通过热壁的热通量作为时间的函数，使用 IFSM 和具有三个不同时间步长的三时间级二阶方案计算（指示为每个振荡周期的时间步数）\end{center}
0.1

62.5 t/ P 125 t/ P 250 t/ P
-0.1
-0.2
-0.3
Q

-0.4
-0.5
-0.6
-0.7
-0.8
180 185 190 195 200
t

\begin{sourceblock}
\begin{tikzpicture}
\node[inner sep=0,anchor=south west] (image) at (0,0) {\includegraphics[page=288,trim={53.00bp 572.04bp 53.87bp 57.71bp},clip,width=0.92\textwidth,max height=0.38\textheight]{Ferziger4th.pdf}};
\begin{scope}[x={(image.south east)},y={(image.north west)}]
\fill[white] (0.39844,0.72685) rectangle (0.44502,0.96702);
\fill[white] (0.81059,0.60181) rectangle (0.88598,0.80929);
\fill[white] (0.88617,0.60181) rectangle (0.90788,0.80929);
\fill[white] (0.79567,0.42209) rectangle (0.87110,0.62957);
\fill[white] (0.87125,0.42209) rectangle (0.90788,0.62957);
\fill[white] (0.81856,0.24210) rectangle (0.87110,0.44957);
\fill[white] (0.87125,0.24210) rectangle (0.90788,0.44957);
\fill[white] (0.41038,0.03298) rectangle (0.44502,0.27315);
\end{scope}
\end{tikzpicture}
\end{sourceblock}
\begin{sourceblock}
\begin{tikzpicture}
\node[inner sep=0,anchor=south west] (image) at (0,0) {\includegraphics[page=288,trim={53.00bp 546.75bp 53.87bp 108.25bp},clip,width=0.92\textwidth,max height=0.38\textheight]{Ferziger4th.pdf}};
\begin{scope}[x={(image.south east)},y={(image.north west)}]
\fill[white] (0.00000,0.07271) rectangle (0.12433,0.97666);
\fill[white] (0.12592,0.07271) rectangle (0.15194,0.97666);
\fill[white] (0.15350,0.07271) rectangle (0.24039,0.97666);
\fill[white] (0.24192,0.07271) rectangle (0.27221,0.97666);
\fill[white] (0.00000,0.00000) rectangle (0.03528,0.08169);
\fill[white] (0.03681,0.00000) rectangle (0.20538,0.08169);
\fill[white] (0.20689,0.00000) rectangle (0.28809,0.08169);
\end{scope}
\end{tikzpicture}
\end{sourceblock}
\begin{sourceblock}
\begin{tikzpicture}
\node[inner sep=0,anchor=south west] (image) at (0,0) {\includegraphics[page=288,trim={53.00bp 521.50bp 53.87bp 108.25bp},clip,width=0.92\textwidth,max height=0.38\textheight]{Ferziger4th.pdf}};
\begin{scope}[x={(image.south east)},y={(image.north west)}]
\fill[white] (0.00000,0.71613) rectangle (0.12433,0.99286);
\fill[white] (0.12592,0.71613) rectangle (0.15194,0.99286);
\fill[white] (0.15350,0.71613) rectangle (0.24039,0.99286);
\fill[white] (0.24192,0.71613) rectangle (0.27221,0.99286);
\fill[white] (0.00000,0.44215) rectangle (0.03528,0.71888);
\fill[white] (0.03681,0.44215) rectangle (0.20538,0.71888);
\fill[white] (0.20689,0.44215) rectangle (0.28809,0.71888);
\fill[white] (0.00000,0.16845) rectangle (0.05083,0.44518);
\fill[white] (0.05239,0.16845) rectangle (0.10250,0.44518);
\fill[white] (0.10403,0.16845) rectangle (0.13997,0.44518);
\fill[white] (0.14153,0.16845) rectangle (0.21380,0.44518);
\fill[white] (0.21537,0.16845) rectangle (0.26544,0.44518);
\fill[white] (0.35750,0.30613) rectangle (0.38737,0.51553);
\fill[white] (0.41038,0.03298) rectangle (0.44502,0.27315);
\fill[white] (0.00000,0.00000) rectangle (0.04517,0.17148);
\fill[white] (0.04671,0.00000) rectangle (0.10454,0.17148);
\fill[white] (0.10614,0.00000) rectangle (0.15624,0.17148);
\fill[white] (0.15777,0.00000) rectangle (0.21350,0.17148);
\fill[white] (0.21507,0.00000) rectangle (0.25239,0.17148);
\end{scope}
\end{tikzpicture}
\end{sourceblock}
\begin{sourceblock}
\begin{tikzpicture}
\node[inner sep=0,anchor=south west] (image) at (0,0) {\includegraphics[page=288,trim={53.00bp 496.22bp 53.87bp 158.78bp},clip,width=0.92\textwidth,max height=0.38\textheight]{Ferziger4th.pdf}};
\begin{scope}[x={(image.south east)},y={(image.north west)}]
\fill[white] (0.39844,0.10772) rectangle (0.44502,0.89228);
\end{scope}
\end{tikzpicture}
\end{sourceblock}
\begin{sourceblock}
\begin{tikzpicture}
\node[inner sep=0,anchor=south west] (image) at (0,0) {\includegraphics[page=288,trim={53.00bp 454.26bp 53.87bp 182.65bp},clip,width=0.92\textwidth,max height=0.38\textheight]{Ferziger4th.pdf}};
\begin{scope}[x={(image.south east)},y={(image.north west)}]
\fill[white] (0.41038,0.65994) rectangle (0.44502,0.95895);
\fill[white] (0.43329,0.40096) rectangle (0.47392,0.69997);
\fill[white] (0.56487,0.40096) rectangle (0.60550,0.69997);
\fill[white] (0.69657,0.40096) rectangle (0.73717,0.69997);
\fill[white] (0.82815,0.40096) rectangle (0.86878,0.69997);
\fill[white] (0.95973,0.40096) rectangle (1.00000,0.69997);
\end{scope}
\end{tikzpicture}
\end{sourceblock}
到目前为止，我们只处理了矩形解域和笛卡尔网格。真正的工程问题很少如此简单。在大多数情况下，流动和传热发生在相当复杂的几何形状中。我们将在下一章中解释如何将迄今为止描述的方法扩展到非笛卡尔结构化或任意多面体非结构化网格，这可以应用于更复杂的解域形状。

\chapter{复杂几何}
工程实践中的大多数流程涉及复杂的几何形状，这些几何形状不容易与笛卡尔网格拟合。尽管前面描述的代数系统的离散化原理和求解方法仍然有效，但它们的实现有许多不同的可能性。求解算法的属性取决于网格、向量和张量分量的选择，以及网格上变量的排列。本章将讨论这些问题。

\section*{9.1 网格的选择}\phantomsection\addcontentsline{toc}{section}{9.1 网格的选择}
当几何形状是规则的（例如，矩形或圆形）时，选择网格很简单：网格线通常遵循边界定义的方向，该方向与适当的坐标方向一致。在复杂的几何形状中，选择绝非易事。网格受到离散化方法施加的约束。如果算法是针对曲线正交网格设计的，则不能使用非正交网格；如果要求CV是四边形或六面体，则不能使用三角形和四面体组成的网格等。当几何形状复杂且无法满足约束时，必须做出妥协。然而，如果离散化和求解方法都是针对任意多面体控制体积设计的，则可以使用任何网格类型。

\subsection*{9.1.1 曲线边界的逐步逼近}\phantomsection\addcontentsline{toc}{subsection}{9.1.1 曲线边界的逐步逼近}
最简单的计算方法是使用正交网格（笛卡尔或极圆柱）的方法。然而，为了将这样的网格应用于具有倾斜或弯曲边界的解域，需要额外的近似：要么必须通过类似楼梯的步骤来近似边界，要么网格延伸

\begin{sourceblock}
\includegraphics[page=290,trim={53.00bp 444.25bp 53.87bp 52.00bp},clip,width=\textwidth,max height=0.38\textheight]{Ferziger4th.pdf}
\par\vspace{0.45\baselineskip}
\small 图 9.1 使用弯曲壁边界的逐步近似的网格示例，并进行局部细化以减少台阶高度
\end{sourceblock}
超出边界及其附近的单元格必须使用特殊的近似值。

第一种方法有时仍然被使用，但它会带来两类问题：

\begin{itemize}
\item 每条网格线的网格点（或CV）数量不是恒定的，因为它是在完全规则的网格中。因此，尽管单元是规则的，但由于必须使用某种间接寻址或限制每条网格线上索引范围的特殊数组，规则网格的主要优点就丧失了。对于每个新问题，可能需要更改计算机代码。
\item 逐步逼近光滑墙会在解中引入误差，尤其是当网格较粗糙时。阶梯壁边界条件的处理也需要特别注意，因为阶梯近似影响的面积比阶梯的直接面积大得多。
\end{itemize}
这种网格的一个例子如图 9.1 所示。当壁面的剪切力或传热发挥重要作用时（例如机翼、飞机或船体、空气动力体、涡轮叶片等），这种方法可能会导致较大的误差。不建议这样做，除非求解算法允许在壁附近进行局部网格细化（有关局部网格细化方法的详细信息，请参阅第 12 章）并且压力占主导地位。对于钝体尤其如此，并且当许多物体占据的空间中的阻塞效应对于流动分析比单个物体上的剪切应力更重要时。这可能是土木和环境工程中的情况（例如，当研究建筑物内部和周围、河流和湖泊、大气等中的流动时）。理论上，如果墙壁上的网格间距与自然墙壁粗糙度相当，则逐步近似将产生准确的结果。

使用此方法的另一个可能原因是，当包含许多物理模型（例如，燃烧、多相流、相变等）的现有求解方法无法适应更好地拟合边界的网格时，或者当壁边界的精确表示并不重要时（例如，当表面非常粗糙时）。事实上，一些作者使用局部细化的逐步近似来模拟壁粗糙度效应。例如，Manhart 和 Wengle (1994) 对壁挂式半球上流动的大涡流模拟以及 Xing-Kaeding (2006) 对圆柱体进出水的模拟。

如上所述，局部细化的使用消除了笛卡尔网格最吸引人的特征之一，即邻域连接的简单性。要么必须将细化界面处的非细化单元视为多面体，要么必须使用所谓的“幽灵节点”（也称为“悬挂节点”）。这些的处理方式与非共形或滑块接口相同，这将在第 1 节中详细解释。 9.6.1.

\subsection*{9.1.2 浸入边界方法}\phantomsection\addcontentsline{toc}{subsection}{9.1.2 浸入边界方法}
浸入边界方法使用规则（主要是笛卡尔）网格，并通过部分阻挡被墙壁切割的单元来处理不规则的墙壁边界。这种方法的首次发表被认为是 Peskin (1972) 的论文；近年来，已经开发了许多浸入边界方法的变体。动机之一是研究移动或变形物体周围的流动，这种方法被认为比移动和变形的贴身网格更灵活。

该方法有多种风格，并且由于细节仅与壁边界条件的处理有关，而该方法的其余部分通常对应于此处描述的方法之一，因此我们不会深入讨论这些细节。基本上，首先必须识别被墙边界（也可能是它们的直接邻居）切割的单元。为此，通常使用边界的三角测量（立体光刻，STL）描述，而不是使用 CAD 表示。最大的问题是在正确的位置强制执行正确的墙边界条件；这导致将速度固定在切割单元中或主体内部的虚影单元中，使得适合一定数量节点的轮廓在穿过壁边界时给出指定的壁速度。当墙壁移动时，必须在每个时间步骤执行切割单元和解域外部单元的识别。

当网格足够细时，这些方法会产生准确的结果。然而，如果研究机翼、涡轮叶片或其他平滑弯曲壁周围的高雷诺数流动，沿壁具有棱镜层的边界拟合网格会更准确。特别是当使用具有所谓“低雷诺数”公式的湍流模型时，需要在壁法线方向上非常小的网格间距，浸入边界方法变得效率低下，因为它们最终也会在壁切向方向上细化网格。

有关浸入边界方法的更多详细信息，请参阅 Peskin (2002) 的评论文章和其他出版物，例如 Tseng 和 Ferziger (2003)、Mittal 和 Iaccarino
( 2005 )、Taira 和 Colonius ( 2007 )、Lundquist 等人。 ( 2012 ) 等。论文中有非常详细的描述，例如 Peller ( 2010 ) 和 Hylla ( 2013 ) 等。

\subsection*{9.1.3 重叠网格}\phantomsection\addcontentsline{toc}{subsection}{9.1.3 重叠网格}
当几何形状复杂时，结构化或块结构、边界拟合网格的问题是，网格质量通常在靠近墙壁的地方恶化，而墙壁附近的精度应该是最高的。通过使用最佳边界拟合、靠近墙壁的棱柱网格和远离墙壁的简单（通常是笛卡尔）网格，重叠网格可以帮助实现靠近墙壁的高网格质量。这类方法在文献中经常与奇美拉网格（奇美拉是一种具有狮头、山羊身体和蛇尾的神话生物）、重叠网格或复合网格等名称联系在一起。当研究物体周围的流动时，无论是在固定位置还是在移动，这种方法尤其有吸引力。通常，人们会创建一个不考虑实体的背景网格，该背景网格仅适合通常不移动的外部边界（入口、出口、对称或远场边界等）。围绕主体创建的网格距主体一定距离，并且与背景网格重叠。如果物体在移动，附着在其上的网格也会移动（请参阅 NASA Chimera Grid Tools 用户手册 NASA CGTUM 2010）。

被主体覆盖的背景网格部分和重叠网格被停用，除了两个网格都处于活动状态的狭窄重叠区域之外。离散化和求解方法通常是此处描述的方法之一；只有单个网格上解的耦合是特定于重叠网格的。一种可能性是将重叠网格的外表面和通过停用背景网格中的一些单元而产生的孔的表面视为边界。在流进入网格的边界部分，规定了入口（狄利克雷）条件，而在流离开网格的部分，规定了压力。要施加的边界值是通过从另一个网格插值获得的。这些边界条件的更新可以在每次内迭代之后或每次外迭代之后进行；前者更加隐式（如在并行计算中）并导致更好的收敛特性。

另一种选择是为每个求解的方程将网格耦合到线性方程组的矩阵 A 中，以便在所有网格上同时获得解。如何做到这一点将参考图 9.2 及其中的符号进行解释。该图显示了两个重叠网格上的活动单元格。定义接受器（幽灵、悬挂）节点，以便可以执行离散化，就好像相邻小区也处于活动状态一样。受体细胞的变量值是通过从另一个网格中选定的一组供体细胞插值获得的。这如图 9.2 所示：标记为 P 的单元的离散方程是指通过通过 P 和 Na 之间的面的通量以空心圆标记的受体相邻单元 N a。无论在离散通量中何处引用受体单元中心的变量值，它都会被插值表达式替换：

\begin{center}\small 图9.2 重叠网格上受体细胞和供体细胞的定义\end{center}
N N
N N
a

N

P

N N
1 2

\begin{sourceblock}
\includegraphics[page=293,trim={53.00bp 424.21bp 53.87bp 225.43bp},clip,width=0.92\textwidth,max height=0.38\textheight]{Ferziger4th.pdf}
\end{sourceblock}
这里 α 4 、α 5 和 α 6 是加起来为 1 的插值因子，通过将线性形状函数拟合到另一个网格的三个最近邻居来获得。因此，在单元格 P 的方程中，线性方程矩阵中将有 6 个非对角系数：三个邻居来自同一网格（N 1 、N 2 和 N 3 ），三个邻居来自重叠网格（N 4 、N 5 和 N 6 ）。根据所使用的线性方程求解器，可能需要对算法进行一些其他修改。此外，供体细胞之间的插值和要使用的供体细胞的数量有许多不同的选项。一种可能性也是仅使用变量值和来自最近邻居的梯度。然后需要采用延迟校正方法来考虑梯度的贡献，但这并不代表一个大的复杂性。

重叠网格的主要优点是：

\begin{itemize}
\item 与使用单个网格相比，可以更好地优化靠近墙壁的网格质量；
\item 如果身体移动，靠近墙壁的网格质量将保持不变；
\item 无需更改网格或边界条件即可轻松实现不同攻角物体的参数化研究——只需旋转重叠网格即可。
\item 通过重叠网格，可以考虑任意身体运动，而这在其他方法中是不可能实现的。
\end{itemize}
请注意，对于重叠网格，特殊插值是在远离墙壁、变量变化不太强的区域执行的，而浸没边界方法直接在需要最高精度的墙壁上应用特殊插值。

过去经常使用重叠网格，最近此功能也可在商业代码中使用。自1992年以来，每两年举办一次关于重叠网格的专门研讨会；参见官方网站：www. oversetgridsymposium.org 了解更多信息。一些方法的详细描述可以在 Hadži ́c (2005) 和 Hanaoka (2013) 的论文中找到；这两篇论文都可以从互联网上下载。重叠网格应用的示例将在下列各节中介绍。 13.4 和 13.10.1。

\subsection*{9.1.4 边界拟合非正交网格}\phantomsection\addcontentsline{toc}{subsection}{9.1.4 边界拟合非正交网格}
边界拟合非正交网格最常用于计算复杂几何形状中的流量（大多数商业代码使用此类网格）。它们可以是结构化的、块结构的或非结构化的。这种网格的优点是（1）它们可以适应任何几何形状，（2）比正交曲线网格更容易实现最佳属性。由于 CV 面落在解域边界上，因此与曲线边界的逐步近似相比，边界条件更容易实现。网格还可以适应流动，即可以选择一组网格线来跟随流线（这提高了精度），并且在变化强烈的区域可以使间距更小，特别是如果使用块结构或非结构化网格。

非正交网格有几个缺点。如果使用有限差分方法，则变换后的方程包含更多需要近似的项，从而增加了编程的难度和每个单元求解方程的计算量。如果使用有限体积方法，则必须创建质量可接受的控制体积，这是一项艰巨的任务。网格的非正交性在某些情况下可能会导致收敛问题或非物理解。速度分量的选择以及网格上变量的排列会影响算法的准确性和效率。这些问题将在下面进一步讨论。

在本书的其余部分中，我们将假设网格是非正交且非结构化的。我们将介绍的离散化原理和求解方法也适用于正交网格，因为它们可以被视为非正交网格的特殊情况。其中一节讨论块结构网格的处理，因为相同的方法也用于非结构网格中的非共形界面，例如滑动界面。

\section*{9.2 网格生成}\phantomsection\addcontentsline{toc}{section}{9.2 网格生成}
复杂几何形状的网格生成是一个需要太多空间才能在此详细处理的问题。我们将仅介绍一些基本思想和网格应具有的属性。有关网格生成的各种方法的更多详细信息可以在专门讨论该主题的书籍和会议记录中找到，例如 Thompson 等人。 ( 1985 ) 和 Arcilla 等人。 （1991）。

尽管必要性要求在复杂的几何形状中网格是非正交的，但使其尽可能接近正交也很重要。在 FV 方法中，CV 顶点处网格线的正交性并不重要，重要的是单元面表面法向量与连接其两侧 CV 中心的线之间的角度。因此，由等边三角形组成的二维网格相当于正交网格，因为连接单元中心的线与单元面正交。这将在第二节中进一步讨论。 9.7.2.

单元拓扑也很重要。如果使用中点法则积分逼近、线性插值和中心差分来离散方程，那么如果CV在2D中为四边形，在3D中为六面体，则分别比使用三角形和四面体的精度更高。原因是，在四边形和六面体 CV 上离散化扩散项时，在相对单元面上产生的部分误差会部分抵消（如果单元面平行且面积相等，则它们会完全抵消）。为了在三角形和四面体上获得相同的精度，必须使用更复杂的插值和梯度近似。特别是在固体边界附近，最好有四边形或六面体（棱镜层；参见第 9.2.3 节等），因为所有量在那里变化很大，并且在此区域中精度尤其重要。

如果一组网格线紧密遵循流动的流线，尤其是对流项，精度也会提高。如果使用三角形或四面体则无法实现这一点，但使用四边形和六面体则可以实现。

当处理复杂的几何形状时，非均匀网格是规则而不是例外。然而，网格质量成为一个重要问题，我们在第 1 章中专门讨论了这一问题。 12.

有经验的用户可能知道哪里会出现速度、压力、温度等的强烈变化；这些区域的网格应该很好，因为那里的误差很可能很大。然而，即使是经验丰富的用户也会偶尔遇到意外，并且无论如何更复杂的方法都是有用的。错误在整个域中进行对流和扩散，如第 1 节中所述。 3.9，因此必须尽可能实现截断误差的均匀分布。然而，可以从粗网格开始，然后根据离散化误差的估计对其进行局部细化；执行此操作的方法称为解自适应网格方法，将在第 1 章中进行描述。 12.

最后，还有网格生成的问题。当几何形状很复杂时，此任务通常会消耗迄今为止最多的用户时间；设计人员花费一周时间生成一个网格并不罕见，而解在并行计算机上只需要几个小时。由于解的精度与用于方程离散化的近似值一样（如果不是更多）取决于网格质量，因此网格优化是一项值得投入的时间。

存在许多用于电网生成的商业和公共领域代码。网格生成过程的自动化，旨在减少用户时间并加快过程，是该领域的主要目标。有时，生成一组高质量的重叠网格比单个块结构化或非结构化网格更容易，但在真正复杂的几何形状中，由于存在太多不规则块并且需要使用太多网格块，因此这种方法很困难。然而，人们可以在出版物中找到使用超过一百个网格块的示例（例如，用于计算军用飞机周围的流量）。

三角形和四面体网格的生成更容易自动化，这是它们受欢迎的原因之一。然而，在大多数商业代码中，使用修剪的六面体或多面体网格。生成多面体网格时，第一步通常是生成四面体网格，然后将其转换为多面体网格。工业应用中的网格生成过程涉及其他几个步骤，将在以下小节中进行描述。

\subsection*{9.2.1 流域定义}\phantomsection\addcontentsline{toc}{subsection}{9.2.1 流域定义}
对于自动网格生成，需要解域边界的闭合表面作为起点。在大多数情况下，曲面是以一种 CAD 格式定义的。然而，人们经常获取实体零件的 CAD 数据，并需要提取实体内部、外部或之间的流体所占据的体积。大多数用于网格生成的商业软件包都包含允许对实体进行布尔运算和提取流量的工具；有时需要额外的步骤，例如为入口和出口创建边界表面。图 9.3 显示了一个示例：CAD 数据包含球阀、手柄和外壳。如果只计算阀门周围的流量，则需要提取阀门内部流体所占的体积，并通过提供入口和出口管道横截面来关闭阀门；然后可以丢弃固体部分。还希望将出口边界移至阀门的更下游；如果CAD模型中不包含出口管道，则可以沿轴向拉伸出口边界，从而创建足够的管道截面，如图9.3所示。

如果对从管道中的热流体到环境的热传递感兴趣，则可以保留固体，并在阀门组件周围添加一个更大的盒子，以便计算从外壁到环境的热传递。这可以通过从盒子中减去阀门组件来实现。然后需要在流体域和所有固体部分内生成网格。固液界面处最好有共形网格；这将在第 1 章中更详细地讨论。 12 讨论网格质量时。

在包含数百甚至数千个固体部件的非常复杂的几何形状的情况下（例如，在研究车辆的热管理时，需要计算车辆周围、发动机舱内、客舱内等的流动以及通过固体部件的热传导），通过布尔运算和固体压印创建流域边界可能太复杂。当零件装配包含需要手动修复的间隙、交叉或重叠表面时，或者当导入的几何体包含太多对流程不重要的细节并且包含这些细节会不必要地增加网格中的单元数时，尤其如此。在这种情况下，所谓的“表面包裹”工具可能是首选。这些工具将表面包裹在所有固体部件周围，就像吹气球一样，直到它紧紧地粘在所有墙壁上。所得的闭合表面以离散形式（通常是三角测量）提供。

\begin{sourceblock}
\includegraphics[page=297,trim={53.00bp 364.25bp 53.87bp 52.00bp},clip,width=\textwidth,max height=0.38\textheight]{Ferziger4th.pdf}
\par\vspace{0.45\baselineskip}
\small 图 9.3 模拟阀门周围流动时流体体积提取和边界挤压的示例：导入的 CAD 零件（上）和带有挤压出口边界的提取流域（下）
\end{sourceblock}
表面包裹工具通常为用户提供控制生成的封闭表面的细节保真度的可能性（例如，保留锐边或材料界面等特征线、表面曲率的分辨率、保留零件之间的间隙等）。因此，可以对几何体进行“去特征化”，即去除不需要的细节（例如，封闭孔和间隙、去除小部件等）。通过严格的公差，可以获得几何形状的精确表示，而不会显著损失细节保真度。

\subsection*{9.2.2 表面网格的生成}\phantomsection\addcontentsline{toc}{subsection}{9.2.2 表面网格的生成}
通过对几何图形的 CAD 表示进行细分而获得的离散表面表示不适合作为网格生成的起点：三角形准确地表示了几何图形，但它们的形状和尺寸变化通常远远超过用于模拟流体流动的网格可接受的范围。因此，网格划分工具通常首先重新创建一个合适的表面网格，该网格充分代表几何形状，同时提供适度的增长率并尽可能接近最佳等边三角形（至少 Delaunay 条件应该

\begin{center}\small 图 9.4 图 9.3 中阀门流域 CAD 描述的初始细分示例（上）和通过三角测量重新划分表面网格，作为体积网格生成的起点（下）\end{center}
满足，即与每个三角形的顶点拟合的圆不应包含任何其他顶点——这是最大化所有三角形的最小角度的条件）。这里我们不讨论表面网格生成的细节；我们鼓励感兴趣的读者查阅有关该主题的文献（Frey 和 George 2008，第 7 章）。我们只想强调这样一个事实，即许多用于生成体积网格的算法需要由足够的三角形表面网格离散化的解域的封闭表面作为起点。图 9.4 显示了初始表面三角剖分和重新划分网格的表面的示例。

\subsection*{9.2.3 体网格的生成}\phantomsection\addcontentsline{toc}{subsection}{9.2.3 体网格的生成}
大多数用于生成四面体网格的工具都是从解域边界上的三角形表面网格开始的。首先生成表面有一个底面的四面体，然后向内继续该过程，形成行进前沿；整个过程有点像通过行进过程求解方程，实际上，某些方法是基于椭圆或双曲偏微分方程的解。

如果需要解析边界层，则四面体单元不适合靠近墙壁，因为第一个网格点必须非常靠近墙壁，而在平行于墙壁的方向上可以使用相对较大的网格尺寸。这些要求导致形成细长的四面体，从而在扩散通量的近似中产生问题。因此，大多数网格生成方法首先在实体边界附近生成棱柱层，从表面的三角形离散化开始并挤压

\begin{center}\small 图 9.5 解域剩余部分中由靠近墙壁的棱柱和四面体组成的网格示例\end{center}
与边界正交的方向上的三角形。通常，对棱镜层进行一些扩展；应避免大于 1.5 的因子，而高达 1.2 的值通常会提供良好的结果。在该层的顶部，在域的剩余部分中自动生成四面体网格。这种网格的一个例子如图 9.5 所示。应创建多少个棱镜层取决于边界层的建模；更多细节将在下一章中给出。

这种方法提高了靠近墙壁的网格质量，并导致更准确的解和数值求解方法更好的收敛性；但是，只有当求解方法允许混合控制体积类型时才能使用它。原则上，任何类型的方法（FD、FV、FE）都可以适用于这种网格。

自动网格生成的另一种方法是用笛卡尔网格覆盖解域，并调整域边界切割的单元以适应边界。这可以通过将落在解域之外的切割单元的顶点投影到边界，或按原样接受切割单元（多面体 CV）来实现。后一种方法要求离散化和求解方法可以处理任意多面体单元。这种方法的问题在于，会生成质量最差的近壁不规则单元，而实际上需要最高的网格质量。如果这是在非常粗略的水平上完成的，然后对网格进行多次细化，则不规则性仅限于少数位置，并且可能不会对精度产生太大影响；然而，复杂的几何形状通常会限制初始网格的粗细程度。

为了使不规则的修剪单元远离墙壁，可以首先沿着墙壁创建棱镜层；然后规则笛卡尔网格被近壁棱镜层的外表面切割。这种网格的一个例子如图 9.6 所示。这种方法允许快速生成网格，但需要一个求解器来处理通过用任意表面切割规则单元而创建的多面体单元。此外，网格质量也是一个问题，因为切割的单元可能很小而且形状相当不规则。较大单元之间的非常小的单元是不利的（膨胀率过高）。同样，原则上所有类型的方法都可以适用于这种类型的网格。

如果离散化和求解方法允许使用具有任意拓扑单元（一般多面体）的非结构化网格，则网格生成程序受到的约束很少。如今，大多数用于流动模拟的商业代码

\begin{center}\small 图 9.6 通过将靠近墙壁的棱柱层与大部分解域中的规则笛卡尔网格相结合而创建的修剪网格示例，以及沿着修剪表面的不规则多面体单元\end{center}
\begin{sourceblock}
\includegraphics[page=300,trim={53.00bp 470.77bp 53.87bp 140.71bp},clip,width=\textwidth,max height=0.38\textheight]{Ferziger4th.pdf}
\par\vspace{0.45\baselineskip}
\small 图 9.7 由靠近壁的多边形棱柱层和解域主体中的任意多面体组成的网格示例
\end{sourceblock}
可以生成和使用任意多面体网格。如上所述，此类网格可以从四面体网格创建；这通常是生成多面体网格的第一步。然而，由于多面体网格不会对拓扑施加任何限制（即单元面和相邻单元的数量是任意的），因此通常会遵循几个优化步骤。可以通过多种方式修改单个单元格或一组多面体单元格，以提高网格质量；这包括移动顶点、分割或合并面或单元等。多面体网格的示例如图 9.7 所示。

多面体单元也是通过单元局部网格细化创建的。未细化的相邻单元虽然保留其原始形状（例如六面体），但变成逻辑多面体，因为一个或多个面被一组子面替换。

求解域也可以首先被划分为多个块，在这些块中，与整个求解域的情况相比，可以更容易地创建具有良好属性的网格。人们可以自由地为每个块选择最佳的网格拓扑（结构化 H、O 或 C 网格，或非结构化四面体、六面体或多面体网格）。然后，块接口处的网格必须相互印记；位于界面中的原始面被几个较小的面所取代，这些较小的面被唯一地定义为界面两侧的两个相邻单元的公共表面。块界面上的单元具有不规则形状的面并且必须

\begin{sourceblock}
\includegraphics[page=301,trim={53.00bp 399.25bp 53.87bp 52.00bp},clip,width=\textwidth,max height=0.38\textheight]{Ferziger4th.pdf}
\par\vspace{0.45\baselineskip}
\small 图 9.8 通过组合两个具有非共形界面和单元格局部细化的块创建的网格；注意带有不规则细胞面的压印界面
\end{sourceblock}
被视为多面体。图 9.8 中可以看到一个例子，它显示了从具有矩形横截面的通道到具有圆形横截面的管道的突然过渡。显然，笛卡尔网格代表矩形通道的最佳选择，而边界拟合的 O 形网格最适合圆形管道。图 9.8 显示了块界面，其中两个网格相互压印，导致不规则的单元面。耦合非共形网格块的替代策略也是可能的；有些将在第 1 节中描述。 9.6.1.

如果要在参数研究（例如形状优化）中改变部分几何形状，则在复杂几何形状中按块生成非共形网格尤其有吸引力。在这种情况下，几何形状中未发生变化的部分的网格可以保持不变，并且仅需要在一小部分区域而不是在整个解域中重新生成网格。与生成整体网格相比，这种方法通常可以更好地优化网格属性。

如果求解方法要求所有 CV 都属于特定类型，则生成复杂几何形状的网格可能会很困难。在 CFD 的某些工业应用中，网格生成过程需要数天甚至数周的时间并不罕见。高效、自动的网格生成器是工业环境中任何 CFD 软件的重要组成部分。自动网格生成工具的存在可能是 CFD 在许多行业如此广泛传播的主要原因，因为这些工具可以实现整个过程的自动化，并将处理时间减少到可承受的范围。

\begin{center}\small 图 9.9 显示了复杂工业 CFD 应用的示例，其中修剪网格的剖面应用于车辆热工分析。\end{center}
\begin{sourceblock}
\includegraphics[page=302,trim={53.00bp 436.25bp 53.87bp 52.00bp},clip,width=\textwidth,max height=0.38\textheight]{Ferziger4th.pdf}
\par\vspace{0.45\baselineskip}
\small 图 9.9 一个经过工业应用局部细化的修剪笛卡尔网格的示例 - 车辆热管理分析；通过体积网格的水平截面。来源：戴姆勒股份公司
\end{sourceblock}
显示年龄。汽车工程中的很大一部分计算工作都花在确保车辆部件在运行中不会过热，因为这可能会导致故障。虽然汽车空气动力学更多地受到设计师的影响而不是工程师的影响，但发动机舱和车身底部是工程师可以在空间限制内自由优化形状和位置的区域。图 9.9 创建于 2014 年，表明了 CFD 工程应用中解域的几何复杂性。平均细胞数约为 20-5000 万个，所有几何相关细节都可以解析。通常，用户会为特定应用创建网格设计模板（包括局部网格细化），从而减少第一个网格精心准备后的后续分析的预处理时间。

\section*{9.3 速度分量的选择}\phantomsection\addcontentsline{toc}{section}{9.3 速度分量的选择}
在第 1 章中 我们讨论了与动量组成部分的选择相关的问题。如图 1.1 所示，只有在固定基础上选择分量才能得到完全保守形式的动量方程。为了确保动量守恒，最好使用这样的基，最简单的基是笛卡尔基。当流是三维的时，使用任何其他基础（例如，面向网格、协变或逆变）没有优势。只有当选择另一个向量基可以简化问题时（例如，通过减小问题的维度）才值得放弃使用笛卡尔分量。这种情况的一个例子是管道或其他轴对称几何形状中的流动。如果流量在圆周方向上不变化，则速度矢量在极柱基上仅具有两个非零分量，但具有三个非零笛卡尔分量。因此，该问题在笛卡尔分量方面具有三个因变量，但如果使用极柱分量，则只有两个因变量，这是一个实质性的简化。如果流量沿圆周方向变化，则无论如何，问题都是三维的。虽然在轴对称几何中使用极柱组件具有明显的优势，因为坐标与流动边界对齐，但笛卡尔选项更优越，因为方程采用强保守形式，并且我们已经在例如图 9.5 中展示了如何处理这种类型的几何。

\subsection*{9.3.1 面向网格的速度分量}\phantomsection\addcontentsline{toc}{subsection}{9.3.1 面向网格的速度分量}
如果使用面向网格的速度分量，则动量方程中会出现非保守源项。这些解释了组件之间动量的重新分配。例如，如果使用极柱分量，对流张量 ρ vv 的散度会导致两个这样的源项：

\begin{itemize}
\item 在r 分量的动量方程中，有一项ρ v 2 θ / r，它代表表观离心力。这不是在旋转坐标系（例如泵或涡轮通道）中分析的流动中发现的离心力，它完全是由于从笛卡尔坐标到极柱坐标的转换而产生的。该术语描述了由于 v θ 方向的变化而将 θ 动量转换为 r 动量。
\item 在θ 分量的动量方程中，有一项− ρ v r v θ / r，它代表表观科里奥利力。该项是 θ 动量的源或汇，具体取决于速度分量的符号。
\end{itemize}
动量方程中曲率项作用的示例如图 9.10 所示。它显示了即使速度场是均匀的，即速度矢量在整个域中相同，速度矢量的径向和切向分量如何沿网格线变化。速度分量的这种变化只是由于网格线方向的变化而发生，与流动物理学无关。虽然所有形式的控制方程在数学上都是等价的，但有些控制方程在尝试数值求解时显然会带来更多困难。很明显，如果使用极柱坐标或球坐标，数值误差会扰乱图 9.10 所示示例中速度场的均匀性，而使用笛卡尔分量时则很难维持精确解。

在一般曲线坐标中，有更多这样的源术语（参见 Sedov 1971 的书籍；Truesdell 1991 等）。它们涉及 Christoffel 符号（曲率项、高阶坐标导数），其离散化通常代表数值误差的重要来源。网格要求平滑——点与点之间网格方向的变化必须小且连续。特别是在非结构化网格上，网格线与坐标方向不相关，这种基础很难使用。

\begin{sourceblock}
\includegraphics[page=304,trim={53.00bp 459.25bp 53.87bp 52.00bp},clip,width=\textwidth,max height=0.38\textheight]{Ferziger4th.pdf}
\par\vspace{0.45\baselineskip}
\small 图 9.10 均匀速度场和极坐标网格的示例，显示径向和周向分量随网格中位置的变化
\end{sourceblock}
\subsection*{9.3.2 笛卡尔速度分量}\phantomsection\addcontentsline{toc}{subsection}{9.3.2 笛卡尔速度分量}
回想一下，我们专门使用笛卡尔向量和张量分量。如果使用其他组件，则离散化和求解技术保持不变，但有更多项需要近似。第 1 章给出了笛卡尔分量的守恒方程。 1.

如果使用FD方法，则只需对非正交坐标采用适当形式的散度和梯度算子（或者将关于笛卡尔坐标的所有导数变换到非正交坐标）。这会导致项数增加，但方程的守恒性质与笛卡尔坐标中的相同，如下所示。

在 FV 方法中不需要坐标变换。当逼近 CV 表面法线的梯度时，可以使用局部坐标变换或在与单元面正交的线上使用辅助节点，如下所示。

\section*{9.4 变量排列的选择}\phantomsection\addcontentsline{toc}{section}{9.4 变量排列的选择}
在第 7 章中 我们提到，除了并置变量排列之外，各种交错排列也是可能的。虽然笛卡尔网格中的一种或另一种没有明显的优势，但当使用非正交网格时情况会发生很大变化。

\subsection*{9.4.1 交错排列}\phantomsection\addcontentsline{toc}{subsection}{9.4.1 交错排列}
\begin{sourceblock}
\includegraphics[page=305,trim={53.00bp 495.25bp 53.87bp 52.00bp},clip,width=\textwidth,max height=0.38\textheight]{Ferziger4th.pdf}
\par\vspace{0.45\baselineskip}
\small 图9.11 非正交网格上的变量排列：a——逆变速度分量交错排列，b——笛卡尔速度分量交错排列，c——笛卡尔速度分量同位排列
\end{sourceblock}
\begin{center}\small 图 9.11 显示了此类网格的部分，其中网格线方向改变了 90°。在一种情况下，逆变速度分量和在另一种情况下笛卡尔速度分量显示在交错位置。回想一下，引入交错排列是为了实现速度和压力梯度之间的强耦合。目标是使垂直于单元面的速度分量位于该面两侧的压力节点之间，见图 8.1。对于逆变或协变面向网格的组件，这个目标也在非正交网格上实现，见图9.11a。对于笛卡尔分量，当网格线改变方向 90° 时，会出现如图 9.11b 所示的情况：存储在单元表面的速度分量对通过该面的质量通量没有贡献，因为它平行于该面。为了计算通过此类 CV 面的质量通量，必须使用周围单元面的插值速度。这使得压力校正方程的推导变得困难，并且不能确保速度和压力的正确耦合——两者都可能导致振荡。\end{center}
因为，在工程流程中，网格线经常会改变方向180°或更多，特别是如果使用非结构化网格，交错排列很难使用。如果所有笛卡尔分量都存储在每个 CV 面上，则可以克服其中一些问题。然而，这在 3D 中变得复杂，特别是如果允许任意形状的 CV。要了解如何做到这一点，感兴趣的读者可能想看看 Maliska 和 Raithby 的论文（1984）。

\subsection*{9.4.2 并置安排}\phantomsection\addcontentsline{toc}{subsection}{9.4.2 并置安排}
它已在第 1 章中展示。从图 7 可以看出，并置排列是最简单的一种，因为所有变量共享相同的 CV，但需要更多插值。当网格非正交时，它并不比其他排列复杂，如图 2 所示。 9.11 c.通过任何CV面的质量通量可以通过对面两侧的两个节点处的速度进行插值来计算；该过程与常规笛卡尔网格相同。大多数商业 CFD 代码使用笛卡尔速度分量和变量的并置排列。我们将重点关注这一安排。

接下来，我们将在前面章节中针对笛卡尔网格所做的工作的基础上，描述非正交网格离散化的新特征。

\section*{9.5 有限差分法}\phantomsection\addcontentsline{toc}{section}{9.5 有限差分法}
有限差分法在工程应用的CFD中并未得到广泛应用；然而，它们在一些地球物理应用中发挥了作用（参见第 13.7 节）。对于这些情况，可以使用第 1 章的方法。 3个很有用。一般来说，使用有限差分方法比有限体积方法更容易开发高阶方法（如许多地球物理问题中所需的那样），但当使用高达二阶的近似时，后者非常适合任意多面体控制体积。对于高于二阶的阶数，复杂性显著增加，本章后面的单独部分将对此进行说明。因此，我们在这里仅简要介绍非正交网格的 FD 方法。

\subsection*{9.5.1 基于坐标变换的方法}\phantomsection\addcontentsline{toc}{subsection}{9.5.1 基于坐标变换的方法}
FD 方法通常仅与结构化网格结合使用，在这种情况下，每条网格线都是一条坐标为 xi i 的线。坐标由变换 x i = x i (xi j ), j = 1 , 2 , 3 定义，其特征在于雅可比行列式 J：

\begin{sourceblock}
\begin{tikzpicture}
\node[inner sep=0,anchor=south west] (image) at (0,0) {\includegraphics[page=306,trim={53.00bp 185.68bp 53.87bp 399.35bp},clip,width=0.92\textwidth,max height=0.38\textheight]{Ferziger4th.pdf}};
\begin{scope}[x={(image.south east)},y={(image.north west)}]
\fill[white] (0.00000,0.93441) rectangle (0.05573,1.00000);
\fill[white] (0.06123,0.93441) rectangle (0.09128,1.00000);
\fill[white] (0.64087,0.50820) rectangle (0.68725,0.66502);
\fill[white] (0.39874,0.34743) rectangle (0.43245,0.48995);
\fill[white] (0.43110,0.33411) rectangle (0.44171,0.43977);
\fill[white] (0.64202,0.33485) rectangle (0.68623,0.49168);
\fill[white] (0.52217,0.18802) rectangle (0.56857,0.34484);
\fill[white] (0.58153,0.18802) rectangle (0.62794,0.34484);
\fill[white] (0.52328,0.01479) rectangle (0.56749,0.17150);
\fill[white] (0.58259,0.01479) rectangle (0.62680,0.17150);
\end{scope}
\end{tikzpicture}
\end{sourceblock}
因为我们使用的是笛卡尔向量分量，所以我们只需要将笛卡尔坐标的导数转换为广义坐标即可：

\begin{sourceblock}
\includegraphics[page=306,trim={53.00bp 114.91bp 53.87bp 521.10bp},clip,width=0.92\textwidth,max height=0.38\textheight]{Ferziger4th.pdf}
\end{sourceblock}
其中 β i j 表示雅可比行列式 J 中 ∂ x i /∂Σ j 的余因子。在 2D 中，这会导致：

\begin{sourceblock}
\begin{tikzpicture}
\node[inner sep=0,anchor=south west] (image) at (0,0) {\includegraphics[page=307,trim={53.00bp 582.67bp 53.87bp 50.00bp},clip,width=0.92\textwidth,max height=0.38\textheight]{Ferziger4th.pdf}};
\begin{scope}[x={(image.south east)},y={(image.north west)}]
\fill[white] (0.29681,0.49059) rectangle (0.33504,0.83597);
\fill[white] (0.29272,0.03585) rectangle (0.33913,0.41589);
\end{scope}
\end{tikzpicture}
\end{sourceblock}
通用守恒定律，在笛卡尔坐标系中为：

\begin{sourceblock}
\begin{tikzpicture}
\node[inner sep=0,anchor=south west] (image) at (0,0) {\includegraphics[page=307,trim={53.00bp 524.17bp 53.87bp 112.86bp},clip,width=0.92\textwidth,max height=0.38\textheight]{Ferziger4th.pdf}};
\begin{scope}[x={(image.south east)},y={(image.north west)}]
\fill[white] (0.26250,0.56132) rectangle (0.34271,0.95878);
\fill[white] (0.39567,0.56132) rectangle (0.41423,0.95878);
\fill[white] (0.34815,0.32154) rectangle (0.37633,0.71865);
\fill[white] (0.28662,0.06939) rectangle (0.31591,0.47578);
\fill[white] (0.38177,0.06939) rectangle (0.41684,0.47578);
\fill[white] (0.41630,0.04122) rectangle (0.42695,0.33597);
\end{scope}
\end{tikzpicture}
\end{sourceblock}
变换为：

\begin{sourceblock}
\begin{tikzpicture}
\node[inner sep=0,anchor=south west] (image) at (0,0) {\includegraphics[page=307,trim={53.00bp 467.25bp 53.87bp 162.78bp},clip,width=0.92\textwidth,max height=0.38\textheight]{Ferziger4th.pdf}};
\begin{scope}[x={(image.south east)},y={(image.north west)}]
\fill[white] (0.21266,0.45251) rectangle (0.29284,0.77264);
\fill[white] (0.34472,0.45251) rectangle (0.36325,0.77264);
\fill[white] (0.19215,0.25173) rectangle (0.21029,0.57186);
\fill[white] (0.29829,0.25893) rectangle (0.32647,0.57934);
\fill[white] (0.23675,0.05566) rectangle (0.26605,0.38327);
\end{scope}
\end{tikzpicture}
\end{sourceblock}
\begin{sourceblock}
\begin{tikzpicture}
\node[inner sep=0,anchor=south west] (image) at (0,0) {\includegraphics[page=307,trim={53.00bp 432.44bp 53.87bp 214.13bp},clip,width=0.92\textwidth,max height=0.38\textheight]{Ferziger4th.pdf}};
\begin{scope}[x={(image.south east)},y={(image.north west)}]
\fill[white] (0.00000,0.72969) rectangle (0.07732,1.00000);
\end{scope}
\end{tikzpicture}
\end{sourceblock}
与垂直于坐标表面的速度分量成正比 xi j = const。系数 B mj 定义为：

\begin{sourceblock}
\includegraphics[page=307,trim={53.00bp 373.79bp 53.87bp 273.91bp},clip,width=0.92\textwidth,max height=0.38\textheight]{Ferziger4th.pdf}
\end{sourceblock}
变换后的动量方程包含几个附加项，这些附加项是由于动量方程中的扩散项包含通用守恒方程中未找到的导数而产生的，请参见方程 1。 （1.16）、（1.18）和（1.19）。这些术语的形式与上面所示的相同，这里不再列出。

方程（9.6）与方程（9.6）具有相同的形式。 ( 9.5 )，但后者的每一项都被前者的三项之和所取代。如上所示，这些项包含坐标的一阶导数作为系数。这些并不难通过数值计算（与二阶导数不同）。非正交网格的不寻常特征是混合导数出现在扩散项中。为了清楚地表明这一点，我们重写了方程。 ( 9.6 ) 的展开形式：

\begin{sourceblock}
\begin{tikzpicture}
\node[inner sep=0,anchor=south west] (image) at (0,0) {\includegraphics[page=307,trim={53.00bp 131.14bp 53.87bp 425.10bp},clip,width=0.92\textwidth,max height=0.38\textheight]{Ferziger4th.pdf}};
\begin{scope}[x={(image.south east)},y={(image.north west)}]
\fill[white] (0.08195,0.82011) rectangle (0.16217,0.92539);
\fill[white] (0.21359,0.82020) rectangle (0.23212,0.92539);
\fill[white] (0.06147,0.75414) rectangle (0.07958,0.85942);
\fill[white] (0.16761,0.75660) rectangle (0.19579,0.86178);
\fill[white] (0.10608,0.68981) rectangle (0.13537,0.79745);
\fill[white] (0.19474,0.48471) rectangle (0.21326,0.58990);
\fill[white] (0.18232,0.34631) rectangle (0.22653,0.46206);
\fill[white] (0.18364,0.14932) rectangle (0.20217,0.25450);
\fill[white] (0.17128,0.01092) rectangle (0.21546,0.12657);
\end{scope}
\end{tikzpicture}
\end{sourceblock}
源自梯度算子的 φ 的所有三个导数都出现在每个源自散度算子的外导数内部，请参见等式 1。 （1.27）。 φ的混合导数乘以指数不等的系数B mj，当网格正交时，无论是直线还是曲线，该系数都为零。如果

\begin{sourceblock}
\begin{tikzpicture}
\node[inner sep=0,anchor=south west] (image) at (0,0) {\includegraphics[page=308,trim={53.00bp 486.92bp 53.87bp 58.42bp},clip,width=0.92\textwidth,max height=0.38\textheight]{Ferziger4th.pdf}};
\begin{scope}[x={(image.south east)},y={(image.north west)}]
\fill[white] (0.96126,0.00993) rectangle (0.99137,0.11498);
\end{scope}
\end{tikzpicture}
\end{sourceblock}
网格是非正交的，它们相对于对角线元素 B ii 的大小取决于网格线之间的角度和网格纵横比。当网格线之间的角度较小且纵横比较大时，混合导数相乘的系数可能大于对角线系数，这可能会导致数值问题（收敛性差、解振荡等）。如果非正交性和纵横比适中，这些项比对角项小得多并且不会引起问题。混合导数项通常被显式处理，因为它们包含在隐式计算分子中将使后者变得更大并且解更昂贵。显式处理通常会增加外部迭代的数量，但从更简单且更便宜的内部迭代中获得的节省要显著得多。

方程中的导数。 ( 9.6 ) 可以使用第 1 章中描述的 FD 方法之一来近似。 3、见图9.12。沿曲线坐标的导数的近似方法与沿直线的导数的近似方法相同。

坐标变换通常被认为是一种将复杂的非正交网格转换为简单、均匀的笛卡尔网格的方法（变换空间中的间距是任意的，但通常取 xi i = 1）。一些作者声称离散化变得更简单，因为变换后的空间中的网格看起来更简单。然而，这种简化是虚幻的：流动确实发生在复杂的几何形状中，并且这一事实无法通过巧妙的坐标变换来隐藏。尽管变换后的网格看起来确实比未变换的网格简单，但有关复杂性的信息包含在度量系数中。虽然均匀变换网格上的离散化简单且准确，但雅可比行列式和其他几何信息的计算并不简单，并且会引入额外的离散化误差；也就是说，真正的困难被隐藏在这里。

网格间距 xi i 不需要明确指定。物理空间中的体积 V 定义为：

\begin{sourceblock}
\begin{tikzpicture}
\node[inner sep=0,anchor=south west] (image) at (0,0) {\includegraphics[page=308,trim={53.00bp 123.20bp 53.87bp 527.48bp},clip,width=0.92\textwidth,max height=0.38\textheight]{Ferziger4th.pdf}};
\begin{scope}[x={(image.south east)},y={(image.north west)}]
\fill[white] (0.00000,0.92367) rectangle (0.07814,1.00000);
\fill[white] (0.10577,0.93014) rectangle (0.12890,1.00000);
\fill[white] (0.12863,0.92367) rectangle (0.14093,1.00000);
\fill[white] (0.14364,0.92367) rectangle (0.16842,1.00000);
\fill[white] (0.17110,0.92367) rectangle (0.26409,1.00000);
\fill[white] (0.26683,0.92367) rectangle (0.30490,1.00000);
\end{scope}
\end{tikzpicture}
\end{sourceblock}
如果我们将整个方程乘以 Σ 1 Σ 2 Σ 3，并将 J Σ 1 Σ 2 Σ 3 处处替换为 V，则网格间距 Σ i 消失在所有项中。如果使用中心差来近似系数 β i j，例如，在 2D 中（见图 9.12）：

\begin{sourceblock}
\begin{tikzpicture}
\node[inner sep=0,anchor=south west] (image) at (0,0) {\includegraphics[page=309,trim={53.00bp 554.60bp 53.87bp 50.00bp},clip,width=0.92\textwidth,max height=0.38\textheight]{Ferziger4th.pdf}};
\begin{scope}[x={(image.south east)},y={(image.north west)}]
\fill[white] (0.28917,0.60936) rectangle (0.33392,0.81378);
\fill[white] (0.28917,0.17355) rectangle (0.33392,0.37797);
\end{scope}
\end{tikzpicture}
\end{sourceblock}
最终的离散项将仅涉及相邻节点之间的笛卡尔坐标的差异以及每个节点周围的假想单元的体积。因此，我们所需要的只是在每个网格节点周围构造这样的不重叠的单元并计算它们的体积——坐标 xi i 不需要赋任何值并且坐标变换被隐藏。

\subsection*{9.5.2 基于形函数的方法}\phantomsection\addcontentsline{toc}{subsection}{9.5.2 基于形函数的方法}
尽管我们不知道是否有人尝试过，但FD方法也可以应用于任意非结构化网格。人们必须指定一个可微的形状函数（可能是多项式）来描述特定网格点附近变量 的变化。多项式的系数可以通过将形状函数拟合到许多周围节点处的 φ 值来获得。不需要变换方程中的任何项，因此可以使用最简单的笛卡尔形式。可以对形状函数进行解析微分，以根据周围节点处的变量值和几何参数提供关于网格点处的笛卡尔坐标的一阶和二阶导数的表达式。所得系数矩阵将是稀疏的，但除非网格是结构化的，否则它不会具有对角结构。

还可以允许使用不同的形状函数，具体取决于本地网格拓扑。这将导致计算分子中的邻居数量不同，但是可以轻松设计出可以处理这种复杂性的求解器（例如代数多重网格或共轭梯度类型求解器）。

人们还可以设计一种根本不需要网格的有限差分方法；所需要的只是一组充分分布在解域上的离散点。然后，我们可以找到每个点的一定数量的近邻，可以拟合合适的形状函数；然后可以对形状函数进行微分以获得该点的导数的近似值。最适合使用微分方程的非保守形式（参见方程 1.20），因为它更容易离散化。该方法在任何情况下都不能完全保守，但如果需要时点的间距足够密集，这不是问题。

在空间中分布点似乎比创建合适的控制体积或高质量的元素更容易。 FV 方法的最大问题之一是创建解域的封闭表面，如上所述。 FD 方法的宽容度更高；两个网格点之间的任何间隙都是不可见的。此外，单元应满足某些最低质量要求，否则可能会出现收敛问题或非物理解，或两者兼而有之。 CV 形状的优化很复杂，并且许多操作（例如使用非共形网格界面压印表面）对公差很敏感。另一方面，在空间中移动单个点或引入新的点更容易，因为只有它们之间的距离起作用——公差不是问题。

第一步是在表面上放置点，然后在垂直于表面的方向上添加一小段距离的点，就像在 FV 网格中创建棱镜层时一样。第二组点可以规则地分布在求解域（例如，笛卡尔网格）中，边界附近具有更高的密度（例如在单元格局部细化中）。然后可以检查两组点是否重叠，如果点彼此太接近，则可以移动、删除或合并它们。同样，如果存在间隙，则可以轻松插入新点。局部细化非常容易，只需在现有点之间插入更多点即可。如果解域的边界发生移动，则距边界一定距离内的网格点将随之移动，而其余点则保持原样。

唯一棘手的事情是推导合适的压力或压力校正方程；但是，这可以按照以下部分中介绍的方法来实现。我们希望在这项工作的未来版本中看到这种方法。

上述原理适用于所有方程。这里不会详细讨论导出压力或压力校正方程或在非正交网格上实现 FD 方法中的边界条件的特殊功能，因为到目前为止给出的技术的扩展是简单的。

\section*{9.6 有限体积方法}\phantomsection\addcontentsline{toc}{section}{9.6 有限体积方法}
\begin{sourceblock}
\begin{tikzpicture}
\node[inner sep=0,anchor=south west] (image) at (0,0) {\includegraphics[page=310,trim={53.00bp 211.40bp 53.87bp 425.28bp},clip,width=0.92\textwidth,max height=0.38\textheight]{Ferziger4th.pdf}};
\begin{scope}[x={(image.south east)},y={(image.north west)}]
\fill[white] (0.07570,0.56687) rectangle (0.09426,0.95927);
\fill[white] (0.14481,0.32960) rectangle (0.18502,0.72234);
\fill[white] (0.19173,0.31738) rectangle (0.23128,0.71317);
\fill[white] (0.23768,0.32960) rectangle (0.26586,0.72234);
\fill[white] (0.07020,0.08045) rectangle (0.09946,0.48201);
\fill[white] (0.12235,0.04073) rectangle (0.13997,0.33164);
\end{scope}
\end{tikzpicture}
\end{sourceblock}
FV 方法的原理已在第 1 章中介绍。如图4所示，其中使用矩形控制体积（CV）来说明各种近似的常用方法。在这里，我们指的是将这些近似应用于任意多面体形状的 CV 的特殊之处。

当使用笛卡尔基向量时，控制方程不包含曲率项，因此 CV 可以通过直线连接的顶点来定义。细胞表面的实际形状并不重要；重要的是细胞表面的形状。因为它以直线段为界，所以无论形状如何，它在笛卡尔坐标曲面（表示曲面矢量分量）上的投影都是相同的。

我们将考虑块结构和非结构化网格，并首先描述如何组织必要的网格数据。

\subsection*{9.6.1 块结构网格}\phantomsection\addcontentsline{toc}{subsection}{9.6.1 块结构网格}
对于复杂的几何形状构建结构化网格很困难，有时甚至是不可能的。例如，要计算自由流中圆柱体周围的流量，可以轻松地在其周围生成结构化的 O 型网格，但如果圆柱体位于狭窄的管道中，则这不再可能。在这种情况下，块结构网格在（i）可用于结构化网格的简单性和多种求解器与（ii）处理非结构化网格允许的复杂几何形状的能力之间提供了有用的折衷。

这个想法是在构建块时在每个块内使用规则的数据结构（字典顺序），以便通过非常粗糙的非结构化网格（每个块代表一个单元）来填充不规则域。

许多方法都是可能的。有些使用重叠块（例如，Hinatsu 和 Ferziger 1991；Perng 和 Street 1991；Zang 和 Street 1995；Hubbard 和 Chen 1994；1995）。其他人依赖于非重叠块（例如，Coelho 等人，1991；Lilek 等人，1997b）。我们将描述一种使用非重叠块的方法。它也非常适合在并行计算机上使用（参见第 12 章）；通常，每个块的计算被分配给单独的处理器。

首先将解域细分为几个子域，使得每个子域都可以拟合具有良好属性的结构化网格（不要太非正交，单个CV纵横比不要太大）。一个例子如图2.2所示。在每个块中，索引 i 和 j 用于标识 CV，但我们还需要一个块标识符。数据存储在一维数组中。该一维数组中块 3 中的节点 ( i , j ) 的索引为（参见表 3.2 ）： l = O 3 + ( i − 1 ) N 3

\begin{sourceblock}
\begin{tikzpicture}
\node[inner sep=0,anchor=south west] (image) at (0,0) {\includegraphics[page=311,trim={53.00bp 289.83bp 53.87bp 358.20bp},clip,width=0.92\textwidth,max height=0.38\textheight]{Ferziger4th.pdf}};
\begin{scope}[x={(image.south east)},y={(image.north west)}]
\fill[white] (0.00000,0.85478) rectangle (0.05991,1.00000);
\end{scope}
\end{tikzpicture}
\end{sourceblock}
其中 O 3 是块 3 的偏移量（所有先前块中的节点数，即
j ) and N m
i and N m
N 1
j + N 2
i N 1
我N 2
j 是块 m 中 i 和 j 方向的节点数。

两个相邻块的网格在界面处不需要匹配；示例如图2.3所示。非共形接口的详细信息如图 1 和 2 所示。 9.13 和 9.14。造成这种情况的原因有两个：(i) 由于精度要求，一个块中的网格比另一块中的网格更细； (ii) 一个块中的网格相对于另一个块移动（所谓的滑动界面，通常在旋转部件周围的流动模拟中发现）。

处理这种情况的一种可能性是在界面的另一侧使用辅助节点，就好像网格在界面上保留相同的结构一样（也称为“幽灵”或“悬挂节点”），见图 9.13。然后使用相邻网格的周围节点的插值来计算这些节点处的变量值。例如，图 9.13 中块 A 中单元格的幽灵节点 E 处的变量值可以通过使用块 B 中最近邻居节点的变量值和梯度来获得：

\begin{sourceblock}
\includegraphics[page=312,trim={53.00bp 457.26bp 53.87bp 52.00bp},clip,width=\textwidth,max height=0.38\textheight]{Ferziger4th.pdf}
\par\vspace{0.45\baselineskip}
\small 图 9.13 具有非共形网格的两个块之间的界面：使用鬼（悬挂）节点的非保守处理
\end{sourceblock}
\begin{sourceblock}
\includegraphics[page=312,trim={53.00bp 396.50bp 53.87bp 253.14bp},clip,width=0.92\textwidth,max height=0.38\textheight]{Ferziger4th.pdf}
\end{sourceblock}
对于来自块 B 的单元，可以在幽灵节点 W 处应用相同的处理。这些值在每次内部迭代之后或每次外部迭代之后更新；该处理类似于使用域分解处理并行计算中的子域接口；见节。详细内容请参见第 12 章 12.6.2。另一种选择是更新等式右侧的节点值。 ( 9.13 ) 在每次内部迭代之后以及每次外部迭代之后的剩余部分。还有额外的簿记需要实现，但我们不会在这里处理它。

这种方法的唯一问题是它不是完全保守的：通过界面一侧的块 A 中的面的通量总和可能不等于通过另一侧的块 B 中的面的通量总和。然而，人们可以设计一种修正来强制保护（例如，Zang 和 Street 1995）。

另一种完全保守的方法是允许沿界面的 CV 具有四个以上（2D 中）或六个以上（3D 中）面，即将它们视为任意多边形或多面体单元。为此，必须识别界面两侧的两个单元所共有的界面表面部分。位于界面中的原始单元面将不会用于通量近似——它们仅用于执行创建界面所需的映射。例如，在块 A 中工作时，图 9.14 中所示的块 A 中阴影单元的东面不包括在内。因此，该 CV 的系数矩阵和源项将不完整，因为缺少其东侧的贡献；特别是，系数 A E 将为零。

因为图 9.14 的 A 块中的阴影 CV 在其东面（分为三个界面）上有三个邻居，所以我们不能在这里使用结构化网格的常用符号。为了处理在块界面处发现的不规则单元面，我们必须使用另一种数据结构——类似于整个网格非结构化时使用的数据结构。两个 CV 共用的每个接口

\begin{sourceblock}
\includegraphics[page=313,trim={53.00bp 452.25bp 53.87bp 52.00bp},clip,width=\textwidth,max height=0.38\textheight]{Ferziger4th.pdf}
\par\vspace{0.45\baselineskip}
\small 图 9.14 具有非共形网格的两个块之间的接口：使用新数据结构的保守处理
\end{sourceblock}
必须被识别（通过预处理工具）并与近似表面积分所需的所有信息一起放置在接口列表中：

\begin{itemize}
\item 左(L) 和右(R) 相邻单元的索引，
\item 表面矢量（从L 指向R），以及
\item 单元面中心的坐标。
\end{itemize}
有了这些信息，人们就可以使用每个块内部使用的方法来近似通过这些面的通量。同样的方法可以用于 O 型和 C 型网格中发生的“切割”；在这种情况下，我们正在处理同一块两侧之间的界面（即，A 和 B 是同一块，但网格在界面处可能不共形）。

每个界面单元面都对相邻 CV 的源项有贡献（对通过延迟校正处理的对流和扩散通量的显式贡献）、对这些 CV 的主对角系数 ( A P ) 以及两个非对角系数：节点 R 的 A L 和节点 L 的 A R。因此，通过将全局系数矩阵的贡献视为属于界面单元面（始终有两个相邻单元），可以克服由于具有三个东邻居而导致的数据结构不规则问题。而不是简历。那么块相对于彼此如何排序是无关紧要的（一个块的东侧可以连接到另一个块的任何一侧）：只需向接口单元面提供相邻 CV 的索引。

来自界面单元面（即 A L 和 A R ）的贡献使全局系数矩阵 A 不规则：每行的元素数量和带宽都不是恒定的。不过，这很容易处理。我们需要做的就是修改迭代矩阵 M（参见第 5 章），使其不包含由于块界面上的面而产生的元素。这与并行计算中子域接口使用的方法相同——接口面的贡献滞后于一次迭代。

我们将描述一种基于 ILU 类型求解器的求解算法；它很容易适应其他线性方程求解器。

1. 在每个块中组装矩阵 A 的元素和源项 Q，忽略块接口的贡献。 2. 循环遍历界面单元面列表，更新节点 L 和 R 处的 A P 和 Q P，并计算存储在单元面处的矩阵元素 A L 和 A R。 3. 计算每个块中矩阵 L 和 U 的元素，忽略相邻块，即，就好像它们是独立的一样。 4. 使用矩阵 A 的正则部分（ A E 、 A W 、 A N 、 A S 、 A P 和 Q P ）计算每个块中的残差；沿着块接口的 CV 的残差是不完整的，因为引用相邻块的系数为零。 5. 循环界面单元面列表，并通过分别添加乘积 A R φ R 和 A L φ L 来更新节点 L 和 R 处的残差；一旦访问完所有面，所有残差就完成了。 6. 计算每个块中每个节点的变量更新，然后返回步骤 1.7. 重复直到满足收敛标准。

由于引用相邻块中的节点的矩阵元素对迭代矩阵 M 没有贡献，因此预计收敛所需的迭代次数将大于单块情况下的迭代次数。这种效应可以通过人为地将结构化网格分成几个子域并将每个子域视为一个块来研究。正如已经提到的，这就是通过使用空间中的域分解来并行隐式方法时所做的事情（参见第 12 章）；然后通过数值效率来衡量线性方程求解器收敛速度的下降。

Schreck 和 Peri c (1993) 以及 Seidl 等人。 （1996）等人进行了大量测试，发现即使对于相对大量的子域，性能（尤其是使用共轭梯度和多重网格求解器时）仍然非常好。当可以构建良好的结构化网格时，就应该使用它。块结构确实增加了计算工作量，但它允许解决更复杂的问题，并且当然需要更复杂的算法。

第 1 节介绍了在 O 网格中的共形界面上应用此方法的示例。 9.12. O 和 C 型网格的算法实现可以在目录 2dgl 中的代码 caffa.f 中找到；参见附录A.1。有关块结构非共形网格实现的更多详细信息，请参见 Lilek 等人。 （1997b）。

当一个块中的网格移动时（即，我们正在处理滑动界面），一个块在一个时间步内旋转 θ 而另一个块固定的事实会导致邻域连接发生变化，无论每个块中的网格是结构化还是非结构化。如果使用上述保守方法，则必须在每个时间步执行界面处的面映射；上一时间步的接口列表将被清除，并创建新的接口列表。如果使用悬挂节点方法，则还需要在每个时间步重新定义供体单元和插值模板。

\subsection*{9.6.2 非结构化网格}\phantomsection\addcontentsline{toc}{subsection}{9.6.2 非结构化网格}
非结构化网格在使网格适应域边界方面具有很大的灵活性。一般来说，可以使用任意形状的控制体积，即具有任意数量的单元面。 CFD 中使用的非结构化网格的早期版本由最多有六个面的单元组成，即四面体、棱柱、金字塔和六面体。它们都可以被视为六面体的特殊情况，因此名义上的六面体网格可能包括少于六个面的 CV。每个 CV 由八个顶点定义，因此 CV 列表还包含关联顶点的列表。列表中顶点的顺序代表单元面的相对位置；例如，前四个顶点定义底面，后四个顶点定义顶面，见图 9.15。六个相邻 CV 的位置也被隐式定义；例如，由顶点 1、2、3 和 4 定义的底面对于相邻 CV 号 1 等是公共的。采用这种方法是为了减少定义 CV 之间的连接性所需的数组数量。

20 世纪 90 年代末，CFD 中引入了多面体网格；这需要更改数据结构，因为每个 CV 的面数和面上的角点数量都不受限制。现代 CFD 代码通常对网格中的所有单元使用类似于下面描述的数据结构，即使它们是笛卡尔的。所有数据都按顶点、面和体积列表进行组织。

首先，列出所有顶点及其索引和三个笛卡尔坐标。然后，创建一个面列表，其中每个面由其索引和定义闭合多边形的顶点索引来定义（顶点按列出的顺序通过直线段连接，最后一个顶点与第一个顶点连接以闭合多边形）。最后，创建一个单元列表，其中包含单元索引和包围它的面列表。

人脸列表中存储的信息还包括：

\begin{itemize}
\item 表面矢量分量（笛卡尔坐标表面上的面投影）；
\item 面质心坐标；
\item 面两侧单元的索引（惯例通常是表面向量从第一个列出的单元指向第二个单元）；
\item 矩阵A 每侧单元的系数，与相邻单元变量值相乘。
\end{itemize}
\begin{sourceblock}
\includegraphics[page=315,trim={53.00bp 94.85bp 53.87bp 464.65bp},clip,width=\textwidth,max height=0.38\textheight]{Ferziger4th.pdf}
\par\vspace{0.45\baselineskip}
\small 图 9.15 通过八个顶点列表定义名义上六面体 CV
\end{sourceblock}
\begin{sourceblock}
\includegraphics[page=316,trim={53.00bp 445.25bp 53.87bp 52.00bp},clip,width=\textwidth,max height=0.38\textheight]{Ferziger4th.pdf}
\par\vspace{0.45\baselineskip}
\small 图 9.16 从四面体网格构造多面体控制体（2D 插图示例）
\end{sourceblock}
存储在单元格列表中的信息包括：

\begin{itemize}
\item 细胞体积；
\item 细胞质心的坐标；
\item 变量值和流体特性；
\item 矩阵A 的系数A P。
\end{itemize}
近年来，多面体网格变得流行。多面体通常通过在四面体网格的顶点周围创建控制体积来获得，这是首先创建的（但也有其他方法）。这在图 9.16 右侧的二维图中得到了演示。通过将边上的中点与面和体积的质心连接起来，将四面体分成四个六面体。一个四面体顶点周围的控制体积是通过简单地连接通过分割包含该顶点的四面体创建的所有六面体部分来定义的；在这个过程中，两个六面体共有的面消失在多面体内部。由相同的两个控制体积共享的所有子面被合并在一起，创建一个多面体面。计算点放置在CVcentroid；然后丢弃四面体网格。这最大限度地减少了多面体 CV 的面数。通常，一些优化步骤如下：可以移动多面体 CV 的顶点，可以分割或合并面甚至单元格以创建具有更好属性的单元格。

将四面体网格转换为多面体网格的另一种方法是用与其正交的平面切割每个四面体的边缘；切割点可能是也可能不是线中心。各种 CV 类型的优点和缺点将在第 1 章中讨论。 12 电网质量问题将得到解决。

\subsection*{9.6.3 基于控制体积的有限元方法的网格}\phantomsection\addcontentsline{toc}{subsection}{9.6.3 基于控制体积的有限元方法的网格}
我们在这里仅对使用三角形元素和线性形状函数的混合 FE/FV 方法进行简短描述。有关适当的有限元方法及其在纳维-斯托克斯方程中的应用的详细信息，请参阅 Oden (2006)、Zienkiewicz 等人的书籍。 （2005）或弗莱彻（1991）。

对于此方法，使用由三角形（2D）或四面体（3D）组成的非结构化网格，但相同的原理也适用于由四边形和六面体组成的网格（包括作为特殊情况的棱柱和金字塔）。网格表示用于描述变量变化的元素，即定义形状函数。计算节点位于单元顶点。通常，假设变量 φ 在单元内线性变化，即其形状函数为（二维）：

\begin{sourceblock}
\begin{tikzpicture}
\node[inner sep=0,anchor=south west] (image) at (0,0) {\includegraphics[page=317,trim={53.00bp 420.54bp 53.87bp 230.14bp},clip,width=0.92\textwidth,max height=0.38\textheight]{Ferziger4th.pdf}};
\begin{scope}[x={(image.south east)},y={(image.north west)}]
\fill[white] (0.00000,0.92367) rectangle (0.03744,1.00000);
\fill[white] (0.04009,0.92367) rectangle (0.09982,1.00000);
\end{scope}
\end{tikzpicture}
\end{sourceblock}
通过将函数拟合到顶点处的节点值来确定系数a 0 、a 1 和a 2。因此它们是节点处的坐标和变量值的函数。

在二维中，通过连接单元的质心和单元边上的中点，在每个单元节点周围形成控制体积，从而在每个计算节点周围创建四边形子单元，如图 9.16 左侧所示。在 3D 中，六面体子元素是通过使用体积、面和边的中点分割四面体来创建的，如上所述，用于将四面体网格转换为多面体网格。主要区别在于，计算点仍然是起始网格的顶点，而不是像上一节介绍的方法中那样的控制体积的质心，参见图 9.16 的右侧。此外，保留在控制体积表面上的分裂四面体产生的每个面都被保留，表面积分的积分点位于其质心。

积分形式的守恒方程应用于这些具有许多面（2D 中大约 10 个，3D 大约 50 个）的不规则多面体形状的 CV。使用为每个元素定义的形状函数，按元素计算表面和体积积分：对于图 9.16 中所示的 2D CV，CV 表面由 12 个子面组成，其体积由六个子体积组成（来自共享相同顶点并因此对 CV 做出贡献的六个元素）。由于元素上变量的变化是以解析函数的形式规定的，因此可以轻松计算积分。通常，形状函数仅用于计算子面或子体积质心处的变量值，并且中点规则近似用于评估表面和体积积分，如下所述。

CV 的代数方程涉及节点 P 及其直接邻居（图 9.16 中的 N 1 到 N 6 ）。尽管该图中所示的 2D 网格仅由三角形组成，但一个 CV 与另一个 CV 之间的邻居数量通常会有所不同，具体取决于有多少个三角形共享一个顶点；这导致不规则的矩阵结构。这限制了可以使用的求解器的范围；通常单独使用共轭梯度和高斯-赛德尔求解器，或者在代数多重网格方法中使用。

因此，这种方法使用多面体控制体积，尽管它们从未显示出来；结果显示在初始网格定义的元素上。与应用于多面体控制体积的经典 FV 方法的主要区别（我们将在本章剩余部分讨论）是：

即使在两种情况下都使用表面和体积积分的中点法则近似，积分点的数量是不同的：经典方法在两个相邻单元之间有一个面，并且在CV质心中具有单个体积积分点，而本方法生成由相同的两个相邻单元共享的多个面，并评估每个子体积的体积积分。因此，经典 FV 方法的每次迭代的计算量和存储要求较低。 • 人们可能认为上述方法的精度可能会稍高一些，因为积分是使用大量较小的区域和体积来执行的，但这是值得怀疑的，因为数据是使用相同数量的计算点和同类近似值进行插值的。 • 在本方法中，还在位于解域边界上的原始网格的顶点周围创建控制体积。这些 CV 需要特殊处理，因为指定边界值时没有需要求解的方程。

Baliga 和 Patankar (1983)、Schneider 和 Raw (1987)、Masson 等人遵循了这种方法，尽管只是在二维中并使用二阶近似。 ( 1994 )、巴利加 ( 1997 ) 等。 Raw (1985) 推出了 3D 版本。

\subsection*{9.6.4 网格参数计算}\phantomsection\addcontentsline{toc}{subsection}{9.6.4 网格参数计算}
在 3D 中，单元面不一定是平面的。为了计算细胞体积和细胞面表面向量，需要合适的近似值。一种简单的方法是用一组平面三角形来表示单元面。对于结构化网格中使用的六面体，Kordula 和 Vinokur（1983）建议将每个 CV 分解为八个四面体（每个 CV 面被细分为两个三角形），以便不会发生重叠。

计算任意 CV 的细胞体积的另一种方法是基于高斯定理。通过使用恒等式 1 = div ( x i )，可以将体积计算为：

\begin{sourceblock}
\begin{tikzpicture}
\node[inner sep=0,anchor=south west] (image) at (0,0) {\includegraphics[page=318,trim={53.00bp 87.59bp 53.87bp 552.00bp},clip,width=0.92\textwidth,max height=0.38\textheight]{Ferziger4th.pdf}};
\begin{scope}[x={(image.south east)},y={(image.north west)}]
\fill[white] (0.11931,0.45122) rectangle (0.14244,0.88663);
\fill[white] (0.15050,0.46139) rectangle (0.17868,0.89680);
\fill[white] (0.22608,0.44746) rectangle (0.26562,0.88663);
\fill[white] (0.27371,0.46139) rectangle (0.30189,0.89680);
\fill[white] (0.19859,0.14049) rectangle (0.21624,0.46365);
\end{scope}
\end{tikzpicture}
\end{sourceblock}
\begin{sourceblock}
\includegraphics[page=319,trim={53.00bp 447.25bp 53.87bp 52.00bp},clip,width=\textwidth,max height=0.38\textheight]{Ferziger4th.pdf}
\par\vspace{0.45\baselineskip}
\small 图 9.17 关于任意控制体积的细胞体积和表面向量的计算：两种替代方法
\end{sourceblock}
其中 k 表示单元面，S x
k 是单元面表面向量的 x 分量（见图 9.17）：

\begin{sourceblock}
\begin{tikzpicture}
\node[inner sep=0,anchor=south west] (image) at (0,0) {\includegraphics[page=319,trim={53.00bp 362.02bp 53.87bp 284.55bp},clip,width=0.92\textwidth,max height=0.38\textheight]{Ferziger4th.pdf}};
\begin{scope}[x={(image.south east)},y={(image.north west)}]
\fill[white] (0.00000,0.81196) rectangle (0.05239,1.00000);
\fill[white] (0.05504,0.81196) rectangle (0.18195,1.00000);
\end{scope}
\end{tikzpicture}
\end{sourceblock}
除了 x i 之外，还可以使用 y j 或 z k，在这种情况下，必须对 y k S y 的乘积求和
k or z k S z
k。如果每个单元面对于其共用的两个 CV 以相同的方式定义，则该过程可确保不会发生重叠，并且所有 CV 体积的总和等于解域的体积。

一个重要的问题是单元表面的表面向量的定义。最简单的方法是将它们分解为具有一个公共顶点的三角形；参见图9.17的左侧部分。三角形的面积和表面法向量很容易计算。整个单元面的表面法向量就是所有三角形的表面向量之和（参见图 9.17 中的面 k）：

\begin{sourceblock}
\begin{tikzpicture}
\node[inner sep=0,anchor=south west] (image) at (0,0) {\includegraphics[page=319,trim={53.00bp 199.21bp 53.87bp 429.23bp},clip,width=0.92\textwidth,max height=0.38\textheight]{Ferziger4th.pdf}};
\begin{scope}[x={(image.south east)},y={(image.north west)}]
\fill[white] (0.36944,0.71485) rectangle (0.38824,0.94218);
\fill[white] (0.38376,0.67241) rectangle (0.39525,0.84695);
\fill[white] (0.38586,0.79363) rectangle (0.39817,0.96817);
\end{scope}
\end{tikzpicture}
\end{sourceblock}
where N v
k 是单元面中的顶点数量，ri 是顶点 i 的位置向量。注意有 N v
k − 2 个三角形。即使单元面是扭曲的或凸出的，上述表达式也是正确的。公共顶点的选择并不重要。

单元面中心可以通过按其面积加权的每个三角形中心坐标（其本身就是其顶点坐标的平均值）求平均值来找到。单元面的面积等于其表面矢量的大小，例如：

\begin{sourceblock}
\begin{tikzpicture}
\node[inner sep=0,anchor=south west] (image) at (0,0) {\includegraphics[page=319,trim={53.00bp 75.23bp 53.87bp 571.22bp},clip,width=0.92\textwidth,max height=0.38\textheight]{Ferziger4th.pdf}};
\begin{scope}[x={(image.south east)},y={(image.north west)}]
\fill[white] (0.26794,0.15033) rectangle (0.29660,0.77908);
\fill[white] (0.30289,0.20518) rectangle (0.33107,0.79279);
\fill[white] (0.33459,0.15033) rectangle (0.38526,0.79279);
\fill[white] (0.38878,0.20518) rectangle (0.41696,0.79279);
\end{scope}
\end{tikzpicture}
\end{sourceblock}
请注意，仅需要将单元面投影到笛卡尔坐标平面上。当 CV 边缘是直的时（正如它们所假设的那样），这些是准确的。

如果除了流体流动之外，还需要计算粒子的运动，则必须确保以所有三角形的表面矢量都指向外的方式执行单元面的三角测量（即，所有单个表面矢量的标量积必须为正）。上述方法不能保证这一点。参考图9.17的右侧部分提出了一种更鲁棒的计算细胞面数据的方法，但它创建了更多的三角形，因此导致更大的计算量。

在该方法中，在每个面 k 处定义辅助中心点 h；例如，它的坐标可以是定义面的所有顶点的平均坐标：

\begin{sourceblock}
\includegraphics[page=320,trim={53.00bp 429.47bp 53.87bp 198.98bp},clip,width=0.92\textwidth,max height=0.38\textheight]{Ferziger4th.pdf}
\end{sourceblock}
其中 r v i 表示定义面的顶点列表中面 k 上的第 i 个顶点。

然后将单元面细分为 N v
通过将每个顶点与中心点 h 连接来形成 k 个三角形。第 m 个三角形的表面向量可以表示为：

\begin{sourceblock}
\includegraphics[page=320,trim={53.00bp 348.25bp 53.87bp 289.87bp},clip,width=0.92\textwidth,max height=0.38\textheight]{Ferziger4th.pdf}
\end{sourceblock}
整个面的表面向量等于所有三角形的表面向量之和：

\begin{sourceblock}
\includegraphics[page=320,trim={53.00bp 282.76bp 53.87bp 345.74bp},clip,width=0.92\textwidth,max height=0.38\textheight]{Ferziger4th.pdf}
\end{sourceblock}
每个三角形 m 、 r k 、 m 的质心坐标定义如下：

\begin{sourceblock}
\includegraphics[page=320,trim={53.00bp 225.21bp 53.87bp 412.91bp},clip,width=0.92\textwidth,max height=0.38\textheight]{Ferziger4th.pdf}
\end{sourceblock}
存储用于离散化过程的面质心坐标为：

\begin{sourceblock}
\begin{tikzpicture}
\node[inner sep=0,anchor=south west] (image) at (0,0) {\includegraphics[page=320,trim={53.00bp 160.83bp 53.87bp 464.39bp},clip,width=0.92\textwidth,max height=0.38\textheight]{Ferziger4th.pdf}};
\begin{scope}[x={(image.south east)},y={(image.north west)}]
\fill[white] (0.00000,0.83284) rectangle (0.10147,1.00000);
\fill[white] (0.10418,0.83284) rectangle (0.15389,1.00000);
\end{scope}
\end{tikzpicture}
\end{sourceblock}
为了计算单元体积，定义另一个辅助中心点 H 是有用的；该点的坐标可以定义为所有单元顶点坐标的平均值。现在可以将每个面中每个三角形的顶点与中心 H 连接起来，从而创建一个四面体，其体积可以轻松计算：

\begin{sourceblock}
\includegraphics[page=320,trim={53.00bp 69.17bp 53.87bp 568.95bp},clip,width=0.92\textwidth,max height=0.38\textheight]{Ferziger4th.pdf}
\end{sourceblock}
现在可以通过将附加到所有单元面中的三角形的所有上述定义的四面体的体积相加来计算单元体积：

\begin{sourceblock}
\includegraphics[page=321,trim={53.00bp 536.13bp 53.87bp 92.32bp},clip,width=0.92\textwidth,max height=0.38\textheight]{Ferziger4th.pdf}
\end{sourceblock}
where N f
P 表示围绕节点 P 包围单元的面数。单元质心 P 的坐标可以通过对每个四面体的坐标与其体积进行加权来计算：

\begin{sourceblock}
\includegraphics[page=321,trim={53.00bp 449.66bp 53.87bp 179.78bp},clip,width=0.92\textwidth,max height=0.38\textheight]{Ferziger4th.pdf}
\end{sourceblock}
其中四面体质心的坐标定义为其顶点的平均坐标：

\begin{sourceblock}
\begin{tikzpicture}
\node[inner sep=0,anchor=south west] (image) at (0,0) {\includegraphics[page=321,trim={53.00bp 395.80bp 53.87bp 242.32bp},clip,width=0.92\textwidth,max height=0.38\textheight]{Ferziger4th.pdf}};
\begin{scope}[x={(image.south east)},y={(image.north west)}]
\fill[white] (0.00000,0.81585) rectangle (0.02908,1.00000);
\fill[white] (0.03179,0.81585) rectangle (0.06493,1.00000);
\fill[white] (0.06758,0.81585) rectangle (0.17913,1.00000);
\end{scope}
\end{tikzpicture}
\end{sourceblock}
\section*{9.7 通量和源项的近似}\phantomsection\addcontentsline{toc}{section}{9.7 通量和源项的近似}
\subsection*{9.7.1 对流通量的近似}\phantomsection\addcontentsline{toc}{subsection}{9.7.1 对流通量的近似}
我们将专门使用表面和体积积分的中点法则近似；它是唯一适用于任意形状的积分域的二阶近似。只需要知道人脸质心的坐标即可；上一节描述了如何计算任意多边形的这些值。开发更高阶方法的必要步骤将在第 4 节中描述。 9.10。

我们首先看一下质量通量的计算。仅考虑一张脸，在图 9.18 所示的 CV 中用索引 k 表示；同样的方法适用于其他面——只需替换索引。 CV 可以有任意数量的面孔。

通过面 k 的质量通量的中点法则近似可得出：

\begin{sourceblock}
\begin{tikzpicture}
\node[inner sep=0,anchor=south west] (image) at (0,0) {\includegraphics[page=321,trim={53.00bp 154.05bp 53.87bp 489.02bp},clip,width=0.92\textwidth,max height=0.38\textheight]{Ferziger4th.pdf}};
\begin{scope}[x={(image.south east)},y={(image.north west)}]
\fill[white] (0.28394,0.40009) rectangle (0.32147,0.94798);
\fill[white] (0.32773,0.44690) rectangle (0.35591,0.94798);
\end{scope}
\end{tikzpicture}
\end{sourceblock}
面 k 处的单位法线矢量由表面矢量 S k 和面面积定义
S k，其在笛卡尔坐标平面上的投影为 S i
k :

\begin{sourceblock}
\includegraphics[page=321,trim={53.00bp 94.61bp 53.87bp 553.08bp},clip,width=0.92\textwidth,max height=0.38\textheight]{Ferziger4th.pdf}
\end{sourceblock}
表面积 S k 为：

\begin{sourceblock}
\includegraphics[page=322,trim={53.00bp 454.25bp 53.87bp 52.00bp},clip,width=\textwidth,max height=0.38\textheight]{Ferziger4th.pdf}
\par\vspace{0.45\baselineskip}
\small 图 9.18 一般 2D 和 3D 控制体积及使用的符号
\end{sourceblock}
\begin{sourceblock}
\begin{tikzpicture}
\node[inner sep=0,anchor=south west] (image) at (0,0) {\includegraphics[page=322,trim={53.00bp 389.46bp 53.87bp 248.32bp},clip,width=0.92\textwidth,max height=0.38\textheight]{Ferziger4th.pdf}};
\begin{scope}[x={(image.south east)},y={(image.north west)}]
\fill[white] (0.39230,0.41008) rectangle (0.42096,0.84661);
\fill[white] (0.42725,0.44852) rectangle (0.45543,0.85614);
\end{scope}
\end{tikzpicture}
\end{sourceblock}
见节。有关如何计算任意多面体 CV 的表面矢量分量的详细信息，请参见 9.6.4。

根据这些定义，质量通量的表达式变为：

\begin{sourceblock}
\begin{tikzpicture}
\node[inner sep=0,anchor=south west] (image) at (0,0) {\includegraphics[page=322,trim={53.00bp 305.19bp 53.87bp 333.94bp},clip,width=0.92\textwidth,max height=0.38\textheight]{Ferziger4th.pdf}};
\begin{scope}[x={(image.south east)},y={(image.north west)}]
\fill[white] (0.32057,0.43058) rectangle (0.35811,0.89856);
\fill[white] (0.36436,0.47057) rectangle (0.39254,0.89856);
\fill[white] (0.39606,0.43058) rectangle (0.42650,0.89856);
\end{scope}
\end{tikzpicture}
\end{sourceblock}
where u i
k 是插值到面质心 k 和 v n 的笛卡尔速度分量
k 是与面 k 正交的速度分量。笛卡尔网格和任意多面体网格之间的区别在于，在后一种情况下，表面矢量在多个笛卡尔方向上具有分量，并且所有笛卡尔速度分量都对质量通量有贡献。每个笛卡尔速度分量乘以相应的表面矢量分量（单元面到笛卡尔坐标平面上的投影），请参见方程 1。 （9.31）。

任何输送量的对流通量通常是通过假设质量通量已知来计算的，通过中点法则近似，可以得出：

\begin{sourceblock}
\begin{tikzpicture}
\node[inner sep=0,anchor=south west] (image) at (0,0) {\includegraphics[page=322,trim={53.00bp 148.75bp 53.87bp 492.85bp},clip,width=0.92\textwidth,max height=0.38\textheight]{Ferziger4th.pdf}};
\begin{scope}[x={(image.south east)},y={(image.north west)}]
\fill[white] (0.31098,0.40913) rectangle (0.34617,0.95110);
\end{scope}
\end{tikzpicture}
\end{sourceblock}
其中 φ k 是面质心 k 处的 φ 值。这表示基于皮卡德迭代的线性化，如第 1 章所述。 5.

面质心处的变量值必须通过适当的插值通过节点值来表示。第 1 章中提出了几种可能性。如图4所示，但并不是所有的都可以轻松扩展到任意多面体网格，因为计算点（CV质心）并不位于可以轻松进行插值的线上。对于高阶插值，需要多维形状函数。

然而，通过使用简单的几何数据和向量运算，可以轻松构建任意多面体 CV 的简单二阶方法。考虑图 9.18 中的两个任意控制体积 (CV)。以下向量是根据细胞和面质心的坐标定义的：

\begin{sourceblock}
\includegraphics[page=323,trim={53.00bp 509.16bp 53.87bp 140.48bp},clip,width=0.92\textwidth,max height=0.38\textheight]{Ferziger4th.pdf}
\end{sourceblock}
我们还沿着连接单元中心 P 和相邻单元 N k 的中心的线引入线性插值因子，如下所示：

\begin{sourceblock}
\includegraphics[page=323,trim={53.00bp 439.55bp 53.87bp 198.30bp},clip,width=0.92\textwidth,max height=0.38\textheight]{Ferziger4th.pdf}
\end{sourceblock}
将向量 ( rk − r P ) 投影到连接细胞中心的线上，该线上的辅助点 k ′ 定义为：

\begin{sourceblock}
\includegraphics[page=323,trim={53.00bp 382.48bp 53.87bp 267.50bp},clip,width=0.92\textwidth,max height=0.38\textheight]{Ferziger4th.pdf}
\end{sourceblock}
面部两侧细胞中心之间的线性插值是最简单的二阶近似，但它导致位置 k ′ 处的变量值，该位置不是面部质心：

\begin{sourceblock}
\begin{tikzpicture}
\node[inner sep=0,anchor=south west] (image) at (0,0) {\includegraphics[page=323,trim={53.00bp 322.70bp 53.87bp 327.28bp},clip,width=0.92\textwidth,max height=0.38\textheight]{Ferziger4th.pdf}};
\begin{scope}[x={(image.south east)},y={(image.north west)}]
\fill[white] (0.00000,0.89542) rectangle (0.02412,1.00000);
\fill[white] (0.02680,0.89542) rectangle (0.06989,1.00000);
\fill[white] (0.07257,0.89542) rectangle (0.11398,1.00000);
\fill[white] (0.11669,0.89542) rectangle (0.17110,1.00000);
\fill[white] (0.17377,0.89542) rectangle (0.28505,1.00000);
\end{scope}
\end{tikzpicture}
\end{sourceblock}
如果用上述插值得到的值来表示面质心处的值，除非k′和k点几乎重合，否则结果不会是二阶精确的。

通过使用变量值及其从位置 k ′ 的梯度可以实现对人脸质心位置 k 的二阶精确插值，如下所示：

\begin{sourceblock}
\includegraphics[page=323,trim={53.00bp 227.24bp 53.87bp 422.92bp},clip,width=0.92\textwidth,max height=0.38\textheight]{Ferziger4th.pdf}
\end{sourceblock}
右侧第一项，通过等式表达时。 （9.36），提供对系数矩阵A的贡献；第二项被明确视为延迟修正。 k '处的梯度可以根据表达式(9.36)使用单元中心值的插值来计算。如果向量 d 通过面心 k，则涉及梯度的延迟校正将为零，因为 k ′ 和 k 重合；然后，面值仅根据单元中心值计算。

另一种选择是从两侧执行从像元中心到面中心的外推，然后根据距离（相当于中心差分近似）或流向（二阶迎风）或其他一些标准，在两者之间使用加权：

\begin{sourceblock}
\includegraphics[page=323,trim={53.00bp 70.78bp 53.87bp 578.86bp},clip,width=0.92\textwidth,max height=0.38\textheight]{Ferziger4th.pdf}
\end{sourceblock}
对于距离加权，我们有：

\begin{sourceblock}
\includegraphics[page=324,trim={53.00bp 569.98bp 53.87bp 80.71bp},clip,width=0.92\textwidth,max height=0.38\textheight]{Ferziger4th.pdf}
\end{sourceblock}
等式中的各个值 φ k 1 和 φ k 1 ( 9.38 ) 代表二阶近似值。如果采用从上游侧外推得到的值，则得到线性迎风方案；大多数商业 CFD 代码均提供该功能。

我们还可以在垂直于单元面并穿过其质心的线上定义两个附加辅助点。首先，我们确定连接单元中心与公共面质心到面法线的向量投影中较短的一个：

\begin{sourceblock}
\begin{tikzpicture}
\node[inner sep=0,anchor=south west] (image) at (0,0) {\includegraphics[page=324,trim={53.00bp 449.39bp 53.87bp 200.25bp},clip,width=0.92\textwidth,max height=0.38\textheight]{Ferziger4th.pdf}};
\begin{scope}[x={(image.south east)},y={(image.north west)}]
\fill[white] (0.00000,0.92848) rectangle (0.09735,1.00000);
\end{scope}
\end{tikzpicture}
\end{sourceblock}
辅助点 P ′ 和 N ′
k 现在被定义为距面质心 k 的距离 a（一个位于单元中心位置到面法线的投影处，另一个靠近面，以便两者距面中心的距离相同，见图 9.19）：

\begin{sourceblock}
\begin{tikzpicture}
\node[inner sep=0,anchor=south west] (image) at (0,0) {\includegraphics[page=324,trim={53.00bp 376.58bp 53.87bp 271.47bp},clip,width=0.92\textwidth,max height=0.38\textheight]{Ferziger4th.pdf}};
\begin{scope}[x={(image.south east)},y={(image.north west)}]
\fill[white] (0.00000,0.90658) rectangle (0.12620,1.00000);
\end{scope}
\end{tikzpicture}
\end{sourceblock}
现在可以将面质心处的变量值计算为辅助节点 P ' 和 N ' 处的值的平均值
k（因为它们距面中心的距离相等）；这些可以计算如下：

\begin{sourceblock}
\includegraphics[page=324,trim={53.00bp 270.80bp 53.87bp 343.20bp},clip,width=0.92\textwidth,max height=0.38\textheight]{Ferziger4th.pdf}
\end{sourceblock}
所有三个选项都提供二阶近似，其中包含参考面两侧单元中心处的变量值的部分（这可以对矩阵系数做出贡献）和依赖于单元中心处的梯度的部分（通常使用延迟校正方法来处理）。第一个选项简化为笛卡尔网格上的标准中心差分近似（线性插值）。如果两个细胞中心的梯度相同，第二个和第三个选项将减少到相同的近似值；否则，即使网格是笛卡尔网格，与梯度差成比例的项也将保留。

二次（三阶）和三次（四阶）插值可以通过使用节点 P 和 N k 处的变量值和梯度轻松构建，如第 1 章中所述。 4，但它们只会导致位置 k' 处的更准确的近似。当网格足够细时，仍然会得到更准确的解，但由于表达式（9.37）中的修正和积分近似都是二阶的，对流通量近似的整体阶数无法提高到二阶以上；见第 4.7.1 节 有关此问题的更多详细信息。

一阶迎风方案很容易在任何电网上实施，请参见第 1 节。 4.4.1；然而，由于其极端的数值扩散，通常仅在二阶近似因网格特性极差而失败时使用。可以使用一阶迎风格式的另一个领域是在产生振荡解的情况下稳定高阶格式。这是通过混合第 2 节中描述的两种方案来实现的。 5.6；混合系数可以由用户指定，也可以由代码根据某些标准进行评估。

\subsection*{9.7.2 扩散通量的近似}\phantomsection\addcontentsline{toc}{subsection}{9.7.2 扩散通量的近似}
\begin{sourceblock}
\begin{tikzpicture}
\node[inner sep=0,anchor=south west] (image) at (0,0) {\includegraphics[page=325,trim={53.00bp 412.96bp 53.87bp 222.83bp},clip,width=0.92\textwidth,max height=0.38\textheight]{Ferziger4th.pdf}};
\begin{scope}[x={(image.south east)},y={(image.north west)}]
\fill[white] (0.15654,0.34036) rectangle (0.19290,0.77858);
\fill[white] (0.17341,0.27446) rectangle (0.18752,0.55684);
\fill[white] (0.19835,0.34926) rectangle (0.22653,0.73015);
\end{scope}
\end{tikzpicture}
\end{sourceblock}
单元面中心处的 φ 梯度可以用相对于全局笛卡尔坐标 x i 或局部正交坐标 ( n , t , s ) 的导数来表示：

\begin{sourceblock}
\begin{tikzpicture}
\node[inner sep=0,anchor=south west] (image) at (0,0) {\includegraphics[page=325,trim={53.00bp 344.31bp 53.87bp 292.73bp},clip,width=0.92\textwidth,max height=0.38\textheight]{Ferziger4th.pdf}};
\begin{scope}[x={(image.south east)},y={(image.north west)}]
\fill[white] (0.00000,0.87560) rectangle (0.04156,1.00000);
\fill[white] (0.04171,0.87560) rectangle (0.06719,1.00000);
\fill[white] (0.06734,0.87216) rectangle (0.10511,1.00000);
\end{scope}
\end{tikzpicture}
\end{sourceblock}
其中 n 、 t 和 s 分别表示与表面垂直和相切的坐标方向； n 、 t 和 s 是每个坐标方向对应的单位向量。请注意，坐标系 ( n 、 t 、 s ) 也是笛卡尔坐标系，但它仅经过旋转，使其一个坐标垂直，另外两个坐标与单元面相切。

有很多方法可以逼近细胞面法线的导数或细胞中心的梯度向量；我们只描述其中的一小部分。如果单元面附近的 φ 变化由形状函数描述，则可以在位置 k 处对该函数进行微分，以找到相对于笛卡尔坐标的导数。则扩散通量为：

\begin{sourceblock}
\begin{tikzpicture}
\node[inner sep=0,anchor=south west] (image) at (0,0) {\includegraphics[page=325,trim={53.00bp 177.75bp 53.87bp 448.93bp},clip,width=0.92\textwidth,max height=0.38\textheight]{Ferziger4th.pdf}};
\begin{scope}[x={(image.south east)},y={(image.north west)}]
\fill[white] (0.34562,0.31526) rectangle (0.38198,0.65231);
\fill[white] (0.36250,0.26457) rectangle (0.37663,0.48175);
\fill[white] (0.38746,0.32210) rectangle (0.41564,0.61505);
\fill[white] (0.43838,0.29473) rectangle (0.45248,0.51191);
\end{scope}
\end{tikzpicture}
\end{sourceblock}
这很容易显式地实现；隐式版本可能很复杂，具体取决于形状函数的阶数和涉及的节点数量。

当使用辅助节点 P ' 和 N ' 的变量值时，可以获得一种导出扩散通量二阶近似的简单方法
k 在上一节中介绍过；见图9.19。关于 n 的导数近似为中心差，简单地为： ∂φ

\begin{sourceblock}
\begin{tikzpicture}
\node[inner sep=0,anchor=south west] (image) at (0,0) {\includegraphics[page=325,trim={53.00bp 68.30bp 53.87bp 559.20bp},clip,width=0.92\textwidth,max height=0.38\textheight]{Ferziger4th.pdf}};
\begin{scope}[x={(image.south east)},y={(image.north west)}]
\fill[white] (0.00000,0.75362) rectangle (0.12484,1.00000);
\fill[white] (0.12752,0.75362) rectangle (0.15233,1.00000);
\fill[white] (0.15498,0.75362) rectangle (0.24971,1.00000);
\fill[white] (0.38259,0.48835) rectangle (0.42081,0.78753);
\fill[white] (0.38379,0.11749) rectangle (0.41973,0.42365);
\end{scope}
\end{tikzpicture}
\par\vspace{0.45\baselineskip}
\small 图 9.19 任意多面体 CV 的扩散通量的近似
\end{sourceblock}
P

k

P

k

d

n

N

k

N

k

n

参照式（9.42），可以将上式改写如下：

\begin{sourceblock}
\begin{tikzpicture}
\node[inner sep=0,anchor=south west] (image) at (0,0) {\includegraphics[page=326,trim={53.00bp 372.06bp 53.87bp 255.45bp},clip,width=0.92\textwidth,max height=0.38\textheight]{Ferziger4th.pdf}};
\begin{scope}[x={(image.south east)},y={(image.north west)}]
\fill[white] (0.07203,0.48848) rectangle (0.11026,0.78773);
\end{scope}
\end{tikzpicture}
\end{sourceblock}
右侧的第一项表示隐式部分（对系数矩阵有贡献），而第二项表示显式部分，使用先前迭代的值（延迟校正）计算。

如果面法线穿过两侧的单元中心并且它们距面中心的距离相等，则涉及梯度的延迟校正将消失；然后，将仅使用单元中心值计算法向导数，就像均匀笛卡尔网格上的情况一样。然而，如果网格不均匀，则上述近似值 ( 9.48 ) 将包含一个校正项，使其在面质心 k 处具有二阶精度。

计算单元面导数的另一种方法是首先在 CV 中心获得它们，然后按照确定 φ k 的方式将它们插值到单元面。然而，这可能会导致振荡解，并且需要应用类似于用于避免共置电网上的振荡压力的校正。让我们首先看看如何计算细胞中心的梯度。

高斯定理提供了一种简单的方法；我们用单元上的平均值来近似 CV 中心的导数：

\begin{sourceblock}
\begin{tikzpicture}
\node[inner sep=0,anchor=south west] (image) at (0,0) {\includegraphics[page=326,trim={53.00bp 121.73bp 53.87bp 500.86bp},clip,width=0.92\textwidth,max height=0.38\textheight]{Ferziger4th.pdf}};
\begin{scope}[x={(image.south east)},y={(image.north west)}]
\fill[white] (0.37660,0.40551) rectangle (0.41483,0.67118);
\fill[white] (0.37389,0.05786) rectangle (0.41549,0.34834);
\end{scope}
\end{tikzpicture}
\end{sourceblock}
那么我们可以将导数 ∂φ/∂ x i 视为向量 φ i i 的散度，并利用高斯定理将上式中的体积积分转化为曲面积分：

\begin{sourceblock}
\begin{tikzpicture}
\node[inner sep=0,anchor=south west] (image) at (0,0) {\includegraphics[page=327,trim={53.00bp 580.17bp 53.87bp 53.51bp},clip,width=0.92\textwidth,max height=0.38\textheight]{Ferziger4th.pdf}};
\begin{scope}[x={(image.south east)},y={(image.north west)}]
\fill[white] (0.29823,0.60690) rectangle (0.33645,0.96303);
\fill[white] (0.26941,0.12939) rectangle (0.28704,0.39341);
\fill[white] (0.29549,0.14048) rectangle (0.33711,0.52988);
\end{scope}
\end{tikzpicture}
\end{sourceblock}
这表明，可以通过将 φ 与 CV 所有面上的表面矢量的 x 分量的乘积相加，并将总和除以 CV 体积，来计算 φ 相对于 CV 中心处的 x 的导数：

\begin{sourceblock}
\begin{tikzpicture}
\node[inner sep=0,anchor=south west] (image) at (0,0) {\includegraphics[page=327,trim={53.00bp 495.02bp 53.87bp 133.69bp},clip,width=0.92\textwidth,max height=0.38\textheight]{Ferziger4th.pdf}};
\begin{scope}[x={(image.south east)},y={(image.north west)}]
\fill[white] (0.39155,0.47208) rectangle (0.42977,0.78092);
\fill[white] (0.38884,0.06759) rectangle (0.43044,0.40529);
\end{scope}
\end{tikzpicture}
\end{sourceblock}
对于 φ k，我们可以使用用于计算对流通量的值，尽管不一定需要对两项使用相同的近似值。对于笛卡尔网格和线性插值，得到标准的中心差近似（只有东面和西面对x方向的导数有贡献，因为其他面的表面向量的x分量为零；单元体积可以表示为V =

\begin{sourceblock}
\begin{tikzpicture}
\node[inner sep=0,anchor=south west] (image) at (0,0) {\includegraphics[page=327,trim={53.00bp 388.68bp 53.87bp 236.81bp},clip,width=0.92\textwidth,max height=0.38\textheight]{Ferziger4th.pdf}};
\begin{scope}[x={(image.south east)},y={(image.north west)}]
\fill[white] (0.00000,0.97097) rectangle (0.02908,1.00000);
\fill[white] (0.03083,0.97097) rectangle (0.12180,1.00000);
\fill[white] (0.12358,0.97097) rectangle (0.21441,1.00000);
\fill[white] (0.21618,0.97097) rectangle (0.24262,1.00000);
\fill[white] (0.24436,0.97097) rectangle (0.31071,1.00000);
\fill[white] (0.31248,0.97097) rectangle (0.37853,1.00000);
\fill[white] (0.38027,0.97097) rectangle (0.42165,1.00000);
\fill[white] (0.42337,0.97097) rectangle (0.48806,1.00000);
\fill[white] (0.48983,0.97097) rectangle (0.53125,1.00000);
\fill[white] (0.53299,0.97097) rectangle (0.58105,1.00000);
\fill[white] (0.58280,0.97097) rectangle (0.67684,1.00000);
\fill[white] (0.67865,0.97097) rectangle (0.72502,1.00000);
\fill[white] (0.72677,0.97097) rectangle (0.75988,1.00000);
\fill[white] (0.76159,0.97097) rectangle (0.88403,1.00000);
\fill[white] (0.88580,0.97097) rectangle (0.91558,1.00000);
\fill[white] (0.94223,0.97343) rectangle (0.96535,1.00000);
\fill[white] (0.97341,0.98007) rectangle (1.00000,1.00000);
\fill[white] (0.00084,0.65929) rectangle (0.02950,0.96384);
\fill[white] (0.05152,0.67700) rectangle (0.09059,0.96384);
\fill[white] (0.39068,0.43419) rectangle (0.42890,0.71882);
\end{scope}
\end{tikzpicture}
\end{sourceblock}
细胞中心梯度也可以通过使用线性形状函数在二阶内近似；如果我们假设两个相邻单元中心（例如 P 和 N k ）之间 φ 的线性变化，我们可以写：

\begin{sourceblock}
\includegraphics[page=327,trim={53.00bp 318.43bp 53.87bp 331.21bp},clip,width=0.92\textwidth,max height=0.38\textheight]{Ferziger4th.pdf}
\end{sourceblock}
我们可以写出与节点 P 周围的单元格的邻居一样多的方程；然而，我们只需要计算三个导数 ∂φ/∂ x i。借助最小二乘法，可以从此类超定系统中针对任意 CV 形状显式计算导数；有关详细信息，请参阅 Demirdži ́c 和 Muzaferija (1995)。

以这种方式计算的导数可以插值到单元表面，并且扩散通量可以从方程（1）计算。 （9.46）。这种方法的问题在于，在迭代过程中可能会生成振荡解，并且振荡不会被检测到。下面描述了如何避免这种情况。

对于显式方法，这种方法实现起来非常简单。然而，它不适合以隐式方法实现，因为它会产生大量的计算分子。第 2 节中描述的延迟修正方法。 5.6提供了解决这个问题的方法，但它不一定有助于消除振荡。

过去 20 年来，已经开发出大量解决这一问题的方法； Demirdži ́c ( 2015 ) 给出了一个概述并指出了最好的一个，这类似于在非正交坐标中离散化扩散项得到的结果：

\begin{sourceblock}
\begin{tikzpicture}
\node[inner sep=0,anchor=south west] (image) at (0,0) {\includegraphics[page=328,trim={53.00bp 581.29bp 53.87bp 54.29bp},clip,width=0.92\textwidth,max height=0.38\textheight]{Ferziger4th.pdf}};
\begin{scope}[x={(image.south east)},y={(image.north west)}]
\fill[white] (0.02463,0.34522) rectangle (0.06099,0.78010);
\fill[white] (0.04150,0.27945) rectangle (0.05564,0.56021);
\fill[white] (0.06644,0.35373) rectangle (0.09462,0.73233);
\fill[white] (0.11738,0.31839) rectangle (0.18027,0.72349);
\fill[white] (0.92433,0.00000) rectangle (1.00000,0.07297);
\end{scope}
\end{tikzpicture}
\end{sourceblock}
(9.54) 右侧第一项被隐式处理，即它对系数矩阵 A 做出贡献；下划线项是使用变量的现行值计算的，并被视为另一个递延修正。其中，( ∇ φ) k 是根据两个 CV 中心与人脸的距离对梯度进行插值得到的，而 ( ∇ φ) k 表示 P 和 N k 处梯度的平均值。这种区别的原因是，右侧的第一项是中心差分近似，在节点 P 和 N k 之间的中点处是二阶精确的，并且当 φ 的变化平滑时，右侧的最后一项应该抵消第一项，因此梯度被插值到中点而不是单元面中心。

请注意，上式右侧的所有项实际上提供了位置 k ' 而不是 k 处的近似值。因此，如果从 k ' 到 k 的距离很大，则通量近似将不是二阶精确的。另请注意，方程的法向导数的近似值。 ( 9.54 ) 类似于使用辅助节点 P ' 和 N ' 获得的结果
k，方程。 ( 9.48 )，但它是使用不同的方法得出的。等式的近似值。即使连接单元中心 P 和 N k 的线不穿过面质心，( 9.48 ) 在面质心 k 处也是二阶精确的。

在动量方程中，扩散通量包含的项比通用守恒方程中的相应项多一些，例如，对于 u i：

\begin{sourceblock}
\begin{tikzpicture}
\node[inner sep=0,anchor=south west] (image) at (0,0) {\includegraphics[page=328,trim={53.00bp 294.13bp 53.87bp 341.95bp},clip,width=0.92\textwidth,max height=0.38\textheight]{Ferziger4th.pdf}};
\begin{scope}[x={(image.south east)},y={(image.north west)}]
\fill[white] (0.22394,0.33367) rectangle (0.26030,0.77611);
\fill[white] (0.24081,0.26713) rectangle (0.25492,0.55256);
\fill[white] (0.26571,0.34265) rectangle (0.29392,0.72721);
\end{scope}
\end{tikzpicture}
\end{sourceblock}
通用守恒方程中不存在下划线项。如果 ρ 和 μ 恒定，则根据连续性方程，所有 CV 面上的下划线项之和为零，请参见第 1 节。 7.1.如果 ρ 和 μ 不是常数，则它们（除了近激波之外）平滑变化，并且整个 CV 表面上划线项的积分小于主项的积分。因此，通常会明确处理带下划线的术语。如上所示，使用来自 CV 中心的梯度向量可以轻松计算单元表面的导数。该项不会引起振荡，并且可以使用插值梯度来计算。

\subsection*{9.7.3 源项的近似}\phantomsection\addcontentsline{toc}{subsection}{9.7.3 源项的近似}
\begin{sourceblock}
\begin{tikzpicture}
\node[inner sep=0,anchor=south west] (image) at (0,0) {\includegraphics[page=328,trim={53.00bp 66.78bp 53.87bp 574.26bp},clip,width=0.92\textwidth,max height=0.38\textheight]{Ferziger4th.pdf}};
\begin{scope}[x={(image.south east)},y={(image.north west)}]
\fill[white] (0.32448,0.37649) rectangle (0.36469,0.95219);
\end{scope}
\end{tikzpicture}
\end{sourceblock}
该近似与 CV 形状无关，并且具有二阶精度。

现在让我们看看动量方程中的压力项。它们可以被视为 CV 表面上的保守力，也可以被视为非保守体力。在第一种情况下，我们有（在 u i 的等式中）：

\begin{sourceblock}
\begin{tikzpicture}
\node[inner sep=0,anchor=south west] (image) at (0,0) {\includegraphics[page=329,trim={53.00bp 519.33bp 53.87bp 118.83bp},clip,width=0.92\textwidth,max height=0.38\textheight]{Ferziger4th.pdf}};
\begin{scope}[x={(image.south east)},y={(image.north west)}]
\fill[white] (0.28770,0.44496) rectangle (0.31417,0.85811);
\fill[white] (0.31098,0.36741) rectangle (0.32746,0.67370);
\fill[white] (0.31338,0.65046) rectangle (0.32869,0.95711);
\fill[white] (0.33435,0.45461) rectangle (0.36253,0.86776);
\fill[white] (0.36605,0.45461) rectangle (0.39423,0.86776);
\end{scope}
\end{tikzpicture}
\end{sourceblock}
在第二种情况下我们得到：

\begin{sourceblock}
\begin{tikzpicture}
\node[inner sep=0,anchor=south west] (image) at (0,0) {\includegraphics[page=329,trim={53.00bp 459.13bp 53.87bp 169.58bp},clip,width=0.92\textwidth,max height=0.38\textheight]{Ferziger4th.pdf}};
\begin{scope}[x={(image.south east)},y={(image.north west)}]
\fill[white] (0.28653,0.43174) rectangle (0.30180,0.66070);
\fill[white] (0.26084,0.27785) rectangle (0.28728,0.58696);
\fill[white] (0.28412,0.21988) rectangle (0.30057,0.44910);
\fill[white] (0.30746,0.28507) rectangle (0.33564,0.59391);
\fill[white] (0.33916,0.28507) rectangle (0.36734,0.59391);
\end{scope}
\end{tikzpicture}
\end{sourceblock}
第一种方法是完全保守的。如果使用高斯定理计算导数 ∂ p /∂ x i，则第二个是保守的（并且相当于第一个）。如果通过对形函数求导来计算 CV 中心的压力导数，则这种方法通常并不保守。

其他体积源项以相同的方式近似：首先计算单元质心处的源项，然后将其简单地乘以单元体积。非线性源项需要线性化；这是使用第 2 节中讨论的方法完成的。 5.5.

\section*{9.8 压力校正方程}\phantomsection\addcontentsline{toc}{section}{9.8 压力校正方程}
当网格是非正交和/或非结构化时，需要修改 SIMPLE 算法（参见第 8.2.1 节）。本节介绍了该方法。前一章介绍的隐式分数步法的扩展遵循相同的思路，这里不再介绍，因为两种方法之间的差异很小，并且在第 1 节中详细描述了。 8.2.3.

对于任何网格类型，离散动量方程具有以下形式：

\begin{sourceblock}
\begin{tikzpicture}
\node[inner sep=0,anchor=south west] (image) at (0,0) {\includegraphics[page=329,trim={53.00bp 194.94bp 53.87bp 443.52bp},clip,width=0.92\textwidth,max height=0.38\textheight]{Ferziger4th.pdf}};
\begin{scope}[x={(image.south east)},y={(image.north west)}]
\fill[white] (0.31964,0.44942) rectangle (0.35901,0.95665);
\end{scope}
\end{tikzpicture}
\end{sourceblock}
注意，这里的下标 k 表示单元中心 N k 而不是面质心。

源项 Q i , P 包含离散压力梯度项。无论该术语如何近似，都可以写成：

\begin{sourceblock}
\begin{tikzpicture}
\node[inner sep=0,anchor=south west] (image) at (0,0) {\includegraphics[page=329,trim={53.00bp 110.74bp 53.87bp 517.97bp},clip,width=0.92\textwidth,max height=0.38\textheight]{Ferziger4th.pdf}};
\begin{scope}[x={(image.south east)},y={(image.north west)}]
\fill[white] (0.23344,0.25461) rectangle (0.28671,0.58723);
\fill[white] (0.29203,0.28533) rectangle (0.32021,0.59418);
\fill[white] (0.32586,0.27812) rectangle (0.36490,0.64040);
\end{scope}
\end{tikzpicture}
\end{sourceblock}
其中 δ p /δ x i 表示相对于笛卡尔坐标 x i 的离散压力导数。如果压力项以保守的方式近似（作为表面力的总和），则 CV 上的平均压力梯度可表示为：

\begin{sourceblock}
\begin{tikzpicture}
\node[inner sep=0,anchor=south west] (image) at (0,0) {\includegraphics[page=330,trim={53.00bp 579.70bp 53.87bp 50.00bp},clip,width=0.92\textwidth,max height=0.38\textheight]{Ferziger4th.pdf}};
\begin{scope}[x={(image.south east)},y={(image.north west)}]
\fill[white] (0.08626,0.44347) rectangle (0.10153,0.67892);
\fill[white] (0.06054,0.28568) rectangle (0.08701,0.60291);
\fill[white] (0.08370,0.22420) rectangle (0.11380,0.46652);
\fill[white] (0.11916,0.29281) rectangle (0.14734,0.61032);
\fill[white] (0.15086,0.29281) rectangle (0.17904,0.61032);
\end{scope}
\end{tikzpicture}
\end{sourceblock}
与往常一样，校正采用压力梯度的形式，并且压力是从通过施加连续性约束而获得的类泊松方程导出的。目标是满足连续性，即进入每个 CV 的净质量通量必须为零。为了计算质量通量，我们需要单元面中心的速度。在交错排列中，这些都是可用的。在共置网格上，它们是通过插值获得的。

它已在第 1 章中展示。如图 7 所示，当使用单元表面的插值速度来导出压力校正方程时，会产生较大的计算分子，从而导致压力和/或速度的振荡。我们描述了一种修改插值速度的方法，该方法产生紧凑的压力校正方程并避免振荡解。我们将简要描述第 2 节中提出的方法的扩展。 8.2.1 非正交网格。下面描述的方法对于动量方程中压力项的保守和非保守处理均有效，并且只需稍加修改即可应用于非正交网格上的 FD 方案。它对于任意形状的简历也有效；下面我们将考虑面 k，见图 9.19。

遵循第 2 节中描述的方法。 8.2.1，通过减去单元表面计算的压力梯度与插值梯度之间的差来校正插值单元面速度：

\begin{sourceblock}
\begin{tikzpicture}
\node[inner sep=0,anchor=south west] (image) at (0,0) {\includegraphics[page=330,trim={53.00bp 300.35bp 53.87bp 325.06bp},clip,width=0.92\textwidth,max height=0.38\textheight]{Ferziger4th.pdf}};
\begin{scope}[x={(image.south east)},y={(image.north west)}]
\fill[white] (0.43146,0.43236) rectangle (0.45125,0.71643);
\fill[white] (0.17386,0.21434) rectangle (0.20162,0.42941);
\fill[white] (0.20788,0.27007) rectangle (0.23606,0.55389);
\fill[white] (0.23958,0.26344) rectangle (0.28307,0.58507);
\fill[white] (0.26719,0.21041) rectangle (0.27783,0.42107);
\fill[white] (0.27976,0.19912) rectangle (0.30505,0.55389);
\fill[white] (0.30965,0.27007) rectangle (0.33783,0.55389);
\fill[white] (0.36460,0.24331) rectangle (0.39618,0.54726);
\end{scope}
\end{tikzpicture}
\end{sourceblock}
其中 * 表示通过使用先前外迭代的压力求解动量方程来预测的外迭代 m 处的速度。它在该节中有所体现。 8.2.1 对于二维均匀网格，应用于插值速度的修正对应于压力三阶导数乘以 ( x ) 2 的中心差近似；它检测振荡并消除它们。如果 A P 太大，则校正项可能很小并且无法发挥其作用。当使用非常小的时间步解决不稳态问题时，可能会发生这种情况，因为 A P 包含 V / t，但这种问题很少发生。校正项可以乘以常数而不影响近似的一致性。这种在同位网格上进行压力-速度耦合的方法是在 20 世纪 80 年代初开发的，通常归功于 Rhie 和 Chow (1983)。它被广泛使用并用于大多数商业 CFD 代码中。

只有法向速度分量对通过单元表面的质量通量有贡献。它取决于法线方向上的压力梯度。这允许我们为单元表面的法向速度分量 v n = v · n 编写以下表达式（尽管我们没有求解该分量的方程）：

\begin{sourceblock}
\begin{tikzpicture}
\node[inner sep=0,anchor=south west] (image) at (0,0) {\includegraphics[page=331,trim={53.00bp 576.19bp 53.87bp 50.00bp},clip,width=0.92\textwidth,max height=0.38\textheight]{Ferziger4th.pdf}};
\begin{scope}[x={(image.south east)},y={(image.north west)}]
\fill[white] (0.43949,0.44105) rectangle (0.45928,0.73041);
\fill[white] (0.17931,0.21852) rectangle (0.21179,0.43780);
\fill[white] (0.21808,0.27534) rectangle (0.24626,0.56471);
\fill[white] (0.24977,0.20726) rectangle (0.31287,0.57697);
\fill[white] (0.31744,0.27534) rectangle (0.34562,0.56471);
\fill[white] (0.37242,0.24831) rectangle (0.40400,0.55795);
\fill[white] (0.42875,0.09161) rectangle (0.46926,0.44305);
\end{scope}
\end{tikzpicture}
\end{sourceblock}
Because A u i
P 对于给定 CV 中的所有速度分量都是相同的（除了某些边界附近），可以替换 A v n
P by A u i
P。我们可以计算相邻 CV 中心处垂直于面 k 的方向上的压力导数，并将其插值到单元面中心。直接计算单元表面的法向导数需要坐标变换，这是结构化网格上的常用过程。当使用任意形状的 CV 时，我们希望避免使用坐标变换。使用形状函数是一种可能性，但它会产生复杂的压力校正方程。可以使用延迟校正方法来降低复杂性。

可以构建另一种方法，使用图 9.19 中所示的面法线上的辅助节点，如第 9 节中所述。 9.7.2 扩散通量。压力对 n 的导数可以通过中心差来近似，如下所示：

\begin{sourceblock}
\begin{tikzpicture}
\node[inner sep=0,anchor=south west] (image) at (0,0) {\includegraphics[page=331,trim={53.00bp 380.87bp 53.87bp 246.64bp},clip,width=0.92\textwidth,max height=0.38\textheight]{Ferziger4th.pdf}};
\begin{scope}[x={(image.south east)},y={(image.north west)}]
\fill[white] (0.38466,0.48149) rectangle (0.41931,0.78773);
\end{scope}
\end{tikzpicture}
\end{sourceblock}
两个辅助节点处的压力值可以使用单元中心值和梯度来计算：

\begin{sourceblock}
\includegraphics[page=331,trim={53.00bp 310.24bp 53.87bp 325.85bp},clip,width=0.92\textwidth,max height=0.38\textheight]{Ferziger4th.pdf}
\end{sourceblock}
通过这些表达式，方程。 ( 9.64 ) 变为：

\begin{sourceblock}
\begin{tikzpicture}
\node[inner sep=0,anchor=south west] (image) at (0,0) {\includegraphics[page=331,trim={53.00bp 250.19bp 53.87bp 377.31bp},clip,width=0.92\textwidth,max height=0.38\textheight]{Ferziger4th.pdf}};
\begin{scope}[x={(image.south east)},y={(image.north west)}]
\fill[white] (0.06238,0.48137) rectangle (0.09699,0.78753);
\end{scope}
\end{tikzpicture}
\end{sourceblock}
当连接节点 P 和 N k 的线与单元面正交并通过其中心时，即当 P 和 P ′ 以及 N k 和 N ′ 时，右侧第二项消失
k 一致。如果目标只是防止并置网格上的压力振荡，则仅使用等式右侧的第一项就足够了。 （9.66），即可以近似方程1 （9.63）为：

\begin{sourceblock}
\begin{tikzpicture}
\node[inner sep=0,anchor=south west] (image) at (0,0) {\includegraphics[page=331,trim={53.00bp 140.21bp 53.87bp 487.29bp},clip,width=0.92\textwidth,max height=0.38\textheight]{Ferziger4th.pdf}};
\begin{scope}[x={(image.south east)},y={(image.north west)}]
\fill[white] (0.27383,0.46014) rectangle (0.30544,0.78054);
\fill[white] (0.01729,0.30072) rectangle (0.04722,0.65166);
\fill[white] (0.03059,0.24896) rectangle (0.06307,0.47567);
\fill[white] (0.06938,0.30745) rectangle (0.09756,0.60688);
\fill[white] (0.10108,0.23706) rectangle (0.16418,0.61931);
\fill[white] (0.16878,0.30745) rectangle (0.19696,0.60688);
\fill[white] (0.20241,0.09291) rectangle (0.25979,0.42365);
\fill[white] (0.25059,0.05331) rectangle (0.26208,0.22360);
\fill[white] (0.26782,0.09446) rectangle (0.35474,0.42365);
\end{scope}
\end{tikzpicture}
\end{sourceblock}
请注意，表面处的插值压力梯度应根据两个单元中心处的值（按 1/2 加权）计算，而不是根据到单元表面的距离进行插值。原因是在面部计算的梯度是二阶准确的，不是在面部而是在细胞中心之间的中间；为了确保在压力变化平滑的情况下校正项相互抵消，应将单元中心的梯度插值到同一位置，即简单平均。

使用插值速度计算的质量通量，

\begin{sourceblock}
\includegraphics[page=332,trim={53.00bp 497.06bp 53.87bp 150.71bp},clip,width=0.92\textwidth,max height=0.38\textheight]{Ferziger4th.pdf}
\end{sourceblock}
不满足连续性要求，因此它们在 CV 所有面上的总和会产生质量源：

\begin{sourceblock}
\begin{tikzpicture}
\node[inner sep=0,anchor=south west] (image) at (0,0) {\includegraphics[page=332,trim={53.00bp 440.41bp 53.87bp 198.53bp},clip,width=0.92\textwidth,max height=0.38\textheight]{Ferziger4th.pdf}};
\begin{scope}[x={(image.south east)},y={(image.north west)}]
\fill[white] (0.00000,0.89338) rectangle (0.02746,1.00000);
\fill[white] (0.03011,0.89338) rectangle (0.04824,1.00000);
\fill[white] (0.05092,0.89338) rectangle (0.11564,1.00000);
\fill[white] (0.11832,0.89338) rectangle (0.20965,1.00000);
\end{scope}
\end{tikzpicture}
\end{sourceblock}
必须减少到零。必须校正速度，以便在每个 CV 中满足质量守恒。在隐式方法中，不必在每次外迭代结束时精确地满足质量守恒。按照前面描述的方法，我们通过压力校正的梯度来表达速度校正来校正质量通量，因此：

\begin{sourceblock}
\begin{tikzpicture}
\node[inner sep=0,anchor=south west] (image) at (0,0) {\includegraphics[page=332,trim={53.00bp 258.18bp 53.87bp 308.41bp},clip,width=0.92\textwidth,max height=0.38\textheight]{Ferziger4th.pdf}};
\begin{scope}[x={(image.south east)},y={(image.north west)}]
\fill[white] (0.10917,0.72858) rectangle (0.14307,0.86469);
\fill[white] (0.13257,0.70849) rectangle (0.14671,0.79458);
\fill[white] (0.15296,0.73129) rectangle (0.18114,0.84741);
\fill[white] (0.18466,0.72858) rectangle (0.23862,0.86469);
\fill[white] (0.51534,0.48599) rectangle (0.53513,0.60211);
\fill[white] (0.31964,0.41949) rectangle (0.38198,0.53571);
\fill[white] (0.40890,0.41688) rectangle (0.43200,0.53300);
\fill[white] (0.43826,0.40864) rectangle (0.47988,0.53571);
\fill[white] (0.32517,0.18081) rectangle (0.42472,0.32607);
\fill[white] (0.43047,0.19769) rectangle (0.44361,0.31381);
\fill[white] (0.44547,0.18543) rectangle (0.49423,0.31381);
\end{scope}
\end{tikzpicture}
\end{sourceblock}
如果在其他 CV 面上应用相同的近似，则要求校正后的质量通量满足连续性方程：

\begin{sourceblock}
\includegraphics[page=332,trim={53.00bp 189.52bp 53.87bp 449.42bp},clip,width=0.92\textwidth,max height=0.38\textheight]{Ferziger4th.pdf}
\end{sourceblock}
我们得到压力校正方程。

方程（9.70）右侧的最后一项导致压力校正方程中的扩展计算分子。由于当非正交性不严重时此项很小，因此通常的做法是忽略它。当解收敛时，压力修正变为零，因此省略此项不会影响解；然而，它确实影响收敛速度。对于基本上非正交的网格，必须使用较小的欠松弛参数α p，参见方程(7.84)。

当使用上述近似时，压力校正方程具有通常的形式；此外，它的系数矩阵是对称的，因此可以使用对称矩阵的特殊求解器（例如，共轭梯度系列的 ICCG 求解器；参见第 5 章和代码存储库中的目录求解器；参见附录 A.1）。

在求解压力校正方程后，使用方程校正通过单元表面的质量通量。 ( 9.70 )，导致外迭代 m 的最终值：

\begin{sourceblock}
\includegraphics[page=333,trim={53.00bp 497.06bp 53.87bp 150.71bp},clip,width=0.92\textwidth,max height=0.38\textheight]{Ferziger4th.pdf}
\end{sourceblock}
细胞中心速度和压力通过以下方式校正：

\begin{sourceblock}
\begin{tikzpicture}
\node[inner sep=0,anchor=south west] (image) at (0,0) {\includegraphics[page=333,trim={53.00bp 436.76bp 53.87bp 191.92bp},clip,width=0.92\textwidth,max height=0.38\textheight]{Ferziger4th.pdf}};
\begin{scope}[x={(image.south east)},y={(image.north west)}]
\fill[white] (0.15934,0.27896) rectangle (0.19585,0.63561);
\fill[white] (0.33263,0.08996) rectangle (0.37200,0.46476);
\end{scope}
\end{tikzpicture}
\end{sourceblock}
可以在压力校正方程中迭代地考虑网格非正交性，即通过使用预测器校正器方法。首先求解 p ' 的方程，其中方程中的非正交项。 ( 9.70 ) 被忽略。在第二步中，通过添加另一个更正来纠正第一步中所犯的错误：

\begin{sourceblock}
\begin{tikzpicture}
\node[inner sep=0,anchor=south west] (image) at (0,0) {\includegraphics[page=333,trim={53.00bp 341.52bp 53.87bp 286.88bp},clip,width=0.92\textwidth,max height=0.38\textheight]{Ferziger4th.pdf}};
\begin{scope}[x={(image.south east)},y={(image.north west)}]
\fill[white] (0.00000,0.75146) rectangle (0.13386,1.00000);
\fill[white] (0.19041,0.28431) rectangle (0.22430,0.64335);
\fill[white] (0.21380,0.23105) rectangle (0.22794,0.45840);
\fill[white] (0.23251,0.29120) rectangle (0.26069,0.59777);
\fill[white] (0.26256,0.28431) rectangle (0.30211,0.64335);
\end{scope}
\end{tikzpicture}
\end{sourceblock}
通过忽略第二个校正 p '' 中的非正交项，但将它们考虑到第一个校正 p ' 中，因为 p ' 现在可用，得出第二个质量通量校正的以下表达式：

\begin{sourceblock}
\begin{tikzpicture}
\node[inner sep=0,anchor=south west] (image) at (0,0) {\includegraphics[page=333,trim={53.00bp 224.40bp 53.87bp 373.24bp},clip,width=0.92\textwidth,max height=0.38\textheight]{Ferziger4th.pdf}};
\begin{scope}[x={(image.south east)},y={(image.north west)}]
\fill[white] (0.17480,0.60555) rectangle (0.21432,0.80336);
\fill[white] (0.25955,0.26263) rectangle (0.35910,0.47372);
\fill[white] (0.36484,0.28701) rectangle (0.37798,0.45591);
\fill[white] (0.37985,0.26934) rectangle (0.42863,0.45591);
\end{scope}
\end{tikzpicture}
\end{sourceblock}
现在可以显式计算右侧第二项，因为 p ′
可用。

\begin{sourceblock}
\begin{tikzpicture}
\node[inner sep=0,anchor=south west] (image) at (0,0) {\includegraphics[page=333,trim={53.00bp 165.32bp 53.87bp 482.80bp},clip,width=0.92\textwidth,max height=0.38\textheight]{Ferziger4th.pdf}};
\begin{scope}[x={(image.south east)},y={(image.north west)}]
\fill[white] (0.03528,0.86737) rectangle (0.13988,1.00000);
\fill[white] (0.14027,0.86737) rectangle (0.18168,1.00000);
\fill[white] (0.18208,0.86737) rectangle (0.29826,1.00000);
\fill[white] (0.29865,0.86737) rectangle (0.37453,1.00000);
\fill[white] (0.37498,0.87292) rectangle (0.44562,1.00000);
\fill[white] (0.44749,0.87292) rectangle (0.48138,1.00000);
\fill[white] (0.48319,0.86737) rectangle (0.54623,1.00000);
\fill[white] (0.54659,0.86737) rectangle (0.63958,1.00000);
\fill[white] (0.63997,0.86737) rectangle (0.72129,1.00000);
\fill[white] (0.72165,0.86737) rectangle (0.74977,1.00000);
\fill[white] (0.75014,0.86737) rectangle (0.83317,1.00000);
\fill[white] (0.83350,0.86737) rectangle (0.87492,1.00000);
\fill[white] (0.87528,0.86737) rectangle (1.00000,1.00000);
\fill[white] (0.00000,0.20422) rectangle (0.27844,0.84573);
\fill[white] (0.34896,0.09822) rectangle (0.36310,0.57436);
\fill[white] (0.37011,0.22420) rectangle (0.39829,0.86626);
\fill[white] (0.40180,0.20422) rectangle (1.00000,0.84573);
\fill[white] (0.00000,0.00000) rectangle (0.03907,0.18202);
\fill[white] (0.04180,0.00000) rectangle (0.08322,0.18202);
\fill[white] (0.08595,0.00000) rectangle (0.17395,0.18202);
\fill[white] (0.17672,0.00000) rectangle (0.28132,0.18202);
\fill[white] (0.28406,0.00000) rectangle (0.41032,0.18202);
\fill[white] (0.41651,0.00000) rectangle (0.45753,0.27026);
\fill[white] (0.46027,0.00000) rectangle (0.53829,0.18202);
\fill[white] (0.54102,0.00000) rectangle (0.58577,0.18202);
\fill[white] (0.58854,0.00000) rectangle (0.62995,0.18202);
\fill[white] (0.63266,0.00000) rectangle (0.69907,0.18202);
\fill[white] (0.70177,0.00000) rectangle (0.78478,0.18202);
\fill[white] (0.79053,0.00000) rectangle (0.81365,0.18757);
\fill[white] (0.81693,0.00000) rectangle (0.84671,0.18202);
\fill[white] (0.84944,0.00000) rectangle (0.89086,0.18202);
\fill[white] (0.89359,0.00000) rectangle (1.00000,0.18202);
\end{scope}
\end{tikzpicture}
\end{sourceblock}
现在应该强制执行 c = 0。这导致第二个压力修正 p ′′ 的方程，其具有与 p ′ 方程相同的矩阵 A，但右侧不同。这可以在某些求解器中被利用。第二次压力修正的源项包含 ̇ m ′′ 显式部分的散度。

可以通过引入第三、第四等校正来继续校正过程。附加修正趋于零；很少有必要超出 SIMPLE 中已经描述为压力校正方程的两个方程

\begin{sourceblock}
\includegraphics[page=334,trim={53.00bp 493.93bp 53.87bp 52.00bp},clip,width=\textwidth,max height=0.38\textheight]{Ferziger4th.pdf}
\par\vspace{0.45\baselineskip}
\small 图 9.20 侧壁倾斜 45°、Re = 1000（左）的盖驱动腔体中的几何形状和预测流线，以及使用 α u = 0 时的迭代次数。 8 和一个或两个压力校正步骤，作为 α p 的函数（右）
\end{sourceblock}
算法包括比网格非正交性影响的非精确处理更严格的近似。

如果网格接近正交，则包含第二次压力校正对算法的性能影响较小。然而，如果 n 和 d 之间的角度（见图 9.19）在大部分域中大于 45°，则仅进行一次校正收敛可能会很慢。强烈的欠松弛（仅增加 5-10\%
p ' 到 pm − 1 ）和速度欠松弛因子的减少可能会有所帮助，但会以效率为代价。通过两个压力校正步骤，非正交网格也可以获得正交网格上的性能。

\begin{center}\small 图 9.20 显示了未经第二次校正的非正交网格上性能下降的示例。侧壁倾斜 45° 的盖驱动腔体中的流动是在 Re = 1000 时计算的；图 9.20 还显示了几何形状和计算的流线。网格线与墙壁平行。通过第二次压力校正，收敛所需的迭代次数及其对压力欠松弛因子 α p 的依赖性与正交网格的结果相似；见图8.13。如果不包括第二次校正，则可用参数α p 的范围非常窄并且需要更多迭代。对于速度 α u 的欠松弛因子的其他值，获得了类似的结果，对于较大的 α u 值，差异更大。当网格线之间的角度减小并且仅计算一次压力校正时，获得收敛的α p 的范围变得更窄。\end{center}
这里描述的方法是在目录 2dgl 中找到的代码中实现的（参见附录）。在结构化网格上，可以将单元面法向压力导数转化为沿网格线方向的导数组合，得到混合导数的压力修正方程；见第 9.5 节.如果隐式处理交叉导数，则压力校正方程的计算分子在 2D 中至少包含 9 个节点，在 3D 中包含 19 个节点。上述两步过程产生与使用隐式离散交叉导数类似的收敛特性（参见 Peri ́c 1990），但计算效率更高，尤其是在 3D 中。

\section*{9.9 轴对称问题}\phantomsection\addcontentsline{toc}{section}{9.9 轴对称问题}
轴对称流动相对于笛卡尔坐标是三维的，即速度分量是所有三个坐标的函数；然而，它们在圆柱坐标系中只是二维的（相对于圆周方向的所有导数均为零，并且所有三个速度分量仅是轴向和径向坐标 z 和 r 的函数）。在没有涡流的情况下，圆周速度分量到处都为零。由于使用两个自变量比使用三个自变量容易得多，因此对于轴对称流动，在圆柱坐标系而不是笛卡尔坐标系中工作更有意义。

在微分形式中，以圆柱坐标系编写的质量和动量的二维守恒方程如下（例如，参见 Bird 等人，2006 年）：

\begin{sourceblock}
\includegraphics[page=335,trim={53.00bp 388.95bp 53.87bp 214.25bp},clip,width=0.92\textwidth,max height=0.38\textheight]{Ferziger4th.pdf}
\end{sourceblock}
\begin{sourceblock}
\includegraphics[page=335,trim={53.00bp 323.08bp 53.87bp 287.64bp},clip,width=0.92\textwidth,max height=0.38\textheight]{Ferziger4th.pdf}
\end{sourceblock}
\begin{sourceblock}
\includegraphics[page=335,trim={53.00bp 257.27bp 53.87bp 353.50bp},clip,width=0.92\textwidth,max height=0.38\textheight]{Ferziger4th.pdf}
\end{sourceblock}
其中非零应力张量分量为：

\begin{sourceblock}
\begin{tikzpicture}
\node[inner sep=0,anchor=south west] (image) at (0,0) {\includegraphics[page=335,trim={53.00bp 148.36bp 53.87bp 436.08bp},clip,width=0.92\textwidth,max height=0.38\textheight]{Ferziger4th.pdf}};
\begin{scope}[x={(image.south east)},y={(image.north west)}]
\fill[white] (0.13035,0.42534) rectangle (0.16866,0.57785);
\fill[white] (0.17642,0.43635) rectangle (0.20460,0.57785);
\fill[white] (0.20881,0.43182) rectangle (0.29726,0.66010);
\fill[white] (0.32968,0.18984) rectangle (0.34821,0.33146);
\fill[white] (0.13182,0.09131) rectangle (0.16863,0.24602);
\fill[white] (0.17642,0.10453) rectangle (0.20460,0.24602);
\fill[white] (0.20881,0.09131) rectangle (0.24562,0.24602);
\fill[white] (0.25275,0.10453) rectangle (0.28093,0.24602);
\fill[white] (0.28445,0.10122) rectangle (0.32084,0.24602);
\end{scope}
\end{tikzpicture}
\end{sourceblock}
正如第 3 节中所讨论的那样。 9.3.1，上述方程包含两项在直角坐标系中没有类似物的项：表观离心力 ρ v 2
v r 方程中的 θ / r，以及 v θ 方程中的视在科里奥利力 ρ v r v θ / r。这些术语来自坐标变换，不应与旋转坐标系中出现的离心力和科里奥利力混淆。如果涡流速度 v θ 为零，则视在力为零，并且第三个方程变得多余。

当使用FD方法时，关于轴向和径向坐标的导数以与笛卡尔坐标中相同的方式近似；第 1 章中描述的任何方法。可以使用3个。

有限体积方法需要小心。先前给出的积分形式的守恒方程（例如，（8.1）和（8.2））保持不变，并添加了表观力作为源项。这些被集成在卷上，如第 1 节中所述。 4.3. θ方向上的CV大小是统一的，即一个弧度。需要注意压力术语。如果将这些视为体积力，并且 z 和 r 方向上的压力导数在体积上积分，如方程式所示： ( 9.58 )，不需要额外的步骤。然而，如果压力在 CV 表面上积分，如式（1）所示。 ( 9.57 )，仅在北、南、西、东单元面上进行积分是不够的，就像平面二维问题的情况一样，必须考虑前表面和后表面上的压力的径向分量。

因此，我们必须将这些在平面二维问题中没有对应项的项添加到 v r 的动量方程中：

\begin{sourceblock}
\begin{tikzpicture}
\node[inner sep=0,anchor=south west] (image) at (0,0) {\includegraphics[page=336,trim={53.00bp 363.89bp 53.87bp 264.47bp},clip,width=0.92\textwidth,max height=0.38\textheight]{Ferziger4th.pdf}};
\begin{scope}[x={(image.south east)},y={(image.north west)}]
\fill[white] (0.38689,0.46983) rectangle (0.41002,0.77581);
\fill[white] (0.21973,0.28507) rectangle (0.25555,0.63870);
\fill[white] (0.26265,0.29195) rectangle (0.29083,0.59820);
\fill[white] (0.29435,0.29195) rectangle (0.36180,0.78295);
\end{scope}
\end{tikzpicture}
\end{sourceblock}
其中 S 是正面面积。

在 v θ 的方程中，如果需要求解，则必须包含源项（表观科里奥利力）：

\begin{sourceblock}
\begin{tikzpicture}
\node[inner sep=0,anchor=south west] (image) at (0,0) {\includegraphics[page=336,trim={53.00bp 285.97bp 53.87bp 347.78bp},clip,width=0.92\textwidth,max height=0.38\textheight]{Ferziger4th.pdf}};
\begin{scope}[x={(image.south east)},y={(image.north west)}]
\fill[white] (0.35110,0.25780) rectangle (0.38842,0.67644);
\fill[white] (0.39552,0.26613) rectangle (0.42370,0.62303);
\fill[white] (0.42719,0.26613) rectangle (0.45537,0.62303);
\end{scope}
\end{tikzpicture}
\end{sourceblock}
与平面二维问题相比，唯一的区别是单元面面积和体积的计算。单元面“n”、“e”、“w”和“s”的面积按照平面几何计算，参见方程 1。 ( 9.29 )，其中包含因子 r k （其中 k 表示单元面中心）。正面和背面的面积的计算方式与平面几何中的体积相同（其中第三维是统一的）。具有任意数量面的轴对称 CV 的体积为：

\begin{sourceblock}
\begin{tikzpicture}
\node[inner sep=0,anchor=south west] (image) at (0,0) {\includegraphics[page=336,trim={53.00bp 161.50bp 53.87bp 469.00bp},clip,width=0.92\textwidth,max height=0.38\textheight]{Ferziger4th.pdf}};
\begin{scope}[x={(image.south east)},y={(image.north west)}]
\fill[white] (0.53429,0.29293) rectangle (0.57380,0.54097);
\fill[white] (0.57714,0.35887) rectangle (0.60532,0.68322);
\fill[white] (0.60653,0.35129) rectangle (0.63654,0.72419);
\end{scope}
\end{tikzpicture}
\end{sourceblock}
其中 N v 表示顶点数，逆时针计数，i = 0 对应于 i = N v。

轴对称旋流中的一个重要问题是径向和周向速度分量的耦合。 v r 的方程包含 v 2
θ，v θ 的方程包含 v r 和 v θ 的乘积作为源项；见上文。顺序（解耦）求解过程和皮卡德线性化的组合可能效率低下。可以通过在外部迭代中使用多重网格方法来改善耦合，请参见在第 12 章中，通过耦合求解方法，或者通过使用更隐式的线性化方案，参见第12节。 5.5.

如果将柱坐标系的坐标 z 和 r 替换为 x 和 y，则与笛卡尔坐标系中的方程的类比变得明显。事实上，如果 r 设置为 1，且 v θ 和 τ θθ 设置为零，这些方程将与笛卡尔坐标中的方程相同，其中 v z = u x 且 v r = u y。因此，相同的计算机代码可用于平面和轴对称二维流动；对于轴对称问题，设置 r = y 并包括 τ θθ，如果旋流分量非零，则包括 v θ 方程。

\section*{9.10 高阶有限体积方法}\phantomsection\addcontentsline{toc}{section}{9.10 高阶有限体积方法}
值得注意的是，高阶FV方法的推导比高阶FD方法的构造更加困难。在FD方法中，我们只需要用高阶近似来近似网格点处的一阶和二阶导数，这在结构化网格上相对容易做到（参见第3章）。在 FV 方法中，存在三种近似：

\begin{itemize}
\item 表面和体积积分的近似；
\item 将变量值插值到 CV 中心以外的位置；
\item 单元中心和所有单元面的一阶导数的近似。
\end{itemize}
中点法则的二阶精度是单点近似所能达到的最高精度。任何高阶 FV 方案都需要对多个单元面位置进行高阶插值，以及更复杂的多点积分近似来计算对流通量。此外，对于扩散通量，还需要以更高阶逼近单元面内多个位置处的导数。这在结构化网格上是可以管理的，但在非结构化网格上相当困难，尤其是由任意多面体 CV 组成的网格。为了实现、扩展、调试和维护的简单性，二阶精度似乎是精度和效率之间的最佳折衷。

仅当需要非常高的精度（离散化误差低于 1\%）时，高阶方法才具有成本效益。人们还必须记住，只有当网格足够细时，高阶方法才能比二阶方法产生更准确的结果。如果网格不够精细，高阶方法可能会产生振荡解，并且平均误差可能高于二阶方案。与二阶方案相比，高阶方案还需要每个网格点更多的内存和计算时间。对于可接受 1\% 左右误差的工业应用，二阶方案与局部网格细化相结合，提供了准确性、编程和代码维护简单性、稳健性和效率的最佳组合。

正如已经提到的，使用 FD 或 FE 方法比使用 FV 方法更容易实现高阶方法。高阶有限元方法在结构力学中非常标准，特别是对于线性问题。不可压缩流的连续性方程会带来麻烦，这就是为什么 CFD 的有限元方法通常使用动量和连续性方程的不等阶近似值。用于非结构化网格的 FD 方法如今也不常见，但我们预计在不久的将来会看到这种高阶方法。

\section*{9.11 边界条件的实现}\phantomsection\addcontentsline{toc}{section}{9.11 边界条件的实现}
在非正交网格上实施边界条件需要特别注意，因为边界通常不与笛卡尔速度分量对齐。 FV 方法要求边界通量已知或以已知量和内部节点值表示。当然，CV 的数量必须与未知数的数量相匹配。

我们经常会提到局部坐标系 ( n , t , s )，它是一个旋转的笛卡尔坐标系，其中 n 是边界的向外法线，t 和 s 与边界相切。

\subsection*{9.11.1 入口}\phantomsection\addcontentsline{toc}{subsection}{9.11.1 入口}
通常，在入口边界处，必须规定所有数量。如果入口处的条件不太清楚并且需要对可变轮廓进行近似，则将边界向上游移动尽可能远离感兴趣的区域是有用的。由于速度和其他变量已给定，因此可以直接计算所有对流通量。扩散通量通常是未知的，但可以使用已知的变量边界值和梯度的单侧有限差分近似来近似。

如果在边界处指定速度，则不需要在迭代期间进行校正。如果边界处的速度修正为零，则转换为 SIMPLE 算法中压力修正的零梯度条件；参见等式。 （9.70）。因此，压力校正方程在指定速度的所有边界处都具有诺伊曼边界条件。

\subsection*{9.11.2 插座}\phantomsection\addcontentsline{toc}{subsection}{9.11.2 插座}
在出口处，我们通常对流量知之甚少。因此，这些边界应尽可能位于感兴趣区域的下游。否则，错误可能会向上游传播。流动应在整个出口横截面上被引导出域，并且如果可能的话，与出口边界平行和正交。在高雷诺数流中，误差的上游传播（至少在稳态流中）很弱，因此很容易找到边界条件的合适近似值。通常，人们沿着网格线从内部到边界（或者更好的是沿着流线）进行推断。最简单的近似是沿网格线的零梯度。对于对流通量，这意味着使用一阶迎风近似。沿网格线的零梯度条件可以很容易地隐式实现。例如，在结构化二维网格的东面，一阶向后近似给出 φ E = φ P。当我们将此表达式插入到边界旁边的 CV 的离散方程中时，我们有：

\begin{sourceblock}
\includegraphics[page=339,trim={53.00bp 462.38bp 53.87bp 188.30bp},clip,width=0.92\textwidth,max height=0.38\textheight]{Ferziger4th.pdf}
\end{sourceblock}
因此方程中没有出现边界值 φ E。这并不意味着扩散通量在出口边界处为零，除非网格与边界正交。

如果需要更高的精度，则必须使用出口边界处导数的高阶单侧有限差分近似。对流通量和扩散通量都必须用内部节点的变量值来表示。

如果将速度外推到出口边界，则通常会对其进行修正（假设流动不可压缩），以便流出质量通量与流入质量通量相匹配。然后，连续性方程得到全局满足，即，当所有 CV 的质量守恒方程求和时，所有内部单元表面上的质量通量相互抵消，并且当边界通量匹配时，净质量通量为零。以这种方式校正边界速度的结果是，它们可以被视为当前外迭代的固定值，从而导致压力校正方程的零梯度条件。如果诺依曼条件适用于所有边界，则必须确保压力校正方程中源项的代数和等于零，否则该公式不成立。上述出口速度的修正保证了该条件的满足。然而，当诺依曼条件适用于所有边界时，压力校正方程的解不是唯一的，可以向所有值添加一个常数，方程仍然满足。因此，通常将压力固定在参考位置，并通过添加在给定网格点计算的压力校正与在参考位置计算的压力校正之间的差值来校正压力。

当流动不稳定时，特别是直接模拟湍流时，需要注意避免在出口边界反映误差。这些问题见第 10.2 节 和 13.6。

\subsection*{9.11.3 不透水墙}\phantomsection\addcontentsline{toc}{subsection}{9.11.3 不透水墙}
\begin{sourceblock}
\includegraphics[page=340,trim={53.00bp 542.09bp 53.87bp 108.60bp},clip,width=0.92\textwidth,max height=0.38\textheight]{Ferziger4th.pdf}
\end{sourceblock}
该条件源自粘性流体粘附到固体边界（无滑移条件）这一事实。

由于没有流过壁面，所有量的对流通量均为零。扩散通量需要一些注意。对于诸如热能之类的标量，它们可以为零（绝热壁），可以指定它们（指定热通量），或者可以指定标量值（等温壁）。如果通量已知，则可以将其插入近壁 CV 的守恒方程中，例如，对于标记为“s”的边界面：

\begin{sourceblock}
\begin{tikzpicture}
\node[inner sep=0,anchor=south west] (image) at (0,0) {\includegraphics[page=340,trim={53.00bp 397.26bp 53.87bp 231.77bp},clip,width=0.92\textwidth,max height=0.38\textheight]{Ferziger4th.pdf}};
\begin{scope}[x={(image.south east)},y={(image.north west)}]
\fill[white] (0.08941,0.27189) rectangle (0.12577,0.63029);
\fill[white] (0.10629,0.21611) rectangle (0.11925,0.44705);
\fill[white] (0.13122,0.27917) rectangle (0.15940,0.59068);
\fill[white] (0.17982,0.03234) rectangle (0.20060,0.28079);
\end{scope}
\end{tikzpicture}
\end{sourceblock}
其中 f 是每单位面积的规定通量。如果 φ 的值是在壁上指定的，我们需要使用单侧差分来近似 φ 的法向梯度。根据这样的近似，我们还可以计算出指定通量时壁面处的 φ 值。存在多种可能性；一是计算位于法线n上的辅助点P′处的φ值，见图9.21，并使用近似值：

\begin{sourceblock}
\begin{tikzpicture}
\node[inner sep=0,anchor=south west] (image) at (0,0) {\includegraphics[page=340,trim={53.00bp 289.97bp 53.87bp 338.74bp},clip,width=0.92\textwidth,max height=0.38\textheight]{Ferziger4th.pdf}};
\begin{scope}[x={(image.south east)},y={(image.north west)}]
\fill[white] (0.39380,0.47181) rectangle (0.43206,0.78066);
\end{scope}
\end{tikzpicture}
\end{sourceblock}
其中，δ n = ( r S − r P ′ ) · n 为点 P ′ 和 S 之间的距离。如果非正交性不严重，可以使用 φ P 代替 φ P ′。也可以使用形状函数或从细胞中心推断的梯度。然后可以使用中点规则将通量近似为：

\begin{sourceblock}
\begin{tikzpicture}
\node[inner sep=0,anchor=south west] (image) at (0,0) {\includegraphics[page=340,trim={53.00bp 194.44bp 53.87bp 434.28bp},clip,width=0.92\textwidth,max height=0.38\textheight]{Ferziger4th.pdf}};
\begin{scope}[x={(image.south east)},y={(image.north west)}]
\fill[white] (0.27302,0.27793) rectangle (0.30938,0.63335);
\end{scope}
\end{tikzpicture}
\end{sourceblock}
动量方程中的扩散通量需要特别注意。如果我们要求解速度分量 v n 、 v t 和 v s，我们可以使用第 1 节中描述的方法。 7.1.6.壁上的粘性应力为：

\begin{sourceblock}
\begin{tikzpicture}
\node[inner sep=0,anchor=south west] (image) at (0,0) {\includegraphics[page=340,trim={53.00bp 115.90bp 53.87bp 514.95bp},clip,width=0.92\textwidth,max height=0.38\textheight]{Ferziger4th.pdf}};
\begin{scope}[x={(image.south east)},y={(image.north west)}]
\fill[white] (0.34391,0.40946) rectangle (0.38863,0.76764);
\fill[white] (0.20208,0.21139) rectangle (0.24174,0.56985);
\fill[white] (0.24773,0.24199) rectangle (0.27591,0.56985);
\fill[white] (0.27943,0.23179) rectangle (0.31991,0.56985);
\fill[white] (0.78580,0.24199) rectangle (0.79886,0.56985);
\fill[white] (0.92439,0.23179) rectangle (1.00000,0.55937);
\fill[white] (0.73489,0.00000) rectangle (0.77558,0.21678);
\end{scope}
\end{tikzpicture}
\end{sourceblock}
这里我们假设坐标 t 沿壁面剪切力的方向，因此 τ ns = 0。该力平行于速度矢量在壁面上的投影（ s 与其正交）。这相当于假设速度矢量

\begin{center}\small 图 9.21 关于墙体边界条件的实现\end{center}
\begin{sourceblock}
\begin{tikzpicture}
\node[inner sep=0,anchor=south west] (image) at (0,0) {\includegraphics[page=341,trim={53.00bp 470.62bp 53.87bp 80.66bp},clip,width=0.92\textwidth,max height=0.38\textheight]{Ferziger4th.pdf}};
\begin{scope}[x={(image.south east)},y={(image.north west)}]
\fill[white] (0.00000,0.93000) rectangle (0.11026,1.00000);
\fill[white] (0.11185,0.93000) rectangle (0.13504,1.00000);
\fill[white] (0.13657,0.93000) rectangle (0.15269,1.00000);
\fill[white] (0.15426,0.93000) rectangle (0.20268,1.00000);
\end{scope}
\end{tikzpicture}
\end{sourceblock}
不会改变第一个网格点和墙壁之间的方向，这并不完全正确，但却是一个合理的近似值，随着第一个网格点到墙壁的距离减小，它会变得更加准确。

v t 和 v n 都可以在节点 P 处轻松计算。在二维中，单位向量 t 很容易从角点“se”和“sw”的坐标获得，见图 9.21。在 3D 中，我们必须确定向量 t 的方向。根据平行于壁的速度，我们可以定义单位向量 t 如下：

\begin{sourceblock}
\includegraphics[page=341,trim={53.00bp 330.68bp 53.87bp 306.62bp},clip,width=0.92\textwidth,max height=0.38\textheight]{Ferziger4th.pdf}
\end{sourceblock}
近似应力所需的速度分量为：

\begin{sourceblock}
\includegraphics[page=341,trim={53.00bp 284.93bp 53.87bp 365.75bp},clip,width=0.92\textwidth,max height=0.38\textheight]{Ferziger4th.pdf}
\end{sourceblock}
导数可以按照方程式计算。 （9.87）。

可以变换应力τ nt 以获得τ xx 、τ xy 等，但这不是必需的。 τ nt 的表面积分产生力：

\begin{sourceblock}
\begin{tikzpicture}
\node[inner sep=0,anchor=south west] (image) at (0,0) {\includegraphics[page=341,trim={53.00bp 200.76bp 53.87bp 442.25bp},clip,width=0.92\textwidth,max height=0.38\textheight]{Ferziger4th.pdf}};
\begin{scope}[x={(image.south east)},y={(image.north west)}]
\fill[white] (0.31690,0.39818) rectangle (0.36758,0.93645);
\fill[white] (0.37281,0.44790) rectangle (0.40099,0.94812);
\end{scope}
\end{tikzpicture}
\end{sourceblock}
其 x 、 y 和 z 分量对应于离散动量方程中所需的积分；例如，在 u x 的方程中：

\begin{sourceblock}
\begin{tikzpicture}
\node[inner sep=0,anchor=south west] (image) at (0,0) {\includegraphics[page=341,trim={53.00bp 130.62bp 53.87bp 512.40bp},clip,width=0.92\textwidth,max height=0.38\textheight]{Ferziger4th.pdf}};
\begin{scope}[x={(image.south east)},y={(image.north west)}]
\fill[white] (0.11417,0.40095) rectangle (0.13832,0.93642);
\fill[white] (0.14520,0.44810) rectangle (0.17338,0.94810);
\end{scope}
\end{tikzpicture}
\end{sourceblock}
或者，我们可以使用单元中心的速度梯度（例如，使用高斯定理计算，参见方程（9.49）），将它们外推到壁单元面的中心，计算剪应力 τ xx 、 τ xy 等，并根据上述表达式计算剪力分量。

因此，我们用剪切力代替壁动量方程中的扩散通量。如果使用先前迭代的值显式计算该力，则可能会损害收敛性。如果将力写为节点 P 处笛卡尔速度分量的函数，则可以隐式处理其中的一部分。在这种情况下，所有速度分量的系数 A P 都不相同（内部单元的情况就是如此）。这是不希望的，因为压力校正方程中需要系数 A P，如果它们不同，我们就必须存储所有三个值。因此，最好使用延迟校正方法，如在室内；我们近似

\begin{sourceblock}
\begin{tikzpicture}
\node[inner sep=0,anchor=south west] (image) at (0,0) {\includegraphics[page=342,trim={53.00bp 475.69bp 53.87bp 162.17bp},clip,width=0.92\textwidth,max height=0.38\textheight]{Ferziger4th.pdf}};
\begin{scope}[x={(image.south east)},y={(image.north west)}]
\fill[white] (0.00000,0.87977) rectangle (0.03907,1.00000);
\fill[white] (0.04174,0.87977) rectangle (0.19627,1.00000);
\end{scope}
\end{tikzpicture}
\end{sourceblock}
隐式并将隐式近似与使用上述方法之一计算的力之间的差值添加到方程的右侧。这里δn是节点P到墙的距离。那么系数 A P 对于所有速度分量都是相同的，并且显式项部分抵消。收敛速度几乎不受影响。

由于速度是在壁处指定的，因此压力校正方程在壁边界处具有诺伊曼条件。

\subsection*{9.11.4 对称平面}\phantomsection\addcontentsline{toc}{subsection}{9.11.4 对称平面}
在许多流动中，存在一个或多个对称平面。当流动稳定时，存在相对于该平面对称的解（在许多情况下，例如扩散器或突然膨胀的通道，也存在非对称稳态解，通常比对称解更稳定）。对称解可以通过使用对称条件仅在部分解域中求解问题来获得。

在对称平面上，所有量的对流通量均为零。此外，平行于对称平面的速度分量和所有标量的法向梯度在那里为零。因此，所有标量的扩散通量在对称平面处为零。法向速度分量为零，但其法向方向的导数不为零；因此，法向应力 τ nn 不为零。 τ nn 的表面积分产生一个力：

\begin{sourceblock}
\begin{tikzpicture}
\node[inner sep=0,anchor=south west] (image) at (0,0) {\includegraphics[page=342,trim={53.00bp 144.02bp 53.87bp 499.00bp},clip,width=0.92\textwidth,max height=0.38\textheight]{Ferziger4th.pdf}};
\begin{scope}[x={(image.south east)},y={(image.north west)}]
\fill[white] (0.30767,0.39792) rectangle (0.35738,0.93642);
\fill[white] (0.36298,0.44810) rectangle (0.39116,0.94810);
\end{scope}
\end{tikzpicture}
\end{sourceblock}
当对称边界与笛卡尔坐标平面不重合时，所有三个笛卡尔速度分量的扩散通量将不为零。这些通量可以通过首先从 (9.95) 获得合成法向力和法向导数的近似值（如前一节所述）并将该力分解为其笛卡尔分量来计算。或者，可以推断从内部到边界的速度梯度，并使用类似于 (9.93) 的表达式，例如，对于面“s”处的 u x 分量（见图 9.21）：

\begin{sourceblock}
\begin{tikzpicture}
\node[inner sep=0,anchor=south west] (image) at (0,0) {\includegraphics[page=343,trim={53.00bp 545.60bp 53.87bp 97.42bp},clip,width=0.92\textwidth,max height=0.38\textheight]{Ferziger4th.pdf}};
\begin{scope}[x={(image.south east)},y={(image.north west)}]
\fill[white] (0.11107,0.40138) rectangle (0.13525,0.93642);
\fill[white] (0.14214,0.44810) rectangle (0.17032,0.94810);
\end{scope}
\end{tikzpicture}
\end{sourceblock}
与壁边界的情况一样，可以将对称边界处的扩散通量分解为隐式部分，其中涉及 CV 中心处的速度分量（对系数 A P 做出贡献），或者使用延迟校正方法使所有速度分量的 A P 保持相同。

同样，由于规定了法向速度分量，压力校正方程在对称边界处具有诺伊曼边界条件。

\subsection*{9.11.5 规定压力}\phantomsection\addcontentsline{toc}{subsection}{9.11.5 规定压力}
在不可压缩流中，通常指定入口处的质量流量并在出口处使用外推法。然而，在某些情况下，质量流量未知，但入口和出口之间的压降是规定的。此外，压力有时是在远场边界处指定的。

当在边界处指定压力时，无法指定速度 - 必须使用与两个 CV 之间的单元面相同的方法从内部推断，请参见方程 1。 （9.63）；唯一的区别是，现在单元面和一个邻居节点的位置重合。边界处的压力梯度使用单侧差分来近似；例如，在“e”面上，可以使用以下表达式，这是一阶向后差分：

\begin{sourceblock}
\begin{tikzpicture}
\node[inner sep=0,anchor=south west] (image) at (0,0) {\includegraphics[page=343,trim={53.00bp 247.17bp 53.87bp 381.61bp},clip,width=0.92\textwidth,max height=0.38\textheight]{Ferziger4th.pdf}};
\begin{scope}[x={(image.south east)},y={(image.north west)}]
\fill[white] (0.37062,0.46360) rectangle (0.40728,0.78025);
\end{scope}
\end{tikzpicture}
\end{sourceblock}
这样确定的边界速度需要进行修正以满足质量守恒；在指定压力的边界处，质量通量修正 ̇ m ′ 不为零。然而，边界压力未修正，即边界处 p ′ = 0。这在压力校正方程中用作狄利克雷边界条件。有关在边界指定静压时边界条件实施的更多详细信息，请参阅第 1 章。 11.

如果雷诺数较高，在指定入口和出口压力的情况下应用上述方法，求解过程将缓慢收敛。另一种可能性是首先猜测入口处的质量流量并将其视为一次外部迭代的规定，并考虑仅在出口处指定压力。然后应通过尝试将入口边界处的推断压力与指定压力相匹配来校正入口速度。使用迭代校正程序将两个压力之间的差值降至零。

\section*{9.12 例子}\phantomsection\addcontentsline{toc}{section}{9.12 例子}
在本节中，我们将介绍计算需要贴体网格的几何形状层流的示例。两个例子涉及稳态流，一个例子涉及非稳态流。在一种情况下，使用结构化网格和可以从互联网下载的代码；对于另外两个示例，使用商业 CFD 软件。这些示例的目的是演示如何解决此类流问题、如何分析解的准确性以及不同网格类型如何影响计算工作量和结果质量。

\subsection*{9.12.1 Re = 20 时圆柱周围的流动}\phantomsection\addcontentsline{toc}{subsection}{9.12.1 Re = 20 时圆柱周围的流动}
作为一个例子，我们考虑在无限环境中围绕圆柱体的第一个层流 2D 流，该环境暴露于雷诺数 Re = 20 的均匀横流。雷诺数基于均匀流速度 U ∞ 、流体粘度 μ 和圆柱体直径 D。解域是有限的，并在圆柱体的上游和下游以及其上方和下方延伸 16 D，见图 9.22，它显示了整个解域和用于这些计算的边界拟合、结构化 O 型网格。所使用的两种计算机代码（一种基于 SIMPLE 算法，另一种基于隐式分数步法（标记为 IFSM））可在 Internet 上获取；详见附录A.1。

使用五个系统细化的网格（即，每个较粗网格 CV 分为四个较细网格 CV），以便能够估计离散化误差；最粗的网格有 24 × 16 CV，最细的网格有 384 × 256 CV（周围的单元格数量）

\begin{sourceblock}
\includegraphics[page=344,trim={53.00bp 101.35bp 53.87bp 374.99bp},clip,width=\textwidth,max height=0.38\textheight]{Ferziger4th.pdf}
\par\vspace{0.45\baselineskip}
\small 图 9.22 三级网格（左）和圆柱体周围四级网格的细节（右），用于计算圆柱体周围的稳态和非稳态二维流动
\end{sourceblock}
\begin{center}\small 图 9.23 Re = 20 时圆柱体附近的预测流线\end{center}
柱面×径向细胞数）。细胞均匀分布在圆柱体周围，同时沿径向扩张；最粗网格上的扩展因子为 1.25，对于每个更细的网格，扩展因子等于前一个网格上的扩展因子的平方根。因此，在最精细的网格上，扩展因子仅为 1.014044，因为现在 256 个单元格都安装在从圆柱体到外部边界的相同距离内。由于网格是O型的，所以它们在东西边界相交处有一个接缝：它是圆柱体后面的水平中心线，如图9.22中四级网格的较粗线所示。在左边界，指定速度 U ∞ = 1 m/s（圆柱直径为 D = 1 m）；在下游侧，速度是用零梯度条件推断的；在顶部和底部边界，指定了对称条件。请注意，在使用的 O 网格中，所有这些线段都是南边界的一部分，而圆柱体表面代表北边界。二阶 CDS 用于空间离散化。

在此雷诺数下的流量是稳定的；它与圆柱体表面分离，并在圆柱体后面形成两个弱循环涡，如图9.23中的流线图和图9.24中的速度矢量图所示。均匀的流入气流被圆柱体在其周围的大区域内偏转；圆柱体到解域外边界16D的距离可能不足以代表均匀流的真正无扰动远场条件，但预计外边界对圆柱体周围流动的影响不会太强。无论如何，当我们确定下面的离散化误差时，它们仅对受指定边界条件影响的流动有效。

\begin{center}\small 图 9.24 还显示了压力等值线；它们清楚地表明，即使在相对较大的距离处也能感觉到圆柱体的存在，因为压力梯度不为零。正如预期的那样，最高压力出现在前驻点。密集的等压线表明圆柱表面上部和下部的压力快速下降，直到赤道后不久达到最小值。从那时起，压力再次开始增加，这种不利的压力梯度导致流动分离并在气缸后面形成再循环区。\end{center}
\begin{sourceblock}
\includegraphics[page=346,trim={53.00bp 452.25bp 53.87bp 52.00bp},clip,width=\textwidth,max height=0.38\textheight]{Ferziger4th.pdf}
\par\vspace{0.45\baselineskip}
\small 图 9.24 Re = 20 时气缸附近的预测压力等值线（左）和速度矢量（右）；仅绘制每四个向量。最大速度矢量显示在插图中的比例尺上
\end{sourceblock}
\begin{sourceblock}
\includegraphics[page=346,trim={53.00bp 260.97bp 53.87bp 247.37bp},clip,width=\textwidth,max height=0.38\textheight]{Ferziger4th.pdf}
\par\vspace{0.45\baselineskip}
\small 图 9.25 Re = 20 时圆柱体上的预测压力（左）和剪切力（右）
\end{sourceblock}
在无粘流中，压力轮廓将围绕垂直中心线完全对称，后驻点处的压力与前部相同，从而导致零阻力并且无再循环。由于流体粘度，压力（垂直于壁）和剪切力（垂直于壁）在圆柱体表面上积分时，会导致流动方向上的非零净力分量。由于稳态流绕水平中心线对称，因此升力为零。压力和剪切阻力分量向网格无关解的收敛如图 9.25 所示。

在每个网格上，进行外迭代，直到残差范数（所有CV处的残差绝对值之和）减少七个数量级；这是非常必要的，但我们希望确保在确定离散化误差时迭代误差可以忽略不计。 SIMPLE 和 IFSM 的结果

\begin{sourceblock}
\includegraphics[page=347,trim={53.00bp 453.25bp 53.87bp 52.00bp},clip,width=\textwidth,max height=0.38\textheight]{Ferziger4th.pdf}
\par\vspace{0.45\baselineskip}
\small 图 9.26 Re = 20 时圆柱体上的压力（左）和剪切力（右）的离散化误差，使用 Richardson 外推技术估计
\end{sourceblock}
方法收敛到相同的网格独立解，但使用 SIMPLE 算法时，粗网格上的误差更小。有趣的是，在图 9.25 中，粗网格上的压力被高估，而剪切力被低估。因此，总力的相对误差低于分力的误差，因为它们的符号相反，因此部分抵消。不同来源的错误相互抵消的情况并不罕见，但它们也可能会增加。因此，检查解的网格依赖性始终很重要。

\begin{center}\small 图 9.26 显示了使用 Richardson 外推法估计的离散化误差（详细信息请参见第 3.9 节）。正如预期的那样，获得了向独立于网格的解的二阶收敛。在最粗糙的网格上，压力和剪切力都有很大的误差（分别约为 10\% 和 3\%）；在最精细的网格上，误差相当低（分别为 0.03\% 和 0.007\%）。四级网格提供的解误差约为 0.1\%，这对于大多数应用来说已经足够小了。\end{center}
对于物体周围的流动，阻力和升力系数定义为：

\begin{sourceblock}
\includegraphics[page=347,trim={53.00bp 175.97bp 53.87bp 458.52bp},clip,width=0.92\textwidth,max height=0.38\textheight]{Ferziger4th.pdf}
\end{sourceblock}
其中 F x 和 F y 是流体对身体施加的力的 x 和 y 分量，S 是垂直于流动方向的身体横截面积。在二维流动模拟中，z 方向的尺寸被假定为统一，并且因为在我们的模拟中使用了 D = 1，所以面积也等于 1m 2。未扰动流的速度也取为 U ∞ = 1m/s。因此，如果流体密度ρ = 1 kg/m 3，则只需将计算出的力乘以2即可获得阻力和升力系数。由此由理查森外推得到的总力得到的阻力系数为2.083，这与实验和数值研究文献中的数据非常一致。

达到规定的残差水平所需的外部迭代次数（在 SIMPLE 中）或时间步长（在 IFSM 中）对速度欠松弛因子（在 SIMPLE 中）或时间步长（在 IFSM 中）的依赖性与上一章中在稳态流中观察到的情况非常相似，参见图 1 和 2.8.14 和 8.18。因此，我们没有针对这种情况显示此类图表，但可以说，在最佳设置下，SIMPLE 和 IFSM 中所需的计算工作是相似的。

当雷诺数增大时，圆柱体后面的两个涡流变得更长、更强；保持两个涡流完全相等并且在 Re = 45 左右变得很困难，即使是最小的扰动也会导致一个涡流变得更大并且流动失去对称性。一旦对称性被破坏，流动就会变得不稳定；涡流开始相互生长，并从圆柱体的每一侧交替分离，形成著名的冯卡门涡街。实验和模拟都表明，在 Re = 200 左右，流动变得三维，不再是完美的周期性；最后，当雷诺数更大时，尾流变得湍急。当然，这只能在 3D 模拟中看到。在下一节中，我们将仔细研究 Re = 200 时圆柱体周围的非定常流动。

\subsection*{9.12.2 Re = 200 时圆柱体周围的流动}\phantomsection\addcontentsline{toc}{subsection}{9.12.2 Re = 200 时圆柱体周围的流动}
我们在这里展示了使用商用流动求解器 STAR-CCM+ 和三个系统细化网格（在每个细化步骤中，网格间距在两个方向上减半）在 Re = 200 下进行的二维模拟的一些结果。在圆柱体周围创建 10 个棱镜层；棱镜层的外表面切割笛卡尔网格，填充剩余空间。域大小与前面的示例相同：它是矩形，并在正负 x 和 y 方向上远离圆柱体中心延伸 16 D。笛卡尔网格分 4 个步骤进行局部细化，使得圆柱体周围及其尾流的矩形区域中的单元尺寸为“基本尺寸”的 6.25\%；图 9.27 显示了包含最细区域的最粗网格部分。最粗的网格有 6,196 个 CV；中网格有 21,744 CV；最细的网格有 109,320 CV。棱镜层的厚度与基底尺寸成正比，因此每次细化网格时棱镜层的厚度都会减半。在最精细的网格上，沿着圆柱体周边大约有 220 个单元。模拟文件和更详细的模拟报告可以在互联网上找到；详情请参见附录。

仅使用二阶时间离散化（三时间级方案；参见第6.3.2.4节）；它在该节中得到了证明。 6.4 当需要时间精确解时，一阶时间推进方案不适合模拟非定常问题。使用四个时间步长来测试解的时间步依赖性：0.04 s、0.02 s、0.01 s 和 0.005 s，对应于约。每个阻力振荡周期分别有 63、126、253 和 506 个时间步长。

\begin{sourceblock}
\includegraphics[page=349,trim={53.00bp 428.48bp 53.87bp 52.00bp},clip,width=\textwidth,max height=0.38\textheight]{Ferziger4th.pdf}
\par\vspace{0.45\baselineskip}
\small 图 9.27 用于模拟 Re = 200 时流动的圆柱体周围最粗糙网格的细节
\end{sourceblock}
Re = 200 时圆柱体周围的流动高度不稳定，强烈的涡流以规则的频率从圆柱体中脱落。在最精细的网格上，该频率被确定为 f = 0.1977 年，对应于无量纲斯特劳哈尔数（因为在我们的例子中 D 和 U ∞ 都等于 1）：

\begin{sourceblock}
\begin{tikzpicture}
\node[inner sep=0,anchor=south west] (image) at (0,0) {\includegraphics[page=349,trim={53.00bp 307.40bp 53.87bp 330.09bp},clip,width=0.92\textwidth,max height=0.38\textheight]{Ferziger4th.pdf}};
\begin{scope}[x={(image.south east)},y={(image.north west)}]
\fill[white] (0.44791,0.55428) rectangle (0.46105,0.95812);
\fill[white] (0.46475,0.55428) rectangle (0.49122,0.95812);
\fill[white] (0.37185,0.30716) rectangle (0.40162,0.71065);
\fill[white] (0.40602,0.32007) rectangle (0.43420,0.72356);
\end{scope}
\end{tikzpicture}
\end{sourceblock}
其中 P 是升力的振荡周期，发现等于 5.059 s。该值与文献中的数据非常吻合。阻力以双频率振荡，因为每次涡旋脱落都有一个最大值和一个最小值，而当涡旋在一侧脱落时，升力最大，而当下一个涡旋在另一侧脱落时，升力最小。

\begin{center}\small 图 9.28 显示了使用最小时间步长在最精细网格上计算的瞬时速度矢量和压力等值线。可以看到一个大漩涡刚刚从下侧脱落；另一个涡流开始在圆柱体的上侧形成，当另一个涡流移开时，该涡流会增大。\end{center}
圆柱体前侧的等压线分布与 Re = 20 时的分布相似，参见。图9.24，只不过最小压力的位置现在已向上游移动并位于赤道之前，并且圆柱表面上侧和下侧的分布不对称。此外，前驻点（发现最大压力的位置）和最小压力位置之间有更多的轮廓，表明压力梯度更高。事实上，通过比较图 2 和 3 中的速度矢量图可以看出。如图 9.24 和 9.28 所示，Re = 200 时的流体加速度比 Re = 20 时强得多。而 Re = 20 时，场中的最大速度为 1.17 m/s（比自由流高 17\%）

\begin{sourceblock}
\includegraphics[page=350,trim={53.00bp 302.25bp 53.87bp 52.00bp},clip,width=\textwidth,max height=0.38\textheight]{Ferziger4th.pdf}
\par\vspace{0.45\baselineskip}
\small 图 9.28 Re = 200 时圆柱体周围非定常流动的预测瞬时速度矢量（上）和等压线（下）；插值速度矢量显示在统一的呈现网格上
\end{sourceblock}
速度），在 Re = 200 时，最大速度为 1.47 m/s（比自由流速度高 47\%）。

SIMPLE 算法用于求解纳维-斯托克斯方程。如前所述，SIMPLE 需要选择两个欠松弛因子。对于稳态流动，通常选择 0.8 的速度和 0.2 的压力，但当流动不稳定并且使用较小的时间步长时，两个欠松弛因子都可以增加。为了安全起见，我们在这里使用 0.8 表示速度，使用 0.5 表示所有网格和所有时间步上的压力，并强制代码每个时间步执行 10 次迭代。对于两个最小的时间步长，可以将速度的欠松弛因子增加到 0.9，而不是每个时间步长的固定数量的外部迭代，可以使用合适的标准在满足标准时停止外部迭代（例如，通过指定应达到的残差水平）。图 9.29 显示了连续性和动量方程中的残差范数如何在具有最小时间步长的最精细网格上的 4 个时间步长内随外部迭代而变化。注意事项

\begin{sourceblock}
\includegraphics[page=351,trim={53.00bp 449.25bp 53.87bp 52.00bp},clip,width=\textwidth,max height=0.38\textheight]{Ferziger4th.pdf}
\par\vspace{0.45\baselineskip}
\small 图 9.29 几个时间步外迭代过程中残差的减少
\end{sourceblock}
动量方程中的残差在 10 次迭代中下降了超过 3 个数量级，而连续性方程中的残差范数下降了大约 2 个数量级。原因是动量方程是非线性的，并且当前进到下一个时间水平时，比线性质量守恒方程对平衡产生更大的扰动。

计算非稳态流时的离散误差分析比稳态流的情况更困难。当边界条件造成不稳定时，如前一章研究的情况（参见第 8.4.2 节），情况会更简单，因为振荡周期是规定的。在本例中，边界条件是稳定的，流动不稳定仅由固有的不稳定性引起；此外，圆柱体表面是光滑的，并且流动分离点不是固定的，就像具有矩形横截面的圆柱体的情况一样。因此，时间步长和网格大小的变化会导致所有流特征的变化。

\begin{center}\small 图 9.30 显示了预测阻力和升力对网格细度的依赖性。显示了所有三个网格的结果；时间步长保持恒定为 0.01 秒（每个阻力振荡周期约 253 个时间步长）。从图中可以明显看出，粗网格和中网格解之间的差异远大于中网格和细网格解之间的差异。人们预计，通过二阶离散化，如果网格间距减半，连续网格上的解之间的差异将减少 4 倍（这里就是这种情况）。\end{center}
\begin{center}\small 图 9.31 显示了解对固定（精细）网格的时间步长大小的依赖性。以最大时间步长（0.04 s）运行模拟，直到达到周期性状态；然后将其保存并用作随后模拟 10 个升力振荡周期（所有四个时间步长）的起点。经过几个周期，解会调整到新的时间步长，并再次达到周期性状态。该图清楚地表明，获得了向时间步独立解的二阶收敛：对应于两个最小时间步的两条曲线无法彼此区分，其差异\end{center}
\begin{sourceblock}
\includegraphics[page=352,trim={53.00bp 331.26bp 53.87bp 52.00bp},clip,width=\textwidth,max height=0.38\textheight]{Ferziger4th.pdf}
\par\vspace{0.45\baselineskip}
\small 图 9.30 网格相关性分析：两个升力振荡周期内阻力（上）和升力（下）的变化，在三个网格上以相同的时间步长 t = 0 计算。 01s
\end{sourceblock}
到下一个较粗的网格增加了 4 倍，正如二阶时间离散方案所预期的那样。

非定常流模拟中的离散误差分析需要选择适当的初始网格和时间步长（基于一次改变一个），然后同时细化网格和时间步长。

\subsection*{9.12.3 Re = 200 时通道内圆柱体周围的流动}\phantomsection\addcontentsline{toc}{subsection}{9.12.3 Re = 200 时通道内圆柱体周围的流动}
我们还对安装在方形横截面管道的两壁之间的圆柱体周围的 3D 层流进行了计算。尽管仍然可以为该几何图形生成一个像样的块结构网格，但我们使用此示例来演示非结构化网格和商业软件在网格生成和流计算中的应用。创建三种网格类型并将其用于流动分析：修剪笛卡尔网格（如前面的示例所示）、四面体网格

\begin{sourceblock}
\includegraphics[page=353,trim={53.00bp 322.26bp 53.87bp 52.00bp},clip,width=\textwidth,max height=0.38\textheight]{Ferziger4th.pdf}
\par\vspace{0.45\baselineskip}
\small 图 9.31 时间步依赖性分析：两个升力振荡周期内阻力（上）和升力（下）的变化，在四个不同时间步的最精细网格上计算
\end{sourceblock}
德拉尔和多面体。在所有三种情况下，沿圆柱体和通道壁的棱镜层均以类似的方式创建。管道轴指向 x 方向，圆柱轴指向 y 方向（水平）。

\begin{center}\small 图 9.32 显示了解域的几何结构和边界上的粗多面体网格；图 9.33 显示了 y = 0 处穿过所有三种网格的纵向剖面。管道延伸长度为（以米为单位）： - 1.75≤x≤3.25，-0.5≤\end{center}
y≤0.5 和 − 0.5≤z≤0.5；圆柱体位于坐标系原点，直径为 0.4 m。因此，圆柱体阻挡了 40\% 的管道横截面。假设流体的密度为1kg/m 3，粘度为0.005Pa·s。在入口 ( x = − 1 . 75 m) 处，指定 u x = 1 m/s 的均匀速度场，而在出口 ( x = 3 . 25 m) 处，指定恒定压力。基于管道高度的雷诺数为 Re = 200。速度场初始化为 x 方向上等于入口速度的恒定速度。如前面的示例所示，无限环境中的圆柱体周围的流动在相同条件下将是不稳定的，但由于管道的限制，在这种情况下层流仍然是稳定的。

\begin{sourceblock}
\includegraphics[page=354,trim={53.00bp 494.25bp 53.87bp 52.00bp},clip,width=\textwidth,max height=0.38\textheight]{Ferziger4th.pdf}
\par\vspace{0.45\baselineskip}
\small 图 9.32 管道中的圆柱体：解域边界上的几何形状和粗多面体网格
\end{sourceblock}
我们在圆柱体周围的较短管道段中生成了三种类型的网格，并从上游和下游管道横截面沿管道轴执行所谓的网格挤压。这导致在入口下游和出口上游的管道部分中产生棱柱形单元，其中实现了单元在x方向上的逐渐膨胀或收缩，如图9.32中所示，其中创建了具有多边形基底的棱柱。在四面体网格的情况下，挤压部分由具有三角形基底的棱柱制成，而在修剪六面体网格的情况下，挤压区域包含细长的六面体。圆柱体周围以及某种程度上下游的局部网格细化是通过指定需要较小单元尺寸的两个体积形状来调用的：圆柱体周围较大的圆柱体形状和下游侧的块。网格生成软件在规则细化区域内生成规则形状的多面体（十二面体）和四面体，如图9.33所示。

使用西门子的STAR-CCM+软件进行计算；它基于所有积分的 FV 方法和中点规则近似。二阶迎风方案用于对流（使用上游单元中心的变量值和梯度线性外推到单元面中心），而线性形状函数用于梯度近似（导致笛卡尔网格上的中心差异）。使用 SIMPLE 算法，速度的欠松弛因子为 0.8，压力的欠松弛因子为 0.2。图 9.34 显示了中等大小的多面体网格（约 190 万个单元）上残差随外部迭代的变化。残差很快下降一个数量级，但迭代误差通常不会从一开始就跟随残差（参见上一章的图 8.9）。从残差图中估计迭代误差的一种安全方法是从最终状态向后推断。例如，我们在图 9.34 中看到，在迭代 800 处，u x 速度的残差范数约为 1e-6。沿着平均斜率向后投影，我们在 0.01 左右的水平上达到迭代 0，从而表明迭代误差的真正减少约为 4 个数量级。事实上，仔细观察速度和压力的监测值会发现前 4 个有效数字没有变化，这是确保解确实收敛的另一个有用的检查。

\begin{sourceblock}
\includegraphics[page=355,trim={53.00bp 131.48bp 53.87bp 52.00bp},clip,width=\textwidth,max height=0.38\textheight]{Ferziger4th.pdf}
\par\vspace{0.45\baselineskip}
\small 图9.33 纵向对称平面y=0的计算网格：四面体网格（上）、多面体网格（中）和修剪六面体网格（下）
\end{sourceblock}
\begin{sourceblock}
\includegraphics[page=356,trim={53.00bp 398.25bp 53.87bp 52.00bp},clip,width=\textwidth,max height=0.38\textheight]{Ferziger4th.pdf}
\par\vspace{0.45\baselineskip}
\small 图 9.34 管道中圆柱体周围 3D 层流的 SIMPLE 算法中残差随外部迭代次数增加的变化
\end{sourceblock}
\begin{center}\small 图 9.35 显示了两个纵向对称平面中的速度矢量，使用修剪的六面体网格计算得出。在水平面（z = 0）中可以看到速度矢量如何在圆柱体两端向后转向；这表明马蹄涡的形成。垂直截面显示，当流动绕过圆柱体时，流动有很强的加速度；与圆柱体上游的速度矢​​量相比，圆柱体周围的速度矢量长度加倍。在气缸后面形成一个再循环区，其长度几乎是两倍直径。由于流动是稳定的，因此由于几何对称性，速度场在两个截面中是对称的。 1 然而，请注意，气缸后面的再循环区的长度在横向上有所不同：中心最长，随着向侧壁移动而变小，但随后在靠近墙壁时再次增加。因此，三维效果并不局限于与侧壁相交的气缸端部，侧壁的效果在整个气缸跨度上都是可见的。\end{center}
\begin{center}\small 图 9.36 显示了垂直对称平面中的压力分布。圆柱体周围的等压线类似于无限环境中圆柱体的二维计算中的等压线：最高压力出现在前驻点，即流动撞击圆柱体壁的位置，最低压力位于圆柱体侧面。下游侧压力有所恢复，但由于粘性损失，下游驻点处的压力远低于上游驻点。气缸下游的一个闭合压力轮廓指示结束\end{center}
\footnotetext[1]{当几何形状对称时，稳态流动不必对称；在一些对称的几何形状中，例如扩散器和突然的通道扩张，在实验和模拟中都可以获得不对称的稳态流动。}
\begin{sourceblock}
\includegraphics[page=357,trim={53.00bp 314.25bp 53.87bp 52.00bp},clip,width=\textwidth,max height=0.38\textheight]{Ferziger4th.pdf}
\par\vspace{0.45\baselineskip}
\small 图 9.35 使用修剪六面体网格计算的两个纵向对称平面中的速度矢量：y = 0（上）和 z = 0（下）
\end{sourceblock}
\begin{sourceblock}
\includegraphics[page=357,trim={53.00bp 107.09bp 53.87bp 375.40bp},clip,width=\textwidth,max height=0.38\textheight]{Ferziger4th.pdf}
\par\vspace{0.45\baselineskip}
\small 图 9.36 纵向对称平面 y = 0 中的压力分布，使用修剪的六面体网格计算（从左到右流动）
\end{sourceblock}
\begin{sourceblock}
\includegraphics[page=358,trim={53.00bp 258.25bp 53.87bp 52.00bp},clip,width=\textwidth,max height=0.38\textheight]{Ferziger4th.pdf}
\par\vspace{0.45\baselineskip}
\small 图 9.37 圆柱体下游的 u x 速度剖面，在三个系统细化的修剪六面体网格上计算： x = 0.35，y = 0（上）且 x = 0.75，z = 0（下）
\end{sourceblock}
再循环区：此处两股气流（来自气缸上方和下方）相遇，导致局部压力升高。

对于每种网格类型，我们对三个系统细化的网格进行了一系列计算，以检查解的网格依赖性。图 9.37 显示了 u x 速度沿圆柱下游 0.875 D 处垂直对称面中的一条线和沿圆柱下游 1.875 D 处水平对称面中的一条线的剖面。给出了三个系统地细化的修剪六面体网格的结果，所有壁均具有五个棱柱层。粗网格有 377,006 个 CV，中网格有 1,611,904 个 CV，细网格有 8,952,321 个 CV（每次细化网格间距减半）。对于沿垂直线的剖面，三个网格上得到的解之间的差异非常小；高峰

\begin{sourceblock}
\includegraphics[page=359,trim={53.00bp 259.26bp 53.87bp 52.00bp},clip,width=\textwidth,max height=0.38\textheight]{Ferziger4th.pdf}
\par\vspace{0.45\baselineskip}
\small 图 9.38 圆柱体下游 u x 速度的剖面图，根据三种不同的网格类型计算： x = 0.35，y = 0（上）且 x = 0.75，z = 0（下）
\end{sourceblock}
粗网格对值的预测不足 ca。 2\%，而中、细网格的轮廓在图中几乎无法区分。在距离气缸更下游的水平部分中，差异更加明显：此处的峰值比管道中的平均速度小一个数量级，因此甚至更小的差异也可以清楚地区分。正如二阶方法所预期的那样，粗网格和中网格之间的差异大约是中网格和细网格之间的差异的 4 倍。因此可以估计，最细网格上的平均离散化误差约为平均管道速度的 0.1\%。

\begin{center}\small 图 9.38 显示了在不同类型的精细网格上计算的相同速度剖面的比较。其差异小于同类型中、细网格之间的差异，见图9.37。这证实了，当网格足够细时，无论使用哪种类型的计算网格，都将获得相同的与网格无关的解。不同之处在于生成网格和求解纳维-斯托克斯方程所需的工作量。\end{center}
当使用诸如此处使用的商业软件时，通常，当使用修剪的六面体网格时，达到相同水平的离散化误差所需的工作量是最低的。效率最低的是四面体网格：与使用多面体或六面体网格相比，需要更多的单元才能达到相同水平的离散化误差。在具有相同控制体积数量的网格上，相同条件下（求解线性方程组、欠松弛因子等时每次外迭代的内迭代次数）迭代的收敛速度也较慢。

请注意，当对所有网格应用相同类型的近似时，上述陈述是有效的（此处：积分近似的中点规则、线性插值或外推、梯度近似的线性形状函数）。如果应用专门针对一种特定网格类型调整的离散化，则比率可能会有所不同。

我们在这里不提供计算时间的定量比较，因为我们没有改变网格来确保相同水平的离散化误差，并且因为比率取决于问题；上述陈述基于 CFD 的许多工业应用的经验。根据离散化过程中使用的有限近似类型以及所使用的线性方程求解器，比较也会有所不同。

三维流比二维流更难以可视化。速度矢量和流线通常在 2D 中使用，但在 3D 问题中很难绘制和解释。在选定表面（平面、一定数量的等值面、边界面等）上显示轮廓和矢量投影以及从不同方向查看它们的可能性可能是分析 3D 流的最佳方法。非稳态流动需要结果的动画。我们不会在这里进一步讨论这些问题，但想强调它们的重要性。

\chapter{湍流}
\section*{10.1 简介}\phantomsection\addcontentsline{toc}{section}{10.1 简介}
工程实践中遇到的大多数流动都是湍流（Pope 2000 和 Jovanovi c 2004），因此与迄今为止研究的层流相比需要不同的处理。湍流具有以下特性：

\begin{itemize}
\item 湍流非常不稳定。对于不熟悉这些流动的观察者来说，流动中大多数点的速度随时间变化的图看起来是随机的。可以使用“混乱”这个词，但近年来它被赋予了另一个定义。
\item 它们是三维的。时间平均速度可能只是两个坐标的函数，但瞬时场在所有三个空间维度上快速波动。
\item 它们包含大量涡度。事实上，涡旋拉伸是增加湍流强度的主要机制之一。
\item 湍流增加了保守量的搅拌速率。搅拌是使具有至少一种守恒特性的不同浓度的流体块接触的过程。实际的混合是通过扩散完成的。尽管如此，这个过程通常被称为湍流扩散。
\item 通过刚才提到的过程，湍流使不同动量含量的流体接触。由于粘度的作用，速度梯度减小，流动的动能降低；换句话说，混合是一个耗散过程。损失的动能不可逆地转化为流体的内能。
\item 近年来的研究表明，湍流包含相干结构——可重复且本质上是确定性的事件，它们造成了大部分的混合。然而，湍流的随机成分导致这些事件在大小、强度和发生时间间隔上彼此不同，使得对它们的研究变得非常困难。
\item 湍流在很宽的长度和时间尺度上波动。这一特性使得湍流的直接数值模拟变得非常困难。 （见下文。）
\end{itemize}
所有这些属性都很重要。湍流产生的效果可能是理想的，也可能不是理想的，具体取决于应用。当需要化学混合或传热时，强烈混合是有用的；这两者都可能因湍流而增加几个数量级。另一方面，动量混合的增加会导致摩擦力增加，从而增加泵送流体或推进车辆所需的功率；同样，增加一个数量级并不罕见。工程师需要能够理解和预测这些影响才能实现良好的设计。在某些情况下，可以至少部分地控制湍流。

过去，研究湍流的主要方法是实验。诸如时间平均阻力或传热之类的整体参数相对容易测量，但随着工程设备复杂程度的提高，所需的细节和精度水平也随之增加，测量的成本、费用和难度也随之增加。为了优化设计，通常有必要了解不良影响的根源；这需要进行详细的测量，既昂贵又耗时。目前，某些类型的测量几乎不可能进行，例如流量内的波动压力。其他的则无法达到所需的精度。因此，数值方法可以发挥重要作用，而这种作用在商用和军用飞机和船舶的设计和优化中可能最为先进。通常对部件（例如机翼、螺旋桨、涡轮机等）以及全身配置进行 CFD 分析。然而，数值方法的要求取决于人们想要在流程中分析的内容，我们将在第 1 节中回到这一原则。 12.1.1.例如，在时间平均流动中，力、热通量等的确定所提出的要求不如要查看雷诺应力或什至更感兴趣的三重相关性（无论是在准确性还是在模拟持续时间方面）的要求严格。

在继续讨论这些流动的数值方法之前，总结一下预测湍流的方法是有帮助的。巴迪纳等人。 （1980）列出了六个类别，我们的总结如下：

\begin{itemize}
\item 第一种方法涉及使用相关性，例如给出作为雷诺数函数的摩擦因数或作为雷诺数和普朗特数函数的传热努塞尔数。这种方法通常在入门课程中教授，非常有用，但仅限于简单类型的流，即只能通过几个参数来表征的流。由于其使用不需要CFD，这里不再赘述。
\item 第二种方法使用积分方程，可以通过对一个或多个坐标进行积分，从运动方程导出积分方程。通常这会将问题简化为一个或多个易于求解的常微分方程。应用于方程的方法是第 1 章中讨论的常微分方程的方法。 6.
\item 第三种方法基于将运动方程分解为平均分量和波动分量所获得的方程（Pope 2000）。不幸的是，这些分解方程不形成闭集（参见第 10.3.5.1 节），因此这些方法需要引入近似值（湍流模型）。本章稍后将介绍当今常用的一些湍流模型以及与包含湍流模型的方程的数值解相关的问题的讨论。在那里我们将重点关注所谓的单点闭包。 1
\end{itemize}
处理湍流模型的实际方法取决于用于获取平均值和波动方程的过程的性质，导致第三种方法的子类别如下：

– 我们获得一组称为雷诺平均的偏微分方程
纳维-斯托克斯（或 RANS）方程，如果创建平均值的过程是对随时间或实现集合（所有可控因素均保持固定的想象的一组流动）的运动方程求平均值。由此产生的方程可以表示随时间变化的流或稳态流，如下文所述。 – 我们获得一组称为大涡模拟 (LES) 方程的方程
当通过对空间中的有限体积进行平均（或过滤）来获得平均值时。 2 然后，LES 求解出流的最大尺度运动的准确表示，同时对小尺度运动进行近似或建模。它可以被视为RANS（见上文）和直接数值模拟（见下文）之间的一种折衷。

\begin{itemize}
\item 最后，第四种方法是直接数值模拟（DNS），其中针对湍流中的所有运动求解纳维-斯托克斯方程。
\end{itemize}
随着这个列表的进展，越来越多的湍流运动被计算出来，而越来越少的湍流运动被模型近似。这使得接近底部的方法更加精确，但计算时间显著增加。

本章描述的所有方法都需要求解某种形式的质量、动量、能量或化学物质的守恒方程。主要困难在于湍流比层流包含更广泛的长度和时间尺度变化。因此，尽管它们与层流方程相似，但描述湍流的方程通常求解起来更加困难和昂贵。

1 还有两点闭包，它们使用方程来计算两个空间点处速度分量的相关性，或者更常见的是，使用这些方程的傅立叶变换。这些方法在实践中并未广泛使用（Leschziner 2010），并且最常用于纯粹的研究，因此我们不会进一步考虑它们。然而，Lesieur（2010，2011）提出了一篇迷人的评论，对他们的历史和最先进的技术有很好的洞察力。 2 通过对空间中“相对”大的体积进行平均，可以获得超大涡流模拟 (VLES)。我们稍后将讨论（第 10.3.1 节；10.3.7 节）如何定义“相对”。

\section*{10.2 直接数值模拟（DNS）}\phantomsection\addcontentsline{toc}{section}{10.2 直接数值模拟（DNS）}
\subsection*{10.2.1 概述}\phantomsection\addcontentsline{toc}{subsection}{10.2.1 概述}
湍流模拟最准确的方法是求解纳维-斯托克斯方程，除了可以估计和控制误差的数值离散化之外，无需平均或近似。从概念的角度来看，这也是最简单的方法。在这样的模拟中，流中包含的所有运动都被解析。计算得到的流场相当于一次流的实现或一次短期的实验室实验；如上所述，这种方法称为直接数值模拟 (DNS)。

在直接数值模拟中，为了确保捕获湍流的所有重要结构，执行计算的域必须至少与要考虑的物理域或最大湍流涡旋一样大。后一个尺度的一个有用的度量是湍流的积分尺度 (L)，它本质上是速度的脉动分量保持相关的距离。因此，域的每个线性维度必须至少是积分尺度的几倍。有效的模拟还必须捕获所有动能耗散。这种情况发生在最小的尺度上，即粘性活跃的尺度上，因此网格的大小必须在粘性确定的尺度上，称为柯尔莫戈洛夫尺度， η。通常分辨率要求表述为

\begin{sourceblock}
\begin{tikzpicture}
\node[inner sep=0,anchor=south west] (image) at (0,0) {\includegraphics[page=364,trim={53.00bp 312.77bp 53.87bp 325.35bp},clip,width=0.92\textwidth,max height=0.38\textheight]{Ferziger4th.pdf}};
\begin{scope}[x={(image.south east)},y={(image.north west)}]
\fill[white] (0.00000,0.80014) rectangle (0.14890,1.00000);
\fill[white] (0.15161,0.80014) rectangle (0.17642,1.00000);
\fill[white] (0.17907,0.80014) rectangle (0.25377,1.00000);
\fill[white] (0.25648,0.80014) rectangle (0.28626,1.00000);
\end{scope}
\end{tikzpicture}
\end{sourceblock}
所以网格大小 ≤ 2 η。舒马赫等人。 ( 2005 ) 指出耗散发生在一定的尺度范围内（耗散谱的峰值在 k η ∼ 0 . 2 处），并且局部尺度小于 η 出现在流动中；他们的模拟针对 k max η 高达 34 进行。为了捕捉其流程的本质，Kaneda 和 Ishihara (2006) 以及 Bermejo-Moreno 等人。 (2009) 使用的 k max η 值分别高达 2 和 4。

对于均匀各向同性湍流（最简单的湍流类型），没有理由使用除均匀网格之外的任何东西。在这种情况下，刚才给出的参数表明每个方向上的网格点数必须是 L / η 的量级；可以证明（Tennekes 和 Lumley 1976）该比率与 Re 3 / 4 成正比
L。 3 这里 Re L 是基于速度波动大小和积分尺度的雷诺数；该参数通常约为工程师用来描述流动的宏观雷诺数的 0.01 倍。因为必须在三个坐标方向中的每一个上使用这个点数，并且时间步长与网格大小相关，所以模拟的成本为 Re 3
L。就
3 情况可能会更糟！ NSF 基于仿真的工程科学报告（2006 年）实际上指出，尺度的暴政主导着许多领域的仿真工作，包括流体力学。因此，即使对于 Re L ∼ 10 7，长度比例也约为 2 × 10 5。然而，他们指出，对于蛋白质折叠，时间尺度比为 ∼ 10 12，而对于先进材料设计，空间尺度比为 ∼ 10 10！

工程师用来描述流量的雷诺数，成本的比例可能会有所不同。

由于计算中可使用的网格点数量受到执行计算的机器的处理速度和内存的限制，因此直接数值模拟通常在几何简单的域中进行。在现有的机器上，可以使用多达 4096 3 个网格点对湍流雷诺数高达 10 5 量级的均匀流进行直接数值模拟（Ishihara 等人，2009）。如前一段所述，这对应于整体流量雷诺数大约大两个数量级，并允许 DNS 达到工程感兴趣的雷诺数范围的低端，使其在某些情况下成为有用的方法。在其他情况下，可以通过使用某种外推法从模拟的雷诺数外推到实际感兴趣的雷诺数。有关 DNS 的更多详细信息，请参阅 Pope (2000) 以及 Moin 和 Mahesh (1998) 的概述以及 Ishihara 等人的具体论文。 (2009)，他们使用傅里叶谱方法（第 3.11 节）来研究各向同性流，Wu 和 Moin (2009)，他们使用分数步（第 7.2.1 节）有限差分法来研究边界层流。

\subsection*{10.2.2 讨论}\phantomsection\addcontentsline{toc}{subsection}{10.2.2 讨论}
DNS 的结果包含有关流的非常详细的信息。这可能非常有用，但一方面，它提供的信息远远超出任何工程师的需求，另一方面，DNS 价格昂贵且不经常用作设计工具。那么人们必须问一下 DNS 可以用来做什么。有了它，我们可以获得大量网格点上的速度、压力和任何其他感兴趣的变量的详细信息。这些结果相当于实验数据，可用于生成统计信息或创建“数值流可视化”。从后者中，人们可以学到很多东西，例如，流中存在的连贯结构（图 10.1）。这些丰富的信息可用于更深入地了解流动的物理原理，或构建定量模型（可能是 RANS 或 LES 类型），这将允许以较低的成本计算其他类似的流动，从而使此类模型可用作工程设计工具。

DNS 的一些用途示例如下：

\begin{itemize}
\item 了解层流到湍流转变的过程，以及湍流产生、能量传递和耗散的机制；
\item 模拟空气动力噪声的产生；
\item 了解压缩性对湍流的影响；
\item 了解燃烧和湍流之间的相互作用；
\item 控制和减少固体表面上的阻力。
\end{itemize}
DNS 的其他应用已经完成，并且将来无疑还将进行许多其他应用。

\begin{sourceblock}
\includegraphics[page=366,trim={53.00bp 508.25bp 53.87bp 52.00bp},clip,width=\textwidth,max height=0.38\textheight]{Ferziger4th.pdf}
\par\vspace{0.45\baselineskip}
\small 图 10.1 平板边界层完全湍流区域中的发夹森林。所示速度梯度张量第二个不变量的等值面；根据距墙的无量纲距离的局部值对表面进行着色。大于 300 的值呈红色（从左到右流动）。来自：吴和莫因 2011
\end{sourceblock}
计算机速度的提高使得在工作站上以低雷诺数执行简单流的 DNS 成为可能。简单流是指任何均匀的湍流（有很多）、通道流、自由剪切流和其他一些流。在大型并行计算机上，如上所述，DNS 是通过 4096 3 完成的
( ∼ 6 . 9 × 10 10 ) 或更多网格点。计算时间取决于机器和所使用的网格点的数量，因此无法给出有用的估计。事实上，人们通常选择要模拟的流和网格点的数量以适应可用的计算机资源。完整的最先进的模拟可能会在计算机系统上使用数十万个核心并消耗数百万个核心小时。随着计算机变得越来越快和存储器越来越大，将模拟更复杂和更高的雷诺数流。

直接数值模拟和大涡模拟可以采用多种数值方法。几乎本书中描述的任何方法都可以使用；然而，对于并行系统上真正的大规模计算，最常使用基于 CDS 或谱方案的显式代码；我们在下面注意到，在某些情况下，隐式方法用于方程中的某些项。由于这些方法已经在前面的章节中介绍过，因此我们不会在这里给出太多细节。然而，DNS 和 LES 以及稳态流模拟之间存在重要差异，因此讨论这些差异很重要。

DNS 和 LES 数值方法最重要的要求源于需要准确实现包含各种长度和时间尺度的流。由于需要准确的时间历史，因此为稳态流设计的技术效率低下，并且在未经大量修改的情况下不应使用。对精度的要求要求时间步长很小，显然，时间提前方法对于所选的时间步长必须是稳定的。在大多数情况下，可以使用在精度要求所需的时间步长内稳定的显式方法，因此没有理由承担与隐式方法相关的额外费用；因此，大多数模拟都使用显式的时间提前方法。一个值得注意的（但不是唯一的）例外发生在固体表面附近。这些区域中的重要结构尺寸非常小，必须使用非常精细的网格，特别是在垂直于墙壁的方向上。数值不稳定可能是由涉及垂直于壁的导数的粘性项引起的，因此通常会隐式处理这些项。在复杂的几何形状中，可能需要隐式地处理更多术语。

DNS和LES中最常用的时间提前方法具有二阶到四阶精度；最常用的是龙格-库塔方法，但也使用了其他方法，例如 Adams-Bashforth 和蛙跳方法。一般来说，对于给定的精度等级，龙格-库塔方法每个时间步需要更多的计算。尽管如此，它们仍然是首选，因为对于给定的时间步长，它们产生的误差比竞争方法小得多。因此，在实践中，它们在相同的精度下允许更大的时间步长，这足以补偿增加的计算量。克兰克-尼科尔森方法通常应用于必须隐式处理的项，例如粘性项或壁法向对流项。

时间提前方法的一个困难在于，精度高于一阶的方法需要在多个时间步长（包括中间时间步长）存储数据。因此，设计和使用需要相对较少存储的方法是有优势的。 Leonard 和 Wray (1982) 提出了一种三阶 Runge-Kutta 方法，与该精度的标准 Runge-Kutta 方法相比，该方法需要更少的存储空间（参见 Bhaskaran 和 Lele 2010 中该方法的应用）。另一方面，对于可压缩流和计算声学，需要不同的特性，例如低耗散和低色散龙格-库塔方案（Hu et al. 1996，应用于 Bhagatwala 和 Lele 2011）。

DNS 中的另一个重要问题是需要处理各种长度范围。这需要改变人们对离散化方法的思考方式。空间离散化方法精度的最常见描述是其阶数，该数字描述了当网格尺寸减小时离散化误差减小的速率。回到第一章的讨论。如图3所示，根据速度场的傅里叶分解来思考是有用的。我们证明（第 3.11 节），在均匀网格上，速度场可以用傅立叶级数表示：

\begin{sourceblock}
\begin{tikzpicture}
\node[inner sep=0,anchor=south west] (image) at (0,0) {\includegraphics[page=367,trim={53.00bp 248.32bp 53.87bp 399.38bp},clip,width=0.92\textwidth,max height=0.38\textheight]{Ferziger4th.pdf}};
\begin{scope}[x={(image.south east)},y={(image.north west)}]
\fill[white] (0.00000,0.71312) rectangle (0.01744,1.00000);
\fill[white] (0.02012,0.71312) rectangle (0.11266,1.00000);
\fill[white] (0.11540,0.71312) rectangle (0.19678,1.00000);
\fill[white] (0.36701,0.07050) rectangle (0.42752,0.71204);
\fill[white] (0.43107,0.08460) rectangle (0.45925,0.71204);
\end{scope}
\end{tikzpicture}
\end{sourceblock}
在大小为 x 的网格上可以解析的最高波数 k 是 π/ x，因此我们仅考虑 0 < k < π/ x。系列（10.2）可以逐项区分。 e i kx 的精确导数，ike e i kx 被 i k eff e i kx 替代，其中 k eff 是第 2 节中定义的有效波数。 3.11 当使用有限差分近似时。图 3.12 中给出的 k eff 绘图表明，中心差异仅在 k < π/ 2 x （感兴趣的波数范围的前半部分）时准确。

稳态流模拟中没有遇到的湍流模拟困难在于，湍流谱（湍流能量在波数或反长度尺度上的分布）通常在波数范围 \{ 0 , π/ x \} 的很大一部分上很大，因此图 3.12 中的方法的阶数及其行为可能变得很重要。在该图中，我们看到理想情况下数值方案的有效波数接近精确的谱线。四阶 CDS 无疑是一种改进，但无论如何都不是理想的。紧凑方案（Lele 1992 和 Mahesh 1998；参见第 3.3.3 节）可以产生类似谱的高阶方案。在可压缩流模拟和计算气动声学中，此类方案很流行，例如，参见 Kim (2007) 的四阶紧凑方案或 Kawai 和 Lele (2010) 的六阶方案的实现。然而，由于网格分辨率对于定义 DNS 中可解析的最小结构至关重要，因此如果网格足够细，使用二阶 CDS 也可以获得令人印象深刻的结果（Wu 和 Moin 2009）。

重申使用节能空间差分方案的重要性也很有用。许多方法，包括所有迎风方法，都是耗散的。也就是说，它们包括作为截断误差一部分的扩散项，该扩散项在与时间相关的计算中耗散能量。它们的使用受到提倡，因为它们引入的耗散通常可以稳定数值方法。当这些方法应用于稳态问题时，稳态结果中的耗散误差可能不会太大（尽管我们在前面的章节中表明这些误差可能相当大）。当这些方法用于 DNS 时，产生的耗散通常远大于物理粘性造成的耗散，并且获得的结果可能与问题的物理性质几乎没有联系。有关能量守恒的讨论，请参阅第 2 节。 7.1.3.另外，如第 1 章所示。如图7所示，能量守恒防止速度无限制地增长，从而保持稳定。

时间和空间上的方法和步长需要相关联。空间和时间离散化中产生的误差应尽可能接近相等，即它们应该是平衡的。对于每个时间步长来说，这是不可能的，但是，如果在平均意义上不满足该条件，则说明在自变量之一中使用了太细的步长，并且可以以较低的成本进行模拟，而精度损失很小。

DNS 和 LES 的准确性很难衡量。原因在于湍流的本质。湍流初始状态的微小变化会随时间呈指数放大，并且在相对较短的时间之后，扰动流几乎与原始流相似。这是一种物理现象，与数值方法无关。由于任何数值方法都会引入一些误差，并且方法或参数的任何变化都会改变该误差，因此不可能直接比较两个解来确定误差。相反，可以使用不同的网格（应该与原始网格有很大不同）重复模拟，并且可以比较两种解的统计特性。根据差异，可以找到误差的估计值。不幸的是，很难知道误差如何随网格大小变化，因此这种估计只能是近似值。大多数计算简单湍流的人都使用一种更简单的方法，即查看湍流的频谱。如果最小尺度的能量比能谱峰值的能量小得多，则可以安全地假设流动已经得到很好的解析。

在域配置和边界条件允许的情况下，精度要求利用 DNS 和 LES 中常见的频谱方法。这些方法已在前面的第 2 节中进行了简要描述。 3.11.本质上，他们使用傅里叶级数作为计算导数的手段。傅里叶变换的使用之所以可行，只是因为快速傅里叶变换算法（Cooley 和 Tukey 1965；Brigham
1988）将计算傅立叶变换的成本降低为 n log 2 n 运算。不幸的是，该算法仅适用于等距网格和其他一些特殊情况。已经开发出许多此类专门方法来求解纳维-斯托克斯方程；光谱方法的更多细节在第 3 节中给出。 3.11 和卡努托等人。 （2007）。

谱方法的一个有趣应用不是直接逼近纳维-斯托克斯方程，而是将它们乘以一系列“测试或基函数”，在整个域上积分，然后找到满足所得方程的解（参见第 3.11.2.1 节）。满足这种形式的方程的函数被称为“弱解”。人们可以将纳维-斯托克斯方程的解表示为一系列向量函数，每个向量函数的散度为零。这种选择消除了方程积分形式的压力，从而减少了需要计算和存储的因变量的数量。通过注意到如果函数的散度为零，则可以根据其他两个分量计算其第三个分量，可以进一步减少因变量集。结果是只需要计算两组因变量，从而将内存需求减少一半。由于这些方法比较专业，开发需要相当大的篇幅，这里不再详细介绍；有兴趣的读者可以参考 Moser 等人的论文。 ( 1983 ) 或 Canuto 等人的 3.4.2 节。 （2007）。

\subsection*{10.2.3 初始条件和边界条件}\phantomsection\addcontentsline{toc}{subsection}{10.2.3 初始条件和边界条件}
DNS 中的另一个困难是生成初始条件和边界条件。前者必须包含初始三维速度场的所有细节。由于相干结构是流动的重要组成部分，因此很难构建这样的场。此外，初始条件的影响通常会被流动“记住”相当长的时间，通常是几个“涡流转换时间”。涡流转换时间本质上是流动的积分时间尺度或积分长度尺度除以均方根速度（q）。因此，初始条件对结果有显著影响。通常，以人工构造的初始条件开始的模拟的第一部分必须被丢弃，因为它不忠实于物理原理。如何选择初始条件的问题既是一门艺术，也是一门科学，无法给出适用于所有流动的唯一规定，但我们将给出一些例子。

对于均匀各向同性湍流，最简单的情况是使用周期性边界条件，并且最容易在傅里叶空间中构造初始条件，即我们需要创建 ˆ u i ( k )。这是通过给出设置傅立叶模式幅度的频谱来完成的，即 |ˆ u i ( k ) |。连续性 k · ˆ u i ( k ) 的要求对该模式提出了另一个限制。这样就只剩下一个随机数来完全定义 ˆ u i ( k )；它通常是相位角。然后，模拟必须运行大约两个涡流翻转时间，然后才能被认为代表真实的湍流。

其他流动的最佳初始条件是从先前的模拟结果中获得的。例如，对于受到应变的均匀湍流，最佳初始条件取自已发展的各向同性湍流。对于通道流，最佳选择是平均速度、不稳定模式（具有几乎正确的结构）和噪声的混合。对于弯曲通道，可以将平面通道流动完全展开的结果作为初始条件。

类似的考虑适用于流进入域的边界条件（流入条件）。正确的条件必须包含每个时间步的湍流平面（或其他表面）上的完整速度场，这是很难构建的。例如，对于弯曲通道中的流动进行此操作的一种方法是使用平面通道中的流动结果。对平面通道流进行模拟（同时或提前），并且垂直于主流方向的一个平面上的速度分量提供了弯曲通道的流入条件。 Chow 和 Street (2009) 对苏格兰一座山上的气流的流入条件使用了这样的策略，并使用了用于预测大气中尺度气流的代码。 4
正如已经指出的，对于在给定方向上不变化（在统计意义上）的流，可以在该方向上使用周期性边界条件。它们易于使用，特别适合光谱方法，并在标称边界处提供尽可能真实的条件。

流出边界不太难处理。一种可能性是使用外推条件，要求所有量在垂直于边界的方向上的导数为零：

\begin{sourceblock}
\begin{tikzpicture}
\node[inner sep=0,anchor=south west] (image) at (0,0) {\includegraphics[page=370,trim={53.00bp 308.22bp 53.87bp 329.64bp},clip,width=0.92\textwidth,max height=0.38\textheight]{Ferziger4th.pdf}};
\begin{scope}[x={(image.south east)},y={(image.north west)}]
\fill[white] (0.00000,0.88331) rectangle (0.02746,1.00000);
\fill[white] (0.03011,0.88331) rectangle (0.07152,1.00000);
\fill[white] (0.07420,0.88331) rectangle (0.19215,1.00000);
\fill[white] (0.19489,0.88331) rectangle (0.22797,1.00000);
\fill[white] (0.23068,0.88331) rectangle (0.29537,1.00000);
\end{scope}
\end{tikzpicture}
\end{sourceblock}
其中 φ 是任何因变量。该条件常用于稳态流，但在非稳态流中并不令人满意。对于后者，最好用非定常对流条件来代替该条件。已经尝试了许多这样的条件，但似乎效果很好的一个也是最简单的条件之一：

\begin{sourceblock}
\includegraphics[page=370,trim={53.00bp 216.80bp 53.87bp 421.05bp},clip,width=0.92\textwidth,max height=0.38\textheight]{Ferziger4th.pdf}
\end{sourceblock}
其中 U 是与流出表面上的位置无关的速度，其选择是为了保持总体守恒，即，它是使流出质量通量等于流入质量通量所需的速度。此条件似乎避免了由压力扰动从流出边界反射回域内部而引起的问题。另一种选择是使用第 2 节中描述的强制技术。 13.6 抑制在朝向出口边界一定距离内垂直于平均流动方向的速度波动。通过这种方式，所有涡流在到达边界和反射之前就消失了

\footnotetext[4]{“中尺度”是指水平尺寸通常在 5 公里左右到数百公里甚至 10 3 公里之间的天气系统。}
被避免。然而，仍然存在确定最佳强迫参数的问题；这个问题将在第二节中更详细地讨论。 13.6.

在光滑的实体壁上，无滑移边界条件，这已在第 2 章中进行了描述。可以使用图8和图9的实施例。必须记住，在这种类型的边界处，湍流往往会形成小但非常重要的结构（“条纹”），这些结构需要非常精细的网格，特别是在垂直于壁的方向上，以及在较小程度上，在展向方向上（垂直于壁和主要流动方向的方向）。另一种方法是对复杂形状（Kang et al. 2009；Fadlun et al. 2000；Kim et al. 2001；Ye et al. 1999）或粗糙墙壁（Leonardi et al. 2003；Orlandi and Leonardi 2008）采用浸入边界法（IBM）。在该方法中，控制方程在规则网格上离散化和求解，但施加了边界条件，例如，通过实际边界上的体力，这导致在规则网格的相关点处施加适当的力，即使它们是近边界，但不是边界。此外，可以使用从实际边界上的所需条件到规则网格边界的速度插值，并且必须保持质量守恒。大多数这些 IBM 方法都采用分数步法（参见第 7.2.1 和 8.3 节），Kim 和 Ye 的工作使用 FV，而其他方法则使用 FD 方案。

对称边界条件经常用于 RANS 计算中以减小域的大小，但通常不适用于 DNS 或 LES，因为尽管平均流可能关于某个特定平面对称，但瞬时流却不是，并且通过应用此类条件可以消除重要的物理效应。然而，对称条件已被用来表示自由表面。

尽管尽一切努力使初始条件和边界条件尽可能真实，但在流体发展出物理流的所有正确特征之前，必须运行模拟一段时间。这种情况源于湍流的物理原理，因此人们几乎无法加速这一过程。下面提到了一种可能性。正如我们所指出的，涡流翻转时间尺度是问题的关键时间尺度。在许多流中，它可以与作为整体的流的时间尺度特征（即，平均流时间尺度）相关。然而，在分离的流中，有些区域在很长的时间范围内与流的其余部分进行通信，并且开发过程可能非常慢，从而需要很长的运行时间。

确定流量发展是否完成的最佳方法是监测一些量，最好是对流量发展缓慢的部分敏感的量；选择取决于所模拟的流量。作为示例，可以测量分离流的再循环区域中的表面摩擦力随时间变化的空间平均值。最初，监测数量通常会系统性增加或减少；当流动充分发展时，该值将随时间呈现统计波动。在该点之后，可以通过对时间和/或流中的统计均匀坐标进行平均来获得统计平均结果（例如，对于平均速度或其波动）。这样做时，重要的是要记住，由于湍流不是纯粹随机的，因此样本大小与平均过程中使用的点数不同。保守估计是假设直径等于积分尺度的每个体积（并且等于积分时间尺度的每个时间段）仅代表单个样本。

最初使用粗网格可以加快开发过程。当流在该网格上发展时，可以引入精细网格。如果这样做，仍然需要一些等待才能在精细网格上发展流动，但它可能比在整个模拟过程中使用精细网格所需的时间要短。

Wu 和 Moin（2009）关于零压力梯度边界层模拟的论文是关于为对它们非常敏感的问题设置初始条件和边界条件以及流动物理如何成为如何进行的良好指南的杰作。

\subsection*{10.2.4 DNS应用举例}\phantomsection\addcontentsline{toc}{subsection}{10.2.4 DNS应用举例}
作为 DNS 可以实现的功能的说明性示例，我们将采用一个看似简单的流，即由大量静态流体中的振荡网格创建的流。网格的振荡会产生湍流，湍流的强度随着距网格的距离而减弱。这种能量从振荡栅格转移的过程通常称为湍流扩散；湍流的能量传递在许多流动中起着重要作用，因此其预测很重要，但建模却出人意料地困难。布里格斯等人。 (1996)对此流动进行了模拟，并与实验确定的湍流随距网格距离的衰减率很好地吻合。能量大约衰减为 x − α
其中 2 < α < 3；指数 α 的确定在实验和计算上都很困难，因为快速衰减无法提供足够大的区域来准确计算其值。

Briggs 等人使用基于该流程模拟的可视化。 (1996)表明，这种流动中湍流扩散的主要机制是高能流体团通过未受干扰的流体的运动。这似乎是一个简单且合乎逻辑的解释，但与之前的提议相反。图 10.2 显示了

\begin{sourceblock}
\includegraphics[page=372,trim={53.00bp 97.70bp 53.87bp 466.65bp},clip,width=\textwidth,max height=0.38\textheight]{Ferziger4th.pdf}
\par\vspace{0.45\baselineskip}
\small 图10.2 静止流体中振荡网格产生的流动平面上的动能等值线；网格位于图的顶部。充满能量的流体包将能量从电网区域转移出去。来自布里格斯等人。 （1996）
\end{sourceblock}
\begin{sourceblock}
\includegraphics[page=373,trim={53.00bp 470.25bp 53.87bp 52.00bp},clip,width=\textwidth,max height=0.38\textheight]{Ferziger4th.pdf}
\par\vspace{0.45\baselineskip}
\small 图 10.3 湍流动能通量 q 的轮廓与一些常用湍流模型的预测相比（Mellor 和 Yamada 1982；Hanjali ́c 和 Launder 1976，1980）；来自布里格斯等人。 （1996）
\end{sourceblock}
该流动中一个平面上的动能轮廓。人们可以看到，在整个流动过程中，大的能量区域的大小大致相同，但远离网格的能量区域较少。原因是那些平行于网格传播的包裹不会在垂直于网格的方向上移动很远，并且高能流体的小“斑点”很快就会被粘性扩散的作用破坏。

结果用于测试湍流模型。这种测试的一个典型例子如图 10.3 所示，其中给出了湍流动能通量的轮廓，并与一些常用湍流模型的预测进行了比较。很明显，即使在像这样简单的流程中，模型也不能很好地工作。可能的原因是，这些模型旨在处理剪切产生的湍流，其特征与振荡网格产生的湍流非常不同。

该模拟使用了专为模拟均匀湍流而设计的代码（Rogallo 1981）。周期性边界条件应用于所有三个方向；这意味着实际上存在周期性的网格阵列，但是只要相邻网格之间的距离足够大于湍流衰减所需的距离，这就不会造成问题。该代码使用了傅立叶谱方法和三阶龙格-库塔时间方法。

这些结果说明了 DNS 的一些重要特征。该方法允许计算可与实验数据进行比较的统计量以验证结果。它还允许计算在实验室中难以测量但对于评估模型有用的数量。同时，该方法产生流动的可视化，可以深入了解湍流的物理原理。在实验室中几乎不可能同时获得同一流程的统计数据和可视化。正如上面的例子所示，这种组合可能非常有价值。

在直接数值模拟中，人们可以以一种在实验室中难以或不可能实现的方式控制外部变量。在很多情况下，DNS 产生的结果与实验结果不一致，但事实证明前者更接近正确。一个例子是通道流中壁面附近的湍流统计分布； Kim 等人的结果。 （1987）被证明比更仔细地重复实验时更准确。 Bardina 等人提供了一个早期的例子。 （1980）解释了旋转对各向同性湍流影响的实验中一些明显反常的结果。

DNS 可以比其他方式更准确地调查某些影响。还可以尝试无法通过实验实现的控制方法。这样做的目的是提供对流动物理的深入了解，从而表明可以实现的可能性（并指出可实现方法的方向）。 Choi 等人进行的平板减阻和控制研究就是一个例子。 （1994）。他们表明，通过使用穿过壁（或脉动壁表面）的受控吹气和吸力，平板的湍流阻力可以减少 30\%。比尤利等人。 (1994) 使用最优控制方法证明了在低雷诺数下可以强制流动重新层化以及在高雷诺数下可以减少表面摩擦的可能性。

10.2.4.2 Re = 5,000 时球体周围的流动
球体是一个形状非常简单的物体，但它周围的流体流动却非常复杂。这种流的 DNS 在中等雷诺数下是可能的；塞德尔等人。 (1998) 提出了使用二阶时间和空间离散化以及具有局部细化的非结构化六面体网格的模拟。在这个阶段我们不会深入讨论流动物理的细节（在描述 LES 方法之后会详细介绍），但我们想介绍一种测试模拟的可实现性并分析流动不同区域中湍流结构变化的方法。

\begin{center}\small 图 10.4 显示了模拟和实验的流动模式。尽管染料分布（来自两个孔，一个位于分离线前面的球体前部，一个位于分离后的球体后侧）与涡度的方位角分量并不完全对应，但很明显，在识别流动的主要特征时，两个图是一致的。流动分离发生在赤道附近；剪切层变得不稳定并卷起，形成涡环；这些最终在下游分解为各向同性湍流。正如实验和模拟所表明的那样，再循环区内的回流是湍流的。\end{center}
分析流结构并同时检查解的可实现性的另一种有趣方法是在 Lumley (1979) 首先引入的图中绘制在流场中各个点计算的各向异性张量的不变量。各向异性张量的分量定义如下：

\begin{sourceblock}
\includegraphics[page=375,trim={53.00bp 418.25bp 53.87bp 52.00bp},clip,width=\textwidth,max height=0.38\textheight]{Ferziger4th.pdf}
\par\vspace{0.45\baselineskip}
\small 图 10.4 DNS 和实验的流型比较：计算的瞬时涡度的方位角分量（上）和实验中染料分布的快照（来自 Seidl 等人，1998）
\end{sourceblock}
\begin{sourceblock}
\begin{tikzpicture}
\node[inner sep=0,anchor=south west] (image) at (0,0) {\includegraphics[page=375,trim={53.00bp 339.10bp 53.87bp 298.90bp},clip,width=0.92\textwidth,max height=0.38\textheight]{Ferziger4th.pdf}};
\begin{scope}[x={(image.south east)},y={(image.north west)}]
\fill[white] (0.37591,0.26866) rectangle (0.40138,0.70860);
\fill[white] (0.39997,0.26866) rectangle (0.41059,0.57321);
\fill[white] (0.41768,0.30704) rectangle (0.44586,0.71819);
\fill[white] (0.45296,0.51741) rectangle (0.49856,0.95736);
\fill[white] (0.49877,0.51741) rectangle (0.50938,0.82232);
\end{scope}
\end{tikzpicture}
\end{sourceblock}
其中 u i 是速度矢量的波动分量，q 代表湍流强度，q = u i u i，δ i j 是克罗内克三角洲。该张量的两个非零不变量 II 和 III 是：

\begin{sourceblock}
\begin{tikzpicture}
\node[inner sep=0,anchor=south west] (image) at (0,0) {\includegraphics[page=375,trim={53.00bp 258.62bp 53.87bp 379.50bp},clip,width=0.92\textwidth,max height=0.38\textheight]{Ferziger4th.pdf}};
\begin{scope}[x={(image.south east)},y={(image.north west)}]
\fill[white] (0.23323,0.29907) rectangle (0.24800,0.71163);
\fill[white] (0.24923,0.29907) rectangle (0.26403,0.71163);
\fill[white] (0.27146,0.30871) rectangle (0.29964,0.72127);
\fill[white] (0.30316,0.30871) rectangle (0.34998,0.95717);
\fill[white] (0.66589,0.26981) rectangle (0.69510,0.71163);
\fill[white] (0.69368,0.26981) rectangle (0.74974,0.71163);
\fill[white] (0.75675,0.30871) rectangle (0.76980,0.72127);
\fill[white] (0.92439,0.29550) rectangle (1.00000,0.70807);
\fill[white] (0.62331,0.04283) rectangle (0.66734,0.71163);
\end{scope}
\end{tikzpicture}
\end{sourceblock}
\begin{sourceblock}
\includegraphics[page=376,trim={53.00bp 287.25bp 53.87bp 52.00bp},clip,width=\textwidth,max height=0.38\textheight]{Ferziger4th.pdf}
\par\vspace{0.45\baselineskip}
\small 图 10.5 平均流动流线上评估各向异性张量不变量的点（上）以及这些点在不变量图中的位置（下）；来自塞德尔等人。 （1998）
\end{sourceblock}
流线上的所有点都正确地位于由限制线界定的三角形内部，这表明模拟没有违反物理定律——解是可实现的。由于球体上游的流动是层流，并且其后面的分离区域不封闭（意味着“新鲜”流体不断进入该区域，而“较旧”流体不断离开该区域），因此出现了一个有趣的问题：来自非湍流上游区域的流体元素如何在地图内进入湍流状态？从图 10.5（下）看来，唯一的可能性是在三角形的顶角：分离流线附近的流体元素在发展为多维湍流之前首先经历一维波动。离开再循环区的流体元素保留在湍流尾流中，其中湍流在下游缓慢衰减。

\subsection*{10.2.5 DNS的其他应用}\phantomsection\addcontentsline{toc}{subsection}{10.2.5 DNS的其他应用}
球体后面的流动是一个具有挑战性的研究课题，因为它包含所有类型的湍流状态：通过扫描湍流尾流，可以找到与不变图中所有位置相对应的点。通常使用 DNS 研究的许多其他流（例如平面通道或管道中完全发育的流）仅覆盖地图的一小部分。然而，由于仅在湍流区域内需要非常精细的网格，因此需要允许局部网格细化的求解方法。此类方法通常基于有限体积方法并使用具有二阶离散化的非结构化网格，因此与用于更简单几何形状的高阶方法相比，需要更精细的网格才能获得满意的精度。使用结构化网格，如 Pal 等人使用的网格。 (2017) 在他们对亚临界 5 雷诺数 3700 和中等弗劳德数下流经球体的分层流的研究中，在不需要精细网格的区域浪费了大量网格点；这在雷诺数较高时变得尤为重要。

针对简单几何形状的专门求解方法的优点在于，它们可以充分求解更高雷诺数的流动。 Lee and Moser (2015)在他们的DNS平面通道流量中达到了Re = 2.5 × 10 5，基于平均速度和通道宽度。研究的流动表现出高雷诺数壁界湍流的特征，并允许对边界层进行更详细的分析，从而为开发新的或调整现有的湍流模型提供宝贵的信息。见第 10.3.5.5 节 讨论依赖于所谓的对数定律的墙函数，其中 DNS 表明某些“定律”参数并不像迄今为止所认为的那样通用。

随着计算技术的进步，DNS 将继续应用于更复杂的几何形状和更高的雷诺数，从而提供无法测量或达不到所需精度的数据，但在科学和工程中都非常有用。

\section*{10.3 用模型模拟湍流}\phantomsection\addcontentsline{toc}{section}{10.3 用模型模拟湍流}
\subsection*{10.3.1 模型类别}\phantomsection\addcontentsline{toc}{subsection}{10.3.1 模型类别}
如上所述，当纳维-斯托克斯方程以某种方式平均时，结果是由于非线性对流项，方程不闭合。这意味着平均方程包含涉及未知波动变量的非线性相关项。人们可以推导出这些相关性的精确方程，但它们包含更复杂的未知相关项，这些项无法从平均变量和方程所针对的项确定。

\footnotetext[5]{经过球体的亚临界流是低雷诺数状态，其中阻力系数不是雷诺数的函数，并且存在层流边界层分离。}
因此，需要为这些未知的相关性构建模型，即近似表达式或方程。 Leschziner (2010) 提出了一系列有用的约束条件和模型所需的特征；我们解释如下：模型应该是：

\begin{itemize}
\item 基于理性原理和物理概念，而不是直觉；
\item 根据适当的数学原理构建，例如维度同质性、一致性和框架不变性；
\item 受限于产生物理上可实现的行为；
\item 广泛适用；
\item 数学上简单；
\item 由具有可访问边界条件的变量构建；
\item 计算稳定。
\end{itemize}
本节和以下部分专门介绍生成模型的方法。我们可以将方法分类如下：

1. RANS、不稳定RANS（URANS）和瞬态RANS（TRANS）：当奥斯本·雷诺兹（Osborne Reynolds，1895）提出后来称为雷诺平均纳维-斯托克斯方程时，它们是在空间体积上而不是时间上进行平均的。然而，文献中的传统平均是随时间或整体的平均。

a.对所有时间求平均值：如果对所有时间求平均值，则所得结果
方程是传统的稳态RANS方程，即由平均值定义的平均流是稳态的。 b.与时间尺度相比，在较长的时间内进行平均，例如
湍流的积分尺度（参见 Chen 和 Jaw 1998）：如果平均值的时间比湍流的时间尺度长，则所得方程将不稳定，成为原始方程的时间低通滤波版本。这种方法很少（如果有的话）实际使用，部分原因是高频湍流的模型必须取决于滤波器比例，如下所述。 c.对系综进行平均：如果方程对系综进行平均
（一组统计上相同的实现），然后求平均会消除所有湍流（随机）涡流，因为它们在系综的不同成员之间具有随机相位。如果系综在统计上是平稳的，那么平均流将是稳定的。然而，如果所研究的运动中存在连贯的确定性元素，则所得的平均流方程可能不稳定，6 即系综不是静止的。一个例子是由边界的周期性运动驱动的流动，即内燃机气缸中的流动，由具有可重复模式的活塞和阀门的运动驱动。然后，这些方程通常被称为非稳态 RANS 或 URANS（Durbin 2002；Iaccarino 等人 2003；Wegner 等人 2004）。 Hanjali ́c (2002) 提出了一个与此类似的构造，他称之为瞬态 RANS (TRANS)（参见第 10.3.7 节）。在任何一种情况下，平均方程仍然是未闭合的，因此需要模型来模拟湍流效应。结果

\footnotetext[6]{Chen 和 Jaw (1998) 在图 1.8 中说明了整体平均值的创建。}
关于平均过程，Wyngaard（2010）提醒我们“在湍流的任何实现中，即使是瞬间，集合平均场也不可能存在。”

文献通常对平均或过滤过程并不精确。事实上，对于这里列出的所有方法，很少明确地进行平均或过滤；平均时间和所需集合的数量均未定义。因此，在所有情况下，平均方程看起来基本相同，并且通常应用相同的湍流模型。这种方法对于 RANS 和集合平均方程（URANS 和 TRANS）可能是正确的，因为在这些情况下，假设所有湍流（随机运动）都已被平均。然后，湍流模型预计能够表示缺失部分对剩余平均流的影响，即使它是不稳定的并且支持例如确定性运动。另一方面，时间滤波 RANS 方程求解流动的低频运动，同时对高频或小时间尺度运动进行建模。因此，在这样的模拟中，湍流模型必须取决于最高解析频率的时间尺度，其模式与 LES 方法一致（如下）；很少有研究（如果有的话）实际上采用这种方法，因此几乎所有模拟都是 RANS（长期平均且稳定）或 URANS/TRANS（集合平均且可能不稳定）模拟，但没有解决湍流问题。 2. LES、VLES：对于大涡模拟 (LES) 和超大涡模拟 (VLES)，纳维-斯托克斯方程和标量方程在空间上进行过滤（平均）。所得方程本质上与 URANS 方程相同，只是未封闭项的模型具有不同的含义和（可能）形式；值得注意的是，虽然 URANS 的大部分湍流能量处于子滤波器尺度，但 LES 可以解析大部分湍流能量。因此，LES 和 VLES 解决了所有大于滤波器尺寸的不稳定特征，同时对小尺度（子滤波器尺度）特征进行建模。关于所涉及的尺度有多种想法。 Pope (2000) 提供了一些指导原则，即如果“滤波器和网格足够精细，可以解析各处 80\% 的能量”，则模拟就是 LES。该定义在墙壁附近有一些警告。如果“滤波器和网格太粗糙，无法解析 80\% 的能量”，则模拟为 VLES。正如想象的那样，这更加重视拥有包含尽可能多物理场的子滤波器模型。布莱恩等人。 (2003) 和 Wyngaard (2004) 提出 LES 的网格尺度应该远小于能谱峰值的长度尺度。 Matheou 和 Chung (2014) 使用 Kolmogorov 能谱来量化大气边界层中流动的这些标准。他们假设，要解析80\%的湍流动能，网格尺度应小于峰值谱尺度的1/12，而要解析90\%的TKE，网格尺度应小于峰值谱尺度的1/32。他们的经验是，合理的收敛需要 90\% 的水平。 3. ILES：这种方法实际上隐式地生成模型，因此被称为隐式大涡模拟。我们将在下一节开始讨论它。

\subsection*{10.3.2 隐式大涡模拟（ILES）}\phantomsection\addcontentsline{toc}{subsection}{10.3.2 隐式大涡模拟（ILES）}
隐式大涡模拟（ILES）的概念是，可以将数值方法直接应用于纳维-斯托克斯方程，然后定制数值方案，使所得方法既具有涡流解析又具有适当的耗散性，即，它是一种大涡模拟，但没有针对亚网格尺度波动的显式模型；这种波动无法用解来表示，因为网格是一个空间滤波器（网格上可以解析的最高波数是 π/ x ）。

理解 ILES 概念的一种简单方法是使用修正方程方法（参见第 6.3.2.2 节），其中我们试图检查数值方法中的截断误差项。 Rider (2007) 对可压缩湍流进行了这种分析，从分析中可以看出截断项是什么以及它们如何构成有效的亚网格尺度模型。该专着由格林斯坦等人编辑。 ( 2007 ) 提供了关于 ILES 实施的广泛概述以及一些具体示例。在那本书中，Smolarkiewicz 和 Margolin (2007) 提供了令人信服的证据，证明他们的 MPDATA 是一个 ILES 代码，可以准确模拟地球物理流，特别是大气环流和边界层运动。 MPDATA 的核心是平流项的迭代有限差分近似；它是二阶精确且保守的。迭代首先使用上游差分项，然后进行第二遍以提高准确性。作者指出 MPDATA 属于非振荡 Lax-Wendroff 方案类别。

阿斯普登等人。 (2008) 基于尺度分析给出了 ILES 的不同且富有洞察力的观点。为了进行计算，他们采用了 Lawrence Berkeley NL 的 CCSE IAMR 代码，该代码是一种不可压缩的变密度 FV 分数阶方法，在空间和时间上具有二阶精度，使用不分裂的 Godunov 方法（Colella 1990）和单调性限制的四阶 CD 斜率近似（Colella 1985）。这种类型的任何方法都应该满足以下要求：如果我们不断细化网格，则应该解决更多的湍流并减少建模，直到我们达到所有湍流都将被解决的网格细度（即，我们有一个 DNS）并且子网格尺度模型的贡献可以忽略不计。 ILES显然满足了这个要求。

\subsection*{10.3.3 大涡模拟（LES）}\phantomsection\addcontentsline{toc}{subsection}{10.3.3 大涡模拟（LES）}
湍流包含广泛的长度和时间尺度；流动中可能出现的涡流尺寸范围示意性地显示在图 10.6 的左侧。该图的右侧显示了流动中某一点的典型速度分量的时间历史；波动发生的尺度范围是显而易见的。

\begin{sourceblock}
\includegraphics[page=381,trim={53.00bp 473.25bp 53.87bp 52.00bp},clip,width=\textwidth,max height=0.38\textheight]{Ferziger4th.pdf}
\par\vspace{0.45\baselineskip}
\small 图 10.6 湍流运动的示意图（左）和一点速度分量的时间依赖性（右）
\end{sourceblock}
大尺度运动通常比小尺度运动更有活力；它们的大小和强度使它们成为迄今为止最有效的保存财产运输者。小鳞片通常要弱得多，并且几乎不能传输这些特性。处理大涡流比处理小涡流更精确的模拟可能是有意义的。大涡模拟就是这样一种方法。大涡模拟是三维的并且与时间相关。虽然价格昂贵，但比相同流量的 DNS 成本低得多。一般来说，只要可行，DNS 就会成为首选方法，因为它更准确。对于雷诺数过高或几何结构过于复杂而无法应用 DNS 的流，LES 是首选方法（参见例如 Rodi 等人 2013 或 Sagaut 2006）。例如，LES 已成为大气科学中的主要工具，涉及云、降水、污染输送、山谷风流（例如，Chen 等人，2004b；Chow 等人，2006；Shi 等人，2018b，a）； LES 处理的行星边界层群落物理过程包括浮力、旋转、夹带、凝结以及与粗糙地面和海洋表面的相互作用（Moeng 和 Sullivan 2015）。 LES 甚至已经进入高压燃气轮机的设计（Bhaskaran 和 Lele 2010），并在气动声学模拟中发挥主导作用（例如 Bodony 和 Lele 2008；Brès 等人 2017）。

必须精确定义要计算的数量。我们需要一个仅包含总场的大尺度分量的速度场。这最好通过过滤速度场 ( ? ) 来产生；在这种方法中，要模拟的大尺度场或解析尺度场本质上是整个场的局部平均值。我们将使用一维符号；推广到三个维度是很简单的。滤波后的速度定义为：

\begin{sourceblock}
\begin{tikzpicture}
\node[inner sep=0,anchor=south west] (image) at (0,0) {\includegraphics[page=381,trim={53.00bp 110.32bp 53.87bp 536.66bp},clip,width=0.92\textwidth,max height=0.38\textheight]{Ferziger4th.pdf}};
\begin{scope}[x={(image.south east)},y={(image.north west)}]
\fill[white] (0.30734,0.06263) rectangle (0.37702,0.72234);
\fill[white] (0.38054,0.11900) rectangle (0.40872,0.72234);
\end{scope}
\end{tikzpicture}
\end{sourceblock}
其中滤波器内核 G ( x , x ′ ) 是局部函数。 7 LES 中应用的滤波器内核包括高斯滤波器、箱式滤波器（简单的局部平均值）和截止滤波器（消除属于截止值以上波数的所有傅里叶系数的滤波器）。每个过滤器都有一个与之相关的长度标度，。粗略地说，尺寸大于 的涡流是大涡流，而尺寸小于 的涡流是需要建模的小涡流。

当对恒定密度（不可压缩流）的纳维-斯托克斯方程进行滤波时，8 得到一组与 URANS 方程非常相似的方程：

\begin{sourceblock}
\begin{tikzpicture}
\node[inner sep=0,anchor=south west] (image) at (0,0) {\includegraphics[page=382,trim={53.00bp 473.98bp 53.87bp 156.05bp},clip,width=0.92\textwidth,max height=0.38\textheight]{Ferziger4th.pdf}};
\begin{scope}[x={(image.south east)},y={(image.north west)}]
\fill[white] (0.11062,0.42260) rectangle (0.19783,0.77264);
\fill[white] (0.23690,0.42315) rectangle (0.32788,0.77319);
\fill[white] (0.32809,0.42315) rectangle (0.35347,0.77319);
\fill[white] (0.20325,0.25921) rectangle (0.23143,0.57934);
\fill[white] (0.13823,0.05594) rectangle (0.16752,0.38355);
\end{scope}
\end{tikzpicture}
\end{sourceblock}
由于连续性方程是线性的，因此滤波不会改变它：

\begin{sourceblock}
\includegraphics[page=382,trim={53.00bp 415.52bp 53.87bp 221.52bp},clip,width=0.92\textwidth,max height=0.38\textheight]{Ferziger4th.pdf}
\end{sourceblock}
值得注意的是，因为

\begin{sourceblock}
\includegraphics[page=382,trim={53.00bp 370.05bp 53.87bp 280.63bp},clip,width=0.92\textwidth,max height=0.38\textheight]{Ferziger4th.pdf}
\end{sourceblock}
并且这个不等式左边的数量不容易计算，这是这个不等式两侧之间差异的建模近似，

\begin{sourceblock}
\includegraphics[page=382,trim={53.00bp 309.09bp 53.87bp 338.95bp},clip,width=0.92\textwidth,max height=0.38\textheight]{Ferziger4th.pdf}
\end{sourceblock}
必须介绍一下。因此，在 LES 中求解的动量方程为：

\begin{sourceblock}
\begin{tikzpicture}
\node[inner sep=0,anchor=south west] (image) at (0,0) {\includegraphics[page=382,trim={53.00bp 249.44bp 53.87bp 380.59bp},clip,width=0.92\textwidth,max height=0.38\textheight]{Ferziger4th.pdf}};
\begin{scope}[x={(image.south east)},y={(image.north west)}]
\fill[white] (0.06902,0.42260) rectangle (0.15624,0.77264);
\fill[white] (0.19531,0.42315) rectangle (0.28632,0.77319);
\fill[white] (0.28653,0.42315) rectangle (0.31191,0.77319);
\fill[white] (0.16168,0.25921) rectangle (0.18986,0.57934);
\fill[white] (0.09666,0.05594) rectangle (0.12592,0.38355);
\end{scope}
\end{tikzpicture}
\end{sourceblock}
在 LES 的背景下，τ s
i j 称为亚网格尺度雷诺应力。 “压力”这个名字源于它的处理方式，而不是它的物理性质。它实际上是由小尺度或未解决尺度的作用引起的大尺度动量通量。 “次网格尺度”这个名称也有点用词不当。滤波器的宽度 不需要与网格大小 h 任何关系，除了前面提到的 ≥ h 的明显条件之外。传统上，作者建立了这样的联系
7 读者可能会注意到这里没有提到网格。过滤器尺寸至少与网格尺寸一样大，并且通常要大得多。在传统的 LES 中，滤波器宽度被假定为网格大小，现在我们将遵循这种做法。见第 10.3.3.4 和 10.3.3.7 节 中，我们将探讨网格和滤波器宽度之间的关系以及它如何影响建模。 8 与上述时间和集合平均方程的情况一样，定义了滤波器操作，但方程通常没有显式滤波。这些方程本质上是传统 LES 中隐式未知滤波器的结果。博斯等人。 (2010)已经展示了显式过滤的价值。

他们的命名法已经固定下来。现在，用于近似亚网格尺度雷诺应力 (10.11) 的模型称为亚网格尺度 (SGS) 或亚滤波器尺度 (SFS) 模型，具体取决于所使用的模型，如下所述。

亚网格尺度雷诺应力包含小尺度场的局部平均值，因此它的模型应该基于局部速度场，或者可能基于局部流体运动的过去历史。后者可以通过使用求解偏微分方程的模型来获得确定 SGS 雷诺应力所需的参数来完成。

10.3.3.2 Smagorinsky 和相关模型
最早且最常用的亚网格尺度模型是 Smagorinsky (1963) 提出的。它是一个涡粘模型。所有此类模型都基于这样的概念：SGS 雷诺应力的主要影响是增加传输和耗散。由于这些现象是由于层流中的粘度造成的，因此可以合理地假设一个可行的模型可能是：

\begin{sourceblock}
\begin{tikzpicture}
\node[inner sep=0,anchor=south west] (image) at (0,0) {\includegraphics[page=383,trim={53.00bp 366.39bp 53.87bp 263.65bp},clip,width=0.92\textwidth,max height=0.38\textheight]{Ferziger4th.pdf}};
\begin{scope}[x={(image.south east)},y={(image.north west)}]
\fill[white] (0.23967,0.38255) rectangle (0.25263,0.61994);
\fill[white] (0.22310,0.25900) rectangle (0.24087,0.57922);
\end{scope}
\end{tikzpicture}
\end{sourceblock}
其中 μ t 是涡流粘度，S i j 是大尺度或解析场的应变率。 Wyngaard ( 2010 ) 展示了 Lilly 在 1967 年如何使用 τ s 的演化方程正式推导出该模型
我 j .类似的模型也经常与 RANS 方程结合使用；见下文。

亚网格尺度涡粘性的形式可以通过量纲参数导出，为：

\begin{sourceblock}
\begin{tikzpicture}
\node[inner sep=0,anchor=south west] (image) at (0,0) {\includegraphics[page=383,trim={53.00bp 268.58bp 53.87bp 379.80bp},clip,width=0.92\textwidth,max height=0.38\textheight]{Ferziger4th.pdf}};
\begin{scope}[x={(image.south east)},y={(image.north west)}]
\fill[white] (0.00000,0.85135) rectangle (0.07570,1.00000);
\fill[white] (0.07841,0.85135) rectangle (0.12647,1.00000);
\fill[white] (0.12917,0.85135) rectangle (0.16229,1.00000);
\end{scope}
\end{tikzpicture}
\end{sourceblock}
其中 C S 是待确定的模型参数， 是滤波器长度尺度，| S | = (S i j S i j ) 1 / 2。涡流粘度的这种形式可以通过多种方式导出。理论提供了参数的估计。大多数这些方法仅适用于各向同性湍流，它们都同意 C S ≈ 0.2. 不幸的是，C S 不是常数；它可以是雷诺数和/或其他无量纲参数的函数并且可以在不同的流中取不同的值。

斯马戈林斯基模型虽然相对成功，但也并非没有问题，并且其使用已经下降，取而代之的是下面描述的更复杂的模型。不建议这样做，但它的使用仍然存在。为了用它模拟通道流动，需要进行一些修改。大部分流中的参数 C S 的值必须从 0 开始减小。 2 至大约 0.065，它使涡流粘度降低了几乎一个数量级。所有剪切流都需要这种幅度的变化。在靠近通道表面的区域，该值必须进一步降低。一种成功的方法是借用 van Driest 阻尼，该阻尼长期以来一直用于降低 RANS 模型中的近壁涡流粘度：

\begin{sourceblock}
\begin{tikzpicture}
\node[inner sep=0,anchor=south west] (image) at (0,0) {\includegraphics[page=384,trim={53.00bp 590.47bp 53.87bp 50.00bp},clip,width=0.92\textwidth,max height=0.38\textheight]{Ferziger4th.pdf}};
\begin{scope}[x={(image.south east)},y={(image.north west)}]
\fill[white] (0.32926,0.05999) rectangle (0.36653,0.54499);
\fill[white] (0.37152,0.10479) rectangle (0.39970,0.55551);
\fill[white] (0.40310,0.05999) rectangle (0.45086,0.54499);
\end{scope}
\end{tikzpicture}
\end{sourceblock}
其中 n + 是距壁的距离（以粘性壁单位表示）（ n + = nu τ /ν，其中 u τ 是剪切速度，u τ = √ τ wall /ρ，τ wall 是壁处的剪切应力），
A + 是一个常量，通常约为 25。尽管这种修改产生了所需的结果，但在 LES 的背景下很难证明其合理性。

另一个问题是，在壁附近，流动结构非常各向异性。创建低速和高速流体区域（条纹）；它们在翼展方向和法向方向上的长度约为 1000 个粘性单位，宽度约为 30-50 个粘性单位。解决条纹问题需要高度各向异性的网格，并且在 SGS 模型中使用的长度尺度 的选择并不明显。通常的选择是 ( 1 2 3 ) 1 / 3 但 ( 2
1 + 2 2 + 2 3 ) 1 / 2 是可能的，其他的很容易构造；这里 i 是与滤波器在第 i 个坐标方向相关的宽度。

在稳定分层的流体中，有必要减小 Smagorinsky 参数。分层在地球物理流中很常见；通常的做法是使参数成为理查森数或弗劳德数的函数。这些是相关的无量纲参数，代表分层和剪切的相对重要性。在旋转和/或曲率起重要作用的流动中也会出现类似的效果。通常，理查森数基于平均流场的属性。在某些代码中，SGS 湍流动能的传递 k = 1

\begin{sourceblock}
\begin{tikzpicture}
\node[inner sep=0,anchor=south west] (image) at (0,0) {\includegraphics[page=384,trim={53.00bp 358.43bp 53.87bp 289.28bp},clip,width=0.92\textwidth,max height=0.38\textheight]{Ferziger4th.pdf}};
\begin{scope}[x={(image.south east)},y={(image.north west)}]
\fill[white] (0.00000,0.86706) rectangle (0.06983,1.00000);
\fill[white] (0.07239,0.86706) rectangle (0.19510,1.00000);
\fill[white] (0.19768,0.86706) rectangle (0.23913,1.00000);
\fill[white] (0.24165,0.86706) rectangle (0.38289,1.00000);
\fill[white] (0.38547,0.86706) rectangle (0.48180,1.00000);
\fill[white] (0.48439,0.86706) rectangle (0.50917,1.00000);
\fill[white] (0.51173,0.86706) rectangle (0.58472,1.00000);
\fill[white] (0.58731,0.86706) rectangle (0.62208,1.00000);
\fill[white] (0.62460,0.86706) rectangle (0.66602,1.00000);
\fill[white] (0.66857,0.86706) rectangle (0.79314,1.00000);
\fill[white] (0.79570,0.86706) rectangle (0.82544,1.00000);
\fill[white] (0.82800,0.86706) rectangle (0.86944,1.00000);
\fill[white] (0.87197,0.86706) rectangle (0.94168,1.00000);
\fill[white] (0.94424,0.86706) rectangle (1.00000,1.00000);
\fill[white] (0.00000,0.21867) rectangle (0.06487,0.84590);
\fill[white] (0.06608,0.21867) rectangle (0.09585,0.84590);
\fill[white] (0.09702,0.21867) rectangle (0.16508,0.84590);
\fill[white] (0.16626,0.21867) rectangle (0.24677,0.84590);
\fill[white] (0.24800,0.21867) rectangle (0.28941,0.84590);
\fill[white] (0.29062,0.21867) rectangle (0.40189,0.84590);
\fill[white] (0.40307,0.21867) rectangle (0.43281,0.84590);
\fill[white] (0.43402,0.21867) rectangle (0.49377,0.84590);
\fill[white] (0.49498,0.21867) rectangle (0.60734,0.84590);
\fill[white] (0.60854,0.21867) rectangle (0.69489,0.84590);
\fill[white] (0.69609,0.21867) rectangle (0.78186,0.84590);
\fill[white] (0.78313,0.22409) rectangle (0.80123,0.85133);
\fill[white] (0.80659,0.23874) rectangle (0.83477,0.86598);
\fill[white] (0.84189,0.47531) rectangle (0.85720,0.94031);
\fill[white] (0.89826,0.10635) rectangle (0.90890,0.57189);
\fill[white] (0.91359,0.21867) rectangle (1.00000,0.84590);
\fill[white] (0.00000,0.00000) rectangle (0.02638,0.28649);
\fill[white] (0.03344,0.00000) rectangle (0.06162,0.21704);
\fill[white] (0.06514,0.00000) rectangle (0.09221,0.20239);
\fill[white] (0.09756,0.00000) rectangle (0.16815,0.21704);
\fill[white] (0.17086,0.00000) rectangle (0.19567,0.19696);
\fill[white] (0.19832,0.00000) rectangle (0.32126,0.19696);
\fill[white] (0.32400,0.00000) rectangle (0.39868,0.19696);
\fill[white] (0.40138,0.00000) rectangle (0.47606,0.19696);
\fill[white] (0.47877,0.00000) rectangle (0.52683,0.19696);
\fill[white] (0.52953,0.00000) rectangle (0.58923,0.19696);
\fill[white] (0.59194,0.00000) rectangle (0.62006,0.19696);
\fill[white] (0.62274,0.00000) rectangle (0.72571,0.19696);
\fill[white] (0.72842,0.00000) rectangle (0.76983,0.19696);
\fill[white] (0.77254,0.00000) rectangle (0.90301,0.19696);
\fill[white] (0.90574,0.00000) rectangle (0.93552,0.19696);
\end{scope}
\end{tikzpicture}
\end{sourceblock}
i = u i − u i ) 被计算（Pope 2000）并用于将系数定义为

\begin{sourceblock}
\includegraphics[page=384,trim={53.00bp 322.05bp 53.87bp 325.87bp},clip,width=0.92\textwidth,max height=0.38\textheight]{Ferziger4th.pdf}
\end{sourceblock}
其中，湍流长度尺度 l 3 的垂直分量根据浮力进行调整； l 1 = l 2 = 或 h （即名义上等于网格或滤波器比例）；参见，例如，Xue (2000)。

因此，斯马戈林斯基模型存在许多困难。如果我们希望模拟更复杂和/或更高的雷诺数流，拥有更准确的模型可能很重要（例如，参见 Shi et al. 2018b , a ）。事实上，基于 DNS 数据结果的详细测试表明，Smagorinsky 模型在表示次网格尺度应力的细节方面相当差。特别是，涡粘模型迫使应力 τ s
i j 与应变率 S i j 一致，但实际上并非如此！

10.3.3.3 动态模型
模拟中解析的最小尺度在许多方面与通过模型处理的更小的尺度相似。这个想法导致了另一种亚网格尺度模型，即尺度相似性模型（Bardina et al. 1980）。主要论点是，已解析尺度和未解析尺度之间的重要相互作用涉及前者的最小涡流和后者的最大涡流，即比与过滤器相关的长度尺度 稍大或稍小的涡流。基于这个概念的论证得出以下模型：

\begin{sourceblock}
\includegraphics[page=385,trim={53.00bp 592.70bp 53.87bp 55.33bp},clip,width=0.92\textwidth,max height=0.38\textheight]{Ferziger4th.pdf}
\end{sourceblock}
其中双上划线表示已过滤两次的数量。该模型与实际的 SGS 雷诺应力相关性非常好，但几乎不耗散任何能量，不能作为“独立”的 SGS 模型。它将能量从较大尺度（前向散射）或多或少均匀地传递到最小尺度，以及从最小解析尺度到较大尺度（后向散射），这是有用的。为了纠正耗散的不足，人们可以结合 Smagorinsky 和尺度相似性模型来生成“混合”模型。该模型提高了模拟的质量。有关更多详细信息，请参阅 Sagaut (2006)。

尺度相似性模型的基本概念，即最小解析尺度运动可以提供可用于建模最大亚网格尺度运动的信息，可以更进一步，导致动态模型或过程（Germano 等人，1991）。该过程基于这样的假设：上述模型之一是小尺度的可接受的表示。

开创性的杰尔马诺程序的本质在于尺度不变性的思想（Meneveau and Katz 2000）：“尺度不变性意味着流动的某些特征在不同的运动尺度中保持相同。”这里，假设SGS模型中的系数和形式在网格尺度和网格尺度的某个倍数的尺度上保持相同。因此，虽然上面的原始过滤产生了

\begin{sourceblock}
\includegraphics[page=385,trim={53.00bp 329.69bp 53.87bp 318.34bp},clip,width=0.92\textwidth,max height=0.38\textheight]{Ferziger4th.pdf}
\end{sourceblock}
并且，使用 Smagorinsky 模型，

\begin{sourceblock}
\includegraphics[page=385,trim={53.00bp 272.32bp 53.87bp 365.80bp},clip,width=0.92\textwidth,max height=0.38\textheight]{Ferziger4th.pdf}
\end{sourceblock}
以“测试过滤器”比例再次过滤过滤后的方程（10.8）产生

\begin{sourceblock}
\includegraphics[page=385,trim={53.00bp 222.06bp 53.87bp 425.19bp},clip,width=0.92\textwidth,max height=0.38\textheight]{Ferziger4th.pdf}
\end{sourceblock}
and

\begin{sourceblock}
\begin{tikzpicture}
\node[inner sep=0,anchor=south west] (image) at (0,0) {\includegraphics[page=385,trim={53.00bp 178.81bp 53.87bp 459.31bp},clip,width=0.92\textwidth,max height=0.38\textheight]{Ferziger4th.pdf}};
\begin{scope}[x={(image.south east)},y={(image.north west)}]
\fill[white] (0.00000,0.81585) rectangle (0.04740,1.00000);
\end{scope}
\end{tikzpicture}
\end{sourceblock}
现在比率 \textasciicircum{} / 成为该方法的可调整常数；已经做出了各种选择，但通常使用--2的比率。杰尔马诺创建了现在被称为杰尔马诺恒等式的东西，L i j = T s

\begin{sourceblock}
\begin{tikzpicture}
\node[inner sep=0,anchor=south west] (image) at (0,0) {\includegraphics[page=385,trim={53.00bp 130.03bp 53.87bp 518.12bp},clip,width=0.92\textwidth,max height=0.38\textheight]{Ferziger4th.pdf}};
\begin{scope}[x={(image.south east)},y={(image.north west)}]
\fill[white] (0.00000,0.86882) rectangle (0.09395,1.00000);
\fill[white] (0.09522,0.86882) rectangle (0.15552,1.00000);
\fill[white] (0.15675,0.86882) rectangle (0.21814,1.00000);
\fill[white] (0.21937,0.86882) rectangle (0.29654,1.00000);
\fill[white] (0.29780,0.86882) rectangle (0.34027,1.00000);
\fill[white] (0.34150,0.86882) rectangle (0.45116,1.00000);
\fill[white] (0.45239,0.86882) rectangle (0.47050,1.00000);
\fill[white] (0.47170,0.86882) rectangle (0.53143,1.00000);
\fill[white] (0.53263,0.86882) rectangle (0.56241,1.00000);
\fill[white] (0.56355,0.86882) rectangle (0.60671,1.00000);
\fill[white] (0.60809,0.86882) rectangle (0.63290,1.00000);
\fill[white] (0.63408,0.86882) rectangle (0.70126,1.00000);
\fill[white] (0.70250,0.86882) rectangle (0.81877,1.00000);
\fill[white] (0.82003,0.86882) rectangle (0.91131,1.00000);
\fill[white] (0.91254,0.86882) rectangle (0.97558,1.00000);
\fill[white] (0.97681,0.86882) rectangle (1.00000,1.00000);
\fill[white] (0.00000,0.20400) rectangle (0.05498,0.84658);
\fill[white] (0.05489,0.20400) rectangle (0.14045,0.84658);
\fill[white] (0.14039,0.20400) rectangle (0.17014,0.84658);
\fill[white] (0.17008,0.20400) rectangle (0.21149,0.84658);
\fill[white] (0.21140,0.20400) rectangle (0.32767,0.84658);
\fill[white] (0.32758,0.20400) rectangle (0.43110,0.84658);
\fill[white] (0.43104,0.16398) rectangle (0.46220,0.86715);
\fill[white] (0.46075,0.16398) rectangle (0.47140,0.64091);
\fill[white] (0.47865,0.22457) rectangle (0.50683,0.86715);
\fill[white] (0.51101,0.20956) rectangle (0.53248,0.85214);
\fill[white] (0.53221,0.45692) rectangle (0.54517,0.93330);
\fill[white] (0.59450,0.08949) rectangle (0.60514,0.56643);
\fill[white] (0.60370,0.08949) rectangle (0.62535,0.84658);
\fill[white] (0.62529,0.20400) rectangle (0.68836,0.84658);
\fill[white] (0.68830,0.20400) rectangle (0.76632,0.84658);
\fill[white] (0.76623,0.20400) rectangle (0.78770,0.84658);
\fill[white] (0.78764,0.20400) rectangle (0.81242,0.84658);
\fill[white] (0.81236,0.20400) rectangle (1.00000,0.84658);
\fill[white] (0.00000,0.00000) rectangle (0.02746,0.18232);
\fill[white] (0.03014,0.00000) rectangle (0.09744,0.18232);
\end{scope}
\end{tikzpicture}
\end{sourceblock}
\begin{sourceblock}
\begin{tikzpicture}
\node[inner sep=0,anchor=south west] (image) at (0,0) {\includegraphics[page=385,trim={53.00bp 86.08bp 53.87bp 552.04bp},clip,width=0.92\textwidth,max height=0.38\textheight]{Ferziger4th.pdf}};
\begin{scope}[x={(image.south east)},y={(image.north west)}]
\fill[white] (0.03856,0.26945) rectangle (0.06971,0.72091);
\fill[white] (0.06827,0.26945) rectangle (0.07892,0.57566);
\fill[white] (0.08598,0.30835) rectangle (0.11417,0.72091);
\fill[white] (0.11768,0.30835) rectangle (0.13883,0.72091);
\fill[white] (0.15714,0.26945) rectangle (0.18421,0.71128);
\fill[white] (0.18791,0.29872) rectangle (0.20770,0.71128);
\fill[white] (0.20791,0.26945) rectangle (0.21856,0.57566);
\fill[white] (0.22397,0.26945) rectangle (0.28105,0.80086);
\fill[white] (0.27979,0.29872) rectangle (0.29958,0.80086);
\fill[white] (0.29979,0.26945) rectangle (0.31044,0.57566);
\end{scope}
\end{tikzpicture}
\end{sourceblock}
其中 L kk 包含各向同性项。 9 在此等式中。 ( 10.17 ) 系数C 2
S 是唯一的未知数，可以通过多种方式确定。各向同性项相互抵消，例如，参见 Sagaut (2006)、Lilly (1992) 或 Germano 等人。 (1991)，因为不可压缩流中S ii = 0，各向同性项在实际求解方法中以S i j δ i j L kk 的形式出现。

动态模型的基本要素是（1）它使用解析场中最小涡流的信息来计算系数，（2）假设具有相同系数值的同一模型适用于实际 LES 和较粗尺度的滤波方程。动态过程将模型系数给出为两个量的比率，并且直接根据当前解析的变量在每个空间网格点和每个时间步计算系数。这里我们不介绍实际的过程；但是，请注意方程式。 ( 10.17 ) 对于一个未知数产生六个方程，因此是超定的。然后动态过程设置方案以最小化误差

\begin{sourceblock}
\begin{tikzpicture}
\node[inner sep=0,anchor=south west] (image) at (0,0) {\includegraphics[page=386,trim={53.00bp 402.97bp 53.87bp 235.15bp},clip,width=0.92\textwidth,max height=0.38\textheight]{Ferziger4th.pdf}};
\begin{scope}[x={(image.south east)},y={(image.north west)}]
\fill[white] (0.12920,0.26945) rectangle (0.15281,0.72091);
\fill[white] (0.15137,0.26945) rectangle (0.16202,0.57566);
\fill[white] (0.16905,0.30835) rectangle (0.19723,0.72091);
\fill[white] (0.20075,0.26945) rectangle (0.23191,0.72091);
\fill[white] (0.23047,0.26945) rectangle (0.24111,0.57566);
\fill[white] (0.24650,0.29515) rectangle (0.32944,0.77302);
\fill[white] (0.31344,0.22734) rectangle (0.34749,0.72091);
\end{scope}
\end{tikzpicture}
\end{sourceblock}
杰尔马诺等人。 ( 1991 ) 将误差与应变率收缩，即 ε i j S i j 生成 C 2 的单个方程
S。对 Germano 等人提出的原始模型的两项重大改进。 (1991) 是 (1) Lilly (1992) 建议的用于评估系数的最小二乘法和 (2) Wong 和 Lilly (1994) 引入新的基本模型来取代 Smagorinsky 模型；他们的模型基于 Kolomogorov 缩放，不需要在动态过程中计算应变率，从而使壁边界条件不太重要。它们的涡流粘度为

\begin{sourceblock}
\begin{tikzpicture}
\node[inner sep=0,anchor=south west] (image) at (0,0) {\includegraphics[page=386,trim={53.00bp 281.89bp 53.87bp 364.68bp},clip,width=0.92\textwidth,max height=0.38\textheight]{Ferziger4th.pdf}};
\begin{scope}[x={(image.south east)},y={(image.north west)}]
\fill[white] (0.00000,0.72969) rectangle (0.07573,1.00000);
\fill[white] (0.07841,0.72969) rectangle (0.10322,1.00000);
\end{scope}
\end{tikzpicture}
\end{sourceblock}
C ε 是动态过程中确定的系数，不要求耗散率等于SGS 能量产生率。

以 Smagorinsky 或 Wong-Lilly 模型为基础的动态过程消除了前面描述的许多困难：

\begin{itemize}
\item 在剪切流中，Smagorinsky 模型参数需要比各向同性湍流中的参数小得多。动态模型自动产生这种变化。
\item 靠近墙壁时模型参数必须进一步减小。动态模型会在壁面附近以正确的方式自动减小参数。
\item 各向异性网格或滤波器的长度尺度的定义尚不清楚。对于动态模型来说，这个问题变得毫无意义，因为该模型通过更改参数值来补偿长度尺度中的任何误差。
\end{itemize}
9 眼光敏锐的读者会注意到，我们悄悄地从方程 1 中的测试滤波器平均值中提取了系数 C 2 S，它是空间和时间的函数。 （10.17）；这相当于假设它在测试过滤器的体积上是恒定的。这是一个方便的选择，但不是唯一可能的选择。

虽然它比 Smagorinsky 模型有相当大的改进，但动态过程存在问题。它产生的模型参数是空间坐标和时间的快速变化函数，因此涡流粘度的两个符号都取很大的值。虽然负涡粘性被认为是表示从小尺度到大尺度的能量传递的一种方式（这个过程称为反向散射），但如果涡粘性在太大的空间区域或太长的时间内为负，则可能并且确实会发生数值不稳定。一种解决方法是将任何涡流粘度 μ t < − μ 重置为 − μ，即等于分子粘度的负值；这称为剪辑。另一种有用的替代方法是采用空间或时间平均。欲知详情，读者可参阅上述引用的论文。这些技术产生了进一步的改进，但仍不完全令人满意；寻找更稳健的子网格尺度模型是当前研究的主题，如下所示。

动态模型所依据的论点并不限于 Smagorinsky 模型。相反，我们可以使用混合斯马戈林斯基尺度相似性模型。 Zang 等人使用了混合模型。 (1993) 以及 Shah 和 Ferziger (1995) 取得了相当大的成功。

最后，我们注意到已经设计了其他版本的动态过程来克服最简单形式的模型的困难。其中较好的模型之一是 Meneveau 等人的拉格朗日动力学模型。 （1996）。在此模型中，动态过程的模型参数表达式的分子和分母中的项沿流动轨迹进行平均。这是通过求解拉格朗日轨道上各项积分的偏微分“松弛-传递”方程来完成的。该模型的应用在第 4 节中给出。 10.3.4.3.

10.3.3.4 使用子滤波器尺度重建的模型
卡拉蒂等人。 (2001) 导出了 LES 的控制方程，这自然导致了 SFS 和 SGS 混合模型，并区分了纳维-斯托克斯方程的显式滤波和数值计算的离散化。虽然 LES 的显式滤波器由等式给出。 ( 10.7 )，离散化的效果也是一个（不幸的是）未知的滤波器（参见第 10.3.2 节的 ILES 方法）。离散网格的控制方程可以通过将显式滤波器 (overbar) 和离散算子 (tilde) 应用于不可压缩且恒定密度的纳维-斯托克斯和连续性方程来获得：

\begin{sourceblock}
\begin{tikzpicture}
\node[inner sep=0,anchor=south west] (image) at (0,0) {\includegraphics[page=387,trim={53.00bp 133.58bp 53.87bp 496.45bp},clip,width=0.92\textwidth,max height=0.38\textheight]{Ferziger4th.pdf}};
\begin{scope}[x={(image.south east)},y={(image.north west)}]
\fill[white] (0.03675,0.42260) rectangle (0.10580,0.77292);
\fill[white] (0.14487,0.45306) rectangle (0.17699,0.77319);
\fill[white] (0.17218,0.42315) rectangle (0.19925,0.77319);
\fill[white] (0.19798,0.44558) rectangle (0.21777,0.77319);
\fill[white] (0.21798,0.42315) rectangle (0.24334,0.77319);
\fill[white] (0.11125,0.25921) rectangle (0.13943,0.57934);
\fill[white] (0.25047,0.25921) rectangle (0.27865,0.57934);
\fill[white] (0.28217,0.25921) rectangle (0.33050,0.76267);
\fill[white] (0.05531,0.05594) rectangle (0.08457,0.38355);
\fill[white] (0.16974,0.05594) rectangle (0.20481,0.38355);
\fill[white] (0.20427,0.03323) rectangle (0.21492,0.27084);
\fill[white] (0.30911,0.06342) rectangle (0.33026,0.38355);
\end{scope}
\end{tikzpicture}
\end{sourceblock}
and

\begin{sourceblock}
\begin{tikzpicture}
\node[inner sep=0,anchor=south west] (image) at (0,0) {\includegraphics[page=387,trim={53.00bp 84.11bp 53.87bp 552.93bp},clip,width=0.92\textwidth,max height=0.38\textheight]{Ferziger4th.pdf}};
\begin{scope}[x={(image.south east)},y={(image.north west)}]
\fill[white] (0.00000,0.90928) rectangle (0.04740,1.00000);
\end{scope}
\end{tikzpicture}
\end{sourceblock}
总 SFS 应力 10 定义为

\begin{sourceblock}
\includegraphics[page=388,trim={53.00bp 569.98bp 53.87bp 80.71bp},clip,width=0.92\textwidth,max height=0.38\textheight]{Ferziger4th.pdf}
\end{sourceblock}
这直接导致了非常有洞察力的分解，即

\begin{sourceblock}
\begin{tikzpicture}
\node[inner sep=0,anchor=south west] (image) at (0,0) {\includegraphics[page=388,trim={53.00bp 519.50bp 53.87bp 131.19bp},clip,width=0.92\textwidth,max height=0.38\textheight]{Ferziger4th.pdf}};
\begin{scope}[x={(image.south east)},y={(image.north west)}]
\fill[white] (0.04791,0.09968) rectangle (0.10069,0.92233);
\fill[white] (0.10571,0.17411) rectangle (0.13389,0.92233);
\fill[white] (0.13738,0.10421) rectangle (0.16084,0.92233);
\fill[white] (0.15940,0.10421) rectangle (0.17005,0.65890);
\fill[white] (0.17714,0.17411) rectangle (0.20532,0.92233);
\fill[white] (0.20884,0.17411) rectangle (0.22998,0.92233);
\fill[white] (0.24827,0.10421) rectangle (0.29386,0.90550);
\fill[white] (0.29408,0.10421) rectangle (0.30472,0.65890);
\fill[white] (0.31011,0.10421) rectangle (0.36722,0.92233);
\fill[white] (0.36592,0.15663) rectangle (0.39041,0.92233);
\fill[white] (0.38592,0.10421) rectangle (0.39657,0.65890);
\end{scope}
\end{tikzpicture}
\end{sourceblock}
我们现在观察到流场有效地分为三个域：解析场 ̃ u i 、解析子滤波器尺度 (RSFS) 场 ( ̃ u i − ̃ u i ) 和次网格尺度 (SGS) 场 ( u i − ̃ u i )。解析的字段是由代码计算的； RSFS 场存在于网格上，因此实际上可以根据已解析的数据将其重建为某种程度的近似值，正如 Stolz 等人所做的那样。 （2001）和Chow等人。 （2005）；并且必须对 SGS 场效应进行建模。我们可以看到这是如何在 SFS 应力项中发挥作用的，一旦定义它们，它们就会成为 LES 中使用的“模型”。

已解决的子滤波器规模 (RSFS) 应力
使用 Stolz 等人的反卷积方法，无需建模即可实际计算 RSFS 应力。 (2001)，它从解析场近似重建原始速度场（另见 Carati 等人 2001；Chow 等人 2005）。重建使用 van Cittert 迭代级数展开方法，并根据滤波速度 ( ̃ u i ) 提供未滤波速度 ( ̃ u i ) 的估计：

\begin{sourceblock}
\includegraphics[page=388,trim={53.00bp 265.14bp 53.87bp 382.63bp},clip,width=0.92\textwidth,max height=0.38\textheight]{Ferziger4th.pdf}
\end{sourceblock}
其中 I 是恒等运算符，G 是等式中的滤波器。 （10.7），*表示卷积。然后，我们观察到近似速度 ̃ u ⋆
i 可以插入 RSFS 应力的第一项，从而无需模型即可计算应力，即：

\begin{sourceblock}
\begin{tikzpicture}
\node[inner sep=0,anchor=south west] (image) at (0,0) {\includegraphics[page=388,trim={53.00bp 191.35bp 53.87bp 456.71bp},clip,width=0.92\textwidth,max height=0.38\textheight]{Ferziger4th.pdf}};
\begin{scope}[x={(image.south east)},y={(image.north west)}]
\fill[white] (0.21164,0.15985) rectangle (0.27838,0.86283);
\fill[white] (0.28337,0.22345) rectangle (0.31155,0.86283);
\fill[white] (0.31507,0.22345) rectangle (0.33621,0.86283);
\fill[white] (0.35453,0.16372) rectangle (0.38162,0.86283);
\fill[white] (0.38030,0.20852) rectangle (0.40481,0.86283);
\fill[white] (0.40030,0.16372) rectangle (0.41095,0.63772);
\fill[white] (0.41636,0.16372) rectangle (0.47347,0.86283);
\fill[white] (0.47218,0.20852) rectangle (0.49197,0.86283);
\fill[white] (0.49218,0.16372) rectangle (0.50283,0.63772);
\end{scope}
\end{tikzpicture}
\end{sourceblock}
\begin{sourceblock}
\begin{tikzpicture}
\node[inner sep=0,anchor=south west] (image) at (0,0) {\includegraphics[page=388,trim={53.00bp 153.37bp 53.87bp 494.69bp},clip,width=0.92\textwidth,max height=0.38\textheight]{Ferziger4th.pdf}};
\begin{scope}[x={(image.south east)},y={(image.north west)}]
\fill[white] (0.00000,0.86449) rectangle (0.05074,1.00000);
\fill[white] (0.05233,0.86449) rectangle (0.17814,1.00000);
\fill[white] (0.17982,0.86449) rectangle (0.28066,1.00000);
\fill[white] (0.28229,0.86449) rectangle (0.32869,1.00000);
\fill[white] (0.33029,0.86449) rectangle (0.38337,1.00000);
\fill[white] (0.38493,0.86449) rectangle (0.41805,1.00000);
\fill[white] (0.41961,0.86449) rectangle (0.55426,1.00000);
\fill[white] (0.55585,0.86449) rectangle (0.60728,1.00000);
\fill[white] (0.60887,0.86449) rectangle (0.65032,1.00000);
\fill[white] (0.65191,0.86449) rectangle (0.71829,1.00000);
\fill[white] (0.71988,0.86449) rectangle (0.79465,1.00000);
\fill[white] (0.79624,0.86449) rectangle (0.86346,1.00000);
\fill[white] (0.86505,0.86449) rectangle (0.90752,1.00000);
\fill[white] (0.90914,0.86449) rectangle (0.93062,1.00000);
\fill[white] (0.93218,0.86449) rectangle (0.95699,1.00000);
\fill[white] (0.95856,0.86449) rectangle (1.00000,1.00000);
\fill[white] (0.00000,0.20299) rectangle (0.10728,0.84292);
\fill[white] (0.10917,0.20299) rectangle (0.13729,0.84292);
\fill[white] (0.13913,0.20299) rectangle (0.17389,0.84292);
\fill[white] (0.17579,0.20299) rectangle (0.21558,0.84292);
\fill[white] (0.21741,0.20299) rectangle (0.27380,0.84292);
\fill[white] (0.27570,0.20299) rectangle (0.32544,0.84292);
\fill[white] (0.32734,0.20299) rectangle (0.36875,0.84292);
\fill[white] (0.37062,0.20299) rectangle (0.44535,0.84292);
\fill[white] (0.44728,0.20299) rectangle (0.50698,0.84292);
\fill[white] (0.50884,0.22345) rectangle (0.52998,0.86283);
\fill[white] (0.52698,0.20852) rectangle (0.55823,0.93363);
\fill[white] (0.57495,0.08905) rectangle (0.58559,0.56305);
\fill[white] (0.59104,0.20299) rectangle (0.65411,0.84292);
\fill[white] (0.65597,0.20299) rectangle (0.75561,0.84292);
\fill[white] (0.75753,0.20299) rectangle (0.79062,0.84292);
\fill[white] (0.79248,0.20299) rectangle (0.90998,0.84292);
\fill[white] (0.91188,0.20299) rectangle (1.00000,0.84292);
\fill[white] (0.00000,0.00000) rectangle (0.02746,0.18142);
\fill[white] (0.02776,0.00000) rectangle (0.06917,0.18142);
\fill[white] (0.06950,0.00000) rectangle (0.20571,0.18142);
\fill[white] (0.20611,0.00000) rectangle (0.29170,0.18142);
\fill[white] (0.29209,0.00000) rectangle (0.31853,0.18142);
\fill[white] (0.31886,0.00000) rectangle (0.35278,0.18142);
\fill[white] (0.35314,0.00000) rectangle (0.42614,0.18142);
\fill[white] (0.42653,0.00000) rectangle (0.55278,0.18142);
\fill[white] (0.55314,0.00000) rectangle (0.60120,0.18142);
\fill[white] (0.60156,0.00000) rectangle (0.67717,0.18142);
\fill[white] (0.67753,0.00000) rectangle (0.75970,0.18142);
\fill[white] (0.76006,0.00000) rectangle (0.81143,0.18142);
\fill[white] (0.81179,0.00000) rectangle (0.86650,0.18142);
\fill[white] (0.86683,0.00000) rectangle (0.91323,0.18142);
\fill[white] (0.91356,0.00000) rectangle (1.00000,0.18142);
\end{scope}
\end{tikzpicture}
\end{sourceblock}
j 实际上必须在计算中显式过滤（Chow 等人，2005 年；Gullbrand 和 Chow，2003 年）。用户可以选择重建的级别。 0 级使用一项并产生基本上相当于 Bardina 等人的尺度相似性模型的近似值。 （1980）； 1 级使用两个术语，对结果有重大影响，并且可能是一个有用的折衷方案

\footnotetext[10]{请注意 Carati 等人。 （2001）和Chow等人。 (2005) 以及许多其他人使用应力定义作为速度乘积，而不包括密度。}
混合模型中的成本和物理之间的关系。 11 斯托尔兹等人。 (2001) 使用此重建作为完整模型，实际上忽略了应力的 SGS 部分；然而，重建项需要通过添加到滤波纳维-斯托克斯方程右侧的能量消耗项来补充。重建项允许能量从子滤波器反向散射到解析尺度（参见 Chow 等人，2005）。

次网格尺度 (SGS) 应力
SGS压力

\begin{sourceblock}
\begin{tikzpicture}
\node[inner sep=0,anchor=south west] (image) at (0,0) {\includegraphics[page=389,trim={53.00bp 462.38bp 53.87bp 188.30bp},clip,width=0.92\textwidth,max height=0.38\textheight]{Ferziger4th.pdf}};
\begin{scope}[x={(image.south east)},y={(image.north west)}]
\fill[white] (0.00000,0.92367) rectangle (0.05074,1.00000);
\fill[white] (0.05344,0.92367) rectangle (0.11317,1.00000);
\fill[white] (0.11588,0.92367) rectangle (0.18728,1.00000);
\fill[white] (0.35017,0.09961) rectangle (0.40638,0.92238);
\fill[white] (0.41140,0.17464) rectangle (0.43958,0.92238);
\fill[white] (0.44310,0.17464) rectangle (0.46424,0.92238);
\end{scope}
\end{tikzpicture}
\end{sourceblock}
包含精确（未过滤）速度的相关性，因此必须建模。可以使用任何 SGS 模型，包括上面提到的模型和下面描述的代数应力模型。

卡拉蒂等人。 (2001) 方法导致总 SFS 应力的混合表达式。周等人。 ( 2005 ) 通过使用重建加上动态 SGS 模型（Wong 和 Lilly 1994，动态模型）以及近壁模型（其作用类似于第 10.3.3.2 节中描述的 van Driest 模型），创建了动态重建模型（DRM）；路德维希等人。 (2009)将该模型的性能与其他模型进行比较，发现 DRM 优于 Smagorinsky 和 ​​TKE 模型。 Zhou和Chow（2011）将DRM策略应用于稳定大气边界层。

另一方面，人们可以忽略 RSFS 或 SGS 术语并创建功能模型，例如 Stolz 等人。 (2001) 仅使用 RSFS 项和能量汇、单独的动态 Smagorinsky 或 Enriquez 等人的代数应力模型来近似反卷积。 (2010) 接下来描述。独立 SGS 模型的基本原理是假设过滤器宽度和网格大小 h 相等，因此 RSFS 区域消失。这并不严格正确，因为显式和（隐式）网格过滤器不同，但它已经足够了。一般来说，对于混合模型，遵循 Chow 和 Moin (2003) 关于误差和 SGS 力大小的建议，设置 2 ≤ / h ≤ 4。

10.3.3.5 显式代数雷诺应力模型 (EARSM)

上述 SGS 模型的问题之一是它们不允许在实验和现场研究中在地面/墙壁附近观察到的法向应力各向异性（参见 Sullivan 等人，2003）。替代方案包括：

\footnotetext[11]{因此，n 层有 n + 1 项展开式 ( 10.24 )；参见 Shi 等人的附录。 (2018b)进行讨论和应用。}
\begin{itemize}
\item 非线性SGS 模型，例如Kosovi ́c (1997) 的模型，它允许能量反向散射和正常SGS 应力各向异性（参见重建模型，它也允许反向散射并且在RSFS 处是各向异性的）。
\item 导出SGS 应力的输运方程，然后保留产生Smagorinsky 模型的附加项。 Wyngaard (2004) 以及 Hatlee 和 Wyngaard (2007) 展示了这些方程，并提出了具有改进性能的新的各向异性 SGS 模型。 Ramachandran 和 Wyngaard (2010) 扩展了这种方法，求解了 SGS 应力的偏微分输运方程的截断版本。他们针对中度对流大气边界层进行了此操作，这与 Sullivan 等人报告的 HATS 实验相匹配。 （2003）。
\item 从SGS 应力传递方程生成代数应力模型，以适中的成本捕获新的物理现象。我们简要评论两个代数应力模型。
\end{itemize}
一般来说，显式代数应力方法旨在生成不需要求解微分方程且无需涡粘性的模型。 Rodi (1976) 为 URANS 创建了这样一个模型，Wallin 和 Johansson (2000) 对此进行了扩展（参见第 10.3.5.4 节）。马斯托普等人。 (2009) 为旋转河道流的 LES 建立了 EARSM。他们沿着与 Wallin 和 Johansson 相同的思路，对 SGS 应力各向异性张量的输运方程中的项进行截断或建模；他们获得了具有两个模型参数（即 SGS 动能和 SGS 时间尺度）的 SGS 应力的显式代数模型。他们提供了动态和非动态程序来评估它们。他们的模拟结果与 DNS 数据非常吻合，并且由于避免使用涡流粘度，解析的雷诺应力各向异性和 SGS 应力各向异性都得到了改善。

最近，Rasam 等人。 (2013) 结合了 Wallin 和 Johansson (2000) 的显式代数标量通量模型和 Marstorp 等人的 LES 的 EARSM。 (2009)，与过滤的 DNS 数据相比，产生了一个能够产生 SGS 标量通量各向异性并相当好地预测标量剖面的模型。

Findikakis 和 Street (1979) 基于 Launder 等人的 RANS 湍流模型的思想，描述了 LES 中 SGS 项的 EARSM。 （1975年）；见节。见下文 10.3.5.4。在 Carati 等人的背景下工作。 ( 2001 ) 分解（如上所述），Enriquez 等人。 ( 2010 ) 建立在 Findikakis 和 Street 的工作以及 Chow 等人的工作之上。 (2005) 创建一个线性代数亚网格尺度应力模型，并结合 RSFS 应力重建，应用于中性大气边界层。 SGS 应力的传输方程已简化为仅表示这些方程中的产生、耗散和压力重新分布项，从而产生一组六个线性代数方程，这些方程在每个时间步长的每个网格点处反转。除了 Launder 等人中隐含的常数之外，没有任何参数需要确定。 （1975）已被采用的模型。该模型嵌入到 ARPS 中尺度 LES 代码中（Xue 等人，2000 年；Chow 等人，2005 年）。 SGS 雷诺应力方程忽略对流、扩散和粘性，但在法向应力方程中对耗散进行建模，并通过求解偏微分输运方程获得湍流动能。报告的结果表明，在不进行重建的情况下使用时，该 EARSM 产生的结果优于 Smagorinsky 模拟并相当于动态 Wong-Lilly 结果，并且它再现了观察到的法向应力各向异性。 Enriquez (2013) 中包含了对中性、稳定和对流大气边界层的应用。

10.3.3.6 LES 的边界条件
LES 使用的边界条件和数值方法与 DNS 中使用的非常相似。最重要的区别是，当 LES 应用于复杂几何形状的流动时，一些数值方法（例如谱方法）变得难以应用。在这些情况下，人们被迫使用有限差分、有限体积或有限元方法。原则上，可以使用本书前面描述的任何方法，但重要的是要记住，挑战网格分辨率的结构几乎可能存在于流程中的任何地方。因此，采用尽可能高精度的方法非常重要。

正如我们在该节中指出的那样。 10.2 和 10.3.5.5 中，壁附近的流动行为可能需要对边界条件进行特殊处理或在近壁区域使用壁函数。 Sagaut (2006) 提出了对替代方案的讨论，Rodi 等人也是如此。 （2013），其中还包括粗糙墙壁的替代品。 Stoll 和 Porté-Agel (2006) 描述了许多粗糙壁模型，但在这里我们描述了最简单的替代模型之一，因为它几乎适用于所有情况。该模型源于大气科学和海洋学，其中固体边界通常是粗糙的并且无滑移边界条件不合适，因为粗糙度无法准确表示和/或粗糙度元素穿透平均粘性子层并且强迫包括压力和粘性效应。然后，应用自由滑移条件，其中速度不受约束，并且在边界处建立速度和动量通量（壁应力）之间的关系。这反映了通常对具有壁模型的 RANS 流所做的操作（参见第 10.3.5.5 节），即假设边界附近存在对数速度剖面（或者如果流不是中性稳定的，则可以使用 Monin-Obukhov 相似剖面（Porté-Agel 等人 2000；Zhou 和 Chow 2011））。因此，解决需要非常精细的垂直于壁的网格的流动问题，12 使用强制这种对数行为的边界条件（Rodi 等人，2013 年；Porté-Agel 等人，2000 年；Chow 等人，2005 年），并且壁应力由下式给出

\begin{sourceblock}
\begin{tikzpicture}
\node[inner sep=0,anchor=south west] (image) at (0,0) {\includegraphics[page=391,trim={53.00bp 134.02bp 53.87bp 514.36bp},clip,width=0.92\textwidth,max height=0.38\textheight]{Ferziger4th.pdf}};
\begin{scope}[x={(image.south east)},y={(image.north west)}]
\fill[white] (0.06198,0.15878) rectangle (0.15122,0.87500);
\fill[white] (0.15642,0.22410) rectangle (0.18460,0.87500);
\fill[white] (0.18812,0.15878) rectangle (0.27792,0.87500);
\end{scope}
\end{tikzpicture}
\end{sourceblock}
其中 C D 是阻力系数，u 1 和 v 1 是速度剖面上离壁第一个网格点处的流向和横向速度分量

\footnotetext[12]{注意：DNS 还需要沿着墙壁的精细网格，但 LES 通常在墙壁或地面附近具有相当高的（水平与正常）网格纵横比。}
\begin{sourceblock}
\begin{tikzpicture}
\node[inner sep=0,anchor=south west] (image) at (0,0) {\includegraphics[page=392,trim={53.00bp 558.46bp 53.87bp 53.54bp},clip,width=0.92\textwidth,max height=0.38\textheight]{Ferziger4th.pdf}};
\begin{scope}[x={(image.south east)},y={(image.north west)}]
\fill[white] (0.72860,0.66568) rectangle (0.74250,0.82397);
\fill[white] (0.74890,0.69745) rectangle (0.79699,0.91097);
\fill[white] (0.79937,0.69745) rectangle (0.84078,0.91097);
\fill[white] (0.84313,0.69745) rectangle (0.93898,0.91097);
\fill[white] (0.94135,0.69745) rectangle (0.98111,0.91097);
\fill[white] (0.98349,0.69745) rectangle (1.00000,0.91097);
\fill[white] (0.00000,0.47654) rectangle (0.07405,0.69006);
\fill[white] (0.07672,0.47654) rectangle (0.13281,0.69006);
\fill[white] (0.13552,0.47654) rectangle (0.16030,0.69006);
\end{scope}
\end{tikzpicture}
\end{sourceblock}
\begin{sourceblock}
\begin{tikzpicture}
\node[inner sep=0,anchor=south west] (image) at (0,0) {\includegraphics[page=392,trim={53.00bp 512.92bp 53.87bp 117.05bp},clip,width=0.92\textwidth,max height=0.38\textheight]{Ferziger4th.pdf}};
\begin{scope}[x={(image.south east)},y={(image.north west)}]
\fill[white] (0.00000,0.73044) rectangle (0.02081,1.00000);
\fill[white] (0.02346,0.73044) rectangle (0.11741,1.00000);
\fill[white] (0.12009,0.73044) rectangle (0.16986,1.00000);
\end{scope}
\end{tikzpicture}
\end{sourceblock}
其中 z 0 是粗糙度长度，z 1 是到离墙第一个点的距离；参见等式。 （10.63）。

在 LES 中，可以将 RANS 建模中使用的那种壁函数（参见第 10.3.5 节）用于平滑壁，以避免原本需要的精细网格。 Piomelli 和 Balaras (2002) 以及 Piomelli (2008) 对一系列墙层模型进行了评论，后者得出的结论是“……没有任何一种方法明显优于其他方法。”

10.3.3.7 截断误差和计算混合
除其他因素外，数值大涡模拟的准确性可能会受到与子网格子滤波器尺度强迫相比的数值截断误差的大小、SFS/SGS 模型的形式、数值算法和计算混合的影响。最后一个是地球物理流模拟的典型特征，其中使用高阶（例如，第五和第六）平流方案；在这种情况下，Xue（2000）指出“大多数数值模型采用数值扩散或计算混合来控制小规模（波长中接近两个网格间隔）噪声，这些噪声可能由数值色散、非线性不稳定性、不连续物理过程和外力引起。”

以下部分提出了一些有关错误、混合效果和模拟质量的指南。

误差和精度：Chow 和 Moin (2003) 对数值误差和 SFS 强迫之间的平衡进行了研究。 Chow 和 Moin (2003) 使用稳定分层剪切流的直接数值模拟 (DNS) 数据集进行先验测试，将几种有限差分方案的数值误差与 SFS 力的大小进行了比较。为了确保 SFS 项大于非线性对流项数值格式中包含的数值误差，他们发现：

1. 对于二阶有限差分格式，滤波器尺寸至少应为网格间距的四倍； 2.对于六阶Padé方案，滤波器尺寸只需至少是网格间距的两倍。

\footnotetext[13]{对于受科里奥利影响的流动，速度随着距边界的距离而旋转，即使流动是远离边界的单向流动，两个分量也可能不为零；见第 10.3.4.3 节.}
然而，Celik 等人。 (2009) 认为“这在工程 LES 应用中几乎不可能实现。”不满足上述标准的结果是，除了模拟的 SGS 应力之外，LES 还包含来自数值误差的未知（但可能很重要）的贡献。

混合效应：计算混合项通常是应用于动量和标量方程的四阶或高阶超粘性项。对于商业或专业开发的规范，指南建议系数值，以便给出稳定性所需的最小混合量。然而，这仍然可能对结果产生重大影响。布莱恩等人。 (2003) 浮力对流测试用例（参见其附录），具有六阶对流和六阶显式滤波器，表明对于能谱，只有波长大于网格间距六倍的信息代表物理解，并且可以通过改变计算混合系数来为小于网格大小六倍的尺度创建任意光谱斜率。 Michioka 和 Chow (2008) 使用中尺度代码在复杂地形上进行了标量传输的高分辨率 LES，并证明：

1.大的计算混合系数可能不会显著改变模拟的中尺度流场（因为计算混合仅在网格尺度噪声附近阻尼）。 2. 计算混合可以显著影响高分辨率模拟中的湍流波动和近地速度剖面，因为小尺度运动是已解析湍流流场的一部分。 3. 使用最小的可行计算混合系数对于实现标量最大地面浓度的模拟和测量数据之间的良好匹配至关重要。

模拟质量：LES 的验证和/或验证很困难，因为例如，随着网格尺寸的减小，数值截断误差和 SGS 或 SFS 模型误差都会减小，并且流量解析的分数也会发生变化。因此，从某种意义上说，LES 向 DNS 靠拢，正如 Celik 等人所言。 ( 2009 ) 指出“不存在独立于电网的 LES 之类的东西。”当然，使用一系列较小的网格尺寸来测试收敛性是合理的（这意味着从一个网格到另一个网格的表示流变化对于所需目的来说足够小）。 Sullivan 和 Patton（2011）对网格尺寸减小的收敛性进行了详细评估，指出某些性能指标的收敛速度比其他指标慢，并且随着网格尺寸的减小，界面和其他锐梯度区域的物理表示可能会发生变化。

迈耶斯等人。 (2007) 对均匀各向同性流进行了误差评估，以演示一种系统地改变 Smagorinsky 系数并使用不同离散化的方法，从而尝试隔离数值和模型误差影响。在一系列先前的论文之后，Celik 等人。 (2009) 提供了工程 LES 应用的评估策略。他们也建议系统化的网格和模型变化，这需要多次代码运行，并应用于许多案例，提供图形证据。然而，沙利文

\begin{sourceblock}
\includegraphics[page=394,trim={53.00bp 506.25bp 53.87bp 52.00bp},clip,width=\textwidth,max height=0.38\textheight]{Ferziger4th.pdf}
\par\vspace{0.45\baselineskip}
\small 图 10.7 安装在通道壁上的立方体上流动的解域；来自 Shah 和 Ferziger (1997)
\end{sourceblock}
和 Patton (2011) 警告说，而 Celik 等人。 (2009)经常依赖DNS作为测试基础，它不适用于高雷诺数行星边界层。

\subsection*{10.3.4 LES应用举例}\phantomsection\addcontentsline{toc}{subsection}{10.3.4 LES应用举例}
作为该方法的示例，我们将使用安装在通道一面墙上的立方体上的流动。几何形状如图10.7所示。对于由 Shah 和 Ferziger (1997) 进行的所示模拟，基于流入处最大速度和立方体高度的雷诺数为 3200。流入是完全发展的通道流，取自该流的单独模拟，出口条件是方程 1 给出的对流条件。 10.4.在翼展方向使用周期性边界条件，在所有壁表面使用无滑移条件。

LES 使用 240 × 128 × 128 控制体积的网格，具有二阶精度。时间推进方法是分数步类型的。对流项采用三阶龙格-库塔时间方法显式处理，而粘性项则采用隐式处理。特别是，后者使用的方法是克兰克-尼科尔森方法的近似因式分解。通过多重网格法求解泊松方程获得压力。

\begin{center}\small 图10.8给出了靠近壁面区域的时间平均流动的流线图；从该图中可以看出有关流量的大量信息。传入的流不会发生传统意义上的分离，而是到达停滞点或鞍点（图中用 A 标记）并绕过身体。下壁上方的一些气流会撞击立方体的正面；大约一半的流动向下流动，并在身体前面形成逆流区域。当流动沿着立方体正面靠近下壁时，在立方体前方存在二次分离和重新附着线（图中用 B 标记）。在每个\end{center}
\begin{sourceblock}
\includegraphics[page=395,trim={53.00bp 333.25bp 53.87bp 52.00bp},clip,width=\textwidth,max height=0.38\textheight]{Ferziger4th.pdf}
\par\vspace{0.45\baselineskip}
\small 图10.8 壁挂式立方体上流动靠近下壁区域的流线；来自 Shah 和 Ferziger (1997)
\end{sourceblock}
在立方体的一侧，可以找到一个会聚流线的区域（标记为 C）和另一个发散流线的区域（标记为 D）；这些是马蹄涡的痕迹（下面会详细介绍）。在身体的后面，人们会发现两个旋流区域（标记为E），它们是拱形涡流的足迹。最后，在主体的更下游有一条重新连接线（标记为 H）。

\begin{center}\small 图 10.9 显示了流动中心平面上的时间平均流动的流线。上述许多特征都清晰可见，包括上游角部的分离区（F），也是马蹄涡的头部，拱形涡的头部（G），重新附着线（H），以及在上表面上不重新附着的主体上方的再循环区（I）。\end{center}
最后，图 10.10 给出了时间平均流动的流线在平行于立方体背面（位于物体下游）的平面上的投影。马蹄涡 (J) 与较小的角涡一样清晰可见。

值得注意的是，瞬时流量看起来与时间平均流动有很大不同。例如，拱形涡旋并不是瞬时存在的；流动中存在涡流，但它们在立方体的两侧几乎总是不对称的。事实上，图 10.8 的近对称表明平均时间（几乎）足够长。

\begin{sourceblock}
\includegraphics[page=396,trim={53.00bp 484.26bp 53.87bp 52.00bp},clip,width=\textwidth,max height=0.38\textheight]{Ferziger4th.pdf}
\par\vspace{0.45\baselineskip}
\small 图10.9 壁挂式立方体上流动的垂直中心平面的流线；来自 Shah 和 Ferziger (1997)
\end{sourceblock}
\begin{sourceblock}
\includegraphics[page=396,trim={53.00bp 305.87bp 53.87bp 205.40bp},clip,width=\textwidth,max height=0.38\textheight]{Ferziger4th.pdf}
\par\vspace{0.45\baselineskip}
\small 图10.10 壁挂式立方体上的流动流线在平行于背面的平面上的投影，立方体后面0.1步高；来自 Shah 和 Ferziger (1997)
\end{sourceblock}
从这些结果可以清楚地看出，LES（或用于更简单流的 DNS）提供了大量有关流的信息。与第 1 节中描述的模拟类型相比，这种模拟的性能与进行实验有更多的共同点。 10.3.5 以及从中获得的定性信息非常有价值。

10.3.4.2 Re = 50,000 时围绕球体的流动
见第 10.2.4.2 节 我们简要讨论了 Re = 5,000 时光滑球体周围流的 DNS；对于 Re = 50,000 时的流量，需要非常精细的网格才能执行适当的 DNS。合乎逻辑的步骤是切换到 LES，它可以解决大尺度问题并对次网格尺度湍流进行建模。即使对于 LES，我们也必须在较低的雷诺数下制作比用于 DNS 的网格更精细的网格：在由笛卡尔单元组成的非结构化网格中使用了大约 4000 万个单元，并在尾迹中进行了局部细化，并通过沿壁的棱镜层进行了修剪。图10.11

\begin{sourceblock}
\includegraphics[page=397,trim={53.00bp 455.48bp 53.87bp 52.00bp},clip,width=\textwidth,max height=0.38\textheight]{Ferziger4th.pdf}
\par\vspace{0.45\baselineskip}
\small 图 10.11 光滑球体 y = 0 处解域的纵向截面，显示球体周围的网格结构（在最细的网格中，所有方向上的间距小两倍）
\end{sourceblock}
显示球体尾流中的网格结构；因此可以看到网格结构，该网格比用于产生下面所示结果的网格粗两倍。

Baki ́c (2002) 对 Re = 50,000 时球体周围的流动进行了实验研究；他使用直径 D = 61 的台球。 4 mm，用直径 d = 8 mm 的棍子固定在横截面为 300 × 300 mm 的矩形小风洞中。他还对一个球体进行了测量，该球体在距前驻点 75° 的位置处沿圆周连接了绊线；绊线的直径为0.5毫米。

模拟中的解域与实验几何形状相匹配，入口边界位于球体中心上游 300 mm 处，出口边界位于球体中心下游 300 mm 处。入口处指定的均匀速度为 12.43 m/s，流体为密度 ρ = 1 的空气。 204 kg/m 3 且粘度μ = 1.837×10−5帕·秒。对光滑球体和带有绊线的球体进行了模拟。光滑球体尾迹的网格间距（见图10.11）在所有三个方向上均为0.265 mm，约为D/232；该间距保持在大约 1 的距离。后驻点下游7 D处，最细网格区直径接近1 . 5D。网格从图10.11所示的最细区域分几步变粗；没有尝试解决沿风洞壁生长的边界层，但沿球体和固定棒表面的棱镜层开始时紧邻壁单元的厚度为 0.03 毫米 (D / 2407)。

对于有绊线的情况，必须改变网格设计。为了捕获绊线表面的层流边界层分离，区域中的单元尺寸延伸约。上游及以上的半个绊线直径和下游的 4 个直径为 0.041667 毫米 (D /1474)。两倍间距的区域沿着球体表面一直延伸，直到球体下游侧发生边界层分离；见图10.12，它显示了详细信息

\begin{sourceblock}
\includegraphics[page=398,trim={53.00bp 400.48bp 53.87bp 52.00bp},clip,width=\textwidth,max height=0.38\textheight]{Ferziger4th.pdf}
\par\vspace{0.45\baselineskip}
\small 图 10.12 带有绊线的球体在 y = 0 处的解域的纵向剖面，显示了球体和绊线周围的网格结构
\end{sourceblock}
绊线周围和直接球体尾流中的实际网格。在球体尾流的最大部分，网格尺寸为0.3333 mm（D/184）。靠近壁的第一棱镜层的厚度为0.02毫米(D/3070)。在有绊线的情况下，尾流要窄得多，因此与光滑球体的情况相比，球体后面的细网格区域相应地变小。

两种情况下的时间步长均为 10 μs，根据平滑球体最细区域的平均速度和网格间距，平均库朗数为 0.5。绊线周围最细的网格中的库朗数更大。在空间和时间上均使用二阶方案（对流和扩散的中心差分以及时间上的二次向后插值）。商业代码 STAR-CCM+ 是求解器（时间上完全隐式，即在新的时间级别计算对流通量、扩散通量和源项；SIMPLE 算法用于压力-速度耦合；假设流动不可压缩）。欠松弛参数的速度为 0.95，压力为 0.75，每个时间步执行 5 次外部迭代以更新非线性项。还在一个较粗糙的网格和三种不同的子网格尺度模型上进行了模拟；这是一项正在进行的研究的一部分，我们不会在这里讨论所有细节。一个重要的信息是，通过使用局部细化的非结构化网格（也可以是多面体），我们可以解决湍流尾流，同时不会浪费流动不湍流且变量在空间中变化不那么快的区域中的单元。对图 10.13 和 10.14 中涡度结构的目视检查表明，所使用的网格对于 LES 来说可能具有足够的分辨率，因为这些结构比细网格间距大得多。

\begin{sourceblock}
\includegraphics[page=399,trim={53.00bp 280.48bp 53.87bp 52.00bp},clip,width=\textwidth,max height=0.38\textheight]{Ferziger4th.pdf}
\par\vspace{0.45\baselineskip}
\small 图 10.13 光滑球体（上）和带有绊线的球体（下）在 y = 0 平面内涡量分量 ω y 的瞬时云图
\end{sourceblock}
使用 RANS 模型无法很好地预测这种亚临界流动；我们不会显示此类计算的任何结果，而只是声明所有测试的模型都显著低估了阻力并高估了球体后面再循环区的长度（显著意味着 25\% 或更多）。然而，使用动态 Smagorinsky 亚网格尺度模型的 LES 预测光滑球体的平均阻力系数约为 0.48，这与文献中发现的实验数据接近。绊线预计会延迟分离并减小球体后面再循环区的大小，从而导致阻力显著降低；情况确实如此。如图 10.15 所示，预测阻力在接近 0.175 的值附近波动，几乎比光滑球体低 3 倍。这证明了 CFD 的强大功能：当使用适当的网格和湍流建模方法时，可以预测微小的几何形状变化对流动的影响。

\begin{sourceblock}
\includegraphics[page=400,trim={53.00bp 427.48bp 53.87bp 52.00bp},clip,width=\textwidth,max height=0.38\textheight]{Ferziger4th.pdf}
\par\vspace{0.45\baselineskip}
\small 图 10.14 光滑球体（左）和带有绊线的球体（右）在 x / D = 1 平面中涡度分量 ω x 的瞬时云图
\end{sourceblock}
\begin{sourceblock}
\includegraphics[page=400,trim={53.00bp 160.86bp 53.87bp 262.41bp},clip,width=\textwidth,max height=0.38\textheight]{Ferziger4th.pdf}
\par\vspace{0.45\baselineskip}
\small 图 10.15 在最后 65,000 个时间步长期间，光滑球体（上）和带有绊线的球体（下）的阻力系数变化
\end{sourceblock}
\begin{sourceblock}
\includegraphics[page=401,trim={53.00bp 279.25bp 53.87bp 52.00bp},clip,width=\textwidth,max height=0.38\textheight]{Ferziger4th.pdf}
\par\vspace{0.45\baselineskip}
\small 图 10.16 光滑球体（上）和带绊线球体（下）的瞬时流型
\end{sourceblock}
LES 提供对流动行为的洞察，这对于可能需要改进设计的工程师来说非常重要。使用 LES 中按指定时间增量创建的图片，可以轻松执行流动画；我们无法在本书中重现动画，但我们在图 1-2 中显示。 10.13、10.14 和 10.16 具有光滑表面和表面安装绊线的球体的瞬时流动模式，揭示了流动分离和尾流行为的细节。这些图清楚地显示了绊线对流动的影响：与绊线分离后重新附着后，边界层变得湍流，导致主分离延迟和尾流窄得多。每个工程师都会立即知道，流动行为的这种变化将导致阻力的显著减少。

工程师通常对其波动的平均量和均方根值（均方根）感兴趣。对于诸如身体阻力之类的积分量，可以轻松获得此信息，例如，通过处理如图 1 所示的信号。

10.15。获得平均速度和压力分布并不那么简单。在渠道或管道流的 LES 中，通常在时间和一个空间方向（渠道的展向和管道的圆周方向）上进行平均，这使得人们可以用较少数量的样本获得稳定的值。在复杂的几何形状中，空间平均通常是不可能的；在本例中，由于球体安装在具有矩形横截面的风洞中，因此流动在整个横截面中肯定不是轴对称的（尽管在球体和固定球体的杆附近可能是轴对称的）。这里遇到三个问题：

\begin{itemize}
\item 如果没有空间平均，获得稳定的时间平均流动所需的样本数量会变得非常大。在当前的模拟中，对最近 65,000 个（光滑球体）或 70,000 个样本（带绊线的球体）进行平均，但时间平均流动不是轴对称的；需要更多的样本才能实现平均速度分布的完美对称。在实验中，人们通常假设对称并沿着从对称轴向外半径的一条线测量轮廓。
\item 当使用非结构化网格时，空间平均不容易执行。在本例中，单元格是笛卡尔坐标系，大小各不相同；为了在圆周方向上平均解，我们必须创建一个结构化的极柱网格，将时间平均解插值到该网格上，根据笛卡尔速度分量计算结构化网格上的轴向、径向和圆周分量，然后在圆周方向上对它们进行平均。
\item 即使几何形状完全轴对称，平均流动轴对称的假设也可能是错误的。事实上，有证据表明——无论是模拟（例如，Constantinescu 和 Squires 2004 年）还是实验（例如，Taneda 1978 年）——球体的尾流往往会相对于流动方向倾斜。通过圆周平均来强制轴对称会伪造结果。
\end{itemize}
我们在图 10.17 中显示了两个球体的时间平均流动模式。虽然光滑球体的平均流场看起来几乎对称，但对于带有绊线的球体，流型仅在平面 z = 0 中接近对称；在平面 y = 0 中，流动高度不对称。正如我们将在该节中展示的那样。 10.3.6，一些RANS模型还预测光滑球体周围的不对称稳态流动，因此这似乎是流动的特征而不是模拟的缺陷。在这两种情况下，再循环区的长度与 Baki ́c 的实验非常吻合，他发现再循环区在 x / D = 1 处结束。 43 表示光滑球体，x/D = 1 表示带有绊线的球体。

我们再次注意到，时间平均流只是通过对不同时间的流实现进行平均而获得的构造；它在现实中从来不存在，哪怕是一瞬间也不存在。虽然它可以通过测量数据和模拟数据的平均来获得，但它在自然界中是无法观察到的。实验中长时间曝光会产生不同的图像，因为平均值周围的扩散不会像数学平均中那样抵消正值和负值。有趣的是，图 10.14 所示的 x / D = 1 平面上的涡量分量 ω x 的瞬时图对于带有绊线的球体来说几乎是对称的，而

\begin{sourceblock}
\includegraphics[page=403,trim={53.00bp 103.48bp 53.87bp 52.00bp},clip,width=\textwidth,max height=0.38\textheight]{Ferziger4th.pdf}
\par\vspace{0.45\baselineskip}
\small 图 10.17 平面 y = 0 中的光滑球体（上）和平面 y = 0（中）和 z = 0（下）中带有绊线的球体的时间平均流动模式
\end{sourceblock}
\begin{sourceblock}
\includegraphics[page=404,trim={53.00bp 371.25bp 53.87bp 52.00bp},clip,width=\textwidth,max height=0.38\textheight]{Ferziger4th.pdf}
\par\vspace{0.45\baselineskip}
\small 图 10.18 x / D = 1 处沿 4 条径向线（正 y 、负 y 、正 z 、负 z ）的时间平均轴向速度分量（上）及其方差（下）的轮廓与 Baki ́c (2002) 的光滑球体（左）和带绊线球体（右）的实验数据的比较
\end{sourceblock}
光滑球体的图像高度不对称——与平均流的结果相反。

预测平均速度和雷诺应力与 Baki ́c (2002) 的实验数据的比较也显示出相当好的一致性。我们在图 10.18 中仅显示了轴向速度分量 U x 及其方差 ( u ′
x ) 2 at x / D = 1。沿 4 条线绘制时间平均值，从杆轴开始，沿正 y 、负 y 、正 z 和负 z 方向延伸。

如果流动是轴对称的并且平均时间足够长，则所有这 4 个轮廓都会崩溃；这里的情况并非如此。光滑球体的平均速度分布相对接近，但方差表明所有 4 个分布在峰值附近存在显著变化；最高值与最低值相差30\%左右。这种差异可能会随着平均时间的延长（更多样本）而减小。在带有绊线的球体的情况下，沿负 z 坐标的平均速度分布与其他三个分布有很大不同，这三个分布几乎彼此重叠。然而，方差分布仅沿 y 坐标方向塌陷（并且与测量数据非常匹配）；沿 z 坐标的两个轮廓彼此不同，并且与其他两个轮廓不同。如上所述，这里的流动似乎具有倾斜的尾流。

这些模拟的更多细节（包括网格和 SGS 模型依赖性以及实验数据）将在未来的出版物中提供。

10.3.4.3 可再生能源模拟：大型风力发电机组
Jacobson 和 Delucchi（2009）概述了“到 2030 年实现可持续能源的道路”。他们的计划需要太阳能、水力、地热能和风能，其中包括 380 万台大型风力涡轮机。事实上，风能已经是一种快速增长的电力来源（例如，2016年，德国所有消耗电能的12.3\%是由风力涡轮机产生的，并且这一趋势还在增加）。斯塔。 Maria和Jacobson（2009）对大型风力发电机组（即风电场）对大气中能量的影响进行了概述和初步分析。由于风电场（来自转子下游的湍流）导致的近地表湍流的增加可能会影响那里的热量和蒸汽通量，并且一些研究表明边界层中的混合增加。此外，在风电场中，涡轮机之间相互作用。

当前和不久的将来的技术表明风力涡轮机的转子直径超过150 m，轮毂高度（支撑塔顶部转子的轴水平）约为200 m，功率输出约为10 MW。 Lu 和 Porté-Agel (2011) 在稳定的大气边界层中对一个超大型风电场进行了三维 LES。它们的转子直径为 112 m，轮毂高度为 119 m。该研究总结如下。

流域：域和物理设置取自著名的稳定边界层 (SBL) 案例，边界层高度约为 175 m，即 Beare 等人的比对研究。 （2006）基于全球能源和水循环实验大气边界层研究（GABLS）。因此，流域和模拟已经过全面测试，并且验证了没有风力涡轮机的流动的数值代码，从而可以轻松观察和量化涡轮机影响的评估。

\begin{center}\small 图 10.19 显示了左侧的域。基本思想是将单个风力涡轮机放置在具有周期性横向边界条件的域中，从而创建具有指定涡轮机放置的无限风电场。域的尺寸为垂直高度 L z = 400 m 和横向宽度 L y = 5 D = 560 m，其中 D = 112 m 是上述转子直径。有两种仿真配置，一种是 L x = 5 D（5D 情况），另一种是 L x = 8 D（8D 情况）。通过这种配置，给定的涡轮机会受到其邻居的影响，反之亦然，就像在真实的风电场中一样。\end{center}
风力涡轮机：图10.19显示了右侧的三叶片涡轮机。它位于距离域左边缘 x c = 80 m 处，高度 z c = 119 m，并沿着域的中心线（y c = 260 m）。涡轮机在流向方向上间隔5或8个直径，在横向上间隔5个直径。这些被认为是典型的风电场间距。涡轮机的转速为 8 rpm，与所选涡轮机、其发电量、叶片等一致。

在计算中，涡轮机采用执行器线方法（ALM）进行参数化，其中考虑了叶片的实际运动（Ivanell 等人，2009）。刀片上的力用沿刀片轴线的线表示。使用叶片单元的局部攻角计算体积力（由沿线每个点、每个叶片上以及每个时间步长的升力和阻力引起）

\begin{sourceblock}
\includegraphics[page=406,trim={53.00bp 469.25bp 53.87bp 52.00bp},clip,width=\textwidth,max height=0.38\textheight]{Ferziger4th.pdf}
\par\vspace{0.45\baselineskip}
\small 图 10.19 Lu 和 Porté-Agel 中的涡轮机域布局（左）和三叶片风力涡轮机、发电机和塔架（右）（2011 年）。 （经AIP出版许可转载）
\end{sourceblock}
这取决于叶片形状、瞬时解析流速以及该叶片的实际测量的翼型性能数据。这会在涡轮机平面上产生不稳定的空间变化力，这些力通过高斯平滑传递到流动动量方程中，以连接到附近的网格点。流动结果反馈至 ALM 并影响沿叶片计算的升力和阻力。详细信息可以在引用的 ALM 论文以及 Lu 和 Porté-Agel 论文及其参考文献中找到。

LES 代码和模型：使用的代码是 Porté-Agel 等人对 LES 代码的修改版本。 (2000)，其在水平方向上是伪谱，并且在垂直方向上使用交错网格上的二阶CDS。连续性方程、带有来自 ALM 的风致强迫项的动量守恒方程以及位温 14 的传输方程在 Boussinesq 近似（第 1.7.5 节）中求解不可压缩流，并包括科里奥利强迫。在域的顶部使用无应力/零梯度边界条件。在地面上，通过应用粗糙壁的 Monin-Obukhov 相似理论，瞬时壁应力与第一个垂直节点处的水平速度相关，该理论适用于分层流动（参见上面的方程（10.27）和（10.29））。类似的边界条件用于热通量。

网格具有均匀间隔的点，网格间距为 ∼ 3.3 m，8D 情况为 270 × 168 × 120 点，5D 情况为 168 × 168 × 120 点。进行了多项测试来验证这些选择。库朗极限 C = 0.06 与二阶精确 Adams-Bashforth 时间步进方案一起使用。根据标准 3/2 规则对光谱结果进行去锯齿。在收集数据之前，运行代码直到获得准统计稳定条件；记得这个
14 位温在气象学中被广泛使用，因为它具有适合研究可压缩介质（空气）中分层流的特性；它是高度为 z 的一团空气在不与周围环境进行热交换（绝热）的情况下移动到地面后的温度；所以 ( z ) = T ( z ) [ p 接地 )/ p ( z ) ] 0.286.

流动是三维且不稳定的，具有解析的湍流运动和由涡轮机引起的相干运动。

动量和热通量的 SGS 模型是 节 动力学模型的一个变体。 10.3.3.3 如上所述，并在 Stoll 和 Porté-Agel (2008) 中进行了描述；波特-阿格尔等人。 ( 2000 ) 及后续论文。添加了两个重要的特征，它们显著改善了动态方法的行为：（1）为了最小化动态过程中的误差，要最小化的总误差是通过沿着拉格朗日轨迹（流动中流体粒子的路径）累积局部误差来生成的（Meneveau et al. 1996）和（2）放宽了尺度相似性，因此模型形式和系数是尺度相关的。基本过滤器/网格宽度计算为 ( x y z ) 1 / 3。网格滤波器和两个测试滤波器均为水平方向的二维低通滤波器；垂直FD方向没有滤波。使用锐截止滤波器在傅里叶空间中对变量进行滤波，从而删除所有大于滤波器尺度的波数。动力学方程的结构与上述基本相同；然而，由于不强制尺度不变性，因此假定系数随网格间距变化的形式，并使用两个测试滤波器，一个为网格尺​​寸的两倍，一个为网格尺​​寸的四倍。然后，使用与传统动态模型中相同的最小化策略，可以确定未知数。当然，这里还有一个SGS位温的动态程序。请注意，因为动态过程使用最小的解析尺度，所以没有必要在 SGS 模型中包括稳定性校正，因为该效应已经存在于小解析尺度中（参见第 10.3.3.2 节）。

模拟结果：在风电场规模下，由于稳定（热分层）条件，速度剖面引起的垂直剪切随高度变化，科里奥利力效应引起的水平剪切也随高度变化（图10.20）。这些会导致涡轮机上出现显著的不对称负载。在这里，我们可以看到 LES 的强大功能，因为它实际上解决了移动涡轮叶片的叶尖涡流。此外，科里奥利力不仅会对风力涡轮机造成额外的横向剪切载荷，而且还会驱使部分湍流能量远离风力涡轮机尾流的中心。风力涡轮机的运动增强了热量的垂直混合，导致风力涡轮机尾流中的空气温度升高，表面热通量降低，从而影响热能预算。

我们展示了该研究的三个数据。图 10.21 清楚地表明，水平速度分量和潜在温度的平均垂直剖面（空间和时间上的平均值）受到涡轮机存在的严重影响。边界层顶部的射流通过涡轮机的能量提取而被移除，并且当涡轮机存在时，科里奥利效应被改变。两种S x = L x / D情况之间的差异并不大。混合层深度的增加在位温上是明显的。

\begin{center}\small 图 10.22 显示了顺风不同 S x 距离处平均流向速度剖面的演变。注意涡轮机位于 S 涡轮机 ∼ 0 处。 7. 当涡轮机靠得更近时，可以从流动中提取更多能量。\end{center}
最后，我们在图 10.23 中看到了涡轮机对流动能谱的影响。 N 是稳定层结中自然振荡的 Brunt–Väisälä 或层结频率。除了顺风处湍流能量明显增加外

\begin{sourceblock}
\includegraphics[page=408,trim={53.00bp 356.48bp 53.87bp 52.00bp},clip,width=\textwidth,max height=0.38\textheight]{Ferziger4th.pdf}
\par\vspace{0.45\baselineskip}
\small 图 10.20 使用涡度等值面 ω(∼ 0 . 3 | ω | ) 可视化 t = 150 s 时由移动涡轮叶片引起的叶尖涡。经 AIP Publishing 许可，转载自 Lu 和 Porté-Agel (2011)
\end{sourceblock}
\begin{sourceblock}
\includegraphics[page=408,trim={53.00bp 163.90bp 53.87bp 343.37bp},clip,width=\textwidth,max height=0.38\textheight]{Ferziger4th.pdf}
\par\vspace{0.45\baselineskip}
\small 图 10.21 速度和位温的垂直平均分布。水平虚线为涡轮轮毂高度；曲线上背景中的光线是 M-O 相似曲线。经 AIP Publishing 许可，转载自 Lu 和 Porté-Agel (2011)
\end{sourceblock}
\begin{sourceblock}
\includegraphics[page=409,trim={53.00bp 488.25bp 53.87bp 52.00bp},clip,width=\textwidth,max height=0.38\textheight]{Ferziger4th.pdf}
\par\vspace{0.45\baselineskip}
\small 图10.22 域中心线上流向（轴向）剖面的顺风变化。经 AIP Publishing 许可，转载自 Lu 和 Porté-Agel (2011)
\end{sourceblock}
\begin{center}\small 图 10.23 风力涡轮机对涡轮机平面下游 x = 3 D 处的湍流能谱的影响，\end{center}
y c，z = z c + D / 2（刀片顶端的高度）。经 AIP Publishing 许可，转载自 Lu 和 Porté-Agel (2011)

本文针对此处所示的涡轮机，深入讨论了涡轮机影响下的湍流通量，并根据气象情况对其进行了解释。

Lu 和 Porté-Agel 表示，他们的结果表明“大涡模拟可以提供有价值的三维高分辨率速度和温度场，用于定量描述风力涡轮机尾流及其对风电场内部和上方的热量和动量湍流通量的影响。”这里，表面动量通量减少了30\%以上，表面浮力通量减少了15\%以上。风电场对动量和热量的垂直湍流通量有很大影响，这可能会影响当地的气象。

\subsection*{10.3.5 雷诺平均纳维-斯托克斯 (RANS) 模拟}\phantomsection\addcontentsline{toc}{subsection}{10.3.5 雷诺平均纳维-斯托克斯 (RANS) 模拟}
传统上，工程师通常只对了解湍流的一些定量特性感兴趣，例如物体上的平均力（也许还有它们的分布）、两股传入流体流之间的混合程度或已发生反应的物质的量。至少可以说，使用上述方法来计算这些数量是矫枉过正的。然而，事情已经发生了变化，今天的问题更加复杂，设计更加严格，流程更加依赖于正在发生的事情的细节。然后，明智地使用上述 DNS 和 LES 方案中包含的方法以及智能使用本节中描述的 RANS 方法就变得至关重要。

如上所述，由于一个多世纪前的奥斯本雷诺兹，本节的方法被称为雷诺平均方法。 Leschziner (2010) 写道“……在撰写本文时，工业流的绝大多数计算预测都是基于 RANS 方程。”在雷诺平均湍流方法中，控制方程以某种方式进行平均，如第 1 节中所述。 10.3.当我们查看下面的各个部分时，我们需要对这些平均策略的影响保持警惕。回想一下，实际策略是：

1. 稳态流：所有不稳定因素均被平均，即所有不稳定因素均被视为湍流的一部分。结果是平均流动方程是稳定的。这些就是雷诺平均纳维-斯托克斯 (RANS) 方程。 2. 非定常流：方程是对一整套统计上相同的流实现（系综）进行平均。 15 因此，所有随机波动都会被平均，因此隐含着“动荡”的一部分。然而，如果流中存在确定性和连贯的结构，它们应该在平均中幸存下来。这些整体平均方程可能具有不稳定的均值。此类流在上面被定义为不稳定的 RANS 或 URANS (Durbin 2002) 或瞬态 RANS 或 TRANS (Hanjali ́c 2002)。

我们再次回顾一下，在平均时，纳维-斯托克斯方程的非线性产生了必须建模的项，就像之前所做的那样。上面简要讨论了湍流的复杂性，使得任何单一的雷诺平均模型都不太可能能够很好地表示所有湍流，因此湍流模型应被视为工程近似而不是科学定律。 Hanjali ́c (2004) 对 RANS 及其湍流模型进行了全面回顾。

10.3.5.1 雷诺平均纳维-斯托克斯 (RANS) 方程
在统计稳态流中，每个变量都可以写成时间平均值 φ 和关于该值 φ ′ 的波动之和：

\begin{sourceblock}
\includegraphics[page=410,trim={53.00bp 143.99bp 53.87bp 505.39bp},clip,width=0.92\textwidth,max height=0.38\textheight]{Ferziger4th.pdf}
\end{sourceblock}
\begin{sourceblock}
\begin{tikzpicture}
\node[inner sep=0,anchor=south west] (image) at (0,0) {\includegraphics[page=410,trim={53.00bp 96.04bp 53.87bp 533.99bp},clip,width=0.92\textwidth,max height=0.38\textheight]{Ferziger4th.pdf}};
\begin{scope}[x={(image.south east)},y={(image.north west)}]
\fill[white] (0.00000,0.71891) rectangle (0.07732,1.00000);
\end{scope}
\end{tikzpicture}
\end{sourceblock}
\begin{sourceblock}
\includegraphics[page=411,trim={53.00bp 456.25bp 53.87bp 52.00bp},clip,width=\textwidth,max height=0.38\textheight]{Ferziger4th.pdf}
\par\vspace{0.45\baselineskip}
\small 图 10.24 统计稳态流的时间平均（左）和非稳态流的整体平均（右）
\end{sourceblock}
这里 t 是时间，T 是平均间隔。与波动的典型时间尺度相比，这个间隔必须很大；因此，我们感兴趣的是 T →∞ 的极限，见图 10.24。如果 T 足够大，则 φ 不依赖于平均开始的时间。

如果流动不稳定，则无法使用时间平均，除非将随附的湍流模型的时间尺度调整为平均时间段 T < ∞。在大多数情况下，不稳态流是通过整体平均来处理的。这个概念之前已经讨论过，如图 10.24 16 所示：

\begin{sourceblock}
\includegraphics[page=411,trim={53.00bp 269.37bp 53.87bp 361.29bp},clip,width=0.92\textwidth,max height=0.38\textheight]{Ferziger4th.pdf}
\end{sourceblock}
其中 N 是系综成员的数量，并且必须足够大以消除湍流（随机）波动的影响。这种类型的平均可以应用于任何流。我们使用术语“雷诺平均”来指代任何这些平均过程；将其应用于纳维-斯托克斯方程可生成稳定情况下的雷诺平均纳维-斯托克斯 (RANS) 方程以及不稳定情况下的 URANS 或 TRANS。

从方程。 ( 10.31 )，因此 φ ′ = 0。因此，对守恒方程中的任何线性项进行平均，只需给出平均量的相同项。从二次非线性项我们得到两项，即平均值和协方差的乘积：

\begin{sourceblock}
\begin{tikzpicture}
\node[inner sep=0,anchor=south west] (image) at (0,0) {\includegraphics[page=411,trim={53.00bp 126.21bp 53.87bp 520.66bp},clip,width=0.92\textwidth,max height=0.38\textheight]{Ferziger4th.pdf}};
\begin{scope}[x={(image.south east)},y={(image.north west)}]
\fill[white] (0.00000,0.78983) rectangle (0.14093,1.00000);
\end{scope}
\end{tikzpicture}
\end{sourceblock}
\footnotetext[16]{对于具有持久结构的不稳定情况，流动可能不会单调变化，如果 LES 线被想象为代表整体平均流动中的相干结构，结果可能看起来很像图 10.6。参见 Chen 和 Jaw (1998) 中的图 1.8。}
仅当两个量不相关时，最后一项才为零；在湍流中很少出现这种情况，因此，守恒方程包含诸如 ρ u ′ 之类的项

\begin{sourceblock}
\begin{tikzpicture}
\node[inner sep=0,anchor=south west] (image) at (0,0) {\includegraphics[page=412,trim={53.00bp 569.05bp 53.87bp 79.43bp},clip,width=0.92\textwidth,max height=0.38\textheight]{Ferziger4th.pdf}};
\begin{scope}[x={(image.south east)},y={(image.north west)}]
\fill[white] (0.00000,0.86127) rectangle (0.02746,1.00000);
\fill[white] (0.03038,0.86127) rectangle (0.14271,1.00000);
\fill[white] (0.14571,0.86127) rectangle (0.21468,1.00000);
\fill[white] (0.21762,0.86127) rectangle (0.27317,1.00000);
\fill[white] (0.27618,0.86127) rectangle (0.30595,1.00000);
\fill[white] (0.30890,0.86127) rectangle (0.32701,1.00000);
\fill[white] (0.32992,0.86127) rectangle (0.40881,1.00000);
\fill[white] (0.41173,0.86127) rectangle (0.45314,1.00000);
\fill[white] (0.45609,0.86127) rectangle (0.61317,1.00000);
\fill[white] (0.61621,0.86127) rectangle (0.73582,1.00000);
\fill[white] (0.73880,0.86127) rectangle (0.83179,1.00000);
\fill[white] (0.83477,0.86127) rectangle (0.90614,1.00000);
\fill[white] (0.90911,0.86127) rectangle (0.96881,1.00000);
\fill[white] (0.97176,0.86127) rectangle (1.00000,1.00000);
\fill[white] (0.00000,0.18969) rectangle (0.04448,0.93205);
\fill[white] (0.06325,0.06795) rectangle (0.08493,0.83918);
\fill[white] (0.08677,0.18403) rectangle (0.16310,0.83918);
\fill[white] (0.16493,0.18403) rectangle (0.20635,0.83918);
\fill[white] (0.20818,0.18969) rectangle (0.32192,0.84485);
\fill[white] (0.32379,0.18403) rectangle (0.44914,0.91450);
\fill[white] (0.45245,0.18403) rectangle (0.50051,0.83918);
\fill[white] (0.50232,0.18969) rectangle (0.54749,0.93205);
\fill[white] (0.53687,0.06795) rectangle (0.58448,0.88732);
\fill[white] (0.58632,0.18403) rectangle (0.67188,0.83918);
\fill[white] (0.67374,0.18403) rectangle (0.70352,0.83918);
\fill[white] (0.70535,0.18403) rectangle (0.74677,0.83918);
\fill[white] (0.74860,0.18969) rectangle (0.86262,0.84485);
\fill[white] (0.86448,0.18969) rectangle (0.94418,0.84485);
\fill[white] (0.94602,0.18403) rectangle (1.00000,0.84485);
\fill[white] (0.00000,0.00000) rectangle (0.08568,0.16195);
\fill[white] (0.08692,0.00000) rectangle (0.17242,0.16195);
\fill[white] (0.17362,0.00000) rectangle (0.24995,0.16195);
\fill[white] (0.25119,0.00000) rectangle (0.33585,0.16195);
\fill[white] (0.33711,0.00000) rectangle (0.37020,0.16195);
\fill[white] (0.37143,0.00000) rectangle (0.51426,0.16195);
\fill[white] (0.51549,0.00000) rectangle (0.62514,0.16195);
\fill[white] (0.62638,0.00000) rectangle (0.65450,0.16195);
\fill[white] (0.65570,0.00000) rectangle (0.72707,0.16195);
\fill[white] (0.72827,0.00000) rectangle (0.75805,0.16195);
\fill[white] (0.75925,0.00000) rectangle (0.80066,0.16195);
\fill[white] (0.80186,0.00000) rectangle (0.87158,0.16195);
\fill[white] (0.87278,0.00000) rectangle (1.00000,0.16195);
\end{scope}
\end{tikzpicture}
\end{sourceblock}
对于没有体积力的不可压缩流动，平均连续性和动量方程可以用张量表示法和笛卡尔坐标写为：

\begin{sourceblock}
\begin{tikzpicture}
\node[inner sep=0,anchor=south west] (image) at (0,0) {\includegraphics[page=412,trim={53.00bp 463.05bp 53.87bp 138.53bp},clip,width=0.92\textwidth,max height=0.38\textheight]{Ferziger4th.pdf}};
\begin{scope}[x={(image.south east)},y={(image.north west)}]
\fill[white] (0.19041,0.23637) rectangle (0.27759,0.43231);
\fill[white] (0.33056,0.25325) rectangle (0.34908,0.43231);
\fill[white] (0.28304,0.14498) rectangle (0.31125,0.32404);
\fill[white] (0.21805,0.03129) rectangle (0.24731,0.21453);
\fill[white] (0.31666,0.03129) rectangle (0.35173,0.21453);
\fill[white] (0.35122,0.01859) rectangle (0.36186,0.15149);
\end{scope}
\end{tikzpicture}
\end{sourceblock}
其中 τ i j 是平均粘性应力张量分量：

\begin{sourceblock}
\begin{tikzpicture}
\node[inner sep=0,anchor=south west] (image) at (0,0) {\includegraphics[page=412,trim={53.00bp 403.15bp 53.87bp 226.88bp},clip,width=0.92\textwidth,max height=0.38\textheight]{Ferziger4th.pdf}};
\begin{scope}[x={(image.south east)},y={(image.north west)}]
\fill[white] (0.34650,0.25921) rectangle (0.36427,0.57934);
\fill[white] (0.36292,0.22902) rectangle (0.37356,0.46663);
\fill[white] (0.37212,0.22902) rectangle (0.38277,0.46663);
\fill[white] (0.38983,0.25921) rectangle (0.41802,0.57934);
\fill[white] (0.42153,0.25921) rectangle (0.44689,0.57934);
\end{scope}
\end{tikzpicture}
\end{sourceblock}
最后，可以写出标量平均值的方程：

\begin{sourceblock}
\begin{tikzpicture}
\node[inner sep=0,anchor=south west] (image) at (0,0) {\includegraphics[page=412,trim={53.00bp 341.09bp 53.87bp 295.95bp},clip,width=0.92\textwidth,max height=0.38\textheight]{Ferziger4th.pdf}};
\begin{scope}[x={(image.south east)},y={(image.north west)}]
\fill[white] (0.19402,0.56151) rectangle (0.27423,0.95876);
\fill[white] (0.32716,0.56151) rectangle (0.34571,0.95876);
\fill[white] (0.27964,0.32131) rectangle (0.30782,0.71890);
\fill[white] (0.38211,0.31237) rectangle (0.42003,0.71890);
\fill[white] (0.42027,0.28419) rectangle (0.45176,0.71890);
\fill[white] (0.45510,0.32131) rectangle (0.48328,0.71890);
\fill[white] (0.48511,0.31237) rectangle (0.53026,0.76254);
\fill[white] (0.21814,0.06907) rectangle (0.24743,0.47560);
\fill[white] (0.31329,0.06907) rectangle (0.34836,0.47560);
\fill[white] (0.34782,0.04124) rectangle (0.35847,0.33574);
\end{scope}
\end{tikzpicture}
\end{sourceblock}
守恒方程中雷诺应力和湍流标量通量的存在意味着这些方程不是封闭的，也就是说，它们包含的变量比方程的数量多。闭合需要使用一些近似值，这些近似值通常采用以平均量来规定雷诺应力张量和湍流标量通量的形式。

可以推导出高阶相关性的方程，例如雷诺应力张量的方程，但这些方程包含更多（和高阶）未知的相关性，需要建模近似。这些方程将在稍后介绍，但重要的一点是不可能导出一组封闭的精确方程。我们在工程中使用的近似湍流模型在地球科学中通常称为参数化。

10.3.5.2 简单湍流模型及其应用
为了闭合方程，我们必须引入湍流模型。为了了解合理的模型可能是什么，我们注意到，正如我们在上一节中所做的那样，在层流中，能量耗散以及垂直于流线的质量、动量和能量的传递是由粘度调节的，因此很自然地假设湍流的影响可以表示为粘度的增加。这导致了雷诺应力的涡粘模型：

\begin{sourceblock}
\begin{tikzpicture}
\node[inner sep=0,anchor=south west] (image) at (0,0) {\includegraphics[page=413,trim={53.00bp 557.67bp 53.87bp 72.36bp},clip,width=0.92\textwidth,max height=0.38\textheight]{Ferziger4th.pdf}};
\begin{scope}[x={(image.south east)},y={(image.north west)}]
\fill[white] (0.24908,0.25145) rectangle (0.32424,0.61451);
\end{scope}
\end{tikzpicture}
\end{sourceblock}
以及标量的涡流扩散模型：

\begin{sourceblock}
\includegraphics[page=413,trim={53.00bp 497.18bp 53.87bp 139.85bp},clip,width=0.92\textwidth,max height=0.38\textheight]{Ferziger4th.pdf}
\end{sourceblock}
在等式中。 ( 10.38 ), k 为湍流动能：

\begin{sourceblock}
\begin{tikzpicture}
\node[inner sep=0,anchor=south west] (image) at (0,0) {\includegraphics[page=413,trim={53.00bp 439.60bp 53.87bp 198.52bp},clip,width=0.92\textwidth,max height=0.38\textheight]{Ferziger4th.pdf}};
\begin{scope}[x={(image.south east)},y={(image.north west)}]
\fill[white] (0.30854,0.54425) rectangle (0.32833,0.95717);
\fill[white] (0.24974,0.29872) rectangle (0.26788,0.71128);
\fill[white] (0.27320,0.30835) rectangle (0.30138,0.72091);
\fill[white] (0.30848,0.04283) rectangle (0.35423,0.76660);
\fill[white] (0.34364,0.22163) rectangle (0.38003,0.76660);
\fill[white] (0.73642,0.30835) rectangle (0.75113,0.72091);
\fill[white] (0.90923,0.29515) rectangle (1.00000,0.70771);
\end{scope}
\end{tikzpicture}
\end{sourceblock}
其中 μ t 是湍流（或涡流）粘度，t 是湍流扩散率。方程仍未闭合，但附加未知数的数量从 9 个（雷诺应力张量的 6 个分量和湍流通量矢量的 3 个分量）减少到两个（ μ t 和 t ）。

方程中的最后一项。 ( 10.38 ) 需要保证，当方程两边都收缩时（两个指数设置相等并求和），方程仍然正确。尽管涡粘性假设在细节上并不正确，但它很容易实现，并且通过仔细应用，可以为许多流动提供相当好的结果。

\begin{sourceblock}
\begin{tikzpicture}
\node[inner sep=0,anchor=south west] (image) at (0,0) {\includegraphics[page=413,trim={53.00bp 309.17bp 53.87bp 343.01bp},clip,width=0.92\textwidth,max height=0.38\textheight]{Ferziger4th.pdf}};
\begin{scope}[x={(image.south east)},y={(image.north west)}]
\fill[white] (0.48974,0.08596) rectangle (0.51991,0.91404);
\fill[white] (0.51507,0.00000) rectangle (0.55762,0.30014);
\fill[white] (0.56346,0.00000) rectangle (0.61152,0.29298);
\fill[white] (0.61735,0.00000) rectangle (0.63546,0.29298);
\fill[white] (0.64129,0.00000) rectangle (0.72099,0.29298);
\fill[white] (0.72683,0.00000) rectangle (0.79901,0.29298);
\fill[white] (0.80653,0.00000) rectangle (0.83783,0.30014);
\fill[white] (0.84364,0.00000) rectangle (1.00000,0.29298);
\end{scope}
\end{tikzpicture}
\end{sourceblock}
2 k 和长度尺度 L。量纲分析表明：

\begin{sourceblock}
\begin{tikzpicture}
\node[inner sep=0,anchor=south west] (image) at (0,0) {\includegraphics[page=413,trim={53.00bp 275.46bp 53.87bp 375.22bp},clip,width=0.92\textwidth,max height=0.38\textheight]{Ferziger4th.pdf}};
\begin{scope}[x={(image.south east)},y={(image.north west)}]
\fill[white] (0.00000,0.92432) rectangle (0.10066,1.00000);
\fill[white] (0.10337,0.92432) rectangle (0.18229,1.00000);
\fill[white] (0.18502,0.92432) rectangle (0.24310,1.00000);
\end{scope}
\end{tikzpicture}
\end{sourceblock}
其中C μ为无量纲常数，稍后给出其值。

在最简单的实际模型（混合长度模型）中，k 是使用近似值 q = L ∂ u /∂ y 从平均速度场确定的，并且 L 是坐标的规定函数。对于简单的流动，L 的精确规定是可能的，但对于分离的或高度三维的流动则不然。因此，混合长度模型只能应用于相对简单的流；它们也称为零方程模型。

规定湍流量的困难表明人们可以使用偏微分方程来计算它们。由于对湍流的最低限度描述至少需要速度尺度和长度尺度，因此从两个此类方程导出所需量的模型是合理的选择。在几乎所有此类模型中，湍流动能方程 k 决定了速度尺度。这个量的精确方程并不难推导：

\begin{sourceblock}
\begin{tikzpicture}
\node[inner sep=0,anchor=south west] (image) at (0,0) {\includegraphics[page=414,trim={53.00bp 552.06bp 53.87bp 50.28bp},clip,width=0.92\textwidth,max height=0.38\textheight]{Ferziger4th.pdf}};
\begin{scope}[x={(image.south east)},y={(image.north west)}]
\fill[white] (0.10192,0.73292) rectangle (0.17844,0.91834);
\fill[white] (0.21750,0.73307) rectangle (0.28271,0.91850);
\fill[white] (0.28292,0.72022) rectangle (0.32340,0.91850);
\fill[white] (0.18388,0.62759) rectangle (0.21206,0.80893);
\fill[white] (0.12418,0.51254) rectangle (0.15347,0.69796);
\fill[white] (0.24608,0.51254) rectangle (0.28114,0.69796);
\fill[white] (0.28060,0.49984) rectangle (0.29125,0.63417);
\end{scope}
\end{tikzpicture}
\end{sourceblock}
有关该方程推导的详细信息，请参见 Pope (2000)、Chen 和 Jaw (1998) 或 Wilcox (2006)。该方程左侧的项和右侧的第一项不需要建模。最后一项表示密度 ρ 和耗散 ε 的乘积，即湍流能量不可逆地转化为内能的速率。我们将给出下面的耗散方程。

右侧第二项代表动能的湍流扩散（实际上是波动本身对速度波动的传递）；它几乎总是使用梯度扩散假设来建模：

\begin{sourceblock}
\begin{tikzpicture}
\node[inner sep=0,anchor=south west] (image) at (0,0) {\includegraphics[page=414,trim={53.00bp 410.08bp 53.87bp 222.94bp},clip,width=0.92\textwidth,max height=0.38\textheight]{Ferziger4th.pdf}};
\begin{scope}[x={(image.south east)},y={(image.north west)}]
\fill[white] (0.29555,0.28261) rectangle (0.32376,0.63164);
\end{scope}
\end{tikzpicture}
\end{sourceblock}
其中 μ t 是上面定义的涡流粘度，σ k 是湍流普朗特数，其值近似为 1。涡粘性的一大弱点是它是一个标量，这极大地限制了它代表一般湍流过程的能力。在更复杂的模型中，这里不会描述，可以将涡粘性制成张量，或者更好的是，可以在不存在涡粘性的情况下创建模型（参见第 10.3.5.4 节和第 10.3.3.5 节中的 EARSM）。

等式右边的第三项。 ( 10.42 ) 表示平均流产生湍流动能的速率，即动能从平均流到湍流的转移。如果我们使用涡粘假设（10.38）来估计雷诺应力，则可以写成：

\begin{sourceblock}
\includegraphics[page=414,trim={53.00bp 243.19bp 53.87bp 386.84bp},clip,width=0.92\textwidth,max height=0.38\textheight]{Ferziger4th.pdf}
\end{sourceblock}
并且，由于该方程的右侧可以根据将要计算的量进行计算，因此湍流动能方程的展开已完成。

如上所述，需要另一个方程来确定湍流的长度尺度。这个选择并不明显，为此目的使用了许多方程。最流行的一种基于这样的观察：能量方程中需要耗散，并且在所谓的平衡湍流中，即湍流的产生速率和破坏速率接近平衡的湍流中，耗散 ε 以及 k 和 L 之间的关系为 18：

\footnotetext[18]{这种关系在使用基于 TKE 的 SGS 模型的 LES 中发挥着重要作用；在这种情况下，L = 并使用比例常数 O (1)。}
\begin{sourceblock}
\includegraphics[page=415,trim={53.00bp 582.79bp 53.87bp 53.84bp},clip,width=0.92\textwidth,max height=0.38\textheight]{Ferziger4th.pdf}
\end{sourceblock}
这个想法基于以下事实：在高雷诺数下，存在从最大尺度到最小尺度的能量级联，并且转移到小尺度的能量被耗散。方程（10.45）基于惯性能量传递的估计。

方程 (10.45) 允许使用耗散方程作为获得 ε 和 L 的一种方法。方程中没有使用常数。 ( 10.45 ) 因为该常数可以与完整模型中的其他常数组合。

尽管可以从纳维-斯托克斯方程导出精确的耗散方程，但应用于该方程的建模非常严格，因此最好将整个方程视为模型。因此，我们不会尝试推导它。该方程最常用的形式是：

\begin{sourceblock}
\begin{tikzpicture}
\node[inner sep=0,anchor=south west] (image) at (0,0) {\includegraphics[page=415,trim={53.00bp 405.47bp 53.87bp 224.56bp},clip,width=0.92\textwidth,max height=0.38\textheight]{Ferziger4th.pdf}};
\begin{scope}[x={(image.south east)},y={(image.north west)}]
\fill[white] (0.09215,0.45251) rectangle (0.16797,0.77264);
\end{scope}
\end{tikzpicture}
\end{sourceblock}
在该模型中，涡流粘度表示为：

\begin{sourceblock}
\includegraphics[page=415,trim={53.00bp 346.70bp 53.87bp 289.60bp},clip,width=0.92\textwidth,max height=0.38\textheight]{Ferziger4th.pdf}
\end{sourceblock}
其中方程。 (10.45) 用于确定 L。

该模型基于方程。 （10.42）和（10.46）被称为k-ε模型并得到了广泛的应用。该模型包含五个参数；它们最常用的值是：

\begin{sourceblock}
\includegraphics[page=415,trim={53.00bp 265.54bp 53.87bp 385.14bp},clip,width=0.92\textwidth,max height=0.38\textheight]{Ferziger4th.pdf}
\end{sourceblock}
该模型在计算机代码中的实现相对简单。 RANS 方程与层流方程具有相同的形式，只要将分子粘度 μ 替换为有效粘度 μ eff = μ + μ t。最重要的区别是需要求解两个新的偏微分方程，并且 μ t 通常在流域内变化几个数量级。由于与湍流相关的时间尺度比与平均流相关的时间尺度短得多，因此 k – ε 模型（或任何其他湍流模型）的方程比层流方程刚性得多。因此，除了下面要讨论的方程之外，这些方程的离散化几乎没有困难，但求解方法必须考虑增加的刚度。

因此，在数值求解过程中，首先执行动量和压力校正方程的外迭代，其中涡粘性值基于前一迭代结束时的 k 和 ε 值。完成此操作后，将进行湍流动能和耗散方程的外部迭代。由于这些方程高度非线性，因此必须在迭代之前对其进行线性化。完成湍流模型方程的迭代后，我们准备重新计算涡粘性并开始新的外部迭代。

使用涡粘模型的方程的刚度需要时间推进或欠松弛才能获得收敛的稳态解。太大的时间步长（或迭代方法中的欠松弛因子）可能导致 ε 的 k 为负值（尤其是在壁附近），从而导致数值不稳定。即使使用时间推进方法，仍可能需要欠松弛来增强数值稳定性。欠松弛参数的典型值与动量方程中使用的参数类似（0.6-0.9，取决于时间步长、网格质量和流动问题；较高的值适用于高质量网格和小时间步长）。

还提出了许多其他二元方程模型；我们只描述其中之一。一个明显的想法是为长度尺度本身写一个微分方程；这已经被尝试过，但并未取得太大成功。第二个最常用的模型是 k – ω 模型，最初由 Saffman 提出，但由 Wilcox 推广。在此模型中，使用反时标 ω 的方程；这个量可以有多种解释，但它们都不是很有启发性，所以这里省略了。 k – ω 模型使用湍流动能方程 ( 10.42 )，但必须稍作修改：

\begin{sourceblock}
\begin{tikzpicture}
\node[inner sep=0,anchor=south west] (image) at (0,0) {\includegraphics[page=416,trim={53.00bp 329.17bp 53.87bp 299.51bp},clip,width=0.92\textwidth,max height=0.38\textheight]{Ferziger4th.pdf}};
\begin{scope}[x={(image.south east)},y={(image.north west)}]
\fill[white] (0.08063,0.46503) rectangle (0.15714,0.78083);
\fill[white] (0.19621,0.46530) rectangle (0.26144,0.78110);
\fill[white] (0.26165,0.44367) rectangle (0.30214,0.78110);
\fill[white] (0.16259,0.28590) rectangle (0.19080,0.59450);
\fill[white] (0.10292,0.08996) rectangle (0.13218,0.40577);
\end{scope}
\end{tikzpicture}
\end{sourceblock}
上面所说的几乎所有内容都适用于此。 Wilcox (2006) 给出的 ω 方程为：

\begin{sourceblock}
\begin{tikzpicture}
\node[inner sep=0,anchor=south west] (image) at (0,0) {\includegraphics[page=416,trim={53.00bp 258.44bp 53.87bp 370.60bp},clip,width=0.92\textwidth,max height=0.38\textheight]{Ferziger4th.pdf}};
\begin{scope}[x={(image.south east)},y={(image.north west)}]
\fill[white] (0.06030,0.46739) rectangle (0.14192,0.77898);
\fill[white] (0.18099,0.46065) rectangle (0.24620,0.77951);
\fill[white] (0.24641,0.43854) rectangle (0.29200,0.77951);
\fill[white] (0.14737,0.27925) rectangle (0.17555,0.59084);
\fill[white] (0.08511,0.08113) rectangle (0.11438,0.40000);
\end{scope}
\end{tikzpicture}
\end{sourceblock}
在该模型中，涡流粘度表示为：

\begin{sourceblock}
\includegraphics[page=416,trim={53.00bp 202.23bp 53.87bp 436.16bp},clip,width=0.92\textwidth,max height=0.38\textheight]{Ferziger4th.pdf}
\end{sourceblock}
该模型中的系数比 k – ε 模型中的系数稍微复杂一些。他们是：

\begin{sourceblock}
\includegraphics[page=416,trim={53.00bp 134.28bp 53.87bp 503.84bp},clip,width=0.92\textwidth,max height=0.38\textheight]{Ferziger4th.pdf}
\end{sourceblock}
该模型的数值行为与 k – ε 模型类似。

有兴趣了解更多关于这些模型的读者可以参考 Wilcox (2006) 的书。 Menter 于 1993 年推出了该模型的一个流行变体
（参见 Menter 1994）；他的剪切应力传输湍流模型（有时与分离涡模拟一起使用；参见第 10.3.7 节）用于飞机、卡车等的空气动力学研究（Menter 等人，2003）。 19
10.3.5.3 v2f 模型
从上面可以清楚地看出，湍流模型的一个主要问题是不知道在壁附近应用的适当条件。困难在于我们根本不知道其中一些量在墙壁附近的表现如何。而且，在壁面附近，湍流动能的变化甚至耗散都非常快。这表明尝试为该地区的这些数量规定条件并不是一个好主意。另一个主要问题是，尽管多年来致力于开发旨在治疗近壁区域的“低雷诺数”模型，但取得的成功相对较小。

Durbin (1991) 认为问题不在于壁附近的湍流雷诺数较低（尽管粘性效应当然很重要）。不渗透条件（壁面法向速度为零）更为重要。这表明，我们不应该尝试寻找低雷诺数模型，而应该使用由于不渗透条件而在墙壁附近变得非常小的数量。这样的量就是法向速度（工程师通常称为 v）及其涨落 ( v ' 2 )，因此杜宾为该量引入了一个方程。人们发现该模型还需要阻尼函数 f，因此命名为 v 2 – f （或 v 2 f ）模型。它似乎以与 k – ε 模型基本相同的成本提供了改进的结果。亚卡里诺等人。 (2003) 在非定常分离流的非定常 RANS (URANS) 模拟中使用了 v 2 f 模型，非常成功。

为了解决靠近墙壁的问题，杜宾建议使用椭圆松弛。这个想法如下。假设 φ i j 是某个建模量。设模型预测值为 φ m
我 j .我们不接受该值作为模型中使用的值，而是求解方程：

\begin{sourceblock}
\includegraphics[page=417,trim={53.00bp 212.65bp 53.87bp 425.57bp},clip,width=0.92\textwidth,max height=0.38\textheight]{Ferziger4th.pdf}
\end{sourceblock}
其中 L 是湍流的长度尺度，通常取 L ≈ k 3 / 2 /ε。这一程序的引入似乎大大缓解了困难。有关这些模型和其他类似模型的更多详细信息，请参阅 Durbin 最近出版的一本书（2009 年）；另请参见 Durbin 和 Pettersson Reif (2011)。

\footnotetext[19]{NASA 湍流建模资源 (NASA TMR 2019 ) 提供 RANS 湍流模型的文档，包括 Spalart-Allmaras、Menter、Wilcox 和其他模型的最新（经常修正）版本以及验证和确认测试用例、网格和数据库。}
10.3.5.4 雷诺应力和代数雷诺应力模型
涡粘模型有明显的缺陷；有些是涡粘性假设的结果，方程。 ( 10.38 )，无效。在二维中，始终可以选择涡流粘度，使该方程能够给出正确的剪切应力分布（τ i j 的 1-2 分量）。在三维流动中，雷诺应力和应变率可能不会以如此简单的方式相关。这意味着涡流粘度可能不再是标量；事实上，测量和模拟都表明它变成了一个张量。

已经提出了基于使用 k 和 ε 方程的各向异性（张量）模型。阿部等人。 (2003) 描述了一种非线性涡粘模型，该模型在捕获壁附近的高度各向异性湍流方面特别成功。 Leschziner (2010) 描述了多种模型，包括线性涡粘性、非线性涡粘性和雷诺应力模型。

\begin{sourceblock}
\begin{tikzpicture}
\node[inner sep=0,anchor=south west] (image) at (0,0) {\includegraphics[page=418,trim={53.00bp 413.63bp 53.87bp 234.85bp},clip,width=0.92\textwidth,max height=0.38\textheight]{Ferziger4th.pdf}};
\begin{scope}[x={(image.south east)},y={(image.north west)}]
\fill[white] (0.03528,0.86127) rectangle (0.08668,1.00000);
\fill[white] (0.09389,0.86127) rectangle (0.15699,1.00000);
\fill[white] (0.16418,0.86127) rectangle (0.27170,1.00000);
\fill[white] (0.27898,0.86127) rectangle (0.37035,1.00000);
\fill[white] (0.37759,0.86127) rectangle (0.40571,1.00000);
\fill[white] (0.41290,0.86127) rectangle (0.52259,1.00000);
\fill[white] (0.52983,0.86127) rectangle (0.57459,1.00000);
\fill[white] (0.58177,0.86127) rectangle (0.65317,1.00000);
\fill[white] (0.66042,0.86127) rectangle (0.70180,1.00000);
\fill[white] (0.70899,0.86127) rectangle (0.90304,1.00000);
\fill[white] (0.91029,0.86127) rectangle (1.00000,1.00000);
\fill[white] (0.00000,0.18460) rectangle (0.07735,0.83918);
\fill[white] (0.08180,0.18460) rectangle (0.12319,0.83918);
\fill[white] (0.12764,0.18460) rectangle (0.20066,0.83918);
\fill[white] (0.20511,0.18460) rectangle (0.23988,0.83918);
\fill[white] (0.24436,0.18460) rectangle (0.35230,0.83918);
\fill[white] (0.35681,0.18460) rectangle (0.47645,0.83918);
\fill[white] (0.48096,0.18460) rectangle (0.52069,0.83918);
\fill[white] (0.52514,0.18460) rectangle (0.56656,0.83918);
\fill[white] (0.57101,0.18460) rectangle (0.76502,0.83918);
\fill[white] (0.76950,0.18460) rectangle (0.84752,0.83918);
\fill[white] (0.85206,0.14439) rectangle (0.87552,0.86014);
\fill[white] (0.87408,0.14439) rectangle (0.88472,0.62967);
\fill[white] (0.89179,0.20555) rectangle (0.91997,0.86014);
\fill[white] (0.92349,0.19026) rectangle (0.96863,0.93205);
\fill[white] (0.98737,0.06795) rectangle (0.99802,0.55323);
\fill[white] (0.00000,0.00000) rectangle (0.07152,0.16195);
\fill[white] (0.07417,0.00000) rectangle (0.15050,0.16195);
\fill[white] (0.15320,0.00000) rectangle (0.27281,0.16195);
\fill[white] (0.27555,0.00000) rectangle (0.32195,0.16195);
\fill[white] (0.32466,0.00000) rectangle (0.35774,0.16195);
\fill[white] (0.36045,0.00000) rectangle (0.45386,0.16195);
\fill[white] (0.45660,0.00000) rectangle (0.51964,0.16195);
\fill[white] (0.52232,0.00000) rectangle (0.56373,0.16195);
\fill[white] (0.56641,0.00000) rectangle (0.74608,0.16195);
\fill[white] (0.74884,0.00000) rectangle (0.86845,0.16195);
\fill[white] (0.87119,0.00000) rectangle (0.91925,0.16195);
\fill[white] (0.92195,0.00000) rectangle (0.97167,0.16195);
\end{scope}
\end{tikzpicture}
\end{sourceblock}
\begin{sourceblock}
\begin{tikzpicture}
\node[inner sep=0,anchor=south west] (image) at (0,0) {\includegraphics[page=418,trim={53.00bp 333.52bp 53.87bp 270.64bp},clip,width=0.92\textwidth,max height=0.38\textheight]{Ferziger4th.pdf}};
\begin{scope}[x={(image.south east)},y={(image.north west)}]
\fill[white] (0.14860,0.77654) rectangle (0.18821,0.98064);
\fill[white] (0.18677,0.77654) rectangle (0.19741,0.91497);
\fill[white] (0.24003,0.77654) rectangle (0.31943,0.98064);
\fill[white] (0.31798,0.77654) rectangle (0.34337,0.98064);
\fill[white] (0.20641,0.68119) rectangle (0.23459,0.86786);
\fill[white] (0.15883,0.56276) rectangle (0.18809,0.75363);
\fill[white] (0.26770,0.54969) rectangle (0.31293,0.75363);
\fill[white] (0.40063,0.26363) rectangle (0.41916,0.45015);
\fill[white] (0.38710,0.01936) rectangle (0.43233,0.22330);
\end{scope}
\end{tikzpicture}
\end{sourceblock}
由于张量是对称的，因此只需要求解六个方程。右侧的前两项是产生式项，不需要近似或建模。

\begin{sourceblock}
\begin{tikzpicture}
\node[inner sep=0,anchor=south west] (image) at (0,0) {\includegraphics[page=418,trim={53.00bp 246.26bp 53.87bp 380.41bp},clip,width=0.92\textwidth,max height=0.38\textheight]{Ferziger4th.pdf}};
\begin{scope}[x={(image.south east)},y={(image.north west)}]
\fill[white] (0.03528,0.79478) rectangle (0.08668,1.00000);
\fill[white] (0.08938,0.79478) rectangle (0.15573,1.00000);
\fill[white] (0.15844,0.79478) rectangle (0.22980,1.00000);
\fill[white] (0.23248,0.79478) rectangle (0.28220,1.00000);
\fill[white] (0.35657,0.03040) rectangle (0.36722,0.24778);
\fill[white] (0.36577,0.03040) rectangle (0.37642,0.24778);
\end{scope}
\end{tikzpicture}
\end{sourceblock}
这通常称为压力应变项。它在雷诺应力张量的各个分量之间重新分配湍流动能，但不改变总动能。下一个术语是：

\begin{sourceblock}
\includegraphics[page=418,trim={53.00bp 158.58bp 53.87bp 475.82bp},clip,width=0.92\textwidth,max height=0.38\textheight]{Ferziger4th.pdf}
\end{sourceblock}
这是耗散张量。最后一项是：

\begin{sourceblock}
\includegraphics[page=418,trim={53.00bp 111.77bp 53.87bp 536.29bp},clip,width=0.92\textwidth,max height=0.38\textheight]{Ferziger4th.pdf}
\end{sourceblock}
通常称为湍流扩散。

耗散项、压力应变项和湍流扩散项无法根据方程中的其他项精确计算，因此必须进行建模。耗散项最简单且最常见的模型将其视为各向同性：

\begin{sourceblock}
\includegraphics[page=419,trim={53.00bp 535.34bp 53.87bp 102.78bp},clip,width=0.92\textwidth,max height=0.38\textheight]{Ferziger4th.pdf}
\end{sourceblock}
这意味着耗散方程必须与雷诺应力方程一起求解。通常，这被视为 k – ε 模型中使用的耗散方程。已经提出了更复杂（因此更复杂）的模型。

压力-应变项最简单的模型是假设该项的功能是尝试使湍流更加各向同性。这种模式并未取得巨大成功。最成功的模型基于将压力应变项分解为“快速”部分，其中涉及湍流和平均流动梯度之间的相互作用，以及仅涉及湍流量之间的相互作用的“慢”部分（该部分通常使用返回各向同性项进行建模）。参见朗德等人。 ( 1975 ) 或 Pope ( 2000 ) 了解更多细节。

湍流扩散项通常使用梯度扩散类型的近似来建模。在最简单的情况下，假设扩散率是各向同性的，并且只是前面讨论的模型中使用的涡流粘度的倍数。近年来，提出了各向异性和非线性模型。同样，这里没有尝试详细讨论它们。

在三维中，雷诺应力模型除了平均流动方程外还需要求解七个偏微分方程。当需要预测标量时，还需要更多方程。这些方程的求解方式与 k – ε 方程类似。唯一的额外问题是，当雷诺平均纳维-斯托克斯方程与雷诺应力模型一起求解时，它们甚至比用 k – ε 方程获得的方程更硬，并且在求解时需要更加小心，并且计算通常收敛得更慢。应用中通常的方法是首先使用 k – ε 湍流模型计算流动，根据涡粘假设估计雷诺应力分量的初始值，然后使用雷诺应力模型继续计算。这通常会有所帮助，因为与使用雷诺应力模型和简单的变量初始化开始计算的情况相比，它提供了更合理的所有变量的起始字段。同时，通过这种方式可以获得两个湍流模型的解，并且两个解的比较也代表了有用的信息。

雷诺应力模型通过消除涡粘性，能够解决各向异性，但需要求解额外的微分方程，这增加了成本。但是，毫无疑问，雷诺应力模型可以比二元方程模型更正确地表示湍流现象（参见 Hadži ́c 1999，了解一些说明性示例）。对于一些 k – ε 模型表现不佳的流动（例如，旋流、具有驻点或停滞线的流动、具有强曲率和与曲面分离的流动等），获得了特别好的结果。

哪种模型最适合哪种流量（预计没有一个模型适合所有流量）尚不清楚；然而，Leschziner（2010）很好地描述了可用的内容，而 Hanjali ́c（2004）则涵盖了许多细节。由于人们无法始终确定湍流模型的选择和/或性能，因此确保解之间的任何差异都是由于模型差异而不是数值误差所致，这一点非常重要。这就是本书强调数值准确性的原因之一。它的重要性无论怎么强调都不过分，需要持续关注。

正如上文第 1 节所述。 10.3.3.5 对于 LES，显式代数雷诺应力模型 (EARSM) 是一种有吸引力的替代方案，其中雷诺应力模型的微分方程负载减少，但保留了解析各向异性和其他重要物理场的基本能力。 Wallin 和 Johansson (2000) 的 EARSM 旨在研究不可压缩和可压缩旋转流，通过平均应变率和平均旋转率张量描述各向异性。它使用 Rodi (1976) 假设，雷诺应力的对流和扩散可以用雷诺应力、湍流能量产生和湍流能量耗散来表示。这导致他们建立了一个非线性隐式模型，并对其进行简化以实现 EARSM。该方法需要湍流动能和耗散率的输运方程，以及线性方程组的解和一个产生与耗散之比的非线性代数方程。该 EARSM 正确地解释了旋转，并且已被证明优于经典的基于涡流粘度的模型。

上述所有模型都有许多版本。这些修改旨在纠正基本模型的各种缺陷，例如没有考虑湍流的各向异性、对停滞或分离的影响建模不充分、不利或有利的压力梯度、流线曲率、固体壁附近的湍流阻尼或从层流到湍流的过渡。这里详细讨论所有这些细节超出了本书的范围；有兴趣的读者可以在上面引用的参考文献中找到足够的信息，特别是 Patel 等人的文献。 （1985）和威尔科克斯（2006）。商业 CFD 代码通常有 20 多个湍流模型版本可供选择。几乎所有模型都有两种变体，具体取决于壁边界条件的处理方式：“高 Re”版本，假设所有计算点都在边界层的湍流部分内；“低 Re”版本，假设近壁网格解析粘性子层并提供该区域的湍流阻尼。墙壁处理将在下一节中更详细地描述。

10.3.5.5 RANS 计算的边界条件
RANS 计算中入口、出口和对称平面边界条件的应用与层流的情况相同，因此这里不再重复任何细节；见第 9.11 节 了解更多信息。然而，值得注意的是，在流入边界处，k 和 ε 通常是未知的；如果可用，则当然应该按照前面章节中针对通用标量变量描述的相同方式使用已知值。如果 k 未知，通常根据假定的湍流强度 I t = u ′ 2 / u 进行估计。例如，通过指定 It = 0.01（低湍流强度 1\%）并假设 u ′ 2 x = u ′ 2 y = u ′ 2 z = I 2 t u 2，可得到 k = 3
2 I 2 t u 2 = 1.5 × 10 − 4 u 2。应选择 ε 的值，以便从方程导出的长度尺度。 ( 10.45 ) 大约是剪切层宽度或域尺寸的十分之一。如果在入口处测量雷诺应力和平均速度，则可以使用局部平衡的假设来估计 ε；这导致（在横截面 x = const 中）：

\begin{sourceblock}
\includegraphics[page=421,trim={53.00bp 471.09bp 53.87bp 166.76bp},clip,width=0.92\textwidth,max height=0.38\textheight]{Ferziger4th.pdf}
\end{sourceblock}
入口处的速度场本身通常并不精确（特别是在内部流动的情况下）。流速通常是已知的，并且当入口横截面已知时，可以计算平均速度。最简单的近似是规定入口处的恒定速度。如果可以对入口边界上的速度变化进行合理的近似，那么就应该这样做。例如，如果入口代表穿过管道、管道或环空的截面，则可以为这种几何形状中完全发展的流动规定速度分布。通过使用单层单元，在入口和出口处应用具有规定流量的周期性边界条件，可以轻松计算出完全发展的流动；然后，这种计算的结果可以用于规定更复杂的解域入口处的变量值。

如果需要在入口处近似变量值的分布，则应尝试将入口边界移动到感兴趣区域的上游尽可能远的地方。当上游流路的几何形状不可用时，可以作为一种近似，在上游方向上拉伸入口横截面，以允许在入口边界处进行的近似在流到达感兴趣区域时发展成合理的轮廓。大多数商业网格生成工具都允许在入口和出口处进行此类挤压。还建议将出口横截面移动到距离感兴趣区域尽可能远的下游，以便最小化在那里进行的近似对解域重要部分中的流动的影响。另外，通常采用较粗的网格或逐渐降低对流项向出口的近似阶数，以增加数值扩散并避免扰动在出口边界的反射。

在描述湍流 RANS 计算的壁边界条件之前，我们首先认识到，紧邻壁的湍流影响相对较小，并且流动基本上是层流。壁边界层的这一部分称为粘性子层，实验和 DNS 都表明，平行于壁的速度分量随壁距离线性变化（我们在这里用壁法线方向的局部坐标 n 表示；见图 9.2.1）。如果该子层由数值网格解析，则动量方程的边界条件与层流中的相同；见第 9.11 节 详情。我们记得法向粘性应力 τ nn 等于 0，因为壁面法向速度分量 v n 对 n 的导数在壁面处必须为零，并且剪切应力等于分子粘度和壁面平行速度分量 v t 对 n 的导数的乘积：

\begin{sourceblock}
\begin{tikzpicture}
\node[inner sep=0,anchor=south west] (image) at (0,0) {\includegraphics[page=422,trim={53.00bp 522.61bp 53.87bp 108.23bp},clip,width=0.92\textwidth,max height=0.38\textheight]{Ferziger4th.pdf}};
\begin{scope}[x={(image.south east)},y={(image.north west)}]
\fill[white] (0.34391,0.40935) rectangle (0.38863,0.76742);
\fill[white] (0.20208,0.21133) rectangle (0.24174,0.56969);
\fill[white] (0.24773,0.24193) rectangle (0.27591,0.56969);
\fill[white] (0.27943,0.23173) rectangle (0.31991,0.56969);
\fill[white] (0.78580,0.24193) rectangle (0.79886,0.56969);
\fill[white] (0.90923,0.23173) rectangle (1.00000,0.55921);
\fill[white] (0.73489,0.00000) rectangle (0.77558,0.21671);
\end{scope}
\end{tikzpicture}
\end{sourceblock}
由于 v t 的分布在壁附近是线性的，因此即使是最简单的半电池单侧向前或向后差异也是准确的。

问题是，如果要解决高雷诺数 3D 流情况下的粘性子层，靠近壁的单元必须非常薄。一方面，这需要靠近壁的许多棱柱形细胞层，另一方面，这使得细胞的纵横比非常高。由于离散拉普拉斯算子产生的系数（传输方程中的扩散项，压力或压力校正方程中的相应项）与纵横比的平方成正比，这意味着壁法线方向上的系数比其他方向上的系数大几个数量级。这使得方程变得僵硬，对计算精度要求更高，并且使离散方程的求解更加困难。此外，如果壁是弯曲的，则靠近壁的薄单元可能会变得过于扭曲，除非网格也在切向方向上也被充分细化。这就是为什么带有解析粘性子层的网格的所谓低Re湍流模型仅用于中等雷诺数的流动（例如，用于与模型规模的实验进行比较）。对于雷诺数非常高的流动，例如船舶、飞机或其他大型物体周围的流动，我们需要一种替代的、更便宜的方法。

工程师总是试图找到广义的标度定律，并且在大约 100 年前就认识到，如果使用所谓的剪切速度 u τ （也称为摩擦速度 ）进行标度计算，则不同雷诺数下边界层上的速度剖面可能会崩溃：

\begin{sourceblock}
\includegraphics[page=422,trim={53.00bp 212.70bp 53.87bp 418.67bp},clip,width=0.92\textwidth,max height=0.38\textheight]{Ferziger4th.pdf}
\end{sourceblock}
其中 τ wall 是壁剪应力的大小（如果局部壁切向坐标 t 与剪应力矢量方向一致，则 τ wall = | τ nt | ）。速度与 u τ 成比例以获得 u + 并根据距壁的无量纲距离 n + 绘制，如下所示：

\begin{sourceblock}
\begin{tikzpicture}
\node[inner sep=0,anchor=south west] (image) at (0,0) {\includegraphics[page=422,trim={53.00bp 130.48bp 53.87bp 506.74bp},clip,width=0.92\textwidth,max height=0.38\textheight]{Ferziger4th.pdf}};
\begin{scope}[x={(image.south east)},y={(image.north west)}]
\fill[white] (0.00000,0.86445) rectangle (0.04580,1.00000);
\fill[white] (0.04848,0.86445) rectangle (0.07826,1.00000);
\fill[white] (0.08093,0.86445) rectangle (0.18322,1.00000);
\end{scope}
\end{tikzpicture}
\end{sourceblock}
\begin{sourceblock}
\includegraphics[page=423,trim={53.00bp 451.25bp 53.87bp 52.00bp},clip,width=\textwidth,max height=0.38\textheight]{Ferziger4th.pdf}
\par\vspace{0.45\baselineskip}
\small 图 10.25 从 Lee 和 Moser (2015)、El Khoury 等人的 DNS 数据中获得的平面通道、圆形横截面管道和平板上边界层中湍流边界层的归一化速度剖面。 ( 2013 ) 以及 Schlatter 和 Örlü ( 2010 )
\end{sourceblock}
\begin{center}\small 图 10.25 显示了不同雷诺数下三种不同湍流的归一化速度分布（这些是 x 方向上的二维流，因此 n + = y + ）；它们确实彼此叠在一起，只是离墙较远。可以识别三个不同的区域：除了已经提到的紧邻壁的粘性子层之外，还有一个具有对数变化的重要部分，以及线性和对数区域之间的缓冲层。当雷诺数增加时，对数范围扩展到更高的 n + 值；参见例如 Wosnik 等人。 (2000) 或 Lee 和 Moser (2015)。所谓的墙的对数定律已通过实验和 DNS 数据得到证实；然而，最新的研究工作表明，该法律并不像人们曾经想象的那么普遍。显然存在重要的变化，具体取决于雷诺数和流域的几何形状，但我们不会详细讨论这些细节；感兴趣的读者可以查阅有关该主题的最新文献，例如 Smits 等人。 (2011) 以及 Smits 和 Marusic (2013)。由于本节的目的是演示如何将湍流模型实现到 CFD 代码中，因此我们将坚持使用经典方法来描述对数范围内的速度分布，如下所示：\end{center}
\begin{sourceblock}
\begin{tikzpicture}
\node[inner sep=0,anchor=south west] (image) at (0,0) {\includegraphics[page=423,trim={53.00bp 157.65bp 53.87bp 480.47bp},clip,width=0.92\textwidth,max height=0.38\textheight]{Ferziger4th.pdf}};
\begin{scope}[x={(image.south east)},y={(image.north west)}]
\fill[white] (0.45260,0.54461) rectangle (0.47239,0.95717);
\fill[white] (0.37350,0.29907) rectangle (0.41122,0.78266);
\fill[white] (0.41624,0.30835) rectangle (0.44442,0.72127);
\end{scope}
\end{tikzpicture}
\end{sourceblock}
其中 κ 称为冯卡门常数，B 是经验常数。通常取κ=0.41 且 B ≈ 5.2 但这些常数也不是严格通用的。使用 Re = 2 处平面通道流的 DNS。 5 × 10 5 （基于平均速度和通道宽度），Lee 和 Moser (2015) 得出较低的值，即 κ = 0.384 且 B = 4.27；对于粗糙的墙壁，获得的 B 值较小，这意味着如图 10.25 所示的轮廓向下移动。

\begin{sourceblock}
\includegraphics[page=424,trim={53.00bp 453.25bp 53.87bp 52.00bp},clip,width=\textwidth,max height=0.38\textheight]{Ferziger4th.pdf}
\par\vspace{0.45\baselineskip}
\small 图10.26 平面通道、管道和平板边界层流中跨边界层湍流动能标准化产生和耗散的变化（数据来源见图10.25）
\end{sourceblock}
如果第一计算点可以放置在对数范围内而不是粘性子层中，则可以大大减少计算量。我们需要壁面的速度梯度来计算壁面剪应力，但无法从高阶多项式中获得足够的精度，就像层流的情况一样。然而，根据对数定律和一些其他假设，可以推导出壁面剪应力与剖面对数部分中某点的速度之间的关系。这就是所谓的墙函数的作用。

Launder 和 Spalding (1974) 提出了所谓的“高 Re 壁函数”。 20 除了对数速度分布之外，还做出了两个假设：

\begin{itemize}
\item 假定流动处于局部平衡状态，这意味着湍流的产生和消散几乎相等。
\item 壁面和第一个计算点之间的总剪应力（即粘性和湍流贡献之和）恒定，并且等于壁面剪应力 τ wall。
\end{itemize}
\begin{center}\small 图 10.26 显示了从 DNS 数据获得的平面通道、管道和平板边界层流中穿过边界层的湍流动能标准化产生和耗散的变化。该数据以及大量测量确实表明，在对数区域内，产生和耗散确实接近平衡，至少在这些相对简单的流动中。第二个假设也得到了 DNS 和测量数据的支持：在速度线性变化的壁面附近，粘性应力显然是恒定的，并且在整个壁面上。\end{center}
\footnotetext[20]{请注意，“高 Re”和“低 Re”的表述与特定流动问题的实际雷诺数无关，它们与计算点距离墙壁的距离有关。}
\begin{sourceblock}
\includegraphics[page=425,trim={53.00bp 520.24bp 53.87bp 52.00bp},clip,width=\textwidth,max height=0.38\textheight]{Ferziger4th.pdf}
\par\vspace{0.45\baselineskip}
\small 图 10.27 雷诺数为 5 × 10 5 的河道流中壁附近的速度矢量（基于平均速度和河道壁之间的距离），使用壁函数和靠近壁的较粗网格（左）和完全解析的边界层（右）计算；向量之间的细线代表 CV 边界（网格线）
\end{sourceblock}
在缓冲层中，粘性部分减少，而湍流部分增加，使总和几乎恒定。

在这些假设下，可以证明以下关系成立：

\begin{sourceblock}
\begin{tikzpicture}
\node[inner sep=0,anchor=south west] (image) at (0,0) {\includegraphics[page=425,trim={53.00bp 391.96bp 53.87bp 249.32bp},clip,width=0.92\textwidth,max height=0.38\textheight]{Ferziger4th.pdf}};
\begin{scope}[x={(image.south east)},y={(image.north west)}]
\fill[white] (0.40728,0.09614) rectangle (0.43774,0.58648);
\fill[white] (0.44526,0.13234) rectangle (0.47344,0.59735);
\fill[white] (0.47684,0.12148) rectangle (0.53507,0.66653);
\end{scope}
\end{tikzpicture}
\end{sourceblock}
从这个方程和方程。 ( 10.63 ) 我们可以推导出一个将壁上方第一个计算点处的速度与壁剪切应力 21 连接起来的表达式：

\begin{sourceblock}
\begin{tikzpicture}
\node[inner sep=0,anchor=south west] (image) at (0,0) {\includegraphics[page=425,trim={53.00bp 321.84bp 53.87bp 314.09bp},clip,width=0.92\textwidth,max height=0.38\textheight]{Ferziger4th.pdf}};
\begin{scope}[x={(image.south east)},y={(image.north west)}]
\fill[white] (0.28962,0.24793) rectangle (0.32256,0.66865);
\fill[white] (0.32761,0.28600) rectangle (0.35579,0.66865);
\fill[white] (0.35931,0.27706) rectangle (0.40926,0.71731);
\end{scope}
\end{tikzpicture}
\end{sourceblock}
其中 E = eκ B。这允许根据近壁 CV 中心的壁平行速度分量和湍流动能的值计算壁剪切应力。壁剪应力和面面积的乘积产生了一个力，当投影到笛卡尔坐标方向上时，该力为近壁 CV 的离散动量方程提供了必要的贡献，与第 2 节中针对层流所描述的方式相同。 9.11。

\begin{center}\small 图 10.27 显示了雷诺数为 5 × 10 5 时平面通道流底壁附近的速度矢量，该矢量是使用较粗网格上的壁函数并使用解析粘性子层的精细网格来计算的。通道宽 100 毫米，在壁函数的情况下，靠近壁的单元厚度为 0.315 毫米（单元中心 n + = 31）。靠近墙壁10毫米的距离有15个棱镜层，扩大了1.1倍；对于“低Re”方法，相同距离上有 40 个棱镜层，膨胀系数也相同，并且靠近壁的第一个单元厚度为 0.013 毫米（比壁函数薄 23.4 倍；在单元中心 n + = 1 . 32）。显然，在平面通道壁的情况下，非常薄的单元没有问题，但如果通道壁是弯曲的，则还需要在切向方向上进行细化以避免单元过度翘曲\end{center}
\footnotetext[21]{如果墙壁是粗糙的，我们可以使用第 21 节中导出的条件。 10.3.3.3 并在等式中给出。 （10.27）和（10.29）。}
（有关电网质量问题的讨论，请参见第 12 章）。该图强调了使用壁函数时壁与第一个计算点之间速度的显著变化。

在 CFD 的许多实际应用中，很难将计算点始终保持在粘性子层内靠近壁面的位置（对于“低 Re 壁面处理”）或保持在对数范围内（对于“高 Re 壁面函数”）。流动分离、停滞和重新附着区域的存在不可避免地导致壁面剪切应力的巨大变化以及实际上为零的地方。壁函数所基于的假设可能无效，或者计算点可能局部落在粘性子层之外。因此，许多研究人员试图开发更通用的湍流壁边界条件。 Jakirlíc和Jovanovíc(2010)提出了一种允许第一个计算点靠近粘性子层边缘的方法，从而放宽了标准“低Re壁处理”对近壁网格细度的要求。其他研究人员使用雷查特定律（Reichardt 1951）的修改来提供所谓的“盟友+壁处理”：如果网格足够细，则假设线速度分布，如果近壁计算点位于对数层内，则应用标准壁函数。当计算点介于两者之间时，两种方法就会混合。我们在这里不再赘述；在第 13 章中 这个问题将从实际应用的角度再次讨论。

湍流动能的分布及其耗散率通常在壁附近比平均速度分布更达到峰值。这些峰很难捕捉；人们可能应该对湍流量使用比平均流动更精细的网格，但很少这样做。如果对所有量使用相同的网格，则湍流量的分辨率可能不足，并且如果使用高阶方案，则解有可能包含摆动，这可能导致该区域中这些量的负值。通过将中心差分格式与 k 和 ε 方程中的对流项的低阶迎风离散化进行局部混合，可以避免这种可能性。当然，这会降低这些量的计算精度，但如果网格不够精细，则这是必要的。

模型方程的壁边界条件也需要特别注意。一种可能性是当网格解析粘性子层时，精确地求解方程直到壁面。在k-ε模型中，在壁面处设置k=0是合适的，但壁面处的耗散不为零；相反，我们可以使用以下条件：

\begin{sourceblock}
\begin{tikzpicture}
\node[inner sep=0,anchor=south west] (image) at (0,0) {\includegraphics[page=426,trim={53.00bp 137.36bp 53.87bp 493.49bp},clip,width=0.92\textwidth,max height=0.38\textheight]{Ferziger4th.pdf}};
\begin{scope}[x={(image.south east)},y={(image.north west)}]
\fill[white] (0.22623,0.24199) rectangle (0.24457,0.56957);
\fill[white] (0.24896,0.24199) rectangle (0.27714,0.56957);
\fill[white] (0.28066,0.24199) rectangle (0.29964,0.56957);
\fill[white] (0.76165,0.24199) rectangle (0.77471,0.56957);
\fill[white] (0.90938,0.23151) rectangle (1.00000,0.55937);
\fill[white] (0.71074,0.00000) rectangle (0.75143,0.21649);
\end{scope}
\end{tikzpicture}
\end{sourceblock}
如上所述，有必要在壁附近修改模型本身，因为它的存在会抑制湍流波动并导致实际层流粘性子层的存在。

当使用壁函数类型边界条件时，通常将 k 通过壁的扩散通量视为零，从而产生 k 的法向导数为零的边界条件。耗散边界条件是通过假设平衡导出的，即近壁区域的产生和耗散平衡。壁区域的产量计算公式为：

\begin{sourceblock}
\includegraphics[page=427,trim={53.00bp 511.32bp 53.87bp 126.53bp},clip,width=0.92\textwidth,max height=0.38\textheight]{Ferziger4th.pdf}
\end{sourceblock}
这是方程主要项的近似值。 （10.44）；它在壁附近有效，因为该区域的剪应力几乎恒定，并且平行于壁的方向上的速度导数远小于垂直于壁的方向上的导数。上式中所需的单元中心的速度导数可以从对数速度剖面（ 10.63 ）导出：

\begin{sourceblock}
\begin{tikzpicture}
\node[inner sep=0,anchor=south west] (image) at (0,0) {\includegraphics[page=427,trim={53.00bp 402.81bp 53.87bp 225.29bp},clip,width=0.92\textwidth,max height=0.38\textheight]{Ferziger4th.pdf}};
\begin{scope}[x={(image.south east)},y={(image.north west)}]
\fill[white] (0.34126,0.43559) rectangle (0.38223,0.76814);
\end{scope}
\end{tikzpicture}
\end{sourceblock}
当使用上述近似时，ε 的离散方程不适用于靠近壁的控制体积；相反，ε 在 CV 中心处设置等于：

\begin{sourceblock}
\includegraphics[page=427,trim={53.00bp 330.95bp 53.87bp 303.34bp},clip,width=0.92\textwidth,max height=0.38\textheight]{Ferziger4th.pdf}
\end{sourceblock}
该表达式源自等式。 ( 10.45 ) 使用长度尺度的近似值

\begin{sourceblock}
\begin{tikzpicture}
\node[inner sep=0,anchor=south west] (image) at (0,0) {\includegraphics[page=427,trim={53.00bp 277.50bp 53.87bp 358.25bp},clip,width=0.92\textwidth,max height=0.38\textheight]{Ferziger4th.pdf}};
\begin{scope}[x={(image.south east)},y={(image.north west)}]
\fill[white] (0.00000,0.73149) rectangle (0.06403,1.00000);
\end{scope}
\end{tikzpicture}
\end{sourceblock}
在用于导出壁函数的条件下，它在壁附近有效。

需要注意的是，当第一个网格点位于对数区域内时，即当 n +
P > 30。分离流中出现问题；在再循环区域内，特别是在分离区域和再附着区域内，不满足上述条件。通常，壁函数在这些区域中无效的可能性被忽略，并且它们被应用到任何地方。然而，如果在大部分实体边界上违反上述条件，可能会导致严重的建模错误。替代的墙功能和所谓的“ally +”模型已经被提出，并且在大多数商业代码中都可用；它们不能完全消除建模错误，但可以帮助最大限度地减少错误。

以下三篇论文可用于构建有用的墙壁处理。首先，Durbin（2009）对壁面应用的湍流建模进行了深入的回顾，包括一些约束、壁面函数和椭圆松弛模型。其次，Popovac 和 Hanjali ́c (2007) 展示了湍流和传热的壁边界条件。第三，比拉德等人。 ( 2015 ) 引入了自适应壁函数的稳健公式，用于近壁椭圆混合涡粘模型背景下的传热计算。

在远离墙壁的计算边界（远场或自由流边界）处，可以使用以下边界条件：

\begin{itemize}
\item 如果周围的流动是湍流：
\end{itemize}
\begin{sourceblock}
\includegraphics[page=428,trim={53.00bp 503.91bp 53.87bp 132.97bp},clip,width=0.92\textwidth,max height=0.38\textheight]{Ferziger4th.pdf}
\end{sourceblock}
\begin{itemize}
\item 在免费流中：
\end{itemize}
\begin{sourceblock}
\includegraphics[page=428,trim={53.00bp 449.24bp 53.87bp 187.91bp},clip,width=0.92\textwidth,max height=0.38\textheight]{Ferziger4th.pdf}
\end{sourceblock}
雷诺应力的边界条件更加复杂。我们不会在这里详细讨论，但请注意，一般来说，模型必须提供在每种边界类型附近求解的每个变量的变化的近似值。在大多数情况下，条件会减少到指定的边界值（狄利克雷条件）或边界法线方向上的指定梯度（诺依曼条件）。这些条件可以使用层流通用标量变量描述的方法实现到边界附近单元的离散方程中；见第 9.11 节 了解更多信息。

\subsection*{10.3.6 RANS 应用示例：Re = 500,000 时围绕球体的流动}\phantomsection\addcontentsline{toc}{subsection}{10.3.6 RANS 应用示例：Re = 500,000 时围绕球体的流动}
我们之前注意到，RANS 模型不能很好地预测亚临界雷诺数下球体周围的流动，其中层流分离之后会出现湍流尾流； LES 非常适合此类流程。在超临界雷诺数下，LES 需要非常精细的网格和小的时间步长，使得模拟成本非常高。对于此类流，通常使用 RANS 或 URANS 方法。我们在此考虑雷诺数 Re = 500,000 时围绕球体的流动。由于雷诺数不太高（就像车辆、轮船或飞机周围的流动一样），我们可以解析边界层，从而创建一个具有 15 层棱镜层的网格，第一层靠近壁的厚度为 0.01 毫米；外部棱镜层尾流中的网格尺寸为 0.4375 mm (D / 140)。使用与 LES 情况相同类型的网格（具有局部细化的修剪笛卡尔），雷诺数低 10 倍，并且使用相同的商业流动求解器 (STAR-CCM+)。球体周围及其尾流的网格如图 10.28 所示。解域几何形状也相同 - 直径 D = 61 的光滑球体。 4 mm 由直径 d = 8 mm 的杆固定在横截面为 300 × 300 mm 的风洞中。在模拟中，使用了与在模拟中相同的速度和粘度。

\begin{sourceblock}
\includegraphics[page=429,trim={53.00bp 485.25bp 53.87bp 52.00bp},clip,width=\textwidth,max height=0.38\textheight]{Ferziger4th.pdf}
\par\vspace{0.45\baselineskip}
\small 图 10.28 球体周围及其尾流的粗网格，用于 Re = 500,000 时流量的 RANS 计算（每次细化网格时，网格间距减小 1.5 倍）
\end{sourceblock}
在之前的情况下，只需将密度增加10倍（从1.204至12.04 kg/m 3 ）即可获得高10倍的雷诺数。

球体周围的流动有几个 RANS 模型通常不能很好处理的特征；最重要的是与光滑、弯曲的表面分离。我们在这里尝试了商业流动求解器中可用的许多版本中的四种 RANS 模型：（i）k - ε 模型的标准低 Re 版本； (ii) k - ω 模型的 SST 版本； (iii) lag-EB k - ε 模型； (iv) 雷诺应力模型。只有lag-EB k - ε模型产生的结果与文献中的实验数据接近；所有其他模型都会导致阻力过高和再循环区过大。

lag-EB k - ε 模型是代码中 k - ε 模型类中相对较新的补充；详细信息可以参见Lardeau（2018）。名称中的“滞后”意味着该模型考虑了平均应力和应变并不总是一致（一个滞后于另一个）这一事实； “EB”代表椭圆混合。除了 k 和 ε 之外，还求解了两个附加变量的传输方程。该模型的计算在三个网格上进行；对于所有三个网格来说，靠近墙壁的第一个棱镜层是相同的，所有棱镜层的总厚度也是如此；粗网格有10个棱柱层，中网格有12个棱柱层，最细网格有15个棱柱层。在棱柱层之外，从细网格到中网格，从中网格到粗网格，网格间距增加了1.5倍；上面给出了细网格值。粗网格有 728,923 个控制体，中网格有 2,131,351 个控制体，细网格有 6,591,260 个控制体。计算得出的阻力值为：粗网格为 0.0679，中网格为 0.0638，细网格为 0.0619。 Richardson 外推法得出与网格无关的估计值为 0.0604，该值接近文献中的数据。其他湍流模型仅与最精细的网格一起使用。

低 Re k - ε 模型不会收敛到稳态解；当在瞬态模式下继续计算时，阻力系数在 0.108 和 0.146 之间振荡，这大大高于文献报道的实验值（例如 Achenbach 1972）。 SST k - ω 和 lag-EB k - ε 模型收敛到接近稳态的解；图 10.29 显示了计算的残差

\begin{sourceblock}
\includegraphics[page=430,trim={53.00bp 464.48bp 53.87bp 52.00bp},clip,width=\textwidth,max height=0.38\textheight]{Ferziger4th.pdf}
\par\vspace{0.45\baselineskip}
\small 图 10.29 使用 lag-EB k - ε 模型在细网格上计算 Re = 500,000 时球体周围流动的残差历史
\end{sourceblock}
使用 lag-EB k - ε 模型。残差下降了近五个数量级，然后保持在该水平，有小幅波动；阻力值在五个最高有效数字上不再发生变化。在瞬态模式下继续计算不会导致流量发生任何显著变化；出于实际目的，可以将其视为收敛到稳态。使用SST k - ω模型计算得到的阻力系数约为实验值的两倍（0.146）。雷诺应力模型的计算是从 lag-EB k - ε 模型获得的解开始的；收敛是振荡的并且非常缓慢，但是阻力系数变化不大并且趋向于平均值约。 0.108，也远高于实验结果。

Leder (1992) 给出了雷诺数在 150,000 到 300,000 之间的球体后面的再循环区的长度，常数约为 0.2 D。所有计算都预测更长的再循环区；最小值是使用 lag-EB k - ε 模型（略长于 0.4 D ）计算得到的。图 10.30 显示了 SST k - ω 和 lag-EB k - ε 模型在 y = 0 和 z = 0 处的剖面中的流型。可以看出，虽然解实际上是稳定的，但流动不是轴对称的；特别是，使用 SST k - ω 模型获得的解是高度不对称的。如前所述，在其他数值和实验研究中也观察到了不对称尾流（Constantinescu 和 Squires 2004；Taneda 1978）。

尽管流域的几何形状非常简单，但使用 RANS 模型来预测球体周围的流动非常困难。正如以下章节所示，在大多数工业相关流程中，实验和 RANS 计算之间的一致性比本测试案例要好得多（例如，船舶阻力通常预测在实验值的 2\% 以内；在许多情况下，一致性甚至更好，并且只有很少情况下差异会大于 2\%）。原因是，在球体的情况下，RANS 模型不太擅长的流特征占主导地位；在更复杂的几何形状中，RANS 模型可以更好地预测其他特征。

\begin{sourceblock}
\includegraphics[page=431,trim={53.00bp 430.48bp 53.87bp 52.00bp},clip,width=\textwidth,max height=0.38\textheight]{Ferziger4th.pdf}
\par\vspace{0.45\baselineskip}
\small 图 10.30 使用 SST k - ω 模型（左）和 lag-EB k - ε 模型（右）计算 y = 0（上）和 z = 0（下）部分的流型
\end{sourceblock}
\subsection*{10.3.7 超大涡模拟/TRANS/DES}\phantomsection\addcontentsline{toc}{subsection}{10.3.7 超大涡模拟/TRANS/DES}
流动模拟的目标通常是以最小的成本获得有关流动的特定属性的信息。明智的做法是使用最简单的工具来提供所需的结果，但提前知道每种方法的效果并不容易。我们正在处理一个层次结构，从简单到复杂，即稳定 RANS、不稳定 RANS、VLES、LES 和 DNS。我们在该集合中添加了 LES 和 URANS 的有趣混合，称为分离涡模拟 (DES)，有多种变体，缩写为 DDES 或 IDDES。

显然，如果 RANS 方法成功，则没有理由使用 LES 等。另一方面，当 RANS 变体不起作用时，尝试 URANS 或 LES 类型的模拟可能是个好主意。然而，请记住，URANS 是一个系综结构，因此所有湍流能量可能都在建模尺度内。另一方面，LES 可以解析大部分湍流能量。

使用模拟方法的一种方法是执行“构建块”流的不稳定 RANS、LES 和/或 DNS，这些流在结构上与实际感兴趣的流类似。根据结果​​，可以验证和改进可应用于更复杂流程的 RANS 模型。 RANS 计算就可以成为日常工具。仅当设计发生重大变化时才需要执行 LES。然而，在某些模拟领域，LES 的使用已占主导地位，例如用于研究和天气预报的中尺度和大气边界层模拟，或工程中的气动声学模拟。

看来我们必须要么使用 RANS（价格实惠），要么使用 LES（更准确但相当昂贵）。人们很自然地会问，是否有一些方法可以同时提供 RANS 和 LES 的优点，同时避免其缺点？

一种直接的方法是检查第 1 节中给出的 LES 域的定义。 10.3.通过增加给定域和流中的网格大小，我们将解析流中的能量减少到标称 80\% 水平以下，并进行所谓的超大涡模拟 (VLES)。因此，在 VLES 中，人们使用 LES 代码来计算非定常流，但采用大网格尺寸和可能更复杂的湍流模型来正确表示建模运动，例如雷诺应力或代数雷诺应力模型（第 10.3.5.4 节）。示例包括：

\begin{itemize}
\item 水力发电厂尾水管中非稳态流动的VLES（Gyllenram 和Nilsson 2006）。
\item 直喷燃烧室中的流动的VLES（Shih 和Liu 2009）；
\item 深层大气中重力波破碎引起的湍流的VLES。 Smolarkiewicz 和 Prusa 2002 采用非振荡时间前移方法 (NFT)（ILES）作为 VLES，因此显式计算可在网格上解析的大相干涡流。
\end{itemize}
另一方面，流过钝体的流动通常会在其尾迹中产生强烈的涡流。这些涡流在流向和翼展方向上对物体产生脉动力，其预测非常重要。其中包括建筑物（风工程）、海洋平台和车辆等上的流动。如果涡流比构成“湍流”的大部分运动足够大，则湍流模型可能只去除较小尺度的运动，并且可以将非周期性流转换为周期性流，这可能会产生重大后果。 Durbin (2002)、Iaccarino 等人。 （2003）和韦格纳等人。 ( 2004 ) 证明，如果其余结构是周期性的，标准集合平均 URANS 方法将实现这一目标；如果不是，长期模拟将导致稳定的流动。亚卡里诺等人。 (2003) 表明 URANS 可以提供与统计上不稳定的流动实验数据的定量和定性一致性。

然而，例如，在存在浮力的主导作用力的情况下，一种称为瞬态 RANS (TRANS) 的结构（Hanjali ́c 和 Kenjereš 2001；Kenjereš 和 Hanjali ́c 1999）也取得了成功。在对热浮力和洛伦兹力驱动的流动中的相干涡流结构进行 TRANS 模拟时，他们使用了将原始运动显式三重分解为 (1) 稳定的长期平均值、(2) 准周期（相干结构）分量和 (3) 随机（随机）波动。大相干涡流结构通过三维系综平均纳维-斯托克斯方程的时间积分得到完全解析。不连贯的部分由 RANS 型闭包建模。在这种特殊情况下，存在强大的浮力，产生了明确的结构，但这些结构只是伪周期性的。然后，作者求助于光谱间隙的概念来允许流的分解。 Hanjali ́c (2002) 显示了瑞利-贝纳德对流的令人信服的结果，Kenjereš 和 Hanjalic (2002) 使用他们的方法显示了地形和热分层对大气边界层昼夜循环的中尺度建模的影响。强大的不稳定力或周期性强迫似乎将 TRANS 模拟与上述 URANS 模拟区分开来（参见 Durbin 2002）。

分离涡模拟 (DES) 创建于 1997 年，旨在预测大规模分离流。最近，斯帕拉特等人。 (2006) 提出了一个名为延迟 DES (DDES) 的新版本，它纠正了原始版本的一些缺陷。在分离流的主要领域，DES 预测的准确性通常优于稳态或非稳态 RANS 方法。在 DES 中，“平均”流动方程在整个流域中是相同的；然而，在近壁区域，湍流模型简化为 RANS 公式，而远离壁区域则应用 LES 子网格模型。过渡距离由算法根据特定规则计算得出。大多数公式使用或源自 Spalart 和 Allmaras (1994) RANS 模型。已经进行了许多有趣的应用，特别是在工程设计方面。举例来说：

1.维斯瓦纳坦等人。 (2008) 使用 DES 来研究攻角下气动体周围的大规模分离流，使用有限体积、非结构化网格代码，在时间和空间上具有二阶精度。允许身体运动。 2. 小南等人。 (2011)将DES和拉格朗日粒子追踪应用于粗糙壁湍流通道流中，以研究壁粗糙度对流动中分散相的影响。

如今，大多数商业 CFD 代码都提供各种计算湍流的方法，从 RANS、URANS 和 DES 到 LES 甚至 DNS。尽管商业代码中的离散化方法通常仅限于空间和时间上的二阶（对流项可能具有更高阶的插值），但在一系列实际应用中已经报告了相当好的结果。

\chapter{可压缩流动}
\section*{11.1 简介}\phantomsection\addcontentsline{toc}{section}{11.1 简介}
可压缩流在空气动力学和涡轮机械等应用中非常重要。在飞机周围的高速流动中，雷诺数极高，湍流效应仅限于薄边界层。阻力由两个部分组成，边界层引起的摩擦阻力和本质上无粘性的压力或形状阻力；还可能存在由于激波而产生的波阻力，只要注意确保遵守热力学第二定律，就可以从无粘性方程计算出波阻力。如果忽略摩擦阻力，则可以使用无粘动量欧拉方程来计算这些流动。

由于可压缩流在民用和军事应用中的重要性，已经开发了许多求解可压缩流方程的方法。其中包括欧拉方程的特殊方法，例如特征方法和许多能够扩展到粘性流的方法。这些方法大多数是专门为可压缩流设计的，当应用于不可压缩流时，效率非常低。其原因可以有多种变化。一是，在可压缩流中，连续性方程包含时间导数，该导数在不可压缩极限处消失。结果，方程在弱压缩性的限制下变得极其僵硬，需要使用非常小的时间步长或隐式方法。该论点的另一个版本是可压缩方程支持具有与其相关的确定速度的声波。当某些信息以流速传播时，两个速度中较大的一个决定了显式方法中允许的时间步长。在低速限制下，对于任何流体速度，人们被迫采取与声速成反比的时间步长；该步长可能比为不可压缩流设计的方法可能允许的步长小得多。

可压缩流动方程的离散化和求解可以用已经描述的方法进行。例如，要求解瞬态方程，可以使用第 1 章中讨论的任何一种时间超前方法。 6.由于雷诺数较高，可压缩流中的扩散效应通常很小，因此流动中可能存在不连续性，例如冲击。已经构建了在激波附近产生平滑解的特殊方法。这些方法包括简单的迎风法、通量混合法、本质无振荡 (ENO) 法和总变差递减 (TVD) 法。这些在该节中有描述。 4.4.6 和 11.3，并且可以在许多其他书籍中找到，例如 Tannehill 等人。 （1977）和赫希（2007）。

我们首先在下一节中描述最初为不可压缩流设计并在第 1 章中介绍的方法。 7 – 9 可以扩展以处理可压缩流。见第 11.3 节 然后我们讨论专门为计算可压缩流而设计的方法的一些方面，包括另一个方向的扩展——从可压缩流到不可压缩流。

\section*{11.2 任意马赫数的压力修正方法}\phantomsection\addcontentsline{toc}{section}{11.2 任意马赫数的压力修正方法}
为了计算可压缩流，不仅需要求解连续性方程和动量方程，还需要求解热能守恒方程（或总能量守恒方程）和状态方程。后者是连接密度、温度和压力的热力学关系。能量方程在第 1 章中给出。 1 ;对于不可压缩流，它简化为温度的标量传输方程，并且只有对流和热传导很重要。在可压缩流中，粘性耗散可能是重要的热源，并且通过流动膨胀将内能转换为动能（反之亦然）也很重要。然后必须保留方程中的所有项。积分形式的能量方程为：

\begin{sourceblock}
\begin{tikzpicture}
\node[inner sep=0,anchor=south west] (image) at (0,0) {\includegraphics[page=435,trim={53.00bp 218.60bp 53.87bp 391.41bp},clip,width=0.92\textwidth,max height=0.38\textheight]{Ferziger4th.pdf}};
\begin{scope}[x={(image.south east)},y={(image.north west)}]
\fill[white] (0.20183,0.77249) rectangle (0.22036,0.97862);
\fill[white] (0.19633,0.51737) rectangle (0.22559,0.72813);
\fill[white] (0.24845,0.49653) rectangle (0.26608,0.64939);
\fill[white] (0.31423,0.16836) rectangle (0.33402,0.37431);
\fill[white] (0.33633,0.17299) rectangle (0.34947,0.37912);
\fill[white] (0.35131,0.16836) rectangle (0.39865,0.37912);
\fill[white] (0.40974,0.17299) rectangle (0.43792,0.37912);
\fill[white] (0.43979,0.16889) rectangle (0.46457,0.37484);
\fill[white] (0.46809,0.17299) rectangle (0.48114,0.37912);
\fill[white] (0.48463,0.16836) rectangle (0.52851,0.37912);
\fill[white] (0.28493,0.02138) rectangle (0.30259,0.17406);
\end{scope}
\end{tikzpicture}
\end{sourceblock}
这里 h 是单位质量的焓，T 是绝对温度 (K)，k 是热导率，S 是应力张量的粘性部分，S = T + p I。对于具有恒定比热 c p 和 c v 的完美气体，其焓变为 h = c p T，从而允许用温度来编写能量方程。此外，在这些假设下，状态方程为：

\begin{sourceblock}
\includegraphics[page=435,trim={53.00bp 126.35bp 53.87bp 525.46bp},clip,width=0.92\textwidth,max height=0.38\textheight]{Ferziger4th.pdf}
\end{sourceblock}
其中 R 是气体常数。这些方程通过添加连续性方程来完成：

\begin{sourceblock}
\begin{tikzpicture}
\node[inner sep=0,anchor=south west] (image) at (0,0) {\includegraphics[page=435,trim={53.00bp 66.14bp 53.87bp 571.15bp},clip,width=0.92\textwidth,max height=0.38\textheight]{Ferziger4th.pdf}};
\begin{scope}[x={(image.south east)},y={(image.north west)}]
\fill[white] (0.00000,0.88562) rectangle (0.11561,1.00000);
\fill[white] (0.31814,0.55737) rectangle (0.33666,0.95841);
\fill[white] (0.31260,0.06101) rectangle (0.34189,0.47106);
\end{scope}
\end{tikzpicture}
\end{sourceblock}
和动量方程：

\begin{sourceblock}
\begin{tikzpicture}
\node[inner sep=0,anchor=south west] (image) at (0,0) {\includegraphics[page=436,trim={53.00bp 558.09bp 53.87bp 78.58bp},clip,width=0.92\textwidth,max height=0.38\textheight]{Ferziger4th.pdf}};
\begin{scope}[x={(image.south east)},y={(image.north west)}]
\fill[white] (0.15663,0.56702) rectangle (0.17516,0.95928);
\fill[white] (0.22574,0.32067) rectangle (0.26364,0.72209);
\fill[white] (0.26430,0.31727) rectangle (0.30385,0.71327);
\fill[white] (0.31026,0.32983) rectangle (0.33844,0.72209);
\fill[white] (0.15110,0.08076) rectangle (0.18036,0.48219);
\fill[white] (0.20325,0.04072) rectangle (0.22087,0.33186);
\end{scope}
\end{tikzpicture}
\end{sourceblock}
其中 T 是应力张量（包括压力项），b 表示每单位质量的体积力；见第 1 章 讨论这些方程的各种形式。

使用连续性方程来计算密度并从能量方程推导出温度是很自然的。这就将确定压力的作用留给了状态方程。因此，我们看到各种方程的作用与它们在不可压缩流中所发挥的作用有很大不同。另请注意，压力的性质完全不同。在不可压缩流动中，只有动态压力，其绝对值无关紧要；对于可压缩流，热力学压力的绝对值至关重要。

方程的离散化可以使用章节中描述的方法进行。 3和4。唯一需要的改变涉及边界条件（需要不同，因为可压缩方程在性质上是双曲线）、密度和压力之间耦合的性质和处理，以及激波这一事实，即许多变量变化极大的非常薄的区域，可能存在于可压缩流中。下面我们将遵循 Demirdži ́c 等人的方法，将压力校正方法扩展到任意马赫数的流动。 （1993）。 Issa 和 Lockwood (1977)、Karki 和 Patankar (1989) 以及 Van Doormal 等人也发表了类似的方法。 （1987）。

如上所述，可压缩动量方程的离散化本质上与不可压缩方程所采用的离散化相同，请参见第 1 章。参见图7和图8，这里不再重复。我们将把讨论限制在第 1 章中描述的隐式压力校正方法。 7，但这些想法也可以应用于其他方案。

为了使用隐式方法获得新时间级别的解，需要执行多次外部迭代；见第 7.2.1 和 7.2.2 节 详细描述了不可压缩流的分数步和 SIMPLE 方案。如果时间步长很小，则每个时间步长只需要几次外部迭代。对于稳态问题，时间步长可能是无限的，并且欠松弛参数的作用类似于伪时间步长。我们仅考虑分离求解方法，其中依次求解速度分量、压力校正、温度和其他标量变量的线性化（围绕先前外部迭代的值）方程。在求解一个变量时，其他变量被视为已知。扩展下列各节中描述的方法所需的步骤。以下两节描述了 7.2.1 和 7.2.2 能够计算任何马赫数下的流量。

\subsection*{11.2.1 所有流速的隐式分数阶方法}\phantomsection\addcontentsline{toc}{subsection}{11.2.1 所有流速的隐式分数阶方法}
在新时间级别的外部迭代 m 处计算的值取自先前的迭代 m − 1。连续性方程仅表示不可压缩流中速度场的约束（它必须始终无散度）；现在它包含时间导数，必须与其他传输方程一致地处理。为了完整起见，我们在此给出所有步骤，包括那些不需要任何更改的步骤。该算法适用于可压缩流和不可压缩流，如下所示：

1. 在新时间步内的第 m 次迭代中，使用二阶、全隐式、三时间级时间积分方案求解以下形式的动量方程，以估计新的时间级解（参见第 6.3.2.4 节）：

\begin{sourceblock}
\includegraphics[page=437,trim={53.00bp 432.47bp 53.87bp 201.67bp},clip,width=0.92\textwidth,max height=0.38\textheight]{Ferziger4th.pdf}
\end{sourceblock}
其中 v ∗ 是 v m 的预测值；需要对其进行校正以强制连续性。因为特定的空间离散化方案在这里并不重要，所以我们使用对流通量 (C)、扩散通量 (L) 和梯度算子 (G) 的符号表示法。请注意，先前外部迭代的密度用于对流项；此外，如果粘度取决于温度或其他变量，我们在粘度术语中使用上一次迭代的值。 2. 要求修正后的速度和压力满足这种形式的动量方程：

\begin{sourceblock}
\includegraphics[page=437,trim={53.00bp 291.41bp 53.87bp 342.73bp},clip,width=0.92\textwidth,max height=0.38\textheight]{Ferziger4th.pdf}
\end{sourceblock}
将式(11.6)减去式(11.5)可得速度与压力修正关系式：

\begin{sourceblock}
\begin{tikzpicture}
\node[inner sep=0,anchor=south west] (image) at (0,0) {\includegraphics[page=437,trim={53.00bp 223.78bp 53.87bp 413.98bp},clip,width=0.92\textwidth,max height=0.38\textheight]{Ferziger4th.pdf}};
\begin{scope}[x={(image.south east)},y={(image.north west)}]
\fill[white] (0.10418,0.53735) rectangle (0.12397,0.94503);
\fill[white] (0.87504,0.29493) rectangle (0.95233,0.77273);
\fill[white] (0.95585,0.30444) rectangle (0.96890,0.71177);
\fill[white] (0.08656,0.04228) rectangle (0.10638,0.44996);
\fill[white] (0.12580,0.04616) rectangle (0.13892,0.45349);
\fill[white] (0.83829,0.04228) rectangle (0.85808,0.44996);
\fill[white] (0.92439,0.00000) rectangle (1.00000,0.17689);
\end{scope}
\end{tikzpicture}
\end{sourceblock}
3G(p′)。 (11.7) 这里 v ′ = v m − v ∗ 且 p ′ = p m − p m − 1.3. 离散连续性方程将不满足 ρ m − 1 和 v ∗ —a 质量不平衡结果：

\begin{sourceblock}
\includegraphics[page=437,trim={53.00bp 144.94bp 53.87bp 489.21bp},clip,width=0.92\textwidth,max height=0.38\textheight]{Ferziger4th.pdf}
\end{sourceblock}
4. 要求校正后的密度 ρ ∗ 和速度 v m 场满足连续性方程：

\begin{sourceblock}
\includegraphics[page=437,trim={53.00bp 75.60bp 53.87bp 558.53bp},clip,width=0.92\textwidth,max height=0.38\textheight]{Ferziger4th.pdf}
\end{sourceblock}
这里 ρ 是迭代 m 时密度的估计；最终值将在计算 T m 后根据状态方程计算出来。我们引入密度校正 ρ ′ = ρ ∗ − ρ m − 1 并将校正后的密度和速度的乘积展开如下：

\begin{sourceblock}
\includegraphics[page=438,trim={53.00bp 535.24bp 53.87bp 112.46bp},clip,width=0.92\textwidth,max height=0.38\textheight]{Ferziger4th.pdf}
\end{sourceblock}
带下划线的项是两次修正的乘积，比其他项趋向于零，因此从这里开始被忽略。使用方程式。 (11.8)和(11.10)，我们可以重写式(11.8)和(11.10) （11.9）为：

\begin{sourceblock}
\includegraphics[page=438,trim={53.00bp 450.91bp 53.87bp 185.62bp},clip,width=0.92\textwidth,max height=0.38\textheight]{Ferziger4th.pdf}
\end{sourceblock}
5. 为了从方程中获得压力校正方程。 (11.11)，我们需要通过压力修正来表达速度和密度修正。对于速度校正，这已经在方程式中完成。 ( 11.7 )：它与压力修正的梯度成正比，就像不可压缩流的情况一样。对于密度修正和压力修正之间的联系，我们需要参考状态方程：

\begin{sourceblock}
\includegraphics[page=438,trim={53.00bp 335.31bp 53.87bp 302.54bp},clip,width=0.92\textwidth,max height=0.38\textheight]{Ferziger4th.pdf}
\end{sourceblock}
利用这些表达式，我们可以重写式（1）。 ( 11.11 ) 作为所有流速的压力校正方程：

\begin{sourceblock}
\includegraphics[page=438,trim={53.00bp 265.18bp 53.87bp 371.34bp},clip,width=0.92\textwidth,max height=0.38\textheight]{Ferziger4th.pdf}
\end{sourceblock}
6. 求解上述压力校正方程，对速度、压力和密度进行校正，得到 v m 、p m 和 ρ ∗；这些值用于下一步求解能量方程，从中获得更新的温度 T m。最后，根据状态方程 ρ m = f ( p m , T m ) 计算新密度。执行足够次数的迭代后，所有校正都可以忽略不计，我们可以设置 v n + 1 = v m 、p n + 1 = pm 、T n + 1 = T m 和 ρ n + 1 = ρ m，并继续下一个时间步骤。

在不可压缩流动的情况下，压力校正方程（11.13）的左侧变为零，我们恢复了之前在第 1 节中引入隐式分数阶方法时所拥有的泊松方程。 7.2.1.在可压缩流的情况下，压力校正方程 ( 11.13 ) 看起来像任何其他传输方程：它的左侧有变化率和对流项，右侧有扩散和源项。在将相同类型的修改引入到 SIMPLE 算法中之后，该方程的性质将在下面进一步讨论。

\subsection*{11.2.2 适用于所有流速的简单方法}\phantomsection\addcontentsline{toc}{subsection}{11.2.2 适用于所有流速的简单方法}
如《小节》所示。如图7.2.2所示，隐式fractionalstep方法与SIMPLE之间唯一重要的区别在于，前者修正了瞬态项中的速度，而后者则对所有对系数矩阵主对角线有贡献的项（瞬态项以及部分对流和扩散项）进行修正。因此，我们仅简要总结扩展到所有流速的简单方法中的步骤。

1. 在新的外部迭代 m 的第一步中，求解 v * 动量方程，其中密度、压力和所有流体属性均取自前一次迭代 m − 1（为了清楚起见，省略了表明这一点的上标，除非必要）。系数矩阵被分成主对角线
A D 和非对角线部分，A OD：

\begin{sourceblock}
\includegraphics[page=439,trim={53.00bp 413.38bp 53.87bp 234.40bp},clip,width=0.92\textwidth,max height=0.38\textheight]{Ferziger4th.pdf}
\end{sourceblock}
其中G表示离散梯度算子；特定的离散化方案并不重要，这就是我们使用符号表示法的原因。通过求解该方程 v m ∗ 获得的速度场通常不满足连续性方程；对于这个密度和压力也需要更新。然而，我们首先引入由于压力校正而引起的速度校正，同时将密度保持在先前的迭代水平：

\begin{sourceblock}
\includegraphics[page=439,trim={53.00bp 305.78bp 53.87bp 339.60bp},clip,width=0.92\textwidth,max height=0.38\textheight]{Ferziger4th.pdf}
\end{sourceblock}
2、速度和压力修正之间的关系是通过要求修正后的速度和压力满足以下方程的简化版本得到的。 （11.14）：

\begin{sourceblock}
\begin{tikzpicture}
\node[inner sep=0,anchor=south west] (image) at (0,0) {\includegraphics[page=439,trim={53.00bp 246.00bp 53.87bp 401.77bp},clip,width=0.92\textwidth,max height=0.38\textheight]{Ferziger4th.pdf}};
\begin{scope}[x={(image.south east)},y={(image.north west)}]
\fill[white] (0.03829,0.84213) rectangle (0.06803,1.00000);
\fill[white] (0.07074,0.84213) rectangle (0.11630,1.00000);
\fill[white] (0.11901,0.84213) rectangle (0.21949,1.00000);
\end{scope}
\end{tikzpicture}
\end{sourceblock}
现在减去方程。 ( 11.14 ) 由式( 11.16 )，我们得到速度和压力修正之间的关系如下：

\begin{sourceblock}
\includegraphics[page=439,trim={53.00bp 186.23bp 53.87bp 461.55bp},clip,width=0.92\textwidth,max height=0.38\textheight]{Ferziger4th.pdf}
\end{sourceblock}
3. 我们现在转向质量守恒方程。以下两个步骤与刚才描述的隐式分步法中的步骤3和4相同（假设使用相同的时间积分方案，这里是具有三个时间级别的完全隐式方案），因此这里不再重复；参见方程式。 （11.8）-（11.11）。 4. 密度和压力修正之间的关系与之前相同，参见。等式。 （11.12）；速度和压力修正之间的关系由方程式给出。 （11.17）。通过将这些关系插入方程。 ( 11.11 ) , 得到压力修正方程如下形式:

\begin{sourceblock}
\includegraphics[page=440,trim={53.00bp 583.58bp 53.87bp 52.95bp},clip,width=0.92\textwidth,max height=0.38\textheight]{Ferziger4th.pdf}
\end{sourceblock}
为了简单起见，我们在这里假设所有三个速度分量的主对角线系数相同；通常是这种情况，但如果不是，也很容易考虑到差异。

通过比较 SIMPLE 的压力校正方程 ( 11.18 ) 与 IFSM 的相应方程 ( 11.13 )，我们发现只有右侧的第一项（类似于离散拉普拉斯算子的项）不同。这是由于速度和压力修正之间的联系的不同表达，参见。等式。 ( 11.7 ) 对于 IFSM 和方程。 ( 11.17 ) 简单。

\subsection*{11.2.3 压力校正方程的性质}\phantomsection\addcontentsline{toc}{subsection}{11.2.3 压力校正方程的性质}
\begin{sourceblock}
\begin{tikzpicture}
\node[inner sep=0,anchor=south west] (image) at (0,0) {\includegraphics[page=440,trim={53.00bp 364.28bp 53.87bp 264.50bp},clip,width=0.92\textwidth,max height=0.38\textheight]{Ferziger4th.pdf}};
\begin{scope}[x={(image.south east)},y={(image.north west)}]
\fill[white] (0.00000,0.74518) rectangle (0.08731,1.00000);
\fill[white] (0.08998,0.74518) rectangle (0.14289,1.00000);
\fill[white] (0.32033,0.25990) rectangle (0.35735,0.58619);
\fill[white] (0.36361,0.28399) rectangle (0.39179,0.59342);
\end{scope}
\end{tikzpicture}
\end{sourceblock}
对于其他气体以及当液体被认为是可压缩的时，可能需要以数值方式计算导数。收敛解与该系数无关，因为所有校正都为零；仅中间结果受到影响。重要的是密度和压力修正之间的联系在质量上是正确的，并且系数当然会影响方法的收敛速度。

压力校正方程中的系数取决于压力校正的梯度和单元面值所使用的近似值。由速度修正产生的部分与不可压缩情况下的部分相同；它需要在垂直于单元表面的方向上对压力校正的导数进行近似，其近似方式应与动量方程中的压力项相同。由密度修正产生的部分对应于其他守恒方程中的对流通量。它需要对单元面中心的压力进行近似校正；参见章节。图 4 和图 7 是各种常用近似值的示例。

尽管在外观上与不可压缩流的压力校正方程相似，但仍存在重要差异。不可压缩方程是离散泊松方程，即系数表示拉普拉斯算子的近似值。在可压缩情况下，有一些贡献表明可压缩流中的压力方程包含对流和不稳定项，即它实际上是对流波动方程。对于不可压缩流，如果在边界处规定质量通量，则压力可能在附加常数内不确定。可压缩压力校正方程中对流项的存在使得该解具有独特性。

速度和密度校正产生的术语的相对重要性取决于流的类型。扩散项相对于对流项的数量级为 1/Ma 2，因此马赫数是决定因素。在低马赫数时，拉普拉斯项占主导地位并覆盖泊松方程。另一方面，在高马赫数（高度可压缩流）时，对流项占主导地位，反映了流动的双曲性质。求解压力校正方程就相当于求解密度连续方程。因此，压力校正方法会自动调整以适应流动的局部性质，并且相同的方法可以应用于整个流动区域，即使它同时包含高马赫数区域和低马赫数区域（例如，在钝体周围的流动中）。

对于拉普拉斯算子的近似，总是应用中心差近似。另一方面，对于对流项的近似，可以使用多种近似，正如动量方程中的对流项的情况一样。如果使用高阶近似，则可以使用“延迟校正”方法。方程左边是基于一阶迎风近似构造的矩阵，右边包含高阶近似与一阶迎风近似之间的差值，保证方法收敛于属于高阶近似的解；见第 5.6 节 详细信息。此外，如果网格严重非正交，则可以使用延迟校正来简化压力校正方程，如第 1 节中所述。 9.8.

这些差异以另一种方式反映在压力校正方程中。因为该方程不再是纯泊松方程，所以中心系数A P 不是相邻系数之和的负数。只有当div v = 0时，才能获得此属性。此外，虽然不可压缩流的压力校正方程具有对称系数矩阵并允许使用一些特殊的求解器，但由于对流项的影响，该方程在可压缩流的情况下失去了这一特性。

\subsection*{11.2.4 边界条件}\phantomsection\addcontentsline{toc}{subsection}{11.2.4 边界条件}
对于不可压缩流动，通常应用以下边界条件：

\begin{itemize}
\item 流入边界上的规定速度和温度；
\item 对于所有标量，垂直于边界的零梯度以及平行于对称平面上的表面的速度分量；垂直于该表面的速度为零；
\item 固体表面上的无滑移（零相对速度）条件、零正应力和规定温度或热通量；
\item 流出表面上所有量的规定梯度（通常为零）。
\end{itemize}
这些边界条件也适用于可压缩流，并且以与不可压缩流相同的方式处理。然而，在可压缩流中，还存在其他可能的边界条件：

\begin{itemize}
\item 规定的总压力；
\item 规定的总温度；
\item 流出边界1 上的规定静压；
\item 在超音速流出边界处，通常指定所有量的零梯度。
\end{itemize}
下面描述这些边界条件的实现。

11.2.4.1 流入边界上的规定总压力
将借助图11.1对二维域的西边界描述这些边界条件的实现。

一种可能性是，对于理想气体的等熵流，总压力定义为：

\begin{sourceblock}
\begin{tikzpicture}
\node[inner sep=0,anchor=south west] (image) at (0,0) {\includegraphics[page=442,trim={53.00bp 407.36bp 53.87bp 220.40bp},clip,width=0.92\textwidth,max height=0.38\textheight]{Ferziger4th.pdf}};
\begin{scope}[x={(image.south east)},y={(image.north west)}]
\fill[white] (0.00000,0.76733) rectangle (0.02412,1.00000);
\fill[white] (0.02680,0.76733) rectangle (0.11979,1.00000);
\fill[white] (0.12250,0.76733) rectangle (0.16060,1.00000);
\fill[white] (0.28373,0.19515) rectangle (0.30938,0.51954);
\fill[white] (0.31501,0.22512) rectangle (0.34319,0.52632);
\fill[white] (0.35017,0.21808) rectangle (0.36995,0.51954);
\end{scope}
\end{tikzpicture}
\end{sourceblock}
其中 p 是静压，γ = c p / c v。必须规定流向；它的定义为：

\begin{sourceblock}
\begin{tikzpicture}
\node[inner sep=0,anchor=south west] (image) at (0,0) {\includegraphics[page=442,trim={53.00bp 348.00bp 53.87bp 289.30bp},clip,width=0.92\textwidth,max height=0.38\textheight]{Ferziger4th.pdf}};
\begin{scope}[x={(image.south east)},y={(image.north west)}]
\fill[white] (0.00000,0.82108) rectangle (0.02081,1.00000);
\fill[white] (0.02346,0.82108) rectangle (0.04827,1.00000);
\fill[white] (0.05092,0.82108) rectangle (0.14391,1.00000);
\fill[white] (0.14668,0.82108) rectangle (0.18977,1.00000);
\end{scope}
\end{tikzpicture}
\end{sourceblock}
这些边界条件可以通过将压力从解域内部外推到边界来实现，然后借助方程式计算那里的速度。 （11.20）和（11.21）。这些速度可以被视为外迭代中已知的。温度可以指定，也可以根据总温度计算：

\begin{sourceblock}
\begin{tikzpicture}
\node[inner sep=0,anchor=south west] (image) at (0,0) {\includegraphics[page=442,trim={53.00bp 240.13bp 53.87bp 395.46bp},clip,width=0.92\textwidth,max height=0.38\textheight]{Ferziger4th.pdf}};
\begin{scope}[x={(image.south east)},y={(image.north west)}]
\fill[white] (0.30057,0.24517) rectangle (0.32671,0.65270);
\fill[white] (0.33236,0.28282) rectangle (0.36054,0.66121);
\fill[white] (0.36475,0.27398) rectangle (0.38623,0.65270);
\end{scope}
\end{tikzpicture}
\end{sourceblock}
这种处理导致迭代方法收敛缓慢，因为有许多压力和速度的组合满足方程（1）。 （11.20）。人们必须隐含地考虑压力对流入速度的影响。下面描述了这样做的一种方法。

在外部迭代开始时，流入边界（图 11.1 中的“w”侧）处的速度必须根据方程式计算。 (11.20)至(11.21)以及压力的现行值；然后，在动量方程的外部迭代期间，它们将被视为固定的。流入处的质量通量取自先前的外部迭代；它们应该满足连续性方程。从解的

\begin{center}\small 图 11.1 具有规定流向的入口边界旁边的控制体积\end{center}
N

\begin{sourceblock}
\begin{tikzpicture}
\node[inner sep=0,anchor=south west] (image) at (0,0) {\includegraphics[page=443,trim={53.00bp 525.97bp 53.87bp 85.67bp},clip,width=0.92\textwidth,max height=0.38\textheight]{Ferziger4th.pdf}};
\begin{scope}[x={(image.south east)},y={(image.north west)}]
\fill[white] (0.00000,0.94440) rectangle (0.04944,1.00000);
\fill[white] (0.05098,0.94440) rectangle (0.06710,1.00000);
\fill[white] (0.06863,0.94440) rectangle (0.17955,1.00000);
\fill[white] (0.18111,0.94440) rectangle (0.23053,1.00000);
\fill[white] (0.00000,0.76165) rectangle (0.09471,0.94642);
\fill[white] (0.79450,0.77853) rectangle (0.81323,0.97798);
\fill[white] (0.83182,0.74018) rectangle (0.85056,0.93963);
\fill[white] (0.61014,0.67982) rectangle (0.63044,0.87927);
\fill[white] (0.71110,0.41706) rectangle (0.73603,0.61651);
\fill[white] (0.87131,0.30991) rectangle (0.88851,0.50936);
\fill[white] (0.79122,0.02202) rectangle (0.80686,0.22147);
\end{scope}
\end{tikzpicture}
\end{sourceblock}
S

j

i

x

动量方程，( u ∗
x , u ∗
y )，计算新的质量通量 ̇ m *。流入边界上的“规定”速度用于计算那里的质量通量。在接下来的校正步骤中，质量通量（包括其在流入边界处的值）被校正并强制质量守恒。边界上的质量通量校正与内部控制体积面上的质量通量校正之间的差异在于，在边界处，仅校正速度（而不校正密度）。速度修正以压力修正（而不是其梯度）表示：

\begin{sourceblock}
\begin{tikzpicture}
\node[inner sep=0,anchor=south west] (image) at (0,0) {\includegraphics[page=443,trim={53.00bp 336.72bp 53.87bp 292.00bp},clip,width=0.92\textwidth,max height=0.38\textheight]{Ferziger4th.pdf}};
\begin{scope}[x={(image.south east)},y={(image.north west)}]
\fill[white] (0.18550,0.27819) rectangle (0.21254,0.64030);
\end{scope}
\end{tikzpicture}
\end{sourceblock}
系数 C u 借助等式确定。 （11.20）：

\begin{sourceblock}
\begin{tikzpicture}
\node[inner sep=0,anchor=south west] (image) at (0,0) {\includegraphics[page=443,trim={53.00bp 260.79bp 53.87bp 357.54bp},clip,width=0.92\textwidth,max height=0.38\textheight]{Ferziger4th.pdf}};
\begin{scope}[x={(image.south east)},y={(image.north west)}]
\fill[white] (0.47296,0.32692) rectangle (0.49281,0.56871);
\fill[white] (0.49946,0.31918) rectangle (0.54926,0.56871);
\fill[white] (0.55997,0.32127) rectangle (0.60346,0.59444);
\fill[white] (0.41768,0.17319) rectangle (0.43747,0.41498);
\fill[white] (0.43934,0.18072) rectangle (0.46752,0.42271);
\fill[white] (0.50114,0.02510) rectangle (0.52093,0.26710);
\end{scope}
\end{tikzpicture}
\end{sourceblock}
流入边界处的质量通量修正表示为：

\begin{sourceblock}
\begin{tikzpicture}
\node[inner sep=0,anchor=south west] (image) at (0,0) {\includegraphics[page=443,trim={53.00bp 211.48bp 53.87bp 432.75bp},clip,width=0.92\textwidth,max height=0.38\textheight]{Ferziger4th.pdf}};
\begin{scope}[x={(image.south east)},y={(image.north west)}]
\fill[white] (0.03035,0.21771) rectangle (0.06424,0.83615);
\fill[white] (0.05377,0.12277) rectangle (0.07371,0.51438);
\fill[white] (0.07874,0.22958) rectangle (0.10692,0.75764);
\fill[white] (0.11044,0.21771) rectangle (0.21032,0.83615);
\end{scope}
\end{tikzpicture}
\end{sourceblock}
边界处的压力校正 ( p ' ) w 通过从相邻控制体积的中心外推来表示，即作为 p ' 的线性组合
P and p ′
E.从上面的方程中，我们获得了对边界附近控制体积的压力校正方程中的系数 A P 和 A E 的贡献。由于密度在流入处未进行校正，因此此处的压力校正方程没有对流贡献，因此系数 A W 为零。

求解压力修正方程后，对包括流入边界在内的整个域中的速度分量和质量通量进行修正。校正后的质量通量在收敛公差内满足连续性方程。这些用于计算下一次外部迭代的所有传输方程中的系数。流入边界处的对流速度由方程式计算： (11.20)至(11.21)。压力自行调节，使得速度满足连续性方程和总压力的边界条件。流入处的温度由方程计算。 ( 11.22 )，以及状态方程中的密度 ( 11.2 )。

11.2.4.2 规定静压
在亚音速流中，静压力通常规定在流出边界上。那么该边界上的压力校正为零（这用作压力校正方程中的边界条件），但质量通量校正不为零。速度分量是通过从相邻控制体积中心外推获得的，其方式类似于计算共置网格上的单元面速度，例如，对于“e”面和第 m 个外部迭代：

\begin{sourceblock}
\begin{tikzpicture}
\node[inner sep=0,anchor=south west] (image) at (0,0) {\includegraphics[page=444,trim={53.00bp 397.55bp 53.87bp 227.78bp},clip,width=0.92\textwidth,max height=0.38\textheight]{Ferziger4th.pdf}};
\begin{scope}[x={(image.south east)},y={(image.north west)}]
\fill[white] (0.37432,0.43347) rectangle (0.39411,0.71674);
\fill[white] (0.10749,0.26464) rectangle (0.13741,0.59667);
\fill[white] (0.12078,0.21392) rectangle (0.15326,0.43029);
\fill[white] (0.15853,0.27126) rectangle (0.18671,0.55452);
\fill[white] (0.19023,0.20289) rectangle (0.25329,0.56653);
\fill[white] (0.25687,0.27126) rectangle (0.28505,0.55452);
\fill[white] (0.31182,0.24308) rectangle (0.34340,0.54815);
\fill[white] (0.36713,0.09140) rectangle (0.40132,0.41044);
\end{scope}
\end{tikzpicture}
\end{sourceblock}
where v ∗
n 是垂直于流出边界方向的速度分量，可以通过求解动量方程 u ∗ 获得的笛卡尔分量轻松计算出来
i，以及单位向外法向量的已知分量 v *
n = v * · n。与计算内部单元表面的速度的唯一区别是，这里的上线表示从内部单元的外推，而不是表面两侧的单元中心之间的插值。在高流速下，如果流出边界位于下游较远的地方，通常可以使用简单的迎风方案，即使用像元中心值（节点P）代替overbar表示的值； W 和 P 的线性外推也可以在结构化网格上轻松实现。

由这些速度构造的质量通量通常不满足连续性方程，因此必须进行校正。通常，如上所述，速度和密度都需要校正。速度修正为：

\begin{sourceblock}
\begin{tikzpicture}
\node[inner sep=0,anchor=south west] (image) at (0,0) {\includegraphics[page=444,trim={53.00bp 201.92bp 53.87bp 426.40bp},clip,width=0.92\textwidth,max height=0.38\textheight]{Ferziger4th.pdf}};
\begin{scope}[x={(image.south east)},y={(image.north west)}]
\fill[white] (0.31146,0.28556) rectangle (0.33612,0.64384);
\fill[white] (0.32475,0.23057) rectangle (0.35723,0.46430);
\fill[white] (0.36250,0.29270) rectangle (0.39068,0.59836);
\fill[white] (0.39420,0.29270) rectangle (0.42238,0.59836);
\fill[white] (0.44250,0.26203) rectangle (0.47408,0.59122);
\end{scope}
\end{tikzpicture}
\end{sourceblock}
\begin{sourceblock}
\begin{tikzpicture}
\node[inner sep=0,anchor=south west] (image) at (0,0) {\includegraphics[page=444,trim={53.00bp 167.40bp 53.87bp 480.88bp},clip,width=0.92\textwidth,max height=0.38\textheight]{Ferziger4th.pdf}};
\begin{scope}[x={(image.south east)},y={(image.north west)}]
\fill[white] (0.00000,0.86562) rectangle (0.05074,1.00000);
\fill[white] (0.05392,0.86562) rectangle (0.18851,1.00000);
\fill[white] (0.19179,0.86562) rectangle (0.30310,1.00000);
\fill[white] (0.30632,0.86562) rectangle (0.45693,1.00000);
\fill[white] (0.46018,0.86562) rectangle (0.48830,1.00000);
\fill[white] (0.49152,0.86562) rectangle (0.53296,1.00000);
\fill[white] (0.53615,0.86562) rectangle (0.60087,1.00000);
\fill[white] (0.60409,0.86562) rectangle (0.65552,1.00000);
\fill[white] (0.65874,0.86562) rectangle (0.78496,1.00000);
\fill[white] (0.78821,0.86562) rectangle (0.86761,1.00000);
\fill[white] (0.87083,0.86562) rectangle (0.92391,1.00000);
\fill[white] (0.92710,0.86562) rectangle (0.97020,1.00000);
\fill[white] (0.97344,0.86562) rectangle (1.00000,1.00000);
\fill[white] (0.00000,0.19597) rectangle (0.03242,0.84378);
\fill[white] (0.03305,0.19597) rectangle (0.08941,0.84378);
\fill[white] (0.09008,0.19597) rectangle (0.19967,0.84378);
\fill[white] (0.20036,0.21669) rectangle (0.22893,0.93281);
\fill[white] (0.44262,0.08623) rectangle (0.45675,0.56607);
\fill[white] (0.46195,0.21669) rectangle (0.49017,0.86394);
\fill[white] (0.49365,0.19597) rectangle (0.51347,0.84378);
\fill[white] (0.51426,0.19597) rectangle (0.61389,0.84378);
\fill[white] (0.61462,0.19597) rectangle (0.65606,0.84378);
\fill[white] (0.65669,0.19597) rectangle (0.76132,0.84378);
\fill[white] (0.76195,0.19597) rectangle (0.78674,0.84378);
\fill[white] (0.78740,0.19597) rectangle (0.93528,0.84378);
\fill[white] (0.93597,0.19597) rectangle (1.00000,0.84378);
\fill[white] (0.00000,0.00000) rectangle (0.06078,0.17413);
\fill[white] (0.06367,0.00000) rectangle (0.17329,0.17413);
\fill[white] (0.17621,0.00000) rectangle (0.28084,0.17413);
\fill[white] (0.28370,0.00000) rectangle (0.30851,0.17413);
\fill[white] (0.31137,0.00000) rectangle (0.44842,0.17413);
\fill[white] (0.45134,0.00000) rectangle (0.49275,0.17413);
\fill[white] (0.49561,0.00000) rectangle (0.64352,0.17413);
\fill[white] (0.64641,0.00000) rectangle (0.67122,0.17413);
\fill[white] (0.67408,0.00000) rectangle (0.71717,0.17413);
\fill[white] (0.72003,0.00000) rectangle (0.78430,0.17413);
\fill[white] (0.78722,0.00000) rectangle (0.81865,0.17413);
\fill[white] (0.82150,0.00000) rectangle (0.84629,0.17413);
\fill[white] (0.84914,0.00000) rectangle (1.00000,0.17413);
\end{scope}
\end{tikzpicture}
\end{sourceblock}
e = 0，因为压力是规定的）；然而，虽然规定了压力，但温度并不固定（从内部推断），因此密度确实需要校正。最简单的近似是一阶迎风近似，即取ρ ′

\begin{sourceblock}
\begin{tikzpicture}
\node[inner sep=0,anchor=south west] (image) at (0,0) {\includegraphics[page=444,trim={53.00bp 107.62bp 53.87bp 516.69bp},clip,width=0.92\textwidth,max height=0.38\textheight]{Ferziger4th.pdf}};
\begin{scope}[x={(image.south east)},y={(image.north west)}]
\fill[white] (0.00000,0.94143) rectangle (0.06238,1.00000);
\fill[white] (0.06349,0.94143) rectangle (0.15732,1.00000);
\fill[white] (0.15844,0.94143) rectangle (0.18989,1.00000);
\fill[white] (0.19101,0.94143) rectangle (0.23242,1.00000);
\fill[white] (0.23353,0.94143) rectangle (0.32487,1.00000);
\fill[white] (0.32598,0.94143) rectangle (0.38571,1.00000);
\fill[white] (0.38686,0.94143) rectangle (0.44824,1.00000);
\fill[white] (0.44938,0.94143) rectangle (0.47750,1.00000);
\fill[white] (0.47862,0.94143) rectangle (0.51170,1.00000);
\fill[white] (0.51284,0.94143) rectangle (0.63654,1.00000);
\fill[white] (0.63771,0.94143) rectangle (0.68911,1.00000);
\fill[white] (0.69023,0.94143) rectangle (0.79492,1.00000);
\fill[white] (0.79606,0.94143) rectangle (0.97555,1.00000);
\fill[white] (0.97675,0.94143) rectangle (1.00000,1.00000);
\fill[white] (0.00000,0.65551) rectangle (0.04075,0.93187);
\fill[white] (0.04298,0.65551) rectangle (0.16758,0.93187);
\fill[white] (0.16983,0.65551) rectangle (0.26454,0.93187);
\fill[white] (0.26680,0.65551) rectangle (0.45383,0.93187);
\fill[white] (0.45615,0.65551) rectangle (0.50505,0.93187);
\fill[white] (0.50731,0.65551) rectangle (0.58701,0.93187);
\fill[white] (0.58926,0.66436) rectangle (0.61786,0.97012);
\fill[white] (0.67296,0.60602) rectangle (0.69874,0.93187);
\fill[white] (0.70099,0.65551) rectangle (0.75239,0.93187);
\fill[white] (0.75465,0.65551) rectangle (0.81937,0.93187);
\fill[white] (0.82162,0.65551) rectangle (0.87302,0.93187);
\fill[white] (0.87534,0.65551) rectangle (1.00000,0.93187);
\fill[white] (0.00000,0.36983) rectangle (0.02412,0.64619);
\fill[white] (0.02680,0.36983) rectangle (0.09152,0.64619);
\end{scope}
\end{tikzpicture}
\end{sourceblock}
但请注意，密度校正并不用于校正流出边界处的密度 - 一旦计算了压力和温度，密度校正始终根据状态方程进行计算。另一方面，质量通量必须使用上述表达式进行校正，因为只有用于导出压力校正方程的校正才能确保质量守恒。因为，在收敛时，所有校正都归零，因此上述密度校正的处理与其他近似一致，并且不影响解的精度，仅影响迭代方案的收敛速度。压力校正方程中边界节点的系数不包含对流项的贡献（由于迎风）——它的贡献归于中心系数 A P。法线方向的压力导数通常近似为：

\begin{sourceblock}
\begin{tikzpicture}
\node[inner sep=0,anchor=south west] (image) at (0,0) {\includegraphics[page=445,trim={53.00bp 460.71bp 53.87bp 168.01bp},clip,width=0.92\textwidth,max height=0.38\textheight]{Ferziger4th.pdf}};
\begin{scope}[x={(image.south east)},y={(image.north west)}]
\fill[white] (0.36653,0.46472) rectangle (0.40776,0.81347);
\end{scope}
\end{tikzpicture}
\end{sourceblock}
遵循第 2 节中描述的方法。 9.8 和等式（9.66）。

因此，与内部 CV 处的系数相比，边界附近控制体积的压力校正方程中的系数 A P 发生了变化。由于压力校正方程中的对流项和指定静压的狄利克雷边界条件，它通常比不可压缩流收敛得更快（其中诺依曼边界条件通常应用于所有边界，并且方程是全椭圆形）。

11.2.4.3 无反射和自由流边界
在边界的某些部分，可能不知道要应用的确切条件，但压力波和/或冲击应该能够穿过边界而不反射。通常，一维理论用于根据规定的自由流压力和温度计算边界处的速度。如果自由流是超音速的，则冲击可能会穿过边界，然后人们会区分平行速度分量（简单地外推到边界）和法向分量（根据理论计算）。后一个条件取决于压缩波或膨胀（普朗特-迈耶）波是否撞击边界。压力通常从内部外推到边界，而法向速度分量是使用外推压力和规定的自由流马赫数计算的。

有许多方案旨在产生无反射和自由流边界。它们的推导依赖于通过一维理论计算的传出特征；实现取决于离散化和求解方法。这些（数值）边界条件的详细讨论可以在 Hirsch (2007) 和 Durran (2010) 中找到。

如果没有激波穿过自由流或出口边界，则可以如上所述规定那里的静压，并应用解强制技术（参见第 13.6 节）以避免压力波从这些边界反射。这对于捕获声压波的弱可压缩流尤其实用（空气声学或水声分析）；必须避免它们在边界处的反射。有关此主题的更多详细信息可以在 Peri ́c (2019) 中找到。

11.2.4.4 超音速流出
如果流出处的流动是超音速的，则边界处的所有变量都必须通过内部外推来获得，即不需要规定边界信息。压力校正方程的处理与规定静压的情况类似。然而，由于边界处的压力不是规定的而是外推的，因此压力校正也需要外推——它不像上述情况那样为零。因为 p ′

\begin{sourceblock}
\begin{tikzpicture}
\node[inner sep=0,anchor=south west] (image) at (0,0) {\includegraphics[page=446,trim={53.00bp 425.23bp 53.87bp 210.93bp},clip,width=0.92\textwidth,max height=0.38\textheight]{Ferziger4th.pdf}};
\begin{scope}[x={(image.south east)},y={(image.north west)}]
\fill[white] (0.00000,0.91828) rectangle (0.02746,1.00000);
\fill[white] (0.03002,0.91828) rectangle (0.06310,1.00000);
\fill[white] (0.06568,0.91828) rectangle (0.26466,1.00000);
\fill[white] (0.26734,0.91828) rectangle (0.29212,1.00000);
\fill[white] (0.29471,0.91828) rectangle (0.33777,1.00000);
\fill[white] (0.34039,0.91828) rectangle (0.39675,1.00000);
\fill[white] (0.39934,0.91828) rectangle (0.42911,1.00000);
\fill[white] (0.43170,0.91828) rectangle (0.45982,1.00000);
\fill[white] (0.46241,0.91828) rectangle (0.50382,1.00000);
\fill[white] (0.50641,0.91828) rectangle (0.58180,1.00000);
\fill[white] (0.58445,0.91828) rectangle (0.64830,1.00000);
\fill[white] (0.65092,0.91828) rectangle (0.75552,1.00000);
\fill[white] (0.76162,0.92161) rectangle (0.78800,1.00000);
\fill[white] (0.77663,0.84890) rectangle (0.79426,1.00000);
\fill[white] (0.79868,0.91828) rectangle (0.82346,1.00000);
\fill[white] (0.82602,0.91828) rectangle (0.94851,1.00000);
\fill[white] (0.95116,0.91828) rectangle (0.98090,1.00000);
\fill[white] (0.98349,0.91828) rectangle (1.00000,1.00000);
\fill[white] (0.00000,0.51935) rectangle (0.07233,0.90494);
\fill[white] (0.07395,0.51935) rectangle (0.22851,0.90494);
\fill[white] (0.23017,0.51935) rectangle (0.25991,0.90494);
\fill[white] (0.26496,0.52268) rectangle (0.29134,0.95997);
\fill[white] (0.36800,0.44997) rectangle (0.39260,0.73616);
\fill[white] (0.39606,0.51935) rectangle (0.42914,0.90494);
\fill[white] (0.43074,0.51935) rectangle (0.47215,0.90494);
\fill[white] (0.47374,0.51935) rectangle (0.57835,0.90494);
\fill[white] (0.57994,0.51935) rectangle (0.68289,0.90494);
\fill[white] (0.68451,0.51935) rectangle (0.73089,0.90494);
\fill[white] (0.73251,0.51935) rectangle (0.76559,0.90494);
\fill[white] (0.76722,0.51935) rectangle (0.89383,0.90494);
\fill[white] (0.89549,0.51935) rectangle (0.94355,0.90494);
\fill[white] (0.94517,0.51935) rectangle (1.00000,0.90494);
\fill[white] (0.00000,0.12041) rectangle (0.05242,0.50634);
\fill[white] (0.05573,0.12041) rectangle (0.09386,0.50634);
\fill[white] (0.10063,0.12375) rectangle (0.12701,0.56104);
\fill[white] (0.18878,0.05137) rectangle (0.22454,0.50634);
\fill[white] (0.22788,0.12041) rectangle (0.26929,0.50634);
\fill[white] (0.27263,0.12041) rectangle (0.33567,0.50634);
\fill[white] (0.33907,0.12041) rectangle (0.36217,0.50634);
\fill[white] (0.36553,0.12041) rectangle (0.42523,0.50634);
\fill[white] (0.42863,0.12041) rectangle (0.47170,0.50634);
\fill[white] (0.47507,0.12041) rectangle (0.54641,0.50634);
\fill[white] (0.54980,0.12041) rectangle (0.57792,0.50634);
\fill[white] (0.58129,0.12041) rectangle (0.62271,0.50634);
\fill[white] (0.62605,0.12041) rectangle (0.74063,0.50634);
\fill[white] (0.74400,0.12041) rectangle (0.85943,0.50634);
\fill[white] (0.86283,0.12041) rectangle (0.89429,0.50634);
\fill[white] (0.90057,0.09440) rectangle (0.93708,0.50967);
\fill[white] (0.94241,0.13276) rectangle (0.97062,0.51835);
\fill[white] (0.97411,0.12041) rectangle (1.00000,0.50634);
\fill[white] (0.00000,0.00000) rectangle (0.05074,0.10740);
\fill[white] (0.05194,0.00000) rectangle (0.19408,0.10740);
\fill[white] (0.19537,0.00000) rectangle (0.22511,0.10740);
\fill[white] (0.22635,0.00000) rectangle (0.30105,0.10740);
\fill[white] (0.30232,0.00000) rectangle (0.42520,0.10740);
\fill[white] (0.42650,0.00000) rectangle (0.45462,0.10740);
\fill[white] (0.45579,0.00000) rectangle (0.49723,0.10740);
\fill[white] (0.49844,0.00000) rectangle (0.67795,0.10740);
\fill[white] (0.67925,0.00000) rectangle (0.70899,0.10740);
\fill[white] (0.71020,0.00000) rectangle (0.75164,0.10740);
\fill[white] (0.75284,0.00000) rectangle (0.87414,0.10740);
\fill[white] (0.87537,0.00000) rectangle (1.00000,0.10740);
\end{scope}
\end{tikzpicture}
\end{sourceblock}
P )，代数方程中没有出现节点 E，因此 A E = 0。通过边界的质量通量校正近似中出现的节点系数与内部区域的节点系数不同。

下面介绍了应用压力校正方案解决可压缩流动问题的一些示例。更多示例可以在 Demirdži ́c 等人中找到。 ( 1993 )、Lilek ( 1995 ) 和 Riahi 等人。 （2018）。

\subsection*{11.2.5 示例}\phantomsection\addcontentsline{toc}{subsection}{11.2.5 示例}
下面我们给出了圆弧凸块上流动的欧拉方程的解结果。图 11.2 显示了亚音速、跨音速和超音速条件下的几何形状和马赫数预测等值线。圆弧的厚度与弦比对于亚音速和跨音速情况为10\%，对于超音速情况为4\%。指定马赫数 Ma = 0.5（亚音速）、0.675（跨音速）和 1.65（超音速）时的均匀入口流量。由于欧拉方程已求解，因此粘度设置为零，并且在壁处指定滑移条件（流动相切，对于对称表面而言）。这些问题是 1981 年研讨会上的测试用例（参见 Rizzi 和 Viviand 1981），通常用于评估数值方案的准确性。

对于亚音速流动，由于几何形状是对称的并且流动是无粘性的，因此流动也是对称的。总压力在整个解域中应保持恒定，这对于评估数值误差很有用。在跨音速情况下，在下壁上获得一次激波。当迎面而来的气流是超音速时，当气流到达凸块时会产生激波。这种冲击力被上壁反射；它穿过另一个冲击，该冲击从凸起的末端发出，在那里遇到了墙坡度的另一个突然变化。

\begin{center}\small 图 11.3 显示了三种情况下马赫数分别沿下壁和上壁的分布。在最细网格和亚音速流上求解误差很小；这可以从网格细化的效果以及从\end{center}
\begin{sourceblock}
\includegraphics[page=447,trim={53.00bp 253.25bp 53.87bp 52.00bp},clip,width=\textwidth,max height=0.38\textheight]{Ferziger4th.pdf}
\par\vspace{0.45\baselineskip}
\small 图 11.2 通过下壁带有圆弧凸块的通道的无粘流的预测马赫数等值线：Ma in = 0 时的亚音速流。 5（上图），Ma in = 0 处的跨音速流。 675（中），Ma in = 1 处的超音速流动。 65（下）；来自利莱克（1995）
\end{sourceblock}
事实上，出口处两壁的马赫数相同并且等于入口值。总压误差低于0.25\%。在跨音速和超音速情况下，网格细化仅影响激波的陡度；它在三个网格点内解决。如果对所有方程中的所有项都使用中心差分，则冲击处的强烈振荡会使求解变得困难。在此介绍的计算中，使用 10\% 的 UDS 和 90\% 的 CDS 来减少振荡；它们仍然存在，如图 11.3 所示，但仅限于激波附近的两个网格点。有趣的是，激波的位置不会随着网格细化而改变，只是陡度得到改善（这在许多应用中都观察到）。所使用的 FV 方法的守恒特性和 CDS 近似的主导作用可能是造成此特征的原因。

\begin{sourceblock}
\includegraphics[page=448,trim={53.00bp 130.25bp 53.87bp 52.00bp},clip,width=\textwidth,max height=0.38\textheight]{Ferziger4th.pdf}
\par\vspace{0.45\baselineskip}
\small 图 11.3 沿着下壁和上壁预测的马赫数分布，用于穿过下壁中带有圆弧凸块的通道的无粘流：Ma in = 0 时的亚音速流动。 5（上图；95\% CDS，5\% UDS），Ma in = 0 处的跨音速流。 675（中；90\% CDS，10\% UDS），以及 Ma in = 1 时的超音速流。 65（底部；90\% CDS，10\% UDS）；来自利莱克（1995）
\end{sourceblock}
\begin{sourceblock}
\includegraphics[page=449,trim={53.00bp 489.25bp 53.87bp 52.00bp},clip,width=\textwidth,max height=0.38\textheight]{Ferziger4th.pdf}
\par\vspace{0.45\baselineskip}
\small 图11.4 下壁有圆弧凸块通道的超音速无粘流的预测马赫数等值线（160×80 CV网格，100\% CDS离散化）；来自利莱克（1995）
\end{sourceblock}
\begin{sourceblock}
\includegraphics[page=449,trim={53.00bp 335.23bp 53.87bp 210.37bp},clip,width=\textwidth,max height=0.38\textheight]{Ferziger4th.pdf}
\par\vspace{0.45\baselineskip}
\small 图 11.5 可压缩通道流的几何形状和边界条件
\end{sourceblock}
在图 11.4 中，显示了对所有单元面量使用纯 CDS 的超音速情况下的马赫数等值线。在这种情况下，在均匀网格上系数 A P 为零；即使方程中存在冲击且不存在扩散项，延迟修正方法也可以获得纯 CDS 的解。该解包含更多的振荡，但冲击得到了更好的解决。

下面介绍压力校正方法应用于高速流动的另一个例子。几何形状和边界条件如图11.5所示。它代表平面、对称会聚/发散通道的上半部分。在入口处，指定总压和焓；在出口处，所有数量均经过推断。将粘度设置为零，即求解欧拉方程。使用了五个网格：最粗的有 42 × 5 CV，最细的有 672 × 80 CV。常数马赫数线如图11.6所示。由于几何形状的变化，流动无法加速，因此在喉部后面会产生激波。激波在通过出口横截面离开之前从壁和对称平面反射两次。

在图 11.7 中，计算出的沿通道壁的压力分布与 Mason 等人的实验数据进行了比较。 （1980）。显示所有网格上的结果。在最粗的网格上，解会振荡；在所有其他网格上它都相当平滑。如在

\begin{sourceblock}
\includegraphics[page=450,trim={53.00bp 500.48bp 53.87bp 52.00bp},clip,width=\textwidth,max height=0.38\textheight]{Ferziger4th.pdf}
\par\vspace{0.45\baselineskip}
\small 图 11.6 可压缩通道流中的马赫数等值线（从入口处的最小 Ma = 0.22 到最大 Ma = 1.46，步长 0.02）；来自利莱克（1995）
\end{sourceblock}
\begin{sourceblock}
\includegraphics[page=450,trim={53.00bp 231.86bp 53.87bp 189.41bp},clip,width=\textwidth,max height=0.38\textheight]{Ferziger4th.pdf}
\par\vspace{0.45\baselineskip}
\small 图 11.7 预测 (Lilek 1995) 和测量 (Mason et al. 1980) 沿通道壁的压力分布比较
\end{sourceblock}
在前面的示例中，冲击的位置不会随着网格细化而改变，但陡度会随着网格的细化而提高。除了出口附近网格相对较粗之外，各处的数值误差都很低；两个最细网格上的结果几乎无法区分。与实验数据的吻合度也相当好。

本节介绍的求解方法往往会随着马赫数的增加而更快地收敛（除非 CDS 贡献太大以致激波处出现强烈振荡；在大多数应用中约为 90-95\%）。在图 11.8 中，求解 Re = 100 时层流不可压缩流和 Ma = 1.65 凸块上超音速流的方法的收敛性（见图 11.4）

\begin{center}\small 图 11.8 Re = 100 时层流和 Ma in = 1 时超音速流压力校正方法的收敛性。 65 超过通道中的凸起（160 × 80 CV 网格）；来自利莱克（1995）\end{center}
在一个频道中显示。两种流动使用相同的网格和欠松弛参数。虽然在可压缩情况下收敛速率几乎恒定，但在低雷诺数的不可压缩情况下，随着公差收紧，收敛速率会降低。在非常高的马赫数下，随着网格的细化，计算时间几乎随着网格点的数量线性增加（指数约为 1.1，而不可压缩流的情况下约为 1.8）。然而，正如我们将在第 1 章中演示的那样。如图12所示，使用多重网格方法可以显著提高椭圆问题方法的收敛性，使得该方法非常有效。该方法的可压缩版本适用于稳态和非稳态流动问题。

为了获得最终的精度，应该在激波附近局部应用网格细化，其中轮廓突然改变斜率。应用局部网格细化的方法以及细化网格的标准将在第 1 章中描述。 12.此外，CDS 和 UDS 的混合应仅在激波附近局部应用，而不是如上述应用中那样全局应用。决定在何处以及多少 UDS 与 CDS 混合的标准可以基于解的单调性要求、总变差递减（TVD，参见下一节）或其他合适的要求。

\section*{11.3 可压缩流动的设计方法}\phantomsection\addcontentsline{toc}{section}{11.3 可压缩流动的设计方法}
\subsection*{11.3.1 简介}\phantomsection\addcontentsline{toc}{subsection}{11.3.1 简介}
上述方法是为计算不可压缩流而设计的方法到可压缩流的处理的改编。本书多次提到有专门为求解可压缩流而设计的方法。特别是，这些方法可以与第 1 章中描述的人工压缩性方法结合使用。 7.在本节中，我们将简要描述其中一些方法。目的是提供有关这些方法的足够信息，以便与上述方法进行比较。我们不会足够详细地介绍它们，以允许读者基于它们开发代码。后一个任务需要一个单独的卷；对这种处理感兴趣的读者可以参考 Hirsch (2007) 和 Tannehill 等人的文章。 （1977）。

从历史上看，可压缩流计算方法的发展是分阶段进行的。最初（直到 1970 年左右），仅求解线性势流方程。后来，随着计算机容量的增加，人们的兴趣逐渐转向非线性势流方程，并在 20 世纪 80 年代转向欧拉方程。粘性流或纳维-斯托克斯方程（在大多数情况下为 RANS 方程，因为高雷诺数确保流动是湍流）的方法已成为过去 20 年的研究焦点。

如果说这些方法有一个主题，那就是明确认识到这些方程是双曲线的，因此具有真实的特征，有关解的信息以有限的速度传播。另一个基本问题（由特征的存在引起）是可压缩流动方程支持解中的激波和其他类型的不连续性。无粘流中的不连续性很尖锐，但当粘度非零时，不连续性的宽度有限。尊重这些属性很重要，因此在大多数方法中都明确考虑到了这一点。

这些方法主要应用于飞机、火箭和涡轮叶片的空气动力学。由于速度很高，显式方法需要使用非常小的时间步长，并且效率非常低。因此，隐式方法将是有用的并且已经被开发出来。然而，使用的许多方法都是明确的。

\subsection*{11.3.2 不连续性}\phantomsection\addcontentsline{toc}{subsection}{11.3.2 不连续性}
处理不连续性的需要引发了另一组问题。我们已经看到，在尝试捕获解中的任何类型的快速变化时，离散化方法可能会产生包含振荡或“摆动”的结果。当使用非耗散离散化（基本上包括所有中心差分方案）时尤其如此。冲击（或任何其他不连续性）代表快速变化解的极端，因此对离散化方法提出了终极挑战。可以证明，当解包含不连续性时，任何高于一阶的离散化方法都不能保证解的单调性（Hirsch 2007，第 8.3 节）。由于通过使用中心差法（或有限体积法中的等效方法）可以最好地获得精度，因此许多现代可压缩流方法在除应用特殊迎风方法的不连续点附近以外的任何地方都使用中心差。

这些都是可压缩流数值方法设计中必须面对的问题。现在，我们以一种非常笼统和肤浅的方式来审视已提出的一些处理这些问题的方法。

最早的方案基于显式方法和中心差分。其中最著名的方法之一是 MacCormack (2003)（1969 年重印）的方法，该方法至今仍在使用。为了避免此类方法中出现冲击时的振荡问题，有必要在方程中引入人工耗散。通常的二阶耗散（相当于普通粘度）会使解各处变得平滑，因此需要一个对冲击时的快速变化更敏感的项。四阶耗散项，即包含速度四阶导数的附加项，是最常见的附加项，但也使用了更高阶项。

第一个有效的隐式方法是由BeamandWarming（1978）开发的。他们的方法基于 Crank-Nicolson 方法的近似因式分解，并且可以被视为第 2 章中提出的 ADI 方法的扩展。 5和8为可压缩流。与 ADI 方法一样，该方法具有收敛到稳态解的最佳时间步长。使用中心差再次需要在方程中添加显式的四阶耗散项。

最近，人们对更复杂的迎风方案产生了兴趣。目标始终是产生明确定义的不连续性，而不会给解的平滑部分引入过多的误差。实现这一点的一种方案是 Steger 和 Warming (1981) 的通量矢量分裂方法，对此提出了许多修改和扩展。这个想法是将通量（因为应用到欧拉方程，这意味着动量的对流通量）局部分割成沿着方程的各种特征流动的分量。一般来说，这些通量沿不同方向流动。然后通过适合其流动方向的迎风方法处理每种通量。由此产生的方法相当复杂，但迎风提供了不连续处的稳定性和平滑度。

\subsection*{11.3.3 限制器}\phantomsection\addcontentsline{toc}{subsection}{11.3.3 限制器}
接下来，我们提到一类使用限制器来提供平滑且准确的解的方案。其中最早的（也是最容易解释的之一）是 Boris 和 Book (1973) 的通量校正输运 (FCT) 方法；另见库兹明等人。 （2012）。在该方法的一维版本中，可以使用简单的一阶迎风法来计算解。可以估计解中的扩散误差（一种方法是使用高阶方案并取差）。然后从解中减去该估计误差（所谓的反扩散步骤），但仅限于不产生振荡的程度。在库兹明等人中。 (2012) 值得注意的是，Zalesak (1979) 将 FCT 扩展到多维形式，而 Parrott 和 Christie (1986) 将 FCT 推广到非结构化网格上的有限元。

更复杂的方法基于类似的想法，通常称为通量限制器；见第 4.4.6 节 我们给出了一些替代方案。提醒一下，这个概念是将守恒量进入控制体积的通量限制到不会在该控制体积中产生该量的分布的局部最大值或最小值的水平。在总变差递减 (TVD) 方案（这些方法中最流行的类型之一）中，其想法是减少数量 q 的总变差，定义为：

\begin{sourceblock}
\begin{tikzpicture}
\node[inner sep=0,anchor=south west] (image) at (0,0) {\includegraphics[page=454,trim={53.00bp 486.56bp 53.87bp 152.57bp},clip,width=0.92\textwidth,max height=0.38\textheight]{Ferziger4th.pdf}};
\begin{scope}[x={(image.south east)},y={(image.north west)}]
\fill[white] (0.32806,0.46094) rectangle (0.34953,0.88893);
\fill[white] (0.35011,0.46094) rectangle (0.37323,0.88893);
\fill[white] (0.37296,0.46094) rectangle (0.43014,0.95557);
\fill[white] (0.43365,0.47094) rectangle (0.46183,0.89893);
\end{scope}
\end{tikzpicture}
\end{sourceblock}
其中 k 是网格点索引，通过限制通过控制体积面的量的通量。

这些方法已被证明能够在一维问题中产生非常干净的冲击。将它们应用于多维问题的明显方法是在每个方向使用一维版本。这并不完全令人满意，原因类似于不可压缩流的迎风方法在多个维度上不准确的原因；这个问题在第三节中讨论过。 4.7.

TVD 方案降低了不连续点附近的近似阶数。它们在不连续性本身处成为一阶，因为这是保证产生单调解的唯一近似。该格式的一阶性质意味着引入了大量的数值耗散。已经开发了另一类方案，称为本质无振荡（ENO）方案（参见第 4.4.6 节中的上下文）。它们不要求单调性，并且不降低近似的阶数，而是在不连续点附近使用不同的计算分子或形状函数；单面模板用于避免跨不连续点进行插值。

在加权 ENO 方案 (WENO) 中，定义了多个模板并检查它们产生的振荡；根据检测到的振荡类型，权重因子用于定义最终的形状函数（通常称为重建多项式）。为了计算效率，模板的数量应该少且紧凑，但是为了避免振荡同时保持高阶近似，需要在方案中使用大量的邻居。对大气边界层流动和冲击的研究表明，由于 WENO 方案在大梯度存在的情况下是耗散的，因此改变平流方案可能会导致结果发生显著变化。傅等人。 （2016，2017）通过所谓的“有针对性的 ENO 方案”的方法来解决这个问题。例如，它们产生优化的高阶（例如，六阶或八阶）方案，具有更低的耗散和更好的近冲击性能。

Abgrall (1994)、Liu 等人描述了用于非结构化自适应网格的复杂非振荡方法。 ( 1994 )、声纳 ( 1997 ) 和弗里德里希 ( 1998 ) 等。这些方案很难用隐式方法实现；在显式方法中，它们会增加每个时间步长的计算时间，但准确性和无振荡通常可以弥补较高的成本。

最后，我们提到，尽管多重网格方法是为求解椭圆方程而设计的，但它已在可压缩流动问题上取得了巨大成功。

\subsection*{11.3.4 预处理}\phantomsection\addcontentsline{toc}{subsection}{11.3.4 预处理}
还应当指出的是，刚刚描述的大多数最近的方法都是明确的。这意味着可以与它们一起使用的时间步长（或有效等效）存在限制。像往常一样，限制采用库朗条件的形式，但由于声波的存在，它具有修改后的形式：

\begin{sourceblock}
\includegraphics[page=455,trim={53.00bp 427.79bp 53.87bp 210.06bp},clip,width=0.92\textwidth,max height=0.38\textheight]{Ferziger4th.pdf}
\end{sourceblock}
其中 c 是气体中的声速，与往常一样，α 是取决于所使用的特定时间提前方法的参数。

对于仅轻微可压缩的流动，即 Ma = u / c ≪ 1，此条件简化为： c t

\begin{sourceblock}
\begin{tikzpicture}
\node[inner sep=0,anchor=south west] (image) at (0,0) {\includegraphics[page=455,trim={53.00bp 350.61bp 53.87bp 287.25bp},clip,width=0.92\textwidth,max height=0.38\textheight]{Ferziger4th.pdf}};
\begin{scope}[x={(image.south east)},y={(image.north west)}]
\fill[white] (0.00000,0.80339) rectangle (0.09561,1.00000);
\fill[white] (0.09832,0.80339) rectangle (0.13477,1.00000);
\end{scope}
\end{tikzpicture}
\end{sourceblock}
这比柯朗条件限制得多：

\begin{sourceblock}
\includegraphics[page=455,trim={53.00bp 297.33bp 53.87bp 340.53bp},clip,width=0.92\textwidth,max height=0.38\textheight]{Ferziger4th.pdf}
\end{sourceblock}
这通常适用于不可压缩流动。因此，用于可压缩流的方法在轻微可压缩流的限制下往往变得非常低效。上述压力校正方法似乎对于不可压缩和可压缩、稳态和非稳态流动都相当有效。这就是为什么它们主要用于通用商业代码，针对从不可压缩流到高度可压缩流的广泛应用。

还有一些方法最初是为可压缩流开发的，后来扩展到处理所有马赫数的流。正如第 1 节中所指出的。 11.1，控制方程在低马赫数时在数值上变得非常僵硬。为了解决这个问题，我们可以使用预处理（这个概念在第 5.3 节中讨论过）。我们简要介绍了一种此类方法的主要思想，该方法在两个商业代码中实现；更多细节可以在 Weiss 和 Smith (1995) 以及 Weiss 等人中找到。 （1999）。该方法使用时间导数的预处理，以便适用于不可压缩流；这种方法首先由 Turkel (1987) 提出。考虑耦合守恒方程组，其中压力
p 、笛卡尔速度分量 u i 和温度 T 作为解变量。预处理矩阵 K 是将耦合方程组中的各项与时间导数相乘，定义为：

\begin{sourceblock}
\includegraphics[page=456,trim={53.00bp 489.06bp 53.87bp 50.00bp},clip,width=0.92\textwidth,max height=0.38\textheight]{Ferziger4th.pdf}
\end{sourceblock}
密度相对于温度的导数是在恒定压力下求得的。对于理想气体，可得

\begin{sourceblock}
\begin{tikzpicture}
\node[inner sep=0,anchor=south west] (image) at (0,0) {\includegraphics[page=456,trim={53.00bp 418.77bp 53.87bp 210.01bp},clip,width=0.92\textwidth,max height=0.38\textheight]{Ferziger4th.pdf}};
\begin{scope}[x={(image.south east)},y={(image.north west)}]
\fill[white] (0.41976,0.47082) rectangle (0.45465,0.78051);
\end{scope}
\end{tikzpicture}
\end{sourceblock}
δ 和 δ 设置为等于 1，而对于不可压缩流体，这两个量都设置为等于 0。最重要的参数是 θ，其定义为：

\begin{sourceblock}
\includegraphics[page=456,trim={53.00bp 347.33bp 53.87bp 288.52bp},clip,width=0.92\textwidth,max height=0.38\textheight]{Ferziger4th.pdf}
\end{sourceblock}
其中 U r 代表参考速度。选择它使得系统相对于对流和扩散时间尺度的特征值保持良好的条件，即，它们被重新调整以消除方程的刚度。这是通过限制 U r 来实现的，使其值不会低于局部对流或扩散速度：

\begin{sourceblock}
\begin{tikzpicture}
\node[inner sep=0,anchor=south west] (image) at (0,0) {\includegraphics[page=456,trim={53.00bp 263.14bp 53.87bp 387.54bp},clip,width=0.92\textwidth,max height=0.38\textheight]{Ferziger4th.pdf}};
\begin{scope}[x={(image.south east)},y={(image.north west)}]
\fill[white] (0.28737,0.10414) rectangle (0.32165,0.90492);
\fill[white] (0.32878,0.17400) rectangle (0.35699,0.92238);
\fill[white] (0.36048,0.15071) rectangle (0.41687,0.89845);
\fill[white] (0.43368,0.15718) rectangle (0.48099,0.92238);
\fill[white] (0.48114,0.17400) rectangle (0.51516,0.92238);
\fill[white] (0.53525,0.15718) rectangle (0.56592,0.92238);
\fill[white] (0.56608,0.17400) rectangle (0.58397,0.92238);
\end{scope}
\end{tikzpicture}
\end{sourceblock}
第三个极限值的选择是出于数值稳定性的原因（特别是因为驻点/线区域；ε通常设置为10 - 3 ）。 x 是基于网格间距的扩散的局部长度尺度。对于可压缩流，U r 还受到局部声速 c 的限制。

使用有限体积方法，进行二阶空间离散。用于将变量插值到单元面质心的单元中心梯度受到限制，以避免振荡；使用 Barthand Jespersen (1989) 建议的方法。

时间导数的预处理会破坏时间精度，因此该方法仅适用于稳态问题，其中仅对最终解感兴趣，即，当时间导数变得等于零并且预处理不会造成任何损害时。如果瞬态问题需要一种时间精确的解，则必须使用双时间步进。然后，我们从非定常流控制方程组开始。它们各自都有一个关于“物理时间”的派生术语；其中添加了相对于伪时间的预处理时间导数。对于每个物理时间步长，在伪时间中采取几个步骤，直到解停止变化（即，达到伪时间的稳态，预处理的时间导数变为零，并且恢复原始方程）。

伪时间中的步骤对应于前面描述的顺序压力校正方案中的外部迭代。对于伪时间积分，使用一阶隐式欧拉方案，因为它允许大的时间步长并且不要求伪时间的精度。求解瞬态流问题时，需要根据精度要求选择物理时间步长的大小；使用的方案通常是二阶的，例如具有三个时间级别的隐式后向方案或克兰克-尼科尔森方案。

稳态不可压缩流动问题，如第 1 章中提出的问题。如果选择伪时间步长使得库朗数非常大（在 1000 到 10000 之间），则可以使用该耦合求解器非常有效地求解 8 和 9 的问题；如果规定较小的时间步长，则该方法不是很有效。商业代码通常提供 1 到 10 之间的库朗数默认值，这很少是最佳的。问题在于库朗数的最大可用值（通常提供最高效率）取决于问题，并且可能会相差几个数量级。

上述耦合求解器的效率（对于不可压缩流和可压缩流）在很大程度上取决于使用代数多重网格方法（参见第 6 节）求解线性化耦合方程。更多细节和一些说明性应用示例可以在 Weiss 等人中找到。 （1999）。

\section*{11.4 对申请的评论}\phantomsection\addcontentsline{toc}{section}{11.4 对申请的评论}
可压缩流动文献中有丰富的应用和论文，例如，汇集了高阶方法、大涡模拟和壁建模等。我们在此对其中的一些进行评论。

Moin 等人是最早研究可压缩流和标量大涡模拟的著作之一。 （1991）。除了使用动态 SGS 模型（参见第 10.3.3.3 节）之外，他们还根据 Favre 过滤（密度加权）变量重新构建了控制方程。正如 Bilger (1975) 指出的那样，这种平均“使连续性方程变得精确，并消除了涉及湍流通量密度涨落的双重相关性。” Farve 过滤变量定义为

\begin{sourceblock}
\includegraphics[page=457,trim={53.00bp 143.12bp 53.87bp 495.00bp},clip,width=0.92\textwidth,max height=0.38\textheight]{Ferziger4th.pdf}
\end{sourceblock}
其中上横线表示 RANS 或 LES 平均。所得到的控制方程看起来非常类似于不可压缩的 RANS 或 LES 方程，除了存在平均密度和具有时间导数的连续性方程之外。标准求解方法适用，并且衰减各向同性湍流和通道流的结果非常好。卡尼尔等人。 ( 2009 ),
见第 2.3.6 节 讨论了 Favre 过滤（注意到大多数作者都使用了这种变量变化），而 Moin 等人。 ( 1991 ) 详细介绍了动量方程和标量方程的发展。

一个活跃的研究领域是预测喷气发动机排气噪声。这引入了声学传播以及对离散化和边界条件的新约束。 Bodony 和 Lele (2008) 使用大涡模拟回顾了喷气噪声的预测，但此后已经做出了改进。布雷斯等人。 ( 2017 ) 使用非结构化网格代码将 LES 应用于超音速喷气机。豪斯曼等人。 ( 2017 ) 使用重叠网格（参见第 9.1.3 节）和混合 RANS/LES 模型来研究喷气式飞机噪声，作为开发安静超音速公务机的重点工作的一部分。布雷姆等人。 (2017) 使用隐式 LES（参见第 10.3.2 节）和改进的六阶冲击捕获 WENO 方案（第 11.3 节）来研究超音速喷气机撞击产生的噪声，作为 NASA 发动机噪声屏蔽项目的一部分，以减少社区噪声。关于流动产生的声音的预测和高速流动的数值方法的概述文章可以在 Wang 等人中找到。 （2016）和皮罗佐利（2011）。

最后，我们提请大家注意 Le Bras 等人。 ( 2017 )，他们研究了壁界可压缩流，结合了高阶数值格式（例如，六阶紧凑格式；参见第 3.3.3 节）、基于涡流粘度的 SGS 模型和壁模型（使用 Reichardt 1951（速度）和 Kader 1981（温度）分析定律；参见第 10.3.5.5 节）。

NASA 湍流建模资源 (NASA TMR 2019) 提供了将湍流模型应用到可压缩 RANS 方程中的指导。

\chapter{效率与精度的改进}
\section*{12.1 简介}\phantomsection\addcontentsline{toc}{section}{12.1 简介}
\subsection*{12.1.1 网格和流特征解析}\phantomsection\addcontentsline{toc}{subsection}{12.1.1 网格和流特征解析}
本章从数值方法的角度讨论计算效率和准确性。作为引言，我们简单看一下流动物理背景下数值解的准确性，即数值解是否准确地代表了流动物理？清楚地了解模拟的目标非常重要；特别是，人们到底想表达或导出什么流动物理学？物理学知识（通常通过理论、流动观察、尺寸分析等获得）有助于我们定义所需的网格分辨率。在第 10 章，我们展示了球体流动的模拟。很明显，给定配置的阻力系数对网格分辨率并不敏感，但添加绊线具有很大的效果（将阻力减少 50\% 以上）。另一方面，随着网格细化为光滑球体，流动分离附近的小涡流的细节发生了变化。例如，如果我们关心传热，那么那里的流动的局部变化可能会引起我们的兴趣。显然，我们需要了解模拟在捕获感兴趣的流动物理方面是否成功。

预测的流动物理会受到网格分辨率的强烈影响。瑞利-贝纳德对流（限制在两个水平板之间的流体从下方加热）就是一个例子。在该流动中，临界瑞利数表示导致从分子传导转变为对流的初始不稳定性，其波长约为实验或直接数值模拟中板之间间距的两倍。然而，在模拟中，如果网格分辨率约为板间距，则无法正确表示初始不稳定性；事实上，使用线性稳定性理论来估计初始波长，我们可能会选择波长十分之一左右的网格间距。周等人。 ( 2014 )研究了对流大气边界层并证明了初始不稳定结构的大小和“临界”湍流瑞利

\begin{sourceblock}
\includegraphics[page=460,trim={53.00bp 103.25bp 53.87bp 52.00bp},clip,width=\textwidth,max height=0.38\textheight]{Ferziger4th.pdf}
\par\vspace{0.45\baselineskip}
\small 图 12.1 垂直速度 (m/s) 5 km AGL（地面以上）和 6 小时进入天气锋模拟的彩色图，显示网格间距 ( △ ) 从 4 km 到 125 m。虚线轮廓标记了表面阵风锋。由乔治·布莱恩提供，国家大气研究中心的数字取决于“网格间距而不是流动的自然状态”。实际上，我们在模拟中看到（或可以从中解释）的流动物理取决于网格解决真实流动现象的程度；只要分辨率不够，我们看到的就会与现实（自然或实验）中看到的不同。
\end{sourceblock}
Bryan (2007) 和 Bryan 等人令人信服地证明了网格尺寸对预测流动物理的影响。 （2003）。他们报告了在 LES 模式下使用非静水压可压缩代码（与 WRF（天气研究和预报）模型中使用的求解器相同）对风暴天气锋面的深层潮湿对流进行的模拟（参见第 10.3.3 节）；疆域沿线长128公里，东西宽512公里，高18公里。对于锋面风速为 20 m/s 的情况，他们在图 12.1 中给出了距地面 5 公里处的垂直速度和模拟 6 小时的结果。结果令人吃惊，因为上升气流和下降气流（即云形式）的数量和大小在很大程度上取决于网格分辨率，直到达到 125 m； Bryan 通过能谱图（图 12.2）表明，流动在 125 m 处得到合理解析。在该图中，再次令人惊讶地注意到，对于较大的网格间距，最具能量的涡流的大小是由网格大小而不是流动物理驱动的；事实上，随着网格尺寸的增加，最有活力的上升气流尺寸渐近地缩放为网格尺寸的六倍！对这种情况的粗略分析，假设这种情况下对流单元的直径为 2 公里（基于现场数据），得出建议的水平

\begin{sourceblock}
\includegraphics[page=461,trim={53.00bp 115.34bp 53.87bp 311.23bp},clip,width=\textwidth,max height=0.38\textheight]{Ferziger4th.pdf}
\par\vspace{0.45\baselineskip}
\small 图 12.2 网格间距（△）范围为 8 km 至 125 m 的天气锋模拟中 5 km AGL 处的垂直速度谱。这里κ是波数。虚线大约终止于栅格间隔波长处。由国家大气研究中心的 George Bryan 提供 解析 LES 的网格分辨率为 O(100–200 m)，这与 Bryan 详细的网格分辨率研究一致；另见Matheou 和Chung (2014)。
\end{sourceblock}
总之，细化和改进模拟网格通常会提高准确性。然而，它也可能导致模拟流动物理学发生根本性的变化。了解预期的流动物理学可以指导正确的模拟。

\subsection*{12.1.2 组织}\phantomsection\addcontentsline{toc}{subsection}{12.1.2 组织}
求解方法效率的最佳衡量标准是达到所需精度所需的计算量。有多种方法可以提高 CFD 方法的效率和准确性；我们将提出三个足够通用的解，可以应用于前面章节中描述的任何解。

\section*{12.2 误差分析与估计}\phantomsection\addcontentsline{toc}{section}{12.2 误差分析与估计}
在流体流动问题的数值求解中不可避免的各种类型的误差已在第 4 节中进行了简要讨论。 2.5.7.在这里，我们对各种类型的误差进行更详细的讨论，并讨论如何估计和消除这些误差。代码和模型验证问题也将得到解决。

\subsection*{12.2.1 错误描述}\phantomsection\addcontentsline{toc}{subsection}{12.2.1 错误描述}
12.2.1.1 建模错误
流体流动和相关过程通常用代表基本守恒定律的积分或偏微分方程来描述。这些方程可以被认为是问题的数学模型。尽管纳维-斯托克斯方程可以被认为是精确的，但对于大多数工程兴趣流程来说，求解它们是不可能的。如果要直接模拟湍流，对计算机资源的要求很高；其他现象，如燃烧、多相流、化学过程等，很难准确描述，不可避免地需要引入建模近似。牛顿定律和傅里叶定律本身只是模型，尽管它们完全基于对许多流体的实验观察。

即使基础数学模型几乎是精确的，流体的某些特性也可能无法被准确地了解。所有流体特性都很大程度上取决于温度、物质浓度以及可能的压力；这种依赖性经常被忽略，从而引入额外的建模错误（例如，使用布辛涅斯​​克近似来计算自然对流、忽略低马赫数流动中的压缩性效应等）。

这些方程需要初始条件和边界条件。这些通常很难准确指定。在其他情况下，由于各种原因，人们被迫近似它们。通常，应将无限解域视为应用有限和人工边界条件。我们经常必须对解域入口处以及横向和出口边界处的流量做出假设。因此，即使控制方程是精确的，在边界处进行的近似也可能会影响解。

最后，几何形状可能难以准确表示；我们经常忽略难以生成网格的细节。如果不简化几何图形，许多使用结构化或块结构网格的代码就无法应用于非常复杂的问题。

因此，即使我们能够准确地求解方程和指定的边界条件，由于模型假设的错误，结果也不会准确地描述流动。因此，我们将建模误差定义为实际流动与指定几何形状、流体特性以及初始和边界条件的数学模型的精确解之间的差异。

12.2.1.2 离散化误差
此外，我们很少能够精确地求解控制方程。每种数值方法都会产生近似解，因为必须进行各种近似才能获得可以在计算机上求解的代数方程组。例如，在 FV 方法中，必须对表面和体积积分、中间位置的变量值以及时间积分采用适当的近似。显然，空间和时间离散元素越小，这些近似值就越准确。使用更好的近似值也可以提高准确性；然而，这并不是一件小事，因为更准确的近似值更难以编程，需要更多的计算时间和存储，并且可能难以应用于复杂的几何形状。通常，人们在编写代码之前选择近似值，因此空间和时间网格分辨率是用户可以控制精度的唯一参数。

同样的近似可能在流的某一部分非常准确，但在其他地方则不准确。均匀的间距（无论是在空间上还是在时间上）很少是最佳的，因为流量可能在空间和时间上局部变化很大；当变量的变化很小时，误差也会很小。因此，使用相同数量的离散元件和相同的近似值，结果中的误差可能会相差一个数量级或更多。由于计算工作量与离散元素的数量成正比，因此它们的适当分布和大小对于计算效率（实现规定精度的成本）至关重要。

我们将离散化误差定义为控制方程的精确解与离散近似的精确解之间的差异。

12.2.1.3 迭代错误
离散化过程通常会产生一组耦合的非线性代数方程。这些通常是线性化的，并且线性化方程也通过迭代方法求解，因为直接求解通常太昂贵。

任何迭代过程都必须在某个阶段停止。因此，我们必须定义一个收敛标准来决定何时停止该过程。通常，迭代会持续进行，直到残差水平减少特定量为止；这可以证明相当于将误差减少了类似的量。

即使求解过程收敛并且迭代足够长的时间，我们也永远无法获得离散方程的精确解；由于计算机的有限算术精度而导致的舍入误差将提供误差的下限。幸运的是，直到解误差接近计算机的算术精度并且比通常需要的精度高得多之前，舍入误差才成为问题。

我们将迭代误差定义为离散方程的精确解和迭代解之间的差异。尽管这种误差与离散化本身无关，但随着离散元素数量的增加，将误差减小到给定大小所需的工作量也会增加。因此，有必要选择最佳的迭代误差水平——与其他误差（否则无法评估）相比足够小，但不能更小（因为成本会大于必要的）。

12.2.1.4 编程和用户错误
人们常说所有计算机代码都有错误——这可能是真的。代码开发人员有责任尝试消除它们；我们将在这里讨论的一个问题。通过研究代码来定位编程错误是很困难的，更好的方法是设计测试问题，其中可能会出现由错误引起的错误。在将代码应用于常规应用之前，必须仔细检查测试计算的结果。人们应该检查代码是否以预期的速率收敛，误差是否以预期的方式随着离散元素的数量而减少，并且解与通过分析或由另一代码生成的可接受的解一致。

代码的关键部分是边界条件的实现。必须检查结果以确定所应用的边界条件是否确实满足；发现它们并非如此并不罕见。 Peri ́c (1993) 讨论了一个与自然对流中的绝热边界相关的问题。问题的另一个常见根源是紧密耦合的项的近似值不一致。例如，在静止的气泡中，自由表面上的压降必须通过表面张力来平衡。已知解析解的简单流程对于计算机代码的验证非常有用。例如，使用移动网格的代码可以通过移动内部网格同时保持边界固定并使用静止流体作为初始条件来检查；流体应保持静止，并且不应受到网格运动的影响。

解的准确性不仅取决于离散化方法和代码，还取决于代码的使用者；即使有好的代码也很容易得到不好的结果！尽管大多数用户错误导致的错误属于上述三类之一，但区分系统错误（方法本身存在）和可避免的错误（由于代码中的错误或代码使用不当或不当造成的）很重要。

许多用户错误是由于输入数据不正确造成的；通常，只有在进行多次计算后才能发现错误，有时甚至从未发现！当使用方程的无量纲形式时，经常出现的错误是由于几何缩放或参数选择造成的。另一种用户错误是由于数值网格质量不佳造成的（网格点分布不充分可能会使误差增加一个数量级或更多，或者根本无法获得解）。

\subsection*{12.2.2 误差估计}\phantomsection\addcontentsline{toc}{subsection}{12.2.2 误差估计}
每个数值解都包含错误；重要的是要知道错误有多大，以及它们的级别在特定应用程序中是否可以接受。可接受的错误程度可能有很大差异。在新产品早期设计阶段的优化研究中，只有定性分析和系统对设计变更的响应很重要，在优化研究中可接受的错误可能在其他应用中可能是灾难性的。

因此，了解解对于特定应用的效果与首先获得解同样重要。特别是在使用商业代码时，用户应尽可能集中精力仔细分析结果并估计错误。这对于初学者来说可能是一个很大的负担，但经验丰富的 CFD 从业者会经常这样做。

错误分析应该按照与上面引入错误的顺序相反的顺序进行。也就是说，应该首先估计迭代误差（可以在一次计算中完成），然后是离散化误差（需要在不同网格上至少进行两次计算），最后是建模误差（需要参考数据，可能还需要许多计算）。其中每一个都应该比前一个小一个数量级，否则后面的误差的估计将不够准确。

12.2.2.1 迭代误差估计
因为纳维-斯托克斯方程是非线性的，所以我们有两个迭代循环（参见图 7.6）：求解特定变量的线性化（可能是解耦）方程组时的内部迭代，以及更新线性方程组和右侧系数的外部迭代。

从计算效率的角度来看，知道何时停止迭代过程至关重要。对于内部迭代，迭代次数过多是没有意义的，因为在正确求解非线性耦合方程组之前，需要多次更新矩阵系数和右侧。大多数情况下，在系数更新之前将残差级别降低一个数量级就足够了；更长的迭代不会减少所需的外部迭代次数，因此只会导致更长的计算时间。另一方面，如果内部迭代停止得太早，则将需要更多的外部迭代，从而再次增加计算量。像往常一样，最优值取决于问题。

控制外部迭代更为关键：当矩阵系数和右侧的更新导致解的变化可以忽略不计时，离散非线性方程就可以正确求解。根据经验，外部迭代误差（有时也称为收敛误差）应至少比离散化误差低一个数量级。迭代到舍入级别是没有意义的；对于大多数工程应用，任何变量中三到四位有效数字的相对精度（与参考值相比的误差）就足够了。

有多种方法可以估计这些误差； Ferziger 和 Peri ́c (1996) 详细分析了其中的三个；另见第 5.7 节.可以看出，除了迭代的初始阶段外，误差的减少率与残差和连续迭代之间的差异的减少率相同。这在图 8.9 中得到了证明：在一些迭代之后，残差范数、连续迭代之间的差异范数、估计误差和实际迭代误差的曲线都是平行的。请注意，对于外部迭代，使用线性化方程的当前解和更新的矩阵系数和右侧计算的残差是相关的（即，在新的内部迭代循环开始时计算的残差）。通过减去两个连续循环的最后一次内部迭代的值，可以获得解的相关差异。

因此，如果知道计算开始时的误差水平（如果从零字段开始，则为解本身，如果进行粗略但合理的猜测，误差水平会稍低），那么如果残差范数（或两次迭代之间的差异）下降 3-4 个数量级，则可以确信误差将下降 2-3 个数量级。这意味着前两个或三个最高有效数字在进一步的迭代中不会改变，因此解的准确度在 0.01-0.1\% 范围内。

上述陈述适用于稳态问题的解决。在解决非稳态问题时，迭代误差的估计有点复杂。在显式方法的情况下，只需确保压力或压力校正方程被求解到足够严格的公差，以确保质量守恒方程得到充分满足；通常将残差减少三个数量级就足够了。在隐式方法的情况下，如果时间步非常小（如 LES 模拟中的情况），则可能不需要对外部迭代要求如此严格的标准，因为解从一个时间步到另一个时间步没有太大变化，因此在每个时间步内将残差水平降低三个数量级可能是一种矫枉过正。在这种情况下，3-5 次外部迭代可能足以更新非线性和耦合效应。建议在任何新的应用领域测试改变收敛标准的效果，以确保迭代误差足够小。

一个常见的错误是查看连续迭代之间差异的大小，并在差异不超过某个小数时停止计算。然而，差异可能很小，因为迭代缓慢收敛，而迭代误差可能很大。为了估计误差的大小，必须对连续迭代之间的差异进行适当的归一化；当收敛缓慢时，归一化因子变大（参见第 5.7 节）。另一方面，要求差异范数下降三到四个数量级通常是一个安全的标准。由于大多数 CFD 方法中的线性方程求解器需要计算残差，因此最简单的做法是监控其范数（绝对值之和或平方和的平方根）。

在粗网格上，离散化误差较大，可以允许较大的迭代误差；精细网格需要更严格的公差。如果收敛标准基于残差总和而不是每个节点的平均残差，则会自动考虑这一点，因为总和随着节点数量的增加而增长，从而收紧收敛标准。

当开发新代码或添加新功能时，必须毫无疑问地证明求解过程确实收敛，直到残差达到舍入水平。通常，缺乏这种收敛表明存在错误，尤其是在边界条件的实现中。有时，该限制低于宣布收敛的阈值，并且问题可能不会被注意到。在其他情况下，该过程可能会更早地停止收敛（甚至发散）。一旦所有新功能都经过彻底测试，就可以返回到通常的收敛标准。

此外，如果试图为本质上不稳定的问题获得稳定的解（例如，以存在冯卡门涡街的雷诺数围绕圆柱体流动），迭代将不会收敛。因为每次迭代都可以解释为伪时间步（参见第 7.2.2.2 节），所以过程很可能不会发散，但残差会无限振荡。如果几何形状是对称的并且稳定对称解不稳定（例如，扩散器或突然膨胀；稳态解（层流和雷诺平均）通常是不对称的，一侧具有较大的分离区域），则通常会发生这种情况。人们可以通过减少雷诺数或计算一半几何形状的流动、使用对称边界条件或通过执行瞬态计算来检查这是否是问题。特别是在复杂的几何形状中，在求解域的一小部分（例如，汽车后视镜后面）中，流动可能局部不稳定。在这种情况下，残差可能会降至通常的收敛水平以下，但如果试图进一步减少残差，它们将在某个阶段开始振荡。通常，不稳定性非常弱，如果继续不稳态计算，积分量（如力、力矩、总热通量等）可能不会随时间发生明显变化，但稳态计算不会收敛。

\begin{center}\small 图 12.3 和 12.4 显示了当尝试计算壁挂式障碍物周围的稳态湍流时可能出现的问题示例。残差在相同（过高）值附近振荡，没有任何减少的迹象\end{center}
\begin{sourceblock}
\includegraphics[page=468,trim={53.00bp 365.25bp 53.87bp 52.00bp},clip,width=\textwidth,max height=0.38\textheight]{Ferziger4th.pdf}
\par\vspace{0.45\baselineskip}
\small 图12.3 壁挂式障碍物周围流动模拟：瞬态模拟得到障碍物和底壁（上）的瞬时压力分布以及障碍物中高处平行于底壁的截面（下）的速度矢量
\end{sourceblock}
当尝试进行稳态计算时。当我们在 2000 次迭代后切换到瞬态模拟时，每个新时间步开始时的残差仍保持在较高水平，并且没有表现出减小的趋势，但在每个时间步内，外部迭代收敛得很好：每个时间步仅进行 5 次外部迭代，动量方程中的残差下降了两个以上数量级。显然，流动没有稳态解，因为障碍物后面的尾流是不稳定的，从速度矢量可以看出（它们不对称）。

稳态和瞬态模拟中障碍物上的阻力和升力如图 1 和 2 所示。 12.5 和 12.6。在稳态计算中，阻力系数在比瞬态模拟情况下低得多的值附近振荡。在这两种情况下，升力都在零附近振荡，但振幅几乎是稳态计算情况下的两倍。显示这些图的原因是警告 CFD 代码的用户：如果稳态计算中的残差开始在高于通常收敛标准的水平上振荡，则没有控制方程的解，并且不应尝试从变量可视化或平均振荡力来解释图片。只有切换到瞬态仿真，才能在每个时间步结束时获得可进行物理解释的控制方程的有效解。现在可以对力或其他积分参数的振荡进行平均（例如，评估平均阻力或传热系数）或进行其他处理（例如，获得振荡频率、平均值周围波动的均方根值等）。

\begin{sourceblock}
\includegraphics[page=469,trim={53.00bp 357.25bp 53.87bp 52.00bp},clip,width=\textwidth,max height=0.38\textheight]{Ferziger4th.pdf}
\par\vspace{0.45\baselineskip}
\small 图 12.4 壁挂式障碍物周围的流动模拟：稳态计算（上）和瞬态模拟（下）中的残差变化
\end{sourceblock}
当细化网格或从低阶离散化方案切换到高阶离散化方案时，遇到收敛问题并不罕见。原因是，当流动的不稳定性很弱时，离散误差可能会引入足够的阻尼（例如，一阶迎风格式的数值扩散）并且迭代可能会收敛到稳态。流动不稳定通常与分离有关，并且只有在网格充分细化后，小的分离区域（例如，在机翼的吸力侧）才可能出现。在任何情况下，如果瞬态计算的每个时间步内的外部迭代收敛并且在稳态计算残差中振荡，则流动本质上是不稳定的并且应该这样计算。如果外部迭代在瞬态模拟的每个时间步长内不收敛，原因可能是 (i) 时间步长太大，(ii) 欠松弛因子太高，或 (iii) 模拟设置中的错误（网格质量、边界条件、流体属性等）。

12.2.2.2 离散化误差估计
只有比较系统细化网格上的解才能估计离散化误差；见第 3.11.1.2 和 3.9 节 了解更多详细信息。如前所述，这些误差是由于对方程中的各项项使用了近似值造成的

\begin{sourceblock}
\includegraphics[page=470,trim={53.00bp 357.25bp 53.87bp 52.00bp},clip,width=\textwidth,max height=0.38\textheight]{Ferziger4th.pdf}
\par\vspace{0.45\baselineskip}
\small 图 12.5 壁挂障碍物周围的流动模拟：稳态计算（上）和瞬态模拟（下）中阻力系数随迭代的变化
\end{sourceblock}
和边界条件。对于平滑解的问题，近似的质量用阶数 来描述，阶 是将近似的截断误差与网格间距的某个幂相关联；如果空间导数的截断误差与 ( x ) p 成正比，我们就说该近似是 p 阶的。阶数并不是误差大小的直接衡量标准；它表明当间距改变时误差如何变化。相同阶数的近似值在给定网格上可能存在误差，误差可达一个数量级；此外，对于特定网格，低阶近似可能比高阶近似具有更小的误差。然而，随着间距变小，高阶近似肯定会变得更加准确。

使用泰勒级数展开很容易找到许多近似值的阶。另一方面，不同的近似值可能用于不同的项，因此求解方法的顺序作为一个整体可能并不明显（通常是方程中重要项的最不精确近似的顺序）。此外，计算机代码中算法的实现中的错误可能会产生与预期不同的顺序。因此，使用实际代码检查每类问题的方法顺序非常重要。

分析结构化网格上的离散化误差的最佳方法是将每个方向的间距减半。然而，这并不总是可能的；在 3D 中，这导致节点数量增加八倍。因此，第三格有 64 倍

\begin{sourceblock}
\includegraphics[page=471,trim={53.00bp 351.25bp 53.87bp 52.00bp},clip,width=\textwidth,max height=0.38\textheight]{Ferziger4th.pdf}
\par\vspace{0.45\baselineskip}
\small 图 12.6 壁挂式障碍物周围流动的模拟：稳态计算（上）和瞬态模拟（下）中升力系数随迭代的变化
\end{sourceblock}
点作为第一个，我们可能无法承担另一个细化级别。另一方面，误差通常不是均匀分布的，因此细化整个网格是没有意义的。此外，当使用具有任意控制体积或元素的非结构化网格时，不存在局部坐标方向，并且元素以不同的方式细化。

重要的是细化是实质性的、系统性的。将一个方向上的节点数量从 112 个增加到 128 个并不是很有用，除非是在具有均匀误差分布和均匀网格的学术问题中；细化后的网格在每个方向上应比原始网格多至少 50\% 的节点（或者，网格间距应至少减少 1.5 倍）。系统细化意味着网格拓扑和网格点的相对空间密度应在所有网格级别上保持可比性。在不改变节点数量的情况下，不同的网格点分布可能会导致离散化误差发生实质性变化。图 8.10 显示了一个例子：在非均匀网格（靠近墙壁更精细）上获得的 ψ max 结果比在具有相同节点数的均匀网格上获得的结果准确一个数量级。两个解都收敛到具有相同阶数（第二个）的相同网格独立解，但误差的大小相差 10 倍或更多（对于具有相同节点数的网格）！具有类似结论的其他例子如图 1 和 2 所示。 6.4、6.6 和 8.17。

上面的例子强调了良好的网格设计的重要性。对于实际工程应用来说，网格生成是最耗时的任务；生成任何网格通常都很困难，更不用说高质量的网格了。一个好的网格应该尽可能接近正交（请注意，正交性在不同的方法中具有不同的含义；在 FV 方法中，单元面法线与连接相邻单元中心的线之间的角度才是重要的 - 从这个意义上说，四面体网格可能是正交的）。在预计会有大截断误差的地方，它应该是密集的——因此网格设计者应该了解一些关于解的信息（如第 12.1.1 节中所建议的）。这是最重要的标准，最好通过使用具有局部细化的非结构化网格来满足。其他质量标准取决于所使用的方法（网格平滑度、纵横比和扩展比等）。有关电网质量测量的更多详细信息将在第 4 节中给出。 12.3.

估计离散化误差的最简单方法是基于理查森外推法，并假设可以在足够精细的网格上进行计算以获得单调收敛。 （如果情况并非如此，则误差很可能比希望的要大。）因此，只有当两个最细的网格都足够细并且已知误差减少的顺序时，该方法才是准确的。如果所有三个网格在上述意义上都足够精细，则可以根据以下公式从三个连续网格上的结果计算顺序（参见 Roache 1994；以及 Ferziger 和 Peri ́c 1996，了解更多详细信息）：

\begin{sourceblock}
\begin{tikzpicture}
\node[inner sep=0,anchor=south west] (image) at (0,0) {\includegraphics[page=472,trim={53.00bp 317.67bp 53.87bp 299.52bp},clip,width=0.92\textwidth,max height=0.38\textheight]{Ferziger4th.pdf}};
\begin{scope}[x={(image.south east)},y={(image.north west)}]
\fill[white] (0.47708,0.57426) rectangle (0.51961,0.83248);
\fill[white] (0.52403,0.57120) rectangle (0.60568,0.83248);
\fill[white] (0.35047,0.17120) rectangle (0.37026,0.40735);
\fill[white] (0.37468,0.17651) rectangle (0.40286,0.41267);
\end{scope}
\end{tikzpicture}
\end{sourceblock}
其中 r 是网格密度增加的因子（如果间距减半，则 r = 2），而 φ h 表示平均间距为 h 的网格上的解。然后离散化误差估计为（详细信息请参见第 3.9 节）：

\begin{sourceblock}
\includegraphics[page=472,trim={53.00bp 236.70bp 53.87bp 401.05bp},clip,width=0.92\textwidth,max height=0.38\textheight]{Ferziger4th.pdf}
\end{sourceblock}
因此，当间距减半时，对于二阶方法，一个网格上的解的误差等于该网格上的解与前一网格上的解之间的差异的三分之一；对于一阶方法，误差等于上述差值。

对于图 8.10 中的示例，应用于均匀网格和非均匀网格的理查森外推法会导致在五个有效数字内对与网格无关的解进行相同的估计，尽管解中的误差存在一个数量级的不同。

请注意，可以计算积分量（阻力、升力等）以及场值的误差，但所有量的收敛顺序可能不同。它通常等于具有平滑解的问题（例如层流）的理论阶数（例如，第二阶）。当复杂模型（湍流、燃烧、两相流等）或使用开关或限制器的方案时，阶数的定义可能很困难。然而，计算诸如 p 阶或 Roache (1994) 所说的网格收敛指数之类的量并不是绝对必要的；足以显示一系列网格（最好是三个）的计算感兴趣量的变化。如果变化是单调的，并且差异随着网格细化而减小，则可以轻松估计与网格无关的解所在的位置。当然，只要有可能，就应该使用理查森外推法来估计与网格无关的解。

还要注意，网格细化不需要扩展到整个域。如果估计表明某些区域的误差比其他区域小得多，则可以使用局部细化。对于物体周围的流动尤其如此，只有在物体附近和尾流中才需要高分辨率。使用局部网格细化策略的方法总是比需要在整个域或网格块中细化的方法更有效。对于非结构化网格，最粗糙的网格通常涉及局部细化。然而，需要小心；如果在误差源很大（即大截断误差）的情况下不对网格进行细化，则细化的效果可能不会很大，因为误差与变量本身一样受到相同的传输过程（对流和扩散）的影响。

估计离散化误差时出现困难的两个原因值得一提。一是与壁函数相关，通常用于计算湍流；见第 10.3.5.5 节。在这种情况下，我们不会指定壁上的唯一边界条件，而是将壁剪切应力与壁旁边的单元中心的速度联系起来，并假设该位置位于边界层的对数范围内。对数定律通常在受复杂壁形状影响的流动中并不严格有效（即，当壁法向和壁切向方向都存在显著梯度时）。当网格细化时，靠近墙壁的单元质心的位置发生移动，因此动量方程中的边界条件也有效地改变；这通常会导致积分量的变化，不遵循所使用的离散化方案的预期行为。如果网格经过细化，使得靠近壁的单元的质心落入缓冲层 (5 < n + < 30)，则尤其如此。如果必须使用壁函数，应确保所有网格中靠近壁的计算点都保持在对数范围内（n + > 30）；另一种方法是使所有网格的靠近壁的第一棱镜层的厚度保持固定，使得近壁单元处的n + 保持相同，并且仅在切向方向上细化紧邻壁的网格。

当边界层由网格解析时（所谓的“低Re”方法，即在近壁单元中心处n + ≈ 1），动量方程中使用无滑移条件，这使得边界条件是唯一的。在这种情况下，积分量（力、力矩、传热等）的变化通常更有利于使用理查森外推法进行误差估计。然而，必须确保所有网格上的所有近壁单元的 n + < 2。如果 n + 值在中间范围内，则使用所谓的“ally +”版本的壁函数比使用“high-Re”版本更好，但最好确保当使用壁函数时，n + 在任何网格的大部分壁表面上都高于 30。

另一个问题与所使用的网格的几何特征的分辨率相关，尤其是壁曲率。一个典型的例子是涡轮机的前沿

\begin{sourceblock}
\includegraphics[page=474,trim={53.00bp 487.25bp 53.87bp 52.00bp},clip,width=\textwidth,max height=0.38\textheight]{Ferziger4th.pdf}
\par\vspace{0.45\baselineskip}
\small 图12.7 计算网格中螺旋桨叶片前缘的分辨率：壁切向网格过粗（左），导致前缘角部粗糙，通过局部细化进行改进（右）
\end{sourceblock}
机器叶片（例如船舶螺旋桨、风扇、水轮机或燃气轮机叶片、风力或潮汐涡轮机叶片等）。曲率通常很高，以至于最大和最小压力位置之间的距离很短，这不仅在壁法线方向上需要精细网格，而且在壁切线方向上也需要精细网格。然而，商业软件中的自动网格生成工具可能无法充分解析前缘曲率，除非在网格生成过程中特别注意（例如，通过指定沿前缘的网格间距，或要求通过圆上一定数量的点解析曲率）。这可能导致前缘在粗网格上变得锋利而不是圆形。在这种情况下，网格细化有效地导致了解域几何形状的变化，这再次使得某些量的变化不符合预期，从而使得理查森外推法的使用变得困难。因此，重要的是初始网格已经尽可能地表示解域的几何形状。一个例子如图12.7所示。

12.2.2.3 建模误差估计
建模误差是最难估计的；为此，我们需要真实流量的数据。在大多数情况下，数据不可用。因此，通常仅针对一些有详细且准确的实验数据或存在准确的模拟数据（例如大涡或直接数值模拟数据）的测试案例来估计建模误差。无论如何，在将计算与实验进行比较之前，应该分析迭代和离散化误差并证明其足够小。在某些情况下，建模和离散化误差会相互抵消，因此粗网格上的结果可能比细网格上获得的结果更符合实验数据。因此，不应使用实验数据来验证代码；必须对结果进行系统分析。只有当毫无疑问地证明结果确实收敛于与网格无关的解并且离散化误差足够小时，才能进行数值解与实验数据的比较。

值得注意的是，网格依赖性研究用于在尝试量化建模误差之前估计离散化误差，因此必须非常仔细地进行。如果网格设计不当，某些特征可能不会在所使用的任何网格中显示出来，并且在不同网格上获得的解之间的微小差异可能会导致离散化误差很小的结论。最精细网格上的模拟结果与实验之间的差异将被错误地解释为建模误差。一个例子将在第 4 节中展示。 13.8：除非在尖端涡流占据的空间内局部高度细化网格，否则不会捕获尖端涡流空化，因为否则涡流核心中的低压会被低估。如果网格依赖性研究中使用的所有网格都不能很好地解析尖端涡流以准确捕获核心中的低压，则推力和扭矩可能仍然看起来很好地收敛，但尖端涡流空化将会丢失。让空化模型对缺失的尖端涡流空化负责是不公平的，因为当网格被充分局部细化时，结果是相当好的。在其他情况下也会遇到类似的问题，例如，湍流模型被归咎于模拟与实验之间的差异，而其中很大一部分可能是由于网格不足造成的。大家都知道，网格在壁法线方向上需要精细才能捕获边界层特征，但弯曲壁、剪切层、旋流或二次流通常需要在某些区域进行局部切向方向细化，这对于没有经验的商业软件用户来说可能并不明显。

重要的是要记住，实验数据也只是近似值，测量和数据处理误差可能很大。它们还可能包含重大的系统错误。然而，它们对于模型的验证是不可或缺的。人们只能将计算结果与高精度的实验数据进行比较。如果要将实验数据用于验证目的，则必须仔细分析实验数据。

还应该注意的是，不同数量的建模误差是不同的；例如，计算出的压力阻力可能与测量值非常一致，但计算出的摩擦阻力可能基本上是错误的。平均速度分布有时可以很好地预测，而湍流量可能会低估或高估两倍。为了确保模型确实准确，将结果与各种数量进行比较非常重要。

12.2.2.4 编程和用户错误的检测
一种难以量化的错误是编程错误。这些可能是简单的“错误”（不会阻止代码编译的键入错误）或严重的算法错误。对迭代和离散化错误的分析通常可以帮助开发人员找到它们，但有些错误可能非常一致，以至于多年来（如果有的话）都未被发现，特别是当没有确切的参考解可供比较时。

对结果进行批判性分析对于发现潜在的用户错误至关重要；因此，用户对一般流体动力学以及要解决的特定问题具有扎实的知识至关重要（参见第 12.1.1 节）。即使正在使用的 CFD 代码已在其他流上得到验证，用户在设置模拟时也可能会出错，从而导致结果可能存在明显错误（例如，由于几何表示、边界条件、流参数等方面的错误）。用户错误可能很难发现（例如，当缩放发生错误并且计算出的流量对应于与预期不同的雷诺数时）；因此，如果可能的话，应该由执行计算的人以外的其他人对结果进行严格评估。

\subsection*{12.2.3 CFD 不确定性分析的推荐做法}\phantomsection\addcontentsline{toc}{subsection}{12.2.3 CFD 不确定性分析的推荐做法}
12.2.3.1 CFD 代码的验证
任何新代码或添加的功能都应经过系统分析，目的是评估离散化误差（空间和时间），定义收敛标准以确保较小的迭代误差，并消除尽可能多的“错误”。为此，必须选择一组代表代码可解决的问题范围的测试用例，并且可以提供足够准确的解（分析或数值）。因为要确保在指定的边界条件下正确求解方程，所以实验数据并不是衡量数值解质量的最佳方法。需要参考解来定位算法或编程中的错误，这些错误可能会通过与迭代和离散化错误估计相关的测试。请注意，设计网格时应避免某些项为零的情况，因为在这种情况下可能不会出现实现错误；因此，即使在简单几何中使用笛卡尔网格时，旋转解域以使网格线不与笛卡尔坐标对齐（避免单元面上的三个表面矢量分量中的两个为零）也是有用的。此外，还必须确保解独立于解域相对于坐标系方向的方向。

人们应该首先分析离散化中使用的近似值，以确定解向网格（或时间步长）独立解的收敛顺序。这是方程中重要项的最低阶截断误差（但请注意，并非所有项都同等重要 - 它们的重要性取决于问题）。在某些情况下，可以在边界处使用比内部低阶的近似值，而不降低总体阶数。一个例子是对边界处的扩散通量使用单侧一阶近似，而在内部使用二阶中心差；总体收敛是二阶的。然而，如果低阶近似与诺依曼型边界条件一起使用，则情况可能并非如此。

接下来应该分析迭代错误；作为第一步，应该进行计算，其中继续迭代，直到其级别降低到双精度舍入级别（这需要残差至少减少 12 个数量级）。必须选择具有已知稳态解的测试用例。否则，迭代可能会在某个阶段停止收敛，因为迭代可以被解释为伪时间步长，并且流的自然不稳定性可能不允许稳定的解。一个例子是雷诺数超过 50 时围绕圆柱体的流动。一旦获得给定网格上离散方程的精确解，就可以将其与中间阶段的解进行比较，从而评估迭代误差。可以将误差与估计进行比较，或者它们的减少可以与残差的减少或连续迭代之间的差异相关，如上所述。这应该有助于建立收敛标准（对于内部迭代，即线性方程求解器，以及外部迭代，即非线性方程的解）。

应通过比较一系列系统化细化网格和时间步长上的解来分析离散化误差。对于结构化或块结构网格来说，系统细化很容易：例如创建三个不同大小的网格。对于非结构化网格，这项任务并不那么简单，但可以创建具有相似的相对网格大小分布但绝对大小不同的网格。系统细化对于高截断误差区域至关重要，这些区域是离散化误差的来源，这些误差以与因变量本身相同的方式进行对流和扩散。根据经验，当解的二阶和高阶导数较大时，网格必须精细且系统地细化。这通常发生在墙壁附近、剪切层和尾流中。

应在至少三个网格上获得迭代误差足够小的解并进行比较；如果网格足够细以便单调收敛，则可以用这种方式估计收敛阶数和离散化误差。如果情况并非如此，则需要进一步细化。如果计算出的顺序不是预期的顺序，则说明离散化或编程中出现了错误，必须予以纠正。估计的离散化误差应与所需的精度进行比较。

必须对与代码所针对的应用程序类似的许多测试用例重复此过程，以便尝试根除尽可能多的错误源。仅当对代码产生的结果进行系统分析并获得网格和时间步独立解（在离散化误差已被可靠估计且足够小的意义上）时，才应将解与解析解或其他参考解进行比较。这是对编程或算法错误的最终检查。将一个网格上获得的解与参考解进行比较是没有意义的，因为通常某些数量可能会意外地吻合，或者某些错误可能会被抵消。

代码验证没有说明数值精确解代表真实流量的准确性。无论我们使用哪种湍流（或其他）模型，我们都必须确保我们正在求解正确包含模型的方程。不同小组使用相同的网格和相同的湍流模型但不同的代码获得的解的比较往往比一组使用相同的代码但不同的湍流模型时显示出更大的差异（这是 20 世纪 90 年代许多研讨会得出的结论）。模型通常有不同的实现方式，边界条件的处理方式也不同，等等。这是一个难题，目前还没有找到令人满意的解。这些差异可能是由于实现的差异造成的，但是如果使用的模型确实相同，实现是正确的并且错误已被评估和消除，那么每个代码都应该产生相同的结果，并且差异应该消失。这就是为什么我们强调验证和错误评估的必要性。

12.2.3.2 CFD 结果验证
CFD结果的验证包括离散化和建模误差分析；我们可以假设使用经过验证的代码并具有适当的收敛标准，以便可以排除迭代错误。

影响CFD结果准确性的最重要因素之一是数值网格的质量。请注意，即使网格很差，如果足够细化，也应该产生正确的解；它只会花费更多。此外，即使是最好的代码也可能在糟糕且不够精细的网格上产生较差的结果，而如果针对要解决的问题调整网格，基于更简单且不太准确的近似的代码可能会产生出色的结果。 （然而，这通常是使各种误差相互抵消的问题。）可以通过网格点的适当分布来减少离散化误差；见图8.10。

许多商业代码已经足够强大，可以在用户可能提供的任何网格上运行。然而，鲁棒性通常是以牺牲准确性为代价来实现的（例如，通过使用迎风近似）。粗心的用户可能不太关注电网质量，从而不费吹灰之力就得到不准确的解。在网格生成、误差估计和优化方面投入的精力应与解所需的准确度水平相关。如果仅寻求流量的定性特征，快速工作可能是可以接受的，但为了以合理的成本获得定量准确的结果，需要高网格质量。

与实验数据的比较需要知道实验的不确定性。最好仅将完全收敛的结果（迭代和离散化误差已被高度消除的结果）与实验数据进行比较，因为这是评估模型效果的唯一方法。实验不确定性条通常延伸到报告值的两侧。如果离散化误差与实验不确定性相比不小，则无法了解所使用模型的价值。建模误差的估计是 CFD 中最困难的任务。

在许多情况下，确切的边界条件是未知的，必须做出假设。例如物体周围流动的远场条件和入口湍流特性。在这种情况下，必须在相当大的范围内改变关键参数（远场边界的位置、湍流量），以估计解对此因素的敏感度。通常，对于合理的参数值可以获得良好的一致性，但是，除非实验数据提供了这些值，否则这只不过是复杂的曲线拟合。这就是为什么必须选择提供所有必要数量的实验数据，并与测量人员讨论获取这些数据的重要性。一些湍流模型对入口和自由流湍流水平非常敏感，导致参数相对较小的变化导致结果发生重大变化。

如果要研究相同几何形状的多个变体，通常可以依赖对典型代表性案例进行的验证。可以合理地假设相同的网格分辨率和相同的模型将产生与测试用例中相同数量级的离散化和建模误差。虽然在许多情况下确实如此，但情况可能并不总是如此，需要小心。几何形状的变化可能导致出现新的流动现象（分离、二次流、不稳定等），而所使用的模型可能无法捕获这些现象。因此，尽管几何形状的变化可能很小，但从一种情况到另一种情况，建模误差可能会急剧增加（例如，对于一个阀门开度，发动机阀门周围的流动计算可能精确到 3\% 以内，而对于稍小的开度，则在质量上是错误的；有关更详细的描述，请参阅 Lilek 等人（1991）。

12.2.3.3 一般建议
定义用于验证 CFD 代码和结果的严格规则很困难，有时甚至不切实际。虽然建议尽可能使用理查森外推法来估计离散化误差，但可能很难获得所有问题的结论性答案（例如，不同数量的顺序可能会有所不同）。正如我们在第 2 节中指出的那样。 10.3.3.7 评估 LES 的模拟质量是一项挑战； Sullivan 和 Patton (2011) 提出了 LES 质量评估的详细示例。

许多期刊和专业组织制定了自己的规则和指南，用于评估和量化 CFD 解的不确定性；例子是：

\begin{itemize}
\item 美国机械工程师协会 (ASME) 制定了验证和确认标准 ( \url{https://www.asme.org/products/codes-standards/v-v-20-2009-standards-verification-validation} ) 并定期组织有关该主题的会议 ( \url{https://event.asme.org/V-V} ) [Celik 等人。 2008 年从 ASME 角度总结了这些程序]；
\item 美国航空航天学会 (AIAA) 还制定了 CFD 模拟的验证和确认标准（指南：计算流体动力学模拟的验证和确认指南 ( \url{https://doi.org/10.2514/4.472855}；AIAA G-077-1998(2002))
\item 国际拖曳水箱会议 (ITTC) 发布了一套指南用于在海事工程中使用 CFD（例如，参见 \url{https://ittc.info/media/4184/75-03-01-01.pdf} ）
\end{itemize}
理查森外推法是所有误差估计程序的主要成分，但许多此类指南更进一步，评估更广泛意义上的不确定性。我们不会详细介绍这些指南，但建议在特定应用领域工作时应遵循这些指南。

当采用多种类型的模型（对于湍流、两相流、自由表面效应等）时，可能很难将不同的效应彼此分开。然而，任何定量 CFD 分析中最重要的步骤可以总结为：

\begin{itemize}
\item 生成适当结构和精细度的网格（在流量和壁曲率快速变化的区域进行局部细化）。
\item 系统地细化网格（非结构化网格可以有选择地细化：如果误差很小，则不需要细化）。
\item 在至少三个网格上计算流并比较解（确保迭代误差很小）；如果收敛不单调，则再次细化网格。估计最精细网格上的离散化误差。
\item 如果可用，将数值解与参考数据进行比较以估计建模误差。
\end{itemize}
对数值误差的任何合理估计总比没有好；并且由于数值解始终是近似解，因此人们必须始终质疑其准确性。

许多教育组织提供关于一般不确定性量化（特别是 CFD）的专业课程或进行研究。例如斯坦福大学的 UQLab (\url{http://web.stanford.edu/group/uq/}) 或冯卡门研究所的讲座系列（VKI 讲座系列 STO-AVT-236，关于计算流体动力学中的不确定性量化）。关于 CFD 不确定性量化的出版物数量也在快速增加；最近的例子是 Bijl 等人编辑的一本书。 ( 2013 ) 和 Rakhimov 等人的论文。 （2018）等。

\section*{12.3 网格质量和优化}\phantomsection\addcontentsline{toc}{section}{12.3 网格质量和优化}
当网格被细化时，离散化误差总是会减少，但是对这些误差的可靠估计需要对每个新应用进行网格细化研究。具有给定数量网格点的网格的优化可以比非优化网格的系统细化减少同样多（或更多）的离散化误差。因此，关注电网质量非常重要。

网格优化旨在提高表面和体积积分的近似精度。这取决于所使用的离散化方法；在本节中，我们讨论影响本书中描述的方法准确性的网格特征。

为了使用线性插值和/或中点规则获得最高精度的对流通量，连接两个相邻 CV 中心的线应该

\begin{sourceblock}
\includegraphics[page=481,trim={53.00bp 486.24bp 53.87bp 52.00bp},clip,width=\textwidth,max height=0.38\textheight]{Ferziger4th.pdf}
\par\vspace{0.45\baselineskip}
\small 图 12.8 由于 k 和 k ′ 之间的距离过大而导致网格质量较差的示例（左）以及通过局部网格细化改进的网格质量（右）
\end{sourceblock}
\begin{sourceblock}
\includegraphics[page=481,trim={53.00bp 332.87bp 53.87bp 206.40bp},clip,width=\textwidth,max height=0.38\textheight]{Ferziger4th.pdf}
\par\vspace{0.45\baselineskip}
\small 图 12.9 当棱镜层环绕尖角（左）和非共形界面（右）时的高网格非正交性示例
\end{sourceblock}
穿过公共面的中心。在某些情况下，特别是当使用块结构网格时，如图12.8所示的情况是不可避免的。大多数自动网格生成器都会在突出的角处创建这种网格，因为它们通常会在边界处创建六面体或棱柱层，如图 12.9 所示。如果平行于界面的单元很薄，则在非共形界面（例如，在块结构网格或滑动界面）处也会产生高非正交性；图12.9也描述了这种情况。为了在不进行自适应的情况下提高精度，应该对网格进行局部细化，如图 12.8 所示。这减小了单元面中心 k 与连接节点 C 和 N k 的直线穿过单元面 k ' 的点之间的距离。这两点之间的距离相对于单元面的大小（例如 √ S k ）是网格质量的度量。该距离太大的单元格应进行细化，直到 k ' 和 k 之间的距离减小到可接受的水平。在非共形块界面处，界面两侧的单元应具有相似的尺寸，并且长宽比不应太大，以便将非正交性限制在可接受的水平。

当连接相邻 CV 中心的线与单元面正交并穿过单元面中心时，可以获得扩散通量的最大精度。正交性提高了单元面法线方向导数的中心差分近似的准确性：

\begin{sourceblock}
\begin{tikzpicture}
\node[inner sep=0,anchor=south west] (image) at (0,0) {\includegraphics[page=482,trim={53.00bp 558.38bp 53.87bp 70.20bp},clip,width=0.92\textwidth,max height=0.38\textheight]{Ferziger4th.pdf}};
\begin{scope}[x={(image.south east)},y={(image.north west)}]
\fill[white] (0.35874,0.47364) rectangle (0.39696,0.78168);
\end{scope}
\end{tikzpicture}
\end{sourceblock}
当连接两个单元中心的线与面正交时，该近似在两个单元中心之间的中点处是二阶精确的；即使 k ' 不在节点之间的中间，也可以使用多项式拟合获得高阶近似值。如果非正交性不可忽略，则法态导数的估计需要使用许多节点。这可能会导致收敛问题。

如果 k ' 不是单元面中心，则 k ' 处的值表示单元面上平均值的假设不再是二阶准确的。尽管可以进行修正或替代近似，但大多数通用 CFD 代码都使用简单的近似，例如 (12.3)，并且如果网格属性不利，精度会大大降低。我们已经在该节中展示过。 9.7.1和9.7.2如何恢复二阶；对于任意网格也可以获得高阶近似，但这增加了代码的复杂性并降低了其鲁棒性。此外，在属性较差的相对粗糙的网格上，高阶方法可能会导致比低阶方法更不准确的近似。

在大多数有限体积方法中，网格线在 CV 角处正交并不重要；仅连接相邻 CV 中心的线与单元面法线之间的角度很重要（参见图 12.10 中的角度 θ）。从这个意义上说，四面体网格可以是正交的。远离 0° 的角度 θ 会导致较大的误差和收敛问题，应避免。在图12.8所示的情况下，连接相邻CV中心的线几乎与单元面正交，因此k'处的梯度可以准确计算，但由于k'和k之间的距离较大，表面上积分的通量的精度很差。

可能会遇到其他类型的不良 CV 失真。两个如图 12.11 所示。在一种情况下，正六面体 CV 的上表面绕其法线旋转，从而使相邻面变形。在另一种情况下，顶面在其自己的平面中被剪切。这两个功能都是不受欢迎的，如果可能的话应该避免。当弯曲壁上存在薄棱柱形电池时，翘曲尤其成问题。细胞质心的位置可能会落在细胞之外，这可能会导致严重的问题。解是增加棱镜层的厚度，或细化壁切向方向的网格，或两者兼而有之。

\begin{sourceblock}
\begin{tikzpicture}
\node[inner sep=0,anchor=south west] (image) at (0,0) {\includegraphics[page=482,trim={53.00bp 109.44bp 53.87bp 506.46bp},clip,width=0.92\textwidth,max height=0.38\textheight]{Ferziger4th.pdf}};
\begin{scope}[x={(image.south east)},y={(image.north west)}]
\fill[white] (0.79922,0.46835) rectangle (0.81817,0.67954);
\end{scope}
\end{tikzpicture}
\end{sourceblock}
\begin{sourceblock}
\includegraphics[page=483,trim={53.00bp 496.25bp 53.87bp 52.00bp},clip,width=\textwidth,max height=0.38\textheight]{Ferziger4th.pdf}
\par\vspace{0.45\baselineskip}
\small 图 12.11 由于单元扭曲（中）和扭曲（右）而导致网格质量较差的示例
\end{sourceblock}
\begin{sourceblock}
\includegraphics[page=483,trim={53.00bp 400.83bp 53.87bp 183.44bp},clip,width=\textwidth,max height=0.38\textheight]{Ferziger4th.pdf}
\par\vspace{0.45\baselineskip}
\small 图 12.12 质量较差的三角形或四面体网格的示例（左）以及通过在边界附近添加棱柱层的补救措施（右）
\end{sourceblock}
求解欧拉方程时，经常使用由 2D 三角形和 3D 四面体组成的网格。当它们用于求解纳维-斯托克斯方程时，如果四面体质量较差，可能会遇到问题。在两个边界之间的拐角处会遇到一种有问题的情况；两个边界相交的边缘附近的某些四面体可能只有 2 个或什至可能只有一个邻居（其他单元面位于解域边界中）。那么仅使用来自直接邻居的数据就不可能准确地计算梯度（即三个坐标方向上的导数）。如果网格中存在此类单元并且使用通常的离散化，则这些单元处的变量很可能会振荡并且可能无法获得收敛解。通常的解是沿着边界生成棱镜层（特别是沿着墙壁，其中在墙壁法线方向上存在高梯度，如果墙壁是弯曲的，也可能在切向方向上）。这确保了邻近边界的单元至少有 4 个邻居，并且还解决了当四面体靠近壁太平坦时导致的高网格非正交性问题，如图 12.12 所示。特别是对于粘性和湍流，如果核心网格是三角形或四面体，则靠近壁的棱柱层是必须的。

平坦、扭曲的四面体可能遇到的另一个有问题的情况是当所有相邻单元质心几乎落入一个平面时，使得难以近似垂直于该平面的方向上的梯度分量。这通常也会导致可变振荡和收敛问题。解决方法是在生成四面体网格期间强制执行 Delaunay 准则（参见第 9.2.2 节），并在单元需要平坦的墙壁附近创建棱柱。

如果计算节点放置在 CV 质心处，则通过中点规则近似的体积积分是二阶精确的。然而，CV 有时可能会变形，以致质心实际上位于 CV 外部。应该避免这种情况。网格生成器应该检查它生成的网格，并向用户指示存在问题的单元格，除非它能够自动纠正它们。然后用户应该尝试修改控制参数以获得更好的网格质量。

其中一些问题可以通过细分有问题的单元格（以及可能的一些周围单元格）来避免。不幸的是，在某些情况下，唯一的解是生成新的网格。

\section*{12.4 流动计算的多重网格方法}\phantomsection\addcontentsline{toc}{section}{12.4 流动计算的多重网格方法}
几乎所有迭代求解方法在更精细的网格上收敛速度都较慢。收敛速度取决于方法；对于许多方法来说，获得收敛解的外迭代次数与一个坐标方向上的节点数成线性比例。这种行为与信息每次迭代传播的距离有限有关，并且为了收敛，信息必须在域中来回传播多次。多重网格方法，其所需的迭代次数与网格点的数量无关，在 20 世纪 90 年代受到了广泛的关注（参见 Brandt 1984；和 Briggs 等人 2000）。许多作者（包括现在的作者）已经证明，用多重网格方法求解纳维-斯托克斯方程是非常有效的。各种稳态层流和湍流的经验表明，实施多重网格思想可以大大减少计算工作量（参见 Wesseling 1990 的评论论文）。我们在这里简要总结了作者使用的方法的一个版本；许多其他变体也是可能的，请参阅专门讨论 CFD 多重网格方法的国际会议记录，例如 McCormick 1987；以及哈克布施和特罗滕伯格（1991）。在第 1 章中。在第 5 章中，我们提出了一种有效求解线性方程组的多重网格方法。我们在那里看到多重网格方法使用网格层次结构；在最简单的情况下，粗粒度是细粒度的子集。当分数步或其他显式时间步长方法应用于非定常流动时，它非常适合求解类泊松压力或压力校正方程，因为需要精确求解压力方程；这通常在复杂几何中的流的 LES 和 DNS 中完成。另一方面，当使用隐式方法时，不需要在每次迭代时非常精确地求解线性方程；将残差水平降低一个数量级就足够了，并且通常可以通过基本求解器（例如 ILU 或 CG）之一的几次迭代来实现。更准确的解不会减少外部迭代的次数，但可能会增加计算时间。因此，如果多重网格方法仅应用于隐式求解方法中线性方程组的求解，则可实现的加速度是有限的。

对于稳态流问题，我们已经看到隐式求解方法是首选，并且外部迭代的加速非常重要。幸运的是，多重网格方法也可以应用于外部迭代。根据多重网格术语，构成一次外部迭代的操作序列被视为“更平滑”。

在结构化网格上稳态流的有限体积方法的多重网格版本中，每个粗网格 CV 由 2D 中下一个更精细网格的 4 个 CV 和 3D 中的 8 个 CV 组成。通常首先生成最粗的网格，然后从求解其上的问题开始求解过程。然后将每个 CV 细分为更精细的 CV。在最粗糙的网格上找到收敛解后，将其插值到下一个更精细的网格以提供起始解。然后开始两网格程序。重复该过程，直到达到最精细的网格并获得其上的解。如前所述，该策略称为完整多重网格程序 (FMG)。在间隔为 h 的网格上进行 m 次外迭代后，短波长误差分量已被去除，中间解满足以下方程：

\begin{sourceblock}
\includegraphics[page=485,trim={53.00bp 401.42bp 53.87bp 246.53bp},clip,width=0.92\textwidth,max height=0.38\textheight]{Ferziger4th.pdf}
\end{sourceblock}
where ρ m
h 是第 m 次迭代后的残差向量。求解过程现在转移到下一个较粗的网格，其间距为 2 h。如前所述，迭代成本和收敛速度在粗网格上都更加有利，从而使该方法更加高效。

在粗网格上求解的方程应该是细网格方程的平滑版本。通过仔细选择定义，可以确保所求解的方程与之前在该网格上求解的方程相同，即系数矩阵相同。然而，方程现在包含一个额外的源项：

\begin{sourceblock}
\includegraphics[page=485,trim={53.00bp 270.21bp 53.87bp 377.35bp},clip,width=0.92\textwidth,max height=0.38\textheight]{Ferziger4th.pdf}
\end{sourceblock}
如果设置为零，则等式的左侧。 ( 12.5 ) 将表示粗网格方程。右侧包含确保解是平滑的细网格解而不是粗网格解本身的校正。附加项是通过对细网格解和残差进行平滑（“限制”）而获得的；它们在粗网格上的迭代过程中保持不变。上式左边所有项的初始值是右边相应项的初始值。如果细网格残差为零，则解为 ˆ φ 2 h = ̃ φ 2 h。

只有当细网格上的残差不为零时，粗网格近似值才会从其初始值发生变化（因为问题是非线性的，系数矩阵和源项也会发生变化，这就是为什么这些项带有 \textasciicircum{} 符号）。一旦获得粗网格上的解（在一定的公差范围内），修正

\begin{sourceblock}
\includegraphics[page=485,trim={53.00bp 103.72bp 53.87bp 542.84bp},clip,width=0.92\textwidth,max height=0.38\textheight]{Ferziger4th.pdf}
\end{sourceblock}
通过插值（“延长”）转移到细网格并添加到现有解 φ m
h.通过这种校正，解中的大部分低频误差

\begin{sourceblock}
\begin{tikzpicture}
\node[inner sep=0,anchor=south west] (image) at (0,0) {\includegraphics[page=486,trim={53.00bp 596.16bp 53.87bp 56.98bp},clip,width=0.92\textwidth,max height=0.38\textheight]{Ferziger4th.pdf}};
\begin{scope}[x={(image.south east)},y={(image.north west)}]
\fill[white] (0.60574,0.09231) rectangle (0.69161,0.85846);
\fill[white] (0.69377,0.09231) rectangle (0.74051,0.85846);
\fill[white] (0.79338,0.14154) rectangle (0.84710,0.90769);
\fill[white] (0.84926,0.14154) rectangle (0.89600,0.90769);
\end{scope}
\end{tikzpicture}
\par\vspace{0.45\baselineskip}
\small 图12.13 变量从细网格到粗网格的传递，反之亦然
\end{sourceblock}
\begin{sourceblock}
\begin{tikzpicture}
\node[inner sep=0,anchor=south west] (image) at (0,0) {\includegraphics[page=486,trim={53.00bp 549.12bp 53.87bp 100.92bp},clip,width=0.92\textwidth,max height=0.38\textheight]{Ferziger4th.pdf}};
\begin{scope}[x={(image.south east)},y={(image.north west)}]
\fill[white] (0.67537,0.15466) rectangle (0.68890,0.61366);
\fill[white] (0.84851,0.07453) rectangle (0.86202,0.53354);
\end{scope}
\end{tikzpicture}
\end{sourceblock}
\begin{sourceblock}
\begin{tikzpicture}
\node[inner sep=0,anchor=south west] (image) at (0,0) {\includegraphics[page=486,trim={53.00bp 525.53bp 53.87bp 127.61bp},clip,width=0.92\textwidth,max height=0.38\textheight]{Ferziger4th.pdf}};
\begin{scope}[x={(image.south east)},y={(image.north west)}]
\fill[white] (0.72953,0.09231) rectangle (0.75242,0.90769);
\end{scope}
\end{tikzpicture}
\end{sourceblock}
\begin{sourceblock}
\begin{tikzpicture}
\node[inner sep=0,anchor=south west] (image) at (0,0) {\includegraphics[page=486,trim={53.00bp 494.45bp 53.87bp 151.09bp},clip,width=0.92\textwidth,max height=0.38\textheight]{Ferziger4th.pdf}};
\begin{scope}[x={(image.south east)},y={(image.north west)}]
\fill[white] (0.59317,0.42718) rectangle (0.61305,0.94175);
\fill[white] (0.76244,0.05825) rectangle (0.78562,0.66117);
\end{scope}
\end{tikzpicture}
\end{sourceblock}
精细网格上的内容被删除，从而节省了精细网格上的大量迭代。继续这个过程直到细网格上的解收敛。然后可以使用理查森外推来获得对下一个更精细网格的改进的起始猜测，并且启动三级V循环，等等。

对于结构化网格，通常使用简单的双线性（2D）或三线性（3D）插值将变量值从精细网格传输到粗网格以及从粗网格到精细网格的校正。尽管可以并且已经使用了更复杂的插值技术，但在大多数情况下，这种简单的技术就足够了。

将变量从一个网格转移到另一个网格的另一种方法是计算该变量在计算该变量的网格的 CV 中心处的梯度（粗略或精细）。第 1 章描述了使用高斯定理计算任意 CV 中心梯度的有效方法。 9.然后很容易使用这个梯度计算附近任何地方的变量值（这对应于线性插值）。对于图12.13所示的情况，我们可以通过对使用细网格CV梯度计算的值进行平均来计算节点C处的粗网格变量值：

\begin{sourceblock}
\begin{tikzpicture}
\node[inner sep=0,anchor=south west] (image) at (0,0) {\includegraphics[page=486,trim={53.00bp 223.97bp 53.87bp 406.54bp},clip,width=0.92\textwidth,max height=0.38\textheight]{Ferziger4th.pdf}};
\begin{scope}[x={(image.south east)},y={(image.north west)}]
\fill[white] (0.00000,0.95481) rectangle (0.12229,1.00000);
\end{scope}
\end{tikzpicture}
\end{sourceblock}
其中 N f 是一个粗网格 CV 中细网格 CV 的数量（在结构化网格上，2D 中有 4 个，3D 中有 8 个）。粗略 CV 不需要知道哪些精细 CV 属于它——它只需要知道有多少个。另一方面，每个细网格 CV（子级）都知道它属于哪个粗网格 CV（父级）；它只有一个父母。

类似地，粗网格校正很容易转移到细网格。计算粗网格CV中心处的修正梯度；位于该 CV 内的细网格节点处的校正计算公式如下：

\begin{sourceblock}
\includegraphics[page=486,trim={53.00bp 91.47bp 53.87bp 555.26bp},clip,width=0.92\textwidth,max height=0.38\textheight]{Ferziger4th.pdf}
\end{sourceblock}
这种插值比简单地将粗略 CV 校正注入其中的所有精细 CV 更准确（这也可以完成，但在延长后需要平滑校正）。在结构化网格上，可以轻松实现其他类型的多项式插值。

在FV方法中，守恒性质可用于将质量通量和残差从细网格转移到粗网格。在二维中，粗网格 CV 由四个细网格 CV 组成，粗网格方程应为其子细网格 CV 方程之和。因此，残差在精细网格 CV 上简单求和，粗 CV 面上的初始质量通量是精细 CV 面上质量通量的总和。在粗网格计算过程中，质量通量不是使用受限速度计算的，而是使用速度校正进行校正的（前者的准确性较低，并且校正比变量本身更平滑）。对于泛型变量，我们可以这样写：

\begin{sourceblock}
\includegraphics[page=487,trim={53.00bp 425.33bp 53.87bp 222.18bp},clip,width=0.92\textwidth,max height=0.38\textheight]{Ferziger4th.pdf}
\end{sourceblock}
where I k − 1
k 是描述从细网格到粗网格转移的算子，I k
k − 1 是描述从粗网格到细网格的转移的算子；等式。 （12.7）和（12.8）提供了这些运算符的示例。

动量方程中压力项的处理值得特别提及。因为最初 ˆ p = ̃ p 并且压力项是线性的，所以我们可以使用差值 p ′ = ˆ p − ̃ p。那么我们就不需要限制从细网格到粗网格的压力。注意，这不是SIMPLE及相关算法的压力修正p′；它是对更精细网格压力的修正，并且基于速度修正 u ′
i = ˆ u i − ̃ u i。如前所述，初始粗网格质量通量 ̃ ̇ m 是通过对相应的细网格质量通量求和获得的。仅当速度改变时这些才会改变；我们假设精细网格质量通量在多重网格循环开始时是质量守恒的；如果不是，则质量不平衡可以包含在粗网格上的压力校正方程中。

重要的是要注意边界条件的实现也满足一致性要求。例如，如果通过将边界值设置为等于近边界节点处的值来实现对称边界条件，则限制算子无法通过对细网格边界值进行插值来计算 ̃ 的边界值；它必须计算所有内部节点处的 ̃ φ，然后对其应用边界条件，即将对称边界处的 ̃ φ 设置为等于近边界节点处的 ̃ φ。如果将边界条件应用于 ˆ φ，则 φ ′ 在边界节点和邻近边界节点处将不相同，且梯度为 φ ′
将被传递到更精细的网格。那么细网格上的解超过一定限度就无法收敛。由于处理其他边界条件的不一致，可能会出现类似的情况，但我们不会在这里列出所有可能性。重要的是要确保迭代误差可以降低到机器精度（即使在代码投入生产时不会使用此标准）；如果这不可能，那就有问题了！

\begin{sourceblock}
\includegraphics[page=488,trim={53.00bp 490.25bp 53.87bp 52.00bp},clip,width=\textwidth,max height=0.38\textheight]{Ferziger4th.pdf}
\par\vspace{0.45\baselineskip}
\small 图 12.14 使用 V 循环的 FMG 方案示意图，显示了一个循环中不同阶段的典型外迭代次数
\end{sourceblock}
其他策略（例如 W 循环）可用于在网格之间循环。根据收敛速度来决定从一个网格切换到另一个网格可以提高效率。最简单的选择是上述 V 循环，在每个网格级别上具有固定的迭代次数。图 12.14 示意性地显示了 V 循环的 FMG 方法的行为以及每个级别的典型迭代次数。参数的最佳选择取决于问题，但它们对性能的影响不如单网格方法那么显著。多重网格方法的详细信息可以在 Hackbusch (2003) 的书中找到。

多重网格方法可以应用于非结构化网格以及结构化网格。在FV方法中，通常将细网格CV连接起来以产生粗网格CV；每个粗 CV 的精细 CV 数量可能不同，具体取决于 CV 的形状（四面体、金字塔、棱柱、六面体等）。如果粗网格和细网格没有通过系统细化或粗化相关联，则甚至可以使用多重网格思想——网格可以是任意的；唯一重要的是，解域和边界条件在所有级别上都相同，并且一个网格比另一个网格显著粗糙（否则计算效率不会提高）。然后，限制和延长算子基于通用插值方法；这种多重网格方法称为代数多重网格方法（参见例如 Raw 1995；和 Weiss 等人 1999）。

对于使用隐式方法和小时间步长计算非定常流，外部迭代通常收敛得非常快（每次外部迭代残差减少一个数量级），因此不需要多重网格加速。对于完全椭圆（扩散主导）问题，节省量最大，而对于对流主导问题（欧拉方程），节省量最小。当使用五个网格级别时，稳态问题的典型加速因子范围为 10 到 100。下面给出一个例子。

当使用 k - ε 湍流模型计算湍流时，插值可能会在多重网格方法的早期循环中产生 k 和/或 ε 的负值；然后必须限制修正以保持积极性。在具有可变属性的问题中，属性在解域内可能会变化几个数量级。方程的这种强非线性耦合可能会导致多重网格

\begin{sourceblock}
\begin{tikzpicture}
\node[inner sep=0,anchor=south west] (image) at (0,0) {\includegraphics[page=489,trim={53.00bp 496.41bp 53.87bp 58.52bp},clip,width=0.92\textwidth,max height=0.38\textheight]{Ferziger4th.pdf}};
\begin{scope}[x={(image.south east)},y={(image.north west)}]
\fill[white] (0.61014,0.90307) rectangle (0.65095,0.98921);
\fill[white] (0.00000,0.81818) rectangle (0.09895,0.90873);
\fill[white] (0.10030,0.81818) rectangle (0.13059,0.90873);
\fill[white] (0.13194,0.81818) rectangle (0.16788,0.90873);
\fill[white] (0.16926,0.81818) rectangle (0.22926,0.90873);
\fill[white] (0.23068,0.81818) rectangle (0.27651,0.90873);
\fill[white] (0.27789,0.81818) rectangle (0.30250,0.90873);
\fill[white] (0.00000,0.72862) rectangle (0.01546,0.81917);
\fill[white] (0.01699,0.72862) rectangle (0.11663,0.81917);
\fill[white] (0.11817,0.72862) rectangle (0.19940,0.81917);
\fill[white] (0.20096,0.72862) rectangle (0.22698,0.81917);
\fill[white] (0.22851,0.72862) rectangle (0.24463,0.81917);
\fill[white] (0.00000,0.63906) rectangle (0.08899,0.72961);
\fill[white] (0.09056,0.63906) rectangle (0.11657,0.72961);
\fill[white] (0.11814,0.63906) rectangle (0.15408,0.72961);
\fill[white] (0.00000,0.54950) rectangle (0.17191,0.64005);
\fill[white] (0.17350,0.54950) rectangle (0.23744,0.64005);
\fill[white] (0.23901,0.54303) rectangle (0.26803,0.64284);
\fill[white] (0.61886,0.52127) rectangle (0.64767,0.60741);
\fill[white] (0.00000,0.45994) rectangle (0.03383,0.55049);
\fill[white] (0.03540,0.45994) rectangle (0.07134,0.55049);
\fill[white] (0.07290,0.45994) rectangle (0.17711,0.55049);
\fill[white] (0.17871,0.45994) rectangle (0.24523,0.55049);
\fill[white] (0.24677,0.45994) rectangle (0.29621,0.55049);
\fill[white] (0.00000,0.37029) rectangle (0.02253,0.46084);
\fill[white] (0.02406,0.37029) rectangle (0.05717,0.46084);
\fill[white] (0.05874,0.37317) rectangle (0.08343,0.46363);
\fill[white] (0.08568,0.37029) rectangle (0.14138,0.46084);
\fill[white] (0.61886,0.32551) rectangle (0.64767,0.41165);
\fill[white] (0.93380,0.01079) rectangle (0.96884,0.11438);
\end{scope}
\end{tikzpicture}
\end{sourceblock}
方法变得不稳定。最好仅在最精细的网格上更新某些量（例如，k - ε 湍流模型中的湍流粘度），并在多重网格循环内保持它们恒定。

在层流多重网格方法中，欠松弛因素相对不重要；与单网格方法相比，这些方法对这些参数的敏感度较低。在图 12.15 中，我们展示了两个网格在 Re = 10,000 时解决盖驱动腔问题所需的外部迭代次数对速度欠松弛因子的依赖性（使用压力修正的最佳欠松弛，参见第 8.4 节）。在 α u 在 0.5 和 0.9 之间的范围内，迭代次数变化约 30\%，而对于单网格方法，变化为 5 到 7 倍（网格越细，迭代次数越高；见图 8.14）。然而，对于湍流、传热等，欠松弛可能会对多重网格方法产生很大影响。

上述完整的多重网格方法提供了所有网格上的解。所有粗网格求解的成本大约是最细网格求解成本的 30\%。如果求解过程在零场的最精细网格上开始，则需要比使用 FMG 方法更多的总体工作量。更准确的初始字段带来的节省通常超​​过获得它们所需的成本。此外，在一系列网格上获得解可以评估离散化误差，如第 1 节中所讨论的。 3.9，并为网格细化提供了基础。当达到期望的精度时，可以停止网格细化过程。另外，可以使用理查森外推法。

上述加速外部迭代的多重网格方法可以应用于纳维-斯托克斯方程的任何求解方法。 Vanka(1986)将其应用到点耦合求解方法中； Hutchinson 和 Raithby (1986) 以及 Hutchinson 等人。 ( 1988 ) 将其与线耦合解技术一起使用。 SIMPLE 类型的方法和分数步方法也非常适合多重网格加速；参见，例如，Hortmann 等人。 （1990），Lilek 等人。 ( 1997a ) 以及 Thompson 和 Ferziger ( 1989 )。平滑器的角色现在由基本算法（例如 SIMPLE）承担；线性方程求解器起着次要作用。

在表 12.1 中，我们比较了使用不同的求解策略在雷诺数 Re = 1000 和 Re = 1000 时求解二维盖驱动腔流问题所需的外部迭代次数。 SG 表示初始场为零的单网格法。 PG表示延长方案，其中下一个解的解

\begin{center}\small 表 12.1 求解盖驱动二维空腔流问题时，不同版本的求解方法将 L 1 残差范数降低 4 个数量级所需的迭代次数和计算时间（α u = 0 . 8，α p = 0 . 2，非均匀网格；对于 PG 和 FMG，CPU 时间包括所有较粗网格上的计算时间）\end{center}
Re Grid No. 外层 Iter. CPU 时间 SG PG MG FMG SG PG MG FMG
100 8 2 58 58 58 58 0.3 0.3 0.3 0.3 16 2 61 51 47 45 0.9 1.2 1.4 1.5 32 2 156 99 41 41 9.1 7.0 4.0 5.0 64 2 555 256 40 40 140.8 71.1 13.0 16.9 128 2 2119 620 40 40 2141.9 702.6 50.9 66.5 256 2 – – 40 40 – – 242.2 293.8 1000 8 2 124 124 124 124 0.5 0.5 0.5 0.5 16 2 156 162 123 132 2.2 2.5 2.8 2.9 32 2 250 288 132 132 14.0 19.2 11.2 13.8 64 2 433 400 93 73 97.0 120.7 32.0 38.5 128 2 1352 725 83 41 1383.4 851.1 121.5 92.4 256 2 – – 83 31 – – 512.9 278.8
较粗的网格用于提供初始场。 MG 表示使用 V 循环的多重网格方法，其中最精细的网格具有零初始场。最后，FMG表示上述多重网格方法，它可以被认为是PG和MG方案的组合。

结果表明，对于 Re = 1000 的情况，MG 和 FMG 在最精细的网格级别上需要大约相同数量的外部迭代 - 高质量的初始猜测并不能节省太多。对于单网格方案，节省的成本是巨大的：在 128 × 128 CV 网格上，迭代次数减少了 3.5 倍。多重网格方法将该网格的迭代次数减少了 15 倍；随着网格的细化，这个因素会增加。从第三个网格开始，MG 和 FMG 方法中的迭代次数保持不变，而 SG 中的迭代次数增加了 4 倍，PG 方案中的迭代次数增加了 2.5 倍。

对于高雷诺数流动，情况会发生一些变化。 SG 需要的迭代次数少于 Re = 1000 时的迭代次数，但粗网格除外，在粗网格上使用 CDS 会减慢收敛速度。 PG 将迭代次数减少了不到两倍。 MG 需要的迭代次数大约是 Re = 100 时的两倍。然而，随着网格的细化，FMG 变得更加高效 — 256 × 256 CV 网格上所需的迭代次数实际上比 Re = 1000 时要少。这是应用于 Navier-Stokes 方程的多重网格方法的典型行为。 FMG 方法通常是最有效的一种。

3D 腔体流动也获得了与表 12.1 类似的结果；参见 Lilek 等人。 （1997a）了解详细信息。

在图 12.16 中，显示了用于计算的湍流动能 k 的剩余范数（所有 CV 的剩余绝对值之和）的减少
0.1

\begin{center}\small 图12.16 管束段流动计算中湍流动能k剩余范数的折减；最好的网格有 176 × 48 CV，使用了五个级别（来自 Lilek 1995）\end{center}
SG

0.01

Residual

1e-3
1e-4
FMG

1e-5
0 200 400 600 1e-6
迭代

使用 k − ε 模型制作的一段管束中的湍流。这些曲线是 MG 和 SG 方法的典型曲线。在实际应用中，将残差减少三到四个数量级通常就足够了。这里残差的减少超过了必要的程度，以表明收敛速度没有恶化。计算时间的节省因应用程序而异：对流主导的流动比扩散主导的流动要少。

\section*{12.5 自适应网格细化（AMR）}\phantomsection\addcontentsline{toc}{section}{12.5 自适应网格细化（AMR）}
\subsection*{12.5.1 自适应网格细化的动机}\phantomsection\addcontentsline{toc}{subsection}{12.5.1 自适应网格细化的动机}
准确性问题从一开始就一直困扰着计算流体动力学。许多已发表的结果存在重大错误。旨在确定模型有效性的测试有时被证明是不确定的，因为数值误差大于模型的影响。不同小组使用不同代码解决同一问题的比较研究表明，使用相同模型的不同代码获得的解之间的差异往往大于使用相同代码和不同模型获得的解之间的差异（Bradshaw et al. 1994；Rodi et al. 1995）。如果模型确实相同，这些差异只能是由于数值错误或用户错误造成的（由于不同的解释、实现或边界处理，假定相同的模型结果不同，这并不罕见）。为了有效地比较模型，估计和减少误差至关重要。

第一个要素是估计误差的方法；前面给出的理查森方法是一个不错的选择。无需在两个网格上执行计算即可估计误差的方法是比较通过 CV 面的通量，这些通量是由解中采用的离散化方法和更准确（高阶）方法产生的。这不如理查森方法准确，但它确实起到了指示误差较大的位置的作用。由于高阶近似通常更复杂，因此仅用于在使用基本方法获得收敛解后计算通量。例如，可以通过特定面两侧的单元中心拟合三阶多项式，并使用这两个单元中心的变量值和梯度来查找多项式的系数（示例请参见第 4.4.4 节）。假设，如果使用这种近似，将获得精确的解，而不是由通常的近似产生的解 φ。使用三次计算的通量之间的差异 ( F
k ) 和线性拟合 ( F φ
k ) 应作为附加源项添加到离散方程中，以恢复“精确”解。如果我们将离散化误差 ε d 和源项 τ （通常称为 tau 误差）定义如下：

\begin{sourceblock}
\begin{tikzpicture}
\node[inner sep=0,anchor=south west] (image) at (0,0) {\includegraphics[page=492,trim={53.00bp 380.63bp 53.87bp 255.92bp},clip,width=0.92\textwidth,max height=0.38\textheight]{Ferziger4th.pdf}};
\begin{scope}[x={(image.south east)},y={(image.north west)}]
\fill[white] (0.00000,0.82122) rectangle (0.10159,1.00000);
\fill[white] (0.24695,0.42987) rectangle (0.27699,0.86989);
\fill[white] (0.28244,0.42987) rectangle (0.31062,0.82055);
\fill[white] (0.34457,0.42987) rectangle (0.39669,0.82055);
\fill[white] (0.41672,0.41737) rectangle (0.46478,0.80804);
\fill[white] (0.48394,0.39067) rectangle (0.51335,0.82055);
\fill[white] (0.51874,0.42987) rectangle (0.54692,0.82055);
\end{scope}
\end{tikzpicture}
\end{sourceblock}
我们得到离散化误差估计和 tau 误差之间的以下联系：

\begin{sourceblock}
\begin{tikzpicture}
\node[inner sep=0,anchor=south west] (image) at (0,0) {\includegraphics[page=492,trim={53.00bp 324.04bp 53.87bp 315.16bp},clip,width=0.92\textwidth,max height=0.38\textheight]{Ferziger4th.pdf}};
\begin{scope}[x={(image.south east)},y={(image.north west)}]
\fill[white] (0.00000,0.90200) rectangle (0.11729,1.00000);
\fill[white] (0.35564,0.42910) rectangle (0.41808,0.95546);
\end{scope}
\end{tikzpicture}
\end{sourceblock}
无需求解该方程组来得到 ε d，通常只需通过 A P 标准化 τ P 并使用该量作为离散化误差的估计就足够了；这对应于在方程组上执行一次雅可比迭代。 ( 12.11 ) 从零初始值开始。原因是上述分析只是近似的，计算量只是一个指示，而不是离散化误差的估计。有关此误差估计方法的更多详细信息和应用示例，请参阅 Muzaferija 和 Gosman (1997)。关键点是我们现在有了一种在整个流域中逐点定义误差的策略。该信息可用于调整网格（即网格）以使误差水平更加均匀。本质上，如果特定网格点的误差估计大于规定水平，则该网格单元将被标记为细化。

\subsection*{12.5.2 AMR 策略}\phantomsection\addcontentsline{toc}{subsection}{12.5.2 AMR 策略}
根据误差估计，可以构建标记为细化的网格单元的地图。为了提供缓冲区，细化区域的边界通常会延伸一定的余量，该余量应该是局部网格尺寸的函数；两到四个单元的宽度通常是明智的选择。

块结构网格要求按块进行细化；如果不是所有块都被细化，则需要不匹配的接口能力。对于非结构化网格，局部细化可以是单元级的。否则，可以对要细化的单元进行聚类并定义新的细化网格块，如稍后将描述的。

目标是使各处的误差小于某个公差 δ，无论是绝对误差 ∥ ε ∥，还是相对误差 ∥ ε/φ ref ∥，其中 φ ref 是用于归一化的代表变量值。这可以通过使用不同精度的方法来实现，这是常微分方程求解器中常用的方法，但很少这样做。我们还可以细化各处的网格，但这很浪费。更灵活的选择是在误差较大的地方进行局部细化网格。 CFD 代码的经验丰富的用户可能会生成必要时精细而其他地方粗糙的网格，从而产生几乎均匀的离散误差分布。当几何形状包含小但重要的突出物时，例如汽车上的镜子、轮船和其他船只上的附件、大室壁上的小入口和出口等，这一点尤其重要。然而，在复杂的几何形状中，很难“手工”进行局部细化；在瞬态流中，需要细化的区域也可能随时间而变化。因此，如果要以最小的努力实现所需的精度，则自动自适应网格细化方法至关重要。例如，C++ 代码的 Combo 包（Adams et al. 2015）支持块结构的 AMR 应用程序。 AMR 在许多论文中都有讨论，首先是 Berger 和 Oliger (1984)；例子包括 Skamarock 等人。 （1989）以及汤普森和汤普森和弗齐格（1989）。

在某些流问题中，网格的哪些地方需要细化是显而易见的，因此自适应细化的自动化很容易。例子有冲击的可压缩流（冲击周围需要精细网格，例如翼型周围和涡轮机等的流动）、自由表面流动（需要精细网格来解析自由表面，例如射流破碎、上升气泡、刚体进水、波浪中的船舶等）、空化流动等。 12.17 和 12.18；下一章将介绍适应自由表面和空化区的示例。合适的误差指示器或估计器将显示初始网格不够精细的其他区域。有时，误差估计可能表明某个区域中某个变量的误差较高，但其他变量可能需要在其他地方进行网格细化；人们可能需要对每个变量和误差水平进行折衷或加权。这表明决定网格的何处细化以及细化到哪个级别的问题绝非微不足道。正是由于这个原因，商业 CFD 代码尚未提供自适应局部网格细化作为标准功能，但原型已经呈现，并且很明显未来的版本将包含它。

一些作者仅对网格的细化部分进行计算，使用从细化界面处的粗网格解获取的边界条件。这被称为被动方法，因为网格未细化部分的解不会重新计算（参见 Berger 和 Oliger 1984）。这一特性使得该方法不适用于椭圆问题，在椭圆问题中，任何区域的条件变化都可能影响各处的解。允许细化网格解的影响遍及整个域的方法称为主动方法。 Caruso 等人开发了此类方法。 ( 1985 ) 和 Muzaferija ( 1994 ) 等。

\begin{sourceblock}
\includegraphics[page=494,trim={53.00bp 256.25bp 53.87bp 52.00bp},clip,width=\textwidth,max height=0.38\textheight]{Ferziger4th.pdf}
\par\vspace{0.45\baselineskip}
\small 图 12.17 Ma = 0.8 时机翼周围湍流的初始网格（上）和计算的马赫数轮廓线（下）
\end{sourceblock}
一种主动方法（例如，Caruso 等人，1985）与被动方法完全相同，但重要的区别在于，当计算出精细网格解时，该过程尚未完成。相反，有必要计算一个新的粗网格解；该解不是在覆盖整个域的粗网格上计算的解，而是细网格解的平滑版本。要查看需要什么，假设：

\begin{sourceblock}
\begin{tikzpicture}
\node[inner sep=0,anchor=south west] (image) at (0,0) {\includegraphics[page=494,trim={53.00bp 127.07bp 53.87bp 524.03bp},clip,width=0.92\textwidth,max height=0.38\textheight]{Ferziger4th.pdf}};
\begin{scope}[x={(image.south east)},y={(image.north west)}]
\fill[white] (0.00000,0.92154) rectangle (0.06238,1.00000);
\fill[white] (0.06505,0.92154) rectangle (0.08986,1.00000);
\fill[white] (0.09254,0.92154) rectangle (0.18968,1.00000);
\fill[white] (0.19242,0.92154) rectangle (0.29374,1.00000);
\fill[white] (0.29645,0.92154) rectangle (0.35453,1.00000);
\fill[white] (0.41543,0.07979) rectangle (0.50689,0.92021);
\fill[white] (0.51041,0.15160) rectangle (0.53859,0.92021);
\fill[white] (0.54424,0.07979) rectangle (0.58292,0.90293);
\fill[white] (0.90926,0.12699) rectangle (1.00000,0.89561);
\end{scope}
\end{tikzpicture}
\end{sourceblock}
是问题在大小为 h 的网格上的离散化； L代表操作员。为了强制解成为已细化区域中细网格解的平滑版本，我们通过以下方式替换粗网格问题：

\begin{sourceblock}
\includegraphics[page=495,trim={53.00bp 244.25bp 53.87bp 52.00bp},clip,width=\textwidth,max height=0.38\textheight]{Ferziger4th.pdf}
\par\vspace{0.45\baselineskip}
\small 图 12.18 Ma = 0.8 时机翼周围湍流的冲击自适应网格（上）和计算的马赫数等值线（下）
\end{sourceblock}
\begin{sourceblock}
\begin{tikzpicture}
\node[inner sep=0,anchor=south west] (image) at (0,0) {\includegraphics[page=495,trim={53.00bp 175.53bp 53.87bp 461.54bp},clip,width=0.92\textwidth,max height=0.38\textheight]{Ferziger4th.pdf}};
\begin{scope}[x={(image.south east)},y={(image.north west)}]
\fill[white] (0.11426,0.25765) rectangle (0.22692,0.69522);
\fill[white] (0.23044,0.29756) rectangle (0.25862,0.69522);
\end{scope}
\end{tikzpicture}
\end{sourceblock}
其中 ̃ φ h 是平滑的细网格解（即其在粗网格上的表示）。然后将解在粗网格和细网格之间迭代，直到迭代误差足够小；通常大约四次迭代就足够了。由于每个网格上的解不需要每次都迭代到最终的容差，因此该方法的成本仅比被动方法多一点。

在另一种主动方法中（Muzaferija 1994；Muzaferija and Gosman 1997），网格被组合成单个全局网格，包括细化网格以及原始网格的非细化部分。这需要一种解决方法，允许

\begin{center}\small 图 12.19 细化界面处的非细化 CV：它具有与六个相邻 CV (N 1 , ..., N 6 ) 共有的六个面 (c 1 , ...,c 6 )\end{center}
\begin{sourceblock}
\begin{tikzpicture}
\node[inner sep=0,anchor=south west] (image) at (0,0) {\includegraphics[page=496,trim={53.00bp 559.63bp 53.87bp 93.89bp},clip,width=0.92\textwidth,max height=0.38\textheight]{Ferziger4th.pdf}};
\begin{scope}[x={(image.south east)},y={(image.north west)}]
\fill[white] (0.67838,0.09509) rectangle (0.70189,0.90491);
\fill[white] (0.69672,0.00713) rectangle (0.71020,0.58875);
\end{scope}
\end{tikzpicture}
\end{sourceblock}
\begin{sourceblock}
\begin{tikzpicture}
\node[inner sep=0,anchor=south west] (image) at (0,0) {\includegraphics[page=496,trim={53.00bp 485.27bp 53.87bp 120.46bp},clip,width=0.92\textwidth,max height=0.38\textheight]{Ferziger4th.pdf}};
\begin{scope}[x={(image.south east)},y={(image.north west)}]
\fill[white] (0.66111,0.80500) rectangle (0.67741,0.97418);
\fill[white] (0.67191,0.79258) rectangle (0.68535,0.91409);
\fill[white] (0.71618,0.76279) rectangle (0.72962,0.88446);
\fill[white] (0.79823,0.79854) rectangle (0.81170,0.92005);
\fill[white] (0.49155,0.54958) rectangle (0.51510,0.71859);
\fill[white] (0.62656,0.60305) rectangle (0.64577,0.77206);
\fill[white] (0.70860,0.45969) rectangle (0.72208,0.58136);
\fill[white] (0.76692,0.46631) rectangle (0.79044,0.63549);
\end{scope}
\end{tikzpicture}
\end{sourceblock}
具有任意数量面孔的简历。细化区域和非细化区域之间的交界处的 CV 比常规 CV 具有更多的面和邻居；见图12.19。为了保留 FV 方法的全局守恒性质，细化边界上的非细化 CV 的面必须被视为两个（2D 中和 3D 中四个）独立的子面，每个子面为两个 CV 所共用。在离散化中，子面的处理方式与两个 CV 之间的任何其他面完全相同。计算机代码需要有一个可以处理这种情况的数据结构，并且求解器需要能够处理由此产生的不规则矩阵结构。共轭梯度型求解器是一个不错的选择；带有 Gauss-Seidel 平滑器的多重网格求解器也可以使用，但有一些限制。可以通过将细胞面和细胞体积相关值存储在单独的数组中来优化数据结构。对于像第一章中描述的那样的简单离散化方案。如图 9 所示，这很容易完成：每个单元面都是两个 CV 所共有的，因此对于每个面，需要存储指向相邻 CV 节点、表面向量分量和矩阵系数的指针。对于局部细化网格和标准网格，用于解决流动问题的计算机代码是相同的；只有预处理器需要进行调整，使其能够处理局部细化网格的数据。这类似于非共形网格块接口的处理，请参见第 1 节。 9.6.1.

如果使用Chimera网格，则代码无需更改，但每次细化后需要重新定义插值系数和插值涉及的节点。

可以根据需要使用尽可能多的网格细化级别；通常至少需要三个级别，但已使用多达八个级别。自适应网格方法的优点在于，由于最精细的网格仅占据域的一小部分，因此网格点总数相对较少，因此计算成本和内存需求都大大降低。此外，它可以被设计成使得用户不必是专家网格设计者。特别是对于汽车、飞机和船舶等钝体周围的流动，在物体附近和尾流中需要非常精细的网格，但在其他地方可以使用粗网格，局部网格细化对于准确和高效的模拟至关重要。第 2 节中介绍了用户定义的逐单元局部网格细化的示例。 10.3.4.2.

最后，这些方法与多重网格方法结合得很好。嵌套网格可以看作多重网格中使用的网格；唯一重要的区别是，由于粗网格在不需要细化的情况下提供了足够的精度，因此最细的网格不会覆盖整个域。在非细化区域中，两个级别的待解方程保持相同。由于多重网格方法的最大成本是由于在最精细网格上进行迭代，因此节省的成本可能非常大，尤其是在 3D 中。欲了解更多详细信息，请参见 Thompson 和 Ferziger (1989) 以及 Muzaferija (1994)。

\section*{12.6 CFD中的并行计算}\phantomsection\addcontentsline{toc}{section}{12.6 CFD中的并行计算}
本书第一版于 1996 年出版时，工作站是单处理器计算机。当时就已经很清楚，计算能力的进一步提高将需要多个处理器，即并行计算机。如今（2020 年左右）所有计算机都有多个处理单元（称为核心）；工作站通常有大约 36 个核心和 120 GB 内存，并且数量有增加的趋势。多个工作站可以连接起来形成集群，集群可能有数千个核心和 TB 内存。与经典矢量超级计算机相比，此类集群的优势在于可扩展性。它们还使用标准芯片，因此生产成本更低。然而，为传统串行处理设计的算法可能无法在并行计算机上有效运行。

如果并行化是在循环级别执行的（就像自动并行化编译器的情况一样），阿姆达尔定律就会发挥作用，该定律本质上说速度是由代码中效率最低的部分决定的。为了实现高效率，不能并行化的代码部分必须非常小。

更好的方法是将求解域细分为子域并将每个子域分配给一个处理器。在这种情况下，相同的代码在所有处理器上运行在其自己的数据集上。由于每个处理器还需要驻留在其他子域中的一些数据，因此处理器之间的数据交换和/或存储重叠是必要的。

显式方案相对容易并行化，因为所有操作都是对先前时间步骤的数据执行的。每一步完成后，只需在相邻子域之间的接口区域交换数据即可。操作顺序和结果在一个和多个处理器上是相同的。问题中最困难的部分通常是求解压力的椭圆形泊松方程（不过，请参见 Sullivan 和 Patton 2008，他们使用 2-D x - y 平面分解来求解行星边界层的不可压缩流动问题）。

隐式方法更难以并行化。虽然系数矩阵和源向量的计算仅使用“旧”数据并且可以有效地并行执行，但线性方程组的求解并不容易并行化。例如，高斯消去法，其中每次计算都需要前一次计算的结果，在并行机上执行起来非常困难。其他一些求解器可以并行化，并在 n 个处理器上执行与单个处理器上相同的操作序列，但它们要么效率不高，要么通信开销非常大。我们将描述两个例子。

\subsection*{12.6.1 线性方程系统迭代求解器的并行化}\phantomsection\addcontentsline{toc}{subsection}{12.6.1 线性方程系统迭代求解器的并行化}
红黑高斯-赛德尔方法非常适合并行处理。它在该节中进行了简要描述。 5.3.8 和由以交替方式对两组点执行雅可比迭代组成。在 2D 中，节点的颜色与棋盘上的颜色相同；因此，对于二维中的五点计算分子，应用于红点的雅可比迭代仅使用来自黑色邻居节点的数据来计算新值，反之亦然。该求解器的收敛特性与 Gauss-Seidel 方法完全相同，该方法因此而得名。

任意一组节点上新值的计算可以并行执行；所需要的只是上一步的结果。结果与在单个处理器上完全相同。每次迭代，在更新每组数据之后，处理相邻子域的处理器之间的通信会发生两次。该本地通信可以与新值的计算重叠。该求解器仅适用于与多重网格方法结合使用，因为它本身效率相当低。

ILU 类型的方法（例如第 5.3.4 节中介绍的 SIP 方法）是递归的，使得并行化不太简单。在 SIP 算法中，L 和 U 矩阵的元素( 5.41 )，取决于 W 和 S 节点处的单元。在从相邻子域获得数据之前，无法开始计算除西南角子域之外的子域上的系数。在 2D 中，最佳策略是将域细分为垂直条纹，即使用 1D 处理器拓扑。 L 和 U 矩阵的计算和迭代可以相当有效地并行执行（参见 Bastian 和 Horton 1989）。子域1的处理器不需要其他处理器的数据，可以立即启动；它沿着底部或最南端的线前进。在计算出最右边节点的元素后，它可以将这些值传递给子域 2 的处理器。当第一个处理器在其下一行开始计算时，第二个处理器可以在其底行进行计算。当第一个处理器到达从底部开始的第 n 行时，所有 n 个处理器都忙。当第一个处理器到达顶部边界时，它必须等待直到后面n行的最后一个处理器完成；见图12.20。在迭代方案中，需要两次迭代。第一个是以刚才描述的方式完成的，而第二个基本上是它的镜像。

算法如下：

对于 j = 2 到 N j − 1 执行：

接收来自西邻的 U E ( i s − 1 , j ), U N ( i s − 1 , j )；对于 i = i s 到 i e 做：

\begin{sourceblock}
\includegraphics[page=499,trim={53.00bp 476.25bp 53.87bp 52.00bp},clip,width=\textwidth,max height=0.38\textheight]{Ferziger4th.pdf}
\par\vspace{0.45\baselineskip}
\small 图 12.20 SIP 求解器中前向循环（左）和后向循环（右）中的并行处理；阴影部分是负载不均匀的区域
\end{sourceblock}
计算UE(i,j)、LW(i,j)、UN(i,j)、LS(i,j)、LP(i,j);结束我;将UE ( i e , j ), UN ( i e , j ) 发送给东邻居；结束 j ;

对于 m = 1 到 M 执行：

对于 j = 2 到 N j − 1 执行：

从西邻接收 R ( i s − 1 , j )；对于 i = i s 到 i e 做：

计算 ρ( i , j ), R ( i , j ) ;结束我;发送 R ( i e , j ) 给东邻；结束 j ;

对于 j = N j − 1 到 2 步骤 -1 执行：

从东邻接收 δ( i e + 1 , j )；对于 i = i e 到 i s 步骤 -1 执行以下操作：

计算 δ( i , j )；更新变量；结束我;发送 δ( i s , j ) 到西邻；结束 j ;结束米。

问题在于，这种并行化技术需要大量（细粒度）通信，并且每次迭代的开始和结束时都有空闲时间；这些降低了效率。此外，该方法仅限于结构化网格。 Bastian 和 Horton (1989) 在基于晶片机的机器上获得了良好的效率，该机器具有良好的通信与计算速度之比。如果比率不太有利，该方法的效率就会降低。

共轭梯度法（无需预处理）可以直接并行化。该算法涉及一些全局通信（收集部分标量乘积并广播最终值），但性能几乎与单个处理器上的性能相同。然而，为了真正有效，共轭梯度法需要一个良好的预处理器。因为最好的预处理器是ILU型的（SIP是一种非常好的预处理器），所以上述问题再次出现。

上述发展表明，并行计算环境需要重新设计算法。在串行机器上表现出色的方法可能几乎不可能在并行机器上使用。此外，必须使用新标准来评估方法的有效性。隐式方法的良好并行化需要修改求解算法。数值运算数量方面的性能可能比串行计算机差，但如果处理器承载的负载均衡并且通信开销和计算时间适当匹配，则修改后的方法总体上可能更高效。

\subsection*{12.6.2 空间域分解}\phantomsection\addcontentsline{toc}{subsection}{12.6.2 空间域分解}
隐式方法的并行化通常基于数据并行或域分解，可以在空间和时间上进行。空间域分解时，将解域划分为一定数量的子域；这类似于网格的块结构。在块结构中，该过程由求解域的几何形状控制，而在域分解中，目标是通过为每个处理器提供相同的工作量来最大化效率。每个子域分配给一个处理器，但一个处理器可以处理多个网格块；如果是这样，我们可以将它们全部视为一个逻辑子域。

如前所述，必须修改并行机的迭代过程。通常的方法是将全局系数矩阵 A 分成对角块系统 A ii，其中包含连接属于第 i 个子域的节点的元素，以及非对角块或耦合矩阵 A i j ( i̸ = j )，表示块 i 和 j 的相互作用。例如，如果将一个方形二维解域分为四个子域，并对CV进行编号，使得每个子域的成员具有连续的索引，则矩阵具有如图12.21所示的结构；该插图中使用了五点分子离散化。下面描述的方法适用于使用较大计算分子的方案；在这种情况下，耦合矩阵更大。

为了提高效率，内部迭代的迭代求解器应具有尽可能小的数据依赖性（由邻居提供的数据）；数据依赖性可能导致较长的通信和/或空闲时间。因此，选择全局迭代矩阵以使块解耦，即，对于 i ̸ = j，M i j = 0。则子域 i 的迭代方案为：

\begin{sourceblock}
\includegraphics[page=501,trim={53.00bp 470.25bp 53.87bp 52.00bp},clip,width=\textwidth,max height=0.38\textheight]{Ferziger4th.pdf}
\par\vspace{0.45\baselineskip}
\small 图12.21 二维方形解域细分为4个子域时全局系数矩阵的结构
\end{sourceblock}
\begin{sourceblock}
\begin{tikzpicture}
\node[inner sep=0,anchor=south west] (image) at (0,0) {\includegraphics[page=501,trim={53.00bp 401.86bp 53.87bp 236.94bp},clip,width=0.92\textwidth,max height=0.38\textheight]{Ferziger4th.pdf}};
\begin{scope}[x={(image.south east)},y={(image.north west)}]
\fill[white] (0.08108,0.42538) rectangle (0.16235,0.94477);
\end{scope}
\end{tikzpicture}
\end{sourceblock}
因此，来自相邻子域的数据取自先前的内部迭代并被视为已知；它在每次内部迭代后更新。 SIP 求解器很容易适应这种方法。每个对角分块矩阵M ii 按常规方式分解为L和U矩阵；全局迭代矩阵 M = LU 不是在单处理器情况下找到的矩阵。对每个子域进行一次迭代后，必须交换未知的 φ m 的更新值，使得残差 ρ m
可以在子域边界附近的节点处计算。

当SIP求解器以这种方式并行化时，随着处理器数量变大，性能会下降；当处理器的数量从 1 增加到 100 时，迭代次数可能会加倍。但是，如果内部迭代不必非常紧密地收敛，就像第 1 章中描述的隐式方案的情况一样。如图7和8所示，并行SIP可以非常高效，因为SIP倾向于在前几次迭代中快速减少误差。特别是如果使用多重网格方法来加速外部迭代，总效率相当高（80-90\%；参见Schreck and Peri ́c 1993；和Lilek et al. 1995，例如）。以这种方式并行化的 2D 流预测代码可通过互联网获得；详情请参见附录。

基于共轭梯度的方法也可以使用上述方法并行化。下面我们展示了预处理 CG 求解器的伪代码。人们发现（Seidl 等人，1996），通过在单处理器或多处理器上每次 CG 迭代执行两次预处理器扫描可以实现最佳性能。使用诺伊曼边界条件求解泊松方程的结果模拟了 CFD 应用中的压力校正方程，如图 12.22 所示。每次 CG 迭代进行一次预处理器扫描时，收敛所需的迭代次数会随着处理器数量的增加而增加。然而，如果每次 CG 迭代进行两次或多次预处理器扫描，迭代次数几乎保持不变。然而，在不同的应用中，人们可能会获得不同的行为。

\begin{sourceblock}
\includegraphics[page=502,trim={53.00bp 422.25bp 53.87bp 52.00bp},clip,width=\textwidth,max height=0.38\textheight]{Ferziger4th.pdf}
\par\vspace{0.45\baselineskip}
\small 图 12.22 ICCG 求解器中的迭代次数与处理器数量的函数关系（具有 64 3 CV 的均匀网格、具有 Neumann 边界条件的泊松方程、每次预处理器扫描后的 LC、每次 CG 迭代的 l 次扫描、残差范数减少了两个数量级；来自 Seidl 1997）
\end{sourceblock}
\begin{itemize}
\item 通过设置进行初始化： k = 0, φ 0 = φ in , ρ 0 = Q − A φ in , p 0 = 0 , s 0 = 10 30
\end{itemize}
\begin{sourceblock}
\begin{tikzpicture}
\node[inner sep=0,anchor=south west] (image) at (0,0) {\includegraphics[page=502,trim={53.00bp 200.71bp 53.87bp 319.69bp},clip,width=0.92\textwidth,max height=0.38\textheight]{Ferziger4th.pdf}};
\begin{scope}[x={(image.south east)},y={(image.north west)}]
\fill[white] (0.00000,0.99451) rectangle (0.01913,1.00000);
\fill[white] (0.02779,0.99197) rectangle (0.13907,1.00000);
\fill[white] (0.14174,0.99197) rectangle (0.17651,1.00000);
\fill[white] (0.17922,0.99197) rectangle (0.27227,1.00000);
\fill[white] (0.27498,0.99266) rectangle (0.29311,1.00000);
\fill[white] (0.29841,0.99451) rectangle (0.32659,1.00000);
\fill[white] (0.33011,0.99197) rectangle (0.35756,1.00000);
\fill[white] (0.36024,0.99451) rectangle (0.39477,1.00000);
\fill[white] (0.39991,0.99451) rectangle (0.42809,1.00000);
\fill[white] (0.43161,0.98072) rectangle (0.48129,1.00000);
\fill[white] (0.48394,0.99451) rectangle (0.51738,1.00000);
\fill[white] (0.52253,0.99451) rectangle (0.55071,1.00000);
\fill[white] (0.55423,0.99266) rectangle (0.58235,1.00000);
\fill[white] (0.58418,0.98072) rectangle (0.68574,1.00000);
\fill[white] (0.68842,0.99266) rectangle (0.72039,1.00000);
\fill[white] (0.72553,0.99451) rectangle (0.75371,1.00000);
\fill[white] (0.75723,0.99197) rectangle (0.78469,1.00000);
\fill[white] (0.78734,0.98655) rectangle (0.81432,1.00000);
\fill[white] (0.81946,0.99451) rectangle (0.84764,1.00000);
\fill[white] (0.85116,0.99197) rectangle (0.90704,1.00000);
\fill[white] (0.07444,0.58179) rectangle (0.12920,0.66111);
\fill[white] (0.13188,0.58179) rectangle (0.21137,0.66111);
\fill[white] (0.21411,0.58179) rectangle (0.26217,0.66111);
\fill[white] (0.26490,0.58179) rectangle (0.34788,0.66111);
\fill[white] (0.35056,0.58248) rectangle (0.37814,0.67078);
\fill[white] (0.00000,0.00000) rectangle (0.01913,0.00741);
\end{scope}
\end{tikzpicture}
\end{sourceblock}
\begin{itemize}
\item 重复直到收敛。
\end{itemize}
更新等式的右侧。 （12.14），来自邻居块的数据是必要的。在图 12.21 的示例中，处理器 1 需要来自处理器 2 和 3 的数据。在具有共享内存的并行计算机上，该数据可由处理器直接访问。当使用具有分布式内存的计算机时（对于大多数集群来说都是如此），处理器之间的通信是必要的。然后，每个处理器需要存储来自接口另一侧的一层或多层单元的数据。区分本地 (LC) 和全局 (GC) 通信非常重要。

本地通信发生在相邻块上运行的处理器之间。它可以在成对的处理器之间同时发生；一个例子是上述问题中内部迭代内的通信。 GC 意味着从“主”处理器中的所有块收集一些信息，并将一些信息广播回其他处理器。一个例子是通过从处理器收集残差并广播收敛检查的结果来计算残差的范数。有一些通信库可用于此目的；如今，消息传递接口（MPI；请参阅 \url{https://www.mpi-forum.org/docs/} ）是事实上的标准，但其他标准，如 PVM (Sunderam 1990 ) 或 TCGMSG (Harrison 1991 ) 过去已经使用过，并且可能仍然可用。这使得代码具有可移植性——在不同的并行计算机上使用 CFD 代码时通常无需对其进行调整。

如果分配给每个处理器的单元或计算点的数量（即每个处理器的负载）在网格细化时保持不变（这意味着使用更多的处理器），则本地通信时间与计算时间的比率将保持不变。我们说LC是完全可扩展的。然而，当处理器数量增加时，GC 时间会增加，与每个处理器的负载无关。随着处理器数量的增加，全局通信时间最终将变得大于计算时间。因此，GC 是不可扩展的，是大规模并行的限制因素。下面讨论测量效率的方法。

\subsection*{12.6.3 域及时分解}\phantomsection\addcontentsline{toc}{subsection}{12.6.3 域及时分解}
隐式方法通常用于解决稳态流问题。尽管人们很容易认为这些方法不太适合并行计算，但它们可以通过使用时间和空间上的域分解来有效地并行化。这意味着多个处理器同时在同一子域上以不同的时间步长执行工作。该技术首先由 Hackbusch (1984) 提出。

由于不需要在外部迭代中精确求解任何方程，因此也可以将离散方程中的“旧”变量（即来自先前时间步的变量）视为未知数。对于二时间级方案，时间步 n 处的解方程可以写成：

\begin{sourceblock}
\includegraphics[page=503,trim={53.00bp 176.59bp 53.87bp 471.09bp},clip,width=0.92\textwidth,max height=0.38\textheight]{Ferziger4th.pdf}
\end{sourceblock}
因为我们正在考虑隐式方案，所以矩阵和源向量可能取决于新的解，这就是它们带有索引 n 的原因。同时求解多个时间步长的最简单迭代方案是解耦每个时间步长的方程，并在必要时使用变量的现有值（即来自先前外部迭代的值）。这允许在当前时间步长的解的第一个估计可用时，即在执行一次外部迭代之后，立即开始下一个时间步长的计算。包含来自先前时间步的信息的额外源项在每个外部之后更新

\begin{sourceblock}
\begin{tikzpicture}
\node[inner sep=0,anchor=south west] (image) at (0,0) {\includegraphics[page=504,trim={53.00bp 581.00bp 53.87bp 58.18bp},clip,width=0.92\textwidth,max height=0.38\textheight]{Ferziger4th.pdf}};
\begin{scope}[x={(image.south east)},y={(image.north west)}]
\fill[white] (0.53949,0.49666) rectangle (0.56403,0.89318);
\fill[white] (0.63314,0.56528) rectangle (0.66848,0.84941);
\fill[white] (0.61591,0.35052) rectangle (0.63895,0.74703);
\fill[white] (0.95820,0.31306) rectangle (0.99353,0.59718);
\fill[white] (0.94205,0.04525) rectangle (0.96662,0.44177);
\end{scope}
\end{tikzpicture}
\end{sourceblock}
\begin{sourceblock}
\begin{tikzpicture}
\node[inner sep=0,anchor=south west] (image) at (0,0) {\includegraphics[page=504,trim={53.00bp 523.39bp 53.87bp 97.19bp},clip,width=0.92\textwidth,max height=0.38\textheight]{Ferziger4th.pdf}};
\begin{scope}[x={(image.south east)},y={(image.north west)}]
\fill[white] (0.00000,0.96773) rectangle (0.03525,1.00000);
\fill[white] (0.03681,0.96773) rectangle (0.14490,1.00000);
\fill[white] (0.14641,0.96773) rectangle (0.17107,1.00000);
\fill[white] (0.17260,0.96773) rectangle (0.25380,1.00000);
\fill[white] (0.95820,0.80553) rectangle (0.97212,0.97366);
\fill[white] (0.91299,0.36435) rectangle (0.93377,0.59899);
\fill[white] (0.96036,0.18503) rectangle (0.99570,0.35316);
\end{scope}
\end{tikzpicture}
\end{sourceblock}
\begin{sourceblock}
\begin{tikzpicture}
\node[inner sep=0,anchor=south west] (image) at (0,0) {\includegraphics[page=504,trim={53.00bp 494.76bp 53.87bp 154.81bp},clip,width=0.92\textwidth,max height=0.38\textheight]{Ferziger4th.pdf}};
\begin{scope}[x={(image.south east)},y={(image.north west)}]
\fill[white] (0.87101,0.44357) rectangle (0.90635,0.90585);
\fill[white] (0.95928,0.46530) rectangle (0.99462,0.92758);
\fill[white] (0.85702,0.09294) rectangle (0.87609,0.73808);
\fill[white] (0.94313,0.07242) rectangle (0.96770,0.71756);
\end{scope}
\end{tikzpicture}
\end{sourceblock}
迭代，而不是像串行处理那样保持不变。当在时间级别 t n 上工作的处理器 k 正在执行其第 m 次外迭代时，在时间级别 t n - 1 上工作的处理器 k − 1 正在执行其第 ( m + 1 ) 次外迭代。那么处理器 k 在第 m 次外迭代中要求解的方程组为：

\begin{sourceblock}
\includegraphics[page=504,trim={53.00bp 385.03bp 53.87bp 262.61bp},clip,width=0.92\textwidth,max height=0.38\textheight]{Ferziger4th.pdf}
\end{sourceblock}
每次外部迭代，处理器只需交换数据一次，即线性方程求解器不受影响。当然，每次传输的数据比基于空间域分解的方法要多得多。如果并行处理的时间步数不大于每个时间步的外部迭代数，则使用滞后的旧值不会导致每个时间步的计算量显著增加。如果并行计算太多时间步，每个时间步所需的外部迭代次数将会增加，效率将会降低。在并行序列的最后一个时间步上，项 (B n φ n − 1 ) k − 1 包含在源项中，因为它在迭代内不会更改。

\begin{center}\small 图 12.23 显示了在四个时间步上同时求解的两时间级方案的矩阵结构。 Burmeister 和 Horton (1991)、Horton (1991) 和 Seidl 等人已使用 CFD 问题的时间并行求解方法。 （1996）等。该方法也可以应用于多级方案；在这种情况下，处理器必须从多个时间级别发送和接收数据。\end{center}
\subsection*{12.6.4 并行计算的效率}\phantomsection\addcontentsline{toc}{subsection}{12.6.4 并行计算的效率}
\begin{sourceblock}
\begin{tikzpicture}
\node[inner sep=0,anchor=south west] (image) at (0,0) {\includegraphics[page=504,trim={53.00bp 75.81bp 53.87bp 561.49bp},clip,width=0.92\textwidth,max height=0.38\textheight]{Ferziger4th.pdf}};
\begin{scope}[x={(image.south east)},y={(image.north west)}]
\fill[white] (0.41624,0.52670) rectangle (0.44472,0.95839);
\fill[white] (0.34286,0.28710) rectangle (0.37269,0.71602);
\fill[white] (0.37865,0.32455) rectangle (0.40683,0.72538);
\end{scope}
\end{tikzpicture}
\end{sourceblock}
\begin{sourceblock}
\includegraphics[page=505,trim={53.00bp 424.25bp 53.87bp 52.00bp},clip,width=\textwidth,max height=0.38\textheight]{Ferziger4th.pdf}
\par\vspace{0.45\baselineskip}
\small 图 12.24 使用 Simcenter STAR-CCM+ 商业代码中的分离和耦合求解算法以及由 10.2 亿个单元组成的网格，计算勒芒赛车周围的湍流时，加速比与所用处理器单元（核心）数量的函数关系
\end{sourceblock}
这里 T s 是单个处理器上最佳串行算法的执行时间，T n 是使用 n 个处理器的并行算法的执行时间。一般情况下T s ̸ = T 1，因为最佳串行算法可能与最佳并行算法不同；人们不应该根据在单个处理器上执行的并行算法的性能来确定效率。

加速比通常小于n（理想值），因此效率通常小于1（或100\%）。然而，在求解耦合非线性方程时，结果可能是在两个或四个处理器上的求解比在 1 个处理器上的求解效率更高，因此原则上，效率高于 100\% 是可能的（效率的提高通常是由于当一个处理器解决较小的问题时更好地使用高速缓存）。一个例子如图12.24所示。

虽然没有必要，但处理器通常在每次迭代开始时进行同步。由于一次迭代的持续时间由具有最多 CV 数量的处理器决定，因此其他处理器会经历一些空闲时间。延迟还可能是由于不同子域中的不同边界条件、不同数量的邻居或更复杂的通信造成的。

计算时间T s 可表示为：

\begin{sourceblock}
\includegraphics[page=505,trim={53.00bp 139.39bp 53.87bp 509.83bp},clip,width=0.92\textwidth,max height=0.38\textheight]{Ferziger4th.pdf}
\end{sourceblock}
其中 N cv 是 CV 总数，τ 是每个浮点运算的时间，i s 是每个 CV 达到收敛所需的浮点运算次数。对于在 n 个处理器上执行的并行算法，总执行时间由计算和通信时间组成：

\begin{sourceblock}
\begin{tikzpicture}
\node[inner sep=0,anchor=south west] (image) at (0,0) {\includegraphics[page=506,trim={53.00bp 592.70bp 53.87bp 55.33bp},clip,width=0.92\textwidth,max height=0.38\textheight]{Ferziger4th.pdf}};
\begin{scope}[x={(image.south east)},y={(image.north west)}]
\fill[white] (0.35155,0.08890) rectangle (0.36683,0.56267);
\end{scope}
\end{tikzpicture}
\end{sourceblock}
其中 N CV
n 是最大子域中的 CV 数量，T com
n 是无法进行计算的总通信时间。将这些表达式插入到总效率的定义中，得出：

\begin{sourceblock}
\begin{tikzpicture}
\node[inner sep=0,anchor=south west] (image) at (0,0) {\includegraphics[page=506,trim={53.00bp 481.78bp 53.87bp 125.52bp},clip,width=0.92\textwidth,max height=0.38\textheight]{Ferziger4th.pdf}};
\begin{scope}[x={(image.south east)},y={(image.north west)}]
\fill[white] (0.19627,0.64361) rectangle (0.24421,0.86948);
\fill[white] (0.40589,0.29164) rectangle (0.42568,0.48827);
\fill[white] (0.53077,0.29351) rectangle (0.55558,0.48997);
\fill[white] (0.55426,0.36523) rectangle (0.57886,0.51088);
\fill[white] (0.39441,0.11370) rectangle (0.43534,0.25952);
\fill[white] (0.47065,0.11370) rectangle (0.50920,0.25952);
\fill[white] (0.71395,0.14072) rectangle (0.72926,0.28637);
\fill[white] (0.74647,0.17471) rectangle (0.78860,0.40075);
\fill[white] (0.32087,0.05286) rectangle (0.34066,0.24932);
\fill[white] (0.34250,0.05897) rectangle (0.37068,0.25561);
\fill[white] (0.37320,0.05455) rectangle (0.39468,0.25102);
\fill[white] (0.38872,0.02039) rectangle (0.40403,0.16621);
\fill[white] (0.43257,0.05455) rectangle (0.47089,0.25561);
\fill[white] (0.46493,0.02039) rectangle (0.48021,0.16621);
\end{scope}
\end{tikzpicture}
\end{sourceblock}
该方程并不精确，因为每个 CV 的浮点运算数量不是恒定的（由于算法中的分支以及边界条件仅影响某些 CV）。然而，确定影响总效率的主要因素就足够了。这些因素的含义是：

\begin{itemize}
\item E num n — 数值效率考虑了由于修改算法以允许并行化而达到收敛所需的每个网格节点的操作数量变化的影响；
\item E par n — 并行效率考虑了无法进行计算的通信所花费的时间；
\item E lb n — 负载平衡效率考虑了某些处理器因负载不均匀而空闲的影响。
\end{itemize}
当同时在时间和空间上进行并行化时，整体效率等于时间效率和空间效率的乘积。

通过测量获得收敛解所需的计算时间，可以轻松确定总效率。并行效率无法直接测量，因为所有外部迭代的内部迭代次数并不相同（除非由用户固定）。然而，如果我们在 1 和 n 个处理器上执行一定数量的外部迭代，并且每个外部迭代执行固定数量的内部迭代，则数值效率是统一的，并且总效率是并行效率和负载平衡效率的乘积。如果通过使所有子域相等而将负载均衡效率降低到统一，我们就得到了并行效率。一些计算机具有允许执行运算计数的工具；然后可以直接测量数值效率。

对于空间域和时间域分解，随着给定网格的处理器数量的增加，所有三种效率通常都会降低。这种下降既是非线性的又是与问题相关的。并行效率尤其受到影响，因为 LC 的时间几乎恒定，GC 的时间增加，并且由于子域大小减小，每个处理器的计算时间减少。对于时间并行化，当针对相同问题规模并行计算更多时间步时，GC 时间会增加，而 LC 和计算时间保持不变。然而，如果处理器数量增加超过一定限制（这取决于每个时间步所需的外部迭代数量），数值效率将不成比例地下降。负载平衡的优化通常很困难，特别是在网格非结构化且采用局部细化的情况下。有优化的算法，但是可能比流计算要花更多的时间！

并行效率可以表示为三个主要参数的函数：

\begin{itemize}
\item 数据传输的建立时间（称为延迟时间）；
\item 数据传输速率（通常以兆字节/秒表示）；
\item 每个浮点运算的计算时间（通常以Mflops 表示）。
\end{itemize}
对于给定的算法和通信模式，可以创建一个模型方程来将并行效率表示为这些参数和域拓扑的函数。 Schreck 和 Peri ́c (1993) 提出了这样一个模型，并表明可以很好地预测并行效率。人们还可以将数值效率建模为求解算法中的替代方案、求解器的选择和子域的耦合的函数。然而，基于类似流问题的经验的经验输入是必要的，因为算法的行为是依赖于问题的。如果解算法允许替代通信模式，这些模型将很有用；人们可以选择最适合所使用的计算机的模式。例如，可以在每次内部迭代之后、在每第二次内部迭代之后或者仅在每次外部迭代之后交换数据。每次共轭梯度迭代可以采用一次、两次或多次预处理器迭代；预调节器迭代可以包括每个步骤之后或仅在最后的本地通信。这些选项会影响数值效率和并行效率；为了找到最佳方案，需要进行权衡。

通信频率显然是并行计算的一个重要因素。耦合求解方法（将速度、压力和温度视为单个未知向量）执行的数据交换 (LC) 比依次求解每个变量的分离方法少五倍。此外，耦合求解方法每次迭代通常需要比分离求解器执行更多的计算操作。尽管交换的数据量是原来的五倍，但节省的延迟时间和更有利的通信与计算时间比率导致耦合求解器通常比分离求解器具有更高的效率。耦合求解器表现出超线性性能的情况也更常见，如图 12.24 所示。在本例中，由于Cell总数超过10亿，单核无法进行计算；使用的处理器数量最少为 1024 个（每个核心约 100 万个单元）。然而，使用耦合求解器对 48 倍核心（每个核心约 20,000 个单元）进行计算的速度快了 50 倍以上，使用分离求解器的计算速度快了约 40 倍。

空间和时间组合并行化比纯空间并行化更高效，因为对于给定的问题规模，效率会随着处理器数量的增加而下降。表 12.2 显示了在最大雷诺数为 10 4 时，使用 32 × 32 × 32 CV 网格和时间步长 t = T / 200 计算具有振荡盖的立方体腔中非定常 3D 流动的结果，其中 T 是盖子振荡的周期。当使用十六个处理器和四个时间步长时

\begin{center}\small 表 12.2 计算具有振荡盖的 3D 腔流时，空间和时间上的各种域分解的数值效率\end{center}
\begin{sourceblock}
\begin{tikzpicture}
\node[inner sep=0,anchor=south west] (image) at (0,0) {\includegraphics[page=508,trim={53.00bp 445.83bp 53.87bp 80.60bp},clip,width=0.92\textwidth,max height=0.38\textheight]{Ferziger4th.pdf}};
\begin{scope}[x={(image.south east)},y={(image.north west)}]
\fill[white] (0.00782,0.91933) rectangle (0.16824,0.99141);
\fill[white] (0.16980,0.91933) rectangle (0.19444,0.99141);
\fill[white] (0.25744,0.91933) rectangle (0.32024,0.99141);
\fill[white] (0.32180,0.91933) rectangle (0.40442,0.99141);
\fill[white] (0.40602,0.91933) rectangle (0.43203,0.99141);
\fill[white] (0.43359,0.91933) rectangle (0.49074,0.99141);
\fill[white] (0.50710,0.91933) rectangle (0.56989,0.99141);
\fill[white] (0.57143,0.91933) rectangle (0.65402,0.99141);
\fill[white] (0.65561,0.91933) rectangle (0.68162,0.99141);
\fill[white] (0.68319,0.91933) rectangle (0.74033,0.99141);
\fill[white] (0.00782,0.84804) rectangle (0.06920,0.92005);
\fill[white] (0.07077,0.84804) rectangle (0.11236,0.92005);
\fill[white] (0.11392,0.84804) rectangle (0.16403,0.92005);
\fill[white] (0.25744,0.84804) rectangle (0.35708,0.92005);
\fill[white] (0.35862,0.84804) rectangle (0.39597,0.92005);
\fill[white] (0.39750,0.84804) rectangle (0.44761,0.92005);
\fill[white] (0.44917,0.84804) rectangle (0.49501,0.92005);
\fill[white] (0.50710,0.84804) rectangle (0.60671,0.92005);
\fill[white] (0.60824,0.84804) rectangle (0.64559,0.92005);
\fill[white] (0.64716,0.84804) rectangle (0.69726,0.92005);
\fill[white] (0.69880,0.84804) rectangle (0.74463,0.92005);
\fill[white] (0.00851,0.77732) rectangle (0.02463,0.84940);
\fill[white] (0.02776,0.77897) rectangle (0.05254,0.85098);
\fill[white] (0.05501,0.77732) rectangle (0.07113,0.84940);
\fill[white] (0.07329,0.77897) rectangle (0.09808,0.85098);
\fill[white] (0.09937,0.77732) rectangle (0.11408,0.84940);
\fill[white] (0.11618,0.77897) rectangle (0.14096,0.85098);
\fill[white] (0.14177,0.77732) rectangle (0.15368,0.84940);
\end{scope}
\end{tikzpicture}
\end{sourceblock}
并行计算，将空间域分解为4个子域，总数值效率为97\%。如果所有处理器仅用于空间或时间分解，则数值效率会降至 70\% 以下。

处理器之间的通信会导致许多机器上的计算停止。然而，如果通信和计算可以同时进行（这在某些新的并行计算机上是可能的），则可以重新安排解决算法的许多部分以利用这一点。例如，虽然 LC 在求解器中进行，但可以在子域内部进行计算，不需要来自相邻子域的数据。通过时间并行性，人们可以在 LC 发生的同时组合新的系数和源矩阵。如果按照 Demmel 等人的建议重新排列算法，即使是共轭梯度求解器中的 GC（这似乎会阻碍执行）也可以与计算重叠。 （1993）。可以在早期阶段跳过稳态计算中的收敛性检查，或者可以重新安排收敛准则来监控先前迭代的残差水平——可以根据基于残差减少率的投影来决定停止迭代。

PericandSchreck(1995)更详细地分析了重叠通信和计算的可能性，发现它可以显著提高并行效率。新的硬件和软件可能允许计算和通信的并发，因此可以预期并行效率可以得到优化。并行隐式 CFD 算法开发人员最关心的问题之一是数值效率。重要的是，对于相同的精度，并行算法不需要比串行算法更多的计算操作。

结果表明，并行计算可以有效地应用于 CFD。在这方面，工作站集群的使用特别有用，因为它们几乎可供所有用户使用，并且大问题并不总是得到解决。现在所有的计算机（个人电脑、工作站和大型机）都是多处理器机器；因此，在开发新的解决方法时必须考虑并行处理。

\subsection*{12.6.5 图形处理单元（GPU）和并行处理}\phantomsection\addcontentsline{toc}{subsection}{12.6.5 图形处理单元（GPU）和并行处理}
图形处理器 (GPU) 专为交互式游戏而设计，但已被证明可以有效解决流体流动问题。由于它们的设计不同，已经开发了一种单独的语言（统一计算设备架构 - CUDA），并且该领域的文献也不断增多。例如，虽然 GPU 在工作站中用作单处理器时可提供令人印象深刻的速度（比 CPU 快 20 倍），但在具有数百个内核的多 GPU 工作站中，计算加速却要大得多。

迄今为止，许多应用程序都来自使用本地软件包链接GPU框架的项目；然而，存在更通用的软件。在此，我们只是想提醒读者注意这个机会并提供一些参考。 Khajeh-Saeed 和 Perot (2013) 使用 GPU 加速的超级计算机对湍流进行了 DNS 计算，演示了如何链接和优化 GPU、MPI 和超级计算机。我们注意到数值方法的基本算法没有改变，大部分精力都花在了 CUDA 的编码以及链接和优化工作上。沙尔克韦克等人。 ( 2012a ) 提出了在具有链接 GPU 的台式 PC 上模拟湍流云的气象示例。他们的侧边栏“移植到 GPU”特别有启发性。通常，应用程序会采用当前代码的缓慢部分，并单独为该部分实现 GPU 求解器，从而实现显著的加速（例如，Williams 等人，2016 年，他们为非结构化网格开发了 GPU 版本的稀疏线性矩阵压力求解器）。一些商业代码还可以链接到 GPU 以加速计算的适当部分。

使用 CPU 和 GPU 在自适应非结构化网格上执行复杂流动（涉及湍流、燃烧、多相流建模等）的计算是一个挑战，但如果要获得高效率，就没有办法解决这个问题。为更简单的流问题和结构化网格开发的方法都不是单独针对复杂问题的最佳方法，必须开发更智能的并行化概念来动态适应当前的问题。解算法和数据结构可能都需要重新组织，以促进在此类应用中有效使用 CPU 和 GPU。例如，面或体积上的某些循环可能需要分解为多个循环，并且某些数据可能需要存储两次（并在特定步骤之间复制），以最小化数据依赖性并允许计算和通信的并发性。预计未来十年将在这一领域投入大量研究工作。

\chapter{专题}
\section*{13.1 简介}\phantomsection\addcontentsline{toc}{section}{13.1 简介}
流体流动可能包括广泛的附加物理现象，这些现象使该主题远远超出了迄今为止本工作的重点的单相非反应流。流动的流体中可能发生许多类型的物理过程。其中每一个都可能与流动相互作用，产生一系列令人惊奇的新现象。几乎所有这些过程都发生在重要的应用程序中。计算方法已应用于它们并取得了不同程度的成功。

可以添加到流中的最简单的元素是标量，例如可溶化学物质的浓度或温度。标量的存在不影响流体性质的情况已在前面的章节中讨论过；在这种情况下，我们称之为被动标量。在更复杂的情况下，流体的密度和粘度可能会因标量的存在而改变，并且我们有一个活动标量。在一个简单的例子中，流体特性是温度或物质浓度的函数。这个领域被称为传热传质。

在其他情况下，溶解标量的存在或流体本身的物理性质会导致流体的行为方式使得应力与应变率不通过简单的牛顿关系（方程 1.9）相关。在某些流体中，粘度成为瞬时应变率的函数，我们称之为剪切稀化或剪切增稠非牛顿流体。在更复杂的流体中，应力由一组附加的非线性偏微分方程确定。然后我们说流体是粘弹性的。许多聚合物材料，包括生物材料，都表现出这种行为，从而引起意想不到的流动现象。这是非牛顿流体力学领域。

流可能包含各种接口。这些可能是由于流体中存在固体所致。在这种简单的情况下，可以转换为随身体移动的坐标系，并且问题被简化为之前处理的类型之一，尽管是在复杂的几何形状中。在其他问题中，可能存在相对运动的物体，因此别无选择，只能引入运动坐标系。此类特别重要且困难的情况是表面可变形的情况。液体表面就是这种类型的例子。

在其他流程中，多个相可以共存。所有可能的组合都很重要。固气情况包括大气中的灰尘、流化床和穿过多孔介质的气流等现象。在固液类别中是浆体（其中液体是连续相），同样，多孔介质流动。气液流包括喷雾流（其中气相是连续的）和气泡流（其中相反）。最后，可能存在三相流。每个案例都有许多子类别。

流体中可能会发生化学反应，而且这种情况也有很多个别情况。当反应物质稀释时，反应速率可以假定为恒定（然而，它们可以取决于温度），并且反应物质相对于它们对流动的影响本质上是被动标量。此类的例子是大气或海洋中的污染物种类。另一种反应涉及主要物种并释放大量能量。这就是燃烧的情况。又一个例子是高速气流；压缩性效应可能导致温度大幅升高以及气体解离或电离的可能性。

地球物理学和天体物理学也需要求解流体运动方程。除了等离子体效应（下面讨论）之外，这些流动中的新元素与工程流动相比规模巨大。在气象学和海洋学中，旋转和分层对流动行为有很大影响。

最后，我们提到在等离子体（电离流体）中，电磁效应起着重要作用。在这个领域，流体运动方程必须与电磁方程（麦克斯韦方程）一起求解，并且现象和特殊情况的数量是巨大的。

在本章的剩余部分，我们将描述处理部分（但不是全部）这些困难的方法。我们应该指出，上述每个主题都是流体力学的一个重要子专业，并且有大量的文献专门讨论它；下面给出了每个领域的教科书参考资料。在此篇幅内不可能对这些主题中的每一个进行公正的阐述。

\section*{13.2 传热传质}\phantomsection\addcontentsline{toc}{section}{13.2 传热传质}
在该主题课程中通常介绍的三种传热机制（传导、辐射和对流）中，最后一种与流体力学关系最为密切。这种联系如此紧密，以至于对流传热可以被视为流体力学的一个子领域。

稳态热传导由拉普拉斯方程（或与其非常相似的方程）描述，而非稳态传导则由热方程控制；这些方程很容易通过章节中提出的方法求解。 3、4和6。当特性依赖于温度时，情况就会变得复杂。在这种情况下，使用当前迭代中的温度计算属性，更新温度，并重复该过程。收敛通常几乎与固定属性情况一样快。

涉及固体表面的辐射与流体力学几乎没有关系（除了多种主动传热机制的问题）。有一些有趣的问题（例如，火箭喷嘴和燃烧室中的流动），其中流体力学和气体中的辐射传热都很重要。这种组合也出现在天体物理学应用和气象学中。我们在这里不讨论此类问题；然而，见第 13.7 节.

在层流对流换热中，主要过程是沿流向的平流（我们之前称为对流！）和沿垂直于流动方向的传导。当流动为湍流时，层流中传导所起的大部分作用被湍流所取代，并由湍流模型表示；这些模型将在第 1 章中讨论。 10.无论哪种情况，人们的兴趣通常集中在与固体表面的热能交换上。

如果温差很小（水中小于 5K，空气中小于 10K）并且雷诺数很高，则流体属性的变化并不重要，并且温度表现为被动标量。此类问题可以通过本书前面描述的方法来处理。由于在这种情况下温度是被动标量，因此可以在速度场计算完全收敛后计算温度，从而使任务变得更加简单。在流动由密度差驱动的情况下，必须考虑后者。这可以借助下面描述的 Boussinesq 近似来完成。

另一个重要的特殊情况是流过光滑形状物体时发生的传热。在这种类型的流动中，我们可以首先计算身体周围的势流，然后使用获得的压力分布作为边界层代码的输入来预测传热。如果边界层不与主体分离，则可以使用纳维-斯托克斯方程的边界层简化来计算这些流动（例如，参见 Kays 和 Crawford 1978；或 Cebeci 和 Bradshaw 1984）。边界层方程是抛物线方程，可以在现代工作站或个人计算机上在几秒钟（对于 2D 情况）或一分钟左右（对于 3D 情况）内求解。本工作并未详细介绍计算这些流量的方法（但一般原则可在第 3 章至第 8 章中找到）；有兴趣的读者可以在 Cebeci 和 Bradshaw（1984）以及 Patankar 和 Spalding（1977）的著作中找到它们。

在一般情况下，温度变化很大。它们以两种方式影响流量。第一个是通过传输特性随温度的变化。这些可能非常大，必须考虑在内，但在数字上处理并不困难。重要的问题是能量和动量方程现在是耦合的，必须同时求解。幸运的是，耦合通常不会太强，以至于无法以顺序方式求解方程。在每次外部迭代中，首先使用从“旧”温度场计算出的传输属性来求解动量方程。在获得新的外部迭代的动量方程的解并且更新属性之后，更新温度场。该技术与第 1 章中描述的使用湍流模型求解动量方程的技术非常相似。 10.

温度变化的另一个影响是密度变化与重力相互作用产生体积力，该体积力可以显著改变流动并且可能是流动中的主要驱动力。在后一种情况下，我们谈论浮力驱动或自然对流。强制对流和浮力效应的相对重要性通过格拉肖夫数和雷诺数之比来衡量。前者定义如式(1.33);如果比率
Gr Re 2 = Ra Pr Re 2 ≪ 1 ,
自然对流的影响可以忽略不计。在纯浮力驱动的流动中，如果密度变化足够小，则可以忽略垂直动量方程中除体积力之外的所有项中的密度变化。这称为 Boussinesq 近似，它允许通过与不可压缩流所使用的方法基本相同的方法来求解方程。 节 中提供了一个例子。 8.4.

浮力很重要的流动计算通常通过上述类型的方法进行，即速度场的迭代先于温度和密度场的迭代。由于场之间的耦合可能相当强，因此该过程可能比等温流收敛得更慢。作为耦合系统的方程的解提高了收敛速度，但代价是增加了编程和存储要求的复杂性；参见 Galpin 和 Raithby (1986) 的例子。耦合的强度还取决于普朗特数。对于普朗特数高的流体，它更强；对于这些流体，耦合求解方法的收敛速度比顺序方法快得多。然而，这种说法仅适用于稳态流的计算。在高瑞利数下，即使边界条件稳定，浮力驱动的流动也会变得不稳定并最终形成湍流；不稳定也可能是由时变边界条件引起的。当需要时间精确的解时，时间步必须足够小；在这种情况下，解从一个时间步长到另一个时间步长的变化不会太大，并且分离求解方法（如上述方法）实际上可能比耦合求解器在计算上更有效。这是因为隐式求解方法每个时间步只需要很少的外部迭代；耦合求解器每次迭代需要更多的计算时间，并且只有当它可以显著减少所需的迭代次数时才有益，这通常仅适用于稳定或弱瞬态流。如果时间步长足够小，甚至是非迭代时间推进方法，例如第 1 章中描述的分数步法。 7、可以使用；参见 Armfield 和 Street (2005) 的例子。

在某些应用中，需要考虑固体中的热传导以及相邻流体中的对流。此类问题称为共轭传热问题，需要通过在描述两种传热类型的方程之间进行迭代来求解。也可以求解能量方程

\begin{sourceblock}
\includegraphics[page=514,trim={53.00bp 453.25bp 53.87bp 52.00bp},clip,width=\textwidth,max height=0.38\textheight]{Ferziger4th.pdf}
\par\vspace{0.45\baselineskip}
\small 图13.1 带涂层的固液界面扩散通量的离散化，假设界面处有非共形网格
\end{sourceblock}
同时流体和固体域；我们描述了下面需要特别注意的离散化步骤，并且还考虑了太薄而无法被网格解析但可能代表对传热的显著阻力的涂层或挡板。

分别求解流体和固体中的能量方程并在每次外部迭代时更新每个连续体的边界条件通常效率不高；希望在所有连续体中同时求解能量方程。这要求将固-液界面处的热通量表示为流体和固体中相邻细胞中心的温度的函数，而不依赖于界面本身的温度知识。我们在下面描述如何实现这一点。考虑图13.1 所示的情况。如果假设在流体侧边界层已分解，则沿垂直于单元面 k 的线的温度分布将如图 13.1 右侧所示。为了一般性，假设固体和流体中的网格在界面处不匹配；网格被绘制为笛卡尔坐标，以指出即使这样，情况也需要特殊处理；如果网格不与界面正交，则适用相同的方法。

通过以节点 C 和 N k 为中心的单元共有的面 k 的每单位面积的热通量，见图 13.1，可以表示为：

\begin{sourceblock}
\begin{tikzpicture}
\node[inner sep=0,anchor=south west] (image) at (0,0) {\includegraphics[page=514,trim={53.00bp 138.26bp 53.87bp 489.88bp},clip,width=0.92\textwidth,max height=0.38\textheight]{Ferziger4th.pdf}};
\begin{scope}[x={(image.south east)},y={(image.north west)}]
\fill[white] (0.32457,0.47289) rectangle (0.36286,0.78421);
\fill[white] (0.20198,0.26763) rectangle (0.23125,0.59342);
\fill[white] (0.23753,0.29605) rectangle (0.26571,0.60026);
\fill[white] (0.26923,0.26579) rectangle (0.29741,0.60026);
\fill[white] (0.57657,0.03158) rectangle (0.59612,0.27421);
\end{scope}
\end{tikzpicture}
\end{sourceblock}
式中，T为温度，λ为热导率，n为垂直于池面的坐标方向，索引“f”、“z”和“s”分别表示界面的流体侧、涂层和固体侧（见图13.1）。涂层通常非常薄，因此可以假设电导率 λ z 以及整个层的温度梯度恒定。对于流体和固体，如果界面温度 T k f 和 T k s 与温度相关，则应分别采用热导率。

紧邻墙壁的流体中的温度变化始终是线性的；如果粘性层未解析（即使用壁函数），则必须在随后的近似中使用有效电导率。对于线性温度分布，热通量的离散近似为：

\begin{sourceblock}
\includegraphics[page=515,trim={53.00bp 497.80bp 53.87bp 138.36bp},clip,width=0.92\textwidth,max height=0.38\textheight]{Ferziger4th.pdf}
\end{sourceblock}
点 C ′ 和 N ′
k 位于面部正常 n；它们距人脸的距离 δ f 和 δ s 分别表示连接细胞中心与人脸中心的向量的投影，即：

\begin{sourceblock}
\includegraphics[page=515,trim={53.00bp 426.77bp 53.87bp 222.87bp},clip,width=0.92\textwidth,max height=0.38\textheight]{Ferziger4th.pdf}
\end{sourceblock}
引入阻力系数α如下：

\begin{sourceblock}
\includegraphics[page=515,trim={53.00bp 368.49bp 53.87bp 268.47bp},clip,width=0.92\textwidth,max height=0.38\textheight]{Ferziger4th.pdf}
\end{sourceblock}
表达式（13.2）可以重写为：

\begin{sourceblock}
\includegraphics[page=515,trim={53.00bp 321.08bp 53.87bp 327.48bp},clip,width=0.92\textwidth,max height=0.38\textheight]{Ferziger4th.pdf}
\end{sourceblock}
从这个表达式中，可以消除界面温度 T k f 和 T k s，得到：

\begin{sourceblock}
\includegraphics[page=515,trim={53.00bp 249.92bp 53.87bp 373.56bp},clip,width=0.92\textwidth,max height=0.38\textheight]{Ferziger4th.pdf}
\end{sourceblock}
通过引入有效电阻系数和有效电导率：

\begin{sourceblock}
\includegraphics[page=515,trim={53.00bp 180.42bp 53.87bp 443.86bp},clip,width=0.92\textwidth,max height=0.38\textheight]{Ferziger4th.pdf}
\end{sourceblock}
热通量可以方便地表示为：

\begin{sourceblock}
\includegraphics[page=515,trim={53.00bp 122.39bp 53.87bp 513.43bp},clip,width=0.92\textwidth,max height=0.38\textheight]{Ferziger4th.pdf}
\end{sourceblock}
如果网格是非正交或非共形的，则辅助节点 C ′ 和 N ′ 处的温度
k 必须通过适当的插值通过节点值来表示。以下近似与离散化过程中使用的其他近似一致，并且是二阶精确的：

\begin{sourceblock}
\includegraphics[page=516,trim={53.00bp 567.86bp 53.87bp 80.18bp},clip,width=0.92\textwidth,max height=0.38\textheight]{Ferziger4th.pdf}
\end{sourceblock}
辅助节点C′和N′的坐标
k 很容易获得（见式（13.3）和图 13.1）：

\begin{sourceblock}
\begin{tikzpicture}
\node[inner sep=0,anchor=south west] (image) at (0,0) {\includegraphics[page=516,trim={53.00bp 520.04bp 53.87bp 128.00bp},clip,width=0.92\textwidth,max height=0.38\textheight]{Ferziger4th.pdf}};
\begin{scope}[x={(image.south east)},y={(image.north west)}]
\fill[white] (0.00000,0.90608) rectangle (0.04740,1.00000);
\fill[white] (0.05008,0.90608) rectangle (0.17696,1.00000);
\end{scope}
\end{tikzpicture}
\end{sourceblock}
最终，单位面积的热通量可表示为（见式（13.8））：

\begin{sourceblock}
\begin{tikzpicture}
\node[inner sep=0,anchor=south west] (image) at (0,0) {\includegraphics[page=516,trim={53.00bp 472.22bp 53.87bp 175.83bp},clip,width=0.92\textwidth,max height=0.38\textheight]{Ferziger4th.pdf}};
\begin{scope}[x={(image.south east)},y={(image.north west)}]
\fill[white] (0.05561,0.20619) rectangle (0.08487,0.88999);
\fill[white] (0.09116,0.26589) rectangle (0.11934,0.90492);
\fill[white] (0.12283,0.17302) rectangle (0.21985,0.90492);
\fill[white] (0.22559,0.20177) rectangle (0.30328,0.90492);
\fill[white] (0.30511,0.26589) rectangle (0.33329,0.90492);
\fill[white] (0.33513,0.20177) rectangle (0.37982,0.90492);
\fill[white] (0.93209,0.26589) rectangle (0.94514,0.90492);
\fill[white] (0.90941,0.00000) rectangle (1.00000,0.13433);
\end{scope}
\end{tikzpicture}
\end{sourceblock}
。 (13.11) 为了获得通过单元表面的总热通量，需要将 q k 乘以表面面积 S k。带下划线的项可以被视为延迟修正，即使用先前迭代的值进行计算。右侧第一项被隐式处理，即它对矩阵方程的系数有贡献。

如果单元中心和辅助节点之间的距离与到单元面的距离相比较小，即当非正交性适中时，则与主项相比，下划线项较小。当通过面中心的面法线穿过单元中心时，它消失。在严重非正交性的情况下，可能会导致两种问题：（i）非物理解（过冲或下冲）或（ii）收敛问题。在这种情况下，可以限制下划线项（或者甚至将其设置为零作为最后的手段）；这会降低近似阶数，但有助于避免无法提高网格质量时出现的问题。例如，通过完全忽略下划线项，可以有效地设置 T C ' = T C，这对应于假设电池内的恒定温度（一阶近似）。

请注意等式。 ( 13.8 ) 类似于经常描述从墙壁到环境的热传递的表达式，

\begin{sourceblock}
\includegraphics[page=516,trim={53.00bp 233.62bp 53.87bp 417.06bp},clip,width=0.92\textwidth,max height=0.38\textheight]{Ferziger4th.pdf}
\end{sourceblock}
其中 q wall 是壁面热通量，T wall 是壁面温度，T ∞ 是环境温度，α 是所谓的传热系数；实验数据通常以这种形式呈现。许多 CFD 代码用户希望计算和可视化内部流动壁面的传热系数，例如涡轮叶片或热发动机表面的传热系数。这很棘手，因为传热系数不是唯一定义的量，只有壁热通量是唯一的（这是人们应该关注的量）。虽然很明显 T wall 应为局部壁温，但尚不清楚内部流动中应使用哪个环境参考温度。因此，α 和 T ∞ 必须始终成对出现。在一些商业代码中，可以从模拟结果中提取几种不同版本的传热系数；当网格被细化时，有些变化会很大，因此在比较解时必须小心。

\begin{center}\small 图 13.2 测试用例的几何形状，显示同时求解能量方程的所有固体和流体区域\end{center}
Insulation

热水

Air

Aluminium

冷水

100 2
Insulation

我们现在提出一个示例，其中主要关注的是外壳中空气的自然对流。我们在该节中讨论了类似的问题。 8.4.1，但我们只求解了流体域中的能量方程并规定了侧壁的温度，同时在顶部和底部边界使用绝热条件。在这里，我们考虑空腔周围的固体结构：顶壁上方和底壁下方的隔热层，以及侧通道中流动的热水和冷水，目的是保持壁温几乎恒定，如图 13.2 所示。在实验中，很难同时达到恒定壁温和零热通量（对应绝热边界条件）；这种边界条件始终是近似值，并且误差通常是未知的。只要有可能，就应该尝试指定尽可能远离感兴趣区域的近似边界条件。

在这个 2D 示例中，空气占据尺寸为 100 × 100mm 的空腔；它由 2 毫米厚的铝墙包围。左壁由在 20mm 宽的通道中以 10m/s 的平均速度向下流动的热水加热，温度为 310 K。在右侧，300K 的冷水在 20mm 宽的通道中以 10 m/s 的平均速度向上流动。充满空气的腔体中的瑞利数约为1.5×10 6，这意味着空气流动是层流的。水道中的雷诺数为20万，即流动为湍流；它是使用 k - ε 湍流模型进行模拟的。假设水不会向外侧壁散失热量；这是一个近似值，但由于流速相当高，即使一些热量通过另一壁逸出，传递到空气腔的热量也不会受到影响。还假设两个隔热层的顶部和底部边界处没有热量损失；由于它们的厚度为 100 毫米，并且材料的热导率非常低（0.036 W/mK，而铝框架为 237 W/mK），因此这种近似值的影响预计也很小。气流的解域上没有施加任何边界条件——驱动气流的温度是共轭传热解的一部分；仅施加两个进水口处的速度和温度。

\begin{center}\small 图 13.3 流体域中的流动模式（上）以及空气、铝框架和绝缘体中的温度轮廓（下）\end{center}
为了简化网格生成，所有区域都使用统一的网格，网格间距为0.5mm。出于演示目的，这在流体区域的靠近墙壁处已经足够了，但是在绝缘区域中的网格太细了。从下面给出的结果中可以明显看出这一点。迭代求解方法遵循之前概述的计算稳态流的步骤：在每次外部迭代中，首先使用来自先前外部迭代的流体区域中的速度场求解所有区域的能量方程来更新温度。然后，在每个流体区域中执行 SIMPLE 算法的一步（即依次求解动量方程和压力校正方程），使用更新的温度来确定空气的流体特性。在水的强制流动中，由于温度变化很小，特性保持恒定，而空气则被视为理想气体。未使用 Boussinesq 近似：使用状态方程在所有项中的每次外部迭代中更新密度，同时使用多项式表达式考虑粘度随温度的变化。欠松弛参数：速度、k 和 ε 为 0.8，压力为 0.3，温度为 0.99。进行了大约 500 次外部迭代；残留水平为

\begin{sourceblock}
\includegraphics[page=519,trim={53.00bp 336.25bp 53.87bp 52.00bp},clip,width=\textwidth,max height=0.38\textheight]{Ferziger4th.pdf}
\par\vspace{0.45\baselineskip}
\small 图 13.4 沿中间高度水平切割的剖面：空气中的垂直速度分量（上）和所有材料的温度（下）；灰色条表示空心壁
\end{sourceblock}
所有方程都减少了 6 个数量级。线性迎风方案（二阶；参见第 4.4.5 节）用于对流和扩散项的中心差。带有模拟文件的详细报告可供下载；详细内容请参见附录。

\begin{center}\small 图 13.3 显示了流体域中的流动模式和温度等值线。空气沿着热壁上升，然后转向冷壁，释放从热壁收集的热量，然后沿着底壁返回。从图13.4所示的速度分布图可以看出，在腔体的中心部分，空气几乎是停滞的。在空腔的中心部分（稳定分层区域），等温线也几乎是水平的。空气中的等温线接近铝框架，以锐角将空气与绝缘材料分开，表明空气和金属框架之间存在显著的热传递；另一方面，在金属和绝缘材料之间，等温线几乎垂直于界面，表明这两种材料之间的热交换非常低。空气与铝在上下界面处的热交换主要沿框架传导。在隔热层中，正如预期的那样，温度从热墙到冷墙呈线性变化；不规则变化仅在型腔拐角附近可见。速度和温度在热壁和冷壁附近都有很高的梯度，需要在壁法线方向上有精细的网格。\end{center}
\begin{sourceblock}
\includegraphics[page=520,trim={53.00bp 329.25bp 53.87bp 52.00bp},clip,width=\textwidth,max height=0.38\textheight]{Ferziger4th.pdf}
\par\vspace{0.45\baselineskip}
\small 图 13.5 沿空气侧热壁（上）和沿热水和铝壁中高处水平切口（下）的温度分布
\end{sourceblock}
观察冷热壁上的温度是否恒定是很有趣的。图 13.5 显示了沿空气侧热壁的温度变化。它不仅不是恒定的，而且在中心部分它比热水温度 (310 K) 还要高！人们可能会认为代码中存在错误；然而，仔细观察整个热水通道的温度变化会发现，水沿着两个通道壁都会稍微升温，而在中心部分，其温度恒定为 310 K，如入口处指定的那样。这是因为求解的能量方程包含由于粘性热产生而产生的源项；当水以 10m/s 的速度在 20mm 宽的通道中流动时，会产生可测量的热量增益。这就是为什么所有流体通过测试部分再循环的实验都需要有热交换器来带走产生的热量并保持工作温度恒定。因此，温度在大约 310K 左右略高。 65\% 的热墙，而它是大约。靠近拐角处降低了 0.3 K。当热壁和冷壁温差为10 K时，与恒温的偏差很大。因此，只要有可能，就应该在模拟中包含尽可能多的实验设置。

在这个例子中，气流是层流；当感兴趣的流动是湍流时，需要解析近壁层（即使用低Re壁处理）。如果这是负担不起的，那么应该确保壁附近的第一个计算点在对数范围内，以便壁函数是合适的。虽然所谓的“盟友+
大多数商业代码中都提供“墙函数”，但只有当一小部分墙面上的网格处于中间范围时它们才有用。当大部分墙面上的网格 y + 介于 5 到 20 之间时，结果可能会很差。在执行网格依赖性研究时需要考虑到这一点：如果更精细的网格落入缓冲范围，则连续网格上的解之间的差异可能会增加而不是减少！

\section*{13.3 具有可变流体属性的流动}\phantomsection\addcontentsline{toc}{section}{13.3 具有可变流体属性的流动}
尽管我们主要处理不可压缩流，但密度、粘度和其他流体属性已保留在微分算子内。这允许使用前面章节中介绍的离散化和求解方法来解决具有可变流体属性的问题。

流体性质的变化通常是由温度变化引起的；压力变化也会影响密度的变化。第 1 章考虑了这种变化。 11，我们处理可压缩流。然而，在许多情况下，压力不会发生显著变化，但温度和/或溶质浓度可能会导致流体性质发生较大变化。例如减压下的气流、液态金属的流动（晶体生长、凝固和熔化问题等）以及由溶解盐分层的流体的环境流动。

密度、粘度、普朗特数和比热的变化会增加方程的非线性。顺序求解方法可以以与应用于可变温度流的方式大致相同的方式应用于这些流。在每次外部迭代之后重新计算流体属性，并在下一次外部迭代期间将它们视为已知的。如果属性变化很大，收敛速度可能会大大减慢。对于稳态流，多重网格方法可以带来显著的加速；参见杜斯特等人。 (1992) 是应用于金属有机化学气相沉积问题的示例，Kadinski 和 Peri ́c (1996) 是应用于涉及热辐射的问题。

对于瞬态流，特别是当时间步相对较小时，前面描述的顺序求解方法可能效果很好，因为从一个时间步到下一时间步的解的变化并不大。无论如何，通常都需要时间步内的外部迭代来更新延迟校正（例如，由于网格非正交性而导致的校正）和非线性项（对流通量、流体特性、湍流模型的贡献等）。当问题变得僵硬时，可能需要更强的欠松弛和每个时间步更多的外部迭代。

大气和海洋中的流动是变密度流动的特殊例子；它们将在稍后讨论。

\section*{13.4 移动网格}\phantomsection\addcontentsline{toc}{section}{13.4 移动网格}
在许多应用领域中，由于边界的移动，求解域会随着时间而变化。运动由外部效应（如活塞驱动流）或作为解一部分的计算确定（例如，在浮动或飞行物体的情况下）。无论哪种情况，网格都必须移动以适应不断变化的边界。如果坐标系保持固定并使用笛卡尔速度分量，则守恒方程中的唯一变化是对流项中相对速度的出现；见第 1.2 节.我们在这里简要描述如何导出移动网格系统的方程。

首先考虑一维连续性方程：

\begin{sourceblock}
\includegraphics[page=522,trim={53.00bp 435.57bp 53.87bp 202.29bp},clip,width=0.92\textwidth,max height=0.38\textheight]{Ferziger4th.pdf}
\end{sourceblock}
通过将该方程积分到其边界随时间移动（即从 x 1 ( t ) 到 x 2 ( t ) ）的控制体积上，我们得到：

\begin{sourceblock}
\includegraphics[page=522,trim={53.00bp 351.86bp 53.87bp 271.39bp},clip,width=0.92\textwidth,max height=0.38\textheight]{Ferziger4th.pdf}
\end{sourceblock}
第二项不会造成任何问题。第一个需要使用莱布尼茨规则，因此，方程。 ( 13.14 ) 变为：

\begin{sourceblock}
\begin{tikzpicture}
\node[inner sep=0,anchor=south west] (image) at (0,0) {\includegraphics[page=522,trim={53.00bp 268.12bp 53.87bp 355.13bp},clip,width=0.92\textwidth,max height=0.38\textheight]{Ferziger4th.pdf}};
\begin{scope}[x={(image.south east)},y={(image.north west)}]
\fill[white] (0.17203,0.75239) rectangle (0.21916,0.97156);
\fill[white] (0.13489,0.54069) rectangle (0.15468,0.81021);
\fill[white] (0.12935,0.21287) rectangle (0.15811,0.48473);
\end{scope}
\end{tikzpicture}
\end{sourceblock}
导数 d x /d t 表示网格（CV 表面）移动的速度；我们用 v s 表示它。因此，方括号中的项的形式类似于涉及流体速度的最后两项，因此我们可以重写方程： （13.14）为：

\begin{sourceblock}
\includegraphics[page=522,trim={53.00bp 170.25bp 53.87bp 452.99bp},clip,width=0.92\textwidth,max height=0.38\textheight]{Ferziger4th.pdf}
\end{sourceblock}
当边界随着流体速度移动时，即 v s = v，第二个积分变为零，我们有拉格朗日质量守恒方程 d m /d t = 0。

方程的三维版本。 ( 13.15 )（使用莱布尼茨法则的 3D 版本获得）给出：

\begin{sourceblock}
\begin{tikzpicture}
\node[inner sep=0,anchor=south west] (image) at (0,0) {\includegraphics[page=522,trim={53.00bp 75.54bp 53.87bp 561.40bp},clip,width=0.92\textwidth,max height=0.38\textheight]{Ferziger4th.pdf}};
\begin{scope}[x={(image.south east)},y={(image.north west)}]
\fill[white] (0.20502,0.55959) rectangle (0.22481,0.95548);
\fill[white] (0.19952,0.07808) rectangle (0.22824,0.47740);
\fill[white] (0.25113,0.04110) rectangle (0.26875,0.33493);
\end{scope}
\end{tikzpicture}
\end{sourceblock}
或者，在上面使用的符号中：

\begin{sourceblock}
\begin{tikzpicture}
\node[inner sep=0,anchor=south west] (image) at (0,0) {\includegraphics[page=523,trim={53.00bp 558.93bp 53.87bp 78.11bp},clip,width=0.92\textwidth,max height=0.38\textheight]{Ferziger4th.pdf}};
\begin{scope}[x={(image.south east)},y={(image.north west)}]
\fill[white] (0.26544,0.56117) rectangle (0.28526,0.95876);
\fill[white] (0.25994,0.07835) rectangle (0.28866,0.47904);
\end{scope}
\end{tikzpicture}
\end{sourceblock}
见第 1.2 节 我们注意到，守恒定律可以通过使用方程（1）从控制质量形式转换为控制体积形式。 （1.4）；这也引出了上面的质量守恒方程。相同的方法可以应用于任何传输方程。

当 CV 表面以速度 v s 移动时，第 i 个动量分量的守恒方程的积分形式采用以下形式：

\begin{sourceblock}
\begin{tikzpicture}
\node[inner sep=0,anchor=south west] (image) at (0,0) {\includegraphics[page=523,trim={53.00bp 453.84bp 53.87bp 183.20bp},clip,width=0.92\textwidth,max height=0.38\textheight]{Ferziger4th.pdf}};
\begin{scope}[x={(image.south east)},y={(image.north west)}]
\fill[white] (0.04526,0.56117) rectangle (0.06505,0.95876);
\fill[white] (0.11389,0.29656) rectangle (0.15907,0.73127);
\fill[white] (0.16280,0.32131) rectangle (0.20235,0.72199);
\fill[white] (0.20875,0.33402) rectangle (0.23693,0.73127);
\fill[white] (0.03976,0.07801) rectangle (0.06851,0.47904);
\fill[white] (0.09140,0.04124) rectangle (0.10905,0.33574);
\fill[white] (0.90941,0.00000) rectangle (1.00000,0.14742);
\end{scope}
\end{tikzpicture}
\end{sourceblock}
(13.19) 通过将对流项中的速度矢量替换为相对速度 v − v s，可以很容易地从固定 CV 的相应方程导出标量守恒方程。

显然，如果边界以与流体相同的速度移动，则通过 CV 面的质量通量将为零。如果所有 CV 面都是如此，则相同的流体保留在 CV 内并成为控制质量；然后我们就有了流体运动的拉格朗日描述。另一方面，如果 CV 不移动，则方程是之前处理过的方程。上述方程中的时间导数在固定网格和移动网格中具有不同的含义，尽管其近似方式相同。如果CV不动，时间导数表示守恒量在固定位置的局部变化，记为 ∂φ/∂ t；当 CV 移动时，我们用 d φ/ d t 来表示在空间移动的位置处 φ 随时间的变化。在表面完全随流体速度移动的 CV 的上述极端情况中，时间导数变为总（材料）导数，因为 CV 始终包含相同的流体，因此代表控制质量。时间导数的意义变化是由对流通量引起的，对流通量也根据 CV 运动而变化。

当网格的运动已知为时间的函数时，纳维-斯托克斯方程的解不会带来新问题：我们只需使用单元表面的相对速度分量来计算对流通量（例如质量通量）。然而，当单元面移动时，如果使用网格速度来计算质量通量，则不一定能确保质量守恒（以及所有其他守恒量）。例如，考虑具有隐式欧拉时间积分的连续性方程；为了简单起见，我们假设 CV 是矩形的，并且流体不可压缩并且以恒定速度移动。图 13.6 显示了新旧时间级别下 CV 的相对大小。我们还假设网格线（CV 面）以恒定但不同的速度移动，因此 CV 的大小随着时间而增长。

\begin{center}\small 图13.6 矩形控制体积，由于其边界处的网格速度不同，其尺寸随时间增加\end{center}
新职位

\begin{sourceblock}
\begin{tikzpicture}
\node[inner sep=0,anchor=south west] (image) at (0,0) {\includegraphics[page=524,trim={53.00bp 580.97bp 53.87bp 67.82bp},clip,width=0.92\textwidth,max height=0.38\textheight]{Ferziger4th.pdf}};
\begin{scope}[x={(image.south east)},y={(image.north west)}]
\fill[white] (0.00000,0.94986) rectangle (0.04547,1.00000);
\fill[white] (0.04704,0.94986) rectangle (0.09639,1.00000);
\fill[white] (0.10929,0.94467) rectangle (0.13251,1.00000);
\fill[white] (0.13408,0.94467) rectangle (0.25344,1.00000);
\fill[white] (0.81871,0.92795) rectangle (0.86737,1.00000);
\fill[white] (0.87474,0.92795) rectangle (0.95513,1.00000);
\fill[white] (0.00000,0.37061) rectangle (0.07627,0.95101);
\fill[white] (0.07786,0.37061) rectangle (0.15856,0.95101);
\fill[white] (0.16012,0.37061) rectangle (0.22998,0.95101);
\fill[white] (0.23158,0.37061) rectangle (0.27600,0.95101);
\fill[white] (0.86220,0.37637) rectangle (0.87907,0.93084);
\fill[white] (0.00000,0.00000) rectangle (0.09750,0.37695);
\fill[white] (0.09904,0.00000) rectangle (0.14911,0.37695);
\fill[white] (0.15068,0.00000) rectangle (0.20078,0.37695);
\fill[white] (0.20232,0.00000) rectangle (0.24391,0.37695);
\fill[white] (0.24550,0.00000) rectangle (0.27014,0.37695);
\fill[white] (0.27167,0.00000) rectangle (0.28779,0.37695);
\fill[white] (0.68752,0.06916) rectangle (0.71645,0.38963);
\end{scope}
\end{tikzpicture}
\end{sourceblock}
旧位置

\begin{sourceblock}
\begin{tikzpicture}
\node[inner sep=0,anchor=south west] (image) at (0,0) {\includegraphics[page=524,trim={53.00bp 551.00bp 53.87bp 95.72bp},clip,width=0.92\textwidth,max height=0.38\textheight]{Ferziger4th.pdf}};
\begin{scope}[x={(image.south east)},y={(image.north west)}]
\fill[white] (0.00000,0.84861) rectangle (0.10671,1.00000);
\fill[white] (0.10830,0.84861) rectangle (0.13293,1.00000);
\fill[white] (0.13447,0.84861) rectangle (0.17041,1.00000);
\fill[white] (0.17197,0.84861) rectangle (0.21780,1.00000);
\fill[white] (0.00000,0.33574) rectangle (0.10138,0.85427);
\fill[white] (0.10289,0.33574) rectangle (0.12608,0.85427);
\fill[white] (0.12764,0.33574) rectangle (0.15654,0.85427);
\fill[white] (0.15805,0.33574) rectangle (0.27465,0.85427);
\fill[white] (0.65540,0.50515) rectangle (0.68433,0.84037);
\fill[white] (0.74824,0.31874) rectangle (0.77047,0.81411);
\fill[white] (0.98030,0.44284) rectangle (0.99585,0.93821);
\fill[white] (0.68647,0.06179) rectangle (0.71540,0.39701);
\end{scope}
\end{tikzpicture}
\end{sourceblock}
\begin{sourceblock}
\begin{tikzpicture}
\node[inner sep=0,anchor=south west] (image) at (0,0) {\includegraphics[page=524,trim={53.00bp 510.12bp 53.87bp 106.23bp},clip,width=0.92\textwidth,max height=0.38\textheight]{Ferziger4th.pdf}};
\begin{scope}[x={(image.south east)},y={(image.north west)}]
\fill[white] (0.00000,0.95200) rectangle (0.10138,1.00000);
\fill[white] (0.10289,0.95200) rectangle (0.12608,1.00000);
\fill[white] (0.12764,0.95200) rectangle (0.15654,1.00000);
\fill[white] (0.15805,0.95200) rectangle (0.27465,1.00000);
\fill[white] (0.68647,0.84515) rectangle (0.71540,0.97590);
\fill[white] (0.82803,0.23177) rectangle (0.84761,0.42498);
\fill[white] (0.94923,0.23177) rectangle (0.96595,0.42498);
\fill[white] (0.85600,0.02410) rectangle (0.87558,0.21731);
\end{scope}
\end{tikzpicture}
\par\vspace{0.45\baselineskip}
\small 图 13.6 所示的 CV 的离散连续性方程采用隐式欧拉方案如下：
\end{sourceblock}
\begin{sourceblock}
\includegraphics[page=524,trim={53.00bp 410.93bp 53.87bp 209.22bp},clip,width=0.92\textwidth,max height=0.38\textheight]{Ferziger4th.pdf}
\end{sourceblock}
其中 u 和 v 是笛卡尔速度分量。因为我们假设流体以恒定速度运动，所以上式中流体速度的贡献被抵消，只剩下网格速度的差异：

\begin{sourceblock}
\includegraphics[page=524,trim={53.00bp 307.30bp 53.87bp 316.46bp},clip,width=0.92\textwidth,max height=0.38\textheight]{Ferziger4th.pdf}
\end{sourceblock}
在上述假设下，CV 两侧网格速度差可表示为（见图 13.6）：

\begin{sourceblock}
\includegraphics[page=524,trim={53.00bp 239.08bp 53.87bp 398.78bp},clip,width=0.92\textwidth,max height=0.38\textheight]{Ferziger4th.pdf}
\end{sourceblock}
将这些表达式代入方程( 13.21 )，并注意到 ( V ) n + 1 = ( x y ) n + 1 和 ( V ) n = [ ( x ) n + 1 − δ x ][ ( y ) n + 1 − δ y ]，我们发现不满足离散质量守恒定律——存在质量源

\begin{sourceblock}
\includegraphics[page=524,trim={53.00bp 158.09bp 53.87bp 479.77bp},clip,width=0.92\textwidth,max height=0.38\textheight]{Ferziger4th.pdf}
\end{sourceblock}
使用显式欧拉方案可以获得相同的误差（具有相反的符号）。对于恒定的网格速度，它与时间步长成正比，即，它是一阶离散化误差。人们可能会认为这不是问题，因为该方案在时间上只是一阶精确的；然而，人造质量源可能会随着时间的推移而积累并导致严重的问题。如果仅一组网格线移动，或者网格速度在相对的 CV 侧相等，则错误消失。

在上述假设下，Crank-Nicolson 格式和三时间级隐式格式都完全满足连续性方程。更一般地，当流体和/或网格速度不恒定时，这些方案还可以产生人造质量源。

质量守恒可以通过执行所谓的空间守恒定律（SCL）来获得，该定律可以被视为零流体速度极限下的连续性方程：

\begin{sourceblock}
\begin{tikzpicture}
\node[inner sep=0,anchor=south west] (image) at (0,0) {\includegraphics[page=525,trim={53.00bp 498.40bp 53.87bp 138.65bp},clip,width=0.92\textwidth,max height=0.38\textheight]{Ferziger4th.pdf}};
\begin{scope}[x={(image.south east)},y={(image.north west)}]
\fill[white] (0.00000,0.89687) rectangle (0.05907,1.00000);
\fill[white] (0.06177,0.89687) rectangle (0.17095,1.00000);
\fill[white] (0.32364,0.56136) rectangle (0.34343,0.95875);
\fill[white] (0.31811,0.07803) rectangle (0.34683,0.47920);
\end{scope}
\end{tikzpicture}
\end{sourceblock}
该方程描述了 CV 随时间改变其形状和/或位置时的空间守恒。

通过考虑恒定密度流体的质量守恒方程（13.18）可以看出为什么遵守 SCL 很重要；那么它可以写成： d d t

\begin{sourceblock}
\begin{tikzpicture}
\node[inner sep=0,anchor=south west] (image) at (0,0) {\includegraphics[page=525,trim={53.00bp 407.43bp 53.87bp 236.23bp},clip,width=0.92\textwidth,max height=0.38\textheight]{Ferziger4th.pdf}};
\begin{scope}[x={(image.south east)},y={(image.north west)}]
\fill[white] (0.24159,0.72642) rectangle (0.26138,1.00000);
\fill[white] (0.42598,0.38078) rectangle (0.45392,0.93461);
\fill[white] (0.45729,0.43194) rectangle (0.47044,0.94662);
\fill[white] (0.47227,0.42037) rectangle (0.49374,0.93461);
\fill[white] (0.49435,0.41548) rectangle (0.53125,0.93461);
\fill[white] (0.53447,0.43194) rectangle (0.56265,0.94662);
\fill[white] (0.23606,0.10098) rectangle (0.26481,0.61966);
\fill[white] (0.28770,0.05338) rectangle (0.30532,0.43461);
\end{scope}
\end{tikzpicture}
\end{sourceblock}
前两项代表 SCL，加起来为零，参见等式。 （13.24）；因此，对于密度恒定的流体，质量守恒方程简化为

\begin{sourceblock}
\includegraphics[page=525,trim={53.00bp 338.22bp 53.87bp 305.44bp},clip,width=0.92\textwidth,max height=0.38\textheight]{Ferziger4th.pdf}
\end{sourceblock}
因此，重要的是确保上述两项在离散方程中也被抵消（即，由于其运动而通过 CV 面的体积通量之和必须等于体积变化率）；否则，引入人工质量源，它们可能会随着时间的推移而积累并破坏解，如 Demirdži ́c 和 Peri ́c (1988) 所证明的那样。

下面使用隐式三时间级时间积分方案和SIMPLE算法进行说明。 Crank-Nicholson 格式以及隐式 Euler 格式的表达式很容易导出；隐式分数步法的实现也很简单。在具有移动边界的瞬态流的情况下，使用一阶方案进行时间积分才有意义，只有当我们知道流正在朝着稳态解发展时，就像某些浮体问题的情况一样（船舶在平静的水中移动；初始位置是船舶和流体都静止的位置，但在恒定速度下，船舶的纵倾和下沉都会由于产生的波浪而改变）。对于空间积分，我们使用中点规则和中心差方案。

离散化的SCL方程可以转化为以下形式（见式（6.25））：

\begin{sourceblock}
\begin{tikzpicture}
\node[inner sep=0,anchor=south west] (image) at (0,0) {\includegraphics[page=525,trim={53.00bp 92.85bp 53.87bp 530.84bp},clip,width=0.92\textwidth,max height=0.38\textheight]{Ferziger4th.pdf}};
\begin{scope}[x={(image.south east)},y={(image.north west)}]
\fill[white] (0.53444,0.29918) rectangle (0.56262,0.57173);
\fill[white] (0.28674,0.12391) rectangle (0.30653,0.39623);
\fill[white] (0.33095,0.12627) rectangle (0.34409,0.39882);
\end{scope}
\end{tikzpicture}
\end{sourceblock}
\begin{sourceblock}
\begin{tikzpicture}
\node[inner sep=0,anchor=south west] (image) at (0,0) {\includegraphics[page=526,trim={53.00bp 560.13bp 53.87bp 58.64bp},clip,width=0.92\textwidth,max height=0.38\textheight]{Ferziger4th.pdf}};
\begin{scope}[x={(image.south east)},y={(image.north west)}]
\fill[white] (0.62171,0.71227) rectangle (0.63558,0.97467);
\fill[white] (0.63889,0.71227) rectangle (0.66568,0.97467);
\fill[white] (0.66902,0.71227) rectangle (0.68289,0.97467);
\fill[white] (0.51074,0.17606) rectangle (0.52460,0.43825);
\fill[white] (0.87817,0.02533) rectangle (0.89639,0.27148);
\end{scope}
\end{tikzpicture}
\end{sourceblock}
t

n −1
其中求和是针对 CV 的所有面进行的。请注意，连续时间级别的 CV 体积之间的差异可以表示为每个 CV 面从旧位置移动到新位置时扫过的体积 δ V k 之和，见图 13.7，即：

\begin{sourceblock}
\begin{tikzpicture}
\node[inner sep=0,anchor=south west] (image) at (0,0) {\includegraphics[page=526,trim={53.00bp 373.15bp 53.87bp 263.40bp},clip,width=0.92\textwidth,max height=0.38\textheight]{Ferziger4th.pdf}};
\begin{scope}[x={(image.south east)},y={(image.north west)}]
\fill[white] (0.00000,0.82122) rectangle (0.04908,1.00000);
\fill[white] (0.30586,0.42987) rectangle (0.32186,0.82055);
\fill[white] (0.34198,0.42075) rectangle (0.36508,0.81142);
\fill[white] (0.36481,0.42987) rectangle (0.41907,0.87868);
\fill[white] (0.42244,0.42987) rectangle (0.46845,0.82055);
\fill[white] (0.48857,0.42075) rectangle (0.51170,0.81142);
\fill[white] (0.51146,0.42987) rectangle (0.53792,0.87259);
\fill[white] (0.54388,0.42987) rectangle (0.57206,0.82055);
\end{scope}
\end{tikzpicture}
\end{sourceblock}
当这个表达式被引入等式中时。 ( 13.27 ), 得到如下表达式:

\begin{sourceblock}
\begin{tikzpicture}
\node[inner sep=0,anchor=south west] (image) at (0,0) {\includegraphics[page=526,trim={53.00bp 308.16bp 53.87bp 315.53bp},clip,width=0.92\textwidth,max height=0.38\textheight]{Ferziger4th.pdf}};
\begin{scope}[x={(image.south east)},y={(image.north west)}]
\fill[white] (0.00000,0.77032) rectangle (0.11561,1.00000);
\fill[white] (0.31880,0.40801) rectangle (0.33290,0.61013);
\fill[white] (0.34376,0.46525) rectangle (0.37194,0.73757);
\end{scope}
\end{tikzpicture}
\end{sourceblock}
尽管在不同的情况下可以满足上述方程，但可以合理地假设左侧和右侧总和的相应部分应该相等（即方程两侧每个面的贡献相等）。在这一假设下，如果通过单元面的体积通量定义为：

\begin{sourceblock}
\includegraphics[page=526,trim={53.00bp 200.72bp 53.87bp 435.13bp},clip,width=0.92\textwidth,max height=0.38\textheight]{Ferziger4th.pdf}
\end{sourceblock}
因此，每个面在一个时间步内扫过的体积 δ V k 是根据两个时间级别的网格位置计算出来的，并用于计算体积通量 ̇ V n + 1
k ;这样就不需要明确定义 CV 表面的速度 v s。

通过一个单元面的质量通量现在可以计算为（参见方程（13.18））：

\begin{sourceblock}
\begin{tikzpicture}
\node[inner sep=0,anchor=south west] (image) at (0,0) {\includegraphics[page=526,trim={53.00bp 104.85bp 53.87bp 524.24bp},clip,width=0.92\textwidth,max height=0.38\textheight]{Ferziger4th.pdf}};
\begin{scope}[x={(image.south east)},y={(image.north west)}]
\fill[white] (0.01438,0.27099) rectangle (0.08087,0.64049);
\fill[white] (0.03780,0.21269) rectangle (0.05191,0.44426);
\fill[white] (0.08589,0.27827) rectangle (0.11408,0.59028);
\end{scope}
\end{tikzpicture}
\end{sourceblock}
这里 (vi) k 和 (Si) k 代表流体速度矢量 v 和面 k 处的表面矢量 S n 的笛卡尔分量。

\begin{center}\small 图 13.8 3D CV 单元面扫过的体积计算；阴影部分是相邻 CV 共有的表面\end{center}
New

position

S f
k

position f S
Old

k

每次 SIMPLE 迭代应满足（在一定公差范围内）的离散质量守恒方程如下：

\begin{sourceblock}
\begin{tikzpicture}
\node[inner sep=0,anchor=south west] (image) at (0,0) {\includegraphics[page=527,trim={53.00bp 369.08bp 53.87bp 260.99bp},clip,width=0.92\textwidth,max height=0.38\textheight]{Ferziger4th.pdf}};
\begin{scope}[x={(image.south east)},y={(image.north west)}]
\fill[white] (0.10195,0.53590) rectangle (0.12177,0.85639);
\fill[white] (0.12192,0.54588) rectangle (0.15423,0.86665);
\fill[white] (0.18117,0.53867) rectangle (0.20430,0.85916);
\fill[white] (0.20403,0.54588) rectangle (0.25826,0.90047);
\fill[white] (0.26165,0.53590) rectangle (0.31146,0.86665);
\fill[white] (0.31170,0.54588) rectangle (0.34400,0.86665);
\fill[white] (0.37092,0.53867) rectangle (0.39405,0.85916);
\fill[white] (0.39380,0.54588) rectangle (0.42027,0.89520);
\fill[white] (0.42457,0.54588) rectangle (0.45275,0.86665);
\fill[white] (0.45459,0.54588) rectangle (0.48692,0.86665);
\fill[white] (0.51380,0.53867) rectangle (0.53693,0.85916);
\fill[white] (0.53669,0.54588) rectangle (0.59092,0.90047);
\end{scope}
\end{tikzpicture}
\end{sourceblock}
在瞬态 SIMPLE 算法中，通过外部迭代接近新时间级别 t n + 1 的值。根据方程（1），在每次外部迭代时计算新质量通量的近似值。 ( 13.31 )，使用通过求解动量方程获得的普遍密度和速度。然后，通过对单元面速度进行校正（与压力校正的梯度成正比），并且在可压缩流的情况下对单元面密度进行校正（与压力校正成正比），对质量通量进行校正以满足质量守恒方程。这些更正遵循下列各节中描述的程序。 7.2.2和11.2.1，此处不再重复。

在 3D 中，计算单元面扫过的体积时必须小心。由于单元面边缘可能会转动，因此计算扫掠体积需要对图 13.8 中的阴影表面进行三角测量。然后可以使用第 1 节中描述的方法计算体积。 9.6.4.然而，由于阴影表面是两个 CV 所共有的，因此必须确保两个 CV 以相同的方式对它们进行三角测量，以确保节省空间。

由质量守恒方程得出的压力修正方程与固定网格的情况（对于可压缩流和不可压缩流）具有相同的形式，除了时间相关项之外。如果新时间水平的网格位置已知（例如，在活塞驱动的流动、围绕旋转机械的流动等），则通过单元面的体积通量 ̇ V n + 1
k 不依赖于外部迭代计数器，并且在每个新时间步开始时计算一次。然而，在网格适应固体结构的位置（流固相互作用）或界面（例如，在自由表面的界面跟踪处理中）的情况下，需要在外部迭代期间与其他校正一起校正体积通量；有关更多详细信息和示例，请参阅 Demirdži ́c 和 Peri ́c (1990)。

在可压缩流的情况下，解域内的质量变化率不为零，因此无法正式证明 SCL 的特定实现保证守恒，正如上面针对不可压缩流所描述的那样。原因是方程（13.32）左侧出现的密度是单元中心值，而右侧（在单元面质量通量中）使用的是密度的面值。然而，三时间级别方案在时间上是二阶的，因此当时间步长减半时，任何时间离散化误差（影响变化率）都会减少四倍。由网格运动引起的附加误差与固定网格的离散化误差具有相同的量级，并且当时间步细化时以相同的速率减少。

可以证明，对于具有不同形状和解域体积的封闭系统（没有入口和出口）中的可压缩流，当使用上述 SCL 离散化时，质量守恒。下面针对三时间级别方案对此进行了演示。

根据离散质量守恒方程 ( 13.32 )，通过对所有 CV 求和即可获得解域中的总质量。通过所有内部单元面的质量通量在该总和中被抵消；只剩下通过边界面的质量通量之和（请注意，质量通量包含流体速度和边界运动的贡献，请参见方程（13.31）；因此，在不渗透壁处，质量通量为零，无论其是否移动）：

\begin{sourceblock}
\includegraphics[page=528,trim={53.00bp 307.22bp 53.87bp 323.97bp},clip,width=0.92\textwidth,max height=0.38\textheight]{Ferziger4th.pdf}
\end{sourceblock}
其中 ̇ m B 是通过边界单元面的质量通量，

\begin{sourceblock}
\begin{tikzpicture}
\node[inner sep=0,anchor=south west] (image) at (0,0) {\includegraphics[page=528,trim={53.00bp 260.23bp 53.87bp 387.45bp},clip,width=0.92\textwidth,max height=0.38\textheight]{Ferziger4th.pdf}};
\begin{scope}[x={(image.south east)},y={(image.north west)}]
\fill[white] (0.37465,0.07096) rectangle (0.41750,0.79523);
\fill[white] (0.42343,0.08505) rectangle (0.45161,0.71181);
\fill[white] (0.52223,0.08505) rectangle (0.55456,0.71181);
\fill[white] (0.57645,0.07096) rectangle (0.59955,0.69718);
\fill[white] (0.59931,0.08505) rectangle (0.62577,0.79523);
\fill[white] (0.90935,0.06501) rectangle (1.00000,0.69177);
\end{scope}
\end{tikzpicture}
\end{sourceblock}
所有细胞

表示时间 t n 时解域中的质量。

在没有流入或流出的封闭系统的情况下，通过边界单元面的质量通量之和等于零。从方程。 ( 13.33 ) 则如下：

\begin{sourceblock}
\includegraphics[page=528,trim={53.00bp 166.40bp 53.87bp 467.73bp},clip,width=0.92\textwidth,max height=0.38\textheight]{Ferziger4th.pdf}
\end{sourceblock}
如果在模拟开始时，质量在两个连续时间水平上相同，即，如果 M n = M n − 1 = M 0，则从上述方程也可以得出 M n + 1 = M 0，即，无论网格运动如何，质量都将在解域内守恒。

在一些隐式时间积分方法（所谓的完全隐式方法）中，仅在最新时间级别计算通量和源项，除了边界附近之外，在任何地方都可以忽略网格运动。此类方法的示例是隐式欧拉方案和三时间级别方案，请参阅第 1 章。 6.因为通量是在时间级别 t n + 1 计算的，所以我们不需要知道前一个时间级别 t n 的网格在哪里（或者 CV 的形状）：我们可以使用空间固定 CV 的常用方程来代替方程（13.19）：

\begin{sourceblock}
\begin{tikzpicture}
\node[inner sep=0,anchor=south west] (image) at (0,0) {\includegraphics[page=529,trim={53.00bp 522.23bp 53.87bp 114.45bp},clip,width=0.92\textwidth,max height=0.38\textheight]{Ferziger4th.pdf}};
\begin{scope}[x={(image.south east)},y={(image.north west)}]
\fill[white] (0.04568,0.56687) rectangle (0.06421,0.95927);
\fill[white] (0.11480,0.29294) rectangle (0.15997,0.72234);
\fill[white] (0.16370,0.31738) rectangle (0.20325,0.71317);
\fill[white] (0.20968,0.32994) rectangle (0.23786,0.72234);
\fill[white] (0.04015,0.08079) rectangle (0.06941,0.48235);
\fill[white] (0.09230,0.04073) rectangle (0.10992,0.33164);
\end{scope}
\end{tikzpicture}
\end{sourceblock}
这两个方程在变化率和对流项的定义上有所不同：对于空间固定的 CV，对流通量是使用流体速度计算的，时间导数表示空间中固定点（例如 CV 中心）的局部变化率。另一方面，对于移动的 CV，使用流体和 CV 表面之间的相对速度来计算对流通量，时间导数表示位置变化的体积的变化率。

由于计算表面和体积积分不需要先前时间步长的解，因此网格不仅可以移动，还可以改变其拓扑结构，即 CV 的数量及其形状都可以从一个时间步长变化到另一个时间步长。旧解中出现的唯一项是不稳定项，它要求必须对某些旧量的新 CV 上的体积积分进行近似。如果为此目的使用中点规则，我们所需要做的就是将旧解插值到新 CV 中心的位置。一种可能性是计算每个旧 CV 中心的梯度向量，然后对于每个新 CV 中心，找到旧 CV 的最近中心并使用线性插值来获取新 CV 中心的旧值：

\begin{sourceblock}
\includegraphics[page=529,trim={53.00bp 296.24bp 53.87bp 351.78bp},clip,width=0.92\textwidth,max height=0.38\textheight]{Ferziger4th.pdf}
\end{sourceblock}
在移动边界附近，我们必须考虑这样一个事实：边界在时间步长期间发生移动，或者取代了流体，或者腾出了空间来填充流体。对于小运动，可以通过在近边界 CV 中规定质量源或汇来考虑这一点；这些的计算方式与使用边界移动的速度作为流体速度计算通过界面的入口或出口质量通量的方式相同。如果 CV 在一个时间步长内沿运动方向移动超过其宽度，则会出现问题，因为新 CV 的中心可能位于旧网格之外。因此，对于靠近移动壁的精细网格，可能需要使用基于近壁区域中的移动控制体积的移动网格和方程，而远离壁的网格运动可以被忽略，从而允许在其属性由于过度变形而恶化时重新生成网格。上述方法的一个例子可以在 Hadži ́c (2005) 的论文中找到；他还对该方法与使用移动控制体积方程的方法进行了比较，发现两种方法的结果与突然膨胀的管道中活塞驱动流动的实验数据之间具有良好的一致性。

许多工程应用需要使用移动网格。然而，不同的问题需要不同的解决方法。一个重要的例子是转子-定子相互作用，这在涡轮机械和混合器中很常见：栅格的一部分连接到定子并且不移动，而另一部分连接到转子并随之移动。移动网格和固定网格之间的界面通常是圆柱面和平面的组合。如果初始时网格在界面处匹配，则可以允许网格的旋转部分移动，同时保持边界点“粘合”到固定网格，直到变形变得相当大（不应超过45°角）；然后，边界点向前“跳跃”一个单元格，并暂时粘在新位置上。这种“点击”网格已在这些应用中与常规的结构化网格结合使用。

另一种可能性是让移动网格沿着界面“滑动”而不变形。在这种情况下，网格在界面处不匹配，因此某些 CV 比其他 CV 具有更多的邻居。然而，这种情况完全类似于具有非共形接口的块结构网格中遇到的情况，并且可以通过第 1 节中描述的方法来处理。 9.6.1；唯一的区别是单元连接随着时间而变化，并且必须在每个时间步之后重新建立。这种方法比上述方法更灵活；网格可以是不同种类和/或细度的，并且界面可以是任意表面。这种方法还可以应用于物体周围相互经过、进入隧道或在已知轨迹的封闭空间中移动的流动。所有商业代码都提供此功能，用于模拟涉及多种涡轮机中的转子-定子相互作用、船舶上的螺旋桨旋转等的流动。

第三种方法是使用重叠（Chimera）网格。同样，一个网格连接到域的固定部分，另一个网格连接到移动体。即使运动体的轨迹事先未知，当运动体的轨迹非常复杂时，或者当周围域的形状复杂时（例如，当运动部件的路径相交而无法构建滑动界面时），也可以使用这种方法。固定网格可以覆盖身体运动的整个“环境”。重叠区域随时间变化，并且在每个时间步之后需要重新建立网格之间的关系。除了难以确保精确守恒（如第 9.1.3 节中所讨论）之外，这种方法的适用性几乎没有任何限制（甚至可以解释物体相互接近的接触）。

如上所述，相同的方程和离散化方法适用于固定网格和移动网格，唯一的区别是在固定网格上，网格速度 v s 显然为零。有时在两个域上使用不同的坐标系可能是有利的；例如，可以在一个部分使用笛卡尔速度分量，在另一网格上使用极坐标分量。这是可能的，前提是：（i）由于框架加速度而添加体积力，并且（ii）在界面或重叠区域将矢量分量从一个系统变换到另一个系统。原则上这两者都很容易做到，但编程可能很乏味。

\begin{center}\small 图 13.9 显示了使用重叠网格将解耦合到附在船体上的固定网格和随螺旋桨旋转的网格上的示例。该船配备了 POD 驱动器，可以绕垂直轴旋转；因此可以改变推力方向，并且船舶不需要舵。在这个\end{center}
\begin{sourceblock}
\includegraphics[page=531,trim={53.00bp 438.25bp 53.87bp 52.00bp},clip,width=\textwidth,max height=0.38\textheight]{Ferziger4th.pdf}
\par\vspace{0.45\baselineskip}
\small 图 13.9 配备所谓 POD 驱动器的全尺寸船舶螺旋桨的纵向剖面，显示连接到螺旋桨的旋转网格和连接到船体的固定网格的重叠
\end{sourceblock}
研究 POD 是固定的，只有螺旋桨以固定速率旋转；船速不断变化，直到阻力与螺旋桨推力相匹配。在此示例中，插入螺旋桨的区域中的固定网格的粗度是旋转网格外部部分的两倍；单元尺寸较大的不匹配是不希望的，而两个网格中的单元尺寸几乎相等将是最佳的。由于空间限制，我们无法讨论特定研究的细节，但除了在图 13.9 中显示两个部分中的重叠网格之外，我们还在图 13.10 中显示速度和压力等值线。这里我们要注意轮廓线：在两个网格的活动单元重叠的区域中，每个轮廓被表示两次：一次在固定网格的单元中，一次在旋转网格的单元中。不可避免地，线条不会完全相同，但如果一切都正确完成，则错误匹配将会很小。在这种特殊情况下，压力和速度轮廓都匹配得很好。有关此模拟的更详细报告可以从本书网站下载；详细内容请参见附录。

一种更简单的方法是使用笛卡尔速度分量，但在固定区域和移动区域使用不同的参考系；当在随零件移动的参考系中观察时，如果流动（例如，旋转）区域中的流动稳定，则通常会这样做。然后，需要通过参考系运动产生的额外项来扩展移动区域中的方程。这种方法无法考虑移动部件和固定部件之间的相互作用（例如涡轮机械中的叶片通过效应），因此引入了额外的建模误差。如果界面附近的两个参考系中的流动稳定且轴对称，则该误差很小，但如果移动部件和固定部件之间的间隙很小，则该误差可能很大。过去，此类方法通常与涡粘性湍流模型和稳态分析结合使用，以节省计算时间；然而，由于移动部件和固定部件之间的相互作用通常是分析的一个重要元素，因此可以预期，未来大多数此类应用将使用移动网格。

\begin{sourceblock}
\includegraphics[page=532,trim={53.00bp 274.25bp 53.87bp 52.00bp},clip,width=\textwidth,max height=0.38\textheight]{Ferziger4th.pdf}
\par\vspace{0.45\baselineskip}
\small 图 13.10 螺旋桨纵向剖面中的压力（上）和轴向速度（下）等值线，使用图 13.9 中所示的重叠网格计算
\end{sourceblock}
\section*{13.5 自由表面流}\phantomsection\addcontentsline{toc}{section}{13.5 自由表面流}
具有自由表面的流动是具有移动边界的特别困难的一类流动。边界的位置仅在初始时刻已知；稍后必须确定其位置作为解的一部分。为了实现这一点，必须使用 SCL 和自由表面的边界条件。

在最常见的情况下，自由表面是空气-水边界，但也会出现其他液-气表面，液-液界面也是如此。如果可以忽略自由表面的相变，则适用以下边界条件：

\begin{itemize}
\item 运动学条件要求自由表面是一个清晰的界面，将两种流体分开，不允许流过它，即：
\end{itemize}
\begin{sourceblock}
\includegraphics[page=533,trim={53.00bp 593.89bp 53.87bp 56.79bp},clip,width=0.92\textwidth,max height=0.38\textheight]{Ferziger4th.pdf}
\end{sourceblock}
其中“fs”表示自由表面。这表明表面流体速度的法向分量是自由表面速度的法向分量，见式（13.18）。 • 动态条件要求作用在自由表面流体上的力处于平衡状态（自由表面动量守恒）。这意味着自由表面两侧的法向力大小相等且方向相反，而切向方向上的力大小相等且方向相反：

\begin{sourceblock}
\includegraphics[page=533,trim={53.00bp 393.09bp 53.87bp 188.30bp},clip,width=0.92\textwidth,max height=0.38\textheight]{Ferziger4th.pdf}
\end{sourceblock}
其中 σ 为表面张力，n、t 和 s 为自由表面局部正交坐标系 ( n , t , s ) 中的单位向量（ n 为从液体侧看自由表面的外法线，另外两个位于切平面内且相互正交）。下标‘l’和‘g’分别表示液体和气体，K是自由表面的曲率，

\begin{sourceblock}
\includegraphics[page=533,trim={53.00bp 289.74bp 53.87bp 347.66bp},clip,width=0.92\textwidth,max height=0.38\textheight]{Ferziger4th.pdf}
\end{sourceblock}
其中R t 和R s 是沿坐标t 和s 的曲率半径，见图13.11。表面张力 σ 是表面元素每单位长度的力，作用于自由表面的切向；在图13.11中，表面张力产生的力f σ 的大小为f σ = σ d l。对于无限小的表面单元 d S，当 σ = const. 时，表面张力的切向分量相互抵消，并且法向分量可以表示为导致整个表面上的压力跃变的局部力，如方程 1 所示。 （13.39）。

表面张力是液体的热力学性质，取决于温度和其他状态变量，例如化学成分和表面清洁度。如果温差很小，则可以将 σ 的温度依赖性线性化，使得 ∂σ/∂ T 为常数；它通常是负数。当温度沿自由表面大幅变化时，表面张力梯度会产生剪切力，导致流体从热区域移动到冷区域。这种现象称为马兰戈尼或毛细管对流，其重要性由无量纲马兰戈尼数来表征：

\begin{center}\small 图13.11 关于自由表面边界条件的描述\end{center}
\begin{sourceblock}
\begin{tikzpicture}
\node[inner sep=0,anchor=south west] (image) at (0,0) {\includegraphics[page=534,trim={53.00bp 553.40bp 53.87bp 97.87bp},clip,width=0.92\textwidth,max height=0.38\textheight]{Ferziger4th.pdf}};
\begin{scope}[x={(image.south east)},y={(image.north west)}]
\fill[white] (0.52036,0.08070) rectangle (0.53669,0.61601);
\end{scope}
\end{tikzpicture}
\end{sourceblock}
\begin{sourceblock}
\begin{tikzpicture}
\node[inner sep=0,anchor=south west] (image) at (0,0) {\includegraphics[page=534,trim={53.00bp 486.94bp 53.87bp 164.33bp},clip,width=0.92\textwidth,max height=0.38\textheight]{Ferziger4th.pdf}};
\begin{scope}[x={(image.south east)},y={(image.north west)}]
\fill[white] (0.59489,0.38399) rectangle (0.61122,0.91930);
\fill[white] (0.93928,0.08070) rectangle (0.95564,0.61601);
\end{scope}
\end{tikzpicture}
\end{sourceblock}
\begin{sourceblock}
\includegraphics[page=534,trim={53.00bp 439.44bp 53.87bp 198.41bp},clip,width=0.92\textwidth,max height=0.38\textheight]{Ferziger4th.pdf}
\end{sourceblock}
其中 T 是域内的整体温差，L 是表面的特征长度，κ 是热扩散率。

在许多应用中，自由表面的剪切应力可以忽略不计（例如，没有明显风的海浪或大型船舶产生的波浪）。法向应力和表面张力的影响也常常被忽略；在这种情况下，动态边界条件简化为 p l = p g。

这些边界条件的实现并不像看起来那么简单。如果自由表面的位置已知，就不会有什么问题了。可以将位于自由表面上的单元面的质量通量设置为零，并且可以计算从外部作用在单元面上的力；如果忽略表面张力，则只剩下压力。问题是自由表面的位置必须作为解的一部分进行计算，并且通常无法提前得知。因此，可以直接在自由表面上仅实现其中一个边界条件；另一个必须用于定位表面。自由表面的定位必须迭代完成，这大大增加了任务的复杂性。通常只有界面一侧的流动（通常是液体流动）是令人感兴趣的；然而，在许多应用中，必须同时计算界面两侧的两种流体的流动（例如，在液体中移动的气泡的流动或在气体中移动的液滴的流动）。

许多方法已被用来寻找自由表面的形状。它们可以分为两大类：

\begin{itemize}
\item 将自由表面视为其运动被跟随的尖锐界面的方法（界面跟踪方法）。在这种类型的方法中，每次移动自由表面时，都会使用边界拟合网格并前进。在必须使用小时间步长的显式方法中，前面讨论的与网格移动相关的问题经常被忽略。作为一个例子，Hodges 和 Street (1999) 提出了一种 FV 分数阶方法，该方法具有用于自由表面流 LES 的移动边界正交网格。
\end{itemize}
\begin{itemize}
\item 不将接口定义为边界的方法（接口捕获方法）。计算在固定网格上进行，该网格覆盖自由表面两侧的流体。自由表面的形状是通过计算每个近界面单元被部分填充的比例来确定的。这可以通过在初始时间在自由表面引入无质量粒子并跟随它们的运动来实现；这称为 Marker-and-Cell 或 MAC 方案，由 Harlow 和 Welsh (1965) 提出。或者，可以求解液相占据的细胞部分的传输方程（流体体积或 VOF 方案，Hirt 和 Nichols 1981）。
\end{itemize}
还有混合方法。所有这些方法也可以应用于某些类型的两相流，如下一节将讨论。

\subsection*{13.5.1 接口捕获方法}\phantomsection\addcontentsline{toc}{subsection}{13.5.1 接口捕获方法}
MAC 方案很有吸引力，因为它可以处理波浪破碎等复杂现象。然而，计算工作量很大，尤其是在三维空间中，因为除了求解控制流体流动的方程之外，还必须跟踪大量标记颗粒的运动。

在 VOF 方法中，除了质量和动量守恒方程外，还求解每个控制体积的填充分数 c 的方程，使得填充 CV 中 c = 1，空 CV 中 c = 0。从连续性方程可以看出，c 的演化受输运方程控制：

\begin{sourceblock}
\begin{tikzpicture}
\node[inner sep=0,anchor=south west] (image) at (0,0) {\includegraphics[page=535,trim={53.00bp 277.15bp 53.87bp 359.53bp},clip,width=0.92\textwidth,max height=0.38\textheight]{Ferziger4th.pdf}};
\begin{scope}[x={(image.south east)},y={(image.north west)}]
\fill[white] (0.11528,0.55771) rectangle (0.14950,0.95927);
\fill[white] (0.41985,0.56687) rectangle (0.43838,0.95927);
\fill[white] (0.15555,0.32994) rectangle (0.18373,0.72234);
\fill[white] (0.18559,0.32994) rectangle (0.22944,0.72234);
\fill[white] (0.23131,0.32077) rectangle (0.28779,0.72234);
\fill[white] (0.29131,0.32994) rectangle (0.31949,0.72234);
\fill[white] (0.32301,0.31738) rectangle (0.34280,0.70978);
\fill[white] (0.36150,0.31738) rectangle (0.39125,0.70978);
\fill[white] (0.11669,0.08079) rectangle (0.14595,0.48235);
\fill[white] (0.41435,0.08079) rectangle (0.44361,0.48235);
\end{scope}
\end{tikzpicture}
\end{sourceblock}
在不可压缩流中，该方程对于 c 和 1 − c 的互换是不变的；为了在数值方法中保证这一点，必须严格执行质量守恒定律。

这种方法比 MAC 方案更有效，并且可以应用于包括碎波在内的复杂自由表面形状。然而，自由表面轮廓并没有明确定义；它通常涂抹在一个或多个单元上（类似于可压缩流中的冲击）。局部网格细化对于自由表面的精确解析非常重要。细化标准很简单：0 < c < 1 的单元格需要细化。 Raad 和他的同事开发了一种称为标记和微细胞方法的此类方法（例如参见 Chen 等人，1997）。

上述方法有多种变体。在最初的 VOF 方法（Hirt 和 Nichols 1981）中在整个域中求解（ 13.42 ）以找到自由表面的位置；仅针对液相求解质量和动量守恒方程。该方法可以计算具有翻转自由表面的流量，但被液相包围的气体不会感受到浮力效应，因此会表现出不切实际的方式。

Kawamura 和 Miyata (1994) 使用方程。 ( 13.42 ) 计算密度函数的分布（密度和体积分数 c 的乘积）并定位自由表面，即 c = 0 的轮廓。 5. 液体和气体流动的运动的计算是分开进行的。自由表面被视为应用运动学和动态边界条件的边界。由于被自由表面切割而变得不规则的细胞需要特殊处理（变量值外推到位于界面另一侧的节点位置）。该方法用于计算船舶和水下物体周围的流量。

或者，可以将界面两侧的流体视为单一流体，其属性根据每一相的体积分数在空间中变化，即：

\begin{sourceblock}
\includegraphics[page=536,trim={53.00bp 450.77bp 53.87bp 200.25bp},clip,width=0.92\textwidth,max height=0.38\textheight]{Ferziger4th.pdf}
\end{sourceblock}
其中下标 1 和 2 表示两种流体（例如液体和气体）。在这种情况下，界面不被视为边界，因此不需要在其上指定边界条件。界面只是流体性质突然变化的位置。然而，方程的解。 ( 13.42 ) 意味着满足运动学条件，并且也隐式考虑了动态条件。如果自由表面的表面张力很大，则其影响可以解释为体积力。

表面张力仅作用在界面区域，即部分填充的单元中。因为在满细胞或空细胞中 c 的梯度为零，表面张力的法向分量可以表示为（连续表面力方法，Brackbill 等人，1992）：

\begin{sourceblock}
\begin{tikzpicture}
\node[inner sep=0,anchor=south west] (image) at (0,0) {\includegraphics[page=536,trim={53.00bp 295.57bp 53.87bp 348.10bp},clip,width=0.92\textwidth,max height=0.38\textheight]{Ferziger4th.pdf}};
\begin{scope}[x={(image.south east)},y={(image.north west)}]
\fill[white] (0.36818,0.38095) rectangle (0.40644,0.93502);
\fill[white] (0.41143,0.43213) rectangle (0.43961,0.94660);
\end{scope}
\end{tikzpicture}
\end{sourceblock}
然而，当表面张力效应占主导地位时，例如在直径为 1 毫米左右且以非常低的速度移动的液滴或气泡的情况下，就会出现问题。在这种情况下，动量方程中有两个非常大的项（压力项和代表表面张力效应的体积力），它们必须相互平衡；如果气泡或液滴是静止的，它们是唯一的非零项。由于界面的曲率也取决于 c（注意界面处 c 的梯度指向界面的法线方向），

\begin{sourceblock}
\begin{tikzpicture}
\node[inner sep=0,anchor=south west] (image) at (0,0) {\includegraphics[page=536,trim={53.00bp 167.27bp 53.87bp 463.59bp},clip,width=0.92\textwidth,max height=0.38\textheight]{Ferziger4th.pdf}};
\begin{scope}[x={(image.south east)},y={(image.north west)}]
\fill[white] (0.00000,0.71967) rectangle (0.08899,1.00000);
\fill[white] (0.09173,0.71967) rectangle (0.11985,1.00000);
\fill[white] (0.12250,0.71967) rectangle (0.24926,1.00000);
\fill[white] (0.32051,0.24206) rectangle (0.34006,0.56973);
\fill[white] (0.34598,0.24206) rectangle (0.37417,0.56973);
\fill[white] (0.37768,0.24206) rectangle (0.42156,0.56973);
\fill[white] (0.42340,0.23441) rectangle (0.44487,0.56207);
\fill[white] (0.44878,0.24206) rectangle (0.47696,0.56973);
\fill[white] (0.48048,0.24206) rectangle (0.54773,0.56973);
\end{scope}
\end{tikzpicture}
\end{sourceblock}
很难确保在任意 3D 网格上这两项相同（一项在 p 中是线性的，另一项在 c 中是非线性的），因此它们的差异可能会导致所谓的寄生电流。这些可以通过使用特殊的二维离散化方法来避免（参见 Scardovelli 和 Zaleski 1999，了解此类特殊方法的一些示例；另请参见 Harvie 等人 2006 的分析）。目前我们还不知道有一种方法可以完全消除 3D 中非结构化任意网格的问题。

\begin{sourceblock}
\includegraphics[page=537,trim={53.00bp 536.25bp 53.87bp 52.00bp},clip,width=\textwidth,max height=0.38\textheight]{Ferziger4th.pdf}
\par\vspace{0.45\baselineskip}
\small 图 13.12 两个连续时间步长的自由表面相对于单元面 k 的可能方向
\end{sourceblock}
此类方法的关键问题是方程（1）中对流项的离散化。 （13.42）。低阶方案（如一阶精确迎风法）会涂抹界面并引入两种流体的人工混合，因此优选高阶方案。因为c必须满足条件

\begin{sourceblock}
\includegraphics[page=537,trim={53.00bp 421.02bp 53.87bp 230.79bp},clip,width=0.92\textwidth,max height=0.38\textheight]{Ferziger4th.pdf}
\end{sourceblock}
重要的是要确保该方法不会产生过冲或下冲。幸运的是，我们可以推导出既保持界面清晰又在界面上产生单调分布的方案；有关单调方案的一些示例，请参阅 Leonard (1997) 和 Lafaurie 等人。 ( 1994 )、Ubbink ( 1997 ) 或 Muzaferija 和 Peri ́c ( 1999 ) 专门为自由表面流中的界面捕获而设计的方法。在过去的二十年中，已经发布了更多界面捕获或混合方法的变体；我们将在下面进一步提及其中一些。我们在此仅对 Muzaferija 和 Peri ́c (1999) 的 HRIC 方案（高分辨率接口捕获）进行简要说明。单元面 k 处的 c 对流值表示为迎风值和顺风值的混合：

\begin{sourceblock}
\includegraphics[page=537,trim={53.00bp 264.21bp 53.87bp 386.20bp},clip,width=0.92\textwidth,max height=0.38\textheight]{Ferziger4th.pdf}
\end{sourceblock}
其中下标U和D表示k面上游侧和下游侧的单元中心，见图13.12。如果 γ k = 0，则获得二阶中心差分格式。 5；否则，该方案在形式上是一阶精确的。然而，这里的目的不是对平滑变化函数的单元面中心进行精确插值（如计算速度或其他标量变量的对流值时的情况），而是保留两种不混溶流体之间的清晰界面；因此，近似阶次并不像其他传输方程那么重要。在 HRIC 方案中，混合因子 γ k 被确定为解的三个属性的函数：

\begin{itemize}
\item 跨接口的c 变化，基于两个上游和一个下游小区中心的c 值；
\item 本地CFL 数，估计界面在一个时间步长内将移动多远；
\item 界面法线（由c 的梯度定义）和细胞面法线之间的角度θ（见图13.12）。
\end{itemize}
接口的位置没有显式计算，但假设值 c = 0.5代表自由表面。如果界面与单元面正交 ( ∇ c · n = 0)，则迎风和顺风近似给出相同的值，因此将哪个值分配给 γ k 并不重要。另一方面，如果界面平行于单元面（向量 ∇ c 和 n 共线），则 CFL 数起着最重要的作用：如果界面在当前时间步长期间不可能到达单元面，则流体穿过界面是其下游侧的界面，因此应使用顺风近似。在界面任意方向的情况下，所有三个因素可能都很重要；详细信息请参见 Muzaferija 和 Peri ́c (1999)。

由于界面上的流体特性可能存在几个数量级的差异，因此在使用插值计算某些量的单元面值时需要小心；压力是最突出的例子。在压力纯流体静力学变化的情况下，斜率与密度成正比，因此它在界面处突然变化。为了正确计算作用在控制体积上的压力，不应跨界面插值，而应使用单侧外推法；参见，例如，Vukˇcevi ́c 等人。 (2017) 详细描述了一种合适的方法。

界面捕获方法的功能示例如图 13.13 所示，它说明了“溃坝”问题的解，这是计算自由表面流方法的标准测试用例。阻挡液体的屏障突然被移除，留下自由的垂直水面。当水沿着地板向右移动时，它会碰到障碍物，流过它并撞击对面的墙壁。当水落到障碍物另一侧的地板上时，受限的空气由于浮力而向上逸出。数值结果与 Koshizuka 等人的实验相比较。 （1995）。此示例说明了计算液相和气相流量的重要性。如果忽略气相，液体就会下落，而不会感受到滞留空气的任何阻力，从而导致明显不同的运动。这在下一个示例中更为重要，其中表面张力也很重要。

当气体被困在液体中时，浮力效应总是很重要，反之亦然；仅当自由表面存在高曲率或表面张力系数因温度或浓度梯度沿自由表面变化时，表面张力效应才显著。浮力和表面张力都很重要的一个例子是小气泡的上升。图 13.14 显示了这样一种情况：空气以 1mm/s 的速度流入一端直径为 40mm 的管道（另一端封闭），并从那里通过直径为 5mm 的连接管进入一个充满液体、深度为 45mm 的较大容器。液体密度为1500kg/m 3，粘度为1Pa·s，空气密度为1.18kg/m 3，粘度为0.0000185Pa·s。表面张力系数是常数，σ = 0.074 牛/米。如果要获得清晰的界面，流体性质的这种巨大变化对于数值方法来说是一个巨大的挑战。这里使用了 Muzaferija 和 Peri ́c (1999) 的 HRIC 方案（高分辨率界面捕获），当界面被认为是尖锐时，它通常会导致只有一个细胞的液体体积分数在 0 到 1 之间。为了更好地解决气泡曲率问题，采用自适应网格细化

\begin{sourceblock}
\includegraphics[page=539,trim={53.00bp 303.25bp 53.87bp 52.00bp},clip,width=\textwidth,max height=0.38\textheight]{Ferziger4th.pdf}
\par\vspace{0.45\baselineskip}
\small 图 13.13 越过障碍物的塌陷水柱流的实验可视化（左）和数值预测（右）比较（Koshizuka 等人的实验，1995 年；Muzaferija 和 Peri ́c 的预测，1999 年）
\end{sourceblock}
使用过；细化标准是界面的存在（即体积分数在 0.01 和 0.99 之间），并且界面周围的另外两个细胞层也被细化。一旦界面到达细化区末端，就需要重新进行网格细化和粗化；因此，每十个时间步执行一次网格自适应。当使用并行计算时，这会导致在处理器之间重新分配单元的额外工作（这是当今的规则而不是例外）。由于使用了二阶时间离散化，因此接口在每个时间步长中不允许移动超过三分之一的单元，以避免过冲和下冲（如果抛物线经过一个步骤，通常会获得这种情况）。如果没有自适应网格细化，则需要在任何时间可能存在自由表面的地方都拥有最精细的网格级别，这将使单元数量（以及计算时间）增加 10 倍。有关此模拟的更详细报告可从本书网站获得；详细内容请参见附录。

\begin{sourceblock}
\includegraphics[page=540,trim={53.00bp 325.25bp 53.87bp 52.00bp},clip,width=\textwidth,max height=0.38\textheight]{Ferziger4th.pdf}
\par\vspace{0.45\baselineskip}
\small 图 13.14 使用 HRIC 方案模拟上升气泡周围的运动和流动：一系列气泡通过自由表面上升和破裂（上）；蒸气体积分数的分布和适应无气泡表面的网格（左下）；气泡轮廓和气泡内部和外部的速度矢量（右下）
\end{sourceblock}
另一类接口捕获方法基于水平集公式，由 Osher 和 Sethian (1988) 引入。该表面被定义为水平集函数 φ = 0 的表面。该函数的其他值没有意义，为了使其成为平滑函数，通常将 φ 初始化为距界面的有符号距离，即，其任意点的值是距表面上最近点的距离，其符号一侧为正，另一侧为负。然后允许该函数演化为传输方程的解：

\begin{sourceblock}
\begin{tikzpicture}
\node[inner sep=0,anchor=south west] (image) at (0,0) {\includegraphics[page=540,trim={53.00bp 139.27bp 53.87bp 498.59bp},clip,width=0.92\textwidth,max height=0.38\textheight]{Ferziger4th.pdf}};
\begin{scope}[x={(image.south east)},y={(image.north west)}]
\fill[white] (0.37438,0.54880) rectangle (0.41263,0.95757);
\end{scope}
\end{tikzpicture}
\end{sourceblock}
其中 v 是局部流体速度，在任何时候，φ = 0 的表面都是界面。如果函数φ变得过于复杂，可以按照上述方式重新初始化。与 VOF 类方法一样，流体属性由 φ 的局部值决定，但在这里，只有 φ 的符号很重要。

相对于 VOF 方法，这种方法的明显优点是 φ 在界面上平滑变化，而体积分数 c 在那里不连续。然而，在求解体积分数 c 时，通常不会维持其在界面上的逐步变化 - 该阶跃被数值近似所模糊。因此，流体特性在界面上经历平滑的变化。在水平集方法中，属性的逐步变化得以保持，因为我们定义

\begin{sourceblock}
\begin{tikzpicture}
\node[inner sep=0,anchor=south west] (image) at (0,0) {\includegraphics[page=541,trim={53.00bp 510.54bp 53.87bp 140.48bp},clip,width=0.92\textwidth,max height=0.38\textheight]{Ferziger4th.pdf}};
\begin{scope}[x={(image.south east)},y={(image.north west)}]
\fill[white] (0.00000,0.92262) rectangle (0.03907,1.00000);
\fill[white] (0.04174,0.92262) rectangle (0.11973,1.00000);
\end{scope}
\end{tikzpicture}
\end{sourceblock}
然而，在计算粘性流时，这通常会引起问题，因此需要引入一个有限厚度 ε 的区域（通常为一到三个单元宽），在该区域上，整个界面上的属性会发生平滑但快速的变化。这使得该解变得与 VOF 方法的解相似。

如上所述，计算出的 φ 需要时不时地重新初始化。苏斯曼等人。 ( 1994 ) 提出这可以通过求解以下方程来完成：

\begin{sourceblock}
\includegraphics[page=541,trim={53.00bp 392.16bp 53.87bp 245.96bp},clip,width=0.92\textwidth,max height=0.38\textheight]{Ferziger4th.pdf}
\end{sourceblock}
直到达到稳态。这保证了 φ 与 φ 0 具有相同的符号和零电平，并满足 | 的条件。 ∇ φ | = 1，使其类似于有符号距离函数。

由于 φ 没有明确出现在任何守恒方程中，因此原始水平集方法并未完全守恒质量。质量守恒可以通过使等式的右侧得到加强。 ( 13.48 ) 局部质量不平衡的函数̇ m 正如Zhang 等人所做的那样。 （1998）。求解该方程的次数越多，达到稳态所需的迭代次数就越少；当然，频繁求解该方程会增加计算成本，因此需要权衡。

已经提出了许多水平集方法；他们在各个步骤的选择上有所不同。张等人。 (1998)描述了一种这样的方法，他们将其应用于气泡合并和模具填充，包括熔体凝固。他们在结构化非正交网格上使用 FV 方法来求解守恒方程，并使用 FD 方法来求解水平集方程。 ENO 方案用于离散后者的对流项。恩莱特等人。 (2002) 描述了拉格朗日标记粒子和水平集方法，以更好地捕获界面。

该方法的另一个版本用于研究火焰传播。在这种情况下，火焰相对于流体传播，这导致表面可能出现尖点，即表面法线不连续的位置。这将在下面进一步讨论。

有关水平集方法的更多详细信息，请参阅 Sethian (1996) 以及 Osher 和 Fedkiw (2003) 的书籍；另见 Smiljanovski 等人。 ( 1997 ) 和 Reinecke 等人。 (1999) 用于火焰跟踪的类似方法的示例。

界面捕获或 VOF 方法是计算自由表面流最广泛使用的方法；它们可用于所有主要的商业和公共规范，并已成功应用于模拟水的进入和物体撞击液体表面、船舶和水下物体周围的流动、初级射流破裂、近壁气泡塌陷、液滴壁相互作用等。下面将进一步简要描述一些应用示例。

\subsection*{13.5.2 接口跟踪方法}\phantomsection\addcontentsline{toc}{subsection}{13.5.2 接口跟踪方法}
在计算水下物体周围的流动时，许多作者对未受扰动的自由表面进行了线性化。这需要引入高度函数，它是相对于未扰动状态的自由表面高程：

\begin{sourceblock}
\includegraphics[page=542,trim={53.00bp 451.56bp 53.87bp 200.25bp},clip,width=0.92\textwidth,max height=0.38\textheight]{Ferziger4th.pdf}
\end{sourceblock}
运动学边界条件 ( 13.38 ) 变为以下方程，描述高度 H 的局部变化：

\begin{sourceblock}
\includegraphics[page=542,trim={53.00bp 379.95bp 53.87bp 257.91bp},clip,width=0.92\textwidth,max height=0.38\textheight]{Ferziger4th.pdf}
\end{sourceblock}
该方程可以使用第 1 章中描述的方法及时积分。 6.自由表面处的流体速度可以通过从内部外推或使用动态边界条件（13.39）来获得。

这种方法通常与结构化网格和显式时间积分结合使用。许多作者使用 FV 方法进行流动计算，使用 FD 方法进行高度方程，并仅在收敛稳态时在自由表面强制执行两种边界条件（参见例如 Farmer 等人，1994）。

Hino (1992) 使用了 FV 方法并强制执行 SCL，从而满足每个时间步长的所有条件并确保体积守恒。 Raithby 等人开发了类似的方法。 ( 1995 )，Thé 等人。 （1994）和利莱克（1995）。这种类型的一种完全保守的 FV 方法包括以下步骤：

\begin{itemize}
\item 使用当前自由表面上的指定压力求解动量方程以获得速度 u ∗
\end{itemize}
我。 • 通过求解压力校正方程，在当前自由表面处使用零压力校正边界条件，强制每个 CV 中的局部质量守恒（参见第 11.2.4 节）。质量在全局和每个 CV 中均守恒，但自由表面处的规定压力会产生速度修正，因此通过自由表面的质量通量不为零。 • 更正自由表面的位置以强制执行运动学边界条件。每个自由表面单元面都会移动，以便由于其移动而产生的体积通量补偿在上一步中获得的通量。 • 迭代直至无需进一步调整且满足连续性方程和动量方程。 • 前进到下一个时间步骤。

影响该方法效率和稳定性的关键问题是自由表面运动的算法。问题是每个自由表面单元面只有一个离散方程，但必须移动大量网格节点。如果要避免波反射和/或不稳定，则必须正确处理自由表面与其他边界（入口、出口、对称性、壁）的交点。我们将简要描述一种这样的方法。仅考虑了两个时间级别的方案，但该方法可以扩展。

通过移动自由表面单元面的质量通量为（参见方程（13.31））：

\begin{sourceblock}
\begin{tikzpicture}
\node[inner sep=0,anchor=south west] (image) at (0,0) {\includegraphics[page=543,trim={53.00bp 473.87bp 53.87bp 167.06bp},clip,width=0.92\textwidth,max height=0.38\textheight]{Ferziger4th.pdf}};
\begin{scope}[x={(image.south east)},y={(image.north west)}]
\fill[white] (0.09841,0.36533) rectangle (0.14177,0.86950);
\fill[white] (0.14677,0.41095) rectangle (0.17495,0.86950);
\end{scope}
\end{tikzpicture}
\end{sourceblock}
上标 τ 表示计算数量的时间 ( t n < t τ < t n + 1 )；对于隐式欧拉方案，t τ = t n + 1，而对于 Crank-Nicolson 方案，t τ = 1
2(tn+tn+1)。由自由表面指定压力的压力修正方程得到的质量通量不为零；我们通过移动自由表面来补偿，即：

\begin{sourceblock}
\begin{tikzpicture}
\node[inner sep=0,anchor=south west] (image) at (0,0) {\includegraphics[page=543,trim={53.00bp 378.84bp 53.87bp 268.56bp},clip,width=0.92\textwidth,max height=0.38\textheight]{Ferziger4th.pdf}};
\begin{scope}[x={(image.south east)},y={(image.north west)}]
\fill[white] (0.00000,0.82551) rectangle (0.09783,1.00000);
\fill[white] (0.10054,0.82551) rectangle (0.15029,1.00000);
\end{scope}
\end{tikzpicture}
\end{sourceblock}
从这个方程我们得到流体的体积 ̇ V ′
由于自由表面运动，fs 必须流入或流出 CV。我们需要从该方程获得位于自由表面上的 CV 顶点的坐标。这必须小心完成，因此需要特别注意。由于每个单元面都有一个体积流量，但 CV 顶点数量较多，因此未知数比方程多。

Thé等人。 (1994) 建议在与自由表面相邻的层中使用交错 CV，但仅在连续性（压力校正）方程中使用。该方法应用于解决多个二维问题并表现出良好的性能。然而，它需要对求解方法进行大量调整，尤其是在 3D 中；参见 Thé 等人。 （1994）了解更多细节。

另一种可能性是不通过顶点而是通过单元面中心来定义自由表面下的 CV；然后通过插值单元面中心位置来定义顶点，如图 13.15 所示的 2D 结构化网格。自由表面单元面扫过的体积为：

\begin{center}\small 图13.15 自由表面下的控制体积，其顶点位于自由表面上，并由单元面中心的坐标（空心符号）定义；时间步长期间细胞表面扫过的体积用阴影表示\end{center}
旧位置

h

V

位置 新 δ

\begin{sourceblock}
\includegraphics[page=544,trim={53.00bp 584.32bp 53.87bp 53.79bp},clip,width=0.92\textwidth,max height=0.38\textheight]{Ferziger4th.pdf}
\end{sourceblock}
其中 h 是自由表面标记在一个时间步长内移动的距离； h n = hi 而h nw 和h ne 是通过对hi 和hi − 1 或hi 和hi + 1 进行线性插值得到的。将 h nw 、 h n 和 h ne 表示为 hi 、 hi − 1 、 hi + 1 并将上述表达式代入式（1）。 ( 13.52 )，我们得到了单元面中心位置 h i 的方程组。在2D中，系统是三对角的，可以直接用节的TDMA方法求解。 5.2.3.在 3D 中，系统是块三对角的，最好通过第 1 章中介绍的迭代求解器之一来求解。 5.必须在自由表面边缘的 CV 顶点处指定“边界条件”。如果不允许边界移动，则h = 0。如果允许自由表面的边缘移动，例如对于开放系统，边界条件应为非反射或“波透射”类型，不会引起波反射；条件（10.4）是一种适当的可能性。

Lilek (1995) 将这种方法应用于结构化网格上的 2D 和 3D 问题。当侧边界面具有不规则形状时（例如船体），体积的表达式为 δ V ′
fs 变得复杂并且需要在每次外部迭代上进行迭代求解。

Muzaferija 和 Peri ́c (1997) 提出了一种更简单的方法。他们指出，没有必要根据自由表面上单元顶点的几何形状来计算扫掠体积；它可以从等式获得。 （13.52）。位于细胞面中心上方的自由表面标记的位移由高度 h 定义，高度 h 是从已知体积和细胞面面积获得的。然后通过插值 h 计算新的顶点位置；所得到的扫过体积并不精确，需要进行迭代修正。该方法适用于隐式方案，无论如何，每个时间步都需要外部迭代。图 13.15 中显示的“旧”和“新”位置现在是当前和先前外部迭代的值；每次外迭代根据方程（13.52）校正扫掠体积。在每个时间步结束时，当外迭代收敛时，所有校正为零。有关此方法及其在任意非结构化 3D 网格上的实现的详细描述，请参阅 Muzaferija 和 Peri ́c (1997, 1999)。

具有自由表面的流动，如明渠流、船舶周围流等，用弗劳德数来表征：

\begin{sourceblock}
\includegraphics[page=544,trim={53.00bp 171.57bp 53.87bp 465.66bp},clip,width=0.92\textwidth,max height=0.38\textheight]{Ferziger4th.pdf}
\end{sourceblock}
其中g为重力加速度，v为参考速度，L为参考长度； √ gL 是深水中长度为 L 的波的速度。当 Fr > 1 时，流体速度大于波速，流动被称为超临界流动，波不能向上游传播（如超音速可压缩流动中的压力波的情况）。当 Fr < 1 时，波可以向各个方向传播。如果计算自由表面形状的方法没有正确实施，可能会产生小波浪形式的扰动，并且可能无法获得稳定的自由表面形状。

\begin{sourceblock}
\includegraphics[page=545,trim={53.00bp 457.26bp 53.87bp 52.00bp},clip,width=\textwidth,max height=0.38\textheight]{Ferziger4th.pdf}
\par\vspace{0.45\baselineskip}
\small 图 13.16 界面跟踪方法中使用的初始（上）和最终（下）最粗糙网格来计算半圆柱体上的临界流动（下图还显示了动态压力分布；二维模型）
\end{sourceblock}
解。一种不会在物理上不应该产生波的地方（例如，在船前面）产生波的方法被认为满足辐射条件。另一种方法在第 2 节中描述。 13.6.

我们在此展示一个示例，其中使用并比较了界面跟踪和界面捕获方法。研究了放置在通道底部的半圆柱体上的湍流，其上游基于深度的弗劳德数低于 1。流在经过气缸时可能会变成超临界。在临界流中，弗劳德数在气缸上方等于 1，并且流量经历从气缸上游的亚临界 (Fr < 1) 到气缸下游的超临界 (Fr > 1) 的转变。上游水位和流速不能同时独立设定；这些量之一必须适应临界条件（增加流速导致上游水深增加）。这里选择上游水位与圆柱半径之比为2.3的情况，并根据Forbes（1988）的结果将入口速度设置为0.275 m/s。图 13.16 显示了界面跟踪方法中使用的粗网格的初始和最终形状。这种方法的问题在于，网格需要移动并适应自由表面的形状，并且这对于任意情况都不容易实现自动化。当表面网格点可以简单地移动（例如垂直方向）时，解相对容易；然而，解域内的网格必须适应表面网格运动，这对于非结构化网格和大变形来说可能会很复杂。在这种特殊情况下，使用了块结构网格；它的设计使得细胞在整个溶液过程中保持合理的质量。使用代数平滑方法使内部网格适应自由表面的运动。与圆柱体上方和下游的初始形状相比，网格经历了很大的变形。

使用界面捕获方法和沿壁具有棱柱层的修剪笛卡尔网格计算了相同的流量。图 13.17 显示了

\begin{sourceblock}
\includegraphics[page=546,trim={53.00bp 420.25bp 53.87bp 52.00bp},clip,width=\textwidth,max height=0.38\textheight]{Ferziger4th.pdf}
\par\vspace{0.45\baselineskip}
\small 图 13.17 介质网格还显示了水的体积分数和界面位置（上）以及使用界面捕获方法计算半圆柱体上的临界流量获得的预测流型（下）
\end{sourceblock}
\begin{sourceblock}
\includegraphics[page=546,trim={53.00bp 243.90bp 53.87bp 279.36bp},clip,width=\textwidth,max height=0.38\textheight]{Ferziger4th.pdf}
\par\vspace{0.45\baselineskip}
\small 图13.18 界面跟踪和界面捕获方法预测的自由表面形状与福布斯实验数据的比较（1988）
\end{sourceblock}
网格和计算的体积分数分布。图 13.18 显示了两种方法使用的最细网格上计算的自由表面轮廓，与 Forbes (1988) 的实验数据进行了比较。协议非常好；实际上，即使使用图13.16所示的粗网格，自由表面形状的预测也相对较好。在界面捕获法的情况下，水的体积分数在1和0之间变化的区域大约是一个单元宽，如图13.17所示；在这种情况下，当网格细度变化时，体积分数 0.5 的等值面位置也不会发生太大变化。然而，在其他应用中，例如，当需要捕获波浪时，结果的网格依赖性可能很大。

\subsection*{13.5.3 混合方法}\phantomsection\addcontentsline{toc}{subsection}{13.5.3 混合方法}
最后，还有一些不属于上述任一类别的计算两相流的方法。这些方法借用了接口捕获和接口跟踪方法的元素，因此我们将它们称为混合方法。其中有一种由 Tryggvason 及其同事开发的方法，已应用于气泡流，请参见 Tryggvason 和 Unverdi (1990) 以及 Bunner 和 Tryggvason (1999)。

在此方法中，流体属性被涂抹在垂直于界面的固定数量的网格点上。然后将两相视为具有可变属性的单一流体，如界面捕获方法中那样。为了防止界面被弄脏，它也像界面跟踪方法一样被跟踪。这是通过使用流解算器生成的速度场移动标记粒子来完成的。为了保持准确性，添加或删除标记颗粒以保持它们之间大致相等的间距。为此目的，建议使用水平集方法作为替代方法；参见 Enright 等人的粒子水平集方法。 (2002) 或 Osher 和 Fedkiw (2003)。每个时间步之后，都会重新计算属性。

Tryggvason 和他的同事用这种方法计算了许多流动，其中包括水中含有数百个水蒸气气泡的流动。相变、表面张力以及气泡的合并和分裂都可以用这种方法处理。

Scardovelli 和 Zaleski (1999) 报道了类似的混合方法，其中使用了一相体积分数的附加方程和界面跟踪。

另一种混合方法以缩写 PLIC（或 VOF-PLIC）而闻名，意思是分段线性界面构造（Youngs 1982）。一相体积分数的方程按照上述界面捕获方法求解，但使用 2D 中的线段和 3D 中的平面对每个单元中的界面进行几何重建，使得单元被分成对应于相的体积分数和界面法线方向（由体积分数的梯度给出）的两部分。 Rider 和 Kothe (1998) 对类似方法进行了回顾。一些特殊的方法是专门针对笛卡尔网格开发的（例如，参见 Weymouth 和 Yue 2010 以及Qin 等人 2015）。还提出了 VOF 和水平集方法的组合；参见例如 Sussman (2003)。处理可压缩相时需要特别注意；参见 Johnsen 和 Ham (2012) 以及 Beig 和 Johnsen (2015)。

与上述方法类似的方法已用于使用 DNS 对详细的自由表面现象进行精确模拟，例如碎波中的空气夹带和混合（例如，Deike 等人，2016 年）和水跃（Mortazavi 等人，2016 年）。文献中提供了大量涉及射流破碎、上升气泡、壁附近气泡破裂等的出版物。这些方法已与湍流的 LES 和 RANS 建模结合使用。

在许多应用中，两种流体之间的界面并不明显。一个例子是碎波或水跃，其中存在水和空气形成泡沫混合物的区域。在这种情况下，有必要添加两种流体的混合模型，类似于单相流中的湍流传输模型。此类模型可以使用 Deike 等人获得的 DNS 数据来开发。 （2016）。在计算沸腾过程中由于热量增加以及其他一些具有自由表面的流动而导致的气泡增长时，也需要进行特殊处理；详细介绍每个特定应用超出了本书的范围，但大多数使用的方法都与上述方法之一相关。

\section*{13.6 强制解}\phantomsection\addcontentsline{toc}{section}{13.6 强制解}
人们可以应用解强制来实现相同的目标，而不是尝试构建一种不在边界处反射波的数值方案。在这种方法中，人们将纳维-斯托克斯（或 RANS）方程的解与参考解混合在一起，参考解通常是平凡解（例如均匀流或平坦自由表面）、理论解或使用另一种方法（例如基于势流理论的方法）或不同边界条件（例如，波在没有任何障碍的无限域中传播）获得的解。混合是通过将源项添加到离散动量方程（如果使用界面捕获方法计算自由表面流，也可能添加到一相的体积分数方程）来实现的，其形式如下：

\begin{sourceblock}
\begin{tikzpicture}
\node[inner sep=0,anchor=south west] (image) at (0,0) {\includegraphics[page=548,trim={53.00bp 319.19bp 53.87bp 319.75bp},clip,width=0.92\textwidth,max height=0.38\textheight]{Ferziger4th.pdf}};
\begin{scope}[x={(image.south east)},y={(image.north west)}]
\fill[white] (0.21844,0.42500) rectangle (0.28460,0.89265);
\fill[white] (0.28827,0.46728) rectangle (0.31645,0.89265);
\fill[white] (0.33423,0.04412) rectangle (0.34833,0.35919);
\end{scope}
\end{tikzpicture}
\end{sourceblock}
\begin{sourceblock}
\begin{tikzpicture}
\node[inner sep=0,anchor=south west] (image) at (0,0) {\includegraphics[page=548,trim={53.00bp 297.33bp 53.87bp 353.20bp},clip,width=0.92\textwidth,max height=0.38\textheight]{Ferziger4th.pdf}};
\begin{scope}[x={(image.south east)},y={(image.north west)}]
\fill[white] (0.09576,0.00000) rectangle (0.11224,0.49263);
\fill[white] (0.11567,0.07687) rectangle (0.19537,0.81807);
\fill[white] (0.19651,0.07687) rectangle (0.23624,0.81807);
\fill[white] (0.23741,0.07687) rectangle (0.27883,0.81807);
\fill[white] (0.27997,0.07687) rectangle (0.39618,0.81807);
\fill[white] (0.39735,0.07687) rectangle (0.49871,0.81807);
\fill[white] (0.49985,0.07687) rectangle (0.52629,0.81807);
\fill[white] (0.52743,0.07687) rectangle (0.56887,0.81807);
\fill[white] (0.57002,0.07687) rectangle (0.61808,0.81807);
\fill[white] (0.61925,0.07687) rectangle (0.72971,0.81807);
\fill[white] (0.73089,0.07687) rectangle (0.78226,0.81807);
\fill[white] (0.78346,0.07687) rectangle (0.86644,0.81807);
\fill[white] (0.86761,0.07687) rectangle (0.92731,0.81807);
\fill[white] (0.92848,0.07687) rectangle (1.00000,0.81807);
\fill[white] (0.00000,0.00000) rectangle (0.02746,0.05125);
\fill[white] (0.02737,0.00000) rectangle (0.06045,0.05125);
\fill[white] (0.06036,0.00000) rectangle (0.16000,0.05125);
\fill[white] (0.15991,0.00000) rectangle (0.18803,0.05125);
\fill[white] (0.18794,0.00000) rectangle (0.31140,0.05125);
\fill[white] (0.31134,0.00000) rectangle (0.38935,0.05125);
\fill[white] (0.38926,0.00000) rectangle (0.41405,0.05125);
\fill[white] (0.41398,0.00000) rectangle (0.58352,0.05125);
\fill[white] (0.58343,0.00000) rectangle (0.61820,0.05125);
\fill[white] (0.61811,0.00000) rectangle (0.71365,0.05125);
\fill[white] (0.71356,0.00000) rectangle (0.75498,0.05125);
\fill[white] (0.75489,0.00000) rectangle (0.84620,0.05125);
\fill[white] (0.84611,0.00000) rectangle (0.97660,0.05125);
\fill[white] (0.97654,0.00000) rectangle (0.99639,0.07495);
\end{scope}
\end{tikzpicture}
\end{sourceblock}
P 代表细胞质心处的参考溶液。源项需要逐渐混合，这是通过在零（在强迫区域的开始处）到其最大值（在解域的边界处）之间改变强迫系数 γ 来实现的。通常希望使用在混合区开始时渐近趋于零的变化，例如指数变化或cos 2 变化。请注意，γ 不是无量纲系数；它的维度是[1/s]，最佳值取决于问题。

这种通用方法已被许多作者用于各种目的，并且可以在不同名称的文献中找到，具体取决于应用领域（松弛区、阻尼区、海绵层等）；我们不会详细讨论各种变体。关键问题是强制系数γ最大值的选择。在大多数出版物中，使用试错法来确定该参数的最佳值。最近，Peri ́c 和 Abdel-Maksoud (2018) 发表了一种方法，该方法允许确定参数范围，当自由表面波向出口边界衰减时，该参数范围会导致反射系数低于目标值。最佳值取决于强制区的波长和宽度；较高的值适用于较短的波长。对于自由表面波，介于 1（非常长的海浪）和 100（实验室实验中的短波长）之间的值以及 1 到 2 个波长宽的强制区域似乎会产生良好的结果。

如果只想将波浪推向平坦的自由表面，则逐渐将垂直速度分量推向零就足够了。然而，如果想在入口边界处施加这种波，也可以迫使所有速度分量，例如，朝向五阶斯托克斯波（例如，参见 Fenton 1985）。使用对某一波长的强制来仅指定入口处的速度和体积分数的理论值的优点在于，前者将避免上游行进扰动的反射，而后者则不会。此外，例如，如果在对围绕船侧边界的流动的研究中将其视为对称平面，则可以将横向速度分量强制为零，从而避免来自这些边界的波反射。在研究波浪中船舶周围的流动时，使用强迫来避免边界处的波浪反射的示例将在第 1 节中介绍。 13.10.1。

同样的方法也可用于阻尼可压缩流（气体和液体）中的声压波，以避免边界反射；在这种情况下，Peri ́c 和 Abdel-Maksoud (2018) 的理论也可用于确定 γ 的最佳值。例如，为了抑制水中500Hz的声压波动，当施力区为两个波长宽时，γ的最佳值为8,000左右。

如果流动带有涡流而难以指定不反映扰动的适当边界条件，则还可以使用强迫方法来避免 DNS 和 LES 模拟中下游边界的扰动。通过迫使流动，例如，沿流向方向的恒定速度和沿交叉流动方向的零速度（这适用于研究自由流中物体周围的流动），可以实现从湍流到均匀流的平滑过渡，对此简单的出口或指定的压力边界条件就足够了。

\section*{13.7 气象和海洋应用}\phantomsection\addcontentsline{toc}{section}{13.7 气象和海洋应用}
大气层和海洋是地球上规模最大的流动场所。速度可能达到每秒数十米，并且长度尺度巨大，因此雷诺数也很大。由于这些流的长宽比非常大（大气和海洋的水平方向有数千公里，垂直方向有几公里），大尺度的流几乎是二维的（尽管垂直运动很重要），而小尺度的流是三维的。地球自转在大尺度上是主要力，但在小尺度上则不太重要。分层或密度的稳定变化很重要，主要是在较小的尺度上。在不同的尺度上，起主导作用的力量和现象是不同的。

此外，还需要对不同时间尺度进行预测。在公众最感兴趣的情况下，人们希望预测未来相对较短的时间内大气或海洋的状况。在天气预报中，时间尺度是几天，而在海洋中，变化更慢，时间尺度是几周到几个月。无论哪种情况，都需要一种及时准确的方法。另一个极端是气候研究，需要预测相对较长时期内大气的平均状态。在这种情况下，可以平均掉短期时间行为并放宽时间精度要求；然而，在这种情况下，必须对海洋和大气进行建模。根据模拟所需的分辨率（例如公里或米），通常会根据如何处理湍流来划分方法。如果流动的 3D 状态很重要，通常会进行大涡模拟。另一方面，对于主要是 2D 但具有垂直混合的流动，使用具有各种湍流建模风格的 RANS 方法。

计算是在各种不同的长度尺度上完成的。最小的感兴趣区域是尺寸为数百米的大气边界层或海洋混合层。下一个尺度可以称为盆地尺度，由城市及其周围环境组成。在区域或中观尺度上，人们认为一个领域是大陆或海洋的重要组成部分。最后是大陆（或海洋）和全球尺度。在每种情况下，计算资源决定了可以使用的网格点的数量，从而决定了网格的大小。小尺度上发生的现象必须用近似模型来表示。即使在最小的关注尺度上，必须进行平均的区域的大小也明显比工程流程中的大得多。因此，用于表示较小尺度的模型比第 1 章中讨论的工程大涡模拟重要得多。 10.

事实上，重要的结构无法在最大尺度的模拟中解析，因此需要在许多不同的尺度上进行计算；在每个尺度上，目的是研究该尺度特有的现象。气象学家区分进行模拟的四到十个尺度（取决于谁进行计数）。正如人们所预料的那样，关于这一主题的文献数量巨大，我们甚至无法开始涵盖所有已完成的工作。

正如已经指出的，在最大尺度上，大气和海洋流动本质上是二维的（尽管垂直运动有重要影响）。在模拟全球大气或整个洋盆时，现有计算机的能力要求水平方向的网格大小在10到100公里之间。因此，在这些类型的模拟中，必须通过近似模型处理重要结构，例如前沿（存在于不同属性的流体块之间的区域），以使其厚度足够大，以便网格可以解析它们。这种模型非常难以构建，并且是预测误差的主要来源。

三维运动仅在大气或海洋流的最小尺度上才重要。同样重要的是要注意，尽管雷诺数很高，但只有最接近地表的大气部分是湍流的；这是大气边界层，通常覆盖约 1-3 公里厚的区域。在边界层之上，大气分层并保持层流状态。同样，只有海洋的顶层是湍流的；厚度为100\textasciitilde{}300 m，称为混合层。对这些层进行建模很重要，因为大气和海洋正是在这些层中相互作用，并且它们对大规模行为的影响非常重要。

这些模拟中使用的数值方法根据模拟执行的规模而有所不同。对于最小大气或海洋尺度的模拟，例如在大气或海洋边界层中，可以使用类似于工程流大涡模拟中使用的方法（参见第 10 章）。例如，沙利文等人。 ( 2016 ) 使用极端分辨率（0.39 m）的非静水力代码研究了夜间稳定层大气 BL；该 LES 代码在水平方向上使用谱公式，在垂直方向上使用有限差分，并具有 RK3 时间推进。 Skyllingstad 和 Samelson (2012) 在 3 米分辨率下的 FV 非静水力 LES 代码的海洋应用允许检查海洋表面边界层中的锋面不稳定性和湍流混合。区域海洋代码是区域海洋建模系统（ROMS；\url{https://www.myroms.org}），它求解静流动方程和其他耦合模型。它是使用预测器（Leapfrog；第 6.3.1.2 节）和校正器（Adams-Moulton；第 6.2.2 节）时间步方案的 FV 公式。这里，静水力是指垂直加速度小并且垂直压力梯度和重力平衡。

对于不涉及云的模拟，标准代码就足够了，但是当存在云及其相关的液态水和水蒸气以及冰等时，有必要添加微观物理包来处理潮湿过程（参见 Morrison 和 Pinto 2005）。这会增加大量需要求解的附加偏微分方程，但不会显著改变数值。 CM1 代码 ( \url{http://www2.mmm.ucar.edu/people/bryan/cm1/} ) 是此类系统的直接实现，是具有 RK3 时间推进的可压缩流体的非静水压 FD 代码。此外，与许多大气代码一样，它使用五阶或六阶精确平流方案。由于 LES 在大气建模中发挥着越来越重要的作用，Shi 等人。 ( 2018b , 2018a ) 探索了云解析 LES 代码（例如 CM1）中使用的子滤波器尺度和子网格尺度模型的准确性（另请参见 Khani 和 Porté-Agel 2017）。

作为另一个例子，Schalkwijk 等人。 ( 2012a， 2015 ) 演示了使用图形处理单元 (GPU) 和荷兰大气大涡模拟 (DALES) 代码（Heus 等人，2010）生成湍流云，这是在 Arakawa C 网格上采用 RK3 时间推进的非静水力 Boussinseq FD 公式。该代码解决了七个主要的预后变量（即通过偏微分方程提前的时间）。一般来说，针对高分辨率和/或涡流/云分辨率的应用是非静水力和 LES，而全球大气和海洋模型是静水力和 RANS；例如，参见 Washington 和 Parkinson (2005)，了解气候模型的观点。在全球尺度上，使用有限体积方法，但更常见的是专门为球体表面设计的光谱方法。该方法使用球谐函数作为基函数。

在选择时间推进方法时，必须考虑对精度的要求，但也必须注意波浪现象在气象学和海洋学中都发挥着重要作用。从气象图和卫星照片中熟悉的大型天气系统可以被视为非常大尺度的行波。数值方法不应放大或消散它们。因此，在这些领域中使用蛙跳法是很常见的（第6.3.1.2节）。

该方法对于波具有二阶精度和中性稳定。不幸的是，它也是无条件不稳定的（它会放大呈指数衰减的解），因此必须使其稳定。因此，龙格-库塔方法（第 6.2.3 节）通常已取代它。特别是，现在普遍使用三阶实时精确 RK3 方案（第 6.2.3 节）。

\section*{13.8 多相流}\phantomsection\addcontentsline{toc}{section}{13.8 多相流}
工程应用通常涉及多相流；例如，气体或液体流（流化床、含尘气体和浆料）携带的固体颗粒、液体中的气泡（起泡流体和锅炉）或气体中的液滴（喷雾）等。另一个复杂之处是燃烧系统中经常出现多相流。在许多燃烧室中，液体燃料或煤粉以喷雾形式喷射。在其他情况下，煤在流化床中燃烧。最后，工程中也经常遇到具有相变（空化、沸腾、冷凝、熔化、凝固）的多相流。

章节中描述的方法。 13.5 可以适用于某些类型的两相流，特别是两相都是流体的两相流。在这些情况下，两种流体之间的界面按照如上所述进行明确处理。有些方法是专门为这种类型的流程设计的。然而，与界面处理相关的计算成本将这些方法限制为界面面积相对较小的流动。

还有其他几种计算两相流的方法。载体或连续相流体始终通过欧拉方法处理，但分散相可以通过拉格朗日或欧拉方法处理。

当分散相的质量载荷不是很大且分散相颗粒较小时，常采用拉格朗日法；含尘气体和一些燃料喷雾是可以应用此方法的流的示例。在这种方法中，分散相由有限数量的颗粒或液滴表示，其运动以拉格朗日方式计算。被跟踪运动的粒子数量通常远小于流体中的实际数量。每个计算粒子代表一定数量（或数据包）的实际粒子；这些称为数据包方法。如果不存在相变和燃烧且负载较轻，则可忽略分散相对载流的影响，可先计算载流。然后注入粒子，并使用预先计算的背景流体速度场计算它们的轨迹。 （这种方法也用于流动可视化；一种使用无质量点粒子并跟随它们的运动来创建条纹。）这种方法需要将速度场插值到粒子位置；插值方案需要至少与用于时间推进的方法一样准确。准确性还要求选择时间步长，以便粒子在一个时间步长内不会穿过多个细胞。

当分散相的质量负载很大时，必须考虑颗粒对流体运动的影响。如果采用分组方法，则粒子轨迹和流体流动的计算必须同时进行，并且需要迭代；每个粒子都向其所在单元中的流体贡献动量（以及能量和质量）。需要对粒子之间（碰撞、团聚和分裂）以及粒子与墙壁之间的相互作用进行建模。对于这些交换，已经使用了基于实验的相关性，但不确定性可能相当大。这些问题需要另一本书来详细描述；参见 Crowe 等人的书。 (1998) 描述了最广泛使用的方法。

对于大质量载荷以及发生相变时，欧拉方法（双流体模型）适用于两个相。在这种情况下，两个相都被视为具有单独速度场和温度场的连续体；两个相通过交换项相互作用，类似于混合欧拉-拉格朗日方法中使用的交换项。函数定义每个相占据每个单元的多少。 Ishii (1975) 以及 Ishii and Hibiki (2011) 详细描述了二流体模型的原理；另见 Crowe 等人。 (1998) 描述了气体粒子和气体液滴流的一些方法。用于计算这些流动的方法与本书前面描述的方法类似，除了添加了相互作用项和边界条件，当然，需要求解两倍的方程。该方程系统也比单相流的情况要硬得多，这需要更强的欠松弛和更小的时间步长。

空化是一种重要的现象，属于两相流类别，需要专门的模型来预测。最广泛使用的方法是使用多相流的均质模型，即相之间没有滑移（它们共享相同的速度、压力和温度场），并且计算具有可变属性的有效流体的流量。各相的分布由蒸气体积分数 c v 决定，需要求解方程：

\begin{sourceblock}
\begin{tikzpicture}
\node[inner sep=0,anchor=south west] (image) at (0,0) {\includegraphics[page=553,trim={53.00bp 223.35bp 53.87bp 413.33bp},clip,width=0.92\textwidth,max height=0.38\textheight]{Ferziger4th.pdf}};
\begin{scope}[x={(image.south east)},y={(image.north west)}]
\fill[white] (0.25495,0.56687) rectangle (0.27350,0.95927);
\fill[white] (0.24944,0.08045) rectangle (0.27871,0.48201);
\fill[white] (0.30159,0.04073) rectangle (0.31922,0.33164);
\end{scope}
\end{tikzpicture}
\end{sourceblock}
与自由表面流的界面捕获方法的相似性是显而易见的，有两个重要的区别：（i）空化不一定会导致网格尺度上相之间出现尖锐的界面，因此不需要对对流项进行特殊的离散化（可以使用其他标量变量的方法），（ii）蒸汽体积分数方程包含一个源项 q v 来模拟蒸汽泡的生长和破裂。空化模型源项的推导通常基于 Rayleigh-Plesset 方程，该方程描述单个蒸汽泡的动力学：

\begin{sourceblock}
\begin{tikzpicture}
\node[inner sep=0,anchor=south west] (image) at (0,0) {\includegraphics[page=553,trim={53.00bp 79.89bp 53.87bp 550.07bp},clip,width=0.92\textwidth,max height=0.38\textheight]{Ferziger4th.pdf}};
\begin{scope}[x={(image.south east)},y={(image.north west)}]
\fill[white] (0.50081,0.03317) rectangle (0.52779,0.38474);
\end{scope}
\end{tikzpicture}
\end{sourceblock}
这里 R 是气泡半径，t 是时间，p s 是给定温度下的饱和压力，
p 是周围液体的局部压力，σ 是表面张力系数，ρ l 是液体密度。

惯性项经常被忽略，因为将其包含在建模中会显著增加复杂性，而在大多数应用中却没有实质性的好处。此外，表面张力效应仅对气泡生长的开始很重要（它限制了可以生长的种子气泡的尺寸 - 较小的气泡会因表面张力而无法膨胀）。简化的瑞利-普莱塞方程（其中忽略惯性、表面张力和粘性项）严格来说仅对气泡生长阶段有意义，它代表渐近解；当周围压力大于饱和压力时，二次方程无法求解，因为此时右侧为负值。大多数作者采用的解决这个问题的方法是简单地对压力差的绝对值求平方根，并将其符号应用于结果：

\begin{sourceblock}
\begin{tikzpicture}
\node[inner sep=0,anchor=south west] (image) at (0,0) {\includegraphics[page=554,trim={53.00bp 397.86bp 53.87bp 239.11bp},clip,width=0.92\textwidth,max height=0.38\textheight]{Ferziger4th.pdf}};
\begin{scope}[x={(image.south east)},y={(image.north west)}]
\fill[white] (0.29540,0.54988) rectangle (0.33654,0.94995);
\end{scope}
\end{tikzpicture}
\end{sourceblock}
尽管存在这一缺陷，但在大多数实际应用中基于此近似的模型仍获得了相当好的结果。

因为通过这种简化，气泡半径变化率仅取决于局部压力，因此不需要明确地跟踪气泡；对于控制体积中存在的每个气泡，无论其直径和其先前的运动历史如何，生长速率仅是流体中局部压力的函数。文献中有几种使用方程式的空化模型。 ( 13.58 ) 确定蒸气体积分数方程中的源项；我们在这里仅简要描述 Sauer (2000) 提出的最广泛使用的模型；参见 Zwart 等人。 （2004）另一个类似的模型。

该模型假设初始半径为 R 0 的球形种子气泡存在并均匀分布在液体中，其特征在于每单位体积液体的气泡数量 n 0。因此，控制体积中的气泡数量由液体量决定。上述模型参数与液体“质量”有关：众所周知，纯液体（没有任何溶解或游离的气体或固体颗粒）可以承受非常高的拉伸应力。大多数工业液体都含有杂质，并且参数n 0 通常取为10 12 的量级。通过过滤和脱气可以减少甚至避免实验中的空化现象；在模拟中，这相当于将 n 0 降低几个数量级。

在上述假设下，任意时刻控制体积内的气泡数量等于

\begin{sourceblock}
\begin{tikzpicture}
\node[inner sep=0,anchor=south west] (image) at (0,0) {\includegraphics[page=554,trim={53.00bp 122.76bp 53.87bp 528.27bp},clip,width=0.92\textwidth,max height=0.38\textheight]{Ferziger4th.pdf}};
\begin{scope}[x={(image.south east)},y={(image.north west)}]
\fill[white] (0.00000,0.92191) rectangle (0.04692,1.00000);
\fill[white] (0.04962,0.92191) rectangle (0.10770,1.00000);
\fill[white] (0.11038,0.92191) rectangle (0.13516,1.00000);
\fill[white] (0.13783,0.92191) rectangle (0.20752,1.00000);
\fill[white] (0.21029,0.92191) rectangle (0.23841,1.00000);
\end{scope}
\end{tikzpicture}
\end{sourceblock}
其中 c l 是控制体积 V 内液体的体积分数。显然，c l + c v = 1，其中c v 是蒸气的体积分数。总蒸气体积等于

\begin{sourceblock}
\includegraphics[page=555,trim={53.00bp 584.32bp 53.87bp 53.79bp},clip,width=0.92\textwidth,max height=0.38\textheight]{Ferziger4th.pdf}
\end{sourceblock}
其中 V b 是一个气泡的体积，R 是局部气泡半径。蒸气体积分数 c v 现在可以用 n 0 、 R 和空化液体体积分数 c l 来表示：

\begin{sourceblock}
\begin{tikzpicture}
\node[inner sep=0,anchor=south west] (image) at (0,0) {\includegraphics[page=555,trim={53.00bp 516.90bp 53.87bp 121.12bp},clip,width=0.92\textwidth,max height=0.38\textheight]{Ferziger4th.pdf}};
\begin{scope}[x={(image.south east)},y={(image.north west)}]
\fill[white] (0.00000,0.89225) rectangle (0.12168,1.00000);
\fill[white] (0.12439,0.89225) rectangle (0.20662,1.00000);
\fill[white] (0.20929,0.86451) rectangle (0.24328,1.00000);
\end{scope}
\end{tikzpicture}
\end{sourceblock}
当蒸汽的体积分数已知时，可以根据以下表达式计算局部气泡半径：

\begin{sourceblock}
\begin{tikzpicture}
\node[inner sep=0,anchor=south west] (image) at (0,0) {\includegraphics[page=555,trim={53.00bp 457.86bp 53.87bp 172.10bp},clip,width=0.92\textwidth,max height=0.38\textheight]{Ferziger4th.pdf}};
\begin{scope}[x={(image.south east)},y={(image.north west)}]
\fill[white] (0.00000,0.78994) rectangle (0.09729,1.00000);
\fill[white] (0.09997,0.78994) rectangle (0.12971,1.00000);
\fill[white] (0.13242,0.78994) rectangle (0.20466,1.00000);
\fill[white] (0.20743,0.78994) rectangle (0.23221,1.00000);
\fill[white] (0.23489,0.78994) rectangle (0.32878,1.00000);
\end{scope}
\end{tikzpicture}
\end{sourceblock}
现在需要定义式（1）中的蒸汽产生或消耗的速率。 （13.56）使用上述建模框架。显然，蒸汽的产生导致气泡生长，反之亦然，因此气泡半径变化率是关键参数。另一个参数是控制体积中可能发生空化的液体量。

蒸气泡随着流动而移动；因此，在任何时刻产生蒸汽的速率可以通过给定时间控制体积中存在的气泡体积变化的速率来近似：

\begin{sourceblock}
\includegraphics[page=555,trim={53.00bp 328.92bp 53.87bp 308.83bp},clip,width=0.92\textwidth,max height=0.38\textheight]{Ferziger4th.pdf}
\end{sourceblock}
气泡半径增长的速率可以根据方程计算： ( 13.58 )，这完成了空化模型。该模型的更详细推导可以在 Sauer (2000) 以及 Schnerr 和 Sauer (2001) 中找到。该模型适用于大多数商业和公共 CFD 代码，并已成功应用于研究船舶螺旋桨周围、泵和涡轮机、燃油喷射器和其他设备中的空化流。一些作者将源项与另一个参数相乘，该参数对于增长和崩溃阶段具有不同的值（正或负源项）。

我们使用开放水域测试条件下螺旋桨周围的流动示例（即螺旋桨在均匀流动中运行，以固定的规定速率旋转）来演示上述空化模型的性能以及误差估计的几个方面以及来自不同来源的误差的相互作用。在开放水域条件下，可以使用圆周方向和旋转参考系的周期性边界条件来计算单个螺旋桨叶片周围的流动。这降低了计算成本，并且允许使用比存在船和舵时更精细的网格（在这种情况下，整个螺旋桨必须包含在模拟中，并且连接到螺旋桨的网格必须随之旋转，如图 13.9 和 13.10 所示）。该螺旋桨是2011年和2015年船用推进器研讨会（\url{www.marinepropulsors.com}）上用作测试用例的螺旋桨；操作条件来自空化测试案例（来自 SVA Potsdam 的螺旋桨 VP1304，测试案例 2.3.1；参见

\begin{sourceblock}
\includegraphics[page=556,trim={53.00bp 247.25bp 53.87bp 52.00bp},clip,width=\textwidth,max height=0.38\textheight]{Ferziger4th.pdf}
\par\vspace{0.45\baselineskip}
\small 图 13.19 螺旋桨纵向剖面中的计算网格：无（上）和针对叶尖涡捕获的特殊局部细化（下）
\end{sourceblock}
测试案例和拖车报告的详细描述的研讨会程序；海因克 2011）。

使用系统细化的网格而不对尖端涡流进行特殊的局部细化（参见图13.19中的上图）会导致空化模式以及推力和扭矩的值在两个最细的网格上变化不大（差异约为0.1％），从而导致两个错误的结论：（i）离散化误差似乎非常小，（ii）建模误差似乎非常大，因为在尖端涡流中看不到任何空化范围（见图13.20左图）。可以使用涡度大小的等值面来识别这些解中的尖端涡，如图 13.20 的右图所示。因此，涡度分布可用于指导尖端涡旋内的局部网格细化，使得

\begin{sourceblock}
\includegraphics[page=557,trim={53.00bp 443.25bp 53.87bp 52.00bp},clip,width=\textwidth,max height=0.38\textheight]{Ferziger4th.pdf}
\par\vspace{0.45\baselineskip}
\small 图 13.20 使用未针对尖端涡捕获和 k - ω 湍流模型进行局部细化的网格计算出的蒸气体积分数 0.05 的等值面（左）以及来自同一解的涡度大小的等值面（右）
\end{sourceblock}
在涡核中，单元尺寸足够小，足以解决那里压力和速度的快速变化（小于螺旋桨直径的 1/1000），如图 13.19 中网格的一个纵向剖面所示。这导致了尖端涡空化的出现，但在非定常RANS模拟中，空化区的结束比局部网格细化结束要早得多，如图13.21所示。这是由于对尖端涡流内的湍流粘度的过高估计，这模糊了压力和速度分布中的峰值，导致压力过快上升到饱和水平以上。通过切换到 LES，空化区持续到局部网格细化的最后，这也如图 13.21 所示。亚网格尺度模型的湍流粘度远低于 RANS 模型的湍流粘度。 LES 解与图 13.22 中所示的实验观察到的模式非常吻合。虽然尖端涡旋和轮毂涡旋的平均形状和位置在 LES 和实验中都保持稳定，但我们也可以观察到平均形状周围的波动。

尽管本研究中最细的网格相对较细（单个螺旋桨叶片约 1500 万个单元），但与下一个较粗网格上的解相比，没有观察到 RANS 解的显著改进。这是这种建模方法的主要缺点之一：一旦离散化误差变得小于建模误差，进一步的网格细化就不会带来任何好处。使用 LES，网格细化不仅可以减少离散化误差，还可以减少需要建模的湍流谱部分；在足够细的网格水平上，它本质上变成 DNS 并且湍流粘度变得可以忽略不计。在这个例子中，粗网格上的 LES 产生脉动尖端涡流空化：它在延伸到局部网格细化结束和收缩到类似于 RANS 解的程度之间周期性变化。然而，一旦网格进一步细化（大约 500 万个单元），空化区就变得稳定。正如预期的那样，下一步的改进暴露了尖端和轮毂涡流中更精细的结构。本书网站上提供了一份包含有关此模拟的更多详细信息的单独报告；详细内容请参见附录。

\begin{sourceblock}
\includegraphics[page=558,trim={53.00bp 456.25bp 53.87bp 52.00bp},clip,width=\textwidth,max height=0.38\textheight]{Ferziger4th.pdf}
\par\vspace{0.45\baselineskip}
\small 图 13.21 使用带有局部细化的尖端涡捕获网格和 k - ω 湍流模型（左）以及使用相同的网格但 LES 模拟（右）计算蒸汽体积分数 0.05 的等值面
\end{sourceblock}
\begin{sourceblock}
\includegraphics[page=558,trim={53.00bp 261.90bp 53.87bp 243.36bp},clip,width=\textwidth,max height=0.38\textheight]{Ferziger4th.pdf}
\par\vspace{0.45\baselineskip}
\small 图13.22 实验中两个时刻观察到的空化模式照片，显示了尖端和轮毂涡流空化的稳定和波动特征（来源：SVA Potsdam）
\end{sourceblock}
这个例子还表明，在评估离散化和建模误差时必须小心。仅仅监控推力和扭矩等一些积分量是不够的，除非人们只对它们感兴趣；当积分量停止显著变化时，尤其是在物体之后（由于高雷诺数时上游影响有限），局部流动特征可能仍然远未收敛。在这种特殊情况下，解对局部网格细化和湍流模型比对空化模型更敏感。考虑到方程给出的气泡增长率和压力之间的关系，后者的表现出人意料地好。 ( 13.58 ) 是瑞利-普莱塞方程的一个非常粗略的近似。

\section*{13.9 燃烧}\phantomsection\addcontentsline{toc}{section}{13.9 燃烧}
另一个重要的问题领域涉及燃烧（即具有显著放热的化学反应）发挥重要作用的流动。某些应用程序对于读者来说应该是显而易见的。一些燃烧器在几乎恒定的压力下运行，因此放热的主要影响是密度的降低。在许多燃烧系统中，绝对温度通过火焰增加五到八倍并不罕见；密度降低相同的因子。在这种情况下，不可能通过前面讨论的 Boussinesq 近似来处理密度差异。在其他系统中（发动机气缸是最常见的例子），压力和密度都有很大的变化。

可以对湍流燃烧流进行直接数值模拟，但仅限于非常简单的情况。值得注意的是，火焰相对于气体的传播速度很少超过 1 m/s（爆炸或爆炸是一个例外）。这个速度远低于气体中的声速；通常，流体速度也远低于声速。那么马赫数远小于单位，我们就会遇到一种奇怪的情况，即温度和密度变化很大的流动基本上是不可压缩的。

通过求解可压缩运动方程可以计算燃烧流。这已经完成了（参见 Poinsot 等人，1991）。问题在于，正如我们之前指出的，大多数为可压缩流设计的方法在应用于低速流时变得非常低效，从而提高了模拟成本。因此，这些模拟的成本非常高；当化学反应简单时尤其如此。然而，当包含更现实（因此更复杂）的化学反应时，与化学反应相关的时间尺度范围几乎总是非常大，这表明要使用较小的时间步长。换句话说，方程非常僵硬。在这种情况下，可以很大程度上消除使用可压缩流方法的损失。

另一种方法是引入低马赫数近似（McMurtry 等人，1986）。我们从描述可压缩流的方程开始，并假设所有要计算的量都可以表示为马赫数的幂级数。这是非奇异微扰理论，因此不需要特别注意。然而，结果有些令人惊讶。到最低（零）阶，动量方程简化为压力
p ( 0 ) 处处恒定。这是热力学压力，气体的密度和温度通过其状态方程相关。连续性方程具有可压缩（可变密度）形式，这并不奇怪。在下一阶，动量方程恢复其通常形式，但它们仅包含一阶压力的梯度 p (1)，这本质上是不可压缩方程中发现的动态水头。这些方程类似于不可压缩的纳维-斯托克斯方程，可以通过本书中给出的方法求解。

在燃烧理论中，区分了两种理想化情况。首先，反应物在发生任何反应之前完全混合，我们有预混合火焰的情况。内燃机已接近这个极限。在预混合燃烧中，反应区或火焰以层流火焰速度相对于流体传播。在另一种情况下，反应物同时混合和反应，即非预混合燃烧。这两种情况截然不同，需要分别对待。当然，有很多情况并不接近这两个极限。它们被称为部分预混。为了完整地阐述燃烧理论，建议读者查阅 Williams (1985) 的著名著作。

反应流中的关键参数是流动时间尺度与化学时间尺度的比率；它被称为 Damköhler 数 Da。当达姆科勒数非常大时，化学反应非常快，几乎在反应物混合后立即发生。在此限制下，火焰非常稀薄，流动被称为混合主导。事实上，如果可以忽略热释放的影响，则可以将极限 Da→∞ 视为涉及被动标量的极限，并且可以使用本章开头讨论的方法。

对于实际燃烧室的计算，其中流动几乎总是湍流，有必要依靠雷诺平均纳维-斯托克斯（RANS）方程的解。第 1 章描述了这种方法以及需要与其结合用于非反应流的湍流模型。 10.当存在燃烧时，有必要求解描述反应物质浓度的附加方程，并包括允许计算反应速率的模型。我们将在本节的其余部分描述其中一些模型。

最明显的方法，雷诺平均反应流方程，不起作用。原因是化学反应速率是温度的强烈函数。例如，物质 A 和物质 B 之间的反应速率可由下式给出：

\begin{sourceblock}
\includegraphics[page=560,trim={53.00bp 235.23bp 53.87bp 411.61bp},clip,width=0.92\textwidth,max height=0.38\textheight]{Ferziger4th.pdf}
\end{sourceblock}
其中 E a 称为活化能，R 是气体常数，Y A 和 Y B 是两种物质的浓度。阿伦尼乌斯因子 e (−E a / RT ) 的存在使问题变得困难。它随温度变化如此之快，以至于用雷诺平均值代替 T 会产生很大的误差。

在高 Damköhler 数非预混湍流火焰中，反应发生在薄的皱纹火焰区域中。对于这种情况，使用了多种方法，我们将简要描述其中两种。尽管他们之间的哲学和外表存在显著差异，但他们比看起来更相似。

在第一种方法中，人们认为，由于混合是较慢的过程，因此反应速率取决于其发生的速度。在这种情况下，两种物质 A 和 B 之间的反应速率由以下形式的表达式给出：

\begin{sourceblock}
\includegraphics[page=561,trim={53.00bp 584.61bp 53.87bp 53.77bp},clip,width=0.92\textwidth,max height=0.38\textheight]{Ferziger4th.pdf}
\end{sourceblock}
其中 τ 是混合的时间尺度。例如，如果使用 k - ε 模型，则 τ = k /ε。在 k - ω 模型中，τ = 1 /ω。已经提出了许多此类模型，其中最著名的可能是斯伯丁（Spalding，1978）的涡流破碎模型。此类模型通常用于预测工业炉的性能。

非预混燃烧的另一种模型是层流火焰模型。在停滞条件下，非预混火焰会随着时间的推移而厚度增加而缓慢衰减。为了防止这种情况发生，火焰上必须有压缩应变。火焰的状态由该应变率或其更常用的替代值标量耗散率 χ 确定。然后假设火焰的局部结构仅由几个参数决定；至少需要反应物的局部浓度和标量耗散率。将火焰结构的数据制成表格。然后将查表获得的反应速率与每单位体积的火焰面积的乘积计算为体积反应速率。已经给出了火焰面积方程的多个版本。我们不会在这里介绍，但应该足以说明这些模型包含描述通过火焰拉伸和火焰面积破坏来增加火焰面积的术语。

对于相对于流动传播的预混合火焰，小火焰模型的等价物是一种水平集方法。如果假设火焰位于某个变量 G = 0 的位置，则 G 满足方程：

\begin{sourceblock}
\includegraphics[page=561,trim={53.00bp 288.86bp 53.87bp 348.19bp},clip,width=0.92\textwidth,max height=0.38\textheight]{Ferziger4th.pdf}
\end{sourceblock}
其中 S L 是层流火焰速度。可以证明反应物的消耗率为 S L |∇ G |，从而完成模型。在模型的更复杂版本中，火焰速度可能是局部应变率的函数，就像它取决于非预混合火焰中的标量耗散率一样。

最后，请注意，有许多效应很难包含在任何燃烧模型中。其中包括点火（火焰的产生）和熄灭（火焰的消灭）。目前，湍流燃烧模型正在快速发展，该领域的任何快照都无法长期保持最新状态。对这个主题感兴趣的读者应该查阅 Peters (2000) 的书。

\section*{13.10 流固耦合}\phantomsection\addcontentsline{toc}{section}{13.10 流固耦合}
流体总是对水下结构施加力（即使没有流动），但到目前为止我们假设固体墙是刚性的。流固耦合 (FSI) 是指暴露于流体流动的固体的运动或变形。因此，流动受到实体壁运动的影响，并且需要对两种运动进行耦合模拟。

FSI 最简单的形式是飞行或漂浮刚体的运动；当固体结构也由于流动引起的力而变形时，就会引入下一个级别的复杂性。我们简要讨论如何模拟这些现象并展示一些说明性示例。

\subsection*{13.10.1 漂浮和飞行体}\phantomsection\addcontentsline{toc}{subsection}{13.10.1 漂浮和飞行体}
如果飞行或漂浮体可以被认为是刚性的，则可以通过求解质心平移和绕其旋转的常微分方程来模拟其运动： d ( m v )

\begin{sourceblock}
\begin{tikzpicture}
\node[inner sep=0,anchor=south west] (image) at (0,0) {\includegraphics[page=562,trim={53.00bp 383.39bp 53.87bp 220.05bp},clip,width=0.92\textwidth,max height=0.38\textheight]{Ferziger4th.pdf}};
\begin{scope}[x={(image.south east)},y={(image.north west)}]
\fill[white] (0.00000,0.91563) rectangle (0.09732,1.00000);
\fill[white] (0.10000,0.91563) rectangle (0.18800,1.00000);
\fill[white] (0.19074,0.91563) rectangle (0.22051,1.00000);
\end{scope}
\end{tikzpicture}
\end{sourceblock}
其中m是物体的质量，v是物体质心的速度，f是作用在物体上的合力，M是物体相对于全局（惯性）坐标系的惯性矩张量，ω是物体的角速度，m是作用在物体上的合力矩。由于在运动过程中，物体不断改变其相对于全局坐标系的方向，因此必须在每个时间步重新计算其惯性矩；对于除了最简单的体型之外的所有体型来说，这是相当不切实际的。因此，通常使用固定在物体质心处的坐标系以修改形式求解角运动方程，相对于该坐标系，转动惯量不会随着物体移动而改变（例如，参见 Shabana 2013）：

\begin{sourceblock}
\includegraphics[page=562,trim={53.00bp 220.19bp 53.87bp 417.56bp},clip,width=0.92\textwidth,max height=0.38\textheight]{Ferziger4th.pdf}
\end{sourceblock}
其中M b 代表物体固定（移动）坐标系中的惯性矩张量。

作用在物体上的力总是包括重力和流动引起的压力（垂直于物体表面）和剪切力（切向于物体表面）；此外，可能存在与流动无关的外力（例如，体内电机产生的推进力）或依赖于流动和身体运动的外力（例如，系泊线产生的力）。除了重力（根据定义，重力作用于质心，因此不产生力矩）外，所有其他力通常都会产生影响物体旋转运动的力矩。

请注意，与载有粒子的两相流的拉格朗日模型中固体粒子的运动方程相反，我们不需要在 f 中包含阻力、升力、虚拟质量或其他力，这些力可以解释流体相互作用产生的特殊效应。这些力用于模拟拉格朗日方法中的流体-颗粒相互作用，因为网格无法解析单个颗粒周围的流动。这里，直接在流固界面考虑物体和流动之间的相互作用；因此，所有影响都包含在压力和剪切力中。

方程 ( 13.67 ) 和 ( 13.69 ) 表示每个速度矢量（线性和角）的三个分量的六个 ODE 系统；因此，如果一个运动体的运动不受任何方式的约束，那么它就有六个自由度 (6 DoF)。这些方程可以使用第 1 章中描述的方法来求解。 6 获取新时间水平下的v和ω；在流动和身体运动的耦合模拟中，通常在两组方程中使用相同的时间推进方法。为了获得身体质心的新位置及其新方向，还必须对以下一组方程进行积分（这只是 v 和 ω 的定义）：

\begin{sourceblock}
\begin{tikzpicture}
\node[inner sep=0,anchor=south west] (image) at (0,0) {\includegraphics[page=563,trim={53.00bp 380.08bp 53.87bp 257.67bp},clip,width=0.92\textwidth,max height=0.38\textheight]{Ferziger4th.pdf}};
\begin{scope}[x={(image.south east)},y={(image.north west)}]
\fill[white] (0.00000,0.87531) rectangle (0.04737,1.00000);
\fill[white] (0.05008,0.87531) rectangle (0.09341,1.00000);
\end{scope}
\end{tikzpicture}
\end{sourceblock}
其中 r 是定义身体质心位置的向量，也是定义身体相对于身体固定坐标系的方向的向量。

由于流动引起的力会影响身体运动，而身体位置和方向的变化会影响流动，因此这两个问题是强耦合的（双向）。因此需要耦合求解两组动量方程。由于运动体周围的流动总是不稳定的，并且时间步长也不太大，因此使用如图 7.6 所示的顺序求解方法。然后很容易将刚体运动方程的解包含在外部迭代循环内。首先，当身体仍处于前一时间步长计算的位置时，估计新时间步长的流场。然后施加作用在身体上的估计的力来确定新的身体位置的估计（在两种情况下都可以使用欠松弛）。然后，流体区域中的网格适应新估计的身体位置，并开始下一个外循环。这些迭代持续进行，直到新时间步的估计流量产生的力和身体的估计位置的变化都不超过规定的容差。即使使用非迭代时间推进方法来计算流体流量（如第 7.2.1 节和第 7.2.2 节中描述的 PISO 或分数步法），仍然需要引入外部迭代循环，以便隐式耦合流体流动和物体运动的解，特别是当轻体在重流体中移动时。

以密度为 1 的轻体浸入密度为 1000 的液体为例；让我们假设身体通过绳子固定在底壁上，并且在模拟开始时，流体和身体都处于静止状态。垂直方向作用在物体上的力为重力 − ρ b Vg 、浮力
ρ l Vg，绳索中的约束力，等于浮力与重力之差。这里ρ b 是物体密度，ρ l 是液体密度，V 是物体体积，g 是垂直方向重力分量的大小。如果我们现在割断绳子并让身体放松，由于浮力大于重力，它就会开始向上移动。在显式求解方法的情况下，流体将在第一个时间步骤中保持静止，因为物体仍处于其旧位置。现在作用在天体上的力只是重力和浮力，导致天体加速度为 999 g（见式（13.67））：

\begin{sourceblock}
\includegraphics[page=564,trim={53.00bp 462.97bp 53.87bp 174.01bp},clip,width=0.92\textwidth,max height=0.38\textheight]{Ferziger4th.pdf}
\end{sourceblock}
这会导致身体在第一步中移动得太远；然后，流体将在第二个时间步长中与向下方向上的过大阻力发生反应。身体运动的第二个时间步将导致运动方向相反，因为流体阻力大于浮力，等等；结果是振荡背离。通过对流体流动和身体运动的隐式耦合和欠松弛，经过几次外部迭代后，可以实现一侧计算的身体加速度与另一侧的重力、浮力和阻力之和之间的平衡；这个平衡方程提供了合理的身体运动。

这种模拟中最大的挑战是流体中计算网格对身体运动的适应。与规定的身体运动的问题相反，网格只需要在每个新时间步开始时适应一次，现在需要在每个外部迭代中执行适应。这增加了对有效解的需求。最广泛使用的方法有两种：重叠网格和网格变形。在第一种方法中，一个网格附着在身体上并随之移动，没有任何变形。网格质量始终保持不变，并且由于网格质量的变化，身体周围的离散化误差不会发生变化，就像其他方法的情况一样。缺点是用于重叠和背景网格上解耦合的插值模板会发生变化，并且需要在每次外部迭代中更新。

网格变形依赖于附加（偏微分或代数）方程的解或网格顶点坐标的某种代数平滑。一种方法是将流体域视为伪固体（具有有利的特性）；由于主体运动而导致的边界顶点的运动被强加为位移的狄利克雷边界条件，并且通过求解伪固体的动量方程，该位移传播到流体域的主体中。通过改变伪固体的属性，可以实现这样的效果，例如，靠近固体壁的棱柱层区域几乎像刚体一样移动，而网格则远离边界变形。然而，初始网格只能通过变形来变形，仅适用于中等身体运动（例如，船舶在波浪中的运动）；极端的运动可能会导致网格扭曲，从而使其无法使用。原则上，当网格质量变差时，可以停止模拟，为给定的身体位置生成一个新网格，将当前和一个或多个旧解插值到新的单元质心上，然后通过将插值解作为初始条件来继续模拟。

我们现在提出两个浮动和飞行体的流动和运动耦合模拟的说明性示例。在第一个示例中，我们想要预测当船舶在头浪中以恒定速度移动时增加的阻力。事实上，船舶以（几乎）恒定的推力推进，当有波浪时，与平静水面的速度相比，前进速度会下降；然而，拖曳水池中的实验通常是通过以恒定速度拖曳船模型来进行的。要回答的主要问题是：（i）与平静水中的阻力相比，由于波浪的存在而导致的平均阻力增加了多少，以及（ii）在船上可用功率的情况下，船舶在波浪中可以达到多少速度？

当出现波浪时，网格设计必须稍微改变：在平静的水中，我们可以在远离船体的所有方向上粗化网格，特别是向前和侧面，但对于波浪，我们必须在波浪传播方向上保持每个波长的单元数量几乎恒定，并且在整个域的垂直方向上每个波浪高度也必须保持几乎恒定的单元数量。网格分辨率的变化（特别是当它们突然变化并导致波形轮廓的分辨率太粗时）可能会引起扰动，然后扰动会在整个域中进一步传播。在图13.23中，保持了波长的最小分辨率，同时网格在翼展方向上变粗，入射的长波峰不受干扰。船体是一艘集装箱船的模型（所谓的 KCS-KRISO 集装箱船；这是一艘从未建造过的船，但其模型已经在许多拖曳池中进行了昂贵的测试），约 100 米。 7.5 m 长。波长等于船长，解域长约 3.3 个波长，宽略高于 1 个波长；模拟中仅包含一半的流域，因为不存在螺旋桨，并且假设流场关于船舶对称平面对称。允许船舶起伏和纵摇；所有其他动议均被抑制。计算是在附在船体上并以恒定速度（此处：1.2 m/s）移动的坐标系中进行的。波高0.2m，波周期2.184s；使用 Fenton (1985) 提出的解将其建模为斯托克斯五阶波。通过使用重叠网格来解释船舶运动，并使用 HRIC 方案捕获自由表面（Muzaferija 和 Peri ́c 1999）。有关此模拟的更多详细信息可以在本书的网页上找到；详细内容请参见附录。

为了使解域尽可能小并避免波从边界反射，第 2 节中提出的强迫方法。使用 13.6：在距入口、侧面和出口边界 5m 的距离上，速度和水的体积分数均被强制朝向 Fenton (1985) 的理论解，使用 cos 2
强迫系数的变化，从强迫开始时的 0 到边界处的 10。船舶引起的来波扰动向各个方向传播，但在强迫区逐渐消失，RANS方程组的解顺利过渡到理论解，见图13.23。仅当参考解足够好地满足纳维-斯托克斯方程时，才能应用这种对所有边界的强制；否则，由于纳维-斯托克斯波传播之间的不匹配，强迫区内会出现扰动

\begin{sourceblock}
\includegraphics[page=566,trim={53.00bp 297.25bp 53.87bp 52.00bp},clip,width=\textwidth,max height=0.38\textheight]{Ferziger4th.pdf}
\par\vspace{0.45\baselineskip}
\small 图13.23 船模型周围未受扰动自由表面的计算网格（上）和计算的瞬时自由表面高程（下）
\end{sourceblock}
解域内的方程和理论解（基于线性理论的理论解就是这种情况，例如基于线性波叠加的长峰不规则波模型）。斯托克斯五阶理论是一个非常精确的波动模型，纳维-斯托克斯方程的解（通过数值网格对波长和波高具有足够的分辨率）与该理论非常吻合。

\begin{center}\small 图 13.24 显示了船舶在波浪中运动的两个阶段：当船头深深潜入来波波峰时，以及当波峰向船中部移动时，船完全露出水面。作用在船舶上的力及其运动如图 13.25 所示。虽然剪切力几乎保持恒定，但压力在平均值附近有很大的变化；振荡幅度比平均值大 5 倍，这导致电阻在超过三分之一的时间段内为负值！船舶纵摇±3°，垂荡±5.5cm；运动周期比波浪周期短——分别为 1.62 秒和 2.184 秒。这是因为船向\end{center}
\begin{sourceblock}
\includegraphics[page=567,trim={53.00bp 524.25bp 53.87bp 52.00bp},clip,width=\textwidth,max height=0.38\textheight]{Ferziger4th.pdf}
\par\vspace{0.45\baselineskip}
\small 图 13.24 船舶在波浪中的运动：潜入波峰（左）并在波峰到达船中部之前不久船头露出水面（右）
\end{sourceblock}
\begin{sourceblock}
\includegraphics[page=567,trim={53.00bp 261.87bp 53.87bp 165.40bp},clip,width=\textwidth,max height=0.38\textheight]{Ferziger4th.pdf}
\par\vspace{0.45\baselineskip}
\small 图 13.25 剪切和压力阻力分量（上）以及船舶升沉和纵摇运动（下）的预测时间变化
\end{sourceblock}
波，这增加了相遇频率。大约 4 个波周期后，解变得接近周期性。

第二个示例涉及救生艇进水的模拟。当救生艇从海上平台或船上释放时，它首先在空中飞行，然后潜入水中，重新浮出水面，并继续像普通船只一样依靠自身的推进力漂浮和移动。这里的关键问题是：(i) 当救生艇触水时，救生艇内的人员会经历多大的减速，(ii) 救生艇结构在入水期间将承受什么载荷，以及 (iii) 救生艇在潜入水中后何时何地重新浮出水面？这些问题的答案对于救生艇内人员的生存至关重要：（i）人体不能在太长的时间内承受太高的减速度； (ii) 入水时负载过高可能会损坏救生艇结构，从而威胁救生艇内人员的生命； (iii) 如果救生艇潜入过深且出水过晚或位置错误，可能会被坠落物体击中或与平台或船舶相撞。

主要问题是，所有问题的答案取决于太多参数，例如：落差高度、波浪传播方向、波长和波高、救生艇入水的波浪剖面上的点、风向和速度等。所有这些因素都相互影响，因此需要考虑的组合数量巨大。许多效果无法在模型实验中实现，因此最近人们对模拟的兴趣变得非常高。 Mørch (2008, 2009) 提出的关于救生艇入水的首批研究之一以及模拟结果与有限实验数据的比较表明，对于相同的波浪传播方向、波长和波高，入水期间救生艇表面的最大压力可能会变化 4 倍，具体取决于救生艇撞击波浪的位置。

救生艇入水的模拟最容易使用重叠网格来完成：在一个域中创建一个网格，例如从圆柱体中减去救生艇体，该圆柱体作为刚体随船移动，并为环境创建一个单独的网格，该网格适合于解析波浪传播，并且可能包括平台、船舶或任何其他物体。通常，救生艇的释放及其在空气中的初始运动是使用更简单的方法来模拟的；使用 CFD 的详细模拟从水面上方一定距离开始，方向、线速度和角速度取自其他模拟方法，并作为初始条件。由于救生艇非常重，其周围的气流会快速适应其运动，而不会引入显著的干扰（与救生艇的重量和初始动量相比，空气阻力相对较低）。图 13.26 展示了救生艇进入平静水域的实验研究中的两张照片，清楚地显示了船周围所产生的两相流的复杂性。幸运的是，我们不需要解决所有的流动特征（例如薄水层、它们向水滴的过渡或重新浮出水面后尾流中的泡沫区域），以便为上述问题提供相当好的答案。

\begin{sourceblock}
\includegraphics[page=568,trim={53.00bp 90.56bp 53.87bp 406.87bp},clip,width=\textwidth,max height=0.38\textheight]{Ferziger4th.pdf}
\par\vspace{0.45\baselineskip}
\small 图13.26 救生艇进水实验研究图片（来自Mørch et al. 2008）
\end{sourceblock}
\begin{sourceblock}
\includegraphics[page=569,trim={53.00bp 441.25bp 53.87bp 52.00bp},clip,width=\textwidth,max height=0.38\textheight]{Ferziger4th.pdf}
\par\vspace{0.45\baselineskip}
\small 图13.27 用于模拟救生艇入水的重叠网格
\end{sourceblock}
当救生艇向前移动时，背景网格应该时不时地在前面进行细化并在后面进行粗化，但在执行模拟时，该功能在代码中不可用。由于这个原因，背景网格被细化为与附接到救生艇的重叠网格的外层中的单元相似的单元尺寸，该重叠网格在救生艇预期在其中移动的较大区域中。此外，它还围绕自由表面进行了细化，以便捕捉救生艇进水和重新铺面时的变形。在本例中，救生艇长约15m，下落高度为36m，下水角度为35°。有关此模拟和类似模拟的更多详细信息，请参阅 Mørch 等人。 （2008）和（2009）。

救生艇在空气中自由落体时，由于空气中阻力不大，几乎以重力加速度g 加速。然而，当救生艇撞击水面时，其运动阻力突然增加，从而导致显著减速。图 13.28 显示了加速度相对于重力的变化，在救生艇内的两个位置进行评估：一处位于前部，一处位于后部（零表示加速度等于重力）。在很短的时间内，前方减速度达到5g；同时，救生艇的后部经历 3 g 的加速度，因为当船头接触到水时，救生艇开始绕着它旋转。在约0.3秒内，后部的加速度从重力方向的3g变为反方向的6g，总共产生9g的减速度。这对于正常人来说已经是一个临界值，但在一些实验中测量到高达 30 g 的减速度！即使船的结构完好无损，里面的人也不会幸存！

从图13.28可以看出，尽管模拟中使用了相对粗糙的网格（半条船大约有100万个单元；因为船只能在x和z方向自由移动并绕y轴旋转）并且水面平静，但应用了对称条件，模拟和实验之间的一致性相当好。虽然该协议在所有细节上并不完美，但很好地捕捉了进水和重新铺面的主要特征。 5秒第二次减速

\begin{sourceblock}
\includegraphics[page=570,trim={53.00bp 372.25bp 53.87bp 52.00bp},clip,width=\textwidth,max height=0.38\textheight]{Ferziger4th.pdf}
\par\vspace{0.45\baselineskip}
\small 图13.28 两时刻救生艇前部和后部加速度（上）以及救生艇表面压力和自由表面变形的预测时间变化（下）
\end{sourceblock}
因为，重新浮出水面后，船首先跳出水面，然后再次着陆，导致约1 g 的减速度。

人们可能认为当船接触水面时，船表面的最高压力会记录在船头，但事实并非如此。从图13.28中可以看出，在未受扰动的自由液面与船面相交处压力最高，并且由于进入阶段船轴与自由液面之间的夹角减小，因此此时压力沿着船从船头向船尾增大。最高压力出现在船底即将开始船尾曲率的位置；在该位置，船底上升角最小。在本例中，测量到的压力高于 5 bar；在不同的条件下可以获得更高的压力。

当船潜入水中时，其后面会形成一个气腔，因为进入水中的身体会将液体转移到侧面。当身体向前移动时，液体会返回并关闭腔体；这会在水面上产生很大的水花，同时也会在船的后部产生高压负载。 Tregde (2015) 研究了气腔的行为，发现在此类模拟中必须考虑空气压缩性的影响；将气相视为不可压缩会导致压力载荷的显著低估，而如果考虑可压缩性，则与实验的比较非常好。贝尔奇什等人。 ( 2015 ) 提出了一项广泛的验证研究，其中将自由落体救生艇发射到规则波浪中的 CFD 模拟与实验数据进行了比较。

\subsection*{13.10.2 可变形体}\phantomsection\addcontentsline{toc}{subsection}{13.10.2 可变形体}
如果与流体流相互作用的物体可能在流引起的力作用下变形，那么复杂性就达到了新的水平。如果解具有稳态（例如，当飞机在平静的空气中以恒定速度巡航时，其机翼向上弯曲并保持稳定在该位置，直到遇到湍流），则可以对流体流动和结构变形进行一系列独立的模拟。首先计算未变形物体的流体流动，直到它（几乎）收敛到稳态；计算出的压力和剪切力作为载荷传递，用于计算结构变形。然后，流体域中的网格适应变形的物体，并继续计算，直到作用在物体上的力再次收敛。重复该过程，直到流动引起的力和主体变形的变化变得小于规定限度。通常需要5到10次迭代才能得到收敛解。流动和结构变形计算的收敛容差最初可以放宽，并在接近收敛状态时收紧。

在刚才描述的场景中，流动和结构变形的模拟通常是由不同的人使用不同的代码进行的；流量是使用内部或商业代码计算的，该代码很可能基于本书中描述的有限体积方法，而结构变形很可能使用商业有限元代码计算。两个模拟团队和代码通过文件交换进行通信：流程团队写出一个文件，其中包含作用在每个墙边界面上的力，而结构团队则发回新的主体形状或一组边界点的位移。由于流体和固体中的网格通常非常不同，因此需要对解进行插值（也称为“映射”）。很多时候，甚至两个代码中的几何形状也不相同。例如，对于流体流动分析，飞机机翼的蒙皮是相关的壁边界；然而，蒙皮与结构分析几乎不相关——它只承受流体载荷，但结构相关的主体是蒙皮下的框架。因此，通过模拟沿着蒙皮的流体流动计算的载荷必须传递到框架，并且框架的变形（一组离散位置处的位移）需要映射到蒙皮表面处的流体网格的顶点。因此，耦合并不像听起来那么微不足道。

当问题是瞬态的时，流和结构计算之间需要更紧密的耦合。至少每个时间步需要发生一次力和位移的交换。然而，正如刚体运动模拟中已经提到的那样，显式耦合（其中流动求解器使用冻结体形状完成一个时间步，然后结构求解器使用冻结流动引起的力计算变形）可能并不总是收敛。理想情况下，如果要实现紧密（隐式）耦合，则交换应该在每个时间步的每次外部迭代之后发生。当使用两种不同的代码时，这是很难实现的，除非它们是为这种协同仿真而设计的（其中一个代码在前，另一个在后）；作者知道只有一套商业代码可以以这种方式进行通信。它们可以并发运行并通过内存套接字交换数据，因此无需编写任何交换文件。

理想的情况是使用相同的代码计算流体流动和结构变形。有一段时间，有限体积方法似乎适用于这两项任务。 Demirdži ́c 和 Muzaferija ( 1994 , 1995 ) 针对流体流动和结构变形提出了一种基于任意多面体控制体积和二阶离散化的有限体积方法；随后发表了许多类似的方法。 Demirdži ́c 等人。 (1997)证明，使用多重网格方法可以大大加速使用有限体积方法计算结构变形。

然而，事实证明，通常用于计算流体流动的二阶离散化（积分近似的中点规则、线性插值和梯度的中心差）不足以解决某些结构问题，特别是在某些条件下薄结构的变形（悬臂梁是一个典型的例子）。 Demirdži ́c (2016) 证明，四阶二维有限体积法（基于表面积分的辛普森规则近似、三次多项式插值和四阶中心差）在求解悬臂问题时既高效又准确。然而，开发任意多面体的高阶有限体积方法并不是一项简单的任务，与成熟的结构分析有限元方法相比可能不会带来优势；后者不对所有结构使用相同的方程，而是对壳、板、膜、梁和通用体积单元使用不同的单元（基于不同的理论）。

商业代码的最新趋势是将有限体积流求解器和有限元结构求解器集成在同一代码中；在这种情况下，可以在外部迭代级别实现无需任何外部通信的紧密耦合，从而允许同时计算流动和结构变形。这种方法的唯一问题是，现在同一个人需要知道如何设置两个模拟，或者耦合模拟的设置必须由两个人组成的团队执行——一个是流体流动专家，另一个是结构方面的专家。一旦正确设置了模拟，单个代码用户（可能是第三人）就可以轻松地执行参数研究，因为为此只需更改几何形状或输入数据。

流固耦合领域广阔；我们不能不谈一些说明性的例子。在工程应用中，人们通常对预测可能导致共振的情况感兴趣。如果流体流动以接近结构共振频率的频率在结构上产生振荡载荷，则结构可能开始以越来越大的幅度振荡，最终导致失效。一些结构对于不同的变形模式（例如弯曲和扭转）可能具有多个共振频率。需要修改流动或结构，以获得明显的流动引起的力振荡的频率和结构的固有频率之间的足够的差距。

我们在此简要描述图 13.29 所示的流体流动和可变形结构运动的耦合模拟示例；更多详细信息请参阅互联网上的另一份报告；详细内容请参见附录。该测试用例旨在用于验证计算方法，Gomes 等人对此进行了详细描述。 （2011）以及戈麦斯和林哈特（2010）。薄的，50毫米长

\begin{sourceblock}
\includegraphics[page=573,trim={53.00bp 343.25bp 53.87bp 52.00bp},clip,width=\textwidth,max height=0.38\textheight]{Ferziger4th.pdf}
\par\vspace{0.45\baselineskip}
\small 图13.29 固体结构中的装配体（上）和网格以及周围流体（下）的草图
\end{sourceblock}
0.04毫米厚的不锈钢板附在直径为22毫米的刚性圆柱体上；在片材的末端，附有长10毫米、宽4毫米的矩形刚体。圆柱体可以绕其轴线旋转。组件展向尺寸为177mm；它被放置在尺寸为 240mm × 180mm 的矩形横截面隧道中测试部分的中心，使得组件跨越整个隧道宽度，从而产生名义上的二维流动。使用以 1.08m/s 的速度流动的非常粘稠的流体（动力粘度 μ = 0 . 1722 Pa·s，密度 ρ = 1050 kg/m 3 ）；流动为层流。

该模拟是使用商业代码 Simcenter STAR-CCM+ 进行的，其中包括计算刚性部件运动和柔性板变形的有限元例程，同时使用通常的有限体积方法计算流量。外部迭代的执行方式与流体流动相同，使用 SIMPLE 算法；仅通过在结构方面添加一次迭代来扩展循环。作用在结构上的流体力和结构位移在每次外部迭代后更新。图13.29显示了所使用的网格；它在流体中是多面体，在刚性固体部分中是四面体，在柔性片材（4层）中是六面体。时间步长为 1 毫秒。

\begin{center}\small 图 13.30 显示了尾体两个极限位移时结构的形状和位置，以及附近流体中的矢量\end{center}
\begin{sourceblock}
\includegraphics[page=574,trim={53.00bp 285.25bp 53.87bp 52.00bp},clip,width=\textwidth,max height=0.38\textheight]{Ferziger4th.pdf}
\par\vspace{0.45\baselineskip}
\small 图13.30 尾体两侧位移最大时结构预测变形及其周围速度矢量
\end{sourceblock}
结构。模拟从结构与流程对齐开始；随着时间的推移，圆柱体两侧的涡流脱落导致尾体产生摆动运动。图 13.31 通过绘制圆柱体的旋转和尾部的侧向位移与时间的关系来显示这种发展。扰动呈指数增长，但 2 秒后会建立周期性状态。振荡周期和振幅与 Gomes 和 Lienhart (2010) 发表的实验数据相当吻合。

非定常流固耦合的实际例子有：风力涡轮机叶片周围的流动（见10.3.4.3节）、帆船帆周围的流动、心脏瓣膜的流动、复合材料制成的螺旋桨周围的流动等。结构的流动引起的振动也会产生噪声，这也是研究的另一个原因

\footnotetext[1]{Gilmanovetal.( 2015 ) 使用分步 LES 报告 FS 的几个问题，包括心脏瓣膜。}
\begin{sourceblock}
\includegraphics[page=575,trim={53.00bp 329.25bp 53.87bp 52.00bp},clip,width=\textwidth,max height=0.38\textheight]{Ferziger4th.pdf}
\par\vspace{0.45\baselineskip}
\small 图13.31 圆柱体旋转预测发展（上）和尾部垂直位移（下）
\end{sourceblock}
流固相互作用。如果结构在运行中发生显著变形（如飞机机翼），那么了解变形有多大就很重要，因为在制造过程中必须考虑到这一点。例如，我们可以使用与流动求解器耦合的优化软件来找到飞机机翼的最佳形状，但​​因为这应该是在运行载荷下的形状，所以我们需要找到要制造的未变形形状，这样，当机翼在流引起的力下变形时，它就达到所需的最佳形状。

\chapter*{补充资料}\addcontentsline{toc}{chapter}{补充资料}
\markboth{补充资料}{补充资料}
\section*{A.1 计算机代码列表以及如何访问它们}\phantomsection\addcontentsline{toc}{section}{A.1 计算机代码列表以及如何访问它们}
读者可以通过互联网获得体现本书中描述的一些方法的许多计算机代码。这些代码可能很有用，但它们也可以作为进一步开发的起点。它们会不时更新，并且可能会添加新代码。所有计算机代码均可从网站 \url{www.cfd-peric.de} 的下载中心获取。

包括用于求解一维和二维通用守恒方程的代码；这些被用来做章节中的例子。 3、4和6。这些代码中使用了几种对流和扩散项离散化以及时间积分的方案。它们可用于研究方案的特征，包括收敛和离散化误差以及求解器的相对效率。它们也可以用作学生作业的基础；例如，可以要求学生修改离散化方案和/或边界条件。

初始包中给出了几个求解器，包括：

\begin{itemize}
\item 用于一维问题的TDMA 求解器；
\item 二维问题的逐行TDMA 求解器（五点分子）；
\item Stone 的 ILU 求解器 (SIP)，适用于 2D 和 3D 问题（五点和七点分子；3D 版本也以矢量化形式给出）；
\item 通过不完全Cholesky 方法(ICCG) 预处理的共轭梯度求解器，适用于2 维和3 维对称矩阵（五点和七点分子）；
\item 改进的SIP 求解器，用于二维九点分子；
\item 用于非对称矩阵和3D 问题的CGSTAB 求解器；
\item 使用Gauss-Seidel、SIP 和ICCG 作为平滑器的2D 问题的多重网格求解器。
\end{itemize}
最后，有几个解决流体流动和传热问题的代码。包含以下源代码：

\begin{itemize}
\item 用于生成笛卡尔二维网格的代码；
\item 用于生成非正交结构化二维网格的代码；
\end{itemize}
\begin{itemize}
\item 用于对笛卡尔和非正交网格上的二维数据进行后处理的代码，可以在x = const 线上绘制网格、速度矢量、任意量的轮廓。或 y = const.，以及任何数量的黑白或彩色轮廓（输出是 postscript 文件）；
\item 具有交错变量排列的笛卡尔二维网格的FV 代码，用于稳态问题；
\item 使用笛卡尔二维网格和并置变量排列的 FV 代码，适用于稳定或不稳态问题；
\item 使用笛卡尔 3D 网格和同位变量排列的 FV 代码，适用于稳定和不稳态问题，并将多重网格应用于外部迭代；
\item 使用边界拟合非正交二维网格和同位变量排列的 FV 代码，适用于层流稳定或非稳态流（包括移动网格）；
\item 上述代码的版本包括带有壁函数的k – ε 和k – ω 湍流模型以及不使用壁函数的版本；
\item 上述层流代码的多重网格版本（多重网格应用于外部迭代）。
\end{itemize}
这些代码采用标准 FORTRAN77 进行编程，并已在许多计算机上进行了测试。对于较大的代码，目录中还有解释文件；每个代码中都包含许多注释行，包括如何使其适应非结构​​化网格上的 3D 问题的建议。

最后，主目录包含一个名为 errata 的文件；在其中，可能会发现的错误将被记录下来（我们希望该文件即使不为空，也将非常小）。

\section*{A.2 示例模拟的扩展报告}\phantomsection\addcontentsline{toc}{section}{A.2 示例模拟的扩展报告}
在章节中。 9 – 13，我们简要介绍了几个测试用例的结果；篇幅不允许详细描述所有相关数据。下载中心包含更详细的 pdf 格式报告（在大多数情况下还提供用于获取结果的模拟文件），读者可以下载。使用西门子的商业软件 Simcenter STAR-CCM+ 进行模拟。模拟文件可用于改变参数（几何形状、流体特性、边界条件、湍流模型等）并研究解的网格或时间步依赖性。

对于某些测试用例，还提供了流动行为的动画。

\section*{A.3 其他免费 CFD 代码}\phantomsection\addcontentsline{toc}{section}{A.3 其他免费 CFD 代码}
广泛使用并拥有庞大的社区基础。它每六个月正式发布一次，包括许多内置模型，用于计算不可压缩和可压缩流以及许多附加物理模型，如多相流、燃烧和固体应力分析。人们可以按原样使用它，但它也适合作为进一步开发数值技术或物理模型的基础。更多信息请访问 \url{www.openfoam.com}。

\chapter*{参考文献}\addcontentsline{toc}{chapter}{参考文献}
\markboth{参考文献}{参考文献}
{\small\raggedright\noindent Abe, K., Jang, Y.-J., \& Leschziner, M. A. (2003). An investigation of wall-anisotropy expressions\par}
{\small\raggedright\noindent and length-scale equations for non-linear eddy-viscosity models. International Journal of Heat Fluid Flow , 24 , 181–198. Abgrall, R. (1994). On essentially non-oscillatory schemes on unstructured meshes: Analysis and\par}
{\small\raggedright\noindent implementation. Journal of Computational Physics , 114 , 45–58. Achenbach, E. (1972). Experiments on the flow past spheres at very high Reynolds numbers. Journal\par}
{\small\raggedright\noindent of Fluid Mechanics , 54 , 565–575. Adams, M., Colell, P., Graves, D. T., Johnson, J. N., Johansen, H. S., Keen, N. D. et al. (2015).\par}
{\small\raggedright\noindent Chombo software package for AMR application design document (Technical Report). Berkeley, CA: Lawrence Berkeley National Laboratory, Applied Numerical Algorithms Group, Computational Research Division. Amezcua, J., Kalnay, E.,\& Williams, P. D. (2011). The effects of the RAW filter on the climatology\par}
{\small\raggedright\noindent and forecast skill of the SPEEDY model. Monthly Weather Review , 13 , 608–619. Arcilla, A. S., Häuser, J., Eiseman, P. R., \& Thompson, J. F. (Eds.). (1991). Numerical grid gener-\par}
{\small\raggedright\noindent ation in computational fluid dynamics and related fields . Amsterdam: North- Holland. Aris, R. (1990). Vectors, tensors and the basic equations of fluid mechanics . New York: Dover\par}
{\small\raggedright\noindent Publications. Armfield, S. (1991). Finite difference solutions of the Navier-Stokes equations on staggered and\par}
{\small\raggedright\noindent non-staggered grids. Computers Fluids , 20 , 1–17. Armfield, S. (1994). Ellipticity, accuracy, and convergence of the discrete Navier-Stokes equations.\par}
{\small\raggedright\noindent Journal of Computational Physics , 114 , 176–184. Armfield, S., \& Street, R. (1999). The fractional-step method for the Navier-Stokes equations on\par}
{\small\raggedright\noindent staggered grids: The accuracy of three variations. Journal of Computational Physics , 153 , 660– 665. Armfield, S., \& Street, R. (2000). Fractional-step methods for the Navier-Stokes equations on non-\par}
{\small\raggedright\noindent staggered grids. ANZIAM Journal , 42 (E), C134–C156. Armfield, S., \& Street, R. (2002). An analysis and comparison of the time accuracy of fractionalstep\par}
{\small\raggedright\noindent methods for the Navier-Stokes equations on staggered grids. International Journal for Numerical Methods in Fluids , 38 , 255–282. Armfield, S., \& Street, R. (2003). The pressure accuracy of fractional-step methods for the Navier-\par}
{\small\raggedright\noindent Stokes equations on staggered grids. ANZIAM Journal , 44 (E), C20–C39. Armfield, S.,\& Street, R. (2004). Modified fractional-step methods for the Navier-Stokes equations.\par}
{\small\raggedright\noindent ANZIAM Journal , 45 (E), C364–C377.\par}
{\small\raggedright\noindent Armfield, S., \& Street, R. (2005). A comparison of staggered and non-staggered grid Navier-Stokes\par}
{\small\raggedright\noindent solutions for the 8:1 cavity natural convection flow. ANZIAM Journal , 46 (E), C918–C934. Armfield, S., Williamson, N., Kirkpatrick, M., \& Street, R. (2010). A divergence free fractionalstep\par}
{\small\raggedright\noindent method for the Navier-Stokes equations on non-staggered grids. ANZIAM Journal , 51 (E), C654– C667. Aspden, A., Nikiforakis, N., Dalziel, S., \& Bell, J. (2008). Analysis of implicit LES methods.\par}
{\small\raggedright\noindent Communications in Applied Mathematics and Computational Science , 3 , 103–126. Asselin, R. (1972). Frequency filter for time integration. Monthly Weather Review , 100 , 487–490. Baki ́c, V. (2002). Experimental investigation of turbulent flows around a sphere (Ph.D. Dissertation).\par}
{\small\raggedright\noindent Germany: Technical University of Hamburg-Harburg. Baliga, B. R. (1997). Control-volume finite element method for fluid flow and heat transfer. In W. J.\par}
{\small\raggedright\noindent Minkowycz \& E. M. Sparrow (Eds.), Advances in numerical heat transfer (Vol. 1, pp. 97–135). New York: Taylor and Francis. Baliga, B. R., \& Patankar, S. V. (1983). Acontrol-volume finite element method for two-dimensional\par}
{\small\raggedright\noindent fluid flow and heat transfer. Numerical Heat Transfer , 6 , 245–261. Bardina, J., Ferziger, J. H., \& Reynolds, W. C. (1980). Improved subgrid models for large-eddy\par}
{\small\raggedright\noindent simulation. In 13th Fluid and Plasma Dynamics Conference (AIAA Paper 80-1357) . Barth, T. J., \& Jespersen, D. C. (1989). The design and application of upwind schemes on unstruc-\par}
{\small\raggedright\noindent tured meshes. In 27th Aerospace Science Meeting (AIAA Paper 89-0366) . Bastian, P., \& Horton, G. (1989). Parallelization of robust multi-grid methods: ILU factorization and\par}
{\small\raggedright\noindent frequency decomposition method. In W. Hackbusch \& R. Rannacher (Eds.), Notes on numerical fluid mechanics (Vol. 30, pp. 24–36). Braunschweig: Vieweg. Beam, R. M., \& Warming, R. F. (1976). An implicit finite-difference algorithm for hyperbolic\par}
{\small\raggedright\noindent systems in conservation-law form. J. Comput. Phys , 22 , 87–110. Beam, R. M., \& Warming, R. F. (1978). An implicit factored scheme for the compressible navier-\par}
{\small\raggedright\noindent stokes equations. AIAA Journal , 16 , 393–402. Beare, R. J., MacVean, M. K., Holtslag, A. A. M., Cuxart, J., Esau, I., Golaz, J.-C., et al. (2006).\par}
{\small\raggedright\noindent An intercomparison of large-eddy simulations of the stable boundary layer. Boundary-Layer Meteorology , 118 , 247–272. Beig, S. A., \& Johnsen, E. (2015). Maintaining interface equilibrium conditions in compressible\par}
{\small\raggedright\noindent multiphase flows using interface capturing. Journal of Computational Physics , 302 , 548–566. Berchiche, N., Östman, A., Hermundstad, O. A., \& Reinholdtsen, S.-A. (2015). Experimental\par}
{\small\raggedright\noindent validation of CFD simulations of free-fall lifeboat launches in regular waves. Ship Technology Research , 62 , 148–158. Berger, M. J., \& Oliger, J. (1984). Adaptive mesh refinement for hyperbolic partial differential\par}
{\small\raggedright\noindent equations. Journal of Computational Physics , 53 , 484–512. Bermejo-Moreno, I., Pullin, D. I., \& Horiuti, K. (2009). Geometry of enstrophy and dissipation, grid\par}
{\small\raggedright\noindent resolution effects and proximity issues in turbulence. Journal of Fluid Mechanics , 620 , 121–166. Bewley, T., Moin, P., \& Temam, R. (1994). Optimal control of turbulent channel flows. In Active\par}
{\small\raggedright\noindent control of vibration and noise (Vol. DE 75, pp. 221–227). New York: American Society of Mechanical Engineers, Design Engineering Division. Bhagatwala, A., \& Lele, S. K. (2011). Interaction of a Taylor blast wave with isotropic turbulence.\par}
{\small\raggedright\noindent Physics of Fluids , 23 , 035103. Bhaskaran, R., \& Lele, S. K. (2010). Large-eddy simulation of free-stream turbulence effects on\par}
{\small\raggedright\noindent heat transfer to a high-pressure turbine cascade. Journal of Turbulence , 11 , N6. Bijl, H., Lucor, D., Mishra, S., \& Schwab, C. (2013). Uncertainty quantification in computational\par}
{\small\raggedright\noindent fluid dynamics . Switzerland: Springer International Publishing. Bilger, R. W. (1975). A note on Favre averaging in variable density flows. Combustion Science and\par}
{\small\raggedright\noindent Technology , 11 , 215–217. Billard, F., Laurence, D., \& Osman, K. (2015). Adaptive wall functions for an elliptic blending\par}
{\small\raggedright\noindent eddy-viscosity model applicable to any mesh topology. Flow, Turbulence and Combustion , 94 , 817–842.\par}
{\small\raggedright\noindent Bini, D. A., \& Meini, D. (2009). The cyclic reduction algorithm: from Poisson equation to stochastic\par}
{\small\raggedright\noindent processes and beyond. Numerical Algorithms , 51 , 23–60. Bird, R. B., Stewart, W. E., \& Lightfoot, E. N. (2006). Transport phenomena (Revised 2 ed.). New\par}
{\small\raggedright\noindent York: Wiley. Bird, R. B., \& Wiest, J. M. (1995). Constitutive equations for polymeric liquids. Annual Review of\par}
{\small\raggedright\noindent Fluid Mechanics , 27 , 169–193. Blumberg, A. F., \& Mellor, G. L. (1987). A description of a three-dimensional coastal ocean cir-\par}
{\small\raggedright\noindent culation model. In N. S. Heaps (Ed.) Three-dimensional Coastal Ocean models, Coastal and Estuarine Science (Vol. 4, pp. 1–16). Washington, D. C.: AGU. Bodony, D. J., \& Lele, S. K. (2008). Current status of jet noise predictions using large-eddy simu-\par}
{\small\raggedright\noindent lation. AIAA J. , 46 , 364–380. Boris, J. P., \& Book, D. L. (1973). Flux-corrected transport. i. SHASTA, a fluid transport algorithm\par}
{\small\raggedright\noindent that works. Journal of Computational Physics , 11 , 38–69. Bose, S. T., Moin, P., \& You, D. (2010). Grid-independent large-eddy simulation using explicit\par}
{\small\raggedright\noindent filtering. Physics of Fluids , 22 , 105103. Boyd, J. P. (2001). Chebyshev and Fourier spectral methods (Revised 2 ed.). Mineola: Dover\par}
{\small\raggedright\noindent Publications. Brackbill, J. U., Kothe, D. B., \& Zemaach, C. (1992). A continuum method for modeling surface\par}
{\small\raggedright\noindent tension. Journal of Computational Physics , 100 , 335–354. Bradshaw, P., Launder, B. E., \& Lumley, J. L. (1994). Collaborative testing of turbulence models.\par}
{\small\raggedright\noindent In K. N. Ghia, U. Ghia, \& D. Goldstein (Eds.), Advances in computational fluid mechanics (Vol. 196). New York: ASME. Brandt, A. (1984). Multigrid techniques: 1984 guide with applications to fluid dynamics. GMDStu-\par}
{\small\raggedright\noindent dien Nr. 85, Gesellschaft für Mathematik und Datenverarbeitung (GMD). Bonn, Germany (see also Multigrid Classics version at \url{http://www.wisdom.weizmann.ac.il/~achi/} ). Brazier, P. H. (1974). An optimum SOR procedure for the solution of elliptic partial differential\par}
{\small\raggedright\noindent equations with any domain or coefficient set. Computer Methods in Applied Mechanics and Engineering , 3 , 335–347. Brehm, C., Housman, J. A., Kiris, C. C., Barad, M. F., \& Hutcheson, F. V. (2017). Four-jet impinge-\par}
{\small\raggedright\noindent ment: Noise characteristics and simplified acoustic model. International Journal of Heat and Fluid flow , 67 , 43–58. Brès, G. A., Ham, F. E., Nichols, J. W., \& Lele, S. K. (2017). Unstructured large-eddy simulations\par}
{\small\raggedright\noindent of supersonic jets. AIAA Journal , 55 , 1164–1184. Briggs, D. R., Ferziger, J. H., Koseff, J. R., \& Monismith, S. G. (1996). Entrainment in a shear free\par}
{\small\raggedright\noindent mixing layer. Journal of Fluid Mechanics , 310 , 215–241. Briggs, W. L., Henson, V. E., \& McCormick, S. F. (2000). A multigrid tutorial (2nd ed.). Philadel-\par}
{\small\raggedright\noindent phia: Society for Industrial and Applied Mathematics (SIAM). Brigham, E. O. (1988). The fast Fourier transform and its applications . Englewood Cliffs: Prentice\par}
{\small\raggedright\noindent Hall. Bryan, G. H. (2007). A comparison of convection resolving-simulations with convection-\par}
{\small\raggedright\noindent permitting simulations. NSSL Colloquium, Norman, OK. \url{http://www2.mmm.ucar.edu/people/-} bryan/Presentations/bryan\_2007\_nssl\_resolution.pdf . Bryan, G. H., Wyngaard, J. C., \& Fritsch, J. M. (2003). Resolution requirements for the simulation\par}
{\small\raggedright\noindent of deep moist convection. Monthly Weather Review , 131 , 2394–2416. Bückle, U., \& Peri ́c, M. (1992). Numerical simulation of buoyant and thermocapillary convection\par}
{\small\raggedright\noindent in a square cavity. Numerical Heat Transfer, Part A (Applications) , 21 , 101–121. Bunner, B., \& Tryggvason, G. (1999). Direct numerical simulations of three-dimensional bubbly\par}
{\small\raggedright\noindent flows. Physics of Fluids , 11 , 1967–1969. Burmeister, J., \& Horton, G. (1991). Time-parallel solution of the Navier-Stokes equations. In\par}
{\small\raggedright\noindent Proceedings of 3rd European Multigrid Conference . Basel: Birkhäuser. Butcher, J. C. (2008). Numerical methods for ordinary differential equations . Chichester: Wiley. Calhoun, D. (2002). A cartesian grid method for solving the two-dimensional streamfunctionvor-\par}
{\small\raggedright\noindent ticity equations in irregular regions. Journal of Computational Physics , 176 , 231–275.\par}
{\small\raggedright\noindent Canuto, C., Hussaini, M. Y., Quarteroni, A.,\& Zang, T. A. (2006). Spectral methods: Fundamentals\par}
{\small\raggedright\noindent in single domains . Berlin: Springer. Canuto, C., Hussaini, M. Y., Quarteroni, A., \& Zang, T. A. (2007). Spectral methods: Evolution to\par}
{\small\raggedright\noindent complex geometries and applications to fluid dynamics . Berlin: Springer. Carati, D., Winckelmans, G. S., \& Jeanmart, H. (2001). On the modelling of the subgrid-scale and\par}
{\small\raggedright\noindent filtered-scale stress tensors in large-eddy simulation. Journal of Fluid Mechanics , 441 , 119–138. Caretto, L. S., Gosman, A. D., Patankar, S. V., \& Spalding, D. B. (1972). Two calculation procedures\par}
{\small\raggedright\noindent for steady three-dimensional flows with recirculation. In Proceedings of the Third International Conference on Numerical Methods in Fluid Mechanics . Paris. Caruso, S. C., Ferziger, J. H., \& Oliger, J. (1985). An adaptive grid method for incompressible flows\par}
{\small\raggedright\noindent (Technical Report No. TF-23). Stanford CA: Department of Mechanical Engineering, Stanford University. Casulli, V., \& Cattani, E. (1994). Stability, accuracy and efficiency of a semi-implicit method for\par}
{\small\raggedright\noindent three-dimensional shallow water flow. Computers \& Mathematics with Applications , 27 , 99–112. Cebeci, T., \& Bradshaw, P. (1984). Physical and computational aspects of convective heat transfer .\par}
{\small\raggedright\noindent New York: Springer. Celik, I., Ghia, U., Roache, P. J., \& Freitas, C. J. (2008). Procedure for estimation and reporting of\par}
{\small\raggedright\noindent uncertainty due to discretization in cfd applications. Journal of Fluids Engineering , 130 , 078001. Celik, I., Klein, M., \& Janicka, J. (2009). Assessment measures for engineering LES applications.\par}
{\small\raggedright\noindent Journal of Fluids Engineering , 131 , 031102. Chang, W., Giraldo, F., \& Perot, B. (2002). Analysis of an exact fractional-step method. Journal of\par}
{\small\raggedright\noindent Computational Physics , 180 , 183–199. Chen, C.-J., \& Jaw, S.-Y. (1998). Fundamentals of turbulence modeling . Washington: Taylor \&\par}
{\small\raggedright\noindent Francis. Chen, S., Johnson, D. B., Raad, P. E., \& Fadda, D. (1997). The surface marker and micro-cell\par}
{\small\raggedright\noindent method. International Journal for Numerical Methods in Fluids , 25 , 749–778. Chen, C., Zhu, J., Zheng, L., Ralph, E., \& Budd, J. W. (2004a). A non-orthogonal primitive equation\par}
{\small\raggedright\noindent coastal ocean circulation model: Application to lake superior. Journal of Great Lakes Research , 30 (Supplement 1), 41–54. Chen, Y., Ludwig, F. L., \& Street, R. L. (2004b). Stably stratified flows near a notched transverse\par}
{\small\raggedright\noindent ridge across the salt lake valley. Journal of Applied Meteorology , 43 , 1308–1328. Choi, H., \& Moin, P. (1994). Effects of the computational time step on numerical solutions of\par}
{\small\raggedright\noindent turbulent flow. Journal of Computational Physics , 113 , 1–4. Choi, H., Moin, P., \& Kim, J. (1994). Active turbulence control for drag reduction in wall-bounded\par}
{\small\raggedright\noindent flows. Journal of Fluid Mechanics , 262 , 75–110. Chorin,A.J.(1967).Anumericalmeth\hspace{0pt}odforsolvingin\hspace{0pt}compressiblevi\hspace{0pt}scousflowprobl\hspace{0pt}ems. Journal\par}
{\small\raggedright\noindent of Computational Physics , 2 , 12–26. Chorin, A. J. (1968). Numerical solution of the Navier-Stokes equations. Mathematics of Compu-\par}
{\small\raggedright\noindent tation , 22 , 745–762. Chow, F. K., \& Moin, P. (2003). A further study of numerical errors in large-eddy simulations.\par}
{\small\raggedright\noindent Journal of Computational Physics , 184 , 366–380. Chow, F. K., \& Street, R. L. (2009). Evaluation of turbulence closure models for large-eddy sim-\par}
{\small\raggedright\noindent ulation over complex terrain: Flow over Askervein Hill. Journal of Applied Meteorology and Climatology , 48 , 1050–1065. Chow, F. K., Street, R. L., Xue, M., \& Ferziger, J. (2005). Explicit filtering and reconstruction\par}
{\small\raggedright\noindent turbulence modeling for large-eddy simulation of neutral boundary layer flow. Journal of the Atmospheric Sciences , 62 , 2058–2077. Chow, F. K., Weigel, A. P., Street, R. L., Rotach, M. W., \& Xue, M. (2006). High-resolution large-\par}
{\small\raggedright\noindent eddy simulations of flow in a steep Alpine valley. part i: Methodology, verification, and sensitivity experiments. Journal of Applied Meteorology and Climatology , 45 , 63–86. Colella, P. (1985). A direct Eulerian MUSCL scheme for gas dynamics. SIAM Journal on Scientific\par}
{\small\raggedright\noindent and Statistical Computing , 6 , 104–117.\par}
{\small\raggedright\noindent Colella, P. (1990). Multidimensional upwind methods for hyperbolic conservation laws. Journol of\par}
{\small\raggedright\noindent Computational Physics , 87 , 171–200. Coelho, P., Pereira, J. C. F., \& Carvalho, M. G. (1991). Calculation of laminar recirculating flows\par}
{\small\raggedright\noindent using a local non-staggered grid refinement system. International Journal for Numerical Methods in Fluids , 12 , 535–557. Constantinescu, G., \& Squires, K. (2004). Numerical investigations of flow over a sphere in the\par}
{\small\raggedright\noindent subcritical and supercritical regimes. Physics Fluids , 16 , 1449–1466. Cooley, J. W., \& Tukey, J. W. (1965). An algorithm for the machine calculation of complex Fourier\par}
{\small\raggedright\noindent series. Mathematics of Computation , 19 , 297–301. Courant, R., Friedrichs, K., \& Lewy, H. (1928). Über die partiellen Differenzengleichungen der\par}
{\small\raggedright\noindent mathematischen Physik. Mathematische Annalen , 100 , 32–74. Crowe, C., Sommerfeld, M., \& Tsuji, Y. (1998). Multiphase flows with droplets and particles . Boca\par}
{\small\raggedright\noindent Raton: CRC Press. Deike, L., Melville, W. K., \& Popinet, S. (2016). Air entrainment and bubble statistics in breaking\par}
{\small\raggedright\noindent waves. Journal of Fluid Mechanics , 801 , 91–129. Demirdži ́c, I. (2015). On the discretization of the diffusion term in finite-volume continuum mechan-\par}
{\small\raggedright\noindent ics. Numerical Heat Transfer , 68 , 1–10. Demirdži ́c, I. (2016). A fourth-order finite volume method for structural analysis. Applied Mathe-\par}
{\small\raggedright\noindent matical Modelling , 40 , 3104–3114. Demirdži ́c, I., \& Peri ́c, M. (1988). Space conservation law in finite volume calculations of fluid\par}
{\small\raggedright\noindent flow. International Journal for Numerical Methods in Fluids , 8 , 1037–1050. Demirdži ́c, I., \& Peri ́c, M. (1990). Finite volume method for prediction of fluid flow in arbitrarily\par}
{\small\raggedright\noindent shaped domains with moving boundaries. International Journal for Numerical Methods in Fluids , 10 , 771–790. Demirdži ́c, I., \& Muzaferija, S. (1994). Finite volume method for stress analysis in complex\par}
{\small\raggedright\noindent domains. International Journal for Numerical Methods in Engineering , 37 , 3751–3766. Demirdži ́c, I., \& Muzaferija, S. (1995). Numerical method for coupled fluid flow, heat transfer\par}
{\small\raggedright\noindent and stress analysis using unstructured moving meshes with cells of arbitrary topology. Computer Methods in Applied Mechanics and Engineering , 125 , 235–255. Demirdži ́c, I., Muzaferija, S., \& Peric, M. (1997). Benchmark solutions of some structural anal-\par}
{\small\raggedright\noindent ysis problems using finite-volume method and multigrid acceleration. International Journal for Numerical Methods in Engineering , 40 , 1893–1908. Demmel, J. W., Heath, M. T., \& van der Vorst, H. A. (1993). Parallel numerical linear algebra. Acta\par}
{\small\raggedright\noindent Numerica , 2 , 111–197. Demirdži ́c, I., \& Lilek, V., \& Peric, M. (1993). A colocated finite volume method for predicting\par}
{\small\raggedright\noindent flows at all speeds. International Journal for Numerical Methods in Fluids , 16 , 1029–1050. Deng, G. B., Piquet, J., Queutey, P., \& Visonneau, M. (1994). Incompressible flow calculations with\par}
{\small\raggedright\noindent a consistent physical interpolation finite volume approach. Computers Fluids , 23 , 1029–1047. de Vahl Davis, G., \& Mallinson, G. D. (1972). False diffusion in numerical fluid mechanics. Univ\par}
{\small\raggedright\noindent (Technical Report No. FM1). New South Wales, Australia: University of New South Wales, School of Mechanical Ind. Engineering Donea, J., \& Huerta, A. (2003). Finite element methods for flow problems . Chichester: Wiley.\par}
{\small\raggedright\noindent (available at Wiley Online Library). Donea, J., Huerta, A., Ponthot, J.-P.,\& Rodríguez-Ferran, A. (2004). Arbitrary Lagrangian-Eulerian\par}
{\small\raggedright\noindent methods. In Encyclopedia of Computational Mechanics. Vol. 1: Fundamentals (pp. 413–437). Donea, J., Quartapelle, L., \& Selmin, V. (1987). An analysis of time discretization in the finite\par}
{\small\raggedright\noindent element solution of hyperbolic problems. Journal of Computational Physics , 70 , 463–499. Durbin, P. A. (1991). Near-wall turbulence closure modeling without ‘damping functions’. Theo-\par}
{\small\raggedright\noindent retical and Computational Fluid Dynamics , 3 , 1–13. Durbin, P. A. (2002). A perspective on recent developments in RANS modeling. In W. Rodi \& N.\par}
{\small\raggedright\noindent Furyo (Eds.), Engineering urbulence modelling and experiments (Vol. 5, pp. 3–16). Durbin, P. A. (2009). Limiters and wall treatments in applied turbulence modeling. Fluid Dynamics\par}
{\small\raggedright\noindent Research , 41 , 012203.\par}
{\small\raggedright\noindent Durbin, P. A., \& Pettersson Reif, B. A. (2011). Statistical theory and modeling for turbulent flows\par}
{\small\raggedright\noindent (2nd ed.). Chichester: Wiley. Durran, D. R. (2010). Numerical methods for fluid dynamics with applications to geophysics (2nd\par}
{\small\raggedright\noindent ed.). Berlin: Springer. Durst, F., Kadinskii, L., Peri ́c, M., \& Schäfer, M. (1992). Numerical study of transport phenomena\par}
{\small\raggedright\noindent in MOCVD reactors using a finite volume multigrid solver. Journal of Crystal Growth , 125 , 612–626. ElKhoury,G.K.,Schlatter,P.,Noorani,A.,Fischer,P.F.,Brethouwer,G.,\&Johansson,A.V.(2013).\par}
{\small\raggedright\noindent Direct numerical simulation of turbulent pipe flow at moderately high Reynolds numbers. Flow, Turbulence and Combustion , 91 , 475–495. Enright, D., Fedkiw, R., Ferziger, J. H., \& Mitchell, I. (2002). A hybrid particle level set method\par}
{\small\raggedright\noindent for improved interface capturing. Journal of Computational Physics , 183 , 83–116. Enriquez, R. M. (2013). Subgrid-scale turbulence modeling for improved large-eddy simulation of\par}
{\small\raggedright\noindent the atmospheric boundary layer (Ph.D. Dissertation) . Stanford: Stanford University. Enriquez, R. M., Chow, F. K., Street, R. L., \& Ludwig, F. L. (2010). Examination of the linear\par}
{\small\raggedright\noindent algebraic subgrid-scale stress [LASS] model, combined with reconstruction of the subfilterscale stress, for large-eddy simulation of the neutral atmospheric boundary layer. In 19th Conference on Boundary Layers and Turbulence, AMS, Paper 3A . (8 pages). Erturk, E. (2009). Discussions on driven cavity flow. International Journal for Numerical Methods\par}
{\small\raggedright\noindent in Fluids , 60 , 275–294. Fadlun, E. A., Verzicco, R., Orlandi, P., \& Mohd-Yusof, J. (2000). Combined immersed-boundary\par}
{\small\raggedright\noindent finite-difference methods for three-dimensional complex flowsimulations. Journal of Computational Physics , 161 , 35–60. Farmer, J., Martinelli, L., \& Jameson, A. (1994). Fast multigrid method for solving incompressible\par}
{\small\raggedright\noindent hydrodynamic problems with free surfaces. AIAA Journol , 32 , 1175–1182. Fenton, J. D. (1985). A fifth-order Stokes theory for steady waves. Journal of Waterway, Port,\par}
{\small\raggedright\noindent Coastal, Ocean Engineering , 111 , 216–234. Ferziger, J. H. (1998). Numerical methods for engineering application (2nd ed.). New York: Wiley-\par}
{\small\raggedright\noindent Interscience. Ferziger, J. H., \& Peri ́c, M. (1996). Further discussion of numerical errors in CFD. International\par}
{\small\raggedright\noindent Journal for Numerical Methods in Fluids , 23 , 1–12. Findikakis, A. N., \& Street, R. L. (1979). An algebraic model for subgrid-scale turbulence in\par}
{\small\raggedright\noindent stratified flows. Journal of the Atmospheric Sciences , 36 , 1934–1949. Fletcher, C. A. J. (1991). Computational techniques for fluid dynamics (2nd ed., Vol. I \& II). Berlin:\par}
{\small\raggedright\noindent Springer. Fletcher, R. (1976). Conjugate gradient methods for indefinite systems. Lecture Notes in Mathe-\par}
{\small\raggedright\noindent matics , 506 , 773–789. Forbes, L. K. (1988). Critical free surface flow over a semicircular obstruction. Journal of Engi-\par}
{\small\raggedright\noindent neering Mathematics , 22 , 3–13. Fornberg, B. (1988). Generation of finite difference formulas on arbitrarily spaced grids. Mathe-\par}
{\small\raggedright\noindent matics of Computation , 51 , 699–706. Fox, D. G., \& Orszag, S. A. (1973). Pseudospectral approximation to two-dimensional turbulence.\par}
{\small\raggedright\noindent Journal of Computational Physics , 11 , 612–619. Freitas, C. J., Street, R. L., Findikakis, A. N., \& Koseff, J. R. (1985). Numerical simulation of\par}
{\small\raggedright\noindent three-dimensional flow in a cavity. International Journal for Numerical Methods in Fluids , 5 , 561–575. Frey, P. J., \& George, P.-L. (2008). Mesh generation: Application to finite elements (2nd ed.). New\par}
{\small\raggedright\noindent Jersey: Wiley-ISTE. Friedrich, O. (1998). Weighted essentially non-oscillatory schemes for the interpolation of mean\par}
{\small\raggedright\noindent values on unstructured grids. Journal of Computational Physics , 144 , 194–212. Fringer, O. B., Armfield, S. W., \& Street, R. L. (2003). A nonstaggered curvilinear grid pressure\par}
{\small\raggedright\noindent correction method applied to interfacial waves. In Second International Conference on Heat Transfer, Fluid Mechanics and Thermodynamics, HEFAT 2003, Paper FO1 . (6 pages).\par}
{\small\raggedright\noindent Fringer, O. B., Armfield, S. W., \& Street, R. L. (2005). Reducing numerical diffusion in interfacial\par}
{\small\raggedright\noindent gravity wave simulations. International Journal for Numerical Methods in Fluids , 49 , 301–329. Fringer, O. B., Gerritsen, M., \& Street, R. L. (2006). An unstructured-grid, finite-volume, nonhy-\par}
{\small\raggedright\noindent drostatic, parallel coastal ocean simulator. Ocean Modelling , 14 , 139–173. Fu, L., Hu, X. Y., \& Adams, N. A. (2016). A family of high-order targeted ENO schemes for\par}
{\small\raggedright\noindent compressible-fluid simulations. Journal of Computational Physics , 305 , 333–359. Fu, L., Hu, X. Y., \& Adams, N. A. (2017). Targeted ENO schemes with tailored resolution property\par}
{\small\raggedright\noindent for hyperbolic conservation laws. Journal of Computational Physics , 349 , 97–121. Galpin, P. F., \& Raithby, G. D. (1986). Numerical solution of problems in incompressible fluid flow:\par}
{\small\raggedright\noindent Treatment of the temperature-velocity coupling. Numerical Heat Transfer , 10 , 105–129. Gamet, L., Ducros, F., Nicoud, F., \& Poinsot, T. (1999). Compact finite difference schemes on non-\par}
{\small\raggedright\noindent uniform meshes. Application to direct numerical simulations of compressible flows. International Journal for Numerical Methods in Fluids , 29 , 159–191. Garnier, E., Adams, N., \& Sagaut, P. (2009). Large-eddy simulation for compressible flows . Berlin:\par}
{\small\raggedright\noindent Springer. Germano, M., Piomelli, U., Moin, P., \& Cabot, W. H. (1991). A dynamic subgrid-scale eddyviscosity\par}
{\small\raggedright\noindent model. Physics of Fluids A , 3 , 1760–1765. Ghia, U., Ghia, K. N., \& Shin, C. T. (1982). High-Re solutions for incompressible flow using the\par}
{\small\raggedright\noindent Navier-Stokes equations and a multigrid method. Journal of Computational Physics , 48 , 387–411. Gibbs, J. W. (1898). Fourier’s series. Nature , 59 , 200. Gibbs, J. W. (1899). Fourier’s series. Nature , 59 , 606. Gilmanov, A., Le, T. B., \& Sotiropoulos, F. (2015). A numerical approach for simulating fluid\par}
{\small\raggedright\noindent structure interaction of flexible thin shells undergoing arbitrarily large deformations in complex domains. Journal of Computational Physics , 300 , 814–843. Glowinski, R., \& Pironneau, O. (1992). Finite element methods for Navier-Stokes equations. Annual\par}
{\small\raggedright\noindent Review of Fluid Mechanics , 24 , 167–204. Göddeke, D., \& Strzodka, R. (2011). Cyclic reduction tridiagonal solvers on GPUs applied to\par}
{\small\raggedright\noindent mixed-precision multigrid. IEEE Transactions on Parallel and Distributed Systems , 22 , 22–32. Golub, G. H., \& van Loan, C. F. (1996). Matrix computations (3rd ed.). Baltimore: Johns Hopkins\par}
{\small\raggedright\noindent University Press. Gomes, J. P., \& Lienhart, H. (2010). Fluid-structure interaction-induced oscillation of flexible\par}
{\small\raggedright\noindent structures in laminar and turbulent flows. Journal of Fluid Mechanics , 715 , 537–572. Gomes, J. P., Yigit, S., Lienhart, H., \& Schäfer, M. (2011). Experimental and numerical study on a\par}
{\small\raggedright\noindent laminar fluid-structure interaction reference test case. Journal of Fluids Structures , 27 , 43–61. Gottlieb, S., Mullen, J. S., \& Ruuth, S. J. (2006). A fifth order flux implicit WENO method. Journal\par}
{\small\raggedright\noindent of Scientific Computing , 27 , 271–287. Gresho, P. M. (1990). On the theory of semi-implicit projection methods for viscous incompressible\par}
{\small\raggedright\noindent flow and its implementation via a finite element method that also introduces a nearly consistent mass matrix. Part 1: Theory. International Journal for Numerical Methods in Fluids , 11 , 587–620. Grinstein, F. F., Margolin, L. G., \& Rider, W. J. (Eds.). (2007). Implicit large-eddy simulation:\par}
{\small\raggedright\noindent Computing turbulent fluid dynamics . Cambridge: Cambridge University Press. Gullbrand, J., \& Chow, F. K. (2003). The effect of numerical errors and turbulence models in large-\par}
{\small\raggedright\noindent eddy simulations of channel flow, with and without explicit filtering. Journal of Fluid Mechanics , 495 , 323–341. Gyllenram,W., \& Nilsson, H. (2006). Very large-eddy simulation of draft tube flow. In 23rd IAHR\par}
{\small\raggedright\noindent Symposium Yokohama . (10 pages). Hackbusch, W. (1984). Parabolic multi-grid methods. In R. Glowinski \& J.-R. Lions (Eds.), Com-\par}
{\small\raggedright\noindent puting methods in applied sciences and engineering . Amsterdam: North Holland. Hackbusch, W. (2003). Multi-grid methods and applications (2nd Printing) . Berlin: Springer. Hackbusch, W., \& Trottenberg, U. (Eds.). (1991). In Proceedings of Third European Multigrid\par}
{\small\raggedright\noindent Conference, International Series of Numerical Mathematics . Basel: Birkhäuser.\par}
{\small\raggedright\noindent Hadži ́c, H. (2005). Development and application of a finite volume method for the computation of\par}
{\small\raggedright\noindent flows around moving bodies on unstructured, overlapping grids (Ph.D. Dissertation). Technische Universität , Hamburg-Harburg. Hadži ́c, I. (1999). Second-moment closure modelling of transitional and unsteady turbulent flows\par}
{\small\raggedright\noindent (Ph.D. Dissertation). Delft University of Technology. Hageman, L. A., \& Young, D. M. (2004). Applied iterative methods . Mineola: Dover Publications. Ham, F., \& Iaccarino, G. (2004). Energy conservation in collocated discretization schemes on\par}
{\small\raggedright\noindent unstructured grids. In Annals Research Briefs . Stanford, CA: Center for Turbulence Research. Hanaoka, A. (2013). An overset grid method coupling an orthogonal curvilinear grid solver and a\par}
{\small\raggedright\noindent Cartesian grid solver (Ph.D. Dissertation) . Iowa City, IA: University of Iowa. Hanjali ́c, K. (2002). One-point closure models for buoyancy-driven turbulent flows. Annual Review\par}
{\small\raggedright\noindent of Fluid Mechanics , 34 , 321–347. Hanjali ́c, K. (2004). Closure models for incompressible turbulent flows (Technical Report). Brus-\par}
{\small\raggedright\noindent sels, Belgium: Lecture Notes at the von Karman Institute for Fluid Dynamics. Hanjali ́c, K., \& Kenjereš, S. (2001). T-RANS simulation of deterministic eddy structure in flows\par}
{\small\raggedright\noindent driven by thermal buoyancy and Lorentz force. Flow Turbulence and Combustion , 66 , 427–451. Hanjali ́c, K., \& Launder, B. E. (1976). Contribution towards a Reynolds-stress closure for low\par}
{\small\raggedright\noindent Reynolds number turbulence. Journal of Fluid Mechanics , 74 , 593–610. Hanjali ́c, K., \& Launder, B. E. (1980). Sensitizing the dissipation equation to irrotational strains.\par}
{\small\raggedright\noindent Journal of Fluids Engineering , 102 , 34–40. Harlow, F. H., \& Welsh, J. E. (1965). Numerical calculation of time dependent viscous Incompress-\par}
{\small\raggedright\noindent ible flow with free surface. Physics of Fluids , 8 , 2182–2189. Harrison, R. J. (1991). Portable tools and applications for parallel computers. International Journal\par}
{\small\raggedright\noindent of Quantum Chemistry , 40 , 847–863. Harten, A. (1983). High resolution schemes for hyperbolic conservation laws. Journal of Compu-\par}
{\small\raggedright\noindent tational Physics , 49 , 357–393. Harvie, D. J. E., Davidson, M. R.,\& Rudman, M. (2006). An analysis of parasitic current generation\par}
{\small\raggedright\noindent in Volume-of-Fluid simulations. Applied Mathematical Modelling , 30 , 1056–1066. Hatlee, S. C., \& Wyngaard, J. C. (2007). Improved subfilter-scale models from the HATS field data.\par}
{\small\raggedright\noindent Journal of the Atmospheric Sciences , 64 , 1694–1705. Heil, M., Hazel, A. L.,\& Boyle, J. (2008). Solvers for large-displacement fluid-structure interaction\par}
{\small\raggedright\noindent problems: segregated versus monolithic approaches. Computational Mechanics , 43 , 91–101. Heinke, H. J. (2011). Potsdam propellet test case (Technical Report No. 3753). Potsdam, Germany:\par}
{\small\raggedright\noindent SVA Potsdam. Hess, J. L. (1990). Panel methods in computational fluid dynamics. Annual Review of Fluid Mechan-\par}
{\small\raggedright\noindent ics , 22 , 255–274. Heus, T., van Heerwaarden, C. C., Jonker, H. J. J., Siebesma, A. P. et al. (2010). Formulation of the\par}
{\small\raggedright\noindent dutch atmospheric large-eddy simulation (DALES) and overview of its applications. Geoscientific Model Development , 3 , 415–444. Hinatsu, M., \& Ferziger, J. H. (1991). Numerical computation of unsteady incompressible flow in\par}
{\small\raggedright\noindent complex geometry using a composite multigrid technique. International Journal for Numerical Methods in Fluids , 13 , 971–997. Hino, T. (1992). Computation of viscous flows with free surface around an advancing ship. In\par}
{\small\raggedright\noindent Proceedings of 2nd Osaka International Colloquium on Viscous Fluid Dynamics in Ship and Ocean Technology . Osaka University. Hirsch,C.(2007). Numericalcompu\hspace{0pt}tationofintern\hspace{0pt}alandexternalf\hspace{0pt}lows (2nded.,Vol.I).Burlington:\par}
{\small\raggedright\noindent Butterworth-Heinemann (Elsevier). Hirt, C. W., Amsden, A. A., \& Cook, J. L. (1997). An arbitrary Lagrangean-Eulerian computing\par}
{\small\raggedright\noindent method for all flow speeds. Journal of Computational Physics , 135 , 203–216. (Reprinted from 14, 1974, 227–253). Hirt, C. W., \& Harlow, F. H. (1967). A general corrective procedure for the numerical solution of\par}
{\small\raggedright\noindent initial-value problems. Journal of Computational Physics , 2 , 114–119.\par}
{\small\raggedright\noindent Hirt, C. W., \& Nichols, B. D. (1981). Volume of fluid (VOF) method for dynamics of free boundaries.\par}
{\small\raggedright\noindent Journal of Computational Physics , 39 , 201–221. Hodges, B. R., \& Street, R. L. (1999). On simulation of turbulent nonlinear free-surface flows.\par}
{\small\raggedright\noindent Journal of Computational Physics , 151 , 425–457. Hortmann, M., Peri ́c, M., \& Scheuerer, G. (1990). Finite volume multigrid prediction of laminar nat-\par}
{\small\raggedright\noindent ural convection: Bench-mark solutions. International Journal for Numerical Methods in Fluids , 11 , 189–207. Horton, G. (1991). Ein zeitparalleles Lösungsverfahren für die Navier-Stokes-Gleichungen (Ph.D.\par}
{\small\raggedright\noindent Dissertation) . Erlangen-Nürnberg: Universität. Housman, J. A., Stich, G.-D., \& Kiris, C. C. (2017). Jet noise prediction using hybrid RANS/LES\par}
{\small\raggedright\noindent with structured overset grids. In 23rd AIAA/CEAS Aeroacoustics Conference, AIAA Paper 2017- 3213 . Hu, F. Q., Hussaini, M. Y., \& Manthey, J. L. (1996). Low-dissipation and low-dispersion Runge-\par}
{\small\raggedright\noindent Kutta schemes for computational acoustics. Journal of Computational Physics , 124 , 177–191. Hubbard, B. J., \& Chen, H. C. (1994). A Chimera scheme for incompressible viscous flows with\par}
{\small\raggedright\noindent applications to submarine hydrodynamics. In 25th AIAA Fluid Dynamics Conference (AIAA Paper 94-2210) . Hubbard, B. J., \& Chen, H. C. (1995). Calculations of unsteady flows around bodies with relative\par}
{\small\raggedright\noindent motion using a Chimera RANS method. In Proceedings of 10th ASCE Engineering Mechanics Conference . Boulder, CO: University of Colorado at Boulder. Hutchinson, B. R., Galpin, P. F., \& Raithby, G. D. (1988). Application of additive correction multi-\par}
{\small\raggedright\noindent grid to the coupled fluid flow equations. Numerical Heat Transfer , 13 , 133–147. Hutchinson, B. R., \& Raithby, G. D. (1986). A multigrid method based on the additive correction\par}
{\small\raggedright\noindent strategy. Numerical Heat Transfer , 9 , 511–537. Hylla,E.A.(2013). EineImmersedBo\hspace{0pt}undaryMethodez\hspace{0pt}urSimulationvo\hspace{0pt}nStrömungeninkomplexen\par}
{\small\raggedright\noindent und bewegten Geometrien (Ph.D. Dissertation) . Berlin, Germany: Technische Universität Berlin. Iaccarino,G.,Ooi,A.,Durbin,P.A.,\&Behnia,M.(2003).Reynoldsaverag\hspace{0pt}edsimulationof\hspace{0pt}unsteady\par}
{\small\raggedright\noindent separated flow. International Journal of Heat and Fluid Flow , 24 , 147–156. Ishihara, T., Gotoh, T., \& Kaneda, Y. (2009). Study of high-Reynolds number isotropic turbulence\par}
{\small\raggedright\noindent by direct numerical simulation. Annual Review of Fluid Mechanics , 41 , 165–180. Ishii, M. (1975). Thermo-fluid dynamic theory of two-phase flow . Paris: Eyrolles. Ishii, M., \& Hibiki, T. (2011). Thermo-fluid dynamics of two-phase flow . New York: Springer. Issa, R. I. (1986). Solution of implicitly discretized fluid flow equations by operator-splitting.\par}
{\small\raggedright\noindent Journal of Computational Physics , 62 , 40–65. Issa, R. I., \& Lockwood, F. C. (1977). On the prediction of two-dimensional supersonic viscous\par}
{\small\raggedright\noindent interaction near walls. AIAA Jorunal , 15 , 182–188. Ivanell,S.,Sørensen,J.,Mikkelsen,R.,\&Henningson,D.(2009).Analysisofnume\hspace{0pt}ricallygenerat\hspace{0pt}ed\par}
{\small\raggedright\noindent wake structures. Wind Energy , 12 , 63–80. Jacobson, M. Z., \& Delucchi, M. A. (2009). A path to sustainable energy by 2030. Scientific\par}
{\small\raggedright\noindent American , 301 , 58–65. Jakirlic, S., \& Jovanovic, J. (2010). On unified boundary conditions for improved predictions of\par}
{\small\raggedright\noindent near-wall turbulence. Journal of Fluid Mechanics , 656 , 530–539. Jakobsen, H. A. (2003). Numerical convection algorithms and their role in Eulerian CFD reactor\par}
{\small\raggedright\noindent simulations. International Journal of Chemical Reactor Engineering , 1 , Art. A1, 15. Johnsen, E., \& Ham, F. (2012). Preventing numerical errors generated by interface-capturing\par}
{\small\raggedright\noindent schemes in compressible multi-material flows. Journal of Computational Physics , 231 , 5705– 5717. Jovanovi ́c, J. (2004). The statistical dynamics of turbulence . Berlin: Springer. Kader, B. A. (1981). Temperature and concentration profiles in fully turbulent boundary layers.\par}
{\small\raggedright\noindent International Journal of Heat Mass Transfer , 24 , 1541–1544. Kadinski, L., \& Peri ́c, M. (1996). Numerical study of grey-body surface radiation coupled with\par}
{\small\raggedright\noindent fluid flow for general geometries using a finite volume multigrid solver. International Journal of Numerical Methods for Heat and Fluid Flow , 6 , 3–18.\par}
{\small\raggedright\noindent Kaneda, Y., \& Ishihara, T. (2006). High-resolution direct numerical simulation of turbulence. Jour-\par}
{\small\raggedright\noindent nal of Turbulence , 7 , N20. Kang, S., Iaccarino, G., Ham, F., \& Moin, P. (2009). Prediction of wall-pressure fluctuation in\par}
{\small\raggedright\noindent turbulent flows with an immersed boundary method. Journal of Computational Physics , 228 , 6753–6772. Karki, K. C., \& Patankar, S. V. (1989). Pressure based calculation procedure for viscous flows at\par}
{\small\raggedright\noindent all speeds in arbitrary configurations. AIAA Journal , 27 , 1167–1174. Kawai, S., \& Lele, S. K. (2010). Large-eddy simulation of jet mixing in supersonic crossflows.\par}
{\small\raggedright\noindent AIAA Journal , 48 , 2063–2083. Kawamura, T., \& Miyata, H. (1994). Simulation of nonlinear ship flows by density-function method.\par}
{\small\raggedright\noindent Journal of the Society of Naval Architects of Japan , 176 , 1–10. Kays, W. M., \& Crawford, M. E. (1978). Convective heat and mass transfer . New York: McGraw-\par}
{\small\raggedright\noindent Hill. Kenjereš, S., \& Hanjali ́c, K. (1999). Transient analysis of Rayleigh-Bénard convection with a RANS\par}
{\small\raggedright\noindent model. International Journal of Heat and Fluid Flow , 20 , 329–340. Kenjereš, S., \& Hanjalic, K. (2002). Combined effects of terrain orography and thermal stratification\par}
{\small\raggedright\noindent on pollutant dispersion in a town valley: A T-RANS simulation. Journal of Turbulence , 3 , N26. Khajeh-Saeed, A., \& Perot, J. B. (2013). Direct numerical simulation of turbulence using GPU\par}
{\small\raggedright\noindent accelerated supercomputers. Journal of Computational Physics , 235 , 241–257. Khani, S., \& Porté-Agel, F. (2017). A modulated-gradient parameterization for the large-eddy\par}
{\small\raggedright\noindent simulation of the atmospheric boundary layer using the weather research and forecasting model. Boundary-Layer Meteorology , 165 , 385–404. Khosla, P. K., \& Rubin, S. G. (1974). A diagonally dominant second-order accurate implicit scheme.\par}
{\small\raggedright\noindent Computers Fluids , 2 , 207–209. Kim, J. W. (2007). Optimised boundary compact finite difference schemes for computational aeroa-\par}
{\small\raggedright\noindent coustics. Journal of Computational Physics , 225 , 995–1019. Kim, D., \& Choi, H. (2000). A second-order time-accurate finite volume method for unsteady\par}
{\small\raggedright\noindent incompressible flow on hybrid unstructured grids. Journal of Computational Physics , 162 , 411– 428. Kim, J., Kim, D., \& Choi, H. (2001). An immersed-boundary finite-volume method for simulations\par}
{\small\raggedright\noindent of flow in complex geometries. Journal of Computational Physics , 171 , 132–150. Kim, J., \& Moin, P. (1985). Application of a fractional-step method to incompressible Navier-Stokes\par}
{\small\raggedright\noindent equations. Journal of Computational Physics , 59 , 308–323. Kim, J., Moin, P., \& Moser, R. D. (1987). Turbulence statistics in fully developed channel flow at\par}
{\small\raggedright\noindent low Reynolds number. Journal of Fluid Mechanics , 177 , 133–166. Kim, S., Kinnas, S. A., \& Du, W. (2018). Panel method for ducted propellers with sharp trailing\par}
{\small\raggedright\noindent edge duct with fully aligned wake on blade and duct. Journal of Marine Science and Engineering , 6 , 6030089. Kirkpatrick, M. P., \& Armfield, S. W. (2008). On the stability and performance of the\par}
{\small\raggedright\noindent projection- 3 method for the time integration of the Navier-Stokes equations. ANZIAM Journal , 49 (EMAC2007), C559–C575. Klemp, J. B., Skamarock, W. C., \& Dudhia, J. (2007). Conservative split-explicit time integration\par}
{\small\raggedright\noindent methods for the compressible nonhydrostatic equations. Monthly Weather Review , 135 , 2897– 2913. Knight, D. D. (2006). Elements of numerical methods for compressible flows . Cambridge: Cam-\par}
{\small\raggedright\noindent bridge University Press. Konan, N. A., Simonin, O., \& Squires, K. D. (2011). Detached-eddy simulations and particle\par}
{\small\raggedright\noindent Lagrangian tracking of horizontal rough wall turbulent channel flow. Journal of Turbulence , 12 , N22. Kordula, W., \& Vinokur, M. (1983). Efficient computation of volume in flow predictions. AIAA\par}
{\small\raggedright\noindent Journal , 21 , 917–918. Koshizuka, S., Tamako, H., \& Oka, Y. (1995). A particle method for incompressible viscous flow\par}
{\small\raggedright\noindent with fluid fragmentation. Computational Fluid Dynamics Journal , 4 , 29–46.\par}
{\small\raggedright\noindent Kosovi ́c, B. (1997). Subgrid-scale modelling for the large-eddy simulation of high-Reynoldsnumber\par}
{\small\raggedright\noindent boundary layers. Journal of Fluid Mechanics , 336 , 151–182. Kundu, P. K., \& Cohen, I. M. (2008). Fluid mechanics (4th ed.). Burlington: Academic (Elsevier). Kuzmin, D., Löhner, R., \& Turek, S. (Eds.). (2012). Flux-corrected transport: Principles, algo-\par}
{\small\raggedright\noindent rithms, and applications (2nd ed.). Dordrecht: Springer. Kwak, D., Chang, J. L. C., Shanks, S. P., \& Chakravarthy, S. R. (1986). A three-dimensional\par}
{\small\raggedright\noindent incompressible Navier-Stokes flow solver using primitive variables. AIAA Journal , 24 , 390–396. Kwak, D., \& Kiris, C. C. (2011). Artificial compressibility method. In Computation of viscous\par}
{\small\raggedright\noindent incompressible flows . Dordrecht: Springer. Lafaurie, B., Nardone, C., Scardovelli, R., Zaleski, S., \& Zanetti, G. (1994). Modelling merging\par}
{\small\raggedright\noindent and fragmentation in multiphase flows with SURFER. Journal of Computational Physics , 113 , 134–147. Lardeau, S. (2018). Consistent strain/stress lag eddy-viscosity model for hybrid RANS/LES. In\par}
{\small\raggedright\noindent Y. Hoarau, S. H. Peng, D. Schwamborn, \& A. Revell (Eds.), Progress in Hybrid RANS-LES Modelling (pp. 39–51). Cham: Springer. Launder, B. E., Reece, G. J., \& Rodi, W. (1975). Progress in the development of a Reynolds-stress\par}
{\small\raggedright\noindent turbulence closure. Journal of Fluid Mechanics , 68 , 537–566. Launder, B. E., \& Spalding, D. B. (1974). The numerical computation of turbulent flows. Computer\par}
{\small\raggedright\noindent Methods in Applied Mechanics and Engineering , 3 , 269–289. Le Bras, S., Deniau, H., Bogey, C., \& Daviller, G. (2017). Development of compressible large-eddy\par}
{\small\raggedright\noindent simulations combining high-order schemes and wall modeling. AIAA Journal , 55 , 1152–1163. Leder,A.(1992). AbgelösteStrömungen.PhysikalischeGrundlagen .Wiesbaden,Germany:Vieweg. Lee, M., \& Moser, R. D. (2015). Direct numerical simulation of turbulent channel flow up to\par}
{\small\raggedright\noindent Re τ ≈ 5200. Journal of Fluid Mechanics , 774 , 395–415. Leister, H.-J., \& Peri ́c, M. (1994). Vectorized strongly implicit solving procedure for seven-diagonal\par}
{\small\raggedright\noindent coefficient matrix. International Journal of Numerical Methods for Heat and Fluid Flow , 4 , 159– 172. Lele, S. J. (1992). Compact finite difference schemes with spectral-like resolution. Journal of\par}
{\small\raggedright\noindent Computational Physics , 3 , 16–42. Leonard, B. P. (1979). A stable and accurate convection modelling procedure based on quadratic\par}
{\small\raggedright\noindent upstream interpolation. Computer Methods in Applied Mechanics and Engineering , 19 , 59–98. Leonard, B. P. (1997). Bounded higher-order upwind multidimensional finite-volume convection-\par}
{\small\raggedright\noindent diffusion algorithms. In W. J. Minkowycz \& E. M. Sparrow (Eds.), Advances in numerical heat transfer (pp. 1–57). New York: Taylor and Francis. Leonard, A.,\& Wray, A. A. (1982). A new numerical method for the simulation of three dimensional\par}
{\small\raggedright\noindent flow in a pipe. In E. Krause (Ed.), Eighth International Conference on Numerical Methods in Fluid Dynamics . Lecture Notes in Physics (Vol. 170). Berlin: Springer. Leonardi, S., Orlandi, P., Smalley, R. J., Djenidi, L., \& Antonia, R. A. (2003). Direct numerical\par}
{\small\raggedright\noindent simulations of turbulent channel flow with transverse square bars on one wall. Journal of Fluid Mechanics , 491 , 229–238. Leschziner, M. A. (2010). Reynolds-averaged Navier-Stokes methods. In R. Blockley \& W. Shyy\par}
{\small\raggedright\noindent (Eds.), Encyclopedia of aerospace engineering (pp. 1–13). Lesieur, M. (2010). Two-point closure based on large-eddy simulations in turbulence, Part 2: Inho-\par}
{\small\raggedright\noindent mogeneous cases. Discrete \& Continuous Dynamical Systems , 28 . Lesieur, M. (2011). Two-point closure based on large-eddy simulations in turbulence, Part 1:\par}
{\small\raggedright\noindent Isotropic turbulence. Discrete \& Continuous Dynamical System Series S , 4 , 155–168. Li, G., \& Xing, Y. (1967). High order finite volume WENO schemes for the Euler equations under\par}
{\small\raggedright\noindent gravitational fields. Journal of Computational Physics , 316 , 145–163. Lilek, Ž. (1995). Ein Finite-Volumen Verfahren zur Berechnung von inkompressiblen und kompress-\par}
{\small\raggedright\noindent iblen Strömungen in komplexen Geometrien mit beweglichen Rändern und freien Oberflächen (Ph.D. Dissertation) . Germany: University of Hamburg. Lilek, Ž., Muzaferija, S., \& Peri ́c, M. (1997a). Efficiency and accuracy aspects of a full-multigrid\par}
{\small\raggedright\noindent SIMPLE algorithm for three-dimensional flows. Numerical Heat Transfer, Part B , 31 , 23–42.\par}
{\small\raggedright\noindent Lilek, Ž., Muzaferija, S., Peri ́c, M., \& Seidl, V. (1997b). An implicit finite-volume method using\par}
{\small\raggedright\noindent non-matching blocks of structured grid. Numerical Heat Transfer, Part B , 32 , 385–401. Lilek, Ž., Nadarajah, S., Peric, M., Tindal, M., \& Yianneskis, M. (1991). Measurement and simula-\par}
{\small\raggedright\noindent tion of the flow around a poppet valve. In Proceedings of 8th Symposium Turbulent Shear Flows (pp. 13.2.1–13.2.6). Lilek, Ž., \& Peri ́c, M. (1995). A fourth-order finite volume method with colocated variable arrange-\par}
{\small\raggedright\noindent ment. Computers Fluids , 24 , 239–252. Lilek, Ž., Schreck, E., \& Peri ́c, M. (1995). Parallelization of implicit methods for flow simulation. In\par}
{\small\raggedright\noindent S. G. Wagner (Ed.), Notes on numerical fluid mechanics (Vol. 50, pp. 135–146). Braunschweig: Vieweg. Lilly, D. K. (1992). A proposed modification of the Germano subgrid-scale closure method. Physics\par}
{\small\raggedright\noindent of Fluids A , 4 , 633–635. Liu, X.-D., Osher, S., \& Chan, T. (1994). Weighted essentially non-oscillatory schemes. Journal of\par}
{\small\raggedright\noindent Computational Physics , 115 , 200–212. Louda, P., Kozel, K., \& Pˇríhoda, J. (2008). Numerical solution of 2D and 3D viscous incompressible\par}
{\small\raggedright\noindent steady and unsteady flows using artificial compressibility method. International Journal for Numerical Methods in Fluids , 56 , 1399–1407. Lu, H., \& Porté-Agel, F. (2011). Large-eddy simulation of a very large wind farm in a stable\par}
{\small\raggedright\noindent atmospheric boundary layer. Physics of Fluids , 23 , 065101. Ludwig, F. L., Chow, F. K., \& Street, R. L. (2009). Effect of turbulence models and spatial resolution\par}
{\small\raggedright\noindent onresolvedvelo\hspace{0pt}citystructurea\hspace{0pt}ndmomentumflux\hspace{0pt}esinlarge-eddysimulation\hspace{0pt}sofneutralboun\hspace{0pt}dary layer flow. Journal of Applied Meteorology and Climatology , 48 , 1161–1180. Lumley, J. L. (1979). Computational modeling of turbulent flows. Advances in Applied Mechanics ,\par}
{\small\raggedright\noindent 18 , 123–176. Lundquist, K. A., Chow, F. K., \& Lundquist, J. K. (2012). An immersed boundary method enabling\par}
{\small\raggedright\noindent large-eddy simulations of flow over complex terrain in the WRF model. Monthly Weather Review , 140 , 3936–3955. MacCormack, R. W. (2003). The effect of viscosity in hypervelocity impact cratering. Journal of\par}
{\small\raggedright\noindent Spacecraft and Rockets , 40 , 757–763. (Reprinted from AIAA Paper 69–354, 1969). Mahesh, K. (1998). A family of high order finite difference schemes with good spectral resolution.\par}
{\small\raggedright\noindent Journal of Computational Physics , 145 , 332–358. Mahesh, K., Constantinescu, G., \& Moin, P. (2004). A numerical method for large-eddy simulation\par}
{\small\raggedright\noindent in complex geometries. Journal of Computational Physics , 197 , 215–240. Malinen, M. (2012). The development of fully coupled simulation software by reusing segregated\par}
{\small\raggedright\noindent solvers. Application of Parallel and Science Computing, Part 1, PARA 2010. LNCS , 7133 , 242– 248. Maliska, C. R., \& Raithby, G. D. (1984). A method for computing three-dimensional flows using\par}
{\small\raggedright\noindent non-orthogonal boundary-fitted coordinates. International Journal for Numerical Methods in Fluids , 4 , 518–537. Manhart, M., \& Wengle, H. (1994). Large-eddy simulation of turbulent boundary layer over a\par}
{\small\raggedright\noindent hemisphere.InP.Voke,L.Kleiser,\&J.P.Chollet(Eds.), Proceedingsof1stERCOFTACWorkshop on Direct and Large Eddy Simulation (pp. 299–310). Dordrecht: Kluwer Academic Publishers. Marstorp,L.,Brethouwer,G.,Grundestam,O.,\&Johansson,A.V.(2009).Explicitalgebr\hspace{0pt}aicsubgrid\par}
{\small\raggedright\noindent stress models with application to rotating channel flow. Journal of Fluid Mechanics , 639 , 403– 432. Mason, M. L., Putnam, L. E.,\& Re, R. J. (1980). The effect of throat contouring on two-dimensional\par}
{\small\raggedright\noindent converging-diverging nozzle at static conditions (p. 1704). Paper No: NASA Techn. Masson, C., Saabas, H. J., \& Baliga, R. B. (1994). Co-located equal-order control-volume finite ele-\par}
{\small\raggedright\noindent ment method for two-dimensional axisymmetric incompressible fluid flow. International Journal for Numerical Methods in Fluids , 18 , 1–26. Matheou, G., \& Chung, D. (2014). Large-eddy simulation of stratified turbulence. Part II: Applica-\par}
{\small\raggedright\noindent tion of the stretched-vortex model to the atmospheric boundary layer. Journal of the Atmospheric Sciences , 71 , 4439–4460.\par}
{\small\raggedright\noindent McCormick, S. F. (Ed.). (1987). Multigrid methods . Philadelphia: Society for Industrial and Applied\par}
{\small\raggedright\noindent Mathematics (SIAM). McMurtry, P. A., Jou, W. H., Riley, J. J., \& Metcalfe, R. W. (1986). Direct numerical simulations\par}
{\small\raggedright\noindent of a reacting mixing layer with chemical heat release. AIAA Journal , 24 , 962–970. Mellor, G. L., \& Yamada, T. (1982). Development of a turbulence closure model for geophysical\par}
{\small\raggedright\noindent fluid problems. Review of Geophysics , 20 , 851–875. Meneveau, C., \& Katz, J. (2000). Scale-invariance and turbulence models for large-eddy simulation.\par}
{\small\raggedright\noindent Annual Review of Fluid Mechanics , 32 , 1–32. Meneveau, C., Lund, T. S., \& Cabot, W. H. (1996). A Lagrangian dynamic subgrid-scale model of\par}
{\small\raggedright\noindent turbulence. Journal of Fluid Mechanics , 319 , 353–385. Menter, F. R. (1994). Two-equation eddy-viscosity turbulence models for engineering applications.\par}
{\small\raggedright\noindent AIAA Journal , 32 , 1598–1605. Menter, F. R., Kuntz, M., \& Langtry, R. (2003). Ten years of industrial experience with the SST\par}
{\small\raggedright\noindent turbulence model. In K. Hanjali ́c, Y. Nagano, \& M. Tummers (Eds.), Turbulence, heat and mass transfer (Vol. 4, pp. 625–632). (Proceedings of 4th International Symposium on Turbulent for Heat and Mass Transfer, Begell House, Inc) Mesinger, F., \& Arakawa, A. (1976). Numerical methods used in atmospheric models . GARP\par}
{\small\raggedright\noindent Publications Series No. 17 (Vol. 1). Geneva: World Meteorological Organization Meyers, J., Geurts, B. J., \& Sagaut, P. (2007). A computational error-assessment of central finitevol-\par}
{\small\raggedright\noindent ume discretizations in large-eddy simulation using a Smagorinsky model. Journal of Computational Physics , 227 , 156–173. Michioka, T., \& Chow, F. K. (2008). High-resolution large-eddy simulations of scalar transport\par}
{\small\raggedright\noindent in atmospheric boundary layer flow over complex terrain. Journal of Applied Meteorology and Climatology , 47 , 3150–3169. Mittal, R., \& Iaccarino, G. (2005). Immersed boundary methods. Annual Review of Fluid Mechanics ,\par}
{\small\raggedright\noindent 37 , 239–261. Moeng, C.-H. (1984). A large-eddy-simulation model for the study of planetary boundary-layer\par}
{\small\raggedright\noindent turbulence. Journal of the Atmospheric Sciences , 41 , 2052–2062. Moeng, C.-H., \& Sullivan, P. P. (2015). Large-eddy simulation. Encyclopedia of atmospheric sci-\par}
{\small\raggedright\noindent ences (2nd ed., Vol. 4, pp. 232–240). Cambridge: Academic. Moin, P. (2010). Fundamentals of engineering numerical analysis (2nd ed.). Cambridge: Cambridge\par}
{\small\raggedright\noindent University Press. Moin, P., \& Kim, J. (1982). Numerical investigation of turbulent channel flow. Journal of Fluid\par}
{\small\raggedright\noindent Mechanics , 118 , 341–377. Moin, P., \& Mahesh, K. (1998). Direct numerical simulation: A tool in turbulence research. Annual\par}
{\small\raggedright\noindent Review of Fluid Mechanics , 30 , 539–578. Moin, P., Squires, K., Cabot, W., \& Lee, S. (1991). Adynamic subgrid-scale model for compressible\par}
{\small\raggedright\noindent turbulence and scalar transport. Physics of Fluids A , 3 , 2746–2757. Mørch, H. J., Enger, S., Peri ́c, M., \& Schreck, E. (2008). Simulation of lifeboat launching under\par}
{\small\raggedright\noindent storm conditions. In 6th International Conference on CFD in Oil and Gas, Metallurgical and Process Industries . Trondheim, Norway. Mørch, H. J., Peri ́c, M., Schreck, E., el Moctar, O., \& Zorn, T. (2009). Simulation of flow and motion\par}
{\small\raggedright\noindent of lifeboats. In ASME 28th International Conference on Ocean, Offshore and Arctic Engineering . Honolulu, Hawaii. Morrison, H., \& Pinto, J. O. (2005). Intercomparison of bulk cloud microphysics schemes in\par}
{\small\raggedright\noindent mesoscale simulations of springtime arctic mixed-phase stratiform clouds. Mon. Wea. Rev. , 134 , 1880–1900. Mortazavi, M., Le Chenadec, V., Moin, P., \& Mani, A. (2016). Direct numerical simulation of a\par}
{\small\raggedright\noindent turbulent hydraulic jump: Turbulence statistics and air entrainment. Journal of Fluid Mechanics , 797 , 60–94. Moser, R. D., Moin, P., \& Leonard, A. (1983). A spectral numerical method for the Navier-Stokes\par}
{\small\raggedright\noindent equations with applications to Taylor-Couette flow. Journal of Computational Physics , 52 , 524– 544.\par}
{\small\raggedright\noindent Muzaferija, S. (1994). Adaptive finite volume method for flow predictions using unstructured\par}
{\small\raggedright\noindent meshes and multigrid approach (Ph.D. Dissertation). University of London. Muzaferija, S., \& Gosman, A. D. (1997). Finite-volume CFD procedure and adaptive error control\par}
{\small\raggedright\noindent strategy for grids of arbitrary topology. Journal of Computational Physics , 138 , 766–787. Muzaferija, S., \& Peri ́c, M. (1997). Computation of free-surface flows using finite volume method\par}
{\small\raggedright\noindent and moving grids. Numerical Heat Transfer, Part B , 32 , 369–384. Muzaferija, S., \& Peri ́c, M. (1999). Computation of free surface flows using interface-tracking and\par}
{\small\raggedright\noindent interface-capturing methods. In O. Mahrenholtz \& M. Markiewicz (Eds.), Nonlinear water wave interaction (pp. 59–100). Southampton: WIT Press. NASA CGTUM. (2010). Chimera grid tools user’s manual, ver. 2.1. NASA Advanced Supercom-\par}
{\small\raggedright\noindent puting Division. Retrieved from \url{https://www.nas.nasa.gov/publications/software/-docs/chimera/} index.html . NASA TMR. (n.d.). Turbulence modeling resource. Langley Research Center. Retrieved from\par}
{\small\raggedright\noindent \url{https://turbmodels.larc.nasa.gov/index.html} . NSF. (2006). Simulation-based engineering science . Retrieved from \url{http://www.nsf.gov/-pubs/}\par}
{\small\raggedright\noindent reports/sbes\_final\_report.pdf . Oden, J. T. (2006). Finite elements of non-linear continua . Mineola: Dover Publications. Orlandi, P., \& Leonardi, S. (2008). Direct numerical simulation of three-dimensional turbulent\par}
{\small\raggedright\noindent rough channels: Parameterization and flow physics. Journal of Fluid Mechanics , 606 , 399–415. Osher, S., \& Fedkiw, R. (2003). Level set methods and dynamic implicit surfaces . New York:\par}
{\small\raggedright\noindent Springer. Osher, S., \& Sethian, J. A. (1988). Fronts propagating with curvature-dependent speed: Algorithms\par}
{\small\raggedright\noindent based on Hamilton-Jacobi formulations. Journal of Computational Physics , 79 , 12–49. Pal, A., Sarkar, S., Posa, A., \& Balaras, E. (2017). Direct numerical simulation of stratified flow\par}
{\small\raggedright\noindent past a sphere at a subcritical Reynolds number of 3700 and moderate Froude number. Journal of Fluid Mechanics , 826 , 5–31. Parrott, A. K., \& Christie, M. A. (1986). FCT applied to the 2-D finite element solution of tracer\par}
{\small\raggedright\noindent transport by single phase flow in a porous medium. In K. W. Morton \& M. J. Baines (Eds.), Proceedings of ICFD-conference on Numerical Methods in Fluid Dynamics (p. 609ff). Oxford: Oxford University Press. Pascau, A. (2011). Cell face velocity alternatives in a structured colocated grid for the unsteady\par}
{\small\raggedright\noindent Navier-Stokes equations. International Journal for Numerical Methods in Fluids , 65 , 812–833. Patankar, S. V. (1980). Numerical heat transfer and fluid flow . New York: McGraw-Hill. Patankar, S. V., \& Spalding, D. B. (1972). A calculation procedure for heat, mass and momentum\par}
{\small\raggedright\noindent transfer in three-dimensional parabolic flows. International Journal of Heat and Mass Transfer , 15 , 1787–1806. Patankar, S. V., \& Spalding, D. B. (1977). Genmix: A general computer program for twodimensional\par}
{\small\raggedright\noindent parabolic phenomena . Oxford: Pergamon Press. Patel, V. C., Rodi, W., \& Scheuerer, G. (1985). Turbulence models for near-wall and low-Reynolds\par}
{\small\raggedright\noindent number flows: A review. AIAA Journal , 23 , 1308–1319. Pekurovsky, D., Yeung, P. K., Donzis, D., Pfeiffer, W., \& Chukkapallli, G. (2006). Scalability of a\par}
{\small\raggedright\noindent pseudospectral DNS turbulence code with 2D domain decomposition on Power41/Federation and Blue Gene systems. In ScicomP12 and SP-XXL . Boulder, CO: International Business Machines. Retrieved from \url{http://www.spscicomp.org/ScicomP12/-Presentations/User/Pekurovsky.pdf} . Peller, N. (2010). Numerische Simulation turbulenter Strömungen mit Immersed Boundaries (Ph.D.\par}
{\small\raggedright\noindent Dissertation) . Fachgebiet Hydromechanik, Mitteilungen: Technische Universität München. Peri ́c, M. (1985). A finite volume method for the prediction of three-dimensional fluid flow in\par}
{\small\raggedright\noindent complex ducts (Ph.D. Dissertation). Imperial College, London. Peri ́c, M. (1987). Efficient semi-implicit solving algorithm for nine-diagonal coefficient matrix.\par}
{\small\raggedright\noindent Numerical Heat Transfer , 11 , 251–279. Peri ́c, M. (1990). Analysis of pressure-velocity coupling on non-orthogonal grids. Numerical Heat\par}
{\small\raggedright\noindent Transfer. Part B (Fundamentals) , 17 , 63–82.\par}
{\small\raggedright\noindent Peri ́c, M. (1993). Natural convection in trapezoidal cavities. Numerical Heat Transfer. Part A (Appli-\par}
{\small\raggedright\noindent cations) , 24 , 213–219. Peri ́c, M., \& Schreck, E. (1995). Analysis of efficiency of implicit CFD methods on MIMD com-\par}
{\small\raggedright\noindent puters. In Proceedings of Parallel CFD ’95 Conference . Peri ́c, R. (2019). Minimierung unerwünschter Wellenreflexionen an den Gebietsrändern bei Strö-\par}
{\small\raggedright\noindent mungssimulationen mit Forcing Zones (Ph.D. Dissertation) . Germany: Technische Universität Hamburg. Peri ́c, R., \& Abdel-Maksoud, M. (2018). Analytical prediction of reflection coefficients for wave\par}
{\small\raggedright\noindent absorbing layers in flow simulations of regular free-surface waves. Ocean Engineering , 47 , 132– 147. Perktold, G., \& K. and Rappitsch.,. (1995). Computer simulation of local blood flow and vessel\par}
{\small\raggedright\noindent mechanics in a compliant carotid artery bifurcation model. J. Biomechanics , 28 , 845–856. Perng, C. Y., \& Street, R. L. (1991). A coupled multigrid-domain-splitting technique for simulating\par}
{\small\raggedright\noindent incompressible flows in geometrically complex domains. International Journal for Numerical Methods in Fluids , 13 , 269–286. Peskin, C. S. (1972). Flow patterns around heart valves: A digital computer method for solving\par}
{\small\raggedright\noindent the equations of motion (Ph.D. Dissertation) . Albert Einstein College of Medicine, Yeshiva University. Peskin, C. S. (2002). The immersed boundary method. Acta Numerica , 11 , 479–517. Peters, N. (2000). Turbulent combustion . Cambridge: Cambridge University Press. Piomelli, U. (2008). Wall-layer models for large-eddy simulations. Progress in Aerospace Sciences ,\par}
{\small\raggedright\noindent 44 , 437–446. Piomelli, U., \& Balaras, E. (2002). Wall-layer models for large-eddy simulations. Annual Review\par}
{\small\raggedright\noindent of Fluid Mechanics , 34 , 349–374. Pirozzoli, S. (2011). Numerical methods for high-speed flows. Annual Review of Fluid Mechanics ,\par}
{\small\raggedright\noindent 43 , 163–194. Poinsot, T., Veynante, D., \& Candel, S. (1991). Quenching processes and premixed turbulent com-\par}
{\small\raggedright\noindent bustion diagrams. Journal of Fluid Mechanics , 228 , 561–605. Pope, S. B. (2000). Turbulent flows . Cambridge: Cambridge University Press. Popovac, M., \& Hanjali ́c, K. (2007). Compound wall treatment for RANS computation of complex\par}
{\small\raggedright\noindent turbulent flows and heat transfer. Flow, Turbulence and Combustion , 78 , 177–202. Porté-Agel, F., Meneveau, C., \& Parlange, M. B. (2000). A scale-dependent dynamic model for\par}
{\small\raggedright\noindent large-eddy simulation: Application to a neutral atmospheric boundary layer. Journal of Fluid Mechanics , 415 , 261–284. Press, W. H., Teukolsky, S. A., Vettering, W. T., \& Flannery, B. P. (2007). Numerical recipes: The\par}
{\small\raggedright\noindent art of scientific computing (3rd ed.). Cambridge: Cambridge University Press. Pritchard, P. J. (2010). Fox and McDonald’s introduction to fluid mechanics (8th ed.). New York:\par}
{\small\raggedright\noindent Wiley. Purser, R. J. (2007). Accuracy considerations of time-splitting methods for models using two-\par}
{\small\raggedright\noindent timelevel schemes. Monthly Weather Review , 135 , 1158–1164. Qin, Z., Delaney, K., Riaz, A., \& Balaras, E. (2015). Topology preserving advection of implicit\par}
{\small\raggedright\noindent interfaces on Cartesian grids. Journal of Computational Physics , 290 , 219–238. Raithby, G. D. (1976). Skew upstream differencing schemes for problems involving fluid flow.\par}
{\small\raggedright\noindent Computer Methods in Applied Mechanics and Engineering , 9 , 153–164. Raithby, G. D., \& Schneider, G. E. (1979). Numerical solution of problems in incompressible fluid\par}
{\small\raggedright\noindent flow: Treatment of the velocity-pressure coupling. Numerical Heat Transfer , 2 , 417–440. Raithby, G. D., Xu, W.-X., \& Stubley, G. D. (1995). Prediction of incompressible free surface\par}
{\small\raggedright\noindent flows with an element-based finite volume method. Computational Fluid Dynamics Journal , 4 , 353–371. Rakhimov, A. C., Visser, D. C., \& Komen, E. M. J. (2018). Uncertainty quantification method for\par}
{\small\raggedright\noindent CFD applied to the turbulent mixing of two water layers. Nuclear Engineering and Design , 333 , 1–15.\par}
{\small\raggedright\noindent Ramachandran, S., \& Wyngaard, J. C. (2010). Subfilter-scale modelling using transport equations:\par}
{\small\raggedright\noindent Large-eddy simulation of the moderately convective atmospheric boundary layer. Boundary- Layer Meteorology ,. \url{https://doi.org/10.1007/s10546-010-9571-3} . Rasam, A., Brethouwer, G., \& Johansson, A. V. (2013). An explicit algebraic model for the subgrid-\par}
{\small\raggedright\noindent scale passive scalar flux. Journal of Fluid Mechanics , 721 , 541–577. Raw, M. J. (1985). A new control-volume-based finite element procedure for the numerical solu-\par}
{\small\raggedright\noindent tion of the fluid flow and scalar transport equations (Ph.D. Dissertation) . Waterloo, Canada: University of Waterloo Raw, M. J. (1995). A coupled algebraic multigrid method for the 3D Navier-Stokes equations. In\par}
{\small\raggedright\noindent W. Hackbusch \& G. Wittum (Eds.), Fast solvers for flow problems, notes on numerical fluid mechanics (Vol. 49, pp. 204–215). Braunschweig: Vieweg. Reichardt, H. (1951). Vollständige Darstellung der turbulenten Geschwindigkei\hspace{0pt}tsverteilung in glat-\par}
{\small\raggedright\noindent ten Leitungen. Z. Angew Mathematics and Mechanics , 31 , 208–219. Reinecke, M., Hillebrandt, W., Niemeyer, J. C., Klein, R., \& Gröbl, A. (1999). A new model for\par}
{\small\raggedright\noindent deflagration fronts in reactive fluids. Astronomy and Astrophysics , 347 , 724–733. Reynolds, O. (1895). On the dynamical theory of incompressible viscous fluids and the determi-\par}
{\small\raggedright\noindent nation of the criterion. Philosophical Transactions of the Royal Society London Series A , 186 , 123–164. Rhie, C. M., \& Chow, W. L. (1983). A numerical study of the turbulent flow past an isolated airfoil\par}
{\small\raggedright\noindent with trailing edge separation. AIAA Journal , 21 , 1525–1532. Riahi, H., Meldi, M., Favier, J., Serre, E., \& Goncalves, E. (2018). A pressure-corrected immersed\par}
{\small\raggedright\noindent boundary method for the numerical simulation of compressible flows. Journal of Computational Physics , 374 , 361–383. Richardson, L. F. (1910). The approximate arithmetical solution by finite differences of physical\par}
{\small\raggedright\noindent problems involving differential equations with an application to the stresses in a masonry dam. Philosophical Transactions of the Royal Society London Series A , 210 , 307–357. Richtmyer, R. D., \& Morton, K. W. (1967). Difference methods for initial value problems . New\par}
{\small\raggedright\noindent York: Wiley. Rider, W. J. (2007). Effective subgrid modeling from the ILES simulation of compressible turbu-\par}
{\small\raggedright\noindent lence. Journal of Fluids Engineering , 129 , 1493–1496. Rider, W. J., \& Kothe, D. B. (1998). Reconstructing volume tracking. Journal of Computational\par}
{\small\raggedright\noindent Physics , 141 , 112–152. Rizzi, A., \& Viviand, H. (Eds.). (1981). Numerical methods for the computation of inviscid transonic\par}
{\small\raggedright\noindent flows with shock waves ., Notes on numerical fluid mechanics (Vol. 3) Braunschweig: Vieweg. Roache, P. J. (1994). Perspective: A method for uniform reporting of grid refinement studies. ASME\par}
{\small\raggedright\noindent Journal of Fluids Engineering , 116 , 405–413. Rodi, W. (1976). A new algebraic relation for calculating the Reynolds stress. ZAMM , 56 , T219–\par}
{\small\raggedright\noindent T221. Rodi, W., Bonnin, J.-C., \& Buchal, T. (Eds.). (1995). In Proceedings of ERCOFTAC Workshop\par}
{\small\raggedright\noindent on Databases and Testing of Calculation Methods for Turbulent Flows, April 3–7 . Germany: University of Karlsruhe. Rodi, W., Constantinescu, G., \& Stoesser, T. (2013). Large-eddy simulation in hydraulics . London:\par}
{\small\raggedright\noindent Taylor \& Francis. Roe, P. L. (1986). Characteristic-based schemes for the Euler equations. Annual Review of Fluid\par}
{\small\raggedright\noindent Mechanics , 18 , 337–365. Rogallo, R. S. (1981). Numerical experiments in homogeneous turbulence (Technical Report No.\par}
{\small\raggedright\noindent 81315). Ames Research Center, CA: NASA. Saad, Y. (2003). Iterative methods for sparse linear systems (2nd ed.). Philadelphia: Society for\par}
{\small\raggedright\noindent Industrial and Applied Mathematics (SIAM). Saad, Y., \& Schultz, M. H. (1986). GMRES: A generalized residual algorithm for solving nonsym-\par}
{\small\raggedright\noindent metric linear systems. SIAM Journal on Scientific Statistical Computing , 7 , 856–869. Sagaut, P. (2006). Large-eddy simulation for incompressible flows: An introduction (3rd ed.). Berlin:\par}
{\small\raggedright\noindent Springer.\par}
{\small\raggedright\noindent Sani, R. L., Shen, J., Pironneau, O., \& Gresho, P. M. (2006). Pressure boundary condition for the\par}
{\small\raggedright\noindent time-dependent incompressible Navier-Stokes equations. International Journal for Numerical Methods in Fluids , 50 , 673–682. Sauer, J. (2000). Instationär kavitierende Strömungen - ein neues Modell, basierend auf Front\par}
{\small\raggedright\noindent Capturing (VoF) und Blasendynamik (Ph.D. Dissertation) . Germany: University of Karlsruhe. Scardovelli, R., \& Zaleski, S. (1999). Direct numerical simulation of free-surface and interfacial\par}
{\small\raggedright\noindent flow. Annual Review of Fluid Mechanics , 31 , 567–603. Schalkwijk, J., Griffith, E., Post, F. H., \& Jonker, H. J. J. (2012). High-performance simulations of\par}
{\small\raggedright\noindent turbulent clouds on a desktop PC: Exploiting the GPU. Bulletin of the American Meteorological Society , 93 , 307–314. Schalkwijk, J., Jonker, H. J. J., Siebesma, A. P., \& van Meijgaard, E. (2015). Weather forecasting\par}
{\small\raggedright\noindent using GPU-based large-eddy simulations. Bulletin of the American Meteorological Society , 96 , 715–723. Schlatter, P., \& Örlü, R. (2010). Assessment of direct numerical simulation data of turbulent bound-\par}
{\small\raggedright\noindent ary layers. Journal of Fluid Mechanics , 659 , 116–126. Schneider, G. E., \& Raw, M. J. (1987). Control-volume finite-element method for heat transfer and\par}
{\small\raggedright\noindent fluid flow using colocated variables. 1. Computational procedure. Numerical Heat Transfer , 11 , 363–390. Schneider, G. E., \& Zedan, M. (1981). A modified strongly implicit procedure for the numerical\par}
{\small\raggedright\noindent solution of field problems. Numerical Heat Transfer , 4 , 1–19. Schnerr, G. H., \& Sauer, J. (2001). Physical and numerical modeling of unsteady Cavitation dynam-\par}
{\small\raggedright\noindent ics. In Fourth International Conference on Multiphase Flow. New Orleans, USA . Schreck, E., \& Peri ́c, M. (1993). Computation of fluid flow with a parallel multigrid solver. Inter-\par}
{\small\raggedright\noindent national Journal for Numerical Methods in Fluids , 16 , 303–327. Schumacher, J., Sreenivasan, K., \& Yeung, P. (2005). Very fine structures in scalar mixing. Journal\par}
{\small\raggedright\noindent of Fluid Mechanics , 531 , 113–122. Sedov, L. (1971). A course in continuum mechanics (Vol. 1). Groningen: Wolters-Noordhoft Pub-\par}
{\small\raggedright\noindent lishing. Seidl, V. (1997). Entwicklung und Anwendung eines parallelen Finite-Volumen-Verfahrens zur Strö-\par}
{\small\raggedright\noindent mungssimulation auf unstrukturierten Gittern mit lokaler Verfeinerung (Ph.D. Dissertation) . Germany: University of Hamburg. Seidl, V., Muzaferija, S., \& Peri ́c, M. (1998). Parallel DNS with local grid refinement. Applied\par}
{\small\raggedright\noindent Scientific Research , 59 , 379–394. Seidl, V., Peri ́c, M., \& Schmidt, S. (1996). Spaceand time-parallel Navier-Stokes solver for 3D\par}
{\small\raggedright\noindent block-adaptive Cartesian grids. In A. Ecer, J. Periaux, N. Satofuka, \& S. Taylor (Eds.), P arallel Computational Fluid Dynamics 1995: Implementations and results using parallel computers (pp. 577–584). Amsterdam: North Holland, Elsevier. Senocak, I., \& Jacobsen, D. (2010). Acceleration of complex terrain wind predictions using many-\par}
{\small\raggedright\noindent core computing hardware. In 5th International Symposium on Computational Wind Engineering. (CWE2010) (Paper 498) . International Association for Wind Engineering. Sethian, J. A. (1996). Level set methods . Cambridge: Cambridge University Press. Shabana, A. A. (2013). Dynamics of multibody systems (4th ed.). New York: Cambridge University\par}
{\small\raggedright\noindent Press. Shah, K. B., \& Ferziger, J. H. (1995). Large-eddy simulations of flow past a cubic obstacle. In\par}
{\small\raggedright\noindent Annual Research Briefs . Stanford, CA: Center for Turbulence Research. Shah, K. B., \& Ferziger, J. H. (1997). A fluid mechanicians view of wind engineering: Large-\par}
{\small\raggedright\noindent eddy simulation of flow over a cubical obstacle. Journal of Wind Engineering \& Industrial Aerodynamics , 67 \& 68 , 211–224. Shen, J. (1993). A remark on the Projection-3 method. International Journal for Numerical Methods\par}
{\small\raggedright\noindent in Fluids , 16 , 249–253. Shewchuk, J. R. (1994). An introduction to the conjugate gradient method without the agonizing\par}
{\small\raggedright\noindent pain. Pitt., PA: School of Computer Science, Carnegie Mellon University. Retrieved from http:// \url{www.cs.cmu.edu/quake-papers/painless-conjugate-gradient.pdf} .\par}
{\small\raggedright\noindent Shi, X., Chow, F. K., Street, R. L., \& Bryan, G. H. (2018a). An evaluation of LES turbulence\par}
{\small\raggedright\noindent models for scalar mixing in the stratocumulus-capped boundary layer. Journal of the Atmospheric Sciences , 75 , 1499–1507. Shi, X., Hagen, H. L., Chow, F. K., Bryan, G. H., \& Street, R. L. (2018b). Large-eddy simulation\par}
{\small\raggedright\noindent of the stratocumulus-capped boundary layer with explicit filtering and reconstruction turbulence modeling. J. Atmos. Sci. , 75 , 611–637. Shih, T.-H., \& Liu, N.-S. (2009). A very-large-eddy simulation of the nonreacting flow in a single\par}
{\small\raggedright\noindent element lean direct injection combustor using PRNS with a nonlinear subscale model (Technical Report No. 2009-21564). Cleveland, OH: NASA Glenn Research Center. Skamarock, W. C., \& Klemp, J. B. (2008). A time-split nonhydrostatic atmospheric model for\par}
{\small\raggedright\noindent weather research and forecasting operations. Journal of Computational Physics , 227 , 3465–3485. Skamarock, W. C., Oliger, J., \& Street, R. L. (1989). Adaptive grid refinement for numerical weather\par}
{\small\raggedright\noindent prediction. Journal of Computational Physics , 80 , 27–60. Skyllingstad, E. D., \& Samelson, R. M. (2012). Baroclinic frontal instabilities and turbulent mixing\par}
{\small\raggedright\noindent in the surface boundary layer. Part I: Unforced simulations. Journal of Physical Oceanography , 42 , 1701–1716. Smagorinsky, J. (1963). General circulation experiments with the primitive equations. Part I: The\par}
{\small\raggedright\noindent basic experiment. Monthly Weather Review , 91 , 99–164. Smiljanovski, V., Moser, V., \& Klein, R. (1997). A capturing-tracking hybrid scheme for deflagration\par}
{\small\raggedright\noindent discontinuities. Combustion Theory and Modelling , 1 , 183–215. Smits, A. J., \& Marusic, I. (2013). Wall-bounded turbulence. Physics Today , 66 , 25–30. Smits, A. J., McKeon, B. J., \& Marusic, I. (2011). High-Reynolds number wall turbulence. Annual\par}
{\small\raggedright\noindent Review of Fluid Mechanics , 43 , 353–375. Smolarkiewicz, P. K.,\& Margolin, L. G. (2007). Studies in geophysics. In F. Grinstein, L. Mar-\par}
{\small\raggedright\noindent golin, \& W. Rider (Eds.), Implicit large-eddy simulation: Computing turbulent fluid dynamics . Cambridge: Cambridge University Press. Smolarkiewicz, P. K., \& Prusa, J. M. (2002). VLES modelling of geophysical fluids with nonoscil-\par}
{\small\raggedright\noindent latory forward-in-time schemes. International Journal for Numerical Methods in Fluids , 39 , 799–819. Sonar, T. (1997). On the construction of essentially non-oscillatory finite volume approximations\par}
{\small\raggedright\noindent to hyperbolic conservation laws on general triangulations: Polynomial recovery, accuracy and stencil selection. Computer Methods in Applied Mechanics and Engineering , 140 , 157. Sonneveld, P. (1989). CGS, a fast Lanczos type solver for non-symmetric linear systems. SIAM\par}
{\small\raggedright\noindent Journal on Scientific Computing , 10 , 36–52. Spalart, P. R., \& Allmaras, S. R. (1994). A one-equation turbulence model for aerodynamic flows.\par}
{\small\raggedright\noindent La Recherche Aerospatiale , 1 , 5–21. Spalart, P. R., Deck, S., Shur, M. L., Squires, K. D., Strelets, M. K., \& Travin, A. (2006). A\par}
{\small\raggedright\noindent new version of detached-eddy simulation, resistant to ambiguous grid densities. Theoretical and Computational Fluid Dynamics , 20 , 181–195. Spalding, D. B. (1972). A novel finite-difference formulation for differential expressions involving\par}
{\small\raggedright\noindent both first and second derivatives. International Journal for Numerical Methods in Fluids , 4 , 551–559. Spalding, D. B. (1978). General theory of turbulent combustion. Journal of the Energy , 2 , 16–23. Spotz, W. (1998). Accuracy and performance of numerical wall boundary conditions for steady, 2D,\par}
{\small\raggedright\noindent incompressible streamfunction vorticity. International Journal for Numerical Methods in Fluids , 28 , 737–757. Spotz, W. F., \& Carey, G. F. (1995). High-order compact scheme for the steady stream-function\par}
{\small\raggedright\noindent vorticity equations. International Journal for Numerical Methods in Fluids , 38 , 3497–3512. Sta. Maria, M., \& Jacobson, M. (2009). Investigating the effect of large wind farms on the energy\par}
{\small\raggedright\noindent in the atmosphere. Energies , 2 , 816–838. Steger, J. L., \& Warming, R. F. (1981). Flux vector splitting of the inviscid gas-dynamic equations\par}
{\small\raggedright\noindent with applications to finite difference methods. Journal of Computational Physics , 40 , 263–293.\par}
{\small\raggedright\noindent Stoll, R.,\&Porté-Agel, F. (2006). Effect of roughness on surface boundary conditions for large-eddy\par}
{\small\raggedright\noindent simulation. Boundary-Layer Meteorology , 118 , 169–187. Stoll, R., \& Porté-Agel, F. (2008). Large-eddy simulation of the stable atmospheric boundary layer\par}
{\small\raggedright\noindent using dynamic models with different averaging schemes. Boundary-Layer Meteorology , 126 , 1–28. Stolz, S., Adams, N. A., \& Kleiser, L. (2001). An approximate deconvolution model for largeeddy\par}
{\small\raggedright\noindent simulation with application to incompressible wall-bounded flows. Physics of Fluids , 13 , 997– 1015. Stone, H. L. (1968). Iterative solution of implicit approximations of multidimensional partial dif-\par}
{\small\raggedright\noindent ferential equations. SIAM Journal on Numerical Analysis , 5 , 530–558. Street, R. L. (1973). Analysis and solution of partial differential equations . Monterey: Brooks/Cole\par}
{\small\raggedright\noindent Publishing Co. Street, R. L., Watters, G. Z., \& Vennard, J. K. (1996). Elementary fluid mechanics (7th ed.). New\par}
{\small\raggedright\noindent York: Wiley. Sullivan, P. P., Weil, J. C., Patton, E. G., Jonker, H. J. J., \& Mironov, D. V. (2016). Turbulent\par}
{\small\raggedright\noindent winds and temperature fronts in large-eddy simulations of the stable atmospheric boundary layer. Journal of the Atmospheric Sciences , 73 , 1815–1840. Sullivan, P. P., Horst, T. W., Lenschow, D. H., Moeng, C.-H., \& Weil, J. C. (2003). Structure of\par}
{\small\raggedright\noindent subfilter-scale fluxes in the atmospheric surface layer with application to large-eddy simulation modelling. Journal of Fluid Mechanics , 482 , 101–139. Sullivan, P. P., \& Patton, E. G. (2008). A highly parallel algorithm for turbulence simulations in\par}
{\small\raggedright\noindent planetary boundary layers: Results with meshes up to 1024 3 . In 18th Conference on Boundary Layers and Turbulence, AMS (pp. Paper 11B.5, 11). Stockholm, Sweden . Sullivan, P. P., \& Patton, E. G. (2011). The effect of mesh resolution on convective boundary layer\par}
{\small\raggedright\noindent statistics and structures generated by large-eddy simulation. Journal of the Atmospheric Sciences , 68 , 2395–2415. Sunderam, V. S. (1990). PVM: A framework for parallel distributed computing. Concurrency and\par}
{\small\raggedright\noindent computaton: Practice and Experience , 2 , 315–339. Sussman, M. (2003). A second-order coupled level set and volume-of-fluid method for computing\par}
{\small\raggedright\noindent growth and collapse of vapor bubbles. Journal of Computational Physics , 187 , 110–136. Sussman, M., Smereka, P., \& Osher, S. (1994). A level set approach for computing solutions to\par}
{\small\raggedright\noindent incompressible two-phase flow. Journal of Computational Physics , 114 , 146–159. Sweby, P. K. (1984). High resolution schemes using flux limiters for hyperbolic conservation laws.\par}
{\small\raggedright\noindent SIAM Journal on Numerical Analysis , 21 , 995–1011. Sweby, P. K. (1985). High resolution TVD schemes using flux limiters. Large-scale computations in\par}
{\small\raggedright\noindent fluid mechanics , Lecture notes in applied mathematics, Part 2 (Vol. 22, pp. 289–309). Providence: American Mathematical Society. Taira, K., \& Colonius, T. (2007). The immersed boundary method: A projection approach. Journal\par}
{\small\raggedright\noindent of Computational Physics , 225 , 2118–2137. Taneda, S. (1978). Visual observations of the flow past a sphere at Reynolds numbers between 10 4\par}
{\small\raggedright\noindent and 10 6 . Journal of Fluid Mechanics , 85 , 187–192. Tannehill, J. C., Anderson, D. A., \& Pletcher, R. H. (1997). Computational fluid mechanics and\par}
{\small\raggedright\noindent heat transfer (Vol. 2). Penn: Taylor \& Francis. Tennekes, H., \& Lumley, J. L. (1976). A first course in turbulence . Cambridge, MA: MIT Press. Thé, J. L., Raithby, G. D., \& Stubley, G. D. (1994). Surface-adaptive finite-volume method for\par}
{\small\raggedright\noindent solving free-surface flows. Numerical Heat Transfer, Part B , 26 , 367–380. Thibault, J. C., \& Senocak, I. (2009). CUDA implementation of a Navier-Stokes solver on multi-\par}
{\small\raggedright\noindent GPU desktop platforms for incompressible flows. In 47th AIAA Aerospace Science Meeting, AIAA Paper 2009-758. Thompson, J. F., Warsi, Z. U. A., \& Mastin, C. W. (1985). Numerical grid generation - foundations\par}
{\small\raggedright\noindent and applications . New York: Elsevier. Thompson, M. C., \& Ferziger, J. H. (1989). A multigrid adaptive method for incompressible flows.\par}
{\small\raggedright\noindent Journal of Computational Physics , 82 , 94–121.\par}
{\small\raggedright\noindent Tokuda, Y., Song, M.-H., Ueda, Y., Usui, A., Akita, T., Yoneyama, S., et al. (2008). Three-\par}
{\small\raggedright\noindent dimensional numerical simulation of blood flow in the aortic arch during cardiopulmonary bypass. European Journal of Cardio-Thoracic Surgery , 33 , 164–167. Tregde, V. (2015). Compressible air effects in CFD simulations of free fall lifeboat drop. In ASME\par}
{\small\raggedright\noindent 34th International Conference on Ocean, Offshore and Arctic Engineering. St John’s, Newfoundland, Canada . Truesdell, C. (1991). A first course in rational continuum mechanics (2nd ed., Vol. 1). Boston:\par}
{\small\raggedright\noindent Academic. Tryggvason, G., \& Unverdi, S. O. (1990). Computations of 3-dimensional Rayleigh-Taylor insta-\par}
{\small\raggedright\noindent bility. Physics of Fluids A , 2 , 656–659. Tseng, Y. H., \& Ferziger, J. H. (2003). A ghost-cell immersed boundary method for flow in complex\par}
{\small\raggedright\noindent geometry. Journal of Computational Physics , 192 , 593–623. Tu, J. Y., \& Fuchs, L. (1992). Overlapping grids and multigrid methods for three-dimensional\par}
{\small\raggedright\noindent unsteady flow calculation in IC engines. International Journal for Numerical Methods in Fluids , 15 , 693–714. Tukovi ́c, Ž., Peri ́c, M., \& Jasak, H. (2018). Consistent second-order time-accurate non-iterative\par}
{\small\raggedright\noindent PISO algorithm. Computers Fluids , 166 , 78–85. Turkel, E. (1987). Preconditioned methods for solving the incompressible and low speed compress-\par}
{\small\raggedright\noindent ible equations. Journal of Computational Physics , 72 , 277–298. Ubbink, O. (1997). Numerical prediction of two fluid systems with sharp interfaces (Ph.D. Disser-\par}
{\small\raggedright\noindent tation) . London: University of London. van der Vorst, H. A. (1992). BI-CGSTAB: a fast and smoothly converging variant of BI-CG for the\par}
{\small\raggedright\noindent solution of non-symmetric linear systems. SIAM Journal on Scientific Computing , 13 , 631–644. van der Vorst, H. A. (2002). Efficient and reliable iterative methods for linear systems. Journal of\par}
{\small\raggedright\noindent Computational and Applied Mathematics , 149 , 251–265. van der Vorst, H. A., \& Sonneveld, P. (1990). CGSTAB, a more smoothly converging variant of\par}
{\small\raggedright\noindent CGS (Technical Report No. 90-50). Delft, NL: Delft University of Technology. van der Wijngaart, R. F. (1990). Composite-grid techniques and adaptive mesh refinement in com-\par}
{\small\raggedright\noindent putational fluid dynamics (Ph.D. Dissertation) . Stanford CA: Stanford University. Van Doormal, J. P., \& Raithby, G. D. (1984). Enhancements of the SIMPLE method for predicting\par}
{\small\raggedright\noindent incompressible fluid flows. Numerical Heat Transfer , 7 , 147–163. VanDoormal,J.P.,Raithby,G.D.,\&McDonald,B.H.(1987).Thesegregateda\hspace{0pt}pproachtopredi\hspace{0pt}cting\par}
{\small\raggedright\noindent viscous compressible fluid flows. ASME Journal of Turbomachinery , 109 , 268–277. Vanka, S. P. (1986). Block-implicit multigrid solution of Navier-Stokes equations in primitive\par}
{\small\raggedright\noindent variables. Journal of Computational Physics , 65 , 138–158. Van Leer, B. (1977). Towards the ultimate conservative difference scheme. IV. A new approach to\par}
{\small\raggedright\noindent numerical convection. Journal of Computational Physics , 23 , 276–299. Van Leer, B. (1985). Upwind-difference methods for aerodynamic problems governed by the Euler\par}
{\small\raggedright\noindent equations. Large-scale computations in fluid mechanics , Lecture notes in Applied Mathematics, Part 2, (Vol. 22, 327–336). Providence: American Mathematical Society. Viswanathan, A. K., Squires, K. D., \& Forsythe, J. R. (2008). Detached-eddy simulation around a\par}
{\small\raggedright\noindent forebody with rotary motion. AIAA Journal , 46 , 2191–2201. Vukˇcevi ́c, V., Jasak, H., \& Gatin, I. (2017). Implementation of the ghost fluid method for free surface\par}
{\small\raggedright\noindent flows in polyhedral finite volume framework. Computers Fluids , 153 , 1–19. Wakashima, S., \& Saitoh, T. S. (2004). Benchmark solutions for natural convection in a cubic cavity\par}
{\small\raggedright\noindent using the high-order time-space method. International Journal of Heat and Mass Transfer , 47 , 853–864. Wallin, S., \& Johansson, A. V. (2000). An explicit algebraic Reynolds stress model for incompress-\par}
{\small\raggedright\noindent ible and compressible turbulent flows. Journal of Fluid Mechanics , 403 , 89–132. Wang, R., Feng, H., \& Huang, C. (2016). A new mapped weighted essentially non-oscillatory\par}
{\small\raggedright\noindent method using rational mapping function. Journal of Scientific Computing , 67 , 540–580. Warming, R. F., \& Hyett, B. J. (1974). The modified equation approach to the stability and accuracy\par}
{\small\raggedright\noindent of finite-difference methods. Journal of Computational Physics , 14 , 159–179.\par}
{\small\raggedright\noindent Washington, W. M., \& Parkinson, C. L. (2005). An introduction to three-dimensional climate mod-\par}
{\small\raggedright\noindent eling (2nd ed.). Sausalito: University Science Books. Waterson, N. P., \& Deconinck, H. (2007). Design principles for bounded higher-order convection\par}
{\small\raggedright\noindent schemes - a unified approach. Journal of Computational Physics , 224 , 182–207. Watkins, D. S. (2010). Fundamentals of matrix computations (3rd ed.). New York: Wiley- Inter-\par}
{\small\raggedright\noindent science. Wegner, B., Maltsev, A., Schneider, C., Sadiki, A., Dreizler, A., \& Janicka, J. (2004). Assessment of\par}
{\small\raggedright\noindent unsteady RANS in predicting swirl flow instability based on LES and experiments. International Journal of Heat and Fluid Flow , 25 , 528–536. Weinan, E., \& Liu, J.-G. (1997). Finite difference methods for 3D viscous incompressible flows\par}
{\small\raggedright\noindent in the vorticity - vector potential formulation on nonstaggered grids. Journal of Computational Physics , 138 , 57–82. Weiss, J. M., Maruszewski, J. P., \& Smith, W. A. (1999). Implicit solution of preconditioned Navier-\par}
{\small\raggedright\noindent Stokes equations using algebraic multigrid. AIAA Journal , 37 , 29–36. Weiss, J. M., \& Smith, W. A. (1995). Preconditioning applied to variable and constant density flows.\par}
{\small\raggedright\noindent AIAA Journal , 33 , 2050–2057. Wesseling, P. (1990). Multigrid methods in computational fluid dynamics. ZAMM - Z. Angew\par}
{\small\raggedright\noindent Mathematics and Mechanics , 70 , T337–T347. Weymouth, G., \& Yue, D. K. P. (2010). Conservative volume-of-fluid method for free-surface\par}
{\small\raggedright\noindent simulations on Cartesian grids. Journal of Computational Physics , 229 , 2853–2865. White, F. M. (2010). Fluid mechanics (7th ed.). New York: McGraw Hill. Wicker, L. J., \& Skamarock, W. C. (2002). Time-splitting methods for elastic models using forward\par}
{\small\raggedright\noindent time schemes. Monthly Weather Review , 130 , 2088–2097. Wie, S. Y., Lee, J. H., Kwon, J. K., \& Lee, D. J. (2010). Far-field boundary condition effects of\par}
{\small\raggedright\noindent CFD and free-wake coupling analysis for helicopter rotor. Journal of Fluids Engineering , 132 , 84501-1–6. Wikstrom, P. M., Wallin, S., \& Johansson, A. V. (2000). Derivation and investigation of a new\par}
{\small\raggedright\noindent explicit algebraic model for the passive scalar flux. Physics of Fluids , 12 , 688–702. Wilcox, D. C. (2006). Turbulence modeling for CFD (3rd ed.). La Cañada: DCW Industries Inc. Williams, F. A. (1985). Combustion theory: The fundamental theory of chemically reacting flow\par}
{\small\raggedright\noindent systems . Menlo Park: Benjamin-Cummings Publishing Co. Williams, J., Sarofeen, C., Shan, H., \& Conley, M. (2016). An accelerated iterative linear solver with\par}
{\small\raggedright\noindent GPUs and CFD calculations of unstructured grids. Procedia Computer Science , 80 , 1291–1300. Williams, P. D. (2009). A proposed modification to the Robert-Asselin time filter. Monthly Weather\par}
{\small\raggedright\noindent Review , 137 , 2538–2546. Wong, V. C., \& Lilly, D. K. (1994). A comparison of two dynamic subgrid closure methods for\par}
{\small\raggedright\noindent turbulent thermal convection. Physics of Fluids , 6 , 1016–1023. Wosnik, M., Castillo, L., \& George, W. K. (2000). A theory for turbulent pipe and channel flows.\par}
{\small\raggedright\noindent Journal of Fluid Mechanics , 412 , 115–145. Wrobel, L. C. (2002). The boundary element method. Vol. 1: Applications in thermo-fluids and\par}
{\small\raggedright\noindent acoustics. New York: Wiley. Wu, X., \& Moin, P. (2009). Direct numerical simulation of turbulence in a nominally zero-\par}
{\small\raggedright\noindent pressuregradient flat-plate boundary layer. Journal of Fluid Mechanics , 630 , 5–41. Wu, X., \& Moin, P. (2011). Evidence for the persistence of hairpin forest in turbulent, zeropressure-\par}
{\small\raggedright\noindent gradient flat-plate boundary layers. In 7th International Symposium for Turbulence and Shear Flow Phenomena (TSFP-7), Paper 6A4P. Ottawa, Canada . Wyngaard, J. C. (2004). Toward numerical modeling in the “Terra Incognita”. Journal of the Atmo-\par}
{\small\raggedright\noindent spheric Sciences , 61 , 1816–1826. Wyngaard, J. C. (2010). Turbulence in the atmosphere . Cambridge: Cambridge University Press. Xing-Kaeding, Y. (2006). Unified approach to ship seakeeping and maneuvering by a RANSE\par}
{\small\raggedright\noindent method (Ph.D. Dissertation, TU Hamburg-Harburg. Hamburg). Retrieved from \url{http://doku.b.tuharburg.de/volltexte/2006/303/pdf/Xing-Kaeding-thesis.pdf} .\par}
{\small\raggedright\noindent Xue, M. (2000). High-order monotonic numerical diffusion and smoothing. Monthly Weather\par}
{\small\raggedright\noindent Review , 128 , 2853–2864. Xue, M., Drogemeier, K. K., \& Wong, V. (2000). The advanced regional prediction system (ARPS) -\par}
{\small\raggedright\noindent a multi-scale nonhydrostatic atmospheric simulation and prediction model. Part I: Model dynamics and verification. Meteorology and Atmospheric Physics , 75 , 161–193. Yang, H. Q., \& Przekwas, A. J. (1992). A comparative study of advanced shock-capturing schemes\par}
{\small\raggedright\noindent applied to Burger’s equation. Journal of Computational Physics , 102 , 139–159. Ye, T., Mittal, R., Udaykumar, H. S., \& Shyy, W. (1999). An accurate Cartesian grid method for\par}
{\small\raggedright\noindent viscous incompressible flows with complex immersed boundaries. Journal of Computational Physics , 156 , 209–240. Youngs, D. L. (1982). Time-dependent multi-material flowwith large fluid distortion. In K. W.\par}
{\small\raggedright\noindent Morton \& M. J. Baines (Eds.), Numerical methods for fluid dynamics (pp. 273–285). New York: Academic. Zalesak, S. T. (1979). Fully multidimensional flux-corrected transport algorithms for fluids. Journal\par}
{\small\raggedright\noindent of Computational Physics , 31 , 335–362. Zang, Y., \& Street, R. L. (1995). A composite multigrid method for calculating unsteady incom-\par}
{\small\raggedright\noindent pressible flows in geometrically complex domains. International Journal for Numerical Methods in Fluids , 20 , 341–361. Zang, Y., Street, R. L., \& Koseff, J. R. (1993). A dynamic mixed subgrid-scale model and its\par}
{\small\raggedright\noindent application to turbulent recirculating flows. Physics of Fluids A , 5 , 3186–3196. Zang, Y., Street, R. L., \& Koseff, J. R. (1994). A non-staggered grid, fractional-step method for\par}
{\small\raggedright\noindent time-dependent incompressible Navier-Stokes equations in curvilinear coordinates. Journal of Computational Physics , 114 , 18–33. Zhang, H., Zheng, L. L., Prasad, V., \& Hou, T. Y. (1998). A curvilinear level set formulation\par}
{\small\raggedright\noindent for highly deformable free surface problems with application to solidification. Numerical Heat Transfer , 34 , 1–20. Zhou, B., \& Chow, F. K. (2011). Large-eddy simulation of the stable boundary layer with explicit\par}
{\small\raggedright\noindent filtering and reconstruction turbulence modeling. Journal of the Atmospheric Sciences , 68 , 2142– 2155. Zhou, B., Simon, J. S., \& Chow, F. K. (2014). The convective boundary layer in the Terra Incognita.\par}
{\small\raggedright\noindent Journal of the Atmospheric Sciences , 71 , 2547–2563. Zienkiewicz, O. C., Taylor, R. L., \& Nithiarasu, P. (2005). The finite element method for fluid\par}
{\small\raggedright\noindent dynamics (6th ed.). Burlington: Butterworth- Heinemann (Elsevier). Zwart, P. J., Gerber, G., \& Belamri, T. (2004). A two-phase flow model for prediction of cavitation\par}
{\small\raggedright\noindent dynamics. In Fifth International Conference on Multiphase Flow Yokohama, Japan\par}

\end{document}
